\xchapter{Prefatory}

\textbf{TO }

\textbf{\emph{HIS HONORED AND BELOVED COLLEAGUES}}

\textbf{IN THE UNION THEOLOGICAL SEMINARY,}

\textbf{Rev. WILLIAM ADAMS, D.D., LL.D.,}

\textbf{Rev. HENRY B. SMITH, D.D., LL.D.,}

\textbf{Rev. ROSWELL D. HITCHCOCK, D.D., LL.D., }

\textbf{Rev. WILLIAM G. T. SHEDD, D.D., LL.D., }

\textbf{Rev. GEORGE L. PRENTISS, D.D., }

\textbf{Rev. CHARLES A. BRIGGS, D.D., }

THIS WORK IS

\textbf{Respectfully Dedicated}

\textbf{BY}

\textbf{THE AUTHOR }

\xsection{Preface}

A 'symbolical library' that contains the creeds and confessions of all Christian denominations fills a vacuum in theological and historical literature. It is surprising that it has not been supplied long ago. Sectarian exclusiveness or doctrinal indifferentism may have prevented it. Other symbolical collections are confined to particular denominations and periods. In this work the reader will find the authentic material for the study of Comparative Theology Symbolics, Polemics, and Irenics. In a country like ours, where people of all creeds meet in daily contact, this study ought to command more attention than it has hitherto received.

The First Volume has expanded into a doctrinal history of the Church, so far as it is embodied in public standards of faith. The most important and fully developed symbolical systems the Vatican Romanism, the Lutheranism of the Formula of Concord, and the Calvinism of the Westminster standards have been subjected to a critical analysis. The author has endeavored to combine the \textgreek{ἀληθεύειν ἐν ἀγάπῃ} and the \textgreek{ἀγαπᾷν ἐν ἀληθείᾳ,} and to be mindful of the golden motto, \emph{In necessariis unitas, in dubiis libertas, in omnibus caritas}. Honest and earnest controversy, conducted in a Christian and catholic spirit, promotes true and lasting union. Polemics looks to Irenics—the aim of war is peace.

The Second Volume contains the Scripture Confessions, the ante-Nicene Rules of Faith, the cumenical, the Greek, and the Latin Creeds, from the Confession of Peter down to the Vatican Decrees. It includes also the best Russian Catechism and the recent Old Catholic Union Propositions of the Bonn Conferences.

The Third Volume is devoted to the Lutheran, Anglican, Calvinistic, and the later Protestant Confessions of Faith. The documents of the Third Part (pp. 707–876) have never been collected before.

The creeds and confessions are given in the original languages from the best editions, and are accompanied by translations for the convenience of the English reader.\footnote{I have used, e.g., the fac-simile of the oldest MS. of the Athanasian Creed from the 'Utrecht Psalter:' the \emph{ed. princeps} of the Lutheran \emph{Concordia }(formerly in the possession of Dr. Meyer, the well-known commentator); the \emph{Corpus et Syntagma Confessionum, }ed. 1654; a copy of the \emph{Harmonia Confessionum, }once owned by Prince Casimir of the Palatinate, who suggested it; the oldest editions of the Westminster Confession and Catechisms, of the Savoy Declaration, etc.}

While these volumes were passing through the press several learned treatises on the ancient creeds by Lumby, Swainson, Hort, Caspari, and others have appeared, though not too late to be noticed in the final revision. The literature has been brought down to the close of 1876. I trust that nothing of importance has escaped my attention.

I take pleasure in acknowledging my obligation to several distinguished divines, in America and England, for valuable information concerning the denominations to which they belong, and for several contributions, which appear under the writers' names.\footnote{The Rev. Drs. Jos. Angus, W. W. Andrews, Chas. A. Briggs, J. R. Brown, E. W. Gilman, G. Haven, A. A. Hodge, Alex. F. Mitchell, E. D. Morris, Chas. P. Krauth, J. R. Lumby, G. D. Matthews, H. Osgood, E. von Schweinitz, H. B. Smith, John Stoughton, E. A. Washburn, W. R. Williams. See Vol. I. pp. 609, 811, 839, 911; Vol. III. pp. 3, 738, 777, 799.} In a history of conflicting creeds it is wise to consult representative men as well as books, in order to secure strict accuracy and impartiality, which are the cardinal virtues of a historian.

May this repository of creeds and confessions promote a better understanding among the Churches of Christ. The divisions of Christendom bring to light the various aspects and phases of revealed truth, and will be overruled at last for a deeper and richer harmony, of which Christ is the key-note. In him and by him all problems of theology and history will be solved. The nearer believers of different creeds approach the Christological centre, the better they will understand and love each other.

\textbf{P. S.}

Bible House, New York,

\emph{December, }1876.

\xsection{Preface to the Second Edition}

The call for a new edition of this work in less than a year after its publication is an agreeable surprise to the author, and fills him with gratitude to the reading public and the many reviewers, known and unknown, who have so kindly and favorably noticed it in American and foreign periodicals and in private letters. One of the foremost divines of Germany (Dr. Dorner, in the \emph{Jahrbüher für Deutsche Theologie, }1877, p. 682) expresses a surprise that the idea of such an œcumenical collection of Christian Creeds should have originated in America, where the Church is divided into so many rival denominations; but he adds also as an explanation that this division creates a desire for unity and co-operation, and a mutual courtesy and kindness unknown among the contending parties and schools under the same roof of state-churches, where outward uniformity is maintained at the expense of inward peace and harmony.

The changes in this edition are very few. The literature in the first volume is brought down to the present date, and at the close of the second volume a fac-simile of the oldest MSS. of the Athanasian Creed and the Apostles' Creed is added.

\textbf{P. S.}

NEW YORK, \emph{April,} 1878.

\xsection{Preface to the Third Edition}

This edition differs from the second in the following particulars:

1. In the first volume several errors have been corrected (e.g., in the statistical table, p. 818), and a list of new works inserted on p. xiv.

2. In the third volume a translation of the Second Helvetic Confession has been added, pp. 831 sqq.

\textbf{P. S.}

New York, \emph{December,} 1880.

\xsection{Preface to the Fourth Edition}

The call for a fourth edition of this work has made it my duty to give the first volume once more a thorough revision and to bring the literature down to the latest date. In this I have been aided by my young friend, the Rev. Samuel M. Jackson, one of the assistant editors of my ``Religious Encyclopædia.'' The additions which could not be conveniently made in the plates have been printed separately after the Table of Contents, pp. xiv–xvii.

The second and third volumes, which embrace the symbolical documents, remain unchanged, except that at the end of the third volume the new Congregational Creed of 1883 has been added.

Creeds will live as long as faith survives, with the duty to confess our faith before men. By and by we shall reach, through the Creeds of Christendom, the one comprehensive, harmonious Creed of Christ.

\textbf{P. S.}

New York,\emph{ May, }1884.

The fifth edition was a reprint of the fourth, without any changes.

\xsection{Preface to the Sixth Edition}

Since the appearance of the \emph{Creeds of Christendom,} 1877, no work has been issued competing with it in scope and comprehensiveness. The valuable collection of W. W. Walker, \emph{The Creeds and Platforms of Congregationalism,} 1893, and W. J. McGlothlin, \emph{Baptist Confessions of Faith,} 1911, are limited to separate Protestant bodies. The extensive collection of Karl Müller, 1903, is confined to the creeds and catechisms of the Reformed Churches. Professor W. A. Curtis of the University of Edinburgh, in his \emph{History of the Creeds and Confessions of Faith in Christendom and Beyond,} gives the contents of creeds and an account of their origins, not their texts. C. Fabricius, in his \emph{Corpus confessionum,} etc., 1928, sqq., proposes in connexion with colaborers to furnish not only the texts of the Christian creeds, but also the texts of hymns, liturgies, books of discipline, and other documents bearing on Christian doctrine, worship, and practice. For example, 250 pages of Volume I are devoted to hymns, and 250 pages to ``The Doctrines and Discipline of the Methodist Episcopal Church, 1924.''

The new material of the present edition is the following:

Volume I. Additions to the literature; notices of the Church of the Disciples and the Universalist and Unitarian Churches; and changes and additions, as, for example, on the primitive creeds and the Russian Church.

Volume II. In the fourth edition Dr. Schaff, in view of the new importance given in Canon Law to papal utterances on doctrine and morals, added one of the important encyclicals of Leo XIII., who was then living. To this encyclical have been added bulls on the Church, by Boniface VIII., 1302, Anglican Orders, by Leo XIII., 1896, ``Americanism'' and ``Modernism'' by Pius X., 1907–10, and Pius XI.'s encyclical on Church Union, 1928.

Volume III. Additions giving Recent Confessional Declarations and Terms of Union between Church organizations. The material on the latter subject, so closely akin to the general topic of the book, makes it quite probable that Dr. Philip Schaff, in view of his pronounced attitude on Church fellowship and union, would have included it, were he himself preparing this edition of the \emph{Creeds of Christendom}.

David S. Schaff.

Union Theological Seminary

New York, \emph{January,} 1931

\xsection{[Original] Table of Contents}

TABLE OF CONTENTS.

(Vol. I.)

HISTORY OF THE CREEDS OF CHRISTENDOM.

FIRST CHAPTER.

ON CREEDS IN GENERAL.

PAGE

§ 1.

Name and Definition

3

§ 2.

Origin of Creeds

4

§ 3.

Authority of Creeds

7

§ 4.

Value and Use of Creeds

8

§ 5.

Classification of Creeds

9

SECOND CHAPTER.

THE ŒCUMENICAL CREEDS.

PAGE

§ 6.

General Character of the Œcumenical Creeds

12

§ 7.

The Apostles' Creed

14

§ 8.

The Nicene Creed

24

§ 9.

The Creed of Chalcedon

29

§ 10.

The Athanasian Creed

34

THIRD CHAPTER.

THE CREEDS OF THE GREEK CHURCH.

PAGE

§ 11.

The Seven Œcumenical Councils

43

§ 12.

The Confessions of Gennadius, A.D. 1453

46

§ 13.

The Answers of the Patriarch Jeremiah to the Lutherans, A.D. 1576

50

§ 14.

The Confession of Metrophanes Critopulus, A.D. 1625

52

§ 15.

The Confession of Cyril Lucar, A.D. 1631

54

§ 16.

The Orthodox Confession of Mogilas, A.D. 1643

58

§ 17.

The Synod of Jerusalem, and the Confession of Dositheus, A.D. 1672

61

§ 18.

The Synods of Constantinople, A.D. 1672 and 1691

67

§ 19.

The Doctrinal Standards of the Russo-Greek Church

68

§ 20.

Anglo-Catholic Correspondence with the Russo-Greek

74

§ 21.

The Eastern Sects: Nestorians, Jacobites, Copts, Armenians

78

FOURTH CHAPTER.

THE CREEDS OF THE ROMAN CHURCH.

PAGE

§ 22.

Catholicism and Romanism

83

§ 23.

Standard Expositions of the Roman Catholic System

85

§ 24.

The Canons and Decrees of the Council of Trent, A.D. 1563

90

§ 25.

The Profession of the Tridentine Faith, A.D. 1564

96

§ 26.

The Roman Catecism, A.D. 1566

100

§ 27.

The Papal Bulls against the Jansenists, A.D. 1653, 1713

102

Note on the Old Catholics in Holland, 107.

§ 28.

The Papal Definition of the Immaculate Conception of the Virgin Mary, A.D. 1854

108

§ 29.

The Argument for the Immaculate Conception

113

§ 30.

The Papal Syllabus, A.D. 1864

128

§ 31.

The Vatican Council, A.D. 1870

134

§ 32.

The Vatican Decrees. The Constitution on the Catholic Faith

147

§ 33.

The Vatican Decrees, Continued. The Papal Infallibility Decree

150

§ 34.

Papal Infallibility Explained, and Tested by Scripture and Tradition

163

§ 35.

The Liturgical Standards of the Roman Church

189

§ 36.

The Old Catholics

191

FIFTH CHAPTER.

THE CREEDS OF THE EVANGELICAL PROTESTANT CHURCHES.

PAGE

§ 37.

The Reformation. Protestantism and Romanism

203

§ 38.

The Evangelical Confessions of Faith

209

§ 39.

The Lutheran and Reformed Confessions

211

SIXTH CHAPTER.

THE CREEDS OF THE EVANGELICAL LUTHERAN CHURCH.

PAGE

§ 40.

The Lutheran Confessions

220

§ 41.

The Augsburg Confession, A.D. 1530

225

§ 42.

The Apology of the Augsburg Confession, A.D. 1530

243

§ 43.

Luther's Catechisms, A.D. 1529

245

§ 44.

The Articles of Smalcald, A.D. 1537

253

§ 45.

The Formula of Concord, A.D. 1577

258

§ 46.

The Formula of Concord, Concluded

312

§ 47.

Superseded Lutheran Symbols. The Saxon Confession, and the Würtemberg Confession, A.D. 1551

340

§ 48.

The Saxon Visitation Articles, A.D. 1592.

345

§ 49.

An Abortive Symbol against Syncretism, A.D. 1655

349

SEVENTH CHAPTER.

THE CREEDS OF THE EVANGELICAL REFORMED CHURCHES.

PAGE

§ 50.

The Reformed Confessions.

354

\emph{I. Reformed Confessions of Switzerland.}

§ 51.

Zwinglian Confessions. The Sixty-seven Articles. The Ten Theses of Berne. The Confession to Charles V. The Confession to Francis I., A.D. 1523-1531

360

§ 52.

Zwingli's Distinctive Doctrines

369

§ 53.

The Confession of Basle, A.D. 1534

385

§ 54.

The First Helvetic Confession, A.D. 1536

388

§ 55.

The Second Helvetic Confession, A.D. 1566

390

§ 56.

John Calvin. His Life and Character

421

§ 57.

444

§ 58.

The Catechism of Geneva, A.D. 1541

467

§ 59.

The Zurich Consensus, A.D. 1549

471

§ 60.

The Geneva Consensus, A.D. 1552

474

§ 61.

The Helvetic Consensus Formula, A.D. 1675

477

\emph{II. Reformed Confessions of France and the Netherlands.}

§ 62.

The Gallican Confession, A.D. 1559

490

§ 63.

The French Declaration of Faith, A.D. 1872

498

§ 64.

The Belgic Confession, A.D. 1561

502

§ 65.

The Arminian Controversy and the Synod of Dort, A.D. 1604-1619.

508

§ 66.

The Remonstrance, A.D. 1610

516

§ 67.

The Canons of Dort, A.D. 1619

519

\emph{III. The Reformed Confessions of Germany.}

§ 68.

The Tetrapolitan Confession, A.D. 1530

524

§ 69.

The Heidelberg Catechism, A.D. 1563

529

§ 70.

The Brandenburg Confessions

554

The Confession of Sigismund (1614), 555.

The Colloquy at Leipzig (1631), 558.

The Declaration of Thorn (1645), 560.

§ 71.

The Minor German Reformed Confessions. . .

563

\emph{IV. The Reformed Confessions of Bohemia, Poland, and Hungary.}

§ 72.

The Bohemian Brethern and the Waldenses before the Reformation

565

§ 73.

The Bohemian Confessions after the Reformation, A.D. 1535 and 1575

576

§ 74.

The Reformation in Poland and the Consensus of Sendomir, A.D. 1570

581

§ 75.

The Reformation in Hungary and the Confession of Czenger, A.D. 1557

589

\emph{V. The Anglican Articles of Religion.}

§ 76.

The English Reformation.

592

§ 77.

The Doctrinal Position of the Anglican Church and her Relation to other Churches

598

§ 78.

The Doctrinal Formularies of Henry VIII.

611

§ 79.

The Edwardine Articles, A.D. 1553.

613

§ 80.

The Elizabethan Articles, A.D. 1563 and 1571

615

§ 81.

Interpretation of the Thirty-nine Articles

622

§ 82.

Revision of the Thirty-nine Articles by the Protestant Episcopal Church in the United States of America, A.D. 1801

650

§ 83.

The Anglican Catechisms, A.D. 1549 and 1662

654

§ 84.

The Lambeth Articles, A.D. 1595

658

§ 85.

The Irish Articles, A.D. 1615

662

§ 86.

The Articles of the Reformed Episcopal Church, A.D. 1875

665

\emph{VI. The Presbyterian Confessions of Scotland.}

§ 87.

The Reformation in Scotland

669

§ 88.

John Knox

673

§ 89.

The Scotch Confession, A.D. 1560.

680

§ 90.

The Scotch Covenants and the Scotch Kirk

685

§ 91.

The Scotch Catechisms

696

\emph{VII. The Westminster Standards.}

§ 92.

The Puritan Conflict

701

§ 93.

The Westminster Assembly.

727

§ 94.

The Westminster Confession .

753

§ 95.

Analysis of the Confession .

760

§ 96.

The Westminster Catechisms .

783

§ 97.

Criticism of the Westminster System of Doctrine

788

§ 98.

The Westminster Standards in America

804

§ 99.

The Westminster Standards among the Cumberland Presbyterians

813

EIGHTH CHAPTER.

THE CREEDS OF MODERN EVANGELICAL DENOMINATIONS.

PAGE

§ 100.

General Survey

817

§ 101.

The Congregationalists

820

§ 102.

English Congregational Creeds

829

§ 103.

American Congregational Creeds

835

§ 104.

Anabaptists and Mennonites

840

§ 105.

The Calvinistic Baptists

844

§ 106.

The Arminian Baptists

856

§ 107.

The Society of Friends (Quakers).

859

§ 108.

The Moravians

874

§ 109.

Methodism

882

§ 110.

Methodist Creeds

890

§ 111.

Arminian Methodism

893

§ 112.

Calvinistic Methodism

901

§ 113.

The Catholic Apostolic Church (Irvingites)

905

§ 114.

The Evangelical Alliance.

915

§ 115.

The Consensus and Dissensus of Creeds

919

§ 116.

The Disciples of Christ

930

§ 117.

The Universalists.

933

§ 118.

The Unitarians

954

\xsection{Additions to the Literature}

\emph{In General}

Kattenbusch:  \emph{Lehrbuch der vergleichenden Confessionskunde}, Freib., 1892.—Gumlich: \emph{Christ. Creeds and Conff.}, Engl. trans., N. Y., 1894.—Callows: \emph{Origin and Development of Creeds}, London, 1899. S. G. Green: \emph{The Christ. Creed and the Creeds of Christendom}, N. Y., 1899.—Skrine: \emph{Creed and the Creeds, their Function in Religion}, London, 1911.—W. A. Curtis: \emph{Hist. of Creeds and Conff. of Faith in Christendom and Beyond}, Aberdeen, 1911. An elaboration of the author's art., ``Confessions,'' in Enc. of Rel. and Ethics; includes the principles of Mormonism, Christian Science and Tolstoy.—Hirsch: Art., ``Creeds,'' in Enc. Brit., 14th ed.—The works on \emph{Symbolics} of Loofs, and Briggs, N. Y., 1914.—Hase:  \emph{Hdbook of the Controversy with Rome}, 2 vols., London, 1906, trans. from Hase's \emph{Polemik}, ed. of 1900.—Plitt:  \emph{Grundriss der Symbolen}, 7th ed., by Victor Schultze, Erl., 1921.—Mulert: \emph{Konfessionskunde}, Giessen, 1929.

\emph{Collections of Creeds}

Hahn, 3rd ed. enlarged, 1897.—C. Fabricius, prof. in Berlin, \emph{Corpus confessionum. Die Bekenntnisse des Christenthums. Sammlung grundlegender Urkunden aus allen Kirchen der Gegenwart}, Berlin, 1928 sqq.—J. T. Müller: \emph{Die symb. Bücker der ev. luth. Kirche, deutsch und latein.}, 12th ed., 1928.— E. F. Karl Müller: \emph{Die Bekenntnisschriften der reform. Kirche}, Leip., 1903.—For papal decrees: \emph{Acta, sedis sanctae}, Rome.—Mirbt:  \emph{Quellen zur Gesch. des Papsttums und des röm. Katholizismus}, 4th ed., 1924.—Denzinger: \emph{Enchiridion symbolorum et definitionum, quae a concillis oecum. et summis pontificibus emanarunt}, 17th ed. by Umberg, 1928.

\emph{Page} 12.

A. E. Burn: \emph{Facsimiles of the Creeds}, etc., London, 1899; \emph{Introd. to the Creeds and Te Deum}, London, 1901.—Mortimer: \emph{The Creeds, App., Nic., Athanas.}, London, 1902.—A. Seeberg: \emph{Katechismus der Urchristenheit}, 1903.—Turner: \emph{Hist. and Use of Creeds in the Early Centuries of the Church}, London, 1906.—Bp. E. C. S. Gibson:  \emph{The Three Creeds}, Oxf., 1908.—Wetzer and Welte: \emph{Enc}. 2nd ed. V., 676–690.—Loofs: \emph{Symbolik}, pp. 1–70.—Briggs: \emph{Theol. Symbolics}, pp. 34–121.—F. J. Badcock: \emph{The Hist. of Creeds, App., Nic., and Athanas.}, London, 1930, pp. 248.

\emph{Page} 14.

The Apostles Creed: Kattenbusch: \emph{Das apostol. Symbol}, 2 vols., Leipsic, 1894–1900.—Zahn: \emph{Das apostol. Symbol}, Erl., 1893, transl. by Burn from 2nd ed., London, 1899.—Harnack, in Herzog Enc., I, 741–55 and separately in Engl. 1901.—H. B. Swete: \emph{The App. Creed. Its Relation to Prim. Christianity}, Cambr., 1894.—Kunze: \emph{Glaubensregel, hl. Schrift und Taufbekenntniss}, Leipsic, 1899; \emph{Das apostol. Glaubensbekenntniss und das N. T.}, Berlin, 1911, Engl. trans. by Gilmore, N. Y., 1912.— Künstle: \emph{Bibliothek der Symbole}, Mainz, 1901.—A. C. McGiffert: \emph{The App. Creed. Its Origin, Purpose}, etc., N. Y., 1902.—Bp. A. MacDonald (R. C.): \emph{The App. Creed. A Vindication of its Apostol. Authority}, 1903, 2nd ed., London, 1925.—\emph{The App. Creed. Questions of Faith}, Lectures by Denney, Marcus Dods, Lindsay, etc., London, 1904.—Popular treatments by Canon Beeching, 1906; W. R. Richards, N. Y., 1906; Barry, N. Y., 1912; Bp. Bell, 1917, 1919; McFadyen, 1927; H. P. Sloan, N. Y., 1930.—Also Bardenhewer: \emph{Gesch. der altchr. Lit.}, 2nd ed., I, 82-90.

\emph{Page} 24.

The Nicene Creed: Hort: \emph{Two Dissertations on the Constan. Creed}, London, 1876.—Lias: \emph{The Nicene Creed}, 1897.—Kunze: \emph{Das nic.-konstant. Symbol}, Leipsic, 1898.—Harnack, in Herzog Enc. XI., 12–27, and Schaff-Herzog, III, 256–260.—Bp. Headlam: \emph{The Nic. Creed}. Noting differences between the Rom. and Angl. Churches.

\emph{Pages} 43–68.

\emph{Die Bekenntnisse und wichtigsten Glaubenszeugnisse der griech.-oriental. Kirche} (thesauros tes orthodoxias) ed., by Michalcescu, with Introd. by Hauck, Leipsic, 1904. Includes creeds and decrees of the first seven œcum. councils.—Loofs: Symbolik, pp. 77–181.—Adeney: \emph{The Gr. and East. Churches}, N. Y., 1908.—Fortescue (R. C.): \emph{The Orthod. East. Church}, last ed., London, 1916.—Langsford-James: \emph{Dict. of the East. Orthod. Church}, London, 1923.—The art. in Herzog, ``Gennadius II,'' ``Jeremias,'' ``Lukaris,'' etc.—Birkbeck:  \emph{The Russ. and Engl. Churches, during the last fifty years}, London, 1895.—Frère: \emph{Links in the Chain of Russ. Ch. Hist.}, London, 1918.

\emph{Page} 69.

Bonewitsch: \emph{Kirchengesch. Russlands}, Leipsic, 1923.—Reyburn: \emph{Story of the Russ. Church}, London, 1924.—Spinka (prof. in Chicago Theol. Seminary): \emph{The Church and the Russ. Revolution}, N. Y., 1927.—Hecker (student of Drew and Union Theol. Seminaries and Prof. of Theol., Moscow): \emph{ Rel. under the Soviets}, N. Y., 1927; \emph{Soviet Russia in the Second Decade}, 1928.—Emhardt: \emph{Rel. in Soviet Russia}, Milwaukee, 1929.—M. Hindus (b. in Russia): \emph{Humanity Uprooted}, N. Y., 1929.—Batsell: \emph{Soviet Rule in Russia}, N. Y., 1930.—The Engl. White Paper, Aug. 12, 1930, which gives a trans. of Soviet regulations ``respecting religion in the Union of Soviet Socialist Republics.''

\emph{Page} 83.

A. Straub (prof. at Innsbruck): \emph{de Ecclesia}, 2 vols., Innsbr., 1912.—Ryan and Millar: \emph{The State and the Church}, N. Y., 1902.—F. Heiler (ex-Cath., prof. in Marburg): \emph{Der Katholizismus, seine Idee und seine Erscheinung}, Munich, 1923.—Döllinger-Reusch: \emph{Selbstbiographie des Kard. Bellarmin}, with notes, 1887.—Card. Gibbons, d. 1921: \emph{The Faith of Our Fathers}, 1875.—The works and biographies of Card. Newman, d. 1890, and Card. Manning, d. 1892.—D. S. Schaff:  \emph{Our Fathers' Faith and Ours}, N. Y., 1928.

\emph{Page} 91.

Buckley, 2 vols., 1852, gives the Reformatory decisions of the council as well as the Decrees and Canons.—Donovan:  \emph{Profession and Catechism of the C. of Trent}, 1920 and since.—Mirbt: \emph{Quellen zur Gesch. des Papsttums}. Gives large excerpts from the Tridentine standards.—Froude: \emph{Lectures on the C. of Trent}, 1896.—Pastor: \emph{Gesch. der Päpste}, vol. vii.—The \emph{Ch. Histories} of Hergenröther-Kirsch, Funk, etc.

\emph{Page} 134.

Mirbt, pp. 456–466.—Shotwell-Loomis: \emph{The See of St. Peter.} Trans. of patristic documents, N. Y., 1927.—Granderath, S.J.: \emph{Gesch. des Vat. Konzils}, ed. by Kirch, 3 vols., Freib. in Breis., 1903.—Döllinger-Friedrich: \emph{Das Papsttum}, 1892.—Lord Acton: \emph{The Vatican Council} in ``Freedom of Thought.''—Pastor: \emph{Hist. of the Popes}, vol. x. for Sixtus V.'s ed. of the Vulgate.—Card. Gibbons (a member of the council): \emph{Retrospect of Fifty Years}, 2 vols., 1906.—The biographies of Manning by Purcell, 2 vols., 1896; Ketteler by Pfulf, 3 vols., Mainz, 1899; Newman by Ward, 4 vols., 1912.— Straub: \emph{de Ecclesia}, vol. ii., 358–394.—Nielsen: \emph{The Papacy in the 19th Cent.}, vol. ii., pp. 290–374.— Koch: \emph{Cyprian und das röm. Primat}, 1910.—Schnitzer: \emph{Hat Jesus das Papstthum gestiftet?} and \emph{Das Papstthum keine Stiftung Jesus}, 1910.—Count von Hoensbroech (was sixteen years a Jesuit, d. 1923): \emph{Das Papstthum in social-kult. Wirsamkeit}, 3 vols., 4th ed., 1903.—Lietzmann: \emph{Petrus und Paulus in Rom}, 2nd ed., 1927.—Koch: \emph{Cathedra Petri} (dedicated to Schnitzer), Giessen, 1930.

\emph{Page} 220.

H. E. Jacobs: \emph{The Book of Concord or the Symbol. Books of the Ev. Luth. Church}, 2 vols., Phil., 1882, 1912.—\emph{The Luth. Cyclopedia} by Jacobs and Haas, Phil., 1899.—\emph{Concordia Cyclopedia}, 3 vols., 1927.— Schmid: \emph{The Doctr. Theol. of the Ev. Luth. Ch.}, trans. by Hay and Jacobs, 3rd ed., Phil., 1899.—\emph{Luther's Primary Writings}, trans. by Buchheim and Wace, 1896.—\emph{Luther's Works}, Engl. trans., 2 vols., Phil., 1915.—\emph{Luther's Correspondence}, trans. by P. Smith and Jacobs, 2 vols., 1913-1918.—\emph{Lives of Luther} by Schaff in ``Hist. of Chr. Church,'' vol. vi.; Jacobs, 1898; Lindsay in ``Hist. of the Reformation,'' 1906; Preserved Smith, 1911; McGiffert, 1914; Boehmer, Engl. trans., 1916; Mackinnon, 4 vols., 1925-1930; Denifle (R. C.), 2 vols., 2nd ed., 1904; Grisar (R. C.), Engl. trans., 3 vols., 1911, 1912.—P. Smith: \emph{Age of the Reformation}, 1920.—Döllinger: \emph{Akad. Vorträge}, vol. i, 1872. Written after his repudiation of the dogma of Infallibility.

\emph{Page} 225.

Editions of the Augsb. Conf. in Latin and German texts by Kolde,Gotha, 1896, 1911 and Wendt, Halle, 1927.—Ficker: \emph{Konfutation des Augsb. Bekenntnisses}, Leipsic, 1892.—A number of publications bearing on the Augsb. Confession were issued in connexion with the quadricentennial of the Confession's appearance, 1930.

\emph{Page} 354.

Zwingli: \emph{Sämmtliche Werke}, ed. by Egli, Köhler, etc., 1904, sqq.—Karl Müller: \emph{Die Bekenntnissschriften der reformirten Kirche}, Leipsic, 1903. Contains documents not given by Schaff, as Calvin's Genevan Catechism, pp. 117–158; Hungar. Conf. of 1562, pp. 376–448; the Larger Westminster Cat., pp. 612–643; the Nassau Cat. of 1578, pp. 720–738, and the Hesse Cat. of 1607, pp. 822–833.—\emph{Lives of Zwingli} by Stähelin, 2 vols., Basel, 1897; S. M. Jackson, N. Y., 1901. Also Selections from Zwingli, Phil., 1901.—S. Simpson, N. Y., 1902; Egli in Herzog Encycl., vol. xxi.—Humbel: \emph{Zwingli im Spiegel der gleichzeit. schweizer. Lit.}, 1912.

\emph{Page} 388.

Art., ``Bullinger,'' by Egli in Herzog Encycl., vol. iii., pp. 536–549.—Bullinger: \emph{Diarium}, ed. by Egli, Basel, 1904, and \emph{Gegensatz der ev. und röm. Lehre}, ed. by Kügelgen, 1906.—Art., ``Helvetische Konfessionen,'' by Karl Müller in Herzog Encycl., vol. vii and ``Helvetische Konfessionsformeln'' by Egli, vol. vii.

\emph{Page} 421.

Choisy: \emph{L’état chr. à Génève au temps de Th. de Bèze}, Paris, 1903.—Borgeau: \emph{Hist. de l’université de Génève}, Paris, 1903.—\emph{Lives of Calvin} by Schaff in ``Hist. of Chr. Ch.,'' vol. vii.; Kampfschulte, ed. by Goetz, 2 vols., 1899; Doumergue, 7 vols., Lausanne, 1899–1927; W. W. Walker, N. Y., 1906; Reyburn, London, 1914; Lindsay in ``Hist. of the Reformation,'' vol. ii.

\emph{Page} 502.

The Works of B. B. Warfield, Oxf., 1928 sqq.

\emph{Page} 565.

Workman and Pope: \emph{Letters of J. Hus}, London, 1904.—\emph{Lives of Huss} by Count Lützow, London, 1909; D. S. Schaff, N. Y., 1915, and Huss' \emph{de Ecclesia}, trans. with Notes, N. Y., 1915.—Kitts: \emph{John XXIII. and J. Hus}, London, 1910. Müller in \emph{Bekenntnisschriften} gives in full the Hungar. Confessions and the Bohem. Conf . of 1609.

\emph{Page} 568.

\emph{The Nobla Leycon}, with Notes, ed., by Stefano, Paris, 1909.—Comba, father and son: \emph{Hist., of the Waldenses in Italy}, Engl. trans. 1889; \emph{Storia dei Valdesi}, 1893.—Jalla:  \emph{Hist. des Vaudois}, Torre Pelice, 1904.

\emph{Page} 589.

Balogh: \emph{Hist. of the Ref. Ch. in Hungary} in Ref . Ch. Rev., July, 1906.

\emph{Page} 592.

\emph{Use of Sarum}, ed. from MSS. by Frère, 2 vols., Cambr., 1898-1901.—Gee and Hardy: \emph{Documents Illustr. of Engl. Ch. Hist.}—Prothero: \emph{Select Statutes of Elizabeth and James I.}—H. E. Jacobs:  \emph{The Luth. Ch. Movement in Engl.}, Phil., 1870, 1891.—Lindsay: \emph{Hist. of the Reformation}, vol. ii., pp. 315–418.—\emph{The Hist. of the Engl. Ch.} from Henry VIII. to Mary's Death by Gairdner and under Elizabeth and James I. by Frère, 1902, 1904.—Pollard: \emph{Henry VIII.}, London, 1902, \emph{Thos. Cranmer}, 1904; \emph{Wolsey}, 1929.

\emph{Page} 650.

Tiffany: \emph{Hist. of the Prot. Bp. Ch.}, N. Y., 1895.—Hodges: \emph{Three Hundred Years of the Ep. Ch. in Am.}, Phil., 1907.—Cross: \emph{The Angl. Episcopate and the Am. Colonies}, N. Y., 1902.

\emph{Page} 669.

\emph{Histories of the Scotch Reformation} by Mitchell, 1900; Fleming, 1904, 1910; MacEwan, 1913.—\emph{Lives of Knox} by Cowan, 1905; P. H. Brown, 1905.—A. Lang: \emph{J. Knox and the Reformation}, 1905.

\emph{Pages} 701, 820, 835.

H. M. Dexter: \emph{The Congregationalists of the Last 300 Years}, N. Y., 1880.—W. W. Walker: \emph{Creeds and Platforms of Congregationalism}, N. Y., 1893; \emph{Hist. of the Cong. Churches in the U. S.}, N. Y., 1894.—J. Brown: \emph{The Engl. Puritans}, London, 1910.—R. C. Usher: \emph{Reconstruction of the Engl. Ch.}, 2 vols., London, 1910.—W. Selbie:  \emph{Engl. Sects. Congregationalism}, London, 1922.—\emph{Orig. Narratives of Early Am. Hist.}, ed. by Jamieson, N. Y., 1908, sqq.—W. E. Barton:  \emph{Congr. Creeds and Covenants}, Chicago, 1917.

\emph{Page} 813.

McDonnold: \emph{Hist. of the Cumber. Presb. Ch.}, Nashville, 1888.—Miller: \emph{Doctr. of the Cumberl. Presb. Ch.}, Nashville, 1892.

\emph{Page} 840.

Vedder: \emph{Balthazar Hübmaier}, N. Y., 1903.—Newman: \emph{Hist. of the Bapt. Chh. in the U. S.}, N. Y., 1894.—Underhill:  \emph{Conff. of Faith of the Bapt. Chh. in England in the 17th Century}, London, 1854.—McGlothlin: \emph{Bapt. Conff. of Faith}, Phil., 1911.—Carroll: \emph{Baptists and their Doctrines}, N. Y., 1913.

\emph{Page} 859.

Thomas: \emph{Hist. of the Soc. of Friends}, in ``Am. Ch. Hist. Series,'' N. Y., 1894.—Sharpless: \emph{Hist. of Quaker Govt. in Pa.}, 2 vols., Phil., 1898.—R. M. Jones: \emph{The Quakers in the Am. Colonies}, London, 1911; The \emph{Faith and Practice of the Quakers}, 1927.—Holder: \emph{The Quakers in Great Britain and Am.}, N. Y., 1913.

\emph{Page} 874.

Hamilton:  \emph{Hist. of the Morav. Ch.}, Bethlehem, 1900. Also in ``Am. Ch. Hist. Series.''

\emph{Page} 882.

\emph{The Journal of John Wesley}, 8 vols., ed. by Curnock, London, 1910.—Buckley: \emph{Hist. of the Methodists in the U. S.}, N. Y., 1896.—E. S. Tipple: \emph{The Heart of Asbury's Journal}, N. Y., 1905.—Simon: \emph{Revival of Rel. in England in the 18th Cent.}, London, 1907.—Lidgett and Reed: \emph{Methodism in the Modern World}, London, 1929.—Rattenburg: \emph{Wesley's Legacy to the World}, London, 1930.—Allen: \emph{Methodism and Modern World Problems}, London, 1930.—Lunn:  \emph{J. Wesley}, London, 1929.—\emph{Lives of Asbury}, by Tipple, N. Y., 1916, and J. Lewis, 1927.

\xchapter{Chapter 1. Of Creeds in General}

General Literature.

Wm. Dunlop (Prof. of Church Hist. at Edinburgh, d. 1720): Account of all the Ends and Uses of Creeds and Confessions of Faith, a Defense of their Justice, Reasonableness, and Necessity as a Public Standard of Orthodoxy, 2d ed. Lond. 1724. Preface to [Dunlop's] Collection of Confessions in the Church of Scotland, Edinb. 1719 sq. Vol. I. pp. v.–cxlv.

J. Caspar Köcher: Bibliotheca theologiæ symbolicæ et catechetiæ itemque liturgicæ, Wolfenb. and Jena, 1761–69, 2 parts, 8vo.

Charles Butler (R.C., d. 1832): An Historical and Literary Account of the Formularies, Confessions of Faith, or Symbolic Books of the Roman Catholic, Greek, and principal Protestant Churches. By the Author of the Horæ Biblicæ, London, 1816 (pp. 200).

Charles Anthony Swainson  (Prof. at Cambridge and Canon of Chichester): The Creeds of The Church in their Relations to the Word of God and to the Conscience of the Individual Christian (Hulsean Lectures for 1857), Cambridge, 1858.

Francis Chaponnière  (University of Geneva): La Question des Confessions de Foi au sein du Protestantisme contemporain, Genève, 1867. (Pt. I. Examen des Faits. Pt II. Discussion des Principes.)

Karl Leohler:  Die Confessionen in ihrem Verhältniss zu Christus, Heilbronn, 1877.

The introductions to the works on Symbolics by  Marheineke,    Winer,    Möhler,    Köllner,    Gunricke,    Matthes,    Hofmann,    Oehler,  contain some account of symbols, as also the Prolegomena to the Collections of the Symbols of the various Churches by  Walch,    Müller,    Niemeyer,    Kimmel,  etc., which will be noticed in their respective places below.

\xsection{Name and Definition}

§ 1. Name and Definition.

A Creed,\footnote{\par From the beginning of the Apostles' Creed (\emph{Credo}, \emph{I believe}), to which the term is applied more particularly.} or Rule of Faith,\footnote{\par \textgreek{ Κανών τῆς πίστεως } or \textgreek{ τῆς άληθείας, }  \emph{regula fidei, regula veritatis.}  These are the oldest terms used by the ante-Nicene fathers, Irenæus, Tertullian, etc.} or Symbol,\footnote{\par  \textgreek{ Σύμβολον, }  \emph{symbolum}  (from \textgreek{ συμβάλλειν, } to throw together, to compare), means a mark, badge, watchword, test. It was first used in a theological sense by Cyprian, A.D. 250 (Ep. 76, al. 69, ad Magnum, where it is said of the schismatic Novatianus, '\emph{eodum }symbolo,\emph{ quo et nos, baptizare}'), and then very generally since the fourth century. It was chiefly applied to the Apostles' Creed as the baptismal confession by which Christians could be known and distinguished from Jews, heathen, and heretics, in the sense of a military signal or watchword (\emph{tessera militaris}); the Christians being regarded as soldiers of Christ fighting under the banner of the cross. Ambrose (d. 397) calls it ' \emph{cordis signaculum et nostræ militiæ sacramentum.} ' Rufinus, in his \emph{Expositio in Symb. Apost.}, uses the word likewise in the military sense, but gives it also the meaning  \emph{collatio,}   \emph{contributio}  (confounding \textgreek{ σύμβολον } with \textgreek{ συμβολη}), with reference to the legend of the origin of the creed from contributions of the twelve apostles (' \emph{quod plures in unum conferunt; id enim fecerunt apostoli,} ' etc.). Others take the word in the sense of a compact, or agreement (so Suicer, Thes. eccl. II. 1084: ' \emph{Dicere possumus, symbolum non a militari, sed a contractuum tessera nomen id accepisse; est enim tessera pacti, quod in baptismo inimus cum Deo} '). Still others derive it (with King, History of the Apostles' Creed, p. 8) from the signs of recognition among the heathen in their mysteries. Luther and Melancthon first applied it to Protestant creeds. A distinction is made sometimes between \emph{Symbol} and \emph{Symbolical Book}, as also between  \emph{symbola publica}  and  \emph{symbola privata.}  The term  \emph{theologia symbolica}  is of more recent origin than the term  \emph{libri symbolici.}} is a confession of faith for public use, or a form of words setting forth with authority certain articles of belief, which are regarded by the framers as necessary for salvation, or at least for the well-being of the Christian Church.

A creed may cover the whole ground of Christian doctrine and practice, or contain only such points as are deemed fundamental and sufficient, or as have been disputed. It may be declarative, or interrogative in form. It may be brief and popular (as the Apostles' and the Nicene Creeds), for general use in catechetical instruction and at baptism; or more elaborate and theological, for ministers and teachers, as a standard of public doctrine (the symbolical books of the Reformation period). In the latter case a confession of faith is always the result of dogmatic controversy, and more or less directly or indirectly polemical against opposing error. Each symbol bears the impress of its age, and the historical situation out of which it arose.

There is a development in the history of symbols. They assume a more definite shape with the progress of biblical and theological knowledge. They are mile-stones and finger-boards in the history of Christian doctrine. They embody the faith of generations, and the most valuable results of religious controversies. They still shape and regulate the theological thinking and public teaching of the churches of Christendom. They keep alive sectarian strifes and antagonisms, but they reveal also the underlying agreement, and foreshadow the possibility of future harmony.

\xsection{Origin of Creeds}

§ 2. Origin of Creeds.

Faith, like all strong conviction, has a desire to utter itself before others—'Out of the abundance of the heart the mouth speaketh;' 'I believe, therefore I confess' (\emph{Credo, ergo confiteor}). There is also an express duty, when we are received into the membership of the Christian Church, and on every proper occasion, to profess the faith within us, to make ourselves known as followers of Christ, and to lead others to him by the influence of our testimony.\footnote{\par Comp. (Matt. x. 32, 33: 'Every one who shall confess me before men, him will I also confess before my Father who is in heaven. But whosoever shall deny me before men, him will I also deny before my Father who is in heaven.' Rom. x. 9, 10: 'If thou shalt \emph{confess} with thy mouth the Lord Jesus [Jesus as Lord], and shalt believe in thine heart that God hath raised him from the dead, then shalt be saved. For with the heart man believeth unto [so as to obtain] righteousness; and with the mouth \emph{confession} is made unto salvation.'}

This is the origin of Christian symbols or creeds. They never precede faith, but presuppose it. They emanate from the inner life of the Church, independently of external occasion. There would have been creeds even if there had been no doctrinal controversies.\footnote{\par Semisch, Das apostolische Glaubensbekenntniss (Berlin, 1872, p. 7): ' \emph{Bekenntnisse, an welchen sich das geistige Leben ganzer Völker auferbaut, welche langen Jahrhunderten die höchsten Ziele und bestimmenden Kräfte ihres Handelns vorzeichnen, sind nicht Noth- und Flickwerke des Augenblicks . . . es sind Thaten des Lebens, Pulsschläge der sich selbst bezeugenden Kirche.} '} In a certain sense it may be said that the Christian Church has never been without a creed (\emph{Ecclesia, sine symbolis nulla}). The baptismal formula and the words of institution of the Lord's Supper are creeds; these and the confession of Peter antedate even the birth of the Christian Church on the day of Pentecost. The Church is, indeed, not founded on symbols, but on Christ; not on any words of man, but on the word of God; yet it is founded on Christ as \emph{confessed} by men, and a creed is man's answer to Christ's question, man's acceptance and interpretation of God's word. Hence it is after the memorable confession of Peter that Christ said, 'Thou art Rock, and upon this rock I shall build my Church,' as if to say, 'Thou art the Confessor of Christ, and on this Confession, as an immovable rock, I shall build my Church.' Where there is faith, there is also profession of faith. As 'faith without works is dead,' so it may be said also that faith without confession is dead.

But this confession need not always be written, much less reduced to a logical formula. If a man can say from his heart, 'I believe in the Lord Jesus Christ,' it is sufficient for his salvation (Acts xvi. 31). The word of God, apprehended by a living faith, which founded the Christian Church, was at first orally preached and transmitted by the apostles, then laid down in the New Testament Scriptures, as a pure and unerring record for all time to come. So the confession of faith, or the creed, was orally taught and transmitted to the catechumens, and professed by them at baptism, long before it was committed to writing. As long as the  \emph{Disciplina arcani}  prevailed, the summary of the apostolic doctrine, called 'the rule of faith,' was kept confidential among Christians, and withheld even from the catechumens till the last stage of instruction; and hence we have only fragmentary accounts of it in the writings of the ante-Nicene fathers. When controversies arose concerning the true meaning of the Scriptures, it became necessary to give formal expression of their true sense, to regulate the public teaching of the Church, and to guard it against error. In this way the creeds were gradually enlarged and multiplied, even to the improper extent of theological treatises and systems of divinity.

The first Christian confession or creed is that of Peter, when Christ asked the apostles, 'Who say ye that I am?' and Peter, in the name of all the rest, exclaimed, as by divine inspiration, 'Thou art the Christ, the Son of the living God' (Matt. xvi. 16).\footnote{\par The similar confession, John vi. 69, is of a previous date. It reads, according to the early authorities, 'Thou art the Holy One of God' (\textgreek{σὺ εἶ ὁ ἅγιος θεοῦ}). A designation of the Messiah. This text coincides with the testimony of the demoniacs, Marc. I. 26, who, with ghostlike intuition, perceived the supernatural character of Jesus.} This became naturally the substance of the baptismal confession, since Christ is the chief object of the Christian faith. Philip required the eunuch simply to profess the belief that 'Jesus was the Son of God.' In conformity with the baptismal formula, however, it soon took a Trinitarian shape, probably in some such simple form as 'I believe in God the Father, the Son, and the Holy Spirit.' Gradually it was expanded, by the addition of other articles, into the various rules of faith, of which the Roman form under the title 'the Apostles' Creed' became the prevailing one, after the fourth century, in the West, and the Nicene Creed in the East. The Protestant Church, as a separate organization, dates from 1517, but it was not till 1530 that its faith was properly formularized in the Augsburg Confession.

A symbol may proceed from the general life of the Church in a particular age without any individual authorship (as the Apostles' Creed); or from an œcumenical Council (the Nicene Creed; the Creed of Chalcedon); or from the Synod of a particular Church (the Decrees of the Council of Trent; the Articles of Dort; the Westminster Confession and Catechisms); or from a number of divines commissioned for such work by ecclesiastical authority (the Thirty-nine Articles of the Church of England; the Heidelberg Catechism; the Form of Concord); or from one individual, who acts in this case as the organ of his church or sect (the Augsburg Confession, and Apology, composed by Melancthon; the Articles of Smalkald, and the Catechisms of Luther; the second Helvetic Confession by Bullinger). What gives them symbolical or authoritative character is the formal sanction or tacit acquiescence of the church or sect which they represent. In Congregational and Baptist churches the custom prevails for each local church to have its own confession of faith or 'covenant,' generally composed by the pastor, and derived from the Westminster Confession, or some other authoritative symbol, or drawn up independently.

\xsection{Authority of Creeds}

§ 3. Authority of Creeds.\footnote{\par On the authority and use of Symbols there are a number of Latin and German treatises by C. U. Hahn (1833), Hoefling (1835), Sartorius (1845), Harless (1846), A. Hahn 1847), Köllner (1847), Genzken (1851), Bretschneider (1830), Johannsen (1833), and others, all with special reference to the Lutheran State Churches in Germany. See the literature in Müller, Die symb. Bücher der evang. luth. Kirche, p. xv., and older works in Winer's Handbuch der theol. Literatur, 3d ed. Vol. I. p. 334. Comp. also Dunlop and Chaponnière (Part II.), cited in § 1.}

1. In the Protestant system, the authority of symbols, as of all human compositions, is relative and limited. It is not co-ordinate with, but always subordinate to, the Bible, as the only infallible rule of the Christian faith and practice. The value of creeds depends upon the measure of their agreement with the Scriptures. In the best case a human creed is only an approximate and relatively correct exposition of revealed truth, and may be improved by the progressive knowledge of the Church, while the Bible remains perfect and infallible. The Bible is of God; the Confession is man's answer to God's word.\footnote{\par For this reason a creed ought to use language different from that of the Bible. A string of Scripture passages would be no creed at all, as little as it would be a prayer or a hymn. A creed is, as it were, a doctrinal poem written under the inspiration of divine truth. This may be said at least of the œcumenical creeds.} The Bible is the \emph{norma normans}; the Confession the \emph{norma normata.} The Bible is the rule of \emph{faith} (\emph{regula fidei}); the Confession the rule of \emph{doctrine} (\emph{regula doctrinæ}). The Bible has, therefore, a divine and absolute, the Confession only an ecclesiastical and relative authority. The Bible regulates the general religious belief and practice of the laity as well as the clergy; the symbols regulate the public teaching of the officers of the Church, as Constitutions and Canons regulate the government, Liturgies and Hymn-books the worship, of the Church.

Any higher view of the authority of symbols is unprotestant and essentially Romanizing. Symbololatry is a species of idolatry, and substitutes the tyranny of a printed book for that of a living pope. It is apt to produce the opposite extreme of a rejection of all creeds, and to promote rationalism and infidelity.

2. The Greek Church, and still more the Roman Church, regarding the Bible and tradition as two co-ordinate sources of truth and rules of faith, claim absolute and infallible authority for their confessions of faith.\footnote{\par Tertullian already speaks of the  \emph{regula fidei immobilis et irreformabilis}  (\emph{De virg. vel.} c. 1); but he applied it only to the simple form which is substantially retained in the Apostles' Creed.}

The Greek Church confines the claim of infallibility to the seven œcumenical Councils, from the first Council of Nicæa, 325, to the second of Nicæa, 787.

The Roman Church extends the same claim to the Council of Trent and all the subsequent official Papal decisions on questions of faith down to the decree of the Immaculate Conception in 1854, and the dogma of Papal Infallibility proclaimed by the Vatican Council in 1870. Since that time the Pope is regarded by orthodox Romanists as the organ of infallibility, and all his official decisions on matters of faith and morals must be accepted as final, without needing the sanction of an œcumenical council.

It is clear that either the Greek or the Roman Church, or both, must be wrong in this claim of infallibility, since they contradict each other on some important points, especially the authority of the pope, which in the Roman Church is an  \emph{articulus stantis et cadentis ecclesiæ,}  and is expressly taught in the Creed of Pius V. and the Vatican Decrees.

\xsection{Value and Use of Creeds}

§ 4. Value and Use of Creeds.

Confessions, in due subordination to the Bible, are of great value and use. They are summaries of the doctrines of the Bible, aids to its sound understanding, bonds of union among their professors, public standards and guards against false doctrine and practice. In the form of Catechisms they are of especial use in the instruction of children, and facilitate a solid and substantial religious education, in distinction from spasmodic and superficial excitement. The first object of creeds was to distinguish the Church from the world, from Jews and heathen, afterwards orthodoxy from heresy, and finally denomination from denomination. In all these respects they are still valuable and indispensable in the present order of things. Every well-regulated society, secular or religious, needs an organization and constitution, and can not prosper without discipline. Catechisms, liturgies, hymn-books are creeds also as far as they embody doctrine.

There has been much controversy about the degree of the binding force of creeds, and the  \emph{quia}  or  \emph{quatenus}  in the form of subscription. The whole authority and use of symbolical books has been opposed and denied, especially by Socinians, Quakers, Unitarians, and Rationalists. It is objected that they obstruct the free interpretation of the Bible and the progress of theology; that they interfere with the liberty of conscience and the right of private judgment; that they engender hypocrisy, intolerance, and bigotry; that they produce division and distraction; that they perpetuate religious animosity and the curse of sectarianism; that, by the law of reaction, they produce dogmatic indifferentism, skepticism, and infidelity; that the symbololatry of the Lutheran and Calvinistic State Churches in the seventeenth century is responsible for the apostasy of the eighteenth.\footnote{\par These objections are noticed and answered at length by Dunlop, in his preface to the Collection of Scotch Confessions, and in the more recent works quoted on p. 7.} The objections have some force in those State Churches which allow no liberty for dissenting organizations, or when the creeds are virtually put above the Scriptures instead of being subordinated to them. But the creeds, as such, are no more responsible for abuses than the Scriptures themselves, of which they profess to be merely a summary or an exposition. Experience teaches that those sects which reject all creeds are as much under the authority of a traditional system or of certain favorite writers, and as much exposed to controversy, division, and change, as churches with formal creeds. Neither creed nor no-creed can be an absolute protection of the purity of faith and practice. The best churches have declined or degenerated; and corrupt churches may be revived and regenerated by the Spirit of God, and the Word of God, which abides forever.

\xsection{Classification of Creeds}

§ 5. Classification of Creeds.

The Creeds of Christendom may be divided into four classes, corresponding to the three main divisions of the Church, the Greek, Latin, and Evangelical, and their common parent. A progressive growth of theology in different directions can be traced in them.

1. The Œcumenical Symbols of the Ancient Catholic Church. They contain chiefly the orthodox doctrine of God and of Christ, or the fundamental dogmas of the Holy Trinity and the Incarnation. They are the common property of all churches, and the common stock from which the later symbolical books have grown.

2. The Symbols of the Greek or Oriental Church, in which the Greek faith is set forth in distinction from that of the Roman Catholic and the evangelical Protestant Churches. They were called forth by the fruitless attempts of the Jesuits to Romanize the Greek Church, and by the opposite efforts of the crypto-Calvinistic Patriarch Cyrillus Lucaris to evangelize the same. They differ from the Roman Creeds mainly in the doctrine of the procession of the Holy Spirit, and the more important doctrine of the Papacy; but in the controversies on the rule of faith, justification by faith, the church and the sacraments, the worship of saints and relics, the hierarchy and the monastic system, they are much more in harmony with Romanism than with Protestantism.

3. The Symbols of the Roman Church, from the Council of Trent to the Council of the Vatican (1563 to 1870). They sanction the distinctive doctrines of Romanism, which were opposed by the Reformers, and condemn the leading principles of evangelical Protestantism, especially the supreme authority of the Scriptures as a sufficient rule of faith and practice, and justification by faith alone. The last dogma, proclaimed by the Vatican Council in 1870, completes the system by making the official infallibility of the Pope an article of the Catholic faith (which it never was before).

4. The Symbols of the Evangelical Protestant Churches. Most of them date from the period of the Reformation (some from the seventeenth century), and thus precede, in part, the specifically Greek and Latin confessions. They agree with the primitive Catholic Symbols, but they ingraft upon them the Augustinian theory of sin and grace, and several doctrines in anthropology and soteriology (\emph{e.g.}, the doctrine of atonement and justification), which had not been previously settled by the Church in a conclusive way. They represent the progress in the development of Christian theology among the Teutonic nations, a profounder understanding of the Holy Scriptures (especially the Pauline Epistles), and of the personal application of Christ's mediatorial work.

The Protestant Symbols, again, are either Lutheran or Reformed. The former were all made in Germany from A.D. 1530 to 1577; the latter arose in different countries—Germany, Switzerland, France, Holland, Hungary, Poland, England, Scotland, wherever the influence of Zwingli and Calvin extended. The Lutheran and Reformed confessions agree almost entirely in their theology, christology, anthropology, soteriology, and eschatology, but they differ in the doctrines of divine decrees and of the nature and efficacy of the sacraments, especially the mode of Christ's presence in the Lord's Supper.

The later evangelical denominations, as the Congregationalists, Baptists, Quakers, Arminians, Methodists, Moravians, acknowledge the leading doctrines of the Reformation, but differ from Lutheranism and Calvinism in a number of articles touching anthropology, the Church, and the sacraments, and especially on Church polity and discipline. Their creeds are modifications and abridgments rather than enlargements of the old Protestant symbols.

The heretical sects connected with Protestantism mostly reject symbolical books altogether, as a yoke of human authority and a new kind of popery. Some of them set aside even the Scriptures, and make their own reason or the spirit of the age the supreme judge and guide in matters of faith; but such loose undenominational denominations have generally no cohesive power, and seldom outlast their founders.

The denominational creed-making period closed with the middle of the seventeenth century, except in the Roman Church, which has quite recently added two dogmas to her creed, viz., the Immaculate Conception of the Virgin Mary (1854), and the Infallibility of the Bishop of Rome (1870).

If we are to look for any new creed, it will be, we trust, a creed, not of disunion and discord, but of union and concord among the different branches of Christ's kingdom.

\xchapter{Chapter 2. The Œcumenical Creeds.}

Literature on the three Œcumenical Creeds.

Gerh. Joan. Voss (Dutch Reformed, b. near Heidelberg 1577, d. at Amsterdam 1649): De tribus Symbolis, Apostolico, Athanasiano, et Constantinopolitano. Three dissertations. Amst. 1642 (and in Vol. VI. of his Opera, Amst. 1701). Voss was the first to dispute and disprove the apostolic authorship of the Apostles', and the Athanasian authorship of the Athanasian Creed.

James Ussher (Lat. Usserius, Protestant Archbishop of Armagh, d. 1655): De Romanæ ecclesiæ Symbolo Apostolico vetere, aliisque fidei formulis, tum ab Occidentalibus tum ab Orientalibus in prima catechesi et baptismo proponi solitis, Lond. 1647 (also Geneva, 1722; pp. 17 fol., and whole works in 16 vols., Dublin, 1847, Vol. VII. pp. 297 sq. I have used the Geneva ed.).

Jos. Bingham (Rector of Havant, near Portsmouth, d. 1723): Origines Ecclesiastici; or the Antiquities of the Christian Church (first publ. 1710–22 in 10 vols., and often since in Engl. and in the Latin transl. of Grischovius), Book X. ch. 4.

C. G. P. Walch (a Lutheran, d. at Göttingen in 1784): Bibliotheca Symbolica vetus, Lemgo, 1770. (A more complete collection than the preceding ones, but defective in the texts.)

E. Köllner: Symbolik aller christlichen Confessionen, Hamburg, 1837 sqq., Vol. I. pp. 1–92.

Aug. Hahn: Bibliothek der Symbole und Glaubensregeln der Apostolisch-katholischen Kirche, Breslau, 1842. A new and revised ed. by Ludwig Hahn, Breslau, 1877 (pp. 300).

W. Harvey: History and Theology of the Three Creeds, Cambridge, 1856, 2 vols.

Charles A. Heurtley (Margaret Prof. of Divinity, Oxford): Harmonia Symbolica: A Collection of Creeds belonging to the Ancient Western Church and to the Mediæval English Church. Oxford, 1858. The same: De fide et Symbolo. Oxon. et Lond. 1869.

C. P. Caspari (Prof. in Christiania): Ungedruckte, unbeachtete und wenig beachtete Quellen zur Geschichte des Taufsymbols und der Glaubensregel. Christiania, 1866 to 1875, 3 vols.

J. Rawson Lumby (Prof. at Cambridge): The History of the Creeds. Cambridge,1873; 2d ed. London,1880.

C. A. Swainson (Prof. of Divinity, Cambridge): The Nicene and Apostles' Creeds. Their Literary History; together with an Account of the Growth and Reception of 'the Creed of St. Athanasius.' Lond. 1875.

F. John Anthony Hort (Prof. in Cambridge): Two Dissertations on \textgreek{μονογενὴς θεός} and on the 'Constantinopolitan' Creed and other Eastern Creeds of the Fourth Century. Cambridge and London, 1876.

\xsection{General Character of the Œcumenical Creeds.}

§ 6. General Character of the Œcumenical Creeds.

By œcumenical or general symbols (\emph{symbola œcumenica, s. catholica})\footnote{The term \textgreek{οἰκουμενικός} (from \textgreek{οἰκουμένη,} sc. \textgreek{γῆ, } \emph{orbis terrarum,} the \emph{inhabited earth}; in a restricted sense, the old \emph{Roman Empire}, as embracing the civilized world) was first used in its ecclesiastical application of the general synods of Nicæa (325), Constantinople (381), Ephesus (431), and Chalcedon (451), also of patriarchs, bishops, and emperors, and, at a later period, of the ancient general symbols, to distinguish them from the confessions of particular churches. In the Protestant Church the term so used occurs first in the Lutheran Book of Concord (\emph{œcumenica seu catholica}).} we understand the doctrinal confessions of ancient Christianity, which are to this day either formally or tacitly acknowledged in the Greek, the Latin, and the Evangelical Protestant Churches, and form a bond of union between them.

They are three in number: the Apostles', the Nicene, and the Athanasian Creed. The first is the simplest; the other two are fuller developments and interpretations of the same. The Apostles' Creed is the most popular in the Western, the Nicene in the Eastern Churches.

To them may be added the christological statement of the œcumenical Council of Chalcedon (451). It has a more undisputed authority than the Athanasian Creed (to which the term œcumenical applies only in a qualified sense), but, as it is seldom used, it is generally omitted from the collections.

These three or four creeds contain, in brief popular outline, the fundamental articles of the Christian faith, as necessary and sufficient for salvation. They embody the results of the great doctrinal controversies of the Nicene and post-Nicene ages. They are a profession of faith in the only true and living God, Father, Son, and Holy Ghost, who made us, redeemed us, and sanctifies us. They follow the order of God's own revelation, beginning with God and the creation, and ending with the resurrection of the body and the life everlasting. They set forth the articles of faith in the form of facts rather than dogmas, and are well suited, especially the Apostles' Creed, for catechetical and liturgical use.

The Lutheran and Anglican Churches have formally recognized and embodied the three œcumenical symbols in their doctrinal and liturgical standards.\footnote{The Lutheran Form of Concord (p. 569) calls them '\emph{catholica et generalia summæ auctoritatis symbola}.' The various editions of the Book of Concord give them the first place among the Lutheran symbols. Luther himself emphasized his agreement with them. The Church of England, in the 8th of her 39 Articles, declares, 'The three Creeds, Nicene Creed, Athanasius's Creed, and that which is commonly called the Apostles' Creed, ought thoroughly to be received and believed, for they may be proved by most certain warrants of Holy Scripture.' The American editions of the Articles and of the Book of Common Prayer omit the Athanasian Creed, and the Protestant Episcopal Church of the United States excludes it from her service. The omission by the Convention of 1789 arose chiefly from opposition to the damnatory clauses, which even Dr. Waterland thought might be left out. But the doctrine of the Athanasian Creed is clearly taught in the first five Articles.} The other Reformed Churches have, in their confessions, adopted the trinitarian and christological doctrines of these creeds, but in practice they confine themselves mostly to the use of the Apostles' Creed.\footnote{The Second Helvetic Confession, art. 11, the Gallican Confession, art. 5, and the Belgic Confession, art. 9, expressly approve the three Creeds, 'as agreeing with the written Word of God.' In 'The Constitution and Liturgy' of the (Dutch) Reformed Church in the United States the Nicene Creed and the Athanasian Creed are printed at the end. The Apostles' Creed is embodied in the Heidelberg Catechism, as containing 'the articles of our catholic undoubted Christian faith.' The Shorter Westminster Catechism gives it merely in an Appendix, as 'a brief sum of the Christian faith, agreeable to the Word of God, and anciently received in the churches of Christ.'} This, together with the Lord's Prayer and the Ten Commandments, was incorporated in the Lutheran, the Genevan, the Heidelberg, and other standard Catechisms.

\xsection{The Apostles' Creed.}

§ 7. The Apostles' Creed.

Literature.

I. See the Gen. Lit. on the Œcum. Creeds, § 6, p. 12, especially  Hahn,   Heurtley,   Lumby,   Swainson,  and  Caspari (the third vol. 1875).

II. Special treatises on the Apostles' Creed:

Rufinus (d. at Aquileja 410, a presbyter and monk, translator and continuator of Eusebius's Church History to A.D. 395, and translator of some works of Origen, with unscrupulous adaptations to the prevailing standard of orthodoxy; at first an intimate friend, afterwards a bitter enemy of St. Jerome): Expositio Symboli (Apostolici), first printed, under the name of Jerome, at Oxford 1468, then at Rome 1470, at Basle 1519, etc.; also in the Appendix to John Fell's ed. of Cyprian's Opera (Oxon. 1682, folio, p. 17 sq.), and in \emph{Rufini Opera}, ed. Vallarsi (Ver. 1745). See the list of edd. in Migne's \emph{Patrol.} xxi. 17–20. The genuineness of this Exposition of the Creed is disputed by Ffoulkes, on the \emph{Athanas. Creed}, p. 11, but without good reason.

Ambrosius (bishop of Milan, d. 397): Tractatus in Symbolum Apostolorum (also sub tit. \emph{De Trinitate}). \emph{Opera}, ed. Bened., Tom. II. 321. This tract is by some scholars assigned to a much later date, because it teaches the double procession of the Holy Spirit; but Hahn, l.c. p. 16, defends the Ambrosian authorship with the exception of the received text of the Symbolum Apostolicum, which is prefixed. Also, Explanatio Symboli ad initiandos, ascribed to St. Ambrose, and edited by Angelo Mai in Scriptorum Veterum Nova Collectio, Rom. 1833, Vol. VII. pp. 156–158, and by Caspari, in the work quoted above, II. 48 sq.

Venant. Fortunatus. (d. about 600): Expositio Symboli (Opera, ed. Aug. Luchi, Rom. 1786).

Augustinus. (bishop of Hippo, d. 430): De Fide et Symbolo liber unus. Opera, ed. Bened., Tom. XI. 505–522. Sermo de Symbolo ad catechumenos, Tom. VIII. 1591–1610. Sermones de traditione Symboli, Tom. VIII. 936 sq.

Mos. Amyraldus (Amyraut, Prof. at Saumur, d. 1664): Exercitationes in Symb. Apost. Salmur. 1663.

Isaac Barrow (Master of Trinity College, Cambridge, d. 1677). Sermons on the Creed (Theolog. Works, 8 vols., Oxf. 1830, Vol. IV.–VI).

John Pearson (Bishop of Chester, d. 1686): An Exposition of the Creed, 1659, 3d ed. 1669 fol. (and several later editions by Dobson, Burton, Nichols, Chevallier). One of the classical works of the Church of England.

Peter King (Lord Chancellor of England, d. 1733): The History of the Apostles' Creed, with Critical Observations, London, 1702. (The same in Latin by \emph{Olearius}, Lips. 1706.)

H. Witsius (Prof. in Leyden, d. 1708): Exercitationes sacræ in Symbolum quod Apostolorum dicitur, Amstel. 1700; Basil. 1739. English translation by Fraser, Edinb. 1823, 2 vols.

J. E. Im. Walch (Professor in Jena, d. 1778): Antiquitates symbolicæ, quibus Symboli Apostolici historia illustratur, Jena, 1772, 8vo.

A. G. Rudelbach (Luth.): Die Bedeutung des apost. Symbolums, Leipz. 1844 (78 pp.).

Peter Meyers (R. C.): De Symboli Apostolici Titulo, Origine et Auctoritate, Treviris, 1849 (pp. 210). Defends the apostolic origin.

J. W. Nevin: The Apostles' Creed, in the 'Mercersburg Review,' Mercersburg, Pa., for 1849, pp. 105, 201, 313, 585. An exposition of the doctrinal system of the Creed.

Michel Nicolas: Le symbole des apôtres, Paris, 1867. Rationalistic.

G. Lisco (jun.): Das apostolische Glaubensbekenntniss, Berlin, 1872. In opposition to its obligatory use in the church.

O. Zöckler: Das apostolische Symbolum, Güterslohe, 1872 (40 pp.). In defense of the Creed.

Carl Semisch (Prof. of Church History in Berlin): Das apostolische Glaubensbekenntniss, Berlin, 1872 (31 pp.).

A. Mücke: Das apostolische Glaubensbekenntniss der ächte Ausdruck apostolischen Glaubens, Berlin, 1873 (160 pp.).

The Apostles' Creed, or Symbolum Apostolicum, is, as to its form, not the production of the apostles, as was formerly believed, but an admirable popular summary of the apostolic teaching, and in full harmony with the spirit and even the letter of the New Testament.

I. Character and Value.—As the Lord's Prayer is the Prayer of prayers, the Decalogue the Law of laws, so the Apostles' Creed is the Creed of creeds. It contains all the fundamental articles of the Christian faith necessary to salvation, in the form of facts, in simple Scripture language, and in the most natural order—the order of revelation— from God and the creation down to the resurrection and life everlasting. It is Trinitarian, and divided into three chief articles, expressing faith—in God the Father, the Maker of heaven and earth, in his only Son, our Lord and Saviour, and in the Holy Spirit (\emph{in Deum Patrem, in Jesum Christum, in Spiritum Sanctum}); the chief stress being laid on the second article, the supernatural birth, death, and resurrection of Christ. Then, changing the language (\emph{credo in} for \emph{credo} with the simple accusative), the Creed professes to believe 'the holy Catholic Church, the communion of saints, the remission of sins, the resurrection of the body, and the life everlasting.'\footnote{This change was observed already by Rufinus (l.c. § 36), who says: '\emph{Non dicit} ``In\emph{ Sanctam Ecclesiam},'' \emph{nec} ``In\emph{ remissionem peccatorum},'' \emph{nec} ``In\emph{ carnis resurrectionem.'' Si enim addidisset} ``in''\emph{ præpositionem, una eademque vis fuisset cum superioribus. . . . Hac præpositionis syllaba Creator a creaturis secernitur, et divina separantur ab humanis.}' The Roman Catechism (P. I. c. 10, qu. 19) also marks this distinction, '\emph{Nunc autem, mutata dicendi forma}, ``\emph{sanctam},'' \emph{et non} ``\emph{in sanctam}'' \emph{ecclesiam credere profitemur}.'} It is by far the best popular summary of the Christian faith ever made within so brief a space. It still surpasses all later symbols for catechetical and liturgical purposes, especially as a profession of candidates for baptism and church membership. It is not a logical statement of abstract doctrines, but a profession of living facts and saving truths. It is a liturgical poem and an act of worship. Like the Lord's Prayer, it loses none of its charm and effect by frequent use, although, by vain and thoughtless repetition, it may be made a martyr and an empty form of words. It is intelligible and edifying to a child, and fresh and rich to the profoundest Christian scholar, who, as he advances in age, delights to go back to primitive foundations and first principles. It has the fragrance of antiquity and the inestimable weight of universal consent. It is a bond of union between all ages and sections of Christendom. It can never be superseded for popular use in church and school.\footnote{Augustine calls the Apostolic Symbol '\emph{regula fidei brevis et grandis; brevis numero verborum, grandis pondere sententiarum}.' Luther says: 'Christian truth could not possibly be put into a shorter and clearer statement.' Calvin (\emph{Inst.}, Lib. II. c. 16, § 18), while doubting its strictly apostolic composition, yet regards it as an admirable and truly scriptural summary of the Christian faith, and follows its order in his \emph{Institutes}, saying: '\emph{Id extra controversiam positum habemus, totam in eo} [\emph{Symbolo Ap.}] \emph{fidei nostræ historiam succincte distinctoque ordine recenseri, nihil autem contineri, quod solidis Scripturæ testimoniis non sit consignatum}.' J. T. Müller (Lutheran, \emph{Die Symb. Bücher der Evang. Luth. K.}, p. xvi.): 'It retains the double significance of being the \emph{bond of union} of the universal Christian Church, and the \emph{seed} from which all other creeds have grown.' Dr. Semisch (Evang. United, successor of Dr. Neander in Berlin) concludes his recent essay on the Creed (p. 28) with the words: 'It is in its primitive form the most genuine Christianity from the mouth of Christ himself (\emph{das ächteste Christenthum aus dem Munde Christi selbst}).' Dr. Nevin (Germ. Reformed, \emph{Mercersb. Rev.} 1849, p. 204): 'The Creed is the substance of Christianity in the form of faith . . . the direct immediate utterance of the faith itself.' Dr. Shedd (Presbyterian, \emph{Hist. Christ. Doctr.}, II. 433): 'The Apostles' Creed is the earliest attempt of the Christian mind to systematize the teachings of the Scripture, and is, consequently, the uninspired foundation upon which the whole after-structure of symbolic literature rests. All creed development proceeds from this germ.' Bishop Browne (Episcopalian, \emph{Exp}. 39 \emph{Art}., p. 222): 'Though this Creed was not drawn up by the apostles themselves, it may well be called Apostolic, both as containing the doctrines taught by the apostles, and as being in substance the same as was used in the Church from the times of the apostles themselves.' It is the only Creed used in the baptismal service of the Latin, Anglican, Lutheran, and the Continental Reformed Churches. In the Protestant Episcopal and Lutheran Churches the Apostles' Creed is a part of the regular Sunday service, and is generally recited between the Scripture lessons and the prayers, expressing assent to the former, and preparing the mind for the latter.}

At the same time, it must be admitted that the very simplicity and brevity of this Creed, which so admirably adapt it for all classes of Christians and for public worship, make it insufficient as a regulator of public doctrine for a more advanced stage of theological knowledge. As it is confined to the fundamental articles, and expresses them in plain Scripture terms, it admits of an indefinite expansion by the scientific mind of the Church. Thus the Nicene Creed gives clearer and stronger expression to the doctrine of Christ's divinity against the Arians, the Athanasian Creed to the whole doctrine of the Trinity and of Christ's person against the various heresies of the post-Nicene age. The Reformation Creeds are more explicit on the authority and inspiration of the Scriptures and the doctrines of sin and grace, which are either passed by or merely implied in the Apostles' Creed.

II. As to the origin of the Apostles' Creed, it no doubt gradually grew out of the confession of Peter, Matt. xvi. 16, which furnished its nucleus (the article on Jesus Christ), and out of the baptismal formula, which determined the trinitarian order and arrangement. It can not be traced to an individual author. It is the product of the Western Catholic Church (as the Nicene Creed is that of the Eastern Church) within the first four centuries. It is not of primary, apostolic, but of secondary, ecclesiastical inspiration. It is not a word of God to men, but a word of men to God, in response to his revelation. It was originally and essentially a \emph{baptismal confession}, growing out of the inner life and practical needs of early Christianity.\footnote{Tertullian, \emph{De corona militum}. c. 3: '\emph{Dehinc ter mergitamur}, amplius aliquid respondentes, \emph{quam Dominus in Evangelio determinavit.}' The \emph{amplius respondentes} refers to the Creed, not as something different from the Gospel, but as a summary of the Gospel. Comp. \emph{De bapt}., c. 6, where Tertullian says that in the baptismal Creed the Church was mentioned after confessing the Father, the Son, and the Spirit.} It was explained to the catechumens at the last stage of their preparation, professed by them at baptism, often repeated, with the Lord's Prayer, for private devotion, and afterwards introduced into public service.\footnote{Augustine (\emph{Op.}, ed. Bened., VI. \emph{Serm.}, 58): '\emph{Quando surgitis, quando vos ad somnum collocatis, reddite Symbolum vestrum; reddite Domino. . . . Ne dicatis, Dixi heri, dixi hodie, quotidie dico, teneo illud bene. Commemora fidem tuam: inspice te. Sit tanquam speculum tibi Symbolum tuum. Ibi te vide si credis omnia quæ te credere confiteris, et gaude quotidie in fide tua.}'} It was called by the ante-Nicene fathers 'the rule of faith,' 'the rule of truth,' 'the apostolic tradition,' 'the apostolic preaching,' afterwards 'the symbol of faith.'\footnote{\textgreek{Κανὼν τῆς πίστεως, κ. τῆς ἀληθείας, παράδοσις ἀποστολική, τό ἀρχαῖον τῆς ἐκκλησίας, σύστημα, }\emph{regula fidei, reg. veritatis, traditio apostolica, prædicatio ap., fides catholica}, etc. Sometimes these terms are used in a wider sense, and embrace the whole course of catechetical instruction.} But this baptismal Creed was at first not precisely the same. It assumed different shapes and forms in different congregations.\footnote{See the older \emph{regulæ fidei} mentioned by Irenæus: \emph{Contra hær.,} lib. I. c. 10, § 1; III. c. 4, § 1, 2; IV. c. 33, § 7; Tertullian: \emph{De velandis virginibus}, c. 1; \emph{Adv. Praxeam}, c. 2; \emph{De præscript. hæret}., c. 13; Novatianus: \emph{De trinitate s. de regula fidei} (\emph{Bibl. P. P.}, ed. Galland. III. 287); Cyprian: \emph{Ep. ad Magnum}, and \emph{Ep. ad Januarium}, etc.; Origen: \emph{De principiis}, I. præf. § 4–10; Const. Apost. VI. 11 and 14. They are given in Vol. II. pp. 11–40; also by Bingham, Walch, Hahn, and Heurtley. I select, as a specimen, the descriptive account of Tertullian, who maintained against the heretics very strongly the unity of the traditional faith, but, on the other hand, also against the Roman Church (as a Montanist), the liberty of discipline and progress in Christian life. \emph{De velandis virginibus}, c. 1: '\emph{Regula quidem fidei una omnino est, sola immobolis et irreformabilis}, credendi\emph{ scilicet }in unicum Deum omnipotentem, \emph{mundi conditorem}, et Filium ejus Jesum Christum, natum ex virgine Maria, crucifixum sub Pontio Pilato, tertia die resuscitatum a mortuis, receptum in cælis, sedentem nunc ad dexteram Patris, venturum judicare vivos et mortuos,\emph{ per }carnis\emph{ etiam}resurrectionem.\emph{ Hac lege fidei manente cætera jam disciplinæ et conversationis admittunt novitatem correctionis, operante scilicet et proficiente usque in finem gratia Dei}.' In his tract against Praxeas (cap. 2) he mentions also, as an object of the rule of faith, '\emph{Spiritum Sanctum, paracletum, sanctificatorem fidei eorum qui credunt in Patrem et Filium et Spiritum Sanctum.}' We may even go further back to the middle and the beginning of the second century. The earliest trace of some of the leading articles of the Creed may be found in Ignatius, \emph{Epistola ad Trallianos}, c. 9 (ed. Hefele, p. 192), where he says of Christ that he was truly born 'of the Virgin Mary' (\textgreek{τοῦ ἐκ Μαρίας, ὃς ἀληθῶς ἐγεννήθη}), 'suffered under Pontius Pilate' \textgreek{(ἀληθῶς ἐδιώχθη ἐπί Ποντίου Πιλάτου),} 'was crucified and died' (\textgreek{ἀληθῶς ἐσταυρώθη καὶ ἀπέθανεν,}) and 'was raised from the dead' (\textgreek{ὃς καὶ ἀληθῶς ἠγέρθη ἀπὸ νεκρῶν, ἐγείραντος αὐτὸν τοῦ πατρὸς, αὐτοῦ.}) The same articles, with a few others, can be traced in Justin Martyr's \emph{Apol.} I. c. 10, 13, 21, 42, 46, 50.} Some were longer, some shorter; some declarative, some interrogative in the form of questions and answers.\footnote{Generally distributed under three heads: 1. \emph{Credis in Deum Patrem omnipotentem}, etc.? Resp. \emph{Credo}. 2. \emph{Credis et in Jesum Christum}, etc.? Resp. \emph{Credo}. 3. \emph{Credis et in Spiritum Sanctum}, etc.? Resp. \emph{Credo}. See the interrogative Creeds in Martene, \emph{De antiquis ecclesiæ ritibus}, 1. I. c. 1, and in Heurtley, l.c. pp. 103–116.} Each of the larger churches adapted the nucleus of the apostolic faith to its peculiar circumstances and wants; but they all agreed in the essential articles of faith, in the general order of arrangement on the basis of the baptismal formula, and in the prominence given to Christ's death and resurrection. We have an illustration in the modern practice of Independent or Congregational and Baptist churches in America, where the same liberty of framing particular congregational creeds ('covenants,' as they are called, or forms of profession and engagement, when members are received into full communion) is exercised to a much larger extent than it was in the primitive ages.

The first accounts we have of these primitive creeds are merely fragmentary. The ante-Nicene fathers give us not the exact and full formula, but only some articles with descriptions, defenses, explications, and applications. The creeds were committed to memory, but not to writing.\footnote{Hieronymus, \emph{Ep.} 61, \emph{ad Pammach.: 'Symbolum fidei et spei nostræ, quod ab apostolis traditum, non scribitur in charta et atramento, sed in tabulis cordis carnalibus}.' Augustine, \emph{Serm.} ccxii, 2: '\emph{Audiendo symbolum discitur, nec in tabulis vel in aliqua materia, sed in corde scribitur}.'} This fact is to be explained from the 'Secret Discipline' of the ante-Nicene Church. From fear of profanation and misconstruction by unbelievers (not, as some suppose, in imitation of the ancient heathen Mysteries), the celebration of the sacraments and the baptismal creed, as a part of the baptismal act, were kept secret among the communicant members until the Church triumphed in the Roman Empire.\footnote{On the \emph{Disciplina arcani} comp. my \emph{Church History}, I. 384 sq., and Semisch, \emph{On the Ap. Creed}, p. 17, who maintains, with others, that the Apostles' Creed existed in full as a part of the Secret Discipline long before it was committed to writing.}

The first writer in the West who gives us the text of the Latin creed, with a commentary, is Rufinus, towards the close of the fourth century.

The most complete or most popular forms of the baptismal creed in use from that time in the West were those of the churches of Rome, Aquileja, Milan, Ravenna, Carthage, and Hippo. They differ but little.\footnote{See these Nicene and post-Nicene Creeds in Hahn, l.c. pp. 3 sqq., and in Heurtley, l.c. 43 sqq. Augustine (and pseudo-Augustine) gives eight expositions of the Symbol, and mentions, besides, single articles in eighteen passages of his works. See Caspari, l.c. II. 264 sq. He follows in the main the (Ambrosian) form of the Church of Milan, which agrees substantially with the Roman. Twice he takes the North African Symbol of Carthage for a basis, which has additions in the first article, and puts the article on the Church to the close (\emph{vitam æternam per sanctam ecclesiam}). We have also, from the Nicene and post-Nicene age, several commentaries on the Creed by Cyril of Jerusalem, Rufinus, Ambrose, and Augustine. They do not give the several articles continuously, but it is easy to collect and to reconstruct them from the comments in which they are expounded. Cyril expounds the Eastern Creed, the others the Western. Rufinus takes that of the Church of Aquileja, of which he was presbyter, as the basis, but notes incidentally the discrepancy between this Creed and that of the Church of Rome, so that we obtain from him the text of the Roman Creed as well. He mentions earlier expositions of the Creed, which were lost (\emph{In Symb. }§ 1).} Among these, again, the Roman formula gradually gained general acceptance in the West for its intrinsic excellence, and on account of the commanding position of the Church of Rome. We know the Latin text from Rufinus (390), and the Greek from Marcellus of Ancyra (336–341). The Greek text is usually regarded as a translation, but is probably older than the Latin, and may date from the second century, when the Greek language prevailed in the Roman congregation.\footnote{See Caspari, Vol. III. pp. 28–161.}

This Roman creed was gradually enlarged by several clauses from older or contemporaneous forms, viz., the article 'descended into Hades' (taken from the Creed of Aquileja), the predicate 'catholic' or 'general,' in the article on the Church (borrowed from Oriental creeds), 'the communion of saints' (from Gallican sources), and the concluding 'life everlasting' (probably from the symbols of the churches of Ravenna and Antioch).\footnote{The last clause occurs in the Greek text of Marcellus and in the baptismal creed of Antioch \textgreek{(καὶ εἰς ἁμαρτιῶν ἄφειν καὶ εἰς νεκρῶν ἀνάστασιν καὶ εἰς ζωὴν αἰώνιον).} See Caspari, Vol. I. pp. 83 sqq.} These additional clauses were no doubt part of the general faith, since they are taught in the Scriptures, but they were first expressed in local creeds, and it was some time before they found a place in the authorized formula.

If we regard, then, the \emph{present }text of the Apostles' Creed as a complete whole, we can hardly trace it beyond the sixth, certainly not beyond the close of the fifth century, and its triumph over all the other forms in the Latin Church was not completed till the eighth century, or about the time when the bishops of Rome strenuously endeavored to conform the liturgies of the Western churches to the Roman order.\footnote{Heurtley says (l.c. p. 126): 'In the course of the seventh century the Creed seems to have been approaching more and more nearly, and more and more generally, to conformity with the formula now in use; and before its close, instances occur of creeds virtually identical with that formula. The earliest creed, however, which I have met with actually and in all respects identical with it, that of Pirminius, does not occur till the eighth century; and even towards the close of the eighth, A.D. 785, there is one remarkable example of a creed, then in use, which retains much of the incompleteness of the formula of earlier times, the Creed of Etherius Uxamensis.' The oldest known copies of our present \emph{textus receptus} are found in manuscripts of works which can not be traced beyond the eighth or ninth century, viz., in a \emph{'Psalterium Græcum Gregorii Magni},' preserved in the Library of Corpus Christi College, Cambridge, and first published by Abp. Usher, 1647 (also by Heurtley, l.c. p. 82), and another in the '\emph{Libellus Pirminii }[who died 758] \emph{de singulis libris canonicis scarapsus}' (=\emph{collectus}), published by Mabillon (\emph{Analecta}, Tom. IV. p. 575). The first contains the Creed in Latin and Greek (both, however, in Roman letters), arranged in two parallel columns; the second gives first the legend of the Creed with the twelve articles assigned to the twelve apostles, and then the Latin Creed as used in the baptismal service. See Heurtley, p. 71.}  But if we look at the several articles of the Creed separately, they are all of Nicene or ante-Nicene origin, while its kernel goes back to the apostolic age. All the facts and doctrines which it contains, are in entire agreement with the New Testament. And this is true even of those articles which have been most assailed in recent times, as the supernatural conception of our Lord (comp. Matt. i. 18; Luke i. 35), the descent into Hades (comp. Luke xxiii. 43; Acts ii. 31; 1 Pet. iii. 19; iv. 6), and the resurrection of the body (1 Cor. xv. 20 sqq., and other places).\footnote{The same view of the origin of the Apostles' Creed is held by the latest writers on the subject, as Hahn, Heurtley, Caspari, Zöckler, Semisch. Zöckler says (l.c. p. 18): '\emph{Das Apostolicum ist hinsichtlich seiner jetzigen Form sowohl nachapostolisch, als selbst nachaugustinisch, aber hinsichtlich seines Inhalts ist es nicht nur voraugustinisch, sondern ganz und gar apostolisch—in diesen einfachen Satz lässt die Summe der einschlägigen kritisch patristischen Forschungsergebnisse sich kurzerhand zusammendrängen. Und die Wahrheit dieses Satzes, soweit er die Apostolicität des Inhalts behauptet, lässt sich bezüglich jedes einzelnen Gliedes oder Sätzchens, die am spätesten hinzugekommenen nicht ausgenommen, mit gleicher Sicherheit erhärten.}' Semisch traces the several articles, separately considered, up to the third and second centuries, and the substance to the first. Fr. Spanheim and Calvin did the same. Calvin says: '\emph{Neque mihi dubium est, quin a prima statim ecclesiæ origine, adeoque ab ipso Apostolorum seculo instar publicæ et omnium calculis receptæ confessionis obtinuerit}' (\emph{Inst.} lib. II. c. 16, § 18). The most elaborate argument for the early origin is given by Caspari, who derives the Creed from Asia Minor in the beginning of the second century (Vol. III. pp. 1–161).}

The rationalistic opposition to the Apostles' Creed and its use in the churches is therefore an indirect attack upon the New Testament itself. But it will no doubt outlive these assaults, and share in the victory of the Bible over all forms of unbelief.\footnote{The discussion of the Apostles' Creed entered a stage of great warmth after Dr. Schaff's death, 1893. The work by Kattenbusch, the most extensive and exhaustive on the subject, was followed by treatments from the pens of Harnack, Cremer, Zahn, Loofs, Kunze, and others in Germany, Burn, and Badcock, 1930, in England and McGiffert in the United States. The early Roman baptismal formula is carried by Harnack and Mirbt to 150 or earlier, and by Kattenbusch and Zahn to 120 or earlier. A. Seeberg found the clauses in the New Testament writings and held that a creedal formula was in use in Apostolic times. McGiffert, who was followed by Krüger, proposed the theory that the formula was a reply to the heresies of Marcion about 160. Badcock opposes the view of Kattenbusch, Harnack, and Burn on the origin of the Apostles' Creed, relying in part upon Irenaeus's recently found treatise, ``The teaching of the Apostles.'' The renewed study of the Apostles' Creed was followed by a new study of the doctrine of the Virgin birth of Christ in view of the omission of the clause ``conceived by the Holy Ghost'' in the forms of the Rule of Faith known to us and the statement of the early Roman baptismal formula, ``born of the Holy Ghost and the Virgin Mary.'' The most recent treatise on the Virgin birth is by Machen, \emph{The Virgin Birth of Christ}, N. Y., 1930.—Ed.}

III. I add a table, with critical notes, to show the difference between the original Roman creed, as given by Rufinus in Latin (about A.D. 390), and by Marcellus in Greek (A.D. 336–341), and the received form of the Apostles' Creed, which came into general use in the seventh or eighth century. The additions are inclosed in brackets.

The old Roman Form.

The Received Form.

1. I believe in God the Father Almighty\footnote{The Creed of Aquileja has, after \emph{Patrem omnipotentem}, the addition: '\emph{invisibilem et impassibilem},' in opposition to Sabellianism and Patripassianism. The Oriental creeds insert \emph{one} before \emph{God}. Marcellus omits \emph{Father}, and reads \textgreek{εἰς θεὸν παντοκράτορα}.}

1. I believe in God the Father Almighty [\emph{Maker of heaven and earth}].\footnote{'\emph{Creatorem cœli et terræ}' appears in the Apostles' Creed from the close of the seventh century, but was extant long before in ante-Nicene rules of faith (Irenæus, \emph{Adv. hœr}. I. c. 10, 1; Tertullian, \emph{De vel. virg.} c. l, '\emph{mundi conditorem};' \emph{De prœscr. hæret.} c. 13), in the Nicene Creed \textgreek{(ποιητὴν οὐρανοῦ καὶ γῆς, κ.τ.λ.),} and all other Eastern creeds, in opposition to the Gnostic schools, which made a distinction between the true God and the Maker of the world (the Demiurge).}

2. And in Jesus Christ, his only Son, our Lord;

2. And in Jesus Christ, his only Son, our Lord;

3. Who was born by the Holy Ghost of the Virgin Mary;\footnote{'\emph{Qui natus est de Spiritu Sancto ex }(or \emph{et}) \emph{Maria virgine}.'}

3. Who was [\emph{conceived}] by the Holy Ghost, born of the Virgin Mary;\footnote{'\emph{Qui }CONCEPTUS\emph{ est de Spiritu Sancto, natus ex Maria virgine.}' The distinction between conception and birth first appears in the \emph{Sermones de Tempore}, falsely attributed to Augustine.}

4. Was crucified under Pontius Pilate and was buried;

4. [\emph{Suffered}]\footnote{'\emph{Passus},' perhaps from the Nicene Creed (\textgreek{παθόντα,} which there implies the crucifixion). In some forms '\emph{crucifixus},' in others '\emph{mortuus}' is omitted.} under Pontius Pilate, was crucified [\emph{dead}], and buried

[\emph{He descended into Hell} (\emph{Hades})];\footnote{From the Aquilejan Creed: '\emph{Descendit ad inferna},' or, as the Athanasian Creed has it, '\emph{ad inferos},' \emph{to the inhabitants of the spirit-world. }Some Eastern (Arian) creeds: \textgreek{κατέβη εἰς τὸν ᾅδην} (also \textgreek{εἰς τὰ καταχθόνια,} or \textgreek{εἰς τὰ κατώτατα).} Augustine says (\emph{Ep.} 99, al. 164, § 3) that unbelievers only deny '\emph{fuisse apud inferos Christum.}' Venantius Fortunatus, A.D. 570, who had Rufinus before him, inserted the clause in his creed. Rufinus himself, however, misunderstood it by making it to mean the same as \emph{buried} (§ 18: '\emph{vis verbi eadem videtur esse in eo quod sepultus dicitur}').}

5. The third day he rose from the dead;

5. The third day he rose from the dead;

6. He ascended into heaven; and sitteth on the right hand of the Father;

6. He ascended into heaven; and sitteth on the right hand of [\emph{God}] the Father [\emph{Almighty}];\footnote{The additions '\emph{Dei}' and '\emph{omnipotentis},' made to conform to article first, are traced to the Spanish version of the Creed as given by Etherius Uxamensis (bishop of Osma), A.D. 785, but occur already in earlier Gallican creeds. See Heurtley, pp. 60, 67.}

7. From thence he shall come to judge the quick and the dead.

7. From thence he shall come to judge the quick and the dead.

8. And in the HOLY GHOST;

8. [\emph{I believe}]\footnote{'\emph{Credo},' in common use from the time of Petrus Chrysologus, d. 450. But \emph{And}, without the repetition of the verb, is no doubt the primitive form, as it grew immediately out of the baptismal formula, and gives clearer and closer expression to the doctrine of the Trinity.} in the Holy Ghost;

9. The Holy Church;

9. The Holy [\emph{Catholic}]\footnote{'\emph{Catholicam}' (universal), in accordance with the Nicene Creed, and older Oriental forms, was received into the Latin Creed before the close of the fourth century (comp. Augustine: \emph{De Fide et Symbolo}, c. 10). The term \emph{catholic}, as applied to the Church, occurs first in the Epistles of Ignatius (\emph{Ad Smyrnæos}, cap. 8: \textgreek{ὥσπερ ὅπου ἂν ᾖ Χριστὸς Ἰησοῦς, ἐκεῖ ἡ καθολικὴ ἐκκλησία} and in the \emph{Martyrium Polycarpi} (inscription, and cap. 8: \textgreek{ἁπάσης τῆς κατὰ τὴν οἰκουμένην καθολικῆς ἐκκλησίας,} comp. c. 19, where Christ is called \textgreek{ποιμὴν τῆς κατὰ οἰκουμένην καθολικῆς ἐκκλησίας}.} Church

[\emph{The communion of saints}];\footnote{The article '\emph{Commumionem sanctorum},' unknown to Augustine (\emph{Enchir.} c. 64, and \emph{Serm.} 213), appears first in the 115th and 118th Sermons \emph{De Tempore}, falsely attributed to him. It is not found in any of the Greek or earlier Latin creeds. See the note of Pearson \emph{On the Creed}, Art. IX. sub '\emph{The Communion of Saints}' (p. 525, ed. Dobson). Heurtley, p. 146, brings it down to the close of the eighth century, since it is wanting in the Creed of Etherius, 785. The oldest commentators understood it of the communion with the saints in heaven, but afterwards it assumed a wider meaning: the fellowship of all true believers, living and departed.}

10. The forgiveness of sins;

10. The forgiveness of sins;

11. The resurrection of the body (flesh).\footnote{The Latin reads \emph{carnis}, the Greek \textgreek{σαρκός,} flesh; the Aquilejan form hujus\emph{ carnis}, \emph{of }this\emph{ flesh} (which is still more realistic, and almost materialistic), '\emph{ut possit caro vel pudica coronari, vel impudica puniri}' (Rufinus, § 43). It should be stated, however, that there are two other forms of the Aquilejan Creed given by Walch (xxxiv. and xxxv.) and by Heurtley (pp. 30–32), which differ from the one of Rufinus, and are nearer the Roman form.}

11. The resurrection of the body (flesh);

12. [\emph{And the life everlasting}].\footnote{Some North African forms (of Carthage and Hippo Regius) put the article of the Church at the close, in this way: '\emph{vitam eternam per sanctam ecclesiam}.' Others: \emph{carnis resurrectionem in vitam æternam.} The Greek Creed of Marcellus, which otherwise agrees with the old Roman form, ends with \textgreek{ζωὴν αἰώνιον}.}

Note on the Legend of the Apostolic Origin of the Creed.—Till the middle of the seventeenth century it was the current belief of Roman Catholic and Protestant Christendom that the Apostles' Creed was '\emph{membratim articulatimque}' composed by the apostles in Jerusalem on the day of Pentecost, or before their separation, to secure unity of teaching, each contributing an article (hence the somewhat arbitrary division into twelve articles).\footnote{The old Roman form has only eleven articles, unless art. 6 be divided into two; while the received text has sixteen articles, if 'Maker of heaven and earth,' 'He descended into Hades,' 'the communion of saints,' and 'the life everlasting,' are counted separately.} Peter, under the inspiration of the Holy Ghost, commenced: 'I believe in God the Father Almighty;' Andrew (according to others, John) continued: 'And in Jesus Christ, his only Son, our Lord;' James the elder went on: 'Who was conceived by the Holy Ghost;' then followed John (or Andrew): 'Suffered under Pontius Pilate;' Philip: 'Descended into Hades;' Thomas: 'The third day he rose again from the dead;' and so on till Matthias completed the work with the words 'life everlasting. Amen.'

The first trace of this legend, though without the distribution alluded to, we find at the close of the fourth century, in the \emph{Expositio Symboli} of Rufinus of Aquileja. He mentions an ancient tradition concerning the apostolic composition of the Creed ('\emph{tradunt majores nostri}'), and falsely derives from this supposed joint authorship the name \emph{symbolon} (from \textgreek{συμβάλλειν,} in the sense \emph{to contribute}); confounding \textgreek{σύμβολον,} sign, with \textgreek{συμβολή,} contribution ('\emph{Symbolum Græce et indicium dici potest et collatio, hoc est, quod plures in unum conferunt}'). The same view is expressed, with various modifications, by Ambrosius of Milan (d. 397), in his \emph{Explanatio Symboli ad initiandos}, where he says: '\emph{Apostoli sancti convenientes fecerunt symbolum breviter};' by John Cassianus (about 424), \emph{De incarnat. Dom.} VI. 3; Leo M., \emph{Ep. 27 ad Pulcheriam}; Venantius Fortunatus, \emph{Expos. brevis Symboli Ap.}; Isidorus of Seville (d. 636). The distribution of the twelve articles among the apostles is of later date, and there is no unanimity in this respect. See this legendary form in the pseudo-Augustinian \emph{Sermones de Symbolo}, in Hahn, l.c. p. 24, and another from a \emph{Sacramentarium Gallicanum} of the seventh century, in Heurtley, p. 67.

The Roman Catechism gives ecclesiastical sanction, as far as the Roman Church is concerned, to the fiction of a direct apostolic authorship.\footnote{Pars prima, cap. 1, qu. 2 (\emph{Libri Symbolici Eccl. Cath.}, ed. Streitwolf and Klener, Tom. I. p. 111): '\emph{Quæ igitur primum Christiani homines tenere debent, illa sunt, quæ fidei duces, doctoresque sancti Apostoli, divino Spiritu afflati, duodecim Symboli articulis distinxerunt. Nam, cum mandatum a Domino accepissent, ut pro ipso legatione fungentes, in universum mundum proficiscerentur, atque omni creaturæ Evangelium prædicarent: Christianæ fidei formulam componendam censuerunt, ut scilicet id omnes sentirent ac dicerent, neque ulla essent inter eos schismata},' etc. Ibid. qu. 3: '\emph{Hanc autem Christianæ fidei et spei professionem a se compositam Apostoli Symbolum appellarunt; sive quia ex variis sententiis, quas singuli in commune contulerunt, conflata est; sive quia ea veluti nota, et tessera quandam uterentur, qua desertores et subintroductos falsos fratres, qui Evangelium adulterabant, ab iis, qui veræ Christi militiæ sacramento se obligarent, facile possent internoscere.}'} Meyers, l.c., advocates it at length, and Abbé Martigny, in his '\emph{Dictionnaire des antiquitées Chrétiennes},' Paris, 1865 (art. \emph{Symbole des apôtres}, p. 623), boldly asserts, without a shadow of proof: '\emph{Fidèlement attaché à la tradition de l’Église catholique, nous tenons, non-seulement qu’il est l’œuvre des apôtres, mais encore qu’il fut composé par eux, alors que réunis à Jérusalem, ils allaient se disperser dans l’univers entier; et qu’ils volurent, avant de séparer, fixer une règle de foi vraiment uniforme et catholique, destinée à être livrée, partout la même, aux catéchumènes.}'

Even among Protestants the old tradition has occasionally found advocates, such as Lessing (1778), Delbrück (1826), Rudelbach (1844), and especially Grundtvig (d. 1872). The last named, a very able but eccentric high-church Lutheran bishop of Denmark, traces the Creed, like the Lord's Prayer, to Christ himself, in the period between the Ascension and Pentecost. The poet Longfellow (a Unitarian) makes poetic use of the legend in his \emph{Divine Tragedy }(1871).

On the other hand, the apostolic origin (after having first been called in question by Laurentius Valla, Erasmus, Calvin\footnote{In his \emph{Catechism}, Calvin says that the formula of the common Christian faith is called \emph{symbolum apostolorum, quod vel ab ore apostolorum excepta fuerit, vel ex eorum scriptis fideliter collecta.}}) has been so clearly disproved long since by Vossius, Rivetus, Voëtius, Usher, Bingham, Pearson, King, Walch, and other scholars, that it ought never to be seriously asserted again.

The arguments against the apostolic authorship are quite conclusive:

1. The intrinsic improbability of such a mechanical composition. It has no analogy in the history of symbols; even when composed by committees or synods, they are mainly the production of one mind. The Apostles' Creed is no piece of mosaic, but an organic unit, an instinctive work of art in the same sense as the \emph{Gloria in Excelsis}, the \emph{Te Deum}, and the classical prayers and hymns of the Church.

2. The silence of the Scriptures. Some advocates, indeed, pretend to find allusions to the Creed in Paul's 'analogy' or 'proportion of faith,' Rom. xii. 7; 'the good deposit,' 2 Tim. i. 14; 'the first principles of the oracles of God,' Heb. v. 12; 'the faith once delivered to the saints,' Jude, ver. 3; and 'the doctrine,' 2 John, ver. 10; but these passages can be easily explained without such assumption.

3. The silence of the apostolic fathers and all the ante-Nicene and Nicene fathers and synods. Even the œcumenical Council of Nicæa knows nothing of a symbol of strictly apostolic composition, and would not have dared to supersede it by another.

4. The variety in form of the various rules of faith in the ante-Nicene churches, and of the Apostolic Symbol itself down to the eighth century. This fact is attested even by Rufinus, who mentions the points in which the Creed of Aquileja differed from that of Rome. 'Such variations in the form of the Creed forbid the supposition of any fixed system of words, recognized and received as the composition of the apostles; for no one, surely, would have felt at liberty to alter any such normal scheme of faith.'\footnote{Dr. Nevin (l.c. p. 107), who otherwise puts the highest estimate on the Creed. See the comparative tables on the gradual growth of the Creed in the second volume of this work.}

5. The fact that the Apostles' Creed never had any general currency in the East, where the Nicene Creed occupies its place, with an almost equal claim to apostolicity as far as the substance is concerned.

\xsection{The Nicene Creed.}

§ 8. The Nicene Creed.

Literature.

I. See the works on the œcumenical Creeds noticed p. 12, and the extensive literature on the Council of Nicæa, mentioned in my Church History, Vol. III. pp. 616, 617, and 622. The acts of the Council are collected in Greek and Latin by Mansi, \emph{Collect. sacr. Concil.}, Tom. II. fol. 635–704. The Council of Nicæa is more or less fully discussed in the historical works, general or particular, of Tillemont, Walch, Schröckh, Gibbon, A. de Broglie, Neander, Gieseler, Baur (Hist. of the Doctrine of the Trinity), Dorner (History of Christology), Hefele (History of Councils, Stanley (History of the Eastern Church).

II. Special treatises on the Nicene symbol:

Ph. Melanchthon: Explicatio Symb. Nicæni, ed. \emph{a J. Sturione}, Viteb. 1561, 8vo.

Casp. Cruciger: Enarrationis Symboli Nicæni articuli duo, etc., Viteb. 1548, 4to, and Symboli Nicæni enarratio cum præfatione Ph. Melanchthonis, acc. priori editioni plures Symboli partes, Basil (without date).

J. H. Heidegger (d. 1698): De Symbolo Nicæne-Constantinopolitano (Tom. II. \emph{Disp. select.} pp. 716 sqq., Turici, 1675–97).

J. G. Baier: De Conc. Nicæni primi et Œcum. auctoritate atque integritate, Jen. 1695 (in \emph{Disputat. theol. decad.} I.).

T. Fecht: Innocentia Concilii et Symboli Nicæni, Rostock, 1711.

T. Caspar Suicer (d. 1684): Symbolum Nicæno-Constant. expositum et ex antiquitate ecclesiastica illustratum, Traj. ad Eh. 1718, 4to.

George Bull (d. 1710): Defensio Fidei Nicænæ, Oxon. 1687, in his Latin works ed. by Grabe, 1703; by Burton, 1827, and again 1846; English translation in the \emph{Anglo-Catholic Library}, Oxf. 1851, 2 vols.

The Nicene Creed, or Symbolum Nicæno-Constantinopolitanum, is the Eastern form of the primitive Creed, but with the distinct impress of the Nicene age, and more definite and explicit than the Apostles' Creed in the statement of the divinity of Christ and the Holy Ghost. The terms 'coessential' or 'coequal' \textgreek{(ὁμοούσιος τῷ πατρί),} 'begotten before all worlds' \textgreek{(πρὸ πάντων τῶν αἰώνων),} 'very God of very God' \textgreek{(θεὸς ἀληθινὸς ἐκ θεοῦ ἀληθινοῦ),} 'begotten, not made' \textgreek{(γεννηθείς, οὐ ποιηθείς),} are so many trophies of orthodoxy in its mighty struggle with the Arian heresy, which agitated the Church for more than half a century. The Nicene Creed is the first which obtained universal authority. It rests on older forms used in different churches of the East, and has undergone again some changes.\footnote{Compare the symbols of the church of Jerusalem, the church of Alexandria, and the creed of Cæsarea, which Eusebius read at the Council of Nicæa, in Usher, l.c. pp. 7, 8; more fully in Vol. II. pp. 11 sqq., and in Hahn, \emph{Bibliothek der Symbole}, pp. 40 sqq., 91 sqq.}

The Eastern creeds arose likewise out of the baptismal formula, and were intended for the baptismal service as a confession of the faith of the catechumen in the Triune God.\footnote{Eusebius, in his Epistle to the people of Cæsarea, says of the creed which he had proposed to the Council of Nicæa for adoption, that he had learned it as a catechumen, professed it at his baptism, taught it in turn as presbyter and bishop, and that it was derived from our Lord's baptismal formula. It resembles the old Nicene Creed very closely; see Vol. II. p. 29. The shorter creed of Jerusalem used at baptism, as given by Cyril, \emph{Catech.} xix. 9, is simply the baptismal formula put interrogatively; see Hahn, pp. 51 sqq.}

We must distinguish two independent or parallel creed formations, an Eastern and a Western; the one resulted in the Nicene Creed as completed by the Synod of Constantinople, the other in the Apostles' Creed in its Roman form. The Eastern creeds were more metaphysical, polemical, flexible, and adapting themselves to the exigencies of the Church in the maintenance of her faith and conflict with heretics; the Western were more simple, practical, and stationary. The former were controlled by synods, and received their final shape and sanction from two œcumenical Councils; the latter were left to the custody of the several churches, each feeling at liberty to make additions or alterations within certain limits, until the Roman form superseded all others, and was quietly, and without formal synodical action, adopted by Western Christendom.

In the Nicene Creed we must distinguish three forms—the original Nicene, the enlarged Constantinopolitan, and the still later Latin.

1. The original Nicene Creed dates from the first œcumenical Council, which was held at Nicæa, A.D. 325, for the settlement of the Arian controversy, and consisted of 318 bishops, all of them from the East (except Hosius of Spain). This Creed abruptly closes with the words 'and in the Holy Ghost,' but adds an anathema against the Arians. This was the authorized form down to the Council of Chalcedon.

2. The Nicæno-Constantinopolitan Creed, besides some minor changes in the first two articles,\footnote{The most remarkable change in the first article is the omission of the words \textgreek{πουτέστιν ἐκ τῆς οὐσίας τοῦ Πατρός, θεὸν ἐκ θεοῦ} on which great stress was laid by the Athanasian party against the Arians, who maintained that the Son was not of the \emph{essence}, but of the \emph{will} of the Father.} adds all the clauses after 'Holy Ghost,' but omits the anathema. It gives the text as now received in the Eastern Church. It is usually traced to the second œcumenical Council, which was convened by Theodosius in Constantinople, A.D. 381, against the Macedonians or Pneumatomachians (so called for denying the deity of the Holy Spirit), and consisted of 150 bishops, all from the East. There is no authentic evidence of an \emph{œcumenical recognition} of this enlarged Creed till the Council at Chalcedon, 451, where it was read by Aëtius (a deacon of Constantinople) as the 'Creed of the 150 fathers,' and accepted as orthodox, together with the old Nicene Creed, or the 'Creed of the 318 fathers.' But the additional clauses existed in 374, seven years before the Constantinopolitan Council, in the two creeds of Epiphanius, a native of Palestine, and most of them as early as 350, in the creed of Cyril of Jerusalem.\footnote{See Vol. II. pp. 31–38, and the Comparative Table, p. 40; Lumby, p. 68; and Hort, pp. 72–150. Dr. Hort tries to prove that the 'Constantinopolitan' or Epiphanian Creed is not a revision of the Nicene Creed at all, but of the Creed of Jerusalem, and that it dates probably from Cyril, about 362–364, when he adopted the Nicene \emph{homoousia}, and may have been read by him at the Council of Constantinople in vindication of his orthodoxy. Ffoulkes (in Smith's \emph{Dict. of Christ. Antiq.} Vol. I. p. 438) conjectures that it was framed at Antioch about 372, and adopted at the supplemental Council of Constantinople, 382.}

The Nicene Creed comes nearest to that of Eusebius of Cæsarea, which likewise abruptly closes with \textgreek{πνεῦμα ἅγιον}; the Constantinopolitan Creed resembles the creeds of Cyril and Epiphanius, which close with 'the resurrection' and 'life everlasting.' We may therefore trace both forms to Palestine, except the Nicene \emph{homoousion.}

3. The Latin or Western form differs from the Greek by the little word \emph{Filioque}, which, next to the authority of the Pope, is the chief source of the greatest schism in Christendom. The Greek Church, adhering to the original text, and emphasizing the \emph{monarchia} of the Father as the only root and cause of the Deity, teaches the \emph{single} procession \textgreek{(ἐκπόρευσις)} of the Spirit from the Father \emph{alone}, which is supposed to be an \emph{eternal} inner-trinitarian process (like the eternal generation of the Son), and not to be confounded with the \emph{temporal mission } \textgreek{(πέμψις)} of the Holy Spirit by the Father and the Son. The Latin Church, in the interest of the co-equality of the Son with the Father, and taking the procession (\emph{processio}) in a wider sense, taught since Augustine the \emph{double} procession of the Spirit from the Father \emph{and the Son}, and, without consulting the East, put it into the Creed.

The first clear trace of the \emph{Filioque} in the Nicene Creed we find at the third Council of Toledo in Spain, A.D. 589, to seal the triumph of orthodoxy over Arianism. During the eighth century it obtained currency in England and in France, but not without opposition. Pope Leo III., when asked by messengers of a council held during the reign of Charlemagne at Aix la Chapelle, A.D. 809, to sanction the \emph{Filioque}, decided in favor of the double procession, but against any change in the Creed. Nevertheless, the clause gained also in Italy from the time of Pope Nicholas I. (858), and was gradually adopted in the entire Latin Church. From this it passed into the Protestant Churches.\footnote{Comp. Vol. II., at the close.}

Another addition in the Latin form, \emph{'Deus de Deo},' in article II., created no difficulty, as it was in the original Nicene Creed, but it is useless on account of the following '\emph{Deus verus de Deo vero},' and hence was omitted in the Constantinopolitan edition.

The Nicene Creed (without these Western additions) is more highly honored in the Greek Church than in any other, and occupies the same position there as the Apostles' Creed in the Latin and Protestant Churches. It is incorporated and expounded in all the orthodox Greek and Russian Catechisms. It is also (with the \emph{Filioque}) in liturgical use in the Roman (since about the sixth century), and in the Anglican and Lutheran Churches.\footnote{In the Reformed Churches, except the Episcopal, the Nicene Creed is little used. Calvin, who had a very high opinion of the Apostles' Creed, depreciates the Nicene Creed, as a '\emph{carmen cantillando magis aptum, quam confessionis formula}' (\emph{De Reform. Eccles.}).} It was adopted by the Council of Trent as the fundamental Symbol, and embodied in the Profession of the Tridentine Faith by Pius IV. It is therefore more strictly an œcumenical Creed than the Apostles' and the Athanasian, which have never been fully naturalized in the Oriental Churches.

. . 'The faith of the Trinity lies,

Shrined for ever and ever, in those grand old words and wise;

A gem in a beautiful setting; still, at matin-time,

The service of Holy Communion rings the ancient chime;

Wherever in marvelous minster, or village churches small,

Men to the Man that is God out of their misery call,

Swelled by the rapture of choirs, or borne on the poor man's word,

Still the glorious Nicene confession unaltered is heard;

Most like the song that the angels are singing around the throne,

With their ``Holy! holy! holy!'' to the great Three in One.'\footnote{From 'A Legend of the Council of Nice,' by Cecil Frances Alexander, in '\emph{The Contemporary Review}' for February, 1867, pp. 176–179.}

The relation of the Nicene Creed to the Apostles' Creed may be seen from the following table:

The Apostles' Creed; Received Text.

The Nicene Creed, as Enlarged A.D. 381.

(The clauses in brackets are the later additions.)

(The words in brackets are Western changes.)

1. I believe in God the Father Almighty,

1. We [I] believe\footnote{The Greek reads the plural \textgreek{(πιστεύομεν),} but the Latin and English versions have substituted for it the singular (\emph{credo}, I believe), in accordance with the Apostles' Creed and the more subjective character of the Western churches.} in one God the Father Almighty,

[Maker of heaven and earth].

Maker of heaven and earth,

And of all things visible and invisible.

2. And in Jesus Christ, his only Son, our Lord;

2. And in one Lord Jesus Christ,

the only-begotten Son of God,

Begotten of the Father before all worlds;

[God of God],

Light of Light.

Very God of very God,

Begotten, not made,

Being of one substance with the Father;

By whom all things were made;

3. Who was [conceived] by the Holy Ghost,

3. Who, for us men, and for our salvation,

Born of the Virgin Mary;

came down from heaven,

And was incarnate by the Holy Ghost of

the Virgin Mary,

And was made man

4. [Suffered] under Pontius Pilate, was crucified [dead], and buried;

4. He was crucified for us under Pontius Pilate;

And suffered and was buried;

[He descended into Hades];

* * * * *

5. The third day he rose again from the dead;

5. And the third day he rose again,

According to the Scriptures;

6. He ascended into heaven,

6. And ascended into heaven,

And sitteth on the right hand of [God] the Father [Almighty];

And sitteth on the right hand of the Father;

7. From thence he shall come to judge the quick and the dead.

7. And he shall come again, with glory, to judge the quick and the dead;

Whose kingdom shall have no end.

8. And [I believe] in the Holy Ghost;

8. And [I believe] in the Holy Ghost,

The Lord, and Giver of life;

Who proceedeth from the Father [and the Son];

Who with the Father and the Son together is worshiped and glorified;

Who spake by the Prophets.

9. The holy [catholic] Church;

9. And [I believe] in\footnote{The Greek reads \textgreek{εἰς μίαν . . . ἐκκλησίαν,} but the Latin and English versions, in conformity with the Apostles' Creed, mostly omit \emph{in} before \emph{ecclesiam}; see p. 15.} one holy catholic and apostolic Church;

[The communion of saints];

* * * * *

10. The forgiveness of sins;

10. We [I] acknowledge\footnote{Here and in art. 11 the singular is substituted in Western translations for \textgreek{ὁμολογοῦμεν} and \textgreek{προςδοκῶμεν}.} one baptism for the remission of sins;

11. The resurrection of the flesh [body];

11. And we [I] look for the resurrection of the dead;

12. [And the life everlasting].

12. And the life of the world to come.

We give also, in parallel columns, the original and the enlarged formulas of the Nicene Creed, italicizing the later additions, and inclosing in brackets the passages which are omitted in the received text:

The Nicene Creed of 325.\footnote{The Greek original is given, together with the similar Palestinian confession, by Eusebius in his \emph{Epistola ad Cæsareenses}, which is preserved by Athanasius at the close of his \emph{Epistola de decretis Synodi Nicænæ} (\emph{Opera}, ed. Montfaucon, I. 239); also, with some variations, in the Acts of the Council of Chalcedon (Act. II. in Mansi, Tom. VII.); in Theoderet, \emph{H. E. }I. 12; Socrates, \emph{H. E.} I. 8; Gelasius, \emph{H. Conc. Nic.} 1. II. c. 35. See the literature and variations in Walch, l.c. pp. 75 and 87 sqq.; also in Hahn, l.c. pp. 105 sqq.}

The Constantinopolitan Creed of 381.\footnote{The Greek text in the acts of the second œcumenical Council (Mansi, Tom. III. p. 565; Hardouin, Vol. I. p. 814), and also in the acts of the fourth œcumenical Council. See Vol. II p. 35; Hahn, l.c. p. 111; and my \emph{Church Hist.} Vol. III. pp. 667 sqq.}

We believe in one God, the Father Almighty, Maker of all things visible and invisible.

We believe in one God, the Father Almighty, Maker of \emph{heaven and earth, and of} all things visible and invisible.

And in one Lord Jesus Christ, the Son of God, begotten of the Father [the only-begotten; that is, of the essence of the Father, God of God], Light of Light, very God of very God, begotten, not made, being of one substance \textgreek{(ὁμοούσιον)} with the Father; by whom all things were made [both in heaven and on earth]; who for us men, and for our salvation, came down and was incarnate and was made man; he suffered, and the third day he rose again, ascended into heaven; from thence he shall come to judge the quick and the dead.

And in one Lord Jesus Christ, the \emph{only-begotten} Son of God, begotten of the Father \emph{before all worlds} (æons), Light of Light, very God of very God, begotten, not made, being of one substance with the Father; by whom all things were made; who for us men, and for our salvation, came down \emph{from heaven}, and was incarnate \emph{by the Holy Ghost of the Virgin Mary}, and was made man; he \emph{was crucified for us under Pontius Pilate}, and suffered, \emph{and was buried}, and the third day he rose again, \emph{according to the Scriptures}, and ascended into heaven, \emph{and sitteth on the right hand of the Father}; from thence he shall come \emph{again, with glory}, to judge the quick and the dead; \emph{whose kingdom shall have no end.}

And in the Holy Ghost.

And in the Holy Ghost, \emph{the Lord and Giver of life, who proceedeth from the Father, who with the Father and the Son together is worshiped and glorified, who spake by the prophets. In one holy catholic and apostolic Church; we acknowledge one baptism for the remission of sins; we look for the resurrection of the dead, and the life of the world to come. Amen.}

[But those who say: 'There was a time when he was not;' and 'He was not before he was made;' and 'He was made out of nothing,' or 'He is of another substance' or 'essence,' or 'The Son of God is created,' or 'changeable,' or 'alterable'—they are condemned by the holy catholic and apostolic Church.]

\xsection{The Creed of Chalcedon.}

§ 9. The Creed of Chalcedon.

Literature.

The Acta Concilii in the collections of Mansi, Tom. VII., and of Hardouin, Tom. II.

Evagrius: Historia eccl. lib. II. c. 2, 4, 18.

Facundus (Bishop of Hermiane, in Africa): Pro defens. trium capitulorum, lib. V. c. 3, 4; lib. VIII. c. 4 (see Gallandi, \emph{Bibl. PP.} Tom. XI. pp. 713 sqq.).

Liberatus (Archdeacon of Carthage): Breviarium causæ Nestorianorum et Eutychianorum, c. 13 (Gallandi, Tom. XII. pp. 142 sqq.).

Baronius: Annal. ad ann. 451, No. 55 sqq.

Edm. Richer: Hist. concil. generalium, Paris, 1680 (Amst. 1686, 3 vols.), lib. I. c. 8.

Tillemont: Mémoires, etc. Tom. XV. pp. 628 sqq. (in the article on Leo the Great).

Natalis Alexander: Hist. eccles. sec. V. Tom. V. pp. 64 sqq. and pp. 209 sqq.

Quesnel: Synopsis actorum Conc. Chalcedon., in his Dissertat. de vita, etc., \emph{S. Leonis} (see the Ballerini edition of the works of Leo the Great, Tom. II. pp. 501 sqq.).

Hulsemann: Exercit. ad Concil. Chalcedon. Lips. 1651.

Cave: Hist. literaria, etc. pp. 311 sqq. ed. Genev. 1705.

Walch: Ketzerhistorie, Vol. VI. p. 329 sq.; and his Historie der Kirchenversammlungen, p. 307 sq.

Arendt: Papst Leo der Grosse, Mainz, 1835, pp. 267–322.

Dorner: History of the Development of the Doctr. of the Person of Christ (2d Germ. ed.), Part II. 99–150.

Hefele: History of the Councils, Freiburg, Vol. II. (1856). p. 392 sq.

Schaff: History of the Christian Church, N. Y. 1867, Vol. III. pp. 740 sqq. Comp. the literature there on pp. 708 sq., 714 sq., 722.

The Creed of Chalcedon was adopted at the fourth and fifth sessions of the fourth œcumenical Council, held at Chalcedon, opposite Constantinople, A.D. 451 (Oct. 22d and 25th). It embraces the Nicæno-Constantinopolitan Creed, and the christological doctrine set forth in the classical \emph{Epistola Dogmatica} of Pope Leo the Great to Flavian, the Patriarch of Constantinople and martyr of diophysitic orthodoxy at the so-called Council of Robbers (held at Ephesus in 449).\footnote{Comp. my \emph{Church Hist.} Vol. III. p. 738.}

While the first Council of Nicæa had established the eternal, pre-existent Godhead of Christ, the Symbol of the fourth œcumenical Council relates to the incarnate Logos, as he walked upon earth and sits on the right hand of the Father. It is directed against the errors of Nestorius and Eutyches, who agreed with the Nicene Creed as opposed to Arianism, but put the Godhead of Christ in a false relation to his humanity. It substantially completes the orthodox Christology of the ancient Church; for the definitions added during the Monophysite and Monothelite controversies are few and comparatively unessential. As the Nicene doctrine of the Trinity stands midway between Tritheism and Sabellianism, so the Chalcedonian formula strikes the true mean between Nestorianism and Eutychianism.

The following are the leading ideas of the Chalcedonian Christology as embodied in this symbol:\footnote{Abridged, in part, from \emph{My Church History}, Vol. III. pp. 747 sqq.}

1. A true incarnation of the Logos, or the second person in the Godhead \textgreek{(ἐνανθρώπησις θεοῦ, ἐνσάρκωσις τοῦ λόγου, } \emph{incarnatio Verbi}).)\footnote{The diametrical opposite of the \textgreek{ἐνανθρώπησις θεοῦ} is the heathen \textgreek{ἀποθέωσις ἀνθρώπου}.} This incarnation is neither a conversion or transmutation of God into man, nor a conversion of man into God, and a consequent absorption of the one, or a confusion \textgreek{(κρᾶσις, σύγχυσις)} of the two; nor, on the other hand, a mere indwelling \textgreek{(ἐνοίκησις, } \emph{inhabitatio}) of the one in the other, nor an outward, transitory connection \textgreek{(συνάφεια, } \emph{conjunctio}) of the two factors, but an actual and abiding union of the two in one personal life.

2. The precise distinction between nature and person. Nature or substance (essence, \textgreek{οὐσία)} denotes the totality of powers and qualities which constitute a being; while person \textgreek{(ὑπόστασις, πρόσωπον)} is the Ego, the self-conscious, self-asserting and acting subject. The Logos assumed, not a human person (else we would have two persons, a divine and a human), but human nature which is common to us all; and hence he redeemed, not a particular man, but all men as partakers of the same nature.

3. The God-Man as the result of the incarnation. Christ is not a (Nestorian) \emph{double} being, with \emph{two} persons, nor a compound (Apollinarian or Monophysite) \emph{middle} being, a \emph{tertium quid}, neither divine \emph{nor} human; but he is \emph{one }person \emph{both }divine \emph{and }human.

4. The duality of the natures. The orthodox doctrine maintains, against Eutychianism, the distinction of nature even after the act of incarnation, without confusion or conversion \textgreek{(ἀσυγχύτως, } \emph{inconfuse},  and \textgreek{ἀτρέπτως, } \emph{immutabiliter}), yet, on the other hand, without division or separation \textgreek{(ἀδιαιρέτως, } \emph{indivise}, and \textgreek{ἀχωρίστως, } \emph{inseparabiliter}), so that the divine will ever remain divine, and the human ever human,\footnote{'\emph{Tenet},' says Leo, in his Epist. 28 ad Flavian., '\emph{sine defectu proprietatem suam utraque natura, et sicut formam servi Dei forma non adimit, ita formam Dei servi forma non minuit. . . . Agit utraque forma cum alterius communione quod proprium est; Verbo scilicet operante quod Verbi est, et carne exsequente quod carnis est. Unum horum coruscat miraculis, aliud succumbit injuriis. Et sicut Verbum ab æqualitate paternæ gloriæ non recedit, ita caro naturam nostri generis non relinquit.}'} and yet the two have continually one common life, and interpenetrate each other, like the persons of the Trinity.\footnote{Here belongs, in further explanation, the scholastic doctrine of the \textgreek{περιχώρησις,} \emph{permeatio, circummeatio, circulatio, circumincessio, intercommunio}, or reciprocal indwelling and pervasion, which has relation, not merely to the Trinity, but also to Christology. The verb \textgreek{περιχωρεῖν} is first applied by Gregory of Nyssa (\emph{Contra Apollinarium}) to the interpenetration and reciprocal pervasion of the two natures in Christ. On this rested also the doctrine of the exchange or communication of attributes, \textgreek{ἀντίδοσις, ἀντιμετάστασις, κοινωνία ἰδιωμάτων,} \emph{communicatio idiomatum}. The \textgreek{ἀντιμετάστασις τῶν ὀνομάτων,} also \textgreek{ἀντιμεδίστασις,} \emph{transmutatio proprietatum}, transmutation of attributes, is, strictly speaking, not identical with \textgreek{ἀντίδοσις,} but a deduction from it, and the rhetorical expression for it.}

5. The unity of the person  \textgreek{(ἕνωσις καθ᾽ ὑπόστασιν, ἕνωσις ὑποστατική, } \emph{unio hypostatica} or \emph{unio personalis}). The union of the divine and human nature in Christ is a permanent state resulting from the incarnation, and is a real, supernatural, personal, and inseparable union—in distinction from an essential absorption or confusion, or from a mere moral union; or from a mystical union such as holds between the believer and Christ. The two natures constitute but one personal life, and yet remain distinct. 'The same who is true God,' says Leo, 'is also true man, and in this unity there is no deceit; for in it the lowliness of man and the majesty of God perfectly pervade one another. . . . Because the two natures make only one person, we read on the one hand: ``The Son of \emph{Man} came down from heaven'' (John iii. 13), while yet the Son of \emph{God} took flesh from the Virgin; and on the other hand: ``The Son of \emph{God} was crucified and buried,''\footnote{Comp. 1 Cor. ii. 8: 'They would not have crucified the Lord of glory.'} while yet he suffered, not in his Godhead as coeternal and consubstantial with the Father, but in the weakness of human nature.' The self-consciousness of Christ is never divided; his person consists in such a union of the human and the divine natures, that the divine nature is the seat of self-consciousness, and pervades and animates the human.

6. The whole work of Christ is to be attributed to his person, and not to the one or the other nature exclusively. The person is the acting subject, the nature the organ or medium. It is the one divine-human person of Christ that wrought miracles by virtue of his divine nature, and that suffered through the sensorium of his human nature. The superhuman effect and infinite merit of the Redeemer's work must be ascribed to his person because of his divinity; while it is his humanity alone that made him capable of, and liable to, toil, temptation, suffering, and death, and renders him an example for our imitation.

7. The anhypostasia, impersonality, or, to speak more accurately, the enhypostasia, of the human nature of Christ;\footnote{\textgreek{Ἀνυπόστατος} is that which has no personality in itself, \textgreek{ἐνυπόστατος} that which subsists in another personality, or partakes of another hypostasis.} for anhypostasia is a purely negative term, and presupposes a fictitious abstraction, since the human nature of Christ did not exist at all before the act of the incarnation, and could therefore be neither personal nor impersonal. The meaning of this doctrine is that Christ's human nature had no independent personality of its own, besides the divine, and that the divine nature is the root and basis of his personality.\footnote{The doctrine of the impersonality of the human nature of Christ may already be found as to its germ in Cyril of Alexandria, and was afterwards more fully developed by John of Damascus (\emph{De orthodoxa fide}, lib. III.), and by the Lutheran scholastics of the seventeenth century, who, however, did not, for all this, conceive Christ as a mere generic being typifying mankind, but as a concrete human individual. Comp. Petavius, \emph{De incarnatione}, lib. V. c. 5–8 (Tom. IV. pp. 421 sqq.); Thomasius, \emph{Christol.} II. 108–110; Rothe, \emph{Dogmatik}, II. 51 and 147.}

There is, no doubt, a serious difficulty in the old orthodox Christology, if we view it in the light of our modern psychology. We can conceive of a human nature without sin (for sin is a corruption, not an essential quality, of man), but we can not conceive of a human nature without personality, or a self-conscious and free Ego; for this distinguishes it from the mere animal nature, and is man's crowning excellency and glory. To an unbiased reader of the Gospel history, moreover, Christ appears as a full human personality, thinking, speaking, acting, suffering like a man (only without sin), distinguishing himself from other men and from his heavenly Father, addressing him in prayer, submitting to him his own will, and commending to him his spirit in the hour of death.\footnote{He calls himself a 'man,' \textgreek{ἄνθρωπος} (John viii. 40; comp. xix. 5), and very often 'the Son of man,' and other men his 'brethren' (John xx. 17).} Yet, on the other hand, be appears just as clearly in the Gospels as a personality in the most intimate, unbroken, mysterious life-union with his heavenly Father, in the full consciousness of a personal pre-existence before the creation, of having been sent by the Father from heaven into this world, of living in heaven even during this earthly abode, and of being ever one with him in will and in essence.\footnote{John viii.58; xvii. 5, 24; iii. 11-13; v. 37; vi. 38, 62; viii. 42; x. 30, and many other passages in the Gospels. Dr. R. Rothe, who rejects the orthodox doctrine of the Trinity and the Incarnation, yet expressly admits (\emph{Dogmatik}, II. 88): '\emph{Ebenso bestimmt, wie seine wahre Menschheit, tritt im Neuen Testament auch die wahre }GOTTHEIT\emph{ des Erlösers hervor}.' To escape the orthodox inference of an incarnation of a divine hypostasis, Rothe must resort (p. 100) to the Socinian interpretation of John xvii. 5, where the Saviour asserts his pre-existence \emph{with the Father } \textgreek{(δόξασόν με σύ, πάτερ, παρὰ σεαυτῷ τῇ δόξῃ, ᾗ εἶχον πρὸ τοῦ τὸν κόσμον εἶναι παρὰ σοί);} thereby distinguishing himself from the hypostasis of the Father, and yet asserting coeternity. The Socinians and Grotius find here merely an ideal glory in the divine counsel; but it must be taken, in analogy with similar passages, of a \emph{real}, personal, self-conscious pre-existence, and a \emph{real} glory attached to it; otherwise it would be nothing peculiar and characteristic of Christ. How absurd would it be for a man to utter such a prayer!} In one word, he makes the impression of a \emph{theanthropic, divine-human} person.\footnote{A \emph{persona σύνθετος,} in the language of the old Protestant divines. \emph{Divina et humana naturæ}' (says Hollaz), '\emph{in una persona} \textgreek{ συνθέτῳ} \emph{ Filii Dei existentes, unam eandemque habent} \textgreek{ ὑπόστασιν,} \emph{ modo tamen habendi diversam. Natura enim divina eam habet primario, per se et independenter, natura autem humana secundario, propter unionem personalem, adeoque participative.} The divine nature, therefore, is, in the orthodox system, that which forms and constitutes the personality (\emph{das personbildende Princip.}).} His human personality was completed and perfected by being so incorporated with the pre-existent Logos-personality as to find in it alone its full self-consciousness, and to be permeated and controlled by it in every stage of its development.

The Chalcedonian Christology has latterly been subjected to a rigorous criticism (by Schleiermacher, Baur, Dorner, Rothe, and others), and has been charged with a defective psychology, and now with dualism, now with docetism, according as its distinction of two natures or of the personal unity has most struck the eye. But these imputations neutralize each other, like the imputations of tritheism and modalism, which may be made against the orthodox doctrine of the Trinity when either the tri-personality or the consubstantiality is taken alone. This, indeed, is the peculiar excellence of the Creed of Chalcedon, that it exhibits so sure a tact and so wise a circumspection in uniting the colossal antithesis in Christ, and seeks to do justice alike to the distinction of the natures and to the unity of the person. In Christ all contradictions are reconciled.

The Chalcedonian Creed is far from exhausting the great mystery of godliness, 'God manifest in flesh.' It leaves much room for a fuller appreciation of the genuine, perfect, and sinless humanity of Christ, of the Pauline doctrine of the \emph{Kenosis}, or self-renunciation and self-limitation of the Divine Logos in the incarnation and during the human life of our Lord, and for the discussion of other questions connected with his relation to the Father and to the world, his person and his work. But it indicates the essential elements of Christological truth, and the boundary-lines of Christological error. It defines the course for the sound development of this central article of the Christian faith so as to avoid both the Scylla of Nestorian dualism and the Charybdis of Eutychian monophysitism, and to save the full idea of the one divine-human personality of our Lord and Saviour. Within these limits theological speculation may safely and freely move, and bring us to clearer conceptions; but in this world, where we 'know only in part \textgreek{(ἐκ μέρους),}' and 'see through a mirror obscurely \textgreek{(δἰ ἐσόπτρου ἐν αἰνίγματι)}' it will never fully comprehend the great central mystery of the theanthropic life of our Lord.

\xsection{The Athanasian Creed.}

§ 10. The Athanasian Creed.

Literature.

I. Comp. the general literature of the Three Creeds noticed p. 12, especially Lumby and Swainson.

II. Special treatises on the Athanasian Creed:

[Venantius Fortunatus (Bishop of Poitiers, d. about A.D. 600)]: Expositio Fidei Catholicae Fortunati. The oldest commentary on the Athanasian Creed, published from a MS. in the Ambrosian Library at Milan by \emph{Muratori}, 1698, in the second vol. of his \emph{Anecdota}, p. 228, and better in an Appendix to \emph{Waterland's} treatise (see below). But the authorship of Ven. Fort. is a mere conjecture of Muratori, from the name \emph{Fortunatus}, and is denied by modern critics.

Dav. Pareus (Ref.): Symbolum Athanasii breviter declaratum. Heidelb. 1618.

J. H. Heidegger (Ref.): De Symbolo Athanasiano. Tur. 1680.

W. E. Tentzel (Luth.): Judicia eruditorum de Symb. Athanasiano. Gothæ, 1687.

Jos. Anthelmi (R. C.): Disquisitio de Symb. Athan. Paris, 1693.

Montfaucon (R. C.): Diatribe de Symbolo Quicunque, in his edition of the works of St. Athanasius. Paris, 1698, Tom. II. pp. 719-735.

Dan. Waterland (Anglican): A Critical History of the Athanasian Creed, etc. Cambridge, 1724, 2d ed. 1728 (in Waterland's works, Vol. III. pp. 97–270, Oxf. ed. 1843), also re-edited by J. R. King. Lond. 1871. The fullest and most learned treatise on the subject, but in part superseded by recent investigations.

Dom. Maria Speroni (R. C.): De Symbolo vulgo S. Athanasii, two dissertations. Patav. 1750 sq.

John Radcliffe: The Creed of St. Athanasius, illustrated from the Old and New Test., Passages of the Fathers, etc. Lond. 1844.

Philip Schaff: The Athanasian Creed, in the 'American Presbyterian Review,' New York, for 1866, pp. 584–625; \emph{Church History}, Vol. III. pp. 689 sqq.

A. P. Stanley (Dean of Westminster): The Athanasian Creed. Lond. 1871.

E. S. Ffoulkes (B. D.): The Athanasian Creed: By whom Written and by whom Published. Lond. 1872.

Ch. A. Heurtley: The Athanasian Creed. Oxford, 1872. (Against Ffoulkes.)

Comp. the fac-simile edition of the \emph{Utrecht Psalter} (Lond. 1875), and Sir Thos. Hardy (Deputy-Keeper of the Public Records), two Reports on the Athanas. Creed in Connection with the Utrecht Psalter. Lond. 1873.

The Athanasian Creed is also called Symbolum Quicunque, from the first word, '\emph{Quicunque vult salvus esse}.'\footnote{It first bears the title, '\emph{Fides sanctæ Trinitatis},' or '\emph{Fides Catholica Sanctæ Trinitatis};' then (in the '\emph{Cod. Usserius secundus}') '\emph{Fides Sancti Athanasii Alexandrini}.' Hincmar of Rheims, about A.D. 852, calls it '\emph{Sermonem Athanasii de fide, cujus initium est: ``Quicunque vult salvus esse}.'''}

I. Its origin is involved in obscurity, like that of the Apostles' Creed, the Gloria in Excelsis, and the Te Deum. It furnishes one of the most remarkable examples of the extraordinary influence which works of unknown or doubtful authorship have exerted. Since the ninth century it has been ascribed to Athanasius, bishop of Alexandria, the chief defender of the divinity of Christ and the orthodox doctrine of the Trinity (d. 373).\footnote{According to the mediæval legend, Athanasius composed it during his exile in Rome, and offered it to Pope Julius as his confession of faith. So Baronius, Petavius, Bellarmin, etc. This tradition was first opposed and refuted by Gerhard Vossius (1642) and Ussher (1647).} The great name of 'the father of orthodoxy' secured for it an almost œcumenical authority, notwithstanding the solemn prohibition of the third and fourth œcumenical Councils to compose or publish any other creed than the Nicene.\footnote{Conc. Ephes. Can. VII. 'The holy Synod has determined that no person shall be allowed to bring forward, or to write, or to compose any other Creed \textgreek{(ἑτέραν πίστιν μηδενὶ ἐξεῖναι προφέρειν ἤγουν συγγράφειν ἢ συντιθέναι),} besides that which was settled by the holy fathers who assembled in the city of Nicæa, with the Holy Spirit. But those who shall dare to compose any other Creed, or to exhibit or produce any such, if they are bishops or clergymen, they shall be deposed, but if they are of the laity, they shall be anathematized.' The Council of Chalcedon (451), although setting forth a new definition of faith, repeated the same prohibition (after the \emph{Defin. Fidei}).}

Since the middle of the seventeenth century the Athanasian authorship has been abandoned by learned Catholics as well as Protestants. The evidence against it is conclusive. The Symbol is nowhere found in the genuine writings of Athanasius or his contemporaries and eulogists. The General Synods of Constantinople (381), Ephesus (431), and Chalcedon (451) make no allusion to it whatever. It seems to presuppose the doctrinal controversies of the fifth century concerning the constitution of Christ's person; at least it teaches substantially the Chalcedonian Christology. And, lastly, it makes its first appearance in the Latin Churches of Gaul, North Africa, and Spain: while the Greeks did not know it till the eleventh century, and afterwards rejected or modified it on account of the Occidental clause on the procession of the Holy Ghost from the Father \emph{and the Son}. The Greek texts, moreover, differ widely, and betray, by strange words and constructions, the hands of unskilled translators.

The pseudo-Athanasian Creed originated in the Latin Church from the school of St. Augustine, probably in Gaul or North Africa. It borrows a number of passages from Augustine and other Latin fathers.\footnote{See the parallel passages in Waterland's treatise and in my \emph{Church History}, Vol. III. pp. 690 sqq.} It appears first in its full form towards the close of the eighth or the beginning of the ninth century. Its structure and the repetition of the damnatory clause in the middle and at the close indicate that it consists of two distinct parts, which may have been composed by two authors, and afterwards welded together by a third hand. The first part, containing the Augustinian doctrine of the Trinity, is fuller and more metaphysical. The second part, containing a summary of the Chalcedonian Christology, has been found separately, as a fragment of a sermon on the Incarnation, at Treves, in a MS. from the middle of the eighth century.\footnote{Now known as the Colbertine MS., in Paris, which is assigned to about A.D. 730–760, but is derived in part from older MSS. This fragment was first published consecutively by Professor Swainson in 1871, and again in his larger work, 1875 (p. 262), also by Lumby, p. 215. It begins thus: '\emph{Est ergo fides recta ut credamus et confitemur quia Dominus ihesus christus Dei filius, deus pariter et homo est},' etc.; and it ends: '\emph{Hæc est fides sancta et Catholica, quam omnes} [\emph{omnis}] \emph{homo qui ad uitam æternam peruenire desiderat scire integræ} [\emph{integre}] \emph{debet, et fideliter custodire.}' The compiler of the two parts intensified the damnatory clause by changing it into '\emph{quam nisi quisque fideliter firmiterque crediderit, salvus esse non poterit}.' The passages quoted by Archbishop Hincmar of Rheims, A.D. 852, are all taken from the first part.} The fact that Athanasius spent some time in exile at Treves may possibly have given rise to the tradition that the great champion of the orthodox doctrine of the Trinity composed the whole.\footnote{The authorship of the Symbolum \emph{Quicunque} is a matter of mere conjecture. The opinions of scholars are divided between Hilary of Arles (420–431), Vigilius of Tapsus (484), Vincentius Lirinensis (450), Venantius Fortunatus of Poitiers (570), Pope Anastasius (398), Victricius of Rouen (401), Patriarch Paulinus of Aquileja (Charlemagne's favorite theologian, d. 804). Waterland learnedly contends for Hilary of Arles; Quesnel, Cave, Bingham, and Neander for Vigilius Tapsensis of North Africa. Gieseler traces the \emph{Quicunque} to the Councils of Toledo in Spain (633, 638, 675, etc.), which used to profess the Nicene Creed with additional articles (like the \emph{Filioque}) against Arianism. Ffoulkes (who seceded to Rome, and returned, a better Protestant, to the Church of England) and Dean Stanley maintain that it arose in France, simultaneously with the forgery of the pseudo-Isidorean Decretals, for controversial purposes against the Greeks, to set up a fictitious antiquity for Latin \emph{doctrine} (the \emph{Filioque}), as the Decretals did for Latin \emph{ polity}. Swainson and Lumby assign the Creed to an unknown writer of the age of Charlemagne (d. 814) and Alcuin (d. 804), or to the period between 813 and 850.  \par The latest investigations since the rediscovery of the oldest (the Cotton) MS. in the 'Utrecht Psalter' (which was exposed for inspection at the British Museum in 1873, and has since been photographed) are unfavorable to an early origin; for this MS., which Ussher and Waterland assigned to the sixth century, dates probably from the ninth century (as the majority of scholars who investigated it, Drs. Vermuelen, Heurtley, Ffoulkes, Lumby, Swainson, contend against Hardy, Westwood, and Baron van Westreenen), since, among other reasons, it contains also the Apostles' Creed in its final form of 750. The authorship of Venantius Fortunatus (570) was simply inferred by Muratori from the common name 'Fortunatus' at the head of a MS. (\emph{Expositio Fidei Catholicæ Fortunati}) which contains a commentary on the Athanasian Creed, but which is not older than the eleventh century, and quotes a passage from Alcuin. Two other MSS. of the same commentary, but without a title, have been found, one at Florence, and one at Vienna (Lumby, p. 208; Swainson, pp. 317 sqq.). The internal evidence for an earlier date is equally inconclusive. The absence of \emph{Mater Dei } \textgreek{(θεοτόκος)} no more proves an ante-Nestorian origin (before 431, as Waterland contended) than the absence of \emph{consubstantialis } \textgreek{(ὁμοούσιος)} proves an ante-Nicene origin. \par So far, then, we have no \emph{proof} that the pseudo-Athanasian Creed in its \emph{present complete shape} existed before the beginning of the ninth century. And yet it \emph{may} have existed earlier. At all events, two separate compositions, which form the groundwork of our \emph{Quicunque}, are of older date, and the doctrinal substance of it, with the most important passages, may be found in the works of St. Augustine and his followers, with the exception of the damnatory clauses, which seem to have had their origin in the fierce contests of the age of Charlemagne. In a Prayer-Book of Charles the Bald, written about A. D. 870, we find the Athanasian Creed very nearly in the words of the received text. \par I may add that the indefatigable investigator, Dr. Caspari, of Christiania, informs me by letter (dated April 29, 1876) that he is still inclined to trace this Creed to the fifth century, between 450 and 600, and that he found, and will publish in due time, some old symbols which bear a resemblance to it, and may cast some light upon its obscure origin. \emph{Adhuc sub judice lis est.}}

II. Character and Contents.—The Symbolum Quicunque is a remarkably clear and precise summary of the doctrinal decisions of the first four œcumenical Councils (from A.D. 325 to A.D. 451), and the Augustinian speculations on the Trinity and the Incarnation. Its brief sentences are artistically arranged and rhythmically expressed. It is a musical creed or dogmatic psalm. Dean Stanley calls it 'a triumphant pæan' of the orthodox faith. It resembles, in this respect, the older \emph{Te Deum}, but it is much more metaphysical and abstruse, and its harmony is disturbed by a threefold anathema.

It consists of two parts.

The first part (ver. 3–28) sets forth the orthodox doctrine of the Holy Trinity, not in the less definite Athanasian or Nicæno-Constantinopolitan, but in its strictest Augustinian form, to the exclusion of every kind of subordination of essence. It is therefore an advance both on the Nicene Creed and the Apostles' Creed; for these do not state the doctrine of the Trinity in form, but only indirectly by teaching the Deity of the Son and of the Holy Spirit, and leave room for a certain subordination of the Son to the Father, and the Holy Spirit to both. The post-Athanasian formula states clearly and unmistakably both the absolute unity of the divine being or essence, and the tri-personality of the Father, the Son, and the Holy Spirit. God is one in three persons or hypostases, each person expressing the whole fullness of the Godhead, with all his attributes. The term \emph{persona} is taken neither in the old sense of a mere personation or form of manifestation (\textgreek{πρόσωπον,} face, mask), nor in the modern sense of an independent, separate being or individual, but in a sense which lies between these two conceptions, and thus avoids Sabellianism on the one hand, and Tritheism on the other. The divine persons are in one another, and form a perpetual intercommunication and motion within the divine essence.\footnote{The later scholastic terms for this indwelling and interpenetration are \textgreek{περιχώρησις}, \emph{inexistentia, permeatio, circumincessio}, etc. See my \emph{Church History}, Vol. III. p. 680.} Each person has all the divine attributes which are inherent in the divine essence, but each has also a characteristic individuality or property,\footnote{Called by the Greeks \textgreek{ἰδιότης} or \textgreek{ἴδιον}, by the Latins \emph{proprietas personalis} or \emph{character hypostaticus}.} which is peculiar to the person, and can not be communicated; the Father is unbegotten, the Son begotten, the Holy Ghost is proceeding. In this Trinity there is no priority or posteriority of time, no superiority or inferiority of rank, but the three persons are coeternal and coequal.

If the mystery of the Trinity can be logically defined, it is done here. But this is just the difficulty: the infinite truth of the Godhead lies far beyond the boundaries of logic, which deals only with finite truths and categories. It is well always to remember the saying of Augustine: 'God is greater and truer in our thoughts than in our words; he is greater and truer in reality than in our thoughts.'\footnote{'\emph{Verius cogitatur Deus quam dicitur, verius est quam cogitatur},' \emph{De Trinitate}, lib. VII. c. 4, § 7. Dr. Isaac Barrow, one of the intellectual giants of the Anglican Church (died 1677), in his \emph{Defense of the Blessed Trinity} (a sermon preached on Trinity Sunday, 1663), humbly acknowledges the transcendent incomprehensibility, while clearly stating the facts, of this great mystery: 'The sacred Trinity may be considered either as it is in itself wrapt up in inexplicable folds of mystery, or as it hath discovered itself operating in wonderful methods of grace towards us. As it is in itself, 'tis an object too bright and dazzling for our weak eye to fasten upon, an abyss too deep for our short reason to fathom; I can only say that we are so bound to mind it as to exercise our faith, and express our humility, in willingly believing, in submissively adoring those high mysteries which are revealed in the holy oracles concerning it by that Spirit itself which searcheth the depths of God. . . . That there is one Divine Nature or Essence, common unto three Persons, incomprehensibly united, and ineffably distinguished—united in essential attributes, distinguished by peculiar idioms and relations; all equally infinite in every divine perfection, each different from the other in order and manner of subsistence; that there is a mutual inexistence of one in all, and all in one, a communication without any deprivation or diminution in the communicant; an eternal generation, and an eternal procession, without precedence or succession, without proper causality or dependence; a Father imparting his own, and the Son receiving his Father's life, and a Spirit issuing from both, without any division or multiplication of essence—these are notions which may well puzzle our reason in conceiving how they agree, but should not stagger our faith in assenting that they are true; upon which we should meditate, not with hope to comprehend, but with dispositions to admire, veiling our faces in the presence, and prostrating our reason at the feet, of Wisdom so far transcending us.'}

The second part (ver. 29–44) contains a succinct statement of the orthodox doctrine concerning the person of Christ, as settled by the general Councils of Ephesus 431 and Chalcedon 451, and in this respect it is a valuable supplement to the Apostles' and Nicene Creeds. It asserts that Christ had a \emph{rational} soul \textgreek{(νοῦς, νεῦμα),} in opposition to the Apollinarian heresy, which limited the extent of his humanity to a mere body with an animal soul inhabited by the divine Logos. It also teaches the proper relation between the divine and human nature of Christ, and excludes the Nestorian and Eutychian or Monophysite heresies, in essential agreement with the Chalcedonian Symbol.\footnote{See the preceding section.}

III. The Damnatory Clauses.—The Athanasian Creed, in strong contrast with the uncontroversial and peaceful tone of the Apostles' Creed, begins and ends with the solemn declaration that the catholic faith in the Trinity and the Incarnation herein set forth is the indispensable condition of salvation, and that those who reject it will be lost forever. The same damnatory clause is also wedged in at the close of the first and at the beginning of the second part. This threefold anathema, in its natural historical sense, is not merely a solemn \emph{warning} against the great danger of heresy,\footnote{So a majority of the 'Ritual Commission of the Church of England,' appointed in 1867: 'The condemnations in this Confession of Faith are to be no otherwise understood than as a solemn \emph{warning of the peril} of those who willfully reject the Catholic faith.' Such a warning would be innocent and unobjectionable, indeed, but fall far short of the spirit of an age which abhorred heresy as the greatest of crimes, to be punished by death.} nor, on the other hand, does it demand, as a condition of salvation, a full knowledge of, and assent to, the \emph{logical statement} of the doctrines set forth (for this would condemn the great mass even of Christian believers); but it does mean to exclude from heaven all who reject the divine \emph{truth} therein taught. It requires every one who would be saved to believe in the only true and living God, Father, Son, and Holy Ghost, one in essence, three in persons, and in one Jesus Christ, very God and very Man in one person.

The damnatory clauses, especially when sung or chanted in public worship, grate harshly on modern Protestant ears, and it may well be doubted whether they are consistent with true Christian charity and humility, and whether they do not transcend the legitimate authority of the Church. They have been defended by an appeal to Mark xvi. 16; but in this passage those only are condemned who reject the \emph{gospel}, i.e., the great facts of Christ's salvation, not any peculiar dogma. Salvation and damnation depend exclusively on the grace of God as apprehended by a living faith, or rejected in ungrateful unbelief. The original Nicene Symbol, it is true, added a damnatory clause against the Arians, but it was afterwards justly omitted. Creeds, like hymns, lose their true force and miss their aim in proportion as they are polemical and partake of the character of manifestoes of war rather than confessions of faith and thanks to God for his mighty works.\footnote{'It seems very hard,' says Bishop Jeremy Taylor, 'to put uncharitableness into a creed, and so to make it become an article of faith.' Chillingworth: 'The damning clauses in St. Athanasius's Creed are most false, and also in a high degree schismatical and presumptuous.'}

IV. Introduction and Use.—The Athanasian Creed acquired great authority in the Latin Church, and during the Middle Ages it was almost daily used in the morning devotions.\footnote{J. Bona, \emph{De divina Psalmodia}, c. 16, § 18, p. 863 (as quoted by Köllner, \emph{Symbolik}, I. 85): '\emph{Illud Symbolum olim, teste Honorio, quotidie est decantatum, jam vero diebus Dominicis in totius cœtus frequentia recitatur, ut sanctæ fidei confessio ea die apertius celebretur}.'}

The Reformers inherited the veneration for this Symbol. It was formally adopted by the Lutheran and several of the Reformed Churches, and is approvingly mentioned in the Augsburg Confession, the Form of Concord, the Thirty-nine Articles, the Second Helvetic, the Belgic, and the Bohemian Confessions.\footnote{It is printed, with the two other œcumenical Creeds, in all the editions of the Lutheran 'Book of Concord,' and as an appendix to the doctrinal formulas of the Reformed Dutch Church in America. It was received into the 'Provisional Liturgy of the German Reformed Church in the United States,' published Philadelphia, 1858, but omitted in the revised edition of 1867.}

Luther was disposed to regard it as 'the most important and glorious composition since the days of the apostles.'\footnote{'\emph{Es ist also gefasset, dass ich nicht weiss, ob seit der Apostel Zeit in der Kirche des Neuen Testamentes etwas Wichtigeres and Herrlicheres geschrieben sei}' (Luther, \emph{Werke}, ed. Walch, VI. 2315).}

Some Reformed divines, especially of the Anglican Church have commended it very highly; even the Puritan Richard Baxter lauded it as 'the best explication [better, statement] of the Trinity,' provided, however, 'that the damnatory sentences be excepted, or modestly expounded.'

In the Church of England it is still sung or recited in the cathedrals and parish churches on several festival days,\footnote{The rubric directs that the Athanasian Creed 'shall be sung or said at Morning Prayer, instead of the Apostles' Creed, on Christmas-day, the Epiphany, St. Matthias, Easter-day, Ascension-day, Whitsunday, St. John the Baptist, St. James, St. Bartholomew, St. Matthew, St. Simon and St. Jude, St. Andrew, and upon Trinity Sunday.'} but this compulsory public use meets with growing opposition, and was almost unanimously condemned in 1867 by the royal commission appointed to consider certain changes in the Anglican Ritual.\footnote{By nineteen out of the twenty-seven members of the Ritual Commission. See their opinions in Stanley, l.c. pp. 73 sqq. Dean Stanley on that occasion urged no less than sixteen reasons against the public use of the Athanasian Creed. On the other hand, Dr. Pusey has openly threatened to leave the Established Church if the Athanasian Creed, and with it the doctrinal status of that Church, should be disturbed. Brewer's defense is rather feeble. Bishop Ellicott proposed, in the Convocation of Canterbury, to relieve the difficulty by a revision of the English translation, e.g. by rendering \emph{vult salvus esse}, 'desires to be in a state of salvation,' instead of 'will be saved.' Others suggest an omission of the damnatory clauses. But the true remedy is either to omit the Athanasian Creed altogether from the Book of Common Prayer, or to leave its public use optional.}

The Protestant Episcopal Church in the United States, when, in consequence of the American Revolution, it set up a separate organization in the Convention of 1785 at Philadelphia, resolved to remodel the Liturgy (in 'the Proposed Book'), and, among other changes, excluded from it both the Nicene and the Athanasian Creeds, and struck out from the Apostles' Creed the clause, 'He descended into hell.' The Archbishops of Canterbury and York, before consenting to ordain bishops for America, requested their brethren to restore the clause of the Apostles' Creed, and 'to give to the other two Creeds a place in their Book of Common Prayer, even though the use of them should be left discretional.'\footnote{Bishop White (of Philadelphia): \emph{Memoirs of the Protestant Episcopal Church in the United States of America}, New York, 2d ed. 1836, pp 305, 306.} In the Convention held at Wilmington Del., October 10, 1786, the request of the English prelates, as to the first two points, was acceded to, but 'the restoration of the Athanasian Creed was negatived.' As the opposition to this Creed was quite determined, especially on account of the damnatory clauses, the mother Church acquiesced in the omission, and granted the desired Episcopal ordination.\footnote{White's \emph{Memoires}, 26, 27. Bishop White himself was decidedly opposed to the Creed, as was Bishop Provost, of New York. The Archbishop of Canterbury told them afterwards: 'Some wish that you had retained the Athanasian Creed; but I can not say that I feel uneasy on the subject, for you have retained the doctrine of it in your Liturgy, and as to the Creed itself, I suppose you thought it not suited to the use of a congregation' (l.c. 117, 118).}

In the Greek Church it never obtained general currency or formal ecclesiastical sanction, and is only used for private devotion, with the omission of the clause on the double procession of the Spirit.\footnote{Additional Lit. on the Athan. Creed.—Swainson: \emph{The Nic. and App. Creeds, with an Account of the Creed of St. Athanasius}, London, 1894.—Burn in Robinson's Texts and Studies, 1896.—Ommanney, London, 1897, is inclined to ascribe it to Vincens of Lerins about 450.—Bp. Gore, Oxf., 1897.—J. B. Smith in \emph{Contemp. Rev.}, Apr., 1901.—Oxenham, London, 1902.—J. A. Robinson, London, 1905.—Bp. Jayne, 1905.—W. S. Bishop: \emph{Devel. of Trin. Doctr. in the Nic. and Athanas. Creeds}, 1910.—H. Brewer (S.J.), Das \emph{sogenannte Athanas. Glaubensbekenntniss}, 1909.—Burkitt, 1912.—Loofs in Herzog, ii, 177–194, who places its probable origin in Southern France, 450–600.—Badcock inclines to the Ambrosian authorship and calls it a hymn to be memorized. The Abp. of Canterbury, following a resolution of the Lambeth Conference, 1908, appointed a commission of seven, including Bp. Wordsworth of Salisbury, Prof. Swete and Dean Kilpatrick, to prepare a revision of the English translation of the Athanas. Creed. Their report proposed thirteen minor changes. The Anglican Book of Common Prayer prescribed that the Creed be said or sung at morning prayer on thirteen feasts, including Christmas, Easter, Ascension day, and Trinity Sunday. By the order of both Convocations it was omitted and a new rubric inserted, making its use optional on Trinity Sunday. In the ``Revised'' Book of Common Prayer, recommended by the House of Bishops and rejected by Parliament, 1928, the following rubrics are printed side by side, making the use of the creed optional: ``may be sung or said at morning or evening prayer'' on the first Sunday after Christmas, the feast of the Annunciation, and Trinity Sunday. 2. On Trinity Sunday, the recitation beginning with clause 3, ``The Catholic faith is this,'' etc., and closing with clause 28. 3. On the Sunday after Christmas and Ascension day, the recitation being from clause 30 to clause 41. 4. On all the thirteen festivals mentioned in the original Book of Common Prayer. A ``revised translation is added'' which differs from the translation of 1909. See the \emph{Translation of 1909 with Latin Text}, by H. Turner, London, 1910, 15 pp. and 1918, 23 pp. Also the \emph{Book of Com. Prayer with the Additions and Deviations Proposed in 1928}, with Pref., Cambr. Press, 1928. By Roman Cath. usage the creed is prescribed for Trinity Sunday and at prime on all Sundays except Easter and such other feasts for which a special service is provided.—Ed.}

\xchapter{Chapter 3. Creeds of the Greek Church}

General Literature.

\emph{Orthodoxa Confessio catholicæ atque apostol. ecclesiæ orientalis a}  Pet. Mogila  \emph{compos., a}  Meletio Syrigo  \emph{aucta et mutata, gr. c. præf.}  Nectarii  \emph{curav.}  Panagiotta, Amst. 1662; \emph{cum interpret. lat. ed.}  Laur. Normann, Leipz. 1695, 8vo; \emph{c. interpret. lat. et vers. german, ed.}  K. Glo. Hofmann, Breslau, 1751, 8vo. Also in Russian: Moscow, 1696; German by J. Leonh. Frisch, Frankfurt and Leipzig, 1727, 4to; Dutch by J. A. Senier, Haarlem, 1722; in Kimmel's \emph{Monumenta}, P. I. 1843.

\emph{Clypeus orthodoxæ fidei, sive Apologia} (\textgreek{Ἀσπἱς ὀρθοδοξίας, ἠ ἀπολογία καὶ ἔλεγχος}) \emph{ab Synodo Hierosolymitana} (A.D. 1672) \emph{sub Hierosolymorum Patriarcha Dositheo composita adversus Calvinistas hæreticos}, etc. Published at Paris, Greek and Latin, 1676 and 1678: then in  Harduini  \emph{Acta Conciliorum}, Par. 1715, Tom. XI. fol. 179–274; also in Kimmel's  \emph{Monum.} P. I. 325–488. Comp. also the Acts of the Synod of \emph{Constantinople}, held in the same year (1672), and publ. in Hard. l.c. 274–284, and in Kimmel, P. II. 214–227.

\emph{Confessio cathol. et apostolica in oriente ecclesiæ, conscripta compendiose per}  Metrophanem Critopulum. \emph{Ed. et. lat. redd.}  J. Hornejus,  Helmst. 1661, 4to (the title-page has erroneously the date 1561).

Cyrilli Lucaris:  \emph{Confessio christ. fidei græca cum additam. Cyrilli}, Geneva, 1633: græc. et lat. (Condemned as heretical.)

\emph{Acta et scripta theologorum Wirtembergensium et patriarchæ Constantinop.}  Hieremiæ, \emph{quæ utrique ab a.} 1576 \emph{usque ad a.} 1581 \emph{de Augustana Confessione inter se miserunt, gr. et lat. ab iisdem theologis edita}, Wittenb. 1584, fol. This work contains the Augsburg Confession in Greek, three epistles of Patriarch Jeremiah, criticising the Augsb. Conf., and the answers of the Tübingen divines, all in Greek and Latin.

E. J. Kimmel  and H. Weissenborn: \emph{Monumenta fidei ecclesiæ orientalis. Primum in unum corpus collegit, variantes lectiones adnotavit, prolegomena addidit}, etc., 2 vols., Jenæ, 1843–1850. The first part contains the two Confessions of Gennadius, the Confession of Cyrillus Lucaris, the Confessio Orthodoxa, and the Acts of the Synod of Jerusalem. The second part, which is added by Weissenborn, contains the Confessio Metrophanis Critopuli, and the Decretum Synodi Constantinopolitanæ, 1672. Kimmel d. 1846.

W. Gass: \emph{Gennadius und Pletho, Aristotelismus und Platonismus in der griechischen Kirche, nebst einer Abhandlung über die Bestreitung des Islam im Mittelalter}, Breslau, 1844, in two parts. The second part contains, among other writings of Gennadius and Pletho, the two Confessions of Gennadius (1453) in Greek. By the same: Symbolik der griechischen Kirche, Berlin, 1872.

H. W. Blackmore: The Doctrine of the Russian Church, being the Primer or Spelling-book, the Shorter and Longer Catechisms, and a Treatise on the Duty of Parish Priests. Translated from the Slavono-Russian Originals, Aberdeen, 1845.

\xsection{The Seven Œcumenical Councils}

§ 11. The Seven Œcumenical Councils.

The entire Orthodox Greek or Oriental Church,\footnote{The full name of the Greek Church is '\emph{the Holy Oriental Orthodox Catholic Apostolic Church.}' The chief stress is laid on the title \emph{orthodox}. The name \textgreek{Γραικός,} used by Polybius and since as equivalent to the Latin \emph{Græcus,} was by the Greeks themselves always regarded as an exotic. Homer has three standing names for the Greeks: \emph{Danaoi, Argeioi}, and \emph{Achaioi}; also \emph{Panthellenes} and \emph{Panachaioi}. The ancient (heathen) Greeks called themselves \emph{Hellenes}, the modern (Slavonic) Greeks, till recently, \emph{Romans}, in distinction from the surrounding Turks. The Greek language, since the founding of the East Roman empire, was called \emph{Romaic.}} including the Greek Church in Turkey, the national Church in the kingdom of Greece, and the national Church of the Russian Empire, and embracing a membership of about eighty millions, adopts, in common with the Roman communion, the doctrinal decisions of the seven oldest œcumenical Councils, laying especial stress on the Nicene Council and Nicene Creed. These Councils were all summoned by Greek emperors, and controlled by Greek patriarchs and bishops. They are as follows:

I. The first Council of Nicæa, A.D. 325; called by Constantine M.

II. The first Council of Constantinople, A.D. 381; called by Theodosius M.

III. The Council of Ephesus, A.D. 431; called by Theodosius II.

IV. The Council of Chalcedon, A.D. 451; called by Emperor Marcian and Pope Leo I.

V. The second Council of Constantinople, A.D. 553; called by Justinian I.

VI. The third Council of Constantinople, A.D. 680; called by Constantine Pogonatus.

VII. The second Council of Nicæa, A.D. 787; called by Irene and her son Constantine.

The first four Councils are by far the most important, as they settled the orthodox faith on the Trinity and the Incarnation. The fifth Council, which condemned the Three (Nestorian) Chapters, is a mere supplement to the third and fourth. The sixth condemned Monothelitism. The seventh sanctioned the use and worship of images.\footnote{Worship in a secondary sense, or \textgreek{δουλεία,} including \textgreek{ἀσπασμὸς καὶ τιμητικὴ προσκύνησις,} but not that adoration or \textgreek{ἀληθινὴ λατρεία,} which belongs only to God. See Hefele, Conciliengeschichte, Bd. III. p. 440.}

To these the Greek Church adds the Concilium Quinisextum,\footnote{This Synod is called \emph{Quinisexta} or \textgreek{πενθέκτη,} because it was to be a supplement to the fifth and sixth œcumenical Councils, which had passed doctrinal decrees, but no canons of discipline. It is also called the second Trullan Synod, because it was held 'in Trullo,' a saloon of the imperial palace in Constantinople. The Greeks regard the canons of this Synod as the canons of the fifth and sixth œcumenical Councils, but the Latins never acknowledged the Quinisexta, and called it mockingly '\emph{erratica}.' As the dates of the Quinisexta are variously given 686, 691, 692, 712. Comp. Baronius, \emph{Annal}. ad ann. 692, No. 7, and Hefele, l.c. III. pp. 298 sqq.} held at Constantinople (in Trullo), A.D. 691 (or 692), and frequently also that held in the same city A.D. 879 under Photius the Patriarch; while the Latins reject these two Synods as schismatic, and count the Synod of 869 (the fourth of Constantinople), which deposed Photius and condemned the Iconoclasts, as the eighth œcumenical Council. But these conflicting Councils refer only to discipline and the rivalry between the Patriarch of Constantinople and the Pope of Rome.

The Greek Church celebrates annually the memory of the seven holy Synods, held during the palmy days of her history, on the first Sunday in Lent, called the 'Sunday of Orthodoxy,' when the service is made to reproduce a dramatic picture of an œcumenical Council, with an emperor, the patriarchs, metropolitans, bishops, priests, and deacons in solemn deliberation on the fundamental articles of faith. She looks forward to an eighth œcumenical Council, which is to settle all the controversies of Christendom subsequent to the great schism between the East and the West.

Since the last of the seven Councils, the doctrinal system of the Greek Church has undergone no essential change, and become almost petrified. But the Reformation, especially the Jesuitical intrigues and the crypto-Calvinistic movement of Cyril Lucar in the seventeenth century, called forth a number of doctrinal manifestoes against Romanism, and still more against Protestantism. We may divide them into three classes:

I. \emph{Primary} Confessions of public authority:

(\emph{a}) The 'Orthodox Confession,' or Catechism of Peter Mogilas, 1643, indorsed by the Eastern Patriarchs and the Synod of Jerusalem.

(\emph{b}) The Decrees of the Synod of Jerusalem, or the Confession of Dositheus, 1672.

To the latter may be added the similar but less important decisions of the Synods of Constantinople, 1672 (\emph{Responsio Dionysii}), and 1691 (on the Eucharist).

(\emph{c}) The Russian Catechisms which have the sanction of the Holy Synod, especially the Longer Catechism of Philaret (Metropolitan of Moscow), published by the synodical press, and generally used in Russia since 1839.

(\emph{d}) The Answers of Jeremiah, Patriarch of Constantinople, to certain Lutheran divines, in condemnation of the doctrines of the Augsburg Confession, 1576 (published at Wittenberg, 1584), were sanctioned by the Synod of Jerusalem, but are devoid of clearness and point, and therefore of little use.

II. \emph{Secondary} Confessions of a mere private character, and hence not to be used as authorities:

(\emph{a}) The two Confessions of Gennadius, Patriarch of Constantinople, 1453. One of them, purporting to give a dialogue between the Patriarch and the Sultan, is spurious, and the other has nothing characteristic of the Greek system.

(\emph{b}) The Confession of Metrophanes Critopulus, subsequently Patriarch of Alexandria, composed during his sojourn in Germany, 1625. It is more liberal than the primary standards.

III. Different from both classes is the Confession of Cyril Lucar, 1629, which was repeatedly condemned as heretical (Calvinistic), but gave occasion for the two most important expositions of Eastern orthodoxy.

We shall notice these documents in their historical order.

\xsection{The Confessions of Gennadius, A.D. 1453.}

§ 12. The Confessions of Gennadius, A.D. 1453.

J. C. T. Otto:  \emph{Des Patriarchen Gennadios von Konstantinopel Confession}, Wien, 1864 (35 pp.).

See also the work of Gass, quoted p. 43, on \emph{Gennadius and Pletho} (1844), and an article of Prof. Otto on the \emph{Dialogue ascribed to Gennadius}, in (Niedner's) \emph{Zeitschrift für historische Theologie} for 1850, III. 399–417.

The one or two Confessions which the Constantinopolitan Patriarch Gennadius handed to the Turkish Sultan Mahmoud or Mahomet II., in 1453, comprise only a very general statement of the ancient Christian doctrines, without entering into the differences which divide the Oriental Church from the Latin Communion; yet they have a historical importance, as reflecting the faith of the Greek Church at that time.

Georgius Scholarius, a lawyer and philosopher, subsequently called Gennadius, was among the companions and advisers of the Greek Emperor John VII., Palæologus, and the Patriarch Joasaph, when they, in compliance with an invitation of Pope Eugenius IV., attended the Council of Ferrara and Florence (A.D. 1438 and '39), to consider the reunion of the Eastern and Western Catholic Churches. Scholarius, though not a member of the Synod (being a layman at the time), strongly advocated the scheme, while his more renowned countryman, Georgius Gemistus, commonly called Pletho (d. 1453), opposed it with as much zeal and eloquence. Both were also antagonists in philosophy, Gennadius being an Aristotelian, Pletho a Platonist. The union party triumphed, especially through the influence of Cardinal Bessarion (Archbishop of Nicæa), who at last acceded to the Latin \emph{Filioque}, as consistent with the Greek \emph{per Filium}.\footnote{See, on the transactions of this Council, Mansi, Tom. XXXI., and Werner: \emph{Geschichte der apologetischen and polemischen Literatur}, Vol. III. pp. 57 sqq.}

But when the results of the Council were submitted to the Greek Church for acceptance, the popular sentiment, backed by a long tradition, almost universally discarded them. Scholarius, who in the mean time had become a monk, was compelled to give up his plans of reunion, and he even wrote violently against it. Some attribute this inconsistency to a change of conviction, some to policy; while others, without good reason, doubt the identity of the anti-Latin monk Scholarius with the Latinizing Gennadius.\footnote{\par Karyophilus, Allatius, and Kimmel deny the identity of the two persons; Robert Creygthon, Renaudot (1704), Richard Simon, Spanheim, and Gass defend it. Spanheim, however, regards the unionistic writings as interpolations. Allatius and Kimmel maintain that Gennadius continued friendly to the union as Patriarch, but Karyophilus supposes that the unionistic Scholarius died before the conquest of Constantinople, and never was Patriarch. See Kimmel, \emph{Monumenta}, etc., \emph{Prolegomena}, p. vi.; Gass, l.c. Vol. I. pp. 5 sqq., and Werner, l.c. Vol. III. pp. 67 sqq. Scholarius was a fertile writer of homilies, hymns, philosophical and theological essays. Four of these are edited in Greek by W. Gass, viz., his Confession, the Dialogue \emph{De via salutis}, the book \emph{Contra Automatistas et Hellenistas}, and the book \emph{De providentia et prædestinatione} (l.c. Vol. II. pp. 3–146).}

Immediately after the conquest in 1453, Scholarius was elected Patriarch of Constantinople, but held this position only a few years, as he is said to have abdicated in 1457 or 1459, and retired to a convent. This elevation is sufficient proof of his Greek orthodoxy, but may have been aided by motives of policy, inspired by the vain hope of securing, through his influence with the Latin church dignitaries, the assistance of the Western nations against the Turkish invasion.

At the request of the Mohammedan conqueror, Gennadius prepared a Confession of the Christian faith. The Sultan received it, invested Gennadius with the patriarchate by the delivery of the crozier or pastoral staff, and authorized him to assure the Greek Christians of freedom in the exercise of their religion.\footnote{An account of the interview is given in the \emph{Historia patriarcharum qui sederunt in hac magna catholicaque ecclesia Constantinopolitanensi postquam cepit eam Sultanus Mechemeta,} written in modern Greek by Emmanuel Malaxas, a Peloponnesian, and sent by him to Prof. M. Crusius, in Tübingen, who translated and published it in his \emph{Turco-Græcia}, 1584. Crusius and Chytræus were prominent in a fruitless effort to convert the Greek Church to Lutheranism.}

This 'Confession' of Gennadius,\footnote{Kimmel calls it the \emph{second} Confession, counting the Dialogue (which is of questionable authenticity; see below) as the first. But Gass more appropriately prints the Confession first, and the Dialogue afterwards, under its own proper title, \emph{De Via Salutis.}} or 'Homily on the true faith of the Christians,' was written in Greek, and translated into the Turko-Arabic (the Turkish with Arabic letters) for the use of the Sultan.\footnote{The title of the Vienna MS. as published by Otto is: \textgreek{Τοῦ αἰδεσιμωτάτου πατριάρχου Κωνσταντινουπόλεως} | \textgreek{ΓΕΝΝΑΔΙΟΥ ΣΧΟΛΑΡΙΟΥ} | \textgreek{Βιβλίον περὶ τινων κεφαλαίων τῆς ἡμετέρας} | \textgreek{πίστεως.} The title as given by Gass from a MS. in Munich reads: \textgreek{Τοῦ ἀγιωτάτου καὶ πατριάρχου καὶ φιλοσόφου} | \textgreek{ΓΕΝΝΑΔΙΟΥ} | \textgreek{ὁμιλία περὶ τῆς ὀρθῆς καὶ ἀληθοῦς} | \textgreek{πίστεως τῶν Χριστιανῶν.} In other titles it is called \textgreek{ὁμλογία} or \textgreek{ὁμολόγησις.} This Confession (together with the Dialogue on the \emph{Way of Life}) was first published in Greek at Vienna by Prof. John Alex. Brassicanus (Kohlburger), in 1530; then in Latin by J. Harold (in his \emph{Hæresiologia}, Basil. 1556, from which it passed into the Patristic Libraries, \emph{Bibl. P. P. Lugdun.} Tom. XXVI. 556, also \emph{B. P. P. Colon.} Tom. XIV. 376, and \emph{B. P. P. Par.} Tom. IV.); then in Greek and Latin by David Chytræus (in his \emph{Oratio de statu ecclesiarum hoc tempore in Græcia, Asia, Bœmia}, etc., Frankf. 1583, pp. 173 sqq.); and soon afterwards in Greek, Latin, and Turkish by Mart. Crusius of Tübingen (in his \emph{Turco-Græcia}, Basil. 1584, lib. II. 109 sqq.). The text of Crusius differs from the preceding editions. He took it from a copy sent to him, together with the Sultan's answer, by Emmanuel Malaxas. Two other editions of the Greek text were published by J. von Fuchten, Helmst. 1611, and by Ch. Daum, Cygneæ (Zwickau), 1677 \emph{(Hieronymi theologi Græci dialogus de Trinitate}, etc.). Kimmel followed the text of Chytræus, compared with that of Crusius and the different readings in the \emph{Bibl. Patr. Lugdun.} See his \emph{Proleg.} p xx. The last and best editions of the Greek text of the Confession are by Gass, l.c. II. 3–15, who used three MSS., and compared older Greek editions and Latin versions; and by Otto (1864), who (like Brassicanus) reproduced the text of the Vienna Codex after a careful re-examination, and added the principal variations of Brassicanus and Gass.} It treats, in twenty brief sections, of the fundamental doctrines on God, the Trinity, the two natures in the person of Christ, his work, the immortality of the soul, and the resurrection of the body. The doctrine of the Trinity is thus stated: 'We believe that there are in the one God three peculiarities (\textgreek{ἰδιώματα τρία}), which are the principles and fountains of all his other peculiarities . . . and these three peculiarities we call the three subsistences (\textgreek{ὑποστάεις}). . . . We believe that out of the nature (\textgreek{ἐκ τῆς φύσεως}) of God spring the Word (\textgreek{λόγος}) and the Spirit (\textgreek{πνεῦμα}), as from the fire the light and the heat (\textgreek{ὥσπερ ἀπὸ τοῦ πυρὸς φῶς καὶ θέρμη}). . . . These three, the Mind, the Word, and the Spirit (\textgreek{νοῦς, λόγος, πνεῦμα}), are one God, as in the one soul of man there is the mind (\textgreek{νοῦς}), the rational word (\textgreek{λόγος νοητός}), and the rational will (\textgreek{θέλησις νοητή}); and yet these three are as to essence but one soul (\textgreek{μία ψυχὴ κατὰ τὴν οὐσίαν}).'\footnote{Compare, on the Trinitarian doctrine of Gennadius and its relation to Latin Scholasticism, the exposition of Gass, I. 82 sqq. Kimmel and Otto (l.c. p. 400) make him a Platonist, but there are also some Aristotelian elements in him.} The difference of the Greek and Latin doctrine on the procession of the Holy Spirit is not touched in this Confession. The relation of the divine and human nature in Christ is illustrated by the relation of the soul and the body in man, both being distinct, and yet inseparably united in one person.

At the end (§ 14–20) are added, for the benefit of the Turks, seven arguments for the truth of the Christian religion, viz.:\footnote{This apologetic appendix is omitted in the editions of Brassicanus and Fuchten, and is rejected by Otto as a later addition (l.c. pp. 5–11).}

1. The concurrence of Jewish prophecies and heathen oracles in the pre-announcement of a Saviour.

2. The internal harmony and mutual agreement of the different parts of the Scriptures.

3. The acceptance of the gospel by the greatest and best men among all nations.

4. The spiritual character and tendency of the Christian faith, aiming at divine and eternal ends.

5. The ennobling effect of Christ's religion on the morals of his followers.

6. The harmony of revealed truth with sound reason, and the refutation of all objections which have been raised against it.

7. The victory of the Church over persecution and its indestructibility.

The other Confession, ascribed to Gennadius, and generally published with the first, is written in the form of a Dialogue ('\emph{Sermocinatio}') between the Sultan and the Patriarch, and entitled '\emph{The Way of Life}.'\footnote{De Via Salutis. The full title, as given by Gass, l.c. II. 16, and Otto, l.c. p. 409, reads:  \par  \textgreek{Τοῦ αἰδεσιμωτάτου πατριάρχου Κονσταντινουπόλεως } \par  \textgreek{ΓΕΝΝΑΔΙΟΥ ΣΧΟΛΑΡΙΟΥ } \par  \textgreek{Βιβλίον σύντομόν τε καὶ σαφὲς περὶ τινων κεφαλαίων τῆς ἡμετέρας πίστεως, περὶ ὦν ἡ διάλεξις γέγονε μετὰ Ἀμοιρᾶ τοῦ Μαχουμέτου, ὃ καὶ ἐπιγέγραπται } \par \textgreek{περὶ τῆς ὀδοῦ τῆς σωτηρίας (τῶν) ἀνθρώπων.} \par The tract was published three times in Greek in the seventeenth century—by Brassicanus, Vienna, 1530; by Joh. von Fuchten, Helmstädt, 1611 (or 1612); and by Daum, Zwickau, 1677; but each of these editions is exceedingly rare. The Latin version was repeated in several patristic collections, but with more or less omissions or additions (occasionally in favor of the Romish system). We have now two correct editions of the Greek text, one by Gass (1844), and another by Otto (1850; the latter was originally intended for an Appendix to Kimmel's collection). Kimmel gives only the Latin version, having been unable to obtain the Greek original (\emph{Proleg.} p. xx.), and seems to confound the special title with the joint title for both Confessions; see \emph{Bibl. P. P. Colon.} XIV. 378; Werner. l.c. III. 68. note. The Dialogue has also found its way into the writings of Athanasius (\emph{Opera}, Tom. II. 280. Patav. 1777, or II. 335, ed. Paris, 1698), but without a name or an allusion to the Sultan, simply as a dialogue between a Christian bishop and a catechumen, and with considerable enlargements and adaptations to the standard of Greek orthodoxy. Comp. Gass, I. pp. 89 sqq., II. pp. 16–30, and Otto, p. 407.} The Sultan is represented as asking a number of short questions, such as: 'What is God?' 'Why is he called God (\textgreek{θεός})' 'How many Gods are there?' 'How, if there is but one God, can you speak of three Divine Persons, Father, Son, and Holy Ghost?' 'Why is the Father called Father?' 'Why is the Son called Son?' 'Why is the Holy Spirit called Spirit?' To these the Patriarch replies at some length, dwelling mainly on the doctrine of the Trinity, and illustrating it by the analogy of the sun, light, and heat, and by the trinity of the human mind.

But there is no external evidence for the authorship of Gennadius; and the internal evidence is against it. There was no need of two Confessions for the same occasion. There is nothing characteristic of a Mohammedan in the questions of the Sultan. The text is more loose and prolix in style than the genuine Confession; it contains some absurd etymologies unworthy of Gennadius;\footnote{The word \textgreek{θεός,} is derived from \textgreek{θεωρεῖν} (\textgreek{ἀπὸ τοῦ θεωρεῖν τὰ πάντα οἱονεὶ θεωρός}), and also from \textgreek{θέειν, } \emph{percurrere} (\textgreek{ὁ γὰρ θεὸς ἀεὶ καὶ πανταχοῦ πάρεστιν}); \textgreek{πατήρ} is derived from \textgreek{τηρεῖν} (\textgreek{ἀπὸ τοῦ τὰ πάντα τηρεῖν}), \textgreek{υἱός} from \textgreek{οἷος,} \emph{talis} (\emph{qualis enim Pater, talis Filius}), \textgreek{πνεῦμα} from \textgreek{νοέω, } \emph{intelligo} (\textgreek{πάντα γὰρ ὀξέως ἐπινοεῖ}).} and it expressly teaches the Latin doctrine of the double procession of the Holy Spirit.\footnote{In the Latin Version (Kimmel, p. 3): '\emph{Quemadmodum substantia solis producit radios, et a sole et radiis procedit lumen: ita Pater generat Filium seu Verbum ejus, et}  a Patre et Filio Procedit Spiritus Sanctus.' In the Greek text (Gass, II. 19): \textgreek{Ὥσπερ ὁ δίσκος ὁ ἡλιακὸς γεννᾷ τὴν ἀκτῖνα, καὶ παρὰ τοῦ ἡλίου καὶ τῶν ἀκτίνων ἐκπορεύεται τὸ φῶς · οὕτω ὁ θεὸς καὶ πατὴρ γεννᾷ τὸν υἱὸν καὶ λόγον αὐτοῦ, καὶ ἐκ τοῦ πατρὸς καὶ υἱοῦ ἐκπορεύεται τὸ πνεῦμα τὸ ἅγιον.} A Greek Patriarch could not have maintained himself with such an open avowal of the Latin doctrine. The text of Pseudo-Athanasius urges the \emph{processio a solo Patre}, and removes all other approaches to the Latin dogma.} For these reasons, we must either deny the authorship of Gennadius, or the integrity of the received text.\footnote{See Gass, I. p. 100, and \emph{Symb. der griech. Kirche}, p. 38; Otto, p. 405. Both reject the authenticity of the Dialogue.} At all events, it can not be regarded in its present form even as a secondary standard of Greek orthodoxy.

\xsection{The Answers of the Patriarch Jeremiah to the Lutherans, A.D. 1576.}

§ 13. The Answers of Patriarch Jeremiah to the Lutherans, A.D. 1576.

\emph{Acta et Scripta theolog. Würtemberg. et Patriarchæ Constant.}  Hieremiæ, quoted p. 43.

Martin Crusius:  \emph{Turco-Græcia}, Basil. 1584.

Mouravieff: History of the Church of Russia, translated by Blackmore, pp. 289–324.

Hefele (now Bishop of Rottenburg): \emph{Ueber die alten und neuen Versuche, den Orient zu protestantisiren}, in the \emph{Tübinger Theol. Quartalschrift}, 1843, p. 544.

Art. \emph{Jeremias II.}, in Herzog's \emph{Encyklop.} 2d ed. Vol. VI. pp. 530–532. Gass: \emph{Symbolik d. gr. K.} pp. 41 sqq.

Melanchthon, who had the reunion of Christendom much at heart, especially in the later part of his life, first opened a Protestant correspondence with the Eastern Church by sending, through the hands of a Greek deacon, a Greek translation (made by Paul Dolscius) of the Augsburg Confession to Patriarch Joasaph II. of Constantinople, but apparently without effect.

Several years afterwards, from 1573–75, two distinguished professors of theology at Tübingen, Jacob Andreæ, one of the authors of the Lutheran 'Form of Concord' (d. 1590), and Martin Crusius, a rare Greek scholar (d. 1607),\footnote{He was able to take Andreæ's sermons down in Greek as they were delivered in German.} on occasion of the ordination of Stephen Gerlach for the Lutheran chaplaincy of the German legation at the Sublime Porte, forwarded to the Patriarch of Constantinople commendatory letters, and soon afterwards several copies of the Augsburg Confession in Greek (printed at Basle, 1559), together with a translation of some sermons of Andreæ, and solicited an official expression of views on the Lutheran doctrines, which they thought were in harmony with those of the Eastern Church.

At that time Jeremiah II. was Patriarch of Constantinople (from 1572–94), a prelate distinguished neither for talent or learning, but for piety and misfortune, and for his connection with the Russian Church at an important epoch of its history. He was twice arbitrarily deposed, saw the old patriarchal church turned into a mosque, and made a collecting tour through Russia, where he was received with great honor, and induced to confer upon the Metropolitan of Moscow the patriarchal dignity over Russia (1589), and thus to lay the foundation of the independence of the Russian Church.\footnote{Mouravieff gives an interesting account of this visit of Jeremiah, who styled himself 'by the grace of God, Archbishop of Constantinople, which is new Rome, and Patriarch of the whole universe.' He made his solemn entry into the Kremlin seated on an ass, and presented to the Czar several rich relics, among which are mentioned 'a gold Panagia [picture of the Virgin Mary], with morsels of the life-giving Cross, of the Robe of the Lord, and of that of the Mother of God, incased within it, as well as portions of the instruments of our Lord's Passion, the Spear, the Reed, the Sponge, and the Crown of Thorns.'}

After considerable delay, Jeremiah replied to the Lutheran divines at length, in 1576, and subjected the Augsburg Confession to an unfavorable criticism, rejecting nearly all its distinctive doctrines, and commending only its indorsement of the early œcumenical Synods and its view on the marriage of priests.\footnote{This third letter of Jeremiah is called \emph{Censura Orientalis Ecclesiæ}, and covers nearly ninety pages folio. His first two letters are brief, and do not enter into doctrinal discussions.} The Tübingen professors sent him an elaborate defense (1577), with other documents, but Jeremiah, two years afterwards, only reaffirmed his former position, and when the Lutherans troubled him with new letters, apologetic and polemic, he declined all further correspondence, and ceased to answer.\footnote{Vitus Myller, in his funeral discourse on Crusius, complains of the Greeks as being prouder and more superstitious than the Papists (\emph{pontificiis longe magis superstitiosi}). Crusius edited also a Greek translation of four volumes of Lutheran sermons (\emph{Corona anni, } \textgreek{στέφανος τοῦ ἐνιαυτοῦ,} Wittemb. 1603) for the benefit of the Greek people, but with no better success.}

The documents of both parties were published at Wittenberg, 1584.

The Answers of Jeremiah received the approval of the Synod of Jerusalem in 1672,\footnote{In Kimmel's \emph{Monumenta}, Vol. I. p. 378.} and may be regarded, therefore, as truly expressing the spirit of the Eastern Communion towards Protestantism. It is evident from the transactions of the Synod of Jerusalem that the Greek Church rejects Lutheranism and Calvinism alike as dangerous heresies.

The Anglican Church has since made several attempts to bring about an intercommunion with the orthodox East, especially with the Russo-Greek Church, during the reign of Peter the Great, and again in our own days, but so far without practical effect beyond the exchange of mutual courtesies and the expression of a desire for the reunion of orthodox Christendom.\footnote{See beyond, § 20.}

\xsection{The Confession of Metrophanes Critopulus, A.D. 1626.}

§ 14. The Confession of Metrophanes Critopulus, A.D. 1625.

Kimmel, Vol. II. pp. 1–213.

Dietelmaier: \emph{De Metrophane Critopulo}, etc., Altdorf, 1769.

Fabricius: \emph{Biblioth. Græca}, ed. Harless, Vol. XI. pp. 597–599.

Gass: Art. M. K. in Herzog's \emph{Encylop.} Vol. 2d ed. Vol. IX pp. 726–729.

Next in chronological order comes the Confession of Metrophanes Critopulus, once Patriarch of Alexandria, which was written in 1625, though not published till 1661.

Metrophanes Critopulus was a native of Berœa, in Macedonia, and educated at Mount Athos. Cyril Lucar, then Patriarch of Alexandria, sent him to England, Germany, and Switzerland (1616), with a recommendation to the Archbishop of Canterbury (George Abbot), that he might be thoroughly educated to counteract, in behalf of the Greek Church, the intrigues of the Jesuits.\footnote{See the letter in Kimmel, Preface to Vol. II. p. vii., and in Colomesii, \emph{Opera}, quoted there. On Cyril Lucar, see the next section.} The Archbishop kindly received him, and, with the consent of King James I., secured him a place in one of the colleges of Oxford. In 1620 Metrophanes visited the Universities of Wittenberg, Tübingen, Altdorf, Strasburg, and Helmstädt. He acquired good testimonials for his learning and character. He entered into close relations with Calixtus and a few like-minded Lutheran divines, who dissented from the exclusive confessionalism and scholastic dogmatism of the seventeenth century, and labored for Catholic union on the basis of the primitive creeds. At their request Metrophanes prepared a work on the faith and worship of the orthodox Greek Church. He also wrote a number of philological essays. After spending some time in Venice as teacher of the Greek language, he returned to the East, and became successor of Cyril Lucar in Alexandria. But he disappointed the hopes of his patron, and, as a member of the Synod of Constantinople, 1638, he even took part in his condemnation. The year of his death is unknown.

The Confession of Metrophanes\footnote{\textgreek{Ὁμολογία τῆς ἀνατολικῆς ἐκκλησίας τῆς καθολικῆς καὶ ἀποστολικῆς, συγγραφεῖσα ἐν ἐπιτομῇ διὰ Μητροφάνους Ἱερομονάχου Πατριαρχικοῦ τε Πρωτοσυγγέλλου τοῦ Κριτοπούλου.} \emph{Confessio catholicæ et apostolicæ in Orienti ecclesiæ, conscripta compendiose per}  Metrophanem Critopulum, \emph{Hieromonachum et Patriarchalem Protosyngellum.} It was first published in Greek, with a Latin translation, by \emph{J. Hornejus}, at Helmstädt. 1661. Kimmel compared with this ed. the MS. which is preserved in the library at Wolfenbüttel, but he died before his edition appeared, with a preface of Weissenborn (1850).} discusses, in twenty-three chapters, all the leading doctrines and usages of the Eastern Church. It is a lengthy theological treatise rather than a Confession of faith. It has never received ecclesiastical sanction, and is ignored by the Synod of Jerusalem; hence it ought not to be quoted as an authority, as is done by Winer and other writers on Symbolics. Nevertheless, as a private exposition of the Greek faith, it is of considerable interest.

Although orthodox in the main, it yet presents the more liberal and progressive aspect of Eastern theology. It was intended to give a truthful account of the Greek faith, but betrays the influence of the Protestant atmosphere in which it was composed. It is strongly opposed to Romanism, but abstains from all direct opposition to Protestantism, and is even respectfully dedicated to the Lutheran theological faculty of Helmstädt, where it was written.\footnote{Nicolaus Comnenus called Metrophanes a \emph{Græco-Lutheranus}, but without good reason.} In this respect it is the counterpart or complement of the Confession of Dositheus, which, in its zeal against Protestantism, almost ignores the difference from Romanism.\footnote{See below, § 17.} Thus Metrophanes excludes the Apocrypha from the canon, denies in name (though maintaining in substance) the doctrine of purgatory, and makes a distinction between sacraments proper, viz., baptism, eucharist, and penance, and a secondary category of sacramental or mystical rites, viz., confirmation (or chrisma), ordination, marriage, and unction.

\xsection{The Confession of Cyril Lucar, A.D. 1631.}

§ 15. The Confession of Cyril Lucar, A.D. 1631.

Literature.

Cyrilli Lucaris \emph{Confessio Christianæ fidei}, Latin, 1629; \emph{c. additam. Cyrilli}, Gr. et Lat., Genev. 1633; (? Amst.) 1645, and often; also in Kimmel's  \emph{Monumenta fidei Ecclesiæ Orient.} P. I. pp. 24–44. Compare \emph{Proleg.} pp. xxi.–l. (\emph{de vita Cyrilli}).

Thom. Smith: \emph{Collectanea de Cyrillo Lucari}, London, 1707. Comp. also, in Th. Smith's \emph{Miscellanea} (Hal. 1724), his \emph{Narratio de vita, studiis, gestis et martyrio C. Lucaris.}

Leo Allatius (d. at Rome, 1669): \emph{De Ecclesiæ Occidentalis atque Orientalis perpetua consensione, libri tres} (III. 11), Gr. et Lat. Colon. 1648. Bitter and slanderous against Cyril.

J. H. Hottinger: \emph{Analecta hist. theol. Dissert. VIII.}, Appendix, Tigur. 1653 (al. 1652). Against him, L. Allatius: \emph{J. H. Hottingerus, fraudis et imposturæ manifestæ convictus}, Rom. 1661.

J. Aymon:  \emph{Lettres anecdotes de Cyrille Lucaris}, Amsterd. 1718.

Bohnstedt:  \emph{De Cyrillo Lucari}, Halle, 1724.

Mohnike: On Cyril, in the \emph{Studien und Kritiken}, 1832, p. 560.

Several articles on Cyril Lucar, in the \emph{British Magazine} for Sept. 1842, Dec. 1843, Jan. and June, 1844.

Twesten: On Cyril, in the \emph{Deutsche Zeitechr. f. christl. Wissensch. u. chr. Leben}, Berl. 1850, No. 39, p. 305.

W. Gass: Article '\emph{Lukaris},' in Herzog's \emph{Encyklop.} 2d ed. Vol. IX. pp. 5 sqq.; and \emph{Symbolik}, pp. 50 sqq.

Aloysius Pichler (Rom. Cath.): \emph{Der Patriarch Cyrillus Lucaris und seine Zeit}, München, 1862, 8vo. (The author has since joined the Greek Church.)

The Confession of Cyril Lucar was never adopted by any branch or party of the Eastern Church, and even repeatedly condemned as heretical; but as it gave rise to the later authentic definitions of the 'Orthodox Faith,' in opposition to the distinctive doctrines of Romanism and Protestantism, it must be noticed here.

Cyrillus Lucaris (Kyrillos Loukaris\footnote{Properly 'the son of Lucar,' hence \textgreek{τοῦ Λουκάρεως.} The word \textgreek{λοῦκαρ} in later Greek is the Latin \emph{lucar,} or \emph{lucrum,} stipend, pay, profit, whence the French and English \emph{lucre.}}), a martyr of Protestantism within the orthodox Greek Church, occupies a remarkable position in the conflict of the three great Confessions to which the Reformation gave rise. He is the counterpart of his more learned and successful, but less noble, antagonist, Leo Allatius (1586–1669), who openly apostatized from the Greek Church to the Roman, and became librarian of the Vatican. His work is a mere episode, and passed away apparently without permanent effect, but (like the attempted reformations of Wyclif, Huss, and Savonarola) it may have a prophetic meaning for the future, and be resumed by Providence in a better form.

Cyril Lucar was born in 1568 or 1572 in Candia (Crete), then under the sovereignty of Venice, and the only remaining seat of Greek learning. He studied and traveled extensively in Europe, and was for a while rector and Greek teacher in the Russian Seminary at Ostrog, in Volhynia. In French Switzerland he became acquainted with the Reformed Church, and embraced its faith. Subsequently he openly professed it in a letter to the Professors of Geneva (1636), through Leger, a minister from Geneva, who had been sent to Constantinople. He conceived the bold plan of ingrafting Protestant doctrines on the old œcumenical creeds of the Eastern Church, and thereby reforming the same. He was unanimously elected Patriarch of Alexandria in 1602 (?), and of Constantinople in 1621. While occupying these high positions he carried on an extensive correspondence with Protestant divines in Switzerland, Holland, and England, sent promising youths to Protestant universities, and imported a press from England (1629) to print his Confession and several Catechisms. But he stood on dangerous ground, between vacillating or ill-informed friends and determined foes. The Jesuits, with the aid of the French embassador at the Sublime Porte, spared no intrigues to counteract and checkmate his Protestant schemes, and to bring about instead a union of the Greek hierarchy with Rome. At their instigation his printing-press was destroyed by the Turkish government. He himself—in this respect another Athanasius '\emph{versus mundum},' though not to be compared in intellectual power to the 'father of orthodoxy'—was five times deposed, and five times reinstated. At last, however—unlike Athanasius, who died in peaceful possession of his patriarchal dignity—he was strangled to death in 1638, having been condemned by the Sultan for alleged high-treason, and his body was thrown into the Bosphorus. His friends surrounded the palace of his successor, Cyril of Berœa, crying, 'Pilate, give us the dead, that we may bury him.'\footnote{\textgreek{Πίλατε, δὸς ἡμῖν τὸν νεκρόν, ἵνα αὐτὸν θάψωμεν.}} The corpse was washed ashore, but it was only obtained by Cyril's adherents after having been once more cast out and returned by the tide. The next Patriarch, Parthenius, granted him finally an honorable burial.

Cyril left no followers able or willing to carry on his work, but the agitation he had produced continued for several years, and called forth defensive measures. His doctrines were anathematized by Patriarch Cyril of Berœa and a Synod of Constantinople (Sept., 1638),\footnote{Cyril of Berœa seemed to assume the authenticity of Cyril's Confession. He was, however, himself afterwards deposed and anathematized on the charge of extortion and embezzlement of ecclesiastical funds, and for the part he took in procuring the death of Cyril Lucar by preferring false accusation against him to the Turks. See Mouravieff, \emph{Hist. of the Church of Russia}, translated by Blackmore, p. 396. Blackmore, however, gives there a wrong date, assigning the death of Cyril to 1628 instead of 1638.} then again by the Synods of Jassy, in Moldavia, 1643, and of Jerusalem, 1672; but on the last two occasions the honor of his name and the patriarchal dignity were saved by boldly denying the authenticity of his Confession, and contradicting it by written documents from his pen.\footnote{The Synods of Jassy and Jerusalem intimate that Cyril's Confession was a Calvinistic forgery, and the Synod of Jerusalem quotes largely from his homilies to prove his orthodoxy. Mouravieff, l.c. p. 189, adopts a middle view, saying: 'Cyril, although he had condemned the new doctrine of Calvin, nevertheless had not stood up decidedly and openly to oppose it, and for his neglect he was himself delivered over to an anathema by his successor, Cyril of Berœa.'}

This Cyril was the same who seat the famous uncial Codex Alexandrinus of the Bible (A) to King Charles I. of England,\footnote{Not to James I. (who died 1625), as Kimmel and Gass wrongly state. Cyril brought the Codex with him from Alexandria, or, according to another report, from Mount Athos, and sent it to England in 1628, where it passed from the king's library into the British Museum, 1753. It dates from the fifth century, and contains the Septuagint Version of the Old Testament, the whole New Testament, with some chasms, and, as an Appendix, the only MS. copy extant of the first Epistle of Clemens Romanus to the Corinthians, with a fragment of a second Epistle. The New Test. has been edited in quasi-fac-simile, by Woide, Lond. 1786, fol., and in ordinary Greek type by Cowper, Lond. 1860.} and who translated the New Testament into the modern Greek language.\footnote{Published at Geneva or Leyden, 1638, and at London, 1703.}

The Confession of Cyril was first written by him in Latin, 1629, and then in Greek, with an addition of four questions and answers, 1631, and published in both languages at Geneva, 1633.\footnote{The Latin edition was first published in 1529, either at the Hague (by the Dutch embassador Cornelius Van der Haga) or at Geneva, or at both places; the authorities I have consulted differ. The subscription to the Græco-Latin edition before me reads: '\emph{Datum Constantinopoli mense Januario} 1631 \emph{ Cyrillus Patriarcha Constantinopoleos.}' Another edition (perhaps by Hugo Grotius) was published 1645, without indication of place (perhaps at Amsterdam). I have used Kimmel's edition, which gives the text of the edition of 1645.} It expresses his own individual faith, which he vainly hoped would become the faith of the Greek Church. It is divided into eighteen brief chapters, each fortified with Scripture references; eight chapters contain the common old Catholic doctrine, while the rest bear a distinctly Protestant character.

In Chapter I. the dogma of the Trinity is plainly stated in agreement with the œcumenical creeds, the procession of the Spirit in the conciliatory terms of the Council of Florence.\footnote{'\emph{Spiritus Sanctus a Patre }Per Filium  \emph{procedens,}' \textgreek{ἐκ τοῦ πατρὸς δἰ υἱοῦ.}} Chapters IV. and V. treat of the doctrines of creation and divine government; Chapter VI., of the fall of man; Chapters VII. and VIII., of the twofold state of Christ, his incarnation and humiliation, and his exaltation and sitting on the right hand of the Father, as the Mediator of mankind and the Ruler of his Church (\emph{status exinanitionis} and \emph{st. exaltationis}); Chapter IX., of faith in general; Chapter XVI., of baptismal regeneration.

The remaining ten chapters breathe the Reformed spirit. Chapter II. asserts that 'the authority of the Scriptures is superior to the authority of the Church,' since the Scriptures alone, being divinely inspired, can not err.\footnote{'\emph{Credimus Scripturam sacram esse } \textgreek{θεοδίδακτον } (i. e., \emph{a Deo traditam}) \emph{habereque auctorem Spiritum Sanctum, non alium, cui habere debemus fidem indubitam. . . . Propterea ejus auctoritatem esse superiorem Ecclesiæ auctoritate; nimis enim differens est, loqui Spiritum Sanctum et linguam humanam, quum ista possit per ignorantiam errare, fallere et falli, Scriptura vero divina nec fallitur, nec errare potest, sed est infallibilis semper et certa.}'} In the appendix to the second (the Greek) edition, Cyril commends the general circulation of the Scriptures, and maintains their perspicuity in matters of faith, but excludes the Apocrypha, and rejects the worship of images. He believes 'that the Church is sanctified and taught by the Holy Spirit in the way of life,' but denies its infallibility, saying: 'The Church is liable to sin (\textgreek{ἁμαρτάνειν)}, and to choose the error instead of the truth (\textgreek{ἀντὶ τῆς ἀληθείας τὸ ψεῦδος ἐκλέγεσθαι)}; from such error we can only be delivered by the teaching and the light of the Holy Spirit, and not of any mortal man' (Ch. XII.). The doctrine of justification (Chapter XIII.) is stated as follows:

'We believe that man is justified by faith, not by works. But when we say ``by faith,'' we understand the correlative of faith, viz., the Righteousness of Christ, which faith, fulfilling the office of the hand, apprehends and applies to us for salvation. And this we understand to be fully consistent with, and in no wise to the prejudice of, works; for the truth itself teaches us that works also are not to be neglected, and that they are necessary means and testimonies of our faith, and a confirmation of our calling. But, as human frailty bears witness, they are of themselves by no means sufficient to save man, and able to appear at the judgment-seat of Christ, so as to merit the reward of salvation. The righteousness of Christ, applied to the penitent, alone justifies and saves the believer.'

The freedom of will before regeneration is denied (Ch. XIV.).\footnote{\textgreek{Πιστεύομεν ἐν τοῖς οὐκ ἀναγεννηθεῖσι τὸ αὐτεξούσιον νεκρὸν εἶναι.} This is in direct opposition to the traditional doctrine of the Greek Church, which emphasizes the liberum arbitrium even more than the Roman, and was never affected by the Augustinian anthropology.} In the doctrine of decrees, Cyril agrees with the Calvinistic system (Ch. III.), and thereby offended Grotius and the Arminians. He accepts, with the Protestants, only two sacraments as being instituted by Christ, instead of seven, and requires faith as a condition of their application (Ch. XV.). He rejects the dogma of transubstantiation and oral manducation, and teaches the Calvinistic theory of a real but spiritual presence and fruition of the body and blood of Christ by believers only (Ch. XVII.). In the last chapter he rejects the doctrine of purgatory and of the possibility of repentance after death.

\xsection{The Orthodox Confession of Mogilas, A.D. 1643.}

§ 16. The Orthodox Confession of Mogilas, A.D. 1643.

The Orthodox Confession of the Catholic and Apostolic Eastern Church\footnote{\par  \textgreek{Ὀρθόδοξος ὁμολογία τῆς καθολικῆς καὶ ἀποστολικῆς ἐκκλησίας τῆς ἀνατολικῆς.} It is uncertain whether it was first written in Greek or in Russ. First published in Greek by Panagiotta, Amst. 1662; then in Greek and Latin by Bishop Normann, of Gothenburg (then Professor at Upsala), Leipz. 1695; in Greek, Latin, and German by C. G. Hofmann, Breslau, 1751; by Patriarch Adrian in Russian, Moscow, 1696, and again in 1839, etc.; in Kimmel's \emph{Momum.} I. 56–324 (Greek and Latin, with the letters of Nectarius and Parthenius). Comp. Kimmel's \emph{Proleg.} pp. lxii. sqq. The Confession must not be confounded with the \emph{Short Russian Catechism} by the same author (Peter Mogilas).} was originally drawn up about the year 1640 by Peter Mogilas (or Mogila), Metropolitan of Kieff, and father of Russian theology (died 1647), in the form of a Catechism for the benefit of the Russian Church.\footnote{The following account of Mogilas is translated from the Russian of Bolchofsky by Blackmore (\emph{The Doctrine of the Russian Church}, p. xviii.): 'Peter Mogila belonged by birth to the family of the Princes of Moldavia, and before he became an ecclesiastic had distinguished himself as a soldier. After having embraced the monastic life, he became first Archimandrite of the Pechersky, and subsequently, in 1632, Metropolitan of Kieff, to which dignity he was ordained by authority of Cyril Lucar [then Patriarch of Constantinople], with the title of Eparch, or Exarch of the Patriarchal See. He sat about fifteen years, and died in 1647. Besides the \emph{Orthodox Confession}, he put out, in 1645, in the dialect of Little Russia, his \emph{Short Catechism}; composed a Preface prefixed to the \emph{Patericon}; corrected, in 1646, from Greek and Slavonic MSS., the \emph{Trebnik}, or Office-book, and added to each Office doctrinal, casuistical, and ceremonial instructions. He also caused translations to be made from the Greek \emph{Lives of the Saints}, by Metaphrastus, though this work remained unfinished at his death; and, lastly, he composed a Short Russian Chronicle, which is preserved in MS., but has never yet been printed. He was the founder of the first Russian Academy at Kieff.' It was called, after him, the Kievo-Mogilian Academy. He also founded a library and a printing-press. See a fuller account of Peter Mogilas in Mouravieff's \emph{History of the Church of Russia}, translated by Blackmore (Oxford, 1842), pp. 186–189. It is there stated that he received his education in the University of Paris. This accounts for the tinge of Latin scholasticism in his Confession.} It was revised and adopted by a Provincial Synod at Kieff for Russia, then again corrected and purged by a Synod of the Greek and Russian clergy at Jassy, in 1643, where it received its present shape by  Meletius Syriga, or Striga, the Metropolitan of Nicæa, and exarch of the Patriarch of Constantinople. As thus improved, it was sent to, and signed by, the four Eastern Patriarchs. The Synod of Jerusalem gave it a new sanction in 1672 (declaring it a \textgreek{ὁμολογία, ἣν ἐδέξατο καὶ δέχεται ἁπαξαπλῶς πᾶσα ἡ ἀνατολικὴ ἐκκλησία}). In this way it became the Creed of the entire Greek and Russian Church. It has been the basis of several later Catechisms prepared by Russian divines.

The Orthodox Confession was a defensive measure against Romanism and Protestantism. It is directed, first, against the Jesuits who, under the protection of the French embassadors in Constantinople, labored to reconcile the Greek Church with the Pope; and, secondly, against the Calvinistic movement, headed by Cyril Lucar, and continued after his death.\footnote{\par See § 15. Mouravieff, in his \emph{Hist. of the Church of Russia}, p. 188, distinctly asserts that the Confession was directed both against the Jesuits and against 'the Calvinistic heresy,' which, 'under the name of Cyril Lucar, Patriarch of Constantinople,' had been disseminated in the East by 'crafty teachers.' As Cyril and the Calvinists are not mentioned by name in the Orthodox Confession, another Russian writer, quoted by Blackmore (\emph{The Doctrine of the Russian Church}, p. xx.), thinks that Mogilas wrote against the Lutherans rather than the Calvinists; adding, however, that it is chiefly directed against the Papists, from whom danger was most apprehended.}

It is preceded by a historical account of its composition and publication, a pastoral letter of Nectarius, Patriarch of Jerusalem, dated Nov. 20, 1662; and by a letter of indorsement of the Greek text from Parthenius, Patriarch of Constantinople, dated March 11, 1643,\footnote{\par This is the date (\textgreek{αχμγ}) given by Kimmel, P. I. p. 53, and the date of the Synod of Jassy, where the Confession was adopted. Butler (\emph{Hist. Acc. of Conf. of Faith}, p. 101) gives the year 1663; but the Confession was already published in 1662 with the letters of the two Patriarchs. See Kimmel, \emph{Proleg.} p. lxii.} followed by the signatures of twenty-six Patriarchs and prelates of the Eastern Church.

The letter of Parthenius is as follows:

'Parthenius, by the mercy of God, Archbishop of Constantinople, New Rome, and Œcumenical Patriarch. Our mediocrity,\footnote{\par  \textgreek{ἡ μετριότης ἡ μῶν,} a title of proud humility, like the papal '\emph{servus servorum Dei,}' which dates from Gregory I.}together with our sacred congregation of chief bishops and clergy present, has diligently perused a small book, transmitted to us from our true sister, the Church of Lesser Russia, entitled ``\emph{The Confession of the Orthodox Faith of the Catholic and Apostolic Church of Christ},'' in which the whole subject is treated under the three heads of \emph{Faith, Love}, and \emph{Hope}, in such a manner that \emph{Faith} is divided into twelve articles, to wit, those of the sacred [Nicene] Symbol; \emph{Love} into the Ten Commandments, and such other necessary precepts as are contained in the sacred and divinely inspired books of the Old and New Testaments; \emph{Hope} into the Lord's Prayer and the nine Beatitudes of the holy Gospel.

'We have found that this book follows faithfully the dogmas of the Church of Christ, and agrees with the sacred canons, and in no respect differs from them. As to the other part of the book, that which is in the Latin tongue, on the side opposite to the Greek text, we have not perused it, so that we only formally confirm that which is in our vernacular tongue. With our common synodical sentence, we decree, and we announce to every pious and orthodox Christian subject to the Eastern and Apostolic Church, that this book is to be diligently read, and not to be rejected. Which, for the perpetual faith and certainty of the fact, we guard by our subscriptions. In the year of salvation 1643, 11th day of March.'

The Confession itself begins with three preliminary questions and answers. Question first: 'What must an orthodox and Catholic Christian man observe in order to inherit eternal life?' Answer: 'Right faith and good works (\textgreek{πίστιν ὀρθὴν καὶ ἔργα καλά}); for he who observes these is a good Christian, and has the hope of eternal salvation, according to the sacred Scriptures (James ii. 24): ``Ye see, then, how that by works a man is justified, and not by faith only;'' and a little after (v. 26): ``For as the body without the spirit is dead, so faith without works is dead also.'' The divine Paul adds the same in another place (1 Tim. i. 19): ``Holding faith and a good conscience; which some having put away, concerning faith have made shipwreck;'' and, in another place, he says (1 Tim. iii. 9): ``Holding the mystery of the faith in a pure conscience.''' This is essentially the same with the Roman Catholic doctrine. It is characteristic that no passage is cited from the Romans and Galatians, which are the bulwark of the evangelical Protestant view of justification by faith. The second Question teaches that faith must precede works, because it is impossible to please God without faith (Heb. xi. 6). The third Question treats of the division of the Catechism according to the three theological virtues, faith, hope, and charity.

The Catechism is therefore divided into three parts.

1. Part first treats of \emph{Faith} (\textgreek{περὶ πίστεως}), and explains the Nicene Creed, which is divided into twelve articles, and declared to contain all things pertaining to our faith so accurately 'that we should believe nothing more and nothing less, nor in any other sense than that in which the fathers [of the Councils of Nicæa and Constantinople] understood it' (Qu. 5). The clause \emph{Filioque} is, of course, rejected as an unwarranted Latin interpolation and corruption (Qu. 72).

2. Part second treats of \emph{Hope} (\textgreek{περὶ ἐλπίδος}), and contains an exposition of the Lord's Prayer and the (nine) Beatitudes (Matt. v. 3–11).

3. Part third treats of \emph{Love} to God and man (\textgreek{περὶ τῆς εἰς θεὸν καὶ τὸν πλησίον ἀγάπης}), and gives an exposition of the Decalogue; but this is preceded by forty-five questions on the three cardinal virtues of prayer, fasting, and almsgiving, and the four general virtues which flow out of them (prudence, justice, fortitude, and temperance), on mortal and venial sins, on the seven general mortal sins (pride, avarice, fornication, envy, gluttony, desire of revenge, and sloth), on the sins against the Holy Ghost (presumption or temerity, despair, persistent opposition to the truth, and renouncing of the Christian faith), and on venial sins. In the division of the Ten Commandments the Greek Confession agrees with the Reformed Church in opposition to the Roman and Lutheran Churches, which follow the less natural division of Augustine by merging the second commandment in the first, and then dividing the tenth.

\xsection{The Synod of Jerusalem, and the Confession of Dositheus, A.D. 1672.}

§ 17. The Synod of Jerusalem and the Confession of Dositheus, A.D. 1672.

Hardouin:  \emph{Acta Conciliorum} (Paris, 1715), Tom. XI. pp. 179–274.

Kimmel: \emph{Monumenta Fidei Ecclesiæ Orientalis}, P. I. pp. 325–488; \emph{Prolegomena}, pp. lxxv.–xcii.

On the Synod of Jerusalem, comp. also Ittig: \emph{Dissert. de Actis Synodi Hieros. a.} 1672 \emph{sub Patr. Hiers. Dositheo adv. Calvinistas habitæ}, Lips. 1696. Aymon: \emph{Monuments authentiques de la religion des Grecs}, à la Haye, 1708. Basnage: \emph{Hist. de la religion des églises réformées}, P. I. ch. xxxii. J. Covel: \emph{Account of the present Greek Church}, Bk. I. ch. v. Schroeckh: \emph{Kirchengeschichte seit der Reformation}, Bd. ix. (by Tzschirner), pp. 90–96. Gass: \emph{Symb. der griech. Kirche}, pp. 79-84.

The Synod convened at Jerusalem in March, 1672, by Patriarch Dositheus, for the consecration of the restored Church of the Holy Nativity in Bethlehem,\footnote{\par Hence it is sometimes called the Synod of Bethlehem, but it was actually held at Jerusalem.} issued a new Defense or Apology of Greek Orthodoxy. It is directed against Calvinism, which was still professed or secretly held by many admirers of Cyril Lucar. It is dated Jerusalem, March 16, 1672, and signed by Dositheus, Patriarch of Jerusalem and Palestine (otherwise little known), and by sixty-eight Eastern bishops and ecclesiastics, including some from Russia.\footnote{Its title is \textgreek{Ασπὶς ὀρθδοξίας ἢ ἀπολογία καὶ ἔλεγχος πρὸς τοὺς διασύροντας τὴν ἀνατολικὴν ἐκκλησίαν αἱρετικῶς φρονεῖν ἐν τοῖς περὶ θεοῦ καὶ τῶν θείων, κ.τ.λ. } \emph{Clypeus orthodoxæ fidei sive Apologia adversus Calvinistas hæreticos, Orientalem ecclesiam de Deo rebusque divinis hæretice cum ipsis sentire mentientes.} The first edition, Greek and Latin, was published at Paris, 1676; then revised, 1678; also by Hardouin, and Kimmel, l.c.}

This Synod is the most important in the modern history of the Eastern Church, and may be compared to the Council of Trent. Both fixed the doctrinal status of the Churches they represent, and both condemned the evangelical doctrines of Protestantism. Both were equally hierarchical and intolerant, and present a strange contrast to the first Synod held in Jerusalem, when 'the apostles \emph{and elders},' in the presence of 'the brethren,' freely discussed and adjusted, in a spirit of love, without anathemas, the great controversy between the Gentile and the Jewish Christians. The Synod of Jerusalem has been charged by Aymon and others with subserviency to the interests of Rome; Dositheus being in correspondence with Nointel, the French embassador at Constantinople. The Synod was held at a time when the Romanists and Calvinists in France fiercely disputed about the Eucharist, and were anxious to secure the support of the Greek Church. But although the Synod was chiefly aimed against Protestantism, and has no direct polemical reference to the Latin Church, it did not give up any of the distinctive Greek doctrines, or make any concessions to the claims of the Papacy.

The acts of the Synod of Jerusalem consist of six chapters, and a confession of Dositheus in eighteen decrees. Both are preceded by a pastoral letter giving an account of the occasion of this public confession in opposition to Calvinism and Lutheranism, which are condemned alike as being essentially the same heresy, notwithstanding some apparent differences.\footnote{\par  \textgreek{Ἄδελφὰ φρονεῖ Λουθῆρος Καλουΐνῳ, εἰ καὶ ἐν τισι διαφέρειν δοκοῦσιν }. '\emph{Non alia est Lutheri hæresis atque Calvini, quamquam nonnihil videtur interesse}' (Kimmel, P. I. p. 335).} The Answers of Patriarch Jeremiah given to Martin Crusius, Professor in Tübingen, and other Lutherans, in 1572, are approved by the Synod of Jerusalem, as they were by the Synod of Jassy, and thus clothed with a semi-symbolical authority. The Orthodox Confession of Peter Mogilas is likewise sanctioned again, but the Confession of Cyril Lucar is disowned as a forgery.

The \emph{Six Chapters} are very prolix, and altogether polemical against the Confession which was circulated under the name of Cyril Lucar, and give large extracts from his homilies preached before the clergy and people of Constantinople to prove his orthodoxy. One anathema is not considered sufficient, and a threefold anathema is hurled against the heretical doctrines.

The \emph{Confessio Dosithei} presents, in eighteen decrees or articles,\footnote{\par  \textgreek{Ὅρος,} decree, decision. It is translated \emph{capitulum} in Hardouin, \emph{decretum} in Kimmel.} a positive statement of the orthodox faith. It follows the order of Cyril's Confession, which it is intended to refute. It is the most authoritative and complete doctrinal deliverance of the modern Greek Church on the controverted articles. It was formally transmitted by the Eastern Patriarchs to the Russian Church in 1721, and through it to certain Bishops of the Church of England, as an ultimatum to be received without further question or conference by all who would be in communion with the Orthodox Church. The eighteen decrees were also published in a Russian version (1838), but with a number of omissions and qualifications,\footnote{Under the title '\emph{Imperial and Patriarchal Letters on the Institution of the Most Holy Synod, with an Exposition of the Orthodox Faith of the Catholic Church of the East}.' See Blackmore, l.c. p. xxviii. Blackmore (pp. xxvi. and xxvii.) gives also two interesting letters of 'the Most Holy Governing Synod of the Russian Church to the Most Reverend the Bishops of the Remnant of the Catholic Church in Great Britain, our Brethren most beloved in the Lord, 'in answer to letters of two Non-Jurors and two Scotch Bishops seeking communion with the Eastern Church. Comp. § 20.} showing that, after all, the Russian branch of the Greek Church reserves to itself a certain freedom of further theological development. We give them here in a condensed summary from the original Greek:

\emph{Article I.—}The doctrine of the Holy Trinity, with the single procession of the Spirit. (\textgreek{Πνεῦμα ἅγιον ἐκ τοῦ πατρὸς ἐκπορευόμενον.} Against the Latins.)

\emph{Article II.}—The Holy Scriptures must be interpreted, not by private judgment, but in accordance with the tradition of the Catholic Church, which can not err, or deceive, or be deceived, and is of equal authority with the Scriptures. (Essentially Romish, but without an infallible, visible head of the Church.)

\emph{Article III.}—God has from eternity predestinated to glory those who would, in his foreknowledge, make good use of their free will in accepting the salvation, and has condemned those who would reject it. The Calvinistic doctrine of unconditional predestination is condemned as abominable, impious, and blasphemous.

\emph{Article IV.}—The doctrine of creation. The triune God made all things, visible and invisible, except sin, which is contrary to his will, and originated in the Devil and in man.

\emph{Article V.}—The doctrine of Providence. God foresees and permits (but does not foreordain) evil, and overrules it for good.

\emph{Article VI.}—The primitive state and fall of man. Christ and the Virgin Mary are exempt from sin.

\emph{Article VII.}—The doctrine of the incarnation of the Son of God, his death, resurrection, ascension, and return to judgment.

\emph{Article VIII.}—The work of Christ. He is the only Mediator and Advocate for our sins; but the saints, and especially the immaculate Mother of our Lord, as also the holy angels, bring our prayers and petitions before him, and give them greater effect.

\emph{Article IX.}—No one can be saved without faith, which is a certain persuasion, and works by love (i.e. the observance of the divine commandments). It justifies before Christ, and without it no one can please God.

\emph{Article X.}—The holy Catholic and Apostolic Church comprehends all true believers in Christ, and is governed by Christ, the only head, through duly ordained bishops in unbroken succession. The doctrine of Calvinists, that bishops are not necessary, or that priests (presbyters) may be ordained by priests, and not by bishops only, is rejected.

\emph{Article XI.}—Members of the Catholic Church are all the faithful, who firmly hold the faith of Christ as delivered by him, the apostles, and the holy synods, although some of them may be subject to various sins.

\emph{Article XII.}—The Catholic Church is taught by the Holy Ghost, through prophets, apostles, holy fathers, and synods, and therefore can not err, or be deceived, or choose a lie for the truth. (Against Cyril; comp. Art. II.)

\emph{Article XIII.}—Man is justified, not by faith alone, but also by works.

\emph{Article XIV.}—Man has been debilitated by the fall, and lost the perfection and freedom from suffering, but not his intellectual and moral nature. He has still the free will (\textgreek{τὸ αὐτεξούσιον}) or the power to choose and do good or to flee and hate evil (Matt. v. 46, 47;  Rom. i. 19;  ii. 14, 15). But good works done without faith can not contribute to our salvation; only the works of the regenerate, done under grace and with grace, are perfect, and render the one who does them worthy of salvation (\textgreek{σωτηρίας ἄξιον ποιεῖται τὸν ἐνεργοῦντα}).

\emph{Article XV.}—Teaches, with the Roman Church, the seven sacraments or mysteries (\textgreek{μυστὴρια}), viz., baptism (\textgreek{τὸ ἅγιον βάπτισμα, }  Matt. xxviii. 19), confirmation (\textgreek{βεβαίωσις } or \textgreek{χρίσμα,}  Luke xxiv. 49;  2 Cor. i. 21; and Dionysius Areop.), ordination (\textgreek{ἱεροσύνη,}  Matt. xviii. 18), the unbloody sacrifice of the altar (\textgreek{ἡ ἀναίμακτος θυσία,}  Matt. xxvi. 26, etc.), matrimony (\textgreek{γάμος,}  Matt. xix. 6;  Eph. v. 32), penance and confession (\textgreek{μετάνοια καὶ ἐξομολόγησις,}  John xx. 23;  Luke xiii. 3, 5), and holy unction (\textgreek{τὸ ἅγιον ἔλαιον} or \textgreek{εὐχέλαιον,}  Mark vi. 13;  James v. 14). Sacraments are not empty signs of divine promises (as circumcision), but they necessarily (\textgreek{ἐξ ἀνάγκης}) confer grace (as \textgreek{ὄργανα δραστικὰ χάριτος}).

\emph{Article XVI.}—Teaches the necessity of baptism for salvation, baptismal regeneration (John iii. 5), infant baptism, and the salvation of baptized infants (Matt. xix. 12). The effect of baptism is the remission of hereditary and previous actual sin, and the gift of the Holy Spirit. It can not be repeated; sins committed after baptism must be forgiven by priestly absolution on repentance and confession.

\emph{Article XVII.}—The Eucharist is both a sacrament and a sacrifice, in which the very body and blood of Christ are truly and really (\textgreek{ἀληθῶς καὶ πραγματικῶς}) present under the figure and type (\textgreek{ἐν εἴδει καὶ τύπῳ}) of bread and wine, are offered to God by the hands of the priest as a real though unbloody sacrifice for all the faithful, whether living or dead (\textgreek{ὑπὲρ πάντων τῶν εὐσεβῶν ζώντων καὶ τεθνεώτων}), and are received by the hand and the mouth of unworthy as well as worthy communicants, though with opposite effects. The Lutheran doctrine is rejected, and the Romish doctrine of transubstantiation (\textgreek{μεταβολή, μετουσίωσις}) is taught as strongly as words can make it;\footnote{\par Decr. 17 (Kimmel, P. I. p. 457): \textgreek{ὥστε μετὰ τὸν ἁγιασμὸν τοῦ ἄρτου καὶ τοῦ οἴνου μεταβάλλεσθαι } (to be translated) \textgreek{μετουσιοῦσθαι} (transubstantiated), \textgreek{μεταποιεῖσθαι} (refashioned, transformed), \textgreek{μεταῤῥυθμίζεσθαι} (changed, reformed), \textgreek{τὸν μὲν ἄρτον εἰς αὐτὸ τὸ ἀληθὲς τοῦ κυρίου σῶμα, ὅπερ ἐγεννήθη ἐν Βηθλεὲμ ἐκ τῆς ἀειπαρθένου, ἐβαπτίσθη ἐν Ἰορδάνῃ, ἔπαθεν, ἐτάφη, ἀνέστη, ἀνελήφθη, κάθηται ἐκ δεζιῶν τοῦ Θεοῦ καὶ πατέρος, μέλλει ἐλθεῖν ἐπὶ τῶν νεφελῶν τοῦ οὐρανοῦ—τὸν δ οἶνον μεταποιεῖσθαι καὶ μετουσιοῦσθαι εἰς αὐτὸ τὸ ἀληθὲς τοῦ κυρίου αἶμα, ὅπερ κρεμαμένου ἐπὶ τοῦ σταυροῦ ἐχύθη ὑπὲρ τῆς τοῦ κόσμου ξωῆς.} Mosheim thinks that the Greeks first adopted in this period the doctrine of transubstantiation, but Kiesling (\emph{Hist. concertat. Græcorum Latinorumque de transsubstantiatione}, pp. 354–480, as quoted by Tzschirner, in Vol. IX. of his continuation of Schroeckh's \emph{Church Hist. since the Reformation}, p. l02) has shown that several Greeks taught this theory long before or ever since the Council of Florence (1439). Yet the opposition to the Calvinistic view of Cyril and his sympathizers brought the Greek Church to a clearer and fuller expression on this point.} but it is disclaimed to give an explanation of the \emph{mode }in which this mysterious and miraculous change of the elements takes place.\footnote{Ibid. (p. 461): \textgreek{ἔτι τῇ μετουσίωσις λέξει οὐ τὸν τρόπον πιστεύομεν δηλοῦσθαι, καθ ὃν ὁ ἄρτος καὶ ὁ οἶνος μεταποιοῦνται εἰς τὸ σῶμα καὶ τὸ αἷμα τοῦ κυρίου—τοῦτο γὰρ ἄληπτον πάντη καὶ ἀδύνατον πλὴν αὐτοῦ τοῦ θεοῦ. } In the Lat. Version: '\emph{Præterea verbo}  Transsubstantiationis  \emph{modum ilium, quo in corpus et sanguinem Domini panis et vinum convertantur, explicari minime credimus—id enim penitus incomprehensibile},' etc. \textgreek{Μετουσίωσις} (not given in the Classical Dict., nor in Sophocles's Byzantine Greek Dict., nor in Suicer's \emph{Thesaurus})—from the classical \textgreek{οὐσιόω,} to call into being (\textgreek{οὐσία}) or existence, and the patristic \textgreek{οὐσίωσις,} a calling into existence—must be equivalent to the Latin \emph{transsubstantiatio,} or change of the elemental substance of bread and wine into the body and blood of Christ.}

\emph{Article XVIII.}—The souls of the departed are either at rest or in torment,\footnote{\par  \textgreek{ἐν ἀνέσει,} lit. in relaxation, recreation, \textgreek{ἢ ἐν ὀδύνῃ,} or in pain, distress.} according to their conduct in life; but their condition will not be perfect till the resurrection of the body. The souls of those who die in a state of penitence (\textgreek{μετανοήσαντες}), without having brought forth fruits of repentance, or satisfactions (\textgreek{ἱκανοποίησις}), depart into Hades (\textgreek{ἀπέρχεσθαι εἰς ᾄδου}), and there they must suffer the punishment for their sins; but they may be delivered by the prayers of the priests and the alms of their kindred, especially by the unbloody sacrifice of the mass (\textgreek{μαγάλα δυναμένης μάλιστα τῆς ἀναιμάκτου θυσίας}), which individuals offer for their departed relatives, and which the Catholic and Apostolic Church daily offers for all alike. The liberation from this intervening state of purification will take place before the resurrection and the general judgment, but the time is unknown.

This is essentially the Romish doctrine of purgatory, although the term is avoided, and nothing is said of material or physical torments.\footnote{The same doctrine is taught in the Longer Russian Catechism of Philaret (on the 11th article of the Nicene Creed). It is often asserted (even by Winer, who is generally very accurate, \emph{Symb. }pp. 158, 159) that the Greek Church rejects the Romish purgatory. Winer quotes the Conf. Metrophanis Critopuli, c. 20; but this has no ecclesiastical authority, and, although it rejects the word \textgreek{πῦρ καθαρτήριον} (\emph{ignis purgatoris}), and all idea of material or physical pain (\textgreek{τὴν ἐκείνων ποινὴν μὴ ὑλικὴν εἶναι, εἴτους ὀργανικήν, μὴ διὰ πυρός, μήτε δἰ ἄλλης ὕλης)}, it asserts, nevertheless, a spiritual pain of conscience in the middle state (\textgreek{ἀλλὰ διὰ θλίψεως καὶ ἀνίας τῆς συνειδήσεως}), from which the sufferers may be released by prayers and the sacrifice of the altar. The Conf. Orthodoxa (P. I. Qu. 66) speaks vaguely of a \textgreek{πρόσκαιρος κόλασις καθαρτικὴ τῶν ψυχῶν,} 'a temporary purifying (disciplinary) punishment of the souls.' The Roman Church, on her part, does not require belief in a \emph{material }fire. The Greek Church has no such minute geography of the spirit world as the Latin, which, besides heaven and hell proper, teaches an intervening region of purgatory for imperfect Christians, and two border regions, the \emph{Limbus Patrum} for the saints of the Old Testament now delivered, and the \emph{Limbus Infantum} for unbaptized children; but it differs much more widely from the Protestant eschatology, which rejects the idea of a third or middle place altogether, and assign all the departed either to a state of bliss or a state of misery; allowing, however, different degrees in both states corresponding to the different degrees of holiness and wickedness.}

To these eighteen decrees are added four questions and answers, with polemic reference to the similar questions at the close of the enlarged edition of Cyril's Confession.\footnote{\par Comp. § 15, p. 57.} The first question discourages and even prohibits the general and indiscriminate reading of the Holy Scriptures, especially certain portions of the Old Testament. The second denies the perspicuity of the Scriptures. The third defines the extent of the canon including the Apocrypha.\footnote{The following Apocrypha are expressly mentioned (Vol. I. p. 467): The Wisdom of Solomon, Judith, Tobit, History of the Dragon, History of Susannah, the books of the Maccabees, the Wisdom of Sirach. The Confession of Mogilas, though not formally sanctioning the Apocrypha, quotes them frequently as authority, e.g.  Tobit xii. 9, in P. III. Qu. 9, on alms. On the other hand, the less important Confession of Metrophanes Critopulus, c. 7 (Kimmel, P. II. p. 104 sq.), mentions only twenty-two canonical books of the Old Test., and excludes from them the Apocrypha, mentioning Tobit, Judith, Wisdom of Solomon, Wisdom of Sirach, Baruch, and the Maccabees. The Russian Catechism of Philaret omits the Apocrypha in enumerating the books of the Old Test., for the reason that 'they do not exist in Hebrew,' but adds that 'they have been appointed by the fathers to be read by proselytes who are preparing for admission into the Church.' (See Vol. II. 451, and Blackmore's translation, pp.38, 39.)} The fourth teaches the worship of saints, especially the Mother of God (who is the object of \emph{hyperdulia, }as distinct from the ordinary \emph{ dulia }of saints, and the \emph{latria }or worship proper due to God), as also the worshipful veneration of the cross, the holy Gospels, the holy vessels, the holy places,\footnote{\textgreek{προσκυνοῦμεν καὶ τιμῶμεν τὸ ξύλον τοῦ τιμίου τοῦ ζωοποιοῦ σταυροῦ, κ.τ.λ.}} and of the images of Christ and of the saints.\footnote{\textgreek{τὴν εἰκόνα τοῦ Κυρίου ἡμῶν Ἰησοῦ Χρ. καὶ τῆς ὑπεραγίας θεοτόκου καὶ πάντων τῶν ἁγίων προσκυνοῦμεν καὶ τιμῶμεν καὶ ἀσπαζόμεθα. }}

In all these important points the Synod of Jerusalem again essentially agrees with the Church of Rome, and radically dissents from Protestantism.

\xsection{The Synods of Constantinople, A.D. 1672 and 1691.}

§ 18. The Synods of Constantinople, A.D. 1672 and 1691.

Three months previous to the Synod of Jerusalem a Synod was held at Constantinople (January, 1672), which adopted a doctrinal statement signed by Dionysius, Patriarch of Constantinople, and forty-three dignitaries belonging to his patriarchate.\footnote{\par It is called Dionysii, \emph{Patr. Const., super Calvinistarum erroribus ac reali imprimis præsentia responsio}, and is published in some editions of the Confession of the Synod of Jerusalem; in Harduini \emph{Acta Conciliorum}, Tom. XI. pp. 274–282; and in the second volume of Kimmel's \emph{Monumenta}, pp. 214–227.} It is less complete than the Confession of Dositheus, but agrees with it on all points, as the authority and infallibility of the Church, the extent of the canon, the seven mysteries (sacraments), the real sacrifice of the altar, and the miraculous transformation\footnote{On this the document teaches (Kimmel, P. II. p. 218) that when the priest prays, 'Make (\textgreek{ποίησον}) this bread the precious blood of thy Christ,' then, by the mysterious and ineffable operation of the Holy Ghost, \textgreek{ὁ μὲν ἄρτος μεταποιεῖται } (\emph{transmutatur}) \textgreek{εἰς αὐτό ἐκεῖνο τὸ ἴδιον σῶμα τοῦ σωτῆρος Χριστοῦ πραγματικῶς καὶ ἀληθῶς καὶ κυρίως} (\emph{realiter}, \emph{vere}, \emph{ac proprie}), \textgreek{ὁ δὲ οἶνος εἰς τὸ ζωοποιὸν αἷμα αὐτοῦ. }} of the elements.

Another Synod was held in Constantinople nineteen years afterwards, in 1691, under Patriarch Callinicus, for the purpose of giving renewed sanction to the orthodox doctrine of the Eucharist, in opposition to Logothet John Caryophylus, who had rejected the Romish theory of transubstantiation, and defended the Calvinistic view of Cyril Lucar. The Synod condemned him, and declared that the Eastern Church had always taught a change (\textgreek{μεταβολή}) of the elements in the sense of a transubstantiation (\textgreek{μετουσίωσις}), or an actual transformation of their essence into the body and blood of Christ.\footnote{\par I have not been able to procure the proceedings of this Synod; they are omitted both by Hardouin and Kimmel. They were first printed at Jassy, 1698; then in Greek and Latin by Eusebius Renaudot, together with some other Greek writings on the Eucharist, Paris, 1709; in German by Heineccius, in his \emph{Abbildung der alten und neuen Griechischen Kirche}, 2 Parts, Leipz. 1711. Appendix. p. 40. etc. So says Rud. Hofmann (in his \emph{Symbolik}, Leipz. 1857, p. 135), who has paid careful attention to the Greek Church.}

\xsection{The Doctrinal Standards of the Russo-Greek Church.}

§ 19. The Doctrinal Standards of the Russo-Greek Church.

Literature.

I. Russian Doctrine and Theology:

The Catechisms of Platon and Philaret (see below).

R. W. Blackmore: \emph{The Doctrine of the Russian Church}, etc., Aberdeen, 1845.

W. Guettée (Russian Priest and Doctor of Divinity): \emph{Exposition de la doctrine de l’église catholique orthodoxe de Russie, }Paris, 1866.

Theophanes Procopowicz:  \emph{Theologia Christiana orthodoxa, }Königsberg, 1773–1775, 5 vols. (abridged, Moscow, 1802).

Hyac. Kirpinski: \emph{Compendium orthodoxæ theologiæ, }Lips. 1786.

II. Worship and Ritual:

\emph{The divine Liturgy of }St. John Chrysostom (the Liturgy used in the Orthodox Eastern Church), Greek ed. by Daniel, \emph{Cod. Liturg.} Tom. IV. P. II. p. 327, etc.; by J. M. Neale, in \emph{Primitive Liturgies}, 2d edition, London, 1868; English translations by King, Neale, Brett, Covel, J. Freeman Young (the last publ. New York, 1865, as No. VI. of the 'Papers of the Russo-Greek Committee'). Comp. also the entire fourth volume of Daniel's  \emph{Codex Liturg.} (which gives the Oriental Liturgies), and Neale's  \emph{Primitive Liturgies}, and his \emph{Introd. to the History of the Holy Eastern Church} (Lond. 1850).

John Glen King (Anglican Chaplain at St. Petersburg): \emph{The Rites and Ceremonies of the Greek Church in Russia}, Lond. 1772. Very instructive.

III. History and Present Condition of the Russian Church:

Alex. de Stourdza: \emph{Considérations sur la doctrine et l’esprit de l’église orthodoxe, }Weimar, 1816.

Strahl:  \emph{Contributions to Russian Church History, }Halle, 1827: and \emph{History of the Russian Church, }Halle, 1830.

Mouravieff:  \emph{History of the Church of Russia, }St. Petersburg, 1840; translated by Blackmore, Oxford, 1842. Comes down to 1721.

Pinkerton:  \emph{Russia}, London, 1833.

Haxthausen: \emph{Researches on Russia}, German and French, 1847–52, 3 vols.

Theiner:  \emph{Die Staats-Kirche Russlands}, 1853.

H. J. Schmitt: \emph{Kritische Geschichte der neugriechischen und der russischen Kirche}, Mainz, 2d ed. 1854.

Prince Aug. Galitzin:  \emph{L’église Græco-Russe}, Paris, 1861.

Dean Stanley:  \emph{Lectures on the History of the Eastern Church, }Lond. and N. Y. 1862, Lect. IX.–XII.

Boissard:  \emph{L’église de Russie, }Paris, 1867, 2 vols.

Philaret (Archbishop of Tschernigow): \emph{Geschichte der Kirche Russlands, }transl. by Blumenthal, 1872.

Basaroff:  \emph{Russische orthodoxe Kirche. Ein Umriss ihrer Entstehung u. ihres Lebens}, Stuttgart, 1873.

Also the \emph{Occasional Papers} of the 'Eastern Church Associations' of the Church of England and the Protestant Episcopal Church in the United States, publ. in Lond. (Rivington's), and N. York, since 1864.

The latest doctrinal standards of Greek Christianity are the authorized Catechisms and Church-books of the orthodox Church of Russia, by far the most important and hopeful branch of the Eastern Communion.

Russia received Christianity from the Byzantine Empire. Cyril and Methodius, two monks of Constantinople, preached the gospel to the Bulgarians on the Danube after the middle of the ninth century, translated the Scriptures\footnote{\par The Psalms and the New Testament, with the exception of the Apocalypse.} into the Slavonic language (creating the Slavonic alphabet in quaint Greek characters), and thus laid the foundation of Slavonic literature and civilization. This event was contemporary with the founding of the Russian Empire by Ruric, of the Norman race (A.D. 862), and succeeded by half a century the founding of the German Empire under Charlemagne, in close connection with Rome (A.D. 800). As the latter was a substitute for the Western Roman Empire, so the former was destined to take the place of the Eastern Roman Empire, and looks forward to the reconquest of Constantinople, as its natural capital. The barbarous Russians submitted, in the tenth century, without resistance, to Christian baptism by immersion, at the command of their Grand Duke, Vladimir, who himself was brought over to Christianity by a picture on the last judgment, and his marriage to a sister of the Greek Emperor Basil. In this wholesale conversion every thing is characteristic: the influence of the picture, the effect of marriage, the power of the civil ruler, the military command, the passive submission of the people.

Since that time the Greek Church has been the national religion of the Slavonic Russians, and identified with all their fortunes and misfortunes. For a long time they were subject to the jurisdiction of the Patriarch of Constantinople. But after the fall of this city (1453) the Metropolitan of Moscow became independent (1461), and a century later (January, 1589) he was raised by Patriarch Jeremiah II. of Constantinople, then on a collecting tour in Russia, to the dignity of a Patriarch of equal rank with the other four (of Constantinople, Alexandria, Antioch, Jerusalem). Moscow was henceforward the holy city, the Rome of Russia.

In the beginning of the eighteenth century, Peter the Great, a second Constantine, founded St. Petersburg (1703), made this city the political and ecclesiastical capital of his Empire, and created, in the place of the Patriarchate of Moscow, the 'Most Holy Governing Synod,' with the Czar as the head (1721). This organic change was sanctioned by the Eastern Patriarchs (1723), who look upon the emperor-pope of Russia as their future deliverer from the intolerable yoke of the Turks.

[Note.—Since the revolution of 1917 and the assassination of the Czar, the position of the Russian Church has undergone a radical change. The Soviet government has passed from a law abolishing the union of Church and State to a relentless war against all religion and religious exercises, the confiscation of Church property, the suppression of religious liberty, the imprisonment and execution of clerical personages, and even to a policy of active atheistic propaganda. Conforming to the new civil order, the Holy Sober—council—met, August, 1917, with 564 delegates present, of whom 278 were laymen, and constituted Tikhon (1866–1925) Most holy Patriarch of Moscow and all Russia, thus re-establishing the patriarchate after an interval of two centuries. Tikhon resisted the Soviet acts instituting civil marriage and disestablishing the Church, and placed the state officials under excommunication. The government replied by further legislation hostile to the Church, and Tikhon was put under arrest and resigned the patriarchate, 1922. In the mean time a 'reforming' organization, calling itself the 'Living Church,' was effected, which acknowledged the Soviet revolution and made the 'white clergy'—in contrast to the monks—eligible to the episcopal office. The Sober of April, 1924, received greetings from Dr. Blake of the Methodist Episcopal Church, disavowed Tikhon's anti-Soviet deliverances, endorsed the separation of Church and State, and granted to widowed and divorced priests the right of remarriage. A third Sober affirmed that supreme ecclesiastical authority resided in itself and not in the patriarch, a declaration accepted by Gregory VII, œcumenical patriarch of Constantinople, other Eastern patriarchs dissenting. The Tikhon wing was continued under Peter, Metropolitan of Krutitsky, whom Tikhon had designated as his successor. Peter was banished for anti-Soviet policies, and his place filled by Abp. Sergius, who himself was imprisoned but released, 1927, after promising to support the existing civil government. The \emph{émigré} bishops, with Serbia as a rallying-place, have favored the restoration of the empire, and June 30, 1930, Sergius deposed Eulogius from the post of so-called supreme bishop of the Russian Church outside of Russia. Soviet legislation, 1930, confirmed all previous acts calculated to blot out religious convictions and ritual. It forbids the teaching of religion to persons under eighteen, the organization of meetings of women and children for purposes of prayer and biblical and literary study or for sewing, the organization under Church influence of libraries and reading-rooms, and even measures intended to give sanitary and medical assistance. It prohibits the teaching of any form of religious belief in educational establishments, and the formation of all boys' and girls' clubs in church buildings. Religious teaching is treated as ``anti-revolutionary activity.'' The secret propaganda of religion among the masses is forbidden, and ministers of religion, including rabbis and nuns, who continue to follow their religion are disfranchised and made ineligible for public office. Bibles and prayer-books are confiscated. Church buildings are put at the State's disposal. Articles of gold and silver and precious stones are to be given up upon the discontinuance of a house of worship, and places of worship having a historic or artistic value pass to the State. Processions on festival days are forbidden, as also is the observance of Christmas, Easter, and other Church feasts. In addition to such laws, the Soviet has carried out its destructive policy by films and posters ridiculing and blaspheming Christianity. By governmental order or the populace, multitudes of icons have been destroyed and pretended bodies of saints dishonored and shown to be made of wax or straw. The treatment of the Russian Church and clergy has called forth from the pope and the Church of England resolutions against the government's policy, and letters of sympathy. Since 1914, friendly gestures have been made from Rome calculated to win favor for the Roman Church. In 1920, Ephraem of Edessa was enrolled among the doctors of the Church. The Oriental College in Rome has been enlarged. In 1921, Benedict XV. addressed the Russians as 'our distant children who, though separated from us by the barriers of centuries, are all the nearer our paternal heart, the greater their misfortunes are.' In 1929, Pius XI. issued an appeal in Italian for prayer for 'our brethren in Russia,' which spoke of 'the sacrilege heaped upon the priests and believers, and the violence done to the conscience by the Soviets.' The pontiff appointed a solemn mass to be celebrated over St. Peter's tomb, March 19, 1930, and called for the help of 'the Immaculate Virgin, Mother of God, her most chaste spouse, St. Joseph, patron of the Church universal, John Chrysostom and other patron saints of the Russians, and of all saints, especially St. Therèse of the Cradle of Jesus, the sweet thaumaturge of Lisieux.' In the form of prayer which Pius added for general use, the petition was made that the Russians may return 'to the one fold and the communion of the Catholic Church,' and an indulgence of 300 days offered to all making the prayer piously. Resolutions passed by the Convocations Canterbury and York, 1930, called for special prayers in the churches at the morning and evening services, March 16.—Ed.]

We have already seen that the 'Orthodox Confession,' or the first systematic and complete exhibition of the modern Greek faith, is the product of a Russian prelate, Peter Mogilas of Kieff. It was followed, and practically superseded, by other catechisms, which are much better adapted to the religious instruction of the young.

1. The Catechism of  Platon, Metropolitan of Moscow (died 1812), one of the very few Russian divines whose name is known beyond their native land.\footnote{\par '\emph{Orthodox Doctrine, or Summary of Christian Divinity};' first published 1762 in Russian, and translated into eight languages: in English, ed. by R. Pinkerton, Edinb. 1814; German ed., Riga, 1770; Latin ed., Moscow, 1774. Blackmore (l.c. p. vii.) speaks of three Catechisms of Platon, which probably differ only in size.} He was the favorite of the Empress Catherine II. (died 1796), and, for a time, of her savage son, the Emperor Paul (assassinated 1801), and at the end of his life he encouraged the Emperor Alexander I. in the terrible year of the French invasion and the destruction of Moscow. When the French atheist Diderot began a conversation with the sneering remark, 'There is no God,' Platon instantly replied, 'The fool says in his heart, There is no God.' He was a great preacher and the leader of a somewhat milder type of Russian orthodoxy, not disinclined to commune with the outside world. His Catechism was originally prepared for his pupil, the Grand Duke Paul Petrovitsch, and shows some influence of the evangelical system by its tendency to go directly to the Bible.

2. The Catechism of  Philaret, revised, authorized, and published by the Holy Synod of St. Petersburg. It is translated into several languages, and since 1839 generally used in the schools and churches of Russia. It was sent to all the Eastern Patriarchs, and unanimously approved by them.\footnote{\par Philaret wrote two Catechisms—a shorter one, called '\emph{Elements of Christian Learning; or, a Short Sacred History and a Short Catechism},' St. Petersburg, at the Synodical Press, 1840 (only about twelve pages), and a longer one under the title, '\emph{A Full Catechism of the Orthodox Catholic Church of the East, examined and approved by the Most Holy Governing Synod, and published for the Use of Schools and of all Orthodox Christians, by order of His Imperial Majesty},' Moscow, at the Synodical Press, 1839 (English translation of Blackmore, Aberdeen, 1845). Most of the German works on Symbolics ignore Philaret altogether. Even Hofmann (p. 136) and Gass (p. 440) barely mention him. We give his Larger Catechism in the second volume.}

Philaret (born 1782, died 1867) was for forty-seven years (1820–67) Metropolitan of Moscow. He was intrusted with the important State secret of the will of Alexander I., and crowned his two successors (Nicholas I. and Alexander II.). He represents, in learning, eloquence, and ascetic piety, the best phase of the Russian State Church in the nineteenth century.\footnote{\par Dean Stanley, who saw him in Moscow in 1857, praises his striking and impressive manner as a preacher, his gentleness, his dignified courtesy and affability, and associates him with a reactionary revival of mediæval sanctity, which had its parallel in the Puseyism of the Church of England. The Scottish Bishop of Moray and Ross, who called on him in behalf of the Eastern Church Association in 1866, describes him as the most venerated and beloved man in the Russian Empire, and as 'gentle, humble, and pious.' Comp. Souchkow, \emph{Memoirs of Philaret}, Moscow, 1868; \emph{Select Sermons of Philaret. transl. from the Russian}, London (Jos. Masters), 1873.}

His longer Catechism (called a \emph{full} catechism) is, upon the whole, the ablest and clearest summary of Eastern orthodoxy, and shows a disposition to support every doctrine by direct Scripture testimony. It follows the plan and division of the Orthodox Confession of Mogilas, and conforms to its general type of teaching, but it is more clear, simple, evangelical, and much better adapted for practical use. In a number of introductory questions it discusses the meaning of a catechism, the nature and necessity of right faith and good works, divine revelation, the holy tradition and Holy Scripture (as the two channels of the divine revelation and the joint rule of faith and discipline), the Canon of the Scriptures (exclusive of the Apocrypha, because 'not written in Hebrew'), with some account of the several books of the Old and New Testaments, and the composition of the Catechism. This is divided into three parts, like the Confession of Mogilas, according to the three cardinal virtues (1 Cor. xiii. 13).

First Part: On Faith. An Exposition of the Nicene Creed, arranged in twelve articles. In the doctrine of the Church the Protestant distinction of the visible and invisible Church is, in a modified sense, adopted; Christ is declared to be the only and ever-abiding Head of the Church, and it is stated that the division of the Church into many particular and independent organizations, as those of Jerusalem, Antioch, Alexandria, Constantinople, Russia (Rome, Wittenberg, Geneva, and Canterbury are ignored); does not hinder them from being spiritually members 'of the one body of the Universal Church, from having one Head, Christ, and one spirit of faith and of grace.'

Second Part: On Hope. An Exposition of the Lord's Prayer (in seven petitions), and of the nine Beatitudes of the Sermon on the Mount.

Third Part: On Love or Charity. An Exposition of the Decalogue as teaching, in two tables, love to God and love to our neighbor. The last question is: 'What caution do we need when we seem to ourselves, to have fulfilled any commandment? A. We must then dispose our hearts according to the words of Jesus Christ: ``When ye have done all those things which are commanded you, say, We are unprofitable servants; we have done that which was our duty to do'' (Luke xvii. 10).'

3. Finally, we may mention, as secondary standards of Russian orthodoxy and discipline, the \emph{Primer} or \emph{Spelling-Book}, and a Treatise on \emph{The Duty of Parish Priests.}\footnote{\par Both translated by Blackmore, l.c.}

The \emph{Primer} contains the rudiments of religious learning for children and the common people, viz., daily prayers (including the Lord's Prayer, and the 'Hail Mary, Virgin Mother of God,' yet without the 'Pray for us' of the Latin formula), the Nicene Creed, the Ten Commandments (the second and fourth abridged), with brief explanations and short moral precepts.

The Treatise on \emph{The Duty of Parish Priests} was composed by George Konissky, Archbishop of Mogileff (died 1795), aided by Parthenius Sopkofsky, Bishop of Smolensk, and first printed at St. Petersburg in 1776. All candidates for holy orders in the Russian Seminaries are examined on the contents of this book. It is mainly disciplinary and pastoral, a manual for the priests, directing them in their duties as teachers, and as administrators of the mysteries or sacraments. But doctrine is incidentally touched, and it is worthy of remark that this Treatise approaches more nearly to the evangelical principle of the supremacy of the Bible in matters of Christian faith and Christian life than any deliverance of the Eastern Church.\footnote{See Part I. No. VIII.–XIII. pp. 160–164 in Blackmore's version: 'All the articles of the faith are contained in the Word of God, that is, in the books of the Old and New Testaments. . . . The Word of God is the source, foundation, and perfect rule, both of our faith and of the good works of the law. . . . The writings of the holy Fathers are of great use . . . but neither the writings of the holy Fathers nor the traditions of the Church are to be confounded or equaled with the Word of God and his Commandments.'}

\xsection{Anglo-Catholic Correspondence with the Russo-Greek Church.}

§ 20. Anglo-Catholic Correspondence with the Russo-Greek Church.

The Reformation of the sixteenth century proceeded entirely from the bosom of Latin or Western Catholicism. The Greek or Eastern Church had no part in the great controversy, and took no notice of it, until it was brought to its attention from without. The antagonism of the Greek Communion to Western innovations, especially to the claims of the Papacy, seemed to open the prospect of possible intercommunion and co-operation. But, so far, all the approaches to this effect on the part of Protestants have failed

1. The first attempt was made by Lutheran divines in the sixteenth century, and ended in the condemnation of the Augsburg Confession.\footnote{\par See above, § 13.}

2. Of a different kind was Cyril's movement, in the seventeenth century, to protestantize the Eastern Church from within, which resulted in a stronger condemnation of Calvinism and Lutheranism.\footnote{\par See §§ 15–18.}

3. The correspondence of the Anglican Non-Jurors with Russia and the East, 1717–1723, had no effect whatever.

Two high-church English Bishops; called 'Non-Jurors' (because they refused to renounce their oath of allegiance to King James II., and to transfer it to the Prince of Orange), in connection with two Scottish Bishops, assumed, October, 1717, the responsibility of corresponding with the Russian Czar, Peter the Great, and the Eastern Patriarchs.\footnote{\par The letters of the four Bishops signing themselves 'Jeremias, \emph{Primus Angliæ Episcopus}; Archibaldus, \emph{Scoto-Britanniæ Episcopus}; Jacobus, \emph{Scoto-Britanniæ Episcopus}; Thomas, \emph{Angliæ Episcopus},' are given by Lathbury, in his \emph{History of the Non-Jurors}, pp. 309–361, as documentary proof of their doctrinal status, but the answers are omitted.} They were prompted to this step by a visit of an Egyptian Bishop to England, who collected money for the impoverished patriarchal see of Alexandria, and probably still more by a desire to get aid and comfort from abroad in their schismatical isolation. They characteristically styled themselves 'The Catholic Remainder in Britain.'

After a delay of several years, the Patriarchs, under date, Constantinople, September, 1723, sent their ultimatum, requiring, as a term of communion, absolute submission of the British to all the dogmas of the Greek Church. 'Those,' they wrote, 'who are disposed to agree with us in the Divine doctrines of the Orthodox faith must necessarily follow and submit to what has been defined and determined by ancient Fathers and the Holy Œcumenical Synods from the time of the Apostles and their Holy Successors, the Fathers of our Church, to this time. We say they must submit to them with sincerity and obedience, and without any scruple or dispute. And this is a sufficient answer to what you have written.' With this answer they forwarded the decrees of the Synod of Jerusalem of 1672.

The Russians were more polite. The 'Most Holy Governing Synod' of St. Petersburg, in transmitting the ultimatum of the Eastern Patriarchs, proposed, in the name of the Czar, 'to the Most Reverend the Bishops of the Remnant of the Catholic Church in Great Britain, our Brethren most beloved in the Lord,' that they should send two delegates to Russia to hold a friendly conference, in the name and spirit of Christ, with two members to be chosen by the Russians, that it may be more easily ascertained what may be yielded and given up by one to the other; what, on the other hand, may and ought for conscience' sake to be absolutely denied.\footnote{\par The two letters of the Holy Synod, the one signed Moscow, February, 1723, the other without date, are given by Blackmore, \emph{Doctrine of the Russian Church}, Pref. pp. xxvi.–xxviii. The anonymous author (probably Dr. Young, now Bishop in Florida) of No. II. of the Papers of 'the Eastern Church Association' supplies the signatures of nine Church dignitaries of Russia from personal inspection of the archives of the Holy Synod, at a visit to St. Petersburg, April, 1864.}

But such a conference was never held. The death of Peter (1725) put an end to negotiations. Archbishop Wake, of Canterbury, wrote a letter to the Patriarch of Jerusalem, in which he exposed the Non-Jurors as disloyal schismatics and pretenders. The Eastern Patriarchs accused the Anglicans of being 'Lutherano-Calvinists,' and the Russian Church historian, Mouravieff, in speaking of the correspondence, represents them as being infected with the same 'German heresy,' which had been previously condemned by the Orthodox Church.\footnote{\par History of the Church of Russia, translated by Blackmore, pp. 286 sq., 407 sqq.}

4. A far more serious and respectable attempt to effect intercommunion between the Anglican and Russo-Greek Churches was begun in 1862, with the high authority of the Convocation of Canterbury, and the General Convention of the Protestant Episcopal Church in the United States. The ostensible occasion was furnished by the multiplication of Russo-Greeks on the Pacific coast, and by the desirableness of securing decent burial for Anglican travelers in the East, but the real cause lies much deeper. It is closely connected with the powerful Anglo-Catholic movement, which arose in Oxford in 1833, and has ever since been aiming to de-protestantize the Anglican Church. Hundreds of her priests and laymen, headed by Dr. John H. Newman, seceded to Rome; while others, less logical or more loyal to the Church of their fathers, are afraid of the charms or corruptions of the Papacy, and look hopefully to intercommunion with the Holy Catholic Orthodox and Apostolic Mother Church of the East to satisfy their longing for Catholic unity, and to strengthen their opposition to Protestantism and Romanism. The writings of the late Dr. John Mason Neale, and Dr. Pusey's \emph{Eirenicon}, contributed not a little towards creating an interest in this direction.

In the General Convention of the Protestant Episcopal Church in the United States, held in New York, October, 1862, a joint committee was appointed 'to consider the expediency of opening communication with the Russo-Greek Church, to collect authentic information upon the subject, and to report to the next General Convention.' Soon afterwards, July 1, 1863, the Convocation of Canterbury appointed a similar committee, looking to 'such ecclesiastical intercommunion with the Orthodox East as should enable the laity and clergy of either Church to join in the sacraments and offices of the other without forfeiting the communion of their own Church.' The Episcopal Church in Scotland likewise fell in with the movement. These committees corresponded with each other, and reported from time to time to their authorities. Two Eastern Church Associations were formed, one in England and one in America, for the publication of interesting information on the doctrines and worship of the Russo-Greek Church. Visits were made to Russia, fraternal letters and Christian courtesies were exchanged, and informal conferences between Anglican and Russian dignitaries were held in London, St. Petersburg, and Moscow.\footnote{\par See the details in the Occasional Papers of the two Eastern Church Associations, published since 1864 in London (Rivington's) and in New York, and the Reports in the \emph{Journal of the General Convention of the Protestant Episcopal Church in the United States}, held in New York, 1868, Append. IV. p. 427, and Append. XI. p. 480, and of the Convention in Baltimore, 1871, Append. VI. pp. 565–85. These reports are signed by Bishops Whittingham, Whitehouse, Odenheimer, Coxe, Young, and others. A curious incident in this correspondence, not mentioned in these documents, was the celebration of Greek mass, by a Russian ex-priest of doubtful antecedents, in the Episcopal Trinity Chapel of New York, on the anniversary of the Czar Alexander II., March 2, 1865.}

The Russo-Greeks could not but receive with kindness and courtesy such flattering approaches from two of the most respectable Churches of Christendom, but they showed no disposition whatever either to forget or to learn or to grant any thing beyond the poor privilege of burial to Anglicans in consecrated ground of the Orthodox (without, however, giving them any right of private property). Some were willing to admit that the Anglican Church, by retaining Episcopacy and respect for Catholic antiquity, 'attached her back by a strong cable to the ship of the Catholic Church; while the other Protestants, having cut this cable, drifted out at sea.' Yet they could not discover any essential doctrinal difference. They found strange novelties in the Thirty-nine Articles; they took especial offense at Art. 19, which asserts that the Churches of Jerusalem, Alexandria, and Antioch have erred; they expressed serious scruples about the validity of Anglican orders, on account of a flaw in Archbishop Barker's ordination, and on account of the second marriage of many Anglican priests and bishops (which they consider a breach of continency, and a flagrant violation of Paul's express prohibition, according to their interpretation of \textgreek{μιᾶς γυναικὸς ἄνδρα, }  1 Tim. iii. 2); they can not even recognize Anglican baptism, because it is not administered by trine immersion.

On the other hand, the Russo-Greeks insist on the expulsion of the \emph{Filioque}, which is their main objection to Rome; the recognition of the seventh œcumenical Council; the invocation of the Holy Virgin and the Saints; the veneration of icons; prayers for the departed; seven sacramental mysteries; trine immersion; a mysterious transformation (\textgreek{μετουσίωσις)} of the eucharistic elements; the eucharistic sacrifice for the living and the dead.\footnote{\par See the documents in the \emph{Journal of the General Convention} for 1871, pp. 567–577, viz., the answers of Gregory, Patriarch of Constantinople, dated Sept. 26, 1869, to a letter of the Archbishop of Canterbury, accompanied by a Greek copy of the English Liturgy; the report of the Greek Archbishop of Syra to the Holy Synod of Greece, concerning his visit to England (1870); also the report of an interesting conference between the Greek Archbishop of Syra and the Anglican bishop of Ely (Dr. Browne, the author of a Commentary on the Thirty-nine Articles), held February 4, 1870, where all the chief points of difference were discussed in a friendly Christian spirit, but without result.}

5. The latest phase of the Anglo-Greek movement is connected with the Old Catholic reunion Conferences in Bonn, 1874 and 1875.\footnote{\par See the results of the Bonn Conferences, at the close of Vol. II. pp. 545–554.} Here the \emph{Filioque} was surrendered as a peace-offering to the Orientals; but the Orientals made no concession on their part. It is not likely that the Anglican Church will sacrifice her own peace, the memory of her reformers and martyrs, and a Protestant history and literature of three centuries to an uncongenial union with the Russo-Greek Church in her present unreformed state.

\xsection{The Eastern Sects: Nestorians, Jacobites, Copts, Armenians.}

§ 21. The Eastern Sects: Nestorians, Jacobites, Copts, Armenians.

Literature.

I. The Nestorians:

Ebedjesu (Nestorian, d. 1318): \emph{Liber Margarita de veritate fidei}, in Angelo Mai's \emph{Script. veter. Nova Collectio}, Vol. XII. p. 317.

Jos. Sim. Assemani (R. C., d. 1678): \emph{De Syris Nestorianis, }in his \emph{Bibl. Or., }Rom. 1719–28, Tom. III. Pt. II.

Gibbon  : \emph{Decline and Fall of the Roman Empire, }chap, xlvii. near the end.

E. Smith and H. G. C. Dwight: \emph{Researches in Armenia, with a Visit to the Nestorian and Chaldean Christians of Oormiah, }etc., 2 vols. Boston, 1833.

Justin Perkins:  \emph{A Residence of Eight Years in Persia, }Andover, 1843.

W. Etheridge  : \emph{The Syrian Churches, their Early History, Liturgies, and Literature, } Lond. 1846.

Geo. Percy Badger:  \emph{The Nestorians and their Rituals, }Illustrated (with colored plates), 2 vols. Lond. 1852.

H. Newcomb:  \emph{A Cyclopædia of Missions, }New York, 1856, p. 553 sq.

Petermann: Article \emph{Nestorianer}, Herzog's \emph{Theol. Encyklop.} Vol. X. (1858), pp. 279–288.

Rufus Anderson (late For. Sec. Am. Board of C. For. Missions: \emph{Republication of the Gospel in Bible Lands; History of the Missions of the Amer. Board of Comm. for For. Miss. to the Oriental Churches, }Boston, 1872, 2 vols.

On the Nestorian controversy which gave rise to the Nestorian sect, see my \emph{Church History, }Vol. III. p. 715 sq., and the works quoted there; also p. 729.

II. The Monophysites (Jacobites, Copts, Abyssinians, Armenians, Maronites):

Euseb. Renaudot (R. C., d. 1720): \emph{Historia Patriarcharum Alexandrinorum Jacobitarum a D. Marco usque ad finem sæc. xiii., }Par. 1713. Also by the same: \emph{Liturgiarum Orientalium Collectio, }Par. 1716, 2 vols. 4to.

Jos. Sim. Assemani (R. C.): \emph{Bibliotheca orientalis}, Rom. 1719 sqq., Tom. II., which treats \emph{De scriptoribus Syris Monophysitis.}

Michael le Quien (R. C., d. 1733): \emph{Oriens Christianus}, Par. 1740, 3 vols. folio (Vols. II. and III.).

Veyssière de la Croze:  \emph{Histoire in Christianisme d'Ethiope et d'Armenie, }La Haye, 1739.

Gibbon: \emph{Decline and Fall of the Roman Empire, }chap, xlvii.

Makrîzi (Mohammedan, an historian and jurist at Cairo, died 1441): \emph{Historia Coptorum Christianorum }(Arabic and Latin), ed. \emph{H. J. Wetzer, }Sulzbach, 1828; a better edition by \emph{F. Wüstenfeld, }with translation and annotations, Göttingen, 1845.

J. E. T. Wiltsch:  \emph{Kirchliche Statistik}, Berlin, 1846, Bd. I. p. 225 sq.

John Mason Neale (Anglican): \emph{The Patriarchate of Alexandria, }London, 1847, 2 vols. Also, \emph{A History of the Holy Eastern Church}, London, 1850, 2 vols. (Vol. II. contains among other things the Armenian and Copto-Jacobite Liturgies.)

E. Dulaurier : \emph{Histoire, dogmes, traditions, et liturgie de l’église Armeniane, }Par. 1859.

Arthur Penrhyn Stanley:  \emph{Lectures on the History of the Eastern Church, }New York, 1862, p. 92.

E. F. K. Fortescue: \emph{The Armenian Church. With Appendix by S. C. Malan}, London, 1872.

Rufus Anderson: \emph{Republication of the Gospel in Bible Lands, }quoted above.

Schaff:  \emph{Church History}, Vol. III. pp. 334 sqq. and 770 sqq.

Compare accounts in numerous works of Eastern travel, and in missionary periodicals, especially the \emph{Missionary Herald}, and the Annual Reports of the American Board of Foreign Missions.

Besides the Orthodox Greek Church there are scattered in the East, mostly under Mohammedan and Russian rule, ancient Christian sects, the Nestorians and Monophysites. They represent petrified chapters of Church history, but at the same time fruitful fields for Roman Catholic and Protestant Missions. They owe their origin to the Christological controversies of the fifth century, and perpetuate, the one the Nestorian, the other the Eutychian heresy, though no more as living issues, but as dead traditions. They show the tenacity of Christological error. The Nestorians protest against the third œcumenical Council (431), the Monophysites against the fourth (451). In these points of dispute the Latin and the orthodox Protestant Churches agree with the Orthodox Greek Church against the schismatics.

In other respects the Nestorians and Monophysites betray their Oriental character and original affinity with the Greek Church. They regard Scripture and tradition as co-ordinate sources of revelation and rules of faith. They accept the Nicene Creed without the \emph{Filioque}; they have an episcopal and patriarchal hierarchy, and a ritualistic form of worship, only less developed than the orthodox. They use in their service their ancient native languages, although these have become obsolete and unintelligible to them, since they mostly speak now the Arabic. They honor pictures and relics of saints, but not to the same extent as the Greeks and Russians. The Bible is not forbidden, but practically almost unknown among the people. Their creeds are mostly contained in their liturgies.

They supported the Arabs and Turks in the overthrow of the Byzantine Empire, and in turn were variously favored by them, and upheld in their separation from the Orthodox Greek Church. They are sunk in ignorance and superstition, but, owing to their prejudice against the Greek Church, they are more accessible to Western influence.

Providence has preserved these Eastern sects, like the Jews, unchanged to this day, doubtless for wise purposes. They may prove entering wedges for the coming regeneration of the East and the conversion of the Mohammedans.

I. The Nestorians, in Turkey and Persia, are called after Nestorius, Patriarch of Constantinople. He was condemned by the Council of Ephesus, 431, for so teaching the doctrine of two natures in Christ as virtually to deny the unity of person, and for refusing to call Mary 'the Mother of God' (\textgreek{θεοτόκος, } \emph{Deipara}), and he died in exile about 440. His followers call themselves Chaldæan or Syrian Christians. They flourished for several centuries, and spread far into Arabia, India, and even to China and Tartary. Mohammed is supposed to have derived his imperfect knowledge of Christianity from a Nestorian monk, Sergius. But by persecution, famine, war, and pestilence, they have been greatly reduced. The Thomas Christians of East India are a branch of them, and so called from the Apostle Thomas, who is supposed to have preached on the coast of Malabar.

The Nestorians hold fast to the dyophysite Christology of their master, and protest against the Council of Ephesus, for teaching virtually the Eutychian heresy, and unjustly condemning Nestorius. They can not conceive of a human nature without a human personality, and infer two independent hypostases from the existence of two natures in Christ. They object to the orthodox view, that it confounds the divine and human, or that it teaches a contradiction, viz., two natures and one person. The only alternative to them seems either two natures and two persons, or one person and one nature. From their Christology it follows that Mary was only the mother of the man Jesus. They therefore repudiate the worship of Mary as the Mother of God; also the use of images (though they retain the sign of the cross), the doctrine of purgatory (though they have prayers for the dead), and transubstantiation (though they hold the real presence of Christ in the eucharist); and they differ from the Greek Church by greater simplicity of worship. They are subject to a peculiar hierarchical organization, with eight orders, from the catholicus or patriarch to the sub-deacon and reader. The five lower orders, including the priests, may marry; in former times even the bishops, archbishops, and patriarchs had this privilege. Their fasts are numerous and strict. Their feast-days begin with sunset, as among the Jews. The patriarch and the bishops eat no flesh. The patriarch is chosen always from the same family; he is ordained by three metropolitans. The ecclesiastical books of the Nestorians are written in the Syriac language.

II. The Monophysites, taken together, outnumber the Nestorians, and are scattered over the mountains, villages, and deserts of Armenia, Syria, Egypt, and Abyssinia. They are divided into four distinct sects: the Jacobites in Syria; the Copts in Egypt, with their ecclesiastical descendants in Abyssinia;\footnote{\par The Abyssinian Church receives its Patriarch (Abuna. i.e. Our Father) from the Copts, but retains some peculiar customs, and presents a strange mixture of Christianity with superstition and barbarism. See my \emph{Church History}, Vol. III. p. 778.} the Armenians, and the ancient  Maronites on Mount Lebanon (who were Monothelites, but have been mostly merged into the Roman Church).

The Armenians (numbering about three millions and a half) excel all the rest in numbers, intelligence, and enterprise, and are most accessible to Protestant missionaries.

The Monophysites have their name from their distinctive doctrine, that Christ had but one nature (\textgreek{μονὴ φύσις}), which was condemned by the fourth œcumenical Council of Chalcedon. They are the antipodes of the Nestorians, whom they call Dyophysites. They agree with the Council of Ephesus (431) which condemned Nestorius, but reject the Council of Chalcedon (451). They differ, however, somewhat from the Eutychean heresy of an \emph{absorption }of the human nature by the divine, as held by Eutyches (a monk of Constantinople, died after 451), and teach that Christ had one \emph{composite} nature (\textgreek{μία φύσις σύνθετος} or \textgreek{μία φύσις διττή}). They make the humanity of Christ a mere accident of the immutable divine substance. Their main argument against the orthodox or Chalcedonian Christology is that the doctrine of two natures necessarily leads to that of two persons, and thereby severs the one Christ into two sons of God. They regarded the nature as something common to all individuals of a species (\textgreek{κοινόν}), yet as never existing simply as such, but only in individuals. Their liturgical shibboleth was, \emph{God has been crucified}, which they introduced into the trisagion, and hence they were also called \emph{Theopaschites.}

With the exception of the Chalcedonian Christology, the Monophysite sects hold most of the doctrines, institutions, and rites of the Orthodox Greek Church, but in simpler and less pronounced form. They reject, or at least do not recognize, the \emph{Filioque}; they hold to the mass, or the eucharistic sacrifice, with a kind of transubstantiation; leavened bread in the Lord's Supper; baptismal regeneration by trine immersion; seven sacraments (yet not explicitly, since they either have no definite term for sacrament, or no settled conception of it); the patriarchal polity; monasticism; pilgrimages and fasting; the requisition of a single marriage for priests and deacons (bishops are not allowed to marry); the prohibition of the eating of blood or of things strangled. On the other hand, they know nothing of purgatory and indulgences, and have a simpler worship than the Greeks and Romans. According to their doctrine, all men after death go into Hades, a place alike without sorrow or joy; after the general judgment they enter into heaven, or are cast into hell; and meanwhile the intercessions and pious works of the living have an influence on the final destiny of the departed.

Note on Russian Schismatics.—The dissenting sects of the Russo-Greek Church are very numerous, but not organized into separate communions like the older Oriental schismatics; the Russian government forbidding them freedom of public worship. They are private individuals or lay-communities, without churches and priests. They have no definite creeds, and differ from the national religion mostly on minor ceremonies. The most important among them are the Raskolniki (i.e. Separatists, Apostates), or, as they call themselves, the Starovers (Old Believers). They date from the time of Nicon, Patriarch of Moscow, and protest against the ritualistic innovations introduced by this remarkable man in the latter part of the seventeenth century, and afterwards by the Czar Peter the Great; they denounce the former as the false prophet, and the latter as the antichrist. They reject the benediction with three fingers instead of two, the pronouncing of the name of Jesus with two syllables instead of three, processions from right to left instead of the opposite course, the use of modern Russ in the service-books, the new mode of chanting, the use of Western pictures, the modern practice of shaving (unknown to the patriarchs, the apostles, and holy fathers), the use of tobacco (though not of whisky), and, till quite recently, also the eating of the potato (as the supposed apple of the devil, the forbidden fruit of paradise). They are again divided into several parties.

For information about these and other Russian Non-conformists, see Strahl: History of Heresies and Schisms in the Greek-Russian Church, and his Contributions to Russian Church History (I. 250 sqq.); Hepworth Dixon: Free Russia (1870), and the literature mentioned in Herzog's \emph{Encyklop}., Art. \emph{Raskolniken}, Vol. XII. p. 533.

\xchapter{Chapter 4. The Creeds of the Roman Church.}

\emph{General Literature.}

I. Collections of Roman Catholic Creeds:

J. Trg. Lbr. Danz: \emph{Libri Symbolici Ecclesiæ Romano-Catholicæ}, Weimar, 1885.

Fr. W. Streitwolf and  R. E. Klener: \emph{Libri Symbolici Ecclesiæ Catholicæ, conjuncti, atque notis, prolegomenis indicibusque instructi}, Götting. 1838, 2 vols. Contains the Conc. Trid., the Prof. Fidei Trid., and the Catech. Rom.

Henr. Denzinger (R. C., d. 1862): \emph{Enchiridion Symbolorum et Definitionum, quæ de rebus fidei et morum a Conciliis Œcumenicis et Summis Pontificibus enumarunt}, edit. quarta, Wirceburgi, 1865 (pp. 548). A convenient collection, including the definition of the Immaculate Conception of the Virgin Mary (1854), and the Papal Syllabus (1864).

II. Roman Catholic Expositions and Defenses of the Roman Catholic System:

Bellarmin's  \emph{Disputationes}, Bossuet's  \emph{Exposition}, Möhler's  \emph{Symbolik}, Perrone's  \emph{Prælectiones Theologicæ.} See § 23.

III. Protestant Expositions of the Roman Catholic system (exclusive of polemical works):

Ph. C. Marheineke (Prof. in Berlin, d. 1846): \emph{Christliche Symbolik oder historisch-kritische und dogmatisch-comparative Darstellung des kathol., luther., reform., und socinian. Lehrbegriffs}, Heidelb. 1810–13. The first 3 vols. (the only ones which appeared) are devoted to Catholicism.

W. H. D. Ed. Köllner (Prof. at Giessen): \emph{Symbolik der heil. apost. kathol. römischen Kirche}, Hamb. 1844. (Part II. of his unfinished \emph{Symbolik aller christlichen Confessionen}.)

A. H. Baier (Prof. at Greifswald): \emph{Symbolik der römisch-katholischen Kirche}, Leipz. 1854. (The first volume of an unfinished \emph{Symbolik der christlichen Religionen und Religionspartheien}.)

\xsection{Catholicism and Romanism.}

§ 22. Catholicism and Romanism.

The Roman Catholic Church embraces over 180 millions of members, or more than one half of nominal Christendom.\footnote{It is estimated that there are about 370 millions of Christians in the world, which is not much more than one fourth of the human family (1,370,000,000). Of these 370 millions the Roman Church may claim about 190, the Greek Church 80, the Protestant Church 100 millions. But the estimates of the Roman Catholic population vary from 180 to 200 millions.} It is spread all over the earth, but chiefly among the Latin races in Southern Europe and America.\footnote{Geographically speaking, the Roman Church may be called the Church of the South, the Greek Church the Church of the East, the Protestant Church the Church of the West.} It reaches in unbroken succession to the days of St. Peter and Paul, who suffered martyrdom in Rome. It is more fully developed and consolidated in doctrine, worship, and polity than any other Church. Its hierarchy is an absolute spiritual monarchy culminating in the Bishop of Rome, who pretends to be nothing less than the infallible Vicar of Jesus Christ on earth. It proudly identifies itself with the whole Church of Christ, and treats all other Christians as schismatics and heretics, who are outside of the pale of ordinary salvation.

But this unproved assumption is the fundamental error of the system. There is a vast difference between Catholicism and Romanism. The former embraces all Christians, whether Roman, Greek, or Protestant; the latter is in its very name local, sectarian, and exclusive. The holy Catholic Church is an article of faith; the Roman Church is not even named in the ancient creeds. Catholicism extends through all Christian centuries; Romanism proper dates from the Council of Trent. Mediæval Catholicism looked towards the Reformation; Romanism excludes and condemns the Reformation. So ancient Judaism, as represented by Abraham, Moses, and the Prophets, down to John the Baptist, prepared the way for Christianity, as its end and fulfillment; while Judaism, after the crucifixion of the Messiah, and the destruction of Jerusalem, has become hostile to Christianity. 'Catholicism is the strength of Romanism; Romanism is the weakness of Catholicism.'

In Romanism, again, a distinction must be made between the Romanism of the Council of Trent, and the Romanism of the Council of the Vatican. The 'Old Catholics' of Holland and Germany adhere to the former, but reject the latter as a new departure. But the papal absolutism has triumphed, and there is no room any longer for a moderate and liberal Romanism within the reign of the Papacy.

The doctrinal standards of the Roman Catholic Church may accordingly be divided into three classes:

1. The Œcumenical Creeds, which the Roman Church holds in common with the Greek, excepting the \emph{Filioque} clause, which the Greek rejects as an unauthorized, heretical, and mischievous innovation.\footnote{The Greek Church is as much opposed to this Latin interpolation as ever. The Encyclical Epistle of the Eastern Patriarchs and other prelates, in reply to the Epistle of Pius IX., dated Jan. 6, 1848, urges no less than fifteen arguments against the \emph{Filioque}, and reminds Pope Pius of the testimony of his predecessors, Leo III. and John VIII., 'those glorious and \emph{last orthodox} Popes.' Leo, when appealed to by the delegates of Charlemagne, in 809, caused the original Nicene Creed to be engraved on two tablets of silver, on the one in Greek, on the other in Latin, and these to be suspended in the Basilica of St. Peter, to bear perpetual witness against the insertion of the \emph{Filioque}. This fact, contrasted with the reverse action of later Popes, is one among the many proofs against papal infallibility.}

2. The Roman or Tridentine Creeds, in opposition to the evangelical doctrines of the Reformation. Here belong the Council of Trent, the Profession of Pius IV., and the Roman Catechism. They sanction a number of doctrines, which were prepared in part by patristic and scholastic theology, papal decrees, and mediæval councils, but had always been more or less controverted, viz., tradition as a joint rule of faith, the extent of the canon including the Apocrypha, the authority of the Vulgate, the doctrine of the primitive state and original sin, justification by works as well as by faith, meritorious works, seven sacraments, transubstantiation, the withdrawal of the cup, the sacrifice of the mass for the living and the dead, auricular confession and priestly absolution, extreme unction, purgatory, indulgences, and obedience to the authority of the Pope as the successor of Peter and vicar of Christ.

3. The modern Papal and Vatican decisions in favor of the immaculate conception of Mary, and the infallibility of the Pope. These were formerly open questions in the Roman Church, but are now binding dogmas of faith.

\xsection{Standard Expositions of the Roman Catholic System.}

§ 23. Standard Expositions of the Roman Catholic System.

Italy, France, and Germany have successively furnished the ablest champions of the doctrinal system of Romanism in opposition to Protestantism. Their authority is, of course, subordinate to that of the official standards. But as faithful expounders of these standards they have great weight. In Romanism, learning is concentrated in a few towering individuals; while in Protestantism it is more widely diffused, and presents greater freedom and variety of opinion.

1. The first commanding work in defense of Romanism, after many weak attempts of a purely ephemeral character, was written towards the close of the sixteenth century, more than fifty years after the beginning of the Protestant controversy, and about thirty years after the Council of Trent, by  Robert Bellarmin (Roberto Bellarmino). He was born 1542, in Tuscany, entered the order of the Jesuits in 1560, became Professor of Theology at Louvain in 1570, and afterwards at Rome, was made a Cardinal in 1599, Archbishop of Capua in 1602, Librarian of the Vatican in 1605, and died at Rome Sept. 17, 1621, nearly eighty years old. Although the greatest controversialist of his age, he had a mild disposition, and was accustomed to say that 'an ounce of peace was worth more than a pound of victory.' His '\emph{Disputations on the Controversies of the Christian Faith}' are the most elaborate polemic theology of the Roman Church against the doctrines of the Protestant Reformation.\footnote{The \emph{Disputationes de controversiis Christianæ fidei adversus hujus temporis hereticos} were first published at Ingolstadt, 1587–90, 3 vols. folio; then at Venice (but with many errors); at Cologne, 1620; at Paris, 1688; at Prague, 1721; again at Venice, 1721–27; at Mayence, 1842, and at Rome, 1832–40, in 4 vols. 4to. They are usually quoted by the titles of the different sections, \emph{De Verbo Dei, De Christo, De Romano Pontifice, De Conciliis et Ecclesia, Die Clericis, De Monachis, De Purgatorio}, etc. The contemporary \emph{Annals} of Baronius (d. 1607) are the most learned \emph{historical} vindication of Romanism in opposition to Protestantism and the 'Magdeburg Centuries.'} They abound in patristic and scholastic learning, logical acumen and dialectical ability. The differences between Romanism and Protestantism are clearly and accurately stated without any attempt to weaken them. And yet the book was placed on the Index Expurgatorius by Sixtus V. for two reasons; first, because Bellarmin introduces the doctrines of the Reformers in their own words, which it was feared might infect Romish readers with dangerous heresies; and, secondly, because he taught merely an indirect, not a direct, authority of the Pope in temporal matters. In France and Venice, on the contrary, even this doctrine of the indirect temporal supremacy was considered too ultramontane, and hence Bellarmin was never a favorite among the Gallicans. After the death of Sixtus V., the inhibition was removed. The work has ever since remained the richest storehouse of Roman controversialists, and can not be ignored by Protestants, although many arguments are now antiquated, and many documents used as genuine are rejected even by Catholics.

2. Nearly a century elapsed before another champion of Romanism appeared, less learned, but more eloquent and popular,  Jacques Bénigne Bossuet. He was born at Dijon, 1627, was educated by the Jesuits, tutor of the Dauphin 1670–81, Bishop of Meaux since 1681, Counselor of State 1697, and died at Paris 1704. The 'Eagle of Meaux' was the greatest theological genius of France, and the oracle of his age, a man of brilliant intellect, untiring industry, magnificent eloquence, and equally distinguished as controversialist, historian, and pulpit orator. He is called 'the last of the fathers of the Church.' While the hypocritical and licentious Louis XIV. tried to suppress Protestantism in his kingdom by cruel persecution, Bossuet betook himself to the nobler and more successful task of convincing the opponents by argument.

This he did in two works, the first apologetic, the second polemical.

(\emph{a}) \emph{Exposition de la doctrine de l’église catholique sur les matières de controverse}.\footnote{First published in Paris 1671, sixth ed. 1686, and often since in French, German, English, and other languages. It was approved and commended by the French clergy, even by Pope and Cardinals at that time, and attained almost the authority of a symbolical book. But the Jesuit father Maimbourg disapproved it.} This book is a luminous, eloquent, idealizing, and plausible defense of the characteristic doctrines of Romanism. It distinguishes between dogmas and theological opinions; presenting the former in a light that is least objectionable to reason, and disowning the latter when especially objectionable to Protestants. 'Bossuet assumes,' says Gibbon, 'with consummate art, the tone of candor and simplicity, and the ten-horned monster is transformed, by his magic touch, into a milk-white hind, who must be loved as soon as seen.'

(\emph{b}) \emph{Histoire des variations des églises protestantes}.\footnote{Paris, 1688, and often since in several languages. Compare also \emph{his Défense de l’histoire des variations contre M. Basnage.} Sir James Stephen says of the \emph{Variations}, that they bring to the religious controversy 'every quality which can render it either formidable or attractive.' The famous historian of the Decline and Fall of Rome was converted by this work to Romanism, but ended afterwards in infidelity. 'Bossuet shows,' says Gibbon in his \emph{Memoirs}, 'by a happy mixture of reasoning and narration, the errors, mistakes, uncertainties, and contradictions of our first Reformers, whose variations, as he learnedly maintains, bear the marks of error, while the uninterrupted unity of the Catholic Church is a sign and testimony of infallible truth. I read, approved, and believed.'} This is an attempt to refute Protestantism, by presenting its history as a constant variation and change; while the Roman Catholic system remained the same, and thus proves itself to be the truth. The argument is plausible, but not conclusive. It would prove more for the Greek Church than for the Latin, which has certainly itself developed from patristic to mediæval, from mediæval to Tridentine, and from Tridentine to Vatican Romanism. Truth in God, or objectively considered, is unchangeable; but truth in man, or the apprehension of it, grows and develops with man and with history. Change, if it be consistent, is not necessarily a mark of heresy, but may be a sign of life and growth, as the want of change, on the other hand, is by no means always an indication of orthodoxy, but still more frequently of stagnation.

Bossuet, with all his strong Roman Catholic convictions, was no infallibilist and no ultramontanist, but a champion of the Gallican liberties. He was the presiding genius of the clerical assembly of 1682, which framed the famous four Gallican propositions; and he wrote a book in their defense, which was, however, not published till some time after his death.\footnote{\emph{Defensio declarationis celeberrimæ, quam de potestate ecclesiaslica sanxit clerus Gallicanus} 1682, \emph{ex speciali jussu Ludovici M. scripta,} Luxemb. 1730, 2 vols.; in French, Paris, 1735, 2 vols.} He carried on a useless correspondence with the great Leibnitz for a reunion of the Catholic and Protestant churches, and proposed to this end a suspension of the anathemas of Trent and a general council in which Protestants should have a deliberative vote. Altogether, although he sanctioned the infamous revocation of the edict of Nantes (as '\emph{le plus bel usage de l’autorité royale}'), and secured the papal condemnation of the noble Fénelon (a man more humble and saint-like than himself), Bossuet can no longer be regarded as sound and orthodox, if judged by the standard of the Vatican Council.\footnote{Döllinger (\emph{Lectures on the Reunion of Churches}, 1872, Engl. translation, p. 90) says: 'Bossuet puts aside the question of infallibility, as a mere scholastic controversy, having no relation to faith; and this was approved at Rome at the time. Now, of course, he is no longer regarded in his own country as the classical theologian and most eminent doctor of modern times; but as a man who devoted his most learned and comprehensive work, the labor of many years, to the establishment and defense of a fundamental error, and spent many years of his life in the perversion of facts and distortion of authorities. For that must be the present verdict of every infallibilist on Bossuet.'}

3. The same may be said of  John Adam Möhler, the greatest German divine of the Roman Church, a man of genius, learning, and earnest piety. He was born 1796, at Igersheim, in the Kingdom of Würtemberg; was Professor of Theology in the University of Tübingen since 1822, at Munich since 1835, where he died in 1838. The great work of his life is his \emph{Symbolics}.\footnote{'\emph{Symbolik, oder Darstellung der dogmatischen Gegensätze der Katholiken und Protestanten nach ihren öffentlichen Bekenntniss-Schriften}.' It appeared first in 1832, at Mayence; the sixth edition in 1843, and was translated into French, English, and Italian. The English translation is by  James Burton Robertson, and bears the title, \emph{Symbolism; or, Exposition of the doctrinal differences between Catholics and Protestants, as evidenced in their symbolical writings} (Lond. 1843, in 2 vols.; republished in 1 vol., New York, 1844). It is preceded by a memoir of Möhler, and a superficial historical sketch of recent German Church history.} It is at once defensive and offensive, a vindication of Romanism and an attack upon Protestantism, and written with much freshness and vigor. It made a profound impression in Germany at a time when Romanism was believed to be intellectually dead or unable to resist the current of Protestant culture. Möhler was well acquainted with Protestant theology, and was influenced by the lectures and writings of Schleiermacher and Neander.\footnote{Neander told me that Möhler, when a student at Berlin, occasionally called on him, and seemed to him very modest, earnest, and inquiring after the truth. Hase calls him a 'delicate and noble mind,' and relates that when he began his academic career in Tübingen with him, Möhler was filled with youthful ideals, and regarded by Catholics as heterodox. (\emph{Handbuch der Prot. Polemik}, Pref. p. ix.)} He divests Romanism of its gross superstitions, and gives it an ideal and spiritual character. He deals, upon the whole, fairly and respectfully with his opponents, but makes too much argumentative use of the private writings and unguarded utterances of Luther. He ignores the post-Tridentine deliverances of Rome, says not a word about papal infallibility, and, although not a Gallican, he represents the antagonism of the episcopal and papal systems as a wholesome check upon extremes. He recognizes the deep moral earnestness from which the Reformation proceeded, deplores the corruptions in the Church, sends many ungodly popes and priests to hell, and talks of a feast of reconciliation, preceded by a common humiliation and confession that all have sinned and gone astray, the Church alone [meaning the institution] is without spot or wrinkle.\footnote{\emph{Symbolik} (6th edition, p. 353): '\emph{Unstreitig liessen es auch oft genug Priester, Bischöfe und Päpste, gewissenlos und unverantwortlich, selbst dort fehlen, wo es nur von ihnen abhing, ein schöneres Leben zu begründen; oder sie löschten gar noch durch ärgerliches Leben und Streben den glimmenden Docht aus, welchen sie anfachen sollten: die Hölle hat sie verschlungen. . . . Beide} [\emph{Katholiken und Protestanten}] \emph{müssen schuldbewusst ausrufen: Wir Alle haben gefehlt, nur die Kirche ist's, die nicht fehlen kann; wir Alle haben gesündigt, nur sie ist unbefleckt auf Erden.}' Incidentally Möhler denies the papal infallibility, when he says (p. 336): '\emph{Keinem einzelnen als solchen kommt diese Unverirrlichkeit zu.}'} His work called forth some very able Protestant replies, especially from Baur and Nitzsch.\footnote{Baur's \emph{Gegensatz des Katholicismus and Protestantismus} (Tübingen, 1833, 2d ed. 1836), in learning, grasp, and polemical dexterity, is fully equal or superior to Möhler's \emph{Symbolik}, but not orthodox, and elicited a lengthy and rather passionate defense from his Catholic colleague (\emph{Neue Untersuchungen}, Mainz, 1834). Nitzsch's \emph{Protestantische Beantwortung der Möhlerschen Symbolik} (Hamb. 1835) is sound, evangelical, calm, and dignified. It is respectfully mentioned, but not answered, by Möhler. Marheineke and Sartorius wrote, likewise, able replies. A counterpart of Möhler's \emph{Symbolik} is Hase's \emph{Handbuch der Protestantischen Polemik gegen die Römisch-Katholische Kirche}, Leipz. 1862; 3d ed. 1871. Against this work Dr. F. Speil wrote \emph{Die Lehren der Katholischen Kirche, gegenüber der Protestantischen Polemik}, Freiburg, 1865, which, compared with Möhler's book, is a feeble defense.}

4.  Giovanni Perrone, born in Piedmont, 1794, Professor of Theology in the Jesuit College at Rome, wrote a system of dogmatics which is now most widely used in the Roman Church, and which most fully comes up to its present standard of orthodoxy.\footnote{\emph{Prælectiones theologicæ quas in Collegio Romano Societatis Jesu habebat J. P.} They appeared first at Rome, 1835 sqq., in 9 vols. 8vo; also at Turin (31st ed. 1865 sqq. in 9 vols.); at Paris (1870, in 4 vols.); at Brussels, and Ratisbon. His compend, \emph{Prælectiones theologicæ in Compendium redactæ, }has been translated into several languages. Perrone wrote also separate works, \emph{De Jesu Christi Divinitate} (Turin, 1870, 3 vols.); \emph{De virtutibus fidei, spei et caritatis} (Tur. 1867, 2 vols.); \emph{De Matrimonio Christiano} (Lond. 1861), and on the Immaculate Conception of Mary.} Perrone defends the immaculate conception of Mary, and the infallibility of the Pope, and helped to mould the decrees of the Vatican Council. His method is scholastic and traditional, but divested of the wearisome and repulsive features of old scholasticism, and adapted to the modern state of controversy.

Note.—English Works on Romanism.—England and the United States have not produced a classical theological work on Romanism, such as those above mentioned, but a number of compilations and popular defenses. We mention the following: \emph{The Faith of Catholics on certain points of Controversy, confirmed by Scripture and attested by the Fathers of the Church during the five first centuries of the Church, compiled by }Rev.  Jos. Berington  \emph{and }Rev.  John Kirk, Lond. 1812, 1 vol.; 2d ed. 1830; 3d ed., revised and greatly enlarged, by Rev.  James Waterworth, 1846, in 3 vols. \emph{The End of Religious Controversy} (Lond. 1818, and often since), a series of letters by the Rt. Rev.  John Milner (born in London, 1752, d. 1826). \emph{Lectures on the Principal Doctrines and Practices of the Catholic Church, }delivered in London, 1836, by Cardinal Nicholas Wiseman (born in Spain, 1802, died in London, 1865).

At present the ablest champions of Romanism in England are ex-Anglicans, especially Dr.  John H. Newman (born in London, 1801) and Archbishop  Henry Edward Manning (born in London, 1809, Wiseman's successor), who use the weapons of Protestant culture against the Church of their fathers and the faith of their early manhood. Manning is an enthusiastic infallibilist, but Newman acquiesced only reluctantly in the latest dogmatic development.\footnote{The views of the older English Romanists are compiled and classified by  Samuel Capper (a Quaker), in the work, \emph{The Acknowledged Doctrines of the Church of Rome . . . as set forth by esteemed doctors of the said Church}, Lond. 1850 (pp. 608). It consists mostly of extracts from the comments in the Douay version of the Scriptures. Comp. an article in the (N.Y.) \emph{Catholic World} for Dec. 1873, on 'Catholic Literature in England since the Reformation.'}

The principal apologists of the Romish Church in America are Archbishops Kenrick and Spaulding, Bishop England, Dr.  Orestes Brownson (in his \emph{Review}), and more recently the editors, chiefly ex-Protestants, of the monthly '\emph{Catholic World}.' We mention  Francis Patrick Kenrick (Archbishop of Baltimore, born in Dublin 1797, died 1863): \emph{The Primacy of the Apostolic See Vindicated}, 4th ed. Balt. 1855, and \emph{A Vindication of The Catholic Church, in a Series of Letters to the Rt. Rev. J. H. Hopkins}, Balt. 1855. His brother,  Peter Richard Kenrick, Archbishop of St. Louis, was an opponent of the infallibility dogma in the Vatican Council, but has since submitted, like the rest of the bishops. In a lengthy and remarkable speech, which he had prepared for the Vatican Council, but was prevented from delivering by the sudden close of the discussion, June 3, 1870, he shows that the doctrine of papal infallibility was not believed either in Ireland, his former home, or in America; on the contrary, that it was formally and solemnly disowned by British bishops prior to the Catholic Emancipation bill.\footnote{See Kenrick's \emph{Concio habenda, at non habita} in Friedrich's \emph{Documenta}, I. 189–226.}

\xsection{The Canons and Decrees of the Council of Trent, A.D. 1563.}

§ 24. The Canons and Decrees of the Council of Trent.

Literature.

I. Latin Editions.

Paul. Manutius (d. 1574): \emph{Canones et Decreta Œcum. et Generalis Conc. Tridentini, jussu Pontificis Romani}, Rome, 1564, fol., 4to, and 8vo.

\emph{Canones et Decreta Œcum. et Generalis Conc. Trident. . . . Index dogm. et reformationum}, etc., Lovan. 1567, fol.

\emph{Canones et Decreta Œcum. et Generalis Conc. Trident. additis declarationibus cardinal. Ex ultima recognitione}   J. Gallemart  \emph{et citationibus}   J. Sotealli  \emph{et}   Hor. Luth,  \emph{nec non remissionib.}   Agst. Barbosæ (Cologne, 1620; Lyons, 1650, 8vo), \emph{quibus accedunt additiones}   Blo. Andræae, etc., Cologne (1664), 1712, 8vo.

Ph. Chifflet: \emph{S. Concilii Trid. Canones et Decreta cum præfatione}, Antw. 1640, 8vo.

Judov. le Plat (or Leplat; a very learned and moderate Catholic, d. 1810): \emph{Concilii Tridentini Canones et Decreta, juxta exemplar authenticum, Romæ} 1564 \emph{editum, cum variantibus lectionibus, notis Chiffletii}, etc., Antwerp, 1779; Madrid, 1786. The most complete Cath. edition.

Æm. Lud. Richter  \emph{et}   Frid. Schulte:  \emph{Canones et Decreta Concilii Tridentini ex editione Romana a.} 1834, \emph{repetiti}, etc., Leipz. 1853. Best Protestant ed.

\emph{Canones et Decreta sacrosancti Œcumenici Concilii Tridentini}, etc., Romæ, ed. stereotypa VII., Leipz. (Tauchnitz), 1854.

W. Smets: \emph{Concilii Tridentini sacrosancti œcumenici et generalis, Paulo III., Julio III., Pio IV., Pontificibus Maximis, celebrati, Canones et Decreta}. Latin and German, with a German introduction, 5th ed. Bielefeld, 1859.

The doctrinal decrees and canons are also given in Denzinger's \emph{Enchiridion.}

II. English Translations.

J. Waterworth (R.C.): \emph{The Canons and Decrees of the Sacred and Œcumenical Council of Trent} (with \emph{Essays on the External and Internal History of the Council}), London, 1848. (From Le Plat's edition.)

Th. A. Buckley (Chaplain of Christ Church, Oxford): \emph{The Canons and Decrees of the Council of Trent, }London, 1851.

There are also translations in French, German, Greek, Arabic, etc.

III. History of the Council.

Hardouin: \emph{Acta Conciliorum }(Paris, 1714), Tom. X. 1–435.

Jodov. Le Plat:\emph{ Monumentorum ad historiam Concilii Trid. potissimum illustrandum spectantium amplissima collectio, }Lovan. 1781–87, Tom. VII. 4to. The most complete documentary collection.

Fra  Paolo Sarpi (liberal Catholic, d. 1623): \emph{Istoria del concilio Tridentino, nella quale si scoprono tutti gl’artificii della corte di Roma, per impedire, che ne la verita di dogmi si palesasse, ne la riforma del papato e della chiesa si trattasse}, Lond. 1619, fol.; Geneva, 1629, 1660. Latin transl., Lond. 1620; Frankf. 1621; Amst. 1694; Leipz. 1699. French translation by \emph{Peter Francis Courayer}, with valuable historical notes, Lond. 1736, 2 vols. fol.; Amst. 1736, 2 vols. 4to; Amst. 1751, 3 vols. (Courayer was a liberal Roman Catholic divine, but, being persecuted, he fled from France to England, and joined the Anglican Church; d. 1776.) English translation by \emph{Sir Nathaniel Brent, }Lond. 1676, fol. German translations by \emph{Rambach }(with Courayer's notes), Halle, 1761, and by \emph{Winterer}, Mergentheim and Leipz. 2d ed. 1844.

Card.  Sforza Pallavicini (strict Catholic, d. 1667): \emph{Istoria del concilio di Trento}, Roma, 1656–57, 2 vols. fol., and other editions, original and translated. Written in opposition to Paul Sarpi. Comp. Brischar: \emph{Beurtheilung der Controversen Sarpi's und Pallavic.'s}, Tübing. 1843, 2 vols.

L. El. Du Pin (R.C.): \emph{Histoire du concile de Trente}, Brussels, 1721, 2 vols. 4to.

Chr. Aug. Salig (Luth.): \emph{Vollständige Historie des Trident. Conciliums}, Halle, 1741–45, 3 vols. 4to.

Jos. Mendham: \emph{Memoirs of the Council of Trent, principally derived from manuscript and unpublished Records}, etc., Lond. 1834; with a \emph{Supplement}, 1846.

J. Göschl: \emph{Geschichte des Conc. z. Tr.}, Regensburg, 1840, 2 vols.

J. H. von Wessenberg (a liberal R.C. and Bishop of Constance, d. 1860): \emph{Geschichte der grossen Kirchenversammlungen des }15 \emph{und }16 \emph{ten Jahrh}., Constance, 1840, Vol. III. and IV.

Card.  Gabr. Paleotto: \emph{Acta Concilii Trid. ab a }1562 \emph{descr}., ed. \emph{Mendham}, Lond. 1842.

Ed. Köllner: \emph{Symbolik der röm. Kirche}, Hamb. 1844, pp. 7–140.

J. T. L. Danz: \emph{Gesch. des Trid. Conc., }Jena, 1846.

Th. A. Buckley: \emph{History of the Council of Trent, }London, 1852.

Felix Bungener: \emph{Histoire du Concile de Trente}, Paris, 2d edition, 1854. English translation, Edinburgh, 1852, and New York, 1855. Also in German, Stuttg. 1861, 2 vols.

A. Baschet: \emph{Journal du Concile de Trente, redigé par un secrétaire vénitien present aux sessions de }1562 \emph{à} 1563, \emph{avec d'autres documents diplomatiques relatifs à la mission des Ambassadeurs de France}, Par. 1870.

Th. Sickel: \emph{Zur Geschichte des Concils von Trient. Actenstücke aus österreichischen Archiven}, Wien, 1872 (650 pp.). Mostly letters to the German Emperor, in Latin and Italian, from 1559 to 1563.

Augustin Theiner (Priest of the Oratory, d.1874): \emph{Acta genuina SS. Œcumenici Concilii Tridentini . . . nunc primum integra edita}. Zagrabiæ (Croatiæ) et Lipsiæ, 1874, 2 Tom. 4to (pp. 722 and 701).

Jos. von Döllingke: \emph{Ungedruckte Berichte und Tagebücher zur Geschichte des Conc. von Trient}, Nördlingen, 1876.

The principal source and the highest standard of the doctrine and discipline of the Roman Church are the Canons and Decrees of the Council of Trent, first published in 1564, at Rome, by authority of Pius IV.\footnote{The editor of this rare authentic edition was the learned  Paulus Manutius (Paolo Manuzio), Professor of Eloquence and Director of the Printing-Press of the Venetian Academy, settled at Rome 1561, and died there 1574. Not to be confounded with his father, Aldo Manuzio, sen. (1447–1515), the editor of the celebrated editions of the classics; nor with his son, Aldo Manuzio, the younger (1547–1597), likewise a printer and writer, and Professor of Eloquence.}

The Council of Trent (1543–63) is reckoned by the Roman Church as the eighteenth (or twentieth) œcumenical Council.\footnote{There is a dispute about the reformatory Councils of Pisa (1409), Constance (1414–18), and Basle (1431), which are acknowledged by the Gallicans, but rejected by the Ultramontanists, or accepted only in part, i.e., as far as they condemned and punished heretics (Hus and Jerome of Prague). The Council of Ferrara and Florence (1439) is regarded as a continuation of, or a substitute for, the Council of Basle. There is also a dispute among Roman historians about the œcumenical character of the Council of Sardica (343), the Quinisexta (692), the Council of Vienne (1311), and the fifth Lateran (1512–17). See Hefele, \emph{Conciliengeschichte}, Vol. I. 50 sqq.} It is also the last, with the exception of the Vatican Council of 1870, which, having proclaimed the Pope infallible, supersedes the necessity and use of any future councils, except for unmeaning formalities. It was called forth by the Protestant Reformation, and convened for the double purpose of settling the doctrinal controversies, which then agitated and divided Western Christendom, and of reforming discipline, which the more serious Catholics themselves, including even an exceptional Pope (Adrian VI.), desired and declared to be a crying necessity.\footnote{Adrian VI., from Holland, the teacher of Charles V., and the last non-Italian Pope, succeeded Leo X. in 1522, but ruled only one year. 'He died of the papacy.' He was a man of ascetic piety, and openly confessed, through his legate Chieregati, at the Diet of Nurnberg, that the Church was corrupt and diseased, from the Pope and the papal court to the members; but at the same time he demanded the sharpest measures against Luther as a second Mohammed. Twelve years later, Paul III. (1534–49) appointed a reform commission of nine pious Roman prelates, who in a memorial declared that the Pope's absolute dominion over the whole Church was the source of all this corruption; but he found it safer to introduce the Inquisition instead of a reformation.} The Popes, jealous of deliberative assemblies, which might endanger their absolute authority, and afraid of reform movements, which might make concessions to heretics, pursued a policy of evasion and intrigue, and postponed the council again and again, until they were forced to yield to the pressure of public opinion. Pius IV. told the Venetian embassador that his predecessors had professed a wish for a council, but had not really desired it.

In the early stages of the Reformation, Luther himself appealed to a general council, but he came to the conviction that even general councils had erred (e.g., the Council of Constance in condemning Hus), so that he had to trust exclusively to the Word of God and the Spirit of God in history. In deference to the special wish of the Emperor Charles V., the evangelical princes and divines were invited; but being refused a deliberative voice, they declined. 'They could not fail,' they replied, 'to appreciate the efforts of the Emperor, and they themselves were longing for an impartial council to be controlled by the supreme authority of the Scriptures, but they could not acknowledge nor attend a Roman council where their cause was to be judged after papal decrees and scholastic opinions, which had always found opposition in the Church. The council promised by the Pope would be neither free nor Christian, nor œcumenical, nor ruled by the Word of God; it would only confirm the authority of the Pope, on whom it was depending, and prove a new compulsion of conscience.' The result shows that these apprehensions were well founded.\footnote{At the second period of the Council, 1552, a number of Protestant divines from Württemberg, Strasburg, and Saxony, arrived in Trent, or were on the way, but they demanded a revision of the previous decrees and free deliberation, which were refused.}

After long delays the Council was opened by order of Pope Paul III., in the Austrian City of Trent (since 1917, belonging to Italy), on the 13th December, 1545, and lasted, with long interruptions, till the 4th of December, 1563. The attendance varied in the three periods: under Paul III. the number of prelates never exceeded 57, under Julius III. it rose to 62, under Pius IV. it was much larger, but never reached the number of the first œcumenical Council (318). The decrees were signed by 255 members, viz., 4 legates of the Pope, 2 Cardinals, 3 Patriarchs, 25 Archbishops, 168 Bishops, 39 representatives of absent prelates, 7 Abbots, and 7 Generals of different orders. Two thirds of them were Italians. \&gt;From France and Poland only a few dignitaries were present; the greater part of the German Bishops were prevented from attendance by the war between the Emperor and the Protestants in Germany. The theologians who assisted the members of the Synod belonged to the monastic orders most devoted to the Holy See.

The pontifical party controlled the preliminary deliberations as well as the final decisions, in spite of those who maintained the rights of an independent episcopacy.\footnote{The overruling influence of the papal court over the Council rests not only on the authority of Paolo Sarpi, but on many contemporary testimonies, e.g., the reports of Franciscus de Vargas, a zealous Catholic, who was used by Charles V. and Philip II. for the most important missions, who watched the proceedings of the Council at Trent from 1551 to '52 and gave minute information to Granvella. See \emph{Lettres et Mémoires de } Fr. de Vargas,  \emph{de Pierre de Malvenda et des quelques erèques d'Espagne, trad. par Michel le Vassor, }Amst. 1699; also in Latin, by Schramm, Brunswick. 1704. Le Plat pronounced this correspondence fictitious, but its authenticity has been sufficiently established (see Köllner, l.c. pp. 40, 41).}

During a period of nearly twenty years twenty-five public sessions were held, of which about one half were spent in mere formalities. But the principal work was done in the committees or congregations. The articles of dispute were always fixed by the papal legates, who presided. They were then first discussed, often with considerable difference of opinion, in the private sessions of the 'Congregations,' and after being secretly reported to, and approved by, the court of Rome, the Synod, in public session, solemnly proclaimed the decisions. They are generally framed with consummate scholastic skill and prudence.

The decisions of the Council relate partly to doctrine, partly to discipline. The former are divided again into Decrees (\emph{decreta}), which contain the positive statement of the Roman dogma, and into short Canons (\emph{canones}), which condemn the dissenting views with the concluding '\emph{anathema sit}.' The Protestant doctrines, however, are almost always stated in an exaggerated form, in which they would hardly be recognized by a discriminating evangelical divine, or they are mixed up with real heresies, which Protestants condemn as emphatically as the Church of Rome.\footnote{Thus the \emph{Canones de Justificatione} (Sess. VI.) reject Pelagianism and Semi-Pelagianism, as well as Solifidianism and Antinomianism.}

The doctrinal sessions, which alone concern us here, are the following:

SESSIO

III.

Decretum de Symbolo Fidei (accepting the Niceno Constantinopolitan Creed as a basis of the following decrees (Febr. 4, 1546).

``

IV.

Decretum de Canonicis Scripturis (Apr. 8, 1546).

``

V.

De Peccato Originali (June 17, 1546).

``

VI.

De Justificatione (Jan. 13, 1547).

``

VII.

De Sacramentis in genere, and some Canones de Baptismo et Confirmatione (March 3, 1547).

``

VIII.

De Eucharistiæ Sacramento (Oct. 11, 1551).

``

XIV.

De S. Pœnitentiæ et Extreme Unctionis Sacramento (Nov. 25, 1551).

``

XXI.

De Communione sub utraque Specie et Parvulorum (July 16, 1562).

``

XXII.

Doctrina de Sacrificio Missæ (Sept. 17, 1562).

``

XXIII.

Vera et Catholica de Sacramento Ordinis doctrina (July 15, 1563).

``

XXIV.

Doctrina de Sacramento Matrimonii (Nov. 11, 1563).

``

XXV.

Decretum de Purgatorio, Doctrina de Invocatione, Veneratione et Reliquiis Sanctorum, et sacris Imaginibus. Decreta de Indulgentiis, de Delectu Ciborum, Jejuniis et Diebus Festis, de Indice Librorum, Catechismo, Breviario et Missali (Dec. 3 and 4, 1563).

The last act of the Council was a double curse upon all heretics.\footnote{The Cardinal of Lorraine said, '\emph{Anathema cunctis hereticis}.' To this the fathers responded, '\emph{Anathema, Anathema}.'}

The decrees, signed by 255 fathers, were solemnly confirmed by a bull of Pius IV. (\emph{Benedictus Deus et Pater Domini nostri}, etc.) on the 26th January, 1564, with the reservation of the exclusive right of explanation to the Pope.

The Council was acknowledged in Italy, Portugal, Spain, France, the Low Countries, Poland, and the Roman Catholic portion of the German Empire; but mostly with a reservation of the royal prerogatives. In France it was never published in form. No attempt was made to introduce it into England. Pius IV. sent the acts to Queen Mary of Scots, with a letter, dated June 13, 1564, requesting her to publish them in Scotland, but without effect.\footnote{On the reception, see the seventh volume of Le Plat's Collection of Documents, Courayer's  \emph{Histoire de la reception du Concil de Trente, dans les differens états catholiques}, Amst. 1756 (Paris, 1766), and Köllner, l.c. pp. 121–129.}

The Council of Trent, far from being truly œcumenical, as it claimed to be, is simply a Roman Synod, where neither the Protestant nor the Greek Church was represented; the Greeks were never invited, and the Protestants were condemned without a hearing. But in the history of the Latin Church, it is by far the most important clerical assembly, unless the unfinished Vatican Council should dispute with it that honor, as it far exceeded it in numbers. It completed, with the exception of a few controverted articles, the doctrinal system of mediæval Catholicism, and stamped upon it the character of exclusive Romanism. It settled its relation to Protestantism by thrusting it out of its bosom with the terrible solemnities of an anathema. Papal diplomacy and intrigue outmanaged all the more liberal elements. At the same time the Council abolished various crying abuses, and introduced wholesome disciplinary reforms, as regards the sale of indulgences, the education and morals of the clergy, the monastic orders, etc. Thus the Protestant Reformation, after all, had indirectly a wholesome effect upon the Church which condemned it.

The original acts of the Council, as prepared by its general secretary, Bishop Angelo Massarelli, in six large folio volumes, are deposited in the Vatican, and have remained there unpublished for more than three hundred years. But most of the official documents and private reports bearing upon the Council were made known in the sixteenth century, and since. The most complete collection of them is that of Le Plat. New materials were brought to light by Mendham (from the manuscript history of Cardinal Paleotto), by Sickel, and by Döllinger. The genuine acts, but only in part, were edited by Theiner (1874).

The history of the Council was written chiefly by two able and learned Catholics of very different spirit: the liberal, almost semi-Protestant monk Fra  Paolo Sarpi, of Venice (first, 1619); and, in the interest of the papacy, by Cardinal  Sforza Pallavicini (1656), who had access to all the archives of Rome. Both accounts must be compared.

The first learned and comprehensive criticism of the Tridentine doctrine, from a Protestant point of view, was prepared by an eminent Lutheran theologian,  Martin Chemnitz (d. 1586), in his \emph{Examen Concilii Tridentini} (1565–73, 4 Parts), best ed., Frankf., 1707; republished, Berlin, 1861.\footnote{The editor, Ed. Preuss, has since become a Romanist at St. Louis (1871).}

\xsection{The Profession of the Tridentine Faith, A.D. 1564.}

§ 25. The Profession of the Tridentine Faith, 1564.

G. C. F. Mohnike:  \emph{Urkundliche Geschichte der sogenannten Professio Fidei Tridentinæ und einiger andern röm. katholischen Glaubensbekenntnisse}, Greifswald, 1822 (310 pp.).

Streitwolf  \emph{et}   Klener:  \emph{Libri Symbolici Ecclesiæ Catholicæ}, Gött. 1838, Tom. I. pp. xlv.–li. and 98–100.

Köllner: \emph{Symbolik der röm. Kirche}, pp. 141–165.

The older literature see in  Walch:  \emph{Biblotheca theol. sel.}, I. p. 410; and in  Köllner, l.c. p. 141.

Next in authority to the decrees of the Council of Trent, or virtually superior to it, stands the Professio Fidei Tridentinæ, or the Creed of Pius IV.\footnote{The original name was \emph{Forma juramenti professionis fidei}. In the two papal bulls which published and enjoined the creed, it is called \emph{Forma professionis fidei catholicæ}, or \emph{orthodoxæ fidei}. The usual name is \emph{Professio fidei Tridentinæ} (or \emph{P. f. Tridentina}, which is properly a misnomer). See Mohnike, l.c. p. 3, and Köllner, 1.c. p. 150.}

It was suggested by the Synod of Trent, which in its last two sessions declared the necessity of a binding formula of faith (\emph{formula professionis et juramenti}) for all dignitaries and teachers of the Catholic Church.\footnote{Sess. XXV. cap. 2 \emph{De Reformatione} (p. 439, ed. Richter): '\emph{Cogit temporum calamitas et invalescentium hæresum malitia, ut nihil sit prætermittendum, quod ad populorum ædificationem et catholicæ fidei præsidium videatur posse pertinere. Præcipit igitur sancta synodus patriarchis, primatibus, archiepiscopis, episcopis, et omnibus aliis, qui de jure vel consuetudine in concilio provinciali interesse debent, ut in ipsa prima synodo provinciali, post finem præsentis concilii habenda, ea omnia et singula, quæ ab hac sancta synodo definita et statuta sunt, palam recipiant, nec non veram obedientiam summo Romano Pontifici spondeant et profiteantur, simulque hæreses omnes, a sacris canonibus et generalibus conciliis, præsertimque ab hac eadme synodo damnatas, publice detestentur et anathematizent.}' Comp. Sess. XXIV. \emph{De Reformatione}, cap. 12, where an examination and profession (\emph{orthodoxæ fidei publica professio}) is required from the clergy, together with a vow to remain obedient to the Roman Church (\emph{in ecclesiæ Romanæ obedientia se permansuros spondeant ac jurent}).} It was prepared by order of Pope Pius IV., in 1564, by a college of Cardinals.

It consists of twelve articles: the first contains the Nicene Creed in full, the remaining eleven are a clear and precise summary of the specific Roman doctrines as settled by the Council of Trent, together with the important additional declaration that the Roman Church is the mother and teacher of all the rest, and with an oath of obedience to the Pope, as the successor of the Prince of the apostles, and the vicar of Christ.\footnote{'\emph{Sanctum catholicam et apostolicam Romanam ecclesiam omnium ecclesiarum matrem et magistram agnosco, Romanoque Pontifici, beati Petri Apostolorum principis successori ac Jesu Christi vicario, veram obedientiam spondeo ac juro.}' Here the 'catholic' Church is identified with the 'Roman' Church, and true obedience to the Pope is made a test of catholicity. The union decree of the Council of Florence makes a similar assertion (see Hardouin, \emph{Acta Conc.} ix. 423): '\emph{Item definimus, sanctam apostolicam sedem et Romanum Pontificem in universum orbem tenere primatum, et ipsum Pontificem Romanum successorem esse beati Petri principis Apostolorum, et verum Christi vicarium, totiusque ecclesiæ caput et omnium Christianorum patrem et doctorem existere.}' But the integrity of the text of this famous union formula is disputed, and the Greeks and Latins charge each other with corruption. Some Greek copies omit the proud words \textgreek{τὸν Ῥωμαικὸν ἀρχιερέα εἰς πᾶσαν τὴν οἰκουμένην τὸ πρωτεῖον κατέχειν.} Comp. Theod. Frommann: \emph{Zur Kritik des Florentiner Unionsdecrets and seiner dogmatischen Verwerthung beim Vaticanischen Concil}, Leipz. 1870, pp. 40 sqq.} The whole is put in the form of an individual profession ('\emph{Ego}, ——, \emph{firma fide credo et profiteor}'), and of a solemn vow and oath ('\emph{spondeo, voveo ac juro. Sic me Deus adjuvet, et hæc sancta Evangelia}').

This formula was made binding, in a double bull of Nov. 13, 1564 ('\emph{Injunctum noblis}'), and Dec. 9, 1564 ('\emph{In sacrosancta beati Petri, principis apostolorum, cathedra,}' etc.), upon the whole \emph{ecclesia docens,} i.e., upon all Roman Catholic priests and public teachers in Catholic seminaries, colleges, and universities. Besides, it has come to be generally used, without special legislation, as a creed for Protestant converts to Romanism, and hence it is called sometimes the 'Profession of Converts.'\footnote{For converts from the Greek Church the form was afterwards (1575) modified by a reference to the compromise of the Council of Florence. See the \emph{Professio Fidei Græcis præscripta a Gregorio XIII.}, in Denzinger's \emph{Enchir.}, p. 294, and the \emph{Professio Fidei Orientalibus præscripta ab Urbano VIII. et Benedicto XIV.}, ibid., p. 296. For Protestants other forms of abjuration were occasionally used, without official sanction. The infamous Hungarian formula for Protestant converts (\emph{Confessio novorum Catholicorum in Hungaria}, first published 1674) is disowned by liberal Catholics as a foul Protestant forgery, but seems to have been used occasionally by Jesuits during the cruel persecutions of Protestants in Hungary and Bohemia in the 17th century. It contains the most extravagant Jesuit views on the authority of the Pope, the worship of the Virgin, the power of the priesthood, and pronounces awful curses on Protestant parents, teachers, and relations ('\emph{maledictos pronuntiamus parentes nostros,}' etc.), and on the evangelical faith, with the promise to persecute this faith in every possible way, even by the sword ('\emph{Juramus etiam, donec una gutta sanguinis in corpore nostro exstiterit, doctrinam maledictam illam evangelicam nos omnimodo, clam et aperte, violenter et fraudulenter, verbo et facto persecuturos, ense quoque non excluso}'). See the formula in Mohnike, l.c. pp. 88–92, in Streitwolf and Klener, II. pp. 343–346; and an account of the controversies concerning it in Köllner, l.c. pp. 159–165, and especially the monograph of Mohnike: \emph{Zur Geschichte des Ungarischen Fluchformulars} (an Appendix to his History of the Profession of the Tridentine Faith), Greifswald, 1823, 264 pages. A copy of this rare book is in the library of the Union Theological Seminary of New York.} For both purposes it is far better adapted than the Decrees of the Council of Trent, which are too learned and extensive for popular use.

As this Profession of Pius IV. is the most concise and, practically, the most important summary of the doctrinal system of Rome, we give it in full, and arrange it in three parts, so that the difference between the ancient Catholic faith, the later Tridentine faith, and the oath of obedience to the Pope as the vicar of Christ, may be more clearly seen. It should be remembered that the Nicene Creed was regarded by the ancient Church as final, and that the third and fourth œcumenical Councils solemnly, and on the pain of deposition and excommunication, forbade the setting forth of any new creed.\footnote{Conc. Ephes. (431), Canon VII.; Conc. Chalced. (451), after the definition of faith.} To bring the Tridentine formula up to the present standard of Roman orthodoxy, it would require the two additional dogmas of the immaculate conception, and papal infallibility.

TRANSLATION OF THE PROFESSION.\footnote{See the Latin text in the two bulls of Pius IV. above mentioned, also in Mohnike, 1.c. pp. 46 sqq., in Streitwolf and Klener, \emph{Libri Symb.} I. 98–100 (with the various readings), and in Denzinger, \emph{Enchir.}, p. 98. Also Mirbt, pp. 337–40. For additions to the oath, Vol. II. 210.}

I. The Nicene Creed of 381, with the Western Changes.

(See p. 27.)

1. I, ——, with a firm faith, believe and profess all and every one of the things contained in the symbol of faith, which the holy Roman Church makes use of, viz.:

I believe in one God the Father Almighty, Maker of heaven and earth, and of all things visible and invisible.

And in one Lord Jesus Christ, the only-begotten Son of God, begotten of the Father before all worlds; \emph{God of God}, Light of Light, very God of very God, begotten, not made, being of one substance with the Father; by whom all things were made;

Who, for us men, and for our salvation, came down from heaven, and was incarnate by the Holy Ghost of the Virgin Mary, and was made man;

He was crucified for us under Pontius Pilate; suffered and was buried;

And the third day he rose again, according to the Scriptures;

And ascended into heaven; sitteth on the right hand of the Father;

And he shall come again, with glory, to judge the quick and the dead; whose kingdom shall have no end.

And in the Holy Ghost, the Lord, and Giver of life; who proceedeth from the Father \emph{and the Son}; who with the Father and the Son together is worshiped and glorified; who spake by the Prophets.

And one holy catholic and apostolic Church;

I acknowledge one baptism for the remission of sins;

And I look for the resurrection of the dead;

And the life of the world to come. Amen.

II. Summary of the Tridentine Creed (1563).

2. I most steadfastly admit and embrace the apostolic and ecclesiastical traditions, and all other observances and constitutions of the same Church.

3. I also admit the holy Scriptures according to that sense which our holy Mother Church has held, and does hold, to which it belongs to judge of the true sense and interpretation of the Scriptures; neither will I ever take and interpret them otherwise than according to the unanimous consent of the Fathers (\emph{juxta unanimem consensum Patrum}).\footnote{It is characteristic that the Scriptures are put after the traditions, and admitted only in a restricted sense, the Roman Church being made the only interpreter of the Word of God. Protestantism reverses the order, and makes the Bible the rule and corrective of ecclesiastical traditions.}

4. I also profess that there are truly and properly seven sacraments of the new law, instituted by Jesus Christ our Lord, and necessary for the salvation of mankind, though not all for every one, to wit: baptism, confirmation, the eucharist, penance and extreme unction, holy orders, and matrimony; and that they confer grace; and that of these, baptism, confirmation, and ordination can not be reiterated without sacrilege. I also receive and admit the received and approved ceremonies of the Catholic Church used in the solemn administration of the aforesaid sacraments.

5. I embrace and receive all and every one of the things which have been defined and declared in the holy Council of Trent concerning original sin and justification.

6. I profess likewise that in the mass there is offered to God a true, proper, and propitiatory sacrifice for the living and the dead (\emph{verum, proprium, et propitiatorium sacrificium pro vivis et defunctis}); and that in the most holy sacrament of the eucharist there is truly, really, and substantially (\emph{vere, realiter, et substantialiter}) the body and blood, together with the soul and divinity of our Lord Jesus Christ; and that there is made a change of the whole essence (\emph{conversionem totius substantiæ}) of the bread into the body, and of the whole essence of the wine into the blood; which change the Catholic Church calls transubstantiation.

7. I also confess that under either kind alone Christ is received whole and entire, and a true sacrament.

8. I firmly hold that there is a purgatory, and that the souls therein detained are helped by the suffrages of the faithful.

Likewise, that the saints reigning with Christ are to be honored and invoked (\emph{venerandos atque invocandos esse}), and that they offer up prayers to God for us; and that their relics are to be held in veneration (\emph{esse venerandas}).\footnote{This should properly be a separate article, but in the papal bulls it is connected with the eighth article.}

9. I most firmly assert that the images of Christ and of the perpetual Virgin, the Mother of God, and also of other saints, ought to be had and retained, and that due honor and veneration are to be given them.

I also affirm that the power of indulgences was left by Christ in the Church, and that the use of them is most wholesome to Christian people.\footnote{This should likewise be a separate article, but is made a part of article 9.}

III. Additional Articles and Solemn Pledges (1564).

10. I acknowledge the holy Catholic Apostolic Roman Church as the mother and mistress of all churches, and I promise and swear (\emph{spondeo ac juro}) true obedience to the Bishop of Rome, as the successor of St. Peter, prince of the Apostles, and as the vicar of Jesus Christ.

11. I likewise undoubtingly receive and profess all other things delivered, defined, and declared by the sacred Canons and œcumenical Councils, and particularly by the holy Council of Trent; and I condemn, reject, and anathematize all things contrary thereto, and all heresies which the Church has condemned, rejected, and anathematized.

12. I do at this present freely profess and truly hold this true Catholic faith, without which no one can be saved (\emph{extra quam nemo salvus esse potest}); and I promise most constantly to retain and confess the same entire and inviolate,\footnote{For \emph{inviolatam} the Roman Bullaria read \emph{immaculatam.}} with God's assistance, to the end of my life. And I will take care, as far as in me lies, that it shall be held, taught, and preached by my subjects, or by those the care of whom shall appertain to me in my office. This I promise, vow, and swear—so help me God, and these holy Gospels of God.

\xsection{The Roman Catechism, A.D. 1566.}

§ 26. Roman Catechism, 1566.

Latin Editions.

\emph{Catechismus ex decreto Conc. Trident. Pii V. jussu editus}, Romæ ap. Paulum Manutium, 1566, in editions of different sizes, very often reprinted all over Europe.

\emph{Catechismus ad Parochos, ex decreto Concilii Tridentini editus. Ex Pii V. Pont. Max. jussu promulgatus. Syncerus et integer, mendisque iterum repurgatus operâ P. D. L. H. P. A quo est additus apparatus ad Catechismum, in quo ratio, auctores, approbatores, et usus declarantur}, Lugduni, 1659: Paris, 1671; Lovan. 1678; Paris, 1684; Colon. 1689, 1698, 1731; Aug. Vindel. 1762; Lugdun. 1829; Mechlin, 1831; Ratisb. 1856 (730 pp.).

\emph{Catechismus ex decreto Conc. Tridentini ad Parochos Pii Quinti Pont. Max. jussu editus. Ad editionem Romæ A.D.} 1566 \emph{juris publici factam accuratissime expressus}, ed. stereotypa VI., Lipsiæ (Tauchnitz), 1859, 8vo.

Also in Streitwolf  \emph{et}  Klener:  \emph{Libri Symb. eccl. cath.}, Tom. I. pp. 101–712. A critical edition, indicating the different divisions, the quotations from the Scriptures, the Councils, and other documents.

Translations.

\emph{The Catechism for the Curates, composed by the Council of Trent, and published by command of Pope Pius the Fifth. Faithfully translated into English. Permissu superiorum}. London, 1687.

\emph{The Catechism of the Council of Trent, translated into English by J. Donovan}, Baltimore, 1829.

\emph{The Catechism of the Council of Trent, translated into English, with Notes, by T. A. Buckley, B.A.}, London, 1852, 8vo.

German translations, first, by \emph{Paul Hoffäus}, Dillingen, 1568, 1576; another at Wien, 1763; one by \emph{T. W. Bodemann}, Göttingen, 1844; and by \emph{Ad. Buse}, Bielefeld, (with the Lat. text), 3d ed. 1867, 2 vols.

French translations, published at Bordeaux, 1568; Paris, 1578, 1650 (by \emph{P. de la Haye}), 1673, etc.

History.

Julii Pogiani Sunensis (d. 1567): \emph{Epistolæ et Orationes olim collectæ ab Antonio Maria Gratiano, nunc ab Hieronymo Lagomarsinio e Societate Jesu advocationibus illustratæ ac primum editæ}, Rom., Vol. I. 1752; II. 1756; III. 1757; IV. 1758.

\emph{Apparatus ad Catechismum}, etc., mentioned above, by an anonymous author (perhaps Anton. Reginaldus), first published in the edition of the Catechism, Lugd. 1659. The chief source of information.

J. C. Köcher: \emph{Catech. Geschichte der Pübstlichen Kirche}, Jen. 1753.

Köllner: \emph{Symbolik der röm. Kirche}, pp. 166–190. K. gives a list of other works on the subject.

The Roman Catechism was proposed by the Council of Trent, which entered upon some preparatory labors, but at its last session committed the execution to the Pope.\footnote{Sessio XXIV. \emph{De Reformatione}, cap. 7 (ed. Richter, p. 344), the Bishops are directed to provide for the instruction of Catholics, '\emph{juxta formam a sancta synodo in catechesi singulis sacramentis præscribendam, quam episcopi in vulgarem linguam fideliter verti, atque a parochis omnibus populo exponi curabunt.}' According to Sarpi, a draft of the proposed Catechism was laid before the Synod, but rejected. In the 25th and last session (held Dec. 24, 1563), the Synod intrusted the Pope (Pius IV.) with the preparation of an index of prohibited books, a catechism, and an edition of the liturgical books ('\emph{idemque de catechismo a Patribus, quibus illud mandatum fuerat, et de missali, et breviario fieri mandat,}' p. 471).} The object was to regulate the important work of popular religious instruction, and to bring it into harmony with the decisions of the Council.\footnote{Several catechisms, not properly authorized, had appeared before and during the Council of Trent to counteract the Lutheran and Reformed Catechisms, which did so much to spread and popularize the Reformation. See a list of them in Streitwolf and Klener, I. p. i.–iv., and in Köllner, p. 169.} Pius IV. (d. 1565), under the advice of Cardinal Carlo Borromeo (Archbishop of Milan), intrusted the work to four eminent divines, viz.,  Leonardo Marini (afterwards Archbishop of Lanciano),  Egidio Foscarari (Bishop of Modena),  Muzio Calini (Archbishop of Jadera-Zara, in Dalmatia), and  Francesco Fureiro (of Portugal). Three of them were Dominicans (as was the Pope himself). This explains the subsequent hostility of the Jesuits. Borromeo superintended the preparation with great care, and several accomplished Latin scholars, especially Jul. Pogianus, aided in the style of composition.\footnote{Winer, Guericke, Möhler, and others, ascribe the Latinity of the Catechism to Paulus Manutius, the printer of the same; but he himself, in his epistles, where he mentions all his literary labors, says nothing about it.} The Catechism was begun early in 1564, and substantially finished in December of the same year, but subjected for revision to Pogianus in 1565, and again to a commission of able divines and Latinists. It was finally completed in July, 1566, and published by order of Pope Pius V., in September, 1566, and soon translated into all the languages of Europe. Several Popes and Bishops recommended it in the highest terms. The Dominicans and Jansenists often appealed to its authority in the controversies about free will and divine grace, but the Jesuits (Less, Molina, and others) took ground against it, and even charged it with heresy.

The work is intended for teachers (as the title \emph{ad Parochos} indicates), not for pupils. It is a very full popular manual of theology, based upon the decrees of Trent. It answers its purpose very well, by its precise definitions, lucid arrangement, and good style.

The Roman Catechism treats, in four parts: 1, \emph{de Symbolo apostolico}; 2, \emph{de Sacramentis}; 3, \emph{de Decalogo}; 4, \emph{de Oratione Dominica}. It was originally written and printed without divisions.\footnote{The division into four parts, and of these into chapters and questions, appeared first in the edition of Fabricius Lodius, Col. 1572, and Antw. 1574. Other editions vary in the arrangement.} Its theology belongs to the school of Augustine and Thomas Aquinas, and hence it displeased the Jesuits. While it passes by certain features of the Roman system, as the indulgences and the rosary, it treats of others which were not touched upon by the Fathers of Trent, as the \emph{limbus patrum}, the doctrine of the Church, and the authority of the Pope.

Notwithstanding the high character and authority of this production, it did not prevent the composition and use of many other catechisms, especially of a more popular kind and in the service of Jesuitism. The most distinguished of these are two Catechisms of the Jesuit  Peter Canisius (a larger one for teachers, 1554, and a smaller one for pupils, 1566); the Catechism of Cardinal Bellarmin (1603), which Clement VIII. and later Popes commended as an authentic and useful exposition of the Roman Catechism, and which is much used by missionaries; and the Catechism of Bossuet for the diocese of Meaux (1687). The Roman Church allows an endless multiplication of such educational books with adaptations to different nationalities, ages, degrees of culture, local wants and circumstances, provided they agree with the doctrinal system set forth by the Council of Trent. Most of these books, however, must now be remodeled and adjusted to the Council of the Vatican.\footnote{Thus, for instance, in Keenan's \emph{Controversial Catechism}, as published by the 'Catholic Publishing Company,' New Bond Street, London, the pretended doctrine of papal infallibility was expressly denied as 'a Protestant invention; it is no article of the Catholic faith; no decision of the Pope can oblige under pain of heresy, unless it be received and enforced by the teaching body, that is, by the Bishops of the Church.' But since 1871 the leaf containing this question and answer has been canceled and another substituted. So says Oxenham, in his translation of Döllinger on the \emph{Reunion of Churches}, p. 126, note. The same is true of many German and French Catholic Catechisms.}

\xsection{The Papal Bulls Against the Jansenists, A.D. 1653, 1713.}

§ 27. The Papal Bulls against the Jansenists, 1653 and 1713.

Cornelius Jansenius (Episcopi Iprensis, 1585-1638): \emph{Augustinus, seu doctrina Augustini de humanæ naturæ sanitate, ægritudine, et medicina, adv. Pelagianos et Massilienses}, Lovan. 1640, 3 vols.; Paris, 1641; Rouen, 1643 (with a \emph{Synopsis vitæ Jansenii}). Prohibited, together with the Jesuit antitheses, by Pope Urban VIII., 1642.

St. Cyran (Du Vergier, d. 1643): \emph{Aurelius}, 1633: again, Paris, 1646. A companion to Jansen's 'Augustinus', and called after the other name of the great Bishop of Hippo.

Anthony Arnauld (Doctor of the Sorbonne, d. at Brussels, 1694): \emph{Œuvres}, Paris, 1775–81, 49 vols. in 44. Letters, sermons, ascetic treatises, controversial books against Jesuits (Maimbourg, Annat), Protestants (Jurieu, Aubertin), and philosophers (Descartes, Malebranche).

M. Leydecker (Ref. Prof. at Utrecht, d. 1721): \emph{Historia Jansenismi}, Utr. 1695.

Gerberon: \emph{Histoire générale de Jansenisme}, Amst. 1700.

Lucchesini: \emph{Hist. polem. Jansenismi}, Rome, 1711, 3 vols.

Fontaine: \emph{Mémoires pour servir a l’histoire de Port-Royal} (Utrecht), 1738, 2 vols.

\emph{Collectio nova actorum Constit. Unigenitus}, ed. R. J. Dubois, Lugd. 1725.

Dom. de Colonia: \emph{Diction, des livres Jansenistes}, Lyons, 1732, 4 vols.

H. Reuchlin: \emph{Geschichte von Port-Royal}, Hamb. 1839–44, 2 vols. Comp. his monograph on \emph{Pascal}, and his art. \emph{Jansen} and \emph{Jansenismus} in Herzog's \emph{Encyklop.} 2d ed. Vol. VI. pp. 481–493.

C. A. Sainte-Beuve: \emph{Port-Royal}, Paris, 1840–42, 2 vols.

Abbé Guettée: \emph{Jansénisme et Jésuitisme, un examen des accusations de Jans., }etc., Paris, 1857. Compare his \emph{Histoire de l’église de France, composé sur les documents originaux et authentiques}, Paris, 1847–56, 12 vols. Placed on the index of prohibited books, 1852. The author has since passed from the Roman to the Greek Church.

W. Henley Jervis: \emph{The Gallican Church: A History of the Church of France from} 1516 \emph{to the Revolution}, Lond. 1872, 2 vols. On Jansenism, see Vol. I. chaps. xi.–xiv., and Vol. II. chaps. v., vi., and viii.

Frances Martin: \emph{Anglique Arnauld, Abbess of Port-Royal}, London, 1873.

(The controversial literature on Jansenism in the National Library at Paris amounts to more than three thousand volumes.)

On the Jansenists, or Old Catholics, in Holland.

Dupac de Bellegarde: \emph{H. de l’église metropol. d’Utrecht}, Utr. 1784, 3d ed. 1852.

Walch: \emph{Neueste Rel. Geschichte}, Vol. VI. pp. 82 sqq.

Theol. Quartalschrift, Tüb. 1826.

Augusti: \emph{Das Erzbisthum Utrecht}, Bonn, 1838.

S. P. Tregelles: \emph{The Jansenists: their Rise, Persecutions by the Jesuits, and existing Remnant}, London, 1851 (with portraits of Jansenius, St. Cyran, and the Mère Angelique).

J. M. Neale: \emph{A History of the so-called Jansenist Church of Holland}, etc., London, 1857. Neale visited the Old Catholics in Holland in 1851, and predicted for them happier days.

Fr. Nippold: \emph{'Die altkatholische Kirche des Erzbisthums Utrecht. Geschichtl. Parallele zur altkathol. Gemeindebildung in Deutschland}, Heidelberg, 1872.

The remaining doctrinal decrees of the Roman Church relate to internal controversies among different schools of Roman Catholics.

Jansenism, so called after Cornelius Jansenius (or Jansen), Bishop of Ypres, and supported by the genius, learning, and devout piety of some of the noblest minds of France, as St. Cyran, Arnauld, Nicole, Pascal, Tillemont, the Mother Angelique Arnauld, and other nuns of the once celebrated Cistercian convent \emph{Port-Royal des Champs} (a few miles from Versailles), was an earnest attempt at a conservative doctrinal and disciplinary reformation in the Roman Church by reviving the Augustinian views of sin and grace, against the semi-Pelagian doctrines and practices of Jesuitism, and made a near approach to evangelical Protestantism, though remaining sincerely Roman Catholic in its churchly, sacerdotal, and sacramental spirit, and legalistic, ascetic piety. It was most violently opposed and almost totally suppressed by the combined power of Church and State in France, which in return reaped the Revolution. It called forth two Papal condemnations, with which we are here concerned.

I. The bull 'Cum Occasione' of Innocent X. (who personally knew and cared nothing about theology), A.D. 1653. It is purely negative, and condemns the following five propositions from a posthumous work of Jansenius, entitled \emph{Augustinus}.\footnote{The book is called after the great African Church Father, whose doctrines it reproduced, and was published by friends of the author in 1640, two years after his death. On Jansen, comp. the Dutch biography of Heeser:  \emph{Historisch Verhaal van de Geboorte, Leven, etc., van Cornelius Jansenius}, 1727. He was born near Leerdam, in Holland, 1585, studied in Paris, was Professor of Theology in the University of Louvain, Bishop of Ypres 1635, and died 1638. He read Augustine's works against Pelagius thirty times, the other works ten times. His book was finished shortly before his death, and advocates the Augustinian system on total depravity, the loss of free-will, irresistible grace, and predestination. In his will he submitted it to the Holy See. He resembles somewhat his countryman, Pope Adrian VI., who vainly endeavored to reform the Papacy.}

(1.) The fulfillment of some precepts of God is impossible even to just men according to their present ability (\emph{secundum præsentes quas habent vires}), and the grace is also wanting to them by which they could be observed (\emph{deest illis gratia, qua possibilia fiant}).

(2.) Interior grace is never resisted in the state of fallen nature.

(3.) For merit or demerit in the state of fallen nature man need not be exempt from all necessity, but only from coercion or constraint (\emph{Ad merendum et demerendum in statu naturæ lapsæ, non requiritur in homine libertas a necessitate, sed sufficit libertas a coactione}—that is, from violence and natural necessity).

(4.) The Semi-Pelagians admitted the necessity of prevenient interior grace for every action, even for the beginning of faith; but they were heretical (\emph{in eo erant hæretici}) in believing this grace to be such as could be resisted, or obeyed by the human will (\emph{eam gratiam talem esse, cui posset humana voluntas resistere, vel obtemperare}).

(5.) It is semi-Pelagian to say that Christ died and shed his blood wholly (altogether) for all men.\footnote{'\emph{Semipelagianum est dicere, Christum pro omnibus omnino mortuum esse aut sanguinem fudisse.}' This supralapsarian proposition is condemned as \emph{falsa, temeraria, scandalosa, impia, blasphema, et hæretica.} See the five propositions of Jansen in Denzinger's \emph{Enchir}., pp. 316, 317.}

The Jansenists maintained that these propositions were not taught by Jansenius, at least not in the sense in which they were condemned; that this was a historical question of fact (\emph{question de fait}), not a dogmatic question of right (\emph{droit}); and, while conceding to the Pope the right to condemn heretical propositions, they denied his infallibility in deciding a question of fact, about which he might be misinformed, ignorant, prejudiced, or taken by surprise.

But Pope Alexander VII., in a bull of 1665, commanded all the Jansenists to subscribe a formula of submission to the bull of Innocent X., with the declaration that the five propositions were taught in the book of Cornelius Jansen in the sense in which they were condemned by the previous Pope.\footnote{'\emph{Ego N. constitutioni apostolicæ Innocentii X., datæ die} 31. \emph{Maji} 1653, \emph{et constitutioni Alexandri VII., datæ die} 16. \emph{Octobris} 1665, \emph{summorum Pontificum, me subjicio, et quinque propositiones ex Cornelii Jansenii libro, cui nomen Augustinus, excerptas, et in sensu ab eodem auctore intento, prout illas per dictas constitutiones Sedes Apostalica damnavit, sincero animo rejicio ac damno, et ita juro. Sic me Deus adjuvet, et hæc sancta Dei evangelia.}'}

The Jansenists, including the nuns of Port-Royal, refused to submit. Many fled to the Netherlands. The Pope abolished their famous convent (1709), the building was destroyed by order of Louis XIV. (1710), even the corpses of the illustrious Tillemonts, Arnaulds, Nicoles, De Sacys, and others, were disinterred with gross brutality (1711), and the church itself was demolished (1713). No wonder that such barbarous tyranny and cruelty, perpetrated in the holy name of the Church of Christ, bred a generation of skeptics and infidels, who at last banished the Church and religion itself from the territory of France. Cardinal Noailles, who from weakness had lent his high authority to these outrages, made afterwards, in bitter repentance, a pilgrimage to the ruins of Port-Royal, and, looking over the desecrated burial-ground, he exclaimed: 'Oh! all these dismantled stones will rise up against me at the day of judgment! Oh! how shall I ever bear the vast, the heavy load!'\footnote{Gregoire: \emph{Les ruines de Port-Royal}, Par. 1709. \emph{Mémoires sur la déstruction de P. R. des Champs}, 1711. Jervis, l.c. Vol. II. pp.191 sqq. Tregelles says, l.c. p. 47: 'The united acts of Louis XIV. and the Jesuits, in crushing alike Protestants, Quietists, and Jansenists, drove religion well-nigh out of France. What a spectacle! The same monarch, under the influence of the same evil-minded and pharisaical woman (Madame de Maintenon), persecuting not only Protestants, but also such men as Fénelon, among the brightest and holiest of those who owned the authority of Rome. Thus was the train laid which led to the fearful explosion in which altar and throne alike fell, and \emph{atheism} was nationally embraced. How the mind of Voltaire was affected by the abominable deeds of men who \emph{professed} the name of Christ, is shown by his juvenile verses, in which he speaks so indignantly of the destruction of Port-Royal that he was sent for a year to the Bastile.'}

II. The more important bull 'Unigenitus (Dei Filius)', issued by Pope Clement XI., Sept., 1713, condemns one hundred and one sentences of the Jansenist Pasquier Quesnel, (d. 1719), extracted from his moral reflections on the New Testament.\footnote{Pasquier or Paschasius Quesnel was born at Paris, 1634, studied at the Sorbonne, joined the Congregation of the Oratory, and was appointed director of the institution belonging to this order at Paris. He was a profound and devout student of the Scriptures and the Fathers, edited the works of Leo I. (1675, with dissertations) in defense of the Gallican Church against the Ultramontane Papacy (hence the edition was condemned by the Congregation of the Index), was exiled from France 1684, joined Arnauld at Brussels, and died at Amsterdam 1719. After the death of Arnauld he was considered the head of the Jansenists. His commentary is one of the most spiritual and reverent. It is entitled '\emph{Le Nouv. Testament en françois avec des réflexions morales sur chaque vers, et pour en rendre la lecture plus utile, et la méditation plus aisée},' Paris, 1687, 2 vols.; 1694; Amsterd. 1736, 8 vols.; also in Latin and other languages; Engl. ed. London, 1819–25, 4 vols. The Gospels were repeatedly published, with an introductory essay by Bishop Daniel Wilson, London and New York. Comp. \emph{Causa Quesnelliana}, Brussels, 1704.}

This bull is likewise negative, but commits the Church of Rome still more strongly than the former against evangelical doctrines. Several of the passages selected are found almost literally in Augustine and St. Paul; they assert the total depravity of human nature, the loss of liberty, the renewing power of the free grace of God in Christ, the right and duty of all Christians to read the Bible.

The following are the most important of these propositions:\footnote{Denzinger's \emph{Enchir}., pp. 351–361.}

(2.) \emph{Jesu Christi gratia, principium efficax boni cujuscunque generis, necessaria est ad omne opus bonum; absque illa non solum nihil fit, sed nec fieri potest.}

(3.) \emph{In vanum, Domine, præcipis, si tu ipse non das, quod præcipis.} (Compare the similar sentence of Augustine, which was so offensive to Pelagius: \emph{Da quod jubes, et jube quod vis.})

(4.) \emph{Ita, Domine; omnia possibilia sunt ei, cui omnia possibilia facis, eadem operando in illo.}

(10.) \emph{Gratia est operatio manus omnipotentis Dei, quam nihil impedire potest aut retardare.}

(11.) \emph{Gratia non est aliud quam voluntas omnipotentis Dei jubentis et facientis, quod jubet.}

(13.) \emph{Quando Deus vult animam salvam facere, et eam tangit interiori gratiæ suæ manu, nulla voluntas humana ei resistit.}

(18.) \emph{Semen verbi, quod manus Dei irrigat, semper affert fructum suum.}

(21.) \emph{Gratia Jesu Christi est gratia fortis, potens, suprema, invincibilis, utpote quæ est operatio voluntatis omnipotentis, sequela et imitatio operationis Dei incarnantis et resuscitantis Filium suum.}

(27.) \emph{Fides est prima gratia et fons omnium aliarum.} (2 Pet. 1. 3.)

(28.) \emph{Prima gratia, quam Deus concedit peccatori, est peccatorum remissio.}

(29.) \emph{Extra ecclesiam nulla conceditur gratia.}\footnote{The denial of this proposition implies the assertion that there is grace outside of the Church, though not sufficient for salvation; else it would be inconsistent with the Roman Catholic doctrine '\emph{Extra ecclesiam nulla salus.}'}

(30.) \emph{Omnes, quos Deus vult salvare per Christum, salvantur infallibiliter.}

(38.) \emph{Peccator non est liber, nisi ad malum, sine gratia Liberatoris.}

(39.) \emph{Voluntas, quam gratia non prævenit, nihil habet luminis, nisi ad aberrandum, ardoris, nisi ad se præcipitandum, virium nisi ad se vulnerandum; est capax omnis mali et incapax ad omne bonum.}

(40.) \emph{Sine gratia nihil amare possumus, nisi ad nostram condemnationem.}

(58.) \emph{Nec Deus est nec religio, ubi non est charitas.} (1 John iv. 8.)

(59.) \emph{Oratio impiorum est novum peccatum; et quod Deus illis concedit, est novum in eos judicium.}

(69.) \emph{Fides, usus, augmentum et præmium fidei, totum est donum puræ liberalitatis Dei.}

(72.) \emph{Nota ecclesiæ Christianæ est, quod sit catholica, comprehendens et omnes angelos cœli, et omnes electos et justos terræ et omnium sæculorum.}

(75.) \emph{Ecclesia est unus solus homo compositus ex pluribus membris, quorum Christus est caput, vita, subsistentia et persona; unus solus Christus compositus ex pluribus sanctis, quorum est Sanctificator.}

(76.) \emph{Nihil spatiosius Ecclesia Dei; quia omnes electi et justi omnium seculorum illam componunt} (Eph. ii. 22).

(77.) \emph{Qui non ducit vitam dignam filio Dei et membro Christi, cessat interius habere Deum pro Patre et Christum pro capite.}

(79.) \emph{Utile et necessarum est omni tempore, omni loco, et omni personarum generi, studere el cognoscere spiritum, pietatem et mysteria sacræ Scripturæ.}

(80.) \emph{Lectio sacræ Scripturæ est pro omnibus.} (John v. 39; Acts xvii. 11.)

(81.) \emph{Obscuritas sancti verbi Dei non est laicis ratio dispensandi se ipsos ab ejus lectione.}

(82.) \emph{Dies Dominicus a Christianis debet sanctificari lectionibus pietatis et super omnia sanctarum Scripturarum. Damnosum est, velle Christianum ab hac lectione retrahere.}

(84.) \emph{Abripere e Christianorum manibus novum Testamentum seu eis illud clausum tenere auferendo eis modum istud intelligendi, est illis Christi os obturare.}

(85.) \emph{Interdicere Christianis lectionem sacræ Scripturæ, præsertim Evangelii, est interdicere usum luminis filiis lucis et facere, ut patiantur speciem quamdam excommunicationis.}

(92.) \emph{Pati potius in pace excommunicationem et anathema injustum, quam prodere veritatem, est imitari sanctum Paulum; tantum abest, ut sit erigere se contra auctoritatem aut scindere unitatem.}

(100.) \emph{Tempus deplorabile, quo creditur honorari Deus persequendo veritatem ejusque discipulos! . . . Frequenter credimus sacrificare Deo impium, et sacrificamus diabolo Dei servum.}

These and similar propositions, some of them one-sided and exaggerated, many of them clearly patristic and biblical, are indiscriminately condemned by the bull \emph{Unigenitus}, as 'false, captious, ill-sounding, offensive to pious ears, scandalous, rash, injurious, seditious, impious, blasphemous, suspected of heresy and savoring of heresy itself, near akin to heresy, several times condemned, and manifestly renewing various heresies, particularly those which are contained in the infamous propositions of Jansenius!'

A large portion of the French clergy, headed by the Archbishop of Paris, Cardinal de Noailles, who repented of his part in the destruction of Port-Royal, protested against the bull, and appealed from the Pope to a future council. But 'when Rome has spoken, the cause is finished.' The bull \emph{Unigenitus} was repeatedly confirmed by the same Clement XI., A.D. 1718 (in the bull '\emph{Pastoralis Officii}'), Innocent XIII., 1722, Benedict XIII. and a Roman Synod, 1725, Benedict XIV., 1756; it was accepted by the Gallican clergy 1730, and, as Denzinger says, by 'the whole Catholic world' ('\emph{ab universo mundo catholico}'). Even the miracles on the grave of a Jansenist saint (Franois Paris, who died 1727, after the severest self-denial, with a protest against the bull \emph{Unigenitus} in his hand), could not save Jansenism from destruction in France.\footnote{The Jesuits, of course, ascribed the Jansenist miracles, visions, and ecstatic convulsions to the devil.}

But a remnant fled to the more liberal soil of Protestant Holland, and was there preserved as a perpetual testimony against Jesuitism, and, as it now seems, for an important mission in connection with the Old Catholic protest against the decisions of the Vatican Council.

Note on the Jansenists in Holland.—The remnant of the Jansenists or the Old Catholics in Holland date their separate existence from the protest against the bull \emph{Unigenitus}, but are properly the descendants of the original Catholics. They disown the name 'Jansenists,' on the ground of alleged error in the papal bulls concerning the true teaching of Jansen, and call themselves the 'Old Episcopal Clergy of the Netherlands;' but they are strongly opposed to the theology and casuistry of the Jesuits, and incline to the Augustinian views of sin and grace. In other respects they are good Catholics in doctrine, worship, and mode of piety; they acknowledge the decrees and canons of Trent, and even the supremacy of the Pope within the limits of the old Gallican theory. They inform him of the election of every new bishop, which the Pope as regularly declares illegitimate, null, and void. They say that the tyranny of a father does not absolve his children from the duty of obedience, and hope against hope that God will convert the Pope, and turn his heart towards them. They number at present one archbishopric of Utrecht and two bishoprics of Deventer and Haarlem, 25 congregations, and about 6000 members. They live very quietly, surrounded by Romanists and Protestants, and are much respected, like the Moravians, for their character and piety. The Pope, after condemning them over and over again, appointed, in 1853, five new bishoprics in Holland, with a rival archbishop at Utrecht, and thus consolidated and perpetuated the schism. When the decree of the Immaculate Conception was promulgated in 1854, the three Old Catholic Bishops issued a pastoral letter, in which they reject the new dogma as contrary to the Scriptures and early tradition, and as lacking the threefold test of catholicity (\emph{semper, ubique, ab omnibus}). The Vatican decree of Papal Infallibility, and the Old Catholic movement in Germany have brought this long afflicted and persecuted remnant of Jansenism into new notice. The Old Catholics of Germany, holding fast to an unbroken episcopal succession, looked to their brethren in Holland for aid in effecting an organization when it should become necessary. At their invitation, Archbishop Loos, of Utrecht (a venerable and amiable old gentleman), made a tour of visitation in the summer of 1872, and confirmed about five hundred children in several congregations in Germany, blessing God that his little Church was spared for happier days. After his death the Bishop of Deventer consecrated Prof. Reinkens Bishop for the Old Catholics in Germany, Aug. 11, 1873. The Old Catholics of Holland agree with those in Germany: 1. In maintaining the doctrinal basis of Tridentine Romanism; 2. In protesting against all subsequent papal decisions, more particularly the bull \emph{Unigenitus}, the decree of the Immaculate Conception (1854), and the Vatican decree of Papal Infallibility. [The Jansenist Abp. of Utrecht was excommunicated by Leo XIII., Feb. 28, 1893. See Mirbt, p. 488, and also the Old Catholic bishops of Germany and Switzerland.—Ed.]

\xsection{The Papal Definition of the Immaculate Conception of the Virgin Mary, A.D. 1854.}

§ 28. The Papal Definition of the Immaculate Conception of the Virgin Mary, 1854.

Literature.

I. In favor of the Immaculate Conception of Mary:

The papal bull of Pius IX., '\emph{Ineffabilis Deus},' Dec. 8 (10), 1854.

John Perrone (Professor of the Jesuit College in Rome, and one of the chief advisers of Pius IX. in framing his decree): \emph{Can the Immaculate Conception of the Blessed Virgin Mary be defined by a Dogmatic Decree?} In Latin, Rome, 1847, dedicated to Pius IX., with a letter of thanks by the Pope; German translation, by \emph{Dietl} and \emph{ Schels}, Regensburg, 1849. (I used the German edition.) See also Perrone's \emph{ Prælectiones theologicæ}, Append. to Tom. VI., ed. Ratisb. 1854.

C. Passaglia: \emph{De immaculato Deiparæ semper virginis conceptu}, Rom. 1854 sqq., Tom. III. 4to. (The author has since become half heretical, at least as regards the temporal power of the Pope, and was obliged to flee from Rome. See his pamphlet on the subject, 1861, which was placed on the \emph{Index}.)

H. Denzinger (d. 1862): \emph{Die Lehre von der unbefleckten Empfängniss der seligsten Jungfrau}, Würzb. 1868.

Aug. de Roskovány (Episc. Nitriensis): \emph{Beata Virgo Maria in suo conceptu immaculata ex monumentis omnium seculorum demonstrata}, Budapest, 1874, 6 vols.

II. Against the Immaculate Conception:

Juan de Turrecremata: \emph{ Tractatus de veritate conceptionis beatissimæ virginis}, etc., Rome, 1547, 4to; newly edited by Dr. E. B. Pusey, with a preface and notes, London, 1869. Card. Joh. de Turrecremata, or Torquemada (not to be confounded with the Great Inquisitor Thomas de T.), attended as \emph{magister sacri palatii} the General Councils of Basle and Ferrara, and, although a faithful champion of Popery, he opposed, as a Dominican, the Immaculate Conception. He died, 1468, at Rome.

J. de Launoy (or Launoius, a learned Jansenist and Doctor of the Sorbonne, d. 1678): \emph{Præscriptiones de Conceptu B. Mariæ Virginis}, 2d ed. 1677; also in the first volume of his \emph{ Opera omnia}, Colonii Allobrogum, fol. 1731, pp. 9–43, in French and Latin.

G. E. Steitz: Art. \emph{Maria, Mutter des Herrn}, in Herzog's \emph{Encyklop.} Vol. IX. pp. 94 sqq.

E. Preuss: \emph{Die römische Lehre von der unbefleckten Empfägniss. Aus den Quellen dargestellt und aus Gottes Wort widerlegt.} Berlin, 1865. The same, translated into English by \emph{Geo. Gladstone}, Edinburgh, 1867. The author has since become a Romanist, and recalled his book, Dec. 1871.

H. B. Smith (Professor in the Union Theological Seminary, N.Y.): \emph{The Dogma of the Immaculate Conception}, in the \emph{Methodist Quarterly Review}, New York, for 1855, pp. 275–311.

Dr.  Pusey: \emph{Eirenikon}, Part II., Lond. 1867.

Art. In \emph{Christian Remembrancer} for Oct. 1855; Jan. 1866; July, 1868.

K. Hase: \emph{Handbuch der Protest. Polemik gegen die röm. kath. Kirche}, 3d ed. Leipz. 1871, pp. 334–344.

The first step towards the proclamation of the dogma of the Immaculate Conception of the Virgin Mary, which exempts her from all contact with sin and guilt, was taken by Pope Pius IX., himself a most devout worshiper of Mary, during his temporary exile at Gaäta. In an encyclical letter, dated Feb. 2, 1849, he invited the opinion of the Bishops on the alleged ardent desire of the Catholic world that the Apostolic See should, by some solemn judgment, define the Immaculate Conception, and thus secure signal blessings to the Church in these evil times. For, he added, 'You know full well, venerable brethren, that the whole ground of our confidence is placed in the most holy Virgin,' since 'God has vested in her the plenitude of all good, so that henceforth, if there be in us any hope, if there be any grace, if there be any salvation (\emph{si quid spei in nobis est, si quid gratiæ, si quid salutis}), we must receive it solely from her, according to the will of him who would have us possess all through Mary.'

More than six hundred Bishops answered, all of them, with the exception of four, assenting to the Pope's belief, but fifty-two, among them distinguished German and French Bishops, dissenting from the expediency or opportuneness of the proposed dogmatic definition. The Archbishop of Paris (Sibour) apprehended injury to the Catholic faith from the unnecessary definition of the Immaculate Conception, which 'could be proved neither from the Scriptures nor from tradition, and to which reason and science raised insolvable, or at least inextricable, difficulties.' But this opposition was drowned in the general current.\footnote{Perrone says: \emph{Vix quatuor responderunt negative quoad definitionem, et ex hic ipsis tres brevi mutarunt sententiam.} These letters, with others from sovereigns, monastic orders, and Catholic societies, are printed in nine volumes.}

After the preliminary labors of a special commission of Cardinals and theologians, and a consistory of consultation, Pope Pius, in virtue of the authority of Christ and the holy Apostles Peter and Paul, and his own authority, solemnly proclaimed the dogma on the Feast of the Conception, Dec. 8, 1854, in the Church of St. Peter, in the presence of over two hundred Cardinals, Bishops, and other dignitaries, invited by him, not to discuss the doctrine, but simply to give additional solemnity to the ceremony of proclamation. After the mass and the singing of the \emph{Veni Creator Spiritus}, he read with a tremulous voice the concluding formula of the bull '\emph{Ineffabilis Deus},' declaring it to be a divinely revealed fact and dogma, which must be firmly and constantly believed by all the faithful on pain of excommunication, '\emph{that the most blessed Virgin Mary, in the first moment of her conception, by a special grace and privilege of Almighty God, in virtue of the merits of Christ, was preserved immaculate from all stain of original sin}.'\footnote{\emph{Postquam numquam intermisimus in humilitate et jejunio privatas nostras et publicas Ecclesiæ preces Deo Patri per Filium ejus offerre, ut Spiritus Sancti virtute mentem nostram dirigere et confirmare dignaretur, implorato universæ cœlestis curiæ præsidio, et advocato cum genitibus Paraclito Spiritu, eoque sic aspirante, ad honorem Sanctæ et Individuæ Trinitatis, ad decus et ornamentum Virginis Deiparæ, ad exaltationem fidei catholicæ et christianæ religionis augmentum, auctoritate Domini nostri Jesu Christi, beatorum Apostolorum Petri et Pauli, ac nostra declaramus, pronuntiamus et definimus, doctrinam, quæ tenet}, beatissimam Virginem Mariam in primo instanti suæ conceptionis fuisse singulari omnipotentis Dei gratia et privilegio, intuitu meritorium Christi Jesu Salvatoris humani generis, ab omni originalis culpæ labe preservatam immunem,  \emph{esse a Deo revelatam atque idcirco ab omnibus fidelibus firmiter constanterque credendam. Quapropter si qui secus ac a Nobis definitum est, quod Deus avertat, præsumpserint corde sentire, ii noverint ac porro sciant, se proprio judicio condemnatos, naufragium circa fidem passos esse, et ab unitale Ecclesiæ defeciise, ac præterca facto ipso suo semet poenis a jure statutis subjicere, si, quod corde, sentiunt, verbo aut scripto, vel alio quovis externo modo significare ausi fuerint.}'}

The shouts of the assembled multitude, the cannons of St. Angelo, the chime of all the bells, the illumination of St. Peter's dome, the splendor of gorgeous feasts, responded to the decree. Rome was intoxicated with idolatrous enthusiasm, and the whole Roman Catholic world thrilled with joy over the crowning glory of the immaculate queen of heaven, who would now be more gracious and powerful in her intercession than ever, and shower the richest blessings upon the Pope and his Church. To perpetuate the memory of the occasion, the Pope caused a bronze tablet to be placed in the wall of the choir of St. Peter's, with the inscription that, on the 8th of December, 1854, he proclaimed the dogma of the Immaculate Conception of the \emph{Deipara Virgo Maria,} and thereby fulfilled the desire of the whole Catholic world (\emph{totius orbis catholici desideria}), and a pompous marble statue of the Virgin to be erected on the Piazza di Spagnia, facing the palace of the Propaganda, and representing the Virgin in the attitude of blessing, with Moses, David, Isaiah, and Ezekiel, as the prophetic witnesses of her conception, at the foot of the column.\footnote{The statue of the Virgin is said to have come out of the Roman fabric with a hideous crack, which was clumsily patched up. See Hase, \emph{Protest. Polemik}, 3d ed. p. 341, and Preuss, l.c. p. 197 (English edition).} He ordered, also, through the Congregation of Rites, the preparation of a new mass and a new office for the festival of the Conception, which was published Sept. 25, 1863, and contains the prayer: 'O God, who, by the immaculate conception of the Virgin, didst prepare a worthy dwelling for thy Son: grant, we beseech thee, that, as thou didst preserve her from every stain, in anticipation of the death of thy Son, so we also may, through her intercession, appear purified before thy presence.'

The dogma lacks the sanction of an œcumenical Council, and rests solely on the authority of the Pope, who, in its proclamation, virtually anticipated his own infallibility; but it has been generally accepted by subsequent assent, and must be considered as an essential and undoubted part of the Roman faith, especially since the Vatican Council has declared the official infallibility of the Pope.

This extraordinary dogma lifts the Virgin Mary out of the fallen and redeemed race of Adam, and places her on a par with the Saviour. For if she is really free from all hereditary as well as actual sin and guilt, she is above the need of redemption. Repentance, forgiveness, regeneration, conversion, sanctification are as inapplicable to her as to Christ himself. The definition of such a dogma implies nothing less than a Divine revelation; for only the omniscient God can know the fact of the immaculate conception, and only he can reveal it. He did not reveal it to the inspired Apostles, nor to the Fathers. Did he reveal it to Pope Pius IX., in 1854, more than eighteen centuries after it took place?

Viewed from the Roman point of view, the new dogma is the legitimate fruit of the genuine spirit of modern Romanism. It only completes that Mariology, and fortifies that Mariolatry, which is the very soul of its piety and public worship. We may almost call Romanism the Church of the Virgin Mary—not of the real Virgin of the Gospels, who sits humbly and meekly at the feet of her and our Lord and Saviour in heaven, but of the apocryphal Virgin of the imagination, which assigns her a throne high above angels and saints. This mythical Mary is the popular expression of the Romish idea of the Church, and absorbs all the reverence and affection of the heart. Her worship overshadows even the worship of Christ. His perfect humanity, by which he comes much nearer to us than his earthly mother, is almost forgotten. She, the lovely, gentle, compassionate woman, stands in front; her Son, over whom she is supposed still to exercise the rights of her divine maternity, is either the stern Lord behind the clouds, or rests as a smiling infant on her supporting arms. By her powerful intercession she is the fountain of all grace. She is virtually put in the place of the Holy Spirit, and made the mediatrix between Christ and the believer. She is most frequently approached in prayer, and the 'Ave Maria' is to the Catholic what the Lord's Prayer is to the Protestant. If she hears all the petitions which from day to day, and from hour to hour, rise up to her from many millions in every part of the globe, she must, to all intents and purposes, be omnipresent and omniscient. She is the favorite subject of Roman painters, who represent her as blending in harmony the spotless beauty of the Virgin and the tender care of the mother, and as the crowned queen of heaven. Every event of her life, known or unknown, even her alleged bodily assumption to heaven, is celebrated with special zeal by a public festival.\footnote{Why should the fiction of the \emph{Assumption} of Mary to heaven (as it is called in distinction from the \emph{Ascension} of Christ) not be proclaimed a divinely revealed fact and a binding dogma, as well as the Immaculate Conception? The evidence is about the same. If Mary was free from all contact with sin, she can not have been subject to death and corruption, which are the wages of sin. The silence of the Bible concerning her end might be turned to good account. Tradition, also, can be produced in favor of the assumption. St. Jerome was inclined to believe it, and even the great Augustine 'feared to say that the blessed body, in which Christ had been incarnate, could become food for the worms.' The festival of the Assumption, which presupposes the popular superstition, is older than the festival of the Immaculate Conception, and is traced by some to the fifth or sixth century.} It is almost incredible to what extent Romish books of devotion exalt the Virgin. In the Middle Ages the whole Psalter was rewritten and made to sing her praises, as 'The heavens declare thy glory, O Mary;' 'Offer unto our lady, ye sons of God, praise and reverence!' In St. Liguori's much admired and commended '\emph{Glories of Mary},' she is called 'our life,' the 'hope of sinners,' 'an advocate mighty to save all,' a 'peacemaker between sinners and God.' There is scarcely an epithet of Christ which is not applied to her. According to Pope Pius IX., 'Mary has crushed the head of the serpent,' i.e., destroyed the power of Satan, 'with her immaculate foot!' Around her name clusters a multitude of pious and blasphemous legends, superstitions, and impostures of wonder-working pictures, eye-rotations, and other unnatural marvels; even the cottage in which she lived was transported by angels through the air, across land and sea, from Nazareth in Galilee to Loretto in Italy; and such a silly legend was soberly and learnedly defended even in our days by a Roman Archbishop.\footnote{Dr. Kenrick, of St. Louis, in his work on the '\emph{Holy House},' a book which is said to be too little known. See Smith, l.c. p. 279.}

Romanism stands and falls with Mariolatry and Papal Infallibility; while Protestantism stands and falls with the worship of Christ as the only Mediator between God and man, and the all-sufficient Advocate with the Father.

\xsection{The Argument for the Immaculate Conception.}

§ 29. The Argument for the Immaculate Conception.

The importance of the subject justifies and demands a brief examination of the arguments in favor of this novel dogma, which is one of the most characteristic features of modern Romanism, and forms an impassable gulf between it and Protestantism. It is a striking proof of Romish departure from the truth, and of the anti-Christian presumption of the Pope, who declared it to be a primitive divine revelation; while it is in fact a superstitious fiction of the dark ages, contrary alike to the Scriptures and to genuine Catholic tradition.

1. The dogma of the sinlessness of the Virgin Mary is \emph{unscriptural}, and even \emph{anti-scriptural.}

(\emph{a}) The Scripture passages which Perrone and other champions of the Immaculate Conception adduce are, with one exception, all taken from the Old Testament, and based either on false renderings of the Latin Bible, or on fanciful allegorical interpretation.

(1) The main (and, according to Perrone, the only) support is derived from the \emph{protevangelium}, Gen. iii. 15, where Jehovah Elohim says to the serpent, according to the Latin Bible (which the Romish Church has raised to an equality with the original): '\emph{Inimicitias ponam inter te et mulierem, et semen tuum et semen illius}; Ipsa  \emph{conteret caput tuum, et tu insidiaberis calcaneo ejus}' (i.e., \emph{she} shall crush thy head, and thou shalt assail \emph{her} heel). Here the \emph{ipsa} is referred to the woman (\emph{mulier}), and understood of the Virgin Mary.\footnote{Pope Pius IX. has given his infallible sanction to this misapplication of the \emph{protevangelium} to Mary in the gallant phrase already quoted (p. 112) from his Encyclical on the dogma.} And it is inferred that the divinely constituted enmity between Mary and Satan must be unconditional and eternal, which would not be the case if she had ever been subject to hereditary sin.\footnote{Speil, in his defense of Romanism against Hase, argues in this way: The woman, whom God will put in enmity against the devil, must be a future particular woman, over whom the devil never had any power—that is, a woman who, by the grace of God, was free from original sin (\emph{Die Lehren der katholischen Kirche}, 1865, p. 165).} To this corresponds the Romish exegesis of the fight of the woman (i.e., the Church) with the dragon, Rev. xii. 4 sqq.; the woman being falsely understood to mean Mary. Hence Romish art often represents her as crushing the head of the dragon.

But the translation of the Vulgate, on which all this reasoning is based, is contrary to the original Hebrew, which uses the masculine form of the verb, he (or it, the \emph{seed} of the woman), i.e., Christ, \emph{shall bruise}, or \emph{crush}, the serpent's \emph{head}, i.e., destroy the devil's power; it is inconsistent with the last clause, '\emph{and thou shalt bruise}  his (i.e., Christ's) \emph{heel},' which contains a mysterious allusion to the crucifixion of the \emph{seed}, not of the woman; and, finally, the Romish interpretation leads to the blasphemous conclusion that Mary, and not Christ, has destroyed the power of Satan, and saved the human race.\footnote{The Hebrew text admits of no doubt; for the verb \texthebrew{יְשׁוּפְ}, in the disputed clause, is \emph{masculine} (he \emph{shall bruise}, or \emph{crush}), and \texthebrew{הוּא }naturally refers to the preceding \texthebrew{זַרְעָהּ }(\emph{her}  seed), i.e., \texthebrew{זֶרַצ אִשָּׁה }(\emph{the woman's}  seed), and not to the more remote  \texthebrew{אִשָּׁה }(\emph{woman}). In the Pentateuch the personal pronoun \texthebrew{הוּא }(\emph{he}) is indeed \emph{generis communis, }and stands also for the feminine \texthebrew{הִיא }(\emph{she}), which (according to the Masora on Gen. xxxviii. 25) is found but eleven times in the Pentateuch; but in all these cases the masoretic punctuators wrote \texthebrew{הִוא}, to signify that it ought to be read \texthebrew{הִיא }(\emph{she}). The Peshito, the Septuagint (\textgreek{αὐτός σοι τηρήσει κεφαλήν}), and other ancient versions, are all right. Even some MSS. of the Vulgate read \emph{ipse }for \emph{ipsa, }and Jerome himself, the author of the Vulgate, in his '\emph{Hebrew Questions},' and Pope Leo I., condemn the translation \emph{ipsa}. But the blunder was favored by other Fathers (Ambrose, Augustine, Gregory I.), who knew no Hebrew, and by the monastic asceticism and fanciful chivalric Mariolatry of the Middle Ages. To the same influence must be traced the arbitrary change of the Vulgate in the rendering of \texthebrew{שׁוּף }from \emph{conteret} (\emph{shall bruise}) into \emph{insidiaberis} (\emph{shall lie in wait, assail, pursue}), so as to exempt the Virgin from the least injury.}

(2) An unwarranted reference of some poetic descriptions of the fair and spotless bride, in the Song of Solomon, to Mary, instead of the people of Jehovah or the Christian Church, Cant. iv. 7, according to the Vulgate: '\emph{Tota pulchra es, amica mea, et macula non est in te.}' In any case, this is only a description of the present character.

(3) An arbitrary allegorical interpretation of the 'garden inclosed, and fountain sealed,' spoken of the spouse, Cant. iv. 12 (Vulg.: '\emph{hortus conclusus, fons signatus}'), and the closed gate in the east of the temple in the vision of Ezekiel, xliv. 1-3, of which it is said: 'It shall not be opened, and no man shall enter in by it; because Jehovah, the God of Israel, hath entered in by it, therefore it shall be shut. It is for the prince; the prince he shall sit in it, to eat bread before the Lord.' This is a favorite support of the doctrine of the perpetual virginity. Ambrose of Milan (d. 397) was perhaps the first who found here a type of the closed womb of the Virgin, by which Christ entered into the world, and who added to the miracle of a conception \emph{sine viro} the miracle of a birth \emph{clauso utero.}\footnote{Epist. 42 \emph{ad Siricium; De inst. Virg.}, c. 8, and in his hymn \emph{A solis ortus cardine}. The earlier Fathers thought differently on the subject. Tertullian calls Mary 'a virgin as to a man, but not a virgin as to birth' (\emph{non virgo, quantum a partu}); and Epiphanius speaks of Christ as 'opening the mother's womb' (\textgreek{ἀνοίγων μήτραν μητρός}). See my \emph{History of the Christian Church}, Vol. II. p. 417.} Jerome and other Fathers followed, and drew a parallel between the closed womb of the Virgin, from which Christ was born to earthly life, and the sealed tomb from which he arose to heavenly life. But none of the Fathers thought of making this prophecy prove the Immaculate Conception. Such exposition, or imposition rather, is an insult to the Bible, as well as to every principle of hermeneutics.

(4) Sap. i. 4: 'Into a malicious soul wisdom shall not enter; nor dwell in the body that is subject unto sin.' This passage (quoted by Speil and others), besides being from an apocryphal book, has nothing to do with Mary.

(5) Luke i. 28: the angelic greeting, 'Hail (Mary), \emph{full of grace} (\emph{gratia plena}),' according to the Romish versions, says nothing of the origin of Mary, but refers only to her condition at the time of the incarnation, and is besides a mistranslation (see below).

(\emph{b}) All this frivolous allegorical trifling with the Word of God is conclusively set aside by the positive and uniform Scripture doctrine of the universal sinfulness and universal need of redemption, with the single exception of our blessed Saviour, who was conceived by the Holy Ghost without the agency of a human father. It is almost useless to refer to single passages, such as Rom. iii. 10, 23; v. 12, 18; 1 Cor. xv. 22; 2 Cor. v. 14, 15; Gal. iii. 22; Eph. ii. 3; 1 Tim. iv. 10; Psa. li. 5. The doctrine runs through the whole Bible, and underlies the entire scheme of redemption. St. Paul emphasizes the \emph{actual} universality of the curse of Adam, in order to show the \emph{virtual} universality of the salvation of Christ (Rom. v. 12 sqq.; 1 Cor. xv. 22); and to insert an exception in favor of Mary would break the force of the argument, and limit the extent of the atonement as well. Perrone admits the force of these passages, but tries to escape it by saying that, if strictly understood, they would call in question even the immaculate birth of Mary, and her freedom from actual sin as well, which is contrary to the Catholic faith;\footnote{L.c. p. 276. In the same manner he disposes of the innumerable patristic passages which assert the universal sinfulness of men, and make Christ the only exception.} hence the Council of Trent has deprived these passages of all force (\emph{omnem vim ademit}) of application to the blessed Virgin! This is putting tradition above and against the Word of the holy and omniscient God, and amounts to a concession that the dogma is extra-scriptural and anti-scriptural. Unfortunately for Rome, Mary herself has made the application; for she calls God \emph{her Saviour} (Luke i. 47: \textgreek{ἐπὶ τῷ θεῷ τῷ σωτῆρί μου}), and thereby includes herself in the number of the redeemed. With this corresponds also the proper meaning of the predicate applied to her by the angel, Luke i. 28, \textgreek{κεχαριτωμένη,} \emph{highly favored, endued with grace} (\emph{die begnadigte}), the one who received, and therefore needed, grace (\emph{non ut mater gratiæ, sed ut filia gratiæ,} as Bengel well observes); comp. ver. 30, \textgreek{εὗρες χάριν παρὰ τῷ θεῷ,} \emph{thou hast found grace with God}; and Eph. i. 6, \textgreek{ἐχαρίτωσεν ἡμᾶς, } \emph{he bestowed grace upon us}. But the Vulgate changed the passive meaning into the active: \emph{gratia plena,} \emph{full of grace,} and thus furnished a spurious argument for an error.

Nothing can be more truthful, chaste, delicate, and in keeping with womanly humility and modesty than both the words and the silence of the canonical Gospels concerning the blessed among women, whom yet our Lord himself, in prophetic foresight and warning against future Mariolatry, placed on a level with other disciples; emphatically asserting that there is a still higher blessedness of spiritual kinship than that of carnal consanguinity. Great is the glory of Mary—the mother of Jesus, the ideal of womanhood, the type of purity, obedience, meekness, and humility—but greater, infinitely greater is the glory of Christ—the perfect God-man—'the glory of the only-begotten of the Father, full of grace (\textgreek{πλήρης χάριτος} not \textgreek{κεχαριτωμένος}) and of truth.'

2. The dogma of the sinlessness of Mary is also \emph{uncatholic. }It lacks every one of the three marks of true catholicity, according to the canon of Vincentius Lirinensis, which is professedly recognized by Rome herself (the \emph{semper,} the \emph{ubique,} and the \emph{ab omnibus}), and instead of a 'unanimous consent' of the Fathers in its favor, there is a unanimous silence, or even protest, of the Fathers against it. For more than ten centuries after the Apostles it was not dreamed of, and when first broached as a pious opinion, it was strenuously opposed, and continued to be opposed till 1854 by many of the greatest saints and divines of the Roman Church, including St. Bernard and St. Thomas Aquinas, and several Popes.

The ante-Nicene Fathers, far from teaching that Mary was free from hereditary sin, do not even expressly exempt her from actual sin, certainly not from womanly weakness and frailty. Irenæus (d. 202), who first suggested the fruitful parallel of Eve as the mother of disobedience, and Mary as the mother of obedience (not justified by the true Scripture parallel between Adam and Christ), and thus prepared the way for a false Mariology, does yet not hesitate to charge Mary with 'unseasonable haste' or 'urgency,' which the Lord had to rebuke at the wedding of Cana (fcJohn ii. 4);\footnote{\par Iren. \emph{Adv. hœr}. iii. c. 16, § 7: \emph{Dominus, repellens intempestivam festinationem, dixit: 'Quid mihi et tibi est, mulier!}'} and even Chrysostom, at the close of the fourth century, ventured to say that she was immoderately ambitious, and wanting in proper regard for the glory of Christ on that occasion.\footnote{Chrys. \emph{Hom. XXI. al. XX. in Joh. Opera}, ed. Bened. Tom. VIII. p. 122. Compare his \emph{Hom. in Matth. XLIV. al. XLV.}, where he speaks of Mary's ambition (\textgreek{φιλοτιμία}) and thoughtlessness (\textgreek{ἀπόνοια}), when she desired to speak with Christ while he yet talked to the people (Matt. xii. 46 sqq.).} The last charge is hardly just, for in the words, 'Whatsoever he saith unto you, do it,' she shows the true spirit of obedience and absolute trust in her Divine Son. Tertullian implicates her in the unbelief of the brethren of Jesus.\footnote{\emph{De carne Christi}, c. 7: \emph{Fratres Domini non crediderant in illum. Mater æque non demonstratur adhæsisse illi, cum Marthæ et Mariæ aliæ in commercio ejus frequententur.}} Origen thinks that she took offense, like the Apostles, at our Lord's sufferings, else 'he did not die for her sins;' and, according to Basil, she, too, 'wavered at the time of the crucifixion.' Gregory of Nazianzus, and John of Damascus, the last of the great Greek Fathers, teach that she was \emph{sanctified} by the Holy Ghost; which has no meaning for a sinless being.

The first traces of the Romish Mariolatry and Mariology are found in the apocryphal Gospels of Gnostic and Ebionitic origin.\footnote{Compare the convenient digest of this apocryphal history of Mary and the holy family in E. Hoffmann's Leben Jesu nach den Apocryphen, Leipz. 1851, pp. 5–117, and Tischendorf's De evangeliorum apocryphorum origine et usu, Hagæ, 1851.} In marked contrast with the canonical Gospels, they decorate the life of Mary with marvelous fables, most of which have passed into the Roman Church, and some also into the Mohammedan Koran and its commentaries.\footnote{It must be remembered that Mohammed derived his defective knowledge of Christianity from Gnostic and other heretical sources. Gibbon and Stanley trace the Immaculate Conception directly to the Koran, III. pp. 31, 37 (Rodwell's translation, p. 499), where it is said of Mary: 'Remember when the angel said: ``Mary, verily has God chosen thee, and purified thee, and chosen thee above the women of the world.''' \par{} [Pius IX., March 24, 1877, spoke of Mary as \emph{divinarum potentissima conciliatrix gratiarum.} If possible, Leo XIII. in encyclicals on the rosary and other deliverances, and Pius X., went further in exalting Mary. Leo, Sept. 1, 1883, pronounced her 'the safest guide to reach the gracious hand of God,' and, Sept., 1891, affirmed that 'except through the Mother, it is hardly possible for any one to reach Christ.' On the fiftieth anniversary of the dogma of the immaculate conception, Oct. 17, 1904, Pius X. made astounding use of the Old Testament to substantiate her alleged virtues. Calling her the Spouse of the Holy Ghost, he announced that 'already Adam saw her in the distance as the destroyer of the serpent's head, and at the sight of her dried up his tears over the curse which had struck him'; Noah recalled her as he was preparing the ark; Abraham was estopped from sacrificing his son as he thought of her; Jacob saw her in the ladder on which the angels ascended and descended; Moses looked up to her at the burning bush; etc. Pius invoked her aid as the 'glorious helper against all heresies,' as Leo XIII. before had acclaimed her 'the glorious victor over all heretics,' and Pius XI. in his encyclical on Church Union, 1928. Mary, in accordance with the petition of the Provincial Baltimore Council, 1843, has been made by papal decree the 'heavenly guardian of the United States,' as Pius XI. took occasion to remind the world when the Peace Conference met in Washington, 1921. And in his apostolic letter recommending the Catholic University in Washington, he made the petition that 'the immaculate conception may bestow on all America the gifts of wisdom and salvation.' Cardinal Gibbons, \emph{Faith of Our Fathers}, p. 167, Bishop Gilmour in his \emph{Bible History for Catholic Schools}, pp. 11, 130, and also the recent Italian version of the Pentateuch, issued with papal approval, repeat the false translation of Gen. III:15, that Mary should bruise the serpent's head.—ED.]}

Mariolatry preceded the Romish Mariology. Each successive step in the excessive veneration (\emph{hyperdulia}) of the Virgin, and each festival memorializing a certain event in her life, was followed by a progress in the doctrine concerning Mary and her relation to Christ and the believer. The theory only justified and explained a practice already existing.

The Mariology of the Roman Catholic Church has passed through three stages: the \emph{perpetual virginity} of Mary, her freedom from \emph{actual} sin, and her freedom from \emph{hereditary} sin.

This progress in Mariolatry is strikingly reflected in the history of Christian art. 'The first pictures of the early Christian ages simply represent the woman. By-and-by we find outlines of the mother and the child. In an after-age the Son is sitting upon a throne, with the mother crowned, but sitting as yet below him. In an age still later, the crowned mother on a level with the Son. Later still, the mother on a throne \emph{above} the Son. And lastly, a Romish picture represents the eternal Son in wrath, about to destroy the earth, and the Virgin Intercessor interposing, pleading, by significant attitude, her maternal rights, and redeeming the world from his vengeance. Such was, in fact, the progress of Virgin-worship. First the woman reverenced for the Son's sake; then the woman reverenced above the Son, and adored.'

(1) The idea of the \emph{perpetual Virginity} of Mary was already current in the ante-Nicene age, and spread in close connection with the ascetic overestimate of celibacy, and the rise of monasticism. It has a powerful hold even over many Protestant minds, on grounds of religious propriety. Tertullian, who died about 220, still held that Mary bore children to Joseph after the birth of Christ. But towards the close of the fourth century the denial of her perpetual virginity (by the Antidicomarianites, by Helvidius and Jovinian) was already treated as a profane and indecent heresy by Epiphanius in the Greek, and Jerome in the Latin Church. Hence the hypothesis that the brethren and sisters of Jesus, so often mentioned in the Gospels, were either children of Joseph by a former marriage (Epiphanius), or only cousins of Jesus (Jerome). On the other hand, however, the same Epiphanius places among his eighty heresies the Mariolatry of the \emph{Collyridianæ}, a company of women in Arabia, in the last part of the fourth century, who sacrificed to Mary little cakes or loaves of bread (\textgreek{κολλυρίς,} hence the name \textgreek{Κολλυριδιανοί}), and paid her divine honor with festive rites similar to those connected with the cult of Cybele, the \emph{magna mater deûm,} in Arabia and Phrygia.

(2) The freedom of Mary from \emph{actual} sin was first clearly taught in the fifth century by Augustine and Pelagius, who, notwithstanding their antagonism on the doctrines of sin and grace, agreed in this point, as they did also in their high estimate of asceticism and monasticism. Augustine, for the sake of Christ's honor, exempted Mary from willful contact with actual sin;\footnote{\emph{De natura et gratia}, c. 36, § 42 (ed. Bened. Tom. X. p. 144): '\emph{Excepta sancta Virgine Maria,}  de qua propter honorem Domini nullam prorsus, cum de peccatis agitur, haberi volo quæstionem  . . . \emph{hac ergo Virgine excepta, si omnes illos sanctos et sanctas . . . congregare possemus et interrogare, utrum essent sine peccato, quid fuisse responsuros putamus, utrum hoc quod iste} [namely, Pelagius] \emph{dicit, an quod Joannes Apostolus} (1 John i. 8)?' This is the only passage in Augustine which at all favors the Romanists; and the force even of this is partly broken by the parenthetical question: '\emph{Unde enim scimus quid ei} [Mariæ] \emph{plus gratiæ collatum fuerit ad vincendum omni ex parte peccatum quæ concipere ac parere meruit, quem constat nullum habuisse peccatum}? For how do we know what \emph{more of grace for the overcoming of sin in every respect was bestowed upon her} who was found worthy to conceive and give birth to him who, it is certain, was without sin?' This implies that in Mary sin was, if not a developed act, at least a power to be conquered.} but he expressly included her in the fall of Adam and its hereditary consequences.\footnote{\emph{Sermo} 2 \emph{in Psalm}. 34: \emph{Maria ex Adam mortua propter peccatum, et caro Domini ex Maria mortua propter delenda peccata}; i.e., Mary died because of inherited sin, but Christ died for the destruction of sin. In his last great work, \emph{Opus imperf. contra Julian. IV.} c. 122 (ed. Bened. X. 1208), Augustine speaks of the grace of regeneration (\emph{gratia renascendi}) which Mary experienced. He also says explicitly that Christ \emph{alone} was without sin, \emph{De peccat. mer. et remiss}., II. c. 24, § 38 (ed. Bened. X. 61: Solus  \emph{ille, homo factus, manens Deus, peccatum nullum habuit unquam, nec sumpsit carnem peccati, quamvis de materna carne peccati}); ib. c. 35, § 57 (X. 69: \emph{Solus unus est qui sine peccato natus est in similitudine carnis peccati, sine peccato vixit inter aliena peccata, sine peccato mortuus est propter nostra peccata}); \emph{De Genesi ad lit}., c. 18, § 32; c. 20, § 35. These and other passages of Augustine clearly prove, to use the words of Perrone (l.c. pp. 42, 43 of the Germ. ed.), that 'this holy Father evidently teaches that Christ alone must be exempt from the general pollution of sin; but that the blessed Virgin, being conceived by the ordinary cohabitation of parents, partook of the general stain, and her flesh, being descended from sin, was sinful flesh, which Christ purified by assuming it.' The pupils of Augustine were even more explicit. One of them, Fulgentius (\emph{De incarn}. c. 15, § 29, also quoted by Perrone), says: 'The flesh of Mary, which was conceived in unrighteousness in a human way, was truly \emph{sinful} flesh.'} Pelagius, who denied hereditary sin, went further, and exempted Mary (with several other saints of the Old Testament) from sin altogether;\footnote{He says: 'Piety must confess that the mother of our Lord and Saviour was sinless' (as quoted by Augustine, \emph{De nat. et gratia}, cc. 36, § 42: '\emph{quam dicit sine peccato confiteri necesse esse pietati}'). Pelagius also excludes from sin Abel, Enoch, Melchisedek, Abraham, Isaac, Jacob, Noah, Samuel, Nathan, Elijah, Elisha, Daniel, Ezekiel, John the Baptist, Deborah, Anna, Judith, Esther, Elisabeth, and Joseph, the husband of Mary, who 'have not only not sinned, but also lived a righteous life.' Julian, his ablest follower, objected to Augustine that, by his doctrine of hereditary sin and universal depravity, he handed even Mary over to the power of the devil (\emph{ipsam Mariam diabolo nascendi conditione transcribis}); to which Augustine replied (\emph{Opus imperf. contra Jul.} 1. IV. c. 122): '\emph{Non transscribimus diabolo Mariam conditione nascendi, sed ideo quia ipsa conditio solvitur gratia renascendi,}' i.e., because this condition (of sinful birth) is solved or set aside by the grace of the second birth. When this took place, he does not state.} and, if he were not a condemned heretic, he might be quoted as the father of the modern dogma.\footnote{It is characteristic that the Dominicans and Jansenists, who sympathized with the Augustinian anthropology, opposed the Immaculate Conception; while the Franciscans and Jesuits, who advocated it, have a more or less decided inclination towards Pelagianizing theories, and reduce original sin to a loss of supernatural righteousness, i.e., something merely negative, so that it is much easier to make an exception in favor of Mary. The Jesuits, at least, have an intense hatred of Augustinian views on sin and grace, and have shown it in the Jansenist controversy.} The view which came to prevail in the Catholic Church was that Mary, though conceived in sin, like David and all men, was sanctified in the womb, like Jeremiah (i. 5) and John the Baptist (Luke i. 15), and thus prepared to be the spotless receptacle for the Son of God and Saviour of mankind. Many, however, held that she was not fully sanctified till she conceived the Saviour by the Holy Ghost. The extravagant praise lavished on 'the Mother of God' by the Fathers after the defeat of Nestorianism (431), and the frequent epithets \emph{most holy} and \emph{immaculate} (\textgreek{πανάγια,} \emph{immaculata} and \emph{immaculatissima}), refer only to her spotless purity of character after her sanctification, but not to her conception.\footnote{The predicate \emph{immaculate} was sometimes applied to other holy virgins, e.g., to S. Catharine of Siena, who is spoken of as \emph{la immaculata vergine,} in a decree of that city as late as 1462. See Hase, l.c. p.§336.} The Greek Church goes as far as the Roman in the \emph{practice} of Mariolatry, but rejects the dogma of the Immaculate Conception as subversive of the Incarnation.\footnote{See A. V. Mouravieff on the dogma, in Neale's  Voices from the East, 1859, pp. 117–155.}

(3) The third step, which exempts Mary from \emph{original} sin as well, is of much later origin. It meets us first as a pious \emph{opinion} in connection with the festival of the Conception of Mary, which was fixed upon Dec. 8, nine months before the older festival of her birth (celebrated Sept. 8). This festival was introduced by the Canons at Lyons in France, Dec. 8, 1139, and gradually spread into England and other countries. Although it was at first intended to be the festival of the Conception of the \emph{immaculate Mary}, it concealed the doctrine of the \emph{Immaculate Conception}, since every ecclesiastical solemnity acknowledges the sanctity of its object.

For this reason, Bernard of Clairvaux, 'the honey-flowing doctor' \emph{doctor mellifluus}), and greatest saint of his age, who, by a voice mightier than the Pope's, roused Europe to the second crusade, opposed the festival as a false honor to the royal Virgin, which she does not need, and as an unauthorized innovation, which was the mother of temerity, the sister of superstition, and the daughter of levity.\footnote{'\emph{Virgo regia falso non eget honore, veris cumalata honorum titulis. . . . Non est hoc Virginem honorare sed honori detraher. . . . Præsumpta novitas mater temeritatis, soror superstitionis, filia levitatis.}' See his \emph{Epistola} 174, \emph{ad Canonicos Lugdunenses, De conceptione S. Mar.} (\emph{Op.} ed. Migne, I. pp. 332–336). Comp. also Bernard's \emph{Sermo} 78 \emph{in Cant., Op.} Vol. II. pp.1160, 1162.} He urged against it that it was not sanctioned by the Roman Church. He rejected the opinion of the Immaculate Conception of Mary as contrary to tradition and derogatory to the dignity of Christ, the only sinless being, and asked the Canons of Lyons the pertinent question, 'Whence they discovered such a hidden fact? On the same ground they might appoint festivals for the conception of the parents, grandparents, and great-grandparents of Mary, and so on without end.'\footnote{. . . '\emph{et sic tenderetur in infinitum, et festorum non esset numerus}' (\emph{Ep.} 174, p. 334 sq.).} It does not diminish, but rather increases (for the Romish stand-point) the weight of his protest, that he was himself an enthusiastic eulogist of Mary, and a believer in her sinless birth. He put her in this respect on a par with Jeremiah and John the Baptist.\footnote{'\emph{Si igitur ante conceptum sui sanctificari minime potuit, quoniam non erat; sed nec in ipso quidem conceptu, propter peccatum quod inerat: restat ut post conceptum in utero jam existens sanctificationem accepisse credatur, quæ excluso peccato sanctam fecerit nativitatem, non tamen et conceptionem}' (l.c. p. 336).}

The same ground was taken substantially by the greatest schoolmen of the Middle Ages till the beginning of the fourteenth century: Anselm of Canterbury (d. 1109), who closely followed Augustine;\footnote{Anselm, who is sometimes wrongly quoted on the other side, says, \emph{Cur Deus Homo}, ii. 16 (\emph{Op.} ed. Migne, I. p. 416): '\emph{Virgo ipsa . . . est in iniquitatibus concepta, et in peccatis concepit eam mater ejus, et cum originali peccato nata est, quoniam et ipsa in Adam peccavit, in quo omnes peccaverunt.}' To these words of Boso, Anselm replies that 'Christ, though taken from the sinful mass (\emph{de massa peccatrice assumptus}), had no sin.' Then he speaks of Mary twice as being purified from sin (\emph{mundata a peccatis}) by the future death of Christ (c. 16, 17). His pupil and biographer, Eadmer, in his book \emph{De excellent. beatæ Virg. Mariæ}, c. 3 (Ans. \emph{Op.} ed. Migne, II. pp. 560–62), says that the blessed Virgin was freed from all remaining stains of hereditary and actual sin when she consented to the announcement of the mystery of the Incarnation by the angel.' Quoted also by Perrone, pp. 47–49.} Peter the Lombard, 'the Master of Sentences' (d. 1161); Alexander of Hales, 'the irrefragable doctor' (d. 1245); St. Bonaventura, 'the seraphic doctor' (d. 1274); Albertus Magnus, 'the wonderful doctor' (d. 1280); St. Thomas Aquinas, 'the angelic doctor' (d. 1274), and the very champion of orthodoxy, followed by the whole school of Thomists and the order of the Dominicans. St. Thomas taught that Mary was conceived from sinful flesh in the ordinary way, \emph{secundum carnis concupiscentiam ex commixtione maris,} and was sanctified in the womb \emph{after} the infusion of the soul (which is called the \emph{passive} conception); for otherwise she would not have needed the redemption of Christ, and so Christ would not be the Saviour of \emph{all} men. He distinguishes, however, three grades in the sanctification of the Blessed Virgin: first, the \emph{sanctificatio in utero,} by which she was freed from the original guilt (\emph{culpa originalis}); secondly, the \emph{sanctificatio in conceptu Domini,} when the Holy Ghost overshadowed her, whereby she was totally purged (\emph{totaliter mundata}) from the fuel or incentive to sin (\emph{fomes peccati}); and, thirdly, the \emph{sanctificatio in morte,} by which she was freed from all consequences of sin (\emph{liberata ab omni miseria}). Of the festival of the Conception, he says that it was not observed, but tolerated by the Church of Rome, and, like the festival of the Assumption, was not to be \emph{entirely} rejected (\emph{non totaliter reprobanda}).\footnote{\emph{Summa Theologiæ}, Pt. III. Qu. 27 (\emph{De sanctificatione B. Virg.}), Art. 1–5; in \emph{Libr. I. Sentent.} Dist. 44, Qu. 1, Art. 3. Nevertheless, Perrone (pp. 231 sqq.) thinks that St. Bernard and St. Thomas are not in the way of a definition of the new dogma, 'because they wrote at a time when this view was not yet made quite clear, and because they lacked the \emph{principal} support, which subsequently came to its aid; hence they must in this case be regarded as \emph{private} teachers, propounding their own particular opinions, but not as witnesses of the traditional meaning of the Church.' He then goes on to charge these doctors with comparative ignorance of previous Church history. This may be true, but does not help the matter; since the fuller knowledge of the Fathers in modern times reveals a still wider dissent from the dogma of the Immaculate Conception.} The University of Paris, which during the Middle Ages was regarded as the third power in Europe, gave the weight of its authority for a long time to the doctrine of the Maculate Conception. Even seven Popes are quoted on the same side, and among them three of the greatest, viz., Leo I. (who says that Christ alone was free from original sin, and that Mary obtained her purification through her conception of Christ), Gregory I., and Innocent III.\footnote{The other Popes, who taught that Mary was conceived in sin, are Gelasius I., Innocent V., John XXII., and Clement VI. (d. 1352). The proof is furnished by the Jansenist Launoy, \emph{Prœscriptions, Opera I.} pp. 17 sqq., who also shows that the early Franciscans, and even Loyola and the early Jesuits, denied the Immaculate Conception of Mary. Perrone calls him an 'irreligious innovator' (p. 34), and an 'impudent liar' (p. 161), but does not refute his arguments, and evades the force of his quotations from Leo, Gelasius, and Gregory by the futile remark that they would prove too much, viz., that Mary was even \emph{born} in sin, and not purified before the Incarnation, which would be impious!}

But a change in favor of the opposite view was brought about, in the beginning of the fourteenth century, by Duns Scotus, 'the subtle doctor' (d. 1308), who attacked the system of St. Thomas and the Augustinian doctrine of original sin, who delighted in the most abstruse questions and the most intricate problems, to show the skill of his acute dialectics, and who could twist a disagreeable text into its opposite meaning. He was the first schoolman of distinction who advocated the Immaculate Conception, first at Oxford, though very cautiously, as a possible and probable fact.\footnote{Duns Scotus, \emph{Opera}, Lugd. 1639, Tom. VII. Pt. I. pp. 91–100. One of his arguments of probability is that, as God blots out original sin by baptism every day, he can as well do it in the moment of conception. Compare Perrone, pp. 18 sqq.} He refuted, according to a doubtful tradition, the opposite theory, in a public disputation at Paris, with no less than two hundred arguments, and converted the University to his view.\footnote{Related by Wadding, in his \emph{Annal. Minorum}, Lugd. 1635, Tom. III. p. 37, but rejected by Natalis Alexander, in his \emph{Church History}, as a fiction, and doubted even by Perrone (p. 163), who says, however, that Duns Scotus refuted all the arguments of his opponents 'in a truly astounding manner.'} At all events, he made it a distinctive tenet of his order.

Henceforward the Immaculate Conception became an apple of discord between rival schools of Thomists and Scotists, and the rival orders of the Dominicans and Franciscans. They charged each other with heresy, and even with mortal sin for holding the one view or the other. Visions, marvelous fictions, weeping pictures of Mary, and letters from heaven were called in to help the argument for or against a fact which no human being, not even Mary herself, can know without a divine revelation. Four Dominicans, who were discovered in a pious fraud against the Franciscan doctrine, were burned, by order of a papal court, in Berne, on the eve of the Reformation. The Swedish prophetess, St. Birgitte, was assured in a vision by the Mother of God that she was conceived without sin; while St. Catharine of Siena prophesied for the Dominicans that Mary was sanctified in the third hour after her conception. So near came the contending parties that the difference, though very important as a question of principle, was practically narrowed down to a question of a few hours. The Franciscan view gradually gained ground. The University of Paris, the Spanish nation, and the Council of Basle (1439) favored it. Pope Sixtus IV., himself a Franciscan, gave his sanction and blessing to the festival of the Immaculate Conception, but threatened with excommunication all those of \emph{both} parties who branded the one or the other doctrine as a heresy and mortal sin, since the Roman Church had not yet decided the question (1476 and 1483).

The Council of Trent (June 17, 1546) confirmed this neutral position, but with a leaning to the Franciscan side, by adding to the dogma on original sin the caution that it was not intended 'to comprehend in this decree the blessed and immaculate Virgin Mary.'\footnote{Sessio V.: '\emph{Declarat S. Synodus, non esse suæ intentionis, comprehendere in hoc decreto, ubi de peccato originali agitur, beatam et immaculatam Virginem Mariam, Dei genitricem; sed observandas esse constitutiones felicis recordationis Sixti Papæ IV. sub pœnis in eis constitutionibus contentis, quas innovat}.'} Pius V. (1570), a Dominican, condemned Baius (De Bay, Professor at Louvain, and a forerunner of the Jansenists), who held that Mary had actual as well as original sin; but soon afterwards he ordered that the discussion of this delicate question should be confined to scholars in the Latin tongue, and not be brought to the pulpit or among the people. In the mean time the Franciscan doctrine was taken up and advocated with great zeal and energy by the Jesuits. At first they felt their way cautiously. Bellarmin declared the Immaculate Conception to be a pious and probable opinion, more probable than the opposite. In 1593 the fifth general assembly of the order directed its teachers to depart from St. Thomas in this article, and to defend the doctrine of Scotus, 'which was then more common and more accepted among theologians.' It is chiefly through their influence that it gained ground more and more, yet under constant opposition. Paul V. (1616) still left both parties the liberty to advocate their opinion; but a decree of the Congregation of the Holy Inquisition and Gregory XV. (1622) prohibited the publication of the doctrine that Mary was conceived in sin, and removed from the liturgy the word \emph{sanctification} with reference to Mary. Then a new controversy arose as to the meaning of the term \emph{immaculate}; whether it referred to the Virgin or to her conception? To make an end to all dispute, Alexander VII., urged on by the King of Spain, issued a constitution, Dec. 8, 1661, which recommends the Immaculate Conception, defining it almost in the identical words of the dogma of Pius IX.\footnote{'\emph{Ejus} (sc. \emph{Mariæ}),' says Alexander VII., in the bull \emph{Sollicitudo Omnium Ecclesiarum} (\emph{Bullar. Rom.} ed. Coquelines, Tom. VI. p. 182), '\emph{animam in primo instanti creationis atque infusionis in corpus fuisse speciali Dei gratia et privilegio, intuitu meritorum Christi, ejus Filii, humani generis Redemptoris, a macula peccati originalis præservatam immunem.}' Compare the decree of Pius IX. p. 110, which substitutes \emph{suæ conceptionis} for \emph{creationis atque infusionis} (\emph{animæ}) \emph{in corpus,} and \emph{ab omni originalis culpæ labe} for \emph{a macula peccati originalis.}}

Nothing was left but the additional declaration that belief in this doctrine was necessary to salvation. 'From this time,' says Perrone,\footnote{L.c. p. 33.} 'every controversy and opposition to the mystery ceased, and the doctrine of the Immaculate Conception attained to full and quiet possession in the whole Catholic Church. No sincere Catholic ventured hereafter to utter even a sound against it, with the exception of some irreligious innovators, among whom Launoy occupies the first place, and, in these last years, George Hermes.' Thus he disposes of the powerful protest of Launoy, issued in 1676, fifteen years after the bull of Alexander VII., with irrefragable testimonies of Fathers and Popes; to which may be added the anonymous treatise '\emph{Against Superstition},' written by Muratori, 1741, one of the most learned antiquarians and historians of the Roman Church. But Jansenism was crushed; Jesuitism, though suppressed for a while, was restored to greater power; Ultramontanism and Papal Absolutism made headway over the decay of independent learning and research; the voice of the ablest remaining Catholic scholars was unheeded; the submissiveness of the Bishops, and the ignorance, superstition, and indifference of the people united in securing the triumph of the dogma.

3. The only \emph{dogmatic} argument adduced is that of congruity or fitness, in view of the peculiar relations which Mary sustains to the persons of the Holy Trinity. Being eternally chosen by the Father to be 'the bride of the Holy Ghost,' and 'the mother of the Son of God,' it was eminently proper that, from the very beginning of her existence, she should be entirely exempt from contact with sin and the dominion of Satan.\footnote{Perrone, ch. xiv. pp. 102 sqq.}

To this it is sufficient to answer that the Word of God is the highest and only infallible standard of religious propriety; and this standard concludes all men under the power of sin and death, with the only exception of the God-man, the sinless Redeemer of the fallen race. Besides, the argument of congruity can at best only prove the possibility of a fact, not the fact itself. And, finally, it would prove too much in this case; for, if propriety demands a sinless mother for a sinless Son, it demands also (as St. Bernard suggested) a sinless grandmother, great-grandmother, and an unbroken chain of sinless ancestors to the beginning of the race.

On the other hand, the new dogma, viewed even from the stand-point of the Roman Catholic system, involves contradictory elements.

In the first place, it is inconsistent with any proper view of original sin, no matter whether we adopt the theory of traducianism, or that of creationism (which prevails among Roman divines), or that of pre-existence. The bull of 1854 speaks indefinitely of the 'conception' of Mary. But Roman divines usually distinguish between the \emph{active} conception, i.e., the marital act by which the seed of the body is formed by the agency of the parents, and the \emph{passive} conception, i.e., the infusion of the soul into the body by a creative act of God (according to the theory of creationism).\footnote{As to the time of the creation and infusion of the soul, whether it took place simultaneously with the generation of the body, or on the fortieth day (as was formerly supposed), there is no fixed opinion among Roman divines.} The meaning of the new dogma is that Mary, by a special grace and privilege, was exempt from original sin in her \emph{passive} conception, that is, in that moment when her soul was created by God for the animation of her body.\footnote{So the matter is explained by Perrone at the beginning of his Treatise, pp. 1–4; and this accords with the bull of Alexander VII. (\emph{in primo instanti creationis atque infusionis in corpus,} etc.), see p. 125.} Now original sin must come either from the body, or from the soul, or from both combined. If from the body, then Mary must have inherited it from her parents, since the dogma does not exclude these from sin; if from the soul, then God, who creates the soul, is the author of sin, which is blasphemous; if from both, then we have a combination of both these inextricable difficulties. Nor is the matter materially relieved if we take the superficial semi-Pelagian view of hereditary sin, which makes it a mere privation or defect, namely, the absence of the supernatural endowment of original righteousness and holiness (the \emph{similitudo Dei,} as distinct from the \emph{imago Dei}), instead of a positive disorder and sinful disposition.\footnote{The profounder schoolmen, however, represented by St. Thomas, had a deeper view of original sin, nearer to that of Augustine and the Reformers. The same is true of Möhler, who speaks of a 'deep vulneration of the soul in all its powers,' and a 'perverse tendency of the will,' as a necessary consequence of the Fall.} For even in this case the same dilemma returns, that this original defect must have been there from the parents, or must be ordinarily derived from God, as the author of the soul, which alone can be said to possess or to lose righteousness and holiness. Rome must either deny original sin altogether (as Pelagius did), or take the further step of making the Immaculate Conception of Mary a strictly miraculous event, like the conception of Christ by the Holy Ghost, \emph{sine virili complexu} and \emph{sine concupiscentia carnis.}

Secondly, the dogma, by exempting Mary from original sin in consequence of \emph{the merits of Christ},\footnote{. . . '\emph{intuitu meritorum Christi Jesu, Salvatoris humani generis.}'} virtually puts her under the power of sin; for the merits of Christ are only for sinners, and have no bearing upon sinless beings. Perrone, following Bellarmin, virtually concedes this difficulty, and vainly tries to escape it by an unmeaning figure, that Mary was delivered from prison before she was put into it, or that her debt was paid which she never contracted!

Finally, the dogma is inconsistent with the Vatican decree of Papal Infallibility. The hidden fact of Mary's Immaculate Conception must, in the nature of the case, be a matter of divine omniscience and divine revelation, and is so declared in the papal decree.\footnote{. . . '\emph{doctrinam . . . esse a Deo revelatam},' etc.} Now it must have been revealed to the mind of Pius IX., or not. If not, he had no right, in the absence of Scripture proof, and the express dissent of the Fathers and the greatest schoolmen, to declare the Immaculate Conception a divinely revealed fact and doctrine. If it was revealed to him, he had no need of first consulting all the Bishops of the Roman Church, and waiting several years for their opinion on the subject. Or if this consultation was the necessary medium of such revelation, then he is not in himself infallible, and has no authority to define and proclaim any dogma of faith without the advice and consent of the universal Episcopate.

\xsection{The Papal Syllabus, A.D. 1864.}

§ 30. The Papal Syllabus, A.D. 1864.

Literature.

The \emph{Enyclica} and \emph{Syllabus} of Dec. 8, 1864, are published in \emph{Pii IX. Epistola encycl}., etc., Regensb. 1865; in \emph{Officielle Actenstücke zu dem v. Pius IX. nach Rom. berufenen Oekum. Concil}, Berlin, 1869, pp. 1–35, in \emph{Acta et Decreta S. œcum. Conc. Vatic.} Frib. 1871, Pt. I. pp. 1–21, etc.

J. Tosi (R.C.): \emph{Vorlesungen über den Syllabus errorum der päpstl. Encyclica}, Wien, 1865 (251 pp.).

J. Hergenröther (R.C.): \emph{Die Irrthümer der Neuzeit gerichtet durch den heil. Stuhl}, 1865.

\emph{Beleuchtung der päpstlichen Encyclica v.} 8 \emph{Dec.} 1864, \emph{und das Verzeichniss der modernen Irrthümer} (by a R.C.), Leipz. 1865.

\emph{Die Encyclica Papst Pius IX. vom} 8 \emph{Dec.} 1864. \emph{Stimmen aus Maria-Laach} (R.C.), Freib. 1866–69. (By Riess, Schneemann, and others.)

\emph{Der Papst und die modernen Ideen} (R.C.), several numbers, Wien, 1865–67. [By  Cl. Schrader, a Jesuit.]

C. Pronier (Prof. of the Free Theol. Sem. at Geneva, 1873): \emph{La liberté: religieuse et le Syllabus}, Genève, 1870.

W. E. Gladstone: \emph{The Vatican Decrees: a Political Expostulation}, London and New York, 1874; \emph{Vaticanism}, 1875. Comp. the Roman Catholic Replies of Monsign. Capel, J. H. Newman, and Archbishop Manning in defense of the Vatican Decrees; see below, § 31.

On the 8th of December, 1864, just ten years after the proclamation of the sinlessness of the Virgin Mary, Pope Pius IX. issued an encyclical letter '\emph{Quanta cura},' denouncing certain dangerous heresies and errors of the age, which threatened to undermine the foundations of the Catholic religion and of civil society, and exhorting the Bishops to counteract these errors, and to teach that 'kingdoms rest on the foundation of the Catholic faith;' that it is the chief duty of civil government 'to protect the Church;' that 'nothing is more advantageous and glorious for rulers of States than to give free scope to the Catholic Church, and not to allow any encroachment upon her liberty.'\footnote{These and similar sentences are inserted from letters of mediæval Popes, who from their theocratic stand-point claimed supreme jurisdiction over the states and princes of Europe. Popes, like the Stuarts and the Bourbons, never forget and never learn any thing.} In the same letter the Pope offers to all the faithful a complete indulgence for one month during the year 1865,\footnote{. . . '\emph{plenariam indulgentiam ad instar jubilæi concedimus intra unius tantum mensis spatium usque ad totum futurum annum} 1865 \emph{et non ultra.}'} and expresses, in conclusion, his unbounded confidence in the intercession of the immaculate and most holy Mother of God, who has destroyed all the heresies in the whole world, and who, being seated as queen at the right hand of her only begotten Son, can secure any thing she asks from him.\footnote{'\emph{Quo vero facilius Deus Nostris, Vestrisque, et omnium fidelium precibus, votisque annuat, cum omni fiducia deprecatricem apud Eum adhibeamus Immaculatam Sanctissimamque Deiparam Virginem Mariam, quæ cunctas hereses interemit in universo mundo, quæque omnium nostrum amantissima Mater} ``\emph{tota suavis est . . . ac plena misericordiæ . . . omnibus sese exorabilem, omnibus clementissimam prœbet, omnium necessitates amplissimo quodam miseratur affectu}'' [quoted from St. Bernard], \emph{atque utpote Regina adstans a dextris Unigeniti Filii Sui, Domini Nostri Jesu Christi, in vestitu deaurato circumamicta varietate, nihil est quod ab Eo impetrare non valeat. Suffragia quoque petamus Beatissimi Petri Apostolorum Principis, et Coapostoli ejus Pauli, omniumque Sanctorum Cœlitum, qui facti jam amici Dei pervenerunt ad cœlestia regna, et coronati possident palmam, ac de sua immortalitate securi, de nostra sunt salute solliciti.}'}

To this characteristic Encyclical is added the so-called Syllabus, i.e., a catalogue of eighty errors of the age, which had been previously pointed out by Pius IX. in Consistorial Allocutions, Encyclical and other Apostolic Letters, but are here conveniently brought together, and were transmitted by Cardinal Antonelli to all the Bishops of the Roman Catholic Church.

This extraordinary document presents a strange mixture of truth and error. It is a protest against atheism, materialism, and other forms of infidelity which every Christian must abhor; but it is also a declaration of war against modern civilization and the course of history for the last three hundred years. Like the papal bulls against the Jansenists, it is purely negative, but it implies the assertion of doctrines the very opposite to those which are rejected as errors.\footnote{A learned Jesuit, Clemens Schrader, translated them into a positive form.} It expressly condemns religious and civil liberty, the separation of Church and State; and indirectly it asserts the Infallibility of the Pope, the exclusive right of Romanism to recognition by the State, the unlawfulness of all non-Catholic religions, the complete independence of the Roman hierarchy from the civil government (yet without allowing, a separation), the power of the Church to coerce and enforce, and its supreme control over public education, science, and literature.

The number of errors was no doubt suggested by the example of Epiphanius, the venerable father of heresy-hunters (d. 403), who, in his \emph{Panarion}, or \emph{Medicine-Chest}, furnishes antidotes for the poison of no less than eighty heresies (including twenty before Christ), probably with a mystic reference to the \emph{octoginta concubinæ} in the Song of Solomon (vi. 8).

The Pope divides the eighty errors of the nineteenth century into ten sections, as follows:

I. Pantheism, Naturalism, and Absolute Rationalism, No. 1–7.

Under this head are condemned the following errors:

(1.) The denial of the existence of God.

(2.) The denial of his revelation.

(3 and 4.) The sufficiency of human reason to enlighten and to guide men.

(5.) Divine revelation is imperfect, and subject to indefinite progress.

(6.) The Christian faith contradicts human reason, and is an obstacle to progress.

(7.) The prophecies and miracles of the Bible are poetic fictions, and Jesus himself is a myth.\footnote{'\emph{Jesus Christus est mythica fictio.}' I am not aware that any sane infidel has ever gone so far. Strauss and Renan resolve the \emph{miracles} of the gospel history into myths or legends, but admit the historical existence and extraordinary character of Jesus, as the greatest religions genius who ever lived.}

II. Moderate Rationalism, No. 8–14.

Among these errors are:

(12.) The decrees of the Roman See hinder the progress of science.

(13.) The scholastic method of theology is unsuited to our age.\footnote{No. 13. '\emph{Methodus et principia, quibus antiqui Doctores scholastici theologiam excoluerunt, temporum nostrorum necessitatibus scientiarumque progressui minime congruunt.}'}

(14.) Philosophy must be treated without regard to revelation.

III. Indifferentism, Latitudinarianism, No. 15–18.

(15.) Every man may embrace and profess that religion which commends itself to his reason.\footnote{No. 15. '\emph{Liberum cuique homini est eam amplecti ac profiteri religionem, quam rationis lumine quis ductus veram putaverit.}'}

(16.) Men may be saved under any religion.\footnote{No. 16. '\emph{Homines in cujusvis religionis cultu viam æternæ salutis reperire æternamque salutem assequi possunt.}'}

(17.) We may at least be hopeful concerning the eternal salvation of all non-Catholics.\footnote{No. 17. '\emph{Saltem bene sperandum est de æterna illorum omnium salute, qui in vera Christi Ecclesia nequaquam versantur.}'}

(18.) Protestantism is only a different form of the same Christian religion, in which we may please God as well as in the Catholic Church.\footnote{No. 18. '\emph{Protestantismus non aliud est quam diversa veræ ejusdem christianæ religionis forma, in qua æque ac in Ecclesia catholica Deo placere datum est.}'}

IV. Socialism, Communism, Secret Societies, Bible Societies, Clerico-Liberal Societies.

Under this head there are no specifications, but the reader is referred to previous Encyclicals of 1848, 1849, 1854, 1863, in which '\emph{ejusmodi pestes sæpe gravissimisque verborum formulis reprobantur.}' The Bible Societies, therefore, are put on a par with socialism and communism, as pestilential errors worthy of the severest reprobation!

V. Errors respecting the Church and her Rights.

Twenty errors (19–38), such as these: the Church is subject to the State; the Church has no right to exercise her authority without the leave and assent of the State; the Church has not the power to define dogmatically that the religion of the Catholic Church is the only true religion; Roman Pontiffs and œcumenical Councils have exceeded the limits of their power, usurped the rights of princes, and have erred even in matters of faith and morals;\footnote{No. 23. '\emph{Romani pontifices et concilia œcumenica a limitibus suæ potestatis recesserunt, jura principum usurparunt, atque etiam in rebus fidei et morum definiendis errarunt.}'} the Church has no power to avail herself of force, or any temporal power, direct or indirect;\footnote{No. 24. '\emph{Ecclesia vis inferendæ potestatem non habet, neque potestatem ullam temporalem directam vel indirectam.}'} besides the inherent power of the Episcopate, there is another temporal power conceded expressly or tacitly by the civil government, which may be revoked by the same at its pleasure; it does not exclusively belong to the jurisdiction of the Church to direct the teaching of theology; nothing forbids a general council, or the will of the people, to transfer the supreme Pontiff from Rome to some other city; national Churches, independent of the authority of the Roman Pontiff, may be established;\footnote{No. 37. '\emph{Institui possunt nationales Ecclesiæ ab auctoritate Romani Pontificis subductæ planeque divisæ.}'} the Roman Pontiffs have contributed to the Greek schism.\footnote{No. 38. '\emph{Divisioni ecclesiæ in orientalem atque occidentalem nimia Romanorum Pontificum arbitria contulerunt.}'}

VI. Errors concerning Civil Society, considered as well in itself as in its relations to the Church. Seventeen errors (39–55).

(44.) 'Civil authority may meddle in things pertaining to religion, morals, and the spiritual government.'

(45.) 'The whole government of public schools, in which the youth of a Christian commonwealth is trained, with the exception of some Episcopal seminaries, can and must be assigned to the civil authority.'\footnote{No. 45. '\emph{Totum scholarum publicarum regimen, in quibus juventus christianæ alicujus Reipublicæ instituitur, episcopalibus dumtaxat seminariis aliqua ratione exceptis, potest ac debet attribui auctoritati civili,}' etc. Compare Nos. 47 and 48. Hence the irreconcilable hostility of the Romish clergy to public schools, especially where the Protestant Bible is read.}

(46.) 'The method of study even in the seminaries of the clergy is subject to the civil authority.'

(52.) 'The lay government has the right to depose Bishops from the exercise of pastoral functions, and is not bound to obey the Roman Pontiff in those things which pertain to the institution of bishoprics and bishops.'

(55.) 'The Church is to be separated from the State, and the State from the Church.'\footnote{No. 55. '\emph{Ecclesia a Statu, Statusque ab Ecclesia sejungendus est.}' Compare Alloc. \emph{Acerbissimum} 27 Sept. 1852.}

VII. Errors in Natural and Christian Ethics, No. 56–64. Here among other things are condemned the principle of non-intervention, and rebellion against legitimate princes.

VIII. Errors on Christian Matrimony, No. 65–74.

Here the Pope condemns not only loose views on marriage and divorce, but also civil marriage, and any theory which does not admit it to be a sacrament.\footnote{No. 73. '\emph{Vi contractus mere civilis potest inter Christianos constare veri nominis matrimonium; falsumque est, aut contractum matrimonii inter Christianos semper esse sacramentum, aut nullum esse contractum, si sacramentum excludatur.}'}

IX. Errors regarding the Civil Principality of the Roman Pontiff, No. 75, 76.

(75.) Concerning the compatibility of the temporal reign with the spiritual, there is a difference of opinion among the sons of the Christian and Catholic Church.

(76.) The abrogation of the civil government of the Apostolic See would be conducive to the liberty and welfare of the Church.

X. Errors referring to Modern Liberalism, No. 77–80.

Under this head are condemned the principles of religious liberty as they have come to prevail in the most enlightened States of Christendom. The Pope still holds that it is right to forbid and exclude all religions but his own, where he has the power to do so (as he had and exercised in Rome before 1870); and he refuses to make any terms with modern civilization.\footnote{(77.) '\emph{Ætate hoc nostra non amplius expedit, religionem catholicam haberi tamquam unicam status religionem, ceteris quibuscumque cultibus exclusis.}' \par (78.) '\emph{Hinc laudabiliter in quibusdam catholici nominis regionibus lege cautum est, ut hominibus illuc immigrantibus liceat publicum proprii cujusque cultus exercitium habere.}' \par (79.) '\emph{Enimvero falsum est, civilem cujusque cultus libertatem, itemque plenam potestatem omnibus attributam quaslibet opiniones cogitationesque palam publiceque manifestandi conducere ad populorum mores animosque facilius corrumpendos ac indifferentismi pestem propagandam.}' \par (80.) '\emph{Romanus Pontifex potest ac debet cum progressu, cum liberalismo et cum recenti civilitate sese reconciliare et componere.}'}

The Syllabus, though resting solely on the authority of the Pope, must be regarded as an integral portion of the Roman Creed; the Pope having since been declared infallible in his official utterances. The most objectionable as well as the least objectionable parts of it have been formally sanctioned by the Vatican Council. The rest may be similarly sanctioned hereafter. The Syllabus expresses the genuine spirit of Popery, to which may be applied the dictum of the General of the Jesuits: '\emph{Aut sit ut est, aut non sit.}' It can not change without destroying itself.

In the mean time the politico-ecclesiastical doctrines of the Syllabus, together with the Infallibility decree, have provoked a new conflict between the Pope and the Emperor. Pius IX. looks upon the State with the same proud contempt as Gregory VII. 'Persecution of the Church,' he said after the recent expulsion of the Jesuits (1872), 'is folly: a little stone [Dan. ii. 45] will break the colossus [of the new German empire] to pieces.' But Bismarck, who is made of sterner stuff than Henry IV., protests: 'We shall not go to Canossa.'

American Protestants and European Free Churchmen reject all interference of the civil government with the liberty and internal affairs of the Church as much as the Pope, but they do this on the basis of a peaceful separation of Church and State, and an equality of all forms of Christianity before the law; while the Syllabus claims absolute freedom and independence exclusively for the Roman hierarchy, and claims this even in those countries where the State supports the Church, and has therefore a right to a share in its government.

[The Syllabus of Pius IX. was substantially confirmed by Leo XIII., Nov. 1, 1885, June 1, 1889, and Feb. 1, 1890, and Pius XI. in \emph{pascendi gregis,} 1907. It is pronounced infallible by Lehmkuhl, \emph{Theol. Mor.,} II., 780, Straub, \emph{de eccles.,} II., 398–402, and Leitner, \emph{Hdbuch. des kath. Kirchenrechts,} 2nd ed., p. 15. Other documents pronounced by Lehmkuhl, II., 726–88, infallible, are Leo X.'s bull against Luther, 1520, Innocent X.'s against Jansen, Innocent XI.'s against the Laxists, etc.—Ed.]

\xsection{The Vatican Council, A.D. 1864.}

§ 31. The Vatican Council, 1870.

Literature.

I. Works Preceding the Council.

\emph{Officielle Actenstücke zu dem von Sr. Heiligkeit dem Papste Pius IX. nach Rom berufenen Oekumenischen Concil,} Berlin, 1869 (pp. 189). This work contains the Papal Encyclica of 1864, and the various papal letters and official documents preparatory to the Council, in Latin and German.

\emph{Chronique concernant le Prochain Concile. Traduction revue et approuvée de la Civiltà cattolica par la correspondance de Rome,} Vol. I. Avant le Concile. Rome, Deuxième ed. 1869, fol. (pp. 192). Begins with the Papal letter of June 26, 1867.

Henry Edward Manning (Archbishop of Westminster): \emph{The Centenary of St. Peter and the General Council. A Pastoral Letter}. London, 1867. \emph{The Œcumenical Council and the Infallibility of the Roman Pontiff. A Pastoral Letter}. London, 1869. In favor of Infallibility.

C. H. A. Plantier (Bishop of Nîmes): \emph{Sur les Conciles généraux à l’occasion de celui que Sa Sainteté Pie IX. a convoqué pour le 8 décembre prochain,} Nîmes et Paris, 1869. The same in German: \emph{Ueber die allgemeinen Kirchenversammlungen,} translated by Th. von Lamezan, Freiburg im Breisgau, 1869. Infallibilist.

Magr. Vict. Aug. Dechamps (Archbishop of Malines): \emph{L’infaillibilité et le Concile général,} 2d ed., Paris et Malines, 1869. German translation: \emph{Die Unfehlbarkeit des Papstes und das Allgemeine Concil,} Mainz, 1869. Strong Infallibilist.

H. L. C. Maret (Dean of the Theol. Faculty of Paris): \emph{Du Concile général et de la paix religieuse,} Paris, 1869, 2 vols. Against Infallibility. Has since recanted.

W. Emmanuel Freiherr von Ketteler (Bishop of Mayence): \emph{Das Allgemeine Concil und seine Bedeutung für unsere Zeit,} 4th ed. Mainz, 1869. First against, now in favor of Infallibility.

Dr.  Joseph Fessler (Bishop of St. Pölten and Secretary of the Vatican Council, d. 1872): \emph{Das letzte und das nächste Allgemeine Concil,} Freiburg im Breisgau, 1869.

F. Dupanloup (Bishop of Orleans): \emph{Lettre sur le futur Concile Œcuménique,} in French, German, and other languages, 1869. The same on the \emph{Infallibility of the Pope}. First against, then in favor of the new dogma.

\emph{Der Papst und das Concil, von } Janus, Leipzig, 1869 (pseudonymous). The same in English: \emph{The Pope and the Council}, by Janus, London, 1869. In opposition to the Jesuit programme of the Council, from the liberal (old) Catholic stand-point; probably the joint production of Profs. Döllinger, Friedrich, and Huber, of the University of Munich.

Dr.  J. Hergenröther (R.C.): \emph{Anti-Janus}, Freiburg im Breisgau, 1870. Also in English, by  J. B. Robertson, Dublin, 1870.

\emph{Reform der Röm. Kirche in Haupt und Gliedern Aufgabe des bevorstehenden Röm. Concils,} Leipz. 1869. Liberal Catholic.

Felix Bungener (Prot.): \emph{Rome and the Council in the Nineteenth Century. Translated from the French, with additions by the Author}. Edinb. 1870. (Conjectures as to what the Council will be, to judge from the Papal Syllabus and the past history of the Papacy.)

II. Reports During the Council.

The \emph{Civiltà catholica,} of Rome, for 1869 and 1870. Chief organ of the Jesuits and Infallibilists.

Louis Veuillot: \emph{Rome pendant le Concile,} Paris, 1870, 2 vols. Collection of his correspondence to his journal, \emph{l’Univers,} of Paris. Ultra-Infallibilist and utterly unscrupulous.

J. Friedrich (Prof. of Church History in Munich, lib. Cath.): \emph{Tagebuch während des Vaticanischen Concils geführt,} Nödlingen, 1871; 2d ed. 1872. A journal kept during the Council, and noting the facts, projects, and rumors as they came to the surface. The author, a colleague and intimate friend of Döllinger, has since been excommunicated.

Quirinus: \emph{letters from Rome on the Council}, first in the Augsb. \emph{Allgemeine Zeitung,} and then in a separate volume, Munich, 1870; also in English, London, 1870 (pp. 856). Letters of three liberal Catholics, of different nations, who had long resided in Rome, and, during the Council, communicated to each other all the information they could gather from members of the Council, and sent their letters to a friend in Germany for publication in the Augsburg \emph{General Gazette.}

Compare against Quirinus: Die Unwahrheiten der Römischen Briefe vom Concil in der Allg. Zeitung,  Von W. Emmanuel Freiherrn von Ketteler (Bishop of Mayence), 1870.

\emph{Ce qui se passe au Concile.} Dated April 16, 1870. Troisième ed. Paris, 1870. [By Jules Gaillard.]

\emph{La dernière heure du Concile,} Paris, 1870. [By a member of the Council.] The last two works were denounced as a calumny by the presiding Cardinals in the session, July 16,1870.

Also the Reports during the Council in the \emph{Giornale di Roma,} the Turin \emph{Unità catholica,} the London \emph{Times}, the London (R.C.) \emph{Tablet}, the Dublin \emph{Review}, the New York \emph{Tribune}, and other leading periodicals.

III. The Acts and Proceedings of the Council.

(1.) Roman Catholic (Infallibilist) Sources.

\emph{Acta et Decreta sacrosancti et œcumenici Concilii Vaticani die} 8 \emph{Dec.} 1869 \emph{a ss. D. N. Pio IX. inchoati. Cum permissione superiorum}, Friburgi Brisgoviæ, 1871, in 2 Parts. The first part contains the Papal Encyclica with the Syllabus and the acts preparatory to the Council; the second, the public acts of the Council itself, with a list of the dioceses of the Roman Church and the members of the Vatican Council.

\emph{Actes et histoire du Concile œcuménique de Rome, premier du Vatican,} ed. under the auspices of \emph{Victor Frond}, Paris, 1869 sqq. 6 vols. Includes extensive biographies of Pope Pius IX. and his Cardinals, etc., with portraits. Vol. VI. contains the \emph{Actes, decrets et documents reccuillis et mis en ordre par M. Pelletier, chanoine d’Orleans.} Each vol. costs 100 francs.

\emph{Atti ufficialli del Concilio ecumenico,} Turino, pp. 682 (? 1870).

\emph{Officielle Actenstücke zu dem von Sr. Heiligkeit dem Papst Pius IX. nach Rom berufenen Oekumenischen Concil, Zweite Sammlung,} Berlin, 1870.

\emph{Das Oekumemische Concil. Stimmen aus Maria-Laach, Neue Folge.} Freiburg im Breisgau, l870. A series of discussions in defense of the Council by Jesuits (Florian Riess, and K. v. Weber).

Henry Edward Manning (R.C. Archbishop of Westminster): \emph{Petri Privilegium.} \emph{Three Pastoral Letters,} London, 1871. \emph{The True Story of the Vatican Council}, London, 1877.

Bp.  Jos. Fessler (Secretary of the Vatican Council): \emph{Das Vaticanische Concil, dessen äussere Bedeutung und innerer Verlauf,} Wien, 1871.

Eugen Cecconi (Canon at Florence): \emph{Geschichte der allg. Kirchenversammlung im Vatican.} Trans. from the Italian by Dr. W. Molitor. Regensb. 1873 sqq. (Vol. I. contains only the history before the Council.)

The stenographic reports of the speeches of the Council are still locked up in the archives of the Vatican.

(2.) Old Catholic (anti-Infallibilist).

Joh. Friedrich: \emph{Documenta ad illustrandum Concilium Vaticanum anni} 1870, Nördlingen, 1871, In 2 parts. Contains official and unofficial documents bearing on the Council and the various \emph{schemata de fide, de ecclesia,} etc. Compare his \emph{Tagebuch während des Vaticanischen Concils geführt,} above quoted. By the same: \emph{Geschichte des Vaticanischen Concils,} Bonn, 1877. Vol. I. (contains the preparatory history to 1869); Vol. II. 1883.

Joh. Friedrich Ritter von Schulte (Prof. of Canon Law in the University of Prague, now in Bonn): \emph{Das Unfehlbarkeitsdecret vom} 18 \emph{Juli} 1870 . . . \emph{geprüft,} Prag, 1871. Also, \emph{Die Macht der Röm. Päpste über Fürsten, Länder, Völker, Individuen}, etc., Prag, 2d ed. 1871.

\emph{Stimmen aus der katholischen Kirche über die Kirchenfragen der Gegenwart,} München, 1870 sqq. 2 vols. A series of discussions against the Vatican Council, by Döllinger, Huber, Schmitz, Friedrich, Reinkens, and Hötzl.

(3.) Protestant.

Dr.  Emil Friedberg (Prof. of Ecclesiastical Law in Leipzig): \emph{Sammlung der Actenstücke zum ersten Vaticanischen Concil, mit einem Grundriss der Geschichte desselben,} Tübingen, 1872 (pp. 954). Very valuable; contains all the important documents, and a full list of works on the Council.

Theod. Frommann (\emph{Privatdocent} in Berlin): \emph{Geschichte und Kritik des Vaticanischen Concils von} 1869 \emph{und} 1870, Gotha, 1872 (pp. 529).

E. de Pressense (Ref. Pastor in Paris): \emph{Le Concile du Vatican, son histoire et ses conséquences politiques et religieuses,} Paris, 1872. Also in German, by \emph{Fabarius}, Nördlingen, 1872.

L. W. Bacon:  \emph{An Inside View of the Vatican Council}, New York, 1872 (Amer. Tract Society). Contains a translation of Archbishop Kenrick's speech against Infallibility, with a sketch of the Council.

G. Uhlhorn:  \emph{Das Vaticanische Concil} (\emph{Vermischte Vorträge}). Stuttgart, 1875, pp. 235–350.

An extensive criticism on the Infallibility decree in the third edition of Dr.  Hase's  \emph{Handbuch der Protestant. Polemik gegen die römisch-katholische Kirche,} Leipz. 1871, pp. 155–200. Comp. pp. 24–37.

The above are only the most important works of the large and increasing literature, historical, apologetic, and polemic, on the Vatican Council. A. Erlecke, in a pamphlet, \emph{Die Literatur des röm. Concils,} gives a list of over 200 books and pamphlets which appeared in Germany alone before 1871. Friedberg notices 1041 writings on the subject till June 1872. Since then the Gladstone \emph{Expostulation} on the political aspects of the \emph{Vatican Decrees}, Lond. 1874, and his \emph{Vaticanism}, 1875, have called forth a newspaper and pamphlet war, and put Dr. J. H. Newman and Archbishop Manning on the defensive.]

More than three hundred years after the close of the Council of Trent, Pope Pius IX., who had proclaimed the new dogma of the Immaculate Conception, who in the presence of five hundred Bishops had celebrated the eighteenth centennial of the martyrdom of the Apostles Peter and Paul, and who was permitted to survive not only the golden wedding of his priesthood, but even—alone among his more than two hundred and fifty predecessors—the silver wedding of his popedom (thus falsifying the tradition '\emph{non videbit annos Petri}'), resolved to convoke a new œcumenical Council, which was to proclaim his own infallibility in all matters of faith and discipline, and thus to put the top-stone to the pyramid of the Roman hierarchy.

He first intimated his intention, June 26, 1867, in an Allocution to five hundred Bishops who were assembled at the eighteenth centenary of the martyrdom of St. Peter in Rome. The Bishops, in a most humble and obsequious response, July 1, 1867, approved of his heroic courage, to employ, in his old age, an extreme measure for an extreme danger, and predicted a new splendor of the Church, and a new triumph of the kingdom of God.\footnote{'\emph{Summo igitur gauaio,}' said the five hundred Bishops, '\emph{repletus est animus noster, dum sacrato ore Tuo intelleximus, tot inter præsentis temporis discrimina eo Te esse consilio, ut} ``\emph{maximum,}'' \emph{prout aiebat inclitus Tuus prædecessor Paulus III.,} ``\emph{in maximis rei Christianæ periculis remedium,}'' \emph{Concilium œcumenicum convoces. Annuat Deus huic Tuo proposito, cuius ipse Tibi mentem inspiravit; habeantque tandem œvi nostri homines, qui infirmi in fide, semper discentes et nunquam ad veritatis agnitionem pervenientes omni vento doctrinæ circumferuntur, in sacrosancta hac Synodo novam, præsentissimamque occasionem accedendi ad sanctam Ecclesiam columnam ac firmamentum veritatis, cognoscendi salutiferam fidem, perniciosos reiiciendi errores; ac fiat, Deo propitio, et conciliatrice Deipara Immaculata, hæc Synodus grande opus unitatis, sanctificationis et pacis, unde novus in Ecclesiam splendor redundet, novus regni Dei triumphus consequatur. Et hoc ipso Tuæ providentiæ opere denuo exibeatur mundo immensa beneficia, per Pontificatum romanum humanæ societati asserta. Pateat cunctis, Ecclesiam eo quod super solidissima Petra fundetur, tantum valere, ut errores depellat, mores corrigat, barbariem compescat, civilisque humanitatis mater dicatur et sit. Pateat mundo, quod divinæ auctoritatis et debitæ eidem obedientiæ manifestissimo specimine, in divina Pontificatus institutione dato, ea omnia stabilita et sacrata sint, quæ societatum fundamenta ac diuturnitatem solident.}'} Whereupon the Pope announced to them that he would convene the Council under the special auspices of the immaculate Virgin, who had crushed the serpent's head and was mighty to destroy alone all the heresies of the world.\footnote{'\emph{Quod sane votum apertius etiam se prodit in eo communi Concilii œcumenici desiderio, quod omnes non modo perutile, sed et necessarium arbitramini. Superbia enim humana, veterem ansum instauratura, jamdiu per commenticium progressum civitatem et turrem extruere nititur, cujus culmen pertingat ad cœlum, unde demum Deus ipse detrahi possit. At is descendisse videtur inspecturus opus, et ædificantium linguas ita confusurus, ut non audiat unusquisque vocem proximi sui: id enim animo objiciunt Ecclesiæ vexationes, miseranda civilis consortii conditio, perturbatio rerum omnium, in qua versamur. Cui sane gravissimæ calamitati sola certe objici potest divina Ecclesiæ virtus, quæ tunc maxime se prodit, cum Episcopi a Summo Pontifice convocati, eo præside, conveniunt in nomine Domini de Ecclesiæ rebus acturi. Et gaudemus omnino, prœvertisse vos hac in re propositum jamdiu a nobis conceptum, commendandi sacrum hunc cœtum ejus patrocinio, cujus pedi a rerum exordio serpentis caput subjectum fuit, quœque deinde universas hæreses sola interemit. Satisfacturi propterea communi desiderio jam nunc nunciamus, futurum quandocunque Concilium sub auspiciis Deiparæ Virginis ab omni labe immunis esse constituendum, et eo aperiendum die, quo insignis hujus privilegii ipsi collati memoria recolitur. Faxit Deus, faxit Immaculata Virgo, ut amplissimos e saluberrimo isto Concilio fructus percipere valeamus.}' While the Pope complains of the pride of the age in attempting to build another tower of Babel, it did not occur to him that the assumption of infallibility, i.e., a predicate of the Almighty by a mortal man, is the consummation of spiritual pride.}

The call was issued by an Encyclical, commencing \emph{Æterni Patris Unigenitus Filius,} in the twenty-third year of his Pontificate, on the feast of St. Peter and Paul, June 29, 1868. It created at once a universal commotion in the Christian world, and called forth a multitude of books and pamphlets even before the Council convened. The highest expectations were suspended by the Pope and his sympathizers on the coming event. What the Council of Trent had effected against the Protestant Reformation of the sixteenth century, the Council of the Vatican was to accomplish against the more radical and dangerous foes of modern liberalism and rationalism, which threatened to undermine Romanism itself in its own strongholds. It was to crush the power of infidelity, and to settle all that belongs to the doctrine, worship, and discipline of the Church, and the eternal salvation of souls.\footnote{After describing, in the stereotyped phrases of the Roman Court, the great solicitude of the successors of Peter for pure doctrine and good government, and the terrible tempests and calamities by which the Catholic Church and the very foundations of society are shaken in the present age, the Pope's Encyclical comprehensively but vaguely, and with a prudent reserve concerning the desired dogma of Infallibility, defines the objects of the Council in these words: '\emph{In œcumenico hoc Concilio ea omnia accuratissime examine sunt perpendenda ac statuenda, quæ hisce præsertim asperrimis temporibus majorem Dei gloriam, et fidei integritatem, divinique cultus decorem, sempiternamque hominum salutem, et utriusque Cleri disciplinam ejusque salutarem solidamque culturam, atque ecclesiasticarum legum observantiam, morumque emendationem, et christianam juventutis institutionem, et communem omnium pacem et concordiam in primis respiciunt. Atque etiam intentissimo studio curandum est, ut, Deo bene juvante, omnia ab Ecclesia et civili societate amoveantur mala, ut miseri errantes ad rectum veritatis, justitiæ salutisque tramitem reducantur, ut vitiis erroribusque eliminatis, augusta nostra religio ejusque salutifera doctrina ubique terrarum reviviscat, et quotidie magis propagetur et dominetur, atque ita pietas, honestas, probitas, justitia, caritas omnesque Christianæ virtutes cum maxima humanæ societatis utilitate vigeant et efflorescant.}'} It was even hoped that the Council might become a general feast of reconciliation of divided Christendom; and hence the Greek schismatics, and the Protestant heretics and other non-Catholics, were invited by two special letters of the Pope (Sept. 8, and Sept. 13, 1868) to return on this auspicious occasion to 'the only sheepfold of Christ,' for the salvation of their souls.\footnote{'\emph{Omnes Christianos etiam atque etiam hortamur et obsecramus, ut ad unicum Christi ovile redire festinent.}' And at the end again, '\emph{unum ovile et unus pastor};' according to the false and mischievous translation of John x. 16 in the Vulgate (followed by the authorized English Version), instead of '\emph{one flock}' (\textgreek{μία ποίμνη,} not \textgreek{αὐλή}). There may be many folds, and yet one flock under one Shepherd, as there are 'many mansions' in heaven (John xiv. 2).}

But the Eastern Patriarchs spurned the invitation, as an insult to their time-honored rights and traditions, from which they could not depart.\footnote{The Patriarch of Constantinople declined even to receive the Papal letter from the Papal messenger, for the reasons that it had already been published in the \emph{Giornale di Roma}; that it contained principles contrary to the spirit of the Gospel, the doctrines of the œcumenical Councils, and the holy Fathers; that there was no supreme Bishop in the Church except Christ; and that the Bishop of Old Rome had no right to convoke an œcumenical Council without first consulting the Eastern Patriarchs. The other Oriental Bishops either declined or returned the Papal letter of invitation. See the documents in Friedberg, l.c. pp. 233–253; in Officielle Actenstücke, etc., pp. 127–135; and in the Chronique concernant le Prochain Concile, Vol. I. pp. 3 sqq., 103 sqq.} The Protestant communions either ignored or respectfully declined it.\footnote{The Evangelical \emph{Oberkirchenrath} of Berlin, the \emph{Kirchentag} of Stuttgart, 1869, the Paris Branch of the Evangelical Alliance, 'The Venerable Company of Pastors of Geneva,' the Professors of the University of Groningen, the Hungarian Lutherans assembled at Pesth, and the Presbyterians of the United States, took notice of the Papal invitation, all declining it, and reaffirming the principles of the Protestant Reformation. The Presbyterian Dr. Cumming, of London, seemed willing to accept the invitation if the Pope would allow a discussion of the reasons of the separation from Rome, but was informed by the Pope, through Archbishop Manning, in two letters (Sept. 4, and Oct. 30, 1869), that such discussion of questions long settled would be entirely inconsistent with the infallibility of the Church and the supremacy of the Holy See. See the documents in Friedberg, pp. 235–257; comp. pp. 16, 17, and Offic. Actenstücke, pp. 158–176. The Chronique concernant le Prochain Concile, p. 169, criticises at length the American Presbyterian letter signed by Jacobus and Fowler (Moderators of the General Assembly), and sees in its reasons for declining a proof of 'heretical obstinacy and ignorance.'}

Thus the Vatican Council, like that of Trent, turned out to be simply a general Roman Council, and apparently put the prospect of a reunion of Christendom farther off than ever before.

While these sanguine expectations of Pius IX. were doomed to disappointment, the chief object of the Council was attained in spite of the strong opposition of the minority of liberal Catholics. This object, which for reasons of propriety is omitted in the bull of convocation and other preliminary acts, but clearly stated by the organs of the Ultramontane or Jesuitical party, was nothing less than the proclamation of the personal \emph{Infallibility of the Pope}, as a binding article of the Roman Catholic faith for all time to come.\footnote{So the Civiltà cattolica (a monthly Review established 1850, at Rome, the principal organ of the Jesuits, and the \emph{Moniteur} of the Papal Court) defined the programme, Feb. 6, 1869; adding to it also the adoption of the Syllabus of 1864, and, perhaps, the proclamation of the assumption of the Virgin Mary to heaven. The last is reserved for the future. The Archbishop of Westminster (Manning) and the Archbishop of Mechlin (Dechamps) predicted, in pastoral letters of 1867 and 1869, the proclamation of the Papal Infallibility as a certain event. To avert this danger, the Bishop of Orleans (Dupanloup), Père Gratry of the Oratory, Père Hyacinthe, Bishop Maret (Dean of the Theological Faculty of Paris), Montalembert, John Henry Newman, the German Catholic laity (in the Coblenz Address), in part the German Bishops assembled at Fulda, and especially the learned authors of the \emph{Janus}, lifted their voice, though in vain. See the literature on the subject in Friedberg, pp. 17–21.} Herein lies the whole importance of the Council; all the rest dwindles into insignificance, and could never have justified its convocation.

After extensive and careful preparations, the first (and perhaps the last) Vatican Council was solemnly opened amid the sound of innumerable bells and the cannon of St. Angelo, but under frowning skies and a pouring rain, on the festival of the Immaculate Conception of the Virgin Mary, Dec. 8, 1869, in the Basilica of the Vatican.\footnote{Hence the name. The right cross-nave of St. Peter's Church, which itself is a large church, was separated by a painted board wall, and fitted up as the council-hall. See a draught of it in Friedberg, p. 98. The hall was very unsuitable for hearing, and had to be repeatedly altered. The Pope, it is said (Hase, l.c. p. 26), did not care that all the orators should be understood. The Vatican Palace, where the Pope now resides, adjoins the Church of St. Peter. Councils were held there before, but only of a local character. Formerly the Roman œcumenical Councils were held in the Lateran Palace, the ancient residence of the Popes, which is connected with the Church of St. John in the Lateran or Church of the Saviour ('\emph{omnium urbis et orbis ecclesiarum mater et caput}'). There are five Lateran Councils: the first was held, 1123, under Calixtus II.; the second, 1139, under Innocent II.; the third, 1179, under Alexander III.; the fourth and largest, 1215, under Innocent III.; the fifth, 1512–1517, under Leo X., on the eve of the Reformation. The basilica of the Lateran contains the head, the basilica of St. Peter the body, of St. Peter. The Pope expressed the hope that a special inspiration would proceed from the near grave of the prince of the Apostles upon the Fathers of the Council.} It reached its height at the fourth public session, July 18, 1870, when the decree of Papal Infallibility was proclaimed. After this it dragged on a sickly existence till October 20, 1870, when it was adjourned till Nov. 11, 1870, but indefinitely postponed on account of the extraordinary change in the political situation of Europe. For on the second of September the French Empire, which had been the main support of the temporal power of the Pope, collapsed with the surrender of Napoleon III., at the old Huguenot stronghold of Sedan, to the Protestant King William of Prussia, and on the twentieth of September the Italian troops, in the name of King Victor Emanuel, took possession of Rome, as the future capital of united Italy. Whether the Council will ever be convened again to complete its vast labors, like the twice interrupted Council of Trent, remains to be seen. But, in proclaiming the personal Infallibility of the Pope, it made all future œcumenical Councils unnecessary for the definition of dogmas and the regulation of discipline, so that hereafter they will be expensive luxuries and empty ritualistic shows. The acts of the Vatican Council, as far as they go, are irrevocable.

The attendance was larger than that of any of its eighteen predecessors,\footnote{As the œcumenical character of two or three Councils is disputed, the Vatican Council is variously reckoned as the 19th or 20th or 21st œcumenical Council; by strict Romanists (as Manning) as the 19th. Compare note on p. 91.} and presented an imposing array of hierarchical dignity and power such as the world never saw before, and as the Eternal City itself is not likely ever to see again. What a contrast this to the first Council of the apostles, elders, and brethren in an upper chamber in Jerusalem! The whole number of prelates of the Roman Catholic Church, who are entitled to a seat in an œcumenical Council, is one thousand and thirty-seven.\footnote{See a full list, with all the titles, in the \emph{Lexicon geographicum} added to the second part of the Acta et Decreta sacrosancti et œcum. Conc. Vaticani, Friburgi, 1871. The Prelates '\emph{quibus aut jus aut privilegium fuit sedendi in œcumenica synodo Vaticana,}' are arranged as follows:  \par (1.) Eminentissimi et reverendissimi Domini S.E. Rom. Cardinales: (\emph{a}) ordinis Episcoporum, (\emph{b}) ordinis Presbyterorum. (\emph{c}) ordinis diaconorum—51. \par (2.) Reverendissimi Domini Patriarchæ—11. \par (3.) Reverendissimi DD. Primates—10. \par (4.) Reverendissimi DD. Archiepiscopi—166. \par (5.) Reverendissimi DD. Episcopi—740. \par (6.) Abbates nullius dioceseos—6. \par (7.) Abbates Generales ordinum monasticorum—23. \par (8.) Generales et Vicarii Generales congregationum clericorum regularium, ordinum monasticorum, ordinum mendicantium—29. In all, 1037.} Of these there were present at the opening of the Council 719, viz., 49 Cardinals, 9 Patriarchs, 4 Primates, 121 Archbishops, 479 Bishops, 57 Abbots and Generals of monastic orders.\footnote{See the list of names in Friedberg, pp. 376–394.} This number afterwards increased to 764, viz., 49 Cardinals, 10 Patriarchs, 4 Primates, 105 diocesan Archbishops, 22 Archbishops in partibus infidelium, 424 diocesan Bishops, 98 Bishops in partibus, and 52 Abbots, and Generals of monastic orders.\footnote{See the official \emph{Catalogo alfabetico dei Padri presenti al Concilio ecumenico Vaticano}, Roma, 1870.} Distributed according to continents, 541 of these belonged to Europe, 83 to Asia, 14 to Africa, 113 to America, 13 to Oceanica. At the proclamation of the decree of Papal Infallibility, July 18, 1870, the number was reduced to 535, and afterwards it dwindled down to 200 or 180.

Among the many nations represented,\footnote{Manning says, 'some thirty nations'—probably an exaggeration.} the Italians had a vast majority of 276, of whom 143 belonged to the former Papal States alone. France, with a much larger Catholic population, had only 84, Austria and Hungary 48, Spain 41, Great Britain 35, Germany 19, the United States 48, Mexico 10, Switzerland 8, Belgium 6, Holland 4, Portugal 2, Russia 1. The disproportion between the representatives of the different nations and the number of their constituents was overwhelmingly in favor of the Papal influence. Nearly one half of the Fathers were entertained during the Council at the expense of the Pope.

The Romans themselves were remarkably indifferent to the Council, though keenly alive to the financial gain which the dogma of the Infallibility of their sovereign would bring to the Eternal City and the impoverished Papal treasury.\footnote{Quirinus, pp. 480, 481 (English translation).} It is well known, how soon after the Council they voted almost in a body against the temporal power of the Pope, and for their new master.

The strictest secresy was enjoined upon the members of the Council.\footnote{They had to promise and swear to observe '\emph{ inviolabilem secreti fidem}' with regard to the discussions, the opinions, and all matters pertaining to the Council. See the form of the oath in Friedberg, p. 96. In ancient Councils the people are often mentioned as being present during the deliberations, and manifesting their feelings of approval and disapproval.} The stenographic reports of the proceedings were locked up in the archives. The world was only to know the final results as proclaimed in the public sessions, until it should please the Roman court to issue an official history. But the freedom of the press in the nineteenth century, the elements of discord in the Council itself, the enterprise or indiscretion of members and friends of both parties, frustrated the precautions. The principal facts, documents, speeches, plans, and intrigues leaked out in the official \emph{schemata,} the controversial pamphlets of Prelates, and the private reports and letters of outside observers who were in intimate and constant intercourse with their friends in the Council.\footnote{Among the irresponsible but well-informed reporters and correspondents must be mentioned especially the writers in the Civiltà cattolica, and the Paris Univers, on the part of the Infallibilists; and the pseudonymous Quirinus, Prof. Friedrich, and the anonymous French authors of Ce qui se passe au Concile, and of La dernière heure du Concile, on the part of the anti-Infallibilists.}

The subject-matter for deliberation was divided into four parts: on Faith, Discipline, Religious Orders, and on Rites, including Missions. Each part was assigned to a special Commission (\emph{Congregatio} or \emph{Deputatio}), consisting of 24 Prelates elected by ballot for the whole period of the Council, with a presiding Cardinal appointed by the Pope. These Commissions prepared the decrees on the basis of \emph{schemata} previously drawn up by learned divines and canonists, and confidentially submitted to the Bishops in print.\footnote{There were in all forty-five \emph{schemata}, divided into four classes: (1) \emph{circa fidem}, (2) \emph{circa disciplinam ecclesiæ}, (3) \emph{circa ordines regulares}, (4) \emph{circa res ritus orientalis et apostolicas missiones.} See a list in Friedberg, pp. 432–434. Only a part of the \emph{schemata} were submitted, and only the first two \emph{schemata de fide} were acted upon. Friedrich, in the Second Part of his Documenta, gives the \emph{schemata,} as far as they were distributed among the Bishops, together with the revisions and criticisms of the Bishops.} The decrees were then discussed, revised, and adopted in secret sessions by the General Congregation (\emph{Congregationes generales}), including all the Fathers, with five presiding Cardinals appointed by the Pope. The General Congregation held eighty-nine sessions in all. Finally, the decrees thus matured were voted upon by simple \emph{yeas} or \emph{nays} (\emph{Placet} or \emph{Non Placet}), and solemnly promulgated in public sessions in the presence and by the authority of the Pope. A conditional assent (\emph{Placet juxta modum}) was allowed in the secret, but not in the public sessions.

There were only four such public sessions held during the ten months of the Council, viz., the opening session (lasting nearly seven hours), Dec. 8, 1869, which was a mere formality, but of a ritualistic splendor and magnificence such as can be gotten up nowhere on earth but in St. Peter's Cathedral in Rome; the second session, Jan. 6, 1870, when the Fathers simply professed each one before the Pope the Nicene Creed and the Profession of the Tridentine Faith; the third session, April 24, 1870, when the dogmatic constitution on the Catholic faith was unanimously adopted; and the fourth session, July 18, 1870, when the first dogmatic constitution on the Church of Christ and the Infallibility of the Pope was adopted with two dissenting votes.

The management of the Council was entirely in the hands of the Pope and his dependent Cardinals and Jesuitical advisers. He originated the topics which were to be acted on; he selected the preparatory committees of theologians (mostly of the Ultramontane school) who, during the winter of 1868–69, drew up the \emph{schemata}; he appointed the presiding officers of the four Deputations, and of the General Congregation; and he proclaimed the decrees in his own name, 'with the approval of the Council.'\footnote{Under the title: \emph{Pius episcopus, servus servorum Dei, sacro approbante Concilio, ad perpetuam rei memoriam.} The order prescribed for voting was this: The Pope, through the Secretary, asked the members of the Council first in general: \emph{Reverendissimi Patres, placentne vobis Decreta et Canones qui in hac Constitutione continentur?} Then each one was called by name, and must vote either \emph{placet} or \emph{non placet.} When the votes were collected and brought to the Pope, he announced the result by this formula: \emph{Decreta et Canones qui in Constitutione modo lecta continentur, placuerunt omnibus Patribus, nemine dissentiente} [if there were dissenting votes the Pope stated their number]; \emph{Nosque, sacro approbante Concilio, illa} [\emph{sc. decreta}] \emph{et illos} [\emph{canones}], \emph{ita ut lecta sunt, definimus, et Apostolica Auctoritate confirmamus.} See the \emph{Monitum} in the Giornale di Roma, April 18, 1870; Friedberg, pp. 462–464.} He provided, by the bull '\emph{Cum, Romanis Pontificibus,}' of Dec. 4, 1869, for the immediate suspension and adjournment of the Council in case of his death. He even personally interfered during the proceedings in favor of his new dogma by praising Infallibilists, and by ignoring or rebuking anti-Infallibilists.\footnote{See the laudatory letters of Pius to several advocates of Infallibility, in Friedberg, pp. 487–495; comp. pp. 108–111. To Archbishop Dechamps, of Mechlin, he wrote that, in his tract on Papal Infallibility, he had proved the harmony of the Catholic faith with human reason so convincingly as to force even the Rationalists to see the absurdity of the opposite views. He applauded the indefatigable and abusive editor of the Paris \emph{Univers}, Veuillot, who had collected 100,000 francs for the Vicar of Christ (May 30, 1870). On the other hand, he is reported to have rebuked in conversation Cardinal Schwarzenberg by the remark: 'I, John Maria Mastai, believe in the infallibility of the Pope. As Pope I have nothing to ask from the Council. The Holy Ghost will enlighten it.' He even attacked the memory of the eloquent French champion of Catholic interests, the Count Montalembert, who died during the Council (March 13, 1870), by saying, in the presence of three hundred persons: 'He had a great enemy, pride. He was a liberal Catholic, i.e., a half Catholic.' \emph{Ce qui se passe au Concile,} 154 sqq.} The discussion could be virtually arrested by the presiding Cardinals at the request of only ten members; we say virtually, for although it required a vote of the Council, a majority was always sure. The revised order of business, issued Feb. 22, 1870, departed even from the old rule requiring absolute or at least moral unanimity in definitions of faith (according to the celebrated canon \emph{quod semper, quod ubique, quod ab omnibus creditum est}), and substituted for it a mere numerical majority, in order to secure the triumph of the Infallibility decree in spite of a powerful minority. Nothing could be printed in Rome against Infallibility, while the organs of Infallibility had full freedom to print and publish what they pleased.\footnote{Several minority documents, as Kenrick's speech against Infallibility, and the Latin edition of Hefele's tract on Honorius, were printed in Naples; the German in Tübingen. But the \emph{Civiltà cattolica}, the irresponsible organ of the Jesuits and the Pope, was provided with a special building and income, and every facility for obtaining information. See Acton, Quirinus, and Frommann (1.c. p. 13).} Such prominence of the Pope is characteristic of a Council convoked for the very purpose of proclaiming his personal infallibility, but is without precedent in history (except in some mediæval Councils); even the Council of Trent maintained its own dignity and comparative independence by declaring its decrees in its own name.\footnote{' \emph{Sacrosancta Tridentina Synodas, in Spiritu Sancto legitime congregata . . . declarat.}' See the order of the Council of Trent as republished in Friedrich's Documenta, I. pp. 265 sqq.}

This want of freedom of the Council—not to speak of the strict police surveillance over the members—was severely censured by liberal Catholics. More than one hundred Prelates of all nations signed a strong protest (dated Rome, March 1, 1870) against the order of business, especially against the mere majority vote, and expressed the fear that in the end the authority of this Council might be impaired as wanting in truth and liberty—a calamity so direful in these uneasy times, that a greater could not be imagined. But this protest, like all the acts of the minority, was ignored.

The proceedings were, of course, in the official language of the Roman Church, which all Prelates could understand and speak, but very few with sufficient ease to do justice to themselves and their subjects. The acoustic defects of the Council-hall and the difference of pronunciation proved a great inconvenience, and the Continentals complained\footnote{'\emph{Id autem, quod spectat ad numerum suffragiorum requisitum, ut quæstiones dogmaticæ solvantur, in quo quidem rei summa est totiusque Concilii cardo vertitur, ita grave est, ut nisi admitteretur, quod reverenter et enixe postulamus, conscientia nostra intolerabili pondere premeretur: timeremus, ne Concilii œcumenici character in dubium vocari posset; ne ansa hostibus prœberetur Sanctam Sedem et Concilium impetendi, sicque demum apud populum Christianum hujus Concilii auctoritas labefactaretur, quasi veritate et libertate caruerit: quod his turbatissimis temporibus tanta esset calamitas, ut pejor excogitari nulla possit.}' See the remarkable protest in Friedberg, pp. 417–422. Also Dollinger's critique of the order of business, ib. 422–432; Archbishop Kenrick's famous \emph{concio habenda at non habita,} published in Naples, 1870 (and republished in Friedrich's \emph{Docum}.); the work \emph{La libertè du Concile et l’infaillibilité,} which was either written or inspired by Archbishop Darboy, of Paris (in Friedrich's \emph{Docum}. I. pp. 129 sqq.), and the same Prelate's speech in the General Congregation, May 20, 1870 (\emph{ibidem.} II. pp. 415 sqq.). Archbishop Manning, sublimely ignoring all these facts and documents, and referring us to the inaccessible Archives of the Vatican, assures us (\emph{Petri Privil.} III. 32) that the Council was as free as the Congress of the United States, and that the wonder is, not that the opposition failed of its object, but that the Council so long held its peace.} that they could not understand the English Latin. The Council had a full share of ignorance and superstition,\footnote{Some amusing examples are reported by the well-informed Quirinus. Bishop Pie, of Poitiers, supported the Papal Infallibility in a session of the General Congregation (May 13) by an entirely original argument derived from the legend that Peter was crucified downward; for as his head bore the whole weight of the body, so the Pope, as the head, bears the whole Church; but he is infallible who bears, not he who is borne! The Italians and Spaniards applauded enthusiastically. Unfortunately for the argument, the head of Peter did not bear his body, but the cross bore both; consequently the cross must be infallible. A Sicilian Prelate said the Sicilians first doubted the infallibility of Peter when he visited the island, and sent a special deputation of inquiry to the Virgin Mary, but were assured by her that she remembered well having been present when Christ conferred this prerogative on Peter; and this satisfied them completely. Quirinus adds: 'The opposition Bishops see a proof of the insolent contempt of the majority in thus putting up such men as Pie and this Sicilian to speak against them.' \emph{Letter XLVI.} p. 534.} and was disgraced by intrigues and occasional outbursts of intolerance and passion such as are, alas! not unusual in deliberative assemblies even of the Christian Church.\footnote{The following characteristic episode (ignored, of course, in Manning's eulogy) is well authenticated by the concurrent and yet independent reports of Lord Acton (\emph{N. Brit. Rev.}), Quirinus (\emph{Letter XXXII.}), Friedrich (\emph{Tagebuch}, pp. 271, 272), and the author of Ce qui se passe au Concile (p. 69); comp. Friedberg (pp. 104–106). When Bishop Strossmayer, the boldest member of the opposition and an eloquent Latinist, in a session of the General Congregation (March 22), spoke favorably of the great Leibnitz, and paid Protestants the poor compliment of honesty (quoting from St. Augustine: '\emph{Errant, sed bona fide errant}'), he was interrupted by the bell of the President (De Angelis) and his rebuke, 'This is no place for praising Protestants' ('\emph{hicce non est locus laudandi Protestantes}')! Very true, for the Council-hall was only a hundred paces from the Palace of the Inquisition. When, resuming, the speaker ventured to attack the principle of deciding questions of faith by mere majorities, he was more loudly interrupted from all sides by confused exclamations: 'Shame! shame! down with the heretic!' ('\emph{Descendat ab ambone! Descendat! Hæreticus! Hæreticus! Damnamus eum! Damnamus!}') 'Several Bishops sprang from their seats, rushed to the tribune, and shook their fists in the speaker's face' (Quirinus, p. 387). When one Bishop (Place, of Marseilles) interposed, '\emph{Ego non damno!}' the cry was raised with increased fury: '\emph{Omnes, omnes illum damnamus! damnamus!}' Strossmayer was forced by the uproar and the continued ringing of the bell to quit the tribune, but did so with a triple '\emph{Protestor.}' The noise was so great that it could be heard in the interior of St. Peter's. Some thought the Garibaldians had broken in; others that Infallibility had been proclaimed, and shouted, according to their opposite views, either 'Long live the infallible Pope!' or 'Long live the Pope, but not the infallible one' (comp. Quirinus, and \emph{Ce qui se passe}, p. 69). Quirinus says that the scene, 'for dramatic force and theological significance, exceeded almost any thing in the past history of Councils' (p. 386), and that a Bishop of the United States said afterwards, 'not without a sense of patriotic pride, that he knew now of one assembly still rougher than the Congress of his own country' (p. 388). Similar scenes of violence occurred in the œcumenical Councils of Ephesus and Chalcedon, but Christian civilization ought to have made some progress since the fifth century.} But it embraced also much learning and eloquence, especially on the part of the French and German Episcopate. Upon the whole, it compares favorably, as to intellectual ability, moral character, and far-reaching effect, with preceding Roman Councils, and must be regarded as the greatest event in the history of the Papacy since the Council of Trent.

The chief importance of the Council of the Vatican lies in its decree on Papal supremacy and Infallibility. It settled the internal dissensions between Ultramontanism and Gallicanism, which struck at the root of the fundamental principle of authority; it destroyed the independence of the Episcopate, and made it a tool of the Primacy; it crushed liberal Catholicism; it completed the system of Papal absolutism; it raised the hitherto disputed opinion of Papal Infallibility to the dignity of a binding article of faith, which no Catholic can deny without loss of salvation. The Pope may now say not only, 'I am the tradition' (\emph{La tradizione son’ io}), but also, 'I am the Church' (\emph{L’église c’est moi})!

But this very triumph of absolutism marks also a new departure. It gave rise to a secession headed by the ablest divines of the Roman Church. It put the Papacy into direct antagonism to the liberal tendencies of the age. It excited the hostility of civil government in all those countries where Church and State are united on the basis of a concordat with the Roman See. No State with any degree of self-respect can treat with a sovereign who claims infallibility, and therefore unconditional submission in matters of moral duty as well as of faith. In reaching the summit of its power, the Papacy has hastened its downfall.

For Protestants and Greeks the Vatican Council is no more œcumenical than that of Trent, and has only intensified the antagonism. Its œcumenicity is also denied by the Old Catholic scholars—Döllinger, von Schulte, and Reinkens —because it lacked the two fundamental conditions of liberty of discussion, and moral unanimity of suffrage.\footnote{See the Old Catholic protests of the Professors in Munich and Breslau in Friedberg, pp. 152–154, and the literature on the reception of the Council, ib. 53–56; also the discussion of Frommann, pp. 325 sqq. 454 sqq. Döllinger, in his famous censure of the new order of the Council, takes the ground that the œcumenicity of a Council depends upon an authority outside of itself, viz., the public opinion as expressed in the subsequent approval of the whole Church; and Pater Hötzl laid down the principle that no Council is œcumenical which is not approved and adopted as such by the Church. Admitting this, the condition is now fulfilled in the case of the Vatican Council to the whole extent of the Roman Episcopate, which constitutes the \emph{ecclesia docens}, the laity having nothing to do but to submit.} But the subsequent submission of all the Bishops who had voted against Papal Infallibility, supplies the defect as far as the Roman Church is concerned. There was nothing left to them but either to submit or to be expelled. They chose the former, and thus destroyed the legal and moral force of their protest, although not the power of truth and the nature of the facts on which it was based. Henceforward Romanism must stand or fall with the Vatican Council. But (as we have before intimated) Romanism is not to be confounded with Catholicism any more than the Jewish hierarchy which crucified our Saviour, is identical with the people of Israel, from which sprang the Apostles and early converts of Christianity. The destruction of the infallible and irreformable Papacy may be the emancipation of Catholicism, and lead it from its prison-house to the light of a new Reformation.

\xsection{The Vatican Decrees. The Constitution of the Catholic Faith.}

§ 32. The Vatican Decrees. The Constitution on the Catholic Faith.

Three schemes on matters of faith were prepared for the Vatican Council—one against Rationalism, one on the Church of Christ, and one on Christian Matrimony. The first two were revised and adopted; the third was indefinitely postponed. There was also much discussion on the preparation of a small popular Catechism adapted to the present doctrinal status of the Roman Church, and intended to supersede the numerous popular Catechisms now in use; but the draft, which assigned the whole teaching power of the Church to the Pope, to the exclusion of the Episcopate, encountered such opposition (57 \emph{Non Placet}, 24 conditional \emph{Placet}) in the provisional vote of May 4, that it was laid on the table and never called up again.\footnote{Cardinal-Archbishop Matthieu of Besançon, who voted \emph{Non Placet}, is reported by Quirinus to have said on this occasion: '\emph{On veut jeter l’église dans I’abîme, nous y jeterons plutôt nos cadavres.}' Comp. Frommann, l.c. p. 160.}

I. The Dogmatic Constitution on the Catholic Faith (Constitutio Dogmatica de Fide Catholica).

It was unanimously adopted in the third public session, April 24 (\emph{Dominica in albis}), 1870.

The original draft laid before the Council embraced eighteen chapters—on Pantheism, Rationalism, Scripture and tradition, revelation, faith and reason, the Trinity, the two natures of Christ, the primitive state, original sin, the Christian redemption, the supernatural order of grace; but was laid aside.\footnote{Friedrich, \emph{Docum.} II. pp. 3–23.} Archbishop Connolly, of Halifax, recommended that it should be decently buried.\footnote{'\emph{Censeo schema cum honore esse sepeliendum}' (Quirinus, p. 122). Rauscher also spoke against the schema, which made much impression, because he had brought its chief author, the Jesuit Schrader, to the University of Vienna.}

In its present form, the Constitution on the Catholic faith is reduced to four chapters, with a proemium and a conclusion. Chap. I. treats of God as the Creator; Chap. II. of revelation; Chap. III. of faith; Chap. IV. of faith and reason. Then follow 18 canons, in which the errors of Pantheism, Naturalism, and Rationalism are condemned in a manner substantially the same, though more clearly and fully, than had been done in the first two sections of the Syllabus.

The decree asserts, in the old scholastic terminology, the well-known principles of Supernaturalism as held by orthodox Christians in all ages, but it completely ignores the freedom and progress of theological and philosophical science and learning since the Council of Trent, and it forbids (in Chap. II.) all interpretation of the Scriptures which does not agree with the Romish traditions, the Latin Vulgate, and the fictitious 'unanimous consent of the Fathers.' Hence a liberal member of the Council, in the course of discussion, declared the \emph{schema de fide} a work of supererogation. 'What boots it,' he said, 'to condemn errors which have been long condemned, and tempt no Catholic? The false beliefs of mankind are beyond the reach of your decrees. The best defense of Catholicism is religious science. Encourage sound learning, and prove by deeds as well as words that it is the mission of the Church to promote among the nations liberty, light, and true prosperity.'\footnote{Quoted in Latin by Lord Acton in the \emph{North British Review}, Oct. 1870, p. 112, and in Friedberg, p. 102. Acton attributes this speech, not to Strossmayer (as Friedberg says, l.c.; comp. pp. 28 and 102), but to a 'Swiss prelate,' whom he does not name.} On the other hand, the Univers calls the schema a 'masterpiece of clearness and force;' the Civiltà cattolica sees in it 'a reflex of the wisdom of God;'\footnote{'\emph{Un riverbero della sapienza di Dio},' VII. 10, p. 523, quoted by Frommann, l.c. p. 383.} and Archbishop Manning thinks that its importance 'can not be overestimated,' that it is 'the broadest and boldest affirmation of the supernatural and spiritual order ever yet made in the face of the world, which is now more than ever sunk in sense and heavy with Materialism.'\footnote{\emph{Petri Privilegium}, III. pp. 49, 50.} Whatever be the value of the positive principles of the schema, its Popish head and tail reduce it to a \emph{brutum fulmen} outside of the Romish Church, and even the most orthodox Protestants must apply to it the warning, \emph{Timeo Danaos et dona ferentes.}

The preamble, even in its present modified form, derives modern Rationalism and infidelity, as a legitimate fruit, from the heresies condemned by the Council of Trent—that is, from the Protestant Reformation; in the face of the fact, patent to every scholar, that Protestant theology has been in the thickest of the fight with unbelief, and, notwithstanding all its excesses, has produced a far richer exegetical and apologetic literature than Romanism during the last three hundred years.\footnote{The objectionable passage, as finally adopted, reads thus: 'No one is ignorant that the heresies proscribed by the Fathers of Trent, by which the divine magisterium of the Church was rejected, and all matters regarding religion were surrendered to the judgment of each individual, gradually became dissolved into many sects, which disagreed and contended with one another, until at length not a few lost all faith in Christ. Even the Holy Scriptures, which had previously been declared the sole source and judge of Christian doctrine, began to be held no longer as divine, but to be ranked among the fictions of mythology. Then there arose, and too widely overspread the world, that doctrine of Rationalism which opposes itself in every way to the Christian religion as a supernatural institution.' See the different revisions of the \emph{schema de fide} in Friedrich's \emph{Monum.} Pt. II. pp. 3, 65, 73.} The boldest testimony heard in the Council was directed against this preamble by Bishop Strossmayer, from the Turkish frontier (March 22, 1870). He characterized the charge against Protestantism as neither just nor charitable. Protestants, he said, abhorred the errors condemned in the schema as much as Catholics. The germ of Rationalism existed in the Catholic Church before the Reformation, especially in the humanism which was nourished in the very sanctuary by the highest dignitaries,\footnote{Allusion to Pope Leo X.} and bore its worst fruits in the midst of a Catholic nation at the time of Voltaire and the Encyclopedists. Catholics had produced no better refutation of the errors enumerated in the schema than such men as Leibnitz and Guizot. There were multitudes of Protestants in Germany, England, and North America who loved our Lord Jesus Christ, and had inherited from the shipwreck of faith positive truths and monuments of divine grace.\footnote{See the principal part of Strossmayer's speech in Latin in Lord Acton's article in the \emph{North British Review}, Oct. 1870, pp. 115, 116, and in Friedberg, pp. 104–106.} Although this speech was greeted with execrations (see page 145), it had at least the effect that the objectionable preamble was somewhat modified.\footnote{The words in the first revision (Friedr. \emph{Docum.} II. p. 65), \emph{systematum monstra, mythismi, rationalismi, indifferentismi nomine designata, }etc., together with some other offensive expressions, were omitted; but, after all, the substance remained. Lord Acton relates that the German Jesuit Kleutgen hastily drew up the more moderate form. Comp. Quirinus, \emph{Letter XXXIII.} p. 394 sq. Political influence was also brought to bear indirectly upon the Council, as appeared afterwards from Italian papers. Bismarck directed the German Embassador at Rome, Count Arnim, to inform Cardinal Antonelli that, unless the charge against Protestantism was withdrawn, he would not allow the Prussian Bishops on their return to resume their functions in a country whose faith they had insulted. Friedrich, \emph{Tagebuch}, pp. 275, 292; Frommann, \emph{Geschichte des Vat. Concils}, p. 145; Hase, \emph{Polem.} p. 34. The latter overestimates the influence of Prussia on the Papal court when he says: 'If France complains of the Council, Antonelli makes three bows, and all remains as before; but if Prussia comes with her mustache and cavalry boots, Rome understands that the word is quickly followed by the deed, and wisely yields. Strossmayer and von Arnim were in doubt which one of them had been most instrumental in saving the Council from an impropriety.'}

The supplement of the decree binds all Catholics to observe also those constitutions and decrees by which such erroneous opinions as are not here specifically enumerated have been proscribed and condemned by the Holy See. This can be so construed as to include all the eighty errors of the Syllabus. The minority who in the General Congregation had voted \emph{Non Placet} or only a conditional \emph{Placet}, were quieted by the official assurance that the addition involved no new dogma, and had a disciplinary rather than a didactic character. 'Some gave their votes with a heavy heart, conscious of the snare.' Strossmayer stayed away. Thus a unanimous vote of 667 or 668 fathers was secured in the public session, and the Infallibility decree was virtually anticipated. The Pope, after proclaiming the dogma, gave the Bishops his benediction of peace, and gently intimated what he next expected from them.\footnote{'\emph{Videtis,}' he said, '\emph{Fratres carissimi, quam bonum sit et jucundum ambulare in domo Dei cum consensu, ambulare cum pace. Sic ambuletis semper. Et quoniam hac die Dominus Noster Jesus Christus dedit pacem Apostolis suis, et ego, Vicarius ejus indignus, nomine suo do vobis pacem. Pax ista, prout scitis, expellit timorem. Pax ista, prout scitis, claudit aures sermonibus imperitis. Ah! ista pax vos comitetur omnibus diebus vitæ vestræ; sit ista pax vis in morte, sit ista pax vobis gaudium sempiternum in cœlis.}'}

\xsection{The Vatican Decrees, Continued. The Papal Infallibility Decree.}

§ 33. The Vatican Decrees, Continued. The Infallibility Decree.

II. The First Dogmatic Constitution on the Church of Christ (constitutio dogmatica prima de ecclesia Christi).

It was passed, with two dissenting votes, in the fourth public session, July 18, 1870. It treats, in four chapters—(1) on the institution of the Apostolic Primacy in the blessed Peter; (2) on the perpetuity of St. Peter's Primacy in the Roman Pontiff; (3) on the power and nature of the Primacy of the Roman Pontiff; (4) on the Infallibility of the Roman Pontiff.

The new features are contained in the last two chapters, which teach \emph{Papal Absolutism} and \emph{Papal Infallibility}. The third chapter vindicates to the Roman Pontiff a superiority of \emph{ordinary} episcopal (not simply an extraordinary primatial) power over all other Churches, and an \emph{immediate} jurisdiction, to which all Catholics, both pastors and people, are bound to submit in matters not only of faith and morals, but even of discipline and government.\footnote{After quoting, in a mutilated form, the definition of the Council of Florence, whose genuineness is disputed (compare p. 97, note 1), the third chapter goes on: '\emph{Docemus et declaramus, Ecclesiam Romanam, disponente Domino, super omnes alias ordinariæ potestatis obtinere principatum, et hanc Romani Pontificis jurisdictionis potestatem, quæ vere episcopalis est, immediatam esse, erga quam cujuscunque ritus et dignitatis pastores atque fideles, tam seorsum singuli quam simul omnes, officio hierarchicæ subordinationis veræque obedientiæ obstringuntur, non solum in rebus, quæ ad fidem et mores, sed etiam in iis, quæ ad disciplinam et regimen Ecclesiæ per totum orbem diffusæ pertinent; ita ut, custodita cum Romano Pontifice tam communionis quam ejusdem fidei professionis unitate, Ecclesiæ Christi sit unus grex sub uno summo pastore. Hæc est catholicæ veritatis doctrina, a qua deviare salva fide atque salute nemo potest. . . . Si quis itaque dixerit, Romanum Pontificem habere tantummodo officium inspectionis vel directionis, non autem plenam et supremam potestatem jurisdictionis in universam Ecclesiam, non solum in rebus, quæ ad fidem et mores, sed etiam in iis, quæ ad disciplinam et regimem Ecclesiæ per totum orbem diffusæ pertinent; aut eum habere tantum potiores partes, non vero totam plenitudinem hujus supremæ potestatis; aut hanc ejus potestatem non esse ordinariam et immediatam sive in omnes ac singulas ecclesias, sive in omnes et singulos pastores et fideles; anathema sit.}'} He is, therefore, the Bishop of Bishops, over every single Bishop, and over all Bishops put together; he is in the fullest sense the Vicar of Christ, and all Bishops are simply Vicars of the Pope. The fourth chapter teaches and defines, as a divinely revealed dogma, that the Roman Pontiff, when speaking from his chair (\emph{ex cathedra}), i.e., in his official capacity, to the Christian world on subjects relating to faith or morals, is infallible, and that such definitions are irreformable (i.e., final and irreversible) in and of themselves, and not in consequence of the consent of the Church.\footnote{'\emph{ Itaque Nos traditioni a fidei Christianæ exordio perceptæ fideliter inhœrendo, ad Dei Salvatoris nostri gloriam, religionis Catholicæ exaltationem et Christianorum populorum salutem, sacro approbante Concilia, docemus et divinitus revelatum dogma esse declaramus}: Romanum Pontificem, cum ex Cathedra loquitur, id est, cum omnium Christianorum Pastoris et Doctoris munere fungens pro suprema sua Apostolica auctoritate doctrinam de fide vel moribus ab universa Ecclesia tenendam definit, per assistentiam divinam, ipsi in beato Petro promissam, ea infallibilitate pollere, qua divinus Redemptor Ecclesiam suam in definienda doctrina de fide vel moribus instructam esse voluit; ideoque ejusmodi Romani Pontificis definitiones ex sese, non autem ex consensu Ecclesiæ, irreformabiles esse.  \par '\emph{Si quis autem huic Nostræ definitioni contradicere, quod Deus avertat, præsumpserit; anathema sit.}'}

To appreciate the value and bearing of this decree, we must give a brief history of it.

The Infallibility question was suspended over the Council from the very beginning as the question of questions, for good or for evil. The original plan of the Infallibilists, to decide it by acclamation, had to be abandoned in view of a formidable opposition, which was developed inside and outside of the Council. The majority of the Bishops circulated, early in January, a monster petition, signed by 410 names, in favor of Infallibility.\footnote{Friedberg, pp. 465–470. Comp. Frommann, p. 59 sq.} The Italians and the Spaniards circulated similar petitions separately. Archbishop Spalding, of Baltimore, formerly an anti-Infallibilist, prepared an address offering some compromise to the effect that an appeal from the Pope to an œcumenical Council should be reproved.\footnote{Friedberg, pp. 470 sqq.; Frommann, pp. 61–63.} But five counter-petitions, signed by very weighty names, in all 137, representing various degrees of opposition, but agreed as to the \emph{inopportunity} of the definition, were sent in during the same month (Jan. 12 to 18) by German and Austrian, Hungarian, French, American, Oriental, and Italian Bishops.\footnote{Friedberg, pp. 472–478. The American petition against Infallibility was signed by Purcell, of Cincinnati; Kenrick, of St. Louis; McCloskey, of New York; Connolly, of Halifax; Bayley, of Newark (now Archbishop of Baltimore), and several others.}

The Pope received none of these addresses, but referred them to the Deputation on Faith. While in this he showed his impartiality, he did not conceal, in a private way, his real opinion, and gave it the weight of his personal character and influence. 'Faith in his personal infallibility,' says a well-informed Catholic, 'and belief in a constant and special communication with the Holy Ghost, form the basis of the character of Pius IX.'\footnote{Ce qui se passe au Concile, p. 130. The writer adds that some of the predecessors of Pius have held his doctrines, but none has been so ardently convinced, none has professed them '\emph{avec ce mysticisme enthousiaste, ce dédain pour les remontrances des savants et des sages, cette confiance impassible. Quel que soit le jugement de l’histoire, personne ne pourra nier que cette foi profonde ne lui ait créé dans le dix-neuvième siècle une personnalité d’une puissance et d’une majesté incomparables, dont l’éclat grandit encore un pontificat déjà si remarquable par une durée, des vertus et des malheurs vraiment exceptionnels.}' Comp. the Discourses of Pius IX., in 2 vols., Rome, 1873, and the review of Gladstone in the \emph{Quarterly Review }for Jan. 1875.} In the Council itself, Archbishop Manning, the Anglican convert, was the most zealous, devout, and enthusiastic Infallibilist; he urged the definition as the surest means of gaining hesitating Anglo-Catholics and Ritualists longing for \emph{absolute} authority; while his former teacher and friend, Dr. Pusey, feared that the new dogma would make the breach between Oxford and Rome wider than ever. Manning is 'more Catholic than Catholics' to the manor born, as the English settlers in Ireland were more Irish than Irishmen,\footnote{So Archbishop Kenrick, of St. Louis, characterized him in his \emph{Concio habenda at non habita.} Quirinus (Appendix I. p. 832) quotes from a sermon of Manning, preached at Kensington, 1869, in the Pope's name, the following passage: 'I claim to be the Supreme Judge and director of the consciences of men—of the peasant that tills the field, and the prince that sits on the throne; of the household that lives in the shade of privacy, \emph{and the Legislature that makes laws for kingdoms}. I am the sole last Supreme Judge of what is right and wrong.'} and is altogether worthy to be the successor of Pius IX. in the chair of St. Peter. Both these eminent and remarkable persons show how a sincere faith in a dogma, which borders on blasphemy, may, by a strange delusion or hallucination, be combined with rare purity and amiability of character.

Besides the all-powerful aid of the Pope, whom no Bishop can disobey without fatal consequences, the Infallibilists had the great advantage of perfect unity of sentiment and aim; while the anti-Infallibilists were divided among themselves, many of them being simply \emph{inopportunists}. They professed to agree with the majority in principle or practice, and to differ from them only on the subordinate question of definability and opportunity.\footnote{Only the address of the German Bishops took openly the ground that it would be difficult from internal reasons (viz., the contradiction of history and tradition) to proclaim Infallibility as a dogma of revelation. See Friedrich, \emph{Tagebuch}, p. 126; and Frommann, \emph{Geschichte}, p. 62.} This qualified opposition had no weight whatever with the Pope, who was as fully convinced of the opportunity and necessity of the definition as he was of the dogma itself.\footnote{On being asked whether he considered the definition of the dogma \emph{opportune}, Pius IX. resolutely answered, 'No! but \emph{necessary}.' He complained of the opposing Bishops, that, living among Protestants, they were infected by their freedom of thought, and had lost the true traditional feeling. Hase, p. 180.} And even the most advanced anti-Infallibilists, as Kenrick, Hefele, and Strossmayer, were too much hampered by Romish traditionalism to plant their foot firmly on the Scriptures, which after all must decide all questions of faith.

In the mean time a literary war on Infallibility was carried on in the Catholic Church in Germany, France, and England, and added to the commotion in Rome. A large number of pamphlets, written or inspired by prominent members of the Council, appeared for and against Infallibility. Distinguished outsiders, as Döllinger, Gratry, Hyacinthe, Montalembert, and others, mixed in the fight, and strengthened the minority.\footnote{See the literature in the next section, and in Friedberg, pp. 33–44. Comp. Frommann, pp. 66 sqq.} A confidential communication of the intellectual leader of the Anglo-Catholic secession revealed the remarkable fact that some of the most serious minds were at that time oscillating between infallibilism and skepticism, and praying to the spirits of the fathers to deliver the Church from 'the great calamity' of a new dogma.\footnote{Dr. John Henry Newman has, after long silence, retracted in 1875 his letter of 1870, which, though confidential, found its way into public 'by permission,' and has given in his adherence to the Vatican decrees, yet with minimizing qualifications, and in a tone of sadness and complaint against those ultra-zealous infallibilists who 'have stated truths in the most paradoxical forms and stretched principles till they were close upon snapping, and who at length, having done their best to set the house on fire, leave to others the task of putting out the flame.' (See his \emph{Letter to the Duke of Norfolk, on occasion of Gladstone's Expostulation}, Lond. 1875, p. 4.) Nevertheless that document deserves to be remembered for its psychological interest, and as a part of the inner history of the infallibility dogma a few months before its birth. 'Rome,' he wrote to Bishop Ullathorne, 'ought to be a name to lighten the heart at all times, and a Council's proper office is, when some great heresy or other evil impends, to inspire hope and confidence in the faithful; but now we have the greatest meeting which ever has been, and that at Rome, infusing into us by the accredited organs of Rome and of its partisans, such as the \emph{Civiltà }(the \emph{Armonia}), the \emph{Univers}, and the \emph{Tablet}, little else than fear and dismay. When we are all at rest, and have no doubts, and—at least practically, not to say doctrinally—hold the Holy Father to be infallible, suddenly there is thunder in the clearest sky, and we are told to prepare for something, we know not what, to try our faith, we know not how. No impending danger is to be averted, but a great difficulty is to be created. Is this the proper work for an œcumenical Council? As to myself personally, please God, I do not expect any trial at all; but I can not help suffering with the many souls who are suffering, and I look with anxiety at the prospect of having to defend decisions which may not be difficult to my own private judgment, but may be most difficult to maintain logically in the face of historical facts. What have we done to be treated as the faithful never were treated before? When has a definition \emph{de fide} been a luxury of devotion, and not a stern, painful necessity? Why should an aggressive, insolent faction be allowed to ``make the heart of the just sad, whom the Lord hath not made sorrowful?'' Why can not we be let alone when we have pursued peace and thought no evil? I assure you, my lord, some of the truest minds are driven one way and another, and do not know where to rest their feet—one day determining ``to give up all theology as a bad job,'' and recklessly to believe henceforth almost that the Pope is impeccable, at another tempted to ``believe all the worst which a book like \emph{Janus} says;'' others doubting about ``the capacity possessed by Bishops drawn from all corners of the earth to judge what is fitting for European society,'' and then, again, angry with the Holy See for listening to ``the flattery of a clique of Jesuits, Redemptorists, and converts.'' Then, again, think of the store of Pontifical scandals in the history of eighteen centuries, which have partly been poured forth, and partly are still to come. What Murphy [a Protestant traveling preacher] inflicted upon us in one way, Mr.Veuillot is indirectly bringing on us in another. And then, again, the blight which is falling upon the multitude of Anglican Ritualists, etc., who themselves, perhaps—at least their leaders—may never become Catholics, but who are leavening the various English denominations and parties (far beyond their own range) with principles and sentiments tending towards their ultimate absorption into the Catholic Church. With these thoughts ever before me, I am continually asking myself whether I ought not to make my feelings public; but all I do is to pray those early doctors of the Church, whose intercession would decide the matter (Augustine, Ambrose, and Jerome, Athanasius, Chrysostom, and Basil), to avert this great calamity. If it is God's will that the Pope's infallibility be defined, then is it God's will to throw back ``the times and moments'' of that triumph which he has destined for his kingdom, and I shall feel I have but to bow my head to his adorable, inscrutable Providence. You have not touched upon the subject yourself, but I think you will allow me to express to you feelings which, for the most part, I keep to myself. . . .' See an excellent German translation of this letter in Quirinus (p. 274, Germ. ed.) and in Friedberg (p. 131). The English translator of Quirinus has substituted the English original as given here from the \emph{Standard}, April 7, 1870.}

After preliminary skirmishes, the formal discussion began in earnest in the 50th session of the General Congregation, May 13, 1870, and lasted to the 86th General Congregation, July 16. About eighty Latin speeches\footnote{According to Manning, but only 65 according to Friedberg, p. 47.} were delivered in the general discussion on the schema \emph{de Romano Pontifice}, nearly one half of them on the part of the opposition, which embraced less than one fifth of the Council. When the arguments and the patience of the assembly were pretty well exhausted, the President, at the petition of a hundred and fifty Bishops, closed the general discussion on the third day of June. About forty more Bishops, who had entered their names, were thus prevented from speaking; but one of them, Archbishop Kenrick, of St. Louis, published his strong argument against Infallibility in Naples.\footnote{Hence the title '\emph{Concio habenda at non habita}'—\emph{prepared for speaking, but not spoken}. See the prefatory note, dated Rome, June 8, 1870.} Then five special discussions commenced on the proemium and the four chapters. 'For the fifth or last discussion a hundred and twenty Bishops inscribed their names to speak; fifty of them were heard, until on both sides the burden became too heavy to bear; and, by mutual consent, a useless and endless discussion, from mere exhaustion, ceased.'\footnote{Manning, \emph{Petri Privil}. III. pp. 31, 32. He gives this representation to vindicate the liberty of the Council; but the minority complained of an arbitrary close of the discussion. They held an indignation meeting in the residence of Cardinal Rauscher, and protested '\emph{contra violationem nostri juris},' but without effect. See the protest, with eighty-one signatures, in Friedrich, \emph{Doc.} II. p. 379; comp. Frommann, \emph{Geschichte}, p. 174.}

When the vote was taken on the whole four chapters of the Constitution of the Church, July 13, 1870, in the 85th secret session of the General Congregation (601 members being present), 451 voted \emph{Placet}, 88 \emph{Non Placet}, 62 \emph{Placet juxta modum}, over 80 (perhaps 91), though present in Rome or in the neighborhood, abstained for various reasons from voting.\footnote{See the list in Friedberg, pp. 146–149; also in Friedrich, \emph{Docum.} II. pp. 426 sqq.; and Quirinus, \emph{Letter LXVI.} pp. 778 sqq. Quirinus errs in counting the 91 (according to others, 85 or only 70) absentees among the 601. There were in all from 680 to 692 members present in Rome at the time. See Fessler, p. 89 (who states the number of absentees to be 'over 80'), and Frommann, p. 201. The protest of the minority to the Pope, July 17, states the number of voters in the same way, except that 70, instead of 91 or 85, is given as the number of absentees: '\emph{Notum est Sanctitati Vestræ}, 88 \emph{Patres fuisse, qui, conscientia urgente et amore s. Ecclesiæ permoti, suffragium suum per verba}  non placet  \emph{emiserunt;} 62 \emph{alios, qui suffragati sunt per verba}  placet juxta modum,  \emph{denique }70 \emph{circiter qui a congregatione abfuerunt atque a suffragio emittendo abstinuerunt. Hic accedunt et alii, qui, infirmitatibus aut gravioribus rationibus ducti, ad suas diœceses reversi sunt}.'} Among the negative votes were the Prelates most distinguished for learning and position, as Schwarzenberg, Cardinal Prince-Archbishop of Prague; Rauscher, Cardinal Prince-Archbishop of Vienna; Darboy, Archbishop of Paris; Matthieu, Cardinal-Archbishop of Besançon; Ginoulhiac, Archbishop of Lyons; Dupanloup, Bishop of Orleans; Maret, Bishop of Sura (i. p.); Simor, Archbishop of Gran and Primate of Hungary; Haynald, Archbishop of Kalocsa; Förster, Prince-Bishop of Breslau; Scherr, Archbishop of Munich; Ketteler, Bishop of Mayence; Hefele, Bishop of Rottenburg; Strossmayer, Bishop of Bosnia and Sirmium; MacHale, Archbishop of Tuam; Connolly, Archbishop of Halifax; Kenrick, Archbishop of St. Louis.

On the evening of the 13th of July the minority sent a deputation, consisting of Simor, Ginoulhiac, Scherr, Darboy, Ketteler, and Rivet, to the Pope. After waiting an hour, they were admitted at 9 o'clock in the evening. They asked simply for a withdrawal of the addition to the third chapter, which assigns to the Pope the exclusive possession of all ecclesiastical powers, and for the insertion, in the fourth chapter, of a clause limiting his infallibility to those decisions which he pronounces '\emph{innixus testimonio ecclesiarum.}' Pius returned the almost incredible answer: 'I shall do what I can, my dear sons, but I have not yet read the scheme; I do not know what it contains.'\footnote{He spoke in French: '\emph{Te ferai mon possible, mes chers fils, mais je n’ai pas encore lu le schéma; je ne sais pas ce qu’il contient.}' Quirinus, \emph{Letter LXIX.} p. 800.} He requested Darboy, the spokesman of the deputation, to hand him the petition in writing. Darboy promised to do so; and added, not without irony, that he would send with it the schema which the Deputation on Faith and the Legates had with such culpable levity omitted to lay before his Holiness, exposing him to the risk of proclaiming in a few days a decree he was ignorant of. Pius surprised the deputation by the astounding assurance that the whole Church had always taught the unconditional Infallibility of the Pope. Then Bishop Ketteler of Mayence implored the holy Father on his knees to make some concession for the peace and unity of the Church.\footnote{Quirinus, \emph{Letter LXIX}. p. 801, gave, a few days afterwards, from direct information, the following fresh and graphic description of this interesting scene: 'Bishop Ketteler then came forward, flung himself on his knees before the Pope, and entreated for several minutes that the Father of the Catholic world would make some concession to restore peace and her lost unity to the Church and the Episcopate. It was a peculiar spectacle to witness these two men, of kindred and yet widely diverse nature, in such an attitude—the one prostrate on the ground before the other. Pius is ``\emph{totus teres atque rotundus,}'' firm and immovable, smooth and hard as marble, infinitely self-satisfied intellectually, mindless and ignorant; without any understanding of the mental conditions and needs of mankind, without any notion of the character of foreign nations, but as credulous as a nun, and, above all, penetrated through and through with reverence for his own person as the organ of the Holy Ghost, and therefore an absolutist from head to heel, and filled with the thought, ``I, and none beside me.'' He knows and believes that the Holy Virgin, with whom he is on the most intimate terms, will indemnify him for the loss of land and subjects by means of the Infallibility doctrine, and the restoration of the Papal dominion over states and peoples as well as over churches. He also believes firmly in the miraculous emanations from the sepulchre of St. Peter. At the feet of this man the German Bishop flung himself, ``\emph{ipso Papa papalior,}'' a zealot for the ideal greatness and unapproachable dignity of the Papacy, and, at the same time, inspired by the aristocratic feeling of a Westphalian nobleman and the hierarchical self-consciousness of a Bishop and successor of the ancient chancellor of the empire, while yet he is surrounded by the intellectual atmosphere of Germany, and, with all his firmness of belief, is sickly with the pallor of thought, and inwardly struggling with the terrible misgiving that, after all, historical facts are right, and that the ship of the \emph{Curia}, though for the moment it proudly rides the waves with its sails swelled by a favorable wind, will be wrecked on that rock at last.'} This prostration of the proudest of the German prelates made some impression. Pius dismissed the deputation in a hopeful temper. But immediately afterwards Manning and Senestrey (Bishop of Regensburg) strengthened his faith, and frightened him by the warning that, if he made any concession, he would be disgraced in history as a second Honorius.

In the secret session on the 16th of July, on motion of some Spanish Bishops, an addition was inserted '\emph{non autem ex consensu ecclesiæ,} which makes the decree still more obnoxious.\footnote{Quirinus, p. 804: 'Thus the Infallibilist decree, as it is now to be received under anathema by the Catholic world, is an eminently Spanish production, as is fitting for a doctrine which was born and reared under the shadow of the Inquisition.'} On the same day Cardinal Rauscher, in a private audience, made another attempt to induce the Pope to yield, but was told, 'It is too late.'

On the 17th of July fifty-six Bishops sent a written protest to the Pope, declaring that nothing had occurred to change their conviction as expressed in their negative vote; on the contrary, they were confirmed in it; yet filial piety and reverence for the holy Father would not permit them to vote \emph{Non Placet}, openly and in his face, in a matter which so intimately concerned his person, and that therefore they had resolved to return forthwith to their flocks, which had already too long been deprived of their presence, and were now filled with apprehensions of war. Schwarzenberg, Matthieu, Simor, and Darboy head the list of signers.\footnote{See the protest in Friedberg, p. 622. Comp. Frommann, p. 207.} On the evening of the same day not only the fifty-six signers, but sixty additional members of the opposition departed from Rome, promising to each other to make their future conduct dependent on mutual understanding.

This was the turning-point: the opposition broke down by its own act of cowardice. They ought to have stood like men on the post of duty, and repeated their negative vote according to their honest convictions. They could thus have prevented the passage of this momentous decree, or at all events shorn it of its œcumenical weight, and kept it open for future revision and possible reversal. But they left Rome at the very moment when their presence was most needed, and threw an easy victory into the lap of the majority.

When, therefore, the fourth public session was held, on the memorable 18th of July (Monday), there were but 535 Fathers present, and of these all voted \emph{Placet}, with the exception of two, viz., Bishop Riccio, of Cajazzo, in Sicily, and Bishop Fitzgerald, of Little Rock, Arkansas, who had the courage to vote \emph{Non Placet}, but immediately, before the close of the session, submitted to the voice of the Council. In this way a moral unanimity was secured as great as in the first Council of Nicæa, where likewise two refused to subscribe the Nicene Creed, 'What a wise direction of Providence,' exclaimed the \emph{Civiltà cattolica}, '535 yeas against 2 nays. \emph{Only two} nays, therefore almost total unanimity; and yet two \emph{nays}, therefore full liberty of the Council. How vain are all attacks against the œcumenical character of this most beautiful of all Councils!'

After the vote the Pope confirmed the decrees and canons on the Constitution of the Church of Christ, and added from his own inspiration the assurance that the supreme authority of the Roman Pontiff did not suppress but aid, not destroy but build up, and formed the best protection of the rights and interests of the Episcopate.\footnote{'\emph{ Summa ista Romani Pontificis auctoritas, Venerabiles Fratres, non opprimit sed adjuvat, non destruit sed ædificat, et sæpissime confirmat in dignitate, unit in charitate, et Fratrum, scilicet Episcoporum, jura firmat atque tuetur. Ideoque illi, qui nunc judicant in commotione, sciant, non esse in commotione Dominum. Meminerint, quod paucis abhinc annis, oppositam tenentes sententiam, abundaverunt in sensu Nostro, et in sensu majoris partis hujus amplissimi Consessus, sed tunc judicaverunt in spiritu auræ lenis. Numquid in eodem judicio judicando duæ oppositæ possunt existere conscientiæ? Absît. Illuminet ergo Deus sensus et corda; et quoniam Ipse facit mirabilia magna solus, illuminet sensus et corda, ut omnes accedere possint ad sinum Patris, Christi Jesu in terris indigni Vicarii, qui eos amat, eos diligit, et exoptat unum esse cum illis; et ita simul in vinculo charitatis conjuncti prœliare possimus prœlia Domini, ut non solum non irrideant nos inimici nostri, sed timeant potius, et aliquando arma malitiæ cedant in conspectu veritatis, sicque omnes cum D. Augustino dicere valeant: ``Tu vocasti me in admirabile lumen tuum, et ecce video.}'''}

The days of the two most important public sessions of the Vatican Council, namely the first and the last, were the darkest and stormiest which Rome saw from Dec. 8, 1869, to the 18th of July, 1870. The Episcopal votes and the Papal proclamation of the new dogma were accompanied by flashes of lightning and claps of thunder from the skies, and so great was the darkness which spread over the Church of St. Peter, that the Pope could not read the decree of his own Infallibility without the artificial light of a candle.\footnote{Quirinus, \emph{Letter LXIX.} p. 809. A Protestant eye-witness, Prof. Ripley, thus described the scene in a letter from Rome, published in the \emph{New York Tribune} (of which he is one of the editors) for Aug. 11, 1870: 'Rome, July 19.—Before leaving Rome I send you a report of the last scene of that absurd comedy called the Œcumenical Vatican Council. . . . It is at least a remarkable coincidence that the opening and closing sessions of the Council were inaugurated with fearful storms, and that the vigil of the promulgation of the dogma was celebrated with thunder and lightning throughout the whole of the night. On the 8th of last December I was nearly drowned by the floods of rain, which came down in buckets; yesterday morning I went down in rain, and under a frowning sky which menaced terrible storms later in the day. . . . \emph{Kyrie eleison} we heard as soon as the mass was said, and the whole multitude joined in singing the plaintive measure of the Litany of the Saints, and then with equal fervor was sung \emph{Veni Creator}, which was followed by the voice of a secretary reading in a high key the dogma. At its conclusion the names of the Fathers were called over, and \emph{Placet }after \emph{Placet }succeeded \emph{ad nauseam.} But what a storm burst over the church at this moment! The lightning flashed and the thunder pealed as we have not heard it this season before. Every \emph{Placet }seemed to be announced by a flash and terminated by a clap of thunder. Through the cupolas the lightning entered, licking, as it were, the very columns of the Baldachino over the tomb of St. Peter, and lighting up large spaces on the pavement. Sure, God was there—but whether approving or disproving what was going on, no mortal man can say. Enough that it was a remarkable coincidence, and so it struck the minds of all who were present. And thus the roll was called for one hour and a half, with this solemn accompaniment, and then the result of the voting was taken to the Pope. The moment had arrived when he was to declare himself invested with the attributes of God—nay, a God upon earth. Looking from a distance into the hall, which was obscured by the tempest, nothing was visible but the golden mitre of the Pope, and so thick was the darkness that a servitor was compelled to bring a lighted candle and hold it by his side to enable him to read the formula by which he deified himself. And then—what is that indescribable noise? Is it the raging of the storm above?—the pattering of hail-stones? It approaches nearer, and for a minute I most seriously say that I could not understand what that swelling sound was until I saw a cloud of white handkerchiefs waving in the air. The Fathers had begun with clapping—they were the fuglemen to the crowd who took up the notes and signs of rejoicing until the church of God was converted into a theatre for the exhibition of human passions. ``\emph{Viva Pio Nono!'' ``Viva il Papa Infallibile}!'' ``\emph{Viva il trionfo dei Cattolici!}'' were shouted by this priestly assembly; and again another round they had; and yet another was attempted as soon as the \emph{Te Deum} had been sung and the benediction had been given.'} This voice of nature was variously interpreted, either as a condemnation of Gallicanism and liberal Catholicism, or as a divine attestation of the dogma like that which accompanied the promulgation of the law from Mount Sinai, or as an evil omen of impending calamities to the Papacy.

And behold, the day after the proclamation of the dogma, Napoleon III., the political ally and supporter of Pius IX., unchained the furies of war, which in a few weeks swept away the Empire of France and the temporal throne of the infallible Pope. His own subjects forsook him, and almost unanimously voted for a new sovereign, whom he had excommunicated as the worst enemy of the Church. A German Empire arose from victorious battle-fields, and Protestantism sprung to the political and military leadership of Europe. About half a dozen Protestant Churches have since been organized in Rome, where none was tolerated before, except outside of the walls or in the house of some foreign embassador; a branch of the Bible Society was established, which the Pope in his Syllabus denounces as a pest; and a public debate was held in which even the presence of Peter at Rome was called in question. History records no more striking example of swift retribution of criminal ambition. Once before the Papacy was shaken to its base at the very moment when it felt itself most secure: Leo X. had hardly concluded the fifth and last Lateran Council in March, 1517, with a celebration of victory, when an humble monk in the North of Europe sounded the key-note of the great Reformation.

What did the Bishops of the minority do? They all submitted, even those who had been most vigorous in opposing, not only the opportunity of the definition, but the dogma itself. Some hesitated long, but yielded at last to the heavy pressure. Cardinal Rauscher, of Vienna, published the decree already in August, and afterwards withdrew his powerful 'Observations on the Infallibility of the Church' from the market; regarding this as an act of glorious self-denial for the welfare of the Church. Cardinal Schwarzenberg, of Prague, waited with the publication till Jan. 11, 1871, and shifted the responsibility upon his theological advisers. Bishop Hefele, of Rottenburg, who has forgotten more about the history of Councils than the infallible Pope ever knew, after delaying till April 10, 1871, submitted, not because he had changed his conviction, but, as he says, because 'the peace and unity of the Church is so great a good that great and heavy personal sacrifices may be made for it;' i.e., truth must be sacrificed to peace. Bishop Maret, who wrote two learned volumes against Papal Infallibility and in defense of Gallicanism, declared in his retractation that he 'wholly rejects every thing in his work which is opposed to the dogma of the Council,' and 'withdraws it from sale.' Archbishop Kenrick yielded, but has not refuted his \emph{Concio habenda at non habita}, which remains an irrefragable argument against the new dogma. Even Strossmayer, the boldest of the bold in the minority, lost his courage, and keeps his peace. Darboy died a martyr in the revolt of the communists of Paris, in April, 1871. In a conversation with Dr. Michaud, Vicar of St. Madeleine, who since seceded from Rome, he counseled external and official submission, with a mental reservation, and in the hope of better times. His successor, Msgr. Guibert, published the decrees a year later (April, 1872), without asking the permission of the head of the French Republic. Of those opponents who, though not members of the Council, carried as great weight as any Prelate, Montalembert died during the Council; Newman kept silence; Père Gratry, who had declared and proved that the question of Honorius 'is totally gangrened by fraud,' wrote from his death-bed at Montreux, in Switzerland (Feb. 1872), to the new Archbishop of Paris, that he submitted to the Vatican Council, and effaced 'every thing to the contrary he may have written.'\footnote{See details on the reception and publication of the Vatican decrees in Friedberg, pp. 53 sqq., 775 sqq.; Frommann, pp. 215–230; on Gratry, the \emph{Annales de Philosophie Chrétienne}, Sept. 1871, p. 236.}

It is said that the adhesion of the minority Bishops was extorted by the threat of the Pope not to renew their 'quinquennial faculties' (\emph{facultates quinquennales}), that is, the Papal licenses renewed every five years, permitting them to exercise extraordinary episcopal functions which ordinarily belong to the Pope, as the power of absolving from heresy, schism, apostasy, secret crime (except murder), from vows, duties of fasting, the power of permitting the reading of prohibited books (for the purpose of refutation), marrying within prohibited degrees, etc.\footnote{See the article \emph{Facultäten}, in Wetzer und Welte's  \emph{Kirchenlexikon oder Encyklop. der katholischen Theologie}, Vol. III. pp. 879 sqq.}

But, aside from this pressure, the following considerations sufficiently explain the fact of submission.

1. Many of the dissenting Bishops were professedly anti-Infallibilists, not from principle, but only from subordinate considerations of expediency, because they apprehended that the definition would provoke the hostility of secular governments, and inflict great injury on Catholic interests, especially in Protestant countries. Events have since proved that their apprehension was well founded.

2. All Roman Bishops are under an oath of allegiance to the Pope, which binds them 'to preserve, defend, \emph{increase}, and \emph{advance }the rights, honors, privileges, and authority of the holy Roman Church, of our lord the Pope, and his successors.'

3. The minority Bishops defended Episcopal infallibility against Papal infallibility. They claimed for themselves what they denied to the Pope. Admitting the infallibility of an œcumenical Council, and forfeiting by their voluntary absence on the day of voting the right of their protest, they must either on their own theory accept the decision of the Council, or give up their theory, cease to be Roman Catholics, and run the risk of a new schism.

At the same time this submission is an instructive lesson of the fearful spiritual despotism of the Papacy, which overrules the stubborn facts of history and the sacred claims of individual conscience. For the facts so clearly and forcibly brought out before and during the Council by such men as Kenrick, Hefele, Rauscher, Maret, Schwarzenberg, and Dupanloup, have not changed, and can never be undone. On the one hand we find the results of a life-long, conscientious, and thorough study of the most learned divines of the Roman Church, on the other ignorance, prejudice, perversion, and defiance of Scripture and tradition; on the one hand we have history shaping theology, on the other theology ignoring or changing history; on the one hand the just exercise of reason, on the other blind submission, which destroys reason and conscience. But truth must and will prevail at last.

\xsection{Papal Infallibility Explained, and Tested by Scripture and Tradition.}

§ 34. Papal Infallibility Explained, and Tested by Tradition and Scripture.

Literature.

I. For Infallibility.

The older defenders of Infallibility are chiefly Bellarmin, Ballerini, Litta, Alphons de Liguori (whom the Pope raised to the dignity of a \emph{doctor ecclesiæ}, March 11, 1872), Card. Orsi, Perrone, and  Joseph Count du Maistre (Sardinian statesman, d. at Turin Feb. 26, 1821, author of \emph{Du Pape}, 1819; new edition, Paris, 1843, with the Homeric motto: \textgreek{εἶς κοίρανος ἔστω.}

During and after the Vatican Council: the works of Archbishops Manning and Dechamps, already quoted, pp. 134, 135.

Jos. Cardoni (Archbishop of Edessa, in partibus): \emph{Elucubratio de dogmatica Romani Pontificis Infallibilitate ejusque Definibilitate}, Romæ (typis Civilitatis Cattolicæ), 1870 (May, 174 pp.). The chief work on the Papal side, clothed with a semi-official character.

Hermann Rump: \emph{Die Unfehlbarkeit den Papstes und die Stellung der in Deutschland verbreiteten theologischen Lehrbücher zu dieser Lehre}, Münster, 1870 (173 pp.).

Franz Friedhoff (Prof. at Münster): \emph{Gegen-Erwägungen über die päpstliche Unfehlbarkeit}, Münster, 1869 (21 pp.). Superficial.

Flor. Riess and  Karl von Weber (Jesuits): \emph{Das Oekum. Concil. Stimmen aus Maria-Laach, Neue Folge}, No. X. \emph{Die päpstliche Unfehlbarkeit und der alte Glaube der Kirche}, Freiburg im Breisgau, 1870 (110pp.).

G. Bickel: \emph{Gründe fur die Unfehlbarkeit des Kirchenoberhauptes nebst Widerlegung der Einwürfe}, Münster, 1870.

Rev.  P. Weninger (Jesuit): \emph{L’infaillibilité du Pape devant la raison et l’écriture, les papes et les conciles, les pères et les théologiens, les rois et les empereurs. }Translated from the German into French by P. Bélét. (Highly spoken of by Pius IX. in a brief to Abbé Bélét, Nov. 17, 1869; see Friedberg, l.c. p. 487. Weninger wrote besides several pamphlets on Infallibility in German, Innsbruck, 1841; Graz, 1853; in English, New York and Cincinnati, 1868. Archbishop Kenrick, in his \emph{Concio}; speaks of him as 'a pious and extremely zealous but ignorant man,' whom he honored with 'the charity of silence' when requested to recommend one of his books.)

Widerlegung der vier unter die Väter des Concils vertheilten Brochüren gegen die Unfehlbarkeit (transl. of Animadversiones in quatuor contra Romani Pontificis infallibilitatem editos libellos), Münster, 1870.

Bishop  Jos. Fessler: Die wahre und die falsche Unfehlbarkeit der Päpste (against Prof. von Schulte), Wien,1871.

Bishop  Ketteler: Das unfehlbare Lehramt des Papstes, nach der Entscheidung des Vaticanischen Concils, Mainz, 1871, 3te Aufl.

M. J. Scheeben: Schulte und Döllinger, gegen das Concil. Kritische Beleuchtung, etc., Regensburg, 1871.

Amédée de Margerie: Lettre au R. P. Gratry sur le Pape Honorius et le Bréviaire Romain, Nancy, 1870

Paul Bottala (S.J.): Pope Honorius before the Tribunal of Reason and History, London, 1868.

II. Against Infallibility.

(a) By Members of the Council.

Mgr.  H. L. C. Maret (Bishop of Sura, in part., Canon of St. Denis and Dean of the Theological Faculty in Paris): \emph{Du Concile général et de la paix religieuse}, Paris, 1869, 2 Tom. (pp. 554 and 555). An elaborate defense of Gallicanism; since revoked by the author, and withdrawn from sale.

Peter Richard Kenrick (Archbishop of St. Louis): \emph{Concio in Concilio Vaticano habenda at non habita}, Neapoli (typis fratrum de Angelis in via Pellegrini 4), 1870. Reprinted in Friedrich, \emph{Documenta, }I. pp. 187–226. An English translation in L. W. Bacon's \emph{An Inside View of the Vatican Council}, New York, pp. 90–166.

Quæstio (no place or date of publication). A very able Latin dissertation occasioned and distributed (perhaps partly prepared) by Bishop Ketteler, of Mayence, during the Council. It was printed but not published in Switzerland, in 1870, and reprinted in Friedrich, \emph{Documenta}, I. pp. 1–128.

\emph{La liberté du Concile et l’infaillibilité}. Written or inspired by Darboy, Archbishop of Paris. Only flfty copies were printed, for distribution among the Cardinals. Reprinted in Friedrich, \emph{Documenta, }I. pp. 129–186.

Card.  Rauscher: \emph{Observationes quædam de infallibilitatis ecclesiæ subjecto}, Neapoli and Vindobonæ, 1870 (83 pp.).

\emph{De Summi Pontificis infallibilitate personali}, Neapoli, 1870 (32 pp.). Written by Prof.  Salesius Mayer, and distributed in the Council by Cardinal Schwarzenberg.

Jos. de Hefele (Bishop of Rottenburg, formerly Prof. at Tübingen): \emph{Causa Honorii Papæ}, Neap. 1870 (pp. 28). The same: \emph{Honorius und das sechste allgemeine Concil} (with an appendix against Pennachi, 43 pp.), Tübingen, 1870. English translation, with introduction, by Dr.  Henry B. Smith, in the \emph{Presbyterian Quarterly and Princeton Review, }New York, for April, 1872, pp. 273 sqq. Against Hefele comp.  Jos. Pennachi (Prof. of Church History in Rome): \emph{De Honorii I. Pontificis Romani causa in Concilio VI.}

(\emph{b}) By Catholics, not Members of the Council.

Janus: \emph{The Pope and the Council}, 1869. See above, p. 134.

\emph{Erwägungen für die Bischöfe del Conciliums über die Frage der päpstlichen Unfehlbarkeit}, Oct. 1869. Dritte Aufl. München. [By J. von Döllinger.]

J. von Döllinger: Einige Worte über die Unfehlbarkeitsadresse, etc., München, 1870.

Jos. H. Reinkens (Prof. of Church History in Breslau): Ueber päpstliche Unfehlbarkeit, München, 1870.

Clemens Schmitz (Cath. Priest): Ist der Papst unfehlbart? Aus Deutschlands und des P. Deharbe Catechismen beantwortet, München, 1870.

J. Fr. Ritter von Schulte (Prof. in Prague, now in Bonn): \emph{Das Unfehlbarkeits-Decret vom }18 \emph{Juli }1870 \emph{auf seine Verbindlichkeit geprüft}, Prague, 1870. Die Macht der röm. Päpste über Fürsten, Länder, Vöker, etc. seit Gregor VII. zur Würdigung ihrer Unfehlbarkeit beleuchtet, etc., 2d edition, Prague. The same, translated into English (The Power of the Roman Popes over Princes, etc.), by Alfred Somers [a brother of Schulte], Adelaide, 1871.

A. Gratry (Priest of the Oratoire and Member of the French Academy): \emph{ Four Letters to the Bishop of Orleans} (Dupanloup) \emph{and the Archbishop of Malines }(Dechamps), in French, Paris, 1870; several editions, also translated into German, English, etc. These learned and eloquent letters gave rise to violent controversies. They were denounced by several Bishops, and prohibited in their dioceses; approved by others, and by Montalembert. The Pope praised the opponents. Against him wrote Dechamps (Three letters to Gratry, in French; German translation, Mayence, 1870) and A. de Margerie. Gratry recanted on his death-bed.

P. Le Page Renouf: \emph{The Condemnation of Pope Honorius}, London, 1868.

Antonio Magrassi: \emph{Lo Schema sull’ infallibilità personale del Romano Pontefice}, Alessandria, 1870.

\emph{Della pretesa infallibilità personale del Romano Pontefice}, 2d ed. Firenze, 1870 (anonymous, 80 pp.).

J. A. B. Lutterbeck: \emph{Die Clementinen und ihr Verhältniss zum Unfehlbarkeitsdogma}, Giessen, 1872 (pp. 85).

Joseph Langen (Old Catholic Prof. in Bonn): \emph{Das Vaticanische Dogma von dem Universal-Episcopat und der Unfehlbarkeit des Papstes in s. Verh. zur exeg. Ueberlieferung vom} 7 \emph{bis zum }13\emph{ten Jahrh}. 3 Parts. Bonn, 1871–73.

The sinlessness of the Virgin Mary and the personal infallibility of the Pope are the characteristic dogmas of modern Romanism, the two test dogmas which must decide the ultimate fate of this system. Both were enacted under the same Pope, and both faithfully reflect his character. Both have the advantage of logical consistency from certain premises, and seem to be the very perfection of the Romish form of piety and the Romish principle of authority. Both rest on pious fiction and fraud; both present a refined idolatry by clothing a pure humble woman and a mortal sinful man with divine attributes. The dogma of the Immaculate Conception, which exempts the Virgin Mary from sin and guilt, perverts Christianism into Marianism; the dogma of Infallibility, which exempts the Bishop of Rome from error, resolves Catholicism into Papalism, or the Church into the Pope. The worship of a woman is virtually substituted for the worship of Christ, and a man-god in Rome for the God-Man in heaven. This is a severe judgment, but a closer examination will sustain it.

The dogma of the Immaculate Conception, being confined to the sphere of devotion, passed into the modern Roman creed without serious difficulty; but the dogma of Papal Infallibility, which involves a question of absolute power, forms an epoch in the history of Romanism, and created the greatest commotion and a new secession. It is in its very nature the most fundamental and most comprehensive of of all dogmas. It contains the whole system in a nutshell. It constitutes a new rule of faith. It is the article of the standing or falling Church. It is the direct antipode of the Protestant principle of the absolute supremacy and infallibility of the Holy Scriptures. It establishes a perpetual divine oracle in the Vatican. Every Catholic may hereafter say, I believe—not because Christ, or the Bible, or the Church, but—because the infallible Pope has so declared and commanded. Admitting this dogma, we admit not only the whole body of doctrines contained in the Tridentine standards, but all the official Papal bulls, including the mediæval monstrosities of the Syllabus (1864), the condemnation of Jansenism, the bull '\emph{Unam Sanctam}' of Boniface VIII. (1302), which, under pain of damnation, claims for the Pope the double sword, the secular as well as the spiritual, over the whole Christian world, and the power to depose princes and to absolve subjects from their oath of allegiance.\footnote{This bull has been often disowned by Catholics (e.g., by the Universities of Sorbonne, Louvain, Alcala, Salamanca, when officially asked by Mr. Pitt, Prime Minister of Great Britain, 1788, also by Martin John Spalding, Archbishop of Baltimore, in his Lectures on Evidences, 1866), and, to some extent, even by Pius IX. (see Friedberg, p. 718), but it is unquestionably official, and was renewed and approved by the fifth Lateran Council, Dec. 19, 1516. Paul III. and Pius V. acted upon it, the former in excommunicating and deposing Henry VIII. of England, the latter in deposing Queen Elizabeth, exciting her subjects to rebellion, and urging Philip of Spain to declare war against her (see the Bullarium Rom., Camden, Burnet, Froude, etc.). The Papal Syllabus sanctions it by implication, in No. 23, which condemns as an error the opinion that Roman Pontiffs have exceeded the limits of their power.} The past is irreversibly settled, and in all future controversies on faith and morals we must look to the same unerring tribunal in the Vatican. Even œcumenical Councils are superseded hereafter, and would be a mere waste of time and strength.

On the other hand, if the dogma is false, it involves a blasphemous assumption, and makes the nearest approach to the fulfillment of St. Paul's prophecy of the man of sin, who 'as God sitteth in the temple of God, showing himself off that he is God' (2 Thess. ii. 4).

Let us first see what the dogma does not mean, and what it does mean.

It does not mean that the Pope is infallible in his \emph{private} opinions' on theology and religion. As a man, he may be a heretic (as Liberius, Honorius, and John XXII.), or even an unbeliever (as John XXIII., and, perhaps, Leo X.), and yet, at the same time, infallible as Pope, after the fashion of Balaam and Kaiphas.

Nor does it mean that infallibility extends beyond the proper sphere of religion and the Church. The Pope may be ignorant of science and literature, and make grave mistakes in his political administration, or be misinformed on matters of fact (unless necessarily involved in doctrinal decisions), and yet be infallible in defining articles of faith.\footnote{Pope Pius IX. started as a political reformer, and set in motion that revolution which, notwithstanding his subsequent reactionary course, resulted in the unification of Italy and the loss of the States of the Church, against which he now so bitterly protests.}

Infallibility does not imply impeccability. And yet freedom from error and freedom from sin are so nearly connected in men's minds that it seems utterly impossible that such moral monsters as Alexander VI. and those infamous Popes who disgraced humanity during the Roman pornocracy in the tenth and eleventh centuries, should have been vicars of Jesus Christ and infallible organs of the Holy Ghost. If the inherent infallibility of the visible Church logically necessitates the infallibility of the visible head, it is difficult to see why the same logic should not with equal conclusiveness derive the personal holiness of the head from the holiness of the body.

On the other hand, the dogma does mean that all official utterances of the Roman Pontiff addressed to the Catholic Church on matters of Christian faith and duty are infallibly true, and must be accepted with the same faith as the word of the living God. They are not simply final in the sense in which all decisions of an absolute government or a supreme court of justice are final until abolished or superseded by other decisions,\footnote{In this general sense Joseph de Maistre explains infallibility to be the same in the spiritual order that sovereignty means in the civil order: '\emph{L’un et l’autre expriment cette haute puissance qui les domine toutes, dont toutes les autres dérivent, qui gouverne et n’est pas gouvernée, qui juge et n’est pas jugée. Quand nous disons que l’Eglise est infaillible, nous ne demandons pour elle, il est bien essentiel de l’observer, aucun privilége particulier; nous demandons seulement qu’elle jouisse du droit commun à toutes les souverainetés possible qui toutes agissent néssairement comme infaillibles; car tout gouvernement est absolu; et du moment où l’on peut lui résister sous prétexte d’erreur ou d’injustice, il n’existe plus.' Du Pape,} ch. i., pp. 15, 16.} but they are irreformable, and can never be revoked. This infallibility extends over eighteen centuries, and is a special privilege conferred by Christ upon Peter, and through him upon all his legitimate successors. It belongs to every Pope from Clement to Pius IX., and to every Papal bull addressed to the Catholic world. It is personal, i.e., inherent in Peter and the Popes; it is independent, and needs no confirmation from the Church or an œcumenical Council, either preceding or succeeding; its decrees are binding, and can not be rejected without running the risk of eternal damnation.\footnote{Archbishop Manning (\emph{Petri Privil.} III. pp. 112, 113) defines the doctrine of Infallibility in this way:  \par '1. The privilege of infallibility \emph{is personal}, inasmuch as it attaches to the Roman Pontiff, the successor of Peter, as a \emph{public person}, distinct from, but inseparably united to, the Church; but it is not personal, in that it is attached, not to the private person, but to the primacy which he alone possesses. \par '2. It is also \emph{independent}, inasmuch as it does not depend upon either the \emph{Ecclesia docens} or the \emph{Ecclesia discens}; but it is not independent, in that it depends in all things upon the divine head of the Church, upon the institution of the primacy by him, and upon the assistance of the Holy Ghost. \par '3. It is \emph{absolute}, inasmuch as it can be circumscribed by no human or ecclesiastical law; it is not absolute, in that it is circumscribed by the office of guarding, expounding, and defending the deposit of revelation. \par '4. It is \emph{separate }in no sense, nor can be, nor can be so called, without manifold heresy, unless the word be taken to mean \emph{distinct}. In this sense, the Roman Pontiff is distinct from the Episcopate, and is a distinct subject of infallibility; and in the exercise of his supreme doctrinal authority, or magisterium, he does not depend for the infallibility of his definitions upon the consent or consultation of the Episcopate, but only on the divine assistance of the Holy Ghost.'}

Even within the narrow limits of the Vatican decision there is room for controversy on the precise meaning of the figurative term \emph{ex cathedra loqui}, and the extent of faith and \emph{morals}, viz., whether Infallibility includes only the supernatural order of revealed truth and duty, or also natural and political duties, and questions of mere history, such as Peter's residence in Koine, the number of œcumenical Councils, the teaching of Jansen and Quesnel, and other disputed facts closely connected with dogmas. But the main point is clear enough. The Ultramontane theory is established, Gallicanism is dead and buried.

\emph{Ultramontanism and Gallicanism.}

The Vatican dogma is the natural completion of the Papal polity, as the dogma of the Immaculate Conception of Mary is the completion of the Papal cultus.

If we compare the Papal or Ultramontane theory with the Episcopal or Gallican theory, it has the undeniable advantage of logical consistency. The two systems are related to each other like monarchy and aristocracy, or rather like absolute monarchy and limited monarchy. The one starts from the divine institution of the Primacy (Matt. xvi. 18), and teaches the infallibility of the head; the other starts from the divine institution of the Episcopate (Matt. xviii. 18), and teaches the infallibility of the body and the superiority of an œcumenical Council over the Pope. Conceding once the infallibility of the collective Episcopate, we must admit, as a consequence, the infallibility of the Primacy, which represents the Episcopate, and forms its visible and permanent centre. If the body of the teaching Church can never err, the head can not err; and, \emph{vice versa}, if the head is liable to error, the body can not be free from error. The Gallican theory is an untenable \emph{via media}. It secures only a periodic and intermittent infallibility, which reveals itself in an œcumenical Council, and then relapses into a quiescent state; but the Ultramontane theory teaches an unbroken, ever living, and ever active infallibility, which alone can fully answer the demands of an absolute authority.

To refute Papal infallibility is to refute also Episcopal infallibility; for the higher includes the lower. The Vatican Council is the best argument against the infallibility of œcumenical Councils, for it sanctioned a fiction, in open and irreconcilable contradiction to older œcumenical Councils, which not only assumed the possibility of Papal fallibility, but actually condemned a Pope as a heretic. The fifth Lateran Council (1512) declared the decrees of the Council of Pisa (1409) null and void; the Council of Florence denied the validity of the Council of Basle, and this denied the validity of the former. The Council of Constance condemned and burned John Hus for teaching evangelical doctrines; and this fact forced upon Luther, at the disputation with Eck at Leipzig, the conviction that even œcumenical Councils may err. Rome itself has rejected certain canons of Constantinople and Chalcedon, which put the Pope on a par with the Patriarch of Constantinople; and a strict construction of the Papal theory would rule out the old œcumenical Councils, because they were not convened nor controlled by the Pope; while the Greek Church rejects all Councils which were purely Latin.

The Bible makes no provision and has no promise for an œcumenical Council.\footnote{The Synod of Jerusalem, composed of Apostles, Elders, and Brethren, and legislating in favor of Christian liberty, differs very widely from a purely hierarchical Council, which excludes Elders and Brethren, and imposes new burdens upon the conscience.} The Church existed and flourished for more than three hundred years before such a Council was heard of. Large assemblies are often ruled by passion, intrigue, and worldly ambition (remember the complaints of Gregory of Nazianzum on the Synods of the Nicene age). Majorities are not necessarily decisive in matters of faith. Christ promised to be even with two or three who are gathered in his name (Matt. xviii. 20). Elijah and the seven thousand who had not bowed the knee to Baal were right over against the great mass of the people of Israel. Athanasius \emph{versus mundum} represented the truth, and the world \emph{versus Athanasium }was in error during the ascendency of Arianism. In the eighteenth century the Church, both Catholic and Protestant, was under the power of infidelity, and true Christianity had to take refuge in small communities. Augustine maintained that one Council may correct another, and attain to a more perfect knowledge of truth. The history of the Church is unintelligible without the theory of progressive development, which implies many obstructions and temporary diseases. All the attributes of the Church are subject to the law of gradual expansion and growth, and will not be finally complete till the second coming of our Lord.

\emph{Papal Infallibility and Personal Responsibility.}

The Christian Church, as a divine institution, can never fail and never lose the truth. Christ has pledged his Spirit and life-giving presence to his people to the end of time, and even to two or three of his humblest disciples assembled in \emph{his }name; yet they are not on that account infallible. He gave authority in matters of discipline to every local Church (Matt. xviii. 17); and yet no one claims infallibility to every congregation. The Holy Spirit will always guide believers into the truth, and the unerring Word of God can never perish. But local churches, like individuals, may fall into error, and be utterly destroyed from the face of the earth. The true Church of Christ always makes progress, and will go on conquering and to conquer to the end of the world. But the particular churches of Jerusalem, Antioch, Alexandria, Constantinople, Asia Minor, and North Africa, where once the Apostles and St. Augustine taught, have disappeared, or crumbled into ruin, or have been overrun by the false prophet.

The truth will ever be within the reach of the sincere inquirer wherever the gospel is preached and the sacraments are rightly administered. God has revealed himself plainly enough for all purposes of salvation; and yet not so plainly as to supersede the necessity of faith, and to resolve Christianity into a mathematical demonstration. He has given us a rational mind to think and to judge, and a free will to accept or to reject. Christian faith is no blind submission, but an intelligent assent. It implies anxiety to inquire as well as willingness to receive. We are expressly directed to 'prove all things, and to hold fast that which is good' (1 Thess. v. 21); to try the spirits whether they are of God (1 John iv. 1), and to refuse obedience even to an angel from heaven if he preach a different gospel (Gal. i. 8). The Berœan Jews are commended as being more noble than those of Thessalonica, because they received the Word with all readiness of mind, and yet searched the Scriptures daily, whether those things were so (Acts xvii. 11). It was from the infallible Scriptures alone, and not from tradition, that Paul and Apollos reasoned, after the example of Christ, who appeals to Moses and the Prophets, and speaks disparagingly of the traditions of the elders as obscuring the Word of God or destroying its true effect.\footnote{It is remarkable that Christ always uses \textgreek{παράδοσις} in an unfavorable sense: see Matt. xv. 2, 3, 6; Mark vii. 3, 5, 8, 9, 13. So also Paul: Gal. i. 14; Col. ii. 8; while in 1 Cor. xi. 2, and 2 Thess. ii. 15; iii. 6, he uses the term in a good sense, as identical with the gospel he preached.}

In opposition to all this the Vatican dogma requires a wholesale slaughter of the intellect and will, and destroys the sense of personal responsibility. The fundamental error, the \textgreek{πρῶτον ψεῦδος} of Rome is that she identifies the true ideal Church of Christ with the empirical Church, and the empirical Church with the Romish Church, and the Romish Church with the Papacy, and the Papacy with the Pope, and at last substitutes a mortal man for the living Christ, who is the only and ever present head of the Church, 'which is his body, the fullness of him who filleth all in all.' Christ needs no vicar, and the very idea of a vicar implies the absence of the Master.\footnote{I add here what Dr. Hodge, of Princeton, says on the Papal theory of Infallibility (Systematic Theology, New York, 1872, Vol. I. pp. 130, 150): 'There is something simple and grand in this theory. It is wonderfully adapted to the tastes and wants of men. It relieves them of personal responsibility. Every thing is decided for them. Their salvation is secured by merely submitting to be saved by an infallible, sin-pardoning, and grace-imparting Church. Many may be inclined to think that it would have been a great blessing had Christ left on earth a visible representative of himself, clothed with his authority to teach and govern, and an order of men dispersed through the world endowed with the gifts of the original Apostles—men every where accessible, to whom we could resort in all times of difficulty and doubt, and whose decisions could be safely received as the decisions of Christ himself. God's thoughts, however, are not as our thoughts. We know that when Christ was on earth men did not believe or obey him. We know that when the Apostles were still living, and their authority was still confirmed by signs, and wonders, and divers miracles and gifts of the Holy Ghost, the Church was distracted by heresies and schisms. If any in their sluggishness are disposed to think that a perpetual body of infallible teachers would be a blessing, all must admit that the assumption of infallibility by the ignorant, the erring, and the wicked, must be an evil inconceivably great. The Romish theory, if true, might be a blessing; if false, it must be an awful curse. That it is false may be demonstrated to the satisfaction of all who do not wish it to be true, and who, unlike the Oxford tractarian, are not determined to believe it because they love it. . . . If the Church be infallible, its authority is no less absolute in the sphere of social and political life. It is immoral to contract or to continue an unlawful marriage, to keep an unlawful oath, to enact unjust laws, to obey a sovereign hostile to the Church. The Church, therefore, has the right to dissolve marriages, to free men from the obligations of their oaths, and citizens from their allegiance, to abrogate civil laws, and to depose sovereigns. These prerogatives have not only been claimed, but time and again exercised by the Church of Rome. They all of right belong to that Church, if it be infallible. As these claims are enforced by penalties involving the loss of the soul, they can not be resisted by those who admit the Church to be infallible. It is obvious, therefore, that where this doctrine is held there can be no liberty of opinion, no freedom of conscience, no civil or political freedom. As the recent œcumenical Council of the Vatican has decided that this infallibility is vested in the Pope, it is henceforth a matter of faith with Romanists, that the Roman Pontiff is the absolute sovereign of the world. All men are bound, on the penalty of eternal death, to believe what he declares to be true, and to do whatever he decides is obligatory.'}

\emph{Papal Infallibility tested by Tradition.}

The dogma of Papal Infallibility is mainly supported by an inferential dogmatic argument derived from the Primacy of Peter, who, as the Vicar of Christ, must also share in his infallibility; or from the nature and aim of the Church, which is to teach men the way of salvation, and must therefore be endowed with an infallible and ever available organ for that purpose, since God always provides the means together with an end. A full-blooded Infallibilist, whose piety consists in absolute submission and devotion to his lord the Pope, is perfectly satisfied with this reasoning, and cares little or nothing for the Bible and for history, except so far as they suit his purpose. If facts disagree with his dogmas, all the worse for the facts. All you have to do is to ignore or to deny them, or to force them, by unnatural interpretations, into reluctant obedience to the dogmas.\footnote{Archbishop Manning (III. p. 118) speaks of history as 'a wilderness without guide or path,' and says: 'Whensoever any doctrine is contained in the divine revelation of the Church' [the very point which can not be proved in the case before us], 'all difficulties from human history are excluded, as Tertullian lays down, by prescription. The only source of revealed truth is God; the only channel of his revelation is the Church. No human history can declare what is contained in that revelation. The Church alone can determine its limits, and therefore its contents.'} But after all, even according to the Roman Catholic theory, Scripture and history or tradition are the two indispensable tests of the truth of a dogma. It has always been held that the Pope and the Bishops are not the creators and judges, but the trustees and witnesses of the apostolic deposit of faith, and that they can define and proclaim no dogma which is not well founded in primitive tradition, written or unwritten. According to the famous rule of Vincentius Lirinensis, a dogma must have three marks of catholicity: the catholicity of time (\emph{semper}), of space (\emph{ubique}), and of number (\emph{ab omnibus}). The argument from tradition is absolutely essential to orthodoxy in the Roman sense, and, as hitherto held, more essential than Scripture proof.\footnote{This Archbishop Kenrick, in his \emph{Concio}, frankly admits: '\emph{Irenæi, Tertulliani, Augustini, Vincentii Lirinensis exempla secutus, fidei Catholicæ probationes ex traditione potius quam ex Scripturarum interpretatione quærendas duxi; quæ interpretatio, juxta Tertullianum magis apta est ad veritatem obumbitandum quam demonstrandum.}'} The difference between Romanism and Protestantism on this point is this: Romanism requires proof from tradition first, from Scripture next, and makes the former indispensable, the latter simply desirable; while Protestantism reverses the order, and with its theory of the Bible as the only rule of faith and practice, and as an inexhaustible mine of truth that yields precious ore to every successive generation of miners, it may even dispense with traditional testimony altogether, provided that a doctrine can be clearly derived from the Word of God.

Now it can be conclusively proved that the dogma of Papal Infallibility, like the dogma of the Immaculate Conception of Mary, lacks every one of the three marks of catholicity. It is a comparatively modern innovation. It was not dreamed of for more than a thousand years, and is unknown to this day in the Greek Church, the oldest in the world, and in matters of antiquity always an important witness. The whole history of Christianity would have taken a different course, if in all theological controversies an infallible tribunal in Rome could have been invoked.\footnote{'\emph{Die ganze Geschichte des ersten Jahrtausends der Kirche wäre eine andere gewesen, wenn in dem Bischof von Rom das Bewusstsein, in der Kirche auch nur eine Ahnung davon gewesen wäre, dass dort ein Quell unfehlbarer Wahrheit fliesse. Statt all der bittern, verstörenden Kämpfe gegen wirkliche oder vermeintliche Häretiker, gegen die man Bücher schrieb und Synoden aller Art versammelte, würden alle Wohlmeinende sich auf den unfehlbaren Spruch des Papstes berufen haben, und mehr als einst das Orakel des Apollo zu Delphi würde das zu Rom befragt worden sein. Dagegen war es in jenen Jahrhunderten, als alles Christenthum auf die Spitze eines Dogmas gestellt wurde, nichts unerhörtes, dass auch ein Papst vor der subtilen Bestimmung des siegenden Dogma zum Häretiker wurde.}' Hase, \emph{Polemik}, Buch I. c.iv. p. 161.} Ancient Creeds, Councils, Fathers, and Popes can be summoned as witnesses against the Vatican dogma.

1. The four \emph{œcumenical Creeds}, the most authoritative expressions of the old Catholic faith of the Eastern and Western Churches, contain an article on the 'holy Catholic and Apostolic Church,' but not one word about the Bishops of Rome, or any other local Church. How easy and natural, yea, in view of the fundamental importance of the Infallibility dogma, how necessary would have been the insertion of \emph{Roman }after the other predicates of the Church, or the addition of the article: 'The Pope of Rome, the successor of Peter and infallible vicar of Christ.' If it had been believed then as now, it would certainly appear at least in the Roman form of the Apostles' Creed; but this is as silent on this point as the Aquilejan, the African, the Gallican, and other forms.

And this uniform silence of all the œcumenical Creeds is strengthened by the numerous local Creeds of the Nicene age, and by the various ante-Nicene rules of faith up to Tertullian and Irenæus, not one of which contains an allusion to such an article of faith.

2. The \emph{œcumenical Councils }of the first eight centuries, which are recognized by the Greek and Latin Churches alike, are equally silent about, and positively inconsistent with, Papal Infallibility. They were called by Greek Emperors, not by Popes; they were predominantly, and some of them exclusively, Oriental; they issued their decrees in their own name, and in the fullness of authority, without thinking of submitting them to the approval of Rome; they even claimed the right of judging and condemning the Roman Pontiff, as well as any other Bishop or Patriarch.

In the first Nicene Council there was but one representative of the Latin Church (Hosius of Spain); and in the second and the fifth œcumenical Councils there was none at all. The second œcumenical Council (381), in the third canon, put the Patriarch of Constantinople on a par with the Bishop of Rome, assigning to the latter only a primacy of honor; and the fourth œcumenical Council (451) confirmed this canon in spite of the energetic protest of Pope Leo I.

But more than this: the sixth œcumenical Council, held 680, pronounced the anathema on Honorius, 'the former Pope of old Rome,' for teaching officially the Monothelite heresy; and this anathema was signed by all the members of the Council, including the three delegates of the Pope, and was several times repeated by the seventh and eighth Councils, which were presided over by Papal delegates. But we must return to this famous case again in another connection.

3. The \emph{Fathers}, even those who unconsciously did most service to Rome, and laid the foundation for its colossal pretensions, yet had no idea of ascribing absolute supremacy and infallibility to the Pope.

Clement of Rome, the first Roman Bishop of whom we have any authentic account, wrote a letter to the Church at Corinth—not in his name, but in the name of the Roman Congregation; not with an air of superior authority, but as a brother to brethren—barely mentioning Peter, but eulogizing Paul, and with a clear consciousness of the great difference between an Apostle and a Bishop or Elder.

Ignatius of Antioch, who suffered martyrdom in Rome under Trajan, highly as he extols Episcopacy and Church unity in his seven Epistles, one of which is addressed to the Roman Christians, makes no distinction of rank among Bishops, but treats them as equals.

Irenæus of Lyons, the champion of the Catholic faith against the Gnostic heresy at the close of the second century, and the author of the famous and variously understood passage about the \emph{potentior principalitas} (\textgreek{προτεία}) \emph{ecclesiæ Romanæ}, sharply reproved Victor of Rome when he ventured to excommunicate the Asiatic Christians for their different mode of celebrating Easter, and told him that it was contrary to Apostolic doctrine and practice to judge brethren on account of eating and drinking, feasts and new moons. Cyprian, likewise a saint and a martyr, in the middle of the third century, in his zeal for visible and tangible unity against the schismatics of his diocese, first brought out the fertile doctrine of the Roman See as the chair of Peter and the centre of Catholic unity; yet with all his Romanizing tendency he was the great champion of the Episcopal solidarity and equality system, and always addressed the Roman Bishop as his 'brother' and 'colleague;' he even stoutly opposed Pope Stephen's view of the validity of heretical baptism, charging him with error, obstinacy, and presumption. He never yielded, and the African Bishops, at the third Council at Carthage (256), emphatically indorsed his opposition. Firmilian, Bishop of Cæsarea, and Dionysius, Bishop of Alexandria, likewise bitterly condemned the doctrine and conduct of Stephen, and told him that in excommunicating others he only excommunicated himself.

Augustine is often quoted by Infallibilists on account of his famous dictum, \emph{Roma locuta est, causa finita est.}\footnote{Or in a modified form: '\emph{Causa finita est, utinam aliquando finiatur error!}' \emph{Serm.} 131, c. 10. See Janus, Rauscher, von Schulte \emph{versus} Cardoni and Hergenröther, quoted by Frommann, p. 424.} But he simply means that, since the Councils of Mileve and Carthage had spoken, and Pope Innocent I. had acceded to their decision, the Pelagian controversy was finally settled (although it was, after all, not settled till after his death, at the Council of Ephesus). Had he dreamed of the abuse made of this utterance,\footnote{As well as some other of his sententious sayings. His explanation of \emph{coge intrare }was made to justify religious persecutions, from which his heart would have shrunk in horror.} he would have spoken very differently. For the same Augustine apologized for Cyprian's opposition to Pope Stephen on the ground that the controversy had then not yet been decided by a Council, and maintained the view of the liability of Councils to correction and improvement by subsequent Councils. He moreover himself opposed Pope Zosimus, when, deceived by Pelagius, he declared him sound in the faith, although Pope Innocent I. had previously excommunicated him as a dangerous heretic. And so determined were the Africans, under the lead of Augustine (417 and 418), that Zosimus finally saw proper to yield and to condemn Pelagianism in his '\emph{Epistula Tractoria}.'

Gregory I., or the Great, the last of the Latin Fathers, and the first of the mediæval Popes (590–604), stoutly protested against the assumption of the title \emph{œcumenical }or \emph{universal }Bishop on the part of the Patriarchs of Constantinople and Alexandria, and denounced this whole title and claim as \emph{blasphemous}, \emph{anti-Christian}, and \emph{devilish}, since Christ alone was the Head and Bishop of the Church universal, while Peter, Paul, Andrew, and John, were members under the same Head, and heads only of single portions of the whole. Gregory would rather call himself 'the servant of the servants of God,' which, in the mouths of his successors, pretending to be Bishops of bishops and Lords of lords, has become a shameless irony.\footnote{The passages of Gregory on this subject are well known to every scholar. And yet the Vatican decree, in ch. iii., by omitting the principal part, makes him say almost the very opposite.}

As to the Greek Fathers, it would be useless to quote them, for the entire Greek Church in her genuine testimonies has never accepted the doctrine of Papal supremacy, much less of Papal Infallibility.

4. \emph{Heretical Popes}.—We may readily admit the rock-like stability of the Roman Church in the early controversies on the Trinity and the Divinity of Christ, as compared with the motion and changeability of the Greek churches during the same period, when the East was the chief theatre of dogmatic controversy and progress. Without some foundation in history, the Vatican dogma could not well have arisen. It would be impossible to raise the claim of infallibility in behalf of the Patriarchs of Jerusalem, or Antioch, or Alexandria, or Constantinople, among whom were noted Arians, Nestorians, Monophysites, Monothelites, and other heretics. Yet there are not a few exceptions to the rule; and as many Popes, in their lives, flatly contradicted their title of holiness, so many departed, in their views, from Catholic truth. That the Popes after the Reformation condemned and cursed Protestant truths well founded in the Scriptures, we leave here out of sight, and confine our reasoning to facts within the limits of Roman Catholic orthodoxy.

The canon law assumes throughout that a Pope may openly teach heresy, or contumaciously contradict the Catholic doctrine; for it declares that, while he stands above all secular tribunals, yet he can be judged and deposed for the crime of heresy.\footnote{\emph{Decret}. Gratian. Dist. xl. c. 6, in conformity with the sentence of Hadrian II.: '\emph{Cunctos ipsos judicaturus} [\emph{Papa}], \emph{a nemine est judicandus}, nisi deprehendatur a fide devius.' See on this point especially von Schulte, \emph{Concilien}, pp. 188 sqq.} This assumption was so interwoven in the faith of the Middle Ages that even the most powerful of all Popes, Innocent III. (d. 1216), gave expression to it when he said that, though he was only responsible to God, he may sin against the faith, and thus become subject to the judgment of the Church.\footnote{\emph{Serm. II. de consecrat. Pontificis: 'In tantum mihi fides necessaria est, cum de cæteris peccatis Deum judicem habeam, ut propter solum peccatum quod in fidem committitur, possim ab Ecclesia judicari.'}} Innocent IV. (d. 1254) speaks of heretical commands of the Pope, which need not be obeyed. When Boniface VIII. (d. 1303) declared that every creature must obey the Pope at the loss of eternal salvation, he was charged with having a devil, because he presumed to be infallible, which was impossible without witchcraft. Even Hadrian VI., in the sixteenth century, expressed the view, which he did not recant as Pope, that 'if by the Roman Church is understood its head, the Pope, it is certain that he can err even in matters of faith.'

This old Catholic theory of the fallibility of the Pope is abundantly borne out by actual facts, which have been established again and again by Catholic scholars of the highest authority for learning and candor. We need no better proofs than those furnished by them.

Zephyrinus (201–219) and Callistus (219–223) held and taught (according to the 'Philosophumena' of Hippolytus, a martyr and saint) the Patripassian heresy, that God the Father became incarnate and suffered with the Son.

Pope Liberius, in 358, subscribed an Arian creed for the purpose of regaining his episcopate, and condemned Athanasius, 'the father of orthodoxy,' who mentions the fact with indignation.

During the same period, his rival, Felix II., was a decided Arian; but there is a dispute about his legitimacy; some regarding him as an anti-Pope, although he has a place in the Romish Calendar of Saints, and Gregory XIII. (1582) confirmed his claim to sanctity, against which Baronius protested.

In the Pelagian controversy, Pope Zosimus at first indorsed the orthodoxy of Pelagius and Celestius, whom his predecessor, Innocent I., had condemned; but he yielded afterwards to the firm protest of St. Augustine and the African Bishops.

In the Three-Chapter controversy, Pope Vigilius (538–555) showed a contemptible vacillation between two opinions: first indorsing; then, a year afterwards, condemning (in obedience to the Emperor's wishes) the Three Chapters (i.e., the writings of Theodore, Theodoret, and Ibas); then refusing the condemnation; then, tired of exile, submitting to the fifth œcumenical Council (553), which had broken off communion with him; and confessing that he had unfortunately been the tool of Satan, who labors for the destruction of the Church. A long schism in the West was the consequence. Pope Pelagius II. (585) significantly excused this weakness by the inconsistency of St. Peter at Antioch.

John XXII. (d. 1334) maintained, in opposition to Nicholas III. and Clement V. (d. 1314), that the Apostles did not live in perfect poverty, and branded the opposite doctrine of his predecessors as heretical and dangerous. He also held an opinion concerning the middle state of the righteous, which was condemned as heresy by the University of Paris.

Contradictory opinions were taught by different Popes on the sacraments, on the immaculate conception of the Virgin Mary (see p. 123), on matrimony, and on the subjection of the temporal power to the Church.\footnote{See examples under this head in \emph{Janus, }pp. 54 sqq. (\emph{Irrthümer and Widersprüche der Päpste}), p. 51 of the London ed.}

But the most notorious case of an undeniably \emph{official }indorsement of heresy by a Pope is that of Honorius I. (625–638), which alone is sufficient to disprove Papal Infallibility, according to the maxim: \emph{Falsus in uno, falsus in omnibus.}\footnote{Or, as Perrone, himself an Infallibilist, who in his Dogmatic Theology characteristically treats of the Pope before the Holy Scriptures and tradition, puts it: '\emph{Si vel unicus ejusmodi error deprehenderetur, appareret omnes adductas probationes in nihilum redactum iri.}'} This case has been sifted to the very bottom before and during the Council, especially by Bishop Hefele and Père Gratry. The following decisive facts are established by the best documentary evidence:

(1.) Honorius taught \emph{ex cathedra }(in two letters to his heretical colleague, Sergius, Patriarch of Constantinople) the Monothelite heresy, which was condemned by the sixth œcumenical Council, i.e., the doctrine that Christ had only one will, and not two (corresponding to his two natures).\footnote{Honorius prescribed the technical term of the Monothelites as a dogma to the Church (\emph{dogma ecclesiasticum}). In a reply to the Monothelite Patriarch Sergius of Constantinople, which is still extant in Greek and Latin (Mansi, \emph{Coll. Concil.} Tom. XI. pp. 538 sqq.), he approves of his heretical view, and says as clearly as words can make it: 'Therefore we confess also \emph{one will} (\textgreek{ἓν θέλημα}) of our Lord Jesus Christ, since the Godhead has assumed our \emph{nature}, but not our guilt.' In a second letter to Sergius, of which we have two fragments (Mansi, l.c. p. 579), Honorius rejects the orthodox term \emph{two energies} (\textgreek{δύο ἐνέργειαι,} \emph{duæ operationes}), which is used alongside with \emph{two wills} (\textgreek{δύο θελήματα,} \emph{voluntates}). Christ, he reasons, assumed human nature as it was before the fall, when it had not a law in the members which resists the law of the Spirit. He knew only a \emph{sinful }human will. The Catholic Church rejects Monothelitism, or the doctrine of \emph{one }will of Christ, as involving or necessarily leading to Monophysitism, i.e., the doctrine that Christ had but \emph{one }nature; for will is an attribute of \emph{nature, }not of the \emph{person. }The Godhead has three persons, but only one nature, and only one will. Christ has two wills, because he has two natures. The compromise formula of Emperor Heraclius and Patriarch Sergius of Constantinople endeavored to reconcile the Monophysites with the orthodox Church by teaching that Christ had two natures, but only one will and one energy.}

(2.) An œcumenical Council, universally acknowledged in the East and in the West, held in Constantinople, 680, condemned and excommunicated Honorius, 'the former Pope of Old Rome,' as a heretic, who with the help of the old serpent had scattered deadly error.\footnote{Sessio XVI.: '\emph{Sergio hæretico anathema, Cyro hæretico anathema, Honorio hæretico anathema.}' . . . Sessio XVIII.: '\emph{Honorius, qui fuit Papa antiquæ Romæ . . . non vacavit . . . Ecelesiæ erroris scandalum suscitare unius voluntatis, et unius operationis in duabus naturis unius Christi,}' etc. See Mansi, \emph{Conc.} Tom. XI. pp. 622, 635, 655, 666.} The seventh œcumenical Council (787) and the eighth (869) repeated the anathema of the sixth.

(3.) The succeeding Popes down to the eleventh century, in a solemn oath at their accession, indorsed the sixth œcumenical Council, and pronounced 'an eternal anathema' on the authors of the Monothelite heresy, together with Pope Honorius, because he had given aid and comfort to the perverse doctrines of the heretics.\footnote{'\emph{Quia pravis hæreticorum assertionibus fomentum impendit.}' This Papal oath was probably prescribed by Gregory II. (at the beginning of the eighth century), and is found in the Liber Diurnus (the book of formularies of the Roman chancery from the fifth to the eleventh century), edited by Eugène de Rozière, Paris, 1869, No. 84. The Liber Pontificalis agrees with the Liber Diurnus. Editions of the Roman Breviary down to the sixteenth century reiterated the charge against Honorius, since silently dropped.} The Popes themselves, therefore, for more than three centuries, publicly recognized, first, that an œcumenical Council may condemn a Pope for open heresy, and, secondly, that Pope Honorius was justly condemned for heresy. Pope Leo II., in a letter to the Emperor, strongly confirmed the decree of the Council, and denounced his predecessor Honorius as one who 'endeavored by profane treason to overthrow the immaculate faith of the Roman Church.'\footnote{' \emph{Nec non et Honorium }[\emph{anathematizamus}], \emph{ qui hanc apostolicam ecclesiam non apostolicæ traditionis doctrina lustravit, sed profana proditione immaculatam fidem subvertere conatus est.}' Mansi, Tom. XI. p. 731.} The same Pope says, in a letter to the Spanish Bishops: 'With eternal damnation have been punished Theodore, Cyrus, Sergius—\emph{together with Honorius}, who did not extinguish at the very beginning the flame of heretical doctrine, as was becoming to his apostolic authority, but nursed it by his carelessness.'\footnote{'\emph{Cum Honorio, qui flammam hæeretici dogmatis, non ut decuit apostolicam auctoritatem, incipientem extinxit, sed negligendo confovit.}' Mansi, p. 1052.}

This case of Honorius is as clear and strong as any fact in Church history.\footnote{Comp. especially the tract of Bishop Hefele, above quoted. The learned author of the History of the Councils has proved the case as conclusively as a mathematical demonstration.} Infallibilists have been driven to desperate efforts. Some pronounce the acts of the Council, which exist in Greek and Latin, downright forgeries (Baronius); others, admitting the acts, declare the letters of Honoring forgeries, so that he was unjustly condemned by the Council (Bellarmin)—both without a shadow of proof; still others, being forced at last to acknowledge the genuineness of the letters and acts, distort the former into an orthodox sense by a non-natural exegesis, and thus unwillingly fasten upon œcumenical Councils and Popes the charge of either dogmatic ignorance and stupidity, or malignant representation.\footnote{So Perrone, in his \emph{Dogmatics}, and Pennachi, in his \emph{Liber de Honorii I. Rom. Pont. causa}, 1870, which is effectually disposed of by Hefele in an Appendix to the German edition of his tract. Nevertheless, Archbishop Manning, sublimely ignoring all but Infallibilist authorities on Honorius, has the face to assert (III. p. 223) that the case of Honorius is doubtful; that he defined no doctrine whatever; and that his two epistles are entirely orthodox! Is Manning more infallible than the infallible Pope Leo II., who denounced Honorius \emph{ex cathedra} as a heretic?} Yet in every case the decisive fact remains that both Councils and Popes for several hundred years believed in the fallibility of the Pope, in flat contradiction to the Vatican Council. Such acts of violence upon history remind one of King James's short method with Dissenters: 'Only hang them, that's all.'

5. The idea of Papal absolutism and Infallibility, like that of the sinlessness of Mary, can be traced to apocryphal origin. It is found first, in the second century, in the pseudo-Clementine Homilies, which contain a singular system of speculative Ebionism, and represent James of Jerusalem, the brother of the Lord, as the Bishop of Bishops, the centre of Christendom, and the general Vicar of Christ; he is the last arbiter, from whom there is no appeal; to him even Peter must give an account of his labors, and to him the sermons of Peter were sent for safe keeping.\footnote{See my Church History, Vol. I. § 69, p. 219, and the tract of Lutterbeck above quoted.}

In the Catholic Church the same idea, but transferred to the Bishop of Rome, is first clearly expressed in the pseudo-Isidorian Decretals, that huge forgery of Papal letters, which appeared in the middle of the ninth century, and had for its object the completion of the independence of the Episcopal hierarchy from the State, and the absolute power of the Popes, as the legislators and judges of all Christendom. Here the most extravagant claims are put into the mouths of the early Popes, from Clement (91) to Damasus (384), in the barbarous French Latin of the Middle Ages, and with such numerous and glaring anachronisms as to force the conviction of fraud even upon Roman Catholic scholars. One of these sayings is: 'The Roman Church remains to the end free from stain of heresy.' Soon afterwards arose, in the same hierarchical interest, the legend of the donation of Constantine and his baptism by Pope Silvester, interpolations of the writings of the Fathers, especially Cyprian and Augustine, and a variety of fictions embodied in the Gesta Liberii, and the Liber Pontificalis, and sanctioned by Gratianus (about 1150) in his \emph{Decretum}, or collection of canons, which (as the first part of the Corpus juris canonici) became the code of laws for the whole Western Church, and exerted an extraordinary influence. By this series of pious frauds the mediæval Papacy, which was the growth of ages, was represented to the faith of the Church as a primitive institution of Christ, clothed with absolute and perpetual authority.

The Popes since Nicholas I. (858–867), who exceeded all his predecessors in the boldness of his designs, freely used what the spirit of a hierarchical, superstitious, and uncritical age furnished them. They quoted the fictitious letters of their predecessors as genuine, the Sardican canon on appeals as a canon of Nicæa, and the interpolated sixth canon of Nicæa,' the Roman Church always had the primacy,' of which there is not a syllable in the original; and nobody doubted them. Papal absolutism was in full vigor from Gregory VII. to Boniface VIII. Scholastic divines, even Thomas Aquinas, deceived by these literary forgeries, began to defend Papal absolutism over the whole Church, and the Councils of Lyons (1274) and of Florence (1439) sanctioned it, although the Greeks soon afterwards rejected the false union based upon such assumption.

But absolute power, especially of a spiritual kind, is invariably intoxicating and demoralizing to any mortal man who possesses it. God Almighty alone can bear it, and even he allows freedom to his rational creatures. The reminiscence of the monstrous period when the Papacy was a football in the hands of bold and dissolute women (904–962), or when mere boys, like Benedict IX. (1033), polluted the Papal crown with the filth of unnatural vices, could not be quite forgotten. The scandal of the Papal schism (1378 to 1409), when two and even three rival Popes excommunicated and cursed each other, and laid all Western Christendom under the ban, excited the moral indignation of all good men in Christendom, and called forth, in the beginning of the fifteenth century, the three Councils of Pisa, Constance, and Basle, which loudly demanded a reformation of the Church, in the head as well as in the members, and asserted the superiority of a Council over the Pope.

The Council of Constance (1414–1418), the most numerous ever seen in the West, deposed two Popes—John XXIII. (the infamous Balthasar Cossa, who had been recognized by the majority of the Church), on the charge of a series of crimes (May 29, 1415), and Benedict XIII., as a heretic who sinned against the unity of the Church (July 26, 1417),\footnote{The third anti-Pope, Gregory XII., resigned.} and elected a new Pope, Martin V. (Nov. 11, 1517), who had given his adhesion to the Council, though after his accession to power he found ways and means to defeat its real object, i.e., the reformation of the Church.

This Council was a complete triumph of the Episcopal system, and the Papal absolutists and Infallibilists are here forced to the logical dilemma of either admitting the validity of the Council, or invalidating the election of Martin V. and his successors. Either course is fatal to their system. Hence there has never been an \emph{authoritative} decision on the œcumenicity of this Council, and the only subterfuge is to say that the whole case is an extraordinary exception; but this, after all, involves the admission that there is a higher power in the Church over the Papacy.

The Reformation shook the whole Papacy to its foundation, but could not overthrow it. A powerful reaction followed, headed by the Jesuits. Their General, Lainez, strongly advocated Papal Infallibility in the Council of Trent, and declared that the Church could not err only because the Pope could not err. But the Council left the question undecided, and the Roman Catechism ascribes infallibility simply to 'the Catholic Church,' without defining its seat. Bellarmin advocated and formularized the doctrine, stating it as an almost general opinion that the Pope could not publicly teach a heretical dogma, and as a probable and pious opinion that Providence will guard him even against private heresy. Yet the same Bellarmin was witness to the innumerable blunders of the edition of the Latin Vulgate prepared by Sixtus V., corrected by his own hand, and issued by him as the only true and authentic text of the sacred Scriptures, with the stereotyped forms of anathema upon all who should venture to change a single word; and Bellarmin himself gave the advice that all copies should be called in, and a new edition printed with a lying statement in the preface making the printers the scape-goats for the errors of the Pope! This whole business of the Vulgate is sufficient to explode Papal Infallibility; for it touches the very source of divine revelation. Other Italian divines, like Alphonsus Liguori, and Jesuitical text-books, unblushingly use long-exploded mediæval fictions and interpolations as a groundwork of Papal absolutism and Infallibility.

It is not necessary to follow the progress of the controversy between the Episcopal and the Papal systems during the seventeenth and eighteenth centuries. It is sufficient to say that the greatest Catholic divines of France and Germany, including Bossuet and Möhler, together with many from other countries, down to the 88 protesting Bishops in the Vatican Council, were anti-Infallibilists; and that popular Catechisms of the Roman Church, extensively used till 1870, expressly denied the doctrine, which is now set up as an article of faith necessary to eternal salvation.\footnote{So Overberg's \emph{Katechismus}, III. Hauptstück, Fr. 349: '\emph{Müssen wir auch glauben, dass der Papst unfehlbar ist?}  Nein, dies ist kein Glaubensartikel.' Keenan's \emph{Controversial Catechism}, in the editions before 1871, declared Papal Infallibility to be 'a Protestant invention.' The Irish Bishops—Doyle, Murray, Kelly—affirmed under oath, before a Committee of the English Parliament in 1825, that the Papal authority is limited by Councils, that it does not extend to civil affairs and the temporal rights of princes, and that Papal decrees are not binding on Catholics without the consent of the whole Church, either dispersed or assembled in Council. See the original in the Appendix to Archbishop Kenrick's \emph{Concio} in Friedrich's \emph{Documenta}, I. pp. 228–242. But the Irish Catholics, who almost believe in the infallibility of their priests, can be very easily taught to believe in the infallibility of the Pope.}

\emph{Papal Infallibility and the Bible.}

The Old Testament gives no tangible aid to the Infallibilists. The Jewish Church existed as a divine institution, and served all its purposes, from Abraham to John the Baptist, without an infallible tribunal in Jerusalem, save the written law and testimony, made effective from time to time by the living voice of inspired prophecy. Pious Israelites found in the Scriptures the way of life, notwithstanding the contradictory interpretations of rabbinical schools and carnal perversions of Messianic prophecies, fostered by a corrupt hierarchy. The Urim and Thummim\footnote{That is, \textgreek{δήλωσις καὶ ἀλήθεια,} \emph{doctrina et veritas,} Ex. xxviii. 15–30; Deut. xxxiii. 8, 9; 1 Sam. xxviii. 6. The Urim and Thummim were inscribed on the garment of Aaron. Some interpreters identify them with the twelve stones on which the names of the tribes of Israel were engraved; others regard them as a plate of gold with the sacred name of Jehovah; still others as polished diamonds, in form like dice, which, being thrown on the table or Ark of the Covenant, were consulted as an oracle. See the able article of Plumptre, in Smith's \emph{Bible Dictionary}, Vol. IV. pp. 3356 sqq. (Am. ed.).} of the High-Priest has no doubt symbolical reference to some kind of spiritual illumination or oracular consultation, but it is of too uncertain interpretation to furnish an argument.

The passages of the New Testament which are used by Roman divines in support of the doctrine of Infallibility may be divided into two classes: those which seem to favor the Episcopal or Gallican, and those which are made to prove the Papal or Ultramontane theory. It is characteristic that the Papal Infallibilists carefully avoid the former.

1. To the first class belong John xiv. 16 sq.; xvi. 13–16, where Christ promises the Holy Ghost to his disciples that he may 'abide with them forever,' teach them 'all things,' bring to their remembrance all he had said to them,\footnote{The \textgreek{πάντα} implies a strong argument for the completeness of Christ's revelation in the New Testament against the Romish doctrine of addition.} and guide them 'into the whole truth;'\footnote{The phrase \textgreek{εἰς τὴν ἀλήθειαν πᾶσαν} (John xvi. 13), or, according to another reading, \textgreek{ἐν τῇ ἀληθείᾳ πάσῃ} (test. rec. \textgreek{ἐις πᾶσαν τὴν ἀλήθειαν}), expresses the truth as taught by Christ in its completeness—the \emph{whole }truth—and proves likewise the sufficiency of the Scriptures. The A.V. and its predecessors ('\emph{into all truth}'), also Luther (\emph{in alle Wahrheit}, instead of \emph{die ganze} or \emph{volle Wahrheit}), miss the true sense by omitting the article, and conveying the false idea that the Holy Ghost would impart to all the apostles a kind of omniscience. Comp. my annotations to Lange's \emph{John} on the passages (pp. 445, 478, etc.).} John xx. 21: 'As the Father hath sent me, even so send I you. . . . Receive ye the Holy Ghost;'\footnote{Literally: 'Receive Holy Spirit'—\textgreek{λάβετε πνεῦμα ἅγιον.} The absence of the article may indicate a partial or preparatory inspiration as distinct from the full Pentecostal effusion.} Matt. xviii. 18: 'Whatever ye shall bind on earth shall be bound in heaven,' etc.; Matt. xxviii. 19, 20: 'Go and disciple all nations . . . and lo, I am with you alway, even unto the end of the world.'

These passages, which are addressed to all Apostles alike, to doubting Thomas as well as to Peter, prove indeed the unbroken presence of Christ and the Holy Ghost in the Church to the end of time, which is one of the most precious and glorious truths admitted by every true Christian. But, in the first place, the Church, which is here represented by the Apostles, embraces all true believers, laymen as well as Bishops. Secondly, the promise of Christ's presence implies no infallibility, for the same promise is given even to the smallest number of true believers (Matt. xviii. 20). Thirdly, if the passages prove infallibility at all, they would prove individual infallibility by continued inspiration rather than corporate infallibility by official succession; for every Apostle was inspired, and so far infallible; and this no Roman Catholic Bishop, though claiming to be a successor of the Apostles, pretends to be.

2. The passages quoted by the advocates of the Papal theory are three, viz., Luke xxii. 31; Matt. xvi. 18; John xxi. 15.\footnote{Perrone and the Vatican decree on Infallibility confine themselves to these passages.}

We admit, at the outset, that these passages in their obvious meaning which is confirmed by the history of the Apostolic Church, assign to Peter a certain primacy among the Apostles: he was the leader and spokesman of them, and the chief agent of Christ in laying the foundations of his Church among the Jews and the Gentiles. This is significantly prophesied in the new name of Peter given to him. The history of Pentecost (Acts ii.) and the conversion of Cornelius (Acts x.) are the fulfillment of this prophecy, and furnish the key to the interpretation of the passages in the Gospels.

This is the truth which underlies the colossal lie of the Papacy. For there is no Romish error which does not derive its life and force from some truth.\footnote{Augustine says somewhere: '\emph{Nulla falsa doctrina est, quæ non aliquid veri permisceat.}'} But beyond this we have no right to go. The position which Peter occupied no one can occupy after him. The foundation of the Church, once laid, is laid for all time to come, and the gates of Hades can not prevail against it. The New Testament is its own best interpreter. It shows no single example of an exercise of jurisdiction of Peter over the other Apostles, but the very reverse. He himself, in his Epistles, disowns and prophetically warns his fellow-presbyters against the hierarchical spirit; exhorting them, instead of being lords over God's heritage, to be ensamples to his flock (1 Pet. v. 1–4). Paul and John were perfectly independent of him, as the Acts and Epistles prove. Paul even openly administered to him a rebuke at Antioch.\footnote{This fact is so obnoxious to Papists that some of them doubt or deny that the Cephas of Galatians ii. 11 was the Apostle Peter, although the New Testament knows no other. So Perrone, who also asserts, from his own preconceived theory, not from the text, that Paul withstood Peter from respectful love as an inferior to a superior, but not as a superior to an inferior! Let any Bishop try the same experiment against the Pope, and he will soon be sent to perdition.} At the Council of Jerusalem James seems to have presided, at all events he proposed the compromise which was adopted by the Apostles, Elders, and Brethren; Peter was indeed one of the leading speakers, but he significantly advocated the truly evangelical principle of salvation by faith alone, and protested against human bondage (Acts xv.; comp. Gal. ii.).

The great error of the Papacy is that it perverts a primacy of honor into a supremacy of jurisdiction, a personal privilege into an official prerogative, and a priority of time into a permanent superiority of rank. And to make the above passages at all available for such purpose, it must take for granted, as intervening links of the argument, that which can not be proved from the New Testament nor from history, viz., that Peter was Bishop of Rome; that he was there as Paul's superior; that he appointed a successor, and transferred to him his prerogatives.

As to the passages separately considered, Matt. xvi., 'Thou art rock,' and John xxi., 'Feed my flock,' could at best only prove Papal absolutism, but not Papal Infallibility, of which they do not treat.\footnote{For a full discussion of \textgreek{Πέτρος} and \textgreek{πέτρα,} see my edition of Lange's \emph{Comm. on} \emph{Matt. xvi. 18}, pp. 203 sqq.; and on the Romish perversion of the \textgreek{βόσκειν} and \textgreek{ποιμαίνειν τὰ ἀρνία, πρόβατα} and \textgreek{προβάτια} into a \textgreek{κατακυρειύειν,} and even withdrawal of nourishment, see my ed. of Lange on \emph{John, }pp. 638 sqq.} The former teaches the indestructibility of the Church in its totality (not of any individual congregation), but this is a different idea. The Council of Trent lays down 'the unanimous consent of the Fathers' as the norm and rule of all orthodox interpretation, as if exegetical wisdom had begun and ended with the divines of the first six centuries. But of the passage Matt. xvi., which is more frequently quoted by Popes and Papists than any other passage in the Bible, there are no less than five different patristic interpretations; the rock on which Christ built his Church being referred to \emph{Christ }by sixteen Fathers (including Augustine); to the \emph{faith }or \emph{confession }of Peter by forty-four (including Chrysostom, Ambrose, Hilary, Jerome, and Augustine again); \emph{to Peter }professing the faith by seventeen; to \emph{all the Apostles}, whom Peter represented by his primacy, by eight; to \emph{ all the faithful}, who, believing in Christ as the Son of God, are constituted the living stones of the Church.\footnote{This patristic dissensus was brought out during the Council in the \emph{Questio} distributed by Bishop Ketteler with all the proofs; see Friedrich, \emph{Docum. }I. pp. 6 sqq. Kenrick in his speech makes use of it. Comp. also my annotations to Lange's \emph{Comm. on Matthew }in \emph{loco.}} But not one of the Fathers finds Papal Infallibility in this passage, nor in John xxi. The 'unanimous consent of the Fathers' is a pure fiction, except in the most general and fundamental principles held by all Christians; and not to interpret the Bible \emph{except }according to the unanimous consent of the Fathers, would strictly mean not to interpret it at all.\footnote{Even Kenrick confesses that it is doubtful whether any instance of that unanimous consent can be found (in his \emph{Concio}, see Friedr. \emph{Docum.} I. p.195): '\emph{Regula interpetrandi Scripturas nobis imposita, hæc est: eas contra unanimem Patrum consensum non interpetrari. Si unquam detur consensus iste unanimis dubitari possit. Eo tamen deficiente, regula ista videtur nobis legem imponere majorem, qui ad unanimitatem accedere videretur, patrum numerum, in suis Scripturæ interpretationibus sequendi.}'}

There remains, then, only the passage recorded by Luke (xxii. 31, 32) as at all bearing on the disputed question: 'Simon, Simon, behold, Satan desired to have you (or, obtained you by asking), that he may sift you as wheat; but I prayed for thee, that thy faith fail not; and thou, when once thou art converted (or, hast turned again), strengthen thy brethren.' But even this does not prove infallibility, and has not been so understood before Popes Leo I. and Agatho. For (1) the passage refers, as the context shows, to the peculiar personal history of Peter during the dark hour of passion, and is both a warning and a comfort to him. So it is explained by the Fathers, who frequently quote it. (2) Faith here, as nearly always in the New Testament, means personal trust in, and attachment to, Christ, and not, as the Romish Church misinterprets it, orthodoxy, or intellectual assent to dogmas. (3) If the passage refers to the Popes at all, it would prove too much for them, viz., that they, like Peter, denied the Saviour, were converted again, and strengthened their brethren—which may be true enough of some, but certainly not of all.\footnote{This logical inference is also noticed by Archbishop Kenrick (\emph{Concio}, in Friedrich's \emph{Docum.} I. p. 200): '\emph{Præterea singula verba in ista Christi ad Petrum allocutione de Petri successoribus intelligi nequeunt, quin aliquid maxime absurdi exinde sequi videretur.} ``\emph{Tu autem conversus,}'' \emph{respiciunt certe conversionem Petri. Si priora verba}; ``\emph{orari pro te},'' \emph{et posteriora}: ``\emph{confirma fratres tuos,}'' \emph{ad successores Petri cœlestem vim, et munus transiisse probent, non videtur quarenam intermedia verba}: ``\emph{tu autem conversus},'' \emph{ad eos etiam pertinere, et aliquali sensu de eis intelligi, non debeant.}'}

The constant appeal of the Roman Church to Peter suggests a significant parallel. There is a spiritual Peter and a carnal Simon, who are separated, indeed, by regeneration, yet, after all, not so completely that the old nature does not occasionally re-appear in the new man.

It was the spiritual Peter who forsook all to follow Christ; who first confessed him as the Son of God, and hence was called Rock; who after his terrible fall wept bitterly; was re-instated and intrusted with the care of Christ's sheep; who on the birthday of the Church preached the first missionary sermon, and gathered in the three thousand converts; who in the Apostles' Council protested against the narrow bigotry of the Judaizers, and stood up with Paul for the principle of salvation by grace alone through faith in Christ; who, in his Epistles, warns all ministers against hierarchical pride, and exhibits a wonderful meekness, gentleness, and humility of spirit, showing that divine grace had overruled and sanctified to him even his fall; and who followed at last his Master to the cross of martyrdom.

It was the carnal Simon who presumed to divert his Lord from the path of suffering, and drew on him the rebuke, 'Get thee behind me, Satan; thou art a stumbling-block unto me, for thou mindest not the things of God, but the things of men;' the Simon, who in mistaken zeal used the sword and cut off the ear of Malchus; who proudly boasted of his unswerving fidelity to his Master, and yet a few hours afterwards denied him thrice before a servant-woman; who even after the Pentecostal illumination was overcome by his natural weakness, and, from policy or fear of the Judaizing party, was untrue to his better conviction, so as to draw on him the public rebuke of the younger Apostle of the Gentiles. The Romish legend of \emph{Domine quo vadis} makes him relapse into his inconstancy even a day before his martyrdom, and memorializes it in a chapel outside of Rome.

[In 1868, Cardinal Manning and Bishop Senestry of Regensburg, while in Rome, made a vow ``to do all in our power to bring about the definition of papal infallibility,'' the vow being attested by the Jesuit father Liberatore. See Purcell: \emph{Life of Manning}, II., 420. Commer, theological professor in Vienna, in an address on the twenty-fifth anniversary of Leo XIII.'s pontificate, announced that the Roman pontiff had properly been called by Catherine of Siena another Christ—\emph{alter Christus. }The Manual of the Catechism of Pius X. quotes with approval that the pope is Jesus Christ on earth—\emph{il papa è Gesu Cristo sulla terra.}—Ed.]

\xsection{The Liturgical Standards of the Roman Church.}

§ 35. The Liturgical Standards of the Roman Church.

Literature.

Missale Romanum,  \emph{ex decreto sacro-sancti Concilii Tridentini restitutum, S. Pii V., Pontificis Maximi, jussu editum, Clementis VIII. et Urbani VIII. auctoritate recognitum; in quo missæ novissimæ sanctorum accurate sunt dispositæ. }(Innumerable editions.)

Breviarium Romanum,  \emph{ex decreto SS. Concilii Tridentini restitutum, S. Pii V., Pontificis Maximi, jussu editum, Clementis VIII. et Urbani VIII. auctoritate recognitum, cum Officiis Sanctorum novissime per Summos Pontifices usque ad hunc diem concessis. }(The Paris and Lyons edition before me has over 1200 pp., with a Supplement of 127 pp. The Mechlin ed. of 1868 is in 4 vols.)

Pontificale Romanum,  \emph{Clementis VIII. ac Urbani VIII. jussu, editum, inde vero a Benedicto XIV. recognitum et castigatum. Cum Additionibus a Sacra Rituum Congregatione approbatis. }(The Mechlin ed. of 1845 is in three parts, with all the rules and directions printed in red; hence the word \emph{Rubrics}.)

George Lewis:  \emph{The Bible, the Missal, and the Breviary; or. Ritualism self-illustrated in the Liturgical Books of Rome}. Edinburgh, 1863, 2 vols.

A secondary symbolical authority belongs to those Latin liturgical works of the Roman Church which have been sanctioned by the Pope for use in public and private worship. They contain, in the form of devotion, nearly all the articles of faith, especially those referring to the sacraments and the cultus of saints and of the holy Virgin, and are, in a practical point of view, even of greater importance than the doctrinal standards, inasmuch as they are interwoven with the daily religious life of the priests.

Among these works the most important is the Missale Romanum, as issued by Pius V. in 1570, in compliance with a decree of the Council of Trent. It was subsequently revised again under Clement VIII. in 1604, and under Urban VIII. in 1634. The substance goes back to the early eucharistic services of the Latin Church, among which the principal ones are ascribed to Popes Leo I. (Sacramentarium Leonianum, probably from 483–492), Gelasius I. (Sacramentarium Gelasianum), and Gregory I. (Sacramentarium Gregorianum). But considerable diversity and confusion prevailed in provincial and local churches. Hence the Council of Trent ordered a new revision, under the direction of the Pope, with a view to secure uniformity. The Missal consists of three parts, besides Introduction and Appendix, viz.: (\emph{a}) The \emph{Proprium Missarum de Tempore,} or the services for the Sundays of the Christian year, beginning with the first Sunday in Advent, and closing with the last after Whitsuntide, all clustering around the great festivals of Christmas, Easter, and Pentecost. (\emph{b}) The \emph{Proprium Missarum de Sanctis} contains the forms for the celebration of mass on saints' days and other particular feasts, arranged according to the months and days of the civil year; the annually recurring death-days of saints being regarded as their celestial birth-days, (\emph{c}) The Commune Sanctorum is supplementary to the second part, and devoted to the celebration of the days of those saints for whom there is no special service provided in the Proprium. The Appendix to the Missal contains various masses and benedictions.

Next comes the Breviarum Romanum, revised by order of the Council of Trent, under Pius V., 1568, and again under Clement VIII., 1602, and finally brought into its present shape under Urban VIII., 1631. Since that time it has undergone no material changes, but received occasional additions of new festivals. The Breviary\footnote{The term \emph{Breviary} is derived from the \emph{abridgments} of the Scriptures and lives of saints contained therein, as distinct from the \emph{plenarium officium}; by others from the fact that later editions of the work are abridgments of former editions.} contains the prayers, psalms, hymns, Scripture lessons, and patristic comments not only for every Sunday, but for every day of the ecclesiastical year, together with the legends of saints and martyrs, presenting model characters and model devotions for each day, some of them good and harmless, others questionable, superstitious, and childish. The Breviary is a complete thesaurus of Romish piety, the private liturgy of the Romish priest, and to all intents and purposes his Bible. It regulates his whole religious life. It is divided into four parts, according to the four seasons; each part has the same four sections: the \emph{Psalterium,} the \emph{Proprium de Tempore,} the \emph{Proprium Sanctorum,} and the \emph{Commune Sanctorum.} The Introduction contains the ecclesiastical calendar. The office of each day consists of the seven or eight canonical hours of devotion, which are brought into connection with the history of the passion.\footnote{Matins, Lauds (3 A.M.), Prime (6 A.M.), Tierce (9 A.M.), Sext (12 M.), Nones (3 P.M.), Vespers (6 P.M.), and Compline (midnight devotion). The Nocturn is a night service. The custom of saying prayers at these hours goes back to the third century, and partly to Jewish tradition. Tertullian (\emph{De jejun.} c. 10) speaks of the \emph{tertia, sixta,} and \emph{nona} as \emph{apostolical} hours of prayer. On the mystical reference to Christ's passion, comp. the old memorial verse:  \par 'Hæc sunt, septenis propter quæ psallimus horis \par \emph{Matutina }ligat Christum, qui crimina purgat. \par \emph{Prima} replet sputis. Dat causam \emph{tertia} mortis. \par \emph{Sexta} cruci nectit. Latus ejus \emph{nona} bipertit. \par \emph{Vespera} deponit. Tumulo \emph{completa} [\emph{completorium}] reponit.'} The Breviary is the growth of many ages. In the early Church great liberty and diversity prevailed in the forms of devotion, but the Popes Leo I., Gelasius I., Gregory I., Gregory VII., Nicholas  III., and others, labored to unify the priestly devotions, and this work was completed after the Council of Trent.

Besides the Missale Romanum and the Breviarium Romanum, there is a Rituale Romanum, or Book of Priests' Rites; an Episcopale Romanum, containing the Episcopal ceremonies, and a Pontificale Romanum, or the Pontifical. They contain the offices for sacramental and other sacred acts and ceremonies, such as baptism, confirmation, ordination, matrimony, dedication of churches, altars, bells, etc., benediction of crosses, sacred vestures, cemeteries, etc.

\xsection{The Old Catholics.}

§ 36. The Old Catholics.

Literature.

I. By Old Catholic Authors.

The writings of Döllinger, Reinkens, von Schulte, Friedrich, Huber, Reusch, Langen, Michelis, Hyacinthe Loyson, Michaud, bearing on the Vatican Council and the Old Catholic movement since 1870. See Literature in §§ 31 and 34.

The Reports of the Old Catholic Congresses, held at Munich, September, 1871; at Cologne, September, 1872; at Constance, September, 1873; at Freiburg, 1874. Published at Munich, Cologne, Leipzig, and Bonn.

Joseph Hubert Reinkens:  \emph{Katholischer Bischof, den im alten Kathol. Glauben verharrenden Priestern und Laien des deutschen Reiches.} Dated August 11, 1873 (the day of his consecration).

The Letter of the Old Catholic Congress of Constance (signed by Bishop Reinkens, President von Schulte, and the Vice-Presidents Cornelius and Keller) to the General Conference of the Evangelical Alliance, held at New York, October, 1873. In the Proceedings of the Conference, New York, 1874.

P. H. Reusch:  \emph{Bericht über die am} 14, 15, \emph{und} 16 \emph{Sept. }1874, \emph{zu Bonn gehaltenen Unions-Conferenzen, im Auftrag Dr. v. Döllinger herausgegeben}, Bonn, 1875 (75 pp.).

Deutscher Merkur,  \emph{Organ für die Katholische Reformbewegung}, ed. by Hirschwälder, \emph{Weltpriester}. The popular aud official weekly organ since 1871.

Theologisches Literaturblatt, ed. by Prof. Reusch, Bonn. The literary organ of the Old Catholics (10th year, 1875).

II. By Protestant Authors.

Friedberg:  \emph{Sammlung der Actenstücke zum ersten Vatic. Concil. }Tübingen, 1872, pp. 53–63, 625–731, 775–898.

Frommann:  \emph{Geschichte und Kritik des Vatic. Concils.} Gotha, 1872, pp. 250–272.

J. Williamson Nevin (of Lancaster, Pa.): \emph{The Old Catholic Movement}, in the 'Mercersburg Review' for April, 1873, pp. 240–294.

\emph{The Alt-Catholic Movement} (anonymous), in the (Amer. Episc.) 'Church Review,' New York, July, 1873.

W. Krafft (Professor of Church History in Bonn): \emph{The Vatican Council and the Old Catholic Movement}, read before, and published in the Proceedings of, the General Conference of the Evangelical Alliance in New York, October, 1873.

Cæsar Pronier (late Professor of Theology in the Free Church Seminary at Geneva, perished in the shipwreck of the Ville du Havre, Nov. 22, 1873, on his return from the General Conference of the Evangelical Alliance): \emph{Roman Catholicism in Switzerland since the Proclamation of the Syllabus}, 1873 (in the Proceedings of the Alliance Conference, New York, 1874).

III. By Roman Catholics.

Besides many controversial writings since the year 1870 (quoted in part in §§ 31 and 34, and articles in Roman Catholic reviews (as the \emph{Dublin Review}, the \emph{Civiltà Cattolica}, the \emph{Catholic World}) and newspapers (as the Paris \emph{L’Univers}, the London \emph{Tablet}, the Berlin \emph{Germania}, etc.), see especially the Papal Encyclical of Nov. 21, 1873, in condemnation of the 'new heretics,' miscalled 'Old Catholics.'

The Old Catholic movement—the most important in the Latin Church since the Reformation, with the exception, perhaps, of Jansenism—began during the Vatican Council, and was organized into a distinct Church three years afterwards (1873), at Constance, in the very hall where, three hundred and sixty years before, an œcumenical Council was held which, by deposing two rival Popes and electing another, asserted its superiority over the Papacy, but which, by burning John Huss for teaching evangelical doctrines, defeated its own professed object of a 'Reformation of the Church in the head and the members.' This strange coincidence of history brings to mind Luther's poem on the Belgian martyrs:

\emph{'Die Asche will nicht lassen ab},

\emph{Sie stäubt in allen Landen};

\emph{Hier hilft kein Loch, noch Grab, noch Grab},

\emph{Sie macht den Feind zu Schanden.}'

The God of history has his \emph{horas et moras}, but he always carries out his designs at last. The Old Catholic secession would have assumed far more formidable proportions, and cut off from the dominion of the Pope the most intelligent and influential dioceses, if the eighty-eight Bishops who in the Vatican Council voted against Papal Infallibility, had carried out their conviction, instead of making their submission for the sake of a hollow peace. But next to the Pope, Bishops, from an instinctive fear of losing power, have always been most hostile to any serious reform. The old story of the Jewish hierarchy, in dealing with Christ and the Apostles, is repeated again and again in the history of the Church, though also with the honorable exceptions of a Nicodemus and Gamaliel.

Œcumenical Councils are very apt to give rise to secessions. A conscientious minority will not yield, in matters of faith, to a mere majority vote. Thus the Council of Nicæa (325) was only the signal for a new and more serious war between orthodoxy and the Arian heresy, and, even after the triumph of the former at Constantinople (381), the latter lingered for centuries among the newly converted German races. The Council of Ephesus (431) gave rise to the Nestorian schism, and the Council of Chalcedon (451) to the several Monophysite sects, which continue in the East to this day with almost as much tenacity of life as the orthodox Greek Church. From the sixth œcumenical Council (680) dates the Monothelite schism. The Council of Florence (1439) failed to effect a union between the Latin and the Greek communions. The Council of Trent (1563), instead of healing the split caused by the Reformation, only deepened and perpetuated it by consolidating Romanism and anathematizing evangelical doctrines. The nearest parallel to the case in hand is the schism of the Bishops and clergy of Utrecht, which originated in a protest against the implied Papal Infallibility of the anti-Jansenist bull \emph{Unigenitus, }and which recently made common cause with the Old Catholics of Germany by giving them the Episcopal succession.

The Old Catholic Church in Germany and Switzerland arose from a protest, in the name of conscience, reason, and honest learning, against the Papal absolutism and infallibilism of the Vatican Council, and against the obsolete mediævalism of the Papal Syllabus. It lifts its voice against unscrupulous Jesuitical falsifications of history, and against that spiritual despotism which requires, as the highest act of piety, the slaughter of the intellect and will, and thereby destroys the sense of personal responsibility. It has in its favor all the traditions of Gallicanism and liberal Catholicism, which place an œcumenical Council or the whole representative Church above the Pope, the testimony of the ancient Græco-Latin Church, which knew nothing of Papal Infallibility, and even condemned some Popes as heretics, and the current of history, which can not be turned backward.

The leaders of the new Church are eminent for learning, ability, moral character, and position, and were esteemed, before the Vatican Council, pillars and ornaments of the Roman Church—viz., Döllinger,\footnote{Dr. John Jos. Ignat. von Döllinger, of Munich (born 1799), the Nestor of Old Catholicism, is the author of an unfinished \emph{Church History} (Lehrbuch der Kirchengeschichte, Regensburg, second edition, 1843, to Leo X.), a polemic work against the \emph{Reformation} (Die Reformation, ihre innere Entwickelung und ihre Wirkungen, 1846–48, 3 vols.), a Sketch of Luther (1851), Judaism and Heathenism in Relation to Christianity (1857), The Church and the Churches (1860), Fables of Popes and Prophecies of the Middle Ages (1863; English translation, with a Preface by Prof. Henry B. Smith, New York, 1872), and a number of essays and pamphlets. He also edited the miscellaneous writings of Möhler, after whose death he was regarded as the foremost Roman Catholic Church historian. Since his excommunication he delivered, in the great hall of the Museum at Munich, seven interesting lectures \emph{On the Reunion of the Churches} (English translation, with Preface by H. N. Oxenham, of Oxford; republished, New York, 1872). He was Rector of the University of Munich during its Jubilee year, 1871–72, and at the celebration of the Jubilee, in July, 1872, he acquitted himself with marked ability and scholarly dignity, and received from the University, the King of Bavaria, and foreign scholars, the highest honors.  \par{} [Döllinger, d. 1890, unreconciled to the papal government. For his later judgment on Luther and the Prot. Reformation, see his \emph{Akad. Vortr.} I., 76, and Schaff: \emph{Our Fathers' Faith and Ours}, pp. 108, 635. Works not given above: \emph{Beiträge zur Sektengesch. des M. A.}, 1890, \emph{Akad. Vorträge}, 3 vols., pp. 188–91, and, in connection with Prof. Reusch, \emph{Selbstbiographie des Kard. Bellarmin}, 1887.—Ed.]} Reinkens,\footnote{\par Formerly Catholic Professor of Church History in the University of Breslau, now Bishop of the Old Catholic Church in Germany. He resides at Bonn, and is a gentleman of great popular eloquence and winning manners.} Friedrich,\footnote{Professor of Church History in Munich, editor of the \emph{Documenta ad illustrandum Conc. Vaticanum} (2 vols.), and of the \emph{Diary} (\emph{Tagebuch während des Vatic. Concils}), which gives an inside view of the Council from his intimate connection with members.} Huber,\footnote{Professor of Philosophy at Munich, and author of works on \emph{the Philosophy of the Fathers}, on \emph{Jesuitism}, and against the last book of Strauss on \emph{The Old and New Faith.}} Michelis,\footnote{Formerly professor at Braunsberg, and once Catholic member of the Prussian Chamber of Deputies, now pastor of the Old Catholic congregation at Zurich, an elderly gentleman of much learning and eloquence.} Reusch,\footnote{Professor of Theology in Bonn, editor of the literary organ of the Old Catholics, and Acting Secretary of Bishop Reinkens.} Langen,\footnote{Likewise Professor of Theology in Bonn, and author of a learned work on the Vatican decrees examined in the light of Catholic tradition (1873).} von Schulte,\footnote{The first canonist of Europe, the lay leader of Old Catholicism, and able president of its Congresses, formerly Professor of Canon Law in Prague, now in the University of Bonn. Before the Council he received many letters and tokens of respect from Pope Pius IX.} and ex-Père Hyacinthe Loyson.\footnote{Born at Orleans, 1827, priest and monk of the order of the Carmelites, formerly esteemed the most eloquent preacher in France. He broke with his order and with Rome in 1869, and is now settled at Geneva as pastor of an Old Catholic congregation. His marriage to an American widow (1872) created almost as much sensation as Luther's marriage to a nun. He has recently withdrawn from state control, and established an independent Church (1874).}

The centres of Old Catholicism are Munich and Bonn in Germany, and Geneva and Soleure (also Olten) in Switzerland. Beyond these two countries it has many isolated sympathies, but no organized form, and no hold upon the people.\footnote{The German origin of the movement operates against it in France, which, with all its Gallican traditions, has, for political reasons, since the war of 1870, become more Romish than it ever was before. When Völk, at the Old Catholic Congress in Constance, alluded to the uprising of the \emph{Deutschthum} versus the \emph{Welschthum}, and the intrigues of French Jesuits, Hyacinthe and Pressensé left the hall. Yet the Old Catholic priests, who were elected pastors of Geneva by the Catholic part of the population in October, 1873—Loyson, Hurtault, and Charard—are all Frenchmen. Once more Geneva seems to become the centre and starting-point of a new reformation, which sooner or later will react upon France. Abbé Michaud, formerly of the Madeleine in Paris, so far is the only prominent Old Catholic in France. Among the Irish Catholics there is not the least indication of sympathy with Old Catholicism, not even in free America. Spain and Italy ought to sympathize with it, for the Pope is the implacable enemy of Italian unity and the Spanish republic; but they have kept aloof so far from any progressive religious movement; and Spain has once more surrendered herself to the rule of a Bourbon and the Pope (1875). In England, the famous pamphlet of Gladstone on the \emph{Vatican Decrees} (1874) has brought to light the Old Catholic sympathies of Lord Acton and other prominent English Catholics.} In September, 1873, the Old Catholics in the German Empire numbered about one hundred congregations (mostly in Prussia, Baden, and Bavaria), forty priests, and fifty thousand professed members. Since their more complete organization they will probably make more rapid progress. Heretofore the movement in Germany has been more scholastic than popular. It has enlisted the sympathies of the educated, but not to an equal extent the enthusiasm of the people. The question of Papal Infallibility has no such direct practical bearing as the question of personal salvation and peace of conscience, which made the Reformation spread with such irresistible power over all Western Christendom. The masses of Roman Catholics are either too ignorant or too indifferent to care much whether another dogma is added to the large number already adopted, and have no more difficulty to believe blindly in Papal Infallibility than in the daily miracle of transubstantiation and the sacrifice of the mass.\footnote{When in Cologne, July, 1873, I asked a domestic of one of the first hotels where the \emph{Old Catholics} worshiped. He promptly replied, 'You mean the \emph{New Protestants.} I have nothing to do with sects; I am a true Catholic, and mean to die one.' This seemed to me characteristic of the popular feeling in Cologne. The Dome was well filled with worshipers all Sunday, while the Old Catholics had a small though intelligent and respectable congregation in the Garrison Church, and in the small chapel at the City Hall. Dr. Tangermann read Latin mass like a Romish priest, but preached an evangelical sermon in German which would do credit to any Protestant pastor.} On the other hand, however, the Old Catholics are powerfully aided by the widespread indignation against priestcraft, and the serious conflict of the German Empire and the Swiss Republic with the Papacy, which was provoked by the Papal Syllabus and the Vatican Council, and may lead to a thorough revision of the ecclesiastical status of the Continent. Their ultimate success as a Church must chiefly depend upon the continued ascendency of the positive Christian element over the negative and radical (which raised and ruined the 'German Catholic' or Ronge movement of 1844); for only the enthusiasm of faith has constructive power, and that spirit of sacrifice and endurance which is necessary for the establishment of permanent institutions.

The Old Catholic movement was foreshadowed in the liberal Catholic literature preceding the Vatican Council, especially \emph{Janus}; it gathered strength during the Council; it uttered itself in a united protest against the decrees of the Council at a meeting of distinguished Catholic scholars at Nuremberg in August, 1870; and it came to an open rupture with Rome by the excommunication of Döllinger and his sympathizers.

Being called upon by the Archbishop of Munich (his former pupil, and at first an anti-Infallibilist) to submit to the new dogma of Papal absolutism and Infallibility, Dr. Döllinger, in an open answer dated Munich, March 28,1871, declared that, as a Christian, as a theologian, as a historian, and as a citizen, he could not accept the Vatican decrees, for the reasons that they are inconsistent with the spirit of the Gospel and the clear teaching of Christ and the Apostles; that they contradict the whole genuine tradition of the Church; that the attempt to carry out the Papal absolutism had been in times past the cause of endless bloodshed, confusion, and corruption; and that a similar attempt now must lead to an irreconcilable conflict of the Church with the State, and of the clergy with the laity.\footnote{The following is the memorable protest of this aged divine, which reminds one of Luther's more bold and defiant refusal at Worms to recant his writings unless convicted of error from Scripture and reason: 'Als Christ, als Theologe, als Geschichtskundiger, als Bürger kann ich diese Lehre nicht annehmen. \emph{Nicht als}  Christ:  \emph{denn sie ist unverträglich mit dem Geiste des Evangeliums and mit den klaren Aussprüchen Christi und der Apostel; sie will gerade das Imperium dieser Welt aufrichten, welches Christus ablehnte, will die Herrschaft über die Gemeinden, welche Petrus allen und sich selbst verbot. Nicht als}  Theologe:  \emph{denn die gesammte echte Tradition der Kirche steht ihr unversöhnlich entgegen. Nicht als}  Geschichtskenner  \emph{kann ich sie annehmen, denn als solcher weiss ich, dass das beharrliche Streben, diese Theorie der Weltherrschaft zu verwirklichen, Europa Ströme van Blut gekostet, ganze Länder vewirrt und heruntergebracht, den schönen organischen Verfassungsbau der älteren Kirche zerrüttet und die ärgsten Missbräuche in der Kirche erzeugt, genährt und festgehalten hat. Als}  Bürger  \emph{endlich muss ich sie von mir weisen, well sie mit ihren Ansprüchen auf Unterwerfung der Staaten und Monarchen und der ganzen politischen Ordnung unter die päpstliche Gewalt und durch die eximirte Stellung, welche sie für den Klerus fordert, den Grund legt zu endloser verderblicher Zwietracht zwischen Staat und Kirche, zwischen Geistlichen und Laien. Denn das kann ich mir nicht verbergen, dass diese Lehre, an deren Folgen das alte deutsche Reich zu Grunde gegangen ist, falls sie bei dem katholischen Theil der deutschen Nation herrschend würde, sofort auch den Keim eines unheilbaren Siechthums in das eben erbaute neue 'Reich verpflanzen würde.}'—J. von Döllinger's \emph{Erklärung an den Erzbishof von München-Freising}, München, 1871, p. 17 sq.} Whereupon Döllinger was excommunicated April 17, 1871, as being guilty of 'the crime of open and formal heresy.'\footnote{'\emph{Crimen hæreseos externæ et formalis.}'}

His colleague, Professor Friedrich, incurred the same fate. Other Bishops, forgetting their recent change of conviction, proceeded with the same rigor against refractory priests. Cardinal Rauscher suspended the Lent preacher Pederzani; Cardinal Schwarzenberg, Professor Pelleter (who afterwards became a Protestant); Bishop Förster (whose offer to resign was refused by the Pope) suspended Professors Reinkens, Baltzer, and Weber, of Breslau; the Bishop of Ermeland, Professors Michelis and Menzel, and Dr. Wollmann, in Braunsberg; the Archbishop of Cologne deposed the priest Dr. W. Tangermann, of Cologne, and suspended Professors Hilgers, Reusch, Langen, and Knoodt, of Bonn, who, however, supported by the Prussian Government, retained their official positions in the University.

In spite of these summary proceedings of the Bishops, the Old Catholic party, aided by the sympathies of the educated classes, made steady progress, organizing congregations, holding annual meetings, and enlisting the secular and religions press. With great prudence the leaders avoided or postponed reforms, till they could be inaugurated and sanctioned by properly constituted authorities, and moved cautiously between a timid conservatism and a radical liberalism; thus retaining a hold on both wings of the nominal Catholic population.

In the year 1873 the Old Catholics effected a regular Church organization, and secured a legal status in the German Empire, with the prospect of support from the national treasury. Professor Joseph Hubert Reinkens was elected Bishop by the clergy and the representatives of the laity, and was consecrated at Rotterdam by the Old Catholic Bishop Heykamp, of Deventer (Aug. 11, 1873).\footnote{In his Pastoral Letter, Bishop Reinkens disclaims all hierarchical ambition, vain show, and display, and promises to exercise his office in the spirit of apostolic simplicity as a pastor of the flock. He lays great stress on the primitive Catholic mode of his election by the \emph{clergy} and the \emph{people}, as contrasted with the modern election by the \emph{Pope}. He claims to stand in the rank of Cyprian, Hilary, Ambrose, Augustine, and those thousands of Bishops who never were elected by the Pope, or were even known to the Pope, and yet are recognized as truly Catholic Bishops. Consecration by one Bishop is canonically valid, though two or more \emph{assistant} Bishops are usually present. The late Archbishop Loos of Utrecht would have performed the act, had he not died a few months before. Rome, of course, considers this election and consecration by excommunicated priests as a mere farce and a damnable rebellion. See the Pope's Encyclical of Nov. 21, 1872, quoted below.} He was recognized in his new dignity by the King of Prussia, and took the customary oath of allegiance at Berlin (Oct. 7). Other governments of Germany followed this example. (The Empire as such has nothing to do with the Church.) To complete the organization, the Congress at Constance adopted a synodical and parochial constitution, which makes full provision for an equal share of the laity with the clergy in the government of the Church; the synodical representation (\emph{Synodal-Repräsentanz}), or executive committee, being composed of five laymen and five clergymen, including the Bishop.\footnote{See the \emph{Entwurf einer Synodal- und Gemeinde-Ordnung}, Sect. III. §§ 13 and 14: '\emph{In der Leitung des altkatholischen Gemeinwesens steht dem Bischof eine von der Synode gewählte Synodal-Repräsentanz zur Seite. Die Synodal-Repräsentanz besteht aus vier Geistlichen und fünf Laien.}'} This implies the Protestant principle of the general priesthood of believers, and will prevent hierarchical abuses. Certain changes in the cultus, such as the simplification of the mass as a memorial service of the atoning sacrifice of Christ, the substitution of the vernacular language for the Latin, the restoring of the cup to the laity, the introduction of more preaching, and the abolition of various abuses (including the forced celibacy of the clergy), will inevitably follow sooner or later.

The \emph{doctrinal} status of the Old Catholic denomination was at first simply \emph{Tridentine} Romanism versus \emph{Vatican} Romanism, or the Creed of Pius IV. against the Creed of Pius IX.\footnote{Their original programme, adopted at the first Congress at Munich, September 21, 1871, probably drawn up by Döllinger, was very conservative, and included the following articles:  \par 1. We hold fast to the Catholic faith as certified by Scriptures and tradition, and also to the Old Catholic worship. We reject from this stand-point the new dogmas enacted under the pontificate of Pius IX., especially that regarding the infallibility and supreme ordinary and immediate jurisdiction of the Pope. \par 2. We hold fast to the old constitution of the Church, and reject every attempt to deprive the Bishops of their diocesan independence. We acknowledge the primacy of the Bishop of Rome, on the ground of the Fathers and Councils of the undivided Church of antiquity; but we deny the right of the Pope to define any article of faith, except in agreement with the holy Scriptures and the ancient and unanimous tradition of the Church. \par 3. We aim at a reformation of various abuses of the Church, and a restoration of the rights of the laity in ecclesiastical affairs. \par 4. We hope for a reunion with the Greek and Orthodox Russian Church, and for an ultimate fraternal understanding with the other Christian confessions, especially the Episcopal churches of England and America.} This is the ground taken by the Old Catholics in Holland, and adhered to by them to this day. But the logic of the protest against modern Popery will hardly allow the Old Catholics of Germany and Switzerland long to remain in this position. Their friendly attitude towards Protestants, as officially shown in their letter to the General Conference of the Evangelical Alliance, is inconsistent with the Tridentine anathemas. Tridentine Romanism, moreover, is as much an innovation on œcumenical Catholicism as the Vatican Romanism is an innovation on that of Trent, and both are innovations in the same line of consolidation of the one-sided principle of authority. There is no stopping at half-way stations. We must go back to the fountain-head, the Word of God, which is the only final and infallible authority in matters of faith, and furnishes the best corrective against all ecclesiastical abuses.

The leaders of the Old Catholic Church are evidently on this road. They still adhere to Scripture \emph{and tradition}, as the joint rule of faith: but they confine tradition to the unanimous consent of the ancient undivided Church, consequently to the œcumenical creeds, which are held in common by Greeks, Latins, and orthodox Protestants. They have been forced to give up their belief in the infallibility of an œcumenical Council, since the Vatican Council, which is as œcumenical (from the Roman point of view) as that of Trent, has sanctioned what they regard as fatal error. Moreover, Bishop Reinkens, in an eloquent speech before the Old Catholic Congress at Constance, disowned all Romish prohibitions of Bible reading, and earnestly encouraged the laity to read the Book of Life, that they may get into direct and intimate communion with God.\footnote{I give a few extracts from this address, which was delivered in the famous Council Hall of Constance, and received with great applause by the crowded assembly: 'The holy Scripture is the reflection of the sun of righteousness which appeared in Jesus Christ our Lord. I say, therefore, Read the holy Scriptures. I say more: \emph{For the Old Catholics who intrust themselves to my episcopal direction, there exists no prohibition of the reading of the Bible. . . . }Let nothing hinder you from approaching the Gospel, that you may hear the voice of the Bridegroom (John iii. 29). Listen to his voice, and remember that, as the flower turns to the light, and never unfolds all its splendor and beauty except by constantly turning to the light of the sun, thus also the Christian's soul can not represent the full beauty and glory of its divine likeness except by constantly turning to this Gospel, in the rays of which its own fire is kindled. . . . Do not read the Scriptures from curiosity, to find things which are not to be revealed in this world; nor presumptuously, to brood over things which can not be explained by men; nor for the sake of controversy, to refute others; but read the Scriptures to enter into the most intimate communion with God, so that you may be able to say, Nothing shall separate me from the love of Christ. . . . It is not sufficient to have the Bible in every house, and to read it at certain hours in a formal and fragmentary manner, but it ought to be the light of the soul, to which it turns again and again. I repeat it once more: For the Old Catholics, no injunction exists against reading the Bible. On the contrary, I admonish you most earnestly: Read again and again in this holy book, sitting down in humility and joy at the feet of the Lord, \emph{for He alone has words of eternal life.}'} This communion with God through Christ as the only Mediator, and through his Word as the only rule of faith, is the very soul of evangelical Protestantism. The Scripture principle, consistently carried out, must gradually rule out the unscriptural doctrines and usages sanctioned by the Council of Trent.

But it is not necessary on this account that the Old Catholics should ever become Protestants in the historical sense of the term. They may retain those elements of the Catholic system which are not inconsistent with the spirit of the Scriptures, though they may not be expressly sanctioned by the letter. They may occupy a peculiar position of mediation, and in this way contribute their share towards preparing the way for an ultimate reunion of Christendom. And this is their noble aim and desire, openly expressed in a fraternal letter to an assembly of evangelical Christians from nearly all Protestant denominations. They declare: 'We hope and strive for the restoration of the unity of the Christian Church. We frankly acknowledge that no branch of it has exclusively the truth. We hold fast to the ultimate view that upon the foundation of the Gospel, and the doctrines of the Church grounded upon it, and upon the foundation of the ancient, undivided Church, a union of all Christian confessions will be possible through a really œcumenical Council. This is our object and intention in the movement which has led us into close relations with the Evangelical, the Anglican, the Anglo-American, the Russian, and the Greek churches. We know that this goal can not easily be reached, but we see the primary evidences of success in the circumstance that a truly Christian intercourse has already taken place between ourselves and other Christian churches. Therefore we seize with joy the hand of fellowship you have extended to us, and beg you to enter into a more intimate fellowship with us in such a way as may be agreed upon by both parties.'\footnote{Letter of the Congress of Constance, September, 1873, to the General Conference of the Evangelical Alliance in New York. Comp. also Döllinger's \emph{Lectures on the Reunion of the Churches}, and Hyacinthe Loyson's letter to the General Conference in New York.}

On the other hand, the Old Catholics have extended the hand of fellowship to the Greeks and Anglo-Catholics, and adopted, at a Union Conference held in Bonn, Sept., 1874, an agreement of fourteen theses, as a doctrinal basis of intercommunion between those Churches which recognize, besides the holy Scriptures, the binding authority of the tradition of the undivided Church of the first six centuries. In a second Conference, in 1875, they surrendered the doctrine of the double procession of the Spirit as a peace-offering to the Orientals.\footnote{See the documents of the two Bonn Conferences, at the close of Vol. II.}

In the mean time the Pope has cut off all prospect of reconciliation. In his Encyclical of November 21, 1873, addressed to all the dignitaries of the Roman Church, Pius IX., after unsparingly denouncing the governments of Italy, Switzerland, and Germany, for their cruel persecution of the Church, speaks at length of 'those new heretics, who, by a truly ridiculous abuse of the name, call themselves Old Catholics,' and launches at their 'pseudo-bishop' and all his abettors and helpers the sentence of excommunication, as follows:

'The attempts and the aims of these unhappy sons of perdition appear plainly, both from other writings of theirs and most of all from that impious and most impudent of documents which has lately been published by him whom they have set up for themselves as their so-called bishop. For they deny and pervert the true authority of jurisdiction which is in the Roman Pontiff and the Bishops, the successors of the Blessed Peter and the Apostles, and transfer it to the populace, or, as they say, to the community; they stubbornly reject and assail the infallible teaching authority of the Roman Pontiff and of the whole Church; and, contrary to the Holy Spirit, who has been promised by Christ to abide in his Church forever, they audaciously affirm that the Roman Pontiff and the whole of the Bishops, priests, and people who are united with him in one faith and communion, have fallen into heresy by sanctioning and professing the definitions of the œcumenical Vatican Council. Therefore they deny even the indefectibility of the Church, blasphemously saying that it has perished throughout the world, and that its visible head and its Bishops have fallen away; and that for this reason it has been necessary for them to restore the lawful Episcopate in their pseudo-bishop, a man who, entering not by the gate, but coming up by another way, has drawn upon his head the condemnation of Christ.

'Nevertheless, those unhappy men who would undermine the foundations of the Catholic religion, and destroy its character and endowments, who have invented such shameful and manifold errors, or, rather, have collected them together from the old store of heretics, are not ashamed to call themselves Catholics, and Old Catholics; while by their doctrine, their novelty, and their fewness they give up all mark of antiquity and of catholicity. . . .

'But these men, going on more boldly in the way of iniquity and perdition, as by a just judgment of God it happens to heretical sects, have wished also to form to themselves a hierarchy, as we have said, and have chosen and set up for themselves as their pseudo-bishop a certain notorious apostate from the Catholic faith, Joseph Hubert Reinkens; and, that nothing might be wanting to their impudence, for his consecration they have had recourse to those Jansenists of Utrecht whom they themselves, before their falling away from the Church, regarded with other Catholics as heretics and schismatics. Nevertheless this Joseph Hubert Reinkens dares to call himself a bishop, and, incredible as it may seem, the most serene Emperor of Germany has by public decree named and acknowledged him as a Catholic bishop, and exhibited him to all his subjects as one who is to be regarded as a lawful bishop, and as such to be obeyed. But the very rudiments of Catholic teaching declare that no one can be held to be a lawful bishop who is not joined in communion of faith and charity to the rock on which the One Church of Christ is built; who does not adhere to the supreme pastor to whom all the sheep of Christ are committed to be fed; who is not united to the confirmer of the brotherhood which is in the world.' [This cuts off all Greek Bishops as well. Then follow the usual patristic texts for the pretensions of Rome.]

'We therefore, who have been placed, undeserving as we are, in the Supreme See of Peter for the guardianship of the Catholic faith, and for the maintenance of the unity of the universal Church, according to the custom and, example of our predecessors and their holy decrees, by the power given us from on high, not only declare the election of the said Joseph Hubert Reinkens to be contrary to the holy canons, unlawful, and altogether null and void, and denounce and condemn his consecration as sacrilegious; but by the authority of Almighty God we declare the said Joseph Hubert—together with those who have taken part in his election and sacrilegious consecration, and whoever adhere to and follow the same, giving aid, favor, or consents—excommunicated under anathema, separated from the communion of the Church, and to be reckoned among those whose fellowship has been forbidden to the faithful by the Apostle, so that they are not so much as to say to them, God speed you!'

As the Pope's letter of complaint to the Emperor of Germany (August, 1873), in which he claims jurisdiction, in some sense, over all baptized Christians, called forth a courteous and pointed reply from the Emperor disclaiming all intention of persecuting the Catholic Church while defending the rights of the civil government against the encroachments of the hierarchy, and informing his Infallibility that Protestants recognize no other mediator between God and themselves than the Lord Jesus Christ; so this Encyclical was met by an able, dignified, and manly Pastoral from Bishop Reinkens, dated Bonn, December 14, 1873, in which, after refuting the accusations of the Pope, he closes with the following words: 'Brethren in the Lord, what shall we do when Pius IX. exhausts the language of reproach and calumny, and calls us even the most miserable sons of perdition (\emph{miserrimi isti perditionis filii}), to embitter the uninquiring multitude against us? If we are true disciples of Jesus—as we trust—we have that peace which the Lord gives, and not the world, and our ``heart will not be troubled, neither be afraid'' (John xiv. 27). O how sweetly sounds the exhortation: ``Bless them which persecute you: bless, and curse not;'' ``Recompense to no man evil for evil;'' ``If it be possible, as much as lieth in you, live peaceably with all men'' (Rom. xii. 14, 17, 18); ``Love your enemies, bless them that curse you, do good to them that hate you, and pray for them which despitefully use you, and persecute you; that ye may be the children of your Father which is in heaven: for he maketh his sun to rise on the evil and on the good, and sendeth rain on the just and on the unjust'' (Matt. v. 44, 45). Let us look up to Christ, our example, ``who, when he was reviled, reviled not again'' (1 Pet. ii. 21–23). ``The peace of God, which passeth all understanding, keep your hearts and minds through Jesus Christ.'''

The Swiss Federal Government, in answer to the charges raised against it in the same Encyclical, has broken off all diplomatic intercourse with the Papal court. In a new Encyclical of March 23, 1875, addressed to the Bishops of Switzerland, Pious IX. confirmed the condemnation of Nov. 21, 1873, and hurled it with increased severity against the Old Catholics of that country, 'who attack the very foundations of the Catholic religion, boldly reject the dogmatic definitions of the Council of the Vatican, and by every means labor for the ruin of souls.' He calls upon the faithful to 'avoid their religious ceremonies, their instructions, their chairs of doctrinal pestilence, which they have the audacity to set up for the purpose of betraying the sacred doctrines, their writings, and contact with them. Let them have no part, no relation of any kind, with those intruding priests and the apostates who dare exercise the functions of the ecclesiastical ministry, and who have absolutely no jurisdiction and no legitimate mission at all. Let them hold them in horror as strangers and thieves, who come only to steal, assassinate, and destroy.'

The Old Catholic movement in Switzerland is more radical and political than the German, and bears a similar relation to it as the Zwinglian Reformation does to the Lutheran. Edward Herzog, an able and worthy priest of Olten, was elected first bishop by the Swiss Synod, and consecrated by Bishop Reinkens at Rheinfelden, Sept. 18, 1876.

\xchapter{Chapter 5. The Creeds of the Evangelical Churches.}

General Literature.

There are no complete collections of Protestant Creeds, but several separate collections of the Lutheran and of the Reformed Creeds, which will be noticed below under the proper sections. The \emph{Corpus et Syntagma, Confessionum fidei,} Genev. 1654, is chiefly Calvinistic, and the Oxford \emph{Sylloge Confessionum sub tempus reformandæ ecclesiæ editarum,} 1827 (pp. 454), contains only six confessions (including the \emph{Prof. Fidei Trid.} and the \emph{Confessio Saxonica}).

On the general history and principles of the Reformation, the reader is referred to the works, correspondence, and numerous biographies of the Reformers (\emph{e.g. }the \emph{Corpus Reformatorum,} ed. Bretschneider and Bindseil: Luther's \emph{Letters,} by De Wette, supplemented by Seidemann: Calvin's \emph{Works,} new edition by Baum, Cunitz, and Reuss; his \emph{Letters,} by Bonnet; Herminjard's \emph{Correspondance des Réformateurs dans les pays de langue française}; Strype's \emph{Ecclesiastical Memorials,} etc.; the publications of the Parker Society); and the historical works of Sleidan, Seckendorf, Salig, De Thou, Hottinger, Hess, Marheineke, Ranke, Merle D'aubigné, Hagenbach (fourth edition, 1870), Geo. P. Fisher; also Schaff (\emph{Principle of Protestantism,} 1845), Dorner (\emph{Geschichte der Protest. Theologie,} 1867, pp. 77–329, Engl. transl. Edinb. 1871, 2 vols.), Kahnis (\emph{Die Deutsche Reformation,} Leipz. 1872). See lists of literature in Gieseler, \emph{Church History,} Vol. IV. pp. 9 sqq. (Anglo-Amer. edition), and Geo. P. Fisher (of Yale College), \emph{The Reformation,} New York, 1873, Appendix II. pp. 567–591.

\xsection{The Reformation. Protestantism and Romanism.}

§ 37. The Reformation. Protestantism and Romanism.

Protestant Christendom has a nominal membership of about one hundred millions, chiefly in the northern and western parts of Europe and America, and among the most vigorous and hopeful nations of the earth. It represents modern or progressive Christianity, while Romanism is mediæval Christianity in conflict with modern progress, and the Eastern Church ancient Christianity in repose.

We must first of all distinguish between \emph{evangelical }or \emph{orthodox }Protestantism, which agrees with the Greek and Roman Church in accepting the holy Scriptures and the œcumenical faith in the Trinity and Incarnation, and \emph{heretical }or \emph{radical }Protestantism, which dissents from the œcumenical \emph{consensus,} and makes a new departure either in a mystical or in a rationalistic direction. The former constitutes the great body of nominal Protestantism, and is the subject of this chapter. It includes, in the first line, the Lutheran and the Reformed Confessions, or the various national churches of the Reformation in Europe and their descendants in America; and then, in the second line, all those denominations which have proceeded or seceded from them, mostly on questions of government or minor points of doctrine, without departing from the essential articles of their faith, such as the Moravians, Methodists, Mennonites, Baptists, Quakers, Irvingites, and a number of free churches holding to the voluntary principle.

The various Evangelical Protestant churches, viewed as distinct ecclesiastical organizations and creeds, take their rise directly or indirectly from the sixteenth century; but their principles are rooted and grounded in the New Testament, and have been advocated more or less clearly, in part or in full, by spiritual and liberal minded divines in every age of the Church. The stream of Latin or Western Christianity was divided in the sixteenth century; the main current moving cautiously and majestically in the old mediæval channel, the other boldly cutting several new beds for the overflowing waters, and rushing forward, at first with great rapidity and energy, then slacking its speed, and then resuming its forward march with the tide of emigration in a western direction, whither, in the prophetic language of the great English idealist, 'the course of empire takes its way.'

The Reformation of the sixteenth century is, next to the introduction of Christianity, the greatest event in history. It was no sudden revolution; for what has no roots in the past can have no permanent effect upon the future. It was prepared by the deeper tendencies and aspirations of previous centuries, and, when finally matured, it burst forth almost simultaneously in all parts of Western Christendom. It was not a superficial amendment, not a mere restoration, but a regeneration; not a return to the Augustinian, or Nicene, or ante-Nicene age, but a vast progress beyond any previous age or condition of the Church since the death of St. John. It went, through the intervening ages of ecclesiasticism, back to the fountain-head of Christianity itself, as it came from the lips of the Son of God and his inspired Apostles. It was a deeper plunge into the meaning of the Gospel than even St. Augustine had made. It brought out from this fountain a new phase and type of Christianity, which had never as yet been fully understood and appreciated in the Church at large. It was, in fact, a new proclamation of the free Gospel of St. Paul, as laid down in the Epistles to the Romans and Galatians. It was a grand act of emancipation from the bondage of the mediæval hierarchy, and an assertion of that freedom wherewith Christ has made us free. It inaugurated the era of manhood and the general priesthood of believers. It taught the direct communion of the believing soul with Christ. It removed the obstructions of legalism, sacerdotalism, and ceremonialism, which, like the traditions of the Pharisees of old, had obscured the genuine Gospel and made void the Word of God.\footnote{\par It is significant that Christ uses \textgreek{ παράδοσις, } \emph{tradition, }only in an unfavorable sense, as opposed to the Word of God, viz., Matt. xv. 3, 6; Mark vii. 5, 8, 9, 13. Paul employs the term in a bad sense, Gal. i. 14 and Col. ii. 8: in a good sense, of the doctrines of the Gospel, 1 Cor. xi. 2; 2 Thess. ii. 15; iii. 6.}

We do not depreciate mediaeval Catholicism, the womb of the Reformation, the grandmother of modern civilization. It was an inestimable blessing in its time. When we speak of the 'dark ages,' we should never forget that the Church was the light in that darkness. She was the training-school of the Latin, Celtic, and Teutonic (partly also the Sclavonic) races in their childhood and wild youth. She gave them Christianity in the shape of a new theocracy, with a priesthood, minute laws, rites, and ceremonies. She acted as a bulwark against the despotism of the civil and military power, and she defended the moral interests, the ideal pursuits, and the rights of the people. But the discipline of law creates a desire which it can not satisfy, and points beyond itself, to independence and self-government: the law is a schoolmaster to lead men to the freedom of the Gospel. When the mediæval Church had fulfilled her great mission in Christianizing and civilizing (to a certain degree) the Western and Northern barbarians, the time was fulfilled, and Christianity could now enter upon the era of evangelical faith and freedom.

And this is Protestantism. If it were a mere negation of popery, it would have vanished long since, leaving no wreck behind. It is constructive as well as destructive; it protests from the positive basis of the Gospel. It attacks human authority from respect for divine authority; it sets the Word of God over all the wisdom of men.

The Reformation was eminently practical in its motive and aim. It started from a question of conscience: 'How shall a sinner be justified before God?' And this is only another form of the older and broader question: 'What shall I do to be saved?' The answer given by the Reformers (German, Swiss, French, English, and Scotch), with one accord, from deep spiritual struggle and experience, was: 'By faith in the all-sufficient merits of Christ, as exhibited in the holy Scriptures.' And by faith they understood not a mere intellectual assent to the truth, or a blind submission to the outward authority of the Church, but a free obedience, a motion of the will, a trust of the heart, a personal attachment and unconditional surrender of the whole soul to Christ, as the only Saviour from sin and death. The absolute supremacy and sufficiency of Christ and his Gospel in doctrine and life, in faith and practice, is the animating principle, the beating heart of the Reformation, and the essential unity of Protestantism to this day.

Here lies its vitality and constructive power. From this central point the whole theology and Church life was directly or indirectly affected, and a new impulse given to the history of the world in every direction.

The Reformers were baptized, confirmed, and educated, most of them also ordained, in the Catholic Church, and had at first no intention to leave it, but simply to purify it by the Word of God. They shrank from the idea of schism, and continued, like The Apostles, in the communion of their fathers until they were expelled from it. When the Pope refused to satisfy the reasonable demand for a reformation of abuses, and hurled his anathemas on the reformers, they were driven to the necessity of organizing new churches and setting forth new confessions of faith, but they were careful to maintain and express in them their \emph{consensus} with the old Catholic faith as laid down in the Apostles' Creed.

The doctrinal principle of evangelical Protestantism, as distinct from Romanism, is twofold—objective and subjective.

The \emph{objective} (generally called the \emph{formal}\} principle maintains the absolute sovereignty of the Bible, as the only infallible rule of the Christian faith and life, in opposition to the Roman doctrine of the Bible \emph{and tradition,} as co-ordinate rules of faith. Tradition is not set aside altogether, but is subordinated, and its value made to depend upon the measure of its agreement with the Word of God.

The \emph{subjective} (commonly called the \emph{material}) principle is the doctrine of justification by the free grace of God through a living faith in Christ, as the only and sufficient Saviour, in opposition to the Roman doctrine of (progressive) justification by faith \emph{and good works,} as co-ordinate conditions of justification. Good works are held by Protestants to be necessary, not as means and conditions, but as results and evidences, of justification.

To these two principles may be added, as a third, the \emph{social} principle, which affects chiefly the government and discipline of the Church, namely, the \emph{universal} priesthood of \emph{believers,} in opposition to the \emph{exclusive} priesthood of the \emph{clergy}. Protestantism emancipates the laity from slavish dependence on the teaching and governing priesthood, and gives the people a proper share in all that concerns the interests and welfare of the Church; in accordance with the teaching of St. Peter, who applies the term \emph{clergy} (\textgreek{ κλῆρος, } \emph{heritage,} 1 Pet. v. 3) to the congregation, and calls all Christians 'living stones' in the spiritual house of God, to offer up 'spiritual sacrifices,' 'a chosen generation, a royal priesthood, a holy nation, a peculiar people,' setting forth 'the praises of him who called them out of darkness into his marvelous light' (1 Pet. ii. 5, 9; comp. v. 1–4; Rev. i. 6; v. 10; xx. 6).

It is impossible to reduce the fundamental difference between Protestantism and Romanism to a single formula without doing injustice to the one or the other. We should not forget that there are evangelical elements in Romanism, as there are legalistic and Romanizing tendencies in certain schools of Protestantism. But if we look at the prevailing character and the most prominent aspects of the two systems, we may draw the following contrasts:

Protestantism corresponds to the Gentile type of Apostolic Christianity, as represented by Paul; Romanism, to the Jewish type, as represented by James and Peter, though not in Peter's Epistles (where he prophetically warns against the fruitful germ of the Papacy, viz., hierarchical pride and assumption), but in his earlier stage and official position as the Apostle of circumcision. Paul was called afterwards, somewhat irregularly and outside of the visible succession, as the representative of a new and independent apostolate of the Gentiles. The temporary collision of Paul and Peter at Antioch (Gal. ii. 11) foreshadows and anticipates the subsequent antagonism between Protestantism and Catholicism.

Protestantism is the religion of freedom (Gal. v. 1); Romanism, the religion of authority. The former is mainly subjective, and makes religion a personal concern; the latter is objective, and sinks the individual in the body of the Church. The Protestant believes on the ground of his own experience, the Romanist on the testimony of the Church (comp. John iv. 42).

Protestantism is the religion of evangelism and spiritual simplicity; Romanism, the religion of legalism, asceticism, sacerdotalism, and ceremonialism. The one appeals to the intellect and conscience, the other to the senses and the imagination. The one is internal, the other external, and comes with outward observation.

Protestantism is the Christianity of the Bible; Romanism, the Christianity of tradition. The one directs the people to the fountain-head of divine revelation, the other to the teaching priesthood. The former freely circulates the Bible, as a book for the people; the latter keeps it for the use of the clergy, and overrules it by its traditions.

Protestantism is the religion of immediate communion of the soul with Christ through personal faith; Romanism is the religion of mediate communion through the Church, and obstructs the intercourse of the believer with his Saviour by interposing an army of subordinate mediators and advocates. The Protestant prays directly to Christ; the Romanist usually approaches him only through the intercession of the blessed Virgin and the saints.

Protestantism puts Christ before the Church, and makes Christliness the standard of sound churchliness; Romanism virtually puts the Church before Christ, and makes churchliness the condition and measure of piety.\footnote{\par This is no doubt the meaning of Schleiermacher's famous formula (\emph{Der Christliche Glaube,} Vol. I. § 24): 'Protestantism makes the relation of the individual to the Church dependent on his relation to Christ; Catholicism, \emph{vice versa,} makes the relation of the individual to Christ dependent on his relation to the Church.' His pupil and successor, Dr. Twesten, puts the distinction in this way: 'Catholicism emphasizes the first, Protestantism the second, clause of the passage of Irenæus: ``Where the Church is, there is the Spirit of God; and where the Spirit of God is, there is the Church and all grace.'''}

Protestantism claims to be only one, but the most advanced portion of the Church of Christ; Romanism identifies itself with the whole Catholic Church, and the Church with Christianity itself. The former claims to be the safest, the latter the only way to salvation.

Protestantism is the Church of the Christian people; Romanism is the Church of priests, and separates them by education, celibacy, and even by their dress as widely as possible from the laity.

Protestantism is the Christianity of personal conviction and inward experience; Romanism, the Christianity of outward institutions and sacramental observances, and obedience to authority. The one starts from Paul's, the other from James's doctrine of justification. The one lays the main stress on living faith, as the principle of a holy life; the other on good works, as the evidence of faith and the condition of justification.

Protestantism proceeds from the invisible Church to the visible; Rome, \emph{vice versa,} from the visible to the invisible.\footnote{\par This is the distinction made by Möhler, who thereby inconsistently admits the essential truth of the Protestant distinction between the visible and invisible Church, which Bellarmin denies as an empty abstraction.}

Protestantism is progressive and independent; Romanism, conservative and traditional. The one is centrifugal, the other centripetal. The one is exposed to the danger of radicalism and endless division; the other to the opposite danger of stagnation and mechanical and tyrannical uniformity.

The exclusiveness and anti-Christian pretensions of the Papacy, especially since it claims infallibility for its visible head, make it impossible for any Church to live with it on terms of equality and sincere friendship. And yet we should never forget the difference between Popery and Catholicism, nor between the system and its followers. It becomes Protestantism, as the higher form of Christianity, to be liberal and tolerant even towards intolerant Romanism.

\xsection{The Evangelical Confessions of Faith.}

§ 38. The Evangelical Confessions of Faith.

The Evangelical Confessions of faith date mostly from the sixteenth century (1530 to 1577), the productive period of Protestantism, and are nearly contemporaneous with the Tridentine standards of the Church of Rome. They are the work of an intensely theological and polemical age, when religious controversy absorbed the attention of all classes of society. They embody the results of the great conflict with the Papacy. A smaller class of Confessions (as the Articles of Dort and the Westminster Standards) belongs to the seventeenth century, and grew out of internal controversies among Protestants themselves. The eighteenth century witnessed a powerful revival of practical religion and missionary zeal through the labors of the Pietists and Moravians in Germany, and the Methodists in England and North America, but, in its ruling genius, it was irreligious and revolutionary, and undermined the authority of all creeds. In the nineteenth century a new interest in the old creeds was awakened, and several attempts were made to reduce the lengthy confessions to brief popular summaries, or to formularize the doctrinal \emph{consensus} of the different evangelical denominations. The present tendency among Protestants is to diminish rather than to increase the number of articles of faith, and to follow in any new formula the simplicity of the Apostles' Creed; while Romanism pursues the opposite course.

The symbols of the Reformation are very numerous, but several of them were merely provisional, and subsequently superseded by maturer statements of doctrine. Some far exceed the proper limits of a creed, and are complete systems of theology for the use of the clergy. It was a sad mistake and a source of incalculable mischief to incorporate the results of every doctrinal controversy with the confession of faith, and to bind lengthy discussions, with all their metaphysical distinctions and subtleties, upon the conscience of every minister and teacher. There is a vast difference between theological opinions and articles of faith. The development of theology as a science must go on, and will go on in spite of all these shackles.

As to the theology of the confessions of orthodox Protestantism, we may distinguish in them three elements, the œcumenical, the Augustinian, and the evangelical proper.

1. The œcumenical element. In theology and Christology the Protestant symbols agree with the Greek and Roman Churches, and also in the other articles of the Apostles' and Nicene Creeds from the creation of the world to the resurrection of the body.

2. The Augustinian element is found in anthropology, or the doctrines of sin and grace, predestination, and perseverance. Here the Protestant confessions agree with the system of Augustine, who had more influence upon the reformers than any uninspired teacher. The Latin Church during the Middle Ages had gradually fallen into Pelagian and semi-Pelagian doctrines and practices, although these had been condemned in the fifth century. The Calvinistic confessions, however, differ from the Lutheran in the logical conclusions derived from the Augustinian premises, which they hold in common.

3. The Evangelical Protestant and strictly original element is found in soteriology, and in all that pertains to subjective Christianity, or the personal appropriation of salvation. Here belong the doctrines of the rule of faith, of justification by faith, of the nature and office of faith and good works, of the assurance of salvation; here also the protest against all those doctrines of Romanism which are deemed inconsistent with the Scripture principle and with justification by faith. The papacy, the sacrifice of the mass, transubstantiation, purgatory, indulgences, meritorious and hypermeritorious works, the worship of saints, images, and relics are rejected altogether, while the doctrine of the Church and the Sacraments was essentially modified.

\xsection{The Lutheran and Reformed Confessions.}

§ 39. The Lutheran and Reformed Confessions.

Literature.

Max. Göbel : \emph{Die religiöse Eigenthümlichkeit der luther. und reformirten Kirche}. Bonn, 1837. (This book started a good deal of discussion in Germany on the peculiar genius of the two churches.)

C. B. Hundeshagen : \emph{Die Conflicte des Zwinglianismus, Lutherthums, und Calvinismus in der Bernischen Landeskirche von} 1522–1558. Berne, 1843. (The esteemed author died in Bonn, 1872.)

Merle D’aubigné  (d. 1872): \emph{Luther and Calvin,} translated into English, New York, 1846.

Alex. Schweizer : \emph{Glaubenslehre der reformirten Kirche}. Zürich, 1844, Vol. I. pp. 7–83.

M. Schneckenburger : \emph{Vergleichende Darstellung des luther. und reform. Lehrbegriffs}. Stuttgart, 1855, 2 vols. (Very acute and discriminating.) Comp. the introduction by \emph{Güder,} the editor.

Philip Schaff:   \emph{Germany; its Universities, Theology, and Religion}. Philadelphia, 1857, Ch. xviii. and xx., Lutheranism and Reform and the Evang. Union, pp. 167–185.

Essays on the same subject by Lücke, in the \emph{Deutsche Zettschrift,} Berlin, for 1853, Nos. 3 sqq.; Hagenbach, in the \emph{Studien und Kritiken} for 1854, Vol. I. pp. 23–34.

Jul. Müller  (Professor in Halle): \emph{Lutheri et Calvini sententiæ de Sacra Cœna inter se comparatæ,} Halle, 1858. Also in his \emph{Dogmatische Abhandlungen,} Bremen, 1870, pp. 404–467.

Catholicism assumed from the beginning, and retains to this day, two distinct and antagonistic types, the Greek and the Roman, which represent a Christian transformation of the antecedent and underlying nationalities of speculative Greece and world-conquering Rome. In like manner, but to a much larger extent (as may be expected from the greater liberty allowed to national and individual rights and peculiarities), is Protestantism divided since the middle of the sixteenth century into the Lutheran and the Reformed Confessions. To the former belong the established churches in most of the German States, in Denmark, Sweden, and Norway, and all others which call themselves after Luther; the Reformed—in the historical arid Continental sense of the term\footnote{\par As used in all Continental works on Church history and symbolics. It means originally the Catholic Church reformed of abuses, or regenerated by the Word of God.}—embraces the national evangelical churches of Switzerland, France, Holland, some parts of Germany, England, Scotland, with their descendants in America and the British colonies.

The designation \emph{Reformed} is insufficient to cover all the denominations and sects which have sprung directly or indirectly from this family since the Reformation, especially in England during the conflict of the Established Church with Puritanism and nonconformity; and hence in English and American usage it has given way to sectional and specific titles, such as \emph{Presbyterians, Episcopalians, Congregationalists, Baptists, Wesleyans} or \emph{Methodists,} etc. The term \emph{Calvinism} designates not a church, but a theological school in the Reformed Church, which in some sections allows also Arminian views. \emph{Puritanism,} likewise, is not a term for a distinct ecclesiastical organization, but for a tendency and party which exerted a powerful influence in the Anglican and other Reformed Churches on questions of doctrine, government, discipline, and worship.

Among the original Reformed Churches the Anglican stands out in many respects distinctly as a third type of Protestantism: it is the most powerful and the most conservative of all the national or established churches of the Reformation, and retains the entire basis of the mediæval hierarchy, without the papacy; it is a compromise between Catholicism and Protestantism, cemented by the royal supremacy, and leaves room, for Romanizing high-churchism and Puritanic low-churchism, as well as for intervening broad-churchism. But its original doctrinal \emph{status} was moderately Calvinistic, and for a time it made even common cause with the ultra-Calvinistic Synod of Dort.

The doctrinal difference between Lutheranism and Reform was originally confined to two articles, namely, the nature of Christ's presence in the Sacrament of the Eucharist, and the extent of God's sovereignty in the ante-historic and premundane act of predestination. At the Conference held in Marburg, Luther and Zwingli agreed in fourteen and a half articles, and differed only in the other half of the fifteenth article, concerning the real presence.\footnote{\par The fifteenth and last of the Marburg articles treats of the Lord's Supper, and after stating the points of agreement, concludes thus: 'And although at present we can not agree whether the true body and the true blood of Christ be corporeally present in the bread and wine (\emph{ob der wahre Leib und das wahre Blut Christi leiblich im Brode und Weine gegenwärtig sei}), yet each party is to show to the other Christian love, as far as conscience permits (\emph{so weit es das Gewissen jedem gestattet}), and both parties should fervently pray to Almighty God that by his Spirit he may strengthen us in the true understanding. Amen.'} The Swiss reformer saw in this difference no obstacle to fraternal fellowship with the Wittenbergers, with whom, he said, he would rather agree than with any people on earth, and, with tears in his eyes, he extended his hand to Luther; but the great man, otherwise so generous and liberal, who had himself departed from the Catholic Church in much more essential points, felt compelled in his conscience to withhold his hand on account of a general difference of 'spirit,'\footnote{\par ' \emph{Ihr habt einen andern Geist,} ' said Luther to Zwingli.} which revealed itself in subsequent controversies, and defeated many attempts at reunion.

The internal quarrels among Christian brethren, which are found more or less in all denominations and ages,\footnote{\par The feuds between monastic orders and theological schools in the Roman and Greek Churches, and the quarrels even in the œcumenical Councils, from the Nicene down to the Vatican, are fully equal in violence and bitterness to the Protestant controversies in the sixteenth and seventeenth centuries, and are less excusable on account of the boasted doctrinal unity of those churches.} are the most humiliating and heart-sickening chapters in Church history, but they are overruled by Providence for the fuller development of theology, a wider spread of Christianity, and a deeper divine harmony, which will ultimately, in God's own good time, spring out of human discord.

The two great families of Protestantism are united in all essential articles of faith, and their members may and ought to cultivate intimate Christian fellowship without sacrifice of principle or loyalty to their communion. Yet they are distinct ecclesiastical individualities, and Providence has assigned them peculiar fields of labor. Their differences in theology, government, worship, and mode of piety are rooted in diversities of nationality, psychological constitution, education, external circumstances, and gifts of the Spirit.

1. The Lutheran Church arose in monarchical Germany, and bears the impress of the German race, of which Luther was the purest and strongest type. The Reformed Church began, almost simultaneously, in republican Switzerland, and spread in France, Holland, England, and Scotland. The former extended, indeed, to kindred Scandinavia, and, by emigration, to more distant countries. But outside of Germany it is stunted in its normal growth, or undergoes, with the change of language and nationality, an ecclesiastical transformation.\footnote{\par This is the case with the great majority of Anglicized and Americanized Lutherans, who adopt Reformed views on the Sacraments, the observance of Sunday, Church discipline, and other points.} The Reformed Church, on the other hand, while it originated in the German cantons of Switzerland, and found a home in several important parts of Germany, as the Palatinate, the Lower Rhine, and (through the influence of the House of Hohenzollern since the Elector Sigismund, 1614) in Brandenburg and other provinces of Prussia, was yet far more fully and vigorously developed among the maritime and freer nations, especially the Anglo-Saxon race, and follows its onward march to the West and the missionary fields of the East. The modern Protestant movements among the Latin races in the South of Europe likewise mostly assume the Reformed, some even a strictly Calvinistic type. Converts from the excessive ritualism of Rome are apt to swing to the opposite extreme of Puritan simplicity.

Germany occupies the front rank in sacred learning and scientific theology, but the future of evangelical Protestantism is mainly intrusted to the Anglo-American churches, which far surpass all others in wealth, energy, liberality, philanthropy, and a firm hold upon the heart of the two great nations they represent.

2. The Lutheran Church, as its name indicates, was rounded and shaped by the mighty genius of Luther, who gave to the Germans a truly vernacular Bible, Catechism, and hymn-book, and who thus meets them at every step in their public and private devotions. We should, indeed, not forget the gentle, conciliatory, and peaceful genius of Melanchthon, which never died out in the Lutheran Confession, and forms the connecting link between it and the Reformed. He represents the very spirit of evangelical union, and practiced it in his intimate friendship with the stern and uncompromising Calvin, who in turn touchingly alludes to the memory of his friend. But the influence of the '\emph{Præceptor Germaniæ}' was more scholastic and theological than practical and popular. Luther was the originating, commanding reformer, 'born,' as he himself says, 'to tear up the stumps and dead roots, to cut away the thorns, and to act as a rough forester and pioneer;' while 'Melanchthon moved gently and calmly along, with his rich gifts from God's own hand, building and planting, sowing and watering.' Luther was, as Melanchthon called him, the Protestant Elijah. He spoke almost with the inspiration and authority of a prophet and apostle, and his word shook the Church and the Empire to the base. He can be to no nation what he is to the German, as little as Washington can be to any nation what he is to the American.\footnote{\par Luther can only be fully understood by a German, while a Frenchman or an Englishman (with some exceptions, as Coleridge, Hare, Carlyle) is likely to be repelled by some of his writings, \emph{e.g.,} his coarse book against Henry VIII. Hence the unfavorable judgments of such scholars as Hallam, Sir William Hamilton, Pusey; while, on the other hand, even liberal Catholics among German scholars can not but admire him as Germans. Dr. Döllinger, long before his secession from Rome, said (in his book \emph{Kirche und Kirchen}): '\emph{Luther ist der gewaltigste Volksmann, der populärste Charakter, den Deutschland je besessen. In dem Geiste dieses Mannes, des grössten unter den Deutschen seines Zeitalters, ist die protestantische Doctrin entsprungen. Vor der Ueberlegenheit und schöpferischen Energie dieses Geistes bog damals der aufstrebende, thatkräftige Theil der Nation demuthsvoll und gläubig die Kniee.}' The towering greatness of Luther is to the Lutherans a constant temptation to hero-worship, as Napoleon's brilliant military genius is a misfortune and temptation to France. Lessing expressed his satisfaction at the discovery of some defects in Luther's character, since he was, as he says, 'in imminent danger of making him an object of idolatrous veneration. The proofs that in some things he was like other men are to me as precious as the most dazzling of his virtues.' There are not a few Lutherans who have more liking for Luther's faults than for his virtues, and admire his conduct at Marburg as much, if not more, than his conduct at Worms. A very respectable Lutheran professor of theology resolved the difference between Luther and Calvin into this: that the one was human, the other inhuman! Calvin once nobly said, 'Though Luther should call me a devil, I would still revere and love him as an eminent servant of God.' If he was cruel, according to our modern notions, in his treatment of Servetus, he acted in the spirit of his age, and was approved even by the gentle Melanchthon. His followers need fear no comparison with any other Christians as to humanity and liberality.} And yet, strange to say, with all the overpowering influence of Luther, his personal views on the canon\footnote{\par He irreverently called the Epistle of St. James an 'epistle of straw,' and had objections to the Epistle to the Hebrews, the Apocalypse, and the Book of Esther. He was as thoroughly convinced of the inspiration and authority of the Word of God as the most orthodox divine can be, but he had free views on the mode of inspiration and the extent of the traditional canon.} and on predestination\footnote{\par Luther, in his work \emph{De servo arbitrio,} against Erasmus, written in 1525, teaches the slavery of the human will, the dualism in the divine will (secret and revealed), and unconditional predestination to salvation and damnation, in language stronger than even Calvin ever used, who liked the views of that book, but objected to some of its hyperbolical expressions (\emph{Opera,} Tom. VII. p. 142). Melanchthon, who originally held the same Augustinian theory (like all the Reformers), gradually changed it (openly since 1535) in favor of a synergistic theory. But Luther never recalled his tract against Erasmus; on the contrary, he counted it among his best, and among the few of his books which he would not be willing 'to swallow, like Saturn his own children.' He never made this a point of difference from the Swiss. In the Articles of Smalcald, 1537 (III. i. p. 318, ed. Hase), he again denied the freedom of the will, as a scholastic error; and in his commentary on \emph{Genesis} (Ch. vi. 6, 18; xxvi), one of his last works, he taught the same view of the secret will of God as in 1525. Comp. J. Müller: \emph{Lutheri de prædestinatione et libero arbitrio doctrina,} 1832, and his \emph{Dogmat. Abhandlungen,} 1870, pp. 187sqq.; Lütkens: \emph{Luther's Prædestinationslehre im Zusammenhang mit seiner Lehre vom freien Willen,} 1858; Köstlin: \emph{Luther's Theologie in ihrer geschichtl. Entwicklung,} 1863, Vol. II. pp. 32–55, 300–331; Schweizer: \emph{Die protest. Centraldogmen,} 1854, Vol. I. pp. 57 sqq.; Dorner: \emph{Geschichte der protest. Theologie,} 1867, Vol. I. pp. 194 sqq.} were never accepted by his followers; and if we judge him by the standard of the Form of Concord, he is a heretic in his own communion as much as St. Augustine, on account of his doctrines of sin and grace, is a heretic in the Roman Church, revered though he is as the greatest among the Fathers.

The Reformed Church had a large number of leaders, as Zwingli, Œcolampadius, Bullinger, Calvin, Beza, Cranmer, Knox, but not one of them, not even Calvin, could impress his name or his theological system upon her. She is independent of men, and allows full freedom for national and sectional modifications and adaptations of the principles of the Reformation.

3. The Lutheran Confession starts from the wants of sinful man and the personal experience of justification by faith alone, and finds, in this 'article of the standing and falling Church,' comfort and peace of conscience, and the strongest stimulus to a godly life. The Reformed Churches (especially the Calvinistic sections) start from the absolute sovereignty of God and the supreme authority of his holy Word, and endeavor to reconstruct the whole Church on this basis. The one proceeds from anthropology to theology; the other, from theology to anthropology. The one puts the subjective or material principle of the Reformation first, the objective or formal next; the other reverses the order; yet both maintain, in inseparable unity, the subjective and objective principles of the Reformation.

The Augsburg Confession, which is the first and the most important Lutheran symbol, does not mention the Bible principle at all, although it is based upon it throughout;\footnote{\par The Preface of the Augsburg Confession declares that the Confession is 'drawn from the holy Scriptures and the pure Word of God.'} the Articles of Smalcald mention it incidentally;\footnote{\par Part II. (p. 309): 'The Word of God, and no one else, not even an angel, can establish articles of faith.' ('\emph{Regulam aliam habemus, ut videlicet Verbum Dei condat articulos fidei, et præterea nemo, ne angelus quidem.')}} and the Form of Concord more formally.\footnote{\par \emph{Form. Conc.,} Part I. or Epit., at the beginning: 'We believe, teach, and confess that the only rule and standard (\emph{unicam regulam et normam}), according to which all doctrines and teachers alike ought to be tried and judged, are the prophetic and apostolic Scriptures of the Old and New Testaments alone.' Comp. Preface to the Second Part.} But the Reformed Confessions have a separate article \emph{de Scriptura Sacra,} as the only rule of faith and discipline, and put it at the head, sometimes with a full list of the canonical books.\footnote{\par \emph{Conf. Helv.} II. ch. i. (\emph{De Scriptura sancta, vero Dei verbo}): ' \emph{Credimus et confitemur Scripturas canonicas sanctorum Prophetarum et Apostolorum utriusque Testamenti, ipsum verum esse Verbum Dei: et auctoritatem sufficientem ex semetipsis, non ex hominibus habere.} ' \emph{Conf. Helv.} I. (Basil. II.) art. 1; \emph{Conf. Gall.} art. 2–5; \emph{Conf. Scot.} art. 18, 19; \emph{Conf. Belg.} art. 2–7; \emph{art. Angl.} art. 6 (\emph{Scriptura sacra continet omnia quæ ad salutem sunt necessaria,} etc., with a list of the canonical books, from which the Apocrypha are carefully distinguished); \emph{Westminster Conf. of Faith,} ch. i. (more fully), etc. The exception of the first Confession of Basle is only apparent, for it \emph{concludes} with a submission of all its articles to the supreme authority of the Scriptures (\emph{Postremo, hanc nostrum confessionem judicio sacræ biblicæ Scripturæ subjicimus; eoque pollicemur, si ex prædictis Scripturis in melioribus instituamur, nos ommi tempore Deo et sacrosancto ipsius Verbo maxima cum gratiarum actione obsecuturos esse}').}

4. The Lutheran Church has an idealistic and contemplative, the Reformed Church a realistic and practical, spirit and tendency. The former aims to harmonize Church and State, theology and philosophy, worship and art; the latter draws a sharper line of distinction between the Word of God and the traditions of men, the Church and the world, the Church of communicants and the congregation of hearers, the regenerate and the unregenerate, the divine and the human. The one is exposed to the danger of pantheism, which shuts God up within the world; the other to the opposite extreme of deism, which abstractly separates him from the world. Hence the leaning of the Lutheran Christology to Eutychianism, the leaning of the Reformed to Nestorianism.

The most characteristic exponent of this difference between the two confessions is found in their antagonistic doctrines of the Lord's Supper; and hence their controversies clustered around this article, as the Nicene and post-Nicene controversies clustered around the person of Christ. Luther teaches the real presence of Christ's body and blood \emph{in,} \emph{with,} and \emph{under} the elements, the oral manducation by unworthy as well as worthy communicants, and the ubiquity of Christ's body; while Zwingli and Calvin, carefully distinguishing the sacramental sign from the sacramental grace, teach—the one only a symbolical, the other a spiritual real, presence and fruition for believers alone. The Romish doctrine of transubstantiation is equally characteristic of the magical supernaturalism and asceticism of Romanism, which realizes the divine only by a miraculous annihilation of the natural elements. Lutheranism sees the supernatural \emph{in} the natural, Calvinism \emph{above} the natural, Romanism \emph{without} the natural.

5. Viewed in their relations to the mediæval Church, Lutheranism is more conservative and historical, the Reformed Church more progressive and radical, and departs much further from the traditionalism, sacerdotalism, and ceremonialism of Rome. The former proceeded on the principle to retain what was not forbidden by the Bible; the latter, on the principle to abolish what was not commanded.

The Anglican Church, however, though moderately Calvinistic in her Thirty-nine Articles, especially in the doctrine on the Scriptures and the Sacraments, makes an exception from the other Reformed communions, since it retained the body of the episcopal hierarchy and the Catholic worship, though purged of popery. Hence Lutherans like to call it a 'Lutheranizing Church;' but the conservatism of the Church of England was of native growth, and owing to the controlling influence of the English monarchs and bishops in the Reformation period.

6. The Lutheran Confession, moreover, attacked mainly the Judaism of Rome, the Reformed Church its heathenism. 'Away with legal bondage and work righteousness!' was the war-cry of Luther; 'Away with idolatry and moral corruption!' was the motto of Zwingli, Farel, Calvin, and Knox.

7. Luther and Melanchthon were chiefly bent upon the purification of doctrine, and established State churches controlled by princes, theologians, and pastors. Calvin and Knox carried the reform into the sphere of government, discipline, and worship, and labored to found a pure and free church of believers. Lutheran congregations in the old world are almost passive, and most of them enjoy not even the right of electing their pastor; while well-organized Reformed congregations have elders and deacons chosen from the people, and a much larger amount of lay agency, especially in the Sunday-school work. Luther first proclaimed the principle of the general priesthood, but in practice it was confined to the civil rulers, and carried out in a wrong way by making them the supreme bishops of the Church, and reducing the Church to a degrading dependence on the State.

8. Luther and his followers carefully abstained from politics, and intrusted the secular princes friendly to the Reformation with the episcopal rights; Calvin and Knox upheld the sole headship of Christ, and endeavored to renovate the civil state on a theocratic basis. This led to serious conflicts and wars, but they resulted in a great advance of civil and religious liberty in Holland, England, and the United States. The essence of Calvinism is the sense of the absolute sovereignty of God and the absolute dependence of man; and this is the best school of moral self-government, which is true freedom. Those who feel most their dependence on God are most independent of men.\footnote{\par The principles of the Republic of the United States can be traced, through the intervening link of Puritanism, to Calvinism, which, with all its theological rigor, has been the chief educator of manly characters and promoter of constitutional freedom in modern times. The inalienable rights of an American citizen are nothing but the Protestant idea of the general priesthood of believers applied to the civil sphere, or developed into the corresponding idea of the general kingship of free men.}

9. The strength and beauty of the Lutheran Church lies in its profound theology, rich hymnology, simple, childlike, trustful piety; the strength and beauty of the Reformed Churches, in aggressive energy and enterprise, power of self-government, strict discipline, missionary zeal, liberal sacrifice, and faithful devotion, even to martyrdom, for the same divine Lord. From the former have proceeded Pietism and Moravianism, a minutely developed scholastic orthodoxy, speculative systems and critical researches in all departments of sacred learning, but also antinomian tendencies, and various forms of mysticism, rationalism, and hypercriticism. The latter has produced Puritanism, Congregationalism, Methodism, Evangelicalism (in the Church of England), the largest Bible, tract, and missionary societies, has built most churches and benevolent institutions, but is ever in danger of multiplying sectarian divisions, overruling the principle of authority by private judgment, and disregarding the lessons of history.

10. Both churches have accomplished, and are still accomplishing, a great and noble work. Let them wish each other God's speed, and stimulate each other to greater zeal. A noble rivalry is far better than sectarian envy and jealousy. There have been in both churches, at all times, men of love and peace as well as men of war, with corresponding efforts to unite Lutheran and Reformed Christians, from the days of Melanchthon and Bucer, Calixtus and Baxter, down to the Prussian Evangelical Union, the German Church Diet, and the Evangelical Alliance. Even the exclusive Church of England has entered into a sort of alliance with the Evangelical Church of Prussia in jointly founding and maintaining the Bishopric of St. James in Jerusalem.\footnote{\par Chiefly the work of Chevalier Bunsen and his congenial friend, Frederick William IV.}

The time for ecclesiastical amalgamation, or organic union, has not yet come, but Christian recognition and union in essentials is quite consistent with denominational distinctions in non-essentials, and should be cultivated by all who love our common Lord and Saviour, and desire the triumph of his kingdom.

\xchapter{Chapter 6. The Creeds of the Evangelical Lutheran Church.}

\xsection{The Lutheran Confessions.}

§ 40. The Lutheran Confessions.

Literature.

I. Collections of the Lutheran Symbols.

(1.) Latin Editions.

Concordia. \emph{Pia et unanimi con ensu repetita Confessio Fidei et Doctrinæ Electorum, Principum et Ordinum Imperii, atque eorundem Theologorum, qui Augustanam Confessionem amplectuntur et nomina sua huic libro subscripserunt. Cui ex Sacra Scriptura, unica illa veritatis norma et regula quorundam Articulorum, qui post Doctoris Martini Lutheri felicem ex hac vita exitum, in controversiam venerunt, solida accessit Declaratio}, etc. (By Selnecker.) Lips. 1580, 4to; 1584. The second ed. '\emph{communi consilio et mandato Electorum.}' Another edition, Lips. 1602, 8vo, by order and with a Preface of Christian II., Elector of Saxony; republished, Lips. 1606, 1612, 1618, 1626, 8vo; Stettin, 1654, 8vo; Lips. 1669, 8vo; 1677. The second ed. (746 pages) is the authentic Latin \emph{editio princeps.}

The same edition, \emph{cum Appendice tripartita} Dr. Adami Rechenbergii, Lips. first, 1677, 1678, 1698, 1712, 1725; last, 1742. Rechenberg's edition is the standard of reference, followed by the later Latin editions in the paging.

Ecclesiæ Evangelicæ Libri Symbolici, etc. C. M. Pfaffius, \emph{ex editionibus primis et præst. recensuit, varias lectiones adjunxit}, etc. Tubing. 1730, 8vo.

Libri Symbolici Ecclesiæ Evangelico-Lutheranæ \emph{accuratius editi variique generis animadvers. ac disput. illustrati a }Mich. Webero. Viteb. 1809, 8vo.

Libri Symbolici Ecclesiæ Evangelicæ. \emph{Ad fidem optim. exemplorum recens.} J. A. H. Tittmann. Lips. 1817, 8vo; 1827.

Libri Symbolici Ecclesiæ Evangelicæ sive Concordia. \emph{Recens.} C. A. Hase. Lipsiæ, 1827, 8vo; 1837, 1845.

Libri Symbolici Ecclesiæ Lutheranæ \emph{ad editt. principes et ecclesiæ auctoritate probat. rec., præcipuam lectionum diversitatem notavit, Christ. II. ordinumque evangelicor. præfationes, artic. Saxon. visitator. et Confut. A. C. Pontific. adj. } H. A. Guil. Meyer. Gotting. 1830, 8vo.

Concordia. \emph{Libri Symbolici Ecclesiæ Evang. Ad edit. Lipsiensem}, 1584; Berolin. (Schlawitz), 1857, 8vo.

(2.) German Editions.

Concordia. \texthebrew{יהוה }\emph{Christliche, Widerholete, einmütige Bekenntnüs nachbenanter Churfürsten, Fürsten und Stende Augspurgischer Confession, und derselben zu ende des Buchs underschriebener Theologen Lere und Glaubens. Mit angeheffter, in Gottes wort, als der einigen Richtschnur, wohlgegründter erklerung etlicher Artickel, bei welchen nach D. Martin Luther's seligen absterben disputation und streit vorgefallen. Aus einhelliger vergleichung und bevehl obgedachter Churfürsten, Fürsten und Stende, derselben Landen, Kirchen, Schulen und Nachkommen, zum underricht und warnung in Druck verfertiget. Mit Churf. Gnaden zu Sachsen befreihung.} Dresden, 1580, fol. (See the whole title in Corp. \emph{Ref.} Vol. XXVI. p. 443.)

Concordia. Magdeburg, 1580, 4to, two ed.; Tübingen, 1580, fol.; Dresden, 1581, 4to; Frankfurt a. O., 1581, fol.; Magdeburg, 1581, 4to; Heidelberg, 1582, fol., two ed.; Dresden, 1598, fol.; Tübingen, 1599, 4to; Leipzig, 1603, 4to; Stuttgart, 1611, 4to; Leipzig, 1622, 4to; Stuttgart, 1660, 4to; 1681, 4to.

Concordia. \emph{Mit }  Heinr. Pipping's  \emph{Hist. theol. Einl. zu den symb. Schriften der Evang. Luth. Kirchen.} Leipz. 1703, 4to; 2te Ausg. \emph{mit } Christ. Weissen's \emph{Schlussrede.} Leipz. 1739, 4to.

Christliches Concordienbuch, etc., von  Siegm. Jac. Baumgarten. Halle, 1747, 2 vols. 8vo.

Christl. Concordienbuch \emph{mit der Leipziger Theol. Facultaet Vorrede.} Wittenberg, 1760, 8vo; 1766, 1789.

Die Symb. Bücher der Ev. Luth. Kirche, etc., von  J. W. Schöpff. Dresden, 1826–27, 8vo.

Concordia.  \emph{Die Symb. Bücher der ev. luth. Kirche,} etc., von  F. A. Koethe. Leipzig, 1830, 8vo.

Evangel. Concordienbuch, etc., von  J. A. Detzer. Nürnberg, 1830, 1842, 1847.

Evangel. Concordienbuch, etc., von  Fr. W. Bodemann. Hanover, 1843.

Christliches Concordienbuch, New York, 1854.

(3.) German-Latin Editions.

Concordia. \emph{Germanico-Latina ad optima et antiquissima exempla recognita, adjectis fideliter allegator. dictor. S. Scr. capitibus et vers. et testimoniorum P. P. aliorumque Scriptorum locis. . . . cum approbatione Facult. Theol. Lips. Wittenb. et Rostoch. Studio } Ch. Reineccii. Lips. 1708, 4to; 1735.

Christliches Concordienbuch. \emph{Deutsch und Lateinisch mit historischen Einleitungen} J. G. Walch's. Jena, 1750, 8vo.

Die Symbolischen Bücher der Evang. Luther. Kirche, \emph{deutsch und lateinisch}, etc., von  J. F. Müller (of Windsbach, Bavaria), 1847; 3d ed. Stuttgart, 1869. (A very useful edition.)

(4.) Translations.

Dutch: Concordia\emph{of Lutersche Geloofs Belydenis in’t licht gegeven door }  Zach. Dezius. Rotterdam, 1715, 8vo.

Swedish: Libri Concordiæ Versio Suecica, Christeliga, Enhelliga, och Uprepade och Läras, etc. Norköping, 1730, 4to.

English: The Christian Book of Concord, \emph{or Symbolical Book of the Evangelical Lutheran Church, translated by } Ambrose and  Socrates Henkel (two Lutheran clergymen of Virginia), \emph{with the assistance of several other Lutheran clergymen.} Newmarket, Virginia, 1851; 2d ed. revised, 1854. This is the first and only complete English edition of the Book of Concord; but the translation (made from the German) is not sufficiently idiomatic.

II. Historical and Critical Works on the Lutheran Symbols in General.

Jo. Benedict Carpzov: \emph{ Isagoge in libros ecclesiarum Lutheranarum symbolicos. Opus posthumum a } J. Oleario: \emph{Continuatum ed. } J. B. Carpzov (\emph{filius}). Lipsiæ, 1665, 4to; 1675, 1691, 1699, 1725.

Jo. Georg Walch: \emph{Introductio in libros Ecclesiæ Lutheranæ symbolicos, observationibus historicis et theologicis illustrata.} Jenæ, 1732, 4to.

J. Albr. Fabricius: \emph{Centifolium Lutheranum.} Hamb. 1728–30, 2 vols. 8vo.

S. J. Baumgarten: \emph{Erleuterungen der im christlichen Concordienbuch enthaltenen symbolischen Schriften der evang. luth. Kirche, nebst einem Anhange von den übrigen Bekenntnissen und feierlichen Lehrbüchern in gedachter Kirche.} Halle, 1747.

J. Christoph. Kœcher: \emph{ Bibliotheca theologiæ symbolicæ et catecheticæ.} Guelph. et Jenæ, 1751–69, 2 vols.

Jac. W. Feuerlin: \emph{Bibliotheca symb. evang. Lutherana. Accedunt appendices duæ}: I. \emph{Ordinationes et Agenda}; II. \emph{Catechismus ecclesiarum nostrarum.} Gotting. 1752. Another enlarged edition by  J. Barthol. Riederer. Nürnberg, 1768, 2 vols. 8vo.

J. G. Walch: \emph{Bibliotheca theologica selecta.} Jena, 1757–65, 4 vols. 8vo.

Chr. Guil. Fr. Walch: \emph{ Breviarium theol. symb. eccles. luther.} Göttingen, 1765–1781, 8vo.

Eduard Köllner: \emph{Symbolik der lutherischen Kirche.} Hamburg, 1837.

J. F. Müller: \emph{Die symb. Bücher der evang. luth. Kirche.} Stuttgart, 1847; 3d ed. 1869. Introduction pp. cxxiv.

Charles P. Krauth (Dr. and Prof. of Theology in the Evang. Theol. Seminary in Philadelphia): \emph{The Conservative Reformation and its Theology, as represented in the Augsburg Confession and in the History and literature of the Evang. Lutheran Church.} Philadelphia, 1871.

For fuller lists of editions and works, see Feuerlin (ed. Riederer), J. G. Walch, Köllner, l.c., and the 26th and 27th vols. of the \emph{Corpus Reformatorum}, ed. Bindseil.

The Evangelical Lutheran Church, in whole or in part, acknowledges nine symbolical books: three of them are inherited from the Catholic Church, viz., the Apostles' Creed, the Nicene Creed (with the \emph{Filioque}), and the Athanasian Creed; six are original, viz., the Augsburg Confession, drawn up by Melanchthon (1530), the Apology of the Confession, by the same (1530), the Articles of Smalcald, by Luther (1537), the two Catechisms of Luther (1529), and the Form of Concord, prepared by six Lutheran divines (1577).

These nine symbols constitute together the Book of Concord (\emph{Concordia}, or \emph{Liber Concordiæ}, \emph{Concordienbuch}), which was first published by order of Elector Augustus of Saxony in 1580, in German and Latin, and which superseded older collections of a similar character.\footnote{See an account of the various \emph{Corpora Doctrinæ} in Baumgarten, \emph{Erläuterungen}, etc., pp. 247–282; Köllner, \emph{Symbolik}, I. pp. 96 sqq.; and Müller, \emph{Symb. Bücher}, pp. cxxii. sqq. The oldest was the \emph{Corpus Doctrinæ Christianæ Philippicum}, or \emph{Misnicum}, 1560, which contained only Melanchthonian writings, and was followed by several other collections of a more strictly Lutheran character.}

The Lutheran symbols are not of equal authority. Besides the three œcumenical Creeds, the Augsburg Confession is most highly esteemed, and is the only one which is generally recognized. Next to it comes the Shorter Catechism of Luther, which is extensively used in catechetical instruction. His Larger Catechism is only an expansion of the Shorter. The Apology is valuable in a theological point of view, as an authentic commentary on the Augsburg Confession. The Smalcald Articles have an historical significance, as a warlike manifesto against Rome, but are little used. The Form of Concord was never generally received, but decidedly rejected in several countries, and is disowned by the Melanchthonian and unionistic schools in the Lutheran Church.

Originally intended merely as testimonies or confessions of faith, these documents became gradually binding formulas of public doctrine, and subscription to them was rigorously exacted from all clergymen and public teachers in Lutheran State churches.\footnote{\par As early as 1533 a statute was enacted in Wittenberg by Luther, Jonas, and others, which required the doctors of theology, at their promotion, to swear to the incorrupt doctrine of the Gospel as taught in the symbols. It was only a modification of the oath customary in the Roman Catholic Church. After the middle of the sixteenth century, subscription began to be enforced, on pain of deposition and exile. See Köllner, \emph{Symb.}, I. pp. 106 sqq.} The rationalistic apostasy, reacting against the opposite extreme of symbololatry and ultra-orthodoxy, swept away these test-oaths, or reduced them to a hypocritical formality. The revival of evangelical Christianity, since the tercentenary jubilee of the Reformation in 1817, was followed by a partial revival of rigid Lutheran confessionalism, yet not so much in opposition to the Reformed as to the Unionists in Prussia and other German States, where the two Confessions have been amalgamated. The meaning and aim of the Evangelical Union in Prussia, however, was not to set aside the two Confessions, but to accommodate them in one governmental household, allowing them to use either the Lutheran or the Heidelberg Catechism as before. The chief trouble was occasioned by the new liturgy of King Frederick William III., which was forced upon the churches, and gave rise to the Old Lutheran secession. In the other States of Germany, and in Scandinavia and Austria, the Lutheran churches have, with a separate government, also their own liturgies and forms of ordination, with widely differing modes of subscription to the symbolical books.\footnote{\par Köllner, I. pp. 121 sqq., gives a number of \emph{Verpflichtungsformeln} in use in Europe.}

In the United States, the Lutherans, left free from the control of the civil government, yet closely connected with the doctrinal and confessional disputes of their brethren in Germany, are chiefly divided into three distinct organizations, which hold as many different relations to the Symbolical Books, and are, in fact, three denominations under a common name, viz.: the General Synod of the Evangelical Lutheran Church of the United States, organized in 1820; the Synodical Conference of North America, organized in 1872;\footnote{\par{} [The statements must be modified in view of the organic unions and Church federations which have recently been formed within the Lutheran communions—movements encouraged by the 400th anniversaries of the XCV Theses, 1917, and the Augsburg Confession, 1930. To follow a statement furnished by the Rev. G. L. Kieffer, Statistician and Librarian of the National Luth. Council—the Luth. churches of the U. S. and Canada, 1930, had a membership of 2,852,843 communicants. Two-thirds of the number are embraced in three corporate groups, namely, The United Luth. Ch. of Am., formed 1918, with 971,187 members; The Am. Luth. Ch., formed 1930 with 340,809 members; The Evang. Luth. Synod of Missouri, formed 1847, with 702,056 members. Two coöperative federations exist, namely: 1. The Am. Luth. Conference with 926,009 members, formed 1930, consisting of five bodies, The Am. Luth. Ch., The Augustana Synod with 234,434 members, the Norwegian Luth. Ch., with 303,358 members, The Luth. Free Ch. and the United Danish Churches with 47,408 members. 2. The Evang. Synod. Luth. Conference of N. Am., founded 1873, having 873,484 members, and consisting of five groups of which the Missouri Synod is much the largest. In addition to the United Luth. Ch. of Am. and the groups joined in the federations there are seven independent synods with 75,397 members. The groups are not to be regarded ``as separate denominations, their main difference being in a gradual gradation from the freedom of the universal priesthood of believers to a more or less highly developed legalistic control of the individual.'' They coöperate in certain general movements through a National Luth. Council and some of the independent synods support the missions of the larger groups.—Ed.]} and the General Council, which, under the lead of the old Synod of Pennsylvania, seceded from the General Synod, and met first at Fort Wayne, Indiana, Nov. 20, 1867. The first has its theological and literary centre in Gettysburg, the second at St. Louis and Fort Wayne, the third in Philadelphia.

The 'General Synod,' which is composed chiefly of English-speaking descendants of German immigrants, and sympathizes with the surrounding Reformed denominations, adopts simply 'the Augsburg Confession as a correct exhibition of the fundamental doctrines of the divine Word,' without mentioning the other symbolical books at all, and allows a very liberal construction even of the Augsburg Confession, especially the articles on the Sacraments.\footnote{\par 'We receive and hold, with the Evangelical Lutheran Church of our fathers, the Word of God, as contained in the canonical Scriptures of the Old and New Testaments, as the only infallible rule of faith and practice, and the Augsburg Confession, as a correct exhibition of the fundamental doctrines of the Divine Word, and of the faith of our Church founded upon that Word.' (\emph{Constitution of General Synod}, adopted at Washington, 1869, Art. II. Sect. 3.)} With this basis the Lutheran Synod of the Southern States, which was organized during the civil war, is substantially agreed.\footnote{'We receive and hold that the Old and New Testaments are the Word of God, and the only infallible rule of faith and practice. We likewise hold that the Apostles' Creed, the Nicene Creed, and the Augsburg Confession contain the fundamental doctrines of the sacred Scriptures; and we receive and adopt them as the exponents of our faith.'}

The Lutheran Synodical Conference of North America, which is so far almost exclusively German as to language, requires its ministers to subscribe the whole Book of Concord (including the Form of Concord), 'as the pure, unadulterated explanation and exposition of the divine Word and will.'\footnote{\par '\emph{Ich erkenne die drei Hauptsymbole der} [\emph{alten}] \emph{Kirche, die ungeänderte Augsburgische Confession und deren Apologie, die Schmalcaldischen Artikel, die beiden Catechismen Luthers und die Concordienformel für die reine, ungefälschte Erklärung und Darlegung des göttlichen Wortes and Willens, bekenne mich zu denselben als zu meinen eigenen Bekenntnissen und will mein Amt bis an mein Ende treulich und fleissig nach denselben ausrichten. Dazu stärke mich Gott durch seinen heiligen Geist! Amen.}' (Ordination vow in the \emph{Kirchen-Agende}, St. Louis, 1856, p. 173.) Here the Lutheran system of doctrine is almost identified with the Bible, according to the adage:  \par '\emph{Gottes Wort und Luther's Lehr} \par \emph{Vergehet nun und nimmermehr.}'}

With the Missourians are agreed the Buffalo and the Iowa Lutherans, except on the question of the origin and nature of the ministerial office, which has been the subject of much bitter controversy between them.

The 'General Council,' which is nearly equally divided as to language and nationality, stands midway between the General Synod and the Synodical Conference. It accepts, primarily, the 'Unaltered Augsburg Confession in its original sense,' and, in subordinate rank, the other Lutheran symbols, as explanatory of the Augsburg Confession, and as equally pure and Scriptural.\footnote{\par 'We accept and acknowledge the doctrines of the Unaltered Augsburg Confession in its original sense as throughout in conformity with the pure truth, of which God's Word is the only rule. We accept its statements of truth as in perfect accordance with the canonical Scriptures; we reject the errors it condemns, and believe that all which it commits to the liberty of the Church, of right belongs to that liberty. In thus formally accepting and acknowledging the Unaltered Augsburg Confession, we declare our conviction that the other Confessions of the Evangelical Lutheran Church, inasmuch as they set forth none other than its system of doctrine and articles of faith, are of necessity pure and Scriptural. Pre-eminent among such accordant, pure, and Scriptural statements of doctrine, by their intrinsic excellence, by the great and necessary ends for which they were prepared, by their historical position, and by the general judgment of the Church, are these: the Apology of the Augsburg Confession, the Smalcald Articles, the Catechisms of Luther, and the Formula of Concord, all of which are, with the Unaltered Augsburg Confession, in the perfect harmony of one and the same Scriptural faith.' (\emph{Principles of Faith and Church Polity of the Gen. Council}, adopted Nov. 1867, Sections VIII. and IX.)}

\xsection{The Augsburg Confession, 1530.}

§ 41. The Augsburg Confession, 1530.

Literature.

I. Editions, Latin and German. In the general collections of Lutheran Symbols, by Rechenberg, Walch, Hase, Müller, etc. (see § 40).

II. Separate Editions of the Augs. Conf.—in Latin or German, or both—by Twesten (1816), Winer (1825), Tittmann (1830), Spieker (1830), M. Weber (1830), Wiggers (1830), Beyschlag (1830), Funk (1830), Förstemann (1833), Härter (1838). The best critical edition of the Latin and German texts, with all the variations, is contained in the \emph{Corpus Reformatorum}, ed. Bretschneider and Bindseil, Vol. XXVI. (issued, Brunsvigæ, 1858), pp. 263 sqq.

For lists of older editions, see Köllner, \emph{Symbolik}, I. p. 344–353, and Bindseil, in \emph{Corp. Ref.} Vol. XXVI. pp. 211–263.

III. English Translations. In Henkel's \emph{Book of Concord}, 1854, and a better one by Dr. Charles P. Krauth:  \emph{The Augsburg Confession, literally translated from the original Latin, with the most important Additions of the German Text incorporated, together with Introduction and Notes.} Philadelphia, 1869. The same, revised for this work, Vol. II. pp. 1 sqq.

IV. Historical and Critical documents and works on the Augsburg Confession:

Philippi Melanthonis  \emph{Opera} in the second and twenty-sixth volumes of the \emph{Corpus Reformatorum}, ed. Bretschneider and Bindseil. Vol. II. (Halis Saxonum, 1835) contains the Epistles of Melanchthon from Jan. 1, 1530, to Dec. 25, 1535; Vol. XXVI. (Brunsv. 1858, pp. 776), the Augsburg Confession itself, with all the preliminary labors and important documents connected therewith.

Luther's  \emph{Briefe, } in De Wette's ed., Vol. IV. pp. 1–180.

E. Sal. Cyprian:  \emph{Historia der Augsburgischen Confession}, etc. Gotha, 1730, 4to.

Christ. Aug. Salig: \emph{ Vollständige Historie der Augsburg. Confession und derselben Apologie}, etc. 3 Thle. Halle, 1730–35, 4to.

G. G. Weber:  \emph{Kritische Geschichte der Augsb. Conf. aus archivalischen Nachrichten.} Frank. a. M. 1783–84, 2 vols.

K. Pfaff:  \emph{Geschichte des Reichstags zu Augsburg, im Jahr 1530, und des Augsb. Glaubensbekenntnisses bis auf die neueren Zeiten}, Stuttgart, 1830, 8vo; 2 Parts.

Carl Eduard Förstemann:  \emph{Urkundenbuch zur Geschichte des Reichstags zu Augsburg, im Jahr} 1530, etc., 2 vols. Halle, 1833–35, 8vo.

C. Ed. Förstemann:  \emph{Neues Urkundenb. zur Gesch. der ev. Kirchen-Reform.} Hamb. 1842, Vol. I. pp. 357–380. \emph{Die Apologie der Augsburg. Confession in ihrem ersten Entwurfe.}

A. G. Rudelbach:  \emph{Die Augsb. Conf. aus und nach den Quellen}, etc. Leipzig, 1829. \emph{Histor. critische Einleit. in die Augsb. Conf.}, etc. Dresden, 1841.

J. E. Calinich:  \emph{Luther und die Augsb. Confession} (\emph{gekrönte Preisschrift}). Leipz. 1861.

G. Plitt:  \emph{Einleitung in die Augustana.} Erlangen, 1867–68, 2 Parts.

O. Zöckler:  \emph{Die Augsburgische Confession als Lehrgrundlage der deutschen Reformationskirche historisch und exegetisch untersucht.} Frankfurt a. M. 1870.

Comp. also  Ranke:  \emph{Deutsche Geschichte im Zeitalter der Reformation}, III. pp. 186 sqq. (3d ed. 1852), and the relevant sections in Marheineke,  Merle D’Aubigné,  Hagenbach, and Fisher, on the \emph{History of the Reformation.}

See lists of Literature especially in Köllner,  \emph{Symb.} I. pp. 150 sqq., 345 sqq.; also  J. T. Müller,  \emph{Die Symb. Bücher der evang. luth. Kirche}, XVII.;  C. P. Krauth,  \emph{Select Analytical Bibliography of the Augsb. Conf.} (Phila. 1858); and  Zöckler,  \emph{Die Augsb. Conf.} pp. 1, 8, 15, 21, 31, 35, 44, 52, 61, 74, 85–88; and \emph{Corp. Ref.} Vol. XXVI. pp. 102 sqq.

ORIGIN AND HISTORY.

The Augsburg Confession, at first modestly called an \emph{Apology}, after the manner of the early Church in the ages of persecution, was occasioned by the German Emperor Charles V., who commanded the Lutheran Princes to present, at the Diet to be held in the Bavarian city of Augsburg, an explicit statement of their faith, that the religious controversy might be settled, and Catholics and Protestants be united in a war against the common enemies, the Turks.\footnote{The imperial letter convening the diet, dated Bologna, Jan. 21, 1530, was purchased by J. P. Morgan, 1911, for \$25,000 and presented to William II., who, in turn, decorated Mr. Morgan with the order of the Black Eagle.—Ed.} Its deeper cause must be sought in the inner necessity and impulse to confess and formularize the evangelical faith, which had been already attempted before. It was prepared, on the basis of previous drafts, and with conscientious care, by Philip Melanchthon, at the request and in the name of the Lutheran States, during the months of April, May, and June, 1530, at Coburg and Augsburg, with the full approval of Luther. It was signed, August 23, by seven German Princes (the Elector John of Saxony and the Landgrave Philip of Hesse, etc.) and the deputies of two free cities (Nuremberg and Reutlingen). This act required no little moral courage, in view of the immense political and ecclesiastical power of the Roman Church at that time. When warned by Melanchthon of the possible effects of his signature, the Elector John of Saxony nobly replied: 'I will do what is right, unconcerned about my electoral dignity; I will confess my Lord, whose cross I esteem more highly than all the power of the earth.'

On the 25th of June, 1530, the Confession was read aloud, in the German language,\footnote{By Dr. Christian Baier, Vice-Chancellor of the Elector of Saxony, after some introductory remarks of Chancellor Brück, who composed the Preface and the Epilogue; see below. The Emperor at first did not want to have it read at all, but simply presented; yielding this point, he sought to diminish its effect by having it read in Latin, but the Lutheran Princes resisted, and carried their point. 'We are on German soil,' said the Elector John, 'and therefore I hope your Majesty will allow the German language.' He did not allow it, however, to be read in a public session of the Diet in the large City Hall, but merely before a select company of Princes, counselors, and deputies of cities, in the small chapel of the episcopal palace, where he resided.} before the assembled representatives of Church and State, and in the hearing of a monarch in whose dominions the sun never set.

This formed an important epoch in the history of the Reformation. The deputies, and the people who stood outside, listened attentively for two hours to the new creed. The Papists were surprised at its moderation. The Bishop of Augsburg is reported to have said privately that it contained nothing but the pure truth. Duke William of Bavaria censured Dr. Eck for misrepresenting to him the Lutheran opinions; and when the Romish doctor remarked that he could refute them with the Fathers, though not with the Scriptures, the Duke replied, 'I am to understand, then, that the Lutherans are within the Scriptures, and we are on the outside.' The Emperor himself, a bigoted Spaniard, a master in shrewd policy, little acquainted with the German language and spirit, and still less with theology, after respectfully listening for a while, fell asleep during the delivery,\footnote{So Brentius, who was at Augsburg at the time, reports (\emph{cum Confessio legeretur, obdormivit}). Considering the length of the document, this is not inconsistent with the other statement of Jonas and Spalatin, that he, like most of the other Princes, was quite attentive (\emph{satis attentus erat Cæsar}). Nor must his drowsiness be construed as a mark of disrespect to the Lutherans, for he was likewise soundly asleep on the third of August when the Romish Confutation was read before the Diet.} but graciously received the Latin copy for his own use, and handed the German to the Elector of Mayence for safe keeping in the imperial archives, yet prohibited the publication without his permission. Both copies are lost.

The Diet ordered a committee of about twenty Romish theologians, among whom were Eck, Faber, Cochlæus, and Wimpina, to prepare a refutation of the Confession on the spot. Their scholastic \emph{Confutatio}, the result of five successive drafts, was a far inferior production, and made little impression upon the Diet, but it fairly expressed the views of the Emperor and the majority of the States, and was accepted as a satisfactory refutation of the Confession.\footnote{\par The best text, Latin and German, of the \emph{Confutatio Confessionis Augustanæ}, with ample \emph{Prolegomena} and the Summary of Cochlæus, see in the 27th volume of the \emph{Corpus Reformatorum} (1859), pp. 1–243.} Melanchthon answered it by his 'Apology of the Augsburg Confession,' but the Diet refused even to receive the reply; and, after several useless conferences, resolved, Sept. 22 and Nov. 19, 1530, to proceed with violent measures against the Protestants if they should not return to the Catholic faith before the 15th of April of the following year.

The Elector John, justly styled the \emph{Constant}, with all his loyalty to the Emperor and wish for the peace of Germany, refused to compromise his conscience, and, in full view of the possible ruin of his earthly interest, he resolved to stand by 'the imperishable Word of God.'\footnote{See the masterly delineation of this Prince by Ranke, in his \emph{Deutsche Geschichte}, etc., Book V. Ch. 9 (Vol. III. pp. 211 sqq.).} The heroic spirit of the Reformers in these trying times found its noblest expression in the words and tune of Luther's immortal battle-song, based on Psalm. xlvi.:

'A tower of strength our God is still,

A mighty shield and weapon;

He'll help us clear from all the ill

That hath us now o'ertaken.

. . . . . . . . .

'And though they take our life—

Goods, honor, children, wife—

Yet is their profit small;

These things shall vanish all—

The City of God remaineth.'

LUTHER'S SHARE IN THE COMPOSITION.\footnote{Comp. Rückert: \emph{Luther's Verhältniss zum Augsb. Bek.}, Jena, 1854; Calinich:  \emph{Luther und die Augsb. Conf.}, Leipz. 1861 (against Rückert and Heppe); Heppe:  \emph{Entstehung and Fortbildung des Lutherthums}, Cassel, 1863, pp. 234 sqq.; Knaake:  \emph{Luther's Antheil an der Augsb. Conf.}, Berl. 1863; Ratz:  \emph{Was hat Luther durch Melanchthon gewonnen?} in the \emph{Zeitschrift f. hist. Theol.}, Leipz. 1870, No. III.; Zöckler:  l.c. pp. 8 sqq.}

Being under the papal excommunication and the imperial ban since the Diet of Worms (1521), Luther could not safely venture to Augsburg, but he closely watched the proceedings of the Diet from the Castle of Coburg on the Saxon frontier, praying, translating the prophets, writing childlike letters to his children, and manly letters to princes, singing \emph{Ein feste Burg ist unser Gott}, giving his advice at every important step, and encouraging his timid and desponding friend Melanchthon.

He had taken the leading part in the important preparatory labors, namely, the \emph{Fifteen, Articles of the Marburg Conference} (Oct. 3, 1529),\footnote{The German autograph of the Marburg Articles, in the handwriting of the Reformers, was discovered in the archives of Cassel and published by Prof. H. Heppe, of Marburg, Cassel, 1847, and also by Bindseil, in the \emph{Corpus Reform.} Vol. XXVI. pp. 122–127 (in German), with the textual variations. The Articles are signed by Luther, Jonas, Melanchthon, Osiander, Agricola, and Brentius, on the part of the Lutherans, and by Œcolampadius, Zwingli, Bucer, and Hedio on the part of the Reformed. Fourteen of them were fully approved by Zwingli and his friends, and in the 15th, which treats of the Lord's Supper, they agree to disagree as to the mode of Christ's presence.} the \emph{Seventeen Articles of Schwabach} (Oct. 16, 1529),\footnote{The \emph{Articuli XVII. Suobacences} (which must not be confounded with the Twenty-two Articles of a previous convent at Schwabach, near Nuremberg. A.D. 1528, see \emph{Corp. Ref.} Vol. XXVI. pp. 132 sqq.) were composed by Luther, with the aid of Melanchthon, Jonas, Osiander, Brentius, and Agricola. They are only a Lutheran revision and enlargement of the Marburg Articles, and seem to have been drawn up in that town, and then presented before a convent of Lutheran princes and delegates at Schwabach, Oct. 16, and again before a similar convent at Smalcald, Nov. 29. They were first published in February or March, 1530, without the knowledge of Luther, under the title: '\emph{Das Bekenntniss Martini Luthers auf den} \emph{angestellten Reichstag zu Augsburg einzulegen, in} 17 \emph{Artikel verfasst};' then by Luther himself, Wittenb. 1530; and again by Frick, in his edition of Seckendorf's \emph{Ausführl. Historie vom Lutherthum.} See \emph{Corp. Ref.} Vol. XXVI. pp. 129–160.} which correspond to the first or positive part of the Augsburg Confession, and the so-called \emph{Articles of Torgau} (March 20, 1530),\footnote{The Torgau Articles (\emph{Articuli Torgavienses}) were formerly often confounded with the Schwabach Articles, till Förstemann first discovered them in the archives at Weimar, and brought them to light, in 1833, in the first volume of his '\emph{Urkundenbuch},' republished in the \emph{Corp. Ref.} Vol. XXVI. pp. 161–200. They were drawn up by Luther, Melanchthon, Jonas, and Bugenhagen, at the command of the Elector of Saxony (then residing at Torgau), for presentation at the approaching Diet of Augsburg, and discuss the controverted articles on the marriage of priests, the communion of both kinds, the mass, the confession, the episcopal jurisdiction, ordination, monastic vows, invocation of saints, faith and works, etc.} which form the basis of its second or polemical part. But in all respects the Confession, especially the second part, is so much enlarged and improved on these previous labors that it may be called a new work.\footnote{Comp. on the historical details of the sources of the Augs. Conf. the \emph{Corpus Reform.}, Vol. XXVI 1858) pp. 113–200; Plitt:  \emph{Einleitung die Augustana} (1867–68), I. pp. 536 sqq., II. pp.3 sqq.; and Zöckler:  \emph{Die Augsb. Conf.} (1870), pp. 8–15.}

Luther thus produced the doctrinal matter of the Confession, while Melanchthon's scholarly and methodical mind freely reproduced and elaborated it into its final shape and form, and his gentle, peaceful, compromising spirit breathed into it a moderate, conservative tone. In other words, Luther was the primary, Melanchthon the secondary author, of the contents, and the sole author of the style and temper of the Confession.\footnote{\par Kahnis, in his \emph{Luther. Dogmatik}, II. p. 424, says: '\emph{Luther war der Meister des Inhalts, Melanchthon der Meister der Form. . . . Mel. war der Mann, welcher mit Objektivität, Feinheit, Klarheit, Milde zu schreiben verstand. Und wie nie hat er diese Gabe in diesem Falle verwerthet.}' Köllner (Vol. I. p. 178), Rückert, and Heppe give all the credit of authorship to Melanchthon. This is true as far as the spirit and the literary composition are concerned; but as to the doctrines, Luther had a right to say, 'The Catechism, the Exposition of the Ten Commandments, and the \emph{Augsburg Confession}, are \emph{mine.}'}

Luther himself was satisfied that his friend was better adapted for the task, and expressed his entire satisfaction with the execution. When the Confession was sent to him from Augsburg for revision, he wrote to the Elector, May 15, 1530: 'I have read the Apology [Confession] of Master Philip; it pleases me very well, and I know of nothing by which I could better it or change it, nor would it be becoming, for I can not move so softly and gently. May Christ our Lord help, that it may bring forth much and great fruit, as we hope and pray. Amen.'\footnote{'\emph{Ich hab M. Philippsen Apologiam überlesen: die gefället mir fast} (i.e., \emph{sehr}) \emph{wohl, und} \emph{weiss nichts daran zu bessern noch ändern, würde sich auch nicht schicken; denn ich so sanft und leise nicht treten kann. Christus unser Herr helfe, dass sie viel and grosse frucht schaffe, wie wir hoffen bitten. Amen.}' (De Wette's ed. of Luther's \emph{Letters}, IV. p. 17; Luther's \emph{Works}, Erlang. ed. Vol. LIV. p. 145).} After the delivery of the Confession, he wrote to Melanchthon, Sept. 15, in an enthusiastic strain: 'You have confessed Christ, you have offered peace, you have obeyed the Emperor, you have endured injuries, you have been drenched in their revilings, you have not returned evil for evil. In brief, you have worthily done God's holy work as becometh saints. Be glad, then, in the Lord, and exult, ye righteous. Long enough have ye been mourning in the world, look up and lift up your heads, for your redemption draweth nigh. I will canonize you as faithful members of Christ, and what greater glory can you desire! Is it a small thing to have yielded Christ faithful service, and shown yourself a member worthy of him?'\footnote{\par '\emph{Christum confessi estis, pacem obtulistis, Cæsari obedistis, injurias tolerastis, blasphemiis saturati estis, nec malum pro malo reddidistis: summa, opus sanctum Dei, ut sanctos decet, digne tractastis. Lætamini etiam aliquando in Domino et exultate, justi: satis diu tristati} (\emph{al. testati}) \emph{estis in mundo: respicite et levate capita vestra, appropinquat redemtio vestra. Ego canonizabo vos, ut fidelia membra Christi, et quid amplius quæritis gloriæ?}' etc. (\emph{Briefe}, IV. p. 165. Comp. also his letter of July 15 to Jonas, Spalatin, Melanchthon, Agricola, ib. IV. p. 96.)}

The only objection which Luther ever raised to the Augsburg Confession was that it was too gentle, and did not denounce the Pope and the doctrine of purgatory.\footnote{In a letter to Justus Jonas, July 21, 1530: '\emph{Satan adhuc vivit, et bene sensit Apologiam vestram Leisetreterin} [the softly stepping Confession] \emph{dissimulasse articulos de purgatorio, de sanctorum cultu, et maxime de Antichristo Papa}' (\emph{Briefe}, IV. p. 110). Melanchthon himself confessed that he wrote the Confession with more leniency than the malice of the Papists deserved. And yet immediately after the delivery, which marks the height of his usefulness, the good man was in an almost desponding state, and was tormented by scruples whether he had not been conservative enough and taken too much liberty with the venerable Catholic Church. He was, moreover, hard pressed by Romish divines and politicians, and was ready to make serious concessions for the sake of unity and peace. Some of his best friends began unjustly to doubt his loyalty to evangelical truth, and Philip of Hesse, one of the signers of tie Confession, wrote to Zwingli, 'Master Philip goes backward like a crab.'}

CONTENTS.

The Augsburg Confession proper (exclusive of Preface and Epilogue) consists of two parts—one positive and dogmatic, the other negative and polemic, or rather apologetic. The first refers chiefly to doctrines, the second to ceremonies and institutions. The order of subjects is not strictly systematic, though considerably improved upon the arrangement of the Schwabach and Torgau Articles. In the manuscript copies and oldest editions the articles are only numbered; the titles were subsequently added.

I. The first part presents, in twenty-one articles—beginning with the Triune God and ending with the worship of saints—a clear, calm, and condensed statement of the doctrines held by the evangelical Lutherans, (1) in common with the Roman Catholics, (2) in common with the Augustinian school, (3) in opposition to Rome, and (4) in distinction from Zwinglians and Anabaptists.\footnote{For other divisions, see Zöckler, l.c. p. 93 sqq.}

(1.) In theology and Christology, i.e., the doctrines of God's unity and trinity (Art. I.), and of Christ's divine-human personality (III.), the Confession strongly reaffirms the ancient Catholic faith as laid down in the œcumenical Creeds, and condemns (\emph{damnamus}) the old and new forms of Unitarianism and Arianism as heresies.

(2.) In anthropology, i.e., in the articles on the fall and original sin (II.), the slavery of the natural will and necessity of divine grace (XVIII.), the cause and nature of sin (XIX.), the Confession is substantially Augustinian, in opposition to the Pelagian and semi-Pelagian heresies. The Donatists are also condemned (VIII.) for denying the objective virtue of the ministry and the Sacraments, which Augustine defended against them.

(3.) The general Protestant views in opposition to Rome appear in the articles on justification by faith (IV.), new obedience (VI.), the Gospel ministry (V.), the Church (VII., VIII.), repentance (XII.), ordination (XIV.), ecclesiastical rites (XV.), civil government (XVI.), good works (XIX.), the worship of saints, and the exclusive mediatorship of Christ (XX.). Prominence is given to the doctrine of justifition by faith, which, though very briefly stated in its proper place (P. I. Art. IV.), is elsewhere incidentally referred to as the essence of the Gospel.\footnote{Part II. Art. 5 (\emph{De discrimine ciborum}): 'Of this persuasion concerning traditions many disadvantages have followed in the Church. For first the doctrine of grace is obscured by it, and the righteousness of faith, which is the principal part of the Gospel (\emph{doctrina de gratia et justitia fidei, quæ est præcipua pars Evangelii}), and which it behoveth most of all to stand forth and to have the pre-eminence in the Church, that the merit of Christ may be well known, and faith, which believeth that sins are remitted for Christ's sake, may be exalted far above works.'}

(4.) The distinctive Lutheran views—mostly retained from prevailing Catholic tradition, and differing in part from those of other Protestant churches—are contained in the articles on the Sacraments (IX., X., XIII.), on confession and absolution (XI.), and the millennium (XVII.). The tenth article plainly asserts the doctrine of a real \emph{bodily} presence and distribution of Christ in the eucharist to \emph{all} communicants (without determining the \emph{mode} of the presence either by way of consubstantiation or transubstantiation),\footnote{The wording of the article—\emph{quod corpus} (in German, \emph{wahrer Leib}) \emph{et sanguis Christi vere} (\emph{wahrhaftiglich}) \emph{adsint et distribuantur vescentibus in Cæna Domini}—leaves room for both theories. The Papistical Confutation, while objecting to the articles \emph{de utraque specie} and \emph{de missa}, in the second part of the Augsb. Conf., was satisfied with Art. X. of the first part, provided only that it be understood as teaching the presence of the \emph{whole} Christ under the \emph{bread} as well as the wine. ('\emph{Decimus articulus in verbis nihil offendit, quia fatentur, in eucharistia post consecrationem legitime factam corpus et sanguis Christi substantialiter et vere adesse, si modo credant, sub qualibet specie integrum Christum adesse.}') In the Apology of the Confession (Art. X.), Melanchthon asserts the \emph{corporalis præsentia}, and even substitutes for \emph{vere adsint} the stronger terms \emph{vere et}  substantialiter  \emph{adsint.} The Lutheran Church, as represented in Luther's writings and in the Form of Concord (R. 729), rejects transubstantiation, and also the doctrine of impanation, i.e., a \emph{local} inclusion of Christ's body and blood in the elements (\emph{localis inclusio in pane}), or a \emph{permanent} and \emph{extra}-sacramental conjunction of the two substances (\emph{durabilis aliqua conjunctio extra usum sacramenti}); but it teaches consubstantiation in the sense of a sacramental conjunction of the two substances effected by the consecration, or a real presence of Christ's very body and blood \emph{in}, \emph{with}, and \emph{under} (\emph{in}, \emph{cum}, et \emph{sub}) bread and wine. The word \emph{consubstantiation}, however, is not found in the Lutheran symbols, and is rejected by Lutheran theologians if used in the sense of \emph{impanation}. The philosophical foundation of this dogma is the \emph{ubiquity} (either absolute or relative) of Christ's body, which is a part of the Lutheran Christology.} and disapproves of dissenting views (especially the Zwinglian, although it is not named).\footnote{\par \emph{Et improbant secus docentes} (\emph{derhalben wird auch die Gegenlehr verworfen}). The omission of Zwingli's name may be due to regard for his friend, the Landgrave Philip of Hesse, but that he was chiefly intended must be inferred from the antecedent controversies, especially the l5th Article of the Marburg Conference, and from the strong opposition of Melanchthon to Zwingli's theory before 1536 or 1540, when he modified his own view on the Eucharist. See below.} The Anabaptists are expressly condemned (\emph{damnamus}), like heretics, for their views on infant baptism and infant salvation (IX.), the Church (VIII.), civil offices (XVI.), the millennium and final restoration (XVII.). These articles, however, have long ceased to be held by all Lutherans. Melanchthon himself materially changed the tenth article in the edition of 1540. The doctrine of the second advent and the millennium (rejected in Art. XVII.) has found able advocates among sound and orthodox Lutheran divines, especially of the school of Bengel.

II. The second part rejects, in seven articles, those abuses of Rome which were deemed most objectionable, and had been actually corrected in the Lutheran churches, namely, the withdrawal of the communion cup from the laity (I.), the celibacy of the clergy (II.), the sacrifice of the mass (III.), obligatory auricular confession (IV.), ceremonial feasts and fasts (V.), monastic vows (VI.), and the secular power of the bishops, as far as it interferes with the purity and spirituality of the Church (VII.).

The style of the Latin edition is such as may be expected from the classic culture and good taste of Melanchthon, while the order and arrangement might be considerably improved.

The diplomatic Preface to the Emperor is not from his pen, but from that of the Saxon Chancellor Brück.\footnote{Förstemann, \emph{Urkundenbuch}, etc., I. p. 460, and Bindseil, \emph{Corp. Ref.}, Vol. XXVI. p. 205. Chancellor Brück (Pontanus) wrote the Preface in German, and Jonas translated it into Latin. A copy in the Seminary Library at Wittenberg has the remark, probably from the hand of Jonas, after the inscription, '\emph{Præfatio ad Cæs.} Car. V.:' '\emph{Reddita e Germanico Pontani tunc per Justum Jonam.}} It is clumsy, tortuous, dragging, extremely obsequious, and has no other merit than to introduce the reader into the historical situation. The brief conclusion (\emph{Epilogus}) is from the same source, and is followed by the signatures of seven Princes and two magistrates.\footnote{There was considerable controversy as to the genuineness of the signatures of two of seven Princes, viz., John Frederick of Saxony (the son of the Elector John) and Duke Francis of Lüneburg. See Köllner, l.c. pp. 201 sqq.} Several manuscript copies omit both Preface and Epilogue, as not belonging properly to the Confession.

CHARACTER AND VALUE.

The Augsburg Confession breathes throughout an earnest and devout evangelical Christian spirit, and is expressed in clear, mild, dignified language. It professes to be both Scriptural and churchly, and in harmony even with the Roman Church as known from the genuine tradition of antiquity.\footnote{At the conclusion of the first part, the Confession says: '\emph{Hæc fere summa est doctrinæ apud nos, in qua cerni potest, nihil inesse, quod}  discrepet a scripturis, vel ab ecclesia catholica, vel ab ecclesia romana, quatenus ex scriptoribus nota est,' and in the Epilogus: '\emph{Apud nos nihil esse receptum}  contra scripturam, aut ecclesiam catholicam,  \emph{quia manifestum est, nos diligentissime cavisse, ne qua}  nova et impia dogmata  \emph{in ecclesias nostras serperent.}' Hence the Confession frequently appeals not only to the Scriptures, but also to the Fathers (Augustine, Ambrose, Chrysostom, etc.) and the canon law (\emph{Decretum Gratiani, veteres canones}, and the \emph{exemplum ecclesiæ}).} It is remarkably moderate and conciliatory in tone, and free from all harsh or abusive terms. It is not aggressive, but defensive throughout. Hence its original modest name \emph{Apology.}\footnote{Melanchthon wrote to Luther: '\emph{Mittitur tibi Apologia nostra, quanquam verius Confessio est.}' Afterwards it was also frequently called the 'Saxon Confession' and the '\emph{Evangelische Augapfel}' (Prov. vii. 2).} It pleads only for toleration and peace. It condemns the ancient heresies (Arianism, Manicheism, Pelagianism, Donatism), which were punishable according to the laws of the German Empire. It leaves the door open for a possible reconciliation with Rome.\footnote{Ranke, l.c. III. p. 201: '\emph{In diesem Sinne der Annäherung, dem Gefühle des Nochnichtvollkommengetrenntseins, dem Wunsche, eine wie im tieferen Grunde der Dinge waltende, so in einigen Einzelnheiten des Bekenntnisses sichtbare Verwandtschaft geltend zu machen, war die Confession gedacht und abgefasst.}' Zöckler, l.c. p. 318: '\emph{Die Augustana ist in ihren Antithesen, sowohl nach der römischen wie nach der reformirten Seite hin, das mildeste, friedliebendste, gegnerischer seits am leichtesten zu ertragende aller evangelisch-lutherischen Symbole.}} Popery itself, and many of its worst abuses, are not even touched, at least not expressly. The modest and peaceful author wrote under a painful sense of responsibility, with a strong desire for the restoration of the unity of faith, and hence he avoided, all that might give unnecessary offense to the ruling party.\footnote{Comp. the Preface, and the repeated assurances of Melanchthon, e.g., in a letter of May 21, 1530, to Joachim Camerarius (\emph{Corp. Ref.} II. p. 57): '\emph{Ego Apologiam paravi scriptam summa verecundia, neque his de rebus dici mitius posse arbitror.}' And in a letter to the same, dated June 19 (ib. p. 119): '\emph{Non dubitabam quin Apologia nostra videretur futura lenior, quam mereatur improbitas adversarioram.}'}

But the same motive made him unjust toward his fellow-Protestants, who differed from him far less than both differed from the Romanists. The Lutheran divines, after refusing at Marburg all connection with the Zwinglians, yet, being unable to convince the Catholic majority, felt that by protesting against what they regarded as ultra-Protestant radicalism they would better succeed in securing toleration for themselves. One of their leaders, however, Philip of Hesse, openly sympathized with Zwingli, and had to be specially urged by Luther to subscribe the Confession, which he did with a dissent from the tenth article. The majority of the citizens of Augsburg likewise adhered to Zwingli at that time.\footnote{See the remarks of L. Ranke, III. p. 220 sq. Kahnis also (\emph{Luth. Dogm.} II. p. 436) admits that 'the desire for an understanding with the Papists made Melanchthon a very decided opponent of the Swiss, and even of the Strasburgers.'}

The Augsburg Confession is the fundamental and generally received symbol of the Lutheran Church, which also bears the name of 'the Church of the Augsburg Confession.' It is inseparable from the theology and history of that denomination; it best exhibits the prevailing genius of the German Reformation, and will ever be cherished as one of the noblest monuments of faith from the pentecostal period of Protestantism.\footnote{For a hearty estimate of the value of the Confession from the Lutheran stand-point, see Dr. Krauth's introduction to his translation, pp. xlvii. sqq., and his \emph{Conservative Reformation}, pp. 255 sqq.: 'With the Augsburg Confession,' he says in both places, 'begins the clearly recognized life of the Evangelical Protestant Church, the purified Church of the West, on which her enemies fixed the name \emph{Lutheran}. With this Confession her most self-sacrificing struggles and greatest achievements are connected. It is hallowed by the prayers of Luther, among the most ardent that ever burst from the human heart; it is made sacred by the tears of Melanchthon, among the tenderest which ever fell from the eyes of man. It is embalmed in the living, dying, and undying devotion of the long line of the heroes of our faith, who, through the world which was not worthy of them, passed to their eternal rest. The greatest masters in the realm of intellect have defended it with their labors; the greatest Princes have protected it from the sword by the sword; and the blood of its martyrs, speaking better things than vengeance, pleads forever, with the blood of Him whose all-availing love, whose sole and all-atoning sacrifice, is the beginning, middle, and end of its witness.'}

But its influence extends far beyond the Lutheran Church. It struck the key-note to other evangelical confessions, and strengthened the cause of the Reformation every where. It is, to a certain extent, also the Confession of the Reformed and the so-called Union Churches, in Germany, namely, with the explanations and modifications of the author himself in the edition of 1540, which extended, as it were, the hand of fellowship to them (see below). In this qualified sense, either expressed or understood, the Augsburg Confession was frequently signed by Reformed divines and Princes, even by John Calvin, while ministering to the Church at Strasburg, and as delegate to the Conference of Ratisbon, 1541;\footnote{Calvin wrote to Rev. Mart. Schalling, at Ratisbon, 1557: '\emph{Nec vero Augustanam Confessionem repudio, cui pridem volens ac libens subscript, sicut eam auctor ipse interpretatus est}' (\emph{Epp.} p. 437). Similarly in his \emph{Ultima Admonitio ad Joach. Westphalum}, Genev. 1557. It is not quite certain whether it was the Altered or the Unaltered Confession which Calvin subscribed at Ratisbon, but probably it was the former, as he says that it contained nothing contrary to his doctrine, and as he appealed without fear to Melanchthon himself as the best interpreter. The Altered edition had appeared a year before, and had been actually used at the previous Conference at Worms, though Eck protested against it. See Köllner, p. 241; Zöckler, pp. 40, 41; Ebrard, \emph{Dogma vom hell. Abendmahl}, II. p. 450; Stähelin, \emph{Joh. Calvin}, I. p. 236; G. v. Polentz, \emph{Geschichte des französischen Calvinismus}, Vol. I. p. 577; Vol. II. p. 62.} by Farel and Beza at the Conference in Worms, 1557; by the Calvinists at Bremen, 1562; by Frederick III., (the Reformed) Elector of the Palatinate, at the convent of Princes in Naumburg, 1561, and again at the Diet of Augsburg, 1566; by John Sigismund, of Brandenburg, in 1614. It is true that till the close of the Thirty-Years' War (1648) the Reformed were tolerated in the German Empire only as allies of the Augsburg Confession,\footnote{'\emph{Augustanæ Confessioni addicti},' '\emph{Augsburgische Confessionsverwandte.}'} but even afterwards they continued their friendly relation to it, and maintain it to the present day without feeling any more bound by it.\footnote{In the electoral, afterwards royal, house of Brandenburg, the Augsburg Confession and the Heidelberg Catechism have always lived in peace together. The Great Elector, Frederick William, as patron of the German Reformed, professed in their name, when the Westphalian Treaty was concluded, their cordial adherence to the Confession of 1530 (\emph{Profitentur dicti Reformati Augustanam Confessionem augustissimo Imp. Carolo V. anno} 1530 \emph{exhibitam corde et ore}). There are, however, German Reformed congregations of a more strictly Calvinistic type (e.g., in Elberfeld), which would rather adopt the Canons of the Synod of Dort than the Augsburg Confession.}

The last, and the most memorable occasion since 1530, on which this noble Confession was publicly acknowledged, but with a saving clause as to the interpretation of the tenth article relating to the doctrine of the Lord's Supper, was at the German Church Diet of Berlin, 1853, composed of over 1400 clergymen, of four denominations—Lutheran, German Reformed, Evangelical Unionists, and Moravians.\footnote{The unanimous declaration of the Berlin Church Diet reads thus: 'The members of the German Evangelical Church Diet hereby put on record that they hold and profess with heart and mouth the Confession delivered, A.D. 1530, at the Diet of Augsburg, by the evangelical Princes and States to Emperor Charles V., and hereby publicly testify their agreement with it, as the oldest, simplest common document of publicly recognized evangelical doctrine in Germany (\emph{dass sie sich zu der im Jahr} 1530 \emph{auf dem Reichstags zu Augsburg von den evangelischen Fürsten und Ständen Kaiser Karl V. überreichten Confession mit Herz und Mund halten und bekennen, und die Uebereinstimmung mit ihr, als der ältesten, einfachsten gemeinsamen Urkunde öffentlich anerkannter evangelischer Lehre in Deutschland, hiedurch öffentlich bezeugen}).' So far orthodox Lutherans might agree. But now follows a qualification to save the consciences of the Reformed and Unionists: 'With this we connect the declaration that they and each one of them adhere to the particular confessions of their respective churches, and the Unionists to the \emph{consensus} of the same; and that they do not mean to interfere with the different positions which the Lutherans, Reformed, and Unionists sustain to the Tenth Article of the Augsburg Confession, nor with the peculiar relations of those Reformed congregations which never held the Augustana as a symbol (\emph{Hiemit verbinden sie die Erklärung, dass sie jeder insonderheit an den besonderen Bekenntniss-Schriften ihrer Kirchen, und die Unirten an dem Consensus derselben festhalten, und dass der verschiedenen Stellung der Lutheraner, Reformirten und Unirten zu Artikel X. dieser Confession, und den eigenthümlichen Verhältnissen derjenigen Reformirten Gemeinden, welche die Augustana niemals als Symbol gehabt haben, nicht Eintrag geschehen soll}).' See \emph{Evang. Kirchenztg.} of Berlin, for 1853, pp. 775 sqq. While fully recognizing the importance of this testimony in opposition to rationalism and popery, we should remember, first, that it has no official or ecclesiastical character (the German \emph{Kirchentag}, like the Evangelical Alliance, being merely a voluntary association without legislative or disciplinary power); and, secondly, that it is a compromise, which was expressly repudiated by the anti-Union Lutherans (the professors at Erlangen, Leipzig, and Rostock), as 'a frivolous depreciation of the most precious symbol of German Evangelical Christendom.'}

On this fact and the whole history of the Augsburg Confession some German writers of the evangelical Unionist school have based the hope that the Augsburg Confession may one day become the united Confession or œcumenical Creed of all the evangelical churches of Germany.\footnote{So Dr. W. Hoffmann, late Court Chaplain of the Emperor of Germany (\emph{Deutschland Einst und Jetzt im Lichte des Reiches Gottes}, Berlin, 1868, pp. 476 sqq. and 512 sqq.); Consistorialrath Leop. Schultze (\emph{Die Augsb. Confession als Gesammtlbekenntniss unserer evang. Landeskirche}, Bremen, 1869); to some extent also Prof. Zöckler (l.c. p. 330), who proposes that the Augsburg Confession be made, not indeed the Union Symbol, but the Confederation Symbol of German Evangelical Christendom.} This scheme stands and falls with the dream of a united and national Protestant Church of the German Empire. Aside from other difficulties, the Reformed and the majority of Unionists, together with a considerable body of Lutherans, can never conscientiously subscribe to the tenth article as it stands in the proper historical Confession of 1530; while orthodox Lutherans, on the other hand, will repudiate the Altered edition of 1540. The \emph{Invariata} is, after all, a purely Lutheran, that is, a denominational symbol; and the \emph{Variata} is a friendly approach of Lutheranism towards the Reformed communion, which had no share in its original production and subsequent modification, although it responded to it. Neither the one nor the other edition can be the expression of a union, or confederation of two distinct denominations, of which each has its own genius, history, and symbols of faith. Such an expression must proceed from the theological and religious life of both, and meet the wants of the present age. Great as the Augsburg Confession is, the Church will produce something greater still whenever the Spirit of God moves it to a new act of faith in opposition to the unbelief and misbelief of modern times. Every age must do its own work in its own way.

THE TEXT. THE INVARIATA AND THE VARIATA.\footnote{See the details in Weber, Köllner, and Bindseil.}

The Augsburg Confession was first completed in Latin,\footnote{\emph{Corp. Reform.} Vol. XXVI. p. 205.} but the German text was read before the Diet. Both copies were delivered in manuscript to the Emperor, but both disappeared: the German copy, first deposited in the imperial archives at Mayence, was probably sent with other official documents to the Council of Trent (1545), and then to Rome; the Latin copy to Brussels or Spain, and no trace of either has been found. For two hundred years the opinion prevailed that the 'Book of Concord' contained the original text, until Pfaff and Weber, by a thorough investigation on the spot, dispelled this error.\footnote{The Latin text of the Book of Concord is substantially from Melanchthon's quarto edition of 1531, and was supposed to correspond entirely with an imaginary Latin manuscript in Mayence. The German text purports to be a true copy of the original manuscript in Mayence, but is derived from a secondary source, viz., the printed text in the \emph{Corpus Brandenburgicum}, 1572, which, again, was based upon a carelessly written copy of the Confession \emph{before} its final revision. Chancellor Pfaff, of Tübingen, first discovered at Mayence that the original German copy was lost long ago, and he published, in 1730, what was regarded as a true copy of the original; but he was fiercely assailed by Adami, Feuerlin, and others, and his discovery traced to a Jesuitical lie. In 1781 Georg Gottlieb Weber, chief pastor at Weimar, was allowed to make a thorough search in the archives of Mayence, and found to his surprise that the copy shown him as the original was the printed German octavo edition of 1540, bearing on the title-page the words 'Wittenberg, M.D.X.L.' He published the results of his patient investigation in his \emph{Kritische Geschichte der Augsb. Confession aus archival. Nachrichten}, Frankf. a. M. 1783–4, 2 vols.}

The twenty-two manuscript copies, still extant in public or private libraries, are inaccurate, defective, and represent the various stages of revision through which the Confession passed before the 25th of August under the ever-improving hands of the author. There was no time, it seems, to make authentic copies of the final revision.

The printed editions (six in German, one in Latin), which were hastily issued during the Diet by irresponsible, anonymous publishers, are full of errors and omissions, and were condemned by Melanchthon.

Consequently, we must depend entirely upon the author's own printed editions; but even these differ very much among themselves, and the German text differs from the English.\footnote{The various readings in Bindseil's edition, in the \emph{Corpus Reformatorum}, cover as much space as the text itself.} Fortunately the changes are mostly verbal and immaterial, and where they alter the sense (as in the edition of 1540), they can be traced to their proper origin.

By the subscription of the Lutheran Princes and the delivery at the Diet, the Confession had become public property, and should have remained unaltered. But at that time neither editors nor publishers were careful and scrupulous in handling official documents. Luther himself changed the Articles of Smalcald after they had been publicly acknowledged. Melanchthon regarded the Confession, not as a fixed and binding creed, but as a basis for negotiation with the Papists, and as representing a movement still in progress toward clearer light.\footnote{Comp. the concluding words: '\emph{Si quid in hoc confessione desiderabitur, parati sumus latiorem informationem, Deo volente, juxta Scripturas exhibere.}'} He therefore kept improving it, openly and honestly, in every new issue, as he would his own work, and in the edition of 1540 he even embodied some doctrinal modifications in the desire of promoting the cause of truth and peace.

The \emph{editio princeps} was issued by the author in both languages, together with the Latin Apology and a German translation of it by Jonas, at Wittenberg in 1531, in spite of the imperial prohibition, yet with the consent (though not by order) of the Elector of Saxony.\footnote{Under the title: 'Confessio Fidei  | \emph{exhibita invictiss. Imp. Carolo V.} | \emph{Cæsaris Aug. in Comiciis} | \emph{Augustæ}, | \emph{Anno} | \emph{M.D.X.X.X.} | \emph{Addita est Apologia Confessionis.} | Beide, Deutsch | und Latinisch. | Ps. 119. | \emph{Et loquebar de testimoniis tuis in conspectu Regum, et non confundebar.} |  Witebergæ.' (In 4). At the end: '\emph{Impressum per Georgiam Rhau.} | \emph{M.D.X.X.X.I.}' This is the title of the copy in the royal library at Dresden, which Melanchthon gave to Luther, with the words, in his own handwriting (below the title): '\emph{D. Doctari Martino. Et rogo ut legat et emendet.}' See \emph{Corp. Ref.} Vol. XXVI. p. 235. Bindseil (pp. 246 sqq.) shows that the Confession was already printed (but not issued) in November, 1530, and that the whole volume, with the Apology, was finished in April or May, 1531. Some copies of the printed Confession seem to have reached Augsburg before the close of the Diet.} The text was taken, not from Melanchthon's own manuscript copy (which had been delivered to the Emperor), but, as he says, \emph{ex exemplari bonæ fidei} (probably the private copy of the Landgrave Philip of Hesse), and contained already verbal changes and improvements.\footnote{He wrote to Joachim Camerarius, June 26 (a day after the delivery at Augsburg): '\emph{Ego mutabam et refingebam pleraque quotidie, plura etiam mulaturus, si nostri} \textgreek{συμφράδμονες} [counselors] \emph{permisissent.}' \emph{Corp. Ref.} II. p. 140. Kaiser has shown that Melanchthon made a number of changes in the first edition—\emph{Beitrag zu einer Kritischen Literär-Geschichte der Melanchthonischen Original-Ausgabe der lat. und deutsch. Augsb. Conf. und Apologie}, Nürnberg, 1830. Comp. Köllner, l.c. I. p. 340, and \emph{Corp. Ref.} Vol. XXVI. pp. 251 sqq.}

The emendations in subsequent editions before 1540, though quite numerous, do not materially affect the sense, and seem to have been approved; at all events, they were acquiesced in by the Lutherans themselves.\footnote{Luther, who took similar liberty with the Smalcald Articles, expresses no judgment, in his writings, on these variations; but he must have known of them, and tolerated them as unessential, even those of 1540, which appeared six years before his death. The sayings attributed to him on this subject by both parties are apocryphal, at all events unreliable, viz., the word of censure: '\emph{Philippe, Philippe, ihr thut nicht recht, dass ihr Augustanam Confessionem so oft ändert; denn es ist nicht euer, sondern der Kirchen Buch;}' and the word of indirect approval (1546): '\emph{Lieber Philipp, ich muss es bekennen, der Sache vom Abendmahl ist viel zu viel gethan}' (the matter of the Lord's Supper has been much overdone). The latter utterance, however, which Luther is reported to have made shortly before his death, has received a high degree of probability by the discovery of the testimony of Pastor Hardenberg, of Bremen (1547–1550), who publicly and solemnly declared to have heard it, together with another living witness (Canon Herbert von Langen, at Bremen), from \emph{Melanchthon's own lips}. See Erlanger \emph{Reform. Kirchenzeitung} for 1853, No. 40. The first Lutheran divine who publicly censured and condemned the \emph{Variata} was Flacius, at the colloquy of Weimar, 1560. He was followed by Mörlin, Stössel, Wigand, Chytræus, Heshusius, and others.}

But the edition of 1540, which appeared in connection with an improved edition of the Apology,\footnote{Under the title (as given in \emph{Corp. Reform.} l.c. p. 243): 'Confessio | \emph{Fidei exhibita} | \emph{invictiss. imp. Carolo} | \emph{V. Cæsari Aug. in Comiciis} | \emph{Augustæ.} | \emph{Anno. M.D.X.X.X. Addita et Apologia Confessionis diligenter recognita.} | \emph{Psalmo CXIX.} | \emph{Vitebergæ}, 1540.' The words \emph{diligenter recognita} (in the German edition, \emph{mit vleis emendirt}) openly indicate the changes.} differs so widely from the first that it was subsequently called the \emph{Altered} Augsburg Confession (\emph{Variata}), in distinction from the \emph{Unaltered} (\emph{Invariata}) of 1530 or 1531.\footnote{The best text of the \emph{Variata}, with the variations of later editions, is given in \emph{Corp. Reform.} Vol. XXVI. pp. 349 sqq.; the history in Köllner, I. pp. 235–267, and the books there quoted; also in Zöckler, l.c. pp. 35 sqq. In Vol. II. of this Symb. Library the principal changes are noted in foot-notes under the text of the Confession.} It attracted little attention till after the death of Melanchthon (1560), when it created as much trouble as the insertion of the \emph{filioque} clause in the Nicene Creed. The Altered Confession, besides a large number of valuable additions and real improvements in style and the order of subjects,\footnote{In Art. 4, 5, 6, 18, 20, 21, of Part First, and the order of the first five articles in Part Second.} embodies the changes in Melanchthon's theology,\footnote{In Art. 4, 5, 10, 18, 20.} which may be dated from the new edition of his \emph{Loci communes}, 1535, and his personal contact with Bucer and Calvin. He gave up, on the one hand, his views on absolute predestination, and gradually adopted the synergistic theory (which brought him nearer to the Roman Catholic system); while on the other hand (departing further from Romanism and approaching nearer to the Reformed Church), he modified the Lutheran theory of the real presence, at least so far as to allow the Reformed doctrine the same right in the evangelical churches. He never liked the Zwinglian view of a symbolical presence, nor did he openly adopt the Calvinistic view of a spiritual real presence, but he inclined to it, and regarded the difference between this and the Lutheran view as no bar to Christian fellowship and church communion.

Hence in the edition of 1540 he laid greater stress on the necessity of repentance and good works, and softened down the strong expressions against the freedom of will. The other and more important change which gave most offense to orthodox Lutherans, is in the tenth article, concerning the Lord's Supper, where the clause on the real presence, and the disapproval of dissenting views are omitted, and the word \emph{exhibeantur} is substituted for \emph{distribuantur}. In other words, the article is so changed that Calvin could give it his hearty consent, and even Zwingli—with the exception, perhaps, of the word \emph{truly}—might have admitted it.\footnote{Zöckler, l.c. p. 38, thinks that the Calvinistic view would require \emph{credentibus} instead of \emph{vescentibus.} This would be true, if the original \emph{distribuantur} had been retained, and not exchanged for the more indefinite \emph{exhibeantur.} He admits, however, that the tenth article is '\emph{calvinisirend}' and '\emph{bucerianisirend}' in the sense of the Wittenberg Concordia of 1536, whereby Bucer, with Melanchthon's express co-operation, and the subsequent consent of Calvin, endeavored to unite the Lutherans and the Swiss.} The difference will best appear from the following comparison:

Edition 1530. Latin Text.

Edition 1540.

'\emph{De Cæna Domini docent, quod corpus et sanguis Christi } vere adsint,  \emph{et}  distribuantur  \emph{vescentibus in Cæna Domini;}  et improbant secur docentes.'\footnote{The German text of 1530 (1531) differs from the Latin, and is even stronger: '\emph{Vom Abendmahl des Herrn wird also gelehret, dass} wahrer  \emph{Leib} (the \emph{true} body) \emph{und Blut Christi wahrhaftiglich} (corresponding to the \emph{vere} in the Latin text) unter (der) gestalt (under the form) \emph{des Brots und Weins im Abendmahl gegenwärtig sei, und da ausgetheilt und genommen wird} (distributed and received). \emph{Derhalben wird auch die Gegenlehr verworfen.}'}

\emph{'De Cœna Domini docent, quod}  cum pane et vino  \emph{vere}  exhibeantur  \emph{corpus et sanguis Christi vescentibus in Cæna Domini.}'

'Concerning the Lord's Supper, they teach that the body and blood of Christ \emph{are truly present}, and are \emph{distributed} (\emph{communicated}) to those that eat in the Lord's Supper. \emph{And they disapprove of those that teach otherwise}.'

'Concerning the Lord's Supper, they teach that \emph{with bread and wine }are truly\emph{ exhibited }the body and blood of Christ to those that eat in the Lord's Supper.'

The difference between the two editions was first observed, not by Protestants, but by the Roman controversialist, Dr. Eck, at a religious conference in Worms early in the year 1541. Melanchthon and the Saxon theologians made there the altered edition the basis of negotiations, but Eck complained of changes, especially in Art. X., from the original copy of 1530, which he had procured from the archives of Mayence. Nevertheless, the \emph{Variata} was again used, either alone or alongside with the \emph{Invariata}, at several subsequent conferences, probably at Ratisbon, 1541, certainly at Ratisbon in 1546, and at Worms, 1557. It was expressly approved by the Lutheran Princes at a convention in Naumburg, 1561, as an innocent and, in some respects, improved modification and authentic interpretation of the \emph{Invariata.} It was introduced into many Lutheran churches and schools, and printed (with the title and preface of the edition of 1530) in the first collection of Lutheran symbols, called \emph{Corpus Doctrinæ Philippicum}, or \emph{Misnicum} (1559).\footnote{See Weber, l.c. II. pp. 214–336; Köllner, l.c. pp. 248 sqq.}

But after 1560, strict Lutheran divines, such as Flacius and Heshusius, attacked the \emph{Variata}, as heretical and treacherous, and overwhelmed it with coarse abuse. A violent theological war was waged against Melanchthonianism and Crypto-Calvinism, and ended in the triumph of genuine Lutheranisrn and the transition of some Lutheran countries to the Reformed Church. The 'Book of Concord' (1580) gave the text of the \emph{Invariata} in the happy illusion of presenting it, especially the German, in its original form. The Melanchthonians and the Reformed still adhered to the \emph{Variata.} The Westphalian Treaty, in 1648, formally embraced the Reformed, together with the Roman Catholics and Lutherans, in the peace of the German Empire; and henceforth subscription to the Augsburg Confession of 1530 or 1540 ceased to be a necessary condition of toleration.\footnote{\emph{Instrum. Pacis Osnabr.} Art. VII. § 1: '\emph{Unanimi quoque . . . consensu placuit, ut quicquid publica hæc transactio, in eaque decisio gravaminum ceteris Calholicis, et}  Augustanæ Conf. Addictis  \emph{statibus et subditis tribuunt, it etiam iis, qui inter illos}  Reformati  \emph{vocantur, competere debeat.}' Quoted by Jacobson in art. \emph{Westf. Friede}, in Herzog's \emph{Real-Encycl.} XVIII. p. 24. Nevertheless, some interpreted this decree as including only such of the Reformed as subscribed the \emph{Invariata.} All other Christians are expressly excluded by the Treaty; and yet the Popes have always, though vainly, protested in the strongest terms (\emph{damnamus, reprobamus, cassamus, annullamus, vacuamus}) even against this partial concession to the principle of religious freedom; taking the ground that Papists alone have a legal right to exist on German soil. See Gieseler, \emph{Lehrbuch der K. G.} III. I. p. 431 sq.}

The Confession, as delivered before the Diet of Augsburg in 1530, or, in the absence of the original text, the edition of 1531, carefully prepared by Melanchthon himself, is the proper historical Confession of Augsburg, and will always remain so. At the same time, the altered edition of 1540, though not strictly speaking a symbolical book of binding authority any where,\footnote{An attempt was made in the Bavarian Palatinate, in 1853, through the influence of Dr. Ebrard, to raise the \emph{Variata} to the dignity of a symbolical book, but it proved abortive.} is yet far more than a private document, and represents an important element in the public history of the Lutheran Church in the sixteenth century, and the present theological convictions of a very large party in that denomination.

\xsection{The Apology of the Augsburg Confession. A.D. 1530–1531.}

§ 42. The Apology of the Augsburg Confession, A.D. 1530–1531.

The Literature in §§ 40 and 41. The history and literature of the Apology are usually combined with that of the Confession. So in J. G. Walch, Feuerlin-Riederer, Köllner, etc.

The best text of the \emph{Apology}, and of the Roman Catholic \emph{Confutation} of the Confession, in Latin and German, with all the variations, is given in the \emph{Corpus Reformatorum}, Vol. XXVII., ed. Bindseil (Brunsvigæ, 1859), pp. 646, fol. There are few separate editions of the Apology. Feuerlin knew only two, one under the singular title: \emph{Evangelischen Augapfels} (name of the Augsb. Conf.) \emph{Brillen-Butzer}, Leipz. 1629.

The 'Apology of the Augsburg Confession' was prepared by Melanchthon in vindication of the Confession against the Roman Catholic 'Confutation,' which the Emperor and the Diet had ordered and accepted, August 3, 1530, as a satisfactory answer, although, in the eyes of scholars, it did the cause of popery more harm than good.

The Confutation follows the order of the Augsburg Confession, approves eighteen articles of the first part, either in full or with sundry restrictions and qualifications, but rejects entirely the articles on the Church (VII.), on faith and good works (XX.), and on the worship of saints (XXI.), and the whole second part; nevertheless, it acknowledges at the close the existence of various abuses, especially among the clergy, and promises a reformation of discipline. The publication of the document was prohibited, and it did not appear till many years afterwards; but its main contents were known from manuscript copies, and through those who heard it read.\footnote{The Latin text of the \emph{Confutatio} was first published by Fabricius Leodius in \emph{Harmonia Confess.}, 1573; the German, by C. G. Müller, 1808, from a copy of the original in the archives of Mayence, which Weber had previously obtained. Both in the \emph{Corp. Reform.} l.c. Comp. also above, p. 226; Weber's \emph{Krit. Gesch. der A. C.} II. pp. 439 sqq.; and Hugo Lämmer (who afterwards joined the Romish Church): \emph{Die vor-Tridentinisch-Katholische Theologie, des Reformations-Zeitalters}, Berlin, 1858, pp. 33–46.}

The Lutherans urged Melanchthon to prepare at once a Protestant refutation of the Romish refutation, and offered the first draft of it to the Diet, Sept. 22, through Chancellor Brück, but it was refused. On the following day Melanchthon left Augsburg in company with the Elector of Saxony, and re-wrote the Apology on the journey,\footnote{His zeal led him to violate even the law of rest on Sunday when at Altenburg, in Spalatin's house. Luther took the pen from him, and told him to serve God on that day by resting from literary labor. So Salig reports in his \emph{Hist. of the Augsb. Conf.} I. p. 375.} and completed it at Wittenberg in April, 1531.

The Apology is a triumphant vindication of the Confession. It far excels the Confutation both in theological and literary merit, and in Christian tone and spirit. It is written with solid learning, clearness, and moderation, though not without errors in exegesis and patristic quotations. It is seven times as large as the Confession itself. It is the most learned of the Lutheran symbols. It greatly strengthened the confidence of scholars in the cause of Protestantism. Its chief and permanent value consists in its being the oldest and most authentic interpretation of the Augsburg Confession by the author himself.

The Apology, though not signed by the Lutheran Princes at Augsburg, was recognized first in 1532, at a convent in Schweinfurt, as a public confession; it was signed by Lutheran divines at Smalcald, 1537; it was used at the religious conference at Worms, 1540, and embodied in the various symbolical collections, and at last in the Book of Concord.

The text of the Apology has, like that of the Confession, gone through various transformations. The original draft made at Augsburg has no authority.\footnote{Manuscript copies of this '\emph{Apologia prior},' which was based on an imperfect knowledge of the Romish \emph{Confutatio}, still exist. The Latin text of it was published forty-seven years afterwards by Chytræus (from Spalatin's copy), 1578, better by Förstemann, in his \emph{Neues Urkundenbuch} (1842), pp. 357–380 (from a copy written partly by Spalatin and partly by Melanchthon). The best edition is by Bindseil, in the \emph{Corp. Reform.} Vol. XXVII. pp. 275 sqq. in Latin, and in German, pp. 322 sqq.} The first Latin edition was much enlarged and improved, and appeared in April, 1531, at Wittenberg, together with a very free German translation by Justus Jonas, assisted by Melanchthon.\footnote{During the preparation of the \emph{editio princeps} he wrote to Brentius (February, 1531): '\emph{Ego retexo Apologiam et edetur longe auctior et melius munita},' and to Camerarius (March 7): '\emph{Apologia mea nondum absoluta est, crescit enim opus inter scribendum.}' Quoted by Köllner, I, p. 426. Six sheets were reprinted, and a copy of the first print is preserved in the library of Nuremberg. See \emph{Corp. Reform.} Vol. XXVII. pp. 391 sqq.} The second Latin edition of the same year was again much changed, and is called the \emph{Variata}.\footnote{See the titles of the various editions in \emph{Corp. Reform.} Vol. XXVI. pp. 235–242, and the best text of the '\emph{Apologia altera}' of 1531, with the changes of later editions till 1542 (viz., of the ed. II. 1531, ed. III. 1540, ed. IV. 1542), in \emph{Corp. Reform.} Vol. XXVII. pp. 419–646.} The German text was also transformed, especially in the edition of 1533. The Book of Concord took both texts from the first edition.

\xsection{Luther's Catechisms. A.D. 1529.}

§ 43. Luther's Catechisms. A.D. 1529.

Literature.

I. Editions. See § 40. We only mention the critical editions.

C. Mönckeberg:  \emph{Die erste Ausgabe v. Luthers Klein. Katechismus.} Hamburg, 1851. (Reprint of the Low-German translation of the first edition, 1529.)

K. F. Th. Schneider:  \emph{Dr. Martin Luthers Kleiner Katechismus. Nach den Originalausgaben kritisch bearbeitet. Ein Beitrag zur Geschichte der Katechetik.} Berlin, 1853. (Reprint of the standard edition of 1531; with a critical introduction, pp. lxx.)

Theodos. Harnack:  \emph{Der Kleine Katechismus Dr. Martin Luthers in seiner Urgestalt. Kritisch untersucht und herausgegeben.} Stuttgart, 1856, 4to. (Reprint of two editions of 1529, and one of 1539; with lxiv. pp. of introduction, and a table of the principal variations of the text till 1542.)

The popular editions of the Smaller Catechism, especially in German, with or without comments and supplements, are innumerable.

II. Works:

A. Fabricii:  \emph{Axiomata Scripturæ Catechismo Lutheri accommodata}, etc. Isleb. 1583.

C. Dieterici:  \emph{Instit. catech.} Ulm, 1613; often reprinted.

Ph. J. Spener:  \emph{Tabulæ catech.} Frf. 1683, 1687, 1713.

Greg. Langemack:  \emph{Hist. catecheticæ oder Gesammelte Nachrichten zu einer Catech. Historie.} Strals. 1729–1740, 3 vols. Part II., 1733, treats of \emph{Lutheri und anderer evang. Lehrer Catechismis.}

J. C. Köcher:  \emph{Einleitung in die catech. Theol.} Jena, 1752. And \emph{Biblioth. theol. symb. catech.} P. I. 1751; P. II. 1769.

J. C. W. Augusti:  \emph{Versuch einer hist. kritischen Einleitung in die beiden Haupt-Katechismen der Evang. Kirche.} Elberf. 1824.

G. Veesenmeyer:  \emph{Liter. bibliograph. Nachrichten von einigen evang. katechet. Schriften und Katechismen vor und nach Luthers Kat.}, etc. Ulm, 1830.

G. Mohnike:  \emph{Das sechste Hauptstück im Katechismus.} Stralsund, 1830.

C. A. Gerh. \emph{von } Zezschwitz:  \emph{System der christlich kirchlichen Katechetik.} Leipz. 1863–69, 2 vols. Vol. II. P. I. treats of Luther's Catechism very fully.

Comp. the Literature in Fabricius, Feuerlin, Walch, Baumgarten, Köllner,  \emph{Symbolik}, I. p. 473.

CATECHETICAL INSTRUCTION.

Religious instruction preparatory to admission to church membership is as old as Christianity itself, but it assumed very different shapes in different ages and countries. In the first three or four centuries (as also now on missionary ground) it always \emph{preceded baptism}, and was mainly addressed to adult Jews and Gentiles. In length and method it freely adapted itself to various conditions and degrees of culture. The three thousand Jewish converts on the day of Pentecost, having already a knowledge of the Old Testament, were baptized simply on their profession of faith in Christ, after hearing the sermon of St. Peter. Men like Cornelius, the Eunuch, Apollos, Justin Martyr, Tertullian, Cyprian, Jerome, Ambrose, Augustine, needed but little theoretical preparation, and Cyprian and Ambrose were elected bishops even while yet catechumens. At Alexandria and elsewhere there were special catechetical schools of candidates for baptism. The basis of instruction was the traditional rule of faith or Apostles' Creed, but there were no catechisms in our sense of the term; and even the creed which the converts professed at baptism was not committed to writing, but orally communicated as a holy secret. Public worship was accordingly divided into a \emph{missa catechumenorum} for half-Christians in process of preparation for baptism, and a \emph{missa fidelium} for baptized communicants or the Church proper.

With the union of Church and State since Constantine, and the general introduction of infant baptism, catechetical instruction began to be imparted to \emph{baptized} Christians, and served as a preparation for \emph{confirmation} or the first communion. It consisted chiefly of the committal and explanation, (1) of the Ten Commandments, (2) of the Creed (the Apostles' Creed in the Latin, the Nicene Creed in the Greek Church), sometimes also of the Athanasian Creed and the Te Deum; (3) of the Lord's Prayer (Paternoster). To these were added sometimes special chapters on various sins and crimes, on the Sacraments, and prayers. Councils and faithful bishops enjoined upon parents, sponsors, and priests the duty of giving religious instruction, and catechetical manuals were prepared as early as the eighth and ninth centuries, by Kero, monk of St. Gall (about 720); Notker, of St. Gall (d. 912); Otfried, monk of Weissenbourg (d. after 870), and others.\footnote{Otfried's Catechism was newly edited by J. G. Eccard: '\emph{Incerti Monachi Weissenburgensis Catechesis Theotisca Seculo IX. conscripta.}' Hanov. 1713. It contains: 1. The Lord's Prayer, with an explanation; 2. The Deadly Sins; 3. The Apostles' Creed; 4. The Athanasian Creed; 5. The Gloria.} But upon the whole this duty was sadly neglected in the Middle Ages, and the people were allowed to grow up in ignorance and superstition. The anti-papal sects, as the Albigenses, Waldenses, and the Bohemian Brethren, paid special attention to catechetical instruction.\footnote{Comp. J. C. Köcher: \emph{Catechet. Geschichte der Waldenser, Böhmischen Brüder}, etc. Amst. 1768. And C. A. G. von Zezschwitz: \emph{Die Catechismen der Waldenser und Böhmischen Brüder als Documente ihres gegenseitigen Lehraustausches.} Erlangen, 1863.}

The Reformers soon felt the necessity of substituting evangelical Catechisms for the traditional Catholic Catechisms, that the rising generation might grow up in the knowledge of the Scriptures and the true faith. Of all the Protestant Catechisms, those of Luther follow most closely the traditional method, but they are baptized with a new spirit.

LUTHER'S CATECHISMS.

After several preparatory attempts,\footnote{They began in 1518 with a popular evangelical exposition of the Lord's Prayer, and the Ten Commandments. See Schneider, l.c. pp. xvii. sqq., and Zezschwitz, l.c. II. I. pp. 316 sqq. Nor stood he alone in these labors. Urbanus Regius (author of three Catechisms), Lonicer (Strasburg, 1523), Melanchthon (1524), Brentius (1527 or 1528), Lachmann (\emph{Catechesis}, 1528), Rürer, Althamer, Moiban, Corvin, Rhau, Willich, Chytræus (1555), and other Lutherans of the Reformation period, wrote books for the religious instruction of the young.} Luther wrote two Catechisms, in 1529, both in the German language—first the larger, and then the smaller. The former is a continuous exposition rather than a Catechism, and is not divided into questions and answers; moreover, it grew so much under his hands that it became altogether unsuitable for the instruction of the young, which he had in view from the beginning. Hence he prepared soon afterwards a smaller one, or \emph{Enchiridion}, as he called it.\footnote{First in the second edition, whose title (as given by Riederer, but now wanting in the copy rediscovered by Harnack, l.c. p. xxii.) is this: ' \emph{Enchiridion. Der kleine Catechismus für die gemeine Pfarher und Prediger, gemehret and gebessert durch Mart. Luther. Wittenberg, MDXXIX.}' The title of the third edition, 1531, is: '\emph{Enchiridion.} | \emph{Der kleine Catechismus für die gemeine Pfarher und Prediger.} | \emph{Mart. Lu. MDXXXI.}' See Schneider, l.c. p. 1. This is the standard edition, from which the editions of 1539 and 1542 differ very slightly.} It is the ripe flower and fruit of the larger work, and almost superseded it, or confined its use to pastors and teachers and a more advanced class of pupils.\footnote{See, on the relation of the two, Köllner, l.c. p. 490. He says, p. 520: 'The Large Catechism has entirely gone out of use.' Comp. also Zezschwitz, l.c. p. 324. The older view of the priority of the Small Catechism is wrong.}

He was moved to this work by the lamentable state of religious ignorance and immorality among the German people, which he found out during his visitations of the churches in Saxony, 1527–29.\footnote{He says, in his characteristic style (Preface to the Small Catechism): '\emph{Diesen Katechismum oder christliche Lehre in solche kleine, schlechte, einfältige Form zu stellen, hat mich gezwungen und gedrungen die klägliche elende Noth, so ich neulich erfahren habe, da ich auch ein Visitator war. Hilf, lieber Gott, wie manchen Jammer habe ich gesehen, dass der gemeine Mann doch so gar nichts weiss von der christlichen Lehre, sonderlich auf den Dörfern! Und leider viel Pfarrherren ganz ungeschickt und untüchtig sind zu lehren; und sollen doch alle Christen heissen, getauft sein und der heiligen Sacramente geniessen; können weder Vaterunser, noch den Glauben, oder Zehn Gebote; leben dahin, wie das liebe Vieh und unvernünftige Säue; und nun das Evangelium kommen ist, dennoch fein gelernt haben, aller Freiheit meisterlich zu missbrauchen. O ihr Bischöfe, was wollt ihr doch Christo immer mehr antworten, dass ihr das Volk so schändlich habt lassen hingehen, und euer Amt nicht einen Augenblick je bewiesen? Dass euch alles Unglück fliehe! Verbietet einerlei Gestalt und treibet auf eure Menchengesetze, fraget aber derweil nichts danach, ob sie das Vaterunser, Glauben, Zehn Gebote oder einiges Gotteswort können. Ach und wehe über euren Hals ewiglich! Darum bitte ich um Gottes willen euch alle meine lieben Herren und Brüder, so Pfarrherren oder Prediger sind, wollet euch eures Amtes von Herzen annehmen, euch erbarmen über euer Volk, das euch befohlen ist, und uns helfen den Katechismus in die Leute, sonderlich in das junge Volk bringen; und welche es nicht besser vermögen, diese Tafeln und Formen vor sich nehmen, und dem Volke von Wort zu Wort fürbilden}'}

With his conservative instinct, he retained the three parts on the Decalogue (after the Latin division), the Creed, and the Lord's Prayer. To these he added, after the example of the Bohemian Brethren, an instruction on the Sacraments of Baptism and the Lord's Supper.\footnote{The Bohemian Brethren, or Hussites, had Catechisms long before Luther, divided into five parts: 1. The Decalogue; 2. The Creed; 3. The Lord's Prayer; 4. The Sacraments; 5. The House Table. They sent a Latin copy to Luther, 1523. See Köllner, l.c. pp. 485, 469.}

Luther's Catechism proper, therefore, has five parts: 1. Decalogus; 2. Symbolum Apostolicum; 3. Oratio Dominica; 4. De Baptismo; 5. De Sacramento Altaris. So the Large Catechism, as printed in the Book of Concord (without any additions\footnote{Luther says, in the Prolegomena to the Large Catechism, '\emph{Also hätte man überall}  Fünf Stücke der Ganzen Christlichen Lehre,  \emph{die man immerdar treiben kann.}'}), and the Small Catechism in the first two editions (with devotional additions).

THE ADDITIONS IN THE ENCHIRIDION.

In the later editions of the Small Catechism (since 1564) there is a sixth part on \emph{Confession} and \emph{Absolution}, or the \emph{Power of the Keys},\footnote{'\emph{Vom Amt der Schlüssel. De potestate clavium.}' It is usually called '\emph{Das sechste Hauptstück},' although it should properly be an appendix.} which is inserted either as Part V., between Baptism and the Lord's Supper, or added as Part VI., or as an Appendix. The precise authorship of the enlarged form or forms (for they vary) of this Part, with the questions 'What is the Power of the Keys,' etc., is uncertain,\footnote{It is variously traced to Luther, Brentius (who has in his Catechism a sixth part '\emph{On the Keys}'), Bugenhagen, Knipstrov, but with greater probability to the popular Catechetical Sermons prepared for public use in Nuremburg and Brandenburg, 1533 (probably by Brentius), and translated into Latin by Justus Jonas, 1539 (and re-edited by Gerlach, Berlin, 1839). See Francke: \emph{Libri symbolici}, etc. P. II. Proleg. p. xxiv.; Müller: \emph{Die Symbolischen Bücher}, etc. p. xcvii.; Köllner, l.c. pp. 502 sqq.; Zezschwitz, l.c. pp. 327 sqq.} but the substance of it, viz., the questions on private or auricular confession of sin to the minister and absolution by the minister, as given in the 'Book of Concord,' date from Luther himself, and appear first substantially in the third edition of 1531, as introductory to the fifth part on the Lord's Supper.\footnote{See the third edition, as republished by Schneider, l.c. pp. lii. and 45 sqq. Those questions appear under the title '\emph{Wie man die Einfeltigen soll leren beichten.}' An admonition to confession ('\emph{Vermahnung zu der Beicht}') was added also to later editions of the Larger Catechism since 1531, but omitted in the 'Book of Concord,' against the remonstrance of Chemnitz.} He made much account of private confession and absolution, while the Calvinists abolished the same as a mischievous popish invention. 'True absolution,' says Luther, 'or the power of the keys, instituted in the Gospel by Christ, affords comfort and support against sin and an evil conscience. Confession or absolution shall by no means be abolished in the Church, but be retained, especially on account of weak and timid consciences, and also on account of untutored youth, in order that they may be examined and instructed in the Christian doctrine. But the enumeration of sins should be free to every one, to enumerate or not to enumerate such as he wishes.'\footnote{\emph{Art. Smalc.} III. p. 8. The Church of England holds a similar position in regard to the confessional, and hence the recent revival of it by the Ritualists, though under the strong protest of the evangelical party. The 'Book of Common Prayer' of the Church of England contains, besides two different forms of \emph{public} confession and absolution (one for Morning and Evening Prayer, another for the Communion Service), a form of \emph{private} confession and absolution in the Order for the Visitation of the Sick. The first two are retained, the third is omitted in the Prayer Book of the Protestant Episcopal Church of the United States. The third form, in the Visitation Office, retains the traditional form of the Latin Church—\emph{'Absolvo te in Nomine Patris},' etc.—'I absolve thee in the Name,' etc. Blunt, in his \emph{Annotated Book of Common Prayer}, Part II. p. 283, comments largely on this formula, and quotes also a passage from the first exhortation in the Communion Office, which reads as follows: 'Therefore, if there be any one who . . . requireth further comfort and counsel, let him come to me, or to some other discreet and learned minister of God's Word, and open his grief; that by the ministry of God's Holy Word he may receive the benefit of absolution together with ghostly counsel and advice, to the guiding of his conscience, and avoiding of all scruple and doubtfulness.' And after some other quotations, he says: 'Numberless practical writers speak of private confession as a recognized habit in the Church of England since the Reformation as well as before. Nearly all such writers, however, protest against its compulsory injunction, and it does not seem to be proved that frequent and habitual confession has ever been very common in the Church of England since the Reformation.'}

Besides these doctrinal sections, the Smaller Catechism, as edited by Luther in 1531 (partly, also, in the first edition of 1529), has three appendices of a devotional or liturgical character, viz.: 1. A series of short family prayers ('\emph{wie ein Hausvater sein Gesinde soll lehren Morgens und Abends sich segnen}); 2. A table of duties ('\emph{Haustafel}') for the members of a Christian household, consisting of Scripture passages 1 Tim. iii. 2 sqq.; Rom. xiii. 1 sqq.; Col. iii. 19 sqq.; Eph. vi. 1 sqq., etc.); 3. A marriage manual ('\emph{Traubüchlin}'); and 4. A baptismal manual ('\emph{Taufbüchlin}').

The first two appendices, which are devotional, were retained in the 'Book of Concord;' but the third and fourth, which are liturgical and ceremonial, were omitted because of the great diversity in different churches as to exorcism in baptism, and the rite of marriage.

TRANSLATIONS AND INTRODUCTION.

The Smaller Catechism was translated from the German original into the Latin (by Sauermann) and many other languages; even into the Greek, Hebrew, and Syriac. It is asserted by Lutheran writers that no book, except the Bible, has had a wider circulation. Thirty-seven years after its appearance Matthesius spoke of a circulation of over a hundred thousand copies.

It was soon introduced into public schools, churches, and families. It became by common consent a symbolical book, and a sort of 'Lay Bible' for the German people. It is still very extensively used in Lutheran churches, though mostly with supplements or in connection with fuller Catechisms. In Southern Germany the Catechism of Brentius obtained a wide currency.

CHARACTER, VALUE, AND DEFECTS.

Luther's Small Catechism is truly a great little book, with as many thoughts as words, and every word telling and sticking to the heart as well as the memory. It bears the stamp of the religious genius of Luther, who was both its father and its pupil.\footnote{'I am also a doctor and a preacher,' he says in the Preface to his Larger Catechism, 'endowed with no less learning and experience than those who presume so much on their abilities . . . yet I am like a child who is taught the Catechism, and I read and recite word by word, in the morning and when I have leisure, the Ten Commandments, the Articles of the Creed, the Lord's Prayer, the Psalms, etc. . . . and must remain, and do cheerfully remain, a child and pupil of the Catechism.'} It exhibits his almost apostolic gift of expressing the deepest things in the plainest language for the common people. It is strong food for a man, and yet as simple as a child. It marks an epoch in the history of religious instruction: it purged it from popish superstitions, and brought it back to Scriptural purity and simplicity. As it left far behind all former catechetical manuals, it has, in its own order of excellence and usefulness, never been surpassed. To the age of the Reformation it was an incalculable blessing. Luther himself wrote no better book, excepting, of course, his translation of the Bible, and it alone would have immortalized him as one of the great benefactors of the human race. Few books have elicited such enthusiastic praise, and have even to this day such grateful admirers.\footnote{I quote some Lutheran testimonies which show the impressions of early childhood, and seem extravagant to members of other denominations. Matthesius: 'The world can never sufficiently thank and repay Luther for his little Catechism.' Justus Jonas: 'It may be bought for sixpence, but six thousand worlds would not pay for it.' Andr. Fabricius: 'A better book, next to the Bible, the sun never saw; it is the juice and the blood, the aim and the substance of the Bible.' Seckendorf: 'I have received more consolation and a firmer foundation for my salvation from Luther's little Catechism than from the huge volumes of all the Latin and Greek fathers together.' Löhe: 'It is, of all Confessions, that which is most suitable and best adapted to the people. It is a fact, which no one denies, that no other Catechism in the world can be made a prayer of but this. But it is less known that it may be called a real marvel in respect of the extraordinary fullness and great abundance of knowledge expressed in it in so few words.' Leopold Ranke: 'The Catechism published by Luther in 1529, of which he himself says that, old a doctor as he was, he used it himself as a prayer, is as childlike as it is profound, as comprehensible as it is unfathomable, simple, and sublime. Happy he whose soul was fed by it, who clings to it. He possesses at all times an imperishable consolation: under a thin shell, a kernel of truth sufficient for the wisest of the wise.' ('\emph{Der Katechismus, den Luther im Jahr} 1529, \emph{herausgab, von dem er sagt, er bete ihn selbst, so ein alter Doctor er auch sei, ist ebenso kindlich wie tiefsinnig, so fasslich wie unergründlich, einfach and erhaben. Glückselig wer seine Seele damit nährte, wer daran festhält! Er besitzt einen unvergänglichen Trost in jedem Momente: nur hinter einer leichten Hülle den Kern der Wahrheit, der dem Weisesten der Weisen genug that.}' \emph{Deutsche Geschichte im Zeitalter der Reformation}, Vol. II. 3d edition, Berlin, 1852, p. 357.) To add an American testimony, I quote from Dr. Ch. P. Krauth: 'So truly did the Shorter Catechism embody the simple Christian faith, as to become, by the spontaneous acclamation of millions, a Confession. It was a private writing, and yet, beyond all the Confessions, the direct pulsation of the Church's whole heart is felt in it. It was written in the rapture of the purest catholicity, and nothing from Luther's pen presents him more perfectly, simply as a Christian; not as the prince of theologians, but as a lowly believer among believers' (\emph{The Conservative Reformation}, Philadelphia, 1872, p. 285).}

But with all its excellences it has some serious defects. It gives the text of the Ten Commandments in an abridged form (the Larger Catechism likewise), and follows the wrong division of the Romish Church, which omits the second commandment altogether, and cuts the tenth commandment into two, to make up the number.\footnote{The Lutheran and the Roman Catholic Catechisms, following the lead of Augustine, regard the second commandment only as an explanation of the first; the Reformed and the Greek Catechisms, following the division of the Jews (Josephus and Philo) and the early Christians (e. g. Origen), treat it as a separate commandment, which prohibits image worship and enjoins the true worship of God, while the first prohibits idolatry and enjoins monotheism. Hence the different modes of counting from the second to the ninth commandment. The division of the tenth commandment follows as a necessity from the omission of the second, but is decidedly refuted by the intrinsic unity of the tenth commandment, and by a comparison of Exod. xx. 17 with Deut. v. 21; for in the latter passage (as also in the Septuagint version of Exod. xx. 17) the order is transposed, and the neighbor's wife put before the neighbor's house, so that what is the ninth commandment in Exodus, according to the Roman Catholic and Lutheran view, would be the tenth according to Deuteronomy. St. Paul, moreover, in enumerating the commandments of the second table, Rom. xiii. 9 (comp. also vii. 7), alludes to the tenth with the words, 'Thou shalt not covet,' without intimating any such division. Comp. also Mark x. 19. The Decalogue consists of two tables, of five commandments each. The first contains the duties to God (\emph{præcepta pietatis}), the second the duties to men (\emph{præcepta probitatis}); the first is strictly religious, the second moral. The fifth commandment belongs to the first table, since it enjoins reverence to parents as representing God's authority on earth. This view is now taken not only by Reformed, but also by many of the ablest Lutheran divines, e.g., Oehler, \emph{Theologie des Alten Testaments} (Tübingen, 1873), I. pp. 287 sqq.; H. Schultz, \emph{Alttestamentliche Theologie} (Frankf. a. M. 1869), I. p. 429. On the other hand, Kurtz, Kahnis, and Zezschwitz defend the Lutheran division. The main thing, of course, is not the dividing, but the keeping of the commandments.} It allows only three questions and answers to the exposition of the Creed. It gives undue importance to the Sacraments by making them co-ordinate parts with the three great divisions, and elevates even private confession and absolution, as a sort of third sacrament, to equal dignity. It omits many important articles, and contains no express instruction on the Bible, as the inspired record of divine revelation and the infallible rule of faith and practice. Hence it is found necessary, where it is used, to supplement it by a number of preliminary and additional questions and answers.

THE TEXT OF THE ENCHIRIDION.

The critical restoration of the best text of Luther's Small Catechism has only recently been accomplished by Mönckeberg, Schneider, and Harnack. The text of the 'Book of Concord' is unreliable.

The \emph{editio princeps} of 1529 had entirely disappeared until Mönckeberg, 1851, published a Low-German translation from a copy in the Hamburg city library; and five years later (1856) Professor Harnack found an Erfurt reprint of the original (without date), and a Marburg reprint dated 1529.

The second recension, of 1529, which contains several improvements and addenda, was described by Riederer, in 1765, from a copy then in the university library at Altdorf. This copy was supposed to have been transferred to Erlangen, but was discovered by Harnack in the German Museum at Nuremburg, and republished by him, 1856, together with a reprint of the \emph{editio princeps}, and a Wittenberg edition of 1539, a valuable critical introduction, and a table of the principal variations of the text till 1542.

The third recension, of 1531, was brought to light by Dr. Schneider, and accurately republished (but without the woodcuts and the \emph{Traubüchlin} and \emph{Taufbüchlin}), 1853, with a learned introduction and critical apparatus.\footnote{See his description, l.c. pp. l.–liv. It is reprinted in the second volume of this work.} It gives the text of the five parts substantially as it has remained since, also the section on confession ('\emph{Wie man die Einfältigen soll lehren beichten}'), the morning and evening prayers, the \emph{Benedicite} and \emph{Gratias}, the \emph{Haustafel}, the \emph{Traubüchlin} and the \emph{Taufbüchlin.}

In 1535 (and 1536) Luther prepared a new edition, to conform the Scripture texts to his translation of the Bible, which was completed in 1534.

The edition of 1542 ('\emph{aufs neu übersehen und zugericht}') adds the promise to the fourth (fifth) commandment, and enlarges the 'House Table.'

\xsection{The Articles of Smalcald. A.D. 1537.}

§ 44. The Articles of Smalcald. A.D. 1537.

Literature.

Carpzov:  \emph{Isagoge in Libras Symbolicos}, etc., 1675, pp. 767 sqq.

J. C. Bertram:  \emph{Geschichte des symbol. Anhangs der Schmalk. Artikel.} Altdorf, 1770.

M. Meurer:  \emph{Der Tag zu Schmalkalden und die Schmalk. Artikel.} Leipz. 1837.

Köllner:  \emph{Symbolik} (1837), I. pp. 439–472.

G. H. Klippel, in Herzog's \emph{Real-Encykl.} Vol. XIII. (1860), pp. 600 sqq.

Ch. P. Krauth:  \emph{The Conservative Reformation and its Theology}, Phila. 1872, pp. 280–283.

F. Sander:  \emph{Geschichtliche Einleitung zu den Schmalkaldischen Artikeln.} In the \emph{Jahrbücher für Deutsche Theologie}, Gotha, 1875, pp. 475–489.

The older literature, mostly doctrinal and polemical, is given by Fabricius, Walch, Baumgarten, Hase (\emph{Libri Symb. Proleg.} cxl.), and Köllner.

ORIGIN.

Pope Paul III., yielding at last to the request of the German Emperor and the pressure of public opinion, convoked a general Council, to be opened May 23, 1537, at Mantua,\footnote{It did not convene, however, till 1545, in Trent, and then it turned out an exclusive Roman Catholic Council.} and extended, through his legate, Peter Paul Vergerius (who subsequently became a Protestant), an invitation also to the Lutherans.\footnote{Vergerius had a fruitless interview with Luther in the electoral castle at Wittenberg, which was characteristic of both parties. The papal nuncio acted the proud prelate and shrewd Italian diplomatist; the Reformer, the plain, free-spoken German. Luther took the matter in good humor, sent for the barber, and put on his best dress to impress the nuncio with his youth and capacity for even greater mischief to the Pope than he had done already. He scorned his tempting offers, and told him frankly that he cared very little about his master and his Council at Mantua or elsewhere, but promised to attend it, and there to defend his heretical opinions against the whole world. Vergerius, in his report, speaks contemptuously of Luther's poor Latin, rude manners, obstinacy, and impudence; but some years afterwards he renounced Romanism, and became the Reformer of the Grisons in Eastern Switzerland. He died October 4, 1565, at Tübingen, where he spent his last years, without office, but in extensive literary activity and correspondence. See the monograph of Sixt: \emph{Petrus Paulus Vergerius}, Braunschweig, 1855, pp. 35–45.} Though by no means sanguine as to the result, Luther, by order of the Elector of Saxony (Dec. 11, 1536), prepared a Creed as a basis of negotiations at the Council, submitted it to Amsdorf, Agricola, Spalatin, and Melanchthon for approval, and sent it to the Elector, Jan. 3, 1537.

Melanchthon, at the request of the convent assembled at Smalcald, prepared an Appendix on the power and primacy of the Pope, about which the Augsburg Confession and Apology are silent.

SIGNATURE. MELANCHTHON'S POSITION.

The Articles, including the Appendix, were laid before the convent of Lutheran Princes and theologians held in the town of Smalcald (\emph{Schmalkalden}), in Thuringia, which lent its name to the political league of those Princes for mutual protection, and also to this new Creed.\footnote{'\emph{Schmalkaldische Artikel, Articuli Smalcaldici},' so called since 1553. The original title is: 'Artikel christlicher Lehre,  \emph{so da hätten sollen aufs Concilium zu Mantua, oder wo es sonst worden wäre, überantwortet werden von unsers Theils wegen, und was wir annehmen oder nachgeben könnten oder nicht, durch D. Martin Luthern geschrieben, Anno} 1537.'} They were signed by the theologians (but not by the Princes) without being publicly discussed.\footnote{The Princes on that occasion required their theologians to sign also the Augsburg Confession and Apology, but they resolved to have nothing to do with the Pope's Council. The Appendix has thirty-two signatures, the Articles have forty-two, obtained partly at Smalcald and partly on the journey. The principal signers are Luther, Melanchthon, Jonas, Spalatin, Bugenhagen, Amsdorf, Bucer, and Brentius. See Köllner, pp. 445 sqq., and Plitt, \emph{De auctoritate Articuloram Smalcaldicorum} (Erlang. 1862), with the strictures of Heppe, \emph{Entstehung und Fortbildung des Lutherthums} (Cassel, 1863), pp. 252 sqq.}

Melanchthon signed the Articles with the following remarkable qualification: 'I, Philip Melanchthon, approve the foregoing Articles as pious and Christian. But in regard to the Pope, I hold that, if he would admit the Gospel, we might also permit him, for the sake of peace and the common concord of Christendom, to exercise, by \emph{human} right, his present jurisdiction over the bishops, who are now or may hereafter be under his authority.'\footnote{' \emph{De pontifice autem statuo, si evangelium admitteret} (\emph{so er das Evangelium woltte zulassen}), \emph{ei propter pacem et communem tranquillitatem Christianorum, qui iam sub ipso sunt et in posterum sub ipso erunt, superioritatem in episcopos, quam alioqui habet, jure humano etiam a nobis permitti.}' Sander (p. 488) thinks that Melanchthon did not mean this authority to apply to Protestants. But this is inconsistent with the words '\emph{etiam a nobis.}'}

This remarkable concession strongly contrasts with the uncompromising anti-popery spirit of the Articles, and exposed Melanchthon to much suspicion and abuse. It is self-contradictory and impracticable, since the Pope and his hierarchy will never allow the free preaching of the Gospel in the Protestant sense. But the author's motive was a noble desire for a more independent and dignified position of the Church. He feared—and not without good reason—a worse than papal tyranny from rapacious Protestant Princes, who now exercised the power of supreme bishops and little popes in their territories. He sincerely regretted the loss, not of the episcopal domination, but of the episcopal administration, as a check upon secular despotism.\footnote{'\emph{ Utinam, utinam}'—he wrote to his friend, Joach. Camerarius, Aug. 31, 1530—'\emph{possim non quidem dominationem confirmare, sed administrationem restituere episcoporum. Video enim, qualem simus habituri Ecclesiam, dissoluta } \textgreek{πολιτείᾳ } \emph{ecclesiastica. Video postea multo intolerabiliorem futuram tyrannidem, quam antea unquam fuit}' (\emph{Corp. Reform.} Vol. II. p. 334. Comp. his letter of Sept. 4, 1530, to the same, p. 341). Köllner defends Melanchthon's course.}

CONTENTS.

The Articles of Smalcald consist of three parts.

The first reaffirms, very briefly in four articles, the doctrines of the Apostles' and Athanasian Creeds, about which there was no dispute with the Papists. It corresponds to Articles I. and III. of the Augsburg Confession.

The second and principal part, concerning 'the office and work of Christ, or our redemption,' is polemical against the mass, purgatory, the invocation of saints, monasticism, and popery, which interfere and set aside the true doctrine of redemption. Justification by faith alone is emphasized as the chief article of faith, 'upon which depends all that we teach and do against the Pope, the devil, and all the world. We must, therefore, be entirely certain of this, and not doubt it, otherwise all will be lost, and the Pope, and the devil, and our opponents will prevail and obtain the victory.' The mass is denounced as 'the greatest and most horrible abomination,'\footnote{Luther calls it also 'the dragon's tail (\emph{Drachenschwanz}), which has produced a multiplicity of abominations and idolatries' (\emph{multiplices abominationes et idololatrias.} In German: \emph{viel Ungeziefers und Geschmeiss mancherlei Abgötterei}), P. II Art. 2. He says that the mass will be the chief thing in the proposed Council, and will never be yielded by the Papists. Cardinal Campeius had told him at Augsburg he would rather be torn to pieces than allow the mass to be discontinued. So would he (Luther) rather be reduced to ashes than allow a performer of the mass to be equal to our Lord and Saviour.} purgatory as a 'satanic delusion,' the Pope as 'the true Antichrist' predicted by Paul (2 Thess. ii. 4), because 'he will not permit Christians to be saved without his power.'

The third part treats, in fifteen articles, of sin, of the law, of repentance, of the sacraments, and other doctrines and ordinances, concerning which Protestants may dispute either among themselves or with 'learned and sensible men' (i.e., Catholics in the Council, but not with the Pope, who is said to have no conscience, and to care only about 'gold, honor, and power'). In the article on the Lord's Supper, transubstantiation is expressly excluded, but otherwise the Lutheran doctrine is asserted even in stronger terms than in the Augsburg Confession (viz. that 'the true body and blood of Christ are administered and received, not only by pious, \emph{but also by impious} Christians.'\footnote{Heppe (l.c. p. 253 sq.) says that Luther in his first draft used simpler language, viz., that 'the body and blood of Christ are offered \emph{with} the bread and \emph{with} the wine;' but that Amsdorf insisted on a stronger, anti-Melanchthonian statement.} Luther concludes with spicy remarks against the juggling tricks of the Pope.

The Appendix of Melanchthon is a theological masterpiece for his age, written in a calm, moderate, and scholarly tone; and refutes, from the Bible and from the history of the early Church, these three assumptions of the Pope, as 'false, impious, tyrannical, and pernicious in the extreme,' viz.: 1. That the Pope, as the Vicar of Christ, has by divine right supreme authority over the bishops and pastors of the whole Christian world; 2. That he has by divine right both swords, that is, the power to enthrone and dethrone kings, and to regulate civil affairs; 3. That Christians are bound to believe this at the risk of eternal salvation. He also shows from Scripture and from Jerome that the power and jurisdiction of bishops, as far as it differs from that of other ministers, is of human origin, and has been grossly abused in connection with the papal tyranny.

CHARACTER AND AUTHORITY.

It is clear from this outline that the Articles of Smalcald mark a considerable advance in the final separation of the Lutheran body from the Church of Rome. Luther left Smalcald in bad health (he suffered much of the stone), with the prayer that God may fill his associates with hatred of the Pope, and wrote as his epitaph,

'\emph{Pestis eram vivus, moriens tua mors ero, Papa.}'

The Articles themselves differ from the Augsburg Confession as much as Luther differs from Melanchthon. They are more fresh, vigorous, and original, but less cautions, wise, circumspect, and symmetrical. They are not defensive, but aggressive; not an overture of peace, but a declaration of war. They scorn all compromises, and made a reconciliation impossible. They were, therefore, poorly calculated to be a basis of negotiation at a general Council, and were, in fact, never used for that purpose. The Convent at Smalcald resolved not to send any delegates to the Council. But the Smalcald Articles define the position of Lutheranism towards the Papacy, and give the strongest expression to the doctrine of justification by faith. They accordingly took their place, together with the Appendix, among the symbolical books of the Lutheran Church, and were received into various \emph{Corpora Doctrinæ}, and at last into the 'Book of Concord.'\footnote{Comp. Plitt and Heppe, above quoted (p. 254).}

TEXT.\footnote{\par See the minor particulars in Bertram, l.c., and Köllner, pp. 454 sqq.}

Luther prepared the Smalcald Articles at Wittenberg in the German language, and edited them, in 1538, with a preface and considerable changes and additions, but without the signatures, and without the Appendix of Melanchthon. In 1543 and 1545 he issued new editions with slight changes. The first draft, as copied by Spalatin, and signed at Smalcald, was published from the archives of Weimar in 1553, together with Luther's additions and Melanchthon's Appendix, and embodied in the 'Book of Concord.'\footnote{The original MS. of Luther, from which Spalatin made his copy before Luther added his changes, was discovered in the Palatinate Library at Heidelberg in 1817, and edited by Marheineke, with notes, Berlin, 1817.}

The Latin text, as it appeared in the first edition of the 'Book of Concord,' was a poor translation, but was much improved in the edition of 1584.

Melanchthon wrote the Appendix at Smalcald in Latin, but a German translation by Dietrich was signed there, and passed, as the supposed original, into the works of Luther and the first edition of the 'Book of Concord' (1580). The corrected Latin edition of 1584 gave the Latin original, but as the work of all the theologians convened at Smalcald.\footnote{Under the title '\emph{De Potestate et Primatu Papæ. Tractatus per Theologos Smalcaldiæ congregatos conscriptus.}'} This error prevailed nearly two hundred years, until the careful researches of Bertram dispelled it.

\xsection{The Form of Concord. A.D. 1577.}

§ 45. The Form of Concord. A.D. 1577.

Literature.

I. The text of the 'Form of Concord' is found in all the editions of the 'Book of Concord' (\emph{Concordienbuch}), see p. 220.

Heinr. Ludw. Jul. Heppe (Professor in Marburg, an indefatigable investigator of the early history of German Protestantism In the interest of Melanchthonianism): \emph{Der Text der Bergischen Concordienformel, verglichen mit dem Text der Schäbischen Concordie, der Schwäbisch-Sächsischen Concordie und des Torgauer Buches.} Marburg, 1857, 2d ed. 1860.

II. Jacob Andrkæ:  \emph{Sechs christlicher Predig von den Spaltungen, so sich zwischen den Theologen Augspurgischer Confession von Anno} 1548 \emph{bis auf diess }1573 \emph{Jar, nach und nach erhoben}, etc. Tübingen, 1573. Republished by Professor Heppe in Appendix I. to the third volume of his \emph{History of German Protestantism} (see below). In the same volume Heppe published also 'the Swabian and Saxon Form of Concord,' the 'Maulbronn Formula,' and other important documents.

\emph{Apologia oder Verantwortung des christl. Concordienbuchs}, etc. (usually called the \emph{Erfurt Book}, an official Apology, prepared at Erfurt and Quedlinburg by Kirchner, Selnecker, Chemnitz,  and other Lutheran divines). Heidelb. 1583; Dresden, 1584, etc.

Rud. Hospinian (Reformed, d. at Zurich 1626): Concordia Discors;  \emph{h. e. de origine et progressu Formulæ Bergensis}, etc., \emph{ex actis tum publicis, tum privatis} . . . Tig. 1607; Genev. 1678, folio. (The chief work against the 'Form of Concord.')

Leonh. Hutter (Lutheran, d. at Wittenberg 1616): Concordia Concors;  \emph{de origine et progressu Formulæ Concordiæ ecclesiarum Conf. Aug.} . . . \emph{in quo eius } orthodoxia  . . . \emph{demonstratur: et Rud. Hospiniani Tigurini Helvetii convitia, mendacia, et manifesta crimina falsi deteguntur ac solide refutantur . . . ex actis publicis.} Vitemb. 1614; Francof. and Lips. 1690. (This is the most elaborate defense of the 'Form of Concord' called forth by Hospinian's \emph{Conc. discors}, and covers 1460 pp., exclusive of Proleg.)

J. Musæus:  \emph{Prælectiones in Epitom. Form. Conc.} Jen. 1701.

Val. Löscher:  \emph{Historia motuum}, etc. Leipz. 1723, Tom. III. Lib. VI. c. 5 and 9.

Jac. H. Balthasar:  \emph{Historie des Torgisehen Buchs als des nächsten Entwurfs des Bergischen Concordienbuchs}, etc. Greifswald, 1741–56. (In nine parts or dissertations.)

J. Nic. Anton:  \emph{Geschichte der Concordienformel.} Leipz. 1779.

G. J. Planck:  \emph{Geschichte der Entstehung}, etc., \emph{unseres Protest. Lehrbegrifs} . . . \emph{bis zur Einführung der Concordienformel.} Leipz. 1791–1800. Vols. IV.–VI. A work of thorough learning, independent judgment, but without, proper appreciation of the doctrinal differences.

Gottfr. Thomasius (Lutheran): \emph{Das Bekenntniss der evangel. luther. Kirche in der Consequenz seines Princips.} Nürnberg, 1848.

K. Fr. Göschel (Lutheran): \emph{Die Concordienformel nach ihrer Geschichte, Lehre und kirchlichen Bedeutung.} Leipz. 1858.

H. L. J. Heppe (Reformed): \emph{Geschichte des deutschen Protestantismus in den Jahren} 1555–81. Marburg, 1852–58. 4 vols. (The last two volumes contain the history of the 'Form of Concord' and of the 'Book of Concord,' and are also published under the separate title '\emph{Geschichte der lutherischen Concordienformel und Concordie.}')

Gieseler:  \emph{Text-Book of Church History.} American edition, by H. B. Smith, Vol. IV. (New York, 1862), pp. 423–490; German edition, Vol. III. P. II. (Bonn, 1853), pp. 187–310. (A condensed, careful, and impartial statement of the controversies, with citations from the original authorities.)

Dan. Schenkel:  Art. \emph{Concordienformel}, in Herzog's \emph{Real-Encykl.}, Vol. III. (1855), pp. 87–105.

Wilh. Gass:  \emph{Geschichte der Protest. Dogmatik in ihrem Zurammenhang mit der Theologie überhaupt.} Berlin, 1854–67, 4 vols. Vol. I. pp. 21–80.

Gustav Frank (of Jena): \emph{Geschichte der Protest. Theologie.} Leipz. 1862. Vol. I. pp. 94–290.

F. H. R. Frank (Lutheran): \emph{Die Theologie der Concordienformel hist. dogmatisch entwichelt und beleuchtet.} Erlangen, 1858–65. 4 vols. (Chiefly doctrinal.)

H. F. A. Kahnis (Lutheran): \emph{Die Luther. Dogmatik}, Vol. II. (Leipzig, 1864), pp. 515–560.

Is. A. Dorner:  \emph{Geschichte der protestantischen Theologie} (München, 1867), pp. 330–374.

Chas. P. Krauth (Lutheran). \emph{The Conservative Reformation and its Theology} (Phila. 1872), pp. 288–328.

\emph{Denkmal der dritten Jubelfeier der Concordienformel}, 1877. St. Louis, 1877.

NAME. ORIGIN AND OCCASION.

The Form of Concord (\emph{Formula Concordiæ}), the last of the Lutheran Confessions, completed in 1577 and first published in 1580, is named from its aim to give doctrinal unity and peace to the Lutheran Church, after long and bitter contention.\footnote{The name was chosen after older formularies (e.g., the \emph{Henoticon} of Emperor Zeno, the \emph{Formula Concordiæ Wittenbergensis}, 1536, the \emph{Formula Concordiæ inter Suevicas et Saxonicas ecclesias}, 1576, etc.), and occurs first in the edition of Heidelberg, 1582. In the \emph{editio princeps} (1580) the book is called '\emph{Das Buch der Concordien},' but this title was afterwards reserved for the collection of all the Lutheran symbols ('\emph{Concordia},' or '\emph{Liber Concordiæ},' '\emph{Book of Concord}'). It was also called the \emph{Bergische-Buch}, from the place of its composition.} The work was occasioned by a series of doctrinal controversies, which raged in the Lutheran Church for thirty years with as much passion and violence as the trinitarian and christological controversies in the Nicene age. They form a humiliating and unrefreshing, yet instructive and important chapter in the history of Protestantism. The free spirit of the Reformation, which had fought the battles against the tyranny of the Papacy and brought to light the pure doctrines of the Gospel, gave way to bigotry and intolerance among Protestants themselves. Calumny, abuse, intrigue, deposition, and exile were unsparingly employed as means to achieve victory. Religion was confounded with theology, piety with orthodoxy, and orthodoxy with an exclusive confessionalism. Doctrine was overrated, and the practice of Christianity neglected. The contending parties were terribly in earnest, and as honest and pious in their curses as in their blessings; they fought as if the salvation of the world depended on their disputes. Yet these controversies were unavoidable in that age, and resulted in the consolidation and completion of the Lutheran system of doctrine. All phases and types of Christianity must develop themselves, and God overrules the wrath of theologians for the advancement of truth.

LUTHER AND MELANCHTHON.

The seeds of these controversies lay partly and chiefly in the theological differences between Luther and Melanchthon in their later years, partly in the relations of Lutheranism to Romanism and Calvinism.

Luther the Reformer, and Melanchthon the Teacher of Germany, essentially one and inseparable in mind and heart, in doctrine and life, represented in their later period, which may be dated from the year 1533, two types of Lutheranism, the one the conclusive and exclusive, the other the expansive and unionistic type. Luther, at first more heroic and progressive, became more cautions and conservative; while Melanchthon, at first following the lead of the older and stronger Luther, became more independent and liberal.

Luther, as the Reformer of the Romish Church, acted in the general interest of evangelical religion, and enjoys the admiration and gratitude of all Protestants; Luther, as the leader of a particular denomination, assumed a hostile attitude towards other churches, even such as rested on the same foundation of the renewed gospel. After his bold destructive and constructive movements, which resulted step by step in the emancipation from popery, he felt disposed to rest in his achievements. His disgust with the radicalism and fanaticism of Carlstadt and Münzer, his increasing bodily infirmities, and his dissatisfaction with affairs in Wittenberg (which he threatened to leave permanently in 1544), cast a cloud over his declining years. He had so strongly committed himself, and was so firm in his convictions, that he was averse to all further changes and to all compromises. He was equally hostile to the Pope, whom he hated as the very antichrist, and to Zwingli, whom he regarded as little better than an infidel.\footnote{The deepest ground of Luther's aversion to Zwingli must be sought in his mysticism and veneration for what he conceived to be the unbroken faith of the Church. He strikingly expressed this in his letter to Duke Albrecht of Prussia (which might easily be turned into a powerful argument against the Reformation itself). He went so far as to call Zwingli a non-Christian (\emph{Unchrist}), and ten times worse than a papist (March, 1528, in his \emph{Great Confession on the Lords Supper}). His personal interview with him at Marburg (October, 1529) produced no change, but rather intensified his dislike. He saw in the heroic death of Zwingli and the defeat of the Zurichers at Cappel (1531) a righteous judgment of God, and found fault with the victorious Papists for not exterminating his heresy (\emph{Wider etliche Rottengeister}, Letter to Albrecht of Prussia, April, 1532, in De Wette's edition of \emph{L. Briefe}, Vol. IV. pp. 352, 353). And even shortly before his death, unnecessarily offended by a new publication of Zwingli's works, he renewed the eucharistic controversy in his \emph{Short Confession on the Lord's Supper} (1544, in Walch's edition, Vol. XX. p. 2195), in which he abused Zwingli and Oecolampadius as heretics, liars, and murderers of souls, and calls the Reformed generally '\emph{eingeteufelte} [\textgreek{ἐνδιαβολισθέντες}], \emph{durchteufelte, überteufelte lästerliche Herzen und Lügenmäuler.}' No wonder that even the gentle Melanchthon called this a 'most atrocious book,' and gave up all hope for union (letter to Bullinger, Aug. 30, 1544, in \emph{Corp. Reform.} Vol. V. p. 475: '\emph{Atrocissimum Lutheri scriptum, in quo bellum} \textgreek{ περὶ δείπνου κυριακοῦ } \emph{instaurat};' comp. also his letter to Bucer, Aug. 28, 1544, in \emph{Corp. Reform.} Vol. V. p. 474, both quoted also by Gieseler, Vol. IV. p. 412, note 38, and p. 434, note 37). But it should in justice be added, first, that Luther's heart was better than his temper, and, secondly, that he never said a word against Calvin; on the contrary, he seems to have had great regard for him, to judge from his scanty utterances concerning him (quoted by Gieseler, Vol. IV. p. 414, note 43). Calvin behaved admirably on that occasion; he warned Bullinger (Nov. 25, 1544) not to forget the extraordinary gifts and services of Luther, and said: 'Even if he should call me a devil, I would nevertheless honor him as a chosen servant of God.' And to Melanchthon he wrote (June 28, 1545): 'I confess that we all owe the greatest thanks to Luther, and I should cheerfully concede to him the highest authority, if he only knew how to control himself. Good God! what jubilee we prepare for the Papists, and what sad example do we set to posterity!' Melanchthon entirely agreed with him.}

Melanchthon, on the other hand, with less genius but more learning, with less force but more elasticity, with less intuition but more logic and system than Luther, and with a most delicate and conscientious regard for truth and peace, yet not free from the weakness of a compromising and temporizing disposition, continued to progress in theology, and modified his views on two points—the freedom of the will and the presence of Christ in the Eucharist; exchanging his Augustinianism for Synergism, and relaxing his Lutheranism in favor of Calvinism; in both instances he followed the ethical, practical, and unionistic bent of his mind. A minor difference on the \emph{human} right of the papacy and episcopacy appeared in private letters and in his qualified subscription to the Smalcald Articles (1537), but never assumed a serious, practical aspect, except indirectly in the adiaphoristic controversy.\footnote{Kahnis \{\emph{Luth. Dogm.} Vol. II. p. 520) traces the changes of Melanchthon to 'a truly evangelical search after truth, to a practical trait, which easily breaks off the theological edges to bring the doctrine nearer to life, and to the endeavor to reconcile opposites.' Krauth (\emph{Conservative Reformation}, p. 289), who sympathizes with strict Lutheranism, says: 'Melanchthon's vacillations were due to his timidity and gentleness of character, tinged as it was with melancholy; his aversion to controversy; his philosophical, humanistic, and classical cast of thought, and his extreme delicacy in matters of style; his excessive reverence for the testimony of the Church, and of her ancient writers; his anxiety that the whole communion of the West should be restored to harmony; or that, if this were impossible, the Protestant elements, at least, should be at peace.' Comp. on this whole subject the works of Galle:  \emph{Characteristik Melanchthon's als Theologen und Entwicklung seines Lehrbegriffs} (Halle, 1840), pp. 247 sqq. and 363 sqq.; Matthes:  \emph{Phil. Melanchthon} (Altenb. 1841); Ebrard:  \emph{Das Dogma vom heil. Abendmahl} (Frankf. 1846), Vol. II. pp. 434 sqq.; Gieseler:  \emph{Church History}, Vol. IV. pp. 423 sqq.; Heppe:  \emph{Die confessionelle Entwicklung der altprotestantischen Kirche Deutschlands} (Marburg, 1854), pp. 95 sqq.; Carl Schmidt:  \emph{Philipp Melanchthon.} (Elberfeld, 1861), pp. 300 sqq.; Kahnis, l.c. pp. 515 sqq.}

These changes were neither sudden nor arbitrary, but the result of profound and constant study, and represented a legitimate and necessary phase in the development of Protestant theology, which was publicly recognized in various ways before the formation of the 'Form of Concord.' If there ever was a modest, cautious, and scrupulously conscientious scholar, it was Melanchthon. 'There is not a day nor a night for the last ten years,' he assures an intimate friend, 'that I did not meditate upon the doctrine of the Lord's Supper.'\footnote{\emph{Ep. ad Vitum Theodorum}, May 24, 1538 (in \emph{Corp. Reform.} Vol. III. p. 537): '\emph{Scias, amplius decennio nullum diem, nullam noctem abiisse, quin hac de re cogitarim.}'}

As to human freedom, Melanchthon at first denied it altogether, like Luther and the other Reformers, and derived all events and actions, good and bad, from the absolute will of God.\footnote{\emph{Loci theol.} first ed. 1521, A. 7: '\emph{Quandoquidem omnia, quæ eveniunt, necessario juxta divinam prædestinationem eveniunt, nulla est voluntatis nostræ libertas.}' In the edition of 1525 he says: '\emph{Omnia necessario evenire Scripturæ docent. . . . Nec in externis nec in internis operibus ulla est libertas, sed eveniunt omnia juxta destinationem divinam. . . . Tollit omnem libertatem voluntatis nostræ prædestinatio divina.}' (Mel. \emph{Opera} in \emph{Corp. Reform.} Vol. XXI. pp. 88, 93, 95.) In his Commentary on the Romans, published 1524 (cap. 8), Melanchthon calls the power of choice a '\emph{ridiculum commentum},' and derives all things, '\emph{tam bona quam mala},' from the absolute will of God, even the adultery of David ('\emph{Davidis adulterium}') and the treason of Judas ('\emph{Judæ proditio}'), which are the proper work of God ('\emph{ejus proprium opus}') as much as the vocation of Paul; for he does all things not '\emph{permissive, sed potenter.}' He saw this doctrine so clearly in the Epistle to the Romans and other portions of Scripture that passages like  1 Tim. ii. 4 (all men, e.g., all sorts of men) must be adjusted to it. See Galle, pp. 252 sqq., and Heppe, \emph{Dogmatik des deutschen Protestantismus in 16ten Jahrh.} (Gotha, 1857) Vol. I. pp. 434 sqq. In December, 1525, Luther expressed the same views in his book against Erasmus, which he long afterwards (1537) pronounced one of his best works. Comp. p. 215, and Köstlin, \emph{Luther's Theol.} Vol. II. pp. 37, 323. But on Melanchthon the reply of Erasmus (1526) had some effect (as we may infer from the tone of his letter to Luther, Oct. 2, 1527, \emph{Corp. Reform.} Vol. I. p. 893).} Then he avoided the doctrine of predestination, as an inscrutable mystery, and admitted freedom in the sphere of natural life and morality, but still denied it in the spiritual sphere or the order of grace.\footnote{So in the Augsburg Confession (1530), Art. XVIII.: '\emph{De libero arbitrio docent, quod humana voluntas habeat aliquam libertatem ad efficiendam civilem justitiam et diligendas res rationi subjectas. Sed non habet vim sine Spiritu Sancto efficiendæ justitiæ spiritualis, quia animalis homo non percipit ea, quæ sunt Spiritus Dei.}' In Art. XIX. the cause of sin is traced to the will of man and the devil.} At last (after 1535) he openly renounced determinism or necessitarianism, as a Stoic and Manichæan error, and taught a certain subordinate co-operation of the human will in the work of conversion; maintaining that conversion is not a mechanical or magical, but a moral process, and is brought about by the Holy Spirit through the Word of God, with the consent, yet without any merit of man. The Spirit of God is the primary, the Word of God the secondary or instrumental agent of conversion, and the human will allows this action, and freely yields to it.\footnote{First in a new edition of his Commentary to the Romans, 1532, and then in the edition of the '\emph{Loci communes theologici recogniti},' 1535. Here he declares that God is not the cause of sin, but the '\emph{voluntas Diaboli}' and the '\emph{voluntas hominis sunt causæ peccati;}' that we should keep clear of the '\emph{deliramenta de Stoico fato aut} \textgreek{ περὶ τῆς ἀνάγκης,}' that the human will can '\emph{suis viribus sine renovatione aliquo modo externa legis opera facere,}' but that it can not '\emph{sine Spiritu Sancto efficere spirituales affectus, quos Deus requirit. . . . Deus antevertit nos, vocat, movet, adjuvat; sed nos viderimus ne repugnemus. Constat enim peccatum oriri a nobis, non a voluntate Dei. Chrysostomus inquit}:  \textgreek{ὁ δὲ ἕλκων τὸν βουλόμενον ἕλκει. } \emph{Id apte dicitur auspicanti a verbo, ne adversetur, ne repugnet verbo.}' (See Mel. \emph{Opera} in \emph{Corp. Reform.} Vol. XXI. pp. 371–376.) In a new revision of his \emph{Loci}, which appeared in 1548, two years after Luther's death, and in all subsequent editions, he traces conversion to three concurrent causes—the Spirit of God, the Word of God, and the will of man; and states that the will may accept or reject God's grace. '\emph{Veteres aliqui},' he says (\emph{Corp. Reform.} Vol. XXI. pp. 567, 659), '\emph{sic dixerunt: Liberum arbitrium in homine}  facultatem  \emph{esse}  applicandi se ad gr tiam, i.e., \emph{audit promissionem et assentiri conatur et abjicit peccata contra conscientiam. . . . Cum promissio sit universalis, nec sint in Deo contradictoriæ voluntates, necesse est in nobis esse aliquam discriminis causam, cur Saul abjiciatur, David recipiatur}, i.e., \emph{necesse est, aliquam esse actionem dissimilem in his duobus. Hæe dextre intellecta vera sunt, et usus in exercitiis fidei et in vera consolatione, cum æquiescunt animi in Filio Dei monstrato in promissione, illustrabit hanc}  copulationem causarum, verbi dei, spiritus sancti, et voluntatis.' This is the chief passage, which was afterwards (1553) assailed as synergistic. Comp. Galle, pp. 314 sqq.; Gieseler, Vol. IV. pp. 426 and 434; Heppe, l.c. pp. 434 sqq., and \emph{Die confessionelle Entwicklung der alt protest. Kirche Deutschlands,} pp. 107 and 130; Kahnis, l.c. Vol. II. p. 505.}

This is the amount of his Synergism, so called by his opponents. It resembles, indeed, semi-Pelagianism in maintaining a remnant of freedom after the fall, and furnished a basis for negotiations with moderate Romanists, but it differs from it materially in ascribing the initiative and the whole merit of conversion to God's grace. He never gave up the doctrine of justification by the free grace and sole merit of Christ through faith, but in his later years he laid greater stress on the responsibility of man in accepting or rejecting the gospel, and on the necessity of good works as evidences of justifying faith.

As to the Lord's Supper, he at first fully agreed with Luther's view, under the impression that it was substantially the old Catholic doctrine held by the fathers, for whom he had great regard, especially in matters of uncertain exegesis.\footnote{He says (1559): '\emph{Existimo ad confirmandas mentes consensum Vetustatis plurimum conducere}' (quoted by Galle, p. 452). He endeavored to prove the agreement of the fathers with Luther in \emph{Sententiæ Patrum de Cæna Domini}, March, 1530. He there quotes Cyril, Chrysostom, Theophylactus, Hilary, Cyprian, Irenæus, Ambrose, and John of Damascus, and labors also to bring Augustine on his side, but with difficulty (as he says that the body of Christ \emph{in uno loco esse}), and he admits that some passages of Jerome, Gregory of Nazianzum, and Basil might be quoted against Luther. See Galle, pp. 390 sqq.} He also shared his dislike of Zwingli's theological radicalism, and was disposed to trace it to a certain insanity.\footnote{He wrote to Luther from Augsburg, July 14,1530 (\emph{Corp. Reform.} Vol. II. p. 193): '\emph{Zwinglius misit huc confessionem impressam typis. Dicas simpliciter mente captum esse. De peccato originali, de usu sacramentorum veteres errores palam renovat. De ceremoniis loquitur valde helvetice, hoc est barbarissime, velle se omnes ceremonias esse abolitas. Suam causam de sacra cœna vehementer urget. Episcopos omnes vult deletes esse.}'} But his deeper and long-continued study of the subject, and his correspondence and personal intercourse with Bucer and Calvin, gradually convinced him that St. Augustine and other fathers favored rather a figurative or symbolical interpretation of the words of institution,\footnote{In this respect the learned \emph{Dialogus} of Oecolampadius (1530), directed against his \emph{Sententiæ}, made a decided impression on his mind. See Galle, p. 407, and Gieseler, Vol. IV. p. 428. He found a great diversity of views among the fathers ('\emph{mira dissimilitudo,}' see letter to Bucer, 1535, \emph{Corp. Reform.} Vol. II. p. 842), but strong proofs for the figurative interpretation in Augustine, Tertullian, Origen, and all those who speak of the eucharistic elements as \emph{figures}, \emph{symbols}, \emph{types}, and \emph{antitypes} of the body and blood of Christ (see his letter to Crato of Breslau, 1559, quoted by Galle, p. 452).} and that the Scriptures taught a more simple, spiritual, and practical doctrine than either transubstantiation or consubstantiation. Owing to his characteristic modesty and caution, and his deep sense of the difficulties surrounding the problem, he did not set forth a fully developed theory or definition of the mode of Christ's presence, but he substantially agreed with Bucer and Calvin. He gave up the peculiar features of Luther's doctrine, viz., the literal interpretation of the words of institution, and the oral manducation of the body of Christ.\footnote{He first renounced Luther's view, after an interview with Bucer at Cassel, in a letter to Camerarius, Jan. 10, 1535 (\emph{Corp. Reform.} Vol. II. p. 822: '\emph{Meam sententiam noli nunc requirere, fui enim nuncius aliæ,}' i.e., Luther's), and in a confidential letter to Brentius, Jan. 12, 1535 (\emph{Ib.} Vol. II. p. 824, where he speaks in a Greek sentence of the typical interpretation of many of the ancients). Then more fully in the revision of his \emph{Loci Theol.}, 1585 (\emph{de cæna Domini}, in \emph{Corp. Reform.} Vol. XXI. p. 478 sq.). In the \emph{Wittenberg Concordia} (1536) he and Bucer yielded too much to Luther for the sake of peace (compare, however, Dorner, p. 325), but in 1540 he introduced his new conviction into the tenth article of the Augsburg Confession (see above, p. 241), and adhered to it. In his subsequent deliverances he protested against ubiquity and \textgreek{ἀρτολατρεία,} and the fanatical intolerance of the ultra-Lutherans, who denounced him as a traitor. Calvin publicly declared that he and Melanchthon were inseparably united on this point: '\emph{Confirmo, non magis a me Philippum quam a propriis visceribus in hac causa posse divelli}' (\emph{Admonitio ultima ad Westphalum, Opp.} VIII. p. 687). Galle maintains that Melanchthon stood entirely on Calvin's side (l.c. p. 445). So does Ebrard, who says: '\emph{Melanchthon kam, ohne auf Calvin Rücksicht zu nehmen, ja ohne von dessen Lehre wissen zu können, auf selbständigem Wege zu derselben Ansicht, welche bei Calvin sich ausgebildet hatte}' (\emph{Das Dogma u. heil. Abendmahl}, Vol. II. p. 437). Yet in the doctrine of predestination they were wide apart. A beautiful specimen of harmony of spirit with diversity in theology! After his death Calvin appealed to the sainted spirit of Melanchthon now resting with Christ: '\emph{Dixisti centies, cum fessus laboribus et molestiis oppressus caput familiariter in sinum meum deponeres: Utinam, utinam moriar in hoc sinu! Ego vero millies postea optavi nobis contingere, ut simul essemus}' (\emph{Opp.} VIII. p. 724).} He also repeatedly rejected (as, in fact, he never taught) the Lutheran dogma of the ubiquity of Christ's body, as being inconsistent with the nature of a body and with the fact of Christ's ascension to heaven and sitting in heaven, whence he shall return to judgment.\footnote{Dorner, l.c. p. 354: '\emph{Melanchthon hat Luther's christologische Ansichten aus der Zeit des Abendmahlsstreites nie getheilt. Die Menschwerdung besteht ihm in der Aufnahme der menschlichen Natur in die}  Person  \emph{des Logos, nicht aber in der Einigung } (\emph{unio}) \emph{der}  Natur  \emph{des Logos mit der Menschheit in realer Mittheilung der Prädicate der ersteren an die letztere. Die communicatio idiomatum ist ihm nur eine dialektische, verbale: die Person des Logos ist Person des ganzen Christus und trägt die Menschheit als ihr Organon.}'} But he never became a Zwinglian; he held fast to a spiritual real presence of the person (rather than the body) of Christ, and a fruition of his life and benefits by faith. In one of his last utterances, shortly before his death, he represented the idea of a vital union and communion with the person of Christ as the one and only essential thing in this sacred ordinance.\footnote{'\emph{Responsio Phil. Mel. ad quæstionem de controversia Heidelbergensi} (\emph{Corp. Reform.} Vol. IX. p. 961): \emph{Non difficile, sed periculosum est respondere. . . . In hac controversia optimum esset  retinere verba Pauli: ``Panis, quem frangimus,} \textgreek{ κοινωνία ἐστὶ τοῦ σώματος.}'' \emph{ Et copiose de fructu Cænæ dicendum est, ut invitentur homines ad amorem hujus pignoris et crebrum usum. Et vocabulum } \textgreek{κοινωνία} \emph{ declarandum est. Non dicit, mutari naturam panis, ut Papistæ dicunt; non dicit, ut Bremenses, panem esse substantiale Corpus Christi; non dicit, ut Heshusius, panem esse verum corpus Christi: sed esse } \textgreek{κοινωνίαν,} i.e., \emph{ hoc, quo fit consociatio cum corpore Christi, quæ fit in usu, et quidem non sine cogitatione, ut cum mures panem rodunt. . . . Adest Filius Dei in ministerio Evangelii, et ibi certo est efficax in credentibus, ac adest non propter panem, sed propter hominem, sicut inguit: ``Manete in me, et ego in vobis.''}' Comp. on the whole eucharistic doctrine of Melanchthon the learned exposition of Heppe, in the third volume of his \emph{Dogmatik des deutschen Protestantismus im} 16\emph{ten Jahrh.} pp. 143 sqq. He says, p. 150, with reference to the passage just quoted: '\emph{Immer und überall betont es Melanchthon, dass Christi Leib und Blut im Abendmahle mitgetheilt wird, inwiefern daselbst eine Mittheilung des}  Lebendigen  \emph{Leibes, der gottmenschlichen}  Person  \emph{Christi stattfindet, dass die Vereinigung Christi und der Gläubigen, für welche das Abendmahl gestiftet ist, eine persönliche Gemeineschaft, persönliches, lebendiges, wirksames Einwohnen des Gottmenschen in dem Gläubigen ist.}' See also Ebrard, Vol. II. pp. 434 sqq.}

Luther no doubt felt much grieved at these changes, and was strongly pressed by contracted and suspicious minds to denounce them openly, but he was too noble and generous to dissolve a long and invaluable friendship, which forms one of the brightest chapters in his life and in the history of the German Reformation.\footnote{Their friendship was, indeed, seriously endangered, and for some time suspended, but fully restored again; for it rested on their union with Christ. Luther wrote to Melanchthon, June 18, 1540 (\emph{Briefe}, Vol. V. p. 293): '\emph{Nos tecum, et tu nobiscum, et Christus hic et ibi nobiscum.}' He spoke very highly of Melanchthon's \emph{Loci} in March, 1545, and in January, 1546, he called him a true man, who must be retained in Wittenberg, else half the university would go off with him (\emph{Corp. Reform.} Vol. VI. p. 10; Gieseler, Vol. IV. pp. 432–435). Dorner justly remarks (l.c. p. 332 sq.): '\emph{Wenn zu dem Edelsten in Luther auch die ihn zum Reformator befähigende Weitherzigkeit und Demuth gehörte, womit er die eigenthümlichen Gaben Anderer, vor allem Melanchthon's anerkannte, so war es das Bestreben jener engherzigen Freunde, Luthern auf sich selbst zu beschränken, der Ergänzungsbedürftigkeit auch dieser vielleicht grössten nachapostolischen Persönlichkeit zu vergessen und, was ihnen jedoch nicht gelang, auch ihn selbst derselben vergessen zu machen.}' Melanchthon, on his part, although he complained at times of Luther's \textgreek{φιλονεικία} (as a \textgreek{πάθος,} not a \emph{crimen}), and overbearing violence of temper, and thought once (1544) seriously of leaving Wittenberg as a 'prison,' admired and loved him to the end, as the Elijah of the Reformation and as his spiritual father. In announcing to his students the death of Luther (Feb. 18, 1546) on the day following, he paid him this noble and just tribute: '\emph{Obiit auriga et currus Israel, qui rexit ecclesiam in hac ultima senecta mundi,}' and added, '\emph{Amemus igitur hujus viri memoriam et genus doctrinæ ab ipso traditum, et simus modestiores et consideremus ingentes calamitates et mutationes magnas, quæ hunc casum sunt secuturæ.}' Comp. Planck, l.c. Vol. IV. pp. 71–77.} He kept down the rising antagonism by the weight of his personal authority, although he foresaw the troubles to come.\footnote{While sick at Smalcald, 1537, he told the Elector of Saxony that after his death discord would break out in the University of Wittenberg, and his doctrine would be changed. Seckendorf, \emph{Com. de Lutheranismo},' III. p. 165.} After his death (1546) the war broke out with unrestrained violence. Melanchthon was too modest, peaceful, and gentle for the theological leadership, which now devolved upon him; he kept aloof from strife as far as possible, preferring to bear injury and insult with Christian meekness, and longed to be delivered from the 'fury of the theologians' (\emph{a rabie theologorum}), which greatly embittered his declining years.\footnote{'\emph{Ego æquissimo animo,}' he wrote to Camerarius, Feb 24, 1545 (\emph{Corp. Reform.} Vol. V. p.684), '\emph{vel potius } \textgreek{ἀναισθήτως } \emph{fero insolentiam } \textgreek{καὶ ὕβρεις} \emph{ multorum, et dum vivam moderate faciam officium meum.}'} He left the scene of discord April 19, 1560, fourteen years after Luther. His last wish and prayer was 'that the churches might be of one mind in Jesus Christ.' He often repeated the words, 'Let them all be one, even as thou, Father, art in me, and I in thee.' He died with the exclamation, 'O God, have mercy upon me for the sake of thy Son Jesus Christ! In thee, O Lord, have I put my trust; I shall not be confounded forever and ever.' The earthly remains of the ``\emph{Præceptor Germamiæ}'' were deposited beneath the castle church of Wittenberg alongside of Luther's: united in life, they sleep together in death till the morning of the resurrection to everlasting life.

LUTHERANS AND PHILIPPISTS.

The differences between Luther the second and Melanchthon the second, if we may use this term, divided the theologians of the Augsburg Confession into two hostile armies.

The rigid Lutheran party was led by Amsdorf, Flacius, Wigand, Gallus, Judex, Mörlin, Heshus, Timann, and Westphal, and had its headquarters first at Magdeburg, then at the University of Jena, and at last in Wittenberg (after 1574). They held fast with unswerving fidelity to the anti-papal and anti-Zwinglian Luther, as representing the ultimate form of sound orthodoxy. They swore by the letter, but had none of the free spirit of their great master.\footnote{Melanchthon applies to them a saying of Polybius, that '\emph{volentes videri similes magnis viris,}' and being unable to imitate the works (\textgreek{ἔργα}) of Luther, they imitated his by-works (\textgreek{πάρεργα}), '\emph{et producunt in theatrum stultitiam suam.}' Calvin more severely but not unjustly remarks (in his second defense against Westphal, 1556): '\emph{O Luthere, quam paucos tuæ præstantiæ imitatores, quam multas vero sanctæ: tuæ jactantiæ simias reliquisti!}' See Gieseler, Vol. IV. p. 435, and especially Planck, Vol. IV. pp. 79 sqq.} They outluthered Luther, made a virtue of his weakness, constructed his polemic extravagances into dogmas, and contracted the catholic expansiveness of the Reformation into sectarian exclusiveness. They denounced every compromise with Rome, and every approach to the Reformed communion, as cowardly treachery to the cause of evangelical truth.

Among these Lutherans, however, we must distinguish three classes—the older friends of Luther (Jonas, his colleague, and Amsdorf, whom he had consecrated Bishop of Naumburg 'without suet or grease or coals'), the younger and stormy generation headed by Flacius, and the milder framers of the 'Form of Concord' (Andreæ, Chemnitz, Selnecker, and Chytræus), who stood mediating between ultra-Lutheranism and Melanchthonianism.

The Melanchthonians, nicknamed Philppists and Crypto-Calvinists,\footnote{\par The term \emph{Philippists} (from the Christian name of Melanchthon, who was usually called Dr. Philippus) is wider, and embraced the Synergists, while the term \emph{Crypto-Calvinists} applies properly only to those who secretly held the Calvinistic doctrine, on the eucharist, but not on predestination. Some of the strict Lutherans—as Flacius, Amsdorf, and Heshus—held fast to the original views of Luther and Melanchthon on predestination, and taught that man was purely passive and even repugnant (\emph{repugnative}) in the work of conversion. Comp. Landerer in Herzog, Vol. XI. p. 538.} prominent among whom were Camerarius, Bugenhagen, Eber, Crell, Major, Cruciger, Strigel, Pfeffinger, Peucer (physician of the Elector of Saxony, and Melanchthon's son-in-law), had their stronghold in the Universities of Wittenberg and Leipzig (till 1574), and maintained, with less force of will and conviction, but with more liberality and catholicity of spirit, the right of progressive development in theology, and sought to enlarge the doctrinal basis of Lutheranism for a final reconciliation of Christendom, or at least for a union of the evangelical churches.\footnote{Kahnis (Vol. II. p. 520) thus characterizes the two parties: '\emph{Dort} [among the strict Lutherans] \emph{das Princip des Festhaltens, hier} [among the Philippists] \emph{das Princip des Fortschreitens; dort scharfe Ausschliesslichkeit, hier Weite, Milde, Vermittelung, Union; dort fertige, faste Doctrin, hier praktische Elasticität.}'}

Both parties maintained the supreme authority of the Bible, but the Lutherans went with the Bible as understood by Luther, the Philippists with the Bible as explained by Melanchthon; with the additional difference that the former looked up to Luther as an almost inspired apostle, and believed in his interpretation as final, while the latter revered Melanchthon simply as a great teacher, and reserved a larger margin for reason and freedom.\footnote{In the Preface to the \emph{Magdeburg Confession}, 1550, Luther is called 'the third Elijah,' 'the prophet of God,' and Luther's doctrine, without any qualification, 'the doctrine of Christ.' See Heppe: \emph{Die Entstehung and Fortbildung des Lutherthums,} pp. 42, 43. In the \emph{Reussische Confession} of 1567 (Heppe, p. 76) it is said: 'We quote chiefly the writings of Luther as our \emph{prophet} (\emph{als unseres Propheten}), and prefer them to the writings of Philippus and others, who are merely \emph{children} of the prophet (\emph{Prophetenkinder}) and his disciples.' The overestimate of Luther is well expressed in the lines—  \par \emph{'Gottes Wort und Luther's Lehr,} \par \emph{Vergehet nun und nimmermehr.'}}

Both parties set forth new confessions of faith and bulky collections of doctrine (\emph{Corpora Doctrinæ}), which were clothed with symbolical authority in different territories, and increased the confusion and intensified the antagonism.\footnote{Prof. Heppe, in his \emph{Die Entstehung und Fortbildung des Lutherthums und die kirchlichen Bekenntniss-Schriften desselben von} 1548–1576 (Cassel, 1863), gives extracts from twenty Luthern Confessions which appeared during this period of twenty-eight years.}

THE THEOLOGICAL CONTROVERSIES IN THE LUTHERAN CHURCH.

The controversies which preceded the composition of the 'Form of Concord,' centred in the soteriological doctrines of the Reformation, concerning sin and grace, justification by faith, and the use of good works, but they extended also to the eucharist and the person and work of Christ. We notice them in the order of the 'Form of Concord.'

I. THE FLACIAN CONTROVERSY ON ORIGINAL SIN, 1560–1580.\footnote{\emph{Disputatio de originali peccato et libero arbitrio inter}  Matthiam Flacium Illyricum  \emph{et}  Victorinum Strigelium,  1563; Flacius:  \emph{De peccato orig.,} in the second part of his \emph{Clavis Scripturæ Sacræ}, 1567; Til. Heshusius:  \emph{Antidoton contra impium et blasphemum dogma M. Fl. III.} 1572, 3d ed. 1579; Wigand:  \emph{De Manichæismo renovato}, 1587; Schlüsselburg:  \emph{Cat. hær.} 1597, Lib. II.; Planck, Vol. V. pp. 1, 285; Döllinger:  \emph{Die Reformation}, etc. Vol. III. (1848), p. 484; Ed. Schmid:  \emph{Des Flacius Erbsündestreit,} in Niedner's \emph{Zeitschrift für hist. Theol.} 1849, Nos. I. and II.;  Frank:  \emph{Die Theologie der Concordienformel,} Vol. I. p. 60; Dorner,  p. 361, and the monograph of Preger on \emph{Flacius and his Age.} Vol. II. p. 310.}

This controversy involved the question whether original sin is essential or accidental—in other words, whether it is the nature of man itself or merely a corruption of nature. It arose, in close connection with the Synergistic controversy, from a colloquy at Weimar between Flacius and Strigel (1560), extended from Saxony as far as Austria, and continued till the death of Flacius (1575), and even after the completion of the 'Form of Concord.'\footnote{About forty adherents of Flacius, driven to German Austria (Opitz, Irenæus, Cölestin, etc.), issued in 1581 a declaration against the 'Form of Concord,' as inconsistent with Luther's pure doctrine on original sin; but in 1582 they fell out among themselves. As late as 1604 there were large numbers of Flacianists in German Austria. Döllinger, Vol. III. p. 492 sq.}

Matthias Flacius Illyricus, the impetuous and belligerent champion of rigid Lutheranism, a man of vast learning, untiring zeal, unyielding firmness, and fanatical intolerance, renewed apparently the Manichean heresy, and thereby ruined himself.\footnote{This remarkable man, born 1520, at Albona, Istria (in Illyria, hence called \emph{Illyricus}), was a convert from Romanism; studied at Basle, Tübingen, and Wittenberg under Luther and Melanchthon, and became Professor of Hebrew in the University of Wittenberg. Luther attended his wedding, and raised him from a state of mental depression almost bordering on despair. In consequence of his opposition to the Augsburg and Leipzig Interim, Flacius removed to Magdeburg (April, 1549), where he opened his literary batteries against Melanchthon and the Interim, and undertook with several others the first Protestant Church history, under the title of 'The Magdeburg Centuries.' In 1557 he was elected Professor in the newly founded University of Jena, but was deposed (1562), persecuted, and forsaken even by his former friends. He spent the remainder of his life in poverty and exile at Ratisbon, Antwerp, Strasburg, and died in a hospital in Frankfort-on-the-Main, March 11, 1575. Many of his contemporaries, and the learned historian Planck, represent him merely as a violent, pugnacious, obstinate fanatic; but more recently his virtues and merits have been better appreciated by Twesten (\emph{Matthias Flacius Illyricus}, Berlin, 1844), Kling (who calls him one of those witnesses of whom the world was not worthy, in Herzog. Vol. IV. p. 410), and W. Preger (\emph{M. Fl. Illyr. und seine Zeit}, Erlangen. 1859–61, 2 vols.). Heppe, from his Melanchthonian standpoint, judges him more unfavorably, and thus characterizes him (in his \emph{Confessionelle Entwicklung}, etc., p. 138): '\emph{M. Flac. Illyricus war ein fanatischer Verehrer Luther's, der von allen Parteigenossen durch Kraft, Consequenz, Klarheit und Sicherheit seiner theologischen Speculation und durch Energie des Willens wie des Denkens hervorragend, kein Opfer und kein Mittel—auch nicht den schändlichsten Verrath am Vertrauen Melanchthon's—scheute, um sein klar erkanntes Ziel, nämlich die, Vernichtung Melanchthon's and der bisherigen Tradition des Protestantisimus zu erreichen und dem Bekenntniss der Kirche einen ganz anderen Charakter aufzuprägen als der war, in dem es sich bisher entwickelt hatte.}' The library of the Union Theological Seminary, New York, possesses a rare collection of the numerous polemical tracts of Flacius. He has undoubted merits in Church history and exegesis. His best works, besides the 'Magdeburg Centuries,' are his \emph{Catalogus testium veritatis}, Basil. 1556, and his \emph{Clavis Scripturæ Sacræ}, 2 P. Basil. 1567.} From an over-intense conviction of total depravity, he represented original sin as the very substance or essence of the natural man, who after the fall ceased to be in any sense the image of God, and became the very image of Satan. He made, however, a distinction between two substances in man—a physical and ethical—and did not mean to teach an evil matter in the sense of Gnostic and Manichean dualism, but simply an entire moral corruption of the moral nature, which must be replaced by a new and holy nature. He departed not so much from the original Protestant doctrine of sin as from the usual conception of the Aristotelian terms \emph{substance} and \emph{accidens.}\footnote{By \textgreek{τὸ συμβεβηκός} Aristotle means a separable property or quality, which does not essentially belong to a thing. In this sense Flacius denied the accidental character of sin, and maintained that it entered into the inmost constitution, just as holiness is inherent and essential in the regenerate.} He quoted many strong passages from Luther, but he found little favor and bitter opposition even among his friends, and was deposed and exiled with forty-seven adherents. The chief argument against him was the alternative that his doctrine either makes Satan the creator of man, or God the author and preserver of sin.

II. THE SYNERGISTIC CONTROVERSY (1550—1567).\footnote{For fuller information, see Pfeffinger:  \emph{Proposit. de libero arbitrio}, 1555; Flacius;  \emph{De orig. peccato et libero arbitrio}, two disputations, 1558 and 1559; Schüsselburg:  \emph{Catal. Hæret.} 1598 (Lib. V. \emph{de Synergistis}); Planck,  Vol. IV. p. 553; Galle,  p. 326; Döllinger,  Vol. III. p. 437; Gust. Frank:  \emph{Gesch. der Prot. Theol.} Vol. I. p. 125, and his art. \emph{Synergismus} in Herzog, Vol. XV. p. 326; Fr. H. R. Frank:  \emph{Theol. der Conc. F.} Vol. I. p. 113; Dorner,  p. 361; and also the literature on the Flacian controversy, especially Schmid and Preger (quoted p. 268).}

It extended over the difficult subject of man's \emph{freedom} and his relation to the \emph{converting} grace of God. It was a conflict between the original Augustinianism of the Reformers and the later Melanchthonian Synergism, or a refined evangelical modification of semi-Pelagianism.\footnote{See above, p. 262.}

Pfeffinger, Professor in Leipzig, who opened the controversy by an academic dissertation (1550), and then wrote a book on the freedom of the will (1555), Major, Eber, and Crell, in Wittenberg, and Victorin Strigel, in Jena, advocated a limited freedom in fallen man, as a rational and responsible being, namely, the power of accepting the prevenient grace of God,\footnote{'\emph{Facultas se applicandi ad gratiam.}'} with the corresponding power of rejecting it. They accordingly assigned to man a certain though very small share in the work of conversion, which Pfeffinger illustrated by the contribution of a penny towards the discharge of a very large debt.

Amsdorf, Flacius, Wigand, and Heshusius, on the other hand, appealing to the teaching of Luther,\footnote{Especially his book \emph{de servo arbitrio.} Luther calls the \emph{voluntas} of the natural man \emph{noluntas,} and compares him to the column of salt, Lot's wife, a block and stone. Similar terms are used in the 'Form of Concord.'} maintained that man, being totally corrupt, can by nature only resist the Spirit of God, and is converted against and in spite of his perverse will, or must receive a new will before he can accept. God converts a man as the potter moulds the clay, as the sculptor carves a statue of wood or stone. They also advocated, as a logical consequence, Luther's original theory of an unconditional predestination and reprobation. But the 'Form of Concord' rejected it as well as Synergism, without attempting to solve the difficulty.

Both parties erred in not making a proper distinction between regeneration and conversion, and between receptive and spontaneous activity. In regeneration, man is passive, in conversion he is active in turning to God, but in response to the preceding action of divine grace, which Augustine calls the \emph{gratia præveniens.} Conversion certainly is not a compulsory or magical, but an ethical process. God operates upon man, not as upon a machine or a dead stone (as Flacius and also the 'Form of Concord' maintain), but as a responsible, rational, moral, and religiously susceptible though very corrupt being; breaking his natural hostility, making willing the unwilling, and preparing him at every step for corresponding action. So far Melanchthon was right. But the defect of the Synergistic theory is the idea of a partnership between God and man, and a corresponding division of work and merit. Synergism is less objectionable than semi-Pelagianisrn, for it reduces co-operation before conversion to a minimum, but even that minimum is incompatible with the absolute dependence of man on God.

III. THE OSIANDRIC CONTROVERSY (1549–1566).\footnote{Osiander: \emph{Disputationes duæ: una de Lege et Evangelio} (1549), \emph{altera de Justifications} (1550), Regiom. 1550; \emph{De unico Mediatore Jes. Chr. et Justificatione fidei confessio A. Osiandri}, Regiom, 1551; \emph{Schmeckbier}, Königsberg, 1552; \emph{Widerlegung der Antwort Melanchthon's}, 1552. Anton Otto Herzberger:  \emph{Wider die tiefgesuchten und scharfgespitzten, aber doch nichtigen Ursachen Osianders,} Magdeburg, 1552; Gallus:  \emph{Probe des Geistes Osiandri}, Magdeb. 1552; Menius:  \emph{Die Gerechtigkeit, die für Gott gilt, wider die neue alcumistische Theologia Osianders,} Erfurt, 1552; Jo. Wigand:  \emph{De Osiandrismo}, Jena, 1583 and 1586; Schlüsselburg:  \emph{Catal. Hæret.} Lib. VI.; Planck,  Vol. IV. p. 249; Baur:  \emph{Disqu. in Osiandri de justif. doctrinam.} Tüb. 1831; Lehnerdt:  \emph{De Osiandri vita et doctr.} Berol. 1835; H. Wilken:  \emph{Osianders Leben, }Stralsund, 1844; Heberle:  \emph{Os. Lehre in ihrer frühsten Gestalt} (\emph{Studien u. Kritiken}, 1844, p. 386); Ritschl:  \emph{Rechtfertigungslehre des A. Os.} (in \emph{Jahrb. für D. Theol.} 1857, p. 795); R. T. Grau:  \emph{De Os. doctrina}, Marb. 1860; Gieseler, Vol. IV. p. 469; Gass, Vol. I. p. 61; Heppe, Vol. I. p. 81; G. Frank, Vol. I. p. 150; J. H. R. Frank, Vol. II. p. 1–47; Dorner, p. 344. Among Roman Catholic divines, Döllinger in his \emph{Reformation, ihre Entwicklung and Wirkungen,} Vol. III. pp. 397–437, gives the best account of the Osiandric controversy.}

It touched the central doctrine of Evangelical Lutheranism, \emph{justification by faith}, whether it is a mere declaratory, forensic art of acquittal from sin and guilt, or an actual infusion of righteousness.

Luther and the other Reformers made a clear distinction between justification as an external act of God \emph{for} man, and sanctification as an internal act of God \emph{in} man; and yet viewed them as inseparable, sanctification being the necessary effect of justification. Faith was to them an appropriation of the whole Christ, a bond of vital union with his person first, and in consequence of this a participation of his benefits.\footnote{See Köstlin:  \emph{Luther's Theologie}, Vol. II. pp. 444 sqq.}

In the Osiandric controversy, justification and sanctification were either confounded or too abstractedly separated, and the person of Christ was lost sight of in his work or in one of his two natures.

Andrew Osiander (1498–1552), an eminent Lutheran minister and reformer at Nuremberg (since 1522), afterwards Professor at Königsberg (1549), a man of great learning and speculative talent, but conceited and overbearing, created a great commotion by a new doctrine of justification, which he brought out after the death of Luther.\footnote{He thought that 'after the death of the lion he could easily dispose of the hares and foxes.' But the germ of his doctrine was already in his tract, '\emph{Ein gut Unterricht und getreuer Rathschlag aus heil. göttlicher Schrift,}' 1524. At the Diet of Augsburg, 1530, he requested Melanchthon, in the presence of Brentius and Urban Regius, to introduce into the new confession of faith the passage Jer. xxiii. 6, 'The Lord our Righteousness,' which he understood to mean that Christ dwells in us by faith, and works in us both to will and to do. See Wilkens, p. 37; Döllinger, p. 398.} He assailed the forensic conception of justification, and taught instead a medicinal and creative act, whereby the sinner is \emph{made} just by an \emph{infusion} of the \emph{divine nature} of Christ, which is our righteousness. This view was denounced as Romanizing, but it is rather mystical. He did not make justification a gradual process, like the Roman system, but a single and complete act, by which Christ according to his divine nature enters the soul of man through the door of faith.\footnote{'\emph{Christus secundum suam veram divinam essentiam in vere credentibus habitat.}'} He meant justification by faith alone without works, but an effective internal justification in the etymological sense of the term. He was Protestant in this also, that he excluded human merit and represented faith which apprehends Christ, as the gift of God. In connection with this he held peculiar views on the image of God, which he made to consist in the essential union of the human nature with the divine nature, and on the necessity of the incarnation, which in his opinion would have taken place even without the fall, in order that through Christ's humanity we might become partakers of the essential righteousness of God.\footnote{'\emph{Per humanitatem devenit in nos divinitas.}'} He appealed to Luther, but denounced Melanchthon as a heretic and pestilential man.

Osiander was protected by Duke Albrecht of Prussia, whom he had converted, but opposed from every quarter by Mörlin, Staphylus, Stancarus, Melanchthon, Amsdorf, Menius, Flacius, Chemnitz. Between the two parties stood the Swabian divines Brentius and Binder. The controversy was carried on with a good deal of misunderstanding, and with such violence that the Professors in Königsberg carried fire-arms into their academic sessions. It was seriously circulated and believed that the devil wrote Osiander's books, while he enjoyed his meals.

After Osiander's death (1552), his son-in-law, John Funck, chaplain of the Duke, became the leader of his small party; but he was executed on the scaffold (1566) as a heretic and disturber of the public peace. Mörlin was recalled from exile and made Bishop of Samland. The Prussian collection of Confessions (\emph{Corpus Doctrinæ Pruthenicum}, or \emph{Borussicum}, Königsberg, 1567) condemned the doctrines of Osiander.

In close connection with the Osiandric controversy on justification was the  Stancarian dispute, introduced by Francesco Stancaro (or Stancarus), an Italian ex-priest, and for a short time Professor in Königsberg (d. 1574 in Poland). He asserted, against Osiander and in agreement with Peter the Lombard, that Christ was our Mediator and Redeemer according to his \emph{human} nature only (since he, being God himself, could not mediate between God and God).\footnote{'\emph{Nemo potest esse mediator sui ipsius.}' Petrus Lombardus says: '\emph{Christus mediator dicitur secundum humanitatem, non secundum divinitatem.}'} He called his opponents and all the Reformers ignoramuses.\footnote{Wigand:  \emph{De Stancarismo,} Lips. 1583; Schlüsselburg, Lib. IX.; Planck, Vol. IV. p. 449; Gieseler, Vol. IV. p. 480; G. Frank, Vol. I. p. 156.}

Another collateral controversy, concerning the \emph{obedience} of Christ, was raised, A.D. 1563, by Parsimonius, or Karg, a Lutheran minister in Bavaria.\footnote{Georg Karg was born 1512, studied at Wittenberg, was ordained by Luther and Melanchthon, became pastor at Oettingen, afterwards at Ansbach, and died 1576. He was a rigid Lutheran in the Interimistic controversies, but otherwise more a follower of Melanchthon.} He derived our redemption entirely from our Lord's passive obedience, and denied that his active obedience had any vicarious merit, since Christ himself, as man, owed active obedience to God. He also opposed the doctrine of imputation, and resolved justification into the idea of remission of sins.

Karg was opposed by Ketzmann in Ansbach, by Heshusius, and the Wittenberg divines. Left without sympathy, and threatened with deposition and exile, he recanted his theses in 1570, and confessed that the obedience of Christ, his righteousness, merit, and innocence are the ground of our justification and our greatest comfort.\footnote{Thomasius:  \emph{Hist. dogmatis de obedientia Christi activa}, Erl. 1845–46; G. Frank, Vol. I. p. 158; Dorner, p. 345; Döllinger, Vol. III. pp. 564–74 (together with the acts from MS. sources in the Appendix, pp. 15 sqq., the best account). Karg's view was afterwards defended by the Reformed divines John Piscator of Herborn and John Camero of Saumur, perhaps also by Ursinus (according to a letter of Tossanus to Piscator). See Döllinger, Vol. III. p. 573; Schweizer: \emph{Centraldogmen}, Vol. II. p. 16.}

The 'Form of Concord' teaches that Christ as God and man in his one, whole, and perfect obedience, is our righteousness, and that his whole obedience unto death is imputed to us.

IV. THE MAJORISTIC CONTROVERSY (1552–1577.)\footnote{D. G. Major:  \emph{Opera}, Viteb. 1569, 3 vols.; N. von Amsdorf:  \emph{Dass die Propositio:} '\emph{Gute Werke sind zur Seligkeit schädlich},' \emph{eine rechte wahre christliche Propositio sei, durch die heiligen Paulas und Luther gepredigt,} 1559; several tracts of Flacius,  Wigand,  and Responsa and Letters of Melanchthon on this subject from 1553 to 1559, in \emph{Corp. Reform.} Vols. VIII. and IX.; Schlüsselburg, Lib. VII.; Planck, Vol. IV. p. 469; Döllinger, Vol. III. p. 493; Thomasius:  \emph{Das Bek. der ev. luth. Kirche in der Consequenz seines Princips}, p. 100; Heppe, Vol. II. p. 264; G. Frank, Vol. I. p. 122; Fr. H. R. Frank, Vol. II. p. 149; Herzog, Vol. VIII. p. 733; Dorner, p. 339.}

It is closely connected with the Synergistic, Osiandric, and Antinomian controversies, and refers to the use of \emph{good works.}

The Reformers derived salvation solely from the merits of Christ through the medium of faith, as the organ of reception, in accordance with the Scripture, 'Believe in the Lord Jesus Christ, and thou shalt be saved.' But faith was to them a work of God, a living apprehension of Christ, and the fruitful parent of good works. Luther calls faith a 'lively, busy, mighty thing,' which can no more be separated from love than fire from heat and light.\footnote{See his classical description of faith in the Preface to the Epistle to the Romans (Walch, Vol. XIV. p. 114, quoted also in the 'Form of Concord,' p. 626, ed. Müller): '\emph{Der Glaube ist ein göttlich Werk in uns, das uns verwandelt und neu gebiert aus Gott und tödtet den alten Adam, macht uns ganz andere Menschen . . . und bringet den heiligen Geist mit sich. O! es ist ein lebendig, geschäftig, thätig, mächtig Ding um den Glauben, dass es unmöglich ist, dass er nicht ohne Unterlass sollte Gutes wirken; er fragt auch nicht, ob gute Werke zu thun sind, sondern ehe man fragt, hat er sie gethan, und ist immer im Thun. Weraber nicht solche Werke thut, der ist ein glaubloser Mensch. . . . Werke vom Glauben scheiden is so unmöglich als brennen und leuchten vom Feuer mag geschieden werden.}' In another place Luther says: '\emph{So wenig das Feuer ohne Hitze und Rauch ist, so wenig ist der Glaube ohne Liebe.}'} Melanchthon, in his later period, laid greater stress on good works, and taught their necessity as fruits of faith, but not as a condition of salvation, which is a free, unmerited gift of God.\footnote{\emph{Loci theol.} ed. 1535 (the edition dedicated to King Henry VIII.): '\emph{Obedientia nostra, hoc est, justitia bonæ conscientiæ seu operum, quæ Deus nobis præcipit, necessario sequi debet reconciliationem. . . . Si vis in vitam ingredi, serva mandata} (Matt. xix. 17). . . . \emph{Justificamur ut nova et spirituali vita vivamus. . . . Ipsius opus sumus, conditi ad bona opera} (Eph. ii. 10). . . . \emph{Acceptatio ad vitam æternam seu donatio vitæ æternæ conjuncta est cum justificatione, i.e., cum remissione peccatorum et reconciliatione, quæ fide contingit. . . . Itaque non datur vita æterna propter dignitatem bonorum operum, sed gratis propter Christum. Et tamen bona opera ita necessaria sunt ad vitam æternam, quia sequi reconciliationem necessario debent}' (\emph{Corp. Reform.} Vol. XXI. p. 429).}

Georg Major (Professor at Wittenberg since 1539, died 1574), a pupil of Melanchthon, and one of the framers of the Leipzig Interim, declared during his sojourn at Eisleben (1552) that good works are \emph{necessary }to salvation.\footnote{'\emph{Bona opera necessaria esse ad salutem.}'} He pronounced the anathema on every one who taught otherwise, though he were an angel from heaven. He meant, however, the necessity of good works as a negative condition, not as a meritorious cause, and he made, moreover, a distinction between salvation and justification.\footnote{He found it necessary afterwards to qualify his proposition, especially since Melanchthon, to his surprise, did not quite approve it. He assigned to good works a \emph{necessitas debiti,} as commanded by God, a \emph{necessitas conjunctionis,} as connected with faith, but no \emph{necessitas meriti.} Our whole confidence is in Christ. '\emph{Hominem,}' he said, '\emph{sola fide esse justum, sed non sola fide salvum.}'}

This proposition seemed to be inconsistent with Luther's solifidianism, and was all the more obnoxious for its resemblance to a clause in the Romanizing Leipzig Interim (1548).\footnote{Viz., the words, '\emph{Es ist gewisslich wahr, dass die Tugenden Glaube, Liebe, Hoffnung, und andere in uns sein müssen und zur Seligkeit nöthig seien.}' In Pezel's edition of Melanchthon's '\emph{Bedenken}' the words \emph{zur Seligkeit} are omitted. Döllinger, Vol. III. p. 496.}

Hence it was violently opposed from every direction. Nicolas von Amsdorf (1483–1565), appealing to St. Paul and Dr. Luther, condemned it as 'the worst and most pernicious heresy,' and boldly advocated even the counter-proposition, that good works are \emph{dangerous }to salvation (1559).\footnote{'\emph{Bona opera perniciosa} (\emph{noxia}) \emph{esse} [not in themselves, but] \emph{ad salutem.}' Whoever held the opposite view was denounced by Amsdorf as a \emph{Pelagianer, Mameluk, zweifältiger Papist} and \emph{Verläugner Christi.}} Flacius denounced Major's view as popish, godless, and most dangerous, because it destroyed the sinner's comfort on the death-bed and the gallows, made the salvation of children impossible, confounded the gospel with the law, and weakened the power of Christ's death.\footnote{See the extracts from Flacius, in Döllinger, Vol. III. pp. 503 sqq.} Wigand objected that the error of the necessity of good works was already condemned by the Apostles in Jerusalem (Acts xv.), that it was the pillar of popery and a mark of Antichrist, and that it led many dying persons unable to find good works in themselves, to despair. Justus Menius, Superintendent of Gotha, tried to mediate by asserting the necessity of good works for the \emph{preservation }of faith; but this was decidedly rejected as indirectly amounting to the same error. A synod, held at Eisenach in 1556, decided in seven theses that Major's proposition was true only \emph{in abstracto} and \emph{in foro legis,} but not \emph{in foro evangelii,} and should be avoided as liable to be misunderstood in a popish sense. Christ delivered us from the curse of the law, and faith alone is necessary both for justification and salvation, which are identical.\footnote{See the theses in Döllinger, Vol. III. p. 511 sq.} The theses were subscribed by Amsdorf, Strigel, Mörlin, Hugel, Stössel, and even by Menius (although the fifth was directed against him). But now there arose a controversy on the admission of the \emph{abstract }and \emph{legal} necessity of good works, which was defended by Flacius, Wigand, and Mörlin; opposed by Amsdorf and Aurifaber as semi-popish. The former view prevailed.

Melanchthon felt that the necessity of good works for salvation might imply their meritoriousness, and hence proposed to drop the words \emph{for salvation}, and to be contented with the assertion that good works are necessary because God commanded them, and man is bound to obey his Creator.\footnote{See his brief \emph{Judicium} on the Majoristic controversy, 1553, \emph{Corp. Reform.} Vol. VIII. p. 194, and his more lengthy German letter \emph{ad Senatum Northusanum} (Nordhausen), Jan. 13, 1555; \emph{Ibid.}, pp. 410–413. '\emph{Diese Deutung, }'he says (p. 412), '\emph{ist zu fliehen: gute Werke sind}  Verdienst  \emph{der Seligkeit; und muss der Glaub und Trost fest allein auf dem Herrn Christo stehen, dass wir gewisslich durch ihn allein, propter eum et per eum, haben Vergebung der Sünden, Zurchnung der Gerechtigkeit, heiligen Geist, und Erbschaft der ewigen Seligkeit. Dieses Fundament ist gewiss. Es folget auch eben aus diesem Fundament, dass diese andere Proposition recht und nöthig ist: gute Werke oder neuer Gehorsam ist nöthig von wegen göttlicher, unwandelbarer Ordnung, dass die vernünftige Creatur Gott Gehorsam schuldig ist, und dazu erschaffen, und jetzund wiedergeboren ist, dass sie ihm gleichförmig werde.}' Melanchthon heard from an Englishman that this controversy created great astonishment in England, where no one doubted the necessity of good works to salvation, nor failed to see the difference between necessity and merit.} This middle course was adopted by the Wittenberg Professors and by the Diet of Princes at Frankfort (1558), but was rejected by the strict Lutherans.

Major consented (in 1558) no longer to use his phrase, and revoked it in his last will (1570), but he was still assailed, and the Professors at Jena prayed for the conversion of the poor old man (1571) with little hope of success. Flacius prayed that Christ might crush also this serpent. Heshusius publicly confessed that he had committed a horrible sin in accepting the Doctor's degree from Major, who was a disgrace to the theological profession.

The 'Form of Concord' settled the controversy by separating good works both from justification and salvation, yet declaring them necessary as effects of justifying faith.\footnote{In accordance with the word of Augustine: '\emph{Opera sequuntur justificatam, non præcedunt justificandum.}' Three or four of the framers of the 'Form of Concord' were inclined to Major's view, and endeavored at first to prevent its condemnation; but the logic of the Lutheran principle triumphed.}

V. THE ANTINOMIAN CONTROVERSY (1527–1560).\footnote{Luther's  \emph{Werke, }Vol. XX. p. 2014 (ed. Walch); Wigand:  \emph{De antinomia veteri et nova}, Jen. 1571; Schlüsselburg, Lib. IV.; Förstemann:  \emph{Neues Urkundenbuch} (Hamburg, 1842), Vol. I. p. 291; J. G. Schulzius:  \emph{Historia Antinomorum}, Viteb. 1708; Planck, Vol. II. p. 399, Vol. V. I. 1; Thomasius, p. 46; Döllinger, Vol. III. p. 372; Gieseler, Vol. IV. p.397; Heppe, Vol. I. p. 80; Gass, Vol. I. p.57; G. Frank, Vol. I. p. 146; Fr. H. R. Frank, Vol. II. pp. 246, 262; Dorner, p. 336; Elwert:  \emph{De Antinomia Agricolæ Islebii, }Tur. 1836; K. J. Nitzsch:  \emph{Die Gesammterscheinung des Antinomismus}, in the \emph{Studien u. Kritiken}, 1846, Nos. I. and II.}

Protestantism in its joyful enthusiasm for the freedom and all-sufficiency of the gospel was strongly tempted to antinomianism, but restrained by its moral force and the holy character of the gospel itself.\footnote{Gass says (Vol. I. p. 57): '\emph{Die Reformation war selbst Antinomismus, insofern sie mit dem werkheiligen auch das gesetzliche Princip, wenn es die Seligkeit des Menschen bewirken will, verwarf. Melanchthon hatte Gesetz und Evangelium wie Schreck- und Trostmittel einander entgegengestellt und nur auf das letzere die Rechtfertigung gebaut, während er doch unter dem Gesetz den bleibenden Inhalt des göttlichen Willens zusammenfasst.}'} Luther, in opposition to Romish legalism, put the gospel and the law as wide apart as 'heaven and earth,' and said,' Moses is dead.'\footnote{Many of his utterances, as quoted by Döllinger, Vol. III. pp. 45 sqq., sound decidedly antinomian, but must be understood \emph{cum grano salis}, and in connection with his whole teaching. Some of the most objectionable are from his 'Table Talk,' as when he calls Moses 'the master of all hangmen' and 'the worst of heretics.'} Nevertheless he embodied in his Catechism an excellent exposition of the Decalogue before the Creed; and Melanchthon, as we have already seen, laid more and more stress on the moral element and good works in opposition to the abuses of solifidianism and carnal security.

The antinomian controversy has two stages. The first touches the office of the law under the gospel dispensation, and its relation to repentance; the second the necessity of good works, which was the point of dispute between Major and Amsdorf, and has already been discussed.

John Agricola, of Eisleben, misunderstood Luther, as Marcion, the antinomian Gnostic, misunderstood St. Paul.\footnote{Agricola (Schnitter, Kornschneider; Luther called him Grickl) was born at Eisleben, 1492 (hence \emph{Magister Islebius}), and studied at Wittenberg, where he boarded with Luther. He was a popular preacher at Eisleben, and became Professor of Theology at Wittenberg, 1536, and chaplain of Elector Joachim II. at Berlin, 1540. In 1548 he took a leading part in the Augsburg Interim, and denied the essential principles of Protestantism, but protested afterwards from the pulpit against the necessity of good works (1558). He died at Berlin, 1566. Luther was more vexed by him, as he said, than by any pope; he charged him with excessive vanity and ambition, and declared him unfit to teach, and fit only for the profession of a jester (\emph{Briefe}, Vol. V. p. 321). He refused to see him in 1545, and said, '\emph{Grickl wird in alle Ewigkeit Grickl bleiben.}' Bretschneider and Gieseler suppose that Melanchthon incurred Agricola's displeasure by not helping him to a theological chair in Wittenberg. He must have had, however, considerable administrative capacity. Döllinger charges the Reformers with misrepresenting him and his doctrine.} He first uttered antinomian principles in 1527, in opposition to Melanchthon, who in his Articles of Visitation urged the preaching of the law unto repentance.\footnote{'\emph{Prædicatio legis ad pænitentiam.}' \emph{Chursächsische Visitations-Artikel}, 1527 and 1528, Latin and German, ed. by Strobel, 1777.} He was appeased in a conference with the Reformers at Torgau (December, 1527). But when Professor at Wittenberg, he renewed the controversy in 1537, in some arrogant theses, and was defeated by Luther in six public disputations (1538 and 1540). He made a severe attack on Luther, which involved him in a lawsuit, but he removed to Berlin, and sent from there a recantation, Dec. 6, 1540. Long afterwards (1562) he reasserted his views in a published sermon on Luke vii. 37. He was neither clear nor consistent.

Agricola taught with some truth that genuine repentance and remission of sin could only be secured under the gospel by the contemplation of Christ's love. In this Luther (and afterwards Calvin) agreed with him. But he went much further. The law in his opinion was superseded by the gospel, and has nothing to do with repentance and conversion. It works only wrath and death; it leads to unbelief and despair, not to the gospel. He thought the gospel was all-sufficient both for the office of terror and the office of comfort. Luther, on the contrary, maintained, in his disputations, that true repentance consists of two things—knowledge and sorrow of sin, and resolution to lead a better life. The first is produced by the law, the second by the gospel. The law alone would lead to despair and hatred of God; hence the gospel is added to appease and encourage the terrified conscience. The law can not justify, but must nevertheless be taught, that by it the impious may be led to a knowledge of their sin and be humbled, and that the pious may be admonished to crucify their flesh with its sinful lusts, and to guard against security.

The 'Form of Concord' teaches a threefold use of the law: (\emph{a}) A \emph{political} or \emph{civil} use in maintaining outward discipline and order; (\emph{b}) An \emph{elenchtic} or \emph{pedagogic} use in leading men to a knowledge of sin and the need of redemption; (\emph{c}) A \emph{didactic} or \emph{normative} use in regulating the life of the regenerate. The Old and New Testaments are not exclusively related as law and gospel, but the Old contains gospel, and the New is law and gospel complete.

VI. THE CRYPTO-CALVINISTIC OR EUCHARISTIC CONTROVERSY (1549–1574).\footnote{Westphal:  \emph{Farrago confusanearum et inter se dissidentium opiniomum de Cæna Domini ex Sacramentariorum libris congesta}, Magdeb. 1552 (chiefly against Calvin, Bullinger, Peter Martyr, and John à Lasco); \emph{Recta Fides de Cæna Domini ex verbis Ap. Pauli et Evangelistarum demonstrata}, 1553; a tract on \emph{Augustine's} view of the eucharist, 1555; another on \emph{Melanchthon's} view, 1557; then \emph{Justa Defensio} against John à Lasco; and, finally, \emph{Apologia contra corruptelas et calumnias Johannis Calvini}, 1558. Calvin:  \emph{Defensio sanæ et orthodoxæ doctrinæ de sacramentis}, Gen. and Tiguri, 1555; \emph{Secunda Defensio planæ et orthod. de sacram. fidei contra Joach. Westphali calumnias, }1556; \emph{Ultima Admonitio ad Joach. Westphalum}, 1557; \emph{Dilucida Explicatio sanæ doctr. de vera participatione carnis et sanguinis Christi in sacra Cæna}, against Heshusius, 1561. (All these tracts of Calvin in his \emph{Opera}, Vol. IX. ed. Baum, Cunitz, and Reuss, Brunsv. 1870.) Minor eucharistic tracts on the Lutheran side by Brenz,  Schnepf,  Alber,  Timann,  Heshusius;  on the Calvinistic side by Bullinger,  Peter Martyr,  Beza,  and Hardenberg.  Wigand:  \emph{De Sacramentariismo}, Lips. 1584; \emph{De Ubiquitate, Regiom.} 1588; Schlüsselburg, Lib. III.; Planck, Vol. V. II. 1; Galle, p. 436; Ebrard:  \emph{Das Dogma vom heil. Abendmahl}, Vol. II. pp. 525–744; Gieseler, Vol. IV. pp. 439, 454; Heppe, Vol. II. p. 384; Stähelin:  \emph{Calvin}, Vol. II. pp. 112, 198; Schmidt:  \emph{Melanchthon}, pp.580, 639; G. Frank, Vol. I. pp. 132, 164; Fr. H. R. Frank, Vol. III. pp. 1–164; Mönckeberg:  \emph{Joach. Westphal und Joh. Calvin}, 1865; Dorner, p. 400; also Art. \emph{Kryptocalvinismus} in Herzog, Vol. VIII. p. 122; and the \emph{Prolegomena} to the ninth volume of the new edition of Calvin's \emph{Opera} (in \emph{Corp. Reform.}).}

The eucharistic controversy between Luther and Zwingli, although it alienated the German and Swiss branches of the Reformation, did not destroy all intercourse, nor discourage new attempts at reconciliation. Calvin's theory, which took a middle course, retaining, on the basis of Zwingli's exegesis, the religious substance of Luther's faith, and giving it a more intellectual and spiritual form, triumphed in Switzerland, gained much favor in Germany, and opened a fair prospect for union. But the controversy of Westphal against Calvin, and the subsequent overthrow of Melanchthonianism, completed and consolidated the separation of the two Confessions.

Melanchthon's later view of the Lord's Supper, which essentially agreed with that of Calvin, was for a number of years entertained by the majority of Lutheran divines even at Wittenberg and Leipzig, and at the court of the Elector of Saxony. It was also in various ways officially recognized with the Augsburg Confession of 1540, which was long regarded as an improved rather than an altered edition.

But the Princes and the people held fast to the heroic name of Luther against any rival authority, and when the alternative was presented to choose between him and Melanchthon or Calvin, the issue could not be doubtful. Besides, the old traditional view of the mysterious power and magical efficacy of the sacraments had a firm hold upon the minds and hearts of German Christians, as it has to this day.

Joachim Westphal, a rigid Lutheran minister at Hamburg, renewed, in 1552, the sacramental war in several tracts against the 'Zurich Consensus' (issued 1549), and against Calvin and Peter Martyr; aiming indirectly against the Philippists, and treating all as sacramentarians and heretics who denied the corporeal presence, the oral manducation, and the literal eating of Christ's body even by unbelievers. He made no distinction between Calvin and Zwingli, spoke of their godless perversion of the Scriptures, and even their satanic blasphemies. About the same time John à Lasco, a Polish nobleman and minister of a foreign Reformed congregation in London, and one hundred and seventy-five Protestants, who were driven from England under the bloody Mary (1553), sought and were refused in cold winter a temporary refuge in Denmark, Rostock, Lübeck, and Hamburg (though they found it at last in East Friesland). Westphal denounced them as martyrs of the devil, enraged the people against them, and gloried in this cruelty as an act of faith.\footnote{See Utenhoven's \emph{Simplex et fidelis narratio}, etc., Bas. 1560, and the extracts from it by Salig, Vol. II. pp. 1090 sqq., and Ebrard, Vol. II. pp. 536 sqq. Mönckeberg attempts to apologize for Westphal, but without effect. Compare the remarks of Dorner, p. 401.}

This intolerance roused the Swiss, who had kept silence for some time, to a defense of their doctrine. Calvin took up his sharp and racy pen, indignantly rebuking 'the no less rude and barbarous than sacrilegious insults' to persecuted members of Christ, and triumphantly vindicating, against misrepresentations and objections, his doctrine of the spiritual real presence of Christ, and the sealing communication of the life-giving virtue of his body in heaven to the believer through the power of the Holy Ghost.\footnote{'\emph{Fatemur},' he says in his \emph{First Defense}, '\emph{Christum, quod panis et vini symbolis figurat, vere præstare, ut animas nostras carnis suæ esu et sanguinis potione alat. . . . Hujus rei non fallacem oculis proponi figuram dicimus, sed pignus nobis porrigi, cui res ipsa et veritas conjuncta est: quod scilicet Christi carne et sanguine animæ nostræ pascantur}' (in the new edition of his \emph{Opera}, Vol. IX. p. 30). In the \emph{Second Defense}: '\emph{Christum corpore absentem doceo nihilominus non tantum divina sua virtute, quæ ubique diffusa est, nobis adesse, sed etiam facere ut nobis vivifica. sit sua caro} (Vol. IX. p. 76). . . . \emph{Cænam plus centies dici sacrum esse vinculum nostræ cum Christo unitatis} (p. 77). . . . \emph{Spiritus sui virtute Christus locorum distantiam superat ad vitam nobis e sua carne inspirandam}' (p. 77). . . . And in his \emph{Last Admonition}: '\emph{Hæc nostræ doctrinæ summa est, carnem Christi panem esse vivificum, quia dum fide in eam coalescimus, vere aninas nostras alit et pascit. Hoc nonnisi spiritualiter fieri docemus, quia hujus sacræ unitatis vinculum arcana est et incomprehensibilis Spiritus Sancti virtus}' (Vol. IX. p. 162).} He claimed to agree with the Augsburg Confession as understood and explained by its author, and appealed to him. Melanchthon, for reasons of prudence and timidity, declined to take an active part in the strife 'on bread-worship,' but never concealed his essential agreement with him.\footnote{He wrote to Calvin, Oct. 14, 1554 (\emph{Corp. Reform.} Vol. VIII. p. 362): '\emph{Quod in proximis literis hortaris, ut reprimam ineruditos clamores illorum, qui renovant certamen } \textgreek{περὶ ἀρτολατρείας, } \emph{scito, quosdam p\&acelig;cipue odio mei eam disputationem movere, ut habeant plausibilem causam ad me opprimendum.}' To Hardenberg, in Bremen, May 9, 1557: '\emph{Crescit, ut vides, non modo certamen, sed etiam rabies in scriptoribus, qui } \textgreek{ἀρτολάτρειαν } \emph{stabiliunt.}' And to Mordeisen, Nov. 15, 1557 (\emph{Corp. Reform.} Vol. IX. p. 374): '\emph{Si mihi concedatis, ut in alia loco vivam, respondebo illis indoctis sycophantis et vere et graviter, et dicam utilia ecclesiæ.}' He gave, however, his views pretty clearly and dispassionately shortly before his death in his \emph{vota} on the Breslau and Heidelberg troubles (1559 and 1560).} His enemies re-published his former views. His followers were now stigmatized as 'Crypto-Calvinists.'

The controversy gradually spread over all Germany, and was conducted with an incredible amount of bigotry and superstition.

In Bremen, John Timann fought for the real presence, and insisted upon the ubiquity of Christ's body as a settled dogma (1555), while Albert Hardenberg opposed it, and was banished (1560); but a reaction took place afterwards in favor of the Reformed Confession.

In Heidelberg, Tilemann Heshusius,\footnote{\par His German name was \emph{Hesshusen.} He was one of the most pugnacious divines of his age; born 1527 at Nieder-Wesel, died 1588 at Helmstädt. See Leuckfeld's biography, \emph{Historia Heshusiana} (1716), and Henke, in Herzog, Vol. VI. p. 49.} General Superintendent since 1558, attacked the Melanchthonian Klebitz openly at the altar by trying to wrest from him the cup. The Elector Frederick III. dismissed both (1559), ordered the preparation of the Heidelberg Catechism, and introduced the Reformed Confession in the Palatinate (1563).

In Würtemberg the ubiquity doctrine triumphed (at a synod in Stuttgart, 1559), chiefly through the influence of Brentius, who had formerly agreed with Melanchthon, but now feared that 'the devil intended through Calvinism to smuggle heathenism, Talmudism, and Mohammedanism into the Church.'\footnote{In his last book against Bullinger (1564). See Hartmann, \emph{Brenz}, p. 252.} A colloquy at Maulbronn (1564) between the Würtemberg and the Palatinate divines on ubiquity led to no result.

Ducal Saxony, under the lead of the Flacianist Professors of Jena, was violently arrayed against Electoral Saxony with the Crypto-Calvinist faculty at Wittenberg. The Elector Augustus, strongly prejudiced against Flacianism, deceived by the \emph{Consensus Dresdensis }(1571), and controlled by his physician, Caspar Peucer, the active and influential lay-leader of the Crypto-Calvinists, unwittingly maintained for some time Calvinism under the disguise of sound Lutheranism. When he became Regent of the Thuringian Principalities (1573), he banished Heshusius and Wigand from Jena, and all the Flacianists of that district.

Thus Philippism triumphed in all Saxony, but it was only for a short season.

Elector Augustus was an enthusiastic admirer of Luther, and would not tolerate a drop of Calvinistic blood in his veins. When he found out the deceptive policy of the Crypto-Calvinists, he suppressed them by force, 1574.\footnote{He was undeceived by a new deception. The crisis was brought about by the discovery of a confidential correspondence with the Reformed in the Palatinate, and especially by the appearance in Leipzig of the anonymous \emph{Exegesis perspicua controversiæ de Cæna Domini, }1574 (newly edited by Scheffer, Marburg, 1853), which openly rejected the \emph{manducatio oralis}, and defended Calvin's view of the eucharist (though without naming him), while the \emph{Consensus Dresdensis} (1571) had concealed it under Lutheran phraseology. This work was generally attributed to Peucer and the Wittenberg Professors, in spite of their steadfast denial, but it was the product of a Silesian physician, Joachim Cureus. See the proof in Heppe, Vol. II. pp. 468 sqq.} The leaders were deposed, imprisoned, and exiled, among them four theological Professors at Wittenberg.\footnote{Cruciger, Moller, Wiedebram, and Pezel (whom the Lutherans called Beelzebub) refused to recant. The first went to Hesse, the second to Hamburg, the other two to Nassau. The old and weak Major yielded to the condemnation of Melanchthon's view. Several other Wittenberg Professors were likewise deposed.} Peucer was confined in prison for twelve years, while his children were wandering about in misery.\footnote{Peucer was released in 1586, at the intercession of the beautiful Princess Agnes Hedwig of Anhalt, and became physician of the Prince of Dessau, where he died, 1602. He wrote the history of his prison life, \emph{Historia carcerum et liberationis divinæ, }ed. by Pezel, Tig. 1605. On his theory of the real presence, see Galle, pp. 460 sqq. He rejected the Lutheran view much more strongly than his father-in-law, Melanchthon, and thought it had no more foundation in the Bible than the popish transubstantiation. Comp. Henke:  \emph{Casp. Peucer und Nic. Crell, }Marburg, 1865.} Thanks were offered in all the churches of Saxony for the triumph of genuine Lutheranism. A memorial coin exhibits the Elector with the sword in one hand, and a balance in the other: one scale bearing the child Jesus; the other, high up, the four Wittenberg Philippists with the devil, and the title 'reason.'

After the death of Augustus (1586), Calvinism again raised its head under Christian I. and the lead of Chancellor Nicolas Crell, but after another change of ruler (1591) it was finally overthrown: the protesting Professors in Wittenberg and Leipzig were deposed and exiled; the leading ministers at Dresden (Salmuth and Pierius) were imprisoned; Crell, who had offended the nobility, after suffering for ten years in prison, was, without an investigation, beheaded as a traitor to his country (Oct. 9, 1601), solemnly protesting his innocence, but forgiving his enemies.\footnote{He was charged with intermeddling in matters of religion, and advising a dangerous treaty with the Reformed Henry IV. of France against Austria. The suit was referred to an Austrian court of appeals at Prague, and decided in the political interest of Austria with a violation of all justice. His confession of guilt before his heavenly Judge was distorted by his fanatical opponents into a confession of guilt before his human judges. It is often stated that he was not beheaded for religion ('\emph{non ob religionem, sed ob perfidiam multiplicem},' as Hutter says, \emph{Concordia concors, }pp. 448 and 1258). But his Calvinism, or rather his Melanchthonianism (for he never read a line of Calvin), was the only crime which could he proved against him; he always acted under the direction and command of the Elector, and he had accepted the chancellorship with a clear confession of his views, and the assurance of his Prince that he should be protected in it, and never be troubled with subscribing to the 'Form of Concord.' As judge, he was admitted, even by his enemies, to have been impartial and just to the poor as well as the rich. Comp. Hasse:  \emph{Ueber den Crell'schen Process}, in Niedner's \emph{Zeitschrift für hist. Theol}. 1848, No. 2; Vogt in Herzog, Vol. III. p. 183; Richard:  \emph{Dr. Nic. Krell. }Dresden, 1859; G. Frank, Vol. I. pp. 296 sqq.; Henke:  \emph{C. Peucer und N. Crell}, Marburg, 1865.} Since that time the name of a Calvinist became more hateful in Saxony than that of a Jew or a Mohammedan.

It is characteristic of the spirit of the age and the doctrine of consubstantiation that they gave rise to all sorts of idle, curious, and unwittingly irreverent speculations about the possible effect of the consecrated elements upon things for which they never were intended. The schoolmen of the Middle Ages, in the interest of transubstantiation, seriously disputed the question whether the eating of the eucharistic bread would kill or sanctify a mouse, or (as the wisest thought) have no effect at all, since the mouse did not receive it \emph{sacramentaliter,} but only \emph{accidentaliter}. Orthodox Lutherans of the sixteenth century went even further. Brentius decidedly favored the opinion that the consecrated bread, if eaten by a mouse, was fully as much the body of Christ as Christ was the Son of God in the mother's womb and on the back of an ass. The sacrament, he admitted, was not intended for animals, but neither was it intended for unbelievers, who nevertheless received the very body and blood of Christ. An eccentric minister in Rostock required the communicants to be shaved to prevent profanation. Licking the blood of Christ from the beard was supposed to be punished with instant death or a monstrous growth of the beard. Sarcerius caused the earth on which a drop of Christ's blood fell, instantly to be dug up and burned. At Hildesheim it was customary to cut off the beard or the piece of a garment which was profaned by a drop of wine; and the Superintendent, Kongius, was expelled from the city, simply because he had taken up from the earth a wafer and given it to a communicant, without first kneeling before it, kissing, and reconsecrating it, as his colleagues thought he should have done. The Lutherans in Ansbach disputed about the question whether the body of Christ were actually swallowed, like other food, and digested in the stomach. When the Rev. John Musculus, in Frankfort-on-the-Oder, inadvertently spilled a little wine at the communion, he was summoned before a Synod, and Elector John Joachim of Brandenburg declared that deposition, prison, and exile were too mild a punishment for such a crime, and that the offender, who had not spared the blood of Christ, must suffer bloody punishment, and have two or three fingers cut off.\footnote{Such details are recorded by Salig, Vol. III. p. 462; Hartmann and Jäger:  \emph{Brenz}, Vol. II. p. 371; Galle:  \emph{Melanchthon}, p. 449 sq.; Ebrard:  \emph{Abendmahl}, Vol. II. pp. 592, 694; Droysen:  \emph{Geschichte der Preuss. Politik}, Vol. II. p. 261; Sudhof: \emph{Olevianus und Ursinus,} p. 239; G. Frank, Vol. I. p. 164.}

There was also a considerable dispute among Lutheran divines about the precise time and duration of the corporeal presence. John Saliger (Beatus) of Lübeck and his friend Fredeland (followers of Flacius, and of his doctrine on original sin) maintained that the bread becomes the body of Christ immediately after the consecration and before the use (\emph{ante usum}), and called those who denied it sacramentarians; while they in turn were charged with the Romish error of transubstantiation. Deposed at Lübeck, Saliger renewed the controversy from the pulpit at Rostock (1568). Chytræus decided that this was a question of idle curiosity rather than piety, and that it was sufficient to attach the blessing of the sacrament to the transaction, without time-splitting distinctions (1569). The usual Lutheran doctrine confines the union of the bread with the body to the time of the use, and hence the term \emph{consubstantiation }was rejected, if thereby be understood a \emph{durabilis inclusio}, or permanent conjunction of the sacramental bread and body of Christ.\footnote{J. Wiggers:  \emph{Der Saligersche Abendmahlsstreit}, in Niedner's \emph{Zeitschrift für hist. Theol.} 1848, No. 4, p. 613.}

VII. THE CHRISTOLOGICAL OR UBIQUITARIAN CONTROVERSY.\footnote{Dorner:  \emph{Entwicklungsgeschichte der Lehre von der Person Christi}, 2d ed. Vol. II. pp. 665 sqq.; Heppe:  \emph{Gesch. des D. Prot.} Vol. II. pp.75 sqq.; G. E. Steitz:  Art. \emph{ Ubiquität,} in Herzog's \emph{Encykl.} Vol. XVI. pp. 558–616, with an addition by Herzog, Vol. XXI. p. 383; Gieseler, Vol. IV. pp. 452, 462; G. Frank, Vol. I. p. 161; Fr. H. R. Frank, Vol. III. pp. 165–396. Comp. also the literature on the eucharistic controversy, p. 279.}

The Lutheran view of the Lord's Supper implies the ubiquity, i.e., the illocal omnipresence, or at all events the multipresence of Christ's body. And this again requires for its support the theory of the \emph{communicatio idiomatum}, or the communication of the attributes of the two natures of Christ, whereby his human nature becomes a partaker of the omnipresence of his divine nature. A considerable amount of interesting speculation was spent on this subject in the sixteenth century.

All Christians believe in the real and abiding omnipresence of Christ's \emph{divine }nature, and of Christ's \emph{person} (which resides in the divine nature or the pre-existing Logos), according to Matt. xxviii. 20; xviii. 20. But the omnipresence of his \emph{human }nature was no article of any creed before the Reformation, and was only held by a few fathers and schoolmen of questionable orthodoxy, as a speculative opinion.\footnote{Origen first taught the ubiquity of the body of Christ, in connection with his docetistic idealism, but without any regard to the eucharist, and was followed by Gregory of Nyssa (\emph{Orat.} 40, and \emph{Adv. Apollinar.} c. 59). They held that Christ's body after the resurrection was so spiritualized and deified as to lay aside all limitations of nature, and to be in all parts of the world as well as in heaven. See Gieseler's \emph{Commentatio qua Clementis Alex. et Origenis doctrinæ de corpore Christi exponuntur}, Gott. 1837, and Neander's \emph{Dogmengeschichte}, Vol. I. pp. 217, 834. Cyril of Alexandria held a similar view (Christ's body is 'every where,' \textgreek{πανταχοῦ}), but in connection with an almost monophysitic Christology. Scotus Erigena revived Origen's ubiquity, gave it a pantheistic turn, and made it subservient to his view of the eucharistic presence, which he regarded merely as a symbol of the every where present Christ. Neander, Vol. II. p. 43.} The prevailing doctrine was that Christ's glorified body, though no more grossly material and sensuous, and not exactly definable in its nature, was still a body, seated on a throne of majesty in heaven, to which it visibly ascended, and from which it will in like manner return to judge the quick and the dead. This was the view even of Gregory Nazianzen and John of Damascus, who otherwise approach very nearly the Lutheran dogma of the \emph{communicatio idiomatum} (the \emph{genus majestaticum}). The mediæval scholastics ascribed \emph{omnipresence} only to the divine nature and the person of Christ, \emph{unipresence }to his human nature in heaven, \emph{multipresence }to his body in the sacrament; but they derived the eucharistic multipresence from the miracle of transubstantiation, and not from an inherent specific quality of the body. Even William Occam (who was inclined to consubstantiation rather than transubstantiation, and had considerable influence upon Luther) ventured only upon the paradox of the hypothetical \emph{possibility }of an absolute ubiquity.

Luther first clearly taught the absolute ubiquity of Christ's body, as a dogmatic support of the real presence in the eucharist.\footnote{On Luther's Christology and ubiquity doctrine, see Heppe (Ref.): \emph{Dogmatik, des D. Protest. im} 16\emph{ten Jahrh.} Vol. II. pp. 93 sqq., and Köstlin (Luth.): \emph{Luther's Theol.} Vol. II. pp. 118, 153, 167, 172, 512. Köstlin, without adopting Luther's views of ubiquity, finds in them '\emph{grossartige, tiefe, geist- und lebensvolle Anschauungen vom göttlichen Sein und Leben}' (Vol. II. p. 154).} He based it exegetically on  Eph. i. 23 ('which is his body, the fullness of him that filleth all in all') and  John iii. 13 ('the Son of \emph{man} who \emph{is} in heaven'), and derived it directly from the personal union of the divine and human natures in Christ (not, as his followers, from the communication of the attributes). He adopted the scholastic distinction of three kinds of presence: 1. \emph{Local} or circumscriptive (material and confined—as water is in the cup); 2. \emph{Definitive} (local, without local inclusion or measurable quantity—as the soul is in the body, Christ's body in the bread, or when it passed through the closed door); 3. \emph{Repletive} (supernatural, divine omnipresence). He ascribed all these to Christ as man, so that in one and the same moment, when he instituted the holy communion, he was \emph{circumscriptive }at the table, \emph{definitive} in the bread and wine, and \emph{repletive }in heaven, i.e., every where.\footnote{In his \emph{Grosse Bekenntniss vom Abendmahl}, published 1528 (in Walch's ed. Vol. XX.; in the Erlangen ed. Vol. XXX.), he says: '\emph{Kann Christus' Leib über Tisch sitzen and dennoch im Brot sein, so kann er auch im Himmel und wo er will sein und dennoch im Brot sein; es ist kein Unterschied fern oder nah bei dem Tische sein, dazu dass er zugleich im Brot sei. . . . es sollte mir ein schlechter Christus bleiben, der nicht mehr, denn an einem einzelnen Orte zugleich eine göttliche and menschliche Person wäre, und an allen anderen Orten müsste er allein ein blosser abgesonderter Gott und göttliche Person sein ohne Menschseit. Nein, Geselle, wo du mir Gott hinsetzest, da must du mir die Menschheit mit hinsetzen. Die lassen sich nicht sondern und von einander trennen; es ist Eine Person worden und scheidet die Menschseit nicht so von sich, wie Meister Hans seinen Rock auszieht und von sich legt, wenn er schlafen geht. Denn, dass ich den Einfältigen ein grob Gleichniss gebe, die Menschheit ist näher vereinigt mit Gott, denn unsere Haut mit unserm Fleische, ja näher denn Leib und Seele.}'} Where God is, there is Christ's humanity, and where Christ's humanity is, there is inseparably joined to it the whole Deity. In connection with this, Luther consistently denied the literal meaning of Christ's ascension to heaven, and understood the right hand of God, at which he sits, to be only a figurative term for the omnipresent power of God (Matt. xxviii. 18).\footnote{He ridicules the popular conception of heaven and the throne of God as childish: '\emph{Die Rechte Gottes},' he says, l.c., '\emph{ist nicht ein sonderlicher Ort, da ein Leib solle oder möge sein, nicht ein Gaukelhimmel, wie man ihn den Kindern pflegt vorzubilden, darin ein gülden Stuhl stehe und Christus neben dem Vater sitze in einer Chorkappen und gülden Krone. . . . Die Rechte Gottes ist an allen Enden, so ist sie gewisslich auch im Brot und Wein über Tische. . . . Wo nun die Rechte Gottes ist, da muss Christi Leib und Blut auch sein; denn die Rechte Gottes ist nicht zu theilen in viele Stücke, sondern ein einiges einfältiges Wesen.}' If this prove any thing, it proves the \emph{absolute} omnipresence of Christ's body. And so Brentius taught.} Here he resorted to a mode of interpretation which he so strongly condemned in Zwingli when applied to the word \emph{is.}

It is very plain that such an absolute omnipresence of the body proves much more than Luther intended or needed for his eucharistic theory; hence he made no further use of it in his later writings, and rested the real presence at last, as he did at first, exclusively on the literal (or rather synecdochical) interpretation of the words, 'This \emph{is} my body.' His earlier Christology was much more natural, and left room for a real development of Christ's humanity.

Melanchthon, in his later period, decidedly opposed the ubiquity of Christ's body, and the introduction of 'scholastic disputations' on this subject into the doctrine of the eucharist. He wished to know only of a \emph{personal} presence of Christ, which does not necessarily involve bodily presence.\footnote{\emph{De inhabitatione Dei in Sanctis ad Osiandrum}, 1551 (\emph{Consil. Lat.} Vol. II. p. 156): '\emph{Tota antiquitas declarans hanc propositionem: Christus est ubique, sic declarat: Christus est ubique}  personaliter.  \emph{Et verissimum est, Filium Dei, Deum et hominem habitare in sanctis. Sed antiquitas hanc propositionem rejicit: Christus}  corporaliter  \emph{est ubique. Quia natura quælibet retinet sua} \textgreek{ ἰδιώματα.} \emph{Unde Augustinus et alii dicunt: Christi corpus est in certo loco. . . Cavendum est, ne ita astruamus divinitatem hominis Christi, ut veritatem corporis auferamus.}' In a new edition of his lectures on the \emph{Colossians} (1556 and 1559), he maintains the literal meaning of the ascension of Christ, 'i.e., \emph{in locum cœlestem. . . . Ascensio fuit visibilis et corporalis, et sæpe ita scripsit tota antiquitas, Christum corporali locatione in aliquo loco esse, ubicunque vult. Corpus localiter alicubi est secundum verum corporis modum, ut Augustinus inquit.}' See Galle, p. 448.} He also rejected the theory of the \emph{communicatio idiomatum} in a real or physical sense, because it leads to a confusion of natures, and admitted with Calvin only a dialectic or verbal communication.\footnote{See on his Christology chiefly Heppe, Vol. II. pp. 99 sqq.} Luther's Christology leaned to the Eutychian confusion, Melanchthon's to the Nestorian separation of the two natures.

The renewal of the eucharistic controversy by Westphal led to a fuller discussion of ubiquity. The orthodox Lutherans insisted upon ubiquity as a necessary result of the real communication of the properties of the two natures in Christ; while the Philippists and Calvinists rejected it as inconsistent with the nature of a body, with the realness of Christ's ascension, and with the general principle that the infinite can not be comprehended or shut up in the finite.\footnote{'\emph{Finitum non capax est infiniti.}'}

The Colloquy at Maulbronn.—These conflicting Christologies met face to face at a Colloquy in the cloister of Maulbronn, in the Duchy of Würtemberg, April 10–15, 1564.\footnote{Both parties published an account—the Lutherans at Frankfort-on-the-Main, the Reformed at Heidelberg. The latter is more full, and bears the title: \emph{Protocollum, h. e. Acta Colloquii inter Palatinos et Wirtebergicos Theologos de Ubiquitate sive Omnipræsentia corporis Christi. . . . A.} 1564 \emph{Maulbrunni habiti} (Heidelb. 1566). See a full résumé of the Colloquy in Ebrard:  \emph{Abendmahl}, Vol. II. pp. 666–685; Sudhoff:  \emph{Olevian und Ursin,} pp.260–290; in Hartmann:  \emph{Joh. Brenz}, pp. 253–256, and in the larger work of Hartmann and Jäger on \emph{Brenz,} 1840–42, Vol. II.} It was arranged by Duke Christopher of Würtemberg and Elector Frederick III. of the Palatinate. Olevianus, Ursinus (the authors of the Heidelberg Catechism), and Boquin defended the Reformed, the Swabian divines, Andreæ, Brenz, Schnepf, Bidenbach, and Lucas Osiander the Lutheran view. Five days were devoted to the discussion of the subject of ubiquity, and one day to the interpretation of the words, 'This \emph{is} my body.' The Lutherans regarded ubiquity as the main pillar of their view of the eucharistic presence. Andreæ proposed three points for the debate—the incarnation, the ascension, and the right hand of God.

The Lutheran reasoning was chiefly dogmatic: The incarnation is the assumption of humanity into the possession of the divine fullness with all its attributes, and the right hand of God means his almighty and omnipresent power; from these premises the absolute ubiquity of Christ's body necessarily follows.\footnote{Andreæ asserted that Christ's body, when in Mary's womb, was omnipresent as to possession (\emph{possessione}), though not as to manifestation (\emph{non patefactione}). Sudhoff, p. 279. This is the Tübingen doctrine of the \textgreek{κρύψις.} See below.}

The Reformed based their argument chiefly on those Scripture passages which imply Christ's presence in a particular \emph{place}, and his absence from other places, as when he says, 'I leave the world;' 'I go to prepare a place for you. . . . I will come again;' 'I have not yet ascended to my Father;' or when the angels say, 'He is not here,' 'Jesus is taken up from you into heaven,' etc. (John xiv. 2–4, 28; xvi. 3, 7, 16; xx. 17; Acts i. 11; iii. 21).\footnote{The same Lutherans, who so strenuously insisted on the literal interpretation of the \textgreek{ἐστί,} outdid the Reformed in the figurative interpretation of all these passages, and explained the ascension and heaven itself out of the Bible.} They urged the difference between the divine and human, and between the state of humiliation and the state of exaltation. In the appeal to the fathers and the Creed of Chalcedon they had also decidedly the advantage. Nevertheless, the Colloquy had no other effect than to confirm the two parties in their opinions.\footnote{Ebrard says (Vol. II. p. 685): '\emph{So endete das Maulbronner Gespräch mit einer vollständigen Niederlage der Lutheraner.}' Sudhoff (p. 290): '\emph{Es kann von niemandem in Abrede gestellt werden, dass die Pfälzer als Sieger aus diesem Streite hervorgegangen},' and he publishes several manuscript letters giving the impressions of the Colloquy on those present. The Swabians returned discontented, but without change of conviction. Dorner, although a Lutheran, and a Swabian by descent, gives the Reformed Christology in many respects the preference before the Lutheran, and says (Vol. II. p. 724): '\emph{Es ist unbestreitbar, dass die reformirte christologische Literatur, die um die Zeit der Concordienformel ihren Blüthepunkt erreicht, durch Geist, Scharfsinn, Gelehrsamkeit und philosophische Bildung der lutherischen Theologie vollkommen ebenbürtig, ja in manchen Beziehungen überlegen ist.}' He then gives a fine analysis of the Christology of Beza, Danæus, Sadeel, and Ursinus.}

The Consensus Dresdensis.—The Wittenberg and Leipzig Professors and other Philippists in Saxony openly rejected ubiquity in the \emph{Consensus Dresdensis} (October, 1571), which satisfied even the Elector Augustus. This document teaches that the human nature of Christ was after the resurrection glorified and transfigured, but not deified, and still remains human nature with its essential properties, flesh of our flesh; that the ascension of Christ must be understood literally, and not as a mere spectacle; that Christ's sitting at the right hand means the elevation of both natures to the priestly and kingly office; that the sacramental presence of the body of Christ must be something special and altogether distinct from omnipresence.\footnote{See Gieseler, Vol. IV. p. 466 sq.}

Absolute and Relative Ubiquity. Brenz and Chemnitz.—There was a very material difference among the advocates of ubiquity themselves as to its nature and extent, viz.: whether it were absolute, or relative, that is to say, an \emph{omni}presence in the strict sense of the term, or merely a \emph{multi}presence depending on the \emph{will} of Christ (hence also called \emph{volipræsentia}, or, by combination, \emph{multivolipræsentia}). The Swabians, under the lead of Brenz and Andreæ, held the former; the Saxon divines, under the lead of Chemnitz, the latter view.

John Brenz, or Brentius (1499–1570), the Reformer of the Duchy of Würtemberg, and after Melanchthon's death the most prominent German divine, developed, since 1559, with considerable speculative talent, a peculiar Christology.\footnote{In a series of tracts: \emph{De personali unione duarum naturarum in Christo}, 1561 (written in 1560); \emph{Sententia de libello Bullingeri}, 1561; \emph{De Divina majestate Domini nostri J. Christi ad dexteram Patris et de vera præsentia corporis et sanguinis ejus in cæna}, 1562; and \emph{Recognitio propheticæ et apost. doctrinæ de vera Majestate Dei}, 1564. In Brentii \emph{Opera}, 1590, T. VIII. pp. 831–1108. Against Brenz wrote Bullinger:  \emph{Tractatio verborum Domini Joh. XIV. 2}, Tiguri, 1561; \emph{Responsio, qua ostenditur, sententiam de cælo et dextera Dei firmiter adhuc perstare}, 1562; also Peter Martyr and Beza. The Roman Catholics sided with the Reformed against the Lutheran ubiquity. On the Christology of Brenz, comp. Dorner:  \emph{Entw. Geschichte der Christologie}, Vol. II. pp. 668 sqq.; Ebrard:  \emph{Abendmahl}, Vol. II. pp. 646 sqq. (\emph{Brenz und die Ubiquität}); and Steitz in Herzog, Vol. XVI. pp. 584 sqq.} It rests on the Chalcedonian distinction between two natures and one person, but implies at the same time, as he felt himself, a considerable departure from it, since he carried the theanthropic perfection of the exalted Saviour to the very beginning of his earthly life. He took up Luther's idea of ubiquity, and developed it to its legitimate consequences in the interest of the eucharistic presence. According to his system, the incarnation is not only a condescension of the eternal Logos to a personal union with human nature, but at the same time a deification of human nature, or an infusion of the divine substance and fullness into the humanity of Christ at the first moment of its existence. Consequently the man Jesus of Nazareth was omnipotent, omniscient, and omnipresent in the Virgin's womb, in the manger, and on the cross, as well as he is now in the state of glory.\footnote{'\emph{Majestatem divinam tempore carnis suæ in hoc seculo dissimulavit seu ea sese} (\emph{ut Paulus loquitur}) \emph{exinanivit, tamen numquam ea caruit. . . . Texit et obduxit suam majestatem forma servi.}'} The only difference is, that these divine attributes were concealed during his earthly life, and were publicly revealed to his disciples at the ascension to the right hand of God, i.e., to the omnipotent and omnipresent power of God.\footnote{'\emph{Eum tunc manifesto spectaculo voluisse testificari et declarare, se verum Deun et hominem, hoc est, una cum divinitate et humanitate sua jam inde ab initio suæ incarnationis omnia implevisse.}'} The states of humiliation and exaltation are not successive states, but co-existed during the earthly life of Christ. While Christ's humanity was poor, weak, suffering, and dying on earth, it was simultaneously almighty and omnipresent in heaven. He ascended in his humanity invisibly to heaven even at his incarnation, and remained there (John iii. 13). The visible ascension from Mount Olivet would have been impossible without the preceding invisible exaltation. Heaven is no particular place, but a state of entire freedom from space, or absolute existence in God. Space and time, with their limitations, belong only to the earthly mode of existence. Wherever the divinity is, there is also Christ's humanity,\footnote{'\emph{Ubicunque est Deitas, ibi etiam est humanitas Christi.}'} i.e., every where, not, indeed, in the way of local extension and diffusion, but in a celestial, supernatural manner, by virtue of the hypostatic union and the real communication of the properties of the divine nature to the human.

This is the most consistent, though also the most objectionable form of the ubiquity dogma. It virtually resolves the earthly life of Christ into a Gnostic delusion, or establishes a double humanity of Christ—one visible and real, and the other invisible and fantastic.\footnote{Brenz was followed by Jacob Andreæ, Schegck, and the Swabians generally, who have shown a good deal of speculative genius (down to Schelling, Hegel, and Baur), and also by a few divines of North Germany, as Andreas Musculus, John Wigand, and for a time by Heshusius, who afterwards opposed absolute ubiquity. Leonhard Hutter and Ægidius Hunnius, who were Swabians by birth, likewise took substantially the Swabian view, though more for the purpose of maintaining the authority of the 'Formula of Concord.' See Dorner, Vol. II. p. 775.}

Martin Chemnitz (1522–1586), the chief author of the 'Formula of Concord,' next to Andreæ, less original and speculative than Brenz, but superior in patristic learning and sound judgment, elaborated a Christology which mediates between Luther and Melanchthon, and taught only a relative or restricted ubiquity, i.e., a multipresence, which depends upon the will of Christ.\footnote{In his important work: \emph{De duabus naturis in Christo, de hypostatica earum unione, de communicatione idiomatum et aliis quæstionibus inde dependentibus}, Jenæ, 1570, and often reprinted. Comp. Steitz, l.c. pp. 592–597; and Dorner, Vol. II. pp. 695 sqq. Heppe says (\emph{Dogm.} Vol. II. p. 131): '\emph{Der Gegensatz der melanchthonischen und der würtembergisch-brenzischen Christologie ist sonnenklar. Jene erbaut sich auf dem Gedanken, dass Gott wirklicher Mensch geworden ist, während diese sich um den Gedanken lagert, dass ein Mensch Gott geworden ist.}'} He was followed by Selnecker, Chytræus, and most of the Saxon divines. He opposes the Swabian doctrine of a physical, natural communication and transfusion of \emph{idiomata}, and of the capacity of the finite for the infinite, except in the sense that God may dwell and reveal himself in man. He calls the absolute ubiquity a monstrosity (\emph{monstrum, portentum}), as Selnecker called it a Satanic fiction (\emph{figmentum Satanæ}). Christ is an incarnate God, not a deified man. But the Logos may temporarily communicate a divine attribute to the human nature in a supernatural manner as a \emph{donum superadditum}, without thereby setting aside the abiding limitations of humanity; just as fire may give heat and brightness to iron without turning the iron into fire. Chemnitz agrees with the Reformed, as he expressly says, in adopting the 'simple, literal, and natural signification' of the ascension of Christ as related by the Evangelists, i.e., that 'he was, by a visible motion, lifted up on high in a circumscribed form and location of the body, and departed further and further from the presence of the Apostles,' and is, consequently, in this sense withdrawn from us who are on earth, until he shall in like manner 'descend from heaven in glory in a visible and circumscribed form.' Even in glory Christ's body is finite and somewhere (\emph{alicubi}). Nevertheless, while seated at the right hand of God, he \emph{may} be present where he chooses to be, and he \emph{is} present where his Word expressly indicates such presence; as in the eucharist (according to the literal interpretation of the words of institution), or when he appeared to dying Stephen, or to Paul on the way to Damascus.\footnote{'\emph{Præsentia hæc assumtæ naturæ in Christo non est naturalis, vel essentialis, sed voluntaria et liberrima, dependens a voluntate et potentia Filii Dei, h. e. ubi se hmnana natura adesse velle certo verbo tradidit, promisit et asseveravit.}'}

Chemnitz escaped some difficulties of the Swabian theory, but by endeavoring to mediate between it and the Melanchthonian and Swiss theory, he incurred the objections to both. Christ's glorified body is indeed not confined to any locality, and may be conceived to move with lightning speed from place to place, but its \emph{simultaneous} presence in \emph{many} places, wherever the eucharist is celebrated, involves the chief difficulty of an omnipresence, and is just as inconsistent with the nature of a body.

Of subordinate interest was the incidental question, disputed mainly between Wigand and Heshusius, whether the flesh of Christ were almighty and adorable only \emph{in concreto}, or also \emph{in abstracto} (\emph{extra, personam}). Chemnitz declared this to be a mere logomachy, and advised the combatants to stop it, but in vain.

The first creed which adopted the ubiquity dogma was the Würtemberg Confession drawn up by Brenz, and adopted by a Synod at Stuttgart, Dec. 19, 1559.\footnote{\emph{Confessio et doctrina theologorum in Ducatu Wurtembergensi de vera præsentia corporis et sanguinis J. Chr. in Cæna dominica.} Here the absolute ubiquity is taught, not, indeed, in the way of a '\emph{diffusio humanæ naturæ}' or '\emph{distractio membroram Christi},' but so that 'homo  \emph{Christus quoque}  implet omnia  \emph{modo cælesti et humanæ naturæ imperscrutabili.}' See the German in Heppe: \emph{Die Entstehung and Fortbildung des Lutherthums und die kirchl. Bekenntniss-Schriften desselben}, p. 63. Melanchthon concealed his grief over this change of Brenz beneath a facetious remark to a friend on the poor Latinity of this confession ('\emph{Hechingense Latinum:}' \emph{Corp. Reform.} Vol. IX. p. 1036; comp. Gieseler, Vol. IV. p. 454; J. Hartmann: \emph{Joh. Brenz}, p. 249).}

The Formula Concordiæ on this subject is a compromise between the Swabian absolute ubiquitarianism represented by Andreæ and expressed in the \emph{Epitome}, and the Saxon hypothetical ubiquitarianism represented by Chemnitz and expressed in the \emph{Solida Declaratio}. The compromise satisfied neither party. The Helmstädt divines—Tilemann Heshusius, Daniel Hoffmann, and Basilius Sattler—who had signed the \emph{written} Formula in 1577, refused to sign the \emph{printed} copy in 1580, because it contained unauthorized concessions to the Swabian view. A colloquy was held in Quedlinburg, 1583, at which the ubiquity question was discussed for several days without result.\footnote{Heshusius wrote concerning this Colloquy: '\emph{Constanter rejicio ubiquitatem. Chemnitzius, Kirchnerus, Chytræus antea rejecerunt eam: nunc in gratiam Tubingensium cum magno ecclesiæ scandalo ejus patrocinium suscipiunt, ipsorum igitur constantia potius accusanda est.}' Comp. \emph{Acta disput.Quedlinb.}; Dorner, Vol. II. p. 773; Heppe, Vol. IV. p. 316; and G. Frank, Vol. I. p. 259 (\emph{Helmstädt und die Ubiquität}).} Chemnitz was in a difficult position, as he nearly agreed with the Helmstädtians, and conceded that certain expressions had been wrested from him, but he signed the Formula for the sake of peace, with the reservation that he understood it in the sense of a hypothetical or limited ubiquity.

The Giessen and Tübingen Controversy about the Kenosis and Krypsis.\footnote{The Saxon \emph{Solida decisio}, 1624, and an \emph{Apologia decisionis}, 1625; Feuerborn:  \emph{Sciagraphia de div. Jes. Christo juxta humanit. communicatæ majestatis usurpatione}, 1621; \textgreek{Κενωσιγραφία χριστολογική,} Marburg, 1627; Mentzer:  \emph{Juxta defensio} against the Tübingen divines, Giss. 1624; Thummius:  \emph{Majestas J. Christi } \textgreek{θεανθρώπου,} Tüb. 1621; \emph{Acta Mentzeriana}, 1625; \textgreek{Ταπεινωσιγραφία } \emph{sacra, h. e. Repetitio sanæ et orthod. doctrinæ de humiliatione Jesu Christi}, Tüb. 1623 (900 pp. 4to). On the Romish side: \emph{Bellum ubiquisticum vetus et novum}, Dilling. 1627; \emph{Alter und neuer lutherischer Katzenkrieg v. d. Ubiquität}, Ingolst. 1629; Cotta:  \emph{Historia doctrinæ de duplici statu Christi} (in his edition of Gerhard's \emph{Loci theologici}, Vol. IV. pp. 60sqq.); Walch:  \emph{Religionsstreitigkeiten}, Vol. I. p. 206; Vol. IV. p. 551; Baur:  \emph{Gesch. der L. v. d. Dreieinigkeit}, Vol. III. p. 450; Thomasius:  \emph{Christi Person und Werk, }Vol. II. pp. 391–450; Dorner, Vol. II. pp. 788–809; G. Frank, Vol. I. p. 336.}—The ubiquity question was revived under a new shape, on the common basis of the 'Formula of Concord' and the dogma of the \emph{communicatio idiomatum}, in the controversy between the \emph{Kenoticism}, of the theologians of Giessen, which followed in the track of Chemnitz, and the \emph{Krypticism} of the theologians of Tübingen, which was based upon the theory of Brenz and Andreæ. The controversy forms the last phase in the development of the orthodox Lutheran Christology; it continued from 1616–1625, and was lost in the Thirty-Years' War.

Both parties agreed that the human nature of Christ from the moment of the incarnation, even in the mother's womb and on the cross, was in \emph{full possession} (\textgreek{κτῆσις}) of the divine attributes of omnipresence, omnipotence, omniscience, etc.; but they differed as to their \emph{use} (\textgreek{χρῆσις}). The Giessen divines—Balthazar Mentzer (d. 1627), his son-in-law, Justus Feuerborn (d. 1656), and John Winckelmann—taught a real self-renunciation (\textgreek{κένωσις, } \emph{evacuatio, exinanitio}),\footnote{\par Hence they were called \emph{Kenotiker, Kenoticists.}} i.e., that Christ voluntarily laid aside the actual use of the divine attributes and functions, except in the working of miracles; while the Tübingen divines—Lucas Osiander II. (d. 1638), Theodor Thumm, or Thummius (d. 1630), and Melchior Nicolai (d. 1659)—taught that he made a \emph{secret use} of them (\textgreek{κρύψις, } \emph{occulta usurpatio}).\footnote{\par Hence their name, \emph{Kryptiker, Krypticists.}}

The Giessen divines, wishing chiefly to avoid the reproach of a \emph{portentosa ubiquitas}, represented the omnipresence of Christ's humanity, not as an all-pervading existence,\footnote{\emph{Indistantia, nuda adessentia ad creaturas, præsentia simplex.}} but as an all-controlling power, or as an element of omnipotence.\footnote{\emph{Actio, operatio, præsentia modificata.} This amounts to pretty much the same thing with the \emph{omimpræsentia energetica} of the Calvinists.} The Tübingen school taught, in consequence of the \emph{unio hypostatica}, an absolute omnipresence of Christ's humanity, as a quiescent quality, which consists in filling all the spaces of the universe, even from the conception to the death on the cross.\footnote{The same applies to omnipotence. The Tübingen divines gave an affirmative answer to the question, '\emph{An }homo  \emph{Christus in Deum assumptus in statu exinanitionis tamquam rex præsens cuncta, licet latenter, gubernarit?}' They made, however, an apparent concession to their opponents by assuming a brief suspension of the \emph{use} of the divine majesty during the agony in Gethsemane and the crucifixion, in order that Christ might really suffer as high-priest. See Dorner, Vol. II. p. 799.}

A theological commission at Dresden, with Hoe von Hoenegg at the head, decided substantially in favor of the Giessen theory (1525), and against the Tübingen doceticism, without, however, advancing the solution of the problem or feeling its real difficulty.

The Giessen theory is more consistent with the realness of Christ's human life, but less consistent with itself, since it admits an \emph{occasional} interruption of it by the use of the inherent powers of the divinity; the Tübingen theory, on the other hand, virtually destroys the distinction between the state of humiliation and the state of exaltation, and resolves the life of Christ into a magical illusion.

The modern Tübingen school of Baur and Strauss forms a strange parallel and contrast to that of the seventeenth century: it starts from the same principle that 'the finite is capable of the infinite,' but extends it pantheistically to humanity at large, and denies its applicability to Christ, on the ground that the divine fullness can not be emptied into a single individual.\footnote{'In an individual,' says Strauss, in the dogmatic conclusion of his first \emph{Leben Jesu} (Vol. II. p. 710\}. 'in one God-man, the properties and functions which the Church doctrine ascribes to Christ contradict themselves; in the idea of the race they agree. \emph{Humanity} is the union of the two natures—the incarnate God—the infinite externalizing itself in the finite, and the finite spirit remembering its infinitude.'} Therefore, while the old Tübingen school in effect, though not in intention, destroys the real humanity of Christ, the modern Tübingen school consistently denies his divinity, and resolves all the supernatural and miraculous elements of the gospel history into a mythic poem or fiction.

In the modern revival of orthodox Lutheranism, the ubiquity of the body of Christ is either avoided, or advocated only in the hypothetical form, and mostly with a leaning towards a more literal acceptation of the \textgreek{κένωσις} (Phil. ii. 7) than the Giessen divines contended for.\footnote{So Thomasius, Liebner, Gess. But the absolute ubiquity also has found an advocate in Philippi (\emph{Kirchl. Glaubenslehre}, Vol. IV. I. pp. 394). Dr. Stahl, the able theological lawyer, in his \emph{Die lutherische Kirche und die Union} (Berlin, 1859, pp.185 sqq.), admits that the ubiquity question has no religious interest except as a speculative basis for the possibility of the eucharistic presence, and approaches Ebrard's view of an 'extra-spacial, central communication of the virtue' of Christ's body to the believer. Dr. Krauth defends Chemnitz's view, and what he would rather style 'the \emph{personal} omnipresence of the human nature of Christ' (l.c. p. 496). But the human nature of Christ is impersonal, and simply taken up into union with the pre-existent personality of the Divine Logos.}

VIII. THE HADES CONTROVERSY.\footnote{Æpinus: \emph{Comment, in Psa. xvi.} Frcf. 1544, and \emph{Enarratio Psalmi lxviii.}, with an appendix \emph{de descensu Christi ad inferna}, Frcf. 1553. A. Grevius:  \emph{Memoria J. Æpini instaurata}, Hamb. 1736; Dietelmaier:  \emph{Historia dogmatis de descensu Christi}, Norimb. 1741, Alt. 1762; Planck, Vol. V. I. pp. 251–264; König:  \emph{Die Lehre von Christi Höllenfahrt}, pp. 152 sqq.; Güder:  \emph{Die Lehre der Erscheinung Christi unter den Todten}, Bern, 1853, pp. 222 sqq.; G. Frank, Vol. I. p. 160 sq.; Fr. H. R. Frank, Vol. III. p. 397 sqq.}

This controversy, which is discussed in the ninth article of the 'Formula of Concord,' referred to the time, manner, extent, and aim of Christ's mysterious descent into the world of departed spirits. It implied the questions whether the descent took place before or after the death on the cross; whether it were confined to the divine nature, or to the soul, or extended to the body; whether it belonged to the state of humiliation, or to the state of exaltation; whether it were a continuation of suffering and a tasting of the second death, or a triumph over hell. The answer to these questions depended in part on the different views of the communication of idiomata and the ubiquity of the body, as also on Hades, or Sheol, itself, which some identified with hell proper (Gehenna), while others more correctly understood it in a wider sense of the whole realm of the dead. Luther himself had at different times very different opinions of the descent, but regarded it chiefly as a victory over the kingdom of Satan.

John Æpinus,\footnote{A Hellenized form (\textgreek{Αἰπεινός, } \emph{high, lofty}) for his German name Höck, or Hoch. He was born, 1499, at Ziegesar, Brandenburg; studied at Wittenberg, became pastor at St. Peter's, Hamburg, 1529, Superintendent in 1532, introduced the Reformation into that city, signed the Articles of Smalcald, 1537, stood in high esteem, and died 1553. He was a colleague of Westphal, and opposed with Flacius the Leipzig Interim.} a Lutheran minister in Hamburg, started the controversy. He taught, first in 1544 and afterwards more fully, that Christ descended with his spirit into the region of the lost, in order to suffer the pains of hell for men, and thus to complete his humiliation or the work of redemption. So he explained Psalm xvi. 10 (comp. Acts ii. 27, 31). Luther himself had at one time (1524) given a similar exposition of this passage. Flacius sided with Æpinus. But this theory was more Reformed than Lutheran, and was opposed by his colleagues, who carried the dispute into the pulpit and excited the people. Matsberger in Augsburg represented the descent, according to the usual view, as a local change, but had to suffer three years' imprisonment for it. Brenz condemned such locomotion as inconsistent with the dignity and ubiquity of Christ, and denied the locality of hell as well as of heaven. This accords with his view of the ascension. Melanchthon, being appealed to by the magistrate of Hamburg, answered with caution, and warned against preaching on subjects not clearly revealed. He referred to a sermon of Luther, preached at Torgau, 1533, in which he graphically describes the descent as a triumphant march of Christ through the dismayed infernal hosts, so that no believer need hereafter be afraid of the devil and damnation. Melanchthon thought this view was more probable than that of Æpinus; at all events, Christ manifested himself as a conqueror in hell, destroyed the power of the devil, raised many dead to life (Matt. xxvii. 53), and proclaimed to them the true doctrine of the Messiah; to ask more is unnecessary. He advised the magistrate to exclude the controversy from the pulpit.\footnote{Sept. 1550, \emph{Corp. Reform.} Vol. VII. p. 665. Comp. Schmidt, \emph{Melanchthon}, p. 554 sq. In his \emph{Loci}, Melanchthon passes by the \emph{descensus} as unessential. In a letter to Spalatin, March 20,1531 (\emph{Corp. Reform.} Vol. II. p. 490), he expresses his inability to explain the dark passage, 1 Pet. iii. 19, 20. He was pleased with Luther's sermon at Torgau, but added, in a private letter to Anton Musa (March 12, 1543, \emph{Corp. Reform.} Vol. V. p. 58), that Christ probably preached the gospel to the heathen in the spirit world, and converted such men as Scipio and Fabius. (Zwingli likewise believed in the salvation of the nobler heathen.) He wrote to Æpinus, April 20, 1546 (\emph{Corp. Reform.} Vol. VI. p. 116), to preach the necessary doctrines of faith, repentance, prayer, good works, rather than speculations on things which even the most learned did not know.} Several of the most violent opponents of Æpinus were deposed and expelled. The dispute was lost in more serious controversies. It was almost confined to Hamburg.

The Formula of Concord sanctioned substantially the view of Luther and Melanchthon, without entering into the minor questions.

IX. THE ADIAPHORISTIC (OR INTERIMISTIC) CONTROVERSY (1548–1555).\footnote{Comp. Flacius:  \emph{Von wahren und falschen Mitteldingen}, etc.; \emph{Entschuldigung geschrieben an die Universität zu Wittenberg der Mittelding halben}, etc.; \emph{Wider ein recht heidnisch, ja Epicurisch Buch der Adiaphoristen, darin das Leipzische Interim vertheidigt wird}, etc.; and other pamphlets, printed at Magdeburg (as the '\emph{Kanzlei Gottes}'), 1549; Wigand:  \emph{De neutralibus et mediis, }Frcf. 1560; Schlüsselburg:  \emph{Cat. Hæret.} Lib. XIII. (\emph{de Adiaphoristis et Interimistis}); Biek:  \emph{Das dreifache Interim}, Leipz. 1725, Planck, Vol. IV. pp. 85–248; H. Rossel:  \emph{Mel. und das Interim} (at the close of Twesten's monograph on Flacius, Berlin, 1844); Ranke:  \emph{Deutsche Gesch.}, etc. Vol. V.; Gieseler, Vol. IV. p. 435; Herzog:  \emph{Encykl.} Vol. I. p. 124; Vol. VIII. p. 288; Schmidt:  \emph{Mel.} pp. 491, 495, 524; G. Frank, Vol. I. pp. 113, 116; Fr. H. R. Frank, Vol. IV. pp. 1–120; Dorner, p. 331.}

This controversy is the subject of the tenth article of the 'Formula of Concord,' but was the first in the order of time among the disputes which occasioned this symbol. It arose, soon after Luther's death, out of the unfortunate Smalcald war, which resulted in the defeat of the Lutheran states, and brought them for a time under the ecclesiastical control of the Emperor Charles V. and his Romish advisers.

Ecclesiastical rites and ceremonies, which are neither commanded nor forbidden in the Word of God, are in themselves indifferent (\textgreek{ἀδιάφορα, } \emph{media, res mediæ, } \emph{Mitteldinge}), but the observance or non-observance of them may, under testing circumstances, become a matter of principle and of conscience. The Augsburg Confession and Apology (Art. VII.) declare that agreement in doctrine and the administration of the sacraments is sufficient for the unity of the Church, and may co-exist with diversity in usages and rites of human origin. Luther himself desired to retain many forms of the Catholic worship which he considered innocent and beautiful, provided only that no merit be attached to them and no burden be imposed upon the conscience.\footnote{See his humorous letter to Buchholzer in Berlin, Dec. 4, 1539 (\emph{Briefe}, Vol. V. p. 235), which might have considerably embarrassed the anti-Adiaphorists had they known it. He advises Elector Joachim II. that in introducing the Reformation he may, if he desired it, put on one or three priestly garments, like Aaron; may hold one or even seven processions, like Joshua before Jericho; and may dance before it, as David danced before the ark, provided only such things were not made necessary for salvation.} But there is a great difference between retaining old forms and restoring them after they have been abolished, as also between a voluntary and a compulsory observance. When circumcision was yet lawful and practiced by Jewish Christians, Paul resisted it, and saved the principle of Christian liberty against the Judaizing error which made circumcision a condition of salvation. Some of the Romish ceremonies, moreover, especially those connected with the canon of the mass, involve doctrine, and affect the whole idea of Christian worship.

When the Emperor, with the aid of the treasonable Elector Maurice of Saxony, had broken up the Lutheran League of Smalcald, he required the Protestants to submit to a doctrinal and ceremonial compromise till the final settlement of the religious controversy by an œcumenical Council.

The first compromise was the so-called \emph{Augsburg Interim}, enacted by the Diet of Augsburg (May, 1548) for the whole empire. It was essentially Romish, and yielded to the Protestants only the marriage of priests and the cup of the laity. It was rigidly executed in the Southern and prevailingly Roman Catholic states, where about four hundred Lutheran preachers were expelled or dismissed for non-conformity.

The second compromise, called the \emph{ Leipzig Interim}, was enacted by the Elector Maurice (December, 1548), with the aid of Melanchthon and other leading Lutheran divines, for his Protestant dominion, where the Augsburg Interim could not be carried out. It was much milder, saved the evangelical creed in its essential features—as justification by the sole merits of Christ through a living faith—but required conformity to the Romish ritual, including confirmation, episcopal ordination, extreme unction, and even the greater part of the canon of the mass, and such ceremonies as fasts, processions, and the use of images in churches.\footnote{See the text of the two Interims in Gieseler,  Vol. IV. pp. 193–196 and 201–203; the Interim Lipsiense, also, in \emph{Corp. Reform.} Vol. VII. The term gave rise to sarcastic conundrums, as \emph{Interimo, interitus, Hinterim, der Schalk ist hinter ihm} (the villain is behind it). On the political aspects of the Interim, see the fifth volume of Ranke.}

The Protestants were forced to the alternative of either submitting to one of these temporary compromises, or risking the fate of martyrs.

Melanchthon, in the desire to protect churches from plunder and ministers from exile, and in the hope of saving the cause of the Reformation for better times, yet not without blamable weakness, gave his sanction to the Leipzig Interim, and undertook to act as a mediator between the Emperor, or his Protestant ally Maurice, and the Protestant conscience.\footnote{To the Augsburg Interim he was decidedly opposed, and he had also sundry objections to the ceremonial part of the Leipzig Interim. He is only responsible for its doctrinal part. See his letters from this period in \emph{Corp. Reform.} Vols. VI. and VII., and Schmidt's \emph{Mel.} pp. 507 and 524.} It was the greatest mistake in his life, yet not without plausible excuses and incidental advantages. He advocated immovable steadfastness in doctrine, but submission in every thing else for the sake of peace. He had the satisfaction that the University of Wittenberg, after temporary suspension, was restored, and soon frequented again by two thousand students; that no serious attempt was made to introduce the Interim there, and that matters remained pretty much as before. But outside of Wittenberg and Saxony his conduct appeared treasonable to the cause of the Reformation, and acted as an encouragement to an unscrupulous and uncompromising enemy. Hence the venerable man was fiercely assailed from every quarter by friend and foe. He afterwards frankly and honorably confessed that he had gone too far in this matter, and ought to have kept aloof from the insidious counsels of politicians.\footnote{In a letter to his enemy, M. Flacius, dated Sept. 5, 1556, he was not ashamed to confess, after some slight reproaches, '\emph{Vincite! Cedo; nihil pugno de ritibus illis, et maxime opto, ut dulcis sit ecclesiarum concordia. Fateor etiam hac in re a me peccatum esse, et a Deo veniam peto, quod non procul fugi insidiosas illas deliberationes. Sed illa quæ mihi falsa a te et a Gallo objiciuntur, refutabo.}' \emph{Corp. Reform.} Vol. VIII. p. 841 sq. And to the Saxon pastors he wrote, Jan. 17, 1557 (Vol. IX. p. 61): '\emph{Pertractus sum ad aularum deliberationes insidiosas. Quare sicubi vel lapsus sum, vel languidius aliquid egi, peto a Deo et ab Ecclesia veniam, et judiciis Ecclesiæ obtemperabo.}'} He fully recovered his manhood in the noble Saxon Confession which he prepared in 1551 for the Council of Trent, and which is not merely a repetition of the Augsburg Confession, but also a refutation of the theology, worship, and government of the papal Church.

Flacius chose the second alternative. Escaping from Wittenberg to the free city of Magdeburg, he opened from this stronghold of rigid Lutheranism, with other 'exiles of Christ,' a fierce and effective war against Melanchthon and the 'dangerous rabble of the Adiaphorists.' He charged his teacher and benefactor with superfluous mildness, weakness, want of faith, treason to truth; and characterized the Leipzig Interim as an undisguised 'union of Christ and Belial, of light and darkness, of sheep and wolf, of Christ and Antichrist,' aiming at the 'reinstatement of popery and Antichrist in the temple of God.'\footnote{Thus he concisely states the case on the long title-page of his \emph{Apology}, or \emph{Entschuldigung}, etc., addressed to the University of Wittenberg, with a letter to Melanchthon, Magdeburg, 1549. The concluding words of the title state the aim of the Interim thus: '\emph{Das Ende ist die Einsetzung des Papstthums und Einstellung des Antichrists in den Tempel Christi, Stärkung der Gottlosen, dass sie über der Kirche Christi stolziren, Betrübung der Gottfürchtigen, item Schwächung, Einführung in Zweifel, Trennung und unzählige Aergerniss.}' He relates of Melanchthon that he derived from an eclipse of the moon in 1548 the vain hope of the near death of the Emperor, which would end these troubles. He also published several confidential letters of Luther to Melanchthon, written during the Diet at Augsburg, 1530, upbraiding him for his philosophy and timidity.} His chief text was  1 Cor. x. 20–23. He had upon the whole the best of the argument, although in form he violated all the laws of courtesy and charity, and continued, even long afterwards, to persecute Melanchthon as an abettor of Antichrist.

In a milder tone the best friends of Melanchthon remonstrated with him. Brenz preferred exile and misery to the \emph{Interim}, which he called \emph{interitus.} Bucer of Strasburg did the same, and accepted a call to England. Calvin on this question sided with the anti-Adiaphorists, and wrote a letter to Melanchthon (June 18, 1550), which is a model of brotherly frankness and reproof. 'My present grief,' he says in substance, 'renders me almost speechless. . . . In openly admonishing you, I am discharging the duty of a true friend; and if I employ a little more severity than usual, do not think it is owing to any diminution of my old affection and esteem for you. . . . I know you love nothing better than open candor. I am truly anxious to approve all your actions, both to myself and to others. But at present I accuse you before yourself, that I may not be forced to join those who condemn you in your absence. This is the sum of your defense: That provided purity of doctrine be retained, externals should not be pertinaciously contended for. . . . But you extend the adiaphora too far. . . . Some of them contradict the Word of God. . . . When we are in the thick of the fight, we must fight all the more manfully; the hesitation of the general brings more disgrace than the flight of a whole herd of common soldiers. All will blame you if you do not set the example of unflinching steadfastness. . . . I had rather die with you a hundred times than see you survive the doctrines surrendered by you. I have no fear for the truth of God, nor do I distrust your steadfastness. . . Pardon me, dear Philip, for loading your breast with these groans. May the Lord continue to guide you by his Spirit and sustain you by his might.'\footnote{\emph{Opera}, Vol. IX. p. 51, and \emph{Letters of Calvin}, by J. Bonnet, English translation, Vol. II. p. 257. A letter of similar spirit and import to Melanchthon, by his friend Anton Corvinus (Räbener), a distinguished reformer in Hesse and Göttingen, who suffered imprisonment for his opposition to the Interim, was recently discovered in the Royal Library at Hanover by Iwan Franz, and published in Kahnis, \emph{Zeitschrift fur die hist. Theol.} 1874, pp. 105 sqq., from which I quote the following passages: '\emph{O Philippe, o inquam Philippe noster, redi per immortalem Christum ad pristinum candorem, ad pristinam tuam sinceritatem! non languefacito ista tua formidine, pusillanimitate et inepta moderatione nostrorum animos tantopere! Non aperito hac ratione ad Papatus recurrentem impietatem ac Idolomanias fenestram ac januam! Non sis tantorum in Ecclesia offendiculorum autor! Ne sinas tua tam egregia scripta, dicta, facta, quibus mirifice de Ecclesia hactenus meritus es, isto condonationis, moderationis, novationis nævo ad eum modum deformari! Cogita, quantum animi ista nostra carnis ac rationis consilia et adversariis addant et nostris adimant.! Perpende, quam placari etiam istis condonationibus adversarii nostri non queant, qui totius Papatus doctrinam et omnes ex cequo impios cultus reposcunt et ex nostra levitate spem concipiunt se hac in re facile voti compotes futuros. Detestatur Dominus apud Jeremiam eos, qui manus pessimormn confortant, ut non convertatur unusquisque a malitia sua. Cur igitur in tam ardua causa non tales nos gerimus ut hujusmodi detestatio competere in nos haud possit? qua perversitate arundo huc illuc ventis agitata dici quam Johannis constantiam imitari malumus! . . . Proinde Te, o noster Philippe, iterum atque iterum per ilium ipsum Christum redemptorem nostrum et brevi futurum judicem rogamus, ut professionis tuæ memor talem te cum reliquis Vitebergensibus jam geras, qualem Te ab initio hujus causæ ad Electoris captivitatem usque gessisti, hoc est, ut ea sentias, dicas, scribas, agas, quæ Philippum, Christianum Doctorem decent, non aulicum Philosophum.}'}

The defeat of the Emperor by Elector Maurice, who now turned against him, as he had turned before against his fellow-Protestants, and the consequent Peace of Augsburg, 1555, made an end to the Interim troubles, and secured freedom to the Lutheran Churches. But among theologians the controversy continued till the death of Melanchthon.

The conduct of Melanchthon weakened his authority and influence, which had been rising higher and higher before and after Luther's death, especially in the University of Wittenberg. Before this unfortunate controversy he was universally regarded as the theological head of the evangelical Church in Germany, but now a large number of Lutherans began to look upon him with distrust.

X. THE STRASBURG CONTROVERSY ON PREDESTINATION BETWEEN ZANCHI AND MARBACH (1561–1563).\footnote{Planck, Vol. VI. pp.809 sqq.; Röhrich:  \emph{Geschichte der Reform. im Elsass, bes. in Strassburg}, 3 Theile, Strasburg, 1830–1882; Schweizer:  \emph{Centraldogmen der Reform. Kirche}, Vol. I. pp. 418–470 (a very full and able account); Heppe:  \emph{Dogmatik des D. Protest}. Vol. II. pp. 44–47; G. Frank, Vol. I. pp. 178–184; Fr. H. E. Frank, Vol. IV. pp. 121–344.}

This is the last specific doctrine discussed in the Formula of Concord (Art. XI.). The German and Swiss Reformers alike renewed, as an impregnable fortress in their war against the Pelagian corruptions of Rome, the Augustinian system, with its two closely connected doctrines of the absolute spiritual slavery or inability of the unregenerate will of man, and the absolute predestination of God; though with the characteristic difference that Luther and Melanchthon emphasized the \emph{servum arbitrium}, Zwingli the \emph{providentia}, Calvin the \emph{prædestinatio}. In other words, the German Reformers started from the anthropological premise, and inferred from it the theological conclusion; while Calvin made the absolute sovereignty of God the cornerstone of his system. Luther firmly adhered to the \emph{servum arbitrium}, but was more cautious, in his later years, on the mystery of the \emph{prædestinatio.}\footnote{The Philippist Lasius first asserted (1568) that Luther had recalled his book \emph{De servo arbitrio} (1525), but this was indignantly characterized by Flacius and Westphal as a wretched lie and an insult to the evangelical church. The fact is that Luther emphatically reaffirmed this book, in a letter to Capito, 1537, as one of his very best ('\emph{nullum enim agnosco meum justum librum nisi forte De servo arbitrio, et Catechismum}'). And, indeed, it is one of his most powerful works. Luthardt (\emph{Die Lehre vom freien Willen}, Leipz. 1863, p. 122) calls it '\emph{eine mächtige Schrift, stoltz, wahrheitsgewiss, kühn in Gedanken und Wort, voll heiligen Eifers, gewaltigen Ernstes, aus innerster Seele herausgeschrieben. . . . Kaum irgendwo sonst ergiesst sich gleich mächtig und reich der Strom seines Geistes.}' Only in regard to predestination Luther may be said to have moderated his view somewhat, although he never recalled it, that is, he still taught in his later writings (in his \emph{Com. on Genesis},  Ch. VI. 6,  18; Ch. XXVI.) the distinction and antagonism between the revealed will of God, which sincerely calls \emph{all} to repentance and salvation, and the inscrutable secret will which saves only a \emph{part} of the race; but he laid the main stress practically on the former and the means of grace, and thus prepared the way for the 11th Article of the Formula of Concord. '\emph{Scripsi},' he wrote in 1536, '\emph{esse omnia absoluta et necessaria, sed simul addidi, quod adspiciendus sit Deus revelatus}' (\emph{Opera exeg.} Vol. VI. p. 300). Luthardt (l.c. p. 146) correctly says (in opposition both to Lütkens and Philippi) that Luther never recalled, but retained his earlier views on predestination and the necessity of all that happens, and only guarded them against abuse. The result of Köstlin's investigation is this, that Luther never attempted a solution of the contradiction between the secret and the revealed will of God. '\emph{Das eben ist seine Lehre, dass unser Erkennen nicht so weit reicht, und dass wir uns auch das Unbegreifliche und Unverständliche gefallen lassen müssen. . . . Er selbst spricht aus, dass ein Widerspruch für uns stehen bleibe, den wir nicht lösen können noch sollen.}' \emph{Luther's Theologie}, Vol. II. p. 328.} Melanchthon gave up both for his synergism and the universality of grace, though he continued in friendly correspondence with Calvin, who on his part put the mildest construction on this departure. The rigid Lutherans all retained Luther's view of total depravity in opposition to synergism, and some of them (namely, Amsdorf, Flacius, Brenz, Wigand, and, for a time, Heshusius) were also strict predestinarians.\footnote{See the proof passages in Frank's \emph{Theol. der Concord. formel}, Vol. IV. pp. 254–261; Luthardt, pp. 240–244; Planck, Vol. IV. pp. 691–712; and Schweizer, l.c.} But the prevailing Lutheran sentiment became gradually averse to a particular predestination, all the more since it was a prominent doctrine of the hated Calvinists. The Formula of Concord sanctioned a compromise between Augustinianism and universalism, or between the original Luther and the later Melanchthon, by teaching both the absolute inability of man and the universality of divine grace, without an attempt to solve these contradictory positions. In regard to the slavery of the human will, the Formula of Concord, following Luther, went even further than Calvin, and compared the natural man with a dead statue, or clod, and stone; while Calvin always (so far agreeing with the later Melanchthon) insisted on the spontaneity and responsibility of the will in sinning, and in accepting or rejecting the grace of God.

The discussion of this subject was opened by the fierce polemic Tilemann Heshusius, who, in his defense of the corporeal presence against the Sacramentarians (Jena, 1560), first attacked also Calvin's doctrine of predestination, as Stoic and fatalistic, although a year afterwards, in opposition to synergism, he returned to his former view of an absolute and particular predestination. Beza answered his attack with superior ability.\footnote{See Schweizer, l.c. pp. 402 sqq. Heshusius and Westphal invented the name \emph{Calvinists}, which henceforth was used by Lutherans for the Reformed, as the term \emph{Zwinglians} had been before. The term \emph{sacramentarians} was applied to both without distinction.}

Of more importance was the controversy between Marbach (a friend of Heshusius) and Zanchi within the Lutheran denomination itself. It decided its position on the question of predestination and perseverance.

The Church of Strasburg had received from its reformer, Martin Bucer (who on account of the Interim followed a call to the University of Cambridge, 1549, and died there, 1551), a unionistic type, and acted as mediator between the Swiss and German churches. The Reformed Tetrapolitan Confession, the Lutheran Augsburg Confession, and the Wittenberg Concordia (a compromise between the Lutheran and Zwinglian views on the eucharist), were held in great esteem. Calvin and Peter Martyr, who preached and taught there, made a deep impression. The celebrated historian Sleidanus, and the learned founder and rector of the academy, John Sturm, labored in the same spirit.

Jerome Zanchi (Zanchius, 1516–1590), a converted Italian, and pupil of Peter Martyr, became his successor as Professor of Theology at Strasburg in 1553. He was one of the most learned Calvinistic divines of the age, and labored for some time with great acceptance. He taught that in the eucharist Christ's true body broken for us, and his blood shed for us, are received in the sacrament, but not with the mouth and teeth, but by faith, and consequently only by believers. This was approved by his superiors, since the communion was not a \emph{cibus ventris sed mentis}, and the same view had been taught by Bucer, Capito, Hedio, Zell, and Martyr. He opposed ubiquity, and the use of images in churches. He taught unconditional predestination, and its consequence, the perseverance of saints, in full harmony, as he believed, with Augustine, Luther, and Bucer. He reduced his ideas to four sentences: 1. The elect receive from God the gift of true saving faith only once; 2. Faith once received can never be totally and finally lost, partly on account of God's promise, partly on account of Christ's intercession; 3. In every elect believer there are two men, the external and the internal—if he sin, he sins according to the external, but against the internal man, consequently he sins not with the \emph{whole} heart and will; 4. When Peter denied Christ, the confession of Christ died in his mouth, but not his faith in his heart.

Several years before Zanchi's call to Strasburg, a Lutheran counter-current had been set in motion, which ultimately prevailed. It was controlled by John Marbach (1521–1581), a little man with a large beard, incessant activity, intolerant and domineering spirit, who had been called from Jena to the pulpit of Strasburg (1545). Inferior in learning,\footnote{Melanchthon called him \emph{mediocriter doctus}, but his own estimate was much higher, and in his inaugural he spoke with such arrogance that Bucer feared he would prove a great misfortune for the Church at Strasburg. See Röhrich and Schweizer, p. 420.} he was superior to Zanchi in executive ability and popular eloquence. He delighted to be called Superintendent, and used his authority to the best advantage. He abolished Bucer's Catechism and introduced Luther's, taught the ubiquity of Christ's body, undermined the authority of the Tetrapolitan Confession, crippled the church of French refugees, to which Calvin had once ministered, weakened discipline, introduced pictures into churches, including those of Luther, and began to republish at Strasburg the fierce polemical book of Heshusius on the eucharist. This brought on the controversy.

Zanchi persuaded the magistrate to suppress the publication of this book, because of its gross abuse of Melanchthon and a noble German Prince, the Elector Frederick III. of the Palatinate, and because it denounced all who differed from his views of the corporeal presence as heretics. From this time Marbach refused to greet Zanchi on the street, and gathered from the notes of his students material for accusation that he taught doctrines contrary to the Augsburg Confession. He objected, however, not so much to predestination itself as to Zanchi's method of teaching it \emph{a priori} rather than \emph{a posteriori.}

The controversy lasted over two years. Zanchi visited and consulted foreign churches and universities. The answers differed not so much on predestination as on perseverance.\footnote{Zanchii \emph{Opera}, Pt. VII. pp. 65 sqq., and Pt. VIII. pp. 114 sqq.; Schweizer, pp. 448–470.}

The theologians of Marburg (Hyperius, Lonicer, Garnier, Orth, Roding, Pincier, and Pistorius), Zurich (Bullinger, Martyr, Gualter, Lavater, Simler, Haller, Zwingli Jr.), and Heidelberg (Boquinus, Tremellius, Olevianus, and Diller) decided in favor of the theses of Zanchi. The ministers of Basel counseled peace and compromise; the divines of Tübingen approved of the doctrine of predestination, but dissented from the theses on perseverance; even Brenz thought the matter might be amicably settled. The divines of Saxony decided according to their different attitudes towards Melanchthon: the Melanchthonians liked Zanchi's doctrine of the eucharist, but disliked his view of predestination; the anti-Melanchthonians hated the former, but were favorable to the latter, because it was so strongly taught by Luther himself (\emph{De servo arbitrio}).

At last the 'Strasburg Formula of Concord' was adopted (1563), which prescribed the Wittenberg Concordia of 1536 as the rule of doctrine on the Lord's Supper, and asserted the possibility of the loss of faith, yet without denying predestination.\footnote{Printed in the Strasburger \emph{Kirchenordnung} of 1598. and in Löscher's \emph{Historia motuum}, Vol. II. p. 229 sq. See Schweizer, pp. 440 sqq.} Calvin judged that it only threw a veil over the truth. Predestination was with Calvin and Luther an independent and central dogma; the later Lutherans assigned it a subordinate and subsidiary position, and denied its logical consequence, the perseverance of saints. This was also the position of Marbach.

Zanchi subscribed the Strasburg Formula with a restriction, but for the sake of peace he soon followed a call to a Reformed Italian church at Chiavenna, and, being driven away by a pestilence to a mountain, he wrote a full account of the Strasburg troubles.\footnote{It is addressed to Philip of Hesse (Oct. 1, 1565), and given by Schweizer, pp. 425–436. Zanchi accepted afterwards a call to a professorship at the Reformed University of Heidelberg, where he died, 1590. He received also calls to England, Lausanne, Geneva, Zurich, and Leyden, and was justly esteemed for his learning and character. A complete edition of his works appeared at Geneva in eight parts, in 3 vols. folio.} He was supported in his position by the worthy Sturm and several professors, but had the disadvantage of being a foreigner unacquainted with the German tongue. The pastors, backed by the people, triumphed over the professors. What Marbach had begun, his pupil Pappus completed. Strasburg was thoroughly Lutheranized, the Tetrapolitan Confession formally abolished as 'Zwinglian,' and the Formula Concordiæ introduced (1597).\footnote{Comp. Heppe, \emph{Gesch. des D. Protest.} Vol. IV. pp.312–315.}

Yet, after all, the spirit of Bucer never died out. From Strasburg proceeded Spener, with his blessed revival of practical piety and a better appreciation of the Reformed Confession;\footnote{Spener was born at Rappoltsweiler, in Upper Alsace, but his parents were from Strasburg, and he was educated there, and called himself a Strasburger. Kliefoth (as quoted by Heppe, Vol. IV. p. 399), from his own rigid Lutheran stand-point, says, not without good reason: '\emph{Mit Spener beginnt jener grosse Eroberungszug der reformirten Kirche gegen die lutherische, der seitdem verschiedene Namen, erst Frömmigkeit, dann Toleranz, dann Union, dann Conföderation auf sein Panier geschrieben hat.}'} and from the theological faculty of Strasburg hail more recently the appreciating biographies of Beza, Bucer, and Capito (by Baum), and Melanchthon (by Carl Schmidt), and the best edition of the works of Calvin (by Baum, Cunitz, and Reuss). Thus history slowly but surely rectifies its own mistakes.

THE PREPARATION OF THE FORMULA OF CONCORD.\footnote{For the fullest account, see the sixth volume of Planck's, and the third volume of Heppe's history.}

These controversies turned the Lutheran churches in Germany into a camp of civil war, exposed them to the ridicule and obloquy of the Papists, and threatened to end in utter confusion and dissolution. The danger was increased by the endless territorial divisions of Germany, where every Prince and magistrate acted a little pope, and 'every fox looked to his own pelt.'\footnote{As Brenz says: '\emph{Es luge ein jeglicher Fuchs seines Balges.}'}

The best men in the Lutheran communion deeply deplored this state of things, and labored for peace and harmony. Augustus, Elector of Saxony (1533—1588), a pious and orthodox, though despotic Prince, controlled the political part, and paid the heavy expenses of the movement.\footnote{80,000 gulden. Augustus was a zealous Lutheran without knowing the difference between Lutheranism and Philippism, and supported or punished the champions of both parties as he happened to be led or misled by his courtiers and the theologians.} Jacob Andreæ, Professor of Theology and Chancellor of the University at Tübingen (1528–1590), a pupil and friend of Brentius, a man of rare energy, learning, eloquence, and diplomatic skill, managed the theological negotiations, made no less than one hundred and twenty-six journeys, and sacrificed the comforts of home and family (he had twelve children) to the pacification of the Lutheran Church.\footnote{On this remarkable man, see Planck, Vol. VI. pp. 372 sqq.; Heppe, Vol. IV. pp. 376 sqq.; G. Frank, Vol. I. p. 219; Hartmann in Herzog, Vol. I. p. 312; Johannsen, \emph{Jacob Andreæ's Concordistische Thätigkeit}, in Niedner's \emph{Zeitschrift für hist. Theol.} 1853, No. 3. Andreæ has often been too unfavorably judged. His contemporary opponents called him 'Schmidlin' (with reference to his father's trade), 'Dr. Jacobellus, the Pope of Saxony, the planet of Swabia, the apostle of ubiquity, \emph{allotrio-episcopus}, a worshiper of Bacchus and Mammon,' etc. He no doubt had a considerable share of vanity, ambition, and theological passion (which he displayed, e.g., against poor Flacius, even after his death). But there is no reason to doubt the general purity of his motives, and, compared with some other orthodox Lutherans of his age, he was even liberal, at least in his earlier years. At a later period he denounced the alterations of the Augsburg Confession, and compared Melanchthon to Solomon, who at first wrote glorious things, but was afterwards so far led astray that the Bible leaves it doubtful whether he were saved ('\emph{ob er zu unserm Herrgott oder zu dem Teufel gefahren sei}'). He seemed to be predestinated for the work of his life. Planck gives a masterly (though not altogether just) analysis of his character, from which I quote a specimen, as it fairly represents the spirit and style of his celebrated history (Vol. VI. p. 274): '\emph{In halb Deutschland herumzureisen, und an jedem neuen Ort mit neuen Menschen zu unterhandlen—hier mit dem Ministerio einer Reichsstadt, und dort mit einer kleinen Synode von Superintendenten, welche die Geistlichkeit einer ganzen Grafschaft oder eines Fürstenthums repräsentiren—heute mit Flacianern und morgen mit Anhängern der Wittenbergischen Schule und Verehrern Melanchthons—jetzt mit den Hauptpersonen, die an dem gelehrten Streit den vorzüglichsten Antheil genommen, und jetzt mit den Schreiern, die bloss den Lärm vermehrt, und dazwischen hinein mit einem oder dem andern Stillen im Lande, die bisher im Verborgenen über den Streit geseufzt hatten—und allen diesen Menschen alles zu werden, um sie zu gewinnen—es gab wirklich kein Geschäft in der Welt, das für ihn so gemacht war, wie dieses, so wie es auch umgekehrt wenige Menschen gab, die für das Geschäft so gemacht waren, wie er. Nimmt man aber noch dies dazu, dass sich auch der gute, Andreæ selbst dazu für gemacht hielt, dass in die natürliche Thätigkeit seines Geistes auch zuweilen ein kleiner Windzug von Ehrgeiz und Eitelkeit hineinblies, dass er auch für den Reiz der bedeutenden Rolle, die er dabei spielen, und des Aufsehens, das er erregen würde, nicht unfühlbar war, ja dass selbst der Gedanke an das} [\emph{den}] \emph{Verkehr, in das er dabei mit so manchen Fürsten und Herrn kommen, an die Ehrenbezeugungen, die man ihm hier und da erweisen, an die Raths-Deputationen, die ihn in so mancher kleinen Reichsstadt bewillkommen, an die Gastpredigten, die man ihm auftragen, und an die Ehrfurcht, womit dann die ehrliche Bürger einer solchen Stadt, die noch keinen Kanzler von Tübingen gesehen hatten, mit Fingern auf ihn weisen würden—dass auch der Gedanke daran den heiteren und offenherzigen Mann, der es mil seinen kleinen Schwachheiten nicht so genau nahm und sie eben so leicht sich selbst as andern vergab, auf gewisse Augenblicke sehr stark anziehen konnte—nimmt man alles diess zusammen, so wird man auch hinreichend erklärt haben, wie es kommen konnte, dass er vor den Schwierigkeiten seines übernommenen Geschäfts nicht erschrak, die sich ihm doch ebenfalls bei seiner Klugheit, bei seiner Weltkenntniss, und bei seiner besondern durch manche Erfahrung erkauften Kenntniss der Menschen, die er dabei zu bearbeiten hatte, lebhafter als hundert andern darstellen mussten. Gewiss standen auch diese Schwierigkeiten lebhaft genug vor seiner Seele, aber der Reiz, durch den er in das Geschäft hineingezogen wurde, war so stark, dass er ihm schwerlich hätte widerstehen können, wenn er nicht nur die Mühe und Arbeit, die es ihn kosten, sondern auch den tausendfachen Verdruss, den es ihm machen, die zahllosen Kränkungen, die es ihm zuziehen, und selbst alle die stechenden Erinnerungen, durch die es ihm sein Alter verbittern sollte, vorausgesehen hätte.}' Andreæ, in connection with Vergerius, founded the first Bible Society, for Sclavonic nations (1555). His grandson, Johann Valentin Andreæ (1586–1654), was a man of genius and more liberal views, and a great admirer of the order and discipline of the Reformed Church in Geneva, which he sadly missed in Germany.} Next to him, and at a later period, Martin Chemnitz (1522–1586), the greatest pupil of Melanchthon and the prince among the Lutheran divines of his age,\footnote{Author of \emph{Loci theologici; Examen Concilii Tridentini; Harmonia Evangeliorum} (completed by Polycarp Leyser and John Gerhard); \emph{De duabus in Christo naturis}, and other works of vast learning. The Romanists called him a second Martin Luther, and said: '\emph{Si posterior non fuisset, prior non stetisset.}' This reminds one of the line, '\emph{Si Lyra non lyrasset, Lutherus non saltasset.}'} and Nicholas Selnecker (1530–1592),\footnote{He prepared the second Latin translation of the Form of Concord, and is best known by one of his hymns ('\emph{Ach bleib bei uns, Herr Jesu Christ},' etc.; although it is only in part from him). His numerous theological writings are forgotten. He was a little man with short legs, at first a Philippist, then a rigid Lutheran ('\emph{parvus Flacius}'); hence in turn attacked by all parties. '\emph{Die Reformirten, gegen die er den Vers wandte: ``Erhalt uns Herr bei deinem wort und steur' der}  Zwinglianer  \emph{Mord!'' und denen er die Schändung seiner Tochter in letzter Instanz zuchreiben zu müssen glaubte, nannten ihn das ``Lutheräfflein;'' bei den strengen Lutheranern hiess er: ``Schelmlecker, Seelhenker, Seelnecator;'' bei den Melanchthonianern: ``Judas alter in suspensus,'' Auch mit seinem Freund Andreæ ist er zuletzt zerfallen. . . . Ein Jahrhundert später wurde er unter die deutschen Propheten gerechnet.}' G. Frank, Vol. I. p. 221.} originally likewise a Melanchthonian, took the most important part in the movement, and formed with Andreæ the theological 'triumvirate,' which finally completed the Form of Concord.\footnote{The remaining three authors were David Chytræus, Professor in Rostock (d. 1600), who remained a faithful Melanchthonian, and met the violent abuse of the zealots with silence; Andreas Musculus, Professor in Frankfort-on-the-Oder (d. 1581), who denounced Melanchthon as a patriarch of all heretics, and praised Luther as the sun among the dim stars of the old fathers; and Christopher Körner, Professor in Frankfort-on-the-Oder, a friend of Chytræus, but unfortunate in his children, who sunk into the lowest vices (G. Frank, Vol. I. p. 222).}

The first attempts at union were made at the conferences in Frankfort, 1558; Naumburg, 1561; Altenburg, 1568; Wittenberg, 1569; Zerbst, 1570; Dresden, 1571; but they utterly failed and increased the dissension.

After the violent suppression of Crypto-Calvinism in Electoral Saxony (1574), and the death of Flacius (1575) and some other untractable extremists, the work was resumed by the Elector and other Princes. Theological conferences were again held at Maulbronn (1575), Lichtenberg (1576), and Torgau (1576). Three forms of agreement were prepared, which, though not satisfactory, served as a basis for the Formula of Concord. The first is the \emph{Swabian and Saxon Formula}, written by Andreæ (1574), and revised by Chemnitz and Chytræus (1575).'\footnote{\emph{Schwabisch-Sächsische Concordie, Formula Suevica et Saxonica}, or \emph{Formula Concordiæ inter Suevicas et Saxonicas Ecclesias}, published from MS., in the original and revised form, by Heppe, \emph{Geschichte des Deutschen Protest.} Vol. III., \emph{Beilagen}, pp. 75–166, and 166–325. They were preceded by six sermons of Andreæ (1573). Likewise republished by Heppe.} The second is the \emph{Maulbronn Formula}, prepared by the Swabian divines Lucas Osiander and Balthasar Bidembach (Nov. 14, 1575), and approved by a convent of Lutheran Princes in the Cloister of Maulbronn (Jan. 19, 1576).\footnote{See Heppe, Vol. III. pp. 76 sqq.} The former was found too lengthy, the latter too brief. Hence on the basis of both a third form was prepared which combined their merits, but omitted the honorable mention of the name of Melanchthon. This is the '\emph{Torgau Book},' consisting of twelve articles.\footnote{The 'Torgische Buch' or '\emph{Torgisch Bedenken, welchergestalt oder massen vermöge Gottes Worts die eingerissene Spaltungen zwischen den Theologen Augsburgischer Confession christlich verglichen und beigelegt werden möchten, anno} 1576.' It was republished by Semler, with Preface and notes, Halle, 1760, but much better by Heppe, Marburg, 1857; second edition, 1866.} It was mainly the work of Andreæ and Chemnitz, and completed by a convention of eighteen Lutheran divines at the Castle of Hartenfels, at Torgau, June 7, 1576. It was sent by the Elector Augustus to all the Lutheran Princes for examination and revision. It was closely scrutinized by twenty conventions of theologians held within three months, and elicited twenty-five vota, mostly favorable; even Heshusius and Wigand, the oracles of orthodoxy, were pleased, except that they wished an express condemnation of Melanchthon and other 'authors and patrons of corruptions.'

At last the present Formula of Concord was completed, on the basis of the Torgau Book, by six learned divines—Andreæ (of Tübingen), Chemnitz (of Brunswick), Selnecker (of Leipzig), Musculus (of Frankfort-on-the-Oder), Cornerus, or Körner (also of Frankfort), and Chytræus (of Rostock)—who met in March and May, 1577, in the Cloister of Bergen, near Magdeburg, by order of the Elector of Saxony. Hence it is also called '\emph{The Bergen Formula.}'\footnote{Or, \emph{Das Bergische Buch.} English writers usually call it 'Form of Concord,' though 'Formula' is more correct.} The Preface was written two years later by the same authors, in the name of the Lutheran Princes, in two conventions at Jüterbock, January and June, 1579. Three years elapsed before the new symbolical book was signed and solemnly published, by order of Augustus, at Dresden, June 25, 1580, the fiftieth anniversary of the Augsburg Confession, together with the other Lutheran symbols, in one volume, called the 'Book of Concord,' which superseded all similar collections.\footnote{See the titles on p. 220, and literary notices in Köllner, pp. 562 sqq. Andreæ directed the editing of the \emph{German} Book of Concord, Glaser and Fuger read the proof. The manuscript was deposited in the library of the chief church at Dresden, and burned up with it July 19, 1700. The first \emph{Latin} Concordia (1580) was superintended and edited, though without proper authority, by Selnecker; the second edition (1584) was issued by authority of the Electors. There are few separate editions of the Formula of Concord, the first by Selnecker, Lipz. 1582. See Köllner, p. 561.} The Elector Augustus celebrated the completion of the work, which cost him so much trouble and money, by a memorial coin representing him in full armor on the storm-tossed ship of the church.\footnote{See a description in Penzel's \emph{Saxon. Numism.} as quoted by Planck, Vol. VI. p. 689. Augustus dismissed Andreæ (1580), ostensibly with great honor and rich presents, but in fact much displeased with the \emph{garrulus Suevus}, who had spoken disrespectfully of his theological ignorance, had fallen out with Chemnitz and Selnecker, and made many enemies. See a full account in Heppe. Vol. IV. pp. 256–270.}

The Formula of Concord, like the three preparatory drafts on which it is based, was first composed in the German language, and published, with the whole Book of Concord, at Dresden, 1580. The Latin text was imperfectly prepared by Lucas Osiander, and appeared in the Latin \emph{Concordia}, at Leipzig, 1580; then it was materially improved by Selnecker for his separate German-Latin edition of the Formula (not the Book) of Concord, Leipzig, 1582; and was again revised by a convent of Lutheran divines at Quedlinburg, 1583, under the direction of Martin Chemnitz. In this last revision it was published in the first \emph{authentic} Latin edition of the Book of Concord, Leipzig, 1584, and has been recognized ever since as the received Latin text. It was also translated into the Dutch, Swedish, and English languages, but seldom separately published.\footnote{See the authorized Latin text of the Epitome, with a new English translation, in Vol. III. pp. 93 sqq. An English Version of the Formula from the German text appeared in \emph{The Christian Book of Concord}; or, \emph{Symbolical Books of the Evangelical Lutheran Church}, New Market, Va., 1851, 2d ed., 1854. It professes to be literal, hut is very stiff and unidiomatic.}

\xsection{The Form of Concord, Concluded.}

§ 46. The Form of Concord, Concluded.

\emph{Analysis and Criticism.}

The Formula of Concord consists of two parts—the \emph{Epitome} and the \emph{Solida Repetitio et Declaratio.} Both treat, in twelve articles, of the same matter—the first briefly, the other extensively. They begin with the anthropological doctrines of original sin and freedom of the will; next pass on to the soteriological questions concerning justification, good works, the law and the gospel, the third use of the law; then to the eucharist and the person of Christ; and end with foreknowledge and election. This order is characteristic of the Lutheran system, as distinct from the Calvinistic, which begins with the Scriptures, or with God and the eternal decrees. The most important articles are those on the Lord's Supper and the Person of Christ, which teach the peculiar features of the Lutheran creed, viz., consubstantiation, the communication of the properties of the divine nature to the human nature of Christ, and the ubiquity of Christ's body.

The \emph{Epitome} contains all that is essential. It first states the controversy (\emph{status controversiæ}), then the true doctrine (\emph{affirmativa}), and, last, it condemns the error (\emph{negativa}). In the \emph{Solid Repetition and Declaration} this division is omitted; but the articles are more fully explained and supported by ample quotations from the Scriptures, the fathers, the older Lutheran Confessions, and the private writings of Dr. Luther, which swell it to about five times the size of the\emph{Epitome.}

Each part is preceded by an important introduction, which lays down the fundamental Protestant principle that the Canonical Scriptures are the only rule of faith and doctrine,\footnote{'\emph{Die einige Regel und Richtschnur} (\emph{unica regula et norma}), \emph{nach welcher alle Lehren and Lehrer gerichtet und geurtheilt werden sollen.}' Comp. Psa. cxix. 15; Gal. i. 8. The extent of the Canon, however, is not defined, as in several Reformed Confessions, and the question of the Apocrypha of the Old Testament is left open.} and fixes the number of (nine) symbolical books to be hereafter acknowledged in the Lutheran Church, not as judges, but as witnesses and expositions of the Christian faith; namely, the three œcumenical Symbols (the Apostles', the Nicene, and the Athanasian), the Unaltered Augsburg Confession,\footnote{\par '\emph{Die erste ungeänderte Augsb. Confession}' (\emph{Augustanam illam primam et non mutatam Confessionem}). The Preface (pp. 13, 14) rejects the Altered Augsburg Confession (of 1540), if it be understood as teaching another doctrine of the Lord's Supper.} the Apology of the Confession, the Articles of Smalcald, the Smaller and Larger Catechisms of Luther,\footnote{These are called the '\emph{Laienbibel}' (\emph{laicorum biblia}, the layman's Bible), '\emph{darin alles begriffen, was in heiliger Schrift weitläuftig gehandelt, und einem Christenmenschen zu wissen vonnöthen ist.}'} and the Formula of Concord. The Scriptures contain the \emph{credenda}, the things to be believed; the Symbols the \emph{credita}, the things that are believed. Yet the second part of the Formula quotes Dr. Luther, '\emph{piæ sanctœque memoriæ},' as freely, and with at least as much deference to his authority, as Roman Catholics quote the fathers. Melanchthon, the author of the fundamental Confession of the Lutheran Church, is never named, but indirectly condemned; and as to poor Zwingli, he is indeed mentioned, but only to be held up to pious horror for his \emph{'blasphemous allæosis.}'\footnote{\emph{Sol. Decl.} Art. VIII. p. 678 (ed. Müller): '\emph{Die gotteslästerliche allæosis Zwinglii},' which Dr. Luther condemned '\emph{als des Teufels Larve, bis in den Abgrund der Höllen.}'} Thus the supremacy of the Bible is maintained in principle, but Luther is regarded as its regulative and almost infallible expounder.

We now proceed to give a summary of the Formula.

Art. I. Of Original Sin.—It is not the moral essence, or substance, or nature of man (as Flacius taught with the old Manichæans), but a radical corruption of that nature, which can never be entirely eradicated in this world (against the Pelagian and semi-Pelagian heresies).

Art. II. Of Free Will.—Man, in consequence of Adam's fall, has lost the divine image, is spiritually blind, disabled, dead, and even hostile to God, and can contribute nothing towards his conversion, which is the work of the Holy Spirit alone, through the means of grace. The Formula, following Luther, uses stronger terms on the slavery of the will and total depravity than the Calvinistic Confessions. It compares the unconverted man to a column of salt, Lot's wife, a statue without mouth or eyes, a dead stone, block and clod,\footnote{\emph{Solida Declaratio}, Art. II. § 24 (p. 662 ed. Rech., p. 534 ed. Müller): '\emph{Autequam homo per Spiritum Sanctum illuminatur, convertitur, regeneratur et trahitur . . . ad conversionem aut regenerationem suam nihil inchoare, operari, aut coöperari potest, nec plus quam lapis, truncus, aut limus} (\emph{so wenig als ein Stein oder Block oder Thon})'. Thomasius und Stahl disapprove of these expressions, and Luthardt (\emph{Lehre v. freien Willen}, p. 272) admits, at least, that they are unfortunately chosen (\emph{unglücklich gewählt}). Fr. H. R. Frank defends them.} and denies to him the least spark of spiritual power.\footnote{\emph{Ibid.} Art. II. § 7 (p. 656 ed. Rech., p. 589 ed. Müller): . . . '\emph{homo ad bonum prorsus corruptus et mortuus sit, ita ut in hominis natura post lapsum ante regenerationem ne scintillula quidem spiritualium virium} (\emph{nicht ein Fünklein der geistlichen Kräfte}) \emph{ reliqua manserit aut restet, quibus ille ex se ad gratiam Dei præparare se aut oblatam gratiam apprehendere, aut eius gratiæ} (\emph{ex sese et per se}) \emph{capax esse possit, aut se ad gratiam applicare aut accommodare, aut viribus suis propriis aliquid ad conversionem suam vel ex toto vel ex dimidia vel ex minima parte conferre, agere, operari aut coöperari} (\emph{ex se ipso tanquam ex semet ipso}) \emph{possit} (\emph{oder aus seinen eigenen Kräften etwas zu seiner Bekehrung, weder zum ganzen noch zum halben oder zu einigem dem wenigsten oder geringsten Theil, helfen, thun, wirken oder mitwirken vermöge, von ihm selbst, als von ihm selbst}). . . . \emph{Inde adeo naturale tiberum arbitrium, ratione corruptarum virium et naturæ suæ depravatæ, duntaxat ad ea, quæ Deo displicent et adversantur, activum et efficax est.}' This and similar statements are followed by quotations from Dr. Luther, where he compares the natural man to 'a column of salt, Lot's wife, a clod and stone, a dead statue without eyes or mouth.' All he said against Erasmus, and later, in his Commentary on Genesis, about free will, is indorsed. Flacius inferred from the same teacher his Manichæan error, which the Formula condemns in Art. I.} He can not even accept the gospel (which is the work of pure grace), but he may \emph{reject} it, and thereby incur damnation.

This article condemns the fatalism of the Stoics and Manichæans, the anthropological heresies of the Pelagians and Semi-Pelagians, but also and especially the Synergism of Melanchthon and the Philippists. The chief framers of the Formula—Andreæ, Chemnitz, Selnecker, and Chytræus—were at first in favor of Synergism, which would have been more consistent with Article XI.; the Swabian-Saxon Concordia, drawn up by Chemnitz and Chytræus, and the Torgau Book actually contained synergistic passages.\footnote{See these passages in Gieseler, Vol. IV. p. 486, note 24; Heppe, \emph{Der Text der Bergischen Concordienformel verglichen, }etc.; Luthardt, \emph{Lehre vom freien Willen, }pp. 262 sqq. Comp. also the remarks of Planck, Vol. VI. pp. 718 sqq.} But they were omitted or exchanged for others, and consistency was sacrificed to veneration for Luther.

There is an obvious and irreconcilable antagonism between Art. II. and Art. XI. They contain not simply opposite truths to be reconciled by theological science, but contradictory assertions, which ought never to be put into a creed. The Formula adopts one part of Luther's book \emph{De servo arbitrio} (1525), and rejects the other, which follows with logical necessity. It is Augustinian—yea, hyper-Augustinian and hyper-Calvinistic in the doctrine of human depravity, and anti-Augustinian in the doctrine of divine predestination. It indorses the anthropological premise, and denies the theological conclusion. If man is by nature like a stone and block, and unable even to accept the grace of God (as Art. II. teaches), he can only be converted by an act of almighty power and \emph{irresistible} grace (which Art. XI. denies). If some men are saved, without any co-operation on their part, while others, with the same inability and the same opportunities, are lost, the difference points to a particular predestination and the inscrutable decree of God. On the other hand, if God sincerely wills the salvation of all men (as Art. XI. teaches), and yet only a part are actually saved, there must be some difference in the attitude of the saved and the lost towards converting grace (which is denied in Art. II.).

The Lutheran system, then, to be consistent, must rectify itself, and develop either from Art. II. in the direction of Augustinianism and Calvinism, or from Art. XI. in the direction of Synergism and Arminianism. The former would be simply returning to Luther's original doctrine, which he never recalled, though he may have modified it a little; the latter is the path pointed out by Melanchthon, and adopted more or less by some of the ablest modern Lutherans.\footnote{As Thomasius, Stahl, Harless, Hoffmann, Luthardt, Kahnis. See Luthardt, \emph{Die Lehre vom freien Willen}, pp. 378 sqq.} In either case the second article needs modification. It uses the language of feeling rather than sober reflection, and gives the rhetorical expressions of subjective experience the dignity of symbolical statement. We can, indeed, not feel too strongly the sinfulness of sin and the awful corruption of our hearts. Nevertheless, God's image in man is not lost or exchanged for Satan's image, but only disfigured, disabled, and lying in ruins. Man is, indeed, in his prevailing inclination, a slave of sin, yet susceptible of the influences of divine grace, and remains moral and responsible in accepting or rejecting the gospel, before as well as after conversion. His reason, his conscience, his sense of sin, his longing for redemption and for peace with God, his prayers, his sacrifices, and all the '\emph{testimonia animæ naturaliter christianæ},' bear witness with one voice to his divine origin, his divine destination, and his adaptation to the Christian salvation.\footnote{Well says Goethe—  \par \emph{'Wär' nicht das Auge sonnenhaft,} \par \emph{Wie könnte en das Licht erblicken! } \par \emph{Lebt' nicht in uns des Gottes eigne Kraft,} \par \emph{Wie könnt' uns Göttliches entzücken?}'} But on the other hand there are innumerable mysteries of Providence in the order of nature as well as of grace, and inequalities in the distribution of gifts and opportunities, which baffle solution in this present world, and can only be traced to the inscrutable wisdom of God. The human mind has not been able as yet satisfactorily to set forth the harmony of God's sovereignty and man's responsibility.

Art. III. Of Justification by Faith.—Christ is our righteousness, not according to the divine nature alone (Andrew Osiander), nor according to the human nature alone (Stancar), but the whole Christ. God justifies us out of pure grace, without regard to antecedent, present, or subsequent works or merit, by imputing to us the righteousness of the obedience of Christ. Faith alone is the medium and instrument by which we apprehend Christ. Justification is a declaratory or forensic act—a sentence of absolution from sin, not an infusion of righteousness (Osiander).

Art. IV. Of Good Works.—Good works must always follow true faith, but they are neither necessary to salvation (Major), nor dangerous or injurious to salvation (Amsdorf). Salvation is of free grace alone, apprehended by faith.

Art. V. Of the Law and the Gospel.—The object of the law is to reprove sin and to preach repentance; the gospel (in its specific sense) is a joyful message, the preaching of Christ's atonement and satisfaction for all sins.

Art. VI. Of the Third Use of the Law—i.e., its obligation to believers, as distinct from its civil or political, and its pædagogic or moral use in maintaining order, and leading to a conviction of sin. Believers, though redeemed from the curse and restraint of the law, are bound to obey the law with a free and willing spirit. Antinomianism is rejected.

Art. VII. Of the Lord's Supper.—The most important controversy and chief occasion of the Formula—hence the length of this Article in the second part. It sets forth clearly and fully the doctrine of \emph{consubstantiation} (as it is usually called, in distinction from the Romish \emph{transubstantiation}), i.e., of the co-existence of two distinct yet inseparable substances in the sacrament. It is the doctrine of the \emph{real} and \emph{substantial} presence of the true body and blood of Christ \emph{in, } \emph{with}, and \emph{under} the elements of bread and wine (\emph{in, cum, et sub pane et vino}), and the \emph{oral} manducation of both substances by \emph{unbelieving} as well as believing communicants, though with opposite effects. The sacramental union of Christ's real body and blood with the elements is not an impanation or local inclusion, nor a mixture of two substances, nor a permanent (extra-sacramental) conjunction, but it is illocal, supernatural, unmixed, and confined to the sacramental transaction or actual use.\footnote{'\emph{Nihil habet rationem sacramenti extra usum, seu actionem divinitus institutam}' (\emph{Sol. Decl.} p. 663). Gerhard and the later Lutheran theologians describe the presence as \emph{sacramentalis, vera et realis, substantialis, mystica, supernaturalis et incomprehensibilis}, and distinguish it from the \emph{præsentia gloriosa} (in heaven), \emph{hypostatica} (of the \textgreek{λόγος} in the human nature), \emph{spiritualis} (\emph{operativa}, or \emph{virtualis}), \emph{figurativa} (\emph{imaginativa, symbolica}). It is a \textgreek{παρουσία}, not an \textgreek{ἀπουσία} (absence), nor \textgreek{ἐνουσία} (inexistence), nor \textgreek{συνουσία} (co-existence in the sense of coalescence), nor \textgreek{μετουσία} (transubstantiation). They reject the term \emph{consubstantiation} in the sense of impanation or incorporation into bread, or physical coalescence and fusion. The Formula itself does not use the term.} Nor is it effected by priestly consecration, but by the omnipotent power of God, and the word and institution of Christ. The body of Christ is eaten with the mouth by all communicants, but the notion of a Capernaitic or physical eating with the teeth is indignantly rejected as a malignant and blasphemous slander of the sacramentarians.\footnote{And yet Dr. Luther himself unequivocally taught the literal \emph{mastication} of Christ's body. He gave it as the sum of his belief, to which he 'would adhere though the world should collapse,' that Christ's body was '\emph{ausgetheilt, gegessen und mit den Zähnen zerbissen}' (\emph{Briefe}, ed. by De Witte, Vol. IV. p. 572, comp. p. 569). He instructed Melanchthon to insist on this in the conference he had with Bucer in Cassel, Dec. 1534; but Melanchthon, though not emancipated from Luther's view at that time, declined to shoulder it as his own, and began to change his ground on the eucharistic question. \emph{Corp. Ref.} Vol. II. p. 822. Comp. Schmidt, \emph{Mel.} p. 319; Ebrard, \emph{Abendmahl}, Vol. II. pp. 375 sqq.}

The Formula condemns the Romish dogma of transubstantiation, the sacrifice of the mass, and the withdrawal of the cup from the laity, but with equal or greater emphasis the Reformed and Melanchthonian (Crypto-Calvinistic) theory of a spiritual real presence and fruition of Christ by faith, or by believers only, without making a distinction between Zwinglians and Calvinists, except that the latter are called 'the most pernicious of all sacramentarians.'\footnote{Planck (Vol. VI. pp. 732 sqq.) charges the Formula with \emph{willful} misrepresentation of Calvin's view, which he had so clearly, distinctly, and repeatedly set forth, especially in his tracts against Westphal, and which had since been embodied in the Confessions of the Reformed churches. Thomasius, Stahl, and other orthodox Lutherans, freely admit the material difference between Calvin and Zwingli in the theory of the eucharist.}

Art. VIII. Of the Person of Christ.—This article gives scholastic support to the preceding article on the eucharistic presence, and contains an addition to the Lutheran creed. It teaches the \emph{communicatio idiomatum} and the ubiquity of Christ's body. It raised the private opinions and speculations of Luther, Brentius, and Chemnitz on these topics to the authority of a dogma. Some regard this as the crowning excellence of the Formula;\footnote{My friend, Dr. Krauth, goes so far as to say (1.c. p. 316): 'The doctrine of the person of Christ presented in the Formula rests upon the sublimest series of inductions in the history of Christian doctrine. In all confessional history there is nothing to be compared with it in the combination of exact exegesis, of dogmatic skill, and of fidelity to historical development. Fifteen centuries of Christian thought culminate in it.' But in his lengthy exposition he does not even mention the important difference between the Swabian and Saxon schools, nor the various forms of the \emph{communicatio idiomatum}, and evades the real difficulty by resolving, apparently (p. 318), the communication of divine properties into an efficacious manifestation of the Godhead in and through the assumed humanity of Christ—which has never been disputed by Reformed divines.} others, even in the Lutheran communion, as its weakest and most assailable point.\footnote{Even Luthardt admits at least the artificial construction of the Christology of the Formula, and its inconsistency with the historical realness of the picture of Christ in the Gospels (\emph{Compend. der Dogmatik}, p. 144; comp. also Kahnis, \emph{Luth. Dogmatik}, Vol. III. p. 338 sq.). The modern Lutheran Kenoticists, Thomasius, Hofmann (Luthardt inclines to them, p. 155)—not to speak of the extreme form to which Gess carried the \textgreek{κένωσις}—virtually depart from the Formula of Concord, which pronounces it a 'blasphemous perversion' to explain  Matt. xxviii. 18 ('all power is \emph{given} to me,' etc.) in the sense that Christ had ever laid aside or abandoned his almighty power in the state of humiliation (\emph{Epit.}, at the close of Art. VIII.).} It was certainly very unwise, as history has shown, to introduce the scholastic subtleties of metaphysical theology into a public confession of faith.

The Formula derives from the personal union of the two natures in Christ (\emph{unio hypostatica}, or \emph{personalis}) the communion of natures (\emph{communio naturarum}), from the communion of natures the communication of properties or attributes (\emph{communicatio idiomatum}, a term used first by the scholastics), and from the communication of properties the omnipresence or ubiquity of Christ's body. The controversy between the Lutheran and Reformed, who both professedly stand on the common theanthropic Christology of Chalcedon, refers to the nature and extent of the communication of properties, and especially to the ubiquity of Christ's body derived therefrom.

The Formula (in the Second Part) distinguishes three kinds of the \emph{communicatio idiomatum}, which were afterwards more fully analyzed, defined, and designated by the Lutheran scholastics of the seventeenth century.\footnote{We anticipate, for the sake of clearness, from the later orthodox writers the names of the three \emph{genera}. The substance is already in the Formula, and in the treatise of Chemnitz, \emph{De duabus naturis in Christo}, 1580. For a fuller exposition, with ample quotations from Chemnitz, John Gerhard, Hafenreffer, Hutter, Calov, Quenstedt, König, Baier, Hollaz, see Heinrich Schmid's \emph{Dogmatik der evang. lutherischen Kirche} (2d ed. 1847), pp. 252 sqq.; comp. also Luthardt, pp. 144 sqq., and Kahnis, Vol. II. pp. 335 sqq.}

1. The \emph{genus idiomaticum}, by which the attributes of one or the other nature are communicated to the whole person. Thus it is said that 'the Son of God was made of the seed of David, according to the flesh' (Rom. i. 3), that 'Christ was put to death in the flesh,' and that 'he suffered in the flesh' (1 Pet. iii. 18;  iv. 1).\footnote{This genus was subsequently subdivided into three species, corresponding to the \emph{concretum} of the divine nature, the \emph{concretum} of the human nature, and the \emph{concretum} of both natures, of which the \emph{idiomata} are predicated, viz., (\emph{a}) \textgreek{ἰδιοποίησις,} or \textgreek{οἰκείσις,} i.e., '\emph{appropriatio, quando idiomata humana de concreto divinæ naturæ enuntiantur},' Acts iii. 15; xx. 28; 1 Cor. ii. 8; Gal. ii. 20; Psa. xlv. 8. (\emph{b}) \textgreek{Κοινωνία τῶν θείων,} '\emph{divinorum idiomatum, quando de persona verbi incarnati, ab humana natura denominata, idiomata divina ob unionem personalem enuntiantur},' John vi. 62; viii. 58; 1 Cor. xv. 47. (\emph{c}) \textgreek{Ἀντίδοσις,} or \textgreek{συναμφοτερισμός,} '\emph{alternatio s. reciprocatio, qua tam divina quam humana idiomata de concreto personæ sive de Christo, ab utraque natura denominato, prædicantur},' Heb. xiii. 8; Rom. ix. 5; 2 Cor. xiii. 4; 1 Pet. iii. 18. See Schmid, p. 258.} Here Luther's warning is quoted against Zwingli's \emph{allœosis}, as 'a mask of the devil.'

2. The \emph{genus apotelesmaticum}, or the \textgreek{κοινωνία ἀποτελεσμάτων},\footnote{\par 'The expression is borrowed from John of Damascus. \textgreek{ἀποτέλεσμα} means properly \emph{completion} of the work (\emph{consummatio operis}), \emph{effect}, \emph{result}; but it is here used for each action in the threefold office of Christ.} which has reference to the execution of the office of Christ: the communication of redeeming acts of the whole person to one of the two natures. Christ always operates in and through both. Thus Christ, neither as God nor man alone, but as God-man, is our Mediator, Redeemer, King, High-Priest, Shepherd, etc. He shed his blood according to his human nature, but the divine nature gave it infinite value (1 Cor. xv. 3: 'Christ died for our sins;' Gal. i. 4; iii. 17; 1 John iii. 8; Luke ix. 56).

3. The \emph{genus majestaticum}, or \emph{auchematicum},\footnote{\par From \textgreek{αὔχημα, } \emph{gloria.} This genus is also called \textgreek{βελτίωσις, ὑπερύψωσις, μετάδοσις, θέωσις, ἀποθεοσία, θεοποίησις, } \emph{unctio.}} i.e., the communication of the attributes of the divine nature to the assumed humanity of Christ. 'The human nature of Christ,' says the Formula, 'over and above its natural, essential, and permanent human properties, has also received special, high, great, supernatural, inscrutable, ineffable, heavenly prerogatives and pre-eminence in majesty, glory, power, and might, above all that can be named (Eph. i. 21).'\footnote{\emph{Sol. Decl.} Art. VIII. p. 685 (ed. Müller).}. . . 'This majesty of the human nature was hidden and restrained in the time of the humiliation. But now, since the form of a servant is laid aside, the majesty of Christ appears fully, efficiently, and manifestly before all the saints in heaven and on earth, and we also in the life to come shall see his glory face to face (John xvii. 24). For this reason, there is and remains in Christ only one divine omnipotence, power, majesty, and glory, which is the property of the divine nature alone; but this shines forth, exhibits, and manifests itself fully, yet spontaneously, in, with, and through the assumed, exalted human nature in Christ; precisely as to shine and to burn are not two properties of iron, but the power to shine and to burn is the property of the fire—but since the fire is united with the iron, it exhibits and manifests its power to shine and to burn in, with, and through this red-hot iron; so that also the red-hot iron, through this union, has the power to shine and to burn, without a change of the essence and of the natural properties of the fire or of the iron.'\footnote{P. 689.}

The Lutheran scholastics make here a distinction between the operative attributes (omnipotence, omniscience, omnipresence) and the quiescent attributes (eternity, infinitude): all were communicated to Christ for inhabitation and possession, but only the operative for use—\textgreek{χρῆσις, } \emph{usurpatio} (Matt. xxviii. 18; John xvii. 2, 5, 27; Col. ii. 3).

4. Strict logic would require a fourth genus (\emph{genus } \textgreek{ταπεινωτικόν,} namely, the communication of the attributes of the human nature to the divine nature. But this is rejected by the Formula and the Lutheran scholastics, on the ground that the divine nature is unchangeable, and received no accession nor detraction from the incarnation.\footnote{\emph{Sol. Decl.} p. 684: '\emph{Was die göttliche Natur in Christo anlanget, weil bei Gott keine Veränderung ist} (\emph{Jac.} 1,17), \emph{ist seiner göttlichen Natur durch die Menschwerdung an ihrem Wesen und Eigenschaften nichts ab-oder zugegangen, ist in oder für sich dadurch weder gemindert noch gemehret.}' This raises the question how far the unchangeableness of God is affected by the incarnation, about which Dr. Dorner has written some profound articles in the \emph{Jahrbücher für Deutsche Theologie}, 1856 and 1858.} This is a palpable inconsistency,\footnote{As Thomasius and Kahnis (Vol. III. p. 339) admit.} and is fatal to the third genus. For if there is any real communication of the properties of the two natures, it must be mutual; the one is the necessary counterpart of the other. If the human nature is capable of the divine, the divine nature must be capable of the human; and if, on the other hand, the divine nature is incapable of the human, the human nature must be incapable of the divine. Luther felt this, and boldly uses such expressions as 'God suffered,' 'God died,' which were familiar to the Monophysites.\footnote{'\emph{Weil Gottheit und Menschheit,}' he says (Vol. XXX. p. 204, Erl. ed.), '\emph{Eine Person ist, so giebt die Schrift um solcher persönlichen Einigkeit willen auch alles, was der Menschheit widerfährt, der Gottheit, und wiederum. Und ist auch also in der Wahrheit. Denn da musst du ja sagen: Die Person leidet, stirbt; nun ist die Person wahrhaftiger Gott: durum ist's recht geredet: Gottes Sohn leidet.}'}

The battle-ground between the Lutheran and the Reformed is the \emph{genus majestaticum}, for which John of Damascus had prepared the way. But just here the Formula is neither quite clear nor consistent. It was unable to harmonize the two different Lutheran Christologies represented among the authors by Andreæ and Chemnitz.\footnote{See above, pp. 290–294.} It teaches, on the one hand (to guard against the charge of Eutychianism and Monophysitism), that the attributes of the divine nature (as omnipotence, eternity, infinitude, omnipresence, omniscience) 'can never become (intrinsically and \emph{per se}) the attributes of the human nature,' and that the attributes of the human nature (as corporeality, limitation, circumscription, passibility, mortality, hunger, thirst) 'can never become the attributes of the divine nature.'\footnote{\emph{Epit.} VIII. (p. 545, ed. Müller): '\emph{Wir gläuben, lehren und bekennen, dass die göttliche und menschliche Natur nicht in ein Wesen vermenget, keine in die andere verwandelt, sondern ein jede ihre wesentliche Eigenschaften behalte, } Welche der andern Natur Eigenschaften Nimmermehr Werden.  \emph{Die Eigenschaften göttlicher Natur sind: allmächtig, ewig}, etc., \emph{sein, welche der menschlichen Natur Eigenschaften nimmermehr werden. Die Eigenschaften menschlicher Natur sind: ein leiblich Geschöpf oder Creatur sein}, etc., \emph{welche der göttlichen Natur Eigenschaften nimmermehr werden.}' Comp. the \emph{Sol. Decl.} Art. VIII.} (This quite agrees with the doctrine of Chemnitz and of the Reformed theologians.) But, on the other hand (in opposition to Nestorianism and the 'sacramentarians,' as the Reformed are called), the Formula asserts that, by virtue of the hypostatic or personal union of the two natures and the communion of natures, one nature may, nevertheless (by derivation and dependency), partake of the properties of the other, or at least that the human nature, while retaining its inherent properties, may and does receive (as peculiar prerogatives, or as \emph{dona superaddita}) the attributes of divine glory, majesty, power, omniscience, and omnipresence.\footnote{\emph{Epit. }VIII. (p. 545): '\emph{Sondern hie ist die höchste Gemeinschaft, welche Gott mit dem Menschen wahrhaftig hat, aus welcher persönlichen Vereinigung und der daraus erfolgenden höchsten und unaussprechlichen Gemeinschaft alles herfleusst, was menschlich von Gott, und göttlich vom Menschen Christo gesaget und gegläubet wird; wie solche Vereinigung und Gemeinschaft der Naturen die alten Kirchenlehrer durch die Gleichniss eines feurigen Eisens, wie auch der Vereinigung Leibes und der Seelen im Menschen erkläret haben.}' The \emph{Sol. Decl.} repeats the same at greater length.} Thus God is really man, and man is really God; Mary is truly the mother of God, since she conceived and brought forth the Son of God; the Son of God truly suffered, though according to the property of his human nature; Christ as man, not only as God, knows all things, is able to do all things, is present to all creatures, and was so from the moment of the incarnation. For (as the Solid Declaration expressly states) Christ, according to his humanity, received his divine Majesty 'when he was conceived in the womb and became man, and when the divine and human natures were united with each other.' That is to say, the incarnation of God was at the same time a deification of man in Christ. (This was the Swabian theory of Brentius and Andreæ.)

As regards the ubiquity in particular, the Formula is again inconsistent. The Epitome favors the doctrine of the \emph{ absolute} ubiquity of Christ's body in all creatures (as taught by Luther, Brentius, Andreæ), and says that Christ, 'not only as God, but also as man, is present \emph{to all creatures} . . . is \emph{omnipresent}, and all things are possible and known to him;' the Solid Declaration, on the contrary, asserts only the \emph{relative} ubiquity or multivolipresence (as taught by Chemnitz); but neutralizes this again by quoting, with full approbation, Luther's strongest passages in favor of absolute ubiquity.\footnote{The words ' \emph{dass Christus auch nach und mit seiner assumirten Menschheit gegenwärtig sein } könne \emph{und auch sei, } wo er will,' clearly express the \emph{multivolipræsentia} of Chemnitz and the Saxons. Nevertheless, Chemnitz, to his own regret, could not prevent the wholesale indorsement and quotation of Luther's views—that wherever Christ's divinity is, there is also his humanity; that he may be and is in all places wherever God is; that the ascension is figurative; that the right hand of God is every where, etc. Hence it is scarcely correct when Kahnis says (Vol. II. p. 581) that the compromise of the Formula leans to the side of Chemnitz. Compare the thorough discussion of Dorner, \emph{Entwicklungsgeschichte}, Vol. II. pp. 710 sqq., who clearly shows that Chemnitz made several fatal concessions to the Swabian Christology. Hence the opposition of Heshusius and the Helmstädt Lutherans (see p. 293).} Hence there arose a fruitless controversy on the subject among the orthodox Lutherans themselves, as has been already stated.

The Formula, therefore, is not a real union of the Swabian and Saxon types, but only a series of concessions and counter-concessions, and a mechanical juxtaposition of discordant sentences from both parties.\footnote{Dorner, Vol. II. p. 771, '\emph{Die Vermittlungsversuche des I. Andreæ und Chemnitz erreichten in Betreff des eigentlichen Gegensatzes zwischen den Schwaben und Niederdeutschen keine innere Einigung, sondern nur eine Vereinigung van disharmonischen Sätzen von beiden Seiten her in einem Buch. Die Folge war daher nicht Eintracht, sondern vielseitige Zwietracht.}'} The later orthodoxy did not settle the question, and both theories continued to find their advocates. Moreover, the Formula does not answer and refute, but simply denies the objections of the Reformed divines, and falls back upon the incomprehensibility of the mystery of the hypostatic union, which is declared to be the highest mystery next to the Trinity, and the one 'on which our whole consolation, life, and salvation depend.'

As regards the \emph{states of humiliation} (\emph{exinanitio}) and \emph{exaltation} (\emph{exaltatio}), the Formula, in the passages already quoted, teaches the full possession (\textgreek{κτῆσις}), and a partial or occult use (\textgreek{κρῆσις}), of the divine attributes by Christ from the moment of his existence as a man. His human nature, and not the divine pre-existent Logos, is understood to be the subject of the humiliation in the classical passage Phil. ii. 7, on which the distinction of two states is based. Consequently the two states refer properly only to the human nature, and consist in a difference of outward condition and visible manifestation. The humiliation is a partial concealment of the actual use (a \textgreek{κρύψις χρήσεως}) of the divine attributes communicated to the human nature at the incarnation; the exaltation is a full manifestation of the same. As to the extent of the concealment or actual use, there arose afterwards, as we have seen already, a controversy between the Giessen and Tübingen divines, but was never properly settled, nor can it be settled on the christological basis of the Formula.\footnote{The Formula teaches the \textgreek{κτῆσις} with a partial \textgreek{κένωσις χρήσεως,} and so far seems to favor the later Giessen view, although the issue was not yet fairly before the authors. \emph{Sol. Decl.} Art. VIII. (p. 767 ed. Rech., p. 680 ed. Müller): '\emph{Eam vero majestatem statim in sua conceptione etiam in utero matris habuit, sed ut apostolus loquitur} (Phil. ii. 7), \emph{se ipsum exinanivit, eamque, ut D. Lutherus docet, in statu suæ humiliationis } secreto  \emph{habuit, neque eam semper, sed } quoties ipsi visum fuit,  \emph{usurpavit.}' An \emph{occasional} use of the divine attributes during the state of humiliation was expressly conceded by the Giessen divines; they only denied the constant and full (though secret) use contended for by the Tübingen school. See above, p. 295. The Lutheran scholastics were more on the side of the Giessen divines.} The modern school of Lutheran Kenoticists depart from it by assuming a real self-renunciation (\textgreek{κένωσις}) of the divine Logos in the incarnation, but thereby they endanger the immutability of the Deity, and interrupt the continuity of the divine government of the world through the Logos during the state of humiliation.

We add some general remarks on the Christology of the Formula, as far as it differs from the Reformed Christology. After renewed investigation of this difficult problem, I have been confirmed in the conviction that the exegetical argument, which must ultimately decide the case, is in favor of the Reformed and against the Lutheran theory; but I cheerfully admit that the latter represents a certain mystical and speculative element, which is not properly appreciated in the Calvinistic theology, and may act as a check upon Nestorian tendencies.

1. The scholastic refinements of the doctrine of the \emph{communicatio idiomatum}, and especially the ubiquity of the body, have no intrinsic religious importance, and owe their origin to the Lutheran hypothesis of the corporeal presence.\footnote{This is admitted, in part at least, by Dr. Stahl, one of the ablest and most clear-headed modern champions of orthodox Lutheranism, when he says: '\emph{Die Lehre von der Allgegenwart des Leibes Christi ist, abgesehen von der Anwendung auf das Abendmahl, } von gar keinem religiösen Interesse' (\emph{Die lutherische Kirche und die Union}, Berlin, 1859, p. 185).} They should, therefore, never have been made an article of faith. A surplus of orthodoxy provokes skepticism.

2. The great and central mystery of the union of the divine and human in Christ, which the Formula desires to uphold, is overstated and endangered by its doctrine of the \emph{genus majestaticum}, or the communication of the \emph{divine} attributes to the \emph{human nature} of Christ. This doctrine runs contrary to the \textgreek{ἀσυγχύτως} and \textgreek{ἀτρέπτως} of the Chalcedonian Creed. It leads necessarily—notwithstanding the solemn protest of the Formula—to a Eutychian confusion and æquation of natures; for, according to all sound philosophy, the attributes are not an outside appendix to the nature and independent of it, but inherent qualities, and together constitute the nature itself. Or else it involves the impossible conception of a double set of divine attributes—one that is original, and one that is derived or transferred.

3. The \emph{genus majestaticum} can not be carried out, and breaks down half-way. The divine attributes form a unit, and can not be separated. If one is communicated, all are communicated. But how can eternity \emph{ab ante} (\emph{anfangslose Existenz}), which is a necessary attribute of the divine nature of Christ, be really communicated to a being born in time, as Jesus of Nazareth undoubtedly was? How can immensity be transferred to a finite man? The thing is impossible and contradictory. An appeal to God's omnipotence is idle, for God can not sin, nor err, nor die, nor do any thing that is inconsistent with his rational and holy nature.

4. The doctrine has no support in the Scriptures; for the passages quoted in its favor speak of the divine human \emph{person}, not of the human \emph{nature} of Christ; as, '\emph{I} am with you alway;' 'all power is given to \emph{me;}'\footnote{It is objected that omnipotence could not be \emph{given} to the \emph{divine person} of Christ, who had it from eternity essentially and of necessity, but only to his \emph{human nature.} But this reasoning implies a virtual denial of the \textgreek{κένωσις,} or laying aside of the pre-existent glory which Christ had as God, and was going to take possession of again as God-man at his exaltation, John xvii. 5 (\textgreek{δόξασον μὲ . . . τῇ δόξῃ ᾖ εἶχον πρὸ τοῦ τὸν κόσμον εἶναι παρὰ σοί). }} 'in \emph{Christ} are hid all the treasures of wisdom and knowledge;' 'in \emph{Christ }dwelleth all the fullness of the Godhead bodily.' And as to the state of humiliation, such passages as Luke ii. 52; Mark xiii. 32; Heb. v. 8, 9, are inconsistent with the teaching of the Formula that he was omniscient as \emph{man} from the mother's womb.

5. The Christology of the Formula makes it impossible to construct a truly human life of our Lord on earth, and turns it into a delusive Christophany, or substitutes a crypto-pantheistic Christ for a personal, historical Christ.

6. The familiar illustrations of the iron and fire, and body and soul, used by the Formula, favor the Reformed rather than the Lutheran theory; for the iron does not transfer its properties to the fire, nor the fire to the iron; neither are the spiritual qualities of the soul, as cognition and volition, communicated to the body, nor the material properties and functions of the body, as weight and extension, eating and drinking, to the soul: both are indeed most intimately and inseparably connected—the soul dwells in the body, and the body is the organ of the soul—but both remain essentially distinct. The same is the case with the other illustration which is borrowed from the intercommunication or inhabitation (\textgreek{περιχώρησις, } \emph{immanentia, permeatio, circumincessio}) of the persons of the Holy Trinity; for the peculiar properties (\textgreek{ἴδια, ἰδιότητες}) of the persons are not communicated or transferred—paternity and being unbegotten (\textgreek{ἀγεννησία}) belongs to the Father alone, sonship (\textgreek{γεννησία, } \emph{filiatio}) to the Son alone, and procession (\textgreek{ἐκπόρευσις, } \emph{processio}) to the Holy Ghost alone.

7. The ubiquity of the body is logically necessary for the hypothesis of consubstantiation, and both stand and fall together. For the eucharistic multipresence must be derived either from a perpetual miracle (performed through the priestly consecration, or by the power of the Holy Ghost, both of which the Lutherans reject),\footnote{According to the Romish liturgy, the elements are literally changed or transubstantiated into the very body and blood of Christ by the consecration of the priest when he repeats the words of institution, \emph{Hoc est corpus meum}; and hence the priest is blasphemously said to create the body of Christ. But, according to the Oriental and Greek liturgies, the presence of the body and blood of Christ is effected by the Benediction or Invocation of the Holy Ghost, which follows the recital of the words of institution. Calvin and the Reformed liturgies likewise bring in the agency of the Holy Ghost, but simply for conveying the energy or the power and effect of the body and blood of Christ in heaven to the believing communicant.} or from an inherent quality of the body itself, which enables it to be present wherever and whenever it is actually partaken of by the mouth of the communicants.

8. But ubiquity proves too much for consubstantiation by extending the eating of Christ to every meal (though this is inconsistently denied), and depriving the eucharistic presence of all specific value. Yea, it is fatal to it, and leads, we will not say to the Calvinistic, but rather to a crypto-pantheistic theory of the eucharist;\footnote{The Roman Catholic Bellarmin (see below) and Reformed polemics (also Steitz on \emph{Ubiquity}, in Herzog's \emph{Encykl.}) argue that the ubiquity dogma destroys the Lutheran corporeal presence, and logically ends in the Calvinistic theory of the spiritual real presence. But we would rather say that it ends in a crypto-panchristism, which is quite foreign to Calvin. The doctrine of ubiquity was, before Luther, always connected with a leaning to Gnosticism and Pantheism, as in Origen and Scotus Erigena.} for a body which is intrinsically and perpetually omnipresent must be so spiritual that it can only be spiritually present and spiritually be partaken of by faith.\footnote{The Lutherans exclude all ideas of local extension or expansion from the body of Christ, and describe it just as the scholastics and the ancient philosophers (Plato, Aristotle, Philo) describe the presence of incorporeal substances, and especially of the \emph{Deity} itself, which is 'unextended,' 'indistant,' 'devoid of magnitude,' not part of it here and part of it there, but whole and undivided every where and nowhere. See Cudworth's \emph{Intellectual System of the Universe}, Harrison's ed. (Lond. 1845), Vol. III. p. 248.}

9. Ubiquity is not only unscriptural, but antiscriptural, and conflicts with the facts of Christ's local limitations while on earth, his descent into Hades, his forty days after the resurrection, his ascension to heaven, his visible return to judgment. We freely admit that Christ's glorified body is not subject to the laws of \emph{earthly} substances or \emph{confined} to a particular locality; it is a 'spiritual' body (comp.  1 Cor. xv.), with its own laws of rest and locomotion, which transcend our present knowledge; nevertheless it is and ever remains a body, as real as the resurrection body of saints which will be fashioned like unto it (\textgreek{σύμμορφον τῷ σώματι τῆς δόξης αὐτοῦ}), and as heaven itself is real, from which Christ will return 'in like manner' as the apostles 'saw him go into heaven.' The ubiquitarian exegesis here runs into an ultra-Zwinglian spiritualism to save the literalism with which it started. But, feeling its own weakness, it falls back again at last upon the literal understanding of the \textgreek{ἐστί} in the words of institution.

10. This first and last resort of consubstantiation is given up by the ablest modern exegetes,\footnote{Including such unbiased philological commentators as De Wette and Meyer. See especially Meyer on  \emph{Matthew} xxvi. 26 (pp. 548 sqq. of the 5th ed.), and my annotations to Lange on \emph{Matthew}, Am. ed., pp. 470–474. Kahnis, who formerly wrote an elaborate historical work in defense of the Lutheran doctrine (\emph{Die Lehre vom Abendmahl}, Lipz. 1851), has more recently (1861) arrived at the conclusion that 'the Lutheran interpretation of \emph{the words of institution} must be given up,' though he thinks that this affects only the Lutheran theology, not the Lutheran faith.} who agree in the following decisive results: (\emph{a}) That the disputed word \textgreek{ἐστί} was not even spoken by our Lord in Aramaic, and can have no conclusive weight, (\emph{b}) That the substantive verb may designate a symbolical as well as a real relation between the subject and the predicate, as is evident from the nature of the case and from innumerable passages of Scripture, (\emph{c}) That in this case the literal interpretation would lead to transubstantiation rather than the semi-figurative (synecdochical) consubstantiation; since Christ does not say what the Lutheran hypothesis would require: 'This is my body \emph{and bread},' 'This is my blood \emph{and wine} (or in, with, and under the bread and wine).' (\emph{d}) That the figurative or metaphorical interpretation (whether in the Zwinglian or Calvinistic sense) is made necessary in connection with the \textgreek{τοῦτο} for \textgreek{οὗτος, ποτήριον} for \textgreek{οἶνος,} or \textgreek{αἶμα,} as well as by the surroundings of the institution of the Lord's Supper, viz.: the nature of the typical passover, the living, personal presence of our Lord, with his body still unbroken and his blood still unshed, which could not be literally eaten and drunk by his disciples.

This, of course, only settles the exegetical basis, and still leaves room for different doctrinal views of this sacred ordinance, into which we can not here enter.\footnote{I have briefly expressed my own view in \emph{Com. on Matthew}, p. 471: . . . 'But we firmly believe that the Lutheran and Reformed views can be essentially reconciled, if subordinate differences and scholastic subtleties are yielded. The chief elements of reconciliation are at hand in the Melanchthonian-Calvinistic theory. The Lord's Supper is: (1.) A commemorative ordinance, a memorial of Christ's atoning death, and a renewed application of the virtue of his broken body and shed blood. (This is the truth of the Zwinglian view, which no one can deny in the face of the words of the Saviour: '\emph{Do this in remembrance of me.}') (2.) A feast of living union of believers with the ever-living, exalted Saviour, whereby we truly, though spiritually, receive Christ with all his benefits, and are nourished by his life unto life eternal. (This was the substance for which Luther contended against Zwingli, and which Calvin retained, though in a different scientific form, and in a sense rightly confined to believers.) (3.) A communion of believers with one another as members of the same mystical body of Christ. . . . It is a sad reflection that the ordinance of the Lord's Supper—this feast of the \emph{unio mystica} and \emph{communio sanctorum}, which should bind all pious hearts to Christ and each other, and fill them with the holiest and tenderest affections—has been the innocent occasion of the bitterest and most violent passions and the most uncharitable abuse. The eucharistic controversies are among the most unrefreshing and apparently fruitless in church history. Theologians will have much to answer for at the judgment-day for having perverted the sacred feast of divine love into an apple of discord. No wonder that Melanchthon's last wish and prayer was to be delivered from the \emph{rabies theologorum}. Fortunately, the blessing of the holy communion does not depend upon the scientific interpretation and understanding of the words of institution, but upon the promise of the Lord, and upon childlike faith which receives it, though it may not fully understand the mystery of the ordinance. Christians celebrated it with most devotion and profit before they contended about the true meaning of those words, and obscured their vision by all sorts of scholastic theories and speculations. Fortunately, even now Christians of different denominations and holding different opinions can unite around the table of their common Lord and Saviour, and feel one with him and in him who died for them all, and feeds them with his life once sacrificed on the cross, but now living forever. Let them hold fast to what they agree in, and charitably judge of their differences; looking hopefully forward to the marriage supper of the Lamb in the kingdom of glory, when we shall understand and adore, in perfect harmony, the infinite mystery of the love of God in his Son our Saviour.'}

11. The Lutheran doctrine of the eucharist overlooks the omnipresence of the Holy Spirit, and substitutes for it the corporeal presence of Christ. It is the Holy Spirit who brings the believer in and out of the sacrament into a living union and communion with the whole Christ, and makes the perpetual virtue and efficacy of his crucified body on the cross, \emph{i.e.}, his atoning sacrifice, and of his glorified body in heaven, available for our spiritual benefit.

12. Finally, as regards the two states of Christ, the Reformed Christology is right in making the pre-existent Logos \textgreek{(Λόγος ἄσαρκος)} the subject of the \textgreek{κένωσις,} or self-humiliation, instead of the human nature (or the \textgreek{Λόγος ἔνσαρκος),} which was never before \textgreek{ἐν μορφῇ θεοῦ,} and consequently could not renounce it in any way. The incarnation itself is the beginning of the humiliation. In this interpretation of Phil. ii. 7 the Reformed Church is sustained not only by Chrysostom and other fathers, but also by the best modern exegetes of all denominations, including Lutherans.\footnote{See, especially, Meyer (who ably defends the patristic and Reformed exegesis against the objections of De Wette and Philippi), and Braune on Phil. ii. 6 sqq. (Am. ed. of Lange). The latter says: '\textgreek{ὅς} of has for its antecedent \textgreek{Χριστῷ Ιησοῦ,} and points to his ante-mundane state, as verses 7 and 8 refer to his earthly existence, and verses 9–11 refer to his subsequent glorified condition. The subject is the Ego of the Lord, which is active in all the three modes of existence. It is the entire summary of the history of Jesus, including his ante-human state.' Among the dogmatic theologians of the Lutheran Church, Liebner, Thomasius, Kahnis, Gess, and others, give up the old Lutheran exegesis of the passage. Kahnis (in the third volume of his \emph{Luth. Dogmatik}, 1868, p. 341) makes, as the result of his earnest investigation, the following clear and honest statement: '(a) \emph{Dass Paulus in der Offenbarungsgeschichte Jesu Christi drei Stadien unterscheidet: das Stadium der Gottesgestalt, da der Logos beim Vater war; das Stadium der Knechtsgestalt, das mit der Selbstverleugnung Christi in der Menschwerdung begann und zur Erniedrigung am Kreuze fortging; das Stadium der Erhöhung, da im Namen Christi sich alle Knie beugen und ihn als Herrn bekennen.} (b) \emph{Dass das Subjekt der Erniedrigung der} \textgreek{ λόγος ἄσαρκος} \emph{ist, wie schon die alte Kirche in ihren namhaftesten Lehrern sah, die reformirten Theologen richtig erkannten und auch die bedeutendsten neueren Ausleqer aller Confessionen zugestehen, das Subjekt der Erhöhung aber der } \textgreek{λόγος ἔνσαρκος.} (c) \emph{Dass die Entäusserung} (\textgreek{ἐαυτόν ἐκέύωσε}) \emph{darin besteht, dass der Logos sich der Gottesgestalt } (\textgreek{μορφὴ θεοῦ}) \emph{ d. h. des Herrlichkeitsstandes beim Vater begab, um Knechtsgestalt } (\textgreek{μορφὴ δούλου}) \emph{ anzunehmen, d.h. ein Mensch wie wir zu werden, ja als Mensch sich zum Kreuzestode zu erniedrigen} (\textgreek{ἐταπείνωσεν ἐαυτόν}): \emph{Entäusserung also gleich Menschwerdung ist. Darnach fordert dieses Lehrstück eine andere Fassung, als die alte }[\emph{Luther.}] \emph{Dogmatik ihm gab.}'}

Art. IX. Of Christ's Descent into Hell..—The fact of a real descent of the whole person of Christ, the God-man, after his death, into the real hell (not a metaphorical hell, nor the grave, nor the \emph{limbus patrum}) is affirmed, and its object defined to be the defeat of Satan and the deliverance of believers from the power of death and the devil; but all curious questions about the mode are deprecated and left for the world to come.

Art. X. Of Church Usages and Ceremonies, called Adiaphora.—The observance of ceremonies and usages neither commanded nor forbidden in the Word of God, should be left to Christian freedom, but should be firmly resisted when they are forced upon us as a part of divine service (Gal. ii. 4, 5; v. 1; Acts xvi. 3; Rom. xiv. 6; 1 Cor. vii. 18;  Col. ii. 16).

This article was a virtual condemnation of Melanchthon's course in the Interim controversy.

Art. XI. Of God's Foreknowledge and Election.—No serious controversy took place on this doctrine in the Lutheran Church, except at Strasburg between Zanchi and Marbach (1561). The rigid predestinarianism of Luther and the Flacianists quietly gave way to the doctrine of the universality of divine grace, while yet the anthropological premises of the Augustinian system were retained (in Art. I. and II.).

The Formula teaches that there is a distinction between foreknowledge (\emph{præscientia, prævisio, } \emph{Vorsehung, } Matt. x. 29; Psa. cxxxix. 16; Isa. xxxvii. 28) and foreordination (\emph{prædestinatio, electio, } \emph{ewige Wahl, } Eph. i. 5); that foreknowledge pertains alike to the good and the evil, and is not the cause of sin and destruction; that foreordination refers only to the children of God; that this predestination of the elect is 'eternal, infallible, and unchangeable,' and is the ultimate and unconditional cause of their salvation; that God, though he elects only a portion, sincerely desires \emph{all} men to be saved, and invites them by his Word to the salvation in Christ; that the impenitent perish by their own guilt in rejecting the gospel; that Christians should seek the eternal election, not in the secret but in the revealed will of God, and avoid presumptuous and curious questions.

Thus the particularism of election and the universalism of vocation, the absolute inability of fallen man (Art. II.), and the guilt of the unbeliever for rejecting what he can not accept, are illogically combined. The obvious contradiction between this article and the second has already been pointed out.\footnote{See above, p. 314. Comp. also Dorner, \emph{Gesch. der Prot. Theol.} pp.366 sqq. Planck (Vol. VI. p. 814) charges this article with a confusion not found in the other parts of the Formula, and Gieseler (Vol. IV. p. 488) with putting together contradictory positions; while, on the other hand, Thomasius (\emph{Das Bekenntniss der ev. luth. Kirche}, etc. p. 222) sees here only supplementary truths to be reconciled by theological science, and Guericke (in his \emph{Kirchengeschichte}, Vol. III. p. 419) calls the logical inconsistency of the Formula 'divinely necessitated' (\emph{eine göttlich nothwendige Verstandes-Inconsequenz}).}

The authors felt the speculative difficulties of this dogma, and emphasized the practical side, which amounts to this: that believers are saved by the free grace of God, while unbelievers are lost by their own guilt in rejecting the grace sincerely offered to them. Later Lutheran divines, like John Gerhard, labored hard to show that God not only sincerely desires the salvation of all men alike, but that he also actually gives an opportunity to \emph{all} men even in this \emph{present} life.\footnote{\emph{Loc. Theol.} Tom. IV. pp. 189 sqq. (\emph{de Electione et Reprob.} § 7; \emph{de Universalitate Vocationis}, § 135\}.} But the argument fails with regard to the heathen, who form the greatest part of the race even to this day (not to speak of the world before Christ): and hence the Lutheran view of the \emph{actual} universality of the offer of grace necessitates an essential change of the orthodox doctrine of the middle state, as far as those are concerned who never heard of the gospel in this world.

Art. XII. Of Several Heresies and Sects.—This article rejects the peculiar tenets of the Anabaptists, Schwenkfeldians, New Arians, and Antitrinitarians, who never embraced the Augsburg Confession.

To the second part of the Formula there is added a \emph{Catalogue of Testimonies} from the Scriptures and the fathers (Athanasius, Gregory Nazianzen, Cyril of Alexandria, John of Damascus) concerning the divine majesty of the human nature of Christ, in support of the doctrine of the \emph{communicatio idiomatum}, as taught in Art. VIII. This Appendix was prepared by Andreæ and Chemnitz; but it has no symbolical authority, and is often omitted from the Book of Concord.\footnote{Tittmann and Hase omit it; Müller gives it (pp. 731–767).}

RECEPTION, AUTHORITY, AND INTRODUCTION.\footnote{Comp. among recent works especially the third volume of Heppe's \emph{Geschichte des D. Protest}, pp. 215–322, and the whole fourth volume. The chief data are also given by Gieseler, Vol. IV. pp. 489–493, and by Köllner, 1.c. pp. 573–583.}

The Form of Concord, as it is the last, is also the most disputed of the Lutheran symbols. It never attained general authority, like the Augsburg Confession or Luther's Catechism, although far greater exertions were made for its introduction.

It was adopted by the majority of the Lutheran principalities and state churches in Germany;\footnote{The Preface of the Book of Concord is signed by eighty-six names representing the Lutheran state churches in the German empire; among them are three Electors (Louis of the Palatinate, Augustus of Saxony, and John George of Brandenburg), twenty Dukes and Princes, twenty-four Counts, thirty-five burgomasters and counselors of imperial cities. The Formula was also signed by about 8000 pastors and teachers under their jurisdiction, including a large number of ex-Philippists and Crypto-Calvinists, who preferred their livings to their theology; hence Hutter was no doubt right when he admitted that many subscribed \emph{mala conscientia.} Yet no \emph{direct} compulsion seems to have been used. See Köllner, p. 551, and Johannsen, \emph{Ueber die Unterschriften des Concordienbuches}, in Niedner's \emph{Zeitschrift für histor. Theologie}, 1847, No. 1.} also by the state church of Sweden, the Lutherans in Hungary, and several Lutheran synods in the United States.\footnote{It was adopted in Sweden at a Council of Upsala, 1593; in Hungary, 1597. In America it is held by the Lutheran Synodical Conference, and by the General Council, but rejected by the General Synod (see p. 224).}

On the other hand, it was rejected by a number of Lutheran Princes and cities of the empire,\footnote{The Landgrave of Hesse, the Palatinate John Casimir, the Prince of Anhalt, the Duke of Pomerania (where, however, the symbol afterwards came into authority), the Duke of Holstein, the Duke of Saxe-Luneburg, the Counts of Nassau and Hanau, the cities of Strasburg, Frankfort-on-the-Main, Spires, Worms, Nuremberg, Magdeburg, Bremen, Danzig, Nordhausen.} and by King Frederick II. of Denmark.\footnote{Frederick II. strictly prohibited, on pain of confiscation and deposition, the importation and publication of the Form of Concord in Denmark (July 24, 1580), and threw the two superbly bound copies sent to him by his sister, the wife of Augustus of Saxony, unceremoniously into the chimney-fire. See Köllner, p. 575 sq.; Gieseler, Vol. IV. p. 493, note 54; and Heppe, Vol. IV. pp. 275 sqq. Nevertheless the document afterwards gained considerable currency in Denmark.}

Some countries of Germany, where it had been first introduced, rejected it afterwards, but remained Lutheran;\footnote{So the Duchy of Brunswick recalled the subscription in 1583. Duke Julius, one of the most zealous promoters of the Form of Concord, became alienated for personal reasons, because he was severely blamed by Chemnitz and several Princes for allowing one of his sons to receive Romish consecration (Dec. 5, 1578), and two others the tonsure, to the great scandal of Protestantism. He was afterwards strengthened by the doctrinal opposition of Heshusius and the Helmstädt Professors, who rejected the Formula for teaching absolute ubiquity. The \emph{Corpus doctrinæ Julium} was retained in Brunswick and Wolfenbüttel. See Planck, Vol. VI. pp. 667 sqq., and especially Heppe, Vol. IV. pp. 203 sqq. These Brunswick troubles brought about an alienation between Andreæ (who labored to reconcile the Duke) and Chemnitz (who was deposed by the Duke). In a widely circulated letter of April 8, 1580, Chemnitz compared Andreæ to a fawning and scratching cat ('\emph{cum coram longe aliud mihi dicas, } \emph{wie die Katzen, die vorne lecken und hinten kratzen}'). Heppe, p. 214.} while others, in consequence of the doctrinal innovations and exclusiveness of the Formula, passed over to the Reformed Confession.\footnote{So the Palatinate, which, after a short Lutheran interregnum of Louis, readopted the Heidelberg Catechism under John Casimir (1583), Anhalt (1588), Zweibrücken (1588), Hanau (1596), Hesse (1604), and especially Brandenburg under John Sigismund (1614). In this respect the Formula of Concord inflicted great territorial loss upon the Lutheran denomination. The greatest loss was the Palatinate and the Electoral, afterwards the royal house of Brandenburg and Prussia.} It is a significant fact, that the successors of the three Electors, who were the chief patrons and signers of the Formula, left the Lutheran Church: two became Reformed, and one (the King of Saxony) a Roman Catholic.

OPPOSITION AND DEFENSE..\footnote{See lists of controversial works for and against the Formula of Concord in Walch, Feuerlin, and Köllner. Comp. also Hutter, \emph{Conc. conc.} Ch. XXXVII. (p. 958), Ch. XLI. (p. 976), Ch. XLV. (p. 1033), and Ch. XLV. (p. 1038); Heppe, Vol. IV. pp. 270 sqq.; and G. Frank, Vol. I. pp. 251–266. Hutter sees in the general attack of 'the devil and his organs, the heretics,' against the Formula, a clear proof that it was composed \emph{instinctu Spiritus Sancti}, and is in full harmony with the infallible Word of God (p. 976).}

The Formula gave rise to much controversy. It was assailed from different quarters by discontented Lutherans and Philippists,\footnote{The rigidly orthodox Heshusius and the Helmstädt divines (in the Quedlinburg Colloquium, 1583), Christopher Irenæus (an exiled Flacianist, formerly court chaplain at Weimar, 1581), Ambrosius Wolff (or Cyriacus Herdesianus, of Nuremberg, 1580), the Bremen preachers (1581), the Anhalt theologians (1580, 1581), and the Margrave of Baden (in the \emph{Stafford Book}, 1599).} Calvinists,\footnote{Ursinus (in connection with Zanchius, Tossanus, and other deposed Heidelberg Professors, who, under John Casimir and during the rule of Lutheranism in Heidelberg, founded and conducted a flourishing theological school at Neustadt an der Hardt, 1576 to 1583): \emph{Admonitio Christiana de libro Concordiæ} (or \emph{Christliche Erinnerung vom Concordienbuch}), Neostadadii in Palatinatu, Latin and German, 1581 (also in Urs. \emph{Opera}, Heidelberg, 1612, Vol. II. pp.486 sqq.). It consists of twelve chapters, and is very able. Extract in Sudhoff, \emph{Olevianus und Ursinus}, pp. 432–452; comp. Schweizer in Herzog, Vol. X. pp. 263–265. Ursinus and some of his pupils defended this work against the Lutheran 'Apology,' in \emph{Defensio Admonitionis Neost. contra Apologiæ Erfordensis sophismata}, Neost. 1584. Beza wrote \emph{Refutatio dogmatis de ficticia carnis Christi omnipræsentia}; Dansæus an \emph{Examen} of Chemnitz's book \emph{De duabus in Christo naturis}, Genev. 1581; Sadeel, a very able tract, \emph{De veritate humanæ naturæ Christi}, 1585 (in his \emph{Opera}, Genev. 1592). Of later Reformed writings must be mentioned the \emph{Emdensche Buch} (1591), and especially Hospinian's \emph{Concordia discors} (1607), which called forth Hutter's \emph{Concordia concors} (1614).} and Romanists.\footnote{The ablest Roman assailant was Robert Bellarmin: \emph{Judicium de libro quem Lutherani vocant Concordiæ}, Ingolst. 1587, 1589, etc. (in his \emph{Opera}, Col. Ag. 1620, Vol. VII. p. 576). Against him Hoe ab Hœnegg wrote \emph{Apol. contra R. B. impium et stolidum judicium}, Fref. 1605. Bellarmin also repeatedly notices the Christology of the Formula in his great controversial work against Protestantism. See below.}

The chief objection was to the new dogma of ubiquity.

The Lutherans attacked, according to their stand-point, either the concessions to the Swabian scheme of absolute ubiquity, or the absence of a direct condemnation of Melanchthon and other heretics, or the rejection of the Flacian theory of original sin, or the condemnation of Synergism. The last point could be made very plausible, since the chief authors of the Formula, Andreæ, Chemnitz, and Selnecker, had at first been decided synergists. Chytræus remained true at least to his love and admiration for Melanchthon, which subjected him to the suspicion of Crypto-Philippism and Calvinism.\footnote{See Schütz, \emph{Vita Chytræi}, and Heppe, Vol. IV. pp. 395 sqq.}

The Reformed, led by Ursinus (chief author of the Heidelberg Catechism), justly complained of the misrepresentations and unfair condemnation of their doctrine under the indiscriminate charge of sacramentarianism,\footnote{This complaint the Erfurt Apology of the Formula of Concord admitted to be just, at least in part. The Formula makes no distinction between Zwingli and Calvin; condemns Zwingli's '\emph{allæosis}' (by which he meant only to guard against a \emph{confusio} and \emph{æquatio naturarum}) as a mask of the devil; charges the Reformed generally with a Nestorian separation of the two natures in Christ, and a denial of all communion between them; with childish literalism concerning the right hand of God and the throne of glory; with shutting Christ up in heaven, as if he had no more to do with us, etc.} and explained the qualified sense in which the Reformed signed the Augsburg Confession in the sense of its author, with wholesome strictures on the unprotestant overestimate of the authority of Luther. They exposed with rigid logic the doctrinal contradiction between Arts. II. and XI., quoted Luther's views on predestination against the Formula, and refuted with clear and strong arguments the new dogma of ubiquity, which is contrary to the Scriptures, the œcumenical creeds, and sound reason, and destructive of the very nature of the sacrament as a communion of the \emph{body} of Christ; for if the body is omnipresent, and there can be but \emph{one} omnipresence, it must be present like God himself, i.e. like a spirit, every where whole and complete, without parts and members, and thus the lineaments and concrete image of Christ are lost. Sadeel pointed out the palpable inconsistency between the hyperphysical and ultrasupernatural outfit of Christ's body for the eucharistic presence, on the one hand, and the emphasizing of a \emph{corporeal} presence and \emph{oral} manducation on the other, as if this were the main thing in the sacrament, while the communion of the believing \emph{soul} with the person of Christ was almost lost sight of.\footnote{Dorner, in his \emph{History of Christology} (Vol. II. pp. 718–750), gives an admirable and impartial summary of the Reformed argument. Dr. Kahnis, of Leipzig, from his Lutheran standpoint, thus fairly and liberally characterizes the Reformed opposition to the Form of Concord (\emph{Luth. Dogm.} Vol. II. p. 590): '\emph{Die Reformirten vertraten den Standpunkt des Verstandes, welcher zwischen Endlichem und Unendlichem abstract}(?) \emph{scheidend} (\emph{finitum non est capax infiniti}) \emph{der menschlichen Natur Christi keinen Antheil an den göttlichen Eigenschaften einräumt; den Standpunkt der Realität, welcher in der Betrachtung der Person Christi, von dem Wandel auf Erden ausgehead, der rein menschlichen Entwicklung Christi freien Raum schaffen will; den Standpunkt des Praktischen, der bei den sicheren Thatsachen der persönlichen Vereinigung Beruhigung fasste, ohne sich in gnostisch-scholastische Theorien verspinnen zu wollen.}'}

Strange to say, the Roman Catholics were just as decidedly opposed to ubiquity, though otherwise much nearer the Lutheran doctrine of the sacraments. Bellarmin, the greatest controversialist of Rome, exposes the absurdity of a dogma which would destroy the human nature of Christ, and involve the presence of his body \emph{in uteris omnium feminarum, imo etiam virorum,} and the presence \emph{extra uterum} from the moment of conception, and \emph{in utero} after the nativity. In his polemic work against Protestantism he urges five arguments against ubiquity,\footnote{Lib. III. \emph{de Sacramento Eucharistiæ}, cap. 17. Comp. also cap. 7, and Lib. III. \emph{de Christo} (where he refers to the views of Luther, Brentius, Wigand, Heshusius, and Chemnitz on ubiquity).} viz.: (1.) It abolishes the sacramental character of the eucharist. (2.) It leads to the Calvinistic spiritual presence and spiritual eating by faith—the very error of the sacramentarians which this Lutheran dogma was to overthrow.\footnote{\par His reasoning is curious: '\emph{Quod est ubique, non potest moveri, nec transire de loco ad locum; ergo licet corpus Christi sit in pane, tamen non manducatur, cum panis manducatur, quia non movetur, nec transit cum pane e manu ad os, et ab ore ad stomachmn; nam etiam antea erat in ore et in stomacho, priusquam panis eo veniret. . . . Sequitur aut esse inanem cænam Domini, aut saltem spiritualiter sumi per energiam et per fidem, et solum a piis, qui habent fidem, et hoc est, quod volunt Calvinistæ.}'} (3.) It destroys the specific effect of the eucharist, and makes it useless. (4) It is refuted by the other Lutheran doctrine which confines the presence to the time of the \emph{use} of the sacrament.\footnote{'\emph{Si enim corpus Christi ubique est, erit etiam ante usum in vane.}} (5.) It is a makeshift to evade the power of priestly consecration which creates the eucharistic presence.\footnote{Bellarmin (\emph{De Sacr. Euch.} Lib. III. c. 7), after quoting Augustine against the \emph{sententia ubiquistarum Lutheranorum}, thus defines the Roman view: '\emph{Nos fatemur Christi corpus non esse ubique diffusum; et ubicunque est, habere suam formam et partium situm, ac dispositionem; quamvis hæc figura, forma, dispositio partium in cælo conspiciatur, ubi locum replet; in Sacramento autem sit quidem, sed non repleat locum, nec videri a nobis possit.'}}

Outside of Germany and Switzerland the Formula of Concord excited little or only passing polemical interest. Queen Elizabeth endeavored to prevent its adoption because it condemned the Reformed doctrine, and threatened to split and weaken the Protestants in their opposition to the united power of Rome. She sent delegates to a convention of Reformed Princes and delegates held at Frankfort-on-the-Main, Sept 1577.\footnote{Comp. on Elizabeth's action and the Convent of Frankfort, Hutter's \emph{Concordia concors}, Cap. XVI. and XVII. (pp. 513–523); Planck, Vol. VI. pp. 591–611; Heppe, Vol. IV. pp. 5 sqq., 16 sqq., and 72 sqq.} The Anglican divines of the sixteenth century rejected ubiquity as decidedly as the Continental Calvinists.\footnote{Cranmer was at first inclined to the Lutheran, theory, but gave it up afterwards. His fellow-Reformers held the Zwinglian or Calvinistic view. Bishop Hooper thus speaks of ubiquity: 'Such as say that heaven and the right hand of God is in the articles of our faith taken for God's power and might, which is every where, they do wrong to the Scripture and unto the articles of our faith. They make a confusion of the Scripture, and leave nothing certain. They darken the simple and plain verity thereof with intolerable sophisms. They make heaven hell, and hell heaven, turn upside down and pervert the order of God. If the heaven and God's right hand, whither our Saviour's body is ascended, be every where, and noteth no certain place, as these uncertain men teach, I will believe no ascension. What needeth it?—seeing Christ's body is every where with his Godhead. I will interpret this article of my creed thus: \emph{Christus ascendit ad dextram Patris. Patris dextra est ubique: ergo Christus ascendit ad ubique.} See what erroneous doctrine followeth their imaginations!' \emph{Early Writings of John Hooper, D.D., Lord Bishop of Gloucester and Worcester, Martyr}, 1555; ed. by the Parker Society, Cambridge, 1843, p.66. The '\emph{Declaration of Christ and his Office},' from which this passage is taken, was first published at Zürich. 1547, in the early stage of the ubiquitarian controversy. See also the \emph{ Remains of Archbishop Grindal}, Camb. 1843, p. 46.} Evangelical Episcopalians hold the Reformed view of the sacraments; and as to modern Anglo-Catholic and Ritualistic Episcopalians, they greatly prefer the Romish or Greek dogma of transubstantiation to the Lutheran consubstantiation.\footnote{Comp. the eucharistic works of Pusey (1855), Philip Freeman (1862), Thomas L. Vogan (1871), and John Harrison (against Pusey, 1871).}

The attacks upon the Formula, especially those proceeding from Lutherans and the Palatinate divines, could not be ignored in silence. Chemnitz, Selnecker, and Kirchner, by order of the three electoral patrons of the work, convened at Erfurt,\footnote{In the \emph{Gasthof zum grünen Weinfasse.} This gave rise to some joke and mockery.} Oct. 23,1581 (afterwards at Braunschweig and Quedlinburg), and prepared, with much labor and trouble, an elaborate 'Apology,' called the '\emph{Erfurt Book},' in four parts.\footnote{The first part was directed against the \emph{Neustadt Admonition} of Ursinus and his colleagues, the second against the Bremen pastors, the third against Irenæus, the fourth against Wolf. Timothy Kirchner, of the Palatinate, prepared the first three parts, Selnecker and Chemnitz the last. They were published singly, and then jointly at Dresden, 1584, and distributed by the Elector Augustus among all the churches of Saxony. See Hutter, pp. 978 sqq. and 1038 sqq. (\emph{De Apol. Libri Concord. et de Colloquio Quedlinburgensi}); Heppe, Vol. IV. pp. 284–311.} It called forth new attacks, which it is unnecessary here to follow.

LATER FORTUNES.

During the palmy period of Lutheran scholasticism the Formula of Concord stood in high authority among Lutherans, and was even regarded as inspired.\footnote{Hutter (\emph{Conc. conc.} p. 976), Deutschmann, and others, who called it \textgreek{θεόπνευστος.}} Its first centennial (1680) was celebrated with considerable enthusiasm.\footnote{Anton, 1.c. Ch. X. \emph{Erste Concordien-Jubelfreude}, pp. 134 sqq. J. G. Walch, in his \emph{Introd.} 1732, represents the last stage of orthodox veneration before the revolution of sentiment took place.} But at the close of another century it was dead and buried. The Pietists, and afterwards the Rationalists, rebelled against symbololatry and lifeless orthodoxy. One stone after another was taken down from the old temple, until it was left a venerable ruin. Those very countries where subscription to creeds had been most rigorously enforced, suffered most from the neological revolution.

Then followed a period of patient research and independent criticism, which led to a more impartial estimate. Planck, the ablest Lutheran historian of the Formula, with complete mastery of the sources, followed the leading actors into all the ramifications and recesses of their psychological motives, political intrigues, and theological passions, and represents the work as the fabrication of a theological triumvirate, which upon the whole did more harm than good, and which produced endless confusion and controversy.\footnote{See his judgment, Vol. VI. pp. 690 sqq.; 816 sqq. and \emph{passim.} Planck's history is, even more than Hospinian's \emph{Concordia discors}, a \emph{chronique scandaleuse} of Lutheran pugnacity and bigotry in the second half of the sixteenth century.} Köllner, another learned and impartial Lutheran, concedes to it higher merit for the past, but no dogmatic significance for the present, except in the article on predestination.\footnote{\emph{ Symb. Vol.} I. p. 596: '\emph{Die Concordienformel hat dogmatisch nur insofern noch Werth, als sie mit den früheren Symbolen übereinstimmt. . . . Allein die Lehre von der Prädestination ausgenommen, kann ihr für das Dogma wie für die äusseren Verhältnisse der Kirche nur der wenigste eigenthümliche Werth unter allen Symbolen der Kirche zugestanden werden. Eigenthümlich ist nur die Ausbildung und mehr systematische Gestaltung des Lehrbegriffs der Kirche als eines Systems.}' This is too low an estimate of the whole document, and too high an estimate of Art. XI.} Heppe, the indefatigable historian of the German post-Reformation period, from a vast amount of authentic information, carries out the one-sided idea that the Lutheranism of the Formula is an apostasy from the normal development of German Protestantism, by which he means progressive, semi-Reformed, unionistic Melanchthonianism.\footnote{In his numerous works, so often quoted.} Even Kahnis thinks that the Lutheran theology of the future must be built on the Melanchthonian elements which were condemned by the Formula.\footnote{\emph{Dogm.} Vol. II. p. 517: '\emph{Man darf, . . mit Zuversicht aussprechen, dass die Zukunft der theologischen Forschung an dem Fortschreiten auf dem von Melanchthon eingeschlagenen Wege hängt.}'}

With the modern revival of orthodoxy, the Formula enjoyed a partial resurrection among Lutherans of the high sacramentarian type, who regard it as the model of pure doctrine and the best summary of the Bible. By this class of divines it is all the more highly esteemed, since they make doctrine the corner-stone of the Church and the indispensable condition of Christian fellowship. In America, too, the Formula has recently found at least one able and scholarly advocate in the person of Dr. Krauth, of Philadelphia.\footnote{Dr. Krauth calls the Formula 'the amplest and clearest confession in which the Christian Church has ever embodied her faith,' and he goes so far as to say: 'But for the Formula of Concord, it may be questioned whether Protestantism could have been saved to the world' (\emph{Conservative Reform.} p. 302). And this in full view of the independent Protestantism in Switzerland, France, Holland. England, and Scotland, which materially differs from the distinctive theology of this book, and was in vain condemned by it!}

Yet the great body of the Lutheran Church will never return to the former veneration for this symbol. History never repeats itself. Each age must produce its own theology. Even modern Lutheran orthodoxy in its ablest champions is by no means in full harmony with the Formula, but departs from its anthropology and Christology, and makes concessions to Melanchthon and the Reformed theology, or attempts a new solution of the mighty problems which were once regarded as finally settled.\footnote{We can simply allude to the internal differences of the Erlangen, Leipzig, and Rostock schools of Lutherans; to Luthardt on the freedom of the will; to Thomasius on the Kenosis; to Kahnis on the Lord's Supper, inspiration, and the canon of the Scripture; to the Hofmann and Philippi controversy on the atonement; to Hengstenberg's articles on justification and the Epistle of James; to the disputes on the millenarian question; and to the controversy on Church government and the relation of the ministry to the general priesthood of believers, in which Huschke, Stahl, Kliefoth, Vilmar, and Löhe take High-Church ground against the Low-Church views of Höfling, Harless, Diedrich, etc. Some of these controversies, especially the question of the ministerial office (\emph{Amtsfrage}), are also disturbing the peace of the orthodox Lutherans in America, and divide them into hostile synods (the Missouri Synod \emph{versus} the Grabau Synod, Iowa Synod, and portions of the General Council, not to mention several subdivisions). The eschatological controversy separates the Iowa Synod from Grabau and the Missourians, who denounce millenarianism as a heresy. The smallest doctrinal difference among orthodox Lutherans in America is considered sufficient to justify the formation of a new synod with close-communion principles. And yet all these Lutherans adopt the Formula Concordiæ as the highest standard of pure Scripture orthodoxy. Is this \emph{Concordia concors}, or \emph{Concordia discors?}}

AN IMPARTIAL ESTIMATE.

The Formula of Concord is, next to the Augsburg Confession, the most important theological standard of the Lutheran Church, but differs from it as the \emph{sectarian} symbol of Lutheranism, while the other is its \emph{catholic} symbol. Hence its authority is confined to that communion, and is recognized only by a section of it. It is both conclusive and exclusive, a Formula of Concord and a Formula of Discord, the end of controversy and the beginning of controversy. It completed the separation of the Lutheran and Reformed Churches, it contracted the territory and the theology of Lutheranism, and sowed in it the seed of discord by endeavoring to settle too much, and yet leaving unsettled some of the most characteristic dogmas. It is invaluable as a theological document, but a partial failure as a symbol, just because it contains too much theology and too little charity. It closes the productive period of the Lutheran reformation and opens the era of scholastic formalism.

The Formula is the fullest embodiment of genuine Lutheran orthodoxy, as distinct from other denominations. It represents one of the leading doctrinal types of Christendom. It is for the Lutheran system what the Decrees of Trent are for the Roman Catholic, the Canons of Dort for the Calvinistic. It sums up the results of the theological controversies of a whole generation with great learning, ability, discrimination, acumen, and, we may add, with comparative moderation. It is quite probable that Luther himself would have heartily indorsed it, with the exception, perhaps, of a part of the eleventh article. The Formula itself claims to be merely a \emph{repetition} and \emph{explication} of the genuine sense of the Augsburg Confession, and disclaims originality in the substance of doctrine.\footnote{See the Preface. An able argument for this agreement is presented by Prof. Thomasius, of Erlangen, in his \emph{Das Bekenntniss der evangelisch-lutherischen Kirche in der Consequent seines Princips}, Nürnberg, 1848. He develops the doctrines of the Formula from Luther's doctrine of justification by faith as the organic life-principle of the Lutheran Church. But the Lutheran doctrine of the eucharist with the \emph{communicatio idiomatum} and ubiquity of the body have—as the creeds of the Reformed churches prove—no necessary connection with justification by faith; and on these points, which constitute the peculiar features of the Formula, the author of the Augsburg Confession himself represented, even before Luther's death, a different line of development.} But there were two diverging tendencies proceeding from the same source. The author of the Confession himself understood and explained it differently, and the Formula added new dogmas which he never entertained. It excludes, indeed, certain extravagances of the Flacian wing of Lutheranism, but, upon the whole, it is a condemnation of Philippism and a triumph of exclusive Lutheranism.\footnote{\par Andreæ, in a letter to Heshusius and Wigand, of July 24, 1576, giving an account of the results of the Torgau Convention (quoted by Heppe, Vol. III. p. 111), thus characteristically sets forth the object of the whole movement in which he and the Elector Augustus were the chief leaders: '\emph{Hoc enim sancte vobis affirmare et polliceri ausim, Illust. Electorem Saxoniæ in hoc unice intentum, ut } Lutheri Doctrina  \emph{partim obscurata, partim vitiata, partim aperte vel occulte damnata, pura et sincera in scholis et Ecclesiis restituatur, adeoque } Lutherus, hoc est Christus,  \emph{cuius fidelis minister Lutherus fuit, vivat. Quid vultis amplius? Nihil hic fucatum, nihil palliatum, nihil tectum est, sed juxta } spiritum Lutheri, qui Christi est.' And Chemnitz wrote, June 29, 1576: '\emph{Mentio librorum Philippi expuncta est, et responsione hoc in parte retulimus nos ad Lichtenbergense decretmn.}' Some zealots, like Heshusius, desired that Melanchthon should be condemned, by name, in the Formula, but Andreæ thought it better 'to cover the shame of Noah,' and to be silent about the apostasy of the Lutheran Solomon. Dr. Krauth, too, says (\emph{Conservative Reform.} p. 327): 'The Book of Concord treats Melanchthon as the Bible treats Solomon. It opens wide the view of his wisdom and glory, and draws the veil over the record of his sadder days.' In the Formula itself he is nowhere named, but in the Preface to the 'Book of Concord' his writings are spoken of as '\emph{utilia neque repudianda ac damnanda, quatenus cum ea norma, quæ Concordiæ libro expressa est, per omnia consentiunt.}'}

The spirit of Melanchthon could be silenced, but not destroyed, for it meant theological progress and Christian union. It revived from time to time in various forms, in Calixtus, Spener, Zinzendorf, Neander, and other great and good men, who blessed the Lutheran Church by protesting against bigotry and the overestimate of intellectual orthodoxy, by insisting on personal, practical piety, by widening the horizon of truth, and extending the hand of fellowship to other sections of Christ's kingdom. The minority which at first refused the Formula became a vast majority, and even the recent reaction of Lutheran confessionalism against rationalism, latitudinarianism, and unionism will be unable to undo the work of history, and to restore the Lutheran scholasticism and exclusivism of the seventeenth century. The Lutheran Church is greater and wider than Luther and Melanchthon, and, by its own principle of the absolute supremacy of the Bible as a rule of faith, it is bound to follow the onward march of Biblical learning.

The great length of this section may be justified by the intrinsic importance of the Formula Concordiæ, and the scarcity of reliable information in English works.\footnote{There is no full and satisfactory account of the history and character of the Form of Concord in the English language, except in Dr. Krauth's \emph{Conservative Reformation and its Theology}, pp. 288–328; and this, in accordance with the aim of this learned and able author, is apologetic and polemic rather than historical. Dr. Shedd, in his valuable \emph{History of Christian Doctrine} (Vol. II. p. 458), devotes only a few lines to it. Dr. Fisher, in his excellent work on the \emph{Reformation }(N. Y. 1873), disposes of it in a foot-note (p. 481). In Dr. Blunt's \emph{Dictionary of Sects}, etc. (London, 1874), it has no place among the \emph{Protestant Confessions}, and the brief allusion to it sub '\emph{Lutherans},' p. 269, only exposes the ignorance of the writer. The doctrines of the Form of Concord are frequently, though mostly polemically, noticed in Dr. Hodge's \emph{Systematic Theology} (N.Y. 1873, 3 vols.).}

\xsection{Superseded Lutheran Symbols. The Saxon Confession. The Würtemberg Confession. 1551.}

§ 47. Superseded Lutheran Symbols. The Saxon Confession. The Würtemberg Confession. 1551.

Literature.

Heinrich Heppe:  \emph{Die Bekenntniss-Schriften, der altprotestantischen Kirche Deutschlands}, Cassel, 1855. This collection contains (besides the œcumenical Creeds, the Augsburg Confession of 1530, the Altered Augsburg Confession of 1540) the \emph{Confessio Saxonica}, pp. 407–483, and the \emph{Confessio Würtembergica}, pp. 491–554.

Phil. Melanchthonis  \emph{Opera quæ supersunt omnia}, or \emph{Corpus Reformatorum}, ed. Bretschneider and Bindseil, Vol. XXVIII. (Brunsvigæ, 1860), pp. 329–568. This vol. contains the Latin and German texts of the \emph{Conf. Saxonica} with critical Prolegomena.

The Book of Concord embraces all the Lutheran symbols which are still in force; but two other Confessions deserve mention for their historical importance, viz., the Saxon Confession and the Würtemberg Confession.

Both were written in 1551, twenty-one years after the Confession of Augsburg and twenty-six years before the Formula of Concord, in full agreement with the former as understood by its author, and without the distinctive and exclusive features of the latter. Both were intended (like the Articles of Smalcald) for the Roman Catholic Council, and, although they failed in accomplishing their direct object, they exhibit the doctrinal status of the Lutheran or the entire Evangelical Church of Germany at that period. It is this Protestantism which received legal toleration and recognition in the German Empire by the Treaty of Passau, 1552, and three years afterwards, without the restriction as to time, at the Diet of Augsburg.\footnote{Heppe, 1.c. p. xxix.: \emph{'Der in der Conf. Saxonica und in der Conf. Würtembergica entfaltete Lehrbegriff der Augsburgischen Confession ist es, welcher i. J.} 1555 \emph{zu kirchenstaatsrechtlicher Geltung kam. Dieses erhellt schon aus den Beschlüssen der im Mai} 1554 \emph{zur Vorbereitung der Reichstagsverhandlungen gehaltenen evangelischen Conferenz, in dem die daselbst versammelten chursäschsischen, hessischen und strassburgischen Deputirten erklärten: Auf bevorstehendem Reichstage habe man als einziges Bekenntniss die } Augsburgische Confession  \emph{festzuhalten. Da aber die sächsische und die würtembergische Confession mit derselben durchaus übereinstimmten, so habe man entweder jene oder eine von diesen dem Kaiser zu übergeben.}'} But in the succeeding generation the exclusive and more energetic school of Lutheranism prevailed, and found its expression in the Formula of Concord, which superseded those interimistic Confessions.

1. The Saxon Confession (Confessio Saxonica) was drawn up by Melanchthon for the Council of Trent, which, after a brief transfer to Bologna by Paul III., in March, 1547, was again convened at Trent by Julius III., May 1, 1551. The German Emperor had previously (Feb. 13) invited the Protestant States to send delegates, promising them full protection, and his best endeavor to secure 'a Christian, useful reformation, and abrogation of improper doctrines and abuses.' Melanchthon expected nothing from a conference with Bishops and Cardinals, but considered it wise and politic to accept the Emperor's invitation, provided he would secure to the Protestant delegates a hearing before the Council. His advice was the best that could be given under the circumstances, and was accepted by Elector Maurice of Saxony.\footnote{See several letters from February to April, 1551, in the \emph{Corp. Reform. }Vol. VII. (1840), especially pp. 736–739, where Melanchthon gives his views on the Council of Trent; and Schmidt, \emph{Melanchthon}, pp. 534 sqq.} He was requested to prepare a '\emph{Repetition and Exposition of the Augsburg Confession},' usually called the '\emph{Saxon Confession}.'\footnote{It appeared first in Latin at Basle, 1552, under the title: 'Confessio Do | ctrinæ Saxonicarum | Ecclesiarum  \emph{Synodo Tridentinæ ob }| \emph{lata, A.D.} 1551, \emph{in qua},' etc. The original MS., with the title 'Repetitio Confessionis Augustanæ: \emph{An.} 1551, \emph{Witebergæ scripta},' etc., and with corrections from Melanchthon's own hand, is preserved in the library of the \emph{Thomaskirche} in Leipzig, to which Selnecker presented it in 1580. From this Heppe and Bindseil have derived their text; the latter with a critical apparatus from eight printed editions. It was translated into German by John Maetsperger, 1552, and by Georg Major, 1555. The Latin text was often republished separately at Leipzig, Wittenberg, Frankfort, etc., and in the Melanchthonian \emph{Corpora Doctrinæ}; also in the \emph{Corpus et Syntagma Confessionum}, Genev. 1612 and 1654, in the \emph{Sylloge Confessionum}, Oxf. 1804 and 1827 (pp. 237–323); and more recently by H. Heppe, l.c., and by Bindseil, who gives also Major's German translation, in \emph{Corp. Reform.} Vol. XXVIII. pp. 370 sqq. On the various editions, see Bindseil, pp. 347 sqq.} To finish this work with more leisure, he went with his friend Camerarius to the Prince of Anhalt at Dessau.

The document is not merely a repetition of the Augsburg Confession, but an adaptation of it to the changed condition of affairs. In 1530 Melanchthon still hoped for a reunion with Rome, and wrote in an apologetic tone, avoiding all that might irritate the powerful enemy; now all hope of reunion had departed, and Protestantism had made a decided progress in ecclesiastical consolidation and independence. Although the Confession was composed after the defeat of the Protestant Princes by the Emperor, and in the midst of the Adiaphoristic troubles, it shows no disposition whatever to recede from the doctrinal positions taken at Augsburg; on the contrary, the errors and abuses of Rome, which made separation an imperative duty, are freely exposed and refuted. The Scriptures, as understood by the ancient Church in the œcumenical Creeds, are declared to be the only and unalterable foundation of the Evangelical faith.\footnote{Art I. \emph{De doctrina:} '\emph{Affirmamus clare coram Deo et universa Ecclesia in cælo et in terra, nos vera fide amplecti omnia } Scripta Prophetarum et Apostolorum:  \emph{et quidem in hac ipsa nativa sententia quæ expressa est in Symbolis, } Apostolico, Nicæno  \emph{et}  Athanasiano.'} The distinctive Evangelic doctrines and usages in opposition to Rome are comprehended under the two articles of the Apostles' Creed: 'I believe the forgiveness of sins,' and 'one holy Catholic Church.' The former excludes human merit and justification by works; the latter the political and secular conceptions and corruptions of the Church, which is represented to be a spiritual though visible communion of believers in Christ. The controverted articles are considered in twenty-three sections, in the order of the Augsburg Confession, namely: Original Sin, Forgiveness and Justification, Free Will, Good Works, New Obedience, the Church, the Sacraments, Satisfaction, Marriage, Monastic Life, Invocation of Saints, Civil Magistrate. The Saxon Confession is signed, not by Princes, as the Augsburg Confession was, but, as Melanchthon suggested, only by theologians, viz., Bugenhagen, Pfeffinger, Camerarius, Major, Eber, Melanchthon, and the Superintendents of Electoral Saxony, who convened at Wittenberg, July 9, for the purpose, and unanimously adopted the work of their dear and venerable 'Preceptor,' as the clear expression of their own faith in full harmony with his Confession of 1530. It was a beautiful moment in Melanchthon's life, for which he felt very grateful to God.\footnote{See his letter to Prince George of Anhalt, July 11, 1551, \emph{Corp. Reform.} Vol. VII. p. 806 sq., and the letter of Major to Jonas, July 14, \emph{ibid.} p. 809.} The danger was now much greater than in 1530, for the Elector Maurice was in league with the victorious Emperor. The theologians of Brandenburg, Ansbach, Baireuth, Mansfeld, Pomerania, Palatinate, Hesse, Würtemburg, and Strasburg likewise sent in their consent to this Confession.\footnote{See Heppe, 1.c. p. xxvii., and especially the \emph{Corpus et Syntagma Conf.}, which gives after the subscriptions the assenting judgments of the churches above mentioned.}

The Council convened in May, 1551, was adjourned to October, and again to January next. Melanchthon was ordered to proceed to Trent, but to stop at Nuremberg for further instructions. While at Nuremberg, in January, 1552, he wrote a preface to Luther's Commentary on Genesis, and expressed himself very decidedly against the preceding acts of the Council.\footnote{Jan. 25, 1552, \emph{Corp. Reform.} Vol. VII. pp. 918–927.} In the mean time the Saxon and Würtemberg lay-embassadors received a hearing at Trent, not, indeed, before the whole Council in public session, but before a private congregation. They requested that the members of the Council be released of their oath of obedience to the Pope, and be free to decide the questions by the rule of the Scriptures alone. A few prelates were inclined to accede, but the majority would never have sacrificed the principle of tradition, nor reconsidered the decrees already adopted. The Saxon embassadors urged Melanchthon to proceed on his journey, but he delayed on account of the rumors of war. The treacherous Elector Maurice of Saxony cut the Gordian knot by making war upon his ally, the Emperor, in the spring, 1552, drove him from Innspruck, scared the fathers of Trent to their homes, and achieved, in the Treaty of Passau (Aug. 2, 1552), ratified at Augsburg (1555), the first victory for liberty of conscience to Protestants, to which the Emperor reluctantly yielded, and against which the Pope never ceases to protest.

II. The Würtemberg Confession (Confessio Würtembergica)\footnote{The full title, as given by Heppe and Bindseil, is 'Confes | sio Piæ Doctri | næ,  \emph{quæ nomine illu }| \emph{strissimi Principis ac Domini }  Chri | stophori  \emph{Ducis Wirtembergen }| \emph{sis et Teccensis, ac Comitis Montisbe }| \emph{ligardi, per legatos ejus Die XXIIII. } | \emph{ mensis Januarij, Anno MDLII. Con } | \emph{gregationi Tridentini Conci }| \emph{lii proposita est.}' It was first printed at Tübingen, 1551; then in 1556, 1559, 1561, etc. It is also embodied in the \emph{Opera Brentii}, Tübingen, 1590, Tom. VIII. pp. 1–34, in \emph{Corpus et Syntagma Conf.} (from a Frankfort ed. of 1561), and in Heppe, 1.c. pp. 491–554. It is frequently quoted in part under different heads, together with the Saxon Confession, in the Reformed \emph{Harmonia Confessionum}, Genev. 1581. Comp. Pfaff, \emph{Acta et scripta publica Ecclesiæ Wirtembergicæ}, Tüb. 1720; Salig, \emph{Historie der Augsb. Conf.} Tom. I. pp. 673 sqq.; and Hartmann, \emph{Johannes Brentz. Leben und ausgewählte Schriften} (Elberfeld, 1862), pp. 211–221.} Was prepared for the same purpose, at the same time and in the same spirit, by Brentius, the Reformer of the Duchy of Würtemberg, in the name of his Prince, Duke Christopher, who likewise resolved to send delegates to the Council of Trent. For Brentius, like Melanchthon, had no confidence in this partial popish Council, but advised, nevertheless, compliance with the Emperor's request, since a refusal might be construed as disobedience and contempt, or as an act of cowardice. The Confession was approved by a commission of ten Swabian divines, and by the City of Strasburg. It was also approved at Wittenberg, as agreeing with Melanchthon's Confession. It was found best to send two Confessions, one representing the Evangelical Churches of the North, the other those of the South of Germany, to avoid the appearance of a conspiracy.

The Würtemberg Confession contains a preface of Duke Christopher, and restates, in thirty-five articles, the doctrines of the Augsburg Confession and other controverted points, for the purpose of showing that the Evangelical Churches agree with the pure doctrine of the apostles, and of the catholic and orthodox Church.\footnote{Prefat.: '\emph{In nostris ecclesiis non nisi veræ apostolicæ, catholicæ, et orthodoxæ doctrinæ locum datum esse.}'} On the Lord's Supper this Confession goes a little beyond the Saxon; but there is no trace of the ubiquity of Christ's body, of which Brentius, ten years afterwards, became a zealous advocate.

Brentius was among the Würtemberg and Strasburg delegates to Trent, and actually arrived there, March 18, 1552, but only to return in April without accomplishing any thing.\footnote{See Sleidanus, \emph{De statu relig. et reipublicæ Carolo V. Cæsare commentar.} Tom. III. pp. 317–333; \emph{Corp. Reform.} Vol. XXVIII. p. 334, and Hartmann, 1.c. p. 215. The other theological delegates to Trent were Beurlin, Heerbrand, Vannius (Wanner), of Würtemberg, and Marbach and Sellius, of Strasburg. Sleidanus was one of the lay-delegates from Strasburg.} It is very doubtful whether he and Melanchthon would have made a deep impression upon the Council, which was already committed to the cause of popery and had sanctioned some of its most obnoxious doctrines.

\xsection{The Saxon Visitation Articles, A.D. 1592.}

§ 48. The Saxon Visitation Articles, 1592.

Literature.

Articuli Visitatorii,  \emph{Anno Christi} 1592 \emph{in Electoratu et Provinciis superioris Saxoniæ publicati, et Judicibus Con istoriorum, Superintendentibus, Ministris ecclesiarum et scholarum, nec non Administratoribus bonorum ecclesiasticorum, quin et ipsis Patronis et Collatoribus ad subscribendum et servandum propositi et demandati.} They are printed in \emph{Corp. juris eccles. Saxonici}, Dresden, 1773, p. 256, and added to Hase's edition of the Lutheran Symbols, pp. 862–866, the Berlin edition of the \emph{Concordia} (1857), pp. 849–854, and Müller's \emph{Symb. Bücher}, pp. 779–784.

\emph{Gründliche Verantwortung der vier streitigen Artikel}, etc. Leipzig, 1593.

A. Hunnius:  \emph{Widerlegung des Calvinischen Büchleins wider die vier Artikel}, 1593.

Comp. Schroekh:  \emph{Kirchengeschichte seit der Reformation}, Vol. IV. pp. 660 sqq.; Henke:  \emph{Art. Hunnius} in Herzog, Vol. VI. pp. 316–321; Müller:  \emph{Symb. Bücher}, pp.cxxi. (Introd.) sqq.; G.Frank:  \emph{Geschichte der Protest. Theologie} (1864), Vol. I. pp. 290 sqq.

The  Four Articles of Visitation of Electoral Saxony owe their origin to the revival and second overthrow of Crypto-Calvinism, and reflect the fierceness and bitterness of this contest.\footnote{See above, p. 283.} They continued in force till the present century, but never extended their authority beyond Saxony. They are strongly anti-Calvinistic, and may be regarded as an Appendix to the Formula of Concord, with which they fully agree.

They were written in 1592, and first published in German in 1593.\footnote{Under the title: '\emph{Visitation-Artikel im gantzen Churkreiss Sachsen. Sampt derer Calvinisten Negativa und Gegenlehr, und die Form der Subscription, welchergestalt dieselbe beyden Partheien sich zu unterschreiben sind vorgelegt worden.}'} Their object was to perpetuate the reign of exclusive Lutheranism. They are based on the articles of a Colloquy between Andreæ and Beza at Mömpelgard (1586). The chief author was Dr. Aegidius Hunnius, one of the foremost Lutheran divines of his age, a native of Winnenden in the Duchy of Würtemberg, professor of theology at Marburg (1576–1592), and afterwards at Wittenberg (d. 1603).\footnote{He was aided in the composition by Mart. Mirus, George Mylius, and Joshua Lonnerus. Mirus was called by Hospinian '\emph{Inquisitor Saxoniæ}, because, as the Lutherans explained this term of reproach, he cleaned the Lord's vineyard of cunning foxes and wild hogs. His last wish was to die an enemy of Calvinists and Papists. Frank, l.c. Vol. I. p. 296.} He was commissioned with several others to visit the churches and schools of Saxony for the purpose of suppressing every trace of Crypto-Calvinism. All clergymen and teachers, and even the civil officers, were required to subscribe the four Articles or lose their places. A great feast of thanksgiving closed the visitation.

The hardest fate was reserved for Chancellor Crell, who, after ten years' imprisonment, was executed (1601), ostensibly for political offenses, but really for opinions which were once honored by the name of Philip Melanchthon. The preachers who attended this auto-da-fé of hyper-Lutheran orthodoxy told Crell that by his wicked Calvinism he had caused in many cases a dangerous delay of infant baptism, undermined the authority of the ministry, and deserved the fire of hell. They laughed at his prayer on the scaffold; whereupon he prayed to God not to change their laughter into weeping. The executioner, holding the severed head high up in the air, said: 'This was a Calvinistic stroke.'\footnote{See Frank, Vol. I. p. 297, and Henke's monograph on \emph{Casp. Peucer und Nic. Crell, }1865.}

The four Articles give a very clear and explicit summary of those peculiar doctrines which distinguish the Lutheran creed from those of all other Protestant churches. The first refers to the Lord's Supper, and teaches the \emph{real} presence and \emph{oral} fruition of the \emph{true} and \emph{natural} body of Christ by \emph{all} communicants. The second treats of the Person of Christ, and teaches, in support of the eucharistic omnipresence, the communication of the attributes whereby the \emph{human nature} of Christ became partaker of the whole majesty, honor, power, and glory of his divine nature. The third teaches \emph{baptismal regeneration} and the ordinary necessity of baptism for salvation.\footnote{Baptism was performed with exorcism in Lutheran churches, and it was counted one of the chief crimes of the Crypto-Calvinists that they abolished this rite. A Saxon pastor who baptized without exorcism gave great offense to the peasants, who cried after him: 'The naughty priest has not expelled the devil' (\emph{Der lose Pfaffe hat den Teufel nicht ausgetrieben}).} The fourth teaches the \emph{universal atonement}, and the \emph{vocation} of \emph{all} men to salvation, with the possibility of a \emph{total} and \emph{final} fall from grace.

In the negative part the opposite doctrines of the Calvinists are rejected. These were henceforth held in perfect abhorrence in Saxony, and it was a common proverb, 'Rather a Papist than a Calvinist.'\footnote{It is almost incredible to what extent the Lutheran bigotry of those days carried its hatred of Zwinglianism and Calvinism. We give a few characteristic specimens. Schlüsselburg (Superintendent of Ratzeburg), one of the most learned champions of Lutheran orthodoxy, in his \emph{Theologiæ Calvinistarum Libri Tres}, Francoforti ad Mænum, 1592, tries to prove that the Calvinists are unsound in almost every article of the Christian faith ('\emph{Sacramentarios de nullo fere doctrinæ Christianæ articulo recte sentire}'), and has a special chapter to show that the Calvinistic writings overflow with \emph{ mendaciis, calumniis, conviciis, maledictis, et contumeliis.} He regards many of their doctrines as downright blasphemy. Philip Nikolai, a pious Lutheran pastor at Unna, afterwards at Hamburg, and author of two of the finest German hymns ('\emph{Wie schön leuchtet der Morgenstern},' and '\emph{Wachet auf! ruft uns die Stimme}'), called the God of the Calvinists 'a roaring bull (\emph{Wucherstier und Brüllochs}), a bloodthirsty Moloch, a hellish Behemoth and Leviathan, a fiend of men!' (\emph{Kurtzer Bericht von der Calvinisten Gott und ihrer Religion}, Frkf. 1597; \emph{Die erst Victoria, Triumph und Freudenjubel über des Calvin Geistes Niederlag}, 1600; \emph{Calvinischer Vitzliputzli}, etc. See Frank, Vol. I. p. 280. Provost Magirus, of Stuttgart, thought that the Calvinists imitated at times the language of Luther, as the hyena the human voice, for the destruction of men. John Modest wrote a book to prove that the Sacramentarians are no Christians, but baptized Jews and Mohammedans ('\emph{Beweis aus der heiligen Schrift dass die Sacramentirer nicht Christen sind, sondern getaufte Juden und Mahometisten}, Jena, 1586). John Prätorius, in a satire (\emph{Calvinisch Gasthaus zur Narrenkayffen}, etc.), distinguishes open Calvinists, who have no more sense than a horse or an ass; secret Calvinists, who fish in the dark; and several other classes (see Frank, Vol. I. p. 282 sq.). The second Psalm, speaking of the rebellion against Jehovah and his Anointed, was applied to the Calvinists, and their condemnation was embodied in catechisms, hymns, and popular rhymes, of which the following are fair specimens:  \par '\emph{Erhalt uns, Herr, bei deinem Wort} \par \emph{ Und wehr der Calvinisten Mord.}' \par '\emph{Wenn ein Calvinist spricht, Gott grüss dich,} \par \emph{ So wünscht sein Herz, der Tod hol dich.}' \par '\emph{Gottes Wort und Luther's Lehr} \par \emph{ Vergehet nun und nimmermehr,} \par \emph{ Und ob's gleich bisse noch so sehr} \par \emph{ Die Calvinisten an ihrer Ehr.}' \par '\emph{Gottes Wort und Lutheri Schrift } \par \emph{ Sind des Papsts und Calvini Gift.}'}

As the Articles are a very clear and succinct statement of the specific doctrines of Lutheranism as opposed to Calvinism, and not easy of access, they are here given in full:

Articulus I.

\emph{De Sacra Cœna.}

pura et vera doctrina nostrarum ecclesiarum de sacra cœna.

I. Quod, verba Christi: '\emph{Accipite et comedite, hoc est corpus meum: Bibite, hic est sanguis meus},' simpliciter, et secundum literam, sicut sonant, intelligenda sint.

II. Quod in Sacramento duæ res sint, quæ exhibentur et simul accipiuntur: una terrena, quæ est panis et vinum; et una cœlestis, quæ est corpus et sanguis Christi.

III. Quod hæc Unio, Exhibitio et Sumptio fiat hic inferius in terris, non superius in cœlis.

IV. Quod exhibeatur et accipiatur verum et naturale corpus Christi, quod in cruce pependit, et verus ac naturalis sanguis, qui ex Christi latere fluxit.

V. Quod corpus et sanguis Christi non fide tantum spiritualiter, quod etiam extra Cœnam fieri potest, sed cum pane et vino oraliter, modo tamen imperscrutabili et supernaturali, illic in Cœna accipiantur, idque in pignus et certificationem resurrectionis nostrorum corporum ex mortuis.

VI. Quod oralis perceptio corporis et sanguinis Christi non solum fiat a dignis, verum etiam ab indignis, qui sine pœnitentia et vera fide accedunt; eventu tamen diverso. A dignis enim percipitur ad salutem, ab indignis autem ad judicium.

Articulus II.

\emph{De Persona Christi.}

pura et vera doctrina nostrarum ecclesiarum de hoc articulo, de persona christi.

I. In Christo sunt duæ distinctæ Naturæ, divina et humana. Hæ manent in æternum inconfusæ et inseparabiles (seu indivisæ).

II. Hæ duæ Naturæ personaliter ita sunt invicem unitæ, ut unus tantum sit Christus, et una Persona.

III. Propter hanc personalem Unionem recte dicitur, atque in re et veritate ita se habet, quod Deus Homo, et Homo Deus sit, quod Maria Filium Dei genuerit, et quod Deus nos per proprium suum sanguinem redemerit.

IV. Per hanc Unionem personalem, et quæ eam secuta est, exaltationem, Christus secundum carnem ad dexteram Dei collocatus est, et accepit omnem potestatem in cœlo et in terra, factusque est particeps omnis divinæ majestatis, honoris, potentiæ et gloriæ.

Articulus III.

\emph{De Baptismo.}

pura et vera doctrina nostrarum ecclesiarum de hoc articulo s. baptismatis.

I. Quod unum tantum Baptisma sit, et una ablutio, non quæ sordes corporis tollere solet, sed quæ nos a peccatis abluit.

II. Per Baptismum tanquam lavacrum illud regenerationis et renovationis Spiritus Sancti salvos nos facit Deus et operatur in nobis talem justitiam et purgationem a peccatis, ut qui in eo fœdere et fiducia usque ad finem perseverat, non pereat, sed habeat vitam æternam.

III. Omnes, qui in Christum Jesum baptizati sunt, in mortem ejus baptizati sunt, et per Baptismum cum ipso in mortem ejus consepulti sunt, et Christum induerunt.

IV. Baptismus est lavacrum illud regenerationis, propterea, quia in eo renascimur denuo et Spiritu Adoptionis obsignamur ex gratia (sive gratis).

V. Nisi quis renatus fuerit ex aqua et Spiritu, non potest introire in regnum cœlorum. Casus tamen necessitatis hoc ipso non intenditur.

VI. Quicquid de carne nascitur, caro est, et natura sumus omnes filii iræ divinæ: quia ex semine peccaminoso sumus geniti, et in peccatis concipimur omnes.

Articulus IV.

\emph{De Prædestinatione et æterna Providentia Dei.}

pura et vera doctrina nostrarum ecclesiarum de hoc articulo.

I. Quod Christus pro omnibus hominibus mortuus sit, et ceu Agnus Dei totius mundi peccata sustulerit.

II. Quod Deus neminem ad condemnationem condiderit, sed velit, ut omnes homines salvi fiant et ad agnitionem veritatis perveniant, propterea omnibus mandat, ut Filium suum Christum in Evangelio audiant, et per hunc auditum promittit virtutem et operationem Spiritus Sancti ad conversionem et salutem.

III. Quod multi homines propria culpa pereant: alii, qui Evangelium de Christo nolunt audire, alii, qui iterum excidunt gratia, sive per errores contra fundamentum, sive per peccata contra conscientiam.

IV. Quod omnes peccatores pœnitentiam agentes in gratiam recipiantur, et nemo excludatur. etsi peccata ejus rubeant ut sanguis; quandoquidem Dei misericordia major est, quam peccata totius mundi, et Deus omnium suorum operum miseretur.

Sequitur Falsa et Erronea Doctrina Calvinistarum.

\emph{De Sacra Cœna.}

I. Quod supra posita verba Christi figurate intelligenda sint, et non secundum literam, sicut sonant.

II. Quod in Cœna tantum nuda signa sint, corpus autem Christi tam procul a pane, quam supremum cœlum a terra.

III. Quod Christus illic præsens sit tantum virtute et operatione sua, et non corpore suo. Quemadmodum sol splendore et operatione sua in terris præsens et efficax est, corpus autem solare superius in cœlo existit.

IV. Corpus Christi esse typicum corpus, quod pane et vino tantam significetur et præfiguretur.

V. Quod sola fide, quæ in cœlum se elevet, et non ore, accipiatur.

VI. Quod soli digni illud accipiant, indigni autem, qui talem fidem evolantem sursum in cœlos non habent, nihil præter panem et vinum accipiant.

Falsa et Erronea Doctrina Calvinistarum.

\emph{De Persona Christi,}

quæ potissimum iii. et iv. articulo purioris doctrinæ repugnat.

I. Quod Deus Homo, et Homo Deus est, esse figuratam locutionem.

II. Quod humana Natura cum divina non in re et veritate, sed tantum nomine et verbis communionem habeat.

III. Quod Deo impossibile sit ex tota omnipotentia sua præstare, ut corpus Christi naturale simul et instantanee in pluribus, quam in unico loco sit.

IV. Quod Christus secundum humanam Naturam per exaltationem suam tantnm creata dona et finitam potentiam acceperit, non omnia sciat aut possit.

V. Quod Christus secundum Humanitatem absens regnet, sicut Rex Hispaniæ novas Insulas regit.

VI. Quod damnabilis idololatria sit, si fiducia et fides cordis in Christum non solum secundum divinam, sed etiam secundum humanam ipsius Naturam collocetar, et honor adorationis ad utramque dirigatur.

Falsa et Erronea Doctrina Calvinistarum.

\emph{De Sacro Baptismo.}

I. Baptismum esse externum lavacrum aquæ, per quod interna quædam ablutio a peccatis tantum significetur.

II. Baptisimum non operari neque conferre regenerationem, fidem, gratiam Dei et salutem, sed tantum significare et obsignare ista.

III. Non omnes, qui aqua baptizantur, consequi eo ipso gratiam Christi aut donum fidei sed tantum electos.

IV. Regenerationem non fieri in, vel cum Baptismo, sed postea demum crescente aetate, imo et multis in senectute demum contingere.

V. Salutem non dependere a Baptismo, atque ideo Baptismum in causa necessitatis non permittendum esse in Ecclesia, sed in defectu ordinarii Ministri Ecclesiæ permittendum esse, ut infans sine Baptismo moriatur.

VI. Christianorum infantes jam ante Baptismum esse sanctos, ab utero matris, imo adhuc in utero materno constitutes esse in fœdere vitae æternæ cæteroqui Sacrum Baptisma ipsis conferri non posse.

Falsa et Erronea Doctrina Calvinistarum.

\emph{De Prædestinatione et Providentia Dei.}

I. Christum non pro omnibus hominibus, sed pro solis electis mortuum esse.

II. Deum potissimam partem hominum ad damnationem æternam creasse, et nolle, ut potissima pars convertatur et vivat.

III. Electos et regenitos non posse fidem et Spiritum Sanctum amittere, aut damnari, quamvis omnis generis grandia peccata et flagitia committant.

IV. Eos vero, qui electi non sunt, necessario damnari, nec posse pervenire ad salutem, etiamsi millies baptizarentur, et quotidie ad Eucharistiam accederent, præterea vitam tam sancte atque inculpate ducerent, quantum unquam fieri potest.

\xsection{An Abortive Symbol against Syncretism, 1655.}

§ 49. AN ABORTIVE SYMBOL AGAINST SYNCRETISM, 1655.

Finally, we must briefly notice an unsuccessful attempt to increase the number of Lutheran symbols which was made during the Syncretistic controversies in the middle of the seventeenth century.\footnote{H. Schmid:  \emph{Geschichte der Synkretistischen Streitigkeiten in der Zeit des Georg Calixt}, Erlangen, 1846. W. Gass:  \emph{G. Calixt und der Synkretismus}, Breslau, 1846; and his \emph{Geschichte der Protest. Dogmatik}, Vol. II. p. 68. Baur:  \emph{Ueber den Charakter und die Bedeutung des calixtin. Synkretismus}, in the \emph{Theol. Jahrbücher} for 1848, p. 163. E. L. Th. Henke:  \emph{G. Calixtus und seine Zeit}, Halle, 1853–1860, 2 vols.; and his Art. \emph{Synkretismus} and \emph{Synkretistische Streitigkeiten}, in Herzog, Vol. XV. (1862), pp. 342 and 346. G. Frank:  \emph{Geschichte der Protest. Theologie.} Leipz. Vol. II. 1865, p. 4.}

George Calixtus (1586 to 1656), Professor of Theology in the University of Helmstädt (since 1614), which had previously protested against the ubiquity dogma of the Formula of Concord, was disgusted with the exclusive and pugnacious orthodoxy of his day, and advocated, in the liberal and catholic spirit of Melanchthon, peace and conciliation among the three great Confessions—the Lutheran, Catholic, and Reformed. He went back to the Apostles' Creed and the œcumenical \emph{consensus} of the first five centuries (\emph{consensus quinquesecularis}) as a common basis for all, claiming for the Lutheran Church only a superior purity of doctrine, and surrendering as unessential its distinctive peculiarities. This reaction against sectarian exclusiveness and in favor of Catholic expansion within the Lutheran communion was denounced by the orthodox divines of Wittenberg and Leipzig as \emph{Syncretism}, i.e., as a Babylonian mixture of all sorts of religions, or a Samaritan compound of Popish, Calvinistic, Synergistic, Arminian, and even atheistic errors. A war to the knife was waged against it, and lasted from 1645 to 1686. Calixtus had expressed a hope to meet many Calvinists in heaven, but this was traced directly to an inspiration of the devil.

The chief opponent of Syncretism was Abraham Calovius, the fearless champion of an infallible orthodoxy, admired by some as the Lutheran Athanasius, abhorred by others as the Lutheran Torquemada; in his own estimation a \emph{strenuus Christi athleta}, certainly a veritable \emph{malleus hæreticorum;} of vast learning and a herculean working power, which no amount of domestic affliction could break down.\footnote{Abraham Calov (properly Kalau) was born in 1612 at Mohrungen, Prussia (the birthplace of the great Herder—'Esau and Jacob from one womb'), and labored with untiring industry as Professor and General Superintendent at Wittenberg from 1650 to his death, 1686. He stood in high esteem, and controlled the whole faculty, except Meisner, who fell out with him in 1675, so that they no more greeted each other, not even at the communion altar. The Elector, George II., always stayed at his house when he was at Wittenberg. Calovius wrote a system of theology, in twelve volumes (\emph{Sytstema locorum theolog.} 1655–1677), a Commentary on the whole Bible against Grotius, in four folios (\emph{Biblia illustrata}, 1672), and an endless number of polemical works against ancient and modern heretics, some of which had to be prohibited. His domestic history is perhaps without a parallel. He buried no less than thirteen children and five wives in succession. At the death-bed of the fourth he sang with all his might the hymn, '\emph{Wie schön leuchtet der Morgenstern},' especially (as Tholuck relates) the last stanza, '\emph{Wie bin ich doch so herzlich froh.}' etc. He asked her whether she were willing to go to her Lord; she replied: '\emph{Herr Jesu, dir leb' ich, Herr Jesu, dir sterb' ich.}' A few months after the death of his fifth partner, when seventy-two years of age ('\emph{senili amore, morbo nequaquam senili, vehementer laborans},' and '\emph{maxima cum multorum offensione}'), he led to the altar the youthful daughter of his colleague, Quenstädt. A friend of Spener wrote to the latter, May 10, 1684 (as quoted by Tholuck): 'The septuagenarian \emph{senex consularis} has prostituted himself strongly \emph{intra} and \emph{extra ecclesiam.} What is the use of all learning, if one can not control his appetites? He is said to be so debilitated that he can not walk five steps \emph{sine lassitudine.}' Calovius enjoyed his sixth marriage only two years. For a full account of him, see Tholuck,  \emph{Wittenberger Theologen}, 1852, pp. 185–211, and his Art. \emph{Calov}, in Herzog, Vol. II. p. 506; also Gass,  \emph{Geschichte der protest. Dogm.} Vol. I. p. 332; and G. Frank, Vol. II. p. 26. Tholuck characterizes him thus (\emph{W. Theol.} p. 207): '\emph{ Gemüthlose Zähigkeit bei innerlich kochender Leidenschaftlichkeit erscheint als Grundzug dieses theologischen Charakters; weder auf der Kanzel, noch in vertraulichen Briefen, noch in den theologischen Schriften ein Lebenshauch christlicher, selten auch nur menschlicher Wärme. Die Menschen erscheinen ihm wie Zahlen, und unter den dogmatischen Problemen bewegt er sich wie unter Rechenexempeln.}'} His daily prayer was, '\emph{Reple me, Deus, odio hœreticorum.}' He excluded Calixtus, as well as Bellarmin, Calvin, and Socinus, from heaven. As the best means of suppressing this complex syncretistic heresy, and of preventing a schism in the Lutheran Church, he prepared in 1655 a \emph{Repeated Consensus of the truly Lutheran Faith}, which was finally published in Latin and German at Wittenberg in 1664.\footnote{'\emph{Consensus repetitus fidei vere Lutheranæ in illis doctrinæ capitibus, quæ contra puram et invariatam Augustanam Confessionem aliosgue libros symbolicos in Libro Concordiæ comprehensos, scriptis publicis impugnant D. G. Calixtus, ejusque complices.}' First published in the \emph{Consilia Theologica Wittebergensia}, 1664, then often separately by Calovius. A new edition by the late Prof. Henke of Marburg: \emph{Consensus repetitus fidei veræ Lutheranæ, MDCLV. Librorum ecclesiæ evangelicæ symbolicorum supplementum}, Marburg, 1847 (pp. viii. and 70). For a summary, see H. Schmid, l.c. pp. 376 sqq., and Frank, l.c. Vol. II. pp. 12 sqq. Calovius wrote no less than twenty-eight books against the Syncretists, the principal of which are \emph{Syncretismus Calixtinus}, 1653; \emph{Synopsis controversiarum . . . cum hæreticis et schismaticis modernis Socinianis, Anabaptistis, Weigelianis, Remonstrantibus, Pontificiis, Calvinianis, Calixtinis}, etc. 1652; and \emph{Harmonia Calixtino-hæretica}, etc., 1655. See H. Schmid, l.c. p. 237, who with all his orthodox sympathies complains of the endless repetitions and prolixity of these controversial writings. They are almost unreadable. I have before me a defense of the \emph{Consensus Repetitus}, by Aegidius Straucher. Wittenb. 1668 (551 pp.), the mere title of which covers twenty-nine lines.}

This creed first professes and teaches, in the order of the Augsburg Confession, the orthodox doctrine, and then rejects and condemns no less than eighty-eight syncretistic heresies, proved from the writings of Calixtus, Hornejus, Latermann, and Dreier. The first fundamental section anathematizes the Calixtine concession of the imperfection of the Lutheran Church, the relative recognition of Catholics and Calvinists as Christian brethren, and the assertion of the necessity of Church tradition alongside of the Scriptures. The following doctrines are rejected, not simply as doubtful, erroneous, or dangerous opinions (which some of them are), but as downright heresies: That the article of the Trinity is not clearly revealed in the Old Testament; that the Holy Spirit dwells in believers as a gift, not as an essence; that theology need not prove the existence of God, since it is already certain from philosophy; that Jews and Mohammedans are not idolaters; that original sin is simply a \emph{carentia justitiæ}; that souls are created by God (creationism); that Christ's body is not omnipresent; that sanctification enters in any way into the idea of justification; that the true Church embraces also Calvinists, Papists, and Greeks; that infants have no faith; that John vi. treats of the Lord's Supper; that man is active in his conversion; that symbolical books are to be only conditionally subscribed \emph{quatenus Scripturæ S. consentiunt}; that the symbols contain many things as necessary to salvation, which God has not fixed as such; that unbaptized infants are only negatively punished; that good works are necessary to obtain eternal life. A prayer that God may avert all innovations and corruptions from the Orthodox Church, and preserve it in this repeated consensus, forms the conclusion.

This new symbol goes far beyond the Formula of Concord, and would have so contracted Lutheranism as to exclude from it all independent thought and theological progress. It prolonged and intensified the controversy, but nowhere attained ecclesiastical authority. It was subscribed only by the theological faculties of Wittenberg and Leipzig, and rejected by the theologians of Jena, who were pupils of the celebrated John Gerhard, and occupied a milder position. With the death of Calovius the controversy died out, and his symbol was buried beyond the hope of a resurrection. Orthodoxy triumphed, but it was only a partial victory, and the last which it achieved.

During these violent controversies and the awful devastations of the Thirty-Years' War, there arose among a few divines in the Lutheran, Reformed, and Catholic Churches an intense desire for the reunion of Christendom, which found its expression in the famous adage so often erroneously attributed to St. Augustine: '\emph{In necessariis unitas, in dubiis libertas, in omnibus caritas.}\footnote{Dr. Lücke (in a special treatise, Göttingen, 1850) traces the authorship with some degree of certainty to Rupert Meldenius, who belonged to the irenical school of the seventeenth century. Comp. Klose, in Herzog, Vol. IX. p. 304.} It had no practical effect, but sounds like a prophecy of better times.

Soon afterwards arose a second and more successful reaction in the Pietism of Spener and Francke, which insisted on the claims of practical piety against a dead orthodoxy in the Lutheran Church, just as the school of Coccejus did in the Reformed Church of Holland, and the Methodism of Wesley and Whitefield in the Church of England. Then followed, toward the close of the eighteenth century, the far more radical reaction of Rationalism, which broke down, stone by stone, the venerable building of Lutheran orthodoxy, and the whole traditional system of Christian doctrine. Rationalism, in its various forms and phases, laid waste whole sections of Germany, especially those where once a rigorous orthodoxy had most prevailed; it affected also the Reformed churches of the Continent, and, in a less degree, those of England and America. Fortunately the power of this great modern apostasy has been broken, in the nineteenth century, by an extensive revival of the principles of the Reformation, with a better appreciation of its Confessions of Faith, not so much in their subordinate differences as in their essential harmony.

\xchapter{Chapter 7. The Creeds of the Evangelical Reformed Churches}

\xsection{The Reformed Confessions.}

§ 50. The Reformed Confessions.

Literature.

I. Collections of Reformed Symbols.

Harmonia | Confession | Fidei | Orthodixarum, et Reformatarum Ecclesiarum, | \emph{quæ in præcipuis quibusque Europæ Regnis, Nationibus, et Provinciis, sacrum Evangelii doctrinam pure profitentur: quarum catalogum et ordinem sequentes paginæ indicabunt. | Additæ sunt ad calcem brevissimæ observationes: quibus turn illustrantur obscura, tum quæ in speciem pugnare inter se videri possunt, perspicue atque modestissime conciliantur: et si quæ adhuc contraversa manent, syncere indiciantur. | Quæ omnia, Ecclesiarum Gallicarum, et Belgicarum nomine, subjiciuntur libero et prudenti reliquarum omnium judicio.} Genevæ apud Petrum Santandreanum. MDLXXXI. (4to).

This is the first attempt at comparative Dogmatics or Symbolics. It grew out of a desire for one common Creed, which was modified into the idea of a selected harmony. In this shape it was proposed by the Protestants of Zurich and Geneva, intrusted to Beza, Daneau, and Salnar (or Salnard, or Salvart, minister of the Church of Castres), and chiefly executed by the last of the three. It was intended as a defense of Protestant, and particularly Reformed, doctrine against the constant attacks of Romanists and Lutherans. It does not give the Confessions in full, but extracts from them on the chief articles of faith, which are classified under nineteen sections. It anticipates Winer's method, but for harmonistic purposes. Besides the principal Reformed Confessions, three Lutheran Confessions are also used, viz., the Augsburg, the Saxon, and the Würtemberg Confessions. The work appeared almost simultaneously with the Lutheran Formula of Concord, and may be called a Reformed Formula of Concord, though differing from the former in being a mere compilation from previous symbols. (I imported a well-bound copy, which seems to have been the property of the Elector John Casimir, whose likeness and escutcheon are impressed on the cover. He suggested the preparation of such a work.)

An English translation of this irenic work appeared first at Cambridge, 1586 (12mo), and then again in London, 1643 (4to), under the title: 'An Harmony of the Confessions of Faith of the Christian and Reformed Churches,  \emph{which purely profess the holy doctrine of the Gospel, in all the chief kingdoms, nations, and provinces of Europe,} etc. \emph{All which things, in the names of the Churches of France and Belgia, are submitted to the free and discreet judgment of all the Churches. Newly translated out of Latin into English,} etc. \emph{Allowed by public authority.}' According to Strype (\emph{Annals of the Reformation,} ad a. 1586), Archbishop Whitgift, owing to some jealousy among publishers, first forbade the publication of the \emph{Harmony,} but afterwards allowed it.

A new edition by Rev. Peter Hall (Rector of Milston, Wilts), under the modified title: The Harmony of Protestant Confessions:  \emph{exhibiting the Faith of the Churches of Christ, Reformed after the pure and holy doctrine of the Gospel, throughout Europe. Translated from the Latin. A new edition, revised and considerably enlarged.} London, 1842 (640 pages, large 8vo).

Corpus et Syntagma | Confessionum | Fidei,  | \emph{quæ in diversis regnis et nationibus, ecclesiarum nomine fuerunt authentice editæ: in celeberrimis conventibus exhibitæ, publicaque auctoritate comprobatæ,} etc. (first ed. Aureliæ Allobrog. 1612). \emph{Editio nova, Genevæ, sumptibus Petri Chouët, }1654.

The first edition of this rare and valuable book was probably compiled by \emph{Gaspar Laurentius,} who is not named on the title-page, but who signs himself in the dedicatory Epistle to Elector Frederick III. of the Palatinate, before the 'Orthodox Consensus' (in Part III.), and says, in the 'General Preface,' that he edited this Consensus a. 1595, and now (1612) in a much improved form. His object was the same as that of the \emph{Harmony,} viz., to show the essential unity of the evangelical faith in the multiplicity and variety of Confessions which, as the Preface says, in the absence of conspiracy, only strengthen the harmony, and mutually illustrate and supplement each other, like many orthodox expositions of the Scriptures. The second edition, of which I have a copy, is a large quarto volume, consisting of three main parts, the several documents being paged separately. It contains the principal Reformed Confessions down to the Synod of Dort, three Lutheran Confessions, and several other documents, as follows: 1. The \emph{Harmonia sive Concordantia Confessionum Fidei per} (xiii.) \emph{Articulos digesta,} with the \emph{Symbolum Apostolicum,} as the basis of a general \emph{consensus,} supported by Scripture texts and references to the various Confessions of the collection (8 pp.); 2. \emph{Confessio Helvetica posterior,} reprinted from a Zurich edition of 1651: 3. \emph{Confessio Helvetica prior} (or \emph{Basileensis II.}), 1536; 4. \emph{Confessio Basileensis I.} (or \emph{Mylhusiana}), 1532; 5. \emph{Confessio Gallica,} from the Latin edition of 1566; 6. \emph{Confessio Anglicana,} 1562; 7. \emph{Confessio Scotica} of 1560, and the second of 1580; 8. \emph{Confessio Ecclesiarum Belgicarum,} 1559; 9. \emph{Confessio Czengerina,} the Hungarian Confession, 1570; 10. \emph{Confessio Polonica,} or \emph{Consensus Poloniæ} (\emph{Sendomirensis}) 1570; 11. Confessio Argentinensis S. Tetrapolitana, 1531; 12. \emph{Confessio Angustana,} from the Wittenberg edition of 1540; 13. \emph{Confessio Saxonica, s. Misnica,} 1551; 14. \emph{Confessio Wirtembergica,} 1552; 15. \emph{Confessio Illustrissimi Electoris Palatini, Friderici} \emph{III.,} 1576; 16. \emph{Confessio Bohemica} (the first of the two Bohemian Confessions, which was presented to King Ferdinand in 1535. It contains a Preface by Luther. The second was compiled 1575); 17. \emph{Consensus Ecclesiarum Majoris el Minoris Poloniæ, Lithuaniæ,} etc., 1583. Appended: \emph{Acta et Conclusiones Synodi Generalis Thoruniensis}; 18. \emph{Articuli Confessionis Basileensis} of the year 1647; 19. \emph{Canones Synodi Dordrechtanæ,} 1619; 20. \emph{Confessio Cyrilli Patriarchæ Constantinop.,} 1631; 21. \emph{Catholicus Consensus,} viz., A Harmony of Christian Doctrine, compiled from the Scriptures and the writings of the Fathers, under the following heads: (\emph{a}) On the Word of God as the Rule of Faith; (\emph{b}) On God, the Trinitarian and Christological Doctrines; (\emph{c}) On Divine Providence; (\emph{d}) On the Head of the Church; (\emph{e}) On Justification; (\emph{f}) On Free Will, Original Sin, Election and Predestination; (\emph{g}) On the Sacraments; (\emph{h}) On Idolatry, the Worship of Images, etc.; (\emph{i}) On the True Way of Worshiping and Serving God; (\emph{k}) On the Church and the Ministry; (\emph{l}) Resurrection and the Future State.

Confessiones Fidei Ecclesiarum Reformatarum. Græce et Lat. \emph{Ecclesiarum Belgicarum Confessio, interpr.} Jac. Revio, \emph{et Catechesis interpr.} F. Sylburgio. Lugd. Bat. Elzev. 1635, 12mo; Amstel. 1638, 12mo. Ultrajecti, 1660, and often. (This little volume contains a Greek translation of the Belgic Confession by Revius, and a Greek translation of the Heidelberg Catechism by Sylburg, both with the Latin text in the second Column, for the use of schools in Holland.)

A Collection of Confessions of Faith, Catechisms, Directories, books of Discipline, \emph{etc., of Publick Authority in the Church of Scotland. Together with all the Acts of the Assembly which are Standing Rules concerning the Doctrine, Worship, Government, and Discipline of the Church of Scotland.} [By William Dunlop.] Edinburgh, 1719, 1722, in 2 vols. (A third volume was promised, but never appeared, as far as I know.) This rare and valuable collection contains, in the first volume, the Westminster Standards; in the second volume, the Confession of Faith of the English Congregation at Geneva, the Scotch Confession of 1560, the Scotch Confession of 1580, the National Covenant of 1638, Calvin's Catechism, the Heidelberg, and some other Catechisms and Books of Discipline. The first volume has also a long Preface (153 pp.) on the Purpose and Use of Creeds.

Sylloge Confessionum  \emph{sub tempus Reformandæ Ecclesiæ editarum.} Oxon. 1804. Ed. altera et auctior (under the revision of Bishop Lloyd). Oxon. 1827. No editor mentioned. This Collection (suggested by Bishop Cleaver) is very elegantly printed in the Clarendon Press, but has no critical value, and is incomplete. It contains: The Profession of the Tridentine Faith, the Second Helvetic Confession, the Basle Confession (1532), the Altered Augsburg Confession of 1540 (to which, in the second edition only, was added the Augustana of 1530), the Saxon Confession, the Belgic Confession, the Heidelberg Catechism, and the Canons of the Synod of Dort, all in Latin, and without a translation or introduction.

Corpus Librorum Symbolicorum  \emph{qui in Ecclesia Reformatorum auctoritatem publicam obtinuerunt,} Ed. J. Chr. G. Augusti. Elberfeldi, 1827, 8vo. Contains three Helvetic, the Gallic, the Anglican, the Scotch, the Belgic, the Hungarian, Polish, and Bohemian Confessions, the Canons of Dort, the Consensus Helveticus, and the Geneva and Heidelberg Catechisms, with an historical and literary dissertation.

Die Symbolischen Bücher der evangelisch-reformierten Kirche. \emph{Zum ersten Male aus dem Lateinischen vollständig übersetzt und mit histor. Einleitungen und Anmerkungen begleitet. . . . Für Freunde der Union und für alle, die über Entstehung, Inhalt und Zweck der Bekenntniss-Schriften sich zu belehren wünschen.} (By Friedrich Adolph Beck.) 2 Theile. Neustadt a. d. Orla, 1830; 2te wohlfeile Ausg. 1845. A good edition, with brief introductions and notes. The Augsburg Confession and the Creed of Pius IV. are appended to the Second Vol., pp. 350–410.

Sammlung Symbolischer Bücher der evang.-reformirten Kirche  \emph{für Presbyterien, Schullehrer, Confirmanden, und alle welche eine Union auf dem Grunde der heilsamen Lehre und in der Einheit der alten wahren Kirche Christi wünschen. Herausgeg. von} J. J. Mess. 3 Theile. Neuwied, 1828, 1830, and 1846, 8vo.

H. A. Niemeyer: Collectio Confessionum in Ecclesiis Reformatis  \emph{publicatarum}. Lips. 1840 (851 pages large octavo, with 88 pages of Introductory Preface), and \emph{Collectionis Confessionum Appendix, qua continentur Puritanorum Libri Symbolici.} Lipsiæ, 1840 (pp. 113). This is the most complete Latin collection of Reformed Symbols, and contains thirty-one in all, including the Zwinglian and early Swiss Confessions. It is, however, poorly edited, without an index and table of contents. Niemeyer had completed the large volume before he had seen a single copy of the Westminster Standards, and he published them nine months afterwards in an Appendix.

Die Bekenntniss-schriften der evangelisch-reformirten Kirche.  \emph{Mit Einleitungen und Anmerkungen, herausgegeben von } E. G. Adolf Böckel (Oberhofprediger and General Superintendent in Oldenburg). Leipzig, 1847 (884 large octavo pages). The best German collection, containing thirty-two Reformed Symbols, including the Anglican Catechism and the Arminian Confessions, which Niemeyer omits.

Die Bekenntniss-schriften der Reformirten Kirchen Deutschlands. \emph{Herausgegeben von} Dr. Heinrich Heppe. Elberfeld, 1860 (310 pp.). Contains the \emph{Confession} of Elector Frederick III. of the Palatinate (1577), the \emph{Repetitio Anhaltina} (1581), \emph{Anfrichtige Rechenschaft von Lehr und Ceremonien} (1593), \emph{Consensus Ministerii Bremensis Ecclesiæ} (1595), the \emph{Confession of the General Synod held at Cassel} (1608), a \emph{Report on the Faith of the Reformed Churches in Germany} (1607), the \emph{Confession of John Sigismund} of Brandenburg (1614), another \emph{Confession} of the same (1615), and the \emph{Emden Catechism} (1554), all in German.

J. Rawson Lumby (Cambridge): \emph{The Confessions of the Sixteenth Century, with Special Reference to the Articles of the Church of England} (in preparation; to be published in Cambridge and London, 1875).

II. Historical and Doctrinal Works Bearing on the Reformed Confessions.

1. The doctrinal works of Zwingli, Calvin, Beza, Œcolampadius, Bullinger, Ursinus, Olevianus, Knox, Cranmer, Ridley, Latimer, Hooper, Grindal, Jewell, Hooker, and other Reformers and standard divines of the sixteenth century.

2. \emph{Leben und ausgewählte Schriften der Väter und Begründer der reformirten Kirche.} Biographies of Zwingli, Calvin, Œcolampadius, and the other Reformers, by Baum, Christoffel, Hagenbach, Heppe, Pestalozzi, Schmidt, Stähelin, Sudhoff, etc. Elberfeld, 1857–1862. Ten Parts. One volume of this series—Christoffel's \emph{Life of Zwingli}—is translated into English, but without the extracts from his writings.

3. Older Controversial Works of Reformed Divines:

J. Hoornbeek:  \emph{Summa controversiarum religionis cum infidelibus, hæreticis, schismaticis.} Utrecht, 1658. 1676, 1689; Francf. a. O. 1697, 8vo.

Fr. Turretin:  \emph{Inst. theologiæ elenchticæ.} Geneva, 1682, 1688, 3 vols. 4to; Utrecht, 1701, 4 vols. 4to, etc.

B. Pictet:  \emph{De consensu et dissensu inter Reformatos et Augustanæ Confessionis fratres.} Genev. 1700.

F. Spanheim:  \emph{Controversiarum de religione cum dissidentibus elenchus hist. theol.} Leyd. 1687; fifth edition, Leyd. 1757, 4to.

Du Gerdes:  \emph{Elenchus veritatum, circa quas defendendas versatur theol. elenchthica.} Gröningen, 1740, 4to.

J. F. Stapfer:  \emph{Institutiones theologicæ polem.} Zurich, 1743–47, 5 vols. 8vo.

Du Wyttenbach:  \emph{Theol. elenchticæ initia.} Francf. a. M. 1763, 1765, 2 vols. 8vo.

Comp. also the list of older dogmatic works of the Reformed Church in Heppe's  \emph{Dogmatik der evang.-reform. Kirche,} at the end of Preface, and in Schweizer's  \emph{Glaubenslehre der evang.-reform. Kirche,} Vol. I. pp. xxi.-xxiii.

4. Recent Historico-Dogmatic Works:

H. Heppe (Marburg): \emph{Dogmatik der evang.-reform. Kirche dargestellt und aus den Quellen belegt,} Elberfeld, 1861; and his \emph{Dogmatik des Deutschen Protestantismus im} 16\emph{ten Jahrh.} Gotha, 1857, 3 vols.

Alex. Schweizer (Zurich): \emph{Die Protestantischen Centraldogmen in ihrer Entwicklung innerhalb der Reformirten Kirche.} Zurich, 1854–56, 2 vols. Also his \emph{Glaubenslehre der evang.-reform. Kirche dargestellt und aus den Quellen belegt.} Zurich, 1844–47, 2 vols.

Aug. Ebrard (Erlangen): \emph{Das Dogma vom heil. Abendmahl und seine Geschichte} (Frankfurt a. M. 1846), the second vol.; and also his \emph{Christliche Dogmatik.} Königsberg, 1851, 1852, 2 vols.

Charles Hodge (Princeton): \emph{Systematic Theology.} New York, 1873, 3 vols.

J. J. van Oosterzee (Utrecht): \emph{Christian Dogmatics.} Translated from the Dutch by \emph{Watson} and \emph{Evans.} London and New York, 1874, 2 vols.

The Reformed Confessions are much more numerous than the Lutheran, because they represent a larger territory and several nationalities—Swiss, German, French, Dutch, English, and Scotch—each of which produced its own doctrinal and disciplinary standards, since the geographical and political divisions and the close relations to the civil government determined also the number of ecclesiastical organizations. The productive period of the Reformed movement, moreover, extended far into the seventeenth century, especially in England, and some of the most important confessions, as the Canons of Dort and the Westminster Standards, were made long after the symbolic development of the Lutheran Church had reached its culmination and rest in the Formula of Concord. Finally the Reformed Church departs further from the authority of ecclesiastical traditionalism than the Lutheran, and allows more freedom for the development of various types of doctrine and schools of theology within the limits of the Word of God, to which it more rigidly adheres.

But with all this variety, the Reformed symbols are as much agreed in the essential articles of faith as the Lutheran, and differ even less than the Augsburg Confession, as explained by its author and his school, differs from the Formula of Concord.\footnote{This doctrinal consensus of the Reformed Creeds has been shown as early as 1581 in the \emph{Harmonia Confessionum} above quoted.} They exhibit substantially the same system of doctrine, and are only variations of one theme according to the wants of the national Churches for which they were intended. The Reformed Churches were never organically united under one form of government, and even every little canton in Switzerland (as every Lutheran principality in Germany) has its own ecclesiastical establishment;\footnote{In this respect the Churches of the United States, being free from government control, are much better organized, according to creeds, without allowing the State boundaries to interfere with their organic unity.} but they recognized each other as branches of the same family, and kept up a lively intercommunion. Even the leading divines and dignitaries of the Episcopal Church of England, during the sixteenth century, freely corresponded with the Reformed Churches of Switzerland, France, and Holland, and the difference in church polity was no bar to church fellowship.

There are in all over thirty Reformed creeds. But many of them had never more than local authority, or were superseded by later and maturer forms. None of them has the same commanding position as the Augsburg Confession in the Lutheran Church. Those which have been most widely accepted and are still most in use are the Heidelberg or Palatinate Catechism, the Thirty-nine Articles, and the Westminster Confession. The second Helvetic Confession and the Canons of Dort are equal to them in authority and theological importance, but less adapted for popular use. All the rest have now little more than historical significance.

As to origin and theological character, the Reformed Confessions may be divided into Zwinglian and Calvinistic. The earlier were the product of Zwingli and his Swiss coadjutors, the later date from Calvin or his pupils and successors, and exhibit a more advanced and matured state of doctrine, with a difference, however, as to the extent to which they are committed to the Calvinistic system; some accepting it in full, while others maintain a reserve in regard to its angular points and rigorous logical consequences.

As to the country in which they originated and for which they were chiefly intended, we may divide them into Swiss, German, French, Dutch, English, and Scotch Confessions.

To the Swiss family belong the Confessions which proceeded from the Churches of Zurich, Basle, Berne, and Geneva, partly of Zwinglian and partly of Calvinistic origin.

The German family embraces the Tetrapolitan Confession, the Heidelberg Catechism, the Brandenburg and Anhalt Confessions, and a few others. They are less pronounced in their Calvinism, and mediate between it and the Lutheran Creed.

To France and the Netherlands belong the French and the Belgic Confessions, the Canons of the Synod of Dort, and also the Arminian Articles, which differ from the Calvinistic creeds in five points.

The English family embraces the Thirty-nine Articles, the old Scotch Confessions, and the later Westminster Standards.

Besides, there are Bohemian, Polish, and Hungarian Confessions of lesser importance.

Note.—We take the term \emph{Reformed} here in its catholic and historical sense for all those Churches which were founded by Zwingli and Calvin and their fellow-reformers in the sixteenth century on the Continent, and in England and Scotland, and which agreed with the Lutheran Church in opposition to the Roman Catholic, but differed from it in the doctrine of the real presence, afterward also in the doctrine of predestination. By their opponents they were first called in derision \emph{Zwinglians} and \emph{Calvinists,} also \emph{Sacramentarians} or \emph{Sacramentschwärmer} (by Luther and in the Formula of Concord), and in France \emph{Huguenots.} But they justly repudiated all such sectarian names, and used instead the designations \emph{Christian} or \emph{Evangelical} or \emph{Reformed,} or \emph{Evangelical Reformed} or \emph{Reformed Catholic.} The term \emph{Reformed} assumed the ascendency in Switzerland, France, and elsewhere. Beza, e.g., uses it constantly. Queen Elizabeth, in sundry letters to the Protestant courts of Germany in 1577, speaks throughout of \emph{ecclesiæ reformatæ,} and once calls the non-Lutheran Churches \emph{ecclesiæ reformatiores, more Reformed,} implying that the Lutheran is Reformed also.

The Lutherans, before the last quarter of the sixteenth century, called themselves likewise \emph{Christian} and \emph{Evangelical,} sometimes \emph{Reformed,} and since 1530 the \emph{Church} or \emph{Churches of the Augsburg Confession,} or \emph{Verwandte der Augsburgischen Confession.} For a long time they disowned the terms \emph{Lutheranus, Luthericus, Lutheranismus,} which were first used by Dr. Eck, Cochlæus, Erasmus, and other Romanists with the view to stigmatize their religion as a recent innovation and human invention. (A Papist once asked a Lutheran, 'Where was your Church before Luther?' The Lutheran answered by asking another question, 'Where was your face this morning before it was washed?') Erasmus speaks of \emph{Lutherana tragædia, negotium Lutheranum, factio Lutherana.} Hence the Lutheran symbols never use the term \emph{Lutheran,} except once, and then by way of complaint that the 'dear, holy Gospel should be called \emph{Lutheran.}'\footnote{\emph{Apology of the Augsburg Confession,} Art. XV. (VIII. p. 213 ed. Müller): '\emph{Das liebe, heilige Evangelium nennen sie} [the Papists] \emph{Lutherisch.}' The name of Luther, however, is often honorably mentioned, especially in the Formula of Concord.} Luther himself complained of this use of his name; nevertheless he had no objection that it should be duly honored in connection with the Word of God, and thought that his followers need not be ashamed of him.\footnote{'\emph{Wahr ist's,}' he says (\emph{Works,} Erl. ed. Vol. XXVIII. p. 316), '\emph{dass du bei Leib und Seele nicht sol1st sagen: ich bin } Lutherisch  \emph{oder} Päpstisch;  \emph{denn derselben ist keiner für dich gestorben, noch dein Meister, sondern allein Christus, und sollst dich} (\emph{als}) Christen  \emph{bekennen. Aber wenn du es dafür hältst, dass des Luthers Lehre evangelisch und des Papstes unevangelisch sei, so musst du den Luther nicht so gar hinwerfen. Du wirfst sonst seine Lehre auch mit hin, die du doch für Christi Lehre erkennest; sondern also musst du sagen: der Luther sei ein Bube oder heilig, da liegt mir nichts an; seine Lehre aber ist nicht sein, sondern Christi selbst.}' And in another place (Vol. XL. p. 127): '\emph{Und wiewohl ich’s nicht gern habe, dass man die Lehre und Leute } Lutherisch  \emph{nennt, und muss von ihnen leiden, dass sie Gottes Wort mit meinem Namen also schänden, so sollen sie doch den Luther, die Lutherischen Lehre und Leute lassen bleiben und zu Ehren kommen.}'} They thought so, too; and, forgetting St. Paul's warning against sectarian names, they gradually themselves appropriated the term \emph{Lutheran,} or \emph{Evangelical Lutheran,} as the official title of their Church, since about 1585, under the influence of Jacob Andreæ, the chief author of the Formula of Concord, and Ægidius Hunnius, and in connection with the faith in Luther as a special messenger of God for the restoration of Christianity in its doctrinal purity. See the proof in the little book of Dr. Heinrich Heppe, \emph{Ursprung und Geschichte der Bezeichnungen} '\emph{reformirte}' \emph{und} '\emph{lutherische}' \emph{Kirche,} Gotha, 1859, pp. 28, 35, 55.

The negative term \emph{Protestant }was used after 1529 for both Confessions by friend and foe, and is so used to this day; but it must be explained from the historical occasion which gave rise to it, and be connected with the \emph{positive} faith in the Word of God, on the ground of which the evangelical members of the Diet of Spires protested against the decision of the papal majority, as an encroachment on the rights of conscience and an enforcement of the traditions of men.

On the Continent of Europe it is still customary to divide orthodox Christendom into three Confessions or Creeds—the Catholic (Greek and Roman), the Lutheran, and the Reformed—and to embrace under the Reformed all other Protestant bodies, such as Methodists and Baptists, or to speak of them as mere sects. But this will not do in England and America, where these sects, so called, have become powerful Churches. \emph{Reformed} is sometimes used among us in a more general sense of all Protestant Churches, sometimes in a restricted sense of a particular branch of the Reformed Church. The Continental terminology suits the ecclesiastical statistics of the sixteenth century, but must be considerably enlarged and modified in view of the greater number of Anglo-American Churches. We shall devote a separate chapter to those Protestant evangelical bodies which have taken their rise since the Reformation.

\xsection{I. Reformed Confessions of Switzerland.}

\xsubsection{Zwinglian Confessions.}

§ 51. Zwinglian Confessions.

Literature.

H. Zwinglii  \emph{Opera} ed. \emph{Gualther }(Zwingli's son-in-law), Tig. 1545 and 1581, 4 Tom.; ed. \emph{M. Schuler u. J. Schulthess,} Tig. 1828–42, 8 Tom. The last and only complete edition contains the German and Latin works, with a supplemental volume of tracts and letters, published 1861. A judicious selection from his writings, in German, for popular use, was edited by \emph{Christoffel,} Zurich, 1843–46, in fifteen small volumes, also in the second part of his biography of Zwingli.

Biographies of Zwingli by Myconius, Nüscheler, Hess, Rotermund, Schuler, Hottinger, Röder, Tichler, Christoffel (Elberfeld, 1857), and especially Mörikofer:  \emph{Ulrich Zwingli nach den urkundlichen Quellen,} Leipzig, 1867–69, 2 vols. Hottinger and Christoffel are translated into English, but the latter without the valuable extracts from Zwingli's writings. Güder's art. on Zwingli, in Herzog's \emph{Encykl.} Vol. XVIII. pp. 701–766, is a condensed biography. Robbins, \emph{Life of Zwingli,} in \emph{Bibliotheca Sacra,} 1851.

Also A. Ebrard:  \emph{Das Dogma vom heil. Abendmahl und seine Geschichte} (Francf. 1846), Vol. II. pp. 1–112 (an able vindication of Zwingli against misrepresentations). Ed. Zeller:  \emph{Das theologische System Zwingli's,} Tüb. 1853. Ch. Sigwart:  \emph{Ulrich Zwingli, der Charakter seiner Theologie, mit besonderer Rücksicht auf Picus von Mirandula,} Stuttg. 1855. H. Spörri:  \emph{Zwinglistudien,} Leipz. 1866. Merle d’Aubigné:  \emph{History of the Reformation,} 4th vol. (French, English, and German). Hagenbach:  \emph{Geschichte der Reform.,} 4th ed. Leipz. 1870, pp. 183 sqq. G. P. Fisher:  \emph{The Reformation,} New York, 1873, pp. 137 sqq.

Zwingli (1484–1531) represents the first stage of the Reformed Church in Switzerland. He began what Calvin and others completed. He died in the prime of life, a patriot and martyr, on the battle-field, when his work seemed to be but half done. His importance is historical rather than doctrinal. He was the most clear-headed and liberal among the reformers, but lacked the genius, depth, and vigor of Luther and Calvin. He held opinions on the sacraments, original sin (as a disorder rather than a state of guilt), and on the salvation of all infants (unbaptized as well as baptized) and the nobler heathen, which then appeared radical, dangerous, and profane. He could conceive of a broad and free Christian union, consistent with doctrinal differences and denominational distinctions. He was a patriotic republican, frank, honorable, incorruptible, cheerful, courteous, and affable. He took an active part in all the public affairs of Switzerland, and labored to free it from foreign influence, misgovernment and immorality. He began at Einsiedeln (1516), and more effectively at Zurich (since 1519), to preach Christ from the pure fountain of the New Testament, and to set him forth as the only Mediator and all-sufficient Saviour. Then followed his attacks upon the corruptions of Rome, and the Reformation was introduced step by step in Zurich, where he exercised a controlling influence, and in the greater part of German Switzerland, until its progress was suddenly checked by the catastrophe at Cappel, 1531.

Zwingli was scarcely two months younger than Luther, who survived him fifteen years. Both were educated and ordained in the Roman Church, and became innocently and providentially reformers of that Church. Both were men of strong mind, heroic character, fervent piety, and commanding influence over the people. Both were good scholars, great divines, and fond of poetry and music.\footnote{See Zwingli's poems, written during the pestilence, in Hagenbach, 1.c. p. 216, and another, p. 404. He published a moral poem, under the title \emph{The Labyrinth,} as early as 1510, while priest at Glarus (\emph{Opera,} Tom. II. B. pp. 243 sqq.; Mörikofer, Vol. I. pp. 13 sqq.). His preference for Puritanic simplicity in public worship gave rise to the fiction of his hostility to music. He was, on the contrary, singularly skilled in that art, and was called in derision by the Papists 'the evangelical lute-player.' A contemporary says that he never knew a man who could play on so many musical instruments—the lute, the harp, the violin, etc. \par{} [Zwingli's copy of the N. T. was confined to Paul's epistles and \emph{Hebrews}.—Ed.]} Both labored independently for the same great cause of evangelical Protestantism—the one on a smaller, the other on a larger field. But their endowment, training, and conversion were different. Zwingli had less prejudice, more practical common-sense, clear discrimination, sober judgment, self-control, courtesy, and polish—Luther more productive genius, poetic imagination, overpowering eloquence, mystic; depth, fire, and passion; and was in every way a richer and stronger, though rougher and wilder nature. Zwingli's eyes were opened by the reading of the Greek Testament, which he carefully copied with his own hand, and the humanistic learning of his friend Erasmus; while Luther passed through the ascetic struggles of monastic life, till he found peace of conscience in the doctrine of justification by faith alone. Zwingli broke more rapidly and more radically with the Roman Church than Luther. He boldly abolished all doctrines and usages not taught in the Scriptures; Luther piously retained what was not clearly forbidden. He aimed at a reformation of government and discipline as well as theology; Luther confined himself to such changes as were directly connected with doctrine. He was a Swiss and a republican; Luther, a German and a monarchist. He was a statesman as well as a theologian; Luther kept aloof from all political complications, and preached the doctrine of passive obedience to established authority. They met but once in this world, and then as antagonists, at Marburg, two years before Zwingli's death. They could not but respect each other personally, though Luther approached the Swiss with the strongest prejudice, looking upon him as a fanatic and semi-infidel.\footnote{Once, at least, Luther speaks kindly of Zwingli, in a letter to Bullinger, of Zurich, May 14, 1588 (De Wette, Vol. V. p. 112): '\emph{Libere enim dicam: Zwinglium, postquam Marpurgi mihi visus et auditus est, virum optimum esse judicavi, sicut et Œcolampadium.}' In the same letter he says that Zwingli's death caused him much pain. But this personal respect did not prevent him from using the most violent language against his doctrine of the Lord's Supper, which he held in utter abhorrence to the last, and this all the more because his fanatical colleague Carlstadt, who gave him infinite trouble, had first proposed and defended it by an untenable exegesis. This accounts also for his absurd charge of fanaticism against the clear, sober-minded, jejune Zwingli. '\emph{Es ist fast lächerlich,}' says the mild and impartial Hagenbach (p. 280), '\emph{wenn Luther mitten in seiner schwärmerisch tobenden Leidenschaft den ehrlichen Zwingli einen Schärmer nennt, ihn, der von aller Schwärmerei so fern war. Es sei denn, dass man den idealistischen Zug in ihm }(\emph{und der war allerdings dem derben Realismus Luthers zuwider}) \emph{mit diesem Namen bezeichnen wolle. Man betrachte auch nur sein Bildniss! Dieser energische, feste, satte Kopf, diese in Stein gehauene, markante Physiognomie, diese breite Stirn, dieses volle klare Auge, diesen geschlossenen Mund mit runden Lippen—genug! ich überlasse einem Lavater die vollendete Deutung des Bildes} (\emph{der in ihm} ``\emph{Ernst, Nachdenken, männliche Entschlossenheit, eine sich zusammenziehende Thatkraft, einen schauenden, durchdringenden Verstand}'' \emph{erkennt}), \emph{und berufe mich allein auf die Geschichte, welche den lebendigen Commentar zu diesem Bildniss ausmacht.}'} They came to an agreement on every article of faith except the real presence in the eucharist. Zwingli proposed, with tears, peace and union, notwithstanding this difference, but Luther refused the hand of Christian fellowship, because he made doctrinal agreement the boundary-line of brotherhood.\footnote{\par On the relation of Luther and Zwingli, see Ebrard, Vol. II. pp. 214 sqq.; Hagenbach, pp. 278 sqq.; and an essay of Hundeshagen in the \emph{Studien und Kritiken} for 1862. Zwingli himself thus described his relation to Luther in 1523, when the German Papists began to denounce his doctrine as a Lutheran heresy: '\emph{Ich habe, ehe noch ein Mensch in unserer Gegend etwas von Luther's Namen gewusst hat, angefangen das Evangelium Christi zu predigen, im Jahr} 1516. \emph{Wer schalt mich damals lutherisch? . . . Luther's Name ist mir noch zwei Jahre unbekannt gewesen, nachdem ich mich allein an die Bibel gehalten habe. Aber es ist, wie gesagt, nur ihre Schlauheit, dass die Päpstler mich und Andere mit solchem Namen beladen. Sprechen sie: Du musst wohl lutherisch sein, du predigest ja, wie Luther schreibt; so ist meine Antwort: Ich predige ja auch wie Paulus; warum nennst du mich nicht vielmehr einen Paulisten? . . . Meines Erachtens ist Luther ein trefflicher Streiter Gottes, der da mit so grossem Ernste die Schrift durchforscht, als seit tausend Jahren irgend einer auf Erden gewesen ist. Mit dem männlichen, unbewegten Gemüthe, womit er den Papst von Rom angegriffen hat, ist ihm keiner nie gleich geworden, so lange das Papstthum gewähret hat, alle Andern ungescholten. Wessen aber ist solche That? Gottes oder Luthers? Frage den Luther selbst, gewiss sagt er dir: Gottes. Warum schreibst du denn anderer Menschen Lehre dem Luther zu, da er sie selbst Gott zuschreibt, und nichts Neues hervorbringt, sondern was in dem ewigen, unveränderlichen Worte Gottes enthalten ist? Fromme Christen! gebet nicht zu, dass der ehrliche Name Christi verwandelt werde in den Namen Luthers; denn Luther ist für uns nicht gestorben, sondern er lehrt uns den erkennen, von dem wir allein alles Heil haben. Predigt Luther Christum, so thut er's grade wie ich; wiewohl, Gott sei Dank! durch ihn eine unzählbare Menge mehr als durch mich und Andere, denen Gott ihr Mass grösser oder kleiner macht, zu Gott geführt wird. Ich will keinen Namen tragen, als meines Hauptmannes Jesu Christi, dessen Streiter ich bin. . . . Es kann kein Mensch sein, der Luther höher achtet, als ich. Dennoch bezeuge ich vor Gott und allen Menschen, dass ich all’ meine Tage nie einen Buchstaben an ihn geschrieben habe, noch er an mich, noch verschafft, dass geschrieben werde. Ich habe es unterlassen, nicht dass ich jemand desswegen gefürchtet, sondern weil ich damit allen Menschen habe zeigen wollen, wie gleichförmig der Geist Gottes sei, da wir so welt von einander entfernt und doch einmüthig sind, aber ohne alle Verabredung, wiewohl ich ihm nicht zuzuzählen bin; denn jeder thut, soviel ihm Gott weiset.}'}

Zwingli wrote four dogmatic works of a semi-symbolic character, which are closely interwoven with the history of the Reformation in German Switzerland, and present a clear exhibition of the Reformed faith in the first stage of its development. These are the Sixty-seven Articles of Zurich (A.D. 1523), the Ten Theses of Berne (1528), the Confession of Faith to the German Emperor Charles V. (1530), and the Exposition of the Christian Faith to King Francis I. of France (1531).\footnote{They are all embodied in the Collections of Niemeyer and Böckel. Niemeyer (\emph{Collectio,} pp. 3–77) gives the first two in Swiss-German and in Latin, the last two in Latin only. Böckel \{\emph{Bekenntniss-Schriften,} pp. 5–107) gives them in High-German, and adds the 'Brief Christian Instruction' which Zwingli wrote in the name of the Magistrate of Zurich, Sept. 1523, for the preachers and pastors, treating of the Gospel and the Law, of Images, and of the Mass (pp. 13–34).}

1. The Sixty-seven Articles, or Conclusions.\footnote{\par Articuli sive Conclusiones LXVII. H. Zwinglii, a. 1523. They were published by Zwingli himself before the disputation, with the title: 'The following 67 Articles and opinions I, Ulrich Zwingli, confess to have preached in the honorable city of Zurich, on the ground of the Scripture which is called theopneustos [i.e. inspired by God], and I offer to defend them. And should I not correctly understand the said Scripture, I am ready to be instructed and corrected, but only by the Scripture.' On the different editions, see the notices of Niemeyer, \emph{Præfatio,} pp. xvi sqq.}

They were prepared for a public disputation held January 29, 1523, in the city of Zurich, where Zwingli was chief pastor from 1519, and were victoriously defended by him, in the presence of the civil magistrate and about six hundred persons, against Dr. Faber, the General Vicar of the Bishop of Constance, who appeared to superintend the meeting rather than to defend the old doctrines, and was unwilling or unable to answer the arguments of a learned and powerful opponent. The magistrate passed a resolution on the same day approving of Zwingli's position, and requiring all the ministers of the canton to preach nothing but what they could prove from the holy gospel. A second disputation followed in October, on the use of images and the mass, before about nine hundred persons, including three hundred priests and delegates from different cantons; a third disputation took place in January, 1524. The result was the emancipation from popery, and the orderly and permanent establishment of the Reformed Church in the city and canton of Zurich.

These Articles resemble the Ninety-five Theses of Luther, which opened the drama of the Reformation in Germany, October 31, 1517, but they mark a considerable advance in Protestant conviction. They are full of Christ, as the only Saviour and Mediator, and clearly recognize the Word of God as the only rule of faith. They attack the primacy of the Pope, the mass, the invocation of saints, the meritoriousness of human works, fasts, pilgrimages, celibacy, and purgatory, as unscriptural traditions of men. They are short, and, in this respect, like the Thirty-nine Articles of the Church of England, better adapted for a creed than the lengthy confessions of that age. But they never had more than local authority. We give a few specimens:

1. All who say that the gospel is nothing without the approbation of the Church, err and cast reproach upon God.

2. The sum of the gospel is that our Lord Jesus Christ, the true Son of God, has made known to us the will of his heavenly Father, and redeemed us by his innocence from eternal death, and reconciled us to God.

3. Therefore Christ is the only way to salvation for all who were, who are, and who shall be.

4. Whosoever seeks or shows another door, errs—yea, is a murderer of souls and a robber.

7. Christ is the Head of all believers.

8. All who live in this Head are his members and children of God. And this is the true Catholic Church, the communion of saints.

15. Who believes the gospel shall be saved; who believeth not shall be damned. For in the gospel the whole truth is clearly contained.

16. From the gospel we learn that the doctrines and traditions of men are of no use to salvation.

17. Christ is the one eternal high-priest.

18. Christ, who offered himself once on the cross, is the sufficient and perpetual sacrifice for the sins of all believers. Therefore the mass is no sacrifice, but a commemoration of the one sacrifice of the cross and a seal of the redemption through Christ.

19. Christ is the only Mediator between God and us.

22. Christ is our righteousness. From this it follows that our works are good so far as they are Christ's, but not good so far as they are our own.

24. Christians are not bound to any works which Christ has not commanded.

26. Nothing is more displeasing to God than hypocrisy.

27. All Christians are brethren.

34. The power of the Pope and the Bishops has no foundation in the Holy Scriptures and the doctrine of Christ.

49. I know of no greater scandal than the prohibition of lawful marriage to priests, while they are permitted for money to have concubines. Shame! (\emph{Pfui der Schande!})

50. God alone forgives sins, through Jesus Christ our Lord alone.

57. The Holy Scripture knows nothing of a purgatory after this life.

2. The Ten Theses of Berne.

After the Conference between the Reformed and the Roman divines (headed by Dr. Eck), held at Baden, in Aargau, May, 1526, which formed a turning-point in the history of the Swiss Reformation (more decided than the similar disputation between Luther and Eck in Leipzig, 1519), the Reformation triumphed in Berne, the most conservative and aristocratic as well as most influential canton of the confederacy. Three ministers, Berthold Haller, Francis Kolb, and Sebastian Meyer, friends of Zwingli, and a gifted layman, Nicolas Manuel, who was a statesman, poet, and painter, had previously prepared the way under great opposition. The magistrate convened a convocation of the clergy and laity, which continued nineteen days, from January 6 to 26, 1528, discussing ten theses which Zwingli had revised and published at the request of Haller. Delegates appeared from other cantons (except the Roman Catholic), and the South German cities of Constance, Ulm, Lindau, and Strasburg. The Bishops of Constance, Basle, Lausanne, and Sion were also invited, but declined to attend, except the Bishop of Lausanne, who sent a few doctors. Dr. Eck, who had figured as the champion of Romanism in Baden (as well as previously at Leipzig), prudently disdained at this time to follow 'the heretics into their corners and dens.' The principal champions of the Reformed cause were Zwingli (who also preached two very effective sermons on the Apostles' Creed, and against the mass), Œcolampadius, Haller, Kolb, Pellican, Megander, Bucer, and Capito. They carried a complete victory, and hereafter Berne, Zurich, and Basle—the three most enlightened and influential German cantons—were closely linked together in the Reformed faith.\footnote{See Samuel Fischer, \emph{Geschichte der Disputation zu Bern,} Berne, 1828; Melch. Kirchhofer, \emph{Berthold Haller, oder die Reformation in Bern,} Zurich, 1828; C. Pestalozzi, \emph{B. Haller, nach handschriftlichen und gleichzeitigen Quellen,} Elberfeld, 1861, pp. 35 sqq. (in Vol. IX. of the \emph{Lives and Writings of the Fathers and Founders of the Reformed Church}); Zwingli's \emph{Werke,} ed. Schuler and Schulthess, Vol. II. I. pp. 630 sqq. Luther was not well pleased with this triumph of Zwinglianism, and wrote to Gabriel Zwilling, March 7 (De Wette, Vol. III. No. 959): '\emph{Bernæ in Helvetiis finita disputatio est; nihil factum, nisi quod missa abrogata et pueri in plateis cantent, se esse a Deo pisto liberatos.}' He also prophesied an evil end to Zwingli.}

The Bernese Theses are as follows:

1. The holy Christian Church, whose only Head is Christ, is born of the Word of God, and abides in the same, and listens not to the voice of a stranger.

2. The Church of Christ makes no laws and commandments without the Word of God. Hence human traditions are no more binding on us than they are founded in the Word of God.

3. Christ is the only wisdom, righteousness, redemption, and satisfaction for the sins of the whole world. Hence it is a denial of Christ when we confess another ground of salvation and satisfaction.

4. The essential and corporeal presence of the body and blood of Christ can not be demonstrated from the Holy Scripture.

5. The mass as now in use, in which Christ is offered to God the Father for the sins of the living and the dead, is contrary to the Scripture, a blasphemy against the most holy sacrifice, passion, and death of Christ, and on account of its abuses an abomination before God.

6. As Christ alone died for us, so he is also to be adored as the only Mediator and Advocate between God the Father and the believers. Therefore it is contrary to the Word of God to propose and invoke other mediators.

7. Scripture knows nothing of a purgatory after this life. Hence all masses and other offices for the dead are useless.

8. The worship of images is contrary to the Scripture. Therefore images should be abolished when they are set up as objects of adoration.

9. Matrimony is not forbidden in the Scripture to any class of men, but permitted to all.

10. Since, according to the Scripture, an open fornicator must be excommunicated, it follows that unchastity and impure celibacy are more pernicious to the clergy than to any other class.\footnote{The German copy adds: '\emph{Alles Gott und seinem heiligen Wort zu Ehren.}'}

In his farewell sermon, Zwingli thus addressed the Bernese: 'Victory has declared for the truth, but perseverance alone can complete the triumph. Christ persevered unto death. \emph{Ferendo vincitur fortuna.} Behold these idols, behold them conquered, mute, and scattered before us. The gold you have spent upon these foolish images must henceforth be devoted to the comfort of the living images of God in their poverty. In conclusion, stand fast in the liberty wherewith Christ has made us free, and be not entangled again with the yoke of bondage (Gal. v. 1). Fear not! the God who has enlightened you, will enlighten also your confederates; and Switzerland, regenerated by the Holy Ghost, shall flourish in righteousness and peace.'

3. The Confession of Faith to Emperor Charles V.\footnote{\emph{Ad Carolum Rom. Imperatorem Germaniæ comitia Augustæ celebrantem fidei Huldrychi Zwinglii ratio} (\emph{Rechenschaft}). \emph{Anno MDXXX. Mense Julio. Vincat veritas} (Zurich). In the same year a German translation appeared in Zurich, and in 1543 an English translation. See Niemeyer, p. xxvi. Comp. also Böckel, pp. 40 sqq.; Mörikofer, Vol. II. pp. 297 sqq.; and Christoffel, Vol. II. pp. 237 sqq.}

Zwingli took advantage of the meeting of the famous Diet at Augsburg, held A.D. 1530, to send a Confession of his faith addressed to the German Emperor Charles V., shortly after the Lutheran Princes had presented theirs (June 25). It is dated Zurich, July 3, and was delivered by his messenger at Augsburg on the 8th of the same month, but it shared the same fate as the 'Tetrapolitan Confession' of Bucer and Capito: it was never laid before the Diet, and was treated with undeserved contempt. Dr. Eck wrote in three days a refutation,\footnote{\emph{Repulsio Articulorum Zwinglii.} Zwingli replied in \emph{Ad illustrissimos Germaniæ principes Augustæ congregatos, de convitiis Eckii} (\emph{Opera,} Vol. IV. pp. 19 sqq.).} slanderously charging Zwingli that for ten years he had labored to root out from the people of Switzerland all faith and all religion, and to stir them, up against the magistrate; that he had caused greater devastation among them than the Turks, Tartars, and Huns; that he had turned the churches and convents founded by the Hapsburgers (the Emperor's ancestors) into temples of Venus and Bacchus; and that he now completed his crime by daring to appear before the Emperor with such an impudent piece of writing. The Lutherans (with the exception of Philip of Hesse, who sympathized with Zwingli) were scarcely less indignant, and much more anxious to conciliate the Catholics than to appear in league with Zwinglians and Anabaptists. They felt especially offended that the Swiss Reformer took strong ground against the corporeal presence, and incidentally alluded to them as persons who 'were looking back to the flesh-pots of Egypt.'\footnote{'\emph{Quod Christi corpus,}' says Zwingli, '\emph{per essentiam et realiter, hoc est corpus ipsum naturale in cœna aut adsit aut ore dentibusque nostris manducatur, quemadmodum Papistæ, et }quidam qui ad ollas Egyptiacas respectant, \emph{perhibent, id non tantum negamus, sed errorem esse qui verbo Dei adversatur, constanter asseveramus.}'} Melanchthon, who was at that time not yet emancipated from the Catholic tradition on that article, judged him insane.\footnote{See his letter to Luther of July 14, 1530, quoted on p. 263.}

Zwingli, having had no time to consult with his confederates, offered the Confession in his own name, and submitted it to the judgment of the whole Church of Christ, under the guidance of the Word of God and the Holy Spirit.

In the first sections he declares, as clearly and even more explicitly than the Lutheran Confession, his faith in the orthodox doctrines of the Trinity and the Person of Christ, as laid down in the Nicene and Athanasian Creeds (which are expressly named). He teaches the election by free grace, the sole and sufficient satisfaction of Christ, and justification by faith, in opposition to all human mediators and meritorious works. He distinguishes between the internal or invisible, and the external or visible Church; the former is the company of the elect believers and their children, and is the bride of Christ; the latter embraces all nominal Christians and their children, and is beautifully described in the parable of the Ten Virgins, of whom five were foolish. Church may also designate a single congregation, as the church in Rome, in Augsburg, in Leyden. The true Church can never err in the foundation of faith. Purgatory he rejects as an injurious fiction which sets Christ's merits at naught. On original sin, the salvation of unbaptized infants, and the sacraments, he departs much further from the traditional theology than the Lutherans. He goes into a lengthy argument against the corporeal presence in the eucharist. On the other hand, however, he protests against being confounded with the Anabaptists, and rejects their views on infant baptism, civil offices, the sleep of the soul, and universal salvation.

The document is frank and bold, yet dignified and courteous, and concludes thus: 'Hinder not, ye children of men, the spread and growth of the Word of God—ye can not forbid the grass to grow. Ye must see that this plant is richly blessed with rain from heaven. Consider not your own wishes, but the demands of the age concerning the free course of the gospel. Take these words kindly, and show by your deeds that you are children of God.'

4. The Exposition of the Christian Faith to King Francis I.\footnote{Christianæ Fidei \emph{ab}  H. Zwinglio\emph{ predicatæ brevis et clara }Expositio \emph{ab ipso Zwinglio paulo ante mortem ejus ad Regem Christianum scripta.} Under this title Bullinger edited the work, with some omissions and changes, from the author's MS., with a preface, 1536. He calls Zwingli \emph{fidelissimus evangelii præco et Christianæ libertatis assertor constantissimus.} Leo Judæ prepared a free German translation: \emph{Eine kurze, klare Summe and Erklärung des christl. Glaubens,} etc., Zurich (no date). Niemeyer took his text directly from a copy of the manuscript made by Bibliander, in the library at Zurich (pp. xxviii. and 36 sqq.). Christoffel (Vol. I. p. 368) states that the original MS. of Zwingli is still in the public library of Paris. A High-German translation in Böckel, pp. 63 sqq., and Christoffel, Vol. II. pp. 262 sqq.}

This is, as Bullinger says, the swan song of Zwingli, in which he surpassed himself. He wrote it in July, 1531, three months before his death, at the request of his friend Maigret, the French ambassador to Switzerland, and sent it in manuscript to Francis I., King of France (1515–1547), who, from political motives, showed himself favorable to the Protestants in Germany and Switzerland, while he persecuted them at home. A few years before he had dedicated to him his 'Commentary on the true and false Religion' (1525), and a few years afterwards (1536) Calvin dedicated to him his Institutes, with a most eloquent and powerful letter; but the frivolous monarch probably never read these voices of warning, which, if properly heeded, might have changed the whole history of France.

This last document of Zwingli is clear, bold, spirited, full of faith and hope. In a brief preface he warns the most Christian King of France against the lies and slanders circulated against the Protestants. He first treats of God, the ultimate ground of our faith and only object of worship. We do not despise the saints and sacraments, we only guard them against abuse; we honor Mary as the perpetual Virgin and Mother of God,\footnote{Zwingli retained this term, but with a restriction to the human nature united to the Logos.} but we do not worship her in the proper sense of the term, which we know she herself would never tolerate. The sacraments we honor as signs or symbols of holy things, but not as the holy things themselves. Then he speaks of the holy Trinity, and the incarnation of the eternal Son of God for our salvation, who made a full satisfaction for all our sins. He gives an able exposition of the two natures in the one person of Christ, his death, resurrection, ascension, and return to judgment. He rejects purgatory as a papal fiction. He dwells very fully on the doctrine of the Sacraments, especially the eucharistic presence (rejecting ubiquity). The remaining chapters are devoted to the Church, the Magistrate, the remission of sins, faith and works, eternal life, and an attack on the Anabaptists, with whom the Protestants were often confounded in France. In conclusion, he entreats the king to give the gospel free course in his kingdom; to imitate the example of some pious princes in Germany; to judge by the fruits of the Reformed faith wherever it was fairly established; and to forgive the boldness with which he approached his majesty. The urgency of the case demanded it. An appendix is devoted to the mass, with proofs from the fathers, especially from Augustine, in favor of his view on the Lord's Supper.

\xsubsection{Zwingli's Distinctive Doctrines.}

§ 52. Zwingli's Distinctive Doctrines.

Zwingli's doctrines are laid down chiefly in his two Confessions to Charles V. and Francis I. (§ 51), his \emph{Commentarius de vera et falsa religione} (1525), and his sermon \emph{De Providentia Dei} (1530).

Of secondary doctrinal importance are the \emph{ Explanation of his Articles and Conclusions} (1523); his \emph{Shepherd} (a sort of pastoral theology); several tracts and letters on the Lord's Supper, on Baptism and re-Baptism; and his Commentaries on Genesis, Exodus, the Gospels, the Romans, and Corinthians (edited, from his lectures and sermons, by Leo Judä, Megander, and others).

Zwingli's theological system contains, in germ, the main features of the Reformed Creed, as distinct from the Lutheran, and must be here briefly considered.

1. Zwingli begins with the objective (or formal) principle of Protestantism, namely, the exclusive and absolute authority of the Bible in all matters of Christian faith and practice. The Reformed Confessions do the same; while the Lutheran Confessions start with the subjective (or material) principle of justification by faith alone, and make this 'the article of the standing and falling Church.' This difference, however, is more a matter of logical order and relative importance. Word and faith are inseparable, and proceed from the same Holy Spirit. In both denominations a living faith in Christ is the first and last principle. Without this faith the Bible may be esteemed as the best book, but not as the inspired word of God and rule of faith.

2. Zwingli teaches the doctrine of unconditional \emph{election} or predestination to salvation (\emph{constitutio de beandis,} as he defines it), and finds in it the ultimate ground of our justification and salvation; faith being only the organ of appropriation. God is the infinite being of beings, in whom and through whom all other beings exist; the supreme cause, including as dependent organs the finite or middle causes; the infinite and only good (Luke xviii. 18), and every thing else is good (Gen. i. 31) only through and in him. It is a fundamental canon that God by his providence, or perpetual and unchangeable rule and administration,\footnote{Zwingli defines \emph{providentia} to be \emph{perpetuum et immutabile rerum universarum regnum et administratio.}} controls and disposes all events, the will and the action; otherwise he would not be omnipotent and omnipresent. There can be no accident. The fall, with its consequences, likewise comes under his foreknowledge and fore-ordination, which can be as little separated as intellect and will. But God's agency in respect to sin is free from sin, since he is not bound by law, and has no bad motive or affection; so the magistrate may take a man's life without committing murder.\footnote{This illustration is used by Myconius in defending the Zwinglian view of Providence. See Schweizer, \emph{Centraldogmen,} Vol. I. p. 133. The illustration of Zwingli, \emph{Opp.} IV. p. 112, concerning the \emph{adulterium Davidis} and the \emph{taurus,} is less happy.} But only those who hear the Gospel and reject it in unbelief are foreordained to eternal punishment. Of those without the reach of Christian doctrine we can not judge, as we know not their relation to election. There may be and are elect persons among the heathen; and the fate of Socrates and Seneca is no doubt better than that of many popes.

Zwingli, however, dwells mainly on the positive aspect of God's providence—the election to salvation. Election is free and independent. It embraces also infants before they have any faith. It does not follow faith, but precedes it. Faith is itself the work of free grace and the sign and fruit of election (Rom. viii. 29, 30;  Acts xiii. 48). We are elected in order that we may believe in Christ and bring forth the fruits of holiness. Faith is trust and confidence in Christ, the union of the soul with him, and full of good works. Hence it is preposterous to charge this doctrine with dangerous tendency to carnal security and immorality.\footnote{As a matter of history, it is an undeniable fact that the strongest predestinarians (whether Augustinians or Calvinists or Puritans) have been the most earnest, energetic, and persevering Christians. Edward Zeller (a cool philosopher and critic of the Tübingen school) clearly explains this connection in his book on the \emph{Theological System of Zwingli,} pp. 17–19: '\emph{Gerade die Lehre von der Erwählung, der man so oft vorgeworfen hat, dass sie die sittliche Kraft lähme, dass sie zu Trägheit and Sorglosigheit hinführe, gerade diese Lehre ist es, aus welcher der Reformirte jene rücksichts- und zweifellose, bis zur Härte und Leidenschaftlichkeit durchgreifende praktische Energie schöpft, wie wir sie an den Helden dieses Glaubens, einem Zwingli, einem Calvin, einem Farel, einem Knox, einem Cromwell, bewundern, welche ihn vor den Zweifeln und Anfechtungen bewahrt, die dem weicheren, tiefer mit sich selbst beschäftigten Gemüth so viel zu schaffen machen, von denen selbst der grosse deutsche Glaubensheld Luther noch in späten Jahren heimgesucht wurde. Die wesentliche religiöse Bedeutung dieser Lehre, ihre Bedeutung für das innere Leben der Gläubigen, liegt nicht in der Ueberzeuzung von der Unbedingtheit des göttlichen Wirkens als solchen, sondern in dem Glauben an seine Unbedingtheit }in seiner Richtung auf dieses bestimmte Subjekt\emph{ in jener}  persönlichen Gewissheit \emph{der Erwählung, welche den Unterschied der reformirten Erwählungslehre von der augustinischen ausmacht, und eben darauf beruht es auch, dass die theoretisch ganz richtigen Konsequenzen des Prädestinatianismus in Beziehung auf die Nutzlosigkeit und Gleichgültigkeit des eigenen Thuns den Reformirten nicht blos nicht stören, sondern gar nicht für ihn vorhanden sind. Was er in den Sätzen von der ewigen Vorherbestimmung aller Dinge, von dem unwandelbaren Rathschluss der Erwählung und der Verwerfung, für sich selbst findet, das ist nur die unzweifelhafte Gewissheit, persönlich zum Dienst Gottes berufen zu sein, und vermöge dieser Berufung in allen seinen Angelegenheiten unter dem unmittelbarsten Schutz Gottes zu stehen, als Werkzeug Gottes zu handeln, der Seligkeit gewiss zu sein. Die Heilsgewissheit ist hier von der sittlich religiösen Anforderung nicht getrennt, der Einzelne hat das Bewusstsein seiner Berufung nur in seinem Glauben, und den Glauben nur in der Kräftigkeit seines gottbeseelten Willens, er ist sich nicht seiner Erwählung zur Seligkeit ohne alle weitere Bestimmung, sondern wesentlich nur seiner Erwählung zu der Seligkeit des christlichen Lebens bewusst; die Erwählung ist hier nur die Unterlage für das praktische Verhalten des Frommen, der Mensch verzichtet nur desshalb im Dogma auf die Kraft und Freiheit seines Willens, um sie für das wirkliche Leben und Handeln von der Gottheit, an die er sich ihrer entäussert hat, als eine absolute, als die Kraft des göttlichen Geistes, als die unerschütterliche Selbstgewissheit des Erwählten zurückzuerhalten.}'}

This is substantially Zwingli's doctrine, as he preached it during the Conference in Marburg (1529), and taught it in his book on \emph{Providence}\footnote{\par Zwingli, being requested by Philip of Hesse (Jan. 25, 1530) to send him a copy of his sermon, which he had preached without manuscript, reproduced the substance of it, and sent it to him, Aug. 20, 1530. under the title, \emph{Ad illustrissimum Cattorum principem Philippum sermonis de Providentia Dei anamnema. Opera} IV. pp. 79–144. See a full extract in Schweizer's \emph{Centraldogmen,} Vol. I. pp. 102 sqq. Ebrard makes too little account of this tract.} It was afterwards more fully and clearly developed by the powerful intellect of Calvin,\footnote{In the later editions of his \emph{Institutes}; for in the first edition he confines himself to a very brief and indefinite statement of this doctrine.} who made it the prominent pillar of his theology and impressed it upon the majority of the Reformed Confessions, although several of them simply teach a free election to salvation, without saying a word of the decree of reprobation.

On this subject, however, as previously stated, there was no controversy among the early Reformers. They were all Augustinians. Luther heard Zwingli's sermon on Providence in Marburg, and made no objection to it, except that he quoted Greek and Hebrew in the pulpit. He had expressed himself much more strongly on the subject in his famous book against Erasmus (1525). There was, however, this difference, that Luther, like Augustine, from his denial of the freedom of the human will, was driven to the doctrine of absolute predestination, as a logical consequence; while Zwingli, and still more Calvin, started from the absolute sovereignty of God, and inferred from it the dependence of the human will; yet all of them were controlled by their strong sense of sin and free grace much more than by speculative principles. The Lutheran Church afterwards dropped the theological inference in part—namely, the decree of reprobation—and taught instead the universality of the offer of saving grace; but she retained the anthropological premise of total depravity and inability, and also the doctrine of a free election of the saints, or predestination to salvation; and this after all is the chief point in the Calvinistic system, and the only one which is made the subject of popular instruction. In the Lutheran Church, morever, the election theory is moderated by the sacramental principle of baptismal regeneration (as was the case with Augustine), while in the Reformed Church the doctrine of election controls and modifies the sacramental principle, so that the efficacy of baptism is made to depend upon the preceding election.

3. The most original and prominent doctrine of Zwingli is that of the \emph{sacraments,} and especially of the \emph{Lord's Supper.}

He adopts the general definition that the sacrament is the visible sign of an invisible grace, but draws a sharp distinction between the sacramental sign (\emph{signum}) and the thing signified (\emph{res sacramenti}), and allows no necessary and internal connection between them. The baptism by water may take place without the baptism of the Spirit (as in the case of Ananias and Simon Magus), and the baptism by the Spirit, or regeneration, without the baptism by water (for the apostles received only John's baptism; the penitent thief was not baptized at all, and Cornelius was baptized after regeneration). Communion with Christ is not confined to the Lord's Supper, neither do all who partake of this ordinance really commune with Christ. The Spirit of God is free and independent of all outward ceremonies and observances.

As to the effect of the sacraments, Zwingli rejects the whole scholastic theory of the \emph{opus operatum,} and makes faith the necessary medium of sacramental efficacy. He differs here not only from the Romish, but also from the Lutheran theory. He regards the sacraments only as signs and seals, and not strictly as means or instrumentalities of grace, except in so far as they strengthen it. They do not originate and confer grace, but presuppose it, and set it forth to our senses, and confirm it to our faith. As circumcision sealed the righteousness of the faith of Abraham, which he had before in a state of uncircumcision (Rom. iv. 11), so baptism seals the remission of sin by the cleansing blood of Christ, and our incorporation in Christ by faith, which is produced by the Holy Spirit. In infant baptism (which he strongly defended against the Anabaptists, not indeed as necessary to salvation, but as proper and expedient), we have the divine promise which extends to the offspring, and the profession of the faith of the parents with their pledge to bring up their children in the same. The Lord's Supper signifies and seals the fact that Christ died for us and shed his blood for our sins, that he is ours and we are his, and that we are partakers of all his benefits. Zwingli compares the sacrament also to a wedding-ring which seals the marriage union.

He fully admits, however, that the sacraments are divinely instituted and necessary for our twofold constitution; that they are significant and efficacious, not empty, signs; that they aid and strengthen our faith ('\emph{auxilium opemque adferunt fidei}'), and so far confer spiritual blessing through the medium of appropriating faith. In this wider sense they may be called means of grace. He also gives them the character of public testimonies, by which we openly profess our faith before God and the world, pledge our obedience to him, and express our gratitude for mercies received. Hence the name \emph{eucharist,} or \emph{gratiarum actio.}

Concerning the Lord's Supper, Zwingli teaches, in opposition to the Romish mass, that it is a \emph{commemoration,} not a repetition, of the atoning sacrifice of Christ, who offered himself once for all time, and can not be offered by any other; that bread and wine signify or represent, but are not really, the broken body and shed blood of our Lord; that he is present only according to his divine nature and by his Spirit to the eye of faith (\emph{fidei contemplatione}), but not according to his human nature, which is in heaven at the right hand of God, and can not be present every where or in many places at the same time; that to eat his flesh and to drink his blood is a spiritual manducation, or the same as to believe in him (John vi.), and no physical manducation by mouth and teeth, which, even if it were possible, would be useless and unworthy and would establish two ways of salvation—one by faith, the other by literal eating in the sacrament; finally, that the blessing of the ordinance consists in a renewed application of the benefits of the atonement by the worthy or believing communicants, while the unworthy receive only the outward signs to their own judgment.

He therefore rejects every form of a local or corporeal presence, whether by transubstantiation, impanation, or consubstantiation, as contrary to the Bible, to the nature of faith, and to sound reason. He supports the figurative interpretation of the words of institution\footnote{That is, of the verbal copula \textgreek{ἐστί, } \emph{est}=\emph{significat,} not of \textgreek{τοῦτο} (Carlstadt), nor \textgreek{σῶμα}= \emph{figura corporis} (Œcolampadius, on the ground that Christ probably did not use the verb at all in the original Aramaic). Zwingli was always inclined to a tropical interpretation, and averse to the notion of a carnal presence, but was led to his exegesis in 1522 by a tract of Honius (Hoen), a lawyer of Holland, \emph{De eucharistia,} which taught him \emph{in qua voce tropus lateret.} See Ebrard, Vol. II. p. 97. His controversy with Luther began when he wrote a letter to Matth. Alber, at Reutlingen, Nov. 16, 1524.} by a large number of passages, where Christ is said to \emph{be} the door, the lamb, the rock, the vine, etc.; also by such passages as Gen. xli. 26, 27 (the seven good kine \emph{are} seven years), Matt. xiii. 31–37 (the field \emph{is} the world; the tares \emph{are} the children of the wicked one; the reapers \emph{are} the angels), and especially Luke xxii. 20; 1 Cor. xi. 25 (the \emph{cup is} the New Testament in my blood). He proves the local absence of Christ's body by the fact of his ascension to heaven, his future visible return to judgment, and by such passages as, 'I go to prepare a place for you;' 'The poor you have always with you, but me you have not always;' 'I go to my Father;' 'The heaven must receive him until the times of restitution of all things.' He also points out the inconsistency of Luther in maintaining the literal presence of Christ in the sacrament, and yet refusing the adoration; for wherever Christ is he must be adored.

I add his last words on the subject from the Confession sent to King Francis I. shortly before his death: 'We believe that Christ is truly present in the Lord's Supper; yea, we believe that there is no communion without the presence of Christ.\footnote{\par '\emph{Christum credimus vere esse in cœna, immo non credimus esse Domini cœnam nisi Christus adsit.}' Niemeyer, p. 71.} This is the proof: ``Where two or three are gathered together in my name, there am I in the midst of them'' (Matt. xviii. 20). How much more is he present where the whole congregation is assembled to his honor! But that his body is literally eaten is far from the truth and the nature of faith. It is contrary to the truth, because he himself says: 'I am no more in the world'' (John xvii. 11), and 'The flesh profiteth nothing'' (John vi. 63), that is to eat, as the Jews then believed and the Papists still believe. It is contrary to the nature of faith (I mean the holy and true faith), because faith embraces love, fear of God, and reverence, which abhor such carnal and gross eating, as much as any one would shrink from eating his beloved son. . . . We believe that the true body of Christ is eaten in the communion in a sacramental and spiritual manner by the religions, believing, and pious heart (as also St. Chrysostom taught). And this is in brief the substance of what we maintain in this controversy, and what not we, but the truth itself teaches.' To this he adds the communion service, which he introduced in Zurich, that his Majesty may see how devoutly the sacrament is celebrated there in accordance with the institution of Christ. This service is much more liturgical than the later Calvinistic formulas, and includes the 'Gloria in Excelsis,' the Apostles' Creed, and responses.

Closely connected with the eucharistic controversy are certain christological differences concerning ubiquity and the \emph{communicatio idiomatum,} which we have already discussed in the section on the Formula of Concord.

Zwingli's doctrine of the Eucharist is unquestionably the simplest, clearest, and most intelligible theory. It removes the supernatural mystery from the ordinance, and presents no obstacles to the understanding. Exegetically, it is admissible, and advocated even by some of the ablest Lutheran scholars, who freely concede that the literal interpretation of the words of institution, to which Luther appealed first and last against the arguments of Zwingli, is impossible, or, if consistently carried out, must lead to the Romish dogma.\footnote{See above, p. 327.} Philosophically and dogmatically, it labors under none of the difficulties of transubstantiation and consubstantiation, both of which imply the simultaneous multipresence of a corporeal substance, and a physical manducation of Christ's crucified body and blood—in direct contradiction to the essential properties of a body, and the testimony of four of our senses. It has been adopted by the Arminians, and it extensively prevails at present even among orthodox Protestants of all denominations, especially in England and America.\footnote{Dr. Hodge, e.g., does not rise above the Zwinglian view. He denies that Christ is present in any other way than spiritually, and that believers receive any other benefit than 'the sacrificial virtue and effects of the death of Christ on the cross,' which he maintains was received already by the saints of the Old Testament and the disciples at the first Supper, before the glorified body of Christ had any existence. 'The efficacy of this sacrament, as a means of grace, is not in the signs, nor in the service, nor in the minister, nor in the word, but in the attending influence of the Holy Spirit.'—\emph{System. Theol.} Vol. III. pp. 646, 647, 650.}

Zwingli is no doubt right in his protest against every form, however refined and subtle, of the old Capernaitic conception of a carnal presence and carnal appropriation (John vi. 63). He is also right in his positive assertion that the holy communion is a commemoration of the all-sufficient sacrifice of Christ on the cross, and a spiritual feeding on Christ by faith. But he falls short of the \emph{whole} truth; he does not do justice to the strong language of our Lord, especially in John vi. 53–58, concerning the eating of the flesh of the Son of Man (whether this be referred directly or indirectly to the Lord's Supper, or not). After all deduction of carnal misconceptions, there remains the mystery of a vital union of the believer with the \emph{whole} Christ, including his \emph{humanity,} viewed not, indeed, as material substance, but as a principle of life and power.

This Calvin felt. Hence he endeavored to find a \emph{via media} between Zwingli and Luther, and assumed, besides the admitted real presence of the Divine Lord, a dynamic presence and influence of his glorified and ever-living humanity, and an actual communication of its life-giving power (not the \emph{matter} of the body and blood) by the Holy Ghost to the worthy communicant through the medium of faith—as the sun is in the heavens, and yet with his light and heat present on earth. This theory passed substantially into the most authoritative confessions of the sixteenth century, and must therefore be regarded as the orthodox doctrine of the Reformed Church.

On three other points—namely, original sin, the salvation of infants, and the salvation of the heathen—Zwingli had peculiar views, which were in advance of his age, and gave great offense to some of his friends as well as to Luther, but were afterwards adopted by the Arminians.

4. The Reformation was born of an intense conviction of the sinfulness of man and the absolute need of a radical regeneration. Zwingli makes no exception, and describes the corruption and slavery of the natural man almost as strongly as Luther, although he never passed through such terrors of conscience as the monk in Erfurt, nor had he such hand-to-hand fights with the devil.\footnote{Dorner (in his \emph{History of German Theology,} p. 287) says that Zwingli retained from his humanistic culture a certain disposition to 'an aesthetic consideration of sin,' i.e., to see in it something disgraceful, unworthy, bestial rather than diabolical.} He derives sin from the fall of Adam, brought about by the instigation of the devil, and finds its essence in selfishness as opposed to the love of God. He goes beyond the Augustinian infralapsarianism, which seems to condition the eternal counsel of God by the first self-determination of man, and he boldly takes the supralapsarian position that God not only foresaw, but foreordained the fall, together with the redemption, that is, as a means to an end, or as the negative condition for the revelation of the plan of salvation. He fully admits the distinction between original or hereditary sin and actual transgression, but he describes the former as a moral disease, or natural defect, rather than punishable sin and guilt.\footnote{\emph{Defectus naturalis,} or, as he often calls it in his Swiss-German, a \emph{Brest,} i.e. \emph{Gebrechen.} '\emph{Die Erbsünd,}' he says in his book on Baptism,' \emph{ist nüts} (\emph{nichts}) \emph{anders weder} (\emph{als}) \emph{der Brest von Adam her. . . . Wir verstond} (\emph{verstehen}) \emph{durch das Wort Brest einen Mangel, den einer ohn sin Schuld von der Geburt her hat oder sust} (\emph{sonst}) \emph{von Zufällen.}' He distinguishes it from \emph{Laster} and \emph{Frevel,} vice and crime. He explains his view more fully in his tract \emph{De peccato originali ad Urbanum Regium,} 1526, and also in his Confession to Charles V., 1530.} It is a miserable condition (\emph{conditio misera}). He compares it to the misfortune of one born in slavery.\footnote{\par '\emph{Peccatum originale non proprie peccatum est, non enim est facinus contra legem. Morbus igitur est proprie et conditio.}' \emph{Fidei Ratio ad Carol. V.} Cap. IV. (Niemeyer, p. 20).} But if not sin in the proper sense of the term, it is an inclination or propensity to sin (\emph{propensio ad peccandum}), and the fruitful germ of sin, which will surely develop itself in actual transgression. Thus the young wolf is a rapacious animal before he actually tears the sheep.

5. Zwingli was the first to emancipate the \emph{salvation of children} dying in infancy from the supposed indispensable condition of water-baptism, and to extend it beyond the boundaries of the visible Church. This is a matter of very great interest, since the unbaptized children far outnumber the baptized, and constitute nearly one half of the race.

He teaches repeatedly that all elect children are saved whether baptized or not, whether of Christian or heathen parentage, not on the ground of their innocence (which would be Pelagian), but on the ground of Christ's atonement. He is inclined to the belief that all children dying in infancy belong to the elect; their early death being a token of God's mercy, and hence of their election. A part of the elect are led to salvation by a holy life, another part by an early death. The children of Christian parents belong to the Church, and it would be 'impious' to condemn them. But from the parallel between the first and the second Adam, he infers that all children are saved from the ruin of sin, else what Paul says would not be true, that 'as in Adam all die, even so in Christ shall all be made alive' (1 Cor. xv. 22). At all events, it is wrong to condemn the children of the heathen, both on account of the restoration of Christ and of the eternal election of God, which precedes faith, and produces faith in due time; hence the absence of faith in children is no ground for their condemnation.\footnote{\par \emph{Fidei Ratio,} Cap. V. (Niemeyer, p. 21): '\emph{Hinc constat, si in Christo secundo Adam vitæ restituimur, quemadmodum in primo Adam sumus morti traditi, quod temere damnamus Christianis parentibus natos pueros, imo} gentium \emph{quoque pueros. Adam enim si perderere universum genus peccando potuit, et Christus moriendo non vivificavit et redemit universum genus a clade per istum data, jam non est par salus reddita per Christum, et perinde} (\emph{quod absit!}) \emph{ nec verum,} ``\emph{Sicut in Adam omnes moriuntur, ita in Christo omnes vitæ restituuntur.}'' \emph{Verum quomodocunque de gentilium infantibus statuendum sit, hoc certe adseveramus, propter virtutem salutis per Christum præstitæ, præter rem pronunciare qui eos æternæ maledictioni addicunt, cum propter dictam reparationis causam, tum propter electionem Dei liberam, quæ non sequitur fidem, sed fides electionem sequitur.}' In another passage against the \emph{Catabaptists} he says: '\emph{Electi eligebantur antequam in utero conciperentur: mox igitur ut sunt, filii Dei sunt, etiamsi moriantur antequam credant aut ad fidem vocentur.} Comp. Zeller, 1.c. p. 162.} As he believed in the salvation of many adult heathen, he had the less difficulty in believing that heathen children are saved; for they have not yet committed actual transgression, and of hereditary sin they have been redeemed by Christ. We have therefore much greater certainty of the salvation of departed infants than of any adults.

This view was a bold step beyond the traditional orthodoxy. The Roman Catholic Church, in keeping with her doctrine of original sin and guilt, and the necessity of water-baptism for salvation (based upon  Mark xvi. 16 and John iii. 5), teaches the salvation of all baptized, and the \emph{condemnation} of all \emph{unbaptized} children; assigning the latter to the \emph{limbus infantum} on the border of hell, where they suffer the mildest kind of punishment, namely, the negative penalty of loss (\emph{pœna damni} or \emph{carentia beatificæ visionis}), but not the positive pain of feeling (\emph{pœna sensus}).\footnote{The \emph{limbus infantum} is, so to speak, the nursery of hell, on the top floor and away from the fire, as Bellarmin says, \emph{in loco inferni altiori, ita ut ad eum ignis non perveniat.} In a still higher region was the \emph{limbus patrum,} the temporary abode of the saints of the Old Testament, but this was vacated at the descent of Christ into Hades, when those saints were freed from prison and translated into Paradise.} St. Augustine first clearly introduced this wholesale exclusion of all unbaptized infants from heaven—though Christ expressly says that to children emphatically belongs the kingdom of heaven. He ought consistently to have made the salvation of infants, like that of adults, depend upon their election; but the churchly and sacramental principle checked and moderated his predestination theory, and his Christian heart induced him to soften the frightful dogma as much as possible.\footnote{'\emph{Parvulos non baptizatos in damnatione omnium lenissima futuros}' (\emph{Contra Jul.} lib. V. c. 11); '\emph{Infantes non baptizati } lenissime  \emph{quidem, sed tamen damnantur. Potest proinde recte dici, parvulos sine baptismo de corpore exeuntes in damnatione omnium }mitissima \emph{futuros}' (\emph{De pecc. mer. et rem.} cap. 16). Pelagius was more liberal, and assumed a middle state of half-blessedness for unbaptized infants between the heaven of the baptized and the hell of the ungodly. See particulars in my \emph{Church History,} Vol. III. pp. 835 sqq.} As he did not extend election beyond the boundaries of the Catholic Church (although he could not help seeing the significance of such holy outsiders as Melchizedek and Job under the old dispensation), he secured at least, by his high view of the regenerative efficacy of water-baptism, the salvation of all baptized infants dying in infancy. To harmonize this view with his system, he must have counted them all among the elect.

The Lutheran Creed retains substantially the Catholic view of baptismal regeneration, and hence limits infant salvation to those who enjoy this means of grace;\footnote{\emph{Conf. August.} Art. IX.: '\emph{Damnant Anabaptistas qui . . . affirmant pueros sine baptismo salvos fieri.}' In the German edition the last clause is omitted.} allowing, however, some exceptions within the sphere of the Christian Church, and making the damnation of unbaptized infants as mild as the case will permit.\footnote{Calovius (in the \emph{consensus repetitus fidei vere Lutheranæ,} 1655), in the name of the strict Lutherans, rejected the milder view of a merely privative punishment of unbaptized infants, as held by Calixtus (see Henke, \emph{Georg Calixtus,} Vol. II. Pt. II. p. 295), but it was defended by others. Fr. Buddæus, one of the most liberal among the orthodox Lutherans, describes the condition of heathen infants as \emph{admodum tolerabilis,} though they are \emph{exclusi a beatitudine} (\emph{Instit. Theol. dogm.} Lips. 1723, p. 631). Others leave the children to the mercy of God. See V. E. Löscher's \emph{Auserlesene Sammlung der besten neueren Schriften vom Zustand der Seele nach dem Tode,} 1735; republished by Hubert Becker, 1835.} At present, however, there is scarcely a Lutheran divine of weight who would be willing to confine salvation to \emph{baptized} infants.

The Reformed Church teaches the salvation of all \emph{elect} infants dying in infancy, whether \emph{baptized or not,} and assumes that they are regenerated before their death, which, according to Calvinistic principles, is possible without water-baptism.\footnote{\emph{Westminster Conf.} chap. x. § 3: 'Elect infants, dying in infancy, are regenerated and saved by Christ through the Spirit, who worketh when and where and how he pleaseth. So also are all other elect persons who are incapable of being outwardly called by the ministry of the word.' The last sentence may be fairly interpreted as teaching the election and salvation of a portion of heathen \emph{adults.}} The second Scotch Confession, of 1580, expressly rejects, among other errors of popery,' the cruel judgment against infants departing without the sacrament.'\footnote{'\emph{Abhorremus et detestamur . . . crudele judicium contra infantes sine baptismo morientes, bapitismi absolutam quant asserit necessitatem.}' Niemeyer, pp. 357, 358.} Beyond this the Confessions do not go, and leave the mysterious subject to private opinion. Some of the older and more rigid Calvinistic divines of the supralapsarian type carried the distinction between the elect and the reprobate into the infant world, though always securing salvation to the offspring of Christian parents, on the ground of inherited Church membership before and independent of the baptismal ratification; while others more wisely and charitably kept silence, or left the non-elect infants—if there are such, which nobody knows—to the uncovenanted mercies of God. But we may still go a step further, within the strict limits of the Reformed Creed, and maintain, as a pious opinion, that all departed infants belong to the number of the elect. Their early removal from a world of sin and temptation may be taken as an indication of God's special favor. From this it would follow that the majority of the human race will be saved. The very doctrine of election, which is unlimitable and free of all ordinary means, at all events widens the possibility and strengthens the probability of general infant salvation; while those Churches which hold to the necessity of baptismal regeneration must either consistently exclude from heaven all unbaptized infants (even those of Christian Baptists and Quakers), or, yielding to the instinct of Christian charity, they must make exceptions so innumerable that these would become, in fact, the rule, and overthrow the principle altogether.

In the seventeenth century the Arminians resumed the position of Zwingli, and with their mild theory of original sin (which they do not regard as responsible and punishable before and independent of actual transgression), they could consistently teach the general salvation of infants. The Methodists and Baptists adopted the same view. Even in the strictly Calvinistic churches it made steady progress, and is now silently or openly held by nearly all Reformed divines.\footnote{Dr. Hodge, the most orthodox Calvinistic divine of the age, very positively teaches (\emph{Syst. Theol.} Vol. I. p. 26) the salvation of all infants dying in infancy, and represents this as the 'common doctrine of evangelical Protestants.' This may be true of the present generation, and we hope it is, though it is evidently inapplicable to the period of scholastic orthodoxy, both Lutheran and Calvinistic. He supports his view by three arguments: 1. The analogy between Adam and Christ (Rom. v. 18, 19, where we have no right to restrict the free gift of Christ upon all more than the Bible itself restricts it); 2. Christ's conduct towards children; 3. The general nature of God to bless and to save, rather than to curse and destroy.}

Whether consistent or not, the doctrine of infant damnation is certainly cruel and revolting to every nobler and better feeling of our nature. It can not be charged upon the Bible except by logical inference from a few passages (John iii. 5; Mark xvi. 16; Rom. v. 12), which admit of a different interpretation. On the other hand, the general salvation of infants, though not expressly taught, is far more consistent with the love of God, the genius of Christianity, and the spirit and conduct of him who shed his precious blood for all ages of mankind, who held up little children to his own disciples as models of simplicity and trustfulness, and took them to his bosom, blessing them, and saying (unconditionally and before Christian baptism did exist), 'Of such is the kingdom of heaven,' and 'Whosoever shall not receive the kingdom of God as a little child, he shall in nowise enter therein.'

6. \emph{Salvation of adult heathen.} This is a still darker problem. Before Zwingli it was the universal opinion that there can be no salvation outside of the visible Church (\emph{extra ecclesiam nulla satus}). Dante, the poet of mediæval Catholicism, assigns even Homer, Aristotle, Virgil, to hell, which bears the terrible inscription—

'Let those who enter in dismiss all hope.'

But the Swiss Reformer repeatedly expressed his conviction, to which he adhered to the last, that God had his elect among the Gentiles as well as the Jews, and that, together with the saints of the Old Testament from the redeemed Adam down to John the Baptist, we may expect to find in heaven also such sages as Socrates, Plato, Aristides, Pindar, Numa, Cato, Scipio, Seneca; in short, every good and holy man and faithful soul from the beginning of the world to the end.\footnote{His last and fullest utterance on this subject occurs towards the close of his \emph{Expositio Chr. Fidei,} where, speaking of eternal life, he thus addresses the French king: '\emph{Deinde sperandum est tibi visurum esse sanctorum, prudentium, fidelium, canstantium, fortium, virtuosorum omnium, quicunque a condito mundo fuerunt, sodalitatem, cœtum et contubernium. Hic duos Adam, redemptum ac Redemptorem: hic Abelum, Enochum, Noam, Abrahamum, Isaacum, Jacobum, Judam, Mosen, Josuam, Gedeonem, Samuelem, Pinhen, Heliam, Heliseum, Isaiam, ac deiparam Virginem de qua ille præcinuit, Davidem, Ezekiam, Josiam, Baptistam, Petrum, Paulum: hic } Herculem, Theseum, Socratem, Aristidem, Antigonum, Numam, Camillum, Catones, Scipiones: \emph{hic Ludovichum pium antecessoresque tuos Ludovicos, Philippos, Pipinnos, et quotquot in fide hinc migrarunt maiores tuos videbis. Et summatim, non fuit vir bonus, non erit mens sancta, non est fidelis anima, ab ipso mundi exordia usque ad eius consummationem, quem non sis isthic cum Deo visurus. Quo spectaculo quid lætius, quid amænius, quid denique honorificentius vel cogitari potent? Aut quo iustius omnes animi vires intendimus quam ad huiuscemodi vitæ lucrum?}' See Niemeyer, p. 61. Similar passages occur in his Epistles, Commentaries, and tract on Providence. Comp. Zeller, p. 163.}

For this liberality he was severely censured. The great and good Luther was horrified at the idea that even 'the godless Numa' (!) should be saved, and thought that it falsified the whole gospel, without which there can be no salvation.\footnote{'\emph{Hoc si verum est, totum evangelium falsum est.}' Luther denied the possibility of salvation outside of the Christian Church. In his \emph{Catech. Major,} Pars II. Art. III. (ed. Rechenb. p. 503, ed. Müller, p. 460), he says: '\emph{Quicunque extra Christianitatem} (\emph{ausser der Christenheit}) \emph{sunt, sive Gentiles sive Turcæ sive Judæi aut falsi etiam Christiani et hypocritæ, quanquam unum tantum et verum Deum esse credant et invocent} (\emph{ob sie gleich nur Einen wahrhaftigen Gott gläuben und anbeten}), \emph{neque tamen certum habent, quo erga eos animatus sit animo, neque quidquam favoris aut gratiæ de Deo sibi polliceri audent et possunt, quamobrem } in perpetua manent ira et damnatione (\emph{darum sie im ewigen Zorn und Verdammniss bleiben}).'}

Zwingli, notwithstanding his abhorrence of heathen idolatry and every relic of paganism in worship, retained, from his classical training in the school of Erasmus, a great admiration for the wisdom and the manly virtues of the ancient Greeks and Romans, and was somewhat unguarded in his mode of expression. But he had no idea of sending any one to heaven without the atonement, although he does not state when and how it was applied to those who died before the incarnation. In his mind the eternal election was inseparably connected with the plan of the Christian redemption. He probably assumed an unconscious Christianity among the better heathen, and a secret work of grace in their hearts, which enabled them to exercise a general faith in God and to strive after good works (comp. Rom. ii. 7, 10, 14, 15). All truth, he says, proceeds from the Spirit of God. He might have appealed to Justin Martyr and other ancient fathers, who traced all that was true and good among the Greek philosophers and poets to the working of the Logos before his incarnation (John i. 5, 10).\footnote{Dr. Dorner, with his usual fairness and fine discrimination, vindicates Zwingli against misrepresentations (\emph{Gesch. d. Prot. Theol.} p. 284): '\emph{Man hat daraus eine Gleichgültigkeit gegen den historischen Christus und sein Werk erschliessen wollen, dass er }[\emph{Zwingli}] \emph{auch von Heiden sagt: sie seien selig geworden; was die Heiden Weisheit nennen, das nennen die Christen Glauben. Allein er sieht in allem Wahren vor Christo mit manchen Kirchenvätern eine Wirkung und Offenbarung des Logos, ohne jedoch so weit zu gehen, mit Justin die Weisen des Alterthums, welche nach dem Logos gelebt haben, Christen zu nennen. Er sagt nur, sie seien nach dem Tode selig geworden, ähnlich wie auch die Kirche dasselbe von den Vätern des Alten Testaments annimmt. Er konnte dabei wohl diese Seligkeit als durch Christus gewirkt und erworben denken und hat dieselbe jedenfalls nur als in der Gemeinschaft mit Christus bestehend gedacht. Ist ihm doch durch den ewigen Rathschluss der Versöhnung Christus nicht bloss ewig gewiss, sondern auch gegenwärtig für alle Zeiten. So sind ihm jene Heiden doch selig nur durch Christus. Freilich das sagt er nicht, dass sie erst im Jenseits sich bekehren; auch erschneidet mit dem Diesseits die Bekehrung ab. Er lässt ihre im Diesseits bewährte Treue gegen das ihnen vom Logos anvertraute Pfund wahrer Erkentniss die Stelle des Glaubens vertreten. Aber es ist wohl kein Zweifel, dass er sie im Jenseits zur Erkentniss und Gemeinschaft Christi gelangend denkt. Bei den Frommen Alten Testaments fordert auch die Kirche zu ihrem Heil nicht eine bestimmtere Erkenntniss Christi im Diesseits, die sie höchstens den Propheten' zuschreiben könnte.}' Ebrard (in his \emph{History of the Dogma of the Lord's Supper.} Vol. II. p. 77) fully adopts Zwingli's view: '\emph{Jetzt wird ihm wohl Niemand mehr daraus ein Verbrechen machen. Wir wissen, dass }  \emph{Röm. ii.} 7: ``\emph{Denen, die in Beharrlickkeit des Gutesthuns nach unvergänglichem Wesen }trachten,'' \emph{ ewiges Leben verheissen ist, wir wissen dass nur der positive Unglaube an das angebotene Heil weder hier noch dort vergeben wird, dass nur auf ihn die Strafe des ewigen Todes gesetzt ist; wir wissen, dass auf die erste Auferstehung der in Christo Entschlafenen noch eine zweite der ganzen übrigen Menschheit folgen soll, die alsdann gerichtet werden sollen nach ihren Werken, und dass im neuen Jerusalem selber die Blätter des Lebensbaumes dienen sollen zur Genesung der Heiden} (Apok. xxii. 2). \emph{ Zwingli hat also an der Hand der heiligen Schrift das Heidenthum ebenso wie das Judenthum als an den } \textgreek{στοιχείοις τοῦ κόσμου } \emph{gehörig} (Gal. iv. 1–3) \emph{angesehen, und mit vollem Rechte einen Socrates neben einen Abraham gestellt. Ihm besteht die Seligkeit darin, dass das ganze Wunderwerk der göttlichen Weltpädagogik in seinen Früchten klar und herrlich vor den Blicken der erstaunten Seligen da liegt.}'}

During the period of rigorous scholastic orthodoxy which followed the Reformation in the Reformed and Lutheran Churches, Zwingli's view could not be appreciated, and appeared as a dangerous heresy. In the seventeenth century the Romanists excluded the Protestants, the Lutherans the Calvinists, the Calvinists the Arminians, from the kingdom of heaven; how much more all those who never heard of Christ. This wholesale damnation of the vast majority of the human race should have stirred up a burning zeal for their conversion; and yet during that whole period of intense confessionalism and exclusive orthodoxism there was not a single Protestant missionary in the field except among the Indians in the wilderness of North America.\footnote{John Eliot, the 'Apostle of the Indians,' labored among the Indians in that polemical age. He died 1690, eighty-six years of age, at Roxbury, Massachusetts. David Brainerd (d. at Northampton, Mass., 1747) likewise labored among the Indians before any missionary zeal was kindled in the Protestant churches of Europe.}

But in modern times Zwingli's view has been revived and applauded as a noble testimony of his liberality, especially among evangelical divines in Germany, and partly in connection with a new theory of Hades and the middle state.

This is not the place to discuss a point which, in the absence of clear Scripture authority, does not admit of symbolical statement. The future fate of the heathen is wisely involved in mystery, and it is unsafe and useless to speculate without the light of revelation about matters which lie beyond the reach of our observation and experience. But the Bible consigns no one to final damnation, except for rejecting Christ in unbelief,\footnote{John iii. 18, 36; xii. 48;  Mark xvi. 16.} and gives us at least a ray of hope by significant examples of faith from Melchizedek and Job down to the wise men from the East, and by a number of passages concerning the working of the Logos among the Gentiles (John i. 5, 10; Rom. i. 19; ii. 14, 15, 18, 19; Acts xvii. 23, 28; 1 Pet. iii. 19; iv. 6). We certainly have no right to confine God's election and saving grace to the limits of the visible Church. We are indeed bound to his ordinances and must submit to his terms of salvation; but God himself is free, and can save whomsoever and howsoever he pleases, and he is infinitely more anxious and ready to save than we can conceive.

\xsubsection{The First Confession of Basle. A.D. 1534.}

§ 53. The First Confession of Basle. A.D. 1534.

Literature.

Jac. Christ. Beck:  \emph{Dissertatio historico-theologica de Confessione Fidei Basileensis Ecclesiæ,} Basil. 1744.

Melchior Kirchhofer:  \emph{Oswald Myconius, Antistes der Baslerischen Kirche,} Zürich, 1813.

Burckhardt:  \emph{Reformationsgeschichte von Basel,} Basel, 1818.

K. R. Hagenbach:  \emph{Kritische Geschichte der Entstehung und der Schicksale der ersten Basler Confession und der auf sie gegründeten Kirchenlehre,} Basel, 1827 (title ed. 1828).

J. J. Herzog:  \emph{Leben Joh. Œkolampads und die Reformation der Kirche von Basel,} Basel, 1843, 2 vols.

Hagenbach:  \emph{Leben Œkolampads und Myconius,} Elberfeld, 1859. (Part II. of \emph{Väter und Bergründer der reform. Kirche.})

Escher, in \emph{Ersch und Gruber's Encyklop.} Art. \emph{Helvet. Confess.} Sect. II. Part V.

Beck:  \emph{Symb. Bücher der ev. reform. Kirche,} Vol. I. pp. 28 sqq.

The two Confessions of Basle are published in German and Latin by Niemeyer, \emph{Coll.} pp. 78–122: in German alone by Beck and Böckel in their collections. The best reprint of the First Confession of Basle, in the Swiss dialect, with the Scripture proofs on the margin, is given by Hagenbach at the close of his biography of \emph{Œkolampad und Myconius,} pp. 465–470.

The First and Second Confessions of Basle belong to the Zwinglian family, and preceded the age of Calvin, but are a little nearer the German Lutheran type of Protestantism.

The rich and venerable city of Basle, on the frontier of Switzerland, France, and South Germany, since 1501 a member of the Swiss Confederacy, renowned for the reformatory Œcumenical Council of 1430, and the University founded by Pius II., became a centre of liberal learning before the Reformation. Thomas Wyttenbach, the teacher of Zwingli, attacked the indulgences as early as 1502. In 1516 Erasmus of Rotterdam, at that time esteemed as the greatest scholar of Europe, took up his permanent residence in Basle, and published the first edition of the Greek Testament and other important works, though, after the peasant war and Luther's violent attack on him, he became disgusted with the Reformation, which he did not understand. He desired merely a quiet literary illumination within the Catholic Church, and formed a bridge between two ages. He died, like Moses, in the land of Moab (1536).\footnote{Erasmus turned his keen wit first against the obscurantism of the monks, but afterwards against the light of the Reformation. He said to Frederick the Wise at Cologne, before the Diet of Worms (within the hearing of Spalatin): '\emph{Lutherus peccavit in duobus, nempe quod tetigit coronam pontificis et ventres monachorum.}' But when Luther, Zwingli, Œcolampadius took wives, he called the Reformation a comedy which ended always in a marriage.} Wolfgang Capito (Köpfli), an Alsacian, labored in Basle as preacher and professor from 1512 to 1520, in friendly intercourse with Erasmus, and was followed by Caspar Hedio (Heid), who continued in the same spirit, and corresponded with Luther. Another preacher in Basle, Wilhelm Röublin, carried on the Corpus Christi festival a large Bible through the city, with the inscription, 'This is the true sanctuary; the rest are dead men's bones.'

The principal Reformer of Basle is John Œcolampadius (Hausschein, b. 1482, d. 1531), who stood to Zwingli in a similar relation as Melanchthon to Luther: inferior to him in originality, boldness, and energy, but superior in learning, modesty, and gentleness of spirit. He was his chief support in the defense of his doctrine on the eucharist, and took a prominent part in the Conference with Luther at Marburg. Born at Weinsberg, he studied philology, scholastic philosophy, law, and theology with unusual success at Heilbronn, Bologna, Heidelberg, and Tübingen. When twelve years old he wrote Latin poems, and at fourteen he graduated as bachelor of arts. He excelled especially as a Greek and Hebrew scholar, and published afterwards learned commentaries on the prophets and other books of the Bible. He aided Erasmus in the edition of his Greek Testament, 1516. He was well-read in the fathers, and promoted a critical study of their writings. After having labored as preacher for some time in different places, and taken some part in the reformatory movements of Germany, he settled permanently at Basle, in 1522, as pastor of St. Martin and as professor of theology. Here he introduced, with the consent of the citizens, the German service, the communion under both kinds, and other changes. But it was only after the transition of Berne that Basle came out decidedly for the Reformation. It was formally introduced Feb. 9, 1529, according to the model of Zurich, but in a rather violent style, by the breaking of images and the dissolution of convents, yet without shedding of blood. In other respects the Reformed Church of Basle is conservative, and occupies a middle position between Zwinglianisrn and Lutheranism. Œcolampadius died Nov. 24, 1531, a few weeks after his friend Zwingli. He communed with his family, and took an affecting farewell of his wife, his three children (Eusebius, Irene, and Aletheia), and the ministers of Basle. His last words were: 'Shortly I shall be with the Lord Christ. . . . Lord Jesus save me!'\footnote{See the particulars in Herzog's \emph{Œkolamp.} Vol. II. pp. 248 sqq. He was buried with all the honors of the city in the Minster. But the mouth of slander spread the lie that he had committed suicide, to which even Luther, blinded by dogmatic prejudice, was not ashamed to give ear. Melanchthon had great respect for Œcolampadius, stood in friendly correspondence with him, and derived from him a better knowledge of the patristic doctrine of the eucharist.}

The First Confession of Basle (\emph{Confessio Fidei Basileensis prior}) was prepared in its first draft by Œcolampadius, 1531,\footnote{See Herzog, 1.c. Vol. II. pp. 217–221, and Hagenbach, \emph{Joh. Œkol. und Oswald Mycon.} pp. 350 sqq. Œcolampadius, in his last address to the Synod of Basle, Sept. 26, 1531, added a brief, terse confession of faith, and a paraphrase of the Apostles' Creed. But the assertion that he composed the Confession of Basle in its present shape, and sent it to the Augsburg Diet, 1530, rests on a mistake, and has no foundation in any contemporary report.} brought into its present shape by his successor, Oswald Myconius,\footnote{His proper name was Geisshüssler. He was born at Luzerne, 1488; taught and preached at Zurich; after Zwingli's death he moved to Basle, was elected Antistes or first preacher, died 1552, and was buried in the Minster. He must not be confounded with Friedrich Myconius, or Mecum, the Lutheran reformer of Thuringia, and court chaplain at Gotha (d. 1546).} 1532, and first published by the magistrate with a preface of Adelberg Meyer, burgomaster of Basle, Jan. 21, 1534.\footnote{Under the title, '\emph{Bekanntnuss unseres heiligen Christlichen Glaubens wie es die Kylch (Kirche) zu Basel halt.}' It is signed by '\emph{Heinrich Rhyner, Rathschreiber der Statt Basel.}' See the German text, with the marginal notes, at the close of Hagenbach's biography of Œcolampadius and Myconius. A Latin edition appeared 1561 and 1581, which was reproduced in the \emph{Corpus et Syntagma Confess.,} under the title '\emph{Basiliensis vel Mylhusiana Confessio Fidei. anno M.D.XXXII. Scripta Germanice. Latine excusa} 1561 \emph{et} 1581.' Here the date of composition (1532) is given instead of the date of publication (1534). The more usual spelling is \emph{Basileensis} and \emph{Mühlhusana.} A better Latin edition was issued, 1647, by the Basle Professors—Theod. Zwinger, Sebastian Beck, and John Buxtorf—for the use of academic disputations; and this Niemeyer has reprinted, pp. 85 sqq.} Two or three years afterwards it was adopted and issued by the confederated city of Mühlhausen, in the Alsace; hence it is also called the \emph{Confessio Mühlhusana} (or \emph{Mylhusiana}).

It is very simple and moderate. It briefly expresses, in twelve articles, the orthodox evangelical doctrines of God, the fall of man, the divine providence, the person of Christ, the Church and the sacraments, the Lord's Supper (Christ the food of the soul to everlasting life), Church discipline, the civil magistrate, faith and works, the judgment, ceremonies and celibacy, and against the views of the Anabaptists, who were then generally regarded as dangerous radicals, not only by Luther, but also by the Swiss and English Reformers. This is the only Reformed Confession which does not begin with the assertion of the Bible principle, but it concludes with this noteworthy sentence: '''We submit this our Confession to the judgment of the divine Scriptures, and hold ourselves ready always thankfully to obey God and his Word if we should be corrected out of said holy Scriptures.'\footnote{'\emph{Postremo, hanc nostrum Confessionem judicio }Sacræ Biblicæ Scripturæ \emph{subjicimus: eoque pollicemur, si ex prædictis Scripturis in melioribus instituamur }(\emph{etwas besseren berichtet}), \emph{nos omni tempore} Deo  \emph{et } sacrosancto ipsius verbo,  \emph{maxima cum gratiarum actione, obsecuturos esse.}'}

'This Confession,' says the late Professor Hagenbach of Basle,\footnote{\emph{Joh. Œkolampad und Oswald Myconius,} p. 353; comp. his \emph{History of the Conf.} pp. 190 sqq.} 'has remained the public Confession of the Church of Basle to this day. It is, indeed, no longer annually read before the congregation as formerly (on Maundy-Thursday at the ante-communion service), but ministers at their ordination are still required to promise ``to teach according to the direction of God's Word and the Basle Confession derived therefrom.'' A motion was made in the city government in 1826 to change it, but the Church Council declared such change inexpedient. Another motion in 1859 to abolish it altogether was set aside. But the political significance of the Confession can no longer be sustained, in view of the change of public sentiment in regard to the liberty of faith and conscience.'

\xsubsection{The First Helvetic Confession. A.D. 1533.}

§ 54. The First Helvetic Confession, A.D. 1536.

See \emph{Literature} in § 53. Comp. also Pestalozzi:  \emph{Heinrich Bullinger,} pp. 183 sqq.

The First Helvetic Confession (\emph{Confessio Helvetica prior}), so called to distinguish it from the Second Helvetic Confession of 1566, is the same with the Second Confession of Basle (\emph{Basileensis posterior}), in distinction from the First of 1534.\footnote{Hagenbach, l.c. p. 357: 'Basler  \emph{Confession heisst diese Confession nur weil sie }in,  \emph{nicht weil sie } für \emph{Basel verfasst ist} (\emph{ähnlich wie die Augsburger Confession von dem Ort der Uebergabe den Namen hat}). \emph{Bezeichnender ist daher der Name erste }Helvetische  \emph{Confession, weil sie das Gesammtbekenntniss der reformirten Schweizerkirchen ist.}'} It owes its origin partly to the renewed efforts of the Strasburg Reformers, Bucer and Capito, to bring about a union between the Lutherans and the Swiss, and partly to the papal promise of convening a General Council. A number of Swiss divines were delegated by the magistrates of Zurich, Berne, Basle, Schaffhausen, St. Gall, Mühlhausen, and Biel, to a Conference in the Augustinian convent at Basle, January 30, 1536. Bucer and Capito also appeared. Bullinger, Myconius, Grynæus, Leo Judæ, and Megander were selected to draw up a Confession of the faith of the Helvetic Churches, which might be used before the proposed General Council. It was examined and signed by all the clerical and lay delegates, February, 1536, and first published in Latin.\footnote{Sub titulo: '\emph{Ecclesiarum per Helvetiam Confessio Fidei summaria et generalis,}' etc. The German is inscribed, '\emph{Eine kurze und gemeine Bekenntniss des heiligen, wahren und uralten christlichen Glaubens der Kirchen,} etc., \emph{Zürich, Bern, Basel, Strassburg, Constanz, St. Gallen, Schaffhausen, Mühlhausen, Biel,} etc., 1536, \emph{Februariy.}'} Leo Judæ prepared the German translation, which is fuller than the Latin text, and of equal authority.

Luther, to whom a copy was sent through Bucer, expressed unexpectedly, in two remarkable letters, his satisfaction with the earnest Christian character of this document, and promised to do all he could to promote union and harmony with the Swiss.\footnote{See his letter to Jacob Meyer, burgomaster of Basle, Feb. 17, 1535, and his response to the Reformed Cantons, Dec. 1, 1537 (in De Wette, Vol. V. pp. 54 and 83). Luther kept the peace with the Swiss churches only for a few years. In his book against the Turks, 1541, he calumniated without provocation the memory of Zwingli; in August, 1543, he acknowledged the present of the Zurich translation of the Bible sent to him by Froschauer, the publisher, but scornfully declined to accept any further works from preachers 'with whom neither he nor the Church of God could have any communion, and who were driving people to hell' (see his letter in De Wette, Vol. V. p. 587); in 1544 he violently renewed, to the great grief of Melanchthon, the sacramental war in his '\emph{Short Confession of the Sacrament;}' and shortly before his death he was not ashamed to travesty the first Psalm thus: '\emph{Beatus vir, qui non abiit in consilio Sacramentariorum: nec stetit in via Cinglianorum, nec sedet in cathedra Tigurinorum.}' (See his letter to Jac. Probst of Bremen, Jan. 17, 1546, in De Wette, Vol. V. p. 778. Comp. also on this whole subject Hagenbach, l.c. p. 358, and Pestalozzi, l.c. pp. 216 sqq.). Myconius was not disturbed by these outbursts of passion, and continued to respect Luther without departing from the doctrine of his friend Zwingli. He judged, not without some reason, that the two Reformers never understood each other; that Luther stubbornly believed that Zwingli taught the sacrament to be an empty sign, and Zwingli that Luther taught a gross Capernaitic eating. See his letter of Sept. 7, 1538, to Bibliander, in \emph{Simmler's Collection,} Vol. XLV., and Hagenbach, p. 350.} He was then under the hopeful impressions of the 'Wittenberg Concordia,' which Bucer had brought about by his elastic diplomacy, May, 1536, but which proved after all a hollow peace, and could not be honestly signed by the Swiss.

The Helvetic Confession is the first Reformed Creed of national authority. It consists of twenty-seven articles, is fuller than the first Confession of Basle, but not so full as the second Helvetic Confession, by which it was afterwards superseded. The doctrine of the sacraments and of the Lord's Supper is essentially Zwinglian, yet emphasizes the significance of the sacramental signs and the real spiritual presence of Christ, who gives his body and blood—that is, himself—to believers, so that he more and more lives in them and they in him.

It seems that Bullinger and Leo Judæ wished to add a caution against the binding authority of this or any other confession that might interfere with the supreme authority of the Word of God and with Christian liberty.\footnote{This addition, which is not found in any copy, is thus stated by Hagenbach and Niemeyer (\emph{Proleg.} p. xxxvi.): '\emph{Durch diese Artikel wollen wir keineswegs allen Kirchen eine einzige Glaubensregel vorschreiben. Denn wir erkennen keine andere Glaubensregel an als die heilige Schrift. Wer also mit dieser übereinstimmt, mit dem sind wir einstimmig, obgleich er anders von unserer Confession verschiedene Redensarten brauchte. Denn auf die Sache selbst und die Wahrheit, nicht auf die Worte soll man sehen. Wir stellen also jedem frei, diejeniqen Redensarten zu gebrauchen, welche er für seine Kirche am passendsten glaubt, und werden uns auch dergleichen Freiheit bedienen, gegen Verdrehung des wahren Sinnes dieser Confession uns aber zu vertheidigen wissen. Dieser Ansdrücke haben wir uns jetzt bedient, um unsere Ueberzeugung darzustellen.}' Pestalozzi. p. 186, gives the same declaration more fully.}

\xsubsection{The Second Helvetic Confession. A.D. 1566.}

§ 55. The Second Helvetic Confession, A.D. 1566.

\emph{Literature}

Confessio Helvetica Posterior. The Latin text, Zurich, 1566, 1568, 1608, 1651, etc.; recent editions by \emph{J. P. Kindler,} with Preface of Winer, Sulzbach, 1825; by \emph{Fritzsche,} Turici, 1839; and by \emph{Ed. Böhl,} Vienna, 1866: also in the Collections of \emph{Corpus et Syntag. Confess.,} Oxford \emph{Sylloge, Augusti,} and \emph{Niemeyer.} The German text appeared frequently—Zurich, 1566; Basle, 1654; Berne, 1676, etc., and in the Collections of \emph{Beck, Mess,} and \emph{Böckel.} French ed. Geneva, 1566, etc. English translations in Hall's \emph{Harmony of Protestant Confessions} (not complete); another by Owen Jones:  \emph{The Church of the Living God}; \emph{also the Swiss and Belgian Confessions and Expositions of the Faith, translated into the English language in} 1862. London (Caryl Book Society), 1865 (complete, but inaccurate), and a third by Prof. Jerem. H. Good (of Tiffin, O.) in Bomberger's \emph{Reformed Church Monthly} (Ursinus College, Pa.), for Sept. 1872, to Dec. 1873 (good, but made from the German translation).

Joh. Jak. Hottinger:  \emph{Helvetische Kirchengeschichte,} Zurich, 1708, Part III. pp. 894 sqq.

Hagenbach:  \emph{Kritische Geschichte der Entstehung und Schicksale der ersten Basler Confession.} Basel, 1827 (1828), pp. 85 sqq.

Niemeyer:  \emph{Collect.,} Prolegomena, pp. lxiii.-lxviii.

L. Thomas:  \emph{La Confession Helvétique, études historico-dogmatiques sur le xvi\textsuperscript{e}. siècle.} Genève, 1853.

K. Sudhoff:  Art. \emph{Helvetische Confession,} in Herzog's \emph{Theol. Encyklop.} 2d ed. Vol. V. pp. 749–755.

Carl Pestalozzi:  \emph{Heinrich Bullinger. Leben und ausgewählte Schriften. Nach handschriftlichen und gleichzeitigen Quellen.} Elberfeld, 1858 (5th Part of \emph{Väter und Begründer der reform. Kirche}), pp. 413–421.

Before we proceed to the Calvinistic Confessions, we anticipate the Second Helvetic Confession, the last and the best of the Zwinglian family.

BULLINGER.

It is the work of Henry Bullinger (1504–1575), the pupil, friend, and successor of Zwingli, to whom he stands related as Beza does to Calvin. He was a learned, pious, wise, and faithful man, and the central figure in the second period of the Reformation in German Switzerland. Born at Bremgarten, in Aargau,\footnote{He was one of five sons of Dean Bullinger, who, like many priests of those days, in open violation of the laws of celibacy, lived in regular wedlock, but was much respected and beloved even by his bishop of Constance. He opposed Samson's traffic in indulgences, and became afterwards a Protestant through the influence of his son.} educated in Holland and Cologne, where he studied patristic and scholastic theology, and read with great interest the writings of Luther and the \emph{Loci} of Melanchthon, he became on his return intimately acquainted with Zwingli, accompanied him to the Conference at Berne (1528), and after laboring for some years at Cappel and Bremgarten, he was chosen his successor as chief pastor (Antistes) at Zurich, Dec. 9, 1531. This was shortly after the catastrophe at Cappel, in the darkest period of the Swiss Reformation.

Bullinger proved to be the right man in the right place. He raised the desponding spirits, preserved and completed the work of his predecessor, and exerted, by his example and writings, a commanding influence throughout the Reformed Church inferior only to that of Calvin. He was in friendly correspondence with Calvin, Bucer, Melanchthon, Laski, Beza, Cranmer, Hooper,\footnote{Bishop Hooper wrote from prison shortly before his martyrdom, May and December, 1554, to Bullinger, as 'his revered father and guide,' and the best friend he had ever found, and commended to him his wife and two children (Pestalozzi, l.c. p. 445).} Lady Jane Grey,\footnote{Three letters of this singularly accomplished and pious lady, the great-granddaughter of Henry VII., to Bullinger, full of affection and gratitude, are still preserved as jewels in the City Library of Zurich, but his letters to her are lost. She translated a part of his book on Christian marriage into Greek, and asked his advice about learning Hebrew. Edward VI., against the will of Henry VIII., bequeathed his crown to Lady Jane Grey to save the Protestant religion, and this led to her execution at the Tower of London, Feb. 12, 1554, by order of Queen Mary. She met her fate with the spirit of a martyr, and sent, as a last token of friendship, her gloves to Bullinger, which were long preserved in his family (Pestalozzi, l.c. p. 445).} and the leading Protestant divines and dignitaries of England. Some of them had found an hospitable refuge in his house and with his friends during the bloody reign of Mary (1553–58), and after their return, when raised to bishoprics and other positions of influence under Queen Elizabeth, they asked his counsel, and kept him informed about the progress of reform in their country. This correspondence is an interesting testimony not only to his personal worth, but also to the fraternal communion which then existed between the Anglican and the Swiss Reformed Churches.\footnote{See the \emph{Zurich Letters,} published by 'The Parker Society,' Cambridge, second edition (chronologically arranged in one series), 1846. They contain, mostly from the archives of Zurich (the Simmler Collection), Geneva, and Berne, letters of Bishops John Jewel, John Parkhurst, Edmund Grindal, Edwin Sandys, Horn, John Foxe, Sir A. Cook, and others to Bullinger, as also to Gualter (Zwingli's son-in-law), Peter Martyr, Simmler, Lavater, Calvin, and Beza. The news of Bullinger's death was received in England with great grief. W. Barlow wrote to J. Simmler (Bullinger's son-in-law), March 13, 1576 (p. 494): ' How great a loss your Church has sustained by the death of the elder Bullinger, of most happy memory, yea, and our Church also, towards which I have heard that he always entertained a truly fraternal and affectionate regard, and indeed all the Churches of Christ throughout Europe.' Bishop Cox wrote to Gualter in the same year (p. 496): 'My sorrow was excessive for the death of Henry Bullinger, whom, by his letters and learned and pious writings, I had . . . known intimately for many years, although he was never known personally to me. Who would not be made sorrowful by the loss of such and so great a man, and so excellent a friend? not to mention that the whole Christian Church is disquieted with exceeding regret that so bright a star is forbidden any longer to shine upon earth.'} Episcopacy was then not yet deemed the only valid form of the Christian ministry. He opened his house also to Italian Protestants, and treated even the elder Sozino, who died at Zurich, with great kindness and liberality, endeavoring to restrain his heretical tendency. In the latter years of his life he was severely tried by the death of his best friends (Bibliander, Froschauer, Peter Martyr, Pellican, Gessner, Blaarer, Calvin, Hyperius), and by a fearful pestilence which deprived him of his beloved wife and three daughters, and brought him to the brink of the grave. He bore all with Christian resignation, recovered from disease, and continued faithfully to labor for several years longer, until he was called to his reward, after taking affectionate farewell of all the pastors and professors of Zurich, thanking them for their devotion, assuring them of his love, and giving each one of them the hand with his blessing. He assumed the care of the Church of Zurich when it was in a dangerous crisis; he left it firmly and safely established.

COMPOSITION.

Bullinger was one of the principal authors of the First Helvetic Confession, and the sole author of the Second. In the intervening thirty years Calvin had developed his amazing energy, while Romanism had formularized its dogmas in the Council of Trent.

Bullinger composed the Second Helvetic Confession in 1562, in latin, for his own use, as an abiding testimony of the faith in which he had lived and in which he wished to die. He showed it to Peter Martyr, who fully consented to it, shortly before his death (Nov. 12, 1562). Two years afterwards lie elaborated it more fully during the raging of the pestilence, and added it to his will, which was to be delivered to the magistrate of Zurich after his death, which he then expected every day.\footnote{See Bullinger's notes to the list of his writings; J. H. Hottinger, \emph{Schola Tigurina,} p. 76; J. J. Simmler, \emph{Oratio de historia Confessionis Helveticæ,} in Simmler's Collection, as quoted by Pestalozzi, l.c. pp. 416 sq. and 641. Also J. J. Hottinger, \emph{Helvet. Kirchengesch.} Pt. III. p. 894.}

PUBLICATION.

But events in Germany gave it a public character. The pious Elector of the Palatinate, Frederick III., being threatened by the Lutherans with exclusion from the treaty of peace on account of his secession to the Reformed Church and publication of the Heidelberg Catechism (1563), requested Bullinger (1565) to prepare a clear and full exposition of the Reformed faith, that he might answer the charges of heresy and dissension so constantly brought against the same. Bullinger sent him a manuscript copy of his Confession. The Elector was so much pleased with it that he desired to have it translated and published in Latin and German before the meeting of the Imperial Diet, which was to assemble at Augsburg in 1566, to act on his alleged apostasy.

In the mean time the Swiss felt the need of such a Confession as a closer bond of union. The First Helvetic Confession was deemed too short, and the Zurich Confession of 1545, the Zurich Consensus of 1549, and the Geneva Consensus of 1552 touched only the articles of the Lord's Supper and predestination. Conferences were held, and Beza came in person to Zurich to take part in the work. Bullinger freely consented to a few changes, and prepared also the German version. Geneva, Berne, Schaffhausen, Biel, the Grisons, St. Gall, and Muhlhausen expressed their agreement. Basle alone, which had its own Confession, declined for a long time, but ultimately acceded.

The new Confession appeared at Zurich, March 12, 1566, in both languages, at public expense, and was forwarded to the Elector and to Philip of Hesse.\footnote{The full title is: '\emph{Confessio et Expositio simplex Orthodoxæ Fidei, et Dogmatum Catholicorum synceræ Religionis Christianæ. Concorditer ab Ecclesiæ Christi Ministris, qui sunt in Helvetia, Tiguri, Bernæ} [\emph{Glaronæ, Basileæ}], \emph{Scaphusii} [\emph{Abbatiscellæ}], \emph{Sangalli, Curiæ Rhetorum, et apud Confœderatos, Mylhusii item, et Biennæ: quibus adjunxerunt se et Genevensis} [\emph{et Neocomensis}] \emph{Ecclesiæ Ministri} [\emph{una cum aliis Evangelii Præconibus in Polonia, Hungaria, et Scotia}]; \emph{edita in hoc, ut universis testentur fidelibus, quod in unitate veræ et antiquæ Christi Ecclesiæ perstent, neque ulla nova, aut erronea dogmata spargant, atque ideo etiam nihil consortii cum ullis Sectis aut Hæresibus habeant. Ad Rom. cap. X. vers.} 10. \emph{Corde creditur ad justitiam, ore autem confessio fit ad salutem. Tiguri: Excudebat Christophorus Froschoverus, Mense Martio, MDLXVI.}' Glarus, Basle, Appenzell, Neufchatel, Poland, Hungary, and Scotland, which we have included in brackets, approved the Confession at a later period, and hence are not mentioned in the first edition, but partly in the second edition of 1568, and more fully in those of 1644 and 1651.} A French translation appeared soon afterwards in Geneva under the care of Beza.

In the same month the Elector Frederick made such a manly and noble defense of his faith before the Diet at Augsburg, that even his Lutheran opponents were filled with admiration for his piety, and thought no longer of impeaching him for heresy.

AUTHORITY.

The Helvetic Confession is the most widely adopted, and hence the most authoritative of all the Continental Reformed symbols, with the exception of the Heidelberg Catechism. Besides the Swiss Cantons and the Palatinate, in whose name it was first issued, the Reformed Churches of Neufchatel (1568), Basle, France (at the Synod of La Rochelle, 1571), Hungary (at the Synod of Debreczin, 1567), Poland (1571 and 1578), and Scotland (1566)\footnote{The ministers of Scotland wrote to Beza, September, 1566: '\emph{Subscripsimus omnes, qui in hoc cœtu interfuimus, et hujus Academiæ sigillo publico obsignavimus.}' This is stated after the Preface in the edition of the \emph{Corpus et Syntagma,} and in Niemeyer, p. 465, but without naming the \emph{cœtus} and \emph{Academia.}} gave it their sanction. It was well received also in Holland and England.\footnote{I find no evidence of a formal sanction by the Anglican Church; but that the Confession was well received there may be inferred from the high esteem in which Bullinger was held (see p. 391), and still more from the fact that his \emph{Decades} (a popular compend of theology in five series of sermons, each containing ten sermons) were, next to Calvin's \emph{Institutes,} the highest theological authority in England, and were recommended, as late as 1586, to the study of young curates along with the Bible. See Ch. Hardwick: \emph{A History of the Christian Church during the Reformation} (third edition, London, 1873, p. 241), where the following order of the Southern Convocation is quoted from Wilkins, IV. 321: 'Every minister having cure, and being under the degrees of master of arts and bachelor of law, and not licensed to be a public preacher, shall, before the second day of February next, provide a Bible, and \emph{Bullinger's Decades in Latin and English,} and a paper book,' etc. On Bullinger's \emph{Decades,} and his abridgment of the same in the \emph{Handbook of the Christian Religion} (1556), see Pestalozzi, pp. 386, 469, 505 sqq.}

It was translated not only into German, French, and English, but also into Dutch, Magyar, Polish, Italian, Arabic, and Turkish.\footnote{See Niemeyer, \emph{Proleg.} p. lxvii. sq.}

CHARACTER AND VALUE.

Like most of the Confessions of the sixteenth century, the Helvetic Confession is expanded beyond the limits of a popular creed into a lengthy theological treatise. It is the matured fruit of the preceding symbolical labors of Bullinger and the Swiss Churches. It is in substance a restatement of the First Helvetic Confession, in the same order of topics, but with great improvements in matter and form. It is scriptural and catholic, wise and judicious, full and elaborate, yet simple and clear, uncompromising towards the errors of Rome, moderate in its dissent from the Lutheran dogmas. It proceeds on the conviction that the Reformed faith is in harmony with the true Catholic faith of all ages, especially the ancient Greek and Latin Church.

Hence it is preceded by the Imperial edict of 380 (from the recognized Justinian code), which draws the line between orthodoxy and heresy, and excludes as heresies only the departures from the Apostolic and Nicene faith. It inserts also the brief Trinitarian creed ascribed to the Roman Pope Damasus (from the writings of Jerome), and referred to in said decree as a standard of orthodoxy.\footnote{Several creeds bear the name of Damasus, and are given by Hahn, \emph{Bibliothek der Symbole,} pp. 179–190. The form inserted in the Confession is from a letter to Jerome (\emph{Opera,} ed. Vallarsi, Tom. XI. p. 145), and is thus referred to in the Imperial edict: '\emph{Cunctos populos . . . in ea volumus religione versari quam divinum Petrum Apostolum tradidisse Romanis . . . quamque } Pontificem Damasum  \emph{sequi claret, et Petrum Alexandriæ Episcopum, virum Apostolicæ sanctitatis.}'} Harmony in the fundamental doctrines of the ancient Church is declared sufficient, and brotherly union consistent with variety in unessentials, such as in fact always has existed in the Christian Church. As in former Confessions, so also in this, Bullinger distinctly recognizes, in the spirit of Christian liberty and progress, the constant growth in the knowledge of the Word of God, and the consequent right of improvement in symbolical statements of the Christian faith.

Upon the whole, the Second Helvetic Confession, as to theological merit, occupies the first rank among the Reformed Confessions, while in practical usefulness it is surpassed by the Heidelberg and Westminster Shorter Catechisms, and in logical clearness and precision by the Westminster Confession, which is the product of a later age, and of the combined learning and wisdom of English and Scotch Calvinism.\footnote{I add some testimonies on the Second Helvetic Confession. Hagenbach (l.c. p. 86): '\emph{In ihrer ganzen Anlage and in der Durchführung einzelner Punkte, namentlich in praktischer Beziehung} (\emph{in der Scheidung des Geistlichen and Weltlichen, u.s.w.}) \emph{ist sie ein wahres dogmatisches Kunstwerk zu nennen.}' Pestalozzi (\emph{Bullinger,} p. 422): '\emph{Diese Confession, zu der Bullinger zweimal Angesichts des Todes sich bekannte, erscheint als das reife Ergebniss seines Glaubenslebens, seiner reichen inneren und äusseren Erfahrung, als der Inbegriff seiner theologischen Ueberzeugung wie seiner kirchlichen Grundsätze, als die ächte, wahrhafte Entwicklung und Fortbildung seiner früheren Bekenntnisse, zumal der ersten helvetischen Confession} (\emph{von} 1536). \emph{Sie ist ein Muster von Klarheit und Einfachheit, wie selbst hervorragende Gegner anerkennen, ausgezeichnet durch den Ueberblick, der das Ganze der christlichen Lehre umfasst, der völlige Ausdruck von Bullingers Gesinnung, scharf ausgeprägt gegenüber den Verirrungen des römisch-katholischen Kirchenthums, milde in Bezug auf die lutherischen Besonderheiten, ohne doch der eigenen Ueberzeugung irgend Eintrag zu thun. Was aber vornehmlich beachtenswerth, sie ist durchaus getragen von dem vollen, klaren und ruhigen Bewusstsein, das mit so durchgreifender Kräftigkeit Bullinger beseelte, der ächten apostolischen und katholischen Kirche anzugehören, der wahrhaft berechtigten und rechtgläubigen Kirche Christi. Sie ist fern davon, bloss mit der Bibel in der Hand alles das zu verwerfen, was nicht ausdrücklich in der heiligen Schrift gelehrt und geboten ist, wiewohl ihr diese von höchster Geltung ist, als oberste Richtschnur der christlichen Wahrheit. Sie bricht nicht mit dem geschichtlich Gewordenen} (\emph{der Ueberlieferung}), \emph{ausser sofern dieses der Schrift nicht gemäss ist. Die ganze Entwicklung der christlichen Kirche seit den Tagen der Apostel bis auf die Gegenwart ist ihr von hohem Werthe und findet ihre ernste Berücksichtigung, nur dass sie sich nach der obersten Norm muss richten lassen. Insofern steht sie mit ihrer evangelischen Schwesterkirche lutherischen Bekenntnisses ganz auf demselben Boden und kann ihr stets die Hand reichen zur Annäherung, möglicher Weise auch zu einer Einigung, wenn gleich die Auffassung der christlichen Wahrheit nach gewissen Richtungen hin sich unterscheiden und deshalb die Entscheidung über diese oder jene einzelnen Lehrpunkte und Gebräuche verschieden ausfallen mag.}' Dr. Hodge (\emph{Syst. Theol.} Vol. III. p. 634): 'The Second Helvetic Confession is, on some accounts, to be regarded as the most authoritative symbol of the Reformed Church, as it was more generally received than any other, and was sanctioned by different parties.'}

CONTENTS.

In view of the importance of this Confession, I give here a condensed translation of the original.\footnote{The full Latin text will be found in Vol. III.} It consists of thirty chapters, which cover in natural order all the articles of faith and discipline which then challenged the attention of the Church.

Chap. I. The Holy Scriptures.—This chapter lays down the evangelical rule of faith, or the objective principle of Protestantism.

We believe and confess that the Canonical Scriptures of the Old and New Testaments are the true Word of God, and have sufficient authority in and of themselves, and not from men; since God himself through them still speaks to us, as he did to the Fathers, the Prophets, and Apostles. They contain all that is necessary to a saving faith and a holy life; and hence nothing should be added to or taken from them (Deut. iv. 2;  Rev. xxii. 18, 19).

From the Scriptures must be derived all true wisdom and piety, and also the reformation and government of the Churches, the proof of doctrines, and the refutation of errors (2 Tim. iii. 16, 17; 1 Tim. iii. 15;  1 Thess. ii. 13;  Matt. x. 20). God may illuminate men directly by the Holy Spirit, without the external ministry; yet he has chosen the Scriptures and the preaching of the Word as the usual method of instruction.

The apocryphal books of the Old Testament, though they may be read for edification, are not to be used as an authority in matters of faith.\footnote{This is the first symbolical exclusion of the Apocrypha from the Canon. The Lutheran symbols leave this question open.}

We condemn the doctrines of the Gnostics and Manichæans, and all others who reject the Scriptures in whole or in part.

Chap. II. The Interpretation of the Scriptures; the Fathers, Councils, and Traditions.—We acknowledge only that interpretation as true and correct which is fairly derived from the spirit and language of the Scriptures themselves, in accordance with the circumstances, and in harmony with other and plainer passages (2 Pet. i. 20, 21).

We do not despise the interpretation of the Greek and Latin fathers and the teaching of Councils, but subordinate them to the Scriptures; honoring them as far as they agree with the Scriptures, and modestly dissenting from them when they go beyond or against the Scriptures. In matters of faith we can not admit any other judge than God himself, who through his Word tells us what is true and what is false, what is to be followed, and what is to be avoided.

We reject traditions which contradict the Scriptures, though they may claim to be apostolical. For the Apostles and their disciples could not teach one thing by writing, and another by word of mouth. St. Paul preached the same doctrine to all the churches (1 Cor. iv. 17; 2 Cor. i. 13; xii. 18). The Jews likewise had their traditions of the elders, but they were refuted by our Lord as 'making void the Word of God' (Matt. xv. 8, 9; Mark vii. 6, 7).

Chap III. Of God, his Unity and Trinity.—We believe and teach that God is one in essence (Deut. vi. 4;  Exod. xx. 2, 3, etc.), and three in persons—Father, Son, and Holy Ghost. The Father hath begotten the Son from eternity; the Son is begotten in an unspeakable manner; the Holy Ghost eternally proceeds from both, and is to be worshiped with both as one God. There are not three Gods, but three persons—consubstantial, coeternal, distinct as to person and order, yet without any inequality. The divine essence or nature is the same in the Father, the Son, and the Spirit (Luke i. 35; Matt. iii. 17; xxviii. 19; John i. 32; xiv. 26; xv. 26).

In short, we accept the Apostles' Creed, which delivers to us the true faith.

We therefore condemn the Jews and Mohammedans, and all who blaspheme this holy and adorable Trinity. We also condemn all heretics, who deny the Deity of Christ and the Holy Ghost.

Chap. IV Of Idols, Images of Gods and of Saints.—As God is a spirit, he can not be represented by any image (John iv. 24; Isa. xl. 18; xliv. 9, 10; Jer. xvi. 19; Acts xvii. 29, etc.).

And although Christ assumed man's nature, yet he did so not in order to afford a model for sculptors and painters. He instituted for the instruction of the people the preaching of the Gospel, and the sacraments, but not images. Epiphanius tore down an image of Christ and some saint in a church, because he regarded it contrary to the Scriptures.

Chap. V. The Adoration and Invocation of God through the only Mediator Jesus Christ.—God is the only object of worship (Matt. iv. 10). And he is to be worshiped 'in spirit and in truth' (John iv. 24), and through our only and sufficient Mediator and Advocate Jesus Christ (1 Tim. ii. 5; 1 John ii. 1).

Hence we neither adore nor invoke the departed saints, and give no one else the glory that belongs to God alone (Isa. xlii. 8; Acts iv, 12).

Nevertheless, we neither despise nor undervalue the saints, but honor them as the members of Christ and the friends of God, who have gloriously overcome the flesh and the world; we love them as brethren, and hold them up as examples of faith and virtue, desiring to dwell with them eternally in heaven, and to rejoice with them in Christ.

Much less do we believe that the relics of saints should be worshiped. Nor do we swear by saints, since it is forbidden to swear by the name of strange gods (Exod. xxiii. 13; Deut. x. 20).

Chap. VI. The Providence of God.—We believe that the wise, eternal, and almighty God by his providence preserves and rules all things in heaven and earth (Psa. cxiii. 4–6; cxxxix. 3–4; Acts xvii. 28; Rom. xi. 36).

We therefore condemn the Epicureans, who blasphemously affirm that God neither sees nor cares for men (Psa. xciv. 3–9).

We do not despise as unnecessary the means whereby divine Providence works, but make use of them as far as they are commended to us in the Word of God. We disapprove of the rash words of those who say that our efforts and endeavors are vain.

St. Paul well knew that he was sailing under the providence of God, who had assured him that he must bear witness at Rome (Acts xxiii. 11), and that not a soul should perish (xxvii. 21, 34); nevertheless, when the sailors were seeking flight, he said to the centurion and the soldiers: 'Unless these abide in the ship, ye can not be saved' (ver. 31). For God has appointed the means by which we attain to the end.\footnote{\par Here we have a clear recognition of secondary causes in opposition to fatalism and determinism which has sometimes been charged upon Calvinism. The Westminster Confession (Chap. III.) is still more explicit: 'God from all eternity did by the most wise and holy counsel of his own will freely and unchangeably ordain whatsoever comes to pass; \emph{yet so as thereby neither is God the author of sin; nor is violence offered to the will of the creatures, nor is the liberty or contingency of second causes taken away, but rather established} (Acts ii. 23; iv. 27, 28; xvii. 23, 24, comp. with 36; Matt. xvii. 12; John xix. 11; Prov. xvi. 33).}

Chap. VII. Of the Creation of all Things; of Angels, the Devil, and Man.—This good and almighty God created all things, visible and invisible, by his eternal Word, and preserves them by his coeternal Spirit (Psa. xxxiii. 6; John i. 3). He made all things very good and for the use of man (Gen. i. 31).

We condemn the Manichæans who impiously imagine two coeternal principles, the one good, the other evil, and two antagonistic gods.

Angels and men stand at the head of all creatures. Angels are ministers of God (Psa. civ. 4), and ministering spirits sent for them who shall be heirs of salvation (Heb. i. 14).

The devil was a murderer and liar from the beginning (John viii. 44).

Some angels persevered in obedience, and are ordained unto the faithful service of God and men; but others fell of their own accord and ran into destruction, and have become enemies of God and men.

Man was made in the image and likeness of God, and placed by God in paradise as ruler over the earth (Gen. i. 27, 28; ii. 8). This is celebrated by David in the 8th Psalm. Moreover, God gave him a wife and blessed them (Gen. ii. 22 sqq.).

Man consists of two diverse substances in one person—of an immortal soul, which, when separated from the body neither sleeps nor dies, and of a mortal body, which at the last judgment, shall be raised again from the dead.

We condemn those who deny the immortality, or affirm the sleep of the soul, or teach that it is a part of God.

Chap. VIII. Of Man's Fall, Sin, and the Cause of Sin.—Man was created according to the image of God, in true righteousness and holiness, good and upright. But by the instigation of the serpent, and through his own guilt, he fell from goodness and rectitude, and became, with all his offspring, subject to sin, death, and various calamities.

Sin is that inborn corruption of man, derived and propagated from our first parents, whereby we are immersed in depraved lusts, averse to goodness and prone to all evil, and unable of ourselves to do or think any thing that is good. And as years roll on, we bring forth evil thoughts, words, and deeds, as corrupt trees bring forth corrupt fruits (Matt. xii. 33). Therefore we are all by nature under the wrath of God, and subject to just punishment.

By death we understand not only the dissolution of the body, but also the eternal punishments of sin (Eph. ii. 1, 5; Rom. v. 12).

We therefore acknowledge that there is original sin in all men, and that all other sins, whether mortal or venial, also the unpardonable sin against the Holy Ghost, spring from this same source. We acknowledge also that sins are not equal, but some are more grievous than others (Matt. x. 14, 15; xi. 24; 1 John v. 16, 17).

We condemn the Pelagians, who deny original sin; the Jovinianists, who with the Stoics declare all sins to be equal; and those who make God the author of sin against the express teaching of Scriptures (Psa. v. 5–7; John viii. 44).

When God is said to blind or harden men, or to give them over to a reprobate mind (Exod. vii. 13; John xii. 40), it is to be understood as a righteous judgment. Moreover, God overrules the wickedness of men for good, as he did in the case of the brethren of Joseph.

Chap. IX. Of Free Will and Man's Ability.—The will and moral ability of man must be viewed under a threefold state.

First, before the fall, he had freedom to continue in goodness, or to yield to temptation.

Secondly, after the fall, his understanding was darkened and his will became a slave to sin (1 Cor. ii. 14; 2 Cor. iii. 5; John viii. 34; Rom. viii. 7). But he has not been turned into 'a stone or stock;' nor is his will (\emph{voluntas}) a non-will (\emph{noluntas}).\footnote{\par Expressions used by Luther, Flacius, and the Formula of Concord. The Helvetic and other Reformed Confessions are much more guarded on this point, and teach that man, though totally depraved, remains a moral and responsible being in the act of sinning. Melanchthon, in his later period, came to the same view, but went beyond it into synergism. Comp. above, pp. 262, 270.} He serves sin willingly, not unwillingly (\emph{servit peccato non nolens, sed volens}). In external and worldly matters man retains his freedom even after the fall, under the general providence of God.

Thirdly, in the regenerate state, man is free in the true and proper sense of the term. His intellect is enlightened by the Holy Spirit to understand the mysteries and the will of God; and the will is changed by the Spirit and endowed with the power freely to will and to do what is good (Rom. viii. 5, 6; Jer. xxxi. 33; Ezek. xxxvi. 26; John viii. 36; Phil. i. 6, 29; ii. 13).

In regeneration and conversion men are not merely passive, but also active. They are moved by the Spirit of God to do of themselves what they do. But even in the regenerate there remains some infirmity. The flesh strives against the spirit to the end of life (Rom. vii. 14; Gal. v. 17).

We condemn the Manichæans, who deny that evil originated in the free will of man, and the Pelagians, who teach that fallen man has sufficient freedom to keep God's commandments. The former are refuted by Gen. i. 27; Eccles. vii. 29; the latter by John viii. 36.

Chap. X. The Predestination of God and the Election of Saints.—God has from eternity predestinated or freely chosen, of his mere grace, without any respect of men, the saints whom he will save in Christ (Eph. i. 4; 2 Tim. i. 9, 10).

God elected us in Christ and for Christ's sake, so that those who are already implanted in Christ by faith are chosen, but those out of Christ are rejected (2 Cor. xiii. 5).\footnote{\par '\emph{Ergo non sine medio, licet non propter ullum meritum nostrum, sed in Christo et propter Christum nos elegit Deus, ut qui jam sunt in Christo insiti per fidem, illi ipsi etiam sint electi, reprobi vero, qui sunt extra Christum.}'}

Although God knows who are his, and a 'small number of the elect' is spoken of, yet we ought to hope well of all, and not rashly count any one among the reprobate (2 Tim. ii. 19; Matt. xx. 16; Phil i. 3 sqq.).

We reject those who seek out of Christ whether they are chosen, and what God has decreed concerning them from eternity. We are to hear the gospel and believe it, and be sure that if we believe and are in Christ, we are chosen. We must listen to the Lord's invitation, 'Come unto me' (Matt. xi. 28), and believe in the unbounded love of God, who gave his own Son for the salvation of the world, and will not that 'one of these little ones should perish' (John iii. 16;  Matt. xviii. 14).\footnote{Comp. ver. 10 and 11. A very strong passage for the doctrine of infant salvation, and so understood by Zwingli and Bullinger.}

Let, therefore, Christ be the mirror in which we behold our predestination. We shall have a sufficiently evident and sure testimony of being written in the book of life if we live in communion with him, and if in true faith he is ours and we his.

And if we are tempted concerning our predestination, let this be our comfort—that God's promises are general to believers, as he himself says: 'Seek, and ye shall find, and whosoever asketh shall receive' (Matt. vii. 8 sq.). We pray with the whole Church, 'Our Father which art in heaven;' by baptism we are ingrafted into the body of Christ, and we are often fed in the Church by his flesh and blood unto life everlasting. Thus strengthened, let us 'work out our own salvation with fear and trembling, for it, is God who worketh in us both to will and to do according to his good pleasure' (Phil. ii. 12, 13).\footnote{This Tenth Article is moderately Calvinistic or Augustinian, and neither Arminian nor Melanchthonian (synergistic), as has sometimes been claimed. Comp. Schweizer, \emph{Centraldogmen,} Vol. I. p. 476; also Sudhof's art. in Herzog.}

Chap. XI. Jesus Christ true God and Man, and the only Saviour of the World.—We believe and teach that the Son of God, our Lord Jesus Christ, was from eternity predestinated by the Father to be the Saviour of the world; that he was begotten of the Father from all eternity in an ineffable manner (Isa. liii. 8; Micah v. 2; John i. 1). Therefore the Son, according to his Divinity, is coequal and consubstantial with the Father; true God, not merely by name or adoption or by conferring of a dignity, but in essence and nature (1 John v. 20; Phil. ii. 6;  Heb. i. 2, 3; John v. 18; xvii. 5).

We abhor the blasphemous doctrine of Arius and Servetus in opposition to the Son of God.

We also believe and teach that the same eternal Son of God became the Son of Man, of the seed of Abraham and David, not through the will of man (Ebionites), but he was conceived by the Holy Ghost and born of the ever-Virgin Mary (\emph{ex Maria semper virgine}), as taught in the gospel history and the Epistles (Matt. i. 18; Luke i. 34, 35; 1 John iv. 3; Heb. ii. 16). The body of Christ was therefore neither a mere appearance, nor brought down from heaven (the Gnostics, Valentinus and Marcion). Moreover his soul was not without reason (Apollinaris), nor his flesh without a soul (Eunomius); but he had a rational soul, and a flesh with senses capable of true suffering (Matt. xxvi. 36; John xii. 27).

Hence we acknowledge in one and the same Lord Jesus Christ two natures, a divine and a human, which are conjoined and united in one person without absorption or confusion and mixture.

We worship one Lord Christ, not two; one true God-Man, coequal (or of one substance, \emph{consubstantialis,} \textgreek{ὁμοούσιος}) with the Father as regards his divine nature, and coequal with us men, sin only excepted (Heb. iv. 15), as regards his human nature.

We therefore abominate Nestorianism, which dissolves the unity of person, and Eutychianism, Monothelitism, and Monophysitism, which destroy the proper character of the human nature.

We do not teach that the divine nature of Christ did suffer, nor that the human nature of Christ is every where present. The true body of Christ was not deified so as to put off its properties and to be absorbed into the divine substance. But we believe that our Lord Jesus Christ did truly suffer for us in the flesh (1 Pet. iii. 18; iv. 1), and that the Lord of glory was crucified for us (1 Cor. ii. 8). For we accept believingly and reverently the 'communication of properties,' which is deduced from the Scriptures and employed by the ancient Church in explaining and harmonizing seemingly contradictory passages.\footnote{'\emph{Nam communicationem idiomatum ex Scripturis petitam et ab universa vetustate in explicandis componendisque Scripturarum locis in speciem pugnantibus usurpatam, religiose et reverenter recipimus et usurpamus.}' It is an error, therefore, to charge the Reformed Church with rejecting the \emph{communicatio idiomatum.} It admits the communication of the properties of one nature to the whole person, but denies the communication of the properties of one nature to the other, viz., the \emph{genus majestaticum,} so called, whereby the infinite attributes of the divine nature (as omnipresence and omnipotence) are ascribed to the human nature, and the \emph{genus tapeinoticon,} whereby the finite attributes of the human nature are ascribed to the divine. Either of these forms leads necessarily to a Eutychian confusion of natures. The Lutheran Church teaches the \emph{genus majestaticum,} as a support to its doctrine of the Eucharist, but rejects the \emph{genus tapeinoticon.}}

We believe and teach that Christ, in the same flesh in which he died, rose from the dead (Luke xxiv. 30), and ascended to the right hand of God in the highest heaven (Eph. iv. 10), which signifies his elevation to the divine majesty and power, but also a definite place (John xiv. 2; Acts iii. 21).

The same Christ will come again to judgment, when the wickedness of the world shall have reached the highest point, and Antichrist corrupted the true religion. He will destroy Antichrist, and judge the quick and the dead (2 Thess. ii. 8; Acts xvii. 51, 52;  1 Thess. iv. 17). The believers will enter into the mansions of the blessed; the unbelievers, with the devil and his angels, will be cast into everlasting torment (Matt. xxv. 41; 2 Tim. ii. 11; 2 Pet. iii. 7).

We reject all who deny the real resurrection; who teach the ultimate salvation of all the godless, and even the devil. We also reject the Jewish dream of a millennium, or golden age on earth, before the last judgment.

We believe and teach that Christ is the only Redeemer of the whole world, in whom all are saved that were saved before the law, under the law, and under the gospel, or will yet be saved to the end of the world (John x. 1, 7; Acts iv. 12; xv. 11; 1 Cor. x. 1, 4; Rev. xiii. 8).

We therefore confess and teach with a loud voice: Jesus Christ is the only Saviour of the world, the King and High-priest, the true Messiah, whom all the shadows and types of the Law and the Prophets did prefigure and promise. God did send him to us, and we need not look for another. There remains nothing but that we should give all glory to him, believe in him, and rest in him alone.

And, to say much in a few words, we sincerely believe and loudly confess all that has been determined out of the Holy Scriptures concerning the mystery of the incarnation of our Lord Jesus Christ, and is contained in the creeds and decrees of the first four œcumenical Councils held in Niceæ, Constantinople, Ephesus, and Chalcedon, in the Creed of St. Athanasius, and all similar creeds; and we reject all contrary to the same. In this manner we retain, unchanged and entire, the Christian, orthodox, and catholic faith; knowing that nothing is contained in the aforesaid creeds which does not correspond with the Word of God and aid in setting forth the true faith.\footnote{An express and emphatic indorsement of the œcumenical Creeds, on the ground of their agreement with the Scriptures: '\emph{Et ut paucis multa hujus causæ dicamus, quæcunque de incarnationis Domini nostri Jesu Christi mysterio definita sunt ex Scripturis sanctis, et comprehensa symbolis ac sententiis quatuor primarum et prœstantissimarum Synodorum celebratarum Niceæ, Constantinopoli, Ephesi, et Chalcedone, una cum beati Athanasii Symbolo, et omnibus his similibus symbolis, credimus corde syncero, et ore libero ingenue profitemur, condemnantes omnia his contraria. Atque ad hunc modum retinemus inviolatam sive integram fidem Christianam, orthodoxam atque catholicam: scientes, symbolis prædictis nihil contineri, quod non sit conforme Verbo Dei, et prorsus faciat ad synceram fidei explicationem.}'}

Chap. XII. The Law of God.—The law of God explains the will of God and the difference between what is good and bad, just and unjust. It is therefore good and holy. It is twofold: the law of nature inscribed on the hearts of men (Rom. ii. 15), and the written law of Moses. The latter we divide for perspicuity's sake into the moral law, comprehended in the two tables of the Decalogue (Exod. xx.; Deut. v.); the ceremonial law, concerning worship and sacred rites; and the judicial, concerning polity and economy.

The law of God is complete, and allows no addition nor subtraction (Deut. iv. 2;  Isa. xxx. 21). It is given to us, not that by keeping it we might be justified, but that we may be led to a knowledge of sin and guilt, and, despairing of our own strength, turn by faith to Christ (Rom. iv. 15; iii. 20; viii. 3;  Gal. iii. 21–24). Christ is the end of the law, and redeemed us from the curse of the law (Rom. x. 4; Gal. iii. 13). He enables us to fulfill the law, and his righteousness and obedience are imputed to us through faith.

The law is abolished inasmuch as it no more condemns and works wrath in them that believe, who are under grace, and not under the law. Besides, Christ has fulfilled all the types of the law, and put the substance in the place of the shadows; in him we have all fullness. Nevertheless, the law is useful in showing us all virtues and vices, and in regulating the life of new obedience. Christ did not come to destroy, but to fulfill the law (Matt. v. 17).

We therefore condemn old and modern Antinomianism.

Chap. XIII. The Gospel of Jesus Christ.—The law works wrath and announces the curse (Rom. iv. 15;  Deut. xxvii. 26); the gospel announces grace and blessing (John i. 17). Nevertheless, those who lived before and under the law were not deprived altogether of the gospel, but had great promises (Gen. iii. 15; xxii. 18; xlix. 10). The promises were partly temporal, partly spiritual and eternal. By the gospel promises the fathers obtained salvation in Christ.

In the strict sense of the term the gospel is the glad tidings of salvation by Christ, in whom we have forgiveness, redemption, and everlasting life. Hence the history of Christ recorded by the four Evangelists is justly called the gospel.

Compared with the legalism of the Pharisees the gospel appeared to be a new doctrine, as it is even now called new by the Papists; but in fact it is the oldest doctrine, for God foreordained from eternity to save the world through Christ, and has revealed this plan in the gospel (2 Tim. i. 9, 10). It is therefore a grave error to call our evangelical faith a recent innovation.

Chap. XIV. Of Repentance and Conversion.—Repentance (\textgreek{μετάνοια}) is a change of heart produced in a sinner by the word of the gospel and the Holy Spirit, and includes a knowledge of native and actual depravity, a godly sorrow and hatred of sin, and a determination to live hereafter in virtue and holiness. True repentance is turning to God and all good, and turning away from the devil and all evil. It is the free gift of God, and not the result of our own strength (2 Tim. ii. 25).

We have examples of true repentance in the woman that was a sinner (Luke vii. 38), in Peter after his fall (xxii. 62), in the prodigal son (xv. 18), and the publican in the temple (xviii. 13).

It is sufficient to confess our sins to God in private and in the public service; it is not necessary to confess to a priest, for this is nowhere commanded in the Scriptures; although we may seek counsel and comfort from a minister of the gospel in time of distress and trial (comp. James v. 16).

The keys of the kingdom of heaven, out of which the Papists forge swords, sceptres, and crowns, are given to all legitimate ministers of the Church in the preaching of the gospel and the maintenance of discipline (Matt. xvi. 19; John xx. 23; Mark xvi. 15;  2 Cor. v. 18, 19). We condemn the profitable popish doctrine of penance and of indulgences, and apply to them Peter's word to Simon Magus: ' Thy money perish with thee' (Acts viii. 20).

Chap. XV. Of True Justification of Believers.—'To justify' means, with the Apostle when treating of this subject, to remit sins, to absolve from guilt and punishment, to receive into grace, and to pronounce just (Rom. viii. 33; Acts xiii. 38; Deut. xxv. 1; Isa. v. 23).

By nature we are all sinners and guilty of death before the tribunal of God, and we can be justified only by the merits of Christ crucified and risen again. For his sake God is reconciled, and imputes to us not our sins, but the righteousness of Christ as our own, so that we are purged and absolved from sin, death and damnation, and heirs of eternal life. Properly speaking, God alone justifies and justifies only for Christ's sake, not imputing to us our sins, but the righteousness of Christ.

We therefore teach and believe, with the Apostle, that the sinner is justified by faith alone in Christ (\emph{sola fide in Christum}), not by the law, nor by any works (Rom. iii. 28; iv. 2 sqq.; Eph. ii. 8, 9). Righteousness is imputed to faith because it receives Christ as our righteousness and ascribes all to the grace of God, but not because it is our work: it is the gift of God. As we receive food by eating, so faith appropriates Christ.

We do not divide justification by ascribing it partly to the grace of God or to Christ, and partly to our works or merits, but solely and exclusively to the grace of God in Christ through faith. We must first be justified before we can do good works. Love is derived from faith (1 Tim. i. 5; Gal. v. 6).

Therefore we speak here not of a false, dead faith, but of a living and vivifying faith which lives in Christ, our life, and proves its life by living works. Even James (chap. ii.) does not contradict our doctrine, for he speaks of a dead faith which even demons have, and he shows that Abraham proved his living and justifying faith by works.

Chap. XVI. Faith and Good Works, their Reward and the Merit of Man.—Christian faith is not a human opinion and persuasion, but a most firm confidence and clear and steady assent of the mind, a most certain apprehension of the truth of God as laid down in the Scriptures and the Apostles' Creed, and therefore of God himself as the highest good, and especially of the divine promise and of Christ, who is the crown of all promises. Such a faith is a free gift of God, who of his grace grants it to his elect through his Holy Spirit by means of the preaching of the gospel and believing prayer when and in what measure he pleases. This faith has degrees and is subject to growth; hence the prayer of the Apostles: 'Lord, increase our faith' (Luke xvii. 5). [Then follow a number of Scripture proofs: Heb. xi. 1; 2 Cor. i. 20; Phil. i. 29; Rom. xii. 3; 2 Thess. ii. 3; Rom. x. 16; Acts xiii. 48; Gal. v. 6, etc.]

We teach that good works proceed from a living faith, through the Holy Spirit, and are done by believers according to the will and rule of the Word of God (2 Pet. i. 5 sqq.; 1 Thess. iv. 3, 6, 23).

Good works must be done, not to merit thereby eternal life, which is a free gift of God (Rom. vi. 23), nor for ostentation or from selfishness, which the Lord rejects (Matt. vi. 2; xxiii. 5), but for the glory of God, to adorn our calling and to show our gratitude to God, and for the good of our neighbor (Matt. v. 16; Eph. iv. 1; Col. iii. 17; Phil. ii. 4; Tit. iii. 14). Although we teach that man is justified by faith of Christ and not by any works, we do not condemn good works. Man is created and regenerated by faith in order to work unceasingly what is good and useful. 'Every good tree bringeth forth good fruit' (Matt. vii. 17). 'He that abideth in me, the same bringeth forth much fruit' (John xv. 5). 'We are God's workmanship, created in Christ Jesus unto good works, which God hath before ordained that we should walk in them' (Eph. ii. 10).

We condemn, therefore, all who despise good works or declare them useless; at the same time we do not deem them necessary to salvation, in the sense that without them no one was ever saved; for we are saved by Christ alone; but good works are necessarily born of faith, and improperly salvation may be ascribed to them which properly is ascribed to grace (Rom. xi. 6).

God is well pleased and approves of works which are done by us through faith (Acts x. 35; Col. i. 9, 10). He also richly rewards them (Jer. xxxi. 16; Matt. v. 12; x. 42). But we ascribe this reward not to the merits of man who receives it, but to the goodness and faithfulness of God who promises and grants it, although he owes nothing to his creatures. Even if we have done all, we are unprofitable servants (Luke xvii. 10). We say with Augustine, that God crowns and rewards in us, not our merits, but the gifts of his grace. It is a reward of grace, not of merit. We have nothing but what we have received (comp. 1 Cor. iv. 7).

We therefore condemn those who so defend the merits of men as to set at naught the grace of God.

Chap. XVII. Of the Catholic and Holy Church of God, and of the only Head of the Church.—Since God willed from the beginning that men should be saved and come to the knowledge of truth, it follows of necessity that there always was, and now is, and shall be to the end of time, a Church or an assembly of believers and a communion of saints, called and gathered from the world, who know and worship the true God in Christ our Saviour, and partake by faith of all the benefits freely offered through Christ. They are fellow-citizens of the same household of God (Eph. ii. 19). To this refers the article in the Creed: 'I believe the holy catholic Church, the communion of saints.'

And as there is but one God, one Mediator between God and man, Jesus the Messiah, one pastor of the whole flock, one head of this body, one Spirit, one salvation, one faith, one testament or covenant, there must needs be but one Church, which we call catholic, that is, universal, spread throughout all parts of the world and all ages.

We therefore condemn the Donatists, who confined the Church to some corners of Africa, and also the Roman exclusiveness, which pretends that the Roman Church alone is the catholic Church.

The Church is divided, not in itself, but on account of the diversity of its members. There is a Church militant on earth struggling against the flesh, the world, and the devil, and a Church triumphant in heaven rejoicing in the presence of the Lord; nevertheless there is a communion between the two. The Church militant is again divided into particular Churches. It was differently constituted among the Patriarchs, then under Moses, then under Christ in the gospel dispensation; but there is only one salvation in the one Messiah, in whom all are united as members of one body, partaking of the same spiritual food and drink. We enjoy a greater degree of light and more perfect liberty.

This Church is called the house of the living God (1 Tim. iii. 15), built of lively and spiritual stones (1 Pet. ii. 5), resting on an immovable rock, the only foundation (1 Cor. iii. 11), the ground and pillar of the truth (1 Tim. iii. 15). It can not err as long as it rests on the rock Christ, on the foundation of the Prophets and Apostles; but it errs as often as it departs from him who is the truth.\footnote{\par '\emph{Non errat illa, quamdiu innititur petræ Christo et fundamento Prophetarum et Apostolorum. Nec mirum, si erret, quoties deserit illum, qui solus est veritas.}'} The Church is also called a virgin, the bride of Christ, the only and beloved (2 Cor. xi. 2), and the body of Christ, because the believers are living members of Christ under him the head (Eph. i. 23, etc.).

The Church can have no other head than Christ. He is the one universal pastor of his flock, and has promised his presence to the end of the world. He needs, therefore, no vicar; for this would imply his absence. [Those who introduce a double headship and government in the Church plainly belong to the errorists condemned by the Apostles (2 Pet. ii.; Acts xx.; 2 Cor. xi.; 2 Thess. ii.).]\footnote{\par The passage in brackets, according to the Zurich MS., was substituted by Bullinger on the margin for the following sentence, which he wished to have canceled (see note in Niemeyer, p. 501): 'We reject the Romish fiction concerning an official head and title of the servant of the servants of Christ; for experience proves that this is an empty boast, and that the Pope makes himself an enemy of Christ, and exalts himself above God, sitting in the temple of God, and showing himself that he is God' (2 Thess. ii. 4).}

But by rejecting the Roman head we do not introduce disorder and confusion into the Church of Christ, since we adhere to the government delivered by the Apostles before there was any Pope. The Roman head preserves the tyranny and corruption in the Church, and opposes and destroys all just reformation.

They object that since our separation from Rome all sorts of controversies and divisions have arisen. As if there had never been any sects and dissensions in the Roman Church, in the pulpits, and among the people! God is indeed a God of order and peace (1 Cor. xiv. 33); nevertheless there were parties and divisions even in the Apostles' Church (Acts xv.; 1 Cor. iii.; Gal. ii.). God overrules these divisions for his glory and for the illustration of truth.

Communion with the true Church of Christ we highly esteem, and deny that those who separate from it can live before God. As there was no salvation out of the ark of Noah, so there is no certain salvation out of Christ, who exhibits himself to the elect in the Church for their nourishment.\footnote{'\emph{Ut extra arcam Noë non erat ulla salus, pereunte mundo in diluvio, ita credimus, extra Christum, qui se electis in Ecclesia fruendum prœbet, nullam esse salutem certam: et proinde docemus, vivere volentes non oportere separari a vera Christi Ecclesia.}' This high estimate of the Church reminds one of Cyprian's '\emph{Extra ecclesiam nulla salus,}' of Tertullian's '\emph{Qui ecclesiam non habet matrem, Deum non habet patrem,}' and of Augustine's '\emph{Ego evangelio non crederem, nisi me commoveret ecclesiæ auctoritas.}' Calvin, in his \emph{Institutes} (lib. IV. c. 1), uses similar language. But we must remember that the Calvinistic system does not bind election to the visible means of grace, and admits the possibility of salvation without baptism. Bullinger denies only the \emph{certainty} of salvation (\emph{salutem certam}) outside of the Church (comp. above what follows); and so must be understood the Westminster Confession of Faith, Ch. XXV. 2, when it asserts that out of the visible catholic or universal Church 'there is no \emph{ordinary} possibility of salvation.'}

But we do not so restrict the Church as to exclude those who from unavoidable necessity and unwillingly do not partake of the sacraments, or who are weak in faith, or still have defects and errors. God had friends even outside of the Jewish people. We know what happened to Peter, and to chosen believers from day to day, and we know that the Apostle censured the Christians in Galatia and Corinth for grave offenses, and yet calls them holy churches of Christ. Yea, God may at times by a righteous judgment allow the Church to be so obscured and shaken as to appear almost annihilated, as in the days of Elijah (1 Kings xix. 18; comp. Rev. vii. 4, 9); but even then he has his true worshipers, even seven thousand and more; for 'the foundation of God standeth sure, having this seal, the Lord knoweth them that are his' (2 Tim. ii. 19). Hence the Church may be called \emph{invisible,} not that the men composing it are invisible, but because they are known only to God, while we are often mistaken in our judgment. There are also many hypocrites in the Church, who outwardly conform to all the ordinances, but will ultimately be revealed in their true character and be cut off (1 John ii. 19; Matt. xiii. 24, 47).

The true unity of the Church is not to be sought in ceremonies and rites, but in the truth and in the catholic faith, as laid down in the Scriptures and summed up in the Apostles' Creed. Among the ancients there was a great diversity of rites without dissolving the unity of the Church.

Chap. XVIII. On the Ministers of the Church, their Institution and Offices.—God always used ministers for gathering and governing the Church (Rom. x. 14, 17; John xiii. 20; Acts xvi. 9; 1 Cor. iii. 9, etc.).

God employed the Patriarchs, Moses, and the Prophets as teachers of their age. At last he sent his only-begotten Son, filled with infinite wisdom, to be our infallible guide. Christ chose the Apostles, and these ordained pastors in all the Churches (Acts xiv. 23), whose successors have taught and governed the Church to this day.

The ministers of the New Testament are called Apostles, prophets, evangelists, bishops, presbyters, pastors, and teachers (1 Cor. xii. 28; Eph. iv. 11). In subsequent times other names were introduced, as patriarchs, archbishops, metropolitans, archpresbyters, deacons, and subdeacons, etc. But we are satisfied with the offices instituted by the Apostles for the teaching and governing of the Church.

A minister should be lawfully called and chosen by the Church, and excel in sacred learning, pious eloquence, prudence, and unblemished character (1 Tim. iii. 2; Tit. i. 5). When elected, a minister should be ordained of the elders by public prayer and the laying on of hands. We reject arbitrary intruders and incompetent pastors. But we acknowledge that innocent simplicity may be more useful than haughty learning.

A minister of the New Testament is not a priest, as in the Jewish dispensation, offering sacrifices for the living and the dead. Christ is our eternal High-priest, who fulfilled and abolished typical sacrifices by his one perfect sacrifice on the cross; and all believers are priests offering spiritual sacrifices—namely, thanksgiving and praise to God continually.

All ministers are equal in power and commission. Bishops and presbyters were originally the same in office, and governed the Church by their united services, mindful of the words of the Lord: 'He who will be chief among you, let him be your servant' (Luke xxii. 26). Jerome (\emph{Com. on Titus}) says: 'Before, by the instigation of the devil, party spirit and sectarianism arose, the churches were governed by the common counsel of the presbyters; but afterwards, when every one thought that those whom he had baptized belonged to him, not to Christ, it was decreed that one of the presbyters should by election be placed over the rest, and be intrusted with the care of the whole Church, and thus the seed of schisms be destroyed.' But Jerome does not present this decree as divine, for he soon adds that presbyters and bishops know that this distinction is based on ecclesiastical custom, and not on divine command. Therefore no one can be lawfully forbidden to return from human custom to the ancient constitution of the Church of Christ.

The chief duties of ministers are the preaching of the gospel, the administration of the sacraments, the care of souls, and the maintenance of discipline. To do this effectually they must live in the fear of God, pray constantly, study the Scriptures diligently, be always watchful, and shine before all by purity of life. In the exercise of discipline, they should remember that the power was given to them for edification and not for destruction (2 Cor. x. 8; comp.  Matt. xiii. 29).

We reject the error of the Donatists, who make the efficacy of the preaching and the sacraments to depend on the moral character of the minister. The voice of Christ must be heard and obeyed even out of the mouth of an unworthy servant (Matt. xxiii. 3); and the sacraments are efficacious to the worthy recipient by virtue of their divine appointment and the Word of Christ. On these things St. Augustine has much disputed from the Scriptures against the Donatists.

Nevertheless, proper control and discipline should be exercised over the doctrine and conduct of ministers in synods. False or immoral teachers should not be tolerated, but warned or deposed. We do not disapprove general or œcumenical councils if they are conducted, according to the apostolic example (Acts xv.), for the welfare, and not for the corruption of the Church.

As the laborer is worthy of reward, the minister is entitled to the maintenance of himself and family from the congregation he serves (1 Cor. ix. 9 sqq.;  1 Tim. v. 18, etc.). Against the Anabaptists, who denounce ministers living off their ministry.

Chap. XIX. The Sacraments of the Church of Christ.—With the preaching of the Word are joined sacraments or sacred rites instituted by God as signs and seals of his promises for the strengthening of our faith, and as pledges on our part for our consecration to him.

The sacraments of the Jewish dispensation were circumcision and the paschal lamb; the sacraments of the Christian dispensation are baptism and the Lord's Supper.

The Papists count seven sacraments. Of these we acknowledge repentance, ordination of ministers, and marriage as useful institutions of God, but not as sacraments. Confirmation and extreme unction are inventions of men, which may be abolished without any loss. We abhor all merchandise carried on with the sacraments by Romish priests.

The supreme benefit of the sacraments is Christ the Saviour, that Lamb of God slain for our sins from the foundation of the world, and that Rock of which all our fathers drank. So far the sacraments of the Old and New Testaments are the same. But we have the abiding substance.

Sacraments consist of the Word, the sign, and the thing signified. By the Word of God and institution of Christ they become sacraments and are sanctified. The sign in baptism is water, the thing signified is regeneration or the washing from sins. The sign in the Lord's Supper is bread and wine, the thing signified is the veritable body and blood of Christ sacrificed for us. The signs are not changed into the things signified; for then they would cease to be sacramental signs, representing the things signified; but they are sacred and efficacious signs and seals. For he who instituted baptism and the Supper intended that we should receive not the outward form only, but the inward blessing, that we should be truly washed from all our sins through faith, and be made partakers of Christ.

The truth and power of the sacraments depend neither on the worthiness of the minister nor that of the receiver, but on the faithfulness of God. Unbelievers do not receive the things offered; but the fault is in men, whose unbelief doth not annul the faith of God (Rom. iii. 3).

Chap. XX. Of Holy Baptism.— Baptism is instituted by Christ (Matt. xxviii. 19;  Mark xvi. 15). There is only one baptism in the Church; it lasts for life, and is a perpetual seal of our adoption. To be baptized in the name of Christ is to be enrolled, initiated, and received into the covenant, into the family and the inheritance of the sons of God, that, cleansed from our sins by the blood of Christ, we may lead a new and innocent life. We are internally regenerated by the Holy Ghost, but we receive publicly the seal of these blessings by baptism. Water washes away filth, and refreshes and comforts the body; the grace of God inwardly and invisibly cleanses the soul.

By baptism, we are separated from the world and consecrated to God. In baptism we confess our faith and pledge obedience to God. We are enrolled into the holy army of Christ to fight against the World, the flesh, and the devil.

Later human additions to the primitive form of baptism, such as exorcism, the use of burning light, oil, salt, spittle, we judge to be unnecessary.

Baptism is not to be administered by women or by midwives, but by the ministers of the Church.

We condemn those who deny that children of believers should be baptized. For to children belongs the kingdom of God, and they are in covenant with God—why then should not the sign of the covenant be given to them? We are therefore no Anabaptists, and have no communion with them.

Chap. XXI. Of the Holy Supper of our Lord.—The Lord's Supper, or Eucharist, is a grateful commemoration of the benefits of redemption, and a spiritual feast of believers instituted by Christ, wherein he nourishes us with his own flesh and blood by true faith unto eternal life. It signifies and seals to us the greatest benefit and blessing ever conferred on the race of mortals, that he truly delivered his body and shed his blood for the remission of our sins. In it we eat his flesh which is meat indeed, and drink his blood which is drink indeed (Matt xxvi. 20 sqq.; Luke xxii. 19; 1 Cor. xi. 21 sqq.; John vi. 51 sqq.).

This eating is not corporeal and Capernaitic, by the mouth and the stomach, but spiritual, i.e., by the Holy Ghost through faith. 'The flesh,' corporeally eaten, 'profiteth nothing; it is the spirit that quickeneth' (John vi. 63). 'I am the bread of life; he that cometh unto me shall never hunger; and he that believeth on me shall never thirst' (John vi. 51). So that eating and drinking here means to come unto Christ and to believe in him. As Augustine says: 'Why preparest thou the tooth and the stomach? Believe, and thou hast eaten.'

Besides the spiritual eating, in the daily communion of the soul with Christ, there is also a sacramental eating, whereby the believer not only inwardly partakes of Christ, but also receives the visible signs and seals of his body and blood at the Lord's table.\footnote{'\emph{Præter superiorem manducationem spiritualem est et sacramentalis manducatio corporis Domini, qua fidelis non tantum spiritualiter et interne participat vero corpore et sanguine Domini, sed foris etiam accedendo ad mensam Domini accipit visibile corporis et sanguinis Domini sacramentum.}' This is strangely mistranslated by Owen Jones (l.c. p. 173): 'Moreover, also, the sacramental eating of the body of the Lord is a superior spiritual eating,' etc. Bullinger rightly distinguishes between the purely spiritual communion with Christ's flesh and blood (i.e., his real humanity), spoken of in the sixth chapter of John, and the sacramental communion in the Eucharist, which includes all the benefit of the former with the additional blessing of the visible signs and seals of Christ's body broken for us, and Christ's blood shed for us.} And with the signs he receives the thing itself.\footnote{'\emph{Qui foris vera fide sacramentum percipit, idem ille non signum duntaxat percipit, sed re ipsa quoque, ut diximus, fruitur.}'} He is nourished and strengthened by spiritual food. The signs are also sure pledges that Christ died not only for men in general, but also individually for every believing communicant. Besides, in partaking of this ordinance we obey the command of our Lord, celebrate his atoning death, give thanks for the great redemption, and openly profess our faith before the congregation.

But those who commune unworthily and without faith receive only the visible signs to their own condemnation or judgment (1 Cor. xi. 27 sqq.).

We therefore do not so conjoin the body and blood of Christ with bread and wine as to say that the bread itself is the body (except sacramentally), or that the body of Christ is corporeally hid under the bread, aud should be adored under the form of bread, or that whosoever receives the signs receives also necessarily the thing itself. [Against the Lutheran theory.] The body of Christ is in heaven at the right hand of the Father (Mark xvi. 19; Heb. viii. 1;  xii. 2); and hence we must raise our hearts to heaven.

And yet he is not absent from his people when they celebrate his communion. For as the sun in heaven is efficaciously present with us, so much more is Christ the sun of righteousness with us, not, indeed, corporeally, but spiritually by his enlivening and vivifying operation, even as he in the Last Supper explained that he himself would be present with us (John xiv.-xvi.). Hence we have not a Supper without Christ, but an unbloody and mystical Supper, as universal antiquity called it.

Moreover, the Lord's Supper reminds us that we are members of his body, and should live peaceably with all our brethren, and grow and persevere in holiness of life.

Therefore it is very proper that we should duly prepare ourselves by self-examination in regard to our repentance and faith in Christ (1 Cor. xi. 28).

As to the external celebration, we adhere to the original form, consisting in the annunciation of the Word of God, devout prayers, the Lord's action, and its repetition in breaking bread, and distributing it together with the wine, in eating the body and drinking the blood of our Lord, in grateful remembrance of his death, in thanksgiving, and in holy reunion of the brethren as one body.

We disapprove of the withdrawal of the cup contrary to the express command of our Lord: 'Drink ye \emph{all} of it' (Matt. xxvi. 27).

The mass—whatever it may have been in ancient times—has been turned from a salutary institution into a vain show, and surrounded with various abuses, which justify its abolition.

Chap. XXII. Of Sacred and Ecclesiastical Assemblies.—It is lawful and right for all men privately to read the Scriptures for edification. At the same time the maintenance of religion demands regular public services. These should be conducted decently, in order, and for edification, in the language understood by the people.

Chap. XXIII. Of Church Prayers, Singing, and Canonical Hours.—Public prayers in sacred assemblies should be made in the vulgar tongue understood by all. Every prayer is to be offered to God alone, through the sole mediation of Christ, not to saints or through them. Churches are at liberty to vary from the usual forms. Prayers are not superstitiously to be confined to particular places or hours. Long and tedious prayers in public assemblies should be avoided. Singing is not indispensable, but lawful and desirable. Canonical hours are not prescribed in the Scriptures, and are unknown to antiquity.

Chap. XXIV. Of Feasts, Fasts, and the Choice of Meats.—The Lord's day is consecrated, from the times of the Apostles, to the worship of God and to sacred rest. But we observe it in Christian freedom, not with Jewish superstition, neither do we believe that one day is in itself holier than another.

If congregations in addition commemorate the Lord's nativity, circumcision, crucifixion, resurrection, and ascension, and the outpouring of the Holy Ghost, we greatly approve of it. But feasts instituted by men in honor of saints we reject, though the memory of the saints is profitable, and should be commended to the people with exhortations to follow their virtues.

True Christian fasting consists in temperance, abstinence, watchfulness, self-government, and chastisement of our flesh, that we may the easier obey the Spirit. Such fasting is a help to prayer and all virtues.

There are also public fasts appointed in times of affliction and calamity, when people abstain from food altogether till evening, and spend all time in prayer and repentance. Such fasts are mentioned by the Prophets (Joel ii. 12 sq.), and should be observed when the Church is afflicted and oppressed. Private fasts are observed by each of us as we may judge it profitable to our souls.

All fasts ought to proceed from a free and willing mind, and be observed in a spirit of true humility, in order to vanquish the flesh and to serve God more fervently, but not in order to gain the favor of men or the merit of righteousness.

The fast, of forty days (Lent) has the testimony of antiquity, but is not enjoined in the Scriptures, and ought not to be imposed upon the conscience of the faithful. There was great diversity and freedom in the early Church as to the time of fasting, as we learn from Irenæus, and Socrates the historian.

As to the choice of meats, we hold that in fasts we should abstain from all such food or drink as stimulates the carnal desires. But otherwise we know that all the creatures of God are good (Gen. i. 31), and may be used without distinction, but with moderation and thanksgiving (1 Cor. x. 25;  Tit. i. 15). Paul calls the prohibition of meats a doctrine of the demons (1 Tim. iv. 1 sqq.), and reproves those who by excessive abstinence wish to acquire the fame of sanctity.

Chap. XXV. Of Catechizing, and of the Visitation and Consolation of the Sick.—The greatest care is to be bestowed on the religious instruction of the youth, especially in the Ten Commandments, the Apostles' Creed, the Lord's Prayer, and the nature of the sacraments. Churches should see to it that children receive catechetical instruction.

It is one of the chief duties of Christian pastors to visit, comfort, and strengthen the sick, and pray for them in private and in public. But the extreme unction of the Papists we disapprove.

Chap. XXV. Of the Burial of the Faithful, the Care of the Dead, of Purgatory, and the Apparition of Spirits.—The bodies of believers, which are the temples of the Holy Ghost, and will rise again in the last day, should be honorably committed to the earth, without superstition, and their relatives, widows, and orphans should be tenderly cared for.

We believe that the faithful after death go directly to Christ, and need not the prayers of the living. Unbelievers are cast into hell, from which there is no escape.

The doctrine of purgatory is opposed to the Scriptures, and to the plenary expiation and cleansing through Christ (comp. John v. 24; xiii. 10).

The tales about the souls of the departed appearing to the living and requesting their services for deliverance we judge to be mockeries or deceptions of the devil. The Lord forbids necromancy (Deut. xviii. 10); and the rich man was told that if his brethren on earth hear not Moses and the Prophets, neither will they be persuaded though one rose from the dead (Luke xvi. 30).

Chap. XXVII. Of Rights and Ceremonies.—The ceremonial law of the Jews was a schoolmaster and guardian to lead them to Christ, the true Liberator, who abrogated it so that believers are no more under the law, but under the gospel freedom. The Apostles would not lay the burden of Jewish ceremonies on the new converts (Acts xv. 28). The more of human rites are accumulated in the Church, the more it is drawn away from Christian liberty and from Christ himself, while the ignorant seek in ceremonies what they should seek in Christ through faith. A few pure and moderate rites consistent with the Word of God are sufficient.

Difference in ceremonies, such as existed in the ancient Church, and exists now among us, need not to interfere with union and harmony in doctrine and faith. In things indifferent, which are neither good nor evil, the Church has always used liberty (1 Cor. viii. 10;  x. 27 sqq.).

Chap. XXVIII. Of Church Property.—The wealth of the Church should be used for the maintenance of public worship and schools, the support of ministers and teachers, and especially also for the benefit of the poor.

Misapplication and abuse of Church property through ignorance or avarice is a sacrilege, and calls for reformation.

Chap. XXIX. Of Celibacy, Marriage, and Economy.—Those who have the gift of celibacy from heaven, so as to be pure and continent from their whole heart, may serve the Lord in that vocation in simplicity and humility, without exalting themselves above others. If not, they should remember the apostolic word: 'It is better to marry than to burn' (1 Cor. vii. 9).

Marriage (the remedy for incontinence, and continence itself) was instituted by God, who blessed it richly, and inseparably joined man and woman to live together in intimate love and harmony (Matt. xix. 5). Marriage is honorable in all, and the bed is undefiled (Heb. xiii. 4;  1 Cor. vii. 28).

We condemn polygamy, and those who reject second marriages.

Marriage should be contracted in the fear of the Lord, with the consent of parents or their representatives, and for the end for which it was instituted.

Children should be brought up in the fear of the Lord, properly supported by their parents (1 Tim. v. 8), and be taught honest arts or trades.

We condemn the doctrine which forbids marriage, or indirectly slights it as unholy and unclean (1 Tim. iv. 1). We execrate unclean celibacy, secret and open fornications, and the pretended continency of hypocrites.

Chap. XXX. Of the Magistrate.— The civil magistrate is appointed by God himself (Rom. xiii.) for the peace and tranquillity of the human race. If opposed to the Church, he can do much harm: if friendly, he can do the Church most useful service.

The duty of the magistrate is to preserve peace and public order; to promote and protect religion and good morals; to govern the people by righteous laws; to punish the offenders against society, such as thieves, murderers, oppressors, blasphemers, and incorrigible heretics (if they are really heretics).\footnote{'\emph{Coërceat et hæreticos} (\emph{qui rere hæretici sunt}) \emph{incorrigibiles, Dei majestatem blasphemare et Ecclesiam Dei conturbare, adeoque perdere non desinentes.}' The same view of the right and duty of the civil government to punish heretics is expressed in other Confessions. The Reformers differed from the Roman Catholics, not so much in the principle of persecution as in the definition of heresy and the degree of punishment. Nevertheless, the Reformation inaugurated the era of religious toleration and freedom.}

Wars are justifiable only in self-defense, and after all efforts at peace have been exhausted.

We condemn the Anabaptists, who maintain that a Christian should not hold a civil office, that the magistrate has no right to punish any one by death, or to make war, or to demand an oath.

All citizens owe reverence and obedience to the magistrate as the minister of God in all righteous commands, and even their lives when the public safety and welfare require it. Therefore we condemn the despisers of the magistrate, rebels and enemies of the commonwealth, and all who openly or artfully refuse to perform their duties as citizens.

We pray to God, our merciful heavenly Father, to bestow his blessing upon princes and rulers, upon us, and upon all his people, through Jesus Christ our only Lord and Saviour: to whom be praise, and glory, and thanksgiving, forever and ever. Amen.

\xsubsection{John Calvin. His Life and Character.}

§ 56. John Calvin. His Life and Character.

Literature

I. Works and Correspondence of Calvin.

Joannis Calvini.  \emph{Opera quæ supersunt omnia, ed. G. Baum, E. Cunitz, E. Reuss, theologi Argentoratenses.} Brunsvigæ, 1863 sqq. (in the \emph{Corp. Reform.}). So far (1884) 27 vols. 4to. The most complete and most critical edition.

Older Latin ed., Geneva, 1617, in 12 vols. folio, and Amstelod. 1671, 9 vols. fol.

An English edition of Calvin's Works, by the 'Calvin Translation Society,' Edinburgh, 1842–1853, in 52 vols.

Convenient editions of Calvin's \emph{Institutes,} by Tholuck (Berol. 1834 and 1846); the Commentaries on \emph{Genesis,} by Hengstenberg (Berol. 1838), on the \emph{Psalms} (Berol. 1830–34), on the \emph{New Testament} (except the Apocalypse, 1833–38, in 7 vols.), by Tholuck.

His most important works were also written in French.

A German translation of his \emph{Institutes,} by Fr. Ad. Krummacher (1834), of his \emph{Comment.,} by C. F. L. Matthieu (1859 sqq.).

The extensive correspondence of Calvin was first edited in part by Beza and Jonvillers (Calvin's secretary), Genevæ, 1575, and other editions; by Bretschneider (the Gotha Letters), Lips. 1835; by A. Crottet, Genève, 1850; then much more completely by Jules Bonnet, \emph{Lettres Françaises,} Paris, 1854, 2 vols.; an English translation (from the French and Latin) by D. Constable and M. R. Gilchrist, Edinburgh and Philadelphia (Presbyt. Board of Publ.), 1855 sqq., in four vols. (the 4th with an index), giving the letters in chronological order (till 1558). The last and best edition is by the Strasburg Professors in \emph{Calvini Opera,} Vol. X. Part II. to Vol. XV., with ample \emph{Prolegomena} on the previous editions of Calvin's Letters and the manuscript sources.

Compare, also,  A. L. Herminjard:  \emph{Correspondance des réformateurs dans les pays de langue française,} (beginning with 1512). Genève and Paris, 1866, sqq., 5th vol. 1883. A most important work, with many new letters from and to the Reformers, illustrated by historical and biographical notes; the correspondence of Calvin begins Tome II. p. 278.

II. Biographies of Calvin.

Th. de Bèze: \emph{ Histoire de la vie et la mort de J. Calvin,} Genève, 1564; second French ed. enlarged and improved by Nic. Colladon, 1565, recently republished by A. Franklin, Paris, 1864; Latin ed. by Beza, as an introduction to Calvin's Letters, 1575, reprinted in Tholuck's ed. of Calvin's Commentaries. There are also German, English, and Italian translations. The \emph{second} French and the Latin editions should be consulted. This work of Beza, together with Calvin's Letters and Works, furnishes the chief material for an authentic biography.

Hieron. Bolsec (a Carmelite monk, then physician at Geneva, expelled on account of Pelagian views and opposition to Calvin, 1551, returned to the Roman Church 1563): \emph{Histoire de la vie de Jean Calvin,} Paris, 1577 (Genève, 1835); then in Latin: \emph{De J. Calvini magni quondam Genevensium ministri vita, moribus, rebus gestis, studiis ac denique morte,} Coloniæ, 1580. 'A mean and slanderous libel,' inspired by feelings of hatred and revenge. See Schweizer, \emph{Centraldogmen,} Vol. I. p. 205.

Jacques Le Vasseur (R.C.): \emph{Annales de l’église cathédrale de Noyon,} Paris, 1633. Contains some notices on the youth of Calvin.

Jacques Desmay (R.C.): \emph{Remarques sur la vie de J. Calvin hérésiarque tirées des Registres de Noyon,} Rouen, 1657.

Drelincourt:  \emph{La défense de Calvin contre l’outrage fait à sa mémoire,} Genève, 1667; in German, Hanau, 1671. A refutation of the slanders of Bolsec.

Paul Henry (pastor of a French Reformed Church in Berlin): \emph{Das Leben Johann Calvins des grossen Reformators,} etc., Hamburg, 1835–44, 3 vols.; also abridged in one vol., Hamburg, 1846. English translation by Stebbing, London and New York, 1854, in 2 vols. The large work is a valuable collection rather than digestion of material for a full biography by a sincere admirer.

E. Stähelin (Reformed minister at Basle): \emph{Johannes Calvin; Leben und ausgewählte Schriften,} Elberfeld, 1863, 2 vols. (in \emph{Väter und Begründer der reform. Kirche,} Vol. IV. in two parts). Upon the whole the best biography, though not as complete as Henry's, and in need of modification and additions from more recent researches.

T. H. Dyer: \emph{ Life of Calvin,} London, 1850. 'Valuable and impartial' (Fisher).

Felix Bungener: \emph{ Calvin, sa vie, son œuvre et ses écrits,} Paris, 1862; English translation, Edinb. 1863.

F. W. Kampschulte (a liberal Roman Catholic, Professor of History at Bonn, died an Old Catholic, 1871): \emph{Joh. Calvin, seine Kirche und sein Staat in Genf,} Leipzig, 1869, Vol. I. (Vols. II. and III. have not appeared). A most able, critical, and, for a Catholic, remarkably fair and liberal work, drawn in part from unpublished sources.

Guizot (the great historian and statesman, a descendant of the Huguenots, d. at Val Richer, Sept. 12, 1874): \emph{St. Louis and Calvin,} London, 1868. Comp. also his sketch in the \emph{Musée des protestants célèbres.}

The work of the Roman Catholic  Audin:  \emph{Histoire de la vie,} etc., \emph{de Calvin,} Paris, 1841, 5th ed., 1851, in 2 vols. (also in English and German), is mostly a slanderous caricature, based upon Bolsec.

III. Biographical Sketches and Essays.

H. Mignet: \emph{ Mémoire sur l’établissement de la réforme et sur la constitution du Calvinisme à Genève,} Paris, 1834.

J. J. Herzog: \emph{ Joh. Calvin,} Basel, 1843; and in his \emph{Real-Encyklop.} Vol. II. p. 511.

E. Renan: \emph{ Jean Calvin,} in \emph{Études d’histoire réligieuse,} 5th ed., Paris, 1862; English translation by O. B. Frothingham (\emph{Studies of Religious History and Criticism,} New York, 1864, pp. 285–297).

Philip Schaff: \emph{ John Calvin,} in the \emph{Bibliotheca Sacra,} Andover, 1857, pp. 125–146.

Henry B. Smith: \emph{ John Calvin,} in Appleton's \emph{American Cyclopædia,} New York, Vol. IV. (1859) pp. 281–288.

James Anthony Froude: \emph{ Calvinism, an Address delivered to the Students of St. Andrew's,} March 17, 1871 (In his \emph{Short Studies on Great Subjects,} Second Series, New York, 1873, pp. 9–53).

A. A. Hodge (of Alleghany, son of Dr. Charles Hodge of Princeton): \emph{Calvinism,} in Johnson's \emph{Universal Cyclopædia} (New York, 1875 sqq.), Vol. I. pp. 727–734.

Lyman H. Atwater: \emph{ Calvinism in Doctrine and Life,} in the \emph{Presbyt. Quarterly and Princeton Review,} New York, Jan. 1875, pp. 73–106.

IV. Histories of the Reformation in Geneva.

Abr. Ruchat (Professor in Lausanne): \emph{Histoire de la réformation de la Suisse,} Genève, 1727 sqq. 6 vols.; new edition, with appendices, by Prof. Vulliemin, Nyon, Giral. 1835–1838, 7 vols.

C. B. Hundeshagen (Professor in Berne, afterwards in Bonn, d. 1872): \emph{Die Conflicte des Zwinglianismus, Lutherthums und Calvinismus in der Bernischen Landeskirche von} 1532–1558. \emph{Nach meist ungedruckten Quellen.} Bern, 1842.

J. Gaberel (ancien pasteur): \emph{Histoire de l’église de Genève depuis le commencement de la réforme jusqu’en} 1815. Genève, 1855–63, 3 vols.

P. Charpenne: \emph{ Histoire de la réforme et des réformateurs de Genève.} Paris, 1861.

Amad. Roget: \emph{ L'église et l’état à Genève de vivant Calvin,} 1867; and \emph{Histoire du peuple de Genève depuis la réforme jusqu’à l’escalade.} Genève, 1870.

Merle d’Aubigné (Professor of Church History at Geneva, d. 1872): \emph{History of the Reformation in Europe in the time of Calvin} (from the French), New York, 1863–1879, 8 vols. (the second division of his general history of the Reformation. The last two volumes were edited from the author's MSS. They carry the history down to the middle of the 16th century.

G. P. Fisher: \emph{ The Reformation.} New York, 1873, Ch. VII. pp. 192–241.

For the political history of \emph{Geneva} preceding and during the time of Calvin are to be compared Fr. Bonnivard: \emph{ Les Chroniques de Genève,} edited by Dunant (Gen. 1831, 4 vols.); Galiffe: \emph{ Matériaux pour l’histoire de Genève; } J. P. Bérenger: \emph{ Histoire de Genève jusqu’en} 1761 (1772, 6 vols.); and the \emph{Mémoires et documents publiés par la Société d’histoire et d’archéologie de Genève} (1840 sqq.).

CALVIN'S LIFE

After the death of Zwingli and the treaty of Cappel (1531), the progress of the Reformation was checked in German Switzerland, but only to make a more important conquest in French Switzerland, and from thence with the course of empire to move westward to France, Holland, beyond the Channel, and beyond the seas.

The supremacy passed from Zurich to Geneva. Providence had silently prepared the person and the place. The 'little corner' on the borders of Switzerland and France, known since the days of Julius Caesar, was predestinated, by its location and preceding history, for a great international mission, and has nobly fulfilled it, not only in the period of the Reformation of the Church, but also in the nineteenth century on the field of international law and peaceful arbitration. After varying fortunes, Geneva became an independent asylum of civil and religious freedom, and furnished the best base of operation for John Calvin, who, though a Frenchman by birth and a Swiss by adoption, was a cosmopolitan in spirit, and acted as the connecting link between the Germanic and Latin races in the work of reform. Farel, Viret, and Froment had destroyed the power of Popery, but to Calvin was left the more difficult task of reconstruction and permanent organization.

John Calvin,\footnote{The Latinized form of the French \emph{Chauvin} or \emph{Cauvin.} He sunk, even in name, his nationality in his catholicity.} the greatest theologian and disciplinarian of the giant race of the Reformers, and for commanding intellect, lofty character, and far-reaching influence one of the foremost leaders in the history of Christianity, was born at Noyon, in Picardy, July 10, 1509. His father, Gerard Chauvin, a man of severe morals, was secretary to the Bishop of Noyon; his mother, a beautiful and devout, but otherwise not remarkable woman. He received his first training with the children of a noble family (de Mommor), to which he was gratefully attached. His ambitious father destined him for the clerical profession, and secured him even in his twelfth year the benefice of a chaplaincy of the cathedral—an abuse not infrequent in those days of decay of ecclesiastical discipline. He received the tonsure, but not the ordination for the priesthood; while Zwingli and Knox were once priests, and Luther both priest and doctor, in the Church they were called to reform. His elder brother, Charles, became a priest at Noyon, and died a libertine and an infidel in the same year in which John proclaimed his faith to the world (1536)—as if to repeat the startling contrast of Esau and Jacob, reprobation and election, from the same womb.\footnote{Guizot (pp. 153, 155): 'Evidently Charles Calvin lived and died a dissolute man and an unbeliever, and at the same time remained chaplain of the Catholic church of his native town. The sixteenth century abounds in similar instances. . . . The same thing was going on every where; unbelievers and fervent Christians, libertines and men of the most austere lives, were springing up and living side by side. Two contrary winds were blowing over Europe at that period, one carrying with it skepticism and licentiousness, while the other breathed only Christian faith and the severest morality. One of these arose chiefly from the revival of the ancient literature and philosophy of Greece and Rome; the other sprang from the struggles made in the Church itself, and in its councils, to arrive at a reform which was at the same time greatly desired and fiercely opposed. . . . It was, in short, the age which produced Erasmus and Luther in Germany, and Montaigne and Calvin in France.' Merle d'Aubigné (Vol. V. p. 455) conjectures that Charles Calvin became a convert to Protestantism on his death-bed, for which the infuriated priests had him buried by night between the four pillars of a gibbet.} Another remarkable coincidence is the fact that the Reformer studied scholastic philosophy under the same Spanish instructor of the College de Montaigu at Paris in which a few years afterwards Ignatius Loyola, the famous founder of Jesuitism—the very opposite pole of Calvinism—laid the foundation of his counter-reformation.\footnote{Kampschulte, Vol. I. p. 223.}

Calvin received the best education which France could afford, in the Universities of Orleans, Bourges, and Paris, first for the priesthood, then, at the request of his father, for the law.\footnote{It seems (according to Jacques Le Vasseur, 1.c. 1153 sqq., as quoted by Kampschulte, Vol. I. p. 226) that Gerard Chauvin became involved in difficulty with his ecclesiastical superiors, and was even excommunicated. Kampschulte conjectures that this was probably the reason why he ordered his son to exchange the study of theology for that of law. But Calvin himself (in his Commentary on the Psalms) assigns a different motive: '\emph{Mon père m’avoit destiné à la Théologie; mais puis après, d’autant qu’il considéroit que la science des Loix communément enrichit ceux qui la suyvent, ceste espérance luy fait incontinent changer d’avis.}' The study of the law was of great use to Calvin in the organization and control of Church and State in Geneva.} He early distinguished himself by excessive industry, which undermined his constitution, severe self-discipline, and a certain censoriousness, for which he was called by his fellow-students 'the Accusative Case.'\footnote{A notice of Jacques Le Vasseur, which agrees with Beza's statement that he was '\emph{tenera ætate mirum in modum religiosus}' and '\emph{severus omnium in suis sodalibus vitiorum censor.}'} He made rapid progress. Even as a student of nineteen he was often called to the chair of an absent professor, so that (as Beza says) he was considered a doctor rather than an auditor. When he left the university he was the most promising literary man of the age. He might have attained the highest position in France, had not his religious convictions undergone a radical change.

Protestant ideas were then pervading the atmosphere and agitating the educated classes of France even at the court, which was divided on the question of religion. Two of Calvin's teachers, Cordier (or Corderius, who afterwards followed him to Geneva) and Wolmar, were friendly to reform, and one of his relatives, Olivétan, became soon afterwards (1534) the first Protestant translator of the Bible into French. He seems, however, to have exerted as much influence on them as they exerted on him.\footnote{According to Beza and Stähelin (Vol. I. p. 88), Calvin took part even in the first edition of Olivétan's French New Testament (1534). But this seems to be an error; see Reuss, \emph{Révue de Theologie,} 1866, No. III. p. 318, and Kampschulte, p. 247. He revised, however, the second edition, which included the Old Testament (1535), and wrote the preface (see Stähelin, pp.89 sq.).}

His first work was a commentary on Seneca's book on \emph{Mercy,} which he published at his own expense, April, 1532.\footnote{'\emph{L. Annei Se- }| \emph{necæ, Romani Senato- }| \emph{ris, ac philosophi clarissi- }| \emph{mi, libri duo de Clementia, ad Ne- }| \emph{ronem Cæsarem: }| \emph{Joannis Calvini Noviodunæi commentariis illustrati. }| \emph{Parisiis} . . . 1532.' Reprinted from the \emph{ed. princeps} in the new edition of the \emph{Opera,} Vol. V. (1866), pp. 6–162. The commentary is preceded by a dedicatory epistle, and a sketch of the life of Seneca.} It moves in the circle of classical philology and moral philosophy, and reveals a characteristic love for the nobler type of Stoicism, great familiarity with Greek and Roman literature, masterly Latinity, rare exegetical skill, clear and sound judgment, and a keen insight into the evils of despotism and the defects of the courts of justice, but makes no allusion to Christianity. Hence it is quite improbable that it was an indirect plea for toleration and clemency intended to operate on the King of France in dealing with his Protestant subjects.\footnote{As is asserted by Henry, Herzog, Dorner (p. 375), and also by Guizot (p. 162), but justly denied by Stähelin (Vol. I. pp. 14 sqq.) and Kampschulte (p. 238). The work is not dedicated to Francis I., but to Claude de Hangest, the Abbot of St. Eloy (Eligius), afterwards Bishop of Noyon, his former schoolmate; and the implied comparison of the French king with Nero, and the incidental mention of the Neronian persecution ('\emph{quum Nero diris suppliciis impotenter sæviret in Christianos,}' \emph{Opera,} Vol. V. p. 10), would have been fatal to such an apologetic aim. Calvin sent a copy to 'Erasmus, and called him 'the honor and the chief delight of the world of letters'—\emph{literarum alterum decus ac primæ deliciæ} (see his letter to Claude de Hangest, April 4, 1532, in Herminjard, Tom. II. p. 411).} His earliest letters, from 1530 to 1532, are likewise silent on religious subjects, and refer to humanistic studies, and matters of friendship and business.\footnote{They were recently brought to light by Jules Bonnet and Herminjard. They are chiefly addressed to his fellow-student, Francis Daniel, an advocate of Orleans, who acknowledged the necessity of the Reformation, but remained in the Church of Rome. See the Edinburgh edition of Calvin's \emph{Letters,} by Bonnet, Vol. I. p. 3; Herminjard, Vol. II. pp. 278 sqq.; and \emph{Opera,} Vol. X. Pt. II. pp. 3 sqq. His first letter to Daniel is dated '\emph{Melliani} (i.e. Meillant, south of Bourges, not Meaux, as the Edinburgh edition misunderstands it), 8 \emph{Idus Septembr.,}' and is put by Herminjard and the Strasburg editors in the year 1530 (not 1529).}

His conversion to the cause of the Reformation seems to have taken place in the latter part of 1532, about one year after the death of Zwingli.\footnote{Stähelin puts his conversion in the year 1533 (Vol. I. p. 21). But we have a familiar letter from Calvin to Martin Bucer, dated Noyon, '\emph{pridie nonas Septembres,}' probably of the year 1532, in which he recommends a French refugee, falsely accused of holding the opinions of the Anabaptists, and says: 'I entreat of you, Master Bucer, if my prayers, if my tears are of any avail, that you would compassionate and help him in his wretchedness. The poor is left in a special manner to your care—you are the helper of the orphan. . . . Most learned Sir, farewell; Thine from my heart (\emph{Tuus ex animo}): Calvin' (J. Bonnet's \emph{Letters,} Vol. I. pp. 9–11; the Latin in \emph{Opera,} Vol. X. Pt. II. p. 24). Kampschulte (Vol. I. p. 231) infers even an earlier acquaintance of Calvin with Bucer, from a letter of Bucer to Farel, May 1, 1528, in which he mentions a \emph{juvenis Noviodunensis} studying Greek and Hebrew in Strasburg (Herminjard, Vol. II. p. 131, and \emph{Opera,} Vol. X. Pt. II. p. 1); but this youth was probably his relative Olivétan, who was likewise a native of Noyon (Herminjard, Vol. II. p. 451). Besides, there were several places in France of the name \emph{Noviodunum.} In a letter of Oct., 1533, to Francis Daniel (Bonnet, Vol. I. p. 12, and \emph{Opera,} Vol. X. Pt. II. p. 27), Calvin first speaks openly of the Reformation in Paris, the rage of the Sorbonne, and the satirical comedy against the Queen of Navarre.} The precise date and circumstances are unknown. It was as he himself characterizes it, a sudden change (\emph{subita conversio}) from Papal superstition to the evangelical faith, yet not without previous struggles. He tenaciously adhered to the Catholic Church until he was able to disconnect the true idea and invisible essence of the Church from its outward organization. Like Luther, he strove in vain to attain peace of conscience by the methods of Romanism, and was driven to a deeper sense of sin and guilt. 'Only one haven of salvation is left for our souls,' he says, 'and that is the mercy of God in Christ. We are saved by grace—not by our merits, not by our works.' After deep and earnest study of the Scriptures, the knowledge of the truth, like a bright light from heaven, burst upon his mind with such force that there was nothing left for him but to abjure his sins and errors, and to obey the will of God. He consulted not with flesh and blood, and burned the bridge after him.\footnote{He alludes to his conversion only twice, and briefly, namely, in the remarkable Preface to his Commentary on the Psalms, and in his answer to Cardinal Sadolet (\emph{Opera,} Vol. V. pp. 389–411 sq.). In the latter he describes his mental conflicts and terrors of conscience.}

There never was a change of conviction purer in motive, more radical in character, more fruitful and permanent in result. It bears a striking resemblance to that still greater event near Damascus, which transformed a fanatical Pharisee into an apostle of Jesus Christ. And indeed Calvin was not unlike St. Paul in his intellectual and moral constitution; and the apostle of sovereign grace and evangelical freedom never had a more sympathetic expounder than the Reformer of Geneva.

With this step Calvin renounced all prospects of a brilliant career, upon which he had already entered, and exposed himself to the danger of persecution and death.\footnote{He says (\emph{Ad Sadoleti Epistolam, Opera,} Vol. V. p. 389) that if he had consulted his personal interest he would never have left the Roman Church, where the way to honor would have been very easy to him. Audin, in tracing Calvin's conversion to wounded ambition, exposes (as Kampschulte justly observes, p. 242) his utter ignorance of Calvin's character, whose only ambition was to serve God most faithfully.} Though naturally bashful and retiring, and seeking one quiet hiding-place after another, he was forced to come forward. He exhorted and strengthened the timid believers, usually closing with the words of St. Paul: 'If God be for us, who can be against us?' There is no evidence that he ever was ordained by human hands to the ministry of the gospel; but he had an extraordinary call, like that of the prophets of old, and the Apostle of the Gentiles. This was felt by his brethren, and about a year after his conversion he was the acknowledged leader of the Protestant party in France.

For awhile matters seemed to take a favorable turn at the court. His friend, Nicholas Cop, a learned physician, was even elected Rector of the University of Paris.\footnote{Bulæus, \emph{Historia universitatis Parisiensis,} Vol. VI. p. 238; Kampschulte, Vol. I. p. 243.} At his request Calvin prepared for him an inaugural address on Christian philosophy, which Cop delivered on All-Saints' Day, in 1533, in the Church of the Mathurins, before a large assembly. He embraced this public occasion to advocate the reform of the Church on the basis of the pure gospel.\footnote{The incomplete draft of this address has recently been discovered by J. Bonnet among the manuscripts of the Geneva library. In it Calvin explains the great difference between the law and the gospel, and charged the Sophists, as he called the scholastic theologians, \emph{Nihil de fide, nihil de amore Dei, nihil de remissione peccatorum, nihil de gratia, nihil de justificatione, nihil de veris operibus disserunt; aut si certe disserunt, omnia calumniantur, omnia labefactant, omnia suis legibus, hoc est sophisticis coërcent. Vos rogo, quotquot hic adestis, ut has hæreses, has in Deum contumelias numquam æquo animo feratis.}' See Kampschulte, p. 244.} Such a provocation Catholic France had never before received. The Sorbonne ordered the address to be burned. Cop was warned, and fled to Basle; Calvin—as tradition says—escaped in a basket from a window, and left Paris in the garb of a vine-dresser, scarcely knowing whither he was going. A few months afterwards the king himself took a decided stand against the Reformation, and between Nov. 10, 1534, and May 3, 1535, twenty-four Protestants were burned alive in Paris, while many more were condemned to less cruel sufferings.\footnote{This is recorded with some satisfaction by a Catholic writer in the \emph{Journal du Bourgeois de Paris,} quoted by Guizot, p. 168. That Francis I. was present at these horrible executions is denied by Michelet, Martin, and Guizot.}

For more than two years Calvin wandered a fugitive evangelist, under assumed names, from place to place. We find him at Angonlême with his learned friend, the young canon Louis du Tillet, using his excellent library, and probably preparing his 'Institutes;' then at the court of Queen Margaret of Navarre, the sister of Francis I., where he met Le Fèvre d’Estaples (Faber Stapulensis), the aged patriarch of French Protestantism, and Gérard Roussel, her chaplain, who advised him 'to purify the house of God, but not to destroy it;' at Noyon (May, 1534), where he parted with his ecclesiastical benefices; at Poictiers, where he celebrated, with a few friends, for the first time, the Lord's Supper according to the evangelical rite, in a cave near the town, called to this day 'Calvin's Cave;' at Orleans, where he published his first theological work, a tract against the Anabaptist doctrine of the sleep of the soul between death and the resurrection, using exclusively Scriptural arguments with rare exegetical and polemical skill;\footnote{\emph{Psychopannychia,} in \emph{Opera,} Vol. V. pp. 165–232. The Preface is dated 'Aureliæ, 1534.' The second edition appeared in Basle, 1535. This work forms a contrast to his commentary on Seneca as great as exists between the classics and the Bible. In matters relating to the future world. Calvin allows no weight to reason and philosophy, but only to the Word of God. On the merits of this book, see Stähelin, Vol. I. pp. 36 sqq.} again (towards the close of 1534) at Paris, where he met for the first time the unfortunate Michael Servetus, and challenged him to a disputation on the Trinity. But the persecution then breaking out against the Protestants forced him to forsake the soil of France. With his friend Du Tillet he fled to Strasburg, where he arrived utterly destitute, having been robbed by an unfaithful servant, and formed an intimate friendship with Bucer. Thence he went to Basle, where he quietly studied Hebrew with Capito and Grynæus, and published the first edition of his 'Institutes' (1536). In the spring of 1536 he spent a short time at the court of the Duchess Renée of Ferrara, the daughter of Louis XII., a little, deformed, but highly intelligent, noble, and pious lady, who gathered around her a circle of friends of the Reformation, and continued to correspond with him as her guide of conscience.\footnote{Guizot, speaking at some length of this correspondence, makes the remark (p. 207): 'I do not hesitate to affirm that the great Catholic bishops, who in the seventeenth century directed the consciences of the mightiest men in France, did not fulfill the difficult task with more Christian firmness, intelligent justice, and knowledge of the world than Calvin displayed in his intercourse with the Duchess of Ferrara. And the Duchess was not the only person towards whom he fulfilled this duty of a Christian pastor. His correspondence shows that he exercised a similar influence, in a spirit equally lofty and judicious, over the consciences of many Protestants.'} Returning from Italy, where he was threatened by the Inquisition,\footnote{He took the route of Aosta and the Great St. Bernard. His short labors and persecution in Aosta were, five years later (1541), commemorated by a monumental cross and inscription—'\emph{Calvini fuga}'—which was restored in 1741, and again in 1841, and stands to this day. See Gaberel, Vol. I. p. 100; Stähelin, Vol. I. p. 110; Guizot, p. 209: and Merle d'Aubigné, Vol. V. p. 454.} he paid a flying visit to Noyon, and had the pleasure to gain his only remaining younger brother Anthony and his sister Mary to the Reformed faith. With them he proceeded to Switzerland, intending to settle at Basle or Strasburg, and to lead the quiet life of a scholar and an author, without the slightest inclination to a public career. But God had decreed otherwise.

Passing through Geneva in August, 1536, where he expected to spend only a night, Calvin was held fast by William Farel, the fearless evangelist, who threatened him with the curse of God if he preferred his studies to the work of the Lord. 'These words,' says Calvin (in the Preface to his Commentary on the Psalms), 'terrified and shook me, as if God from on high had stretched out his hand to stop me, so that I renounced the journey which I had undertaken.'\footnote{According to Beza (\emph{Vita}), Farel used these words: '\emph{At ego tibi studia, prætexenti denuntio, omnipotentis Dei nomine, futurum, ut, nisi in opus istud Domini nobiscum incumbas, tibi non tam Christum quam te ipsum quærenti Dominus maledicat.}' Beza adds that Calvin was '\emph{territus hac terribili denuntiatione.}' Merle d'Aubigné gives a very dramatic account of this scene, Vol. V. pp. 456 sqq.} Farel, a French nobleman, twenty years older than Calvin, and like him driven by persecution to Switzerland, where he destroyed the strongholds of idolatry with the zeal of a prophet, did a great work when 'he gave Geneva to the Reformation,' but a still greater one when 'he gave Calvin to Geneva.'

This was the turning-point in Calvin's life. Once resolved to obey the voice from heaven, the timid and delicate youth shrunk from no danger. Geneva was then a city of only twelve or fifteen thousand inhabitants, but within its narrow limits it was to become 'the scene of every crisis and every problem, great or small, which can agitate human society.'\footnote{Guizot, p. 210.} It then represented 'a tottering republic, a wavering faith, a nascent Church.' Calvin felt that a negative state of freedom from the tyranny of Savoy and Popery was far worse than Popery itself, and that positive faith and order alone could save the city from political and religious anarchy. He insisted on the abolition of immoral habits, the adoption of an evangelical confession of faith and catechism, the introduction of a strict discipline, Psalm singing, and monthly celebration of the Lord's Supper, with the right of excluding unworthy communicants.\footnote{\emph{Mémoire de Calvin et Farel sur l’organisation de l’église, de Genève,} recently brought to light by Gaberel (\emph{Hist. de l’église de Genève,} 1858, Tom. I. p. 102), and in the Strasburg edition of the \emph{Opera,} Vol. X. Pt. I. pp. 5–14. See a summary in Kampschulte, Vol. I. pp. 287 sqq.}

The magistrate refused to comply, and forbade Calvin and Farel the pulpit; but they, preferring to obey God rather than men, preached at Easter, 1538, to an armed crowd, and declared their determination not to administer the holy communion, lest it be desecrated. On the following day they were deposed and expelled from the city by the great Council of the Two Hundred.

Calvin, again an exile, though now for the principle of authority and discipline rather than doctrine, spent three quiet and fruitful years (1538–41) with Bucer at Strasburg, as teacher of theology and preacher to a congregation of several hundred French refugees.\footnote{Guizot says fifteen hundred. On Calvin's life and labors in Strasburg, see especially the full accounts of Stähelin, Vol. I. pp. 168–318, and Kampschulte, Vol. I. pp. 320–368.} Here he became acquainted with the German Reformation, for Strasburg was the connecting link between Germany and France, as also between Lutheranism and Zwinglianism. But he was disagreeably impressed with the want of Church discipline, and the slavish dependence of the German clergy on the secular rulers. His French congregation was admired for its activity and order. In Strasburg he wrote his tract on the Lord's Supper, his Commentary on the Romans, his masterly answer to Cardinal Sadolet's letter to the Genevese, and his revision of Olivétan's French translation of the Bible. Some of these books attracted the favorable notice of Luther, whom he never met in this world, but always esteemed, with a full knowledge of his faults, as one of the greatest servants of Christ.\footnote{Luther wrote to Bucer: 'Greet Calvin, whose little works I have read with remarkable pleasure;' and Melanchthon wrote: 'Calvin is in high favor here (\emph{magnam gratiam iniit}).' See. Calvin to Farel, Dec. 12, 1539; Stähelin, Vol. I. p. 226; and De Wette's edition of Luther's \emph{Letters,} Vol. V. p. 210. Calvin wrote to Bullinger, when the latter was provoked by the last rude assault of Luther upon the Zwinglians (1544): 'I implore you never to forget how great a man Luther is, and by what extraordinary gifts he excels. Think with what courage, what constancy, what power and success he has devoted himself to this day to the overthrow of the reign of Antichrist and the spreading of the doctrine of salvation far and near. As for me, I have often said, and I say it again, though he should call me a \emph{devil,} I would still give him due honor, and recognize him, in spite of the great faults which obscure his extraordinary virtues, as a mighty servant of the Lord.' See Henry, Vol. II. p. 351; Stähelin, Vol. I. p. 204; Guizot, p. 243; \emph{Opera,} Vol. XI. p. 774.}

In September, 1540, he married Idelette de Bure (a little town in Gueldres), a grave, pious, modest, amiable, and cultivated widow, with three children, whose first husband he had converted from Anabaptism to the orthodox faith. She was in delicate health, but very devoted to him, and satisfied all his desires. He lived with her in perfect harmony nine years, and she bore him one child, a son who died in infancy. He seldom alludes to her in his correspondence, but always in terms of respect and love; and in informing his friend Viret of her departure, he calls her 'the best companion, who would cheerfully have shared with me exile and poverty, and followed me unto death; during her life she was to me a faithful assistant in all my labors; she never dissented from my wishes even in the smallest things.' Seven years afterwards, in a letter of consolation to a friend (Rev. Richard de Valeville, of Frankfort), he says: 'I know from my own experience how painful and burning is the wound which the death of thy wife must have inflicted upon you. How difficult it was for me to become master of my grief. . . . Our chief comfort, after all, is the wonderful providence of God, which overrules our affliction for our spiritual benefit, and separates us from our beloved only to reunite us in his heavenly kingdom.' His grief at her death, and at the death of his child, reveals a hidden spring of domestic affection which is rare in men of his austerity of character and absorption in public duty. He remained a widower the rest of his life.\footnote{Comp. the beautiful tribute to Idelette de Buren, by Jules Bonnet, in the fourth volume of the \emph{Bulletin pour l'histoire du protestantisme français} (1860), and Stähelin, Vol. I. pp. 274–283.}

From the Strasburg period dates also his intimate friendship with Melanchthon, which was not broken by death, and is the more remarkable in view of their difference of opinion on the subject of predestination and free-will. He met him at religious conferences with Romanists, at Frankfort (1539), at Worms (1540), and at Regensburg (Ratisbon, 1541), which he attended as delegate from Strasburg. Their correspondence is a noble testimony to the mind and heart of these great men, so widely different in nationality, constitution, and temper—the one as firm as a rock, the other as timid as a child—and yet one in their deepest relations to Christ and his salvation. They represent the higher union of the Lutheran and Reformed, the Teutonic and the Romanic types of Protestantism. This truly Christian friendship was touchingly expressed by Calvin a year after the death  of the Preceptor of Germany (1561): 'O Philip Melanchthon! for it is upon thee that I call, upon thee, who now livest with Christ in God, and art there waiting for us, until we shall also be gathered with thee to that blessed rest! A hundred times, worn out with fatigue and overwhelmed with care, thou didst lay thy head upon my breast, and say, ``Would to God that I might die here, on thy breast!'' And I, a thousand times since then, have earnestly desired that it had been granted us to be together. Certainly thou wouldst have been more valiant to face danger, and stronger to despise hatred, and bolder to disregard false accusations. Thus the wickedness of many would have been restrained, whose audacity of insult was increased by what they called thy weakness.'\footnote{This passage occurs on the first page of his book against the fanatical Lutheran, Heshusius (\emph{Opera,} Vol. IX. p. 461): '\emph{O Philippe Melanchthon! Te enim appello, qui apud Deum cum Christo vivis, nosque illic expectas, donec tecum in beatam quietem colligamur. Dixisti centies, quum fessus laboribus et molestiis oppressus caput familiariter in sinum meum deponeres: Utinam, utinam moriar in hoc sinu. Ego vero millies postea optavi nobis contingere, ut simul essemus. Certe animosior fuisses ad obeunda certamina, et ad spernendam invidiam, falsasque criminationes pro nihilo ducendas fortior. Hoc quoque modo cohibita fuisset multorum improbitas, quibus ex tua mollitie, quam vocabant, crevit insultandi audacia.}' Comp. on the relation of Calvin to Melanchthon, the full discussion of Stähelin, Vol. I. pp. 230–254; also Guizot, p. 246.}

'It would be difficult,' says Guizot, 'to reconcile truth, piety, and friendship more tenderly.'

In the mean time the Genevese had been brought by sad experience to repent of the expulsion of the faithful pastors, and to feel that the Reformed faith and discipline alone could put their commonwealth on a firm and enduring foundation. The magistrate and people united in an urgent and repeated recall of Calvin. He reluctantly yielded at last, and in September, 1541, after passing a few days with Farel at Neufchatel, he made a triumphant entry into the beautiful city on lake Leman.\footnote{The date is variously given—Sept. 10 by Roget, Sept. 12 by Guizot, Sept. 13 by Kampschulte (following Beza).} The magistracy provided for him a house and garden near the Cathedral of St. Pierre, broadcloth for a coat, and, in consideration of his generous hospitality to strangers and refugees, an annual salary of five hundred florins,\footnote{'Worth about 3600 francs, or £150 at the present time.'—Guizot, p. 257. A syndic received only one fifth of this sum; but Calvin's house was a home for poor refugees of faith from France and other lands, the widows and orphans of martyrs, so that he had often not a penny left. See Stähelin, Vol. II. p. 391, and Hagenbach, \emph{Kirchengesch.} Vol. III. p. 581.} twelve measures of wheat, and two tubs of wine. The rulers of Strasburg, says Beza, stipulated that he should always remain a burgess of their city, and requested him to retain the revenues of a prebend which had been assigned as the salary of his professorship in theology, but they could not persuade him to accept so much as a single farthing.

This second settlement was final. Geneva was now wedded to Calvin, and had to sink or swim with his principles.\footnote{Well says Kampschulte (Vol. I. pp. 385 sq.): '\emph{Genf war im Herbst} 1541 \emph{den geistlichen Tendenzen Calvins dienstbar geworden, es war an den Siegeswagen des Reformators gefesselt und musste ihm folgen trotz allen Sträubens, trotz aller Auflehnungsversuche, die später nicht ausgeblieben sind. Nicht anders fasste Calvin selbst seine Stellung von vorne herein auf. Für ihn ergab sich sein Herrscherrecht über Genf aus dem wunderbaren Gange der letzten Ereignisse mit der Zweifellosigkeit eines von Gott selbst erklärten Glaubenssatzes. Schimpflich vor drei Jahren vertrieben, sah er sich mit den grössten Ehren auf den Schauplatz zurückgeführt, den ihm Farel einst in ernster Stunde ``im Namen des allmächtigen Gottes'' angewiesen: mit Jubel wurde er von demselben Volke begrüsst, das ihm unversöhnlichen Hass geschworen! . . . Calvin fühlte sich fast nur noch als Werkzeug in der Hand Gottes, durch den ewigen göttlichen Rathschluss, ohne jedes persönliche Zuthun, für Genf bestimmt, um des Herrn Willen, wie er ihn erkannt, auf diesem wichtigen Fleck der Erde ohne Furcht und Scheu zu verkündigen, jenes Programm, welches er in der christlichen Institution niedergelegt, hier zur Ausführung zu bringen, dem Herrn hier ein christliches Geschlecht zu sammeln, das der übrigen Welt als Leuchte diene.}'} He continued to labor there, without interruption, for twenty-three years, till his death, May 27, 1564: fighting a fierce spiritual war against Romanism and superstition, but still more against infidelity and immorality; establishing a model theocracy on the basis of Moses and Christ; preaching and teaching from day to day; writing commentaries, theological and polemical treatises; founding an academy, which in the first year attracted more than eight hundred students, and flourishes to this day; attending the sessions of the consistory and the senate; entertaining and counselling strangers from all parts of the world; and corresponding in every direction. He was, in fact, the spiritual head of the Church and the republic of Geneva, and the leader of the Reformed movement throughout Europe. And yet he lived all the time in the utmost simplicity. It is reported that Cardinal Sadolet, when passing through Geneva \emph{incognito,} and calling on Calvin, was surprised to find him residing, not in an episcopal palace, with a retinue of servants, as he expected, but in a little house, himself opening the door. The story may not be sufficiently authenticated, but it corresponds fully with all we know about his ascetic habits.\footnote{This fact is related by Drelincourt in his \emph{Defense de Calvin} (1667), and Bungener (p. 503), and is believed in Geneva, but doubted by Guizot, p. 237, for chronological reasons which are not conclusive (Sadolet died 1549). '\emph{Se non e vero, e ben trovato.}'} For years he took but one meal a day.\footnote{Beza: '\emph{Per decem minimum annos prandio abstinuit, ut nullum omnino cibum extra statam cœnæ horam sumeret.}' Sometimes he abstained for thirty-six hours from all food.} He refused an increase of salary and presents of every description, except for the poor and the refugees, whom he was always ready to aid. He left, besides his library, only about two hundred dollars, which he bequeathed to his younger brother Anthony and his children.\footnote{See his testament in Beza's \emph{Vita.}} When Pope Pius IV. heard of his death, he paid him this high compliment: 'The strength of that heretic consisted in this, that money never had the slightest charm for him. If I had such servants, my dominions would extend from sea to sea.'\footnote{Quoted by Guizot, p. 361.}

His immense labors and midnight studies,\footnote{'\emph{Somni pene nullius,}' says Beza in his closing remarks.} the care of all the churches, and bodily infirmities—such as headaches, asthma, fever, gravel—gradually wore out his delicate body. He died, in full possession of his mental powers, in the prime of manhood and usefulness, not quite fifty-five years of age, leaving his Church in the best order and in the hands of an able and faithful successor, Theodore Beza. Like a patriarch, he assembled first the syndics of Geneva, and afterwards the ministers, around his dying bed, thanked them for their kindness and devotion, asked humbly their pardon for occasional outbursts of violence and wrath, and affected them to tears by words of wisdom and counsel to persevere in the pure doctrine and discipline of Christ. It was a sublime scene, worthily described by Beza,\footnote{With Beza's account of his parting addresses (in the French and Latin edition of the \emph{Vita}) should be compared the official copy, which Bonnet published in the Appendix to the \emph{French Letters,} Tom. II. p. 573, and the Strasburg editors at the close of the 9th vol. of the \emph{Opera} (\emph{Discours d’adieu aux membres du Petit Conseil,} pp. 887–890, and \emph{Discours d’adieu aux ministres,} pp. 891–894). Comp. also Stähelin. Vol. II. pp. 462–468.} and well represented by a painter's skill.\footnote{Hornung's picture of Calvin on his death-bed.}

The Reformer died with the setting sun. 'Thus,' says Beza, 'God withdrew into heaven that most brilliant light, which was a lamp of the Church. In the following night and day there was immense grief and lamentation in the whole city; for the republic had lost its wisest citizen, the Church its faithful shepherd, the academy an incomparable teacher—all lamented the departure of their common father and best comforter next to God. A multitude of citizens streamed to the death-chamber, and could scarcely be separated from the corpse. Among them also were several foreigners, as the distinguished English embassador to France, who had come to Geneva to make the acquaintance of the celebrated man. On the Lord's day, in the afternoon, the remains were carried to the common graveyard on Plainpalnis, followed by all the patricians, pastors, professors, and teachers, and nearly the whole city, in sincere mourning.'

Calvin expressly forbade the erection of any monument over his grave.\footnote{Beza, however, wrote a suitable poem, in Latin and French, which might have been inscribed on the tomb. See his \emph{Vita,} at the close, and \emph{Opera,} Vol. V. pp. xxvi. sqq. (with three other French sonnets); a German translation in Stähelin, Vol. II. p.470.} The stranger asks in vain even for the spot which covers his mortal remains in the cemetery of Geneva. Like Moses, he was buried out of the reach of idolatry. The Reformed Churches of both hemispheres are his monument, more enduring than marble. On the third centenary of his death (1864), his friends in Geneva, aided by gifts from foreign lands, erected to his memory the \emph{Salle de la Réformation}—a noble building, founded on the principles of the Evangelical Alliance, and dedicated to the preaching of the pure gospel and the advocacy of every good cause.

CALVIN'S PERSONAL CHARACTER.

Calvin was of middle, or rather small stature (like David and Paul), of feeble health, courteous, kind, grave and dignified in deportment. He had a meagre and emaciated frame, a thin, pale, finely chiseled face, a well-formed mouth, a long, pointed beard, black hair, a prominent nose, a lofty forehead, and flaming eyes. He was modest, plain, and scrupulously neat in dress, orderly and methodical in all his habits, temperate and even abstemious, allowing himself scarcely nourishment and sleep enough for vigorous work. His physical tent barely covered the mighty spirit within. Conscience and logic, a commanding mind and will, shone through the thin veil of mortality.\footnote{See different portraits of Calvin—in Henry (small biography), in first volume of the \emph{Opera,} in Stähelin, in first volume of Merle d'Aubigné; also Hornung's Calvin on his death-bed, and the medallion portrait made at the festival of the Geneva Reformation. \par{}  [In technical disregard of Calvin's wish the large mural monument was erected in Geneva, 1917, commemorating the Reformation and containing figures of Calvin, Luther, etc.—Ed.]}

How different Luther and Zwingli, with their strong animal foundation, and their abundance of flesh and blood! Calvin seemed to be all bone and nerve. Beza says he looked in death almost the same as alive in sleep.\footnote{Beza thus tersely describes him (at the close of the \emph{Vita}): '\emph{Statura fuit mediocri, colore subpallido et nigricante, oculis ad mortem usque limpidis, quique ingenii sagacitatem testarentur: cultu corporis neque culto neque sordido, sed qui singularem modestiam deceret: victu sic temperato, ut a sordibus et ab omni luxu longissime abesset: cibi parcissimi, ut qui multos annos semel quotidie cibum sumpserit, vintriculi imbecillitatem causatus: somni pæne nullius: memoriæ incredibilis, ut quos semel aspexisset multis post annis statim agnosceret, et inter dictandum sæpe aliquot horas interturbatus statim ad dictata nullo commonefaciente rediret, et eorum, quæ ipsum nosse muneris sui causa interesset, quantumvis multiplicibus et infinitis negotiis oppressus, nunquam tamen oblivisceretur. Judicii, quibuscunque de rebus consuleretur, tam puri et exacti, ut pæne vaticinari sæpe sit visus, nec aberasse meminerim, qui consilium ipsius esset sequutus. Facundiæ contemptor et verborum parcus, sed minime ineptus scriptor, et quo nullus ad hunc diem theologus (absit verbo invidia) purius, gravius, judiciosius denique scripsit, quum tamen tam multa scripserit, quam nemo vel nostra vel patrum memoria.}'}

His intellectual endowments were of the highest order and thoroughly disciplined. He had more constructive, systematizing, and organizing genius than any other Reformer, and was better adapted to found a solid, compact, and permanent school of theology. He was not a speculative or intuitive philosopher, but a consummate logician and dialectician. Luther and Zwingli cut the stones from the quarry; Calvin gave them shape and polish, and erected a magnificent cathedral of ideas with the skill of a master architect. His precocity and consistency were marvelous. He did not grow before the public, like Luther and Melanchthon, and pass through contradictions and retractations, but when a mere youth of twenty-six he appeared fully armed, like Minerva from the head of Jupiter, and never changed his views on doctrine or discipline. He had an extraordinary and well-stored memory, a profound, acute, and penetrating intellect, a clear, sound, and almost unerring judgment, a perfect mastery over the Latin and French tongues. His Latin is as easy and elegant, and certainly as nervous and forcible, as Cicero's, yet free from the pedantic and affected purism of a Bembo and Castalio.\footnote{Who would substitute \emph{respublica} for \emph{ecclesia, genius} for \emph{angelus, lotio} for \emph{baptismus,} etc.} He is one of the fathers of modern French, as Luther is the father of modern German. His eloquence is logic set on fire by intense conviction. His Preface to the 'Institutes,' addressed to the King of France, is reckoned as one of the three immortal prefaces in literature (to which only that of President De Thou to his French History and that of Casaubon to Polybius can be compared); and his 'Institutes' themselves, as has been well said, are 'in truth a continuous oration, in which the stream of discussion rolls onward with an impetuous current, yet always keeps within its defined channel.'\footnote{Fisher, \emph{The Reformation,} p. 198.}

He surpassed all other Reformers (except Beza) in classical culture and social refinement. He was a patrician by education and taste, and felt more at ease among scholars and men of high rank than among the common people. Yet he was quite free from aristocratic pride, despised all ostentation and display, and esteemed every man according to his real worth.

History furnishes, perhaps, no example of a man who with so little personal popularity had such influence upon the people, and who with such natural timidity and bashfulness combined such strength and control over his age and future generations. Constitutionally a retiring scholar and a man of thought, he became providentially a mighty man of action and an organizer of churches.

His moral and religious character is impressed with a certain majesty which keeps the admirer at a respectful distance.\footnote{This was the judgment of the magistrate of Geneva, expressed in these words (June 8, 1564): '\emph{Dieu, lui avait imprimé un charactère d’une si grande majesté.}'} He has often been compared to an old Roman Censor or Stoic; but he resembles much more a Hebrew Prophet. Severe against others, he was far more severe against himself, and was always guided by a sense of duty. Fear of God, purity of motive, spotless integrity, single devotion to truth and duty, unswerving fidelity, sincere humility are the prominent traits of his character. Soaring high above the earth, he was absorbed in God—who alone is great—and looked down upon man as a fleeting shadow. The glory of the Lord and the reformation of the Church constituted the single passion of his life. His appropriate symbol was a hand offering the sacrifice of a bleeding heart to God.\footnote{'\emph{Cor meum velut mactatum Domino in sacrificium offero.}' Subscribed below his autograph in the frontispiece of Henry's smaller biography.}

It must be admitted that this kind of greatness, while it commands our admiration and respect, does not of itself secure our affection and love. There is a censoriousness and austerity about Calvin and his creed which repelled many good men, even among his contemporaries.\footnote{His ungrateful enemy, Balduin, started the saying among the Genevese, 'Rather with Beza in hell than with Calvin in heaven.' And yet they obeyed and revered him. Beza, it should be remembered, was the perfection of a French gentleman; yet his theological system was even more severe than that of Calvin, and he carried the dogma of predestination to the extreme of supralapsarianism. I have met with not a few French, Scotch, and American Christians who, in the combination of severity and purity, gravity and kindliness of character, reminded me strongly of Calvin and Beza. I may mention Gaussen, Malan, Merle d'Aubigné, Pio ie Adolph Monod, and Guizot.} He looked more to the holiness than to the love of God. His piety bears more the stamp of the Old Testament than that of the New. He represents the majesty and severity of the law rather than the sweetness and loveliness of the gospel, the obedience of a servant of Jehovah rather than the joyfulness of a child of our heavenly Father.

Yet even this must be qualified. He sympathized with the spirit of David and Paul as much as with the spirit of Moses and Elijah, and had the strongest sense of the freedom of the gospel salvation. Moreover, behind his cold marble frame there was beating a noble, loving, and faithful heart, which attracted and retained to the last the friendship of such eminent servants of God as Farel, Viret, Beza, Bucer, Bullinger, Knox, and Melanchthon. 'He obtained,' says Guizot, 'the devoted affection of the best men and the esteem of all, without ever seeking to please them.'\footnote{Page 362.} John Knox, his senior in years, sat at his feet as a humble pupil, and esteemed him the greatest man after the Apostles. Farel, in his old age, hastened on foot from Neufchatel to Geneva to take leave of his sick friend, and desired to die in his place. Beza, who lived sixteen years on terms of personal intimacy with him, revered and loved him as a father. And even Melanchthon wished to repose and to die on his bosom. His familiar correspondence shows him in the most favorable light, and is a sufficient refutation of all the calumnies and slanders of his enemies.

He lacked the good-nature, the genial humor, the German \emph{Gemüthlichkeit,} the overflowing humanity of Luther, who for this reason will always be more popular with the masses; but he surpassed him in culture, refinement, consistency, and moral self-control. Both were equally unselfish and unworldly. Both were headstrong and will-strong; but Calvin was more open to argument and less obstinate. Both had, like St. Paul, a fiery and violent temper, which was the propelling force in their hard work, and in fierce battles with the pope and the devil. Hegel says somewhere that 'nothing great can be done without passion.'\footnote{'\emph{Nichts Grosses geschieht ohne Leidenschaft.}'} It is only men of intense convictions and fearless courage that make deep and lasting impressions upon others. But temper is a force of nature, which must be controlled by reason and regulated by justice and charity. Luther came down like a thunder-storm upon his opponents, and used the crushing sledge-hammer indiscriminately against Eck, Cochlæus, Henry VIII., Erasmus, the Sacramentarians, and Zwinglians; while Calvin wielded the sharp sword of irony, wit, scorn, and contempt in defense of truth, but never from personal hatred and revenge. 'Even a dog barks,' he says, 'when his master is attacked; how could I be silent when the honor of my Lord is assailed?'\footnote{The strongest terms of Calvin against ferocious enemies are \emph{canes, porci, bestial, nebulones} (with reference, no doubt, to Scripture usage— Isa. lvi. 10; Matt. vii. 6; Phil. iii. 2; Rev. xxii. 15); but they are mild compared to the coarse and vulgar epithets with which Luther overwhelmed his opponents, without expressing any regret afterwards, except in the case of Henry VIII., where it was least needed, and made the matter worse.} He confessed, however, in a letter to Bucer, and on his death-bed, that he found it difficult to tame 'the wild beast' of his wrath, and humbly asked forgiveness for his weakness. He had no children to write to, and to play with around the Christmas-tree, like Luther, but he appears to better advantage in his relations with men and women. He treated them, even the much younger Beza, as equals, overlooked minor differences, and in correcting their faults expected the same manly frankness from them in return; while Luther, growing more irritable and overbearing with advancing years, made even Melanchthon tremble and fear. But we should charitably remember that the faults of these truly great and good men were only the long shadows of their extraordinary virtues.\footnote{Calvin, though fully aware of the defects of Luther, often expressed his admiration for him (see p. 430), and in January, 1545 (a year before Luther's death), he sent him a letter (which Melanchthon was afraid to hand to the old lion on account of his excited state of feeling against the Swiss), closing with these touching words: 'If I could only fly to you and enjoy your society, even for a few hours! . . . But since this privilege is not granted to me on earth, I hope I may soon enjoy it in the kingdom above. Farewell, most illustrious man, most excellent minister of Christ and father [\emph{pater,} al. \emph{frater}], forever venerable to me. May the Lord continue to guide you by his Spirit to the end for the common good of his Church.' \emph{Opera,} Vol. XII. p. 8.}

It may be found strange that Calvin never alludes to the paradise of nature by which he was surrounded on the lovely shores of Lake Leman, in sight of the lofty Alps that pierce the skies in silent adoration of their Maker. But we look in vain for descriptions of natural scenery in the whole literature of the sixteenth century; and the proper appreciation of the beauties of Switzerland, as well as of other countries, is of more recent date. Calvin had no special organ nor time for the enjoyment of the beautiful either in nature or in art, but he appreciated poetry and music.\footnote{Guizot says (p. 164): 'Although Calvin was devoted to the severe simplicity of evangelical worship, he did not overlook the inherent love of mankind for poetry and art. He himself had a taste for music, and knew its power. He feared that, in a religious service limited to preaching and prayer only, the congregation, having nothing else to do than to play the part of audience, would remain cold and inattentive. For this reason he attached great importance to the introduction and promotion of the practice of Psalm-singing in public worship. ``If the singing,'' he said, ``is such as befits the reverence which we ought to feel when we sing before God and the angels, it is an ornament which bestows grace and dignity upon our worship; and it is an excellent method of kindling the heart, and making it burn with great ardor in prayer. But we must at all times take heed lest the ear should be more attentive to the harmony of the sound than the soul to the hidden meaning of the words'' (\emph{Instit.} Ch. XX.). With this pious warning, he strongly urged the study of singing, and its adaptation to public worship.' Comp. Gaberel, Vol. I. p. 353.} He insisted on the introduction of congregational singing in Geneva, and wrote himself a few poetic versions of the Psalms, and a hymn of praise to Christ, which are worthy of Clement Marot and reveal an unexpected vein of poetic fervor and tenderness.\footnote{These poetic pieces were recently discovered, and published in the sixth volume of the new edition of his \emph{0pera} (1867), pp. 212–224. His \emph{Salutation à Jésus-Christ} was translated into German by Stähelin, and into English by Mrs. Smith, of New York, for Schaff's \emph{Christ in Song,} London edition, p. 549. His \emph{Epinicion Christo cantatum} is a polemic poem in Latin hexameters and pentameters, composed during the Conference at Worms, 1541, in which he describes the Romish polemics Eck, Cochlæus, Nausea, and Pelargus as dragged after the chariot of the victorious Redeemer. \emph{Opera.} Vol. V. pp. 417–428.} The following specimen must suffice:

'I greet thee, who my sure Redeemer art,

My only trust, and Saviour of my heart!

Who so much toil and woe

And pain didst undergo,

For my poor, worthless sake:

We pray thee, from our hearts,

All idle griefs and smarts

And foolish cares to take.

'Thou art the true and perfect gentleness,

No harshness hast thou, and no bitterness:

Make us to taste and prove,

Make us adore and love,

The sweet grace found in thee;

With longing to abide

Ever at thy dear side,

In thy sweet unity.

'Poor, banished exiles, wretched sons of Eve,

Full of all sorrows, unto thee we grieve;

To thee we bring our sighs,

Our groanings, and our cries:

Thy pity, Lord, we crave;

We take the sinner's place,

And pray thee, of thy grace,

To pardon and to save.'

TRIBUTES TO CALVIN.

I add some estimates of Calvin's character, which represent very different stand-points.\footnote{We omit Henry and Stähelin, from whom it would he difficult to select passages in praise of Calvin. See especially the entire Seventh Book of Stähelin, Vol. II., pp. 365–393: \emph{Calvin als Mensch und als Christ.}}

Beza, who knew Calvin best and watched at his death-bed, concludes his biography with these words:

'Having been an observer of Calvin's life for sixteen years, I may with perfect right testify that we have in this man a most beautiful example of a truly Christian life and death, which it is easy to calumniate but difficult to imitate.'\footnote{'\emph{Ego historiam vitæ et obitus ipsius, cujus spectator sedecim annos fui, bona fide persequutus testari mihi optimo jure posse videor, longe pulcherrimum vere Christianæ tum vitæ tum mortis exemplum in hoc homine cunctis propositum fuisse, quod tam facile sit calumniari, quam difficile fuerit æmulari.}'}

Bungener, a pastor of the national Church of Geneva, and author of several historical works, says:\footnote{\emph{Calvin,} etc. English translation, pp. 338, 349.}

'Let us not give him praise which he would not have accepted. God alone creates; a man is great only because God thinks fit to accomplish great things by his instrumentality. Never did any great man understand this better than Calvin. It cost him no effort to refer all the glory to God; nothing indicates that he was ever tempted to appropriate to himself the smallest portion of it. Luther, in many a passage, complacently dwells on the thought that a petty monk, as he says, has so well made the Pope to tremble, and so well stirred the whole world. Calvin will never say any such thing; he never even seems to say it, even in the deepest recesses of his heart: every where you perceive the man, who applies to all things—to the smallest as to the greatest—the idea that it is God who does all and is all. Read again, from this point of view, the very pages in which he appeared to you the haughtiest and most despotic, and see if, even there, he is any thing other than the workman referring all, and in all sincerity, to his Master. . . . But the man, in spite of all his faults, has not the less remained one of the fairest types of faith, of earnest piety, of devotedness, and of courage. Amid modern laxity, there is no character of whom the contemplation is more instructive; for there is no man of whom it has been said with greater justice, in the words of an apostle, ``\emph{he endured as seeing him who is invisible.}'''

Jules Michelet, the French historian, remarks:\footnote{In his \emph{Histoire de France au seizième siècle,} quoted by Stähelin, Vol. I. p. 276.}

'Among the martyrs, with whom Calvin constantly conversed in spirit, he became a martyr himself; he felt and lived like a man before whom the whole earth disappears, and who tunes his last Psalm, his whole eye fixed upon the eye of God, because he knows that on the following morning he may have to ascend the stake.'

Ernest Renan, once educated for the Romish priesthood, then a skeptic, with all his abhorrence of Calvin's creed, pays the following striking tribute to his character:\footnote{In his article on Jean Calvin, above quoted, pp. 286, etc. The translation is by O. B. Frothingham, a radical Unitarian in New York.}

Calvin was one of those absolute men, cast complete in one mould, who is taken in wholly at a single glance: one letter, one action suffices for a judgment of him. There were no folds in that inflexible soul, which never knew doubt or hesitation. . . . Careless of wealth, of titles, of honors, indifferent to pomp, modest in his life, apparently humble, sacrificing every thing to the desire of making others like himself, I hardly know of a man, save Ignatius Loyola, who could match him in these terrible transports. . . . It is surprising that a man who appears to us in his life and writings so unsympathetic should have been the centre of an immense movement in his generation, and that this harsh and severe tone should have exerted so great an influence on the minds of his contemporaries. How was it, for example, that one of the most distinguished women of her time, Renée of France, in her court at Ferrara, surrounded by the flower of European wits, was captivated by that stern master, and by him drawn into a course that must have been so thickly strewn with thorns? This kind of austere seduction is exercised by those only who work with real conviction. Lacking that vivid, deep, sympathetic ardor which was one of the secrets of Luther's success, lacking the charm, the perilous, languishing tenderness of Francis of Sales, Calvin succeeded, in an age and in a country which called for a reaction towards Christianity, simply because he was the most christian man of his generation.'

Guizot, a very competent judge of historical and moral greatness, thus concludes his biography:\footnote{\emph{St. Louis and Calvin,} pp. 361 and 362.}

'Calvin is great by reason of his marvelous powers, his lasting labors, and the moral height and purity of his character. . . . Earnest in faith, pure in motive, austere in his life, and mighty in his works, Calvin is one of those who deserve their great fame. Three centuries separate us from him, but it is impossible to examine his character and history without feeling, if not affection and sympathy, at least profound respect and admiration for one of the great Reformers of Europe and of the great Christians of France.'

Prof. Kahnis, of Leipzig, whose personal and theological sympathies are with Luther, nevertheless asserts the moral superiority of Calvin above the other Reformers:\footnote{\emph{Die Lutherische Dogmatik,} Vol. II. pp. 490, 491.}

'The fear of God was the soul of his piety, the rock-like certainty of his election before the foundation of the world was his power, and the doing of the will of God his single aim, which he pursued with trembling and fear. . . . No other Reformer has so well demonstrated the truth of Christ's word that, in the kingdom of God, dominion is service. No other had such an energy of self-sacrifice, such an irrefragable conscientiousness in the greatest as well as the smallest things, such a disciplined power. This man, whose dying body was only held together by the will flaming from his eyes, had a majesty of character which commanded the veneration of his contemporaries.'

Prof. Dorner, of Berlin, the first among the theologians of the age, distinguished by profound learning, penetrating thought, rare catholicity of spirit, and nice sense of justice and discrimination, says:

'Calvin was equally great in intellect and character, lovely in social life, full of tender sympathy and faithfulness to friends, yielding and forgiving towards personal offenses, but inexorably severe when he saw the honor of God obstinately and malignantly attacked. He combined French fire and practical good sense with German depth and soberness. He moved as freely in the world of ideas as in the business of Church government. He was an architectonic genius in science and practical life, always with an eye to the holiness and majesty of God.'\footnote{\emph{Geschichte der Protest. Theologie,} pp. 374 and 376. I add his considerate judgment of Calvin in full: '\emph{Die nach Zwingli's and Œcolampad's Tode verwaiste reformirte Kirche erhielt am} Johann Calvin, \emph{gleich gross an Geist und Character, einen festen Mittelpunkt und eine ordnende Seele für Lehre und Kirchenverfassung. Durch ihn wurde Genf statt Zürichs die neue reformirte Metropole; und dieses Gemeinwesen bewies eine wunderbare, weithin erobernde Kraft. . . . Calvin's persönliche Erscheinung war die eines altrömischen Censors; er war von feinem Wuchs, blass, hager, mit dem Ausdruck tiefen Ernstes und einschneidender Schärfe. Der Senat von Genf sagte nach seinem Tode, er sei ein majestätischer Charakter gewesen. Liebenswürdig im socialen Leben, voll zarter Theilnahme und Freundestreue, nachsichtig und versöhnlich bei persönlichen Beleidigungen, war er unerbittlich streng, wo er Gottes Ehre in Hartnäckigkeit oder Bosheit angegriffen sah. Unter seinen Collegen hatte er keine Neider, aber viele begeisterte Verehrer. Französisches Feuer und praktischer Verstand schienen mit deutscher Tiefe und Besonnenheit einen Bund geschlossen zu haben. War er auch nicht spekulativen oder intuitiven Geistes, so war dagegen sein Verstand und sein Urtheil um so eindringender und schärfer, sein Gedächtniss umfassend; und er bewegte sich ebenso leicht in der Welt der Ideen und der Wissenschaft, wie in den Geschäften des Kirchenregiments. Zwar ist er nicht ein Mann des Volkes, wie Luther, sondern in seiner Sprache mehr der Gelehrte, und seine Wirksamkeit als Prediger und Seelsorger kann daher mit der Luthers nicht verglichen werden. Dagegen ist er mehr ein architektonischer Geist und zwar sowohl im Gebiete der Wissenschaft als des Lebens. Beide sind ihm in ihrer Wurzel eins, und seine dogmatischen Constructionen, so kühn sie in der Folgerichtigkeit ihrer Gedanken sind, behalten ihm doch immer zugleich erbaulichen Charakter. Auch wo er verwegen in die göttlichen Geheimnisse der Prädestination einzudringen sucht, immer leitet ihn der praktische Trieb, der Heiligkeit und Majestät Gottes zu dienen, für das Gemüth aber den ewigen Ankergrund zu finden, darin es im Bewusstsein der Erwählung durch freie Gnade sicher ruhen könne.}'}

Prof. Gr. T. Fisher, of Yale College, New Haven, gives the following fair and impartial estimate of Calvin:\footnote{\emph{The Reformation,} pp. 206 and 238.}

'When we look at his extraordinary intellect, at his culture—which opponents, like Bossuet, have been forced to commend—at the invincible energy which made him endure with more than stoical fortitude infirmities of body under which most men would have sank, and to perform, in the midst of them, an incredible amount of mental labor; when we see him, a scholar naturally fond of seclusion, physically timid, and recoiling from notoriety and strife, abjuring the career that was most to his taste, and plunging, with a single-hearted, disinterested zeal and an indomitable will, into a hard, protracted contest; and when we follow his steps, and see what things he effected, we can not deny him the attributes of greatness. . . . His last days were of a piece with his life. His whole course has been compared by Vinet to the growth of one rind of a tree from another, or to a chain of logical sequences, He was endued with a marvelous power of understanding, although the imagination and sentiments were less roundly developed. His systematic spirit fitted him to be the founder of an enduring school of thought. In this characteristic he may be compared with Aquinas. He has been appropriately styled the Aristotle of the Reformation. He was a perfectly honest man. He subjected his will to the eternal rule of right, as far as he could discover it. His motives were pure. He felt that God was near him, and sacrificed every thing to obey the direction of Providence. The fear of God ruled in his soul; not a slavish fear, but a principle such as animated the prophets of the Old Covenant. The combination of his qualities was such that he could not fail to attract profound admiration and reverence from one class of minds, and excite intense antipathy in another. There is no one of the Reformers who is spoken of, at this late day, with so much personal feeling, either of regard or aversion. But whoever studies his life and writings, especially the few passages in which he lets us into his confidence and appears to invite our sympathy, will acquire a growing sense of his intellectual and moral greatness, and a tender consideration for his errors.'

\xsubsection{Calvin's Work. His Theology and Discipline.}

§ 57. Calvin's Work.

Of Calvin it may be said, without exaggeration, that he 'labored more' than all the other Reformers.

He raised the little town of Geneva to the dignity and importance of the Protestant Rome.\footnote{The eminent French historian, H. Martin (in his \emph{Histoire de France depuis les temps les plus reculés jusqu’en} 1789, Tom. VIII. p. 325 of the fourth edition, Par. 1860), thus speaks of what Calvin did for the city of Geneva: '\emph{Calvin ne la sauve pas seulement, mais conquiert à cette petite ville une grandeur, une puissance morale immense. Il en fait la capitale de la Réforme, autant que la Réforme peut avoir une capitale, pour la moitié du monde protestante, avec une vaste influence, acceptée ou subie, sur l’autre moitié. Genève n’est rien par la population, par les armes, par le territoire: elle est tout par l’esprit. Un seul avantage matériel lui garantit tous ses avantages moraux: son admirable position, qui fait d’elle une petite France républicaine et protestante, indépendante de la monarchie catholique de France et à l’abri de l’absorption monarchique et catholique; la Suisse protestante, alliée necessaire de la royauté française contre l’empereur, couvre Genève par la politique vis-à-vis du roi et par l’épée contre la maison d’Autriche et de Savoie.}'}

From this radiating centre he controlled, directly or indirectly, through his writings and his living disciples, the Reformed, yea, we may say, the whole Protestant movement; for, wherever it had not already taken root, as in Germany and Scandinavia, Protestantism assumed a Calvinistic or semi-Calvinistic character.\footnote{Kampschulte, Vol. I. p. xii.: \emph{Der romanische Reformator zählte seine Anhänger in der romanischen, germanischen und slavischen Welt und zeigte sich überall, wo nicht das Lutherthum in dem deutschen Character eine Stütze fand, diesem überlegen.}' He quotes the fact that in Bohemia, which borders on Germany, the Slavonian Protestants nearly all profess Calvinism, while Lutheranism is confined to the Germans. The same is still more the case with the Anglo-Saxon race in England, America, and Australia, and in the mission fields among the heathen. In Italy and Spain, too, the Waldenses and the evangelical Churches are, both in doctrine and discipline, much more Calvinistic than Lutheran; but so far Protestantism has a very feeble hold on the Latin races, which are more apt to swing from popery to infidelity, and from infidelity to popery, than to adopt the \emph{via media} either of Lutheranism or Calvinism or Anglicanism.}

His heart continued, indeed, to beat warmly for his native land, which he reluctantly left to share the fortunes of truth exiled, and he raised the cry which is to this day the motto of his faithful disciples: 'France must be evangelized to be saved.' But his true home was the Church of God. He broke through all national limitations. There was scarcely a monarch or statesman or scholar of his age with whom he did not come in contact. Every people of Europe was represented among his disciples. He helped to shape the religious character of churches and nations yet unborn. The Huguenots of France, the Protestants of Holland and Belgium, the Puritans and Independents of England and New England, the Presbyterians of Scotland and throughout the world, yea, we may say, the whole Anglo-Saxon race, in its prevailing religious character and institutions, bear the impress of his genius, and show the power and tenacity of his doctrines and principles of government.\footnote{'In his vast correspondence we find him conversing familiarly with the Reformers—Farel, Viret, Beza, Bullinger, Bucer, Grynæus, Knox, Melanchthon—on the most important religious and theological questions of his age; counseling and exhorting Prince Condé, Jeanne D'Albret, mother of Henry IV., Admiral Coligny, the Duchess of Ferrara, King Sigismund of Poland, Edward VI. of England, and the Duke of Somerset; respectfully reproving Queen Marguerite of Navarre; withstanding libertines and the pseudo-Protestants; strengthening the martyrs, and directing the Reformation in Switzerland, France, Poland, England, and Scotland. He belongs to the small number of men who have exerted a moulding influence, not only upon their own age and country, but also upon future generations in various parts of the world; and not only upon the Church, but indirectly also upon the political, moral, and social life. The history of Switzerland, Germany, France, the Netherlands, Great Britain, and the United States for the last three hundred years bears upon a thousand pages the impress of his mind and character. He raised the small republic of Geneva to the reputation of a Protestant Rome. He gave the deepest impulse to the Reform movement, which involved France, his native land, in a series of bloody civil wars, which furnished a host of martyrs to the evangelical faith, and which continues to live in that powerful nation in spite of the horrid massacre of St. Bartholomew and the revocation of the Edict of Nantes, the dragoonades and exile of hosts of Huguenots, who, driven from their native soil, carried their piety, virtue, and industry to all parts of Western Europe and North America. He kindled the religious fire which roused the moral and intellectual strength of Holland, and consumed the dungeons of the Inquisition and the fetters of the political despotism of Spain. His genius left a stronger mark on the national character of the Anglo-Saxon race and the Churches of Great Britain than their native Reformers. His theology and piety raised Scotland from a semi-barbarous condition, and made it the classical soil of Presbyterian Christianity, and one of the most enlightened, energetic, and virtuous countries on the face of the globe. His spirit stirred up the Puritan revolution of the seventeenth century, and his blood ran in the veins of Hampden and Cromwell, as well as Baxter and Owen. He may be called, in some sense, the spiritual father of New England and the American republic. Calvinism, in its various modifications and applications, was the controlling agent in the early history of our leading colonies (as Bancroft has shown); and Calvinism is, to this day, the most powerful element in the religious and ecclesiastical life of the Western world.'—From the author's Essay on Calvin, in the \emph{Bibl. Sacra} for 1857.}

From him proceeded the first Protestant missionary colony in the newly discovered American Continent.\footnote{On the interesting French colony in Brazil, 1556, consisting of two clergymen and about two hundred members of the Church of Geneva, see Stähelin, Vol. II. pp. 234 sqq. The colony was broken up by the interference of the French government and by Papal intrigues. But it was a harbinger of the later emigrations of persecuted Huguenots in several parts of North America, who enriched the Presbyterian, Dutch, and German Reformed and other Churches.}

He conceived the idea of a general Evangelical Alliance which, though impracticable in his age, found an echo in Melanchthon and Cranmer, and was revived in the nineteenth century (1846) to be realized at no distant future.\footnote{Comp. Stähelin, Vol. II. pp. 198, 241.}

His work and influence were twofold, theological and ecclesiastical. With him theory and practice, theology and piety, were inseparably united. Even when, soaring beyond the limits of time, he dared to lift the veil of the eternal decrees of the omniscient Jehovah, he aimed at a strong motive for holiness, and a firm foundation of hope and comfort. On the other hand, his moral reforms are all based upon principles and ideas. He was thoroughly consistent in his views and actions.

HIS THEOLOGY.

As a scientific theologian, Calvin stands foremost among the Reformers, and is the peer of Augustine and Thomas Aquinas. He has been styled the Aristotle of Protestantism. Melanchthon, 'the Teacher of Germany,' first called him 'the Theologian,' in the emphatic sense in which this title was given to Gregory of Nazianzen in the Nicene age, and to the inspired Apostle John. The verdict of history has confirmed this judgment. Even Rationalists and Roman Catholics must admit his pre-eminence among the systematic divines and exegetes of all ages.\footnote{The Strasburg editors of Calvin's Works, though belonging to the modern liberal school of theology, thus characterize him as a theologian (\emph{Opera,} Vol. I. p. ix.): '\emph{Si Lutherum virum maximum, si Zwinglium civem Christianum nulli secundum, si Melanthonem præceptorem doctissimum merito appellaris, Calvinum jure vocaris } theologorum principem et antesignanum.  \emph{In hoc enim quis linguarum et literarum præsidia, quis disciplinarum fere omnium non miretur orbem? De cujus copia doctrinæ, rerumque dispositione aptissime concinnata, et argumentorum vi ac validitate in dogmaticis; de ingenii acumine et subtilitate, atque nunc festiva nunc mordaci salsedine in polemicis, de felicissima perspicuitate, sobrietate ac sagacitate in exegeticis, de nervosa eloquentia et libertate in paræneticis; de prudentia sapientiaque legislatoria in ecclesiis constituendis, ordinandis ac regendis incomparabili, inter omnes viros doctos et de rebus evangelicis libere sentientes jam abunde constat. Imo inter ipsos adversarios romanos nullus hodie est, vel mediocri harum rerum cognitione imbutus vel tantilla judicii præditus æquitate, qui argumentorum et sententiarum ubertatem, proprietatem verborum sermonemque castigatum, stili denique, tam latini quam gallici, gravitatem et luciditatem non admiretur. Quæ cuncta quum in singulis fere scriptis, tum præcipue relucent in immortali illa Institutione religionis Christianæ, quæ omnes ejusdem generis expositiones inde ab apostolorum temporibus conscriptas, adeoque ipsos Melanthonis Locos theologicos, absque omni controversia longe antecellit atque eruditum et ingenuum lectorem, etiamsi alicubi secus senserit, hodieque quasi vinetum trahit et vel invitum rapit in admirationem.}' To this we add a remarkable tribute of a liberal Roman Catholic historian who abhors Calvin's doctrine of absolute predestination, and yet becomes eloquent when he speaks of the literary merits of his 'Institutes.' '\emph{Sein Lehrbuch der christlichen Religion,}' says Kampschulte (Vol. I. p. xiv.), '\emph{bringt die kirchliche Revolution in ein System, das durch logische Schärfe, Klarheit des Gedankens, rücksichtslose Consequenz, die vor nichts zurückbebt, noch heute unser Staunen und unsere Bewunderung erregt.}' Ibid. p. 274: '\emph{Calvin's Lehrbuch der christlichen Religion ist ohne Frage das hervorragendste und bedeutendste Erzeugniss, welches die reformatorische Literatur des sechszehnten Jahrhunderts auf dem Gebiete der Dogmatik aufzuweisen hat. Schon ein oberflächlicher Vergleich lässt uns den gewaltigen Fortschritt erkennen, den es gegenüber den bisherigen Leistungen auf diesem Gebiete bezeichnet. Statt der unvollkommenen, nach der einen oder andern Seite unzulänglichen Versuche Melanchthon's, Zwingli's, Farel's erhalten wir aus Calvin's Hand das Kunstwerk eines, wenn auch nicht harmonisch in sich abgeschlossenen, so doch wohlgegliederten, durchgebildeten Systems, das in allen seinen Theilen die leitenden Grundgedanken widerspiegelt und von vollständiger Beherrschung des Stoffes zeugt. Es hatte eine unverkennbare Berechtigung, wenn man den Verfasser der Institution als den Aristoteles der Reformation bezeichnete. Die ausserordentliche Belesenheit in der biblischen und patristischen Literatur, wie sie schon in den früheren Ausgaben des Werkes hervortritt, setzt in Erstaunen. Die Methode ist lichtvoll und klar, der Gedankengang streng logisch, überall durchsichtig, die Eintheilung und Ordnung des Stoffes dem leitenden Grundgedanken entsprechend; die Darstellung schreitet ernst und gemessen vor und nimmt, obschon in den späteren Ausgaben mehr gelehrt als anziehend, mehr auf den Verstand als auf das Gemüth berechnet, doch zuweilen einen höheren Schwung an. Calvin's Institution enthält Abschnitte, die dem Schönsten, was von Pascal und Bossuet geschrieben worden ist, an die Seite gestellt werden können: Stellen, wie jene über die Erhabenheit der heiligen Schrift, über das Elend des gefallenen Menschen, über die Bedeutung des Gebetes, werden nie verfehlen, auf den Leser einen tiefen Eindruck zu machen. Auch von den katholischen Gegnern Calvin's sind diese Vorzüge anerkannt und manche Abschnitte seines Werkes sogar benutzt worden. Man begreift es vollkommen, wenn er selbst mit dem Gefühl der Befriedigung und des Stolzes auf sein Werk blickt und in seinen übrigen Schriften gern auf das ``Lehrbuch'' zurückverweist.}'}

The appearance of his \emph{Institutes of the Christian Religion}\footnote{The full title of the first edition is 'Christia- | næ Religionis Insti-  | \emph{tutio totam fere pietatis summam et quic | quid est in doctrina salutis cognitu ne- | cessarium, complectens: omnibus pie- | tatis studiosis lectu dignissi- | mum opus, ac re- | cens edi- | tum. }|  Præfatio ad Chri- | stianissimum Regem Franciæ,  \emph{qua | hic ei liber pro confessione fidei | offertur. }| Joanne Calvino | \emph{Nouiodunensi authore. }| Basileæ, | M.D.XXXVI.' The dedicatory Preface is dated '\emph{X. Calendas Septembres}' (\emph{i.e.} August 23), without the year; but at the close of the book the month of March, 1536, is given as the date of publication. The first two French editions (1541 and 1545) supplement the date of the Preface correctly: '\emph{De Basle le vingt-troysiesme d’Aoust mil cinq cent trente cinq.}' The manuscript, then, was completed in Aug. 1535, but it took nearly a year to print it. The eighth and last improved edition from the pen of the author bears the title: 'Institutio Chri- | stianæ Religionis,  \emph{in libros qua- }| \emph{tuor nunc primum digesta, certisque distincta capitibus, ad aptissimam }| \emph{methodum: aucta etiam tam magna accessione ut propemodum opus }| \emph{novum haberi possit. }| Joanne Calvino authore. | Oliva Roberti Stephani. | \emph{Genevæ. }| M.D.LIX.'} (first in Latin, then in French) marks an epoch in the history of theology, and has all the significance of an event. This book belongs to those few uninspired compositions which never lose their interest and power. It has not only a literary, but an institutional character. Considering the youth of the author, it is a marvel of intellectual precocity. The first edition even contained, in brief outline, all the essential elements of his system; and the subsequent enlargements to five times the original size were not mechanical additions to a building or changes of conviction,\footnote{'\emph{In doctrina,}' says Beza, towards the close of his \emph{Vita Calv.,} '\emph{quam initio tradidit ad extremum constans nihil prorsus immutavit, quod paucis nostra memoria theologis contigit.}' Bretschneider was quite mistaken when he missed in the first edition the doctrine of predestination, which is clearly though briefly indicated, pp. 91 and 138. See Kampschulte, p. 256.} but the natural growth of a living organism from within.\footnote{The Strasburg editors devote the first four volumes to the different editions of the \emph{Institutes} in both languages. Vol. I. contains the \emph{editio princeps Latina,} of Basle, 1536 (pp. 10–247), and the variations of six editions intervening between the first and the last, viz., the Strasburg editions of 1539, 1543, 1545, and the Geneva editions of 1550, 1553, 1554 (pp. 253–1152); Vol. II. the \emph{editio postrema} of 1559 (pp. 1–1118); Vol. III. and IV. the last edition of the French translation, or free reproduction rather (1560), with the variations of former editions. The question of the priority of the Latin or French text is now settled in favor of the former. See Jules Bonnet, in the \emph{Bulletin de la Société de l’histoire du protestantisme français} for 1858, Vol. VI. pp. 137 sqq., Stähelin, Vol. I. p. 55, and the Strasburg editors of the \emph{Opera,} in the ample \emph{Prolegomena} to Vols. I. and III. Calvin himself says expressly (in the Preface to his French ed. 1541) that he first wrote the \emph{Institutes} in Latin ('\emph{premièrement l’ay mis en latin}') for readers of all nations, and that he translated them afterwards for the special benefit of Frenchmen. In a letter to his friend, Francis Daniel, dated Lausanne, Oct. 13, 1536, he writes that he began the French translation soon after the publication of the Latin (\emph{Letters,} ed. Bonnet, Vol. I. p. 21), but it did not appear till 1541, bearing the title '\emph{Institution de la religion Chrestienne } composée en latin,  \emph{par Jean Calvin, et translatée en français par luymesme.}' The erroneous assertion of a French original, so often repeated (by Bayle, Maimbourg, Basnage, and more recently by Henry, Vol. I. p. 104; III. p. 177; Dorner, \emph{Gesch. der protest. Theol.} p. 375; H. B. Smith, 1.c. p. 283; and Guizot, p. 176, who assumes that the first French ed. was published anonymously), arose from confounding the date of the Preface in the French editions (23 Aug. 1535) with the later date of publication (1536). It is quite possible, however, that the dedication to Francis I. was first written in French, and this would most naturally account for the earlier date in the French editions. On the difference of the several editions, comp. also J. Thomas,  \emph{Histoire de l’instit. chrétienne de J. Calv.,} Strasb. 1859, and Köstlin,  \emph{Calvin's Institutio nach Form und Inhalt,} in the \emph{Studien und Kritiken} for 1868.}

The 'Institutes' are by far the clearest and ablest systematic and scientific exposition and vindication of the ideas of the Reformation in their vernal freshness and pentecostal fire. The book is inspired by a heroic faith ready for the stake, and a glowing enthusiasm for the saving truth of the gospel, raised to a new life from beneath the rubbish of human additions. Though freely using reason and the fathers, especially Augustine, it always appeals to the supreme tribunal of the Word of God, to which all human wisdom must bow in reverent obedience. It abounds in Scripture-learning thoroughly digested, and wrought up into a consecutive chain of exposition and argument. It is severely logical, but perfectly free from the dryness and pedantry of a scholastic treatise, and flows on, like a Swiss river, through green meadows and sublime mountain scenery. It overshadowed all previous attempts at a systematic treatment of Protestant doctrines, not only those of Zwingli and Farel, but even Melanchthon's \emph{Loci theologici,} although Calvin generously edited them twice in a French translation with a complimentary preface (1546).\footnote{See the Preface in \emph{Opera,} Vol. IX. pp. 847–850. It is written in excellent taste, and with profound respect and affection for Melanchthon, whose work, he concludes, '\emph{conduit à la pure verité de Dieu, à laquelle it nous convient tenir, nous servant des hommes pour nous aider à y parvenir.}'}

No wonder that the 'Institutes' were greeted with enthusiastic praises by Protestants, which are not exhausted to this day.\footnote{See the eulogies of Bucer, Beza, Sainte-Marthe, Thurius, Blunt, Salmasius, John von Müller, and others, quoted by Henry and Stähelin (Vol. I. pp. 59 sqq.). To these may be added some more recent testimonies. Guizot says (1.c. p. 173): 'The \emph{Institutes} were and are still the noblest monument of the greatness of mind and originality of idea which distinguished Calvin in his own century. More than that, I believe this book to be the most valuable and enduring of all his labors; for those churches which are specially known as the Reformed Churches of France, Switzerland, Holland, Scotland, and the United States of America received from Calvin's \emph{Institutes} the doctrine, organization, and discipline which, in spite of sharp trials, grave mistakes, and claims which are incompatible with the progress of liberty, have still, for more than three centuries, been the source of all their strength and vitality.' Hase (in his \emph{Kirchengeschichte}) calls the \emph{Institutes} '\emph{die grossartigste wissenschaftliche Rechtfertigung des Augustinismus voll religiösen Tiefsinns in unerbittlicher Folgerichtigkeit der Gedanken.}' G. Frank (\emph{Gesch. der Protest. Theol.} Vol. I. p. 74): '\emph{Wie Melanchthon hat auch Calvin seinen Glauben zusammengefasst in einem besonderen Werke, der Inst. rel. chr., nur methodischer, folgerichtiger, überlegner, die grösste Glaubenslehre des} 16 \emph{Jahrh. ist sie wie ein hochgewölbter, dunkler Dom, darin der Ernst der Religion in andächtigem Schauer sich über die Seele legt.}' H. B. Smith (1.c. p. 288): 'It is the most complete system [of theology] which the 16th century produced, nor has it been supplanted by any single work.' Baur (\emph{Dogmengeschichte,} Vol. III. p. 27) calls it 'in every respect a truly classical work, distinguished in a high degree by originality and acuteness of conception, systematic consistency, and clear, luminous method.' To many editions of the \emph{Institutes} the well-known distich of the Hungarian Paul Thurius is affixed: \par '\emph{Præter apostolicas post Christi tempora chartas, Huic peperere libro sæcula nulla parem.}'} They created dismay among Romanists, were burned at Paris by order of the Sorbonne, and hated and feared as the very 'Talmud' and 'Koran of heresy.'\footnote{Florimond de Ræmond, \emph{Histoire de la naissance, progrez et decadence de l’hérésie de ce siècle,} pp. 838, 883, quoted by Kampschulte (p. 278), who adds: '\emph{Seine Schrift des Reformationszeitalters ist von den Katholiken mehr gefürchtet, eifriger bekämpft und verfolgt worden, als Calvin's Christliche Institution.}' See his own judgment quoted on pp. 446 sq., note.} In spite of severe prohibition, they were translated into all the languages of Europe, and passed through innumerable editions. Among the Protestants of France they acquired almost as much authority as Luther's Bible in Germany, and comforted the martyrs in prison. In England, after the accession of Elizabeth, they were long used as the text-book of theology; and even the moderate and 'judicious' Hooker prized them highly, and pronounced Calvin 'incomparably the wisest man that ever the French Church did enjoy.'

This remarkable work was originally a defense of the evangelical doctrines against ignorant or willful misrepresentation, and a plea for toleration in behalf of his scattered fellow-Protestants in France, who were then violently persecuted as a set of revolutionary fanatics and heretics. Hence the dedicatory Preface to Francis I. As the early Apologists addressed the Roman emperors to convince them that the Christians were innocent of the foul charges of atheism, immorality, and hostility to Caesar, so Calvin appealed to the French monarch in defense of his equally innocent countrymen, with a manly dignity, frankness, force, and pathos never surpassed before or since. It is a sad reflection that such a voice of warning should have had so little effect, and that the noble French nation even this day would rather listen to the revolutionary 'Marseillaise' of Voltaire and Rousseau than to the reformatory trumpet of Calvin.

The 'Institutes,' to which this dedication to the French monarch forms the magnificent portal, consist of four books (each divided into a number of chapters), and treat, after the natural and historical order of the Apostles' Creed, first of the knowledge of God the Creator (theology); secondly, of the knowledge of God the Redeemer (christology); thirdly, of the Holy Spirit and the application of the saving work of Christ (soteriology); fourthly, of the external means of salvation, viz., the Church and the Sacraments.\footnote{The first edition of the \emph{Institutes} contains only six chapters: 1. \emph{De lege,} with an explanation of the Decalogue; 2. \emph{De fide,} with an exposition of the Apostles' Creed; 3. \emph{De oratione,} with an exposition of the Lord's Prayer; 4. Of the Sacraments of Baptism and the Lord's Supper; 5. Of the other so-called Sacraments; 6. Of Christian liberty, Church-government, and discipline.}

The most prominent and original features of Calvin's theological system, which have left their impress upon the Reformed Creed, are the doctrine of Predestination and the doctrine of the Lord's Supper. By the first he widened the breach between the Reformed and the Lutheran Church; by the second he furnished a basis for reconciliation.

THE DOCTRINE OF PREDESTINATION.

All the Reformers of the sixteenth century, including even the gentle Melanchthon and the compromising Bucer, under a controlling sense of human depravity and saving grace, in extreme antagonism to Pelagianism and self-righteousness, and, as they sincerely believed, in full harmony not only with the greatest of the fathers, but also with the inspired St. Paul, came to the same doctrine of a double predestination which decides the eternal destiny of all men. Nor is it possible to evade this conclusion on the two acknowledged premises of Protestant orthodoxy—namely, the wholesale condemnation of men in Adam, and the limitation of saving grace to the present world. If the Lutheran theology, after the Formula of Concord (1577), rejected Synergism and Calvinism alike, and yet continued to teach the total depravity of all men and the unconditional election of some, it could only be done at the expense of logical consistency.\footnote{Schleiermacher, the greatest divine of the nineteenth century, has defended Calvinism as the only consistent system on the basis of the orthodox anthropology and eschatology (though he runs it out into a final, unscriptural universalism); and his pupil, Alexander Schweizer, of Zurich (in his \emph{Glaubenslehre der evang. reform. Kirche,} Vol. I. pp. 79 and 81), thus clearly and sharply states the logical aspect of the case: '\emph{Der reformirte Lehrbegriff, consequent gegründet auf das Materialprincip schlechthiniger Abhängigkeit von Gott und von da aus das menschliche Thun beleuchtend, ohne dessen willensmässige Natur zu verkleinern, ist weniger durch seinen Determinismus anstössig geworden, als durch das dualistisch Particularistische der auf die Prädestination angewandten Weltansicht. Gerade dieses aber gehört der Weltansicht aller damaligen Confessionen gleich sehr an and folgt wirklich aus der Vorstellung, dass unser ewiges Loos beim irdischen Sterben entschieden sei, nur hienieden Erlöste selig werden, alle Andern aber verdammt bleiben. . . . Das Harte am reformirten Lehrbegriff ist der dualistische Particularismus, der aber allen Confessionen gemein durch die reformirte Consequenz nur heller in’s Licht gestellt wird, wodurch allein, falls er irrig wäre, die Förderung zur Wahrheit angebahnt ist.} 1. \emph{Dualistischer Particularismus ist die Idee, dass in der Menschen- und Engelwelt die einen selig werden, die andern ewig verdammt. Diess war die Ansicht aller kirchlichen Confessionen, indem der Universalismus, die Beseligung aller rationalen Kreaturen in allen drei Confessionen, als hæretische Irrlehre abgewiesen wurde.} 2. \emph{Liegt im Particularismus Hartes, die Güte Gottes Beschränkendes, so ist es ungerecht, darüber nur die reformirte Confession anzugehen, die weiter nichts gethan, als gelehrt hat: Das Weltergebniss müsse dem Weltplan entsprechen, somit habe Gott ewig grade diese Welt mit diesem Ergebniss gewollt und eine particularistische Prädestination bei sich beschlossen, wovon nun alle Weltentwicklung einfach die Ansführung sei; denn dass alles anders herauskomme, als Gott es gewollt, heisse Gott von den Kreaturen abhängig machen, die Kreaturen zu Göttern machen, Gott aber zum Ungott.}' Comp. also Baur, \emph{Dogmengeschichte,} Vol. III. (1867), pp. 144 sqq.}

Yet there were some characteristic differences among the Reformers. Luther started from the \emph{servum arbitrium,} Zwingli from the idea of an all-ruling \emph{providentia,} Calvin from the timeless or eternal \emph{decretum absolution.} Calvin elaborated the doctrine of predestination with greater care and precision, and avoided 'the paradoxes' of his predecessors. He made it, moreover, the corner-stone of his system, and gave it undue proportion. He set the absolute sovereignty of God over against the mock sovereignty of the Pope. It was for him the 'article of the standing or falling Church;' while Luther always assigned this position to the doctrine of justification by faith alone. In this estimate, both were mistaken, for the central place in the Christian system belongs only to the person and work of Christ—the incarnation and the atonement. Finally, the Augustinian and Lutheran predestinarianism is moderated by the sacramentarian principle of baptismal regeneration; while the Calvinistic predestinarianism confines the sacramental efficacy to the elect, and turns the baptism of the non-elect into an empty form.

Predestination, according to Calvin, is the eternal and unchangeable decree of God by which he foreordained, for his own glory and the display of his attributes of mercy and justice, a part of the human race, without any merit of their own, to eternal salvation, and another part, in just punishment of their sin, to eternal damnation. The decree is, therefore, twofold—a decree of \emph{election} to holiness and salvation, and a decree of \emph{reprobation} on account of sin and guilt.\footnote{\par 'Præscientiam  \emph{quum tribuimus Deo, significamus omnia semper fuisse ac perpetuo mamere sub ejus oculis; ut ejus notitiæ nihil futurum aut præteritum, sed omnia sint præsentia, et sic quidem præsentia, ut non ex ideis tantum imaginetur (qualiter nobis obversantur ea quorum memoriam mens nostra retinet), sed tanquam ante se posita vere intueatur ac cernat. Atque hæc præscientia ad universum mundi ambitum et ad omnes creaturas extenditur. }Prædestinationem  \emph{vocamus æternum Dei decretum, quo apud se constitutum habuit, quid de unoquoque homine fieri vellet. Non enim pari conditione creantur omnes; sed aliis vita æterna, aliis damnatio æterna præordinatur. Itaque, prout in alterutrum finem quisque conditus est, ita vel ad vitam, vel ad mortem prædestinatum dicimus.}' \emph{Instit.} Lib. III. c. 21, § 5 (\emph{Opera,} Vol. II. pp. 682, 683). Comp. his \emph{Articuli de prædest.,} first published from an autograph of Calvin, Vol. IX. p. 713.} The latter is the negative counterpart, which strict logic seems to demand, but against which our better feelings revolt, especially if it is made to include multitudes of innocent children, for their unconscious connection with Adam's fall. Calvin himself felt this, and characteristically called the decree of reprobation a 'decree horrible, though nevertheless true.'\footnote{'\emph{Iterum quæro, unde factum est ut tot gentes una cum liberis eorum infantibus æternæ morti involveret lapsus Adæ absque remedio, nisi quia Deo ita visum est? Hic obmutescere oportet tam dicaces alioqui linguas. Decretum quidem horribile, fateor; infitiari tamen nemo poterit quin præsciverit Deus, quem exitum esset habiturus homo, antequam ipsum conderet, et ideo præsciverit, quia decreto suo sic ordinarat. In præscientiam Dei si quis hic invehatur, temere et inconsulte impingit. Quid enim, quæso, est cur reus agatur cœlestis judex quia non ignoraverit quod futurum erat? In prædestinationem competit, si quid est vel justæ vel speciosæ querimoniæ. Nec absurdum videri debet quod dico, Deum non modo primi hominis casum, et in eo posterorum ruinam prævidisse, sed arbitrio quoque suo dispensasse. Ut enim ad ejus sapientiam pertinet, omnium quæ futura sunt esse præscium, sic ad potentiam, omnia manu sua regere ac moderari.}' \emph{Instit.} Lib. III. c. 23, § 7 (Vol. II. p. 704).} All he could say was that God's will is inscrutable, but always holy and unblamable. It is the ultimate ground of all things, and the highest rule of justice. Foreordination and foreknowledge are inseparable, and the former is not conditioned by the latter, but God foresees what he foreordains. If election were dependent on man's faith and good works, grace would not be free, and in fact would cease to be grace. Man's holiness is not the cause or condition, but the effect of God's election. The unequal distribution of gospel privileges can be traced only to the secret will of God. All men are alike corrupt and lost in Adam; some are saved by free grace, others, who are no worse by nature, reject the gospel. These are undeniable every-day facts, and admit of no other explanation within the limits of the present life; and as to the future world, we know nothing but what God has revealed to us in the Scriptures.

Calvin carried the doctrine of the divine decrees beyond the Augustinian infralapsarianism, which makes the fall of Adam the object of a permissive or passive decree, and teaches the preterition rather than the reprobation of the wicked, to the very verge of supralapsarianism, which traces even the first sin to an efficient or positive decree, analogous to that of election. But while his inexorable logic pointed to this abyss, his moral and religious sense shrunk from the last inference of making God the author of sin, which would be blasphemous, and involve the absurdity that God abhors and justly punishes what he himself decreed. Hence his phrase, which vacillates between infralapsarianism and supralapsarianism: 'Adam fell, God' providence having so ordained it; yet he fell by his own guilt.'\footnote{'\emph{Lapsus est enim primus homo, quia Dominus ita expedire censuerat; cur censuerit, nos latet. Certum tamen est non aliter censuisse, nisi quia videbat, nominis sui gloriam inde merito illustrari. Unde mentionem gloriæ Dei audis, illic justitiam cogita. Justum enim esse oportet quod laudem meretur. Cadit igitur homo, Dei providentia sic ordinante, sed suo vitio cadit. . . . Propria ergo malitia, quam acceptrat a Domino puram naturam corrupit; sua ruina totam posteritatem in exitium secum attraxit.}' \emph{Instit.} Lib. III. c. 23, § 8 (Vol. II. p. 705). The difference between the supralapsarians and infralapsarians was not agitated at the time of Calvin, but afterwards during the Arminian controversy in Holland. Both schools appealed to him. The difference is more speculative than moral and practical. In creating man free, God created him necessarily temptable and liable to fall, but the fall itself is man's own act and abuse of freedom. God decreed sin not efficiently but permissively, not as an actual fact but as a mere possibility, not for its own sake but for the sake of the good or as a negative condition of redemption. Besides, sin has no positive character, is no created substance, bat it is privative and negative, and consists simply in the abuse of faculties and gifts essentially good.}

Calvin defended this doctrine against all objections with consummate skill, and may be said to have exhausted the subject on his side of the question. His arguments were chiefly drawn from the Scriptures, especially the ninth chapter of the Epistle to the Romans; but he unduly stretched passages which refer to the historical destiny of individuals and nations in this world, into declarations of their eternal fate in the other world; and he escaped the proper force of opposite passages (such as John i. 29; iii. 16; 1 John ii. 2; iv. 14; 1 Tim. ii. 4; 2 Pet. iii. 9) by a distinction between the secret and revealed or declared will of God (\emph{voluntas arcani} and \emph{voluntas beneplaciti}), which carries an intolerable dualism into the divine will.

The motive and aim of this doctrine was not speculative, but practical. It served as a bulwark of free grace, an antidote to Pelagianism and human pride, a stimulus to humility and gratitude, a source of comfort and peace in trial and despondency. The charge of favoring license and carnal security was always indignantly repelled by the Pauline 'God forbid!' It is moreover refuted by history, which connects the strictest Calvinism with the strictest morality.

The doctrine of predestination, in its milder, infralapsarian form, was incorporated into the Geneva Consensus, the Second Helvetic, the French, Belgic, and Scotch Confessions, the Lambeth Articles, the Irish Articles, the Canons of Dort, and the Westminster Standards; while the Thirty-nine Articles,\footnote{There is a dispute about the precise meaning of Art. XVII.; but, as Prof. Fisher says (\emph{The Reform.} p. 335), 'the article can not fairly be interpreted in any other sense than that of unconditional election; and the cautions which are appended, instead of being opposed to this interpretation, demonstrate the correctness of it.''} the Heidelberg Catechism, and other German Reformed Confessions, indorse merely the positive part of the free election of believers, and are wisely silent concerning the decree of reprobation, leaving it to theological science and private opinion.

Supralapsarianism, which makes unfallen man, or man before his creation (\emph{i.e.,} a \emph{non ens,} a mere abstraction of thought), the object of God's double foreordination for the manifestation of his mercy in the elect, and his justice in the reprobate, was ably advocated by Beza in Geneva, Gomarus in Holland, Twisse (the Prolocutor of the Westminster Assembly) in England, Nathaniel Emmons (1745–1840) in New England, but it never received symbolical authority, and was virtually or expressly excluded (though not exactly condemned) by the Synod of Dort, the Westminster Assembly, and even the 'Formula Consensus Helvetica' (1675).\footnote{Can. IV.: '\emph{Ita Deus gloriam suam illustrare constituit, ut decreverit, primo quidem hominem integrum creare. } tum  \emph{ejusdem lapsum } permittere, \emph{ac demum ex lapsis quorundam misereri, adeoque eosdem eligere, alios vero in corrupta massa } Relinquere, \emph{æternoque tandem exitio devovere.}' This does not go beyond the limits of Augustinianism. Van Oosterzee errs when he says (\emph{Christian Dogmatics,} Vol. I. p. 452) that the Form. Cons. Hel. asserts the supralapsarian view; while Hodge errs on the other side when he says (\emph{Syst. Theol.} Vol. II. p. 317) that this document contains 'a formal repudiation of the supralapsarian view.'} All Calvinistic Confessions, without exception, trace the fall to \emph{a permissive} decree, make man responsible and justly punishable for sin, and reject, as a blasphemous slander, the charge that God is the author of sin. And this is the case with all the Calvinistic divines of the present day.\footnote{Dr. Hodge, who best represents the Old School Calvinism in America, rejects supralapsarianism and defends infralapsarianism, which he defines thus (\emph{Syst. Theol.} Vol. II. pp. 319 and 320): 'According to the infralapsarian doctrine, God, with the design to reveal his own glory—that is, the perfections of his own nature—determined to create the world; secondly, to permit the fall of man; thirdly, to elect from the mass of fallen men a multitude whom no man could number as ``vessels of mercy;'' fourthly, to send his Son for their redemption; and, fifthly, to leave the residue of mankind, as he left the fallen angels, to suffer the just punishment of their sins.'}

CALVIN'S DOCTRINE OF THE LORD'S SUPPER.

Calvin's doctrine of the Lord's Supper, on which he spent much deep and earnest thought, is an ingenious compromise between the realism and mysticism of the Lutheran, and the idealism and spiritualism of the Zwinglian theory. It aims to satisfy both the heart and the reason.

He retained the figurative interpretation of the words of institution, and rejected all carnal and materialistic conceptions of the eucharistic mystery; but he very strongly asserted, at the same time, a spiritual real presence and fruition of Christ's body and blood for the nourishment of the soul. He taught that believers, while they receive with their mouths the visible elements, receive also by faith the spiritual realities signified and sealed thereby, namely, the benefit of the atoning sacrifice on the cross, and the life-giving virtue of Christ's glorified humanity in heaven, which the Holy Ghost conveys to the soul in a supernatural manner; while unbelieving or unworthy communicants, having no inward connection with Christ, receive only bread and wine to their own judgment. He thus sought to avoid alike the positive error of Luther and the negative error of Zwingli (whose view of the Eucharist he even characterized as 'profane'), and to unite the elements of truth advocated by both in a one-sided and antagonistic way. Luther and Zwingli always had in mind a corporeal or dimensional presence of the material substance of body and blood, and an oral manducation of the same by all communicants—which the one affirmed, the other denied; Calvin substituted for this the idea of a virtual or dynamic presence of the psychic life-power and efficacy of Christ's humanity, and a spiritual reception and assimilation of the same by the organ of faith, and therefore on the part of believing communicants only, through the secret mediation of the Holy Spirit.\footnote{Calvin taught his view of the Eucharist in the first edition of his \emph{Institutes} (cap. 4, \emph{De Sacramentis, }pp. 236 sqq., in the new ed. of the \emph{Opera,} Vol. I. pp. 118 sqq.; comp. Ebrard, \emph{Das Dogma v. heil. Abendmahl,} Vol. II. p. 412), and in the \emph{Confessio fidei de eucharistia.} (1537); then more fully in the later editions of the \emph{Institutes,} 1.c. Lib. IV. cap. 17, 18; in his two Catechisms (1538 and 1542); in his admirable tract \emph{De Cœna Domini} (first in French, 1541, then in Latin, 1545; see Opera, Vol. V. pp. 429–460); in the \emph{Consensus Tigurinus} (1549); and he defended it in several polemical treatises against Westphal (1555–1557) and Heshusius (1561).}

Calvin's doctrine of the Eucharist was substantially approved by Melanchthon in his later period, although from fear of Luther and the ultra-Lutherans he never fully committed himself. It passed into all the leading Reformed Confessions of the sixteenth and seventeenth centuries, and must be regarded as the orthodox Reformed doctrine. Zwingli's theory, which is more simple and intelligible, has considerable popular currency, but no symbolical authority.\footnote{See, on this whole subject, the very elaborate exposition of Ebrard, \emph{Das Dogma v. heil. Abendmahl,} Vol. II. pp. 402–525; Baur, \emph{Geschichte der christl. Kirche,} Vol. IV. pp. 398–402; and Nevin's article on the \emph{Reformed Doctrine of the Lord's Supper,} in the \emph{Mercersburg Review} for Sept. 1850, pp. 421–548 (in defense of his 'Mystical Presence'). Dr. Nevin has clearly and correctly stated Calvin's doctrine of the Eucharist and abundantly fortified it with quotations from all the symbolical standards, in entire harmony with Ebrard (who indorsed him in the \emph{Studien und Kritiken}). After rejecting both the dogma of transubstantiation and consubstantiation, he says (p. 429): 'In opposition to this view, the Reformed Church taught that the participation of Christ's flesh and blood in the Lord's Supper is \emph{spiritual} only, and in no sense corporal. The idea of a local presence in the case was utterly rejected. The elements can not be said to comprehend or include the body of the Saviour in any sense. It is not \emph{there,} but remains constantly in heaven, according to the Scriptures. It is not handled by the minister and taken into the mouth of the communicant. The manducation of it is not oral, but only by faith. It is present in fruition accordingly to believers only in the exercise of faith; the impenitent and unbelieving receive only the naked symbols, bread and wine, without any spiritual advantage to their own souls. Thus we have the doctrine defined and circumscribed on both sides; with proper distinction from all that may be considered a tendency to Rationalism in one direction, and from all that may be counted a tendency to Romanism in the other. It allows the \emph{presence} of Christ's person in the sacrament, including even his flesh and blood, so far as the actual participation of the believer is concerned. Even the term \emph{real presence} Calvin tells us he was willing to employ, if it were to be understood as synonymous with \emph{true} presence; by which he means a presence that brings Christ truly into communion with the believer in his human nature as well as in his divine nature. The word \emph{real,} however, was understood ordinarily to denote a local, corporal presence, and on this account was not approved. To guard against this, it may be qualified by the word \emph{spiritual;} and the expression will then be quite suitable to the nature of the doctrine as it has been now explained. A \emph{real} presence, in opposition to the notion that Christ's flesh and blood are not made present to the communicant in \emph{any} way. A \emph{spiritual} real presence, in opposition to the idea that Christ's body is in the elements in a local or corporal manner. Not real simply, and not spiritual simply, but real and yet spiritual at the same time. The body of Christ is in heaven, the believer on earth; but by the power of the Holy Ghost, nevertheless, the obstacle of such vast local distance is fully overcome, so that in the sacramental act, while the outward symbols are received in an outward way, the very body and blood of Christ are at the same time inwardly and supernaturally communicated to the worthy receiver, for the real nourishment of his new life. Not that the material particles of Christ's body are supposed to be carried over, by this supernatural process, into the believer's person. The communion is spiritual, not material. It is a participation of the Saviour's life; of his life, however, as human, subsisting in a true bodily form. The living energy, the vivific virtue, as Calvin styles it, of Christ's flesh, is made to flow over into the communicant, making him more and more one with Christ himself, and thus more and more an heir of the same immortality that is brought to light in his person.'}

Calvin thus combined his high predestinarianism with a high view of the Church and the Sacraments. Augustine and Luther did the same to a still greater extent, with more prominence given to the sacramental idea. It is the prerogative of great minds to maintain apparently opposite truths and principles which hold each other in check; while with minds less strong and comprehensive, the one principle is apt to rule out the other. In the Catholic and Lutheran Churches the sacramental principle gradually overruled the doctrine of absolute predestination; in the more rigid Calvinistic school, the sacramental principle yielded to the doctrine of predestination. But the authoritative standards are committed to both.

CALVIN AS AN EXEGETE.

Among the works which have more or less influenced the Reformed Confessions we can not ignore Calvin's commentaries. To expound the Scriptures in books, from the chair, and from the pulpit, was his favorite occupation. His whole theology is scriptural rather than scholastic, and distinguished for the skillful and comprehensive working up of the teaching of the Bible, as the only pure fountain of revealed truth and the infallible rule of the Christian faith. As it is systematically comprehended in his 'Institutes,' and defended in his various polemical tracts against Sadolet, Pighius, the Council of Trent, Caroli, Bolsec, Castallio, Westphal, Heshusius, so it is scattered through his Commentaries on the Gospels, Acts, and Epistles, and the principal books of the Old Testament, especially the Psalms and the Prophets. He opened this important series of works, during his sojourn at Strasburg, by an exposition of the Epistle to the Romans (1539), on which his theological system is chiefly based.

He could assert with truth on his death-bed that he never knowingly twisted or misinterpreted a single passage of the Scriptures, that he always aimed at simplicity, and restrained the temptation to show acuteness and ingenuity. He regarded it as the chief object of a commentator to adhere closely to the text, and to bring out clearly and briefly the meaning of the writer. He detested irrelevant talk and diffuseness, and avoided allegorical fancies, which substitute pious imposition for honest exposition. He combined in a very rare degree all the necessary hermeneutical qualifications, a fair knowledge of Greek and Hebrew, sound grammatical tact, thorough sympathy with the spirit and aim of the Bible, and aptitude for fruitful practical application. He could easily enter into the peculiar situation of the Prophets and Apostles, as though he had been with them in their trials, and shared their varied experience. He is free from pedantry, and his exposition is an easy, continuous flow of reproduction. He never evades difficulties, but frankly meets and tries to solve them.

With all his profound reverence for the Word of God, to which his reason bows in cheerful obedience, he is not swayed by a peculiar theory of inspiration or dogmatic prejudice, but shows often remarkable freedom and sagacity in discovering the direct historical import of prophecies, in distinction from their ulterior Messianic bearing.\footnote{In his exposition of Gen. iii. 15, he understands the '\emph{woman's seed}' collectively of the human family in its perpetual struggle with Satan, which at last culminates in the victory of Christ, the head of the race. Comp. also his remarks on Isa. iv. 2; vi. 3; Psa. xxxiii. 6; Matt. ii. 15;  Heb. ii. 6–8.} He notices the difference of style and argument in the Second Epistle of Peter as compared with the first, and in the Epistle to the Hebrews as compared with the undisputed Pauline Epistles. He never ventured to explain the mysteries of the Apocalypse. Luther, with an equally profound reverence and enthusiasm for the Word of God, was even much bolder, and passed sweeping judgments on whole books of the canon (as the Epistle of James, the book of Esther, and the book of Revelation), because he could not find enough of Christ in them. Calvin and his followers retained the Canon in full, but excluded more rigidly the Apocrypha of the Old Testament.

The scholastic Calvinism and Lutheranism of the seventeenth century departed from the more liberal view of the Reformers on the mode and degree of inspiration, and substituted for it a rigid mechanical theory which ignored the human and historical aspect of the Scriptures, and reduced the sacred writers to mere penmen of the Holy Ghost. This theory found symbolical expression in the 'Formula Consensus Helvetica' (1675), which advocates even the inspiration of the Hebrew-vowel points, and cuts off all textual criticism.

Upon the whole, Calvin is ' beyond all question the greatest exegete of the sixteenth century,'\footnote{Reuss: \emph{Geschichte der H. Schriften des N. T.,} 4th edition, p. 564.} which of all centuries was the most fruitful in this department of sacred learning. Luther was the prince of translators; Calvin, the prince of commentators. Augustine and Luther had occasionally a deeper intuition into the meaning of difficult passages, and seized on the main idea with the instinct of genius; but Calvin was more accurate and precise, and more uniformly excellent. Modern commentators have made great progress in textual criticism and grammatical and historical exegesis, but do not attain to his religious depth and fervor. His commentaries have stood the test of time, and will always be consulted with profit. Scaliger, who was displeased with all men, said that no scholar had penetrated so deeply into the meaning of the Prophets as Calvin; the Roman Catholic critic Richard Simon admitted that his commentaries would be 'useful to the whole world,' if they were free from declamations against popery; and of all older expounders none is more frequently quoted by the best modern critical scholars than John Calvin.\footnote{See the frequent references to him in the Commentaries of Tholuck, Hengstenberg, Lücke, Bleek, De Wette, Meyer, Alford; also the Essay of Tholuck, '\emph{Die Verdienste Calvin's als Ausleger der heil. Schrift,}' 1831 (reprinted in his \emph{Vermischte Schriften,} Vol. II. pp. 330–360); Ed. Reuss, \emph{Calvin considéré comme exégète} (\emph{Revue,} Vol. VI. p. 223); and Stähelin, \emph{Joh. Calvin,} Vol. I. pp. 182 sqq. Stähelin says (p. 198): '\emph{Der altlestamentliche wie der neutestamentliche Bibelerklärer, der Lutheraner, wie der Unirte und Reformirte, der wissenschaftliche Exeget, wie der populäre Schriftausleger alle schöpften und schöpfen immer noch aus der Arbeit Calvins bei weitem das Meiste und Beste, was sie von Schrifterklärung aus dem Reformationszeitalter beibringen.}' Comp. also Kahnis, \emph{Dogmatik,} Vol. II. p. 492, and Herzog, \emph{Encykl.} Vol. II. p. 528.}

CALVIN'S CHURCH POLITY.

The practical and ecclesiastical part of Calvin's work is in some respects even more important than his theology, and must be briefly considered in those features which have affected the Calvinistic Confessions. These are the duty of discipline, the principle of lay-representation, and the autonomy of the Church in its relation to the State. In these points Calvinism differs from Lutheranism, and also from Zwinglianism and Anglicanism. Calvin aimed at a moral and social as well as a doctrinal and religious reformation, and succeeded in establishing a model Church, which excited the admiration not only of sympathizing contemporaries, like Farel and Knox,\footnote{John Knox, the Reformer of Scotland, who studied at the feet of Calvin, though four years his senior, in a letter to his friend Locke, in 1556, called the Church of Geneva 'the most perfect school of Christ that ever was in the earth since the days of the Apostles. In other places I confess Christ to be truly preached; but manners and religion to be so sincerely reformed, I have not yet seen in any other place besides.' Farel wrote, in 1557, that he never saw Geneva in such excellent condition before, and that he wonld rather be the last there than the first any where else. There, it was said, the pure gospel is preached in all temples and houses (Calvin himself preached daily, every other week); there the music of psalms never ceases; there hands are folded and hearts lifted up to heaven from morning till night and from night until morning. The Italian refugee, Bernardino Ochino, gives a most favorable description of the moral condition of Geneva. See his Life by Beurath (1875), p. 169.} but even of visitors of other creeds long after his death.\footnote{Dr. Valentine Andreæ of Würtemberg (a grandson of Jacob Andreæ, the chief author of the Formula of Concord), a great and shining light of the Lutheran Church in Germany during the desolations of the Thirty-Years' War (d. 1654), visited Geneva in the early part of the seventeenth century, and held it up as a model of moral purity well worthy of imitation. '\emph{Als ich in Genf war,}' he says in his \emph{Respublica Christianopolitana,} 1619, '\emph{bemerkte ich etwas Grosses, woran die Erinnerung, ja vielmehr, wonach die Sehnsucht nur mit meinem Leben absterben wird. Nicht nur findet sich hier das vollkommene Institut einer vollkommenen Republik, sondern als eine besondere Zierde und Mittel der Disciplin eine Sittenzucht, nach welcher über die Sitten und selbst die geringsten Ueberschreitungen der Bürger wöchentlich Untersuchung angestellt wird, zuerst durch die Viertelsinspectoren, dann durch die Senioren, endlich durch den Magistrat, je nachdem der Frevel der Sache oder die Verhärtung und Verstockung der Schuldigen es erfordern. In Folge dessen sind denn alle Fluchworte, alles Würfel- und Kartenspiel, Ueppigkeit, Uebermuth, Zank, Hass, Betrug, Luxus, u.s.w., geschweige denn grössere Vergehungen, die fast unerhört sind, untersagt. Welche herrliche Zierde für die christliche Religion solche Sittenreinheit, vor der wir mil allen Thränen beweinen müssen, dass sie uns fehlt und fast ganz venachlässigt wird, und alle Gutgesinnten sich anstrengen, dass sie in’s Leben gerufen werde! Mich, wofern mich die Verschiedenheit der Religion nicht abgehalten, hätte die sittliche Uebereinstimmung hier auf ewig gefesselt, und mit allem Eifer habe ich von da an getrachtet, dass etwas Aehnliches auch unserer Kirche zu Theil würde. Nicht geringer als die öffentliche Zucht war auch die häusliche meines Hausherrn Scarron ausgezeichnet durch stetige Gebetsübungen, Lectüre der heiligen Schrift, Gottesfurcht in Worten und Thaten, Masshalten in Speise und Kleidung, dass ich eine grössere Sittenreinheit selbst im väterlichen Hause nicht gesehen.}'} During the eighteenth century his severe system of theology and discipline gave way to the prevailing spirit of Socinianism and the revolutionary spirit of Jean Jacques Rousseau—the counterpart of Calvin; but revived in the nineteenth century, though in a modified form, so that Geneva has become a second time the centre of evangelistic labors in the French-speaking world.\footnote{The Haldanes repaid the debt of Scotland to Geneva, and, in connection with Cesar Malan, gave the first impulse to a revival which resulted in the establishment of a Free Church, and a school of theology distinguished by the labors of Gaussen, Merle d'Aubigné, Pronier, La Harpe. The old National Church which Calvin founded has likewise undergone a salutary change, though the old rigor can never be restored. In point of literary culture and social refinement, Geneva always retained the first rank among French cities next to Paris.}

1. Discipline.—Calvin's zeal for discipline, especially for the honor of the Lord's table, in excluding unworthy communicants, was the cause of his expulsion from Geneva, the cause of his recall from Strasburg, the condition of his acceptance, the struggle and triumph of his life. He had a long and fierce conflict with the ferocious politico-religious party of the Libertines, or 'Spirituals,' as they called themselves, who combined a pantheistic creed with licentiousness and free-lovism, and anticipated the worst forms of modern infidelity to the extent of declaring the gospel a tissue of lies of less value than Æsop's Fables.\footnote{See Calvin's \emph{Instructio adv. fanaticam et furiosam sectam Libertinorum, qui se Spirituales vocant,} written first in French, 1544, \emph{Opera,} Vol. VII. pp. 145–252. Comp. Trechsel's art. \emph{Libertiner} in Herzog's \emph{Real-Encykl.,} and Stähelin, Vol. I. pp. 383 sqq.} He regarded them as worse enemies of God and the truth than the Pope. They resorted to personal indignities and every device of intimidation; they named the very dogs of the street after him; they one night fired fifty shots before his bedchamber; they threatened him in the pulpit; they approached the communion table as if to seize the sacred elements, when he cried out, 'You may break these limbs and shed my blood, I would rather die than dishonor the table of my God,' whereupon they left the church. On another occasion he walked into the midst of an excited mob and offered his breast to their daggers. It seems incredible that a man constitutionally 'unwarlike and timorous' should have completely overcome at last such a powerful and determined opposition, which reached its height in 1553.

The system of discipline which he established saved Geneva from anarchy, into which the Libertines would have plunged it, and was a training-school of self-government for other Reformed Churches; but it was carried to unwarrantable excesses in the punishment of religious and civil offenses, and even innocent amusements, and entered too much into details of private and domestic life.

2. Presbyterian and Synodical Church Polity.—It rests on the principle of ministerial equality, and the principle of lay-representation by elders or seniors in the government of the Church. This polity, founded by Calvin, was consistently carried out in the Presbyterian Churches of France, Holland, Scotland, England, and the United States; but in German Switzerland and Germany it succeeded only partially, while the Church of England retained the Episcopal hierarchy. Calvin himself, however, was not an exclusive Presbyterian. He allowed modifications of the form of government in different countries. He did not object to Episcopacy or the liturgical worship in England; he only protested against the ecclesiastical supremacy of Henry VIII. and a number of abuses.

3. The Autonomy of the Church.—The German Reformers, including Zwingli, yielded too much authority to the civil rulers in matters of religion. Calvin theoretically made the Church independent in her own sphere, and claimed for her the right of self-government. This leads consistently to a separation of Church and State, where the latter is hostile to the former, as was the case in France and to some extent in Scotland. In recent times the Calvinistic Churches, without changing their creed, tend naturally towards complete freedom, from State control. Yet in practice he had no idea of such a separation. He regarded the civil and the spiritual power as the two arms of God's government in the world, which should co-operate together for the same end—the glory of God and the good of society: the Church by infusing a religious spirit into the State, the State by protecting and promoting the interests of the Church. He established, after the model of the Old Testament, a theocracy at Geneva, and governed it by tacit consent as long as he lived, presiding over the 'Venerable Company' of Pastors, and exerting a molding influence upon the civil legislation of the little republic of about 20,000 inhabitants.\footnote{Kampschulte, Vol. I. p. 471: '\emph{Der Grundgedanke, von dem der Gesetzgeber Genfs ausgeht, ist die Theokratie. Calvin will in Genf den Gottesstaat herstellen. Nur Einer ist ihm König und Herr in Staat und Kirche: Gott im Himmel. In seinem Namen herrscht jede irdische Gewalt. Gottes Herrscherruhm zu verkündigen, seine Majestät zu verherrlichen, seinen heiligen Willen zur Ausführung zu bringen und seine Bekenner zu heiligen, ist die gemeinsame Aufgabe von Staat und Kirche.}' Comp. Stähelin, Vol. I. pp. 319 sqq.}

Bossuet, Möhler, and other Roman Catholic divines saw in this a return to the hierarchy, with Calvin as its pope. He has sometimes been compared to Hildebrand; and Kampschulte remarks that the dominion of the spiritual sovereignty was more thoroughly carried out in Geneva than by the Gregories and Innocenses in the Middle Ages. But Calvin's theocracy differed essentially from the Roman Catholic by its popular (though by no means democratic) basis: it was not priestcraft ruling over statecraft, but a self-governing Christian commonwealth. Geneva was an aristocratic republic, ruled by the clergy and the people in orderly representation and friendly co-operation. The highest civil and executive power was lodged in the 'Little Council' of twenty-four syndics, the highest ecclesiastical power in the 'Consistory,' composed (at first) of six pastors and double that number of lay-elders.\footnote{Guizot says of this ecclesiastical organization (p. 265): 'In its origin it was a profoundly Christian and evangelical system; it was republican in many of its fundamental principles and practices, and at the same time it recognized the necessity of authority and order, and originated general and permanent rules of discipline.' Michelet calls the Geneva of Calvin 'the city of the spirit, founded by Stoicism on the rock of predestination;' and Kampschulte (p. 430), 'the metropolis of a grand, sublime, and terrible idea.'}

RELIGIOUS PERSECUTION AND RELIGIOUS FREEDOM.

Unfortunately Calvin inherited from the Theodosian Code and the Catholic Church the worst feature of the theocratic system, namely, the principle of appeal to the secular arm for the temporal, and, if necessary, capital punishment of spiritual offenses, as being offenses against the order and peace of society. This principle is inconsistent with liberty of conscience (which Beza called a diabolical dogma), and justifies all manner of persecution, as duty or policy may suggest. With his intense antagonism to the papal tyranny, he might have thrown off this relic of the Middle Ages, if it had not been for his conviction of the perpetual validity of the Mosaic civil code and his theocratic theory. He thought that the burning of innocent people by Romanists was no good reason why Protestants should spare the guilty.

It was the misfortune of Calvin that this false theory, which confounds two distinct spheres and ignores the spiritual nature of Christ's kingdom, was brought to its severest test and explosion under his own eye, and to the perpetual injury of his fair fame. ``We mean, of course, the terrible theological tragedy of the Spanish physician Michael Servetus, a restless fanatic, a pantheistic pseudo-reformer, and the most audacious and even blasphemous heretic of the sixteenth century, who attacked the doctrine of the Holy Trinity as tritheistic and atheistic, as the greatest monstrosity, and the source of all corruption in the Church. After being condemned to death, and burned in effigy by the Roman Catholic authorities in France,\footnote{See the acts of the process of Servetus at Vienne and Lyons (first published by the Abbé d'Artigny, Paris, 1749), in Calvin's \emph{Opera,} Vol. VIII. pp. 833 sqq.} he fled to Geneva, was arrested, tried, and executed at the stake, for heresy and blasphemy, by the civil government, with the full consent of Calvin, except that he made an ineffectual plea for a mitigation of the punishment (by a substitution of the sword for the fagot).\footnote{For full discussion of the trial and execution of Servetus and Calvin's part in them, see Schaff: \emph{Ch. Hist.,} VII., 680–793. A large boulder has been placed by Swiss and French Protestants at the spot where the Spaniard suffered, to serve as an expiatory monument, expressing regret for the tragedy and at the same time respect for Calvin. On the one side is the record that Michael Servetus of Villeneuve, Aragon, b., Sept. 29, 1511, was executed there, Oct. 26, 1553, and on the other the inscription: \emph{Fils respectueux et reconnaissants de Calvin, notre grand Réformateur, mais condamnant une erreur qui fut celle de son siècle et fermement attachés à la liberté de conscience selon les vrais principes de la Réformation et de l’Evangile nous avons élevé ce monument expiatoire, le 27 Octobre, 1903.} (We, respectful children of Calvin, our great Reformer, but condemning an error which was the error of his age and firmly attached to liberty of conscience according to the true principles of the Reformation and the Gospel, have erected this expiatory monument, Oct. 27, 1903). Calvin was called by Renan ``the most Christian man of his age.''—Ed.}

Severely as we must condemn the great Reformer, from the standpoint of our modern civilization, for this the saddest mistake of his life, it is evident that even here he acted consistently and conscientiously, and that the blame attaches not to his personal character (for towards sincere and earnest heretics, like Lælius Socinus, he showed marked courtesy and leniency), but to his system, and not to his system alone, but to the inherited system of his age, which had not yet emerged from the traditions of the Romish pseudo-theocracy. The burning of Servetus was fully approved by all the Reformers—Beza, Farel, Bucer, Bullinger, even the mild and gentle Melanchthon.\footnote{It may he questioned whether Zwingli and Luther, had they lived, would have sanctioned the execution; their impulses at least were more liberal. With all his polemic violence in argument, Luther disapproved of the shocking cruelties against the Anabaptists in Germany, and said that 'on this plan, the hangman would be the best theologian.'} If Romanists condemned Calvin, they did it from hatred of the man, and condemned him for following their own example even in this particular case. The public opinion of Christendom at that time and down to the eighteenth century justified the right and duty of civil government not only to protect but to support orthodoxy, and to punish heresy by imprisonment, exile, and death; and this right was exercised, with more or less severity, in all countries of Europe, and even in Puritan New England during the colonial period. Protestants differed from Romanists only in their definition of heresy, and by greater moderation in its punishment. Protestants complained of being \emph{innocently} persecuted in France, Spain, Holland, and under the bloody Mary in England; and Catholics raised the same complaint against the systematic cruelty of the penal code of Queen Elizabeth, which looked to the utter extermination of Romanism and Puritanism alike.

A protest against the principle of persecution, first raised by Justin Martyr and Tertullian in the early Church, but forgotten as soon as the Church ascended the throne of the Cæsars, was revived by heretical Anabaptists and Socinians, who themselves suffered from it, without having a chance to persecute their persecutors, and who thus became martyrs of religious freedom. All honor to them, even to Servetus, for the service they rendered under this view to future generations. Liberty is the sweet fruit of bitter persecution. During the seventeenth century this feeble and isolated protest was considerably strengthened by Arminians, Baptists, and Quakers for the same reason; and during the eighteenth century Christian liberality and philanthropy on the one hand, and religious indifferentism and infidelity on the other, made such progress that the doctrinal foundations of persecution were gradually undermined, and toleration (as it was first patronizingly and condescendingly called, and is still called in despotic countries) became the professed policy of civilized governments. But this is not enough: all Christian governments should legally recognize and protect liberty of conscience, as an inherent and inalienable right of every immortal soul; and this requires for its full realization a peaceful separation of Church and State, or an equality of all denominations before the law.

In view of this radical revolution of public opinion on the subject of persecution, it becomes a practical question whether those sections of the Protestant confessions of faith which treat of the relation of Church and State should not be reconstructed and adapted to the principle of religious freedom, all the more since the Papal Syllabus has consistently condemned it, as being one of the errors of modern times. Such a change, at all events, is necessary in the United States, and has actually been made in the American revision of the Thirty-Nine Articles, and of the Westminster Confession.

The principle of religious liberty does not necessarily, as was formerly supposed, imply indifference to truth or a weakening of intensity of conviction. It follows legitimately from a sharper discrimination between the secular and spiritual sphere, between the Old and the New Testaments, between the law of Moses and the gospel of Christ, and from the spirit and example of Him who said, 'My kingdom is not of this world,' and who commanded the carnal-minded Peter to 'put up his sword into the sheath.' God alone is Lord of the conscience, and allows no one with impunity to interfere with his sovereign right. Religion flourishes best in the atmosphere of freedom, and need not fear error as long as truth is left free to combat it.

It is nevertheless true that Calvinism, by developing the power of self-government and a manly spirit of independence which fears no man, though seated on a throne, because it fears God, the only sovereign, has been one of the chief agencies in bringing about this progress, and that civil and religious liberty triumphed first and most completely in Calvinistic countries. 'Calvin,' says Guizot, 'is undoubtedly one of those who did most towards the establishment of religious liberty.'

\xsubsection{The Catechism of Geneva. A.D. 1536 and 1541.}

§ 58. The Catechism of Geneva. A.D. 1536 AND 1541.

Literature.

Calvini \emph{ Opera, ed. Baum, Cunitz, and Reuss,} Vol. V. (1866), pp. 313–362 (the first draft, or \emph{Catechismus prior,} 1538); Vol. VI. (1867), pp. 1–160 (the second catechism, in French and Latin).

Niemeyer, pp. 123–190 (the Latin text of the Larger Catechism, together with the prayers and liturgical forms); comp. his Proleg. pp. xxxvii.-xli.

The German text of the Larger Catechism in Beck (Vol. I. pp. 208–292), and Böckel (pp. 127–172).

An English translation, probably by the same Marian exiles who prepared the 'Geneva Bible,' appeared first at Geneva, 1556; then in Edinburgh, 1564; and is reprinted in Dunlop's \emph{Confessions,} Vol. II. pp. 139–272; also in Horatius Bonar:  \emph{Catechisms of the Scotch Reformation} (Lond. 1866), pp. 4–88. It is divided into fifty-five Sundays.

Stähelin:  \emph{Joh. Calvin,} Vol. I. pp. 124 sqq.

The commanding influence of Calvin's theology and Church polity is manifest in all the leading confessions of the Reformed Churches, especially the French, Dutch, and Scotch, also in the Lambeth Articles, the Irish Articles, and the Westminster Standards. But the confessions which he himself prepared were intended, like those of Zwingli, for local and temporary rather than general purposes, and possess only a secondary authority. These are the Geneva Catechism, the Zurich Consensus, and the Geneva Consensus.\footnote{They were not included in the \emph{Corpus et Syntagma Confessionum,} which appeared in Geneva.}

Calvin, like Luther and other Reformers, did not consider it beneath his dignity, but rather a duty and a privilege, to utilize his profound learning for the benefit of children by adapting it to their simplicity. He made general education and catechetical instruction the basis of the republic.\footnote{George Bancroft calls Calvin 'the father of popular education, the inventor of the system of free schools.'—\emph{Liter. and Histor. Miscellanies,} p. 406.}

During his first residence at Geneva (1536), he prepared a catechism, in the French language, together with a form of discipline, as a basis of instruction for the newly reformed Church of that city.\footnote{The Latin translation has been recently republished by the Strasburg editors from a Basle edition: '\emph{Catechismus, sive Christianæ Religionis institutio, communibus renatæ nuper in Evangelio Genevensis Ecclesiæ suffragiis recepta et vulgari quidem prius idiomate, nunc vero Latine etiam . . . in lucem edita. Joanne Calvino autore. Basileæ, A. M.D. XXXVIII.}' See the Prolegomena to \emph{Opera,} Vol. V. pp. xli. sqq. The French original, which was probably printed at Geneva, 1537, seems to have been lost.} It is a brief summary of the Christian religion, a popular extract from his 'Institutes.' It treats, in fifty-eight sections (but not in the form of question and answer), of the religious constitution of man, the distinction between false and true religion, the knowledge of God, the original state of man, free-will, sin and death, the way of salvation, the law of God, the Ten Commandments, the sum of the law (Matt. xxii. 37), the aim of the law, faith in Christ, election and predestination, the nature of faith, justification and sanctification, repentance and regeneration, faith and good works, an exposition of the articles of the Apostles' Creed, and the petitions of the Lord's Prayer, the sacraments of Baptism and the Lord's Supper, the Church, human traditions, excommunication, and the civil magistrate. Then follows a short confession of faith, in twenty-one articles, extracted from the Catechism, which was to be binding upon all the citizens of Geneva—probably the first instance of a formal pledge to a symbolical book in the history of the Reformed Church.\footnote{\par '\emph{Confessio Fidei, in quam jurare cives onmes Genevenses et qui sub civitatis ejus ditione agunt, jussi sunt: excerpta e Catechismo quo utitur Ecclesia Genevensis.}' It begins with the Word of God and ends with the magistrate. It seems to have been drawn up before the Catechism, immediately after the disputation at Lausanne, for Beza says: '\emph{Tunc edita est a Calvino Christianæ doctrinæ quædam veluti formula, vixdum emergenti e papatus sordibus Genevensi Ecclesiæ accommodata. Addidit etiam Catechismum,}' etc.}

After his return from Strasburg Calvin rewrote the Catechism on a larger scale, and arranged in questions and answers: the catechist drawing out the information, and the pupil or child seeming to teach the master. It was prepared in great haste, for the printer demanded copy without giving him time to revise it. He often desired to perfect the book, but found no time.\footnote{So he said himself on his death-bed; see Stähelin, Vol. II. p. 467.} It appeared in French, 1541 or 1542,\footnote{'\emph{Le Catechisme de l’Église de Genève, c’est à dire le Formulaire d’instruire les enfans en la Chrestienté fait en manière de dialogue ou le ministre interrogue et l’enfant respond.}' The oldest copy extant was found in the ducal library at Gotha, printed 1545. On other editions, see the Prolegomena to \emph{Opera,} Vol. VI.} in Latin, 1545,\footnote{'Catechismus Ecclesiæ Genevensis,  \emph{hoc est, Formula erudiendi pueros in doctrina Christi. Autore Joanné Calvino.}' The Preface to the Latin edition is dated '\emph{Genevæ,} 4 \emph{Calendas Decembris,} 1545.' The Strasburg editors give the French and Latin texts of 1545 in parallel columns, Vol. VI. pp. 8–159. In many editions Calvin's Liturgy is added.} and very often. It was also translated into Italian (1551 and 1556), Spanish (1550), English (1556), German, Dutch, Hungarian, even into Greek and Hebrew.\footnote{Beza, in \emph{Vita,} ad ann. 1541: '\emph{Scripsit Catechismum Gallice et Latine, ab illo priore minime discrepantem, sed multo auctiorem, et in quæstiones ac responsiones distributum: quem merito nobis liceat admirandum quoddam opus vocare, tantopere plurimis etiam exteris populis probatum, ut non modo vernaculis plurimis linguis, utpote Germanica, Anglica, Scotica, Belgica, Hispanica, sed etiam Hebraice ab Immanuele Tremellio Judæo Christiano, et Græce ab Henrico Stephano legatur elegantissime conversus.}' The title of the Greek translation is, \textgreek{Στοιχείωσις τῆς Χριστιανῶν πίστεως, ἠ Κατηχισμὸς, κατὰ τὴν παλαιὰν ὀνομασίαν. } \emph{Græce et Latine,} 1563.} It was used for a long time in Reformed Churches and schools, especially in France and Scotland, and served a good purpose in promoting an intelligent piety and virtue on the solid basis of systematic Biblical instruction. Educational religion, which grows with our growth, is the most substantial, and must ever be the main reliance of the Church.

The object of this work, as explained in the preface, was to restore the catechetical instruction of the ancient Church, so sadly neglected by the Papists, who substituted for it the ceremony of confirmation, and to secure greater unity of faith and doctrine in the scattered Reformed congregations. Calvin showed his churchly tact in making the Apostles' Creed, the Ten Commandments, and the Lord's Prayer the basis. The leading idea is man's relation to God, and his heavenly destination. The whole is divided into five parts, as follows: 1. Of Faith—an exposition of the Creed (which here, as in the Heidelberg Catechism, precedes the Ten Commandments, while in the earlier Catechism of Calvin the opposite order was observed);\footnote{He made the Apostles' Creed the basis of his 'Catechism' and 'Institutes,' not because he believed it to be literally the product of the Apostles, but because it is a faithful summary of their teaching ('\emph{ex eorum scriptis fideliter collecta,}' '\emph{tiré de la pure doctrine apostolique}'), and a formula which best expresses the common Christian faith ('\emph{formula confessionis, quam inter se communem habent Christiani omnes}').} 2. Of the Law, or the Ten Commandments; 3. Of Prayer; 4. Of the Word of God; 5. Of the Sacraments. In the French edition the Catechism is divided into fifty-five lessons, for the fifty-two Sundays of the year and the three great festivals—a method followed in the later editions of the Heidelberg Catechism.\footnote{The distribution into Sundays appears first in the French edition of 1548, which has a '\emph{Table pour trouver le lieu du Catechisme que le Ministre explique un chascun Dimanche.}' See \emph{Opera,} Vol. VI. \emph{Proleg.} p. x. The First Book of Discipline of Scotland (1560), ch. 11, directs the ministers to teach the children Calvin's Catechism—'the most perfect that ever yet was used in the Kirk'—every Sunday afternoon in the presence of the people. See Bonar, l.c. pp. 3, 4.}

Calvin's Catechism is fuller than Luther's, but less popular and childlike. It prepared the way and furnished material for a number of similar works, which have gradually superseded it, especially the Anglican (Nowell's), the Heidelberg, and the Westminster Catechisms. The Anglican Catechism is much shorter and more churchly in taking its starting-point from Baptism. The first question of the Westminster Catechism makes the glory of God 'the chief \emph{end} of man,' and is a happy condensation of the first three questions of Calvin.\footnote{Calvin's Catechism. Westminster Shorter Catechism.   \par \emph{Min. Quis humanæ vitæ præcipuus est finis?} \par 1\emph{st Ques.} What is the chief end of man?   \par \emph{Puer. }Ut Deum, a quo conditi sunt homines, ipsi noverint.  \par \emph{Ans. Man's chief end is to glorify God, and to enjoy him forever.}   \par \emph{Min. Quid causæ habes, cur hoc dicas?}     \par \emph{Puer.} Quoniam nos ideo creavit et collocavit in hoc mundo, quo glorificetur in nobis. Et sane vitam nostram, cujus ipse est initium, æquum est in ejus gloriam referri.     \par \emph{Min. Quod vero est summum bonum hominis?}     \par \emph{Puer.} Illud ipsum.} The Heidelberg Catechism begins more subjectively with 'the only \emph{comfort} of man in life and in death,' herein betraying its German origin; but this also was suggested by the next questions of Calvin concerning the highest good or felicity of man and the firm foundation of our salvation. Otherwise the Heidelberg Catechism adheres to the order of the Genevan more closely than the Westminster, by retaining, as a basis of the dogmatic section, the Apostles' Creed (which the Westminster Catechism merely adds as an appendix).\footnote{Comp. Karl Sudhof: \emph{Olevianus und Ursinus} (1857), pp. 88 sqq. Calvin is also responsible for the unhistorical interpretation of Christ's descent into Hades, by which he understood the anticipation of the sufferings of hell in Gethsemane and on the Cross. This is quite inconsistent with the position of this article between the burial and the resurrection. Ihe Westminster Catechism falls into another error by making it mean simply, 'He continued in the state of the dead and under the power of death till the third day.'}

Guizot gives the preference to Calvin's Catechism over those modern ones which begin with speculative questions on the nature and existence of God. 'Calvin,' he says,\footnote{\emph{St. Louis and Calvin,} p. 348.} 'proceeds in a very different manner; he does not seek God—he knows him, possesses him, and takes God as his starting-point. God the Creator, man his creature, and the relation of man to God—these form the fundamental facts and natural basis of the history, doctrines, and laws of Christianity. Calvin's Catechism commences thus: ``What is the chief end of human life?'' ``To know God.'' And this first assertion is the mainspring of all the principles and religious duties which are afterwards presented, not as the discoveries of the human mind, but as communications made by God in order to meet man's aspirations, and enable him to regulate his life. It is neither a scientific method, nor is the Catechism a philosophical work; it contains the assertion of a real, immemorial, universal, and historical fact, and explains the consequences of that fact. It is the natural and legitimate method of imparting religions instruction, inherent in the very first principle of all religion; it is especially in harmony with the origin and history of Christianity, and no one has ever recognized its power or proved its efficacy more fully than Calvin.'

\xsubsection{The Consensus of Zurich. A.D. 1549.}

§ 59. The Consensus of Zurich. A.D. 1549.

Literature.

I. Consensio Mutua in re Sacramentaria  \emph{ministrorum }Tigurinæ Ecclesiæ  \emph{et J. Calvini ministri } Genevensis Ecclesiæ  \emph{jam nunc ab ipsis autoribus edita.} Tiguri, 1549. In \emph{Opera,} Vol. VII. pp. 689–748. Comp. Proleg. pp. xliv. sqq. Defensio  \emph{sanæ et orthodoxæ de sacramentis eorumque vi, fine, et usu, et fructu quam pastores et ministri Tigurinæ ecclesiæ et Genevensis antehac brevi Consensionis mutuæ formula complexi sunt. Johanne Calvino autore,} Tiguri, 1555, in \emph{Opera,} Vol. IX. pp. 1–40. The same volume contains the later eucharistic tracts of Calvin against the attacks of Joachim Westphal (1556 and 1557) and Tilemann Heshusius (1561).

The \emph{Consensus Tigurinus} with Calvin's \emph{Exposition} is also reprinted in Niemeyer's \emph{Collect.} pp. 191–217; a German translation (in part) in Beck and Böckel.

II. On the History of the Zurich Consensus, see Calvin's correspondence with Bullinger, 1548 and 1549, \emph{Opera,} Vols. XII. and XIII. Hundeshagen:  \emph{Conflicte des Zwinglianismus,} etc.; Henry:  \emph{Calvin,} Vol. II. pp.128 sqq.; Ebrard:  \emph{Das Dogma vom heil. Abendmahl,} Vol. II. pp. 484–524; Pestalozzi:  \emph{Bullinger,} pp. 373–387; Stähelin:  \emph{Calvin,} Vol. II. pp. 112–124.

In the sacramental controversy—the most violent, distracting, and unprofitable in the history of the Reformation—Calvin stood midway between Luther and Zwingli, and endeavored to unite the elements of truth on both sides, in his theory of a spiritual real presence and fruition of Christ by faith.\footnote{See § 57, pp. 455 sqq.} This satisfied neither the rigid Lutherans nor the rigid Zwinglians. The former could see no material difference between Calvin and Zwingli, since both denied the literal interpretation of 'this \emph{is} my body,' and a corporeal presence and manducation. The latter suspected Calvin of leaning towards Lutheran consubstantiation and working into the hands of Bucer, who had made himself obnoxious by his facile compromises and ill-concealed concessions to the Lutheran view in the Wittenberg Concordia (1536).

The wound was reopened by Luther's fierce attack on the Zwinglians (1545), and their sharp reply. Calvin was displeased with both parties, and counselled moderation. It was very desirable to harmonize the teaching of the Swiss Churches. Bullinger, who first advanced beyond the original Zwinglian ground, and appreciated the deeper theology of Calvin, sent him his book on the Sacraments, in manuscript (1546), with the request to express his opinion. Calvin, did this with great frankness, and a degree of censure which at first irritated Bullinger. Then followed a correspondence and personal conference at Zurich, which resulted in a complete union of the Calvinistic and Zwinglian sections of the Swiss Churches on this vexed subject. The negotiations reflect great credit on both parties, and reveal an admirable spirit of frankness, moderation, forbearance, and patience, which triumphed over all personal sensibilities and irritations.\footnote{See the details in Ebrard, Pestalozzi, and Stähelin, who speak in the highest terms of the truly Christian spirit which characterized the two leaders of the Swiss Reformation.}

The first draft of the Consensus Tigurinus, from November, 1548, consists of twenty-four brief propositions drawn up by Calvin, with annotations by Bullinger, to which Calvin responded in January, 1549. They assert that the Sacraments are not in and of themselves effective and conferring grace, but that God, through the Holy Spirit, acts through them as means; that the internal effect appears only in the elect; that the good of the Sacraments consists in leading us to Christ, and being instruments of the grace of God, which is sincerely offered to all; that in baptism we receive the remission of sins, although this proceeds primarily not from baptism, but from the blood of Christ; that in the Lord's Supper we eat and drink the body and blood of Christ, not, however, by means of a carnal presence of Christ's human nature, which is in heaven, but by the power of the Holy Spirit and the devout elevation of our soul to heaven.\footnote{Opera, Vol. VII. pp. 693 sqq.}

In the month of March Calvin sent twenty Articles to the Synod of Berne,\footnote{Ibid. pp. 717 sqq.} but in this canton there was strong opposition to Calvin's rigorism, which subsided only after his death.\footnote{See Hundeshagen, and Stähelin, Vol. II. pp. 125 sqq. Calvin complained on his deathbed of the ill-treatment he had repeatedly received from the government of Berne.}

In May, 1549, he had, in company with Farel, a personal interview with Bullinger in Zurich at his cordial invitation, and drew up the Consensus as it now stands, in Twenty-six Articles. It was published in 1551 at Zurich and at Geneva.\footnote{\emph{Opera,} Vol. VII. pp. 733 sqq. These Twenty-six Articles alone are given, with Calvin's Exposition of 1554, in Niemeyer's \emph{Collectio,} pp. 191–217.} It contains the Calvinistic doctrine, adjusted as nearly as possible to the Zwinglian in its advanced form, but with a disturbing predestinarian restriction of the sacramental grace to the elect.\footnote{Art. XVI. '\emph{Præterea sedulo docemus, Deum non promiscue vim suam exserere in omnibus qui sacramenta recipiunt: sed tantum in electis. Nam quemadmodum non alios in fidem illuminat, quam quos præordinavit ad vitam, ita arcana Spiritus sui virtute efficit, ut percipiant electi quod offerunt sacramenta.}' Yet this is qualified in Art. XVIII. '\emph{Certum quidem est, offeri communiter omnibus Christum cum suis donis, nec hominum infidelitate labefactari Dei veritatem, quin semper vim suam retineant sacramenta: sed non omnes Christi et donorum ejus sunt capaces. Itaque ex Dei parte nihil mutatur: quantum vero ad homines spectat, quisque pro fidei suæ mensura accipit.}' See the lengthy discussion of Ebrard, 1.c. pp. 503 sqq. He fully adopts the doctrine of the Consensus with the exception of the predestinarian restriction, which, however, is inseparable from the Calvinistic system, as formerly held by Ebrard himself.} The truth of the Zwinglian view is fully acknowledged in opposition to transubstantiation and consubstantiation, but the real life union with Christ in the sacrament is as clearly asserted, and made still more plain in the 'Exposition' of the Consensus which Calvin wrote four years afterwards (1554). 'The Sacraments,' he declares, 'are helps and media (\emph{adminicula et media}), by which we are either inserted into the body of Christ, or being so inserted coalesce with it more and more, till he unites us with himself in full in the heavenly life. . . . The Sacraments are neither empty figures, nor outward badges merely of piety, but seals of the promises of God, attestations of spiritual grace for cherishing and confirming faith, organs also by which God efficaciously works in his elect.'\footnote{'\emph{Sacramenta neque inanes esse figuras neque externa tantum pietatis insignia, sed promissionum Dei sigilla, testimonia spiritualis gratiæ ad fidem fovendam et confirmandam, item organa esse quibus efficaciter agit Deus in suis electis, ideoque, licet a rebus signatis distincta sint signa, non tamen disjungi ac separari,}' etc. Niemeyer, p. 204.}

The Consensus was adopted by the Churches of Zurich, Geneva, St. Gall, Schaffhausen, the Grisons, Neuchatel, and, after some hesitation, by Basle, and was favorably received in France, England, and parts of Germany. Melanchthon declared to Lavater (Bullinger's son-in-law) that he then for the first time understood the Swiss, and would never again write against them; but he erased those passages of the Consensus which made the efficacy of the sacrament depend on election.

While the Consensus brought peace and harmony to the Swiss Churches, it was violently assailed by Joachim Westphal, of Hamburg (1552), in the interest of the ultra-Lutheran party in Germany, and became the innocent occasion of the second sacramental war, which has been noticed in the section on the Formula Concordiæ.\footnote{See pp. 279 sqq. A full account of the controversy of Calvin with Westphal is given by Ebrard, Vol. II. pp. 525 sqq., and by Nevin in the \emph{Mercersburg Review} for 1850, pp. 486 sqq.}

\xsubsection{The Consensus of Geneva. A.D. 1552.}

§ 60. The Consensus of Geneva. A.D. 1552.

Literature.

I. De Æterna Dei prædestinatione  \emph{qua in salutem alios ex hominibus elegit, alios suo exitio reliquit: item de providentia qua res humanas gubernat, } Consensus \emph{ pastorum}  Genevensis \emph{ Ecclesiæ a Jo. Calvino, expositus.} Genevæ, 1552. Reprinted in the \emph{Opera,} Vol. VIII. (187O), pp. 249–366. Also in Niemeyer. pp. 218–310. The German text in Böckel (\emph{Die Genfer Uebereinkunft}), pp. 182–280.

II. Alex. Schweizer:  \emph{Die Protest. Centraldogmen der Reform. Kirche,} Vol. I. (1854), pp.l80–238; Henry Vol. II. p. 285; Vol. III. pp. 40 sqq.; Stähelin, Vol. II. (1863), pp. 271–308, and Vol. I. pp. 411 sqq.

Calvin's doctrine of predestination\footnote{See § 57, pp. 450 sqq.} met with strong opposition, which drew from him some able defenses.

The first assault came from an eminent Roman Catholic divine, Albertus Pighius, 1542, who taught the freedom of will almost to the extent of Pelagianism, and conditioned predestination by foreknowledge.\footnote{Pighius of Campen (d. at Utrecht, Dec. 26, 1542) wrote against Luther and Calvin \emph{De libero hominis arbitrio et divina gratia,} Colon. 1542, dedicated to Cardinal Sadolet. This book was first greatly lauded by the Romanists, but after the Council of Trent had fixed its more cautious doctrine of free-will and condemned semi-Pelagianism, it was put by the Spanish Inquisition on the Index of forbidden books.} Calvin wrote a reply to the first part (1543), and dedicated it to Melanchthon, who in the second article of the Augsburg Confession had expressed the Augustinian doctrine of total depravity.\footnote{\emph{Defensio sanæ et orthodoxæ doctrinæ de servitute et liberatione humani arbitrii adv. calumnias A. Pighii Campensis,} Genevæ, 1543. \emph{Opera,} Vol. VI. pp. 225–404.}

A more troublesome opponent was Jerome Bolsec, formerly a Carmelite monk from Paris, then a fugitive Protestant and physician at Geneva and Lausanne, a restless and turbulent spirit. He denounced Calvin's doctrine of predestination as godless and blasphemous, and tried to break down his influence, but was publicly refuted and admonished, and at last expelled from Geneva (1551) and from Berne (1555). He returned to France and to the Roman Church (1563), and thirteen years after Calvin's death he took cruel revenge by a shameless and malignant libel (1577 and 1588), long since refuted.\footnote{On Bolsec, see Bayle, \emph{Dict.}; Henry, \emph{Calv.} Vol. III. pp. 48 sqq.; Trechsel, \emph{Antitrinitarier,} Vol. I. pp. 185 sqq.; Baum, \emph{Beza,} Vol. I. pp. 160 sqq.; and especially Schweizer, l.c. pp. 205–238. It is a sad fact that the blind zeal of modern Romanism has repeatedly republished the libel of Bolsec, with its wicked and absurd charges of theft, adultery, unnatural crimes, blasphemy, insanity, and invocations of the devil. See Audin's biography of Calvin, which has gone through six editions in French (also translated into German and English), and several popular polemic tracts, published by the Society of St. Francis of Sales, of which Stähelin gives some specimens, Vol. I. p. 414.}

These attacks were the occasion of the \emph{Consensus Genevensis,} which first appeared at Geneva, 1552, in the name of the pastors of that city. Calvin contemptuously alludes in the preface to Bolsec, but without naming him, and directs his attack mainly against Pighius (whose doctrine of predestination he had not noticed in the previous work), and a certain Georgius of Sicily (whom he calls an ignorant monk, more deserving of contempt than persecution). The Consensus is, in fact, the second part of his controversial treatise against Pighius (the first being devoted to free-will). It is an elaborate theological argument for the doctrine of absolute predestination, as the only solid ground of comfort to the believer, but is disfigured by polemical violence, and hence unsuited for a public confession. It received the signatures of the pastors of Geneva on account of the disturbances created by Bolsec, but was not intended to be binding for future generations. Beyond Geneva it acquired no symbolical authority. The attempt to enlist the civil government in favor of this dogma created dissatisfaction and opposition in Berne, Basle, and Zurich. Several of Calvin's old friends withdrew; Bullinger counseled peace and moderation; Fabri, of Neuchatel, declared the decree of reprobation untenable; Melanchthon, who in the mean time had changed his view on free-will and predestination, wrote to Peucer that Geneva attempted to restore Stoic fatalism, and imprisoned men for not agreeing with Zeno.\footnote{Bullinger prepared, March, 1553, for an English friend (Barthol. Traheron), a tract, whose title indicates his partial dissent from Calvin: '\emph{De providentia Dei ejusque prædestinatione, et quod Deus non sit auctor peccati, . . . in quo quæ in Calvini formulis loquendi circa hæc improbet, candide et copiose satis exponit,} 3 \emph{Mart.} 1553.' (Appended by mistake to Peter Martyr's \emph{Loci communes,} Gen. 1626. See the extracts of Schweizer from a MS. copy in Zurich, \emph{Centraldogmen,} Vol. I. pp. 266 sqq.). Bullinger disapproved of the supralapsarian assertion, '\emph{Deum non modo ruinam} (\emph{lapsum}) \emph{prævidisse sed etiam arbitrio suo dispensasse.}' Nevertheless, he called Peter Martyr, who was a strict predestinarian, to Zurich, took sides with Zanchi in the Strasburg controversy, and expressed the infralapsarian view in the Second Helvetic Confession, Art. X. See J. H. Hottinger, \emph{Histor. eccles.} Vol. VIII. p. 723; Schweizer, pp. 237 and 255 sqq.}

The dissatisfaction was increased and the matter complicated by the trial and execution of Servet which soon followed (1553), and by the controversy with Castellio, which involved likewise the doctrine of predestination, together with the question of inspiration and the canon. Sebastian Castellio\footnote{Also written \emph{Castallio} (by Calvin); in French, \emph{Chateillon} and \emph{Chatillon,} probably from his birth-place in Savoy.} (1515–1563), a convert from Romanism, a classical philologist of unusual ability and learning, an advocate of toleration, and a forerunner of Socinianism and Rationalism, was received by Calvin into his house at Strasburg (1510), and called by him to the head of the college at Geneva (Sept., 1541), but was refused admission to the clergy on account of his 'profane view' of the Canticles, which he regarded as a sensual love-song.\footnote{'\emph{Carmen lascivum et obscænum, quo Salomo impudicos suos amores descripserit.}' Castellio doubted the verbal inspiration, and called the Greek of the New Testament impure.} These and other theological differences caused his resignation or dismissal from the school, though with an honorable letter of recommendation from Calvin (Feb. 17, 1545). He removed with his family to Basle, and spent there the remainder of his life—for eight years in great poverty, supporting himself by literary and manual labor, then as professor of Greek in the University (since 1553). His principal work is a Latin translation of the Bible (1551), which was much praised and censured for its pedantic Ciceronian elegance. He attacked Calvin and the Church of Geneva very bitterly in anonymous and pseudonymous books, to which Calvin and Beza replied with equal bitterness. In his 'Dialogue on Predestination,' he charges Calvin with making God the author of sin, and dividing the will of God into two contradictory wills. His own view is that all men are alike created in God's image and for salvation, and are by nature the sons and heirs of God; but that final salvation depends upon faith and perseverance. God loves even his enemies, else he could not command us to love them, and would be worse than the wild beast, which loves its own offspring. God's foreknowledge involves no necessity of human actions: things happen, not because God foreknew them, but God foreknew them because they were to happen. God wills a thing because it is right, and not \emph{vice versa.} He reasons as if there were an established moral order outside and independent of God. He compares God to a musician who unites two tunes because they harmonize. Christ came as a physician to heal all the sick, and if some remain sick it is because they refuse the medicine. The famous passage about Jacob and Esau (Rom. ix.) does not refer to these individuals (for Jacob never served Esau), but to the nations which proceeded from them; and 'to hate' means only 'to love less;' moreover, Esau was not foreordained to sell his birthright, but he did this by his own guilt. Paul himself says that God will have all men to be saved, and that 'he concluded all in unbelief, that he might have mercy upon all.' Castellio died a few months before Calvin, without leaving a school behind him; but his ideas were afterwards more fully developed by the Socinians and Arminians.\footnote{On Castellio, see Schweizer, \emph{Centraldogmen,} Vol. I. pp. 310–373, and his essay, \emph{S. Castellio als Bestreiter der calvinischen Prädestinationslehre,} in the \emph{Theol. Jahrbücher} of Baur and Zeller, 1851.}

Notwithstanding these difficulties, the doctrine of predestination made headway in the Reformed Church. It was strongly advocated in Zurich by Peter Martyr. His opponent, Theodor Bibliander (Buchmann), a distinguished Orientalist, 'the father of exegetical theology in Switzerland,' and a forerunner of Arminianism, was removed from his professorship of Hebrew on account of his advocacy of free-will (1560), though his salary was continued to his death (1564).\footnote{See Schweizer, pp. 276 sqq.} The dogma of predestination consolidated the Calvinistic creed, as the dogma of consubstantiation consolidated the Lutheran creed. Both these distinctive dogmas maintained their hold on the two Churches until the theological revolution towards the close of the eighteenth century began to undermine the whole fabric of Protestant orthodoxy and to clear the way for new creations.

\xsubsection{The Helvetic Consensus Formula. A.D. 1675.}

§ 61. The Helvetic Consensus Formula. A.D. 1675.

Literature.

I. Formula Consensus Ecclesiarum Helveticarum Reformatarum,  \emph{circa Doctrinam de Gratia universali et connexa, aliaque nonnulla capita} (\emph{Einhellige Formul der reform. eidg. Kirchen, betreffend die Lehre von der allgemeinen Gnad und was derselben anhanget, sodann auch etliche andere Religionspunkten}). Composed A.D. 1675; first printed at Zurich, 1714, as an appendix to the Second Helvetic Confession; then 1718, 1722, etc., in Latin and German. The official copy, in both languages, is in the archives of Zurich. The Latin text has a place in Niemeyer's \emph{Collectio,} pp. 729–739; the German text in Böckel, pp. 348–360.

The writings of Amyraut,  Cappel,  and La Place; their friends, Paul Testard,  Jean Daillé, and David Blondel; their opponents, Pierre du Moulin,  Fr. Spanheim, and André Rivet; and the decisions of the Synods of Alençon,  Charenton,  and Loudon (1637–1669). See below.

II. J. Jac. Hottinger (d. 173S): \emph{Succincta et solida ac genuina Formulæ Consensus . . . historia, }Latin and German, 1723. By the same: \emph{ Helvetische Kirchengeschichte, }Zurich, Theil III. pp.1086 sqq.; IV. pp. 258, 268 sqq.

Bayle:  \emph{Dict.} art. \emph{Amyraut.}

Ch. M. Pfaff:  \emph{Dissertatio histor. theologica de Formula Consensus Helv. }Tübingen, 1723.

J. Rud. Salchli: \emph{Stricturæ et observationes in Pfaffi dissertationem de F. C.} Bern, 1723.

(Barnaud:) \emph{ Mémoires pour servir à l’histoire des troubles arrivées en Suisse à l’occasion du Consensus.} Amsterd. 1726.

Walch:  \emph{Religionsstreitigkeiten ausserhalb der luth. Kirche,} Jena, 1733, Vol. I. pp. 454 sqq.; III. pp. 736 sqq.

Hagenbach:  \emph{Kritische Gesch. der ersten Basler Confession.} Basle, 1827, pp. 173 sqq.

Alex. Schweizer:  \emph{Die Protest. Centraldogmen in ihrer Entwicklung innerhalb der Reformirten Kirche. Zweite Häfte} (Zurich, 1856), pp. 439–563. By the same: \emph{Die Enstehung der helvetischen Consensus-Formel,} \emph{aus Zürich’s Specialgeschichte näher beleuchtet,} in Niedner's \emph{Zeitschrift für histor. Theologie} for 1860, pp. 122–148 (gives an extract from the MS. of J. H. Heidegger's \emph{Gründliche und wahrhaftige Historie}). Comp. also Schweizer's art. \emph{Amyraut,} in Herzog's \emph{Real-Encykl.} 2d ed. Vol. I. pp. 356–361; and on the \emph{Life and Writings of Amyraut,} in the \emph{Tübinger Theol Jahrbücher} for 1852.

F. Trechsel:  \emph{Helvetische Consensus-Formel,} in Herzog's \emph{Real-Encyklop.} 2d ed. Vol. V. pp. 755–764 (partly based on MS. sources).

Gust. Frank:  \emph{Geschichte der Protestant. Theologie,} Leipz. 1865, Vol. II. pp. 35 sqq.

Aug. Ebrard:  \emph{Kirchen- und Dogmengeschichte,} Vol. III. (1866), pp. 538 sqq. and 552 sqq. Also his art. on \emph{Amyraldism} (against Schweizer), in the \emph{Reform. Kirchenzeitung} for 1853, No. 27 sqq.

The Helvetic Consensus Formula (\emph{Formula Consensus Helvetica}) is the last doctrinal Confession of the Reformed Church of Switzerland, and closes the period of Calvinistic creeds. It has been called a 'symbolical after-birth.' It was composed in 1675, one hundred and eleven years after Calvin's death, by Professor John Henry Heidegger, of Zurich (1633–1698),\footnote{Author of \emph{Concilii Tridentini Anatome historico-theologica; Enchiridion Biblicum; Historia sacra patriarcharum;} and \emph{Histoire du Papisme.}} at the request and with the co-operation of the Rev. Lucas Gernler, of Basle (d. 1675), and Professor Francis Turretin, of Geneva (1623–1687).\footnote{Author of the \emph{Institutio theologicæ elenchthicæ} (1679–85), which still keeps its place among the best systems of Calvinistic theology. New edition, Edinburgh and New York, 1847, in four volumes. His son, John Alphonsus (1671–1737), Professor of Church History in Geneva, was inclined to Arminianism, and advocated toleration. See Schweizer, \emph{Centraldogmen,} Vol. II. pp. 784 sqq.} It never extended its authority beyond Switzerland, but it is nevertheless a document of considerable importance and interest in the history of Protestant theology. It is a defense of the scholastic Calvinism of the Synod of Dort against the theology of Saumur (\emph{Salmurium}), especially against the universalism of Amyraldus. Hence it may be called a \emph{formula anti-Salmuriensis,} or \emph{anti-Amyraldensis.}

THE SYNOD OF DORT AND THE THEOLOGY OF SAUMUR.

The Twenty-third National Synod of the Reformed Church in France, held at Alais, Oct. 1, 1620, adopted the Canons of Dort (1619), as being in full harmony with the Word of God and the French Confession of 1559, and bound all ministers and elders by a solemn oath to defend them to the last breath. The Twenty-fourth National Synod at Charenton, September, 1623, reaffirmed this adoption.\footnote{Aymon \emph{Tous les Synodes nationaux des églises réformées de France.} A 1a Haye, 1710, Vol. II. pp. 183, 298; Schweizer, 1.c. pp. 229 sqq.}

But in the theological academy at Saumur, founded by the celebrated Reformed statesman Du Plessis Mornay (1604), there arose a more liberal school, headed by three contemporary professors—Josué de la Pace (Placeus, 1596–1655), Louis Cappel (Capellus, 1585–1658), and Moyse Amyraut (Moses Amyraldus, 1596–1664)—which, without sympathizing with Arminianism, departed from the rigid orthodoxy then prevailing in the Lutheran and Reformed Churches on three points—the verbal inspiration of the Scriptures, the particular predestination, and the imputation of Adam's sin.

Saumur acquired under these leaders great celebrity, and attracted many students from Switzerland. It became for the Reformed Church of France what Helmstädt, under the lead of Calixtus, was for the Lutheran Church in Germany; and the Helvetic Consensus Formula of Heidegger may be compared to the '\emph{Consensus repetitus}' of Calovius (1664), which was intended to be a still more rigorous symbolical protest against Syncretism, although it failed to receive any public recognition.\footnote{See p. 851, and Schweizer's comparison of the two documents, Vol. II. pp. 532 sqq.}

The further development of the Saumur theology was arrested by the political oppression which culminated in the cruel revocation of the Edict of Nantes by Louis XIV. (1685), and aimed at the utter annihilation of the Reformed Church in France. But its ideas have silently made progress, and were independently revived in more recent times.

VERBAL INSPIRATION.

Louis Cappel, the most distinguished of an eminent Huguenot family, and one of the first Biblical scholars of the seventeenth century, made the history of the text of the Hebrew Scriptures his special study, and arrived at conclusions which differed from the orthodox theory of a literal inspiration. He discovered and proved that the Hebrew system of vocalization did not date from Adam, nor from Moses, nor from Ezra and the Great Synagogue, but from the Jewish grammarians after the completion of the Babylonian Talmud.\footnote{'\emph{Arcanum punctationis revelatum,}' added to his \emph{Commentarii et notæ criticæ in Vetus Testamentum,} Amst. 1689. Cappel wrote this tract in 1622, and sent the MS. to the elder Buxtorf. of Basle (d. 1629), who returned it with the advice to keep back his view. It was first published anonymously by Erpenius at Leyden, 1624. Twenty years afterwards Buxtorf the younger (d. 1664) attacked it in his \emph{Tractatus de punctorum origine, antiquitate et autoritate,} Basil. 1648. Against this Cappel wrote his \emph{Vindiciæ Arcani punctat. revel.,} but they were not published till 1689, by his son, Jacques C., in an Appendix to his Commentary. His views on the late origin of the Hebrew vowels were anticipated by rabbinical scholars, Abn-Ezra (d. 1174) and Elias Levita (d. 1549).} This view is confirmed by the absence of vowels on Jewish coins, on the Phœnician and Punic monuments, on the inscription of the Moabite stone (discovered 1868), and by the analogy of the other Semitic languages. Cappel unsettled also the traditional view of the literal integrity and sacredness of the Masoretic text, and showed that the different readings (\emph{Keri} and \emph{Ktib}\}, while they had no bearing on faith and morals, and therefore could not undermine the authority of the Scriptures, are not to be traced to willful corruption, but must be consulted, together with the ancient translations, in ascertaining the true text.\footnote{\emph{Critica sacra,} etc., Paris, 1650, folio; another edition, by Vogel, in three volumes, Halle, 1775–86. The work was finished October, 1634, but the printing was delayed by the opposition of the Protestants until his son, Jean Cappel, who seceded to the Roman Church, procured a royal privilege for its publication in Paris.}

These views, which are now generally accepted among Biblical scholars, met with violent opposition. Even the Buxtorfs, father and son, at Basle, who immortalized themselves by their rabbinical learning, advocated the divine inspiration of the Hebrew vowels. The Protestant orthodoxy of the seventeenth century, both Calvinistic and Lutheran, was very sensitive on this point, because it substituted an infallible Bible for an infallible papacy; while the Roman orthodoxy cared much more for the divine authority of the Church than for that of the Scriptures.

UNIVERSAL AND PARTICULAR PREDESTINATION.

Moses Amyraut, originally a lawyer, but converted to the study of theology by the reading of Calvin's 'Institutes,' an able divine and voluminous writer, developed the doctrine of hypothetical or conditional universalism, for which his teacher, John Cameron (1580–1625), a Scotchman, and for two years Professor at Sanmur, had prepared the way. His object was not to set aside, but to moderate and liberalize Calvinism by ingrafting this doctrine upon the particularism of election, and thereby to fortify it against the objections of Romanists, by whom the French Protestants were surrounded and threatened. Being employed by the Reformed Synod in important diplomatic negotiations with the government, he came in frequent contact with bishops, and with Cardinal Richelieu, who esteemed him highly. His system is an approach, not so much to Arminianism, which he decidedly rejected, as to Lutheranism, which likewise teaches a universal atonement and a limited election.\footnote{Amyraut's writings on this subject are: \emph{Traité de la Prédestination} (also in Latin), Saumur, 1634; \emph{Echantillon de la doctrine de Calvin sur la Prédestination,} 1637; \emph{De la justification,} 1638; \emph{De providentia Dei in malo,} 1638; \emph{Defensio doctrinæ Calvini de absoluto reprobationis decreto,} 1641; \emph{Dissertationes theol. quatuor,} 1645; \emph{Exercitatio de gratia universali,} 1646; \emph{Disputatio de libero hominis arbitrio,} 1647; \emph{Sermons sur divers textes de la Ste. Écriture,} 1653; \emph{Irenicum sive de ratione pacis in religionis negotio inter Evangelicos,} 1662. Amyraut wrote besides a system of Christian Ethics (in six volumes), and a number of exegetical and practical works. See a list in Herzog, Vol. I. pp. 296 sq.}

Amyraut maintained the Calvinistic premises of an eternal foreordination and foreknowledge of God, whereby he caused all things inevitably to pass—the good efficiently, the bad permissively.\footnote{'\emph{Ou de permettre tellement les mauvaises, que l’événement soit entièrement undubitable.}'} He also admitted the double decree of election and reprobation. But in addition to this he taught that God foreordained a \emph{universal salvation} through the universal sacrifice of Christ offered to all alike (\emph{également pour tous}), on condition of \emph{faith,} so that on the part of God's will and desire (\emph{voluntas, velleitas, affectus}) the grace is universal, but as regards the condition it is particular, or only for those who do not reject it and thereby make it ineffective. The universal redemption scheme precedes the particular election scheme, and not \emph{vice versa.} He reasons from the benevolence of God towards his creatures; Calvinism reasons from the result, and makes actual facts interpret the decrees. Amyraut distinguished between objective grace which is offered to all, and subjective grace in the heart which is given only to the elect. He also makes a distinction between natural ability and moral ability, or the power to believe and the willingness to believe; man possesses the former, but not the latter, in consequence of inherent depravity.\footnote{The same distinction was a century later made by New England Calvinists under the lead of Jonathan Edwards, who knew of the Saumur theology through the works of Stapfer.} He was disposed, like Zwingli, to extend the grace of God beyond the limits of the visible Church, inasmuch as God by his general providence operates upon the heathen, and may produce in them a sort of unconscious Christianity, a faith without knowledge; while within the Church he operates more fully and clearly through the means of grace. Those who never heard of Christ are condemned if they reject the general grace of providence; but the same persons would also reject Christ if he were offered to them. As regards the result, Amyraut agreed with the particularists. His ideal universalism is unavailable, except for those in whom God previously works the condition of faith, that is, for those who are included in the particular decree of election.\footnote{'\emph{Notre saint éternel depend de cette condition, que nous appellons la foy; cette foy depend de la grace de Dieu et de la puissance de son Esprit; cette grace, cette puissance de l’Esprit depend du conseil de l’election de Dieu, et ce conseil n’ayant autre fondement que sa volonté est constant et irrevocable, l’événement sursuit necessairement. Ce conseil depend de la libre volonté de Dieu.}' Schweizer, pp. 296 sq.}

Amyraut's doctrine created a great commotion in the Reformed Churches of France, Holland, and Switzerland. Jean Daillé (1594–1670),\footnote{Joh. Dallæi: \emph{Apologia pro duabus synodis nationalibus, altera Alensone} 1637, \emph{altera Carentone} 1645 \emph{habitis adv. Fr. Spanhemii Exercitationes de gratia universali.} Amst. 1655 (1227 pages), and \emph{Vindiciæ Apologiæ pro duabus synodis.} Amst. 1657. See extracts in Schweizer, pp. 390 sqq. Daillé is best known by his work \emph{Sur l’usage des Pères} (\emph{De Usu Patrum}).} David Blondel (1591–1655),\footnote{\emph{Actes authentiques touchant la paix et charité fraternelle aves les Protestantes,} etc. Amst. 1655. Blondel is best known by his \emph{De la primauté en église} (1641), and other historical works. He was Secretary of the French Synod, which made him honorary professor, with a salary sufficient to enable him to devote himself without pastoral care to his studies. He had an enormous memory, and when blind in his old age he dictated two folios on difficult points in chronology.} and others considered it innocent and consistent with the decrees of the Synod of Dort, where German Reformed and Anglican delegates professed similar views against the supralapsarianism of Gomarus. But Peter du Moulin (Molinæus, since 1621 Professor of the rival theological school of Sedan), Frederick Spanheim (1600–1649, Professor in Leyden), Andrew Rivet (1572–1651, Professor in Leyden), and the theologians of Geneva opposed it as a departure from the orthodox faith and a compromise between Calvinism and Arminianism.\footnote{See especially Pierre du Moulin:  \emph{Examen de la doctrine des Messieurs Amyraut et Têtard touchant la prédestination et les poins, qui en dependent,} Amsterd. 1638; and \emph{Eclaircissement des controverses Salmuriennes, ou défense de la doctrine des églises réformées sur l’immutabilité des decrets de Dieu,} etc. Leyden, 1648. Spanheim (the elder): \emph{Disputatio de gratia universali,} Lugd. Bat. 1644; and \emph{Exercitationes de gratia universali,} Lugd. Bat. 1646 (1856 pages). André Rivet:  \emph{Opera omnia,} Lugd. Bat. 1651–60, Vol. III. pp. 828–878.}

The friends of Amyraut urged the love, benevolence, and impartial justice of God, and the numerous passages in Scripture which teach that God loves 'the whole world,' that he will have 'all men to be saved,' that Christ died 'not for our sins only, but also for the sins of the whole world,' that 'he shut up \emph{all} in unbelief that he might have mercy upon \emph{all.}' On the other hand, it was objected that God could not really will and intend what is never accomplished; that he could not purpose an end without providing adequate means; that, in point of fact, God did not actually offer salvation to all; and that a universalism based on an impossible condition is an unfruitful abstraction.\footnote{The orthodox Lutherans, as far as they took notice of this controversy, saw in Amyraldism a concealment of Calvinism, a mockery on the part of God, a bridge to Syncretism, and characterized the \emph{gratia Amyraldina} as a \emph{gratia Calvina, non divina.} So Reheboldus, \emph{De natura et gratia Mosi Amyraldo opposita,} Gissæ, 1651 (quoted by G. Frank, Vol. I. p. 43). Among American divines, Dr. Hodge notices this controversy (\emph{Syst. Theology,} Vol. II. p. 322), and says that hypothetical redemption is liable to the objections against both Augustinianism and Arminianism. 'It does not remove the peculiar difficulties of Augustinianism, as it asserts the sovereignty of God in election. Besides, it leaves the case of the heathen out of view. They, having no knowledge of Christ, could not avail themselves of this \emph{decretum hypotheticum,} and must therefore be considered as passed over by a \emph{decretum absolutum.}' But Amyraut does notice the case of the heathen; see above.}

The national Synods at Alençon, 1637; at Charenton, 1645; and at Loudun, 1659 (the last synod permitted by the French Government), decided wisely and moderately, saving the orthodoxy of Amyraut, and guarding only against misconceptions. He gave the assurance that he did not change the doctrine, but only the method of instruction. And his opponents were forced at last to admit that the idea of a universal grace, by which no one was actually saved unless included in the particular, effective decree of election, was quite harmless. In this way universalism and particularism were equally sanctioned, and a schism in the French Church was avoided.\footnote{Schweizer, pp. 307 sqq.; Ebrard, p. 555.} The literary controversy continued for several years longer, and developed a large amount of learning and ability, until it was brought to an abrupt close by the political oppressions of the Reformed Church in France.\footnote{Schweizer gives a very full account of the writings on both sides, pp. 320–439. In modern times the great Schleiermacher has revived Amyraldism on German soil, but in a much bolder form, and at the expense of the Scripture doctrine of eternal punishment. He widens Calvinism (which he very acutely defends against Lutheranism and Arminianism) into a \emph{real} and \emph{effective} universalism of salvation, and makes the particularism of election and reprobation merely a temporary means to this end. Schweizer, one of his ablest pupils, adopts this solution of the problem in his \emph{Christliche Glaubenslehre,} Leipzig, 1872, Vol. II. Part II. pp. 78 sqq. and 444 sqq. But this solution is subject to all the objections of what in America is popularly called the system of Universalism: it turns conversion into a process of nature or necessity; it dulls the edge of warning; freedom implies the continued power of resistance; repentance becomes more and more difficult, and at last impossible, especially in hell and in the case of the devil and diabolized men.}

MEDIATE AND IMMEDIATE IMPUTATION\footnote{\emph{Syntagma thesium theologicarum in academia Salmuriensi disputatarum,} Ed. II. Salmur. 1664. Placeus:  \emph{De statu hominis lapsi ante gratiam,} 1640; his defense, \emph{De imputatione primi peccati Adami,} 1655, in his \emph{Opera omnia,} 1699 and 1702, two vols. Against him, A. Rivet:  \emph{Decretum Synodi nationalis Ecclesiarum Reformatarum Galliæ, A.D.} 1645 \emph{de imputatione primi peccati omnibus Adami posteris, cum Ecclesiarum et doctorum protestantium consensu, ex scriptis eorum collecto,} in the \emph{Opera Theol.} of Rivet, Rotterd. 1660, Tom. III. pp. 798–827, translated in part in the \emph{Princeton Review} for 1839, pp. 553–579. Comp. also Schweizer's art. \emph{Placeus,} in Herzog, Vol. XI. pp. 755–57, and several American treatises on the imputation controversy by Hodge, Baird, Landis, G. P. Fisher, quoted in my annotations to Lange's \emph{Com.} on Rom. v. 12 (pp. 191 sqq.), where the exegetical aspects are fully discussed in connection with the classical passage \textgreek{ἐφ᾽ ᾦ πάντες ἥμαρτον}}

All Augustinians and Calvinists agree in the doctrine of total depravity and original sin in consequence of Adam's fall; but differences arose among them concerning the imputation of Adam's sin and guilt to his posterity. The majority advocated the realistic theory of an actual, though impersonal and unconscious, participation of the whole human race in the fall of Adam as their natural organic head, who by his individual transgression vitiated the generic human nature, and transmitted it in this corrupt state by physical generation to his descendants. This, the old Augustinian view, was renewed by the Reformers. Others, since the seventeenth century, adopted the federal theory of a vicarious legal representation of mankind by Adam, in virtue of an assumed covenant of works made with him by the Sovereign Creator, to the effect that Adam should stand a moral probation in behalf of all his descendants (acting like a guardian for children yet unborn, or like a representative for future constituents), and that his act of obedience or disobedience, with all its consequences, should be judicially imputed to them, or accounted theirs in law.\footnote{\emph{Fœdus operum,} or \emph{fœdus naturæ,} as distinct from \emph{fœdus gratiæ.} The only Scripture passage which the Federalists alleged in favor of this primal covenant is Hos. vi. 7: ' For they, like Adam [\texthebrew{כְאָדָם}], have broken the covenant;' but others translate with the Sept.: 'They [are] like men [who] break a covenant' (\textgreek{ὡς ἄνθρωπος παραβαίνων διαθήκην.})} Still others combined the two theories so as to make imputation rest both on the moral ground of participation and on the legal ground of representation.

In connection with this doctrine of hereditary sin there arose among the Calvinists of the seventeenth century a controversy about immediate or antecedent, and mediate or consequent imputation.\footnote{Turretin (\emph{Instit.} Pars I. pp. 556, Loc. ix. \emph{de peccato,} Qu. X.) charges De la Place with inventing this distinction to evade the force of the synodical decision of Charenton, 1645. Augustine and the Reformers did not use the terms, and hence are quoted on both sides.} The theory of immediate imputation makes all descendants of Adam responsible for his disobedience as participants \emph{in actu,} and condemns them independently of, and prior to, native depravity and personal transgression, so that hereditary guilt precedes hereditary sin. The theory of mediate imputation makes inherent depravity derived from Adam, and this alone, the ground of imputation and condemnation (\emph{vitiositas præcedit imputationem}). The school of Montauban, Rivet of Leyden, the elder Turretin of Geneva, Heidegger of Zurich, Garissol, Maresius, and the supralapsarians and federalists advocated the former, some exclusively, some in connection with mediate imputation. La Place (Placeus) of Saumur denied immediate imputation of a foreign sin as arbitrary and unjust, and allowed only a mediate imputation, but claimed to be nevertheless in full harmony with Calvin's teaching on this subject.

The Reformed national Synod at Charenton, near Paris, in 1645, rejected the theory of La Place (yet without calling him to an account or naming him), at least so far as it restricts the nature of original sin to the mere hereditary corruption of Adam's posterity. In vindication of the decree of the Synod, Rivet prepared a collection of passages on imputation (many of them very general and inconclusive) from Reformed and Lutheran confessions and the writings of Calvin, Beza, Bullinger, Melanchthon, Chemnitz, and others.

THE CONSENSUS FORMULA.

Several years after the leaders of the Saumur theology had passed from the stage of history it was thought desirable by some of the prominent divines of Switzerland to protect their Churches against possible danger from the new doctrines of Saumur, which were imported through writings and students, and met with considerable sympathy, especially in Geneva. It was feared—and not without reason— that, however innocent in themselves, they might lead, by legitimate logical development, to an ultimate abandonment of the system of Calvinism.

Hence the new Formula of orthodoxy which forms the subject of this section, was agreed upon by the ecclesiastical and civil authorities of Zurich, Basle, and Geneva, and adopted in other Reformed cantons as a binding rule of public teaching for ministers and professors. Its authority was confined to Switzerland, and even there it could not maintain itself longer than about half a century. French ministers, who fled to Lausanne after the revocation of the Edict of Nantes, refused to sign it; the great Elector Frederick William of Brandenburg (1686), and afterwards the Kings of Prussia and England, and the \emph{Corpus Evangelicorum} at Ratisbon (1722), urged the Reformed cantons, in the interest of peace and union, to abandon the Formula. It gradually lost its hold upon the Swiss churches, and was allowed to die and be buried without mourners. Nevertheless the theology which it represents continues to be advocated by a respectable school of strict Calvinists in Europe, and especially in America.

The Helvetic Consensus Formula was not so much intended to be a new confession of faith, as an explanatory appendix to the former Confessions (resembling in this respect the Saxon Visitation Articles, which were an appendix to the Lutheran Formula of Concord, to guard the churches of Saxony against the dangers of crypto-Calvinism). The document does not breathe the fresh and bracing air of faith and religious experience which characterize the Confessions of the Reformation period. It is the product of scholasticism, which formularized the faith of Calvin into a stiff doctrinal system, and anxiously surrounded it with high walls to keep out the light of freedom and progress. Nevertheless it is more liberal than is generally represented and than might be expected from the bigotry and polemical violence of the seventeenth century. Heidegger was personally mild and modest; he spoke the truth in love, and resisted the pressure of extremists in Switzerland and Holland, who suspected even him of unsoundness, and desired a formal condemnation of the schools not only of Saumur but also of Cocceius and Cartesius. Instead of this, he speaks in the preface of the Formula, respectfully and kindly, of the Saumur theologians, and calls them venerable brethren in Christ, who built on the same foundation of faith, and whose peculiar doctrines are not condemned as heresies, but simply disapproved.\footnote{'\emph{Salvum enim utrinque per Dei gratiam stat fundamentum fidei. . . . Salva unitas corporis mystici et Sprititus. . . . Salvum denique apud nos semper tenerrimæ caritatis vinculum,}' etc. The original draft of the Formula was even milder and much shorter. Schweizer has, in a purely historical interest, vindicated the memory of Heidegger and the comparatively moderate character of the Consensus Formula. See his extracts from the MS. of Heidegger's Report, in Niedner's \emph{Zeitschrift,} above quoted, and his art. \emph{Heidegger,} in Herzog's \emph{Real. Encykl.}}

The Formula consists of a preface and twenty-six canons or articles, which clearly state the points of difference between strict Calvinism and Salmurianism. They teach the following points:

1. The literal inspiration of the Scriptures and the integrity of the traditional Hebrew text of the Old Testament, including the vowels as well as consonants; so that we need not resort to manuscripts, translations, and conjectures.\footnote{'\emph{In specie autem Hebraicus Veteris Testamenti Codex, quem ex traditione Ecclesiæ Judaicæ, cui olim Oracula Dei commissa sunt, accepimus hodieque retinemus, tum quoad } consonas,  \emph{tum quoad } vocalia,  \emph{sive puncta ipsa, sive punctorum saltem potestatem, et tum } quoad res,  \emph{tum } quoad verba  \textgreek{θεόπνευστος, } \emph{ut fidei et vitæ nostræ, una cum. Codice Novi Testamenti sit } canon  \emph{unicus et illibatus, ad cuius normam, ceu Lydium lapidem, universæ, quæ extant, Versiones, sive orientales, sive occidentales exigendæ, et sicubi deflectunt, revocandæ sunt.}' The same theory of plenary inspiration of words and thoughts, which dates from Rabbinical orthodoxy, but was not held by the Reformers, prevailed in the Lutheran Church since John Gerhard, and is even now extensively held, especially in England and America, by those whose faith in the Word of God is not affected by modern criticism. It was most ably defended by the venerable Dr. Louis Gaussen (1790–1863), Professor in the Free Church Theological School of Geneva, in his works on \emph{Theopneusty} (1840; second edition, 1842), and on the \emph{Canon} (1862, two vols.). Dissent from him led to the resignation of his colleague, Scherer. Gaussen admitted, however, the individualities of the sacred writers, and compares them to the keys of an immense organ, on which the Holy Spirit played.} Art. 1–3. Against Cappel.

This attempt to canonize the Hebrew vowels gave great offense to Claude, Daillé, and other French Calvinists; and Heidegger explained to Turretin that the object of the Formula was only to guard the authority and integrity of the original text, and not to decide grammatical and critical questions. But in its natural effect such a mechanical theory of inspiration, which, to be of any practical use, requires a perpetual literary miracle in the preservation of the text, would supersede all textual criticism, and make the Targums, the Septuagint, the Vulgate, and other ancient versions, worse than useless.

2. God decreed from eternity, first, to create man innocent; second, to permit (\emph{permittere}) the fall; third, to elect some to salvation, and thus to reveal in them his mercy, but to leave the rest in the corrupt mass (\emph{alios vero in corrupta massa relinquere}), and to devote them to eternal perdition. (This is clearly the Augustinian infralapsarianism.) In the gracious decree of election Christ himself is included, as the Mediator and our first-born Brother. The doctrine of an antecedent hypothetical will or intention of God\footnote{Called \emph{voluntas conditionata, velleitas, misericordia prima, desiderium inefficax.}} to save all men on condition of faith is rejected as unscriptural and as involving God in imperfection and contradiction. Art. 4–6. Against Amyraut.

3. The covenant of works made by God with Adam before the fall, promising to him eternal life (symbolized by the tree of life in Paradise), on condition of perfect obedience. Art. 7–9. Against Amyraut.

4. Immediate imputation of Adam's sin to all his posterity who fell in him, their representative head (\emph{in ipso ut capite et stirpe}), and forfeited the promised blessing of the covenant of works. Man is thus doubly condemned, for his participation in the sin of Adam and for his hereditary depravity; to deny the former makes the latter doubtful.\footnote{Art. X. '\emph{Censemus igitur peccatum Adami omnibus ejus posteris judicio Dei arcano et justo imputari}' (Rom. v. 12, 19; 1 Cor. xv. 21, 22). Art. XI. '\emph{Duplici igitur nomine post peccatum homo natura, indeque ab ortu suo, antequam ullum actuale peccatum in se admittat, iræ ac maledictioni divinæ obnoxius est; primum quidem ob } \textgreek{παράπτωμα } \emph{et inobedientiam, quam in Adami lumbis commisit; deinde ab consequentem in ipso conceptu hereditariam corruptionem insitam, qua tota ejus natura depravata et spiritualiter mortua est, adeo quidem, ut recte peccatum originale statuatur duplex . . . imputatum videlicet, et hereditarium inhærens.}'} Art. 10–12. Against La Place, not because he asserted mediate or consequent imputation (which the Formula likewise teaches), but because he excluded the other.

5. Limited atonement. Christ died only for the elect, and not indiscriminately for all men.\footnote{Art. XIII. '\emph{Pro solis electis ex decretorio Patris consilio propriaque intentione diram mortem oppetiit} [\emph{Christus}], \emph{solos illos in sinum paternæ gratiæ restituit, solos Deo Patri offenso reconciliavit et a maledictione legis liberavit.}' Art. XVI. '\emph{Haud probare possumus oppositam doctrinam illorum qui statuunt, Christum propria intentione et consilio tum suo tum Patris ipsum mittentis, mortuum esse pro omnibus et singulis, addita conditione impossibili, si videlicet credant.}' The ablest modern advocate of this limited atonement theory is Dr. Hodge, \emph{Syst. Theol.} Vol. II. pp. 544 sqq.} The infinite value and inherent sufficiency of Christ's satisfaction is not denied, but the divine intention and the practical efficiency are limited, and adjusted to the particularism of the decree of election. Art. 13–16. Against Amyraut.

6. The actual vocation to salvation never was absolutely general (\emph{numquam absolute universalis}), but was confined to Israel in the old dispensation and to Christians in the new (Matt. xi. 25; Eph. i. 9). God's revelation in nature and providence (Rom. i. 19, 20) is insufficient for purposes of salvation, though it leaves the heathen without excuse for rejecting even this remnant of the knowledge of God. The external call of God through his ``Word is always serious, and so far effective that it works salvation in the elect, and makes the unbelief of the reprobate inexcusable.\footnote{Art. XIX. '\emph{Vocatio externa quæ per præconium Evangelicum fit, etiam vocantis Dei respectu, seria et sincera est. . . . Neque voluntas illa respectu eorum, qui vocationi non parent, inefficax est, quia semper Deus id, quod volens intendit, assequitur,}' etc.} Art. 17–20. Against Amyraut, who extended the vocation beyond the limits of the visible Church and the ordinary means of grace.

7. The natural as well as moral inability of man to believe the gospel of himself.\footnote{Art. XXI. 'Moralis  \emph{ea impotentia dici possit, quatenus scilicet circa subjectum et objectum morale versatur: } naturalis  \emph{tamen esse simul et dici debet, quatenus homo } \textgreek{φύσει, } \emph{natura, adeoque nascendi lege, inde ab ortu est filius iræ}' (Eph. ii. 2). Dr. Hodge likewise defends this doctrine against the New School Calvinists, who, with Amyraut, claim for man the natural ability, but admit his moral inability.} This twofold inability has its ground in the depravity of our nature, from which only the omnipotent power of the Holy Spirit can deliver us (1 Cor. ii. 14; 2 Cor. iv. 6). Art. 21, 22. Against Amyraut.

8. A twofold covenant of God with man—the covenant of \emph{works} made with Adam and through him with all men, but set aside by the fall, and the covenant of \emph{grace} made only with the elect in Christ, which is forever valid, and exists under two economies, the Jewish and the Christian. The saints of the Old Testament were saved by the same faith in the Lamb of God as we are (Apoc. xiii. 8; Heb. xiii. 8; John xiv. 1); for out of Christ there is no salvation. The doctrine of the Holy Trinity is revealed in the Old Testament in words, figures, and types, sufficiently for salvation, though not as clearly as in the New. For no one can believe in Christ without the Holy Spirit, the third person in the Trinity. Amyraut's doctrine of three essentially different covenants—natural, legal, and evangelical, with different degrees of knowledge and piety—is disapproved. Art. 23–25.

The concluding article (the 26th) prohibits the teaching of new or doubtful and unauthorized doctrines which are contrary to the Word of God, the Second Helvetic Confession, the Canons of the Synod of Dort, and other Reformed symbols.

\xsection{The Reformed Confessions of France and the Netherlands.}

\xsubsection{The Gallican Confession. A.D. 1559.}

§ 62. The Gallican Confession. A.D. 1559.

Literature.

I. Editions of the Gallican Confessions.

The original French text in  Theod. de Beza:  \emph{Histoire ecclésiastique des églises réformées au royaume de France,} Antw. 1580, Tom. II. pp. 173–190; in Niemeyer's  \emph{Collectio Conf. in eccles. reformatis public}. pp. 311–326; and in the \emph{Zeitschrift für die histor. Theologie} for 1875, pp. 506–544, with an introduction by Dr. Heppe. The shorter recension in the new edition of Calvin's \emph{Opera,} Vol. IX. pp. 739 sqq. The text, as revised by the Synod of Rochelle (1571), was often printed in French Bibles, and separately. Comp. the Toulouse edition of 1864, entitled \emph{Confession de Foi et Discipline ecclésiastique des églises réformées de France} (\emph{Société des livres religieux,} pp. 9–35).

The Latin translation: \emph{Gallicarum ecclesiarum Confessio Christianissimo Carolo IX. regi anno MDLXI. exhibita. Nunc vero in Latinum conversa, ut omnino constet eas ab omnibus hæresibus sive sectis esse prorsus aliena. Anno Domini} 1566—and often reprinted; also in \emph{Corpus et Syntagma Conf.} 1654, pp. 77–88, and in Niemeyer's  \emph{Collectio,} pp. 327–339.

A German translation appeared first at Heidelberg, 1562 (see Niemeyer, \emph{Præfat.} p. 1.); also in Böckel's  \emph{Bekenntniss-Schriften der evang. reform. Kirche,} pp. 461–474.

An English translation in John Quick's  \emph{Synodicon in Gallia Reformata,} Lond. 1692, Vol. I. pp. vi.-xvi.

II. History of the Reformation and the Reformed Church in France.

See partly the Literature on Calvin, quoted p. 421.

Theod. Beza:  \emph{Histoire ecclés. des églises réformées au royaume de France} (1521–63), Antw. 1580, 3 vols.

Jean Crespin (d. 1572): \emph{Livre des martyrs} (\emph{Acta Martyrum}), \emph{depuis Jean Hus jusqu’en} 1554. Geneva, 1560; enlarged edition, Genève, 1617, and Amsterd. 1684.

Serranus (Jean de Serres, historiographer of France, 1540–98): \emph{Commentarius de statu religionis et reipublicæ in regno Galliæ,} 1571–73 (five parts).

Theod. Agrippa d’Aubigné (Albinæus, a Huguenot in the service of Henry IV.; d. at Geneva, 1630): \emph{Histoire universelle de mon temps,} 1616–20, 3 vols.

Du Plessis Mornay:  \emph{Mémoires et correspondance.} Paris, 1824–25.

John Quick (a learned Non-conformist, d. 1706): \emph{Synodicon in Gallia Reformata; or, the Acts, Decisions, Decrees, and Canons of the National Councils of the Reformed Churches in France.} London, 1692, 2 vols. fol. (with a history of the Church till 1685). Much more accurate than Aymon.

Aymon:  \emph{Tous les synodes nationaux des églises réformées de France.} La Haye, 1710, 2 vols. 4to.

E. A. Laval:  \emph{Compendious History of the Reformation in France . . . to the Repealing of the Edict of Nantes.} London, 1737–41, 7 vols.

Smedley:  \emph{History of the Reformed Religion in France.} London, 1832, 3 vols.

G. de Félice:  \emph{Histoire des Protestants en France.} Toulouse, 1851; Engl. translation, by Lobdel, 1851. By the same: \emph{Histoire des synodes nationaux des églises réformées de France.} Paris.

W. G. Soldan:  \emph{Geschichte des Protestantismus in Frankreich bis zum Tode Karl’s IX.} Leipzig, 1855, 2 vols.

G. von Polenz:  \emph{Geschichte des französischen Calvinismus bis zur Nationalversammlung i. J.} 1789, \emph{zum Theil aus handschriftl. Quellen.} Gotha, 1857–64, 4 vols.

E. Stähelin:  \emph{Der Uebertritt Heinrich’s IV.} Basle, 1856.

Ath. Coquerel:  \emph{Histoire des églises du désert.} Paris, 1857, 2 vols.

W. Haag:  \emph{La France protestante.} Paris, 1858 (biographies).

Weiss:  \emph{Histoire des réfugiés protestants de France depuis la révocation de l’édit de Nantes jusqu’à nos jours.} Paris, 1853; English translation, London, 1854, 2 vols.

Much valuable information on the early history of Calvinism and French Protestantism generally is contained in Herminjard's  \emph{Correspondance des réformateurs dans les pays de langue français,} Genève and Paris, 1866 sqq. (so far 4 vols.), and in the \emph{Bulletin de la Société de l’histoire du Protestantisme français. Documents historiques inédits et originaux XVI \textsuperscript{e}, XVII\textsuperscript{ e}, et XVIII\textsuperscript{ e} siécles.} Paris (3, rue Lafitte), 1854–63; so far 22 vols.

III. General Histories of France touching upon the Reformation Period.

Thuanus (Jacques Auguste de Thou—born, 1553; died, 1617): \emph{Historiarum sui temporis} libri 138, from 1546–1607 (several editions in five, seven, and sixteen volumes). The author was a moderate Catholic, witnessed the massacre of St. Bartholomew, and helped to prepare the Edict of Nantes. His history was put in the Index Expurg. 1609.

Lacretelle:  \emph{Histoire de France pendant les guerres de la religion.} Paris, 1822, 4 vols.

Sismondi:  \emph{Histoire des Français,} Par. 1821–44, 31 vols. (from vol. 16th).

Jules Michelet (born, 1798): \emph{Histoire de France,} 1833–62, 14 vols. (vols. 8 and 9).

Sir James Stephen:  \emph{Lectures on the History of France,} 1857, third edition, 2 vols.

Leop. Ranke:  \emph{Französische Geschichte namentlich im} 16. \emph{und} 17. \emph{Jahrh.} 1852–68, 6 vols. (English translation in part, London, 1852, 2 vols.)

Henri Martin:  \emph{Histoire de France depuis les temps les plus reculés jusqu’en} 1789, fourth edition, Paris, 1855–60, 16 Tom. (Vols. VII. to X.).

FRENCH PROTESTANTISM.

In France the Reformation seemed to be better prepared than even in Germany, if we look only at the surface of the situation. The French Church had always maintained a certain independence of Rome, under the name of Gallican rights or liberties. Paris was, it is true, the chief seat of orthodox scholasticism, and the Sorbonne took an early opportunity to condemn Luther and his writings (1521); but it nursed also the spirit of mysticism and disciplinary reform, which led to the Councils of Pisa, Constance, and Basle. In the South a remnant, of the Waldenses had survived the bloody persecutions. The humanistic studies flourished greatly at Paris, Orleans, Bourges, and found favor at the court of Francis I. (1494–1547), who invited classical scholars from Italy, thought of calling Erasmus and even Melanchthon to his capital, and aided, for political, reasons, the Protestants in Germany, while yet he inflicted imprisonment and death upon them in France.

For half a century, and amid bloody civil wars, three conflicting tendencies, represented by Calvin, Rabelais, and Loyola—who happened to be in Paris at about the same period—struggled for the mastery: Calvinism, with its high intelligence and uncompromising virtue; the Renaissance, with its elegant culture and frivolous skepticism; and Jesuitism, with its reactionary and unscrupulous fanaticism. Francis I. wavered between the Renaissance, which suited his natural taste, and Romanism, which was the religion of the masses of Frenchmen; his gifted sister, Queen Margaret, of Navarre (grandmother of Henry IV.), protected the Reformation and the Renaissance, and harbored at one time Calvin, and at another the Libertines. Romanism triumphed first over Protestantism, and afterwards over semi-evangelical Jansenism, and France reaped infidelity and the Revolution.

Calvinism, always in the minority, and too stern and exacting for the national character, after a period of heroic martyrdom, gained for a time a limited legal existence under Henry IV. in the Edict of Nantes (1598), but was expelled under Louis XIV. to fertilize other countries, and reduced to a proscribed sect of the desert at home, where nevertheless, like the burning bush, it could not be consumed, and was providentially preserved for better days.\footnote{On an old seal, the device of which has been preserved, the French [Reformed] Church may be seen represented under the image of the burning bush of Moses, with this motto: ``\emph{Flagror, sed non comburor.}'' These words sum up the tragical history of our Church. This Church has been essentially militant; she has known better, perhaps, than any other what it is to fight for life. . . . Most young Frenchmen are brought up in a holy horror of Protestantism; and traces of this early impression are even found clinging to the minds of men of independent thought—nay, of those whose boast it is that they are free-thinkers.'—A. Decoppet, in his report on the Reformed Church in France, at the General Conference of the Evangelical Alliance in New York, 1873. See \emph{Proceedings,} p. 72. The synodical seal, with the above motto and the date 1559, is reproduced on the title-page of the first volume of Bersier's \emph{Histoire du Synode Général de l’église réform, de France} 1872 (Paris, 1872).}

The father of French Protestantism in its unorganized form is Jacques Lefèvre d’Etaples (Faber Stapulensis, 1455–1537), Professor of the Sorbonne and tutor of the royal princes. He translated the Bible from the Vulgate (completed 1530); he taught, even before Luther and Zwingli,\footnote{His Commentary on the Pauline Epistles appeared in 1512.} the doctrine of justification by faith without human works or merit, and the supremacy of the Bible as a rule of faith, and predicted a reformation, saying to his pupil, Farel, 'God will renovate the world, and you will be a witness of it;' but he had to flee to Strasburg, and afterwards to the court of Queen Margaret.

In the same spirit labored his friends and pupils—Briçonnet, Bishop of Meaux, who fostered evangelical doctrines and practices in his diocese, but afterwards timidly joined in the condemnation of Luther; Melchior Wolmar, a native of Germany, Professor of Greek in Bourges and teacher of Calvin; Louis de Berquin (1489–1599), a royal counselor, who was burned at the stake; Clement Marot (1495–1544), the favorite poet of his age and translator of the Psalms in verse; Peter Robert Olivetan (d. 1538), a relative of Calvin and translator of the Bible into French (printed at Neuchatel, 1535); William Farel (1489–1565), Peter Viret, Anton Froment, Calvin, and Beza—who were driven to French Switzerland. The radical extravagances of Anabaptists and anti-Trinitarians also spread in France, and were confounded by the government with the sound evangelical doctrines, and made a pretext for persecution.

But it was only after Calvin, himself the greatest Protestant of France, had taken up his permanent abode in Geneva, that the Reformation movement was organized into a separate Church, and acquired a national importance. He therefore, and his friend and successor Beza, may be regarded as the fathers of the Reformed Church of France. Geneva became an asylum for their persecuted countrymen, and the nursery of evangelists. Henceforward French Protestantism assumed a Calvinistic type in doctrine and discipline, but, owing to the hostile attitude of the government, it was kept separate and distinct from the state. Although cruelly persecuted, and numbering its martyrs by thousands, it spread rapidly among the middle and higher classes, and in 1558 it embraced four hundred thousand followers.

The first national Synod was held in Paris, May 25–28, 1559, under the moderatorship of François de Morel, then pastor of Paris, a friend and pupil of Calvin.\footnote{An account of this Synod in Polenz, Vol. I. pp. 435 sqq. Owing to the troubles of the times there were only eleven congregations represented—Dieppe, Paris, Angers, Orleans, Tours, etc.} It gave the Reformed Church a compact organization by the adoption of the Gallican Confession of Faith, in connection with a Presbyterian form of government and discipline, which remained the firm basis of the Church as long as she was allowed to exist and to hold national Synods, twenty-nine in all, the last being that at Loudun, 1659.

ANTOINE DE CHANDIEU.

The Gallican Confession is the work of John Calvin, who prepared the first draft, and of his pupil, Antoine de la Roche Chandieu, who, with the Synod of Paris in 1559, brought it into its present enlarged shape.\footnote{Quick, in the \emph{Synod. Gall. Ref.} (London, 1692, Vol. I. p. xv.), says: 'Calvin first drew up the Confession itself.' But Beza, in his History, connects Chandieu prominently with the origin of the Confession, without expressly naming him as the author. It is based, in part at least, on a shorter Confession to the King (\emph{Au Roy}), which Calvin probably prepared, 1557, for the congregation of Paris, in vindication against false charges. See Bonnet, \emph{Lettres de Calvin,} Tom. II. p. 131, and \emph{Opera,} Vol. IX. p. 715 (comp. \emph{Proleg.} p. lix.). Calvin also wrote another French Confession of Faith, in the name of the French Churches, during the war, to be presented to the Emperor Maximilian and the German Diet at Frankfort, 1562. Reprinted in \emph{Opera,} Vol. IX. pp. 753–772.}

Chandieu, or, as he is also called, Sadeel,\footnote{\par The Hebrew name for Chandieu, \emph{i.e. Champ de Dieu,} Field of God.} was born 1534, of a wealthy noble family, in the castle Chabot, in Burgundy, studied law in the University of Toulouse, was converted to Protestantism in Paris, renounced a splendid career, studied theology at Geneva, was ordained 1554, and elected pastor of the small Reformed congregation in Paris. He was imprisoned 1557, escaped under the name Sadeel, was again imprisoned, but delivered by the hand of Anton de Bourbon (the father of Henry IV.), engaged in mission work near Poitiers, and returned to his congregation in Paris, 1559. He presided over the third National Reformed Synod at Orleans, 1562, attended as delegate the seventh National Synod at La Rochelle, 1571, barely escaped the massacre of St. Bartholomew (Aug. 24), fled with his family to Geneva, and taught theology at Lausanne. He received a commission in 1578 to attend a Protestant Union meeting at Frankfort, suggested by the Elector John Casimir, but never carried out. He was called back to France as chaplain of King Henry of Navarre (afterwards Henry IV.), returned to Geneva, 1589, and labored there as pastor and Professor of Hebrew till his death, Feb. 23, 1591. Beza esteemed him very highly. De Thou recommends him for 'noble birth, fine appearance, elegant manners, learning, eloquence, and rare modesty.'\footnote{\emph{Histor. Lib.} XXIX. (on occasion of his election as president of the National Synod of Orleans, 1562): '\emph{Ecclesiæ Parisiensis pastor, adolescens, in quo præter gentis nobilitatem, oris venusta facies, eruditio, eloquentia cum singulari modestia certabant.'}} Sadeel wrote twenty-three books and tracts, mostly in Latin, some in French, relating to Christian doctrines (especially the Word of God; the priesthood and sacrifice of Christ; the human nature of Christ; the spiritual manducation of his body), Church discipline, and the history of martyrs.\footnote{Ant. Sadeelis  \emph{Opera theologia,} edited after his death by his son John, and dedicated to Henry of Navarre, Genev. 1592; fifth edition, 1620. He also wrote three sonnets on Calvin's death, and \emph{Octonaires sur la vanité du monde.} See \emph{France protestante,} s. v. Chandieu, Vol. III. pp. 320–332; \emph{Bulletin de la société de l’histoire du protestantisme français,} 1853, p. 279; G. von Polenz, \emph{Gesch. des franz. Calv.,} Vol. I. p. 435; Borrel (pastor in Nismes), art. \emph{Chandieu} in Herzog, \emph{Real-Encykl.} Vol. XIX. p. 318. On Sadeel's Christology, see Dorner, \emph{Entwicklungsgesch. der Lehre von der Person Christi,} Vol. II. pp. 725, 733 sq., etc.}

THE GALLICAN CONFESSIONS.

On a visit to the mission church of Poitiers, after the holy communion, Chandieu was requested by the brethren to suggest to the church in Paris the importance of preparing a common confession of faith and order of discipline.\footnote{Beza, \emph{Histoire,} etc., Tom. I. pp. 172 sq., quoted in Calv. \emph{Opera,} Vol. IX. p. lvii.} Calvin was consulted, and sent three delegates with a draft of a confession to Paris. This was enlarged and adopted by the Synod at Paris, 1559; presented, with a Preface, to King Francis II. at Amboise, 1560, and afterwards by Beza to Charles IX. at the religious conference in Poissy, 1561. It was revised and ratified at the seventh National Synod held at La Rochelle, 1571, with Beza as moderator, in the presence of the Queen of Navarre and her son (Henry IV.), and Admiral Coligny. Hence it is also called the '\emph{Confession of Rochelle.}' Three copies were written on parchment—one for La Rochelle, one for Geneva, one for Béarn—and signed by the ministers and elders present.\footnote{The Geneva copy has been reproduced in fac-simile by Ed. Delessert. See Heppe, p. 513.}

As to the text, the French is the original, but it exists in two recensions: the shorter contains thirty-five articles, the larger forty articles. The latter was sanctioned by the Synod of La Rochelle.\footnote{'\emph{D’autant que nostre confession de foy est imprimée de differentes manières, le Synode declare que celle-là est la véritable confession de nos Églises reformées de France qui commence par ces paroles}: ``\emph{Nous croyons qu’il y a un seul Dieu,'' etc., laquelle a esté dressée au premier Synode national tenu à Paris, le} 25 \emph{mai de l’an} 1559.' Quoted in Calv. \emph{Opera,} Vol. IX. p. lix., from Aymon. The shorter edition is printed in \emph{Opera,} Vol. IX. p. 739, under the title \emph{Confession de Foy faite d'un commun accord par les Églises qui sont dispersées en France et s'abstienent des idolatries papales.} The larger edition is incorporated in the third volume of this work. It substitutes in the title for '\emph{qui sont},' etc., the words '\emph{qui désirent vivre selon la pureté de l’évangile de nostre Seigneur Jesus-Christ.}' Comp. Heppe, pp. 509 sqq.} It was often printed in different languages, and attached to many French Bibles.

CONTENTS.

The Gallican Confession is a faithful summary of the doctrines of Calvin. It begins with God (art. 1), his revelation (2), and the Scriptures as the Word of God and certain rule of our faith, which is above all customs, edicts, decrees, and councils (3–5). The three œcumenical Symbols are adopted (5), because they agree with the Word of God. The Holy Scripture teaches the unity of essence and tripersonality of God—the Father, who is the first cause, principle, and origin of all things; the Son, his eternal Word and Wisdom, eternally begotten by the Father; the Holy Spirit, his virtue and power eternally proceeding from both (6). God in three co-working persons created all things, visible and invisible (7); and governs all things, even sin and evil, yet without being the author of sin, but so making use of devils and sinners as to turn to good the evil which they do, and of which they alone are guilty (8). Man was created pure and perfect, but fell by his own guilt, and became totally corrupt and a slave of sin, although he can still discern good and evil (9). All posterity of Adam is in bondage to original sin, which is an inherited evil (not an imitation merely), and sufficient for condemnation; even after baptism it is still sin, but the condemnation of it is abolished out of free grace (10, 11). God, according to his eternal and immutable purpose, calls out of this corrupt mass those whom he has chosen in the Lord Jesus Christ, without regard to their merit, to the praise of his glorious grace, leaving the rest in their corruption and condemnation, to the praise of his eternal justice (12).\footnote{'\emph{Laissant les autres en cette même corruption et condamnation, pour démontrer en eux ca justice, comme aux premiers il fait luire les richesses de sa miséricorde.}'}

Jesus Christ is our all-sufficient Saviour, and 'made unto us wisdom, and righteousness, and sanctification, and redemption' (13). He assumed our human nature, being God and man in one person, like unto us in body and soul, yet without sin. We detest all ancient and modern heresies on the person of Christ, especially that of Servetus (14). The two natures in the one person of Christ are inseparably united, and yet remain distinct, so that the divine nature retains its attributes, being uncreated, infinite, and omnipresent, and the human nature continues finite and circumscribed (15). By the one sacrifice of Christ on the cross we are reconciled to God, and have the forgiveness of all our sins (16, 17). Our justification is founded on the remission of sins by the atoning death of Christ, without any merit of our own, and is apprehended and appropriated by faith alone (18–20). By this faith we are regenerated, and receive grace to lead a holy life, according to the Holy Spirit dwelling in us. Faith, then, of necessity produces good works, but these works are not accounted to us for righteousness, which must rest exclusively on the satisfaction of Christ; otherwise we would never have peace (21, 22). Christ is our only Advocate before the Father. We therefore reject the intercession of saints, and all other devices which detract from the all-sufficient sacrifice of Christ, as purgatory, monastic vows, pilgrimages, auricular confession, indulgences. We reject them not only on account of the false idea of merit attached to them, but also because they impose a yoke upon the conscience (23, 24).

The Church, with the ministry and preaching of the Word of God, is a divine institution, and must be respected and obeyed. The true Church is the company of believers who agree to live according to the Word of God, and to advance in holiness. Nevertheless there may be hypocrites and reprobates in it, who can not destroy its character and title. We reject the papacy for its many superstitions, idolatries, and corruptions of the Word and Sacraments. But as some trace of the true Church is left in the papacy, together with the virtue and efficacy of baptism, and as the efficacy of baptism does not depend upon the personal character of the minister, we teach that those who received baptism in the Romish Church do not need a second baptism. The true Church should be governed by pastors, elders, and deacons. All true pastors have the same authority and power under one head, the only sovereign and universal bishop, Jesus Christ; and consequently no Church shall claim any authority or dominion over the other (25–33).\footnote{The National Synod of Gap, 1603, inserted an article (31) declaring the pope to be 'the Antichrist and man of sin,' but the Synod of La Rochelle (1607) struck it out on account of the protest of the king. Heppe, p. 537.} The Sacraments are added to the Word as pledges and seals of the grace of God to aid and comfort our faith. They are external signs through which God operates by the power of his Spirit. Their substance and truth is in Christ; separated from him they are empty shadows. There are but two Sacraments: Baptism and the Lord's Supper. Baptism is the permanent pledge and signature of our adoption; by it we are grafted into the body of Christ, so as to be cleansed by his blood and renewed by the Holy Ghost. The Lord's Supper is the witness of our union with Christ, who truly nourishes us with his broken body and shed blood through the secret and incomprehensible power of his Spirit. We hold that this is done spiritually and by faith, not because we substitute imagination or thought for reality and truth, but because this great mystery surpasses our senses and the order of nature. In Baptism and the Lord's Supper God really gives us what they represent. Those who approach the Lord's table with true faith, as a vessel, receive the body and blood of Christ, which nourish the soul no less than bread and wine nourish the body (34—38).

God has instituted kingdoms, republics, and other forms of government, whether hereditary or elective, for the order and peace of society. He has given the sword to the magistrate for the punishment of sin and crime, and the transgressions of the first as well as the second table of the Decalogue.\footnote{'\emph{Il a mis le glaive en la main des magistrats pour réprimer les pechés commis non seulement contre la seconde table des commandements de Dieu, mais aussi contre la première.}' This clause justifies civil punishment of heresy. It is one of the chief causes why even orthodox members of the National Synod of 1872 were opposed to the re-adoption of this Confession in full.} We must therefore obey the magistrate, pay tribute and taxes with a good and free will, even if the rulers are unbelievers. We therefore detest those who would resist authority, establish community of goods, and overthrow the order of justice (39, 40).

\xsubsection{The Declaration of Faith of the Reformed Church in France. A.D. 1872.}

§ 63. The Declaration of Faith of the Reformed Church in France. A.D. 1872.

Literature.

\emph{XXX\textsuperscript{ e} Synode général de l’Église Réformée de France, Première session tenue à Paris du} 6 \emph{Juin au} 10 \emph{Juillet,} 1872. \emph{Procès verbaux et actes publiés par l’ordre du Synode.} Paris, 1873. (Comp. also the \emph{Compte Rendu} of the secretaries, and the discourses of Laurens, Pécaut, Ath. Coquerel, Fontanès, Colani, and Clamagerau, which appeared during the session.)

\emph{Do. Second session tenue a Paris du} 20 \emph{Novembre au} 3 \emph{Décembre,} 1873. Paris, 1873.

Eugène Bersier:  \emph{Histoire du Synode général de l’Église Réformée de France, Paris,} 6 \emph{Juin au} 10 \emph{Juillet,} 1872. Paris, 1872, 2 vols. E. B. attended the Synod of 1872, as a delegate of the Free Church of France, and gave an account of it in the \emph{Journal de Genève}. He has since joined the National Church.

The thirtieth meeting of the General Synod of the Reformed Church in France forms an epoch in its history. It resumed the series of twenty-nine National Synods after an interruption of two hundred and twelve years.\footnote{See a list of the French National Synods in Bersier, Vol. II. pp. 429 sqq.} The last was held at Loudun (Anjou), and was brought to a close in Jan., 1660, by an order of Louis XIV. prohibiting such synods in future, on the pretext that they were too expensive and troublesome, and that their business could be transacted in provincial synods. Daillé, the moderator, protested in vain. This act of injustice aimed to destroy the force of the Reformed communion by breaking it up into incoherent sections, and was crowned by the sweeping Revocation of the Edict of Nantes (Oct. 22, 1685), which deprived France of a million of her best citizens, and reduced the remnant of Protestants to a forbidden sect The history of this dark period is full of touching and dramatic interest. 'The Reformed Church of the Desert,' under 'the most Christian' King of France, like the primitive Church under the sway of heathen Rome, had to hold its synodical meetings in the open fields, in mountain-passes, and caverns of the earth.\footnote{Eight of these forbidden Synods were held between 1726 and 1763.} In those meetings the Gallican Confession was read, and prayer offered for the persecuting king. The spread of infidelity, which followed as a reaction against the tyranny of superstition and bigotry, brought first an edict of universal toleration under Louis XVI. (1787), and soon afterwards a total overthrow of Christianity and social order, until Napoleon, in 1802, restored the Roman Church as the religion of the majority of Frenchmen, and the Reformed Church as the religion of a small though respectable minority, but both under the pay and control of the State, and without the right of synodical self-government and discipline.\footnote{Napoleon's motive was chiefly of a political character. He needed religion as a basis of society, and Protestantism as a check upon the ambition of popery; yet he professed to a number of Protestant pastors to be a friend of the liberty of conscience, whose 'indefinite empire begins where the empire of law ends,' and he authorized them to brand with the name of Nero any one of his successors who should violate this liberty. Napoleon III. professed the same policy, but threw the weight of his power into the scale of Romanism, and made a distinction between the private liberty of conscience, which nobody can touch, and the public liberty of worship, which requires a recognition by the State.}

This right, denied by the Bourbon, the Napoleon, and the Orleanist dynasties, was at last restored to the Reformed Church by the Republican government under Thiers, who, by an edict of Nov. 29, 1871, authorized the Consistories in France and Algiers to elect delegates to a General Synod. Under these auspices the General Synod convened in the Temple du Saint-Esprit, at Paris, from June 6th to July 10th, 1872. It consisted of one hundred and eight delegates (forty-nine ministers and fifty-nine laymen), the legitimate descendants of those few humble but enthusiastic and heroic pastors and elders who met in the same city, in 1559, with torture and death staring them in the face. It was opened by a sermon of pastor Charles Babut of Nîmes on John viii. 14. Charles Edouard Bastie, pastor of Bergerac (Dordogne), was elected moderator. The object of the Synod was to again effect a complete organization on the basis of a confession of faith and a system of discipline.

But the preparation and adoption of a confession of faith is a more difficult task in the nineteenth century than it was in the sixteenth. For, like all other Protestant denominations, the French Church had during the eighteenth century undergone a theological revolution, and is still in a process of transition. The doctrinal system of the Gallican Confession had lost its hold upon a large portion of the clergy and laity; and even the most orthodox Protestants could not subscribe that article which, in harmony with the general sentiment of the sixteenth century, conceded to the civil government (hostile as it then was to the Huguenots) the power to punish heresy by the sword.\footnote{Art. 39: 'God has put the sword into the hands of magistrates to suppress crimes against the \emph{first} as well as against the second table of his Commandments.' It was on that ground that Servet's execution in Geneva for blasphemy was justified.} On the other hand, that venerable document, which embodied the faith of the fathers and martyrs of the French Church, could not be ignored without ingratitude and want of self-respect. Under these circumstances the General Synod, at its thirteenth session, June 20, 1872, adopted a middle course in the following declaration of faith, proposed by Charles Bois, Professor of Church History at Montauban:

'The Reformed Church of France, on resuming her synodical action, which for so many years had been interrupted, desires, before all things, to offer her thanks to God, and to testify her love to Jesus Christ, her Divine Head, who has sustained and comforted her during her successive trials.

'\emph{Au moment où elle reprend la suite de ses Synodes, interrompus depuis tant d’années, l’Église réformée de France éprouve, avant toutes choses, le besoin de rendre grâces à Dieu, et de témoigner son amour à Jésus-Christ, son divin Chef, qui l’a soutenue et consolée durant le cours de ses épreuves.}

'She declares, through the organ of her representatives, that she remains faithful to her principles of faith and freedom on which she was founded.

'\emph{Elle déclare par l’organe de ses représentants qu’elle reste fidèle aux principes de foi et de liberté sur lesquels elle a été fondée.}

'With her fathers and her martyrs in the Confession of Rochelle,\footnote{That is, the Gallican Confession as revised and adopted by the National Synod of La Rochelle, 1571. See § 62.} and with all the Churches of the Reformation in their respective creeds, she proclaims the sovereign authority of the Holy Scriptures in matters of faith, and salvation by faith in Jesus Christ, the only-begotten Son of God, who died for our sins, and was raised again for our justification.

'\emph{Avec ses pères et ses martyrs dans la} Confession de la Rochelle, \emph{avec toutes les Églises de la Réformation dans leurs symboles, elle proclame } l’autorité souveraine des Saintes Écritures en matière de foi, et le salut par la foi en Jésus-Christ, Fils unique de Dieu, mort pour nos offenses et ressuscité pour notre justification.

'She preserves and maintains, as the basis of her teaching, of her worship and her discipline, the grand Christian facts represented in her religious solemnities, and set forth in her liturgies, especially in the Confession of sins, the Apostles' Creed, and in the order for the administration of the Lord's Supper.'

'\emph{Elle conserve donc et elle maintient, à la base de son enseignement, de son culte et de sa discipline, les grands faits chrétiens représentés dans ses solennités religieuses et exprimés dans ses liturgies, notamment dans la Confession des péchés, dans le} Symbole des Apôtres, \emph{et dans la liturgie de la saint Cène.}'

This moderate Confession was adopted by 61 votes against 45, or a majority of only 16 members.\footnote{Two members were absent. The official report says: '\emph{Le nombre des votants est de} 106. \emph{Majorité absolue} 54. \emph{Le dépouillement du scrutin donne} 61 \emph{bulletins blancs,} 45 \emph{bulletins bleus.}'} Among the affirmative votes are those of Babut, Bois, Breyton, Dhombres, Juillerat, and the venerable octogenarian Guizot, whose last public act was a testimony of faith on the floor of this General Synod of the Church of his fathers, declaring before his retirement that the Church must affirm its faith in the supernatural incarnation, the miracles, the resurrection of Christ, or cease to be a Church. The rationalistic minority, including Colani, Coquerel (Athanase and Etienne), Pécaut, Rivet, protested against the adoption of any creed, and asserted the right of each pastor, elder, and private member of the Church to adhere to whatever creed he may think proper. Nevertheless, they expressed their determination to hold on to the National Reformed Church.

The French Government ratified the decision of the Synod (1873). Subscription to its Confession may be hereafter a qualification of electors. The liberal party abstained from participation in the second session of the General Synod held in Nov. and Dec., 1873, and sent in a request to agree to a peaceful separation; but this request was refused.\footnote{The following action was taken by the Synod in reference to the petition of the minority: 'The Assembly, considering that the General Synod is the high court of the Church, and so acknowledged by the State; considering that the decisions arrived at in reference to the Confession of Faith reproduce the doctrines on which the Reformed Church of France was founded, and that, therefore, all who reject them are \emph{ipso facto} without the pale of the Church; considering that none can be constrained to remain in a Church the creed of which he rejects, and from which he wishes to retire—every man having entire liberty to remain or separate himself, according to the dictation of his conscience; considering that the Synod has taken no resolutions to restrict the liberty of any, especially none to prevent the retirement of any pastors and members in order to found another Church, and none to prevent such persons from obtaining the recognition of the State, the advantages of the concordat, and an equitable share of ecclesiastical temporalities; considering, lastly, that it is not the business of the General Synod itself to inaugurate the formation of a new Church, its mission being to construct, and not to rend asunder, passes to the order of the day.'}

Hence the Rationalists, if they have sufficient interest in positive Christianity, will be obliged to secede and organize a new society similar to the Unitarian body in England and the United States.

A separation is preferable to an unnatural alliance at the expense of truth and charity. And it would be all the more honorable if it be done with an equitable division of Church property.

The acts of the General Synod of the National Church had the double effect of virtually excluding the rationalistic party, and of attracting to a closer fellowship the Free Church, which, like the Free Churches in French Switzerland, represents modern evangelical Calvinism, independent of state support and state control.\footnote{The Free Church, or 'Union of the Evangelical Churches in France' (\emph{l’Union des églises évangéliques de France}), to which Pressensé, Fish, and Bersier belong, owes its existence to the rationalism in the National Church which, at the synodical meeting held after the February Revolution of 1848 (without government sanction, and hence without legislative effect), refused to acknowledge the divinity of Christ. This induced Frederick Monod to secede, while his more distinguished and equally conscientious brother Adolph remained, to the benefit of the National body, which since that time has become more orthodox. The Union manifests a good deal of missionary zeal and literary activity, and reacts favorable on the Established Church. Bersier, in his History of the General Synod, expresses himself satisfied with its results (close of Introduction to Vol. I. p. lvii.): '\emph{Nos sympathies personnelles sont avec la droite dans les trois grandes questions que le Synode a eu à résoudre: celle de l’autorité du Synode, celle de la déclaration de foi, celle enfin des conditions de foi et de doctrine auxquelles les pasteurs et les électeurs devront désormais souscrire. Nous estimons que par ces trois votes la majorité a accompli des actes nécessaires, et que si, par un abus de pouvoir que nous ne voulons pas prévoir, le gouvernement refusait de ratifier son œuvre} [the ratification has since been granted], \emph{elle aurait néanmoins posé les fondations futures sur lesquelles, avec ou sans appui de l’État, l’Église réformée devra désormais s’élever.}'}

\xsubsection{The Belgic Confession. A.D. 1561.}

§ 64. The Belgic Confession. A.D. 1561.

Literature.

I. Editions of the Confession.

La Confession de Foi des églises réformées Wallonnes et Flamandes (\emph{Apoc.} ii. 10, 1 \emph{Pierre} iii. 15). \emph{Reimprimée par décision de la Société Évangélique Belge.} Bruxelles, 1850 (Librairie Chrétienne Évangélique, Rue de l’Impératrice, 33). The authentic French text, as revised by the Synod of Dort, with a brief historical and critical introduction.

The Latin text is found in different recensions, in the \emph{Corpus et Syntagma} (1612 and 1654): in the Acts, of the Synod of Dort; the Oxford \emph{Sylloge}; Augusti's \emph{Collect.} (the text of Dort); Niemeyer's \emph{Collect.} (the translation of Hommius, 1518, with various readings).

English translations, likewise differing in minor details, in the \emph{Harmony of Prot. Conf.}; in the Constitution of the Reformed (Dutch) Church in America (very good); and a new one made in 1862 by Owen Jones: \emph{Church of the Living God,} London, 1865, pp. 203–237 (incomplete and inaccurate).

German translation in Beck (Vol. I. pp. 293 sqq.), and Böckel (pp. 480 sqq.).

A Greek translation by Jac. Revius (Pastor of the Church at Deventer): \textgreek{Ἐκκλησιῶν τῆς Βελγικῆς ἐξομολόγησις,} Ultrajecti, 1660; earlier eds. in 1623 and 1653.

Comp. Herzog: art. \emph{Belgische Confession,} in his \emph{Real-Encyklop.} 2d ed. Vol. II. p. 238; M. Goebel: art. \emph{Guido de Brès,} ibid. Vol. V. p. 465.

II. Historical..

H. Grotius:  \emph{Annales et Hist. de rebus Belgicis} (1556–1609). Amstel. 1658.

H. Venema:  \emph{Institutiones historiæ ecclesiæ V. et N.T.} Tom. VII. p. 252 (ad ann. 1563).

J. le Long:  \emph{Kort historisch Verhaal van den oorsprong der Nederlandschen Gereformeerden Kerken ondert Kruys, beneffens alle derselver Leeren Dienst-Boeken.} Amst. 1741.

Gerh. Brandt (Arminian): \emph{Historie der Reformatie in en omtrent de Nederlanden.} Amst. 1671–74, 4 vols. (Also in French: \emph{Histoire de la Réformation des Pays-Bas,} 1726, and in English by Chamberlayne, London, 1720–23, 4 vols.).

Ypey en Dermout:  \emph{Geschiedenissen der Nederlandsche Hervormde Kerk.} Breda, 1819–27, 4 vols.

Van der Kemp:  \emph{De Eere der Nederlandsche Hervormde Kerk.} Rotterd. 1830.

Gachard:  \emph{Correspondance de Guillaume le Taciturne, Prince d’Orange,} 1847–58, 6 vols.

Groen van Prinsterer:  \emph{Archives ou Correspondance inédite de la maison d’Orange-Nassau} (1552–84), 1857–61, 10 vols.; second series (1584–1688), 6 vols. 1857–61.

Wm. H. Prescott:  \emph{History of the Reign of Philip II., King of Spain.} New York, 1855–58, 3 vols.

A. Henne:  \emph{Hist. du règne de Charles V. en Belgique.} Brux. 1858 sqq. 10 Tom.

J. L. Motley:  \emph{The Rise of the Dutch Republic,} London and New York, 1856, 3 vols. By the same: \emph{History of the United Netherlands,} New York, 1861, 4 vols.

M. Koch:  \emph{Untersuch. über die Empörung der Niederlande.} Leipz. 1860.

F. Holzwarth:  \emph{Abfall der Niederlande.} Schaffhausen, 1865–72, 3 vols.

THE REFORMATION IN THE NETHERLANDS.

The Low Countries, conquered from the sea by indomitable energy—the land of Erasmus, of free cities, of inventions, and flourishing commerce —was flooded, through merchants, soldiers, and books, with Protestant ideas from Germany and France, as with waters from the Rhine and the Meuse. Already in 1521 Charles V., who afterwards regretted that he had not burned Luther at Worms, issued from that city an edict for the suppression of heresy in this the most valuable of his inherited dominions. To Belgium belongs the honor of having furnished the first martyrs of evangelical Protestanism in Henry Voes and John Esch, two Augustinian monks, who were burned at the stake in Brussels, July 1, 1523, reciting the Apostles' Creed and singing the \emph{Te Deum,} and who were celebrated by Luther in a stirring hymn.\footnote{See a part of it, in English and German, quoted by Gieseler, Vol. IV. p. 311 (Am. ed.).} This was the fiery signal of a fearful persecution, which reached its height under Philip II. of Spain, and the executor of his bloody designs, the Duke of Alva, but resulted at last in the establishment of national independdence and of the Reformed Church in a large part of the Netherlands. The number of her martyrs exceeds that of any other Protestant Church during the sixteenth century, and perhaps that of the whole primitive Church under the Roman empire.\footnote{Grotius estimates the number of Protestant martyrs in Holland, under one reign, at one hundred thousand. Gibbon (\emph{History of the Decline,} etc., at the close of Ch. XVI.) confidently asserts that 'the number of Protestants who were executed by the Spaniards in a single province and a single reign, far exceeded that of the primitive martyrs in the space of three centuries, and of the Roman empire.' And Motley (\emph{History of the Rise of the Dutch Republic,} Vol. II. p. 504) says of the terrible reign of Alva: 'The barbarities committed amid the sack and ruin of those blazing and starving cities are almost beyond belief; unborn infants were torn from the living bodies of their mothers; women and children were violated by the thousands, and whole populations burned and hacked to pieces by soldiers in every mode which cruelty in its wanton ingenuity could devise.'} During the ever-memorable conflict under William of Orange, who was assassinated by a fanatical papist in 1584, and his second son Maurice—an able military commander and strict Calvinist (d. 1625)—the Bible, with the Belgic Confession and Heidelberg Catechism, was the spiritual guide and comforter of the Protestants, and fortified them against the assaults of the enemy. Calvinism, which fears God and no body else, inspired that heroic courage which triumphed over the political and religious despotism of Spain, and raised Holland to an extraordinary degree of commercial and literary eminence.\footnote{It is strange that Motley, in his great works on the Rise, and the History of the Dutch Republic, ignores the Belgic Confession, and barely mentions the name of Guido de Brès.}

GUIDO DE BRÈS.

The chief author of the Belgic Confession is Guido (or Guy, Wido) de Brès, a noble evangelist and martyr of the Reformed Church of the Netherlands. He was born about 1523 at Mons, in Hennegau, educated in the Roman Church, and by diligent reading of the Scriptures converted to the evangelical faith. Expelled from his country, he sought refuge in London under Edward VI., where he joined the Belgic fugitives, and prepared himself for the ministry. Afterwards he studied at Lausanne, and became a traveling evangelist in Southwestern Belgium and Northern France—from Dieppe to Sedan, from Valenciennes to Antwerp. After the conquest of French Flanders he was, together with a younger missionary from Geneva, Peregrin de la Grange, taken prisoner, put in chains, and hanged on the last day of May, 1567, for disobedience to the commands of the court at Brussels, and especially for the distribution of the holy communion in the Reformed congregations. From prison the youthful martyr wrote letters of comfort to his brethren, his old mother, his wife, and his children, and met his death as if it were a marriage-feast.\footnote{See, on Guy de Brès, the enlarged edition of Crespin's \emph{Histoire des Martyrs,} Genève, 1617, pp. 731–750, and the Brussels edition of the \emph{Conf. de foi,} p. 19.} In his proper home Protestantism was completely suppressed, but in the neighboring countries of Holland and the Lower Rhine it spread and flourished.

THE BELGIC CONFESSION.

The Belgic Confession was prepared in 1561 by Guido de Brès, with the aid of Adrien de Saravia (professor of theology in Leyden, afterwards at Cambridge, where he died, 1613), H. Modetus (for some time chaplain of William of Orange), and G. Wingen, in the French language, to prove the Reformed faith from the Word of God.\footnote{Saravia, in a letter to Uytenbogardus (Apr. 13, 1612), quoted by Niemeyer (Proleg. p. lii.) and Gieseler (\emph{Ch. Hist.} Vol. IV. p. 314, Am. ed.), says: '\emph{Ego me illius confessionis ex primis unum fuisse auctoribus profiteor, sicut et Hermannus Modetus: nescio an plures sint superstites. Illa primo fuit conscripta Gallico sermone a Christi servo et martyre Guidone de Brès, sed antequam ederetur ministris verbi Dei, quos potuit nancisci, illam communicavit: et emendandum si quid displiceret, addendum, detrahendum proposuit, ut unius opus censeri non debeat. Sed nemo eorum, qui manum apposuerunt, umquam cogitavit fidei canonem edere, verum ex canonicis scriptis fidem suam probare.}'} It was revised by Francis Junius, of Bourges (1545–1602)—a student of Calvin, pastor of a Walloon congregation at Antwerp, and afterwards professor of theology at Leyden—who abridged the sixteenth article, and sent a copy to Geneva and other churches for approval. It was probably printed in 1562, or at all events in 1566, and afterwards translated into Dutch, German, and Latin. It was presented to the bigoted Philip II., 1562, in the vain hope of securing toleration, and with an address which breathes the genuine spirit of martyrdom. The petitioners protest against the charge of being rebels, and declare that notwithstanding they number more than a hundred thousand, and are exposed to the most cruel oppression, they obey the Government in all lawful things; but that rather than deny Christ before men they would 'offer their backs to stripes, their tongues to knives, their mouths to gags, and their whole bodies to the fire, well knowing that those who follow Christ must take his cross and deny themselves.'\footnote{The address is given in full by Böckel, 1.c. pp. 480–484.}

The Confession was publicly adopted by a Synod at Antwerp (1566), then at Wesel (1568), more formally by a Synod at Emden (1571)\footnote{The Brussels ed. (p. viii.) says: '\emph{Le} 8 \emph{Octobre, en} 1571, \emph{il fût statué par le premier synode national des Églises wallonnes et flamandes ténu à Embden, que cette Confession serait signée par tous les membres présents au dit synode et par tous ceux qui seraient admis au saint ministère.}'} by a national Synod at Dort (1574), another at Middelburg (1581), and again by the great Synod of Dort, April 29, 1619. But inasmuch as the Arminians had demanded partial changes, and the text had become confused, the Synod of Dort submitted the French, Latin, and Dutch texts to a careful revision. Since that time the Belgic Confession, together with the Heidelberg Catechism, has been the recognized symbol of the Reformed Churches in Holland and Belgium.\footnote{The \emph{Société évangélique} or \emph{Église Chrétienne missionnaire belge} requires from its ministers a qualified subscription to the Belgic Confession with '\emph{une réserve préalable en repoussant ce qui dans la Confession belge regarde l’exercise du pouvoir civil en matière de foi.}'} It is also the doctrinal standard of the Reformed (Dutch) Church in America, which holds to it even more tenaciously than the mother Church in the Netherlands.\footnote{The following formula of subscription is required from ministers of the Dutch Reformed Church in America: 'We, the underwritten, Ministers of the Word of God, residing within the bounds of the Classis of N. N., do hereby sincerely, and in good conscience before the Lord, declare by this our subscription, that we heartily believe, and are persuaded, that all the articles and points of doctrine contained in the [Belgic] Confession and [Heidelberg] Catechism of the Reformed [Dutch] Church, together with the explanation of some points of the aforesaid doctrine made in the National Synod held at Dordrecht, in the year 1619, do fully agree with the Word of God. We promise, therefore, diligently to teach, and faithfully to defend the aforesaid doctrine, without either directly or indirectly contradicting the same by our public preaching or writings. We declare, moreover, that we not only reject all errors that militate against this doctrine, and particularly those which are condemned in the above-mentioned Synod, but that we are disposed to refute and contradict them, and to exert ourselves in keeping the Church pure from such errors. And if hereafter any difficulties or different sentiments respecting the aforesaid doctrine should arise in our minds, we promise that we will neither publicly nor privately propose, teach, or defend the same, either by preaching or by writing, until we have first revealed such sentiment to the Consistory, Classis, or Synod, that the same may be there examined,' etc.}

CONTENTS.

The Belgic Confession contains thirty-seven Articles, and follows the order of the Gallican Confession, but is less polemical and more full and elaborate, especially on the Trinity, the Incarnation, the Church, and the Sacraments.\footnote{Ebrard (\emph{Handbuch der Kirchen- und Dogmengesch.} Vol. III. p. 319) says that besides the Gallican Confession as the basis, use was made also of the Friesian Confession of Utenhoven, which the English exiles brought with them to Emden, and of the Catechism of Laski.} It is, upon the whole, the best symbolical statement of the Calvinistic system of doctrine, with the exception of the Westminster Confession.

THE TEXT.

The text has undergone several modifications as regards the wording and length, but not as regards the doctrine.

The French text must be considered as the original.\footnote{It is entitled, '\emph{Confession de Foy faicte d’un commun accord pour les fidèles qui conversent ès Pays-Bas, lesquels désirent vivre selon la pureté de l’Évangile de nostre Seigneur Jésus-Christ.}' This title is followed by two mottoes—the one from Apoc. ii. 10: '\emph{Sois fidèle jusques à la mort et je te donneray la couronne de vie;}' the other from 1 Pet. iii. 15: '\emph{Soyez tousjours appareillez à respondre à chacun qui vous demande raison de l’espérance qui est en vous.}' On the second leaf there is over the head of the first article the brief title, '\emph{Confession vrayement Chrétienne contenant le sommaire de la doctrine de Dieu et salut éternel de l’âme.}'} Of the first edition of 1561 or 1562 no copies are known. The Synod of Antwerp, in Sept., 1580, ordered a precise parchment copy of the revised text (of Junius) to be made for its archives, which copy had to be signed by every new minister. This manuscript has always been regarded in the Belgic churches as the authentic document.\footnote{The Brussels ed. says (p. 39): '\emph{C’est probablement d’après la copie de Junius que cette Confession a été imprimée dans le livre des Martyrs de Crespin. Le text de Crespin ne diffère pas de celui du manuscrit authentique.}'} The Synod of Dort ordered a new revision, with a view to bring the Latin, French, and Dutch texts into harmony on the basis of the manuscript copy of 1580. The Leyden edition of 1669 gives in two parallel columns the original text and the revised text of Dort. A Rotterdam edition of the Psalter, 1787, carefully reprints the original text in the old spelling from the manuscript, with the changes of Dort in notes. The Brussels edition of 1850 presents the ancient text of 1580, as revised at Dort, in modern French.\footnote{This careful edition, issued by the Evangelical Society of Belgium, is reproduced in the third volume of this work, together with the English version now used by the Dutch Reformed Church in America. Both agree, sentence for sentence.}

Next in authority is the Latin text, but of this there are likewise several recensions, a shorter and a larger. The first Latin translation was made from the revised French copy of Francis Janius, probably by Beza, or under his direction, for the \emph{Harmonia Confessionum,} Geneva, 1581 (distributed under different heads, with the other Confessions).\footnote{See \emph{Note critique} at the close of the Brussels edition, p. 39: '\emph{Junius envoya une copie de cette révision à Genève. Theodore de Beza la fit imprimer} [\emph{in French?}]. \emph{C'est lui, sans doute, qui la traduisit en latin, comme elle se trouve dans} ``\emph{l’Harmonia Confessionum,}'' \emph{Genevæ,} 1581.' That this was the first Latin translation is stated in the \emph{Harmonia,} p. 3: '\emph{Belgica, Gallice omnium Belgicarum Ecclesiarum nomine anno} 1566 \emph{edita, ac demum anno} 1579 [1571?] \emph{in publica Belgii Synodo repetita et confirmata, Belgiceque versa. Nunc denique a nobis etiam Latine expressa.}'} The same passed into the first edition of the \emph{Corpus et Syntagma Confessionum,} Geneva, 1612. Another translation was prepared, 1618, for the use of the Synod of Dort, by Festus Hommius, pastor in Leyden, and one of the scribes of that Synod.\footnote{'\emph{Confessio ecclesiarum reformatarum in Belgio. . . . in usum futuræ synodi nationalis latine edidit et collegit Festus Hommius.}' Ludg. Batav. 1618. Niemeyer (pp. 360 sqq.) gives this translation, which more nearly agrees with the older version, and he adds some readings from the first edition of the \emph{Corpus et Syntagma.}} This text was revised in the following year by that Synod, and thus approved and incorporated with its acts in the 146th session.\footnote{See the extracts from the Acts of the 144th Session, April 29, 1619, in Niemeyer, p. lv.} The revision of Dort was reproduced in the second edition of the \emph{Corpus et Syntagma Conf.,} 1654.\footnote{Under the title \emph{Ecclesiarum Belgicarum Christiana atque Orthodoxa Confessio, summam doctrinæ de Deo et æterna animarum salute complectens, prout in Synodo Dortrechtana fuit recognita et approbata.} The articles are numbered, but have no titles. The difference between this and the first Latin translation may be judged from the following specimen:    Harmonia Confessionum, 1581 (p. 36).  Corpus et Syntagma Confessionum, ed. II., 1654 (p. 129).    \par Art. I. \emph{Corde credimus, et ore confitemur, unicam esse et simplicem essentiam spiritualem, quam Deum vocamus, æternum, incomprehensibilem, inconspicuum, immutabilem, infinitum, qui totus est sapiens, fonsque omnium bonorum uberrimus.} \par Art. I. \emph{Corde credimus, et ore confitemur } omnes, \emph{ unicam esse et simplicem essentiam spiritualem, quam Deum vocamus, } eumque  \emph{æternum, incomprehensibilem, invisibilem, infinitum, } omnipotentem, summe sapientem, justum et bonum,  \emph{omniumque bonorum fontem uberrimum.}} The excellent English version in use in the Reformed Dutch Church of America is made from the Latin text of the Synod of Dort.

\xsubsection{The Arminian Controversy and the Synod of Dort, A.D. 1604–1619.}

§ 65. The Arminian Controversy. A.D. 1604–1619.

Literature.

I. Arminian Sources.

\emph{Scripta adversaria } Collationis Hagiensis. In Dutch, Gravenhage, 1612; in Latin, by Petrus Bertius, Leyden, 1616. This contains the authentic text of the Remonstrance.

Remonstrantia, or the Five Articles of 1610. A German translation in Böckel's \emph{Evang. Reform. Bekenntniss-Schriften.} Leipzig, 1847, pp. 545–553.

Simon Episcopius (Prof. at Leyden, 1612; expelled by the Synod of Dort, 1618; Prof. at the Remonstrant Seminary, 1634; d. 1643): \emph{Confessio seu Declaratio Pastorum qui Remonstrantes vocantur,} etc. Harderw. 1621 in Dutch, 1622 in Latin (German transl. in Böckel, l.c. pp. 572–640). Also his \emph{Apologia pro Confessione Remonstr.,} 1629. Both are included in the works of Episcopius, 2d ed. London, 1678, Vol. II. Part II. pp. 69 sqq.; 95 sqq.

Acta et Scripta Synodalia Dordracena  \emph{ministrorum } Remonstrantium  \emph{in fœderato Belgio. } 2 Cor. xiii. 8. Harderwiici, 1620. This volume (a copy of which is in the Union Theol. Seminary Library) contains the official acts and dogmatic writings of the Remonstrants in explanation and defense of their five articles against the decisions of the Synod of Dort, including a lengthy exposition of the ninth chapter of Romans and other Scripture passages quoted against them.

Jac. Arminius (1560–1609): \emph{Disputationes publicæ et privatæ.} Ludg. Bat. 1614, 2d ed. (with the \emph{Oratio Petri Bertii de vita et obitu Arminii.}). Armin.  \emph{Opera,} Lugd. Bat. 1629; and other editions. English translation of \emph{The Works of } James Arminius, by \emph{James Nichols and William Nichols.} London, 1825, 1828, and 1875, 3 vols.

Also the writings of Episcopius (d. 1643); Grotius (d. 1645); Limborch (d. 1714); Clericus (d. 1736); Wetstein (d. 1754), and other distinguished Arminian scholars. Comp. A. van Cattenburgh:  \emph{Bibliotheca Scriptorum Remonstrantium.} Amst. 1728.

II. Anti-Arminian or Calvinistic Sources.

The Acts and Proceedings of the National Synod of Dort:  \emph{Acta Synodi Nationalis, in nomine Domini nostri Jesu Christi, autoritate ordinum generalium Fœderati Belgii provinciarum, Dortrechti habitæ anno} 1818 \emph{et} 1619. \emph{Accedunt plenissima de quinque articulis theologorum judicia.} Dord. 1620, 4to. (The \emph{judicia theologorum} are omitted in the Elzevir folio ed. of the same date.)

\emph{The Suffrage of the Divines of Great Britain concerning the Articles of the Synod of Dort, signed by them in} 1619 [? Lond. 1624].

Reports of Breitinger, the Hessian, and other foreign delegates.

III. Historical and Controversial.

P. Molinæus (Calvinist): \emph{Anatome Arminianismi.} Leyden, 1619, etc.

N. Vedel (Calv.): \emph{Arcana Arminianismi.} Leyden, 1632–34, 4 Parts, 4to.

Peltius:  \emph{Harmonia Remonstrantium et Socinianorum.} Ludg. 1633.

Byssen:  \emph{De prædestinatione contra Remonstrantes et Jesuitas.} Gorchum, 1660.

Sam. Rhetorfort:  \emph{Examen Arminianismi.} Utrecht, 1668.

Janus Uytenbogaert (Arminian): \emph{Kerckelijcke Historie,} etc. Rotterdam, 1647.

Jac. Triglandius (Calvinist): \emph{Kerekelijcke Geschiedenissen van de vereen. Nederlanden.} Leyden, 1650.

Jo. Halesii  \emph{Historia Concilii Dordraceni; J. L. Moshemius vertit, variis observationibus et vita Halesii auxit.} Hamburg, 1724. John Hales (1584–1656), Canon of Windsor—called 'the Ever-memorable'—attended the Synod of Dort, by which he became a convert to Arminianism, and wrote \emph{Golden Remains; Letters from the Synod of Dort; Acta Synodi Dordr.; Sententia Arminii;} see \emph{Works,} 1765, 3 vols.

Peter Heylin (a friend of Laud and Arminian, d. 1662): \emph{Historia Quinquarticularis; or, a Declaration of the Judgment of the Western Churches, and more particularly of the Church of England, in the Five Controverted Points, reproached in these last times by the name of Arminianism.} London, 1660, in 3 Parts.

Gerhard Brandt (Remonstrant preacher at Amsterdam, d. 1685): \emph{Historie der Reformatie} (\emph{History of the Reformation in and about the Low Countries, from the Eighth Century down to the Synod of Dort}), Amst. 1677–1704, 4 vols. Very full on the Remonstrant controversy. An English translation, by \emph{Chamberlayne,} London, 1720–23, 4 vols. fol. (The last volume gives the history from 1600 to 1623.) Also in French, 1726.

Zeltner (d. 1738): \emph{Breviarium controversiarum cum Remonstrantibus agitatarum.} Norimb. and Altdorf, 1719.

Jac. Regenboog:  \emph{Hist. der Remonstranten,} in Dutch, Amsterd. 1774 sqq., 3 vols.; in German, Lemgo, 1741–84.

G. S. Franke:  \emph{Historia dogmatum Arminianorum.} Kiel, 1814.

Thomas Scott:  \emph{The Articles of the Synod of Dort; with a History of Events which made way for that Synod,} etc. London, 1818. (Calvinistic.)

James Nichols (Arminian): \emph{Calvinism and Arminianism compared in their Principles and Tendency.} Lond. 1824, 2 vols. (An ill-digested mass of materials.)

M. Graf:  \emph{Beiträge zur Geschichte der Synode von Dordrecht.} Basle, 1825.

D. de Bray:  \emph{L’histoire de l’Église Arminienne.} Strasburg, 1835.

Joannes Tideman (Remonstrant preacher at Rotterdam): \emph{De Remonstrantie en het Remonstrantisme. Historisch onderzoek.} Te Haarlem, 1851 (pp. 131).

H. Heppe (Melanchthonian): \emph{Historia Synodi Nat. Dordr.} in Niedner's \emph{Zeitschrift für hist. Theol.,} 1853, pp. 227–327. Contains the Report of the Hessian deputies to Landgrave Moritz, with Introduction and Notes. The same: Art. \emph{Dortrecht} in Herzog's \emph{Real-Encykl.} Vol. III. p. 486.

Alex. Schweizer:  \emph{Centraldogmen.} Zurich, Vol. II. (1856) pp. 31–201.

G. Frank:  \emph{Geschichte der Protest. Theol.} Leipz. 1862, Vol. I. pp. 403 sqq.

M. Schneckenburger (independent, d. 1848): \emph{Vorlesungen über die Lehrbegriffe der kleineren protest. Kirchenparteien,} ed. by Hundeshagen. Frankf. a. M. 1863, pp. 5–26.

William Cunningham (Calvinist): \emph{Historical Theology.} Edinb. 1864, Vol. II. ch. xxv. pp. 371–513.

E. Böhl (Calvinist): \emph{Blätter der Erinnerung an die Dordrechter Synode,} 250 \emph{Jahre nach ihrem Zusammentritt allen Freunden der reform. Lehre gewidmet.} Detmold, 1868 (41 pp.).

John L. Motley:  \emph{The Life and Death of John of Barneveld, Advocate of Holland.} N.Y. 1874, 2 vols. chs. viii. and xiv. Motley gives the political history of the period, but barely touches on the Synod of Dort, and with strong antipathy to Calvinism.

Comp. also Whedon (Methodist), art. \emph{Arminianism,} and A. A. Hodge (Presbyterian), \emph{Calvinism,} both in Johnson's \emph{Cyclop.} Vol. I. (1874), representing both sides. Also art. \emph{Arminianism,} in M'Clintock and Strong's \emph{Cyclop.} Vol. I. p. 412 (Methodist).

The Arminian controversy is the most important which took place within the Reformed Church. It corresponds to the Pelagian and the Jansenist controversies in the Catholic Church. It involves the problem of ages, which again and again has baffled the ken of theologians and philosophers, and will do so to the end of time: the relation of divine sovereignty and human responsibility. It started with the doctrine of predestination, and turned round five articles or 'knotty points' of Calvinism; hence the term 'quinquarticular' controversy. Calvinism represented the consistent, logical, conservative orthodoxy; Arminianism an elastic, progressive, changing liberalism. Calvinism triumphed in the Synod of Dort, and excluded Arminianism. So, in the preceding generation, strict Lutheranism had triumphed over Melanchthonianism in the Formula of Concord. But in both Churches the spirit of the conquered party rose again from time to time within the ranks of orthodoxy, to exert its moderating and liberalizing influence or to open new issues in the progressive march of theological science.

ORIGIN AND PROGRESS OF ARMINIANISM TILL 1618.

The Arminian controversy arose in Holland towards the close of the heroic conflict with foreign political and ecclesiastical despotism. This very contest of forty-five years' duration, so full of trials and afflictions, stimulated the intellectual and moral energies of an honest, earnest, freedom-loving, and tenacious people, and made the Protestant part of the Netherlands the first country in Christendom for industry, commerce, education, and culture. The Universities of Leyden, founded in 1575, as the city's reward for its heroic resistance to Spain, Franecker (1585), Groningen (1612), Utrecht (1636), and Harderwyk (1648) soon excelled older schools of learning. The general prosperity of the United Provinces excited the admiration of the foreign delegates to the Synod of Dort, where they found clean and stately mansions, generous hospitality, and every comfort and luxury which commerce could bring from all parts of the earth. This was the soil on which the Calvinistic system was brought to its severest test. The controversy was purely theological in its nature, but owing to the intimate connection of Church and State it became inevitably entangled in political issues, and shook the whole country. The Reformed Churches in France, Switzerland, Germany, England, and Scotland took a deep interest in it, and sided, upon the whole, with the Calvinistic party; while the Lutheran Church sympathized to some extent with the Arminian.

The founder of Arminianism, from whom it derives its name, is James Arminius (1560–1609).\footnote{His Dutch name is Jacob van Hermanns or Hermanson, Harmensen.} He studied under Beza at Geneva, was elected minister at Amsterdam (1588), and then professor of theology at Leyden (1603), as successor of Francis Junius, who had taken part in the revision of the Belgic Confession. He was at first a strict Calvinist, but while engaged in investigating and defending the Calvinistic doctrines against the writings of Dirik Volckaerts zoon Koornheert,\footnote{Koornheert was Secretarius at Haarlem, and a forerunner of the Remonstrants (d. 1590). He attacked the doctrine of Calvin and Beza on predestination and the punishment of heretics (1578), wrote against the Heidelberg Catechism (1583), and advocated toleration and a reduction of the number of articles of faith. His works were published at Amsterdam, 1630. See Bayle, art. \emph{Koornheert,} and Schweizer, Vol. II. p. 40. Another forerunner of Arminianism was Caspar Koolhaas, preacher in Leyden, who was protected by the civil magistrate, but excommunicated by a provincial Synod at Haarlem, 1582. It should be remembered also that Erasmus, the advocate of free-will, against Luther, was held in high esteem in his native country, and that the views of Castellio, Bolsec, and Huber had made some impression.} at the request of the magistrate of Amsterdam, he found the arguments of the opponent stronger than his own convictions, and became a convert to the doctrine of universal grace and of the freedom of will. He saw in the seventh chapter of Romans the description of a legalistic conflict of the awakened but unregenerate man, while Augustine and the Reformers referred it to the regenerate. He denied the decree of reprobation, and moderated the doctrine of original sin. He advocated a revision of the Belgic Confession and Heidelberg Catechism. He came into open conflict with his supralapsarian colleague, Francis Gomar (1563–1645), who had conferred on him the degree of doctor of divinity, but now became his chief antagonist. Hence the strict Calvinists were called 'Gomarists.' The controversy soon spread over all Holland. Arminius applied to the Government to convoke a synod (appealing, like the Donatists, to the very power which afterwards condemned him), but died of a painful disorder before it convened.\footnote{In the same year (1609) the Pilgrim Fathers of New England arrived in Leyden, where they enjoyed religions freedom till their departure for America (1620). Arminius was born in the same year in which Melanchthon died (1560).} He was a learned and able divine; and during the controversy which embittered his life he showed a meek, Christian spirit. 'Condemned by others,' said Grotius, 'he condemned none.' His views on anthropology and soteriology approached those of the Melanchthonian school in the Lutheran Church, but the tendency of his theology was towards a latitudinarian liberalism, which developed itself in his followers.\footnote{Caspar Brandt: \emph{Historia vitæ J. Arminii,} ed. by Gerhard Brandt (son of the author), with additions by Mosheim, 1725; Engl. transl. by Guthrie, Lond. 1854. Bangs's \emph{Life of Arminius,} N. York, 1843. Mosheim calls him 'a man whom even his enemies commend for his ingenuity, acuteness, and piety.' His motto was, 'A good conscience is a paradise.' In his testament (see extract in Gieseler, Vol. IV. p. 508, note 7), he affirms that he diligently labored to teach nothing but what he could prove from the Scriptures, and what tended to edification and peace among Christians, excepting popery, 'with which,' he says, 'there can be no unity of faith, no bond of piety and peace.' Grotius was much milder towards the Catholics.}

After his death the learned Simon Episcopius (Bisschop, 1583–1644), his successor in the chair of theology at Leyden, afterwards professor in the Arminian College at Amsterdam,\footnote{Limborch: \emph{Vita Episcopii.} Amst. 1701.} and the eloquent Janus Uytenbogaert (1557–1644), preacher at the Hague, and for some time chaplain of Prince Maurice, became the theological leaders of the Arminian party. The great statesman, John van Olden Barneveldt (1549–1619), Advocate-General of Holland and Friesland, and Hugo Grotius (1583–1645), the most comprehensive scholar of his age, equally distinguished as statesman, jurist, theologian, and exegete, sympathized with the Arminians, gave them the weight of their powerful influence, and advocated peace and toleration; but they favored a republican confederacy of States rather than a federal State tending to monarchy, against the ambitious designs of Maurice, the Stadtholder and military leader of the Republic, who wished to consolidate his power, and by concluding a truce with Spain (1609) they incurred the suspicion of disloyalty.\footnote{On Barneveldt, see the work of Motley; on Hugo Grotius, the monograph of Luden, Berlin, 1806.} The Calvinists were the national and popular party, and embraced the great majority of the clergy. They stood on the solid basis of the recognized standards of doctrine. At the same time they advocated the independent action of the Church against the latitudinarian Erastianism of their opponents.

The Arminians formularized their creed in Five Articles (drawn up by Uytenbogaert), and laid them before the representatives of Holland and West Friesland in 1610 under the name of \emph{Remonstrance,} signed by forty-six ministers. The Calvinists issued a \emph{Counter-Remonstrance.} Hence the party names \emph{Remonstrants} (Protestants against Calvinism), and \emph{Counter-Remonstrants} (Calvinists, or Gomarists). A Conference was held between the two parties at the Hague (\emph{Collatio Hagiensis}) in 1611, but without leading to an agreement. A discussion at Delft, 1613, and the edict of the States of Holland in favor of peace, 1614, prepared by Grotius, had no better result.

THE SYNOD OF DORT.

At last, after a great deal of controversy and complicated preparations, the National Synod of Dort\footnote{In Dutch, Dordrecht or Dordtrecht; in Latin, Dordracum—an old fortified town in which the independence of the United Provinces was declared in 1572.} was convened by the States-General, Nov. 13, 1618, and lasted till May 9, 1619. It consisted of eighty-four members and eighteen secular commissioners. Of these fifty-eight were Dutchmen, the rest foreigners. The foreign Reformed Churches were invited to send at least three or four divines each, with the right to vote.

James I. of England sent Drs. George Carleton, Bishop of Llandaff (afterwards of Chichester); John Davenant, Bishop of Salisbury; Samuel Ward, Professor of Cambridge; the celebrated Joseph Hall, afterwards Bishop of Exeter and Norwich (who, however, had to leave before the close, and was replaced by Thomas Goad), and Walter Balcanquall, a Scotchman, and chaplain of the King. The Palatinate was represented by Drs. Abraham Scultetus, Henry Alting, Professors at Heidelberg, and Paulus Tossanus; Hesse, by Drs. George Cruciger, Paul Stein, Daniel Angelocrator, and Rudolph Goclenius; Switzerland, by Dr. John Jacob Breitinger, Antistes of Zurich, Sebastian Beck and Wolfgang Meyer of Basle, Marcus Rutimeyer of Berne, John Conrad Koch of Schaffhausen, John Deodatus and Theodor Tronchin of Geneva; Bremen, by Matthias Martinius, Henry Isselburg, and Ludwig Crocius. The Elector of Brandenburg chose delegates, but excused their absence on account of age. The national Synod of France elected four delegates—among them the celebrated theologians Chamier and Du Moulin—but the King forbade them to leave the country. King James instructed the English delegates to 'mitigate the heat on both sides,' and to advise the Dutch ministers\footnote{See the nine instructions of James to the delegates, in Fuller, \emph{Ch. H. of Brit.} Vol. V. p. 462.} 'not to deliver in the pulpit to the people those things for ordinary doctrines which are the highest points of schools and not fit for vulgar capacity, but disputable on both sides.'

The Synod was opened and closed with great solemnity, and held one hundred and fifty-four formal sessions, besides a larger number of conferences.\footnote{The Dutch delegates held twenty-two additional sessions on Church government.} The expenses were borne by the States-General on a very liberal scale, and exceeded 100,000 guilders.\footnote{The five English delegates were allowed the largest sum, viz., ten pounds sterling per day—more than any other foreign divines.—Fuller, l.c. p. 465.} The sessions were public, and crowded by spectators. John Bogerman, pastor at Leuwarden, was elected President; Festus Hommius, pastor in Leyden, first Secretary—both strict Calvinists. The former had translated Beza's tract on the punishment of heretics into Dutch; the latter prepared a new Latin version of the Belgic Confession. The whole Dutch delegation was orthodox. Only three delegates from the provincial Synod of Utrecht were Remonstrants, but these had to yield their seats to the three orthodox members elected by the minority in that province. Gomarus represented supralapsarian Calvinism, but the great majority were infralapsarians or sublapsarians.

Thus the fate of the Arminians was decided beforehand. Episcopius and his friends—thirteen in all—were summoned before the Synod simply as defendants, and protested against unconditional submission.

Orthodox Calvinism achieved a complete triumph. The Five Articles of the Remonstrance were unanimously rejected, and five Calvinistic canons adopted, together with the Belgic Confession and the Heidelberg Catechism. A thorough and most excellent revision of the Dutch Bible from the Hebrew and Greek was also ordered, besides other decisions which lie beyond our purpose.

The victory of orthodoxy was obscured by the succeeding deposition of about two hundred Arminian clergymen, and by the preceding though independent arrest of the political leaders of the Remonstrants, at the instigation of Maurice. Grotius was condemned by the States-General to perpetual imprisonment, but escaped through the ingenuity of his wife (1621). Van Olden Barneveldt was unjustly condemned to death for alleged high-treason, and beheaded at the Hague (May 14, 1619). His sons took revenge in a fruitless attempt against the life of Prince Maurice.

The canons of Dort were fully indorsed by the Reformed Church in France, and made binding upon the ministers at the Twenty-third National Synod at Alais, Oct. 1, 1620, and again at the Twenty-fourth Synod at Charenton, Sept., 1623. In other Reformed Churches they were received with respect, but not clothed with proper symbolical authority. In England there arose considerable opposition.\footnote{See Hardwick's \emph{History of the Thirty-nine Articles,} ch. ix., and Heylin's \emph{Historia Quinquarticularis.}} The only Church outside of Holland where they are still recognized as a public standard of doctrine is the Reformed Dutch Church in America.

The Synod of Dort is the only Synod of a quasi-œcumenical character in the history of the Reformed Churches. In this respect it is even more important than the Westminster Assembly of Divines, which was confined to England and Scotland, although it produced superior doctrinal standards. The judgments of the Synod of Dort differ according to the doctrinal stand-point. It was undoubtedly an imposing assembly; and, for learning and piety, as respectable as any ever held since the days of the Apostles. Breitinger, a great light of the Swiss Churches, was astonished at the amount of knowledge and talent displayed by the Dutch delegates, and says that if ever the Holy Spirit were present in a Council, he was present at Dort. Scultetus, of the Palatinate, thanked God that he was a member of that Synod, and placed it high above similar assemblies. Meyer, a delegate of Basle, whenever afterwards he spoke of this Synod, uncovered his head and exclaimed '\emph{Sacrosancta Synodus!} Even Paolo Sarpi, the liberal Catholic historian, in a letter to Heinsius, spoke very highly of it. A century later, the celebrated Dutch divine, Campegius Vitringa, said: 'So much learning was never before assembled in one place, not even at Trent.'\footnote{Schweizer, Vol. II. pp. 26, 143 sq.; also, Graf, and Böhl, 1.c.}

On the other hand, the Remonstrants, who had no fair hearing, abhorred the Synod of Dort on account of its Calvinism and intolerance. The Lutherans were averse to it under the false impression that the condemnation of Arminianism was aimed at their own creed. Some secular historians denounce it as a Calvinistic tribunal of inquisition.\footnote{Motley (\emph{Life and Death of John of Barneveld,} Vol. II. p. 309) caricatures the Synod of Dort in a manner unworthy of an impartial historian. 'It was settled,' he said, 'that one portion of the Netherlands and of the rest of the human race had been expressly \emph{created} by the Deity to be forever damned, and another portion to be eternally blessed. . . . On the 30th April and 1st May the Netherland Confession and the Heidelberg Catechism were declared \emph{infallible.}'}

The Canons of Dort have for Calvinism the same significance which the Formula of Concord has for Lutheranism. Both betray a very high order of theological ability and care. Both are consistent and necessary developments. Both exerted a powerful conservative influence on these Churches. Both prepared the way for a dry scholasticism which runs into subtle abstractions, and resolves the living soul of divinity into a skeleton of formulas and distinctions. Both consolidated orthodoxy at the expense of freedom, sanctioned a narrow confessionalism, and widened the breach between the two branches of the Reformation.

ARMINIANISM AFTER THE SYNOD OF DORT.

The banishment of the Arminians was of short duration. After the death of Prince Maurice of Nassau (1625), and under the reign of his milder brother and successor, Frederick Henry, they were allowed to return and to establish churches and schools in every town of Holland, which became more and more a land of religious toleration and liberty. In this respect their principles triumphed over their opponents.\footnote{Hugo Grotius carried the principle of toleration so far that it was said Socinus, Luther, Calvin, Arminius, the Pope, and Arius contended for his religion as seven cities for the birth of the divine Homer. See the verse of Menage, quoted by G. Frank, \emph{Geschichte der Protest. Theologie,} Vol. I. p. 410.} They founded a famous Theological College at Amsterdam (1630), which exists to this day, and has recently been removed to Leyden.

Peace was not so favorable to their growth as controversy. They gradually diminished in number, and are now a very small sect in Holland, almost confined to Rotterdam and Amsterdam.

But their literary and religious influence has gone far beyond their organization. Their eminent scholars, Hugo Grotius, Episcopius, Limborch, Curcellæus, Clericus (Le Clerc), and Wetstein, have enriched exegetical and critical learning, and liberalized theological opinions, especially on religious toleration and the salvation of unbaptized infants. Arminianism, in some of its advocates, had a leaning towards Socinianism, and prepared the way for Rationalism, which prevailed to a great extent in the Established Churches of Holland, Geneva, and Germany from the end of the last century till the recent reaction in favor of orthodox Calvinism and Lutheranism. But many Arminians adhered to the original position of a moderated semi-Pelagianism.

The distinctive Arminian doctrines of sin and grace, free-will and predestination, have been extensively adopted in the Episcopal Church since the reign of Charles I., and in the last century by the Methodists of Great Britain and America, and thereby have attained a larger territory and influence than they ever had in the land of their birth.\footnote{The Wesleys were Arminians, while Whitefield was a Calvinist. They separated on the question of predestination.} Methodism holds to the essential doctrines of the Reformation, but also to the five points of Arminianism, with some important evangelical modifications.

\xsubsection{The Remonstrance.}

§ 66. The Remonstrance.

The Arminian or quinquarticular controversy started with opposition to the doctrine of absolute decrees, and moved in the sphere of anthropology and soteriology. The peculiar tenets are contained in the five points or articles which the Arminians in their 'Remonstrance' laid before the estates of Holland in 1610. They relate to predestination, the extent of the atonement, the nature of faith, the resistibility of grace, and the perseverance of saints.

The Remonstrance is first negative, and then positive. It rejects five Calvinistic propositions, and then asserts the five Arminian propositions. The doctrines rejected are thus stated:

1. That God has, before the fall, and even before the creation of man, by an unchangeable decree, foreordained some to eternal life and others to eternal damnation, without any regard to righteousness or sin, to obedience or disobedience, and simply because it so pleased him, in order to show the glory of his righteousness to the one class and his mercy to the other. (This is the supralapsarian view.)

2. That God, in view of the fall, and in just condemnation of our first parents and their posterity, ordained to exempt a part of mankind from the consequences of the fall, and to save them by his free grace, but to leave the rest, without regard to age or moral condition, to their condemnation, for the glory of his righteousness. (The sublapsarian view.)

3. That Christ died, not for all men, but only for the elect.

4. That the Holy Spirit works in the elect by irresistible grace, so that they \emph{must} be converted and be saved; while the grace necessary and sufficient for conversion, faith, and salvation is withheld from the rest, although they are externally called and invited by the revealed will of God.

5. That those who have received this irresistible grace can never totally and finally lose it, but are guided and preserved by the same grace to the end.

These doctrines, the Remonstrants declare, are not contained in the Word of God nor in the Heidelberg Catechism, and are unedifying, yea dangerous, and should not be preached to Christian people.

Then the Remonstrance sets forth the five positive articles as follows:

ARTICLE FIRST.

\emph{Conditional Predestination.}—God has immutably decreed, from eternity, to save those men who, by the grace of the Holy Spirit, believe in Jesus Christ, and by the same grace persevere in the obedience of faith to the end; and, on the other hand, to condemn the unbelievers and unconverted (John iii. 36).

Election and condemnation are thus conditioned by foreknowledge, and made dependent on the foreseen faith or unbelief of men.

SECOND ARTICLE.

\emph{Universal Atonement.}—Christ, the Saviour of the world, died for all men and for every man, and his grace is extended to all. His atoning sacrifice is in and of itself sufficient for the redemption of the whole world, and is intended for all by God the Father. But its inherent sufficiency does not necessarily imply its actual efficiency. The grace of God may be resisted, and only those who accept it by faith are actually saved. He who is lost, is lost by his own guilt (John iii. 16; 1 John ii. 2).

The Arminians agree with the orthodox in holding the doctrine of a vicarious or expiatory atonement, in opposition to the Socinians; but they soften it down, and represent its direct effect to be to enable God, consistently with his justice and veracity, to enter into a new covenant with men, under which pardon is conveyed to all men on condition of repentance and faith. The immediate effect of Christ's death was not the salvation, but only the salvability of sinners by the removal of the legal obstacles, and opening the door for pardon and reconciliation. They reject the doctrine of a \emph{limited} atonement, which is connected with the supralapsarian view of predestination, but is disowned by moderate Calvinists, who differ from the Arminians in all other points. Calvin himself says that Christ died \emph{sufficienter pro omnibus, efficaciter pro electis.}

THIRD ARTICLE.

\emph{Saving Faith.}—Man in his fallen state is unable to accomplish any thing really and truly good, and therefore also unable to attain to saving faith, unless he be regenerated and renewed by God in Christ through the Holy Spirit (John xv. 5).

FOURTH ARTICLE.

\emph{Resistible Grace.}—Grace is the beginning, continuation, and end of our spiritual life, so that man can neither think nor do any good or resist sin without prevening, co-operating, and assisting grace. But as for the manner of co-operation, this grace is not irresistible, for many resist the Holy Ghost (Acts vii.).

FIFTH ARTICLE.

\emph{The Uncertainty of Perseverance.}—Although grace is sufficient and abundant to preserve the faithful through all trials and temptations for life everlasting, it has not yet been proved from the Scriptures that grace, once given, can never be lost.

On this point the disciples of Arminius went further, and taught the possibility of a total and final fall of believers from grace. They appealed to such passages where believers are warned against this very danger, and to such examples as Solomon and Judas. They moreover denied, with the Roman Catholics, that any body can have a certainty of salvation except by special revelation.

These five points the Remonstrants declare to be in harmony with the Word of God, edifying and, as far as they go, sufficient for salvation. They protest against the charge of changing the Christian Reformed religion, and claim toleration and legal protection for their doctrine.

\xsubsection{The Canons of Dort.}

§ 67. The Canons of Dort.

The Canons of Dort are likewise confined to five points or 'Heads of Doctrine,' and exhibit what is technically called the Calvinistic system—first positively, then negatively, in the rejection of the Arminian errors.\footnote{The term '\emph{rejectio errorum,}' instead of the condemnation and anathemas of the Greek and Roman Churches in dealing with heresies, indicates that Protestant orthodoxy is more liberal and charitable than the Catholic.} Each Head of Doctrine (subdivided into Articles) is subscribed by the Dutch and foreign delegates.

FIRST HEAD OF DOCTRINE.

\emph{Of Divine Predestination.}—Since all men sinned in Adam and lie under the curse [according to the Augustinian system held by all the Reformers], God would have done no injustice if he had left them to their merited punishment; but in his infinite mercy he provided a salvation through the gospel of Christ, that those who believe in him may not perish, but have eternal life. That some receive the gift of faith from God and others not, proceeds from God's eternal decree of election and reprobation.

Election is the unchangeable purpose of God whereby, before the foundation of the world, he has, out of mere grace, according to the sovereign good pleasure of his own will, chosen from the whole human race, which has fallen through their own fault from their primitive state of rectitude into sin and destruction, a certain number of persons to redemption in Christ, whom he from eternity appointed the Mediator and Head of the elect, and the foundation of salvation. These elect, though neither better nor more deserving than others, God has decreed to give to Christ to be saved by him, and bestow upon them true faith, conversion, justification and sanctification, perseverance to the end, and final glory (Eph. i. 4, 5, 6; Rom. viii. 30).

Election is absolute and unconditional. It is not founded upon foreseen faith and holiness, as the prerequisite condition on which it depended; on the contrary, it is the fountain of faith, holiness, and eternal life itself. God has chosen us, not \emph{because} we are holy, but to the \emph{end} that we should be holy (Eph. i. 4; Rom. ix. 11–13; Acts xiii. 38). As God is unchangeable, so his election is unchangeable, and the elect can neither be cast away nor their number be diminished. The sense and certainty of election is a constant stimulus to humility and gratitude.

The non-elect are simply left to the just condemnation of their own sins. This is the decree of reprobation, which by no means makes God the author of sin (the very thought of which is blasphemy), but declares him to be an awful, irreprehensible, and righteous judge and avenger (\emph{Cat.} Ch. I. Art. 15).

SECOND HEAD OF DOCTRINE.

\emph{Of the Death of Christ.} [Limited Atonement.]—According to the sovereign counsel of God, the saving efficacy of the atoning death of Christ extends to all the elect [and to them only], so as to bring them infallibly to salvation. But, intrinsically, the sacrifice and satisfaction of Christ is of infinite worth and value, abundantly sufficient to expiate the sins of the whole world. This death derives its infinite value and dignity from these considerations; because the person who submitted to it was not only really man and perfectly holy, but also the only-begotten Son of God, of the same eternal and infinite essence with the Father and Holy Spirit, which qualifications were necessary to constitute him a Saviour for us; and because it was attended with a sense of the wrath and curse of God due to us for sin.

Moreover the promise of the gospel is, that whosoever believeth in Christ crucified shall not perish, but have everlasting life. This promise, together with the command to repent and believe, ought to be declared and published to all nations, and to all persons promiscuously and without distinction, to whom God out of his good pleasure sends the gospel.

And, whereas many who are called by the gospel do not repent nor believe in Christ, but perish in unbelief; this is not owing to any defect or insufficiency in the sacrifice offered by Christ upon the cross, but is wholly to be imputed to themselves.\footnote{The advocates of a limited atonement reason from the effect to the cause, and make the divine intention co-extensive with the actual application; but they can give no satisfactory explanation of such passages as John iii. 16 ('God so loved the \emph{world},' which never means the elect only, but all mankind); 1 John ii. 2 ('Christ is the propitiation for our sins, and \emph{not for ours only, but also} for the sins of the \emph{whole world}'); 1 Tim. ii. 4; 2 Pet. iii. 9. All admit, however, with the Articles of Dort, that the intrinsic value of the atonement, being the act of the God-man, is infinite and sufficient to cover the sins of all men. Dr. W. Cunningham says: 'The value or worth of Christ's sacrifice of himself depends upon, and is measured by, the dignity of his person, and is therefore infinite. Though many fewer of the human race had been to be pardoned and saved, an atonement of infinite value would have been necessary, in order to procure for them these blessings; and though many more, yea, all men, had been to be pardoned and saved, the death of Christ, being an atonement of infinite value, would have been amply sufficient, as the ground or basis of their forgiveness or salvation' (\emph{Historical Theol.} Vol. II. p. 331). Similarly, Dr. Hodge, Vol. II. pp. 544 sqq. After such admissions the difference of the two theories is of little practical account. Full logical consistency would require us to measure the value of Christ's atonement by the extent of its actual benefit or availability, and either to expand or to contract it according to the number of the elect; but such an opinion is derogatory to the dignity of Christ, and is held by very few extreme Calvinists of little or no influence. Cunningham says (p. 331): 'There is no doubt that all the most eminent Calvinistic divines hold the infinite worth or value of Christ's atonement—its full sufficiency for expiating all the sins of all men.'}

THIRD AND FOURTH HEADS OF DOCTRINE.

\emph{Of the Corruption of Man, his Conversion to God, and the Manner thereof.}—Man was originally formed after the image of God. His understanding was adorned with a true and saving knowledge of his Creator, and of spiritual things; his heart and will were upright, all his affections pure, and the whole Man was holy; but revolting from God by the instigation of the devil, and abusing the freedom of his own will, he forfeited these excellent gifts, and on the contrary entailed on himself blindness of mind, horrible darkness, vanity, and perverseness of judgment; became wicked, rebellious, and obdurate in heart and will, and impure in [all] his affections.

Man after the fall begat children in his own likeness. A corrupt stock produced a corrupt offspring. Hence all the posterity of Adam, Christ only excepted, have derived corruption from their original parent, not by imitation, as the Pelagians of old asserted, but by the propagation of a vicious nature in consequence of a just judgment of God.

Therefore all men are conceived in sin, and are by nature children of wrath, incapable of any saving good, prone to evil, dead in sin, and in bondage thereto; and, without the regenerating grace of the Holy Spirit, they are neither able nor willing to return to God, to reform the depravity of their nature, nor to dispose themselves to reformation.

What, therefore, neither the light of nature nor the law could do, that God performs by the operation of his Holy Spirit through the word or ministry of reconciliation: which is the glad tidings concerning the Messiah, by means whereof it hath pleased God to save such as believe, as well under the Old as under the New Testament.

As many as are called by the gospel are unfeignedly called; for God hath most earnestly and truly declared in his Word what will be acceptable to him, namely, that all who are called should comply with the invitation. He, moreover, seriously promises eternal life and rest to as many as shall come to him, and believe on him.

It is not the fault of the gospel, nor of Christ offered therein, nor of God, who calls men by the gospel, and confers upon them various gifts, that those who are called by the ministry of the Word refuse to come and be converted. The fault lies in themselves.

But that others who are called by the gospel obey the call must be wholly ascribed to God, who, as he hath chosen his own from eternity in Christ, so he calls them effectually in time, confers upon them faith and repentance, rescues them from the power of darkness, and translates them into the kingdom of his own Son, that they may show forth the praises of him who hath called them out of darkness into his marvelous light; and may glory not in themselves but in the Lord, according to the testimony of the Apostles in various places.

Faith is therefore the gift of God, not on account of its being offered by God to man, to be accepted or rejected at his pleasure, but because it is in reality conferred, breathed, and infused into him; nor even because God bestows the power or ability to believe, and then expects that man should, by the exercise of his own free will, consent to the terms of salvation, and actually believe in Christ; but because he who works in man both to will and to do, and indeed all things in all, produces both the will to believe and the act of believing also.

FIFTH HEAD OF DOCRINE.

\emph{Of the Perseverance of the Saints.}—Whom God calls, according to his purpose, to the communion of his Son our Lord Jesus Christ, and regenerates by the Holy Spirit, he delivers also from the dominion and slavery of sin in this life; though not altogether from the body of sin and from the infirmities of the flesh, so long as they continue in this world.

By reason of these remains of indwelling sin, and the temptations of sin and of the world, those who are converted could not persevere in a state of grace if left to their own strength. But God is faithful, who having conferred grace, mercifully confirms and powerfully preserves them therein, even to the end.

Of this preservation of the elect to salvation, and of their perseverance in the faith, true believers for themselves may and do obtain assurance according to the measure of their faith, whereby they arrive at the certain persuasion that they ever will continue true and living members of the Church; and that they experience forgiveness of sins, and will at last inherit eternal life.

This certainty of perseverance, however, is so far from exciting in believers a spirit of pride, or of rendering them carnally secure, that, on the contrary, it is the real source of humility, filial reverence, true piety, patience in every tribulation, fervent prayers, constancy in suffering and in confessing the truth, and of solid rejoicing in God; so that the consideration of this benefit should serve as an incentive to the serious and constant practice of gratitude and good works, as appears from the testimonies of Scripture and the examples of the saints.

In opposition to the Canons of Dort, Episcopius prepared a lengthy defense of the Arminian Articles and a confession of faith in Dutch, 1621, and in Latin, 1622. It claims no binding symbolical authority, and advocates liberty and toleration.\footnote{A German translation in Böckel's \emph{Bekenntniss-Schriften,} pp. 545–640.}

\xsection{III. Reformed Confessions of Germany.}

\xsubsection{The Tetrapolitan Confession.}

§ 68. The Tetrapolitan Confession.

Literature.

I. Editions of the Confessio Tetrapolitana.

The Latin text was first printed at Strasburg (Argentorati), A.D. 1531, Sept. (21 leaves); then in the \emph{Corpus et Syntagma} (1612 and 1654); in Augusti's  \emph{Corpus libr. symb.} (1827), pp. 327 sqq.; and in Niemeyer's  \emph{Collect. Confess.} (1840), pp. 740–770; comp. Proleg. p. lxxxiii.

The German text appeared first at Strasburg, Aug. 1531 (together with the Apology, 72 leaves); then again 1579, ed. by John Sturm, but suppressed by the magistrate, 1580; at Zweibrücken, 1604; in Beck's \emph{Symbol. Bücher,} Vol. I. pp. 401 sqq.; in Böckel's  \emph{Bekenntniss-Schriften,} pp. 363 sqq.

II. History.

Gottl. Wernsdorff:  \emph{Historia Confessionis Tetrapolitanæ.} Wittenb. 1694, ed. iv. 1721.

J. H. Fels: \emph{ Dissert. de varia Confess. Tetrapolitanæ fortuna præsertim in civitate Lindaviensi.} Götting. 1755.

Planck: \emph{ Geschichte des Protest. Lehrbegrifs,} Vol. III. Part I. (second ed. 1796), pp. 68–94.

J. W. Röhrich: \emph{ Geschichte der evangel. Kirche des Elsasses.} Strassburg, 1855, 3 vols.

J. W. Baum: \emph{ Capito und Butzer} (Elberf. 1860), pp. 466 sqq. and 595.

H. Mallet, in Herzog's \emph{Encykl.} Vol. XV. pp. 574–576.

Comp. also the literature on the Augsburg Diet and the Augsburg Confession, especially Salig and Förstemann, quoted in § 41, p. 225.

THE REFORMED CHURCH IN GERMANY.

The mighty genius of Luther, aided by the learning of Melanchthon, controlled the German Reformation at first to the exclusion of every other influence; and if Lutheranism had not assumed a hostile and uncompromising attitude towards Zwinglianism, Calvinism, and the later theology of Melanchthon, it would probably have prevailed throughout the German empire, as the Reformed creed prevailed in all the Protestant cantons of Switzerland. But the bitter eucharistic controversies and the triumph of rigid Lutheranism in the Formula of Concord over Melanchthonianism drove some of the fairest portions of Germany, especially the Palatinate and Brandenburg, into the Reformed communion.

The German branch of the Reformed family grew up under the combined influences of Zwingli, Calvin, and Melanchthon. Zwingli's reformation extended to the southern portions of Germany bordering on Switzerland, especially the free imperial cities of Strasburg, Constance, Lindau, Memmingen, and Ulm. It is stated that the majority of the Protestant citizens of Augsburg during the Diet of 1530 sympathized with him rather than with Luther. Calvin spent nearly three years at Strasburg (1538–41), and exerted a great influence on scholars through his writings. Melanchthon (who was a native of the Palatinate), in his later period, emancipated himself gradually from the authority of Luther, and sympathized with Calvin in the sacramental question, while in the doctrines of divine sovereignty and human freedom he pursued an independent course. He trained the principal author of the Heidelberg Catechism (Ursinus), reorganized the University of Heidelberg (1557), which became the Wittenberg of the Reformed Church in Germany, and threw on several occasions the weight of his influence against the exclusive type of Lutheranism advocated by such men as Flacius, Heshusius, and Westphal. He impressed upon the German Reformed Church his mild, conciliatory spirit and tendency towards union, which, at a later period, prevailed also in a large part of the Lutheran Church. The German Reformed Church, then, occupies a mediating position between Calvinism and Lutheranism. It adopts substantially the Calvinistic creed, but without the doctrine of reprobation (which is left to private opinion), and without its strict discipline; while it shares with the Lutheran Church the German language, nationality, hymnology, and mystic type of piety.\footnote{Dr. Heppe, in his numerous and learned works on the history and theology of the German Reformation period, endeavors to identify the German Reformed Church with Melanchthonianism (which was only an element in it), and Melanchthonianism with original German Protestantism (which was prevailingly Lutheran in the strict sense of the term), thus overestimating the influence of Melanchthon and underrating the influence of Zwingli and Calvin. His books are very valuable, but one-sided, and must be supplemented by the writings of Alex. Schweizer (\emph{Die Centraldogmen}) and others on the same subject.} The great majority of German Reformed congregations have, since 1817, under the lead of the royal house of Prussia, been absorbed in what is called the Evangelical or United Evangelical Church. The aim of this union was originally to substitute one Church for two, but the result has been to add a third Church to the Lutheran and Reformed, since these still continue their separate existence in Germany and among the German emigrants in other countries.\footnote{The large German Protestant population of the United States is divided among Lutherans (the most numerous), German Reformed, and Evangelicals (or Unionists). A considerable number is connected with English denominations, especially the Methodists and Presbyterians.}

BUCER.

Among the framers of the character of the Reformed Church in Germany, Martin Bucer (Butzer),\footnote{He wrote his name in German \emph{Butzer} (i.e., \emph{Cleanser,} from \emph{putzen, to cleanse}), in Latin \emph{Bucerus,} in Greek \textgreek{Βούκηρος.} See Baum, l.c. p. 88.} Wolfgang Fabricius Capito, and Caspar Hedio occupy the next place after Zwingli, Calvin, and Melanchthon. Bucer (1491–1551), the learned and devoted reformer of Strasburg, and a facile diplomatist, was a personal friend of Zwingli, Luther, and Calvin, and a mediator between the Swiss and the German Reformation, as also between Continental and Anglican Protestantism. He labored with indefatigable zeal for an evangelical union, and hoped to attain it by elastic compromise formulas (like the Wittenberg Concordia of 1536), which concealed the real difference, and in the end satisfied neither party. He drew up with Melanchthon the plan of a reformation in Cologne at the request of the archbishop. During the Interim troubles he accepted a call to England, aided Cranmer in his reforms, and died as Professor of Theology at Cambridge, universally lamented. In the reign of Bloody Mary he was formally condemned as a heretic, his bones were dug up and publicly burned (Feb. 6, 1556); but Elizabeth solemnly restored the 'blessed' memory of 'the dear martyrs Martin Bucer and Paul Fagius.' In attainments and fertility as a writer he was not surpassed in his age.\footnote{See a chronological list of his very numerous printed works in Baum, pp. 586 sqq. Baum says: '\emph{An Fruchtbarkeit kommt ihm} [Bucer] \emph{kaum Luther gleich, trots dem dass er bei weitem mehr als Luther, ja in seiner letzten Lebensperiode beinahe beständig, auf Reisen, Conventen, Reichstagen und Colloquien, in befreundeten Städten und Orten als Organisator der Kirchenreformation abwesend und in Anspruch genommen war. Mit einer beispiellosen Elasticität des Geistes angethan, mit einem fieberhaftigen Thätigkeitstriebe behaftet, schrieb er, vermöge des ungemeinen Reichthums seiner Kenntnisse mit solcher fabelhaften Leichtigkeit und Unleserlichkeit, dass nicht allein zu dem Meisten was von Anderen gelesen werden sollte, ein mit seiner die Worte blos andeutenden Schrift genau vertrauter Amanuensis nothwendig war, sondern dass er auch neben seinen Amtsgeschäften noch bei weitem mehr förderte als zwei der geübtesten Schreiber in’s Reine bringen konnten. Er hat umfangreiche Bücher auf seinen Reisen geschrieben.}' His best amanuensis, Conrad Huber, began a complete edition of his works, of which the first volume only appeared at Basle, 1577 (959 pages, folio). It is called \emph{Tomus Anglicanus,} because it contains mostly the books which Bucer wrote in England. Many of his MSS. are preserved in Strasburg and in England.}

THE CONFESSION OF THE FOUR CITIES.

The oldest Confession of the Reformed Church in Germany is the Tetrapolitan Confession, also called the Strasburg and the Swabian Confession.\footnote{\emph{Confessio Tetrapolitana, C. Quatuor Civitatum, C. Argentinensis} (\emph{Argentorati}), \emph{C. Suevica, die Confession der vier Städte, das Vierstädte-Bekenntniss.}}

It was prepared in great haste, during the sessions of the Diet of Augsburg in 1530, by Bucer, with the aid of Capito and Hedio, in the name of the four imperial cities (hence the name) of Strasburg, Constance, Memmingen, and Lindau which, on account of their sympathy with Zwinglianism, were excluded by the Lutherans from their political and theological conferences, and from the Protestant League. They would greatly have preferred to unite with the Lutherans in a common confession; but at that time even Melanchthon was more anxious to conciliate the Papists than the Zwinglians and Anabaptists; and of the Lutheran princes the Landgrave Philip of Hesse was the only one who, from a broad, statesman-like view of the critical situation, favored a solid union of the Protestants against the common foe, but in vain. Hence after the Lutherans had presented their Confession, June 25, and Zwingli his own, July 8, the Four Cities handed theirs, July 11, to the Emperor, in German and Latin. It was not read before the Diet, but a Confutation full of misrepresentations was prepared by Faber and Cochläus, and read October 24 (or 17). The Strasburg divines were not even favored with a copy of this Confutation, but procured one secretly, and answered it by a 'Vindication and Defense' (as Melanchthon wrote his Apology of the Augsburg Confession during the Diet). The Confession and Apology, after being withheld for a year from print for the sake of peace, were officially published in both languages at Strasburg in the autumn of 1531.\footnote{Under the title, '\textbf{Bekandtnuss der vier Frey und Reichstätt, Strassburg, Constantz, Memmingen und Lindaw, in deren sie keys-Majestat, uff dem Reichstag zu Augspurg im xxx. Jar gehalten, ires glaubens und fürhabens, der Religion halb, rechenschaft gethon haben.—Schriftliche Beschirmung und verthedigung derselbigen Bekandtnuss, gegen der Confutation und Widerlegung, so den gesandten der vier Stätten, uff bemeldtem Reichstage, offentlich fürgelesen, und hie getrewlich eingebracht ist.}' At the end, '\textbf{Getruckt zu Strassburg durch Johann Schweintzer, uff den xxii. Augusti, MDXXXI.}' Shortly after the appearance of the German original there appeared a Latin translation, which, however, did not contain the Apology. The title is as follows: '\emph{Confessio Religionis Ckristianæ Sacratissimo Imperatori Carolo V. Augusto, in Comitiis Augustanis Anno MDXXX. per legatos Civitatum Argentorati, Constantiæ, Memmingæ, et Lindaviæ exhibita. Si quis voluerit voluntati ejus obtemperare, is cognoscet de doctrina utrum ex Deo sit an ego a me ipso loquar Joh. VII.}' At the end, '\emph{Argentorati Georgio Ulrichero Andlano Impressore Anno MDXXXI., mense Septemb.}'—These titles are copied from Baum, l.c. p. 595. Comp. Niemeyer, \emph{Proleg.} pp. lxxxiv. sq. A new German translation from the Latin is given in Walch's edition of Luther's Works, Vol. XX. pp. 1966–2008.}

The Tetrapolitan Confession consists of twenty-three chapters, besides Preface and Conclusion. It is in doctrine and arrangement closely conformed to the Lutheran Confession of Augsburg, and breathes the same spirit of moderation. The Reformed element, however, appears in the first chapter (On the Matter of Preaching), in the declaration that nothing should be taught in the pulpit but what was either expressly contained in the Holy Scriptures or fairly deduced therefrom.\footnote{\par '\emph{Mandavimus iis, qui concionandi apud nos munere fungebantur, ut nihil aliud quam quæ sacris literis aut continentur, aut certe nituntur, e suggestu docerent. Videbatur namque nobis haud indignum, eo in illo tanto discrimine confugere, quo confugerunt olim et semper, non solum sanctissimi Patres, Episcopi, et Principes, sed quilibet etiam privati, nempe ad authoritatem Scripturæ arcanæ. Ad quam nobiliores Thessalonicensium auditum Christi Evangelium explorasse, divus Lucas cum laude illorum memorat, in qua Paulus summo studio versari suum Timotheum voluit, sine cuius authoritate, nulli Pontifices suis decretis obedientiam, nulli patres suis scriptis fidem, nulli denique Principes suis legibus authoritatem unquam postularunt, ex qua demum ducendas sacras conciones, et magnum Sacri Imperii concilium Nurembergæ, anno Christi M.D.XXIII. celebratum sancivit. Si enim verum divus Paulus testatus est, per divinam Scripturam hominem Dei penitus absolvi, atque ad omne opus bonum instrui, nihil poterit is veritatis Christianæ, nihil doctrinæ salutaris desiderare, Scripturam qui consulere religiose studeat.}'} (The Lutheran Confession, probably from prudential and irenical considerations, is silent on the supreme authority of the Scriptures.) The evangelical doctrine of justification is stated in the third and fourth chapters more clearly than by Melanchthon, namely, that we are justified not by works of our own, but solely by the grace of God and the merits of Christ through a living faith, which is active in love and productive of good works. Images are rejected in Ch. XXII. The doctrine of the Lord's Supper (Ch. XVIII.) is couched in dubious language, which was intended to comprehend in substance the Lutheran and the Zwinglian theories, and contains the germ of the view afterwards more clearly and fully developed by Calvin. In this ordinance, it is said, Christ offers to his followers, as truly now as at the institution, his very body and blood as spiritual food and drink, whereby their souls are nourished to everlasting life.\footnote{\par '\emph{De hoc venerando corporis et sanguinis Christi sacramento omnia, quæ de illo Evangelistæ, Paulus et sancti Patres scripta reliquerunt, nostri fide optima docent, commendant, inculcant. Indeque singulari studio hanc Christi in suos bonitatem, semper depredicant, qua is non minus hodie, quam in novissima illa cœna, omnibus qui inter illius discipulos ex animo nomen dederunt, cum hanc cœnam, ut ipse instituit repetunt, verum suum corpus, verumque suum sanguinem, vere edendum et bibendum, in cibum potumque animarum, quo illæ in æternam vitam alantur, dare per sacramenta dignatur, ut jam ipse in illis, et illi in ipso vivant et permaneant, in die novissimo, in novam et immortalem vitam per ipsum resuscitandi, juxta sua illa æternæ veritatis verba}: ``\emph{Accipite et manducate, hoc est corpus meum},'' \emph{etc.} ``\emph{Bibite ex eo omnes, hic calix est sanguis meus},'' \emph{etc. Præcipua vero diligentia populi animos, nostri ecclesiastæ ab omni tum contentione, tum supervacanea et curiosa disquisitione, ad illud revocant, quod solum prodest, solumque a Christo servatore nostro spectatum est, nempe ut ipso pasti, in ipso et per ipsum vivamus, vitam Deo placitam, sanctam, et ideo perennem quoque et beatam, simusque inter nos omnes unus panis, unum corpus, qui de uno pane in sacra cœna participamus. Quo sane factum est, ut divina sacramenta, sacrosancta Christi cœna, quam religiosissime, reverentiaque singulari apud nos et administrentur, et sumantur.}' Ebrard (\emph{Kirchen- und Dogmengeschichte,} Vol. III. p. 93) says of Bucer, that he had the theological elements for a true doctrinal union of the Lutheran and Reformed views of the eucharist. '\emph{In der richtigen exegetisehen Grundlage völlig mit Zwingli einig, brachte er das Element, welches auch in Zwingli keimartig vorhanden gewesen, aber in der Hitze des Streites ganz zurückgetreten war—die Lebensgemeinschaft oder unio mystica mit der } Person  \emph{Christi—im Sinne der Tetrapolitana} (\emph{d.i. im Sinne der nachherigen calvinisch-melanchthonischen Lehre}) \emph{zur Entwicklung.}'} Nothing is said of the oral manducation and the fruition of unbelievers, which are the distinctive features of the Lutheran view. Bucer, who had attended the Conference at Marburg in 1529, labored with great zeal afterwards to bring about a doctrinal compromise between the contending theories, but without effect.

We may regard the Strasburg Confession as the first attempt at an evangelical union symbol. But Bucer's love for union was an obstacle to the success of his confession, which never took deep root; for in the Reformed Churches it was soon superseded by the clearer and more logical confessions of the Calvinistic type, and the four cities afterwards signed the Lutheran Confession to join the Smalcald League. Bucer himself remained true to his creed, and reconfessed it in his last will and testament (1548), and on his death-bed.\footnote{Baum, pp. 569, 572.}

\xsubsection{The Heidelberg Catechism. A.D. 1563.}

§ 69. The Heidelberg Catechism. A.D. 1563.

Literature.

I. Standard Editions of the Catechism.

Official German editions of 1563 (three), 1585, 1595, 1684, 1724, 1863 (American). The original title is '\textbf{Catechismus} | \textbf{Oder} | \textbf{Christlicher Underricht}, | \textbf{wieder in Kirchen und Schu-} | \textbf{len der Churfürstlichen} | \textbf{Pfaltz getrieben} | \textbf{wirdt.} | \textbf{Gedruckt in der Churfürstli-} | \textbf{chen Stad Heydelberg, durch} | \textbf{Johannem Mayer.} | 1563.' With the Electoral arms. 95 pages.

There is but one copy of the first edition known to exist, and this did not come into public notice till 1864. It belonged to Prof. Hermann Wilken, of Heidelberg, whose name it bears, with the date 1563; was bought by Dr. Treviranus, of Bremen, in 1823, given by him to Dr. Menken, bought back after Menken's death, 1832, and is now in the University Library at Utrecht. I examined it in October and November, 1865, at Bremen. It has the remark, '\emph{Diesses ist die allererste Edition, in welcher Pag.} 55 \emph{die} 80\emph{ste Frag und Antwort nicht gefunden wirdt. Auff Churfürstlichen Befehl eingezogen. Liber rarissimus.}' The Scripture texts are quoted in the margin, but only the chapters, since the versicular division (which first appeared in Stephens's Greek Testament of 1551) had not yet come into general use. A quasi fac-simile of this copy was issued by the Rev. Albrecht Wolters, then at Bonn (now at Halle), under the title, '\emph{Der Heidelberger Katechismus in seiner ursprünglichen Gestalt, herausgegeben nebst der Geschichte seines Textes im Jahre} 1563.' Bonn, 1864. Comp. his art. in the \emph{Studien und Kritiken} for 1867, pp. 1, 2.

Niemeyer, in his collection of Reformed Confessions, pp. 390 sqq., gives, besides the Latin text, a faithful reprint of the \emph{third} German edition, with the eightieth question in full.

Philipp Schaff:  \emph{Der Heidelberger Katechismus. Nach der ersten Ausgabe von} 1563 \emph{revidirt und mit kritischen Anmerkungen, sowie einer Geschichte und Charakteristik des Katechismus versehen.} Philadelphia (J. Kohler), 1863; second edition, revised and enlarged, 1866. This edition was prepared for the tercentenary celebration of the Heidelberg Catechism, and gives the received text of the third edition with the readings of the first and second editions, and the Scripture proofs in full.

The Latin translation was published in 1563, and again in 1566, under the title, 'Cate- |  chesis Religio- |  nis Christianæ,  | \emph{ quæ traditur in Ecclesiis | et Scholis Pala- | tinatus. | Heydelbergæ. | Excusum anno post Christum | natum M.D.LXVI.}' I saw a copy of this \emph{ed. Latina,} in the library of the late Dr. Treviranus, in Bremen (1865). On the title-page the words are written, '\emph{Editio rara et originalis};' also the name of G. Menken, the former owner. The Scripture references are marked on the margin, including the verses. The eightieth question is complete (with '\emph{execranda idololatia}'), pp. 62 and 63, and supported by many Scripture texts and the \emph{Can. Missæ.} The questions are divided into fifty-two Sundays. '\emph{Precationes aliquot privatæ et publicæ},' a '\emph{Precatio scholastica},' and some versified prayers of Joachim Camerarius (the friend and biographer of Melanchthon), are added.

The best \emph{English,} or rather \emph{American,} edition of the Catechism is the stately triglot tercentenary edition prepared at the direction of the German Reformed Church in the United States, by a committee consisting of E. V. Gerhart, D.D., John W. Nevin, D.D., Henry Harbaugh, D.D., John S. Kessler, D.D., Daniel Zacharias, D.D., and three laymen, and issued under the title, '\emph{The Heidelberg Catechism, in German, Latin, and English; with an Historical Introduction} (by Dr. Nevin), New York (Charles Scribner), 1863.' 4to. The German text is a reprint of the third edition after Niemeyer, with the German in modern spelling added; the English translation is made directly from the German original, and is far better than the one in popular use, which was made from the Latin. It is the most elegant and complete edition of the Catechism ever published, but it appeared before the discovery of the \emph{editio princeps,} and repeats the error concerning the eightieth question (see Introd. p. 38).

II. Commentaries.

The commentaries and sermons on the Heidelberg Catechism are exceedingly numerous, especially in the German and Dutch languages. The first and most valuable is from the chief author, Zach. Ursinus: \emph{Corpus Doctrinæ orthodoxæ, or Commentary on the Heidelberg Catechism,} ed. by his pupil, David Pareus, and repeatedly published at Heidelberg and elsewhere—1591, 1618, etc.—in Latin, German, Dutch, and English. An American edition, on the basis of the English translation of Bishop Dr. H. Parry, was issued by Dr. Williard (President of Heidelberg College, Tiffin, O.), Columbus, O. 1850. Other standard commentaries are by Coccejus (1671), d’Outrein (1719), Lampe (1721), Stähelin (1724), and van Alpen (1800). See a fuller list by Harbaugh in 'Mercersb. Rev.' for 1860, pp. 601–625, and in Bethune's \emph{Lectures.}

Of more recent works we name—

Karl Sudhoff: \emph{ Theologisches Handbuch zur Auslesung des Heidelberger Catechismus.} Francf. a M. 1862.

Geo. W. Bethune (D.D., and minister of the Ref. Dutch Ch., N.Y.; d. 1862): \emph{Expository Lectures on the Heidelb. Catech.} N. York, 1864, 2 vols., with an alphabet. list of works by Van Nest at close of Vol. II.

Hermann Dalton (Ger. Ref. minister at St. Petersb.): \emph{Immanuel. Der Heidelberger Katechismus als Bekenntniss- und Erbauungsbuch der evangel. Kirche erklärt und an’s Herz gelegt.} Wiesbaden, 1870 (pp. 539).

III. Historical Works on the Catechism.

H. Alting (Prof. of Theology at Heidelberg and Gröningen, d. 1644): \emph{Historia Ecclesiæ Patatinæ.} Frankf. a. M. 1701.

B. G. Struve: \emph{ Pfälzische Kirchenhistorie.} Frankf. 1721, Ch. V. sqq.

D. L. Wundt: \emph{ Grundriss der pfälzischen Kirchengeschichte bis zum Jahr} 1742. Heidelb. 1798.

Jaques Lenfant: \emph{ L’innocence du Catéchisme de Heidelberg.} Heidelb. 1688 (1723).

J. Chr. Köcher: \emph{ Katechetische Geschichte der Reformirten Kirche, sonderlich der Schicksale des Heidelberger Katechismi. }Jena, 1766, pp. 237–444.

G. J. Planck: \emph{ Geschichte der protestantischen Theologie von Luther’s Tode,} etc. Vol. II. Part II. pp. 475–491. (This is Vol. V. of his great work on the \emph{Geschichte der Entstehung,} etc., \emph{unseres protestant. Lehrbegriffs.})

Heinr. Simon van Alpen; \emph{ Geschichte u. Literatur des Heidelb. Katechismus.} Frankf. a. M. 1800. Vol. III. Part II. (The first two volumes and the first part of the third volume of this catechetical work contain explanations and observations on the Catechism, which are, however, semi-rationalistic.)

Joh. Chr. W. Augusti: \emph{ Versuch einer hist.-kritischen Einleitung in die beiden Haupt-Katechismen} (the Luth. and Heidelb.) \emph{der evangelischen Kirche.} Elberfeld, 1824, pp. 96 sqq.

Rienäcker: Article on the Heidelb. Catechism in \emph{Ersch und Gruber, Allgem. Encyklop.} Sect. II. Part IV. pp. 386 sqq.

Ludwig Häusser: \emph{ Geschichte der Rhein-Pfalz.} Heidelb. 1845. Vol. II.

D. Seisen: \emph{ Geschichte der Reformation zu Heidelberg, von ihren ersten Anfängen bis zur Abfassung des Heidelb. Katechismus. Eine Denkschrift zur dreihundertjährigen Jubelfeier daselbst am} 3. \emph{Jan.} 1846. Heidelb. 1846.

Aug.. Ebrard: \emph{ Das Dogma vom heil. Abendmahl und seine Geschichte.} F. a. M. 1846. Vol. II. pp. 575 sqq.

K. Fr. Vierordt: \emph{ Geschichte der Reformation im Grossherzogthum Baden. Nach grossentheils handschriftlichen Quellen.} Karlsruhe, 1847.

John W. Nevin: \emph{ History and Genius of the Heidelberg Catechism.} Chambersburg, Pennsylvania, 1847. (The best work on the Catechism in English.) Comp. Dr. Nevin's able Introduction to the triglot tercentenary edition of the H. C. New York, 1863, pp. 11–127.

Karl Sudhoff: \emph{ C. Olevianus und Z. Ursinus. Leben und ausgewählte Schriften.} Elberfeld, 1857.

G. D. J. Schotel: \emph{ History of the Origin, Introduction, and Fortunes of the Heidelberg Catechism} (in Dutch). Amsterdam, 1863.

Several valuable essays on the Heidelberg Catechism, by Plitt, Sack, and Ullmann, in the \emph{Studien und Kritiken} for 1863, and by Wolters and Trechsel, ibid, for 1867.

Tercentenary Monument. \emph{ In Commemoration of the Three Hundredth Anniversary of the Heidelberg Catechism.} Published by the German Reformed Church of the United States of North America, in English and German. The German ed. by Dr. Schaff, with an historical introduction. Chambersburg and Philadelphia, Pa. 1863. This work contains about twenty essays, by European and American theologians, on the history and theology of the Heidelberg Catechism.

J. I. Dœdes (Prof. at Utrecht): \emph{De Heidelbergsche Catechismus in zijne eerste Levensjaren} 1563–1567. \emph{Historische en Bibliografische Nalezing met} 26 \emph{Facsimiles.} Utrecht, 1867 (pp. 154). Very valuable for the early literary history of the H. C., with fac-similes of the first German, Latin, and Dutch editions.

THE REFORMATION IN THE PALATINATE.

The Palatinate, one of the finest provinces of Germany, on both sides of the upper Rhine, was one of the seven electorates (\emph{Kurfürstenthümer}), whose rulers, in the name of the German people, elected the Emperor of Germany. After the dissolution of the old empire (1806) it ceased to be a politico-geographical name, and its territory is now divided between Baden, Bavaria, Hesse Darmstadt, Nassau, and Prussia. Its capital was Heidelberg (from 1231 till 1720), famous for its charming situation at the foot of the Königsstuhl, on the banks of the Swabian river Neckar, for its picturesque castle, and for its university (founded in 1346).

Luther made a short visit to Heidelberg in 1518, and defended certain evangelical theses. In 1546, the year of Luther's death, the Reformation was introduced under the Elector Frederick II. Melanchthon, who was a native of the Palatinate, and twice received a call to a professorship of theology at Heidelberg (1546 and 1557), but declined, acted as the chief counselor in the work, and aided, on a personal visit in 1557, in reorganizing the university on an evangelical basis under Otto Henry (1556–59). He may therefore be called the Reformer of the Palatinate. He impressed upon it the character of a moderate Lutheranism friendly to Calvinism. The Augsburg Confession was adopted as the doctrinal basis, and the cultus was remodeled (as also in the neighboring Duchy of Würtemberg) after Zwinglian simplicity. Heidelberg now began to attract Protestant scholars from different countries, and became a battle-ground of Lutheran, Philippist, Calvinist, and Zwinglian views. The conflict was enkindled as usual by the zeal for the real presence. Tilemann Heshusius, whom Melanchthon, without knowing his true character, had recommended to a theological chair (1558), introduced, as General Superintendent, exclusive Lutheranism, excommunicated Deacon Klebitz for holding the Zwinglian view, and even fought with him at the altar about the communion cup. This public scandal was the immediate occasion of the Heidelberg Catechism.

FREDERICK III.

During this controversy Frederick III., surnamed the Pious (1515–1576), became Elector of the Palatinate, 1559. He made it the chief object of his reign to carry out the reformation begun by his predecessors. He tried at first to conciliate the parties, and asked the advice of Melanchthon, who, a few months before his death, counseled peace, moderation, and Biblical simplicity, and warned against extreme and scholastic subtleties in the doctrine of the Lord's Supper.\footnote{\emph{Responsio Ph. Mel. ad quæstionem de controversia Heidelbergensi} (Nov. 1, 1559), in \emph{Corp. Reform.} Vol. IX. pp. 960 sqq. It is the last public utterance of Melanchthon on the eucharistic question, and agrees substantially with the doctrine of Calvin, as it was afterwards expressed in the Heidelberg Catechism.} He deposed both Heshusius and Klebitz, arranged a public disputation (June, 1560) on the eucharist, decided in favor of the Melanchthonian or Calvinistic view, called distinguished foreign divines to the university, and intrusted two of them with the composition of the Heidelberg Catechism, which was to secure harmony of teaching and to lay a solid foundation for the religious instruction of the rising generation.

Frederick was one of the purest and noblest characters among the princes of Germany. He was to the Palatinate what King Alfred and Edward VI. were to England, what the Electors Frederick the Wise and John the Constant were to Saxony, and Duke Christopher to Würtemberg. He did more for educational and charitable institutions than all his predecessors. He devoted to them the entire proceeds of the oppressed convents. He lived in great simplicity that he might contribute liberally from his private income to the cause of learning and religion. He was the first German prince who professed the Reformed Creed, as distinct from the Lutheran. For this he suffered much reproach, and was threatened with exclusion from the benefits of the Augsburg Treaty of Peace (concluded in 1555), since Zwinglianism and Calvinism were not yet tolerated on German soil. But at the Diet of Augsburg, in 1566, he made before the Emperor a manly confession of his faith, and declared himself ready to lose his crown rather than violate his conscience. Even his opponents could not but admire his courage, and the Lutheran Elector Augustus of Saxony applauded him, saying, 'Fritz, thou art more pious than all of us.' He praised God on his death-bed that he had been permitted to see such a reformation in Church and school that men were led away from human traditions to Christ and his divine Word. He left in writing a full confession of his faith, which may be regarded as an authentic explanation of the Heidelberg Catechism; it was published after his death by his son, John Casimir (1577).

URSINUS AND OLEVIANUS.

Frederick showed his wisdom by calling two young divines, Ursinus and Olevianus, to Heidelberg to aid in the Reformation and to prepare an evangelical catechism. They belong to the reformers of the second generation. Theirs it was to nurture and to mature rather than to plant. Both were Germans, but well acquainted with the Reformed Churches in Switzerland and France. Both suffered deposition and exile for the Reformed faith.

Zacharias Ursinus (Bär), the chief author of the Heidelberg Catechism, was born at Breslau, July 18, 1534, and studied seven years (1550–1557) at Wittenberg under Melanchthon, who esteemed him as one of his best pupils and friends. He accompanied his teacher to the religious conference at Worms, 1557, and to Heidelberg, and then proceeded on a literary journey to Switzerland and France. He made the personal acquaintance of Bullinger and Peter Martyr at Zurich, of Calvin and Beza at Geneva, and was thoroughly initiated into the Reformed Creed. Calvin presented him with his works, and wrote in them the best wishes for his young friend. On his return to Wittenberg he received a call to the rectorship of the Elizabeth College at Breslau. After the death of Melanchthon he went a second time to Zurich (Oct., 1560), intending to remain there. In the following year he was called to a theological chair at Heidelberg. Here he labored with untiring zeal and success till the death of Frederick III., 1576, when, together with six hundred steadfast Reformed ministers and teachers, he was deposed and exiled by Louis VI., who introduced the Lutheran Creed. Ursinus found a refuge at Neustadt an der Hardt, and established there, with other deposed professors, a flourishing theological school under the protection of John Casimir, the second son of Frederick III. He died in the prime of his life and usefulness, March 6, 1583, leaving a widow and one son. In the same year Casimir succeeded his Lutheran brother in the Electorate, recalled the exiled preachers, and re-established the Reformed Church in the Palatinate.

Ursinus was a man of profound classical, philosophical, and theological learning, poetic taste, rare gift of teaching, and fervent piety. His devotion to Christ is beautifully reflected in the first question of the Heidelberg Catechism, and in his saying that he would not take a thousand worlds for the blessed assurance of being owned by Jesus Christ. He was no orator, and no man of action, but a retired, modest, and industrious student.\footnote{On the door of his study he inscribed the warning, '\emph{Amice, quisquis huc venis, aut agita paucis, aut abi, aut me laborantem adjuva.}'} His principal works, besides the Catechism, are a Commentary on the Catechism (\emph{Corpus doctrinæ orthodoxæ}) and a defense of the Reformed Creed against the attacks of the Lutheran Formula of Concord.

Caspar Olevianus (Olewig), born at Treves Aug. 10,1536, studied the ancient languages at Paris, Bourges, and Orleans, and theology at Geneva and Zurich. He enjoyed, like Ursinus, the personal instruction and friendship of the surviving reformers of Switzerland. He began to preach the evangelical doctrines at Treves, was thrown into prison, but soon released, and called to Heidelberg, 1560, by Frederick III., who felt under personal obligation to him for saving one of his sons from drowning at the risk of his own life. He taught theology and preached at the court. He was the chief counselor of the Elector in all affairs of the Church. In 1576 he was banished on account of his faith, and accepted a call to Herborn, 1584, where he died, Feb. 27, 1585. His last word was a triumphant '\emph{certissimus},' in reply to a friend who asked him whether he were certain of his salvation. Theodore Beza lamented his death in a Latin poem, beginning

'\emph{Eheu, quibus suspiriis,}

\emph{Eheu, quibus te lacrymis}

\emph{Oleviane, planxero?'}

Olevianus was inferior to Ursinus in learning, but his superior in the pulpit and in church government. He wrote an important catechetical work on the covenant of grace, and is regarded as the forerunner of the federal theology of Coccejus and Lampe. He labored earnestly, but only with moderate success, for the introduction of the Presbyterian form of government and a strict discipline, after the model of Geneva. Thomas Erastus (Lieber), Professor of Medicine at Heidelberg, and afterwards of Ethics at Basle (died 1583), opposed excommunication, and defended the supremacy of the state in matters of religion; hence the term 'Erastianism' (equivalent to Cæsaropapism).

PREPARATION AND PUBLICATION OF THE CATECHISM.

The Heidelberg Catechism, as it is called after the city of its birth, or the Palatinate (also Palatine) Catechism, as it is named after the country for which it was intended, was prepared on the basis of two Latin drafts of Ursinus and a German draft of Olevianus. The peculiar gifts of both, the didactic clearness and precision of the one, and the pathetic warmth and unction of the other, were blended in beautiful harmony, and produced a joint work which is far superior to all the separate productions of either. In the Catechism they surpassed themselves. They were in a measure inspired for it. At the same time, they made free and independent use of the Catechisms of Calvin, Lasky, and Bullinger. The Elector took the liveliest interest in the preparation, and even made some corrections.

In December, 1562, Frederick submitted the work to a general synod of the chief ministers and teachers assembled at Heidelberg, for revision and approval. It was published early in 1563, in German, under the title 'Catechismus, Or Christian Instruction, as conducted in the Churches and Schools of the Electoral Palatinate.'\footnote{See the original title in the literature above.} It is preceded by a short Preface of the Elector, dated Tuesday, January 19, 1563, in which he informs the superintendents, clergymen, and schoolmasters of the Palatinate that, with the counsel and co-operation of the theological faculty and leading ministers of the Church, he had caused to be made and set forth a summary instruction or Catechism of our Christian religion from the Word of God, to be used hereafter in churches and schools for the benefit of the rising generation.

THE THIRD EDITION AND THE EIGHTIETH QUESTION.

There appeared, in the year 1563, three official editions of the Catechism with an important variation in the eightieth question, which denounces the Romish mass as 'a denial of the one sacrifice of Christ, and as an accursed idolatry.' In the first edition this question was wanting altogether; the second edition has it in part; the third in full, as it now stands.\footnote{By the discovery of the copy of the first ed., 1864, the origin of the eightieth question was satisfactorily decided. A second copy of the original ed. is in the Imperial Library of Vienna. The Brit. Museum contains a copy of the Engl. trans. by ``William Turner, Doctor of Physick, Imprinted at London, by Richard Jones, 1572.''—Ed.} This question was inserted by the express command of the Elector, perhaps by his own hand, as a Protestant counter-blast to the Romish anathemas of the Council of Trent, which closed its sessions Dec. 4, 1563. Hence the remark at the end of the second and third editions: 'What has been overlooked in the first print, as especially on folio 55 [which contains the eightieth question], has now been added by command of his electoral grace. 1563.'

The same view of the Romish doctrine of transubstantiation and the sacrifice of the mass was generally entertained by the Reformers, and is set forth as strongly in the Articles of Smalcald and other symbolical books, both Lutheran and Reformed. It must be allowed to remain as a solemn protest against idolatry. But the wisdom of inserting controversial matter into a catechism for the instruction of the youth has been justly doubted. The eightieth question disturbs the peaceful harmony of the book, it rewards evil for evil, it countenances intolerance, which is un-Protestant and unevangelical. It provoked much unnecessary hostility, and led even, under the Romish rule of the Elector Charles Philip, in 1719, to the prohibition of the Catechism; but the loud remonstrance of England, Prussia, Holland, and other Protestant states forced the Elector to withdraw the tyrannical decree within a year, under certain conditions, to save appearances.

TRANSLATIONS.

The Heidelberg Catechism was translated into all the European and many Asiatic languages. It has the pentecostal gift of tongues in a rare degree. It is stated that, next to the Bible, the 'Imitation of Christ,' by Thomas à Kempis, and Bunyan's 'Pilgrim's Progress,' no book has been more frequently translated, more widely circulated and used. Whole libraries of paraphrases, commentaries, sermons, attacks, and defenses were written about it. In many Reformed churches, especially in Holland (and also in the United States), it was and is to some extent even now obligatory or customary to explain the Catechism from the pulpit every Sunday afternoon. Hence the division of the questions into fifty-two Sundays, in imitation of the example set by Calvin's Catechism.\footnote{This division was first introduced in the Latin edition of 1566, perhaps earlier. Van Alpen, Niemeyer, and others are wrong in dating it from the German edition of 1573 or 1575.}

A Latin translation, for the use of colleges, was made by order of the Elector, by Joshua Lagus and Lambert Ludolph Pithopœus, and appeared soon after the German, since Olevianus sent a copy of each to Bullinger, in Zurich, as early as April, 1563.\footnote{Dœdes gives a fac-simile of the title-page of the Latin edition of 1563, from a copy in the University Library at Utrecht. It is nearly the same as the title of the edition of 1566, given in the literature above.} It is, however, much inferior to the German in force and unction. The Latin text was often edited separately as well as in the works of Ursinus, in connection with his commentary and other Latin commentaries, and in collections of Reformed symbols.\footnote{Niemeyer (pp. 428 sqq.) reproduces the edition of 1584, which agrees with the \emph{ed. princeps} of 1563 (as far as I can judge from the few fac-simile pages given by Dœdes), and with the text in the Oxford \emph{Sylloge,} while that in the Græco-Latin edition of Sylburg slightly differs. Dr. Louis H. Steiner, of Frederick City, Md., published an elegant and accurate edition under the title '\emph{Catechesis Religionis Christianæ seu Catechismus Heidelbergensis.} Baltimore, 1862.' He gives the variations of three Latin editions: of Cambridge, 1585; of Geneva, 1609 (formerly in the possession of Chevalier Bunsen); and the Oxford \emph{Sylloge,} 1804.}

There are three Dutch translations: the first appeared at Emden, 1563; the second, by Peter Dathenus, in connection with a Dutch version of the Psalter, in 1566, and very often separately.\footnote{On the Dutch translations, see especially the learned work of Professor Dœdes, of Utrecht, pp. 74–128, with fac-similes at the end of the volume.}

A Greek translation was prepared by a distinguished classical scholar, D. Frid. Sylburg, 1597.\footnote{I have before me a Græco-Latin edition of the Catechism \textgreek{κατηχήσεις τῆς χρισπανικῆς θρησκείας,} by Sylburg, and of the Belgic Confession by Jac. Revius, printed at Utrecht, 1660. Earlier editions I see noticed in catalogues.}

Besides these there are editions in modern Greek, in Hebrew, Arabic, etc.\footnote{Niemeyer (\emph{Proleg.} p. lxii.) mentions a Polish translation by \emph{Prasmovius,} a Hungarian by \emph{Scarasius,} an Arabic by \emph{Chelius,} a Singalese by \emph{Konyer,} besides French, Italian, Spanish, English, Bohemian, modern Greek, and Hebrew versions. Dœdes (p. 41) adds a Persian and a Malayan translation. There are no doubt many other versions.}

Three or four English translations were made from the Latin, and obtained a wide circulation in Scotland, England, and America.\footnote{An English edition, without the name of the translator, appeared A.D. 1591 at Edinburgh, 'by publick Authority, for the Use of Scotland,' and also repeatedly in connection with the 'Psalm-Book and the Book of Common Order.' It is embodied in Dunlop's \emph{Collection of Confessions of Faith,} etc., \emph{of publick authority in the Church of Scotland} (Edinburgh, 1719–1722), Vol. II. pp. 273–361, and reproduced by Dr. Horatius Bonar in his \emph{Catechisms of the Scottish Reformation} (London, 1866), pp. 112–170. Dr. Bonar says (p. 171): 'There are several translations of the Heidelberg or Palatine Catechism; and our Church [the Church of Scotland] seems not to have kept to one. In the edition of the Book of Common Order before us (1615), the Catechism is given alone; in that which Dunlop has followed, it has the ``Arguments'' and ``Uses'' of Bastingius.' Another translation by Bishop Henry Parry, of Worcester (d. 1616), appeared (together with the commentary of Ursinus) at Oxford, 1509 and 1601. It was often republished—at Edinburgh, 1615 (with sundry variations, see Bonar, p. 172), again in London, 1633, 1645, 1728, 1851, and quite recently (from the Oxford edition of 1601, with the variations of the edition of 1728) by Dr. Gerhart and Dr. Louis Steiner in the 'Mercersburg Review' for 1861, pp. 74 sqq. The one now in use in the Dutch and German Reformed Churches in America, is traced (by the late Dr. De Witt of New York) to Dr. Laidlie, originally from Scotland, minister at Flushing, Long Island, and was adopted, 1771, by the Synod of the Reformed Dutch Church. These three English translations seem to be only different recensions of one translation compared with the Latin text.} A more correct one from the German original was prepared for the tercentenary celebration of the Catechism, by a learned and able committee appointed by the German Reformed Synod in Pennsylvania, but has not yet come into public use.\footnote{See the tercentenary triglot edition of 1863, noticed in the literature above.}

The merits of the Latin and English translations, and their relation to the German original, may be seen from the following specimens:

The German Original, 1563.

The Latin Version, 1563.

\textbf{Frage 1. Was ist dein einiger Trost im Leben und im Sterben?}

\emph{Qu.} 1. \emph{Quæ est unica tua consolatio in vita et in morte?}

\textbf{Das ich mit Leib und Seele, beides im Leben und im Sterben, nicht mein, sondern meines getreuen Heilandes Jesu Christi eigen bin, der mit seinem theuren Blute für alle meine Sünden vollkommen bezahlet, und mich aus aller Gewalt des Teufels erlöset hat; und also bewahret, dass ohne den Willen meines Vaters im Himmel kein Haar von meinem Haupte kann fallen, ja auch mir alles zu meiner Seligkeit dienen muss. Darum er mich auch durch seinen heiligen Geist des ewigen Lebens versichert, und ihm forthin zu leben von Herzen willig und bereit macht.}

Quod animo pariter et corpore, sive vivam, sive moriar, non meus, sed fidissimi Domini et Servatoris mei Jesus Christi sum proprius, qui pretioso sanguine suo pro omnibus peccatis meis plenissime satisfaciens,\footnote{\par So also the Oxford \emph{Sylloge.} The \emph{ed. Græco-Latina} of Sylburg reads instead: \emph{plenissima solutione facta.}} me ab omni potestate diaboli liberavit, meque ita conservat, ut sine voluntate patris mei cœlestis, ne pilus quidem de meo capite posit cadere: imò verò etiam omnia saluti meæ servire oporteat. Quocirca me quoque suo Spiritu de vita æterna certum facit, utque ipsi deinceps vivam promptum ac paratum reddit.

\textbf{\emph{ Frage 2. Wie viele Stücke sind dir nöthig zu wissen, dass du in diesem Troste seliglich leben und sterben mögest?}}

\emph{Qu.} 2. \emph{Quot sunt tibi scitu necessaria, ut ista}\footnote{Al. edd. \emph{illa.}} \emph{consolatione fruens, beatè vivas et moriaris?}

\textbf{Drei Stücke: Erstlich, wie gross meine Sünde und Elend sei. Zum Andern, wie ich von allen meinen Sünden und Elend erlöset werde. Und zum Dritten, wie ich Gott für solche Erlösung soll dankbar sein.}

Tria. Primum, quanta sit peccati mei et miseriæ meæ magnitudo. Secundum,\footnote{Al. \emph{Alterum.}} quo pacto ab omni peccato et miseria liberer. Tertium, quam gratiam Deo pro ea liberatione debeam.

Scotch Edition of 1591.

Bishop Perry's Translation (1591).

\emph{From Dunlop's Collection }(1722).

\emph{Oxford Edition of} 1601.

\emph{Ques.} 1. \emph{What is thy only comfort in life and in death?}

\emph{Ques.} 1. \emph{What is thy only comfort in lift and death?}

That in soul and body, whether I live or die, I am not mine own, but I belong unto my most faithful Lord and Saviour, Jesus Christ: who by his precious blood, most fully satisfying for all my sins, hath delivered me from the whole power of the Devil; and doth so preserve me, that without the will of my heavenly Father, not so much as a hair can fall from my head: yea, all things are made to serve for my salvation. Wherefore by his Spirit also, he assureth me of everlasting life, and maketh me ready and prepared, that henceforth I may live unto him.

That both in soul and body, whether I live or die, I am not mine own, but belong wholly\footnote{The redundant 'wholly' occurs also in the Edinburgh edition of 1615, which, to judge from the specimens given by Horatius Bonar (in \emph{Catechisms of the Scottish Reformation,} p. 172), is a reprint of Parry's translation with a few variations.} unto my most faithful Lord and Saviour Jesus Christ, who by his precious blood most fully satisfying for all my sins, hath delivered me from all the power of the devil, and so preserveth me, that without the will of my heavenly Father not so much as a hair may fall from my head, yea all things must serve for my safety. Wherefore by his Spirit also he assureth me of everlasting life, and maketh me ready, and prepared, that henceforth I may live to him.

\emph{Ques.} 2. \emph{How many things are needful for thee to know, to the end} [\emph{that}] \emph{thou, enjoying this comfort, mayest live and die an happy man?}

\emph{Ques.} 2. \emph{How many things are necessary for thee to know, that thou enjoying this comfort mayest live and die happily?}

Three things. First, What is the greatness of my sin, and of my misery. Secondly, By what means I may be delivered from all my sin and misery. Thirdly, What thankfulness I owe to God for that deliverance.

Three. The first, what is the greatness of my sin and misery. The second, how I am delivered from all sin and misery. The third, what thanks I owe unto God for this delivery.

The Received American Version, 1771.

The New American Version, 1863.

\emph{Ques.} 1. \emph{What is thy only comfort in life and death?}

\emph{Ques.} 1. \emph{What is thy only comfort in life and in death?}

That I with body and soul, both in life and death, am not my own, but belong unto my faithful Saviour Jesus Christ, who, with his precious blood, hath fully satisfied for all my sins, and delivered me from all the power of the devil; and so preserves me that without the will of my heavenly Father, not a hair can fall from my head; yea, that all things must be subservient to my salvation; and therefore, by his Holy Spirit, he also assures me of eternal life, and makes me sincerely willing and ready henceforth, to live unto him.

That I, with body and soul, both in life and in death, am not my own, but belong to my faithful Saviour Jesus Christ, who with His precious blood has fully satisfied for all my sins, and redeemed me from all the power of the devil; and so preserves me, that without the will of my Father in heaven not a hair can fall from my head; yea, that all things must work together for my salvation. Wherefore, by His Holy Spirit, He also assures me of eternal life, and makes me heartily willing and ready henceforth to live unto Him.

\emph{Ques.} 2. \emph{How many things are necessary for thee to know, that thou, enjoying this comfort, mayest live and die happily?}

\emph{Ques.} 2. \emph{How many things are necessary for thee to know, that thou in this comfort mayest live and die happily?}

Three; the first, how great my sins and miseries are; the second, how I may be delivered from all my sins and miseries; the third, how I shall express my gratitude to God for such deliverance.

Three things: First, the greatness of my sin and misery. Second, how I am redeemed from all my sins and misery. Third, how I am to be thankful to God for such redemption.

Note.—All the English versions, except the last, follow the Latin in its departures from the German, as '\emph{most} faithful Lord' (\emph{fidelissimi Domini}) for 'faithful' (\emph{getreuen}), 'heavenly Father' (\emph{Patris cœlestis}) for 'Father in heaven' (\emph{Vater im Himmel}). The dependence on the Latin may be seen also in the words 'most fully satisfying' (\emph{plenissime satisfaciens}), 'delivered' (\emph{liberavit}) for 'redeemed' (\emph{erlöset}), 'delivery' (\emph{liberatio}) for 'redemption' (\emph{Erlösung}) and in the omission of 'heartily' (\emph{von Herzen}), for which, however, the common American version (which seems to have made use also of the Dutch version) substitutes 'sincerely.'

CHARACTER AND AIM.

The Heidelberg Catechism answers the double purpose of a guide for the religious instruction of the youth and a confession of faith for the Church.

As a catechism it is an acknowledged masterpiece, with few to equal and none to surpass it. Its only defect is that its answers are mostly too long for the capacity and memory of children. It is intended for a riper age. Hence an abridgment was made as early as 1585, but no attempts to simplify and popularize it have been able to supersede it.

As a standard of public doctrine the Heidelberg Catechism is the most catholic and popular of all the Reformed symbols. The German Reformed Church acknowledges no other. The Calvinistic system is herein set forth with wise moderation, and without its sharp, angular points. This may be a defect in logic, but it is an advantage in religion, which is broader and deeper than logic. Children and the mass of the people are unable to appreciate metaphysical distinctions and the transcendent mysteries of eternal decrees. The doctrine of election to holiness and salvation in Christ (or the positive and edifying part of the dogma of predestination) is indeed incidentally set forth as a source of humility, gratitude, and comfort (Ques. 1, 31, 53, 54), but nothing is said of a \emph{double} predestination, or of an eternal decree of \emph{reprobation,} or of a \emph{limited} atonement (comp. Ques. 37). These difficult questions are left to private opinion and theological science. This reserve is the more remarkable since the authors (as well as all other Reformers, except Melanchthon in his later period) were strict predestinarians.

PLAN AND ARRANGEMENT.

The Heidelberg Catechism follows the order of the Epistle to the Romans, and is divided into three parts. The first two questions are introductory. The first part treats of the sin and misery of man (Ques. 3–11; comp. Rom. i. 18-iii. 20); the second of the redemption by Christ (Ques. 12–85; comp. Rom. iii. 21-xi. 36); the third of the thankfulness of the redeemed, or the Christian life (Ques. 86–129; comp. Rom. xii.-xvi.). The second part is the largest, and contains an explanation of all the articles of the Apostles' Creed under the three heads of God the Father, God the Son, and God the Holy Ghost. The doctrine of the sacraments is rightly incorporated in this part, instead of being treated in separate sections, as in the Roman and Lutheran Catechisms. The third part gives an exposition of the Decalogue (as a rule of obedience, viewed in the light of redemption) and of the Lord's Prayer.

This order corresponds to the development of religious life and to the three leading ideas of repentance, faith, and love. The conception of Christian life, as an expression of gratitude for redeeming grace, is truly evangelical. In older catechisms the five or six parts of a catechism—namely, the Creed, the Decalogue, the Lord's Prayer, Baptism, the Lord's Supper—are mechanically co-ordinated; here they are worked up into an organic system.

The execution is admirable throughout. Several answers are acknowledged gems in the history of catechetical literature—e.g., the definition of faith (Ques. 21), on providence (Ques. 27 and 28), on the significance of the Christian name (Ques. 31 and 32), on the benefit of the ascension (Ques. 49), and on justification by faith (Ques. 60).

THE SPIRIT OF THE CATECHISM.

The genius of the Catechism is brought out at once in the first question, which contains the central idea, and strikes the key-note. It is unsurpassed for depth, comfort, and beauty, and, once committed to memory, can never be forgotten. It represents Christianity in its evangelical, practical, cheering aspect, not as a commanding law, not as an intellectual scheme, not as a system of outward observances, but as the best gift of God to man, as a source of peace and comfort in life and in death. What can be more comforting, what at the same time more honoring and stimulating to a holy life than the assurance of being owned wholly by Christ our blessed Lord and Saviour, who sacrificed his own spotless life for us on the cross? The first question and answer of the Heidelberg Catechism is the whole gospel in a nutshell; blessed is he who can repeat it from the heart and hold it fast to the end.\footnote{Dr. Nevan \{\emph{Tercentenary Edition,} Introd. p. 95) says: 'No question in the whole Catechism has been more admired than this, and none surely is more worthy of admiration. Where shall we find, in the same compass, a more beautifully graphic, or a more impressively full and pregnant representation of all that is comprehended for us in the grace of our Lord and Saviour Jesus Christ? For thousands and tens of thousands, during the past three hundred years, it has been as a whole system of theology in the best sense of the term, their pole-star over the sea of life, and the sheet-anchor of their hope amid the waves of death. But what we quote it for now is simply to show the mind that actuates and rules the Catechism throughout. We have here at once its fundamental conception and the reigning law of its construction; the key-note, we may say, which governs its universal sense, and whose grandly solemn tones continue to make themselves heard through all its utterances from beginning to end.'}

It would be difficult to find a more evangelical definition of faith than in Ques. 21: 'Faith is not only a certain knowledge, whereby I hold for truth all that God has revealed to us in his Word; but also a hearty trust, which the Holy Spirit works in me by the gospel, that not only to others, but to me also, forgiveness of sins, everlasting righteousness, and salvation are freely given by God, merely of grace, only for the sake of Christ's merits.' How rich and consoling is the lesson derived from God's all-ruling Providence in Ques. 28! 'That we may be patient in adversity, thankful in prosperity, and for what is future have good confidence in our faithful God and Father, that no creature shall separate us from his love, since all creatures are so in his hand that without his will they can not so much as move.'

The Catechism is a work of religious enthusiasm, based on solid theological learning, and directed by excellent judgment. It is baptized with the pentecostal fire of the great Reformation, yet remarkably free from the polemic zeal and intolerance which characterized that wonderfully excited period—by far the richest and deepest in Church history next to the age of Christ and his inspired apostles. It is the product of the heart as well as the head, full of faith and unction from above. It is fresh, lively, glowing, yet clear, sober, self-sustained. The ideas are Biblical and orthodox, and well fortified by apt Scripture proofs.\footnote{Ques. 44 is hardly an exception; for the idea therein expressed is no error \emph{per se,} but only a false interpretation of the article on Christ's descent into hell (Hades) in the Apostles' Creed, which places it, as an actual fact, between death and the resurrection, in accordance with the Scriptures (Luke xxiii. 43; Acts ii. 27, 31; 1 Pet. iii. 19; iv. 6; Eph. iv. 9, 10); while the Catechism, following Calvin and Lasky, understands it figuratively of Christ's suffering on the cross.} The language is dignified, terse, nervous, popular, and often truly eloquent. It is the language of devotion as well as instruction. Altogether the Heidelberg Catechism is more than a book, it is an institution, and will live as long as the Reformed Church.

COMPARISON WITH THE LUTHERN AND WESTMINSTER CATECHISMS.

The Heidelberg Catechism stands mediating between Luther's Small Catechism, which appeared thirty-four years earlier (1529), and the Shorter Westminster Catechism, which was prepared eighty-four years later (1647).

These are the three most popular and useful catechisms that Protestantism has produced, and have still the strongest hold upon the churches they represent. They have the twofold character of catechisms and symbolical books. They are alike evangelical in spirit and aim; they lead directly to Christ as the one and all-sufficient Saviour, and to the Word of God as the only infallible rule of the Christian's faith and life.

Luther's Catechism is the most churchly of the three, and adheres to the Catholic tradition in its order and arrangement. It assigns a very prominent place to the Sacraments, treating them in separate chapters, co-ordinate with the Decalogue, the Creed, and the Lord's Prayer; while the others incorporate them in the general exposition of the articles of faith. Luther teaches baptismal regeneration and the corporeal presence, and even retains private confession and absolution as a quasi-sacrament. Heidelberg and Westminster are free from all remnants of sacerdotalism and sacramentalism, and teach the Calvinistic theory of the sacraments, which rises, however, much higher than the Zwinglian.

On the other hand, the Lutheran and the Heidelberg Catechisms differ from the Westminster in the following points: 1. They retain the Apostles' Creed as the basis of doctrinal exposition; while the Westminster Catechism puts it in an appendix, and substitutes a new logical scheme of doctrine for the old historical order of the Creed. 2. They are subjective, and address the catechumen as a Church member, who answers from his real or prospective personal experience; while the Westminster Catechism is objective and impersonal, and states the answer in an abstract proposition. 3. They use the warm and direct language of life, the Westminster the scholastic language of dogma; hence the former two are less definite but more expansive and suggestive than the Presbyterian formulary, which, on the other hand, far surpasses them in brevity, terseness, and accuracy of definition.

Upon the whole we prefer the catechetical style and method of the creative Reformation period, because it is more Biblical and fresh, to that of the seventeenth century—the age of scholastic orthodoxy—although we freely concede the relative progress and peculiar excellences of the Westminster standard.\footnote{'It may be questioned,' says Dr. Bonar, of the Free Church of Scotland, 'whether the Church gained any thing by the exchange of the Reformation standards for those of the seventeenth century. The scholastic mold in which the latter are cast has somewhat trenched upon the ease and breadth which mark the former; and the skillful metaphysics employed at Westminster in giving lawyer-like precision to each statement have imparted a local and temporary aspect to the new which did not belong to the more ancient standards. Or, enlarging the remark, we may say that there is something about the theology of the Reformation which renders it less likely to become obsolete than the theology of the covenant. The simpler formulas of the older age are quite as explicit as those of the later; while by the adoption of the Biblical in preference to the scholastic mode of expression they have secured for themselves a buoyancy which will bear them up when the others go down. The old age of that generation is likely to be greener than that of their posterity.' (\emph{Catechisms of the Scottish Reformation,} Preface, p. viii.)}

The Heidelberg Catechism differs from that of Luther—1. By its fullness and thoroughness, and hence it is better adapted to a maturer age; while that of Luther has the advantage of brevity and childlike simplicity, and adaptation to early youth. The one has one hundred and twenty-nine, the other only forty questions and answers, and of these only three are devoted to the exposition of the Apostles' Creed, while the Sacraments receive disproportionate attention. 2. The Heidelberg Catechism gives the words of the Decalogue in full, according to the twentieth chapter of Exodus, and follows the old Jewish and Greek division, which is adopted by the best commentators; while Luther presents merely an abridgment,\footnote{For example, the fourth (third) commandment is thus condensed: '\emph{Du sollst den Feiertag heiligen}' (Thou shalt keep holy the rest-day).} and follows the Roman division by omitting the second commandment and splitting the tenth into two.\footnote{Comp. p. 251, note 2.} 3. The former gives a summary of the law, through which comes the knowledge of sin, in the first part (Ques. 3 and 4), but explains the Decalogue in the third division, viewing it in its Christian aspect as a permanent rule of life; while Luther regards the law in its Jewish or pedagogic aspect, as a schoolmaster leading men to Christ, and hence he puts it as the first head before the Creed. Ursinus correctly says: 'The Decalogue belongs to the first part so far as it is a mirror of our sin and misery, but also to the third part as being the rule of our new obedience and Christian life.'\footnote{The Germans express the different aspects of the law by calling it a \emph{Sündenspiegel, } \emph{Sündenriegel,} and \emph{Lebensregel,} a mirror of sin, a bar of sin, and a rule of life.} 4. In the rendering of the Creed, besides minor verbal differences, the Heidelberg Catechism retains 'the holy catholic Church,' with the addition of 'Christian' (\emph{eine heilige allgemeine christliche Kirche}); while Luther's omits 'catholic,' and substitutes for it 'Christian.'\footnote{Hence in Germany the term 'Catholic' and 'Romanist' are used synonymously, and the proverb '\emph{Das ist um katholisch zu werden}' expresses a desperate condition of things. The English Churches have properly retained the term 'catholic' in its good old sense, instead of allowing Romanists to monopolize it.} 5. In the Lord's Prayer the Heidelberg Catechism uses the modern form 'Our Father' (\emph{Unser Vater}), while Luther in his Catechism (though not in his translation of Matt. vi. 9 and Luke xi. 2) adheres to the Latin and old German form of 'Father our' (\emph{Vater unser}), a difference tenaciously maintained by German Lutherans. The former divides the Prayer into six petitions (with the Greek commentators), and renders \textgreek{ἐκ πονηροῦ} 'from the evil one' (\emph{vom Bösen,} i.e., from the devil); while Luther (with Augustine) numbers seven petitions, and translates (herein agreeing with the English version) 'from evil' (\emph{vom Uebel}).

The difference between the Heidelberg and Westminster Catechisms is chiefly one of nationality. Where the choice is between the two, the former will be used in preference by Germans, the other by Scotch and English Presbyterians. The Westminster Shorter Catechism has the advantage of greater condensation and precision. It is not impossible to make a better one than either by blending the excellences of both. They represent also two types of piety: the one is more emotional and hearty, the other more scholastic and intellectual. This appears at once in the first question. The Heidelberg Catechism asks: 'What is thy only \emph{comfort} in life and in death?' The Westminster: 'What is the chief \emph{end} of man?' The one goes at once into the heart of evangelical piety—the mystical union of the believer with Christ; the other goes back to the creation and the glory of God; but both teach the same God and Christ, and the same way of salvation, whereby God is glorified, and man is raised to everlasting felicity in his enjoyment.

HISTORY OF THE CATECHISM.

1. The Heidelberg Catechism was greeted with great joy, and was at once introduced into the churches and schools of the Lower Palatinate; while the Upper Palatinate, under the governorship of Louis (the eldest son of Frederick III.) remained strictly Lutheran.

But, like every good book, it had to pass through a trial of probation and a fire of martyrdom. Even before it was printed an anonymous writer attacked the Heidelberg Synod which, in December, 1562, had adopted the Catechism in manuscript, together with sundry measures of reform.\footnote{This curious document, which throws light upon that Synod hitherto little known, has been recently recovered and published by Wolters in the \emph{Studien mud Kritiken} for 1867, No. 1, pp. 15 sqq. The Lutheran author, perhaps a dissenting member of the Synod, gives a list of the measures for the introduction of the Catechism and the abolition of various abuses, and accompanies them with bitter marginal comments, such as: 'This is a lie and against God's Word;' 'This is the Anabaptist heresy;' 'To spread Zwinglianism;' '\emph{Friss Vogel oder stirb};' '\emph{Ad spargendam zizaniam};' '\emph{Ut citius imbibant venenum};' '\emph{Evangelii abrogatio};' '\emph{Hispanica inquisitio.}'} After its publication it was violently assailed by strict Lutherans for its alleged Zwinglian and Calvinistic heresies, and by Jesuits on account of the condemnation of the idolatry of the mass in the eightieth question. The first opponents were Lutheran princes (Margrave Charles II. of Baden, Duke Christopher of Würtemberg, the Palatine of Zweibrücken), and Lutheran divines, such as Heshusius, Flacius, Brentius, and Andreæ.\footnote{See on this Lutheran opposition Wolters, l.c., and in his earlier book, \emph{Der Heidelb. Katechismus in seiner Urgestalt} (1864), pp. 141–196; Nevin, Introd. to the \emph{Tercent. Ed.} pp. 42 sqq.; and especially Sudhoff, \emph{Olevianus und Ursinus,} pp. 140 sqq.} Ursinus wrote an able apology of his Catechism, which is embodied in several older editions since 1584. A theological colloquy was held at Maulbronn in April, 1564, where the theological leaders of the Lutheran Duchy of Würtemberg and the Reformed Palatinate, in the presence of their princes, debated for six days in vain on the eucharist and the ubiquity of Christ's body. Both parties were confirmed in their opinions, though the Reformed had the best of the argument.\footnote{See above, pp. 288 sqq.}

Frederick III., notwithstanding his appeal to Melanchthon and the Altered Augsburg Confession, was openly charged with apostasy from the Lutheran faith, and seriously threatened with exclusion from the peace of the empire. Even the liberal Emperor Maximilian II. wrote him a letter of remonstrance. His fate was to be decided at the Diet of Augsburg, 1566. At this critical juncture the pious Elector boldly defended his Catechism, which, he said, was all taken from the Bible, and so well fortified with marginal proof-texts that it could not be overthrown. He declared himself willing to yield to God's truth, if any one could show him any thing better from the Scripture, which was at hand for the purpose. Altogether he made, at the risk of his crown and his life, such a noble and heroic confession as reminds us of Luther's stand at the Diet of Worms. Even his Lutheran opponents were filled with admiration and praise, and left him thereafter in quiet possession of his faith. 'Why do ye persecute this man?' said the Margrave of Baden; 'he has more piety than the whole of us.' The Elector Augustus of Saxony gave similar testimony on this memorable occasion.\footnote{Hundeshagen says of Frederick III.: 'He is acknowledged to be the greatest ruler which the evangelical Palatinate ever had, and as to personal piety and loyalty to his faith the shining model of an evangelical prince.' See his art. on the City and University of Heidelberg, in the \emph{Gedenkbuch der} 300 \emph{jähr. Jubelfeier des Heidelb. Kat.} pp. 58, 59.}

Thus the Catechism had gained a sort of legal existence in the German empire, although it was not till after the Thirty-Years' War, in the Treaty of Westphalia, that the Reformed Church, as distinct from the Lutheran, was formally recognized in Germany.

After the death of Frederick it had to pass through another persecution in the home of its birth. His successor, Louis VI. (1576–1583), exiled its authors, and replaced it by Luther's Catechism and the Formula of Concord. But under the regency of Frederick's second son, Prince John Casimir, the Heidelberg Catechism and the Reformed Church were restored to their former honor, and continued to flourish till the outbreak of the Thirty-Years' War.

This war brought terrible devastation and untold misery upon Heidelberg and the Palatinate, which were laid waste by the merciless Tilly (1622). Then followed the repeated invasions of Turenne, Melac, and Marshal de Lorges, under Louis XIV. The Palatinate fell even into the hands of Roman Catholic rulers (1685), and never again rose to its former glory. Thousands of Protestants emigrated to America, and planted the Catechism in Pennsylvania, so that what it lost in the old world it gained in the new. The indifferentism and rationalism of the eighteenth century allowed all creeds to go into disuse and neglect. In the nineteenth century faith revived, and with it respect for the Heidelberg Catechism; but, owing to the introduction of the union of the Lutheran and Reformed Churches in the Grand Duchy of Baden, to which Heidelberg now belongs, it was merged into a new catechism compiled from it and from that of Luther.\footnote{On the symbolical status of the Evangelical Church in Baden, see two essays of Dr. Hundeshagen, \emph{Die Bekenntnissgrundlage der vereinigten evangelischen Kirche im Grossherzogthum Baden} (1851), and an address delivered before a Pastoral Conference at Durlach, on the same subject, 1851, republished in his \emph{Schriften und Abhandlungen,} ed. by Dr. Christlieb, Gotha, 1875, Vol. II. pp. 119 sqq.}

2. The history of the Palatinate Catechism extends far beyond the land of its birth. It took deeper root and acquired greater influence in other countries. Soon after its appearance it commended itself by its intrinsic excellences to all Reformed Churches of the German tongue. It was introduced in East Friesland, Jülich (Juliers), Cleve (Cleves), Berg, the Wupperthal, Bremen, Hesse Cassel, Anhalt, Brandenburg, East and West Prussia, the free imperial cities, in Hungary, Poland, and in several cantons of Switzerland, as St. Gall, Schaffhausen, and Berne.\footnote{The editions used in the Canton Berne have an anti-supralapsarian addition to Question 27: '\emph{Und obwohl die Sünden durch Gottes Fürsehung werden regiert, so ist doch Gott keine Ursache der Sünde; denn das Ziel unterscheidet die Werke. Siehe Exempel an Joseph und seinen Brüdern, an David und Simei, an Christo und den Juden.}' This addition is found as early as 1697. Noticed by Trechsel in \emph{Studien und Kritiken} for 1867, p. 574.} In the royal house of Prussia it is still used in the instruction of the princes, even after the introduction of the union of the two confessions.\footnote{So I was informed by the late court chaplain, Dr. Snethlage, of Berlin, who was originally Reformed, and who confirmed several members of the royal family.}

It was surrounded with a large number of learned works which fill an important place in the history of Reformed theology. Eminent professors made it the basis of lectures in the University.

In no country was the Catechism more honored than in Holland and her distant colonies in Asia and Africa. It soon replaced the catechisms of Calvin and Lasky. The synods of Wesel, 1568, of Emden, 1571, and of Dort, 1574, recommended and enjoined its use; and ministers were required to explain it to the people in fifty-two lessons throughout the year in the afternoon service of the Lord's day. In the beginning of the sixteenth century the Arminians called for a revision of it, to remove certain features to which they objected. But the famous General Synod of Dort, after a careful examination, opposed any change, and, in its 148th Session, May 1, 1619, it unanimously delivered the judgment that the Heidelberg Catechism 'formed altogether a most accurate compend of the orthodox Christian faith; being, with singular skill, not only adapted to the understanding of the young, but suited also for the advantageous instruction of older persons; so that it could continue to be taught with great edification in the Belgic churches, and ought by all means to be retained.' This judgment was agreed to by all the foreign delegates from Germany, Switzerland, and England, and has thus an œcumenical significance for the Reformed communion.

The Heidelberg Catechism was also clothed with symbolical authority in Scotland, and was repeatedly printed 'by public authority,' even after the Westminster standards had come into use. It seems to have there practically superseded Calvin's Catechism, but it was in turn superseded by Craig's Catechism, and Craig's by that of the Westminster Assembly.

3. From Holland the Heidelberg Catechism crossed the Atlantic to Manhattan Island (1609), with the discoverer of the Hudson River, and was the first Protestant catechism planted on American soil. A hundred years later, German emigrants, driven from the Palatinate by Romish persecution and tyranny, carried it to Pennsylvania and other colonies. It has remained ever since the honored symbol of the Dutch and German Reformed Churches in America, and will continue to be used as long as they retain their separate denominational existence, or even if they should unite with the larger Presbyterian body.

One of the first acts of the reunited Presbyterian Church in the United States, at the session of the General Assembly in Philadelphia, May, 1870, was the formal sanction of the use of the Heidelberg Catechism in any congregation which may desire it.\footnote{A special committee, appointed by the Old School Assembly of 1869, reported to the first reunited Assembly of 1870, after a laudatory description of the Heidelberg Catechism, the following resolutions, which were unanimously adopted: \par 1. \emph{Resolved,} That this General Assembly recognizes in the Heidelberg Catechism a valuable Scriptural compendium of Christian doctrine and duty. \par 2. \emph{Resolved,} That if any churches desire to employ the Heidelberg Catechism in the instruction of their children, they may do so with the approbation of this Assembly. \par See the \emph{Minutes of the General Assembly of the Presbyterian Church in the United States of America for} 1870, p. 120, and the Memorial volume on \emph{Presbyterian Reunion} (New York, 1870), p. 454.}

4. In the year 1863, three centuries after its first publication, the Heidelberg Catechism witnessed its greatest triumph, not only in Germany and Holland, but still more in a land which the authors never saw, and in a language the sound of which they probably never heard. The Reformation was similarly honored in 1817, and the Augsburg Confession in 1830, but no other catechism.

In Germany the tercentenary celebration of the Heidelberg Catechism was left to individual pastors and congregations, and called forth some valuable publications.\footnote{Among these we mention the articles on the Heidelberg Catechism by Ullmann, Sack, Plitt, Hundeshagen, Wolters, and Trechsel, in the \emph{Studien und Kritiken} for 1863, 1864, and 1867, the discovery and reprint of the \emph{ed. princeps} by Wolters (1864), and a collection of excellent sermons by distinguished Reformed pulpit orators, under the title, '\emph{Der einzige Trost im Leben und Sterben},' Elberfeld, 1863.}

The German Reformed Church in the United States took it up as a body, and gave it a wider scope. She made the three-hundredth anniversary of her confession the occasion for a general revival of theological and religious life, the publication of a triglot edition of the Catechism, the endowment of a tercentenary professorship in her seminary, and the collection of large sums of money for churches, missions, and other benevolent objects. All these ends were accomplished. The celebration culminated in a general convention of ministers and laymen in Philadelphia, which lasted a whole week, January 17–23, 1863, in the midst of the raging storm of the civil war. About twenty interesting and instructive essays on the Catechism and connected topics, which had been specially prepared for the occasion by eminent German, Dutch, and American divines, were read in two churches before crowded and attentive assemblies. Luther, Calvin, Zwingli, Melanchthon, Frederick III., Ursinus, and Olevianus were called from their graves to reproduce before an American audience the ideas, trials, and triumphs of the creative and heroic age of the Reformation. Altogether the year 1863 marks an epoch in the history of the Heidelberg Catechism and of the German Reformed Church in America.\footnote{See the \emph{Tercentenary Monument} (574 pages), and the \emph{Gedenkbuch der dreihundert jährigen Jubelfeier des Heidelberger Katechismus} (449 pages), both published at Philadelphia. 1863. The German edition gives the correspondence and essays of Drs. Herzog, Ebrard, Ullmann, Hundeshagen, Lange, and Schotel, in the original German, together with a history of the Catechism by the editor. The Anglo-American essays and addresses of Drs. Nevin, Schaff, Gerhart, Harbaugh, Wolff, Bomberger, Porter, De Witt, Kieffer, Theodor and Thomas Appel, Schneck, Russell, Gans, and Bausmann, are found in full in the English edition.}

OPINIONS ON THE CATECHISM.

We close this chapter with a selection from the many warm commendations which the Heidelberg Catechism has received from distinguished divines of different countries.

Henry Bullinger, the friend and successor of Zwingli, himself the author of a catechism (1559) and of the Second Helvetic Confession (1566), wrote to a friend:

'The order of the book is clear; the matter true, good, and beautiful; the whole is luminous, fruitful, and godly; it comprehends many and great truths in a small compass. I believe that no better catechism has ever been issued.'\footnote{\par '\emph{Arbitror meliorem Catechismum non editum esse. Deo sit gloria qui largiatur successum}' (1563). See Ursinus, \emph{Apol. Catech.} in the \emph{Præfatio.}}

The Hessian divines quoted by David Pareus:

'There is no catechism more thorough, more perfect, and better adapted to the capacity of adults as well as the young.'

The English delegates to the Synod of Dort, George Carleton (Bishop of Llandaff), John Davenant (afterwards Bishop of Salisbury), Archdeacon Samuel Ward, Dr. Thomas Goade, and Walter Balcanqual, said:

'That neither their own nor the French Church had a catechism so suitable and excellent; that those who had compiled it were therein remarkably endowed and assisted by the Spirit of God; that in several of their works they had excelled other theologians, but that in the composition of this Catechism they had outdone themselves.'\footnote{This judgment is quoted on the title-page of the later editions of Bishop Parry's translation, London ed. 1728; reprinted, London, 1851.}

The favorable judgment of the Synod of Dort itself has already been quoted.

Dr. Ullmann (d. 1865), formerly Professor at Heidelberg, and one of the best Church historians of the nineteenth century:\footnote{In Piper's \emph{Evang. Kalender} for 1862, p. 191. Comp. also his art. in the \emph{Studien und Kritiken} for 1863, and in the \emph{Gedenkbuch,} etc.}

'The Heidelberg Catechism, more systematically executed than Luther's, unfolds upon the fundamental thoughts of sin, redemption, and thankfulness, the Reformed doctrine, yet without touching upon predestination, with rare pithiness and clearness, and obtained through these excellences not only speedy and most extended recognition in the Reformed Churches, but is to-day still regarded by all parties as one of the most masterly productions in this department.'

Dr. Aug. Ebrard, one of the ablest and most prolific German Reformed divines:\footnote{\emph{Das Dogma v. heil. Abendmahl,} Vol. II. p. 604.}

'For wonderful union of dogmatic precision and genial heartiness,\footnote{Or, fullness of soul (\emph{gemüthliche Innigkeit}).} of lucid perspicuity and mysterious depth, the Heidelberg Catechism stands alone in its kind. It is at once a system of theology and a book of devotion; every child can understand it at the first reading, and yet the catechist finds in it the richest material for profound investigation.'

Max Göbel, the author of an excellent history of Christian life in the Reformed Church:\footnote{\emph{Geschichte des christl. Lebens,} Vol. I. p. 392.}

'The Heidelberg Catechism may be properly regarded as the flower and fruit of the entire German and French Reformation; it has Lutheran fervor, Melanchthonian clearness, Zwinglian simplicity, and Calvinistic fire blended in one, and therefore—notwithstanding many defects and angles—it has been (together with the Altered Augsburg Confession of 1540), and remains to this day, the only common confession and doctrinal standard of the entire German Reformed Church from the Palatinate to the Netherlands, and to Brandenburg and Prussia.'

Karl Sudhoff, formerly a Roman Catholic priest, then pastor of the German Reformed Church at Frankfort-on-the-Main:\footnote{\emph{Theol. Handbuch zur Auslegung des Heid. Kat.} p. 493.}

'A peculiar power and unction pervades the whole work, which can not easily be mistaken by any one. The book, therefore, speaks with peculiar freshness and animation directly to the soul, because it appears as a confident, joyous confession of the Christian heart assured of salvation. It is addressed to the heart and will as much as to the head. Keen and popular unfolding of ideas is here most beautifully united with the deep feeling of piety, as well as with the earnest spirit of revival and joyous believing confidence. And who that have read this Catechism but once can mistake how indissolubly united with these great excellences is the powerful, dignified, and yet so simple style! What a true-hearted, intelligible, simple, and yet lofty eloquence speaks to us even from the smallest questions!'

Dr. K. B. Hundeshagen, Professor of Theology at Heidelberg, afterwards in Bonn (d. 1873), calls the Heidelberg Catechism a 'witness of Reformed loyalty to the Word of God, of Reformed purity and firmness of faith, of Reformed moderation and sobriety,' and a work 'of eternal youth and never-ceasing value.'\footnote{See his instructive review of Sudhoff's \emph{Handbuch,} in the \emph{Studien und Kritiken} for 1864, pp. 153–180. It is gratifying to me that this distinguished divine fully indorses, on p. 169, the view which I had previously given of the theology of the Heidelberg Catechism and its relation to Calvinism in opposition to Sudhoff on the one hand and Heppe on the other.}

Dr. Plitt, formerly Pastor in Heidelberg, then Professor of Theology in Bonn:\footnote{In the \emph{Studien und Kritiken} for 1863, p. 25.}

'The Heidelberg Catechism still lives; it has not died in three hundred years. It lives in the hearts of Christians. How many catechisms have since then disappeared, how many in the last thirty or forty years, and have been so long sunk in the ``sea of oblivion,'' that one scarcely knows their titles. The Heidelberg Catechism has survived its tercentenary jubilee, and will, God willing, see several such jubilees. It will not die; it will live as long as there is an Evangelical Church.'

Dr. Henry Harbaugh, late Professor of Theology at Mercersburg (d. 1867), a gifted poet and the author of several popular religious works:\footnote{In the \emph{Mercersburg Review} for 1857, p. 102.}

'It is worthy of profound consideration, that the Heidelberg Catechism, which has always ruled the heart, spirit, and body of the Reformed side of the Reformation, has no prototype in any of the Reformers. Zwingli and Calvin can say, It is not of me; it has the suavity but not the compromising spirit of Melanchthon. It has nothing of the dashing terror of Luther. What is stranger than all, it is farthest possible removed from the mechanical scholasticism and rigid logic of Ursinus, its principal author. Though it has the warm, practical, sacred, poetical fervor of Olevianus, it has none of his fire and flame. It is greater than Reformers; it is purer and sounder than theologians.'

Dr. J. W. Nevin, successively Professor of Theology in the Presbyterian Seminary at Alleghany, in the German Reformed Seminary at Mercersburg, and President of Franklin and Marshall College, Lancaster, Pa.:\footnote{\emph{Tercentenary Edition,} Introd. pp. 120–122.}

'In every view, we may say, the Catechism of the Palatinate, now three hundred years old, is a book entitled, in no common degree, to admiration and praise. It comes before us as the ripe product of the proper confessional life of the Reformed Church, in the full bloom of its historical development, as this was reached at the time when the work made its appearance. Its wide-spread and long-continued popularity proclaims its universal significance and worth. It must have been admirably adapted to the wants of the Church at large, as well as admirably true to the inmost sense of its general life, to come in this way into such vast credit. Among all Protestant symbols, whether of earlier or later date, there is no other in which we find the like union of excellent qualities combined and wrought together in the same happy manner. It is at once a creed, a catechism, and a confession; and all this in such a manner, at the same time, as to be often a very liturgy also, instinct with the full spirit of worship and devotion. It is both simple and profound; a fit manual of instruction for the young, and yet a whole system of divinity for the old; a text-book, suited alike for the use of the pulpit and the family, the theological seminary, and the common school. It is pervaded by a scientific spirit, beyond what is common in formularies of this sort; but its science is always earnestly and solemnly practical. In its whole constitution, as we have seen, it is more a great deal than doctrine merely, or a form of sound words for the understanding. It is doctrine apprehended and represented continually in the form of life. It is for the heart every where full as much as for the head. Among its characteristic perfections deserves to be noted always, with particular praise, its catholic spirit, and the rich mystical element that pervades so largely its whole composition. . . . Simple, beautiful, and clear in its logical construction, the symbol moves throughout also in the element of fresh religious feeling. It is full of sensibility and faith and joyous childlike trust. Its utterances rise at times to a sort of heavenly pathos, and breathe forth almost lyrical strains of devotion.'

Dr. Hagenbach, the well-known historian (d. at Basle, 1874):\footnote{\emph{Kirchengeschichte,} Leipz. 1870 (3d edition), Vol. IV. p. 312.}

'The Heidelberg Catechism was greeted not only in the Palatinate but in all Reformed churches as the correct expression of the Reformed faith, and attained the authority of a genuine symbolical standard. It was translated into nearly all languages, and has continued to be the basis of religious instruction to this day. . . . Its tone, notwithstanding the scholastic and dogmatizing or (as Ullmann says) constructive tendency, is truly popular and childlike.'

Then he quotes several questions as models of the catechetical style.

Dr. Dalton, of St. Petersburg:\footnote{\emph{Immanuel. Der Heidelb. Kat.,} etc., 1870, p. 15.}

'The Heidelberg Catechism exhibits the harmonious union of the Calvinistic and the Melanchthonian spirit. It is the ripe fruit of the whole Reformation and the true heir of the treasures gathered, not in ten years, but during that entire period. It is thoroughly Biblical, and represents its particular denominational type with great wisdom and moderation. We feel from beginning to end in the clear and expressive word the warm and sound pulse of a heart that was baptized by the fire and Spirit from above, and knows what it believes.'

It is gratifying that the Lutheran hostility of former days has given way to a sincere appreciation. Drs. Guericke and Kurtz, two prominent champions of Lutheran orthodoxy in the nineteenth century, in almost the same words praise the Heidelberg Catechism for 'its signal wisdom in teaching, its Christian fervor, theological ability, and mediating moderation.'\footnote{Guericke, \emph{Kirchengeschichte,} Vol. III. p. 610 (7th edition), and his \emph{Symbolik.} Kurtz, \emph{Lehrbuch der Kirchengeschichte,} p. 508 (5th edition).} Dr. Julius Stahl, an eminent jurist and the ablest apologist of modern Lutheranism within the Prussian Union, derived the religious revival of the Lutheran Church in his native Bavaria and his own conversion chiefly from the late venerable Reformed pastor and professor, Dr. J. Chr. G. L. Krafft, in Erlangen (died 1845). 'The man,' he said, before the General Synod at Berlin, 1846, 'who built up the Church in my fatherland, \emph{the most apostolic man I ever met in my life,} Pastor Krafft, was a strict adherent of the Reformed creed. Whether he carried the Heidelberg Catechism in his pocket I know not, but this I know, that he caused throughout the whole land a spring to bloom whose fruits will ripen for eternity.'\footnote{See art. \emph{Krafft,} by Goebel, in Herzog's \emph{Encykl.} Vol. VIII. p. 37.}

\xsubsection{The Brandenburg Confessions.}

§ 70. The Brandenburg Confessions.

(\emph{Confessiones Marchicæ.})

Literature.

Hartknoch: \emph{ Preussische Kirchenhistorie.} Frankf. 1686.

Zorn:  \emph{Historia derer zwischen den Lutherischen und Reformirten Theologis gehaltenen Colloquiorum.} Hamburg, 1705.

D. H. Hering:  \emph{Historische Nachricht von dem ersten Anfang der evang.-reformirten Kirche in Brandenburg und Preussen unter dem gottseligen Churfürsten Johann Sigismund, nebst den drei Bekenntniss-Schriften dieser Kirche.} Halle, 1778. The same: \emph{Neue Beiträge zur Geschichte der evangel.-reform. Kirche in den Preuss. Brandenburg. Ländern.} Berlin, 1787.

C. W. Hering:  \emph{Geschichte der kirchlichen Unionsversuche seit der Reformation.} Leipzig, 1836, 1837.

Beck:  \emph{Symbol. Bücher der ev.-reform. Kirche,} Vol. I. pp. 472 sqq.; Vol. II. pp. 110 sqq., 130 sqq.

Niemeyer:  \emph{Collectio,} Proleg. pp. lxxiv. sqq. and 642–689.

Böckel:  \emph{Die Bekenntniss-Schriften,} etc., pp. 425 sqq.

Möller:  \emph{Joh. Sigismund's Uebertritt zum reform. Bekenntniss,} in the \emph{Deutsche Zeitschrift.} Berlin, 1858, pp. 189 sqq.

Alex. Schweizer:  \emph{Die Protest. Centraldogmen,} Vol. II. pp. 6 sqq., 525 sqq., 531 sqq.

Comp. Herzog's \emph{Encyklop.} articles: \emph{Leipziger Colloquium,} Vol. VIII. p. 286; \emph{Joh. Sigismund,} Vol. XIV. p. 364; and \emph{Thorn} (by Henke), Vol. XVI. p. 101.

Brandenburg, the central province of Prussia, with Berlin as its capital, ruled since 1415 by princes of the house of Hohenzollern, at first embraced the Lutheran Reformation, but at the beginning of the seventeenth century the Elector became Calvinistic, drawing with him a few influential ministers and congregations. This Reformed diaspora received an accession of about twenty thousand exiled Huguenots under the liberal policy of the great Elector Frederick William (1620–1688), the proper founder of the Prussian monarchy, who secured the legal recognition of the Reformed Church in the Treaty of Westphalia (1648).

There are three Reformed Confessions of Brandenburg—namely, the Confession of the Elector Sigismund (1614), the Leipzig Colloquy (1631), and the Declaration of Thorn (1645). They bear a moderately Calvinistic, we may say a Unionistic, type, and had a certain symbolical authority in Brandenburg till the introduction of the union of the Lutheran and Reformed Churches in 1817. The great Elector mentions them together in 1664. The Canons of Dort were respectfully received but never adopted by the Brandenburg divines.

THE CONFESSION OF SIGISMUND. A.D. 1614.

See the original German text in the collections of Beck, Niemeyer, Böckel, and also in Heppe's \emph{Bekenntniss-Schriften der reform. Kirchen Deutschlands,} pp. 284–294.

John Sigismund (or Siegmund), Elector of Brandenburg (b. 1572, d. 1619) and ancestor of the royal line of Prussia, was brought up in the rigorous orthodoxy of the Lutheran Formula of Concord, and in his twenty-first year a solemn pledge was exacted from him by his father that he would always adhere to this creed (1593). But religious compulsion had on him an effect directly contrary to that contemplated (as is often the case with independent minds). His social relations with Holland, Cleves, and the Palatinate gave him a favorable impression of the doctrines and discipline of the Calvinistic Churches. In 1608 he succeeded to the throne. At Christmas, 1613, he publicly professed the Reformed faith by receiving the holy communion, according to the Reformed rite, in the Dome of Berlin, together with fifty-four others, including his brother John George, the Count of Nassau, Ernst Casimir, and the English embassador.

This act was the result of conscientious conviction.\footnote{Some writers, including Voltaire, trace the change to political motives—viz., that Sigismund wished to secure the friendship of Holland and England—but without proof. On the contrary, it was bad policy, and in its immediate effect rendered the Elector very unpopular among his German fellow-sovereigns and his own people. '\emph{Kein Wort,}' says Böckel, p. 427, '\emph{keine Handlung des Kurfürsten Johann Sigismund verräth, dass ihn irgend eine unreine Nebenabsicht geleitet habe.}' See also Möller and Hollenberg, l.c.} It was meant to be not so much a change of creed as a further progress in Protestantism, but it created a great sensation, and called forth violent protests from Lutheran princes and pulpits.\footnote{See Hutter's \emph{Calvinista aulico-politicus.}} An edict forbidding public denunciations had little effect. A fanatical mob arose in rebellion against the Reformed preachers, and plundered their houses (1615). The great majority of the Elector's subjects and his own wife remained Lutherans.\footnote{Dr. Tholuck (\emph{Geist der luther. Theologen Wittenbergs,} p. 118, referring to Hartknoch's \emph{Preuss. Kirchenhistorie,} p. 544) mentions the fact that Anna, the wife of Sigismund, in her will and testament ordered her chaplain in the funeral sermon to disown the Calvinistic (?) heresy that Christ's blood and death are merely a \emph{man's} blood and death.}

Nevertheless, his transition was of great prospective importance, for the house of Brandenburg was destined to become, by extraordinary talents and achievements, one of the leading dynasties of Europe, and to take the helm of the new Protestant German empire.

In May, 1614, Sigismund issued a personal confession of faith, which is called after him and also after his country. It was drawn up by himself, with the aid of Dr. Pelargus, General Superintendent at Frankfort-on-the-Oder. It is brief, moderate, conciliatory, and intended to be merely supplementary concerning the controverted articles. The Elector professes faith in the 'true, infallible, and saving Word of God, as the only rule of the pious which is perfect, sufficient for salvation, and abides forever.' Then he accepts, as agreeing with the Bible, the œcumenical creeds (namely, the Apostles', the Nicene, the Athanasian, also the doctrinal decisions of Ephesus, 431, and of Chalcedon, 451), and the Augsburg Confession of 1530, with the later improvements of Melanchthon.

In regard to the controverted articles, Sigismund rejects the Lutheran doctrine of the ubiquity of Christ's body, and exorcism in baptism as a superstitious ceremony, and the use of the wafer instead of the breaking of bread in the communion. He adopts the Reformed doctrine of the sacraments, and of an eternal and unconditional election of grace, yet with the declaration that God sincerely wished the salvation of \emph{all} men, and was not the author of sin and damnation.

In conclusion the Elector expresses his wish and prayer that God may enlighten his faithful subjects with his truth, but disclaims all intention to coerce their conscience, since faith was the free gift of God (John vi. 29; 2 Thess. iii. 2; Phil. i. 29; Eph. iii. 8), and no one should presume to exercise dominion over men's religion (2 Cor. i. 24). He thus freely waived, in relation to his Lutheran subjects, the right of reformation, which was claimed and exercised by other Protestant princes, and established a basis for religious liberty and union.

This wise toleration was in advance of the age, and contrasts favorably with the opposite policy of the Elector Augustus of Saxony, who forced the Formula of Concord upon his people, and answered the Emperor Maximilian II., when he interceded for the release from prison of Peucer (Melanchthon's son-in-law): 'I want only such servants as believe and confess in religion neither more nor less than I myself believe and confess.'\footnote{The Emperor replied: '\emph{Das wage ich von meinen Dienern nicht zu fordern.}' The same Elector Augustus said that 'if he had only one Calvinistic vein in his body, he wished the devil (sic!) would pull it out.'} These times of terrorism over men's consciences are happily passed, and Sigismund's toleration has become the settled policy of his successors to this day.

The conduct of Luther and Zwingli at Marburg gave tone and character to all subsequent union conferences of the two confessions they represent. The Reformed, with a larger charity, were always willing to commune with Lutherans notwithstanding minor doctrinal differences; while the Lutherans, with a narrower conscience and a more compact system of theology, refused the hand of fellowship to the Reformed, and abhorred as a syncretistic heresy all union that was not based upon perfect agreement in dogma; yea, during the seventeenth century they would rather make common cause with Romanists than Calvinists, and went so far as to exclude the Calvinists from heaven.\footnote{Dr. Hülsemann of Wittenberg traced the charitable hope of Calixtus that he would meet many Reformed in heaven to the inspiration of the devil ('\emph{spes dubio procul a diabolo inspirata}'). Calixtus asked, Who inspired this opinion of Hülsemann? Leyser wrote a book to show that communion with Papists was preferable to communion with Calvinists. Another book of that age professed to prove that 'the damned Calvinistic heretics have six hundred and sixty-six theses in common with the Turks.' The French Reformed Synod of Charenton in 1631 sanctioned the admission of Lutheran sponsors in baptism on the ground of essential agreement of the Augsburg Confession with the Reformed doctrine. This resolution was pronounced 'atheistic' by Lutherans as well as Romanists. The spirit of Lutheran bigotry in that classical period of polemic confessionalism and exclusivism is well characterized and illustrated by Dr. Tholuck, in his \emph{Geist der luther. Theologen Wittenbergs im} 17\emph{ten Jahrh.} (1852), pp. 115, 169, 211, etc. Comp. also above, p. 346; Gieseler, \emph{Kirchengeschichte,} Vol. III. Pt. II. (1853), p. 456; Hase, \emph{Kirchengesch.} 9th ed. p. 510.} Fortunately Calixtus and his school, who had the Melanchthonian spirit, formed an honorable exception, and the exception, after much misrepresentation and persecution, has become the rule in the Lutheran Church.

THE COLLOQUY AT LEIPZIG. A.D. 1631.

See the German text of the \emph{Colloquium Lipsiense} in Niemeyer, pp. 653–668, and in Böckel, pp. 443–456.

In the midst of the fierce polemics between the Churches and the horrors of the Thirty-Years' War growing out of it, there arose from time to time a desire for union and peace, which was strengthened by the common danger. In 1629, Ferdinand II., a pupil of the Jesuits, issued an edict aiming at the destruction of Protestantism, which might have been accomplished had not Gustavus Adolphus soon afterwards appeared on German soil. It was during this period that the classical union sentence (often erroneously attributed to Augustine), 'In necessary things unity, in doubtful things liberty, in all things charity,' was first uttered as a prophetic voice in the wilderness by a Lutheran divine of the school of Calixtus, and re-echoed in England by Richard Baxter.\footnote{See Lücke's treatise, \emph{Ueber das Alter, den Verfasser,} etc., \emph{des kirchlichen Friedensspruches,} etc., Göttingen, 1850. He traces it to Rupertus Meldenius, the obscure author of \emph{Parænesis votiva pro pace ecclesiæ ad theologos Augustanæ Confessionis} (before 1635), directed against the \textgreek{φιλοδοξία} and \textgreek{φιλονεικία} of the theologians, and commending humility and love of peace. Here the sentence occurs, '\emph{Si nos servaremus } in necessariis Unitatem, in non necessariis Libertatem, in utrisque Caritatem,  \emph{optimo certe loco essent res nostræ.}' A copy of the first edition of this book, though without date, is preserved in the City Library of Hamburg.}

Under the operation of this feeling and the threatening pressure of Romanism, the Elector Christian William of Brandenburg, accompanied by his chaplain, John Bergius, and the Landgrave William of Hesse, with the theological Professor Crocius and Chaplain Theophilus Neuberger, met at Leipzig with the Elector George of Saxony and the Lutheran divines Matthias Hoë of Hoënegg, Polycarp Leyser, and Henry Höpfner, to confer in a private way about a friendly understanding between the two confessions, hoping to set a good example to other divines of Germany. The conference lasted from March 3 to 23, 1631, and each session continued three hours.

The Augsburg Confession of 1530, with Melanchthon's subsequent explanations, was made the basis of the proceedings, and was discussed article by article. They agreed essentially on all the doctrines except the omnipresence of Christ's human nature, the oral manducation of his body in the eucharist by worthy and unworthy communicants. The Reformed divines were willing, notwithstanding these differences, to treat the Lutherans as brethren, and to make common cause with them against the Papists. But the Lutherans were not prepared to do more than to take this proposal into serious consideration.

The question of election was then also taken up, although it is not expressly mentioned in the Augsburg Confession. They agreed that only a portion of the race was actually saved. The Reformed traced election to the absolute will of God, and reprobation to the unbelief of men; the Lutherans (adhering to the happy inconsistency of the Formula of Concord) brought in God's foreknowledge of the faith of the elect, but they derived faith itself entirely from God's free electing grace. The difference was therefore very immaterial, and simply a matter of logic.

In conclusion, the theologians declared that the conference was intended not to compromise the Churches and sovereigns, but only to find out whether and to what extent both parties agreed in the Twenty-eight Articles of the Augsburg Confession, and whether there was reason to hope for some nearer approach in the future, whereby the true Church might be strengthened against the Papists. In the mean time the proceedings of the conference were to be regarded as strictly private, and not to be published by either party without the consent of the other. The theologians of the two Churches were to show each other Christian love, praying that 'the God of truth and peace grant that we may be one in him, as he is one with the Son (John xvii. 21). Amen, Amen in the name of Jesus Christ, Amen.'

The document is not signed by the princes who arranged the conference, but only by the theologians — namely, Drs. von Hoënegg, Leyser, Höpfner (Lutherans), and Bergius, Crocius, Neuberger (Reformed).\footnote{The proceedings were published by Hoë of Hoënegg, and by Bergius, 1635. See literature in Niemeyer, Proleg. p. lxxix.}

The proceedings were characterized by great theological ability and an excellent Christian temper, and showed a much closer harmony than was expected. They excited considerable sympathy among the Reformed at home and abroad. But the Lutheran members were severely taken to task for favoring syncretism, and in vindicating themselves they became more uncompromising against Calvinism than before. The conference was in advance of the spirit of the age, and left no permanent effect.

THE COLLOQUY OF THORN. A.D. 1645.

The official edition of the Acts: \emph{Acta Conventus Thoruniensis celebrati a.} 1645, etc., Warsaw, 1646 (very incorrect). The Acts, with the two Protestant Confessions (which were excluded from the official Acts), in Calovius, \emph{Historia Syncretistica} (1682), 1685, pp. 199–560. The Reformed \emph{Declaratio Thoruniensis,} Latin, in Niemeyer (pp. 669–689); German, in Böckel (pp. 865–884).

The Colloquy of Thorn, in West Prussia (\emph{Colloquium Thoruniense}), was likewise a well-meant but fruitless union conference in a time of sectarian intolerance and the suicidal folly of the Thirty-Years' War.

In this case the movement proceeded from the Roman Catholic king, Wladislaus IV., of Poland (1632–1648). In this country moderate Lutherans, Calvinists, and Moravians had formed a conservative union in the Consensus of Sendomir (1570), and a treaty of peace secured equal civil rights to Protestants and Romanists (\emph{Pax Dissidentium} in 1573). But this peace was denounced by the Pope as a league of Christ with Belial, and undermined by the Jesuits, who obtained the control of the education of the Polish nobility, and are to a large extent responsible for the ultimate dismemberment and ruin of that unfortunate kingdom.

Wladislaus made a patriotic effort to heal the religious discords of his subjects, and invited Romanists and Dissenters (Protestants) to a charitable colloquy (\emph{colloquium caritativum, fraterna collatio}) in the city of Thorn, which was then under the protection of the King of Poland (since 1454), and had embraced the Lutheran faith (1557). It began April 18, 1645, in the town-hall. There were three parties. The twenty-eight Roman deputies, including eight Jesuits, were determined to defeat the object of peace, and to prevent any concessions to Protestants. The Reformed had twenty-four delegates, chief among them the electoral chaplains John Bergius and Fr. Reichel, of Brandenburg, and the Moravian bishop Amos Comenius. The Lutheran deputation consisted of fifteen, afterwards of twenty-eight members; the most prominent were Calovius of Dantzic and Hülsemann of Wittenberg, the champions of the strictest orthodoxy, and George Calixtus of Helmstädt, the leader of a mild and comprehensive union theology.\footnote{It took Calixtus nearly three weeks to travel from Helmstädt to Thorn.} The sessions were private ('\emph{plebs penitus arcenda}'). The king's chancellor, Prince George Ossolinski, presided.

The first business, called '\emph{liquidatio},' was to be the preparation of a correct statement of the doctrinal system of each party. The Roman Catholic Confession, with a list of rejected misrepresentations, was ready early in September, and read in the second public session, Sept. 16. It was received among the official acts. On the same day the Reformed Confession was read, under the title \emph{Declaratio doctrinæ, ecclesiarum Reformatarum catholicæ.} But the Romanists objected to the word '\emph{catholic},' which they claimed as their monopoly, and to the antithetical part as being offensive to them, and excluded the document from the official acts. The Lutheran Confession was ready the 20th of September, but was even refused a public reading.\footnote{The Latin text in Calovius's \emph{Hist. syncret.} pp. 403–421; the German and Latin texts were separately issued at Leipzig, 1655, and at Dantzic, 1735. See also \emph{Scripta facientia ad Colloquium Thoruniense; accessit G. Calixti consideratio et} \textgreek{ἐπίκρισις,} Helmstädt, 1645, and \emph{Calixti Annotationes et animadversiones in Confessionem Reformatorum,} Wolfenbüttel, 1655.}

The Protestants sent a deputation to the king, who received them and their confessions with courtesy and kindness; but the Romanists demanded more alterations than the Protestants were willing to make, and used every effort to prevent the official publication of heresies. Unfortunately the dissensions among the Lutherans, and between them and the Reformed, strengthened the Romish party. The Colloquy closed Nov. 21, '\emph{mutua valedictione et in fraterna caritate,}' but without accomplishing its end. Calixtus says: 'The Colloquy was no colloquy at all, certainly no \emph{colloquium caritativum,} but \emph{irritativum.}' It left the three confessions where they were before, and added new fuel to the syncretistic controversy in Germany,\footnote{Hence the distich on the Synod of Thorn: \par  '\emph{Quid synodus? nodus: Patrum chorus integer? æger:}  \par \emph{Conventus? ventus: Gloria? stramen. Amen.}'} Calovius and Hülsemann charged Calixtus with aiding the Calvinists in their confession. The city of Thorn, which spent 50,000 guilders for the conference, suffered much from the Thirty-Years' War, also by a plague, and became the scene of a dreadful massacre of Protestants, Dec. 7, 1724, stirred up by the Jesuits in revenge for an attack on their college.

The Declaration of Thorn\footnote{The full title is '\emph{Professio Doctrinæ Ecclesiarum Reformatarum in Regno Poloniæ, Magno Ducatu Lithuaniæ, annexisque Regni Provinciis, in Conventu Thoruniensi, Anni} 1645, \emph{ad liquidationem Controversiarum maturandam, exhibita d.} 1 \emph{Septembris.}' First published at Berlin, 1646, under the title '\emph{Scripta partis Reformatæ in Colloquio Thoruniensi},' etc.} is one of the most careful statements of the Reformed Creed, and the only one among the three confessions of this Colloquy which acquired a practical importance by its adoption among the three Brandenburg Confessions. It is divided into a general part (\emph{generales professio}) and a special declaration (\emph{specialis declaratio}). The former acknowledges the canonical Scriptures of the Old and New Testaments in the original Hebrew and Greek, as the only perfect rule of faith, containing all that is necessary for our salvation. It adopts, also, in a subordinate sense, as explanatory summaries of Scripture doctrine, the œcumenical Creeds, and doctrinal decisions of the ancient undivided Church in opposition to the trinitarian, christological, and Pelagian heresies.\footnote{In the expression of agreement with the ancient Church the Declaration of Thorn is more explicit than any other Protestant confession, Lutheran or Calvinistic or Anglican. After saying that the summary of Scripture doctrine is contained in the Apostles' Creed, the Ten Commandments, the Lord's Prayer, and the Words of Institution of Baptism and the Lord's Supper, the Declaration proceeds: \par '\emph{Si quid vero, in hisce Doctrinæ Christianæ capitibus, dubitationis aut controversiæ de genuino eorum sensu exoriatur, profitemur porro, nos amplecti ceu interpretationem Scripturarum certam et indubitatam, Symbolum Nicænum et Constantinopolitanum, iisdem plane verbis, quibus in Synodi Tridentinæ Sessione tertia, tanquam Principium illud, in quo omnes, qui fidem Christi profitentur, necessario conveniunt, et Fundamentum firmum et unicum, contra quod portæ inferorum nunquam prævalebunt, proponitur.} \par '\emph{Cui etiam consonare Symbolum, quod dicitur Athanasianum, agnoscimus: nec non Ephesinæ primæ, et Chalcedonensis Synodi Confessiones: quinetiam, quæ Quinta et Sexta Synodi, Nestorianorum et Eutychianorum reliquiis opposuere: quæque adversus Pelagianos olim Milevitana Synodus et Arausicana secunda ex Scripturis docuere. Quinimo, quicquid primitiva Ecclesia ab ipsis usque Apostolorum temporibus, unanimi deinceps et notorio consensu, tanquam Articulum fidei necessarium, credidit, docuit, idem nos quoque ex Scripturis credere et docere profitemur.} \par '\emph{Hoc igitur Fidei nostræ professione, tanquam Christiani vere Catholici, ab omnibus veteribus et recentibus Hæresibus, quas prisca universalis Ecclesia unanimi consensu ex Scripturis rejecit atque damnavit, nos nostrasque Ecclesias segregamus.'}} Finally, as regards the controversy with Rome, it accepts the Altered Augsburg Confession and the Consensus of Sendomir (1570) as correct statements of the Scripture doctrines, differing in form, but agreeing in essence.

The 'Special Declaration' states the several articles of the Reformed system, both in its agreement with, and in its departure from, the creeds of Romanists and Lutherans.

The document is signed by a number of noblemen and clergymen from Poland, Lithuania, and Brandenburg.

\xsubsection{Minor German Reformed Confessions.}

§ 71. Minor German Reformed Confessions.

Heinrich Heppe:  \emph{Die Bekenntniss-Schriften der reformirten Kirchen Deutschlands.} Elberfeld, 1860. (Contains nine confessions of secondary importance, most of which are not found in other collections.)

The remaining Confessions of the Reformed Churches in Germany have only a local importance, and may be briefly disposed of.

1. The Confession of Elector Frederick III. of the Palatinate, 1577.—It was his last will and testament, and was published after his death by his son, John Casimir. It may be regarded as an explanatory appendix to the Heidelberg Catechism. It is a clear and strong testimony of his catholic and evangelical faith, and contains some wholesome warnings against the unchristian intolerance of the princes and theologians of his age.\footnote{The German text is given by Heppe, pp. 1–18; a Latin translation in the \emph{Corpus et Syntagma Confessionum,} with a Preface by John Casimir.}

2. The Confession of Anhalt, or Repetitio Anhaltina (i.e., a Repetition of the Augsburg Confession), 1581.\footnote{The German text in Heppe, pp. 19–67, the Latin in Niemeyer, pp. 612–641. Böckel excludes it from his collection because it is not strictly Reformed.}—It was prepared chiefly by Wolfgang Amling, Superintendent of Anhalt, and laid before a conference with Hessian divines held at Cassel, March, 1579.

The duchy of Anhalt, on the banks of the Elbe and Saale (formerly divided into four duchies, called after the principal towns, Anhalt-Dessau, Anhalt-Zerbst, Anhalt-Bernburg, Anhalt-Cöthen, in 1853 united into two, 1863 into one) embraced the Lutheran reformation in 1534, but during the controversies which led to the Formula Concordiæ it adhered to Melanchthon, and finally passed over to the Reformed faith in 1596. Prince John George married a daughter of Prince Casimir of the Palatinate, and introduced the Heidelberg Catechism and a simpler form of worship. At a later period (1644) Lutheranism was partly re-established, but Dessau, Bernburg, and Cöthen remained Reformed.

The 'Anhalt Repetition' can scarcely be numbered among the Reformed Confessions. It belongs to the Melanchthonian transition period, and represents simply a milder type of Lutheranism in opposition to the Flacian party. It recognizes, along with the Altered Augsburg Confession and the \emph{Corpus Doctrinæ} of Melanchthon, the Smalcald Articles and Luther's Catechisms, and professes even the \emph{manducatio oralis} and the \emph{manducatio indignorum}.\footnote{Ebrard (\emph{Kirchen- and Dogmengeschichte,} Vol. III. p. 575) is certainly wrong when he says that the \emph{Repetitio Anhaltina} proves that the Anhalt clergy '\emph{schon damals ganz und gar reformirt über die Person Christi und} das h. Abendmahl \emph{dachte.}' It expressly asserts in Art. vii. that even '\emph{indigne viscentes non quidem nudum aut communem panem calicemque manducant et bibunt, sed ipsum corpus et sanguinem Domini in Sacramento Cœnæ manducantes et bibentes . . . rei fiunt corporis et sanguinis Domini.}' See Niemeyer, p. 628, and Heppe, p. 46.} This is clearly incompatible with the Reformed system of doctrine.

3.  The Confession of Nassau, 1578, prepared, at the request of Count John of Nassau-Dillenburg, by the Rev. Christopher Pezel, who had been expelled from Saxony for Crypto-Calvinism. It was adopted by a general synod of that country, and first printed in 1593. It is Melanchthonian in the sense of the Altered Augsburg Confession and the Confession of Saxony, and rejects the doctrine of ubiquity as an unscriptural innovation and fiction.\footnote{Heppe, pp. 68–146.}

4. The Bremen Confession (\emph{Consensus Ministerii Bremensis}), prepared, 1598, by the same Pezel, who in the mean time had removed to Bremen, and signed by the pastors of that city. It is more decidedly Reformed, and adopts the Calvinistic view of predestination. Among, the books herein approved and recommended to the study of the pastors are also the Geneva Harmonia Confessionum, the Heidelberg Catechism, the Decades of Bullinger, and the Institutes of Calvin, as well as the works of Melanchthon.\footnote{Ibid. pp. 147–243.}

5. The Hessian Confession, adopted by a General Synod at Cassel, A.D. 1607, and published 1608.\footnote{Ibid. pp. 244–249.} It treats only of five articles: the Ten Commandments, the abolition of popish picture idolatry, the Person of Christ (against ubiquity), the eternal election, and the Lord's Supper (against the \emph{manducatio indignorum}). The Heidelberg Catechism and a modification of Luther's Small Catechism were both used in Electoral Hesse.\footnote{Comp. Heppe, \emph{Geschichte der Hessischen Generalsynoden von} 1568–1582, Kassel, 1847, 2 vols. The vexed question whether Hessia is Lutheran or Calvinistic has called forth a large controversial literature, in which the numerous works of this indefatigable investigator of the early history of German Protestantism are very prominent.}

6. The Confession of the Heidelberg Theologians, of 1607, is an exposition of what the Reformed Churches of Germany believe, and what they reject.\footnote{Heppe, pp. 250 sqq.}

7. The Catechism of Emden, 1554, prepared, after the model of Calvin's Catechism, by John a Lasko, or Laski (1499–1560), a converted nobleman and reformer of Poland. It was used in the Reformed Church of East Friesland, where he labored several years. It was afterwards superseded by the Heidelberg Catechism, which is partly based upon it.\footnote{Ibid. pp. 294–310. Comp. Bartels, \emph{Johannes a Lasco,} in the ninth volume of the valuable series of \emph{Väter und Begründer der reformirten Kirche} (1861), pp. 53 sq.}

\xsection{The Reformed Confessions of Bohemia, Poland, and Hungary.}

\xsubsection{The Bohemian Brethren and the Waldenses.}

§ 72. The Bohemian Brethern and the Waldenses.

Literature.

Franz Palacky (Historiographer of the Kingdom of Bohemia): \emph{Geschichte von Böhmen grösstentheils nach Urkunden und Handschriften.} Prag. (1836 sqq.), 3d ed. 1864 sqq. 5 vols. (the 5th vol. comes down to 1526). The same: \emph{Documenta Mag. Joannis Hus, vitam, doctrinam, causam in Constantiensi Concilio actam . . . illustrantia.} Prag. 1869 (mostly from unpublished sources). The same: \emph{Die Vorläufer des Hussitenthums in Böhmen.} Prag. 1869 (new ed.). The same: \emph{Urkundliche Beiträge zur Geschichte des Hussitenkrieges.} 1873, 2 vols. Palacky was a descendant of the Bohemian Brethren, and is the best authority on Bohemian history. He died May 27, 1876.

Jos. Alex. von Helfert:  \emph{Hus und Hieronymus.} Prag. 1853.

Anton Gindely:  \emph{Böhmen und Mähren im Zeitalter der Reformation.} Prag. 1857, 1858, 2 vols. (containing the History of the Bohemian Brethren from 1450–1609). The same: \emph{Quellen zur Geschichte der Böhm. Brüder,} in \emph{Fontes Rerum Austriacarum,} Vol. XIX. Wien, 1859. Gindely is a Roman Catholic, but kindly disposed to the Bohemian Brethren, and thoroughly at home in their literature.

Chr. Ad. Pescheck:  \emph{Geschichte der Gegenreformation in Böhmen.} Leipzig, 1850, 2d ed. 2 vols.

E. H. Gillett (d. 1875, in New York): \emph{Life and Times of John Huss; or, The Bohemian Reformation of the} 15\emph{th Century.} Boston, 1864, 2d ed. 2 vols., 3d ed. 1871.

W. Berger:  \emph{Joh. Hus und Kaiser Sigmund.} Augsb. 1871.

L. Krummel:  \emph{Utraquisten und Taboriten.} Gotha, 1871.

Fr. von Bezold:  \emph{König Sigmund und die Reichskriege gegen die Husiten.} 1872. By the same: \emph{Zur Geschichte des Husitenthums.} München, 1874.

Jaroslav Goll:  \emph{Quellen und Untersuchwngen zur Geschichte der Böhmischen Brüder.} Prag, 1878 (I.).

HUS\footnote{\par Hus (i.e., \emph{Goose}) and \emph{Hussites} (from the Bohemian genitive \emph{Husses}) is the correct spelling, followed by Palacky and Gindely, instead of \emph{Huss} and \emph{Husites.}} AND THE HUSSITES.

The reformation in the Kingdom of Bohemia (now a political division of the Austro-Hungarian Empire), began with John Hus and Jerome of Prague, who were influenced by the doctrines of Wycliffe, and who carried with them the greater part of the population, the Slavic Czechs. They were condemned by the œcumenical Council of Constance as heretics, and burned at the stake, the former July 6, 1415, the latter May 30, 1416. But their martyrdom provoked the Husite wars which would have resulted in the triumph of the Husites, had not internal divisions broken their strength.

The followers of Hus were, from 1420, divided into two parties, the conservative Calixtines, so called from their zeal for the chalice (\emph{calix}) of the laity, or Utraquists (\emph{communio sub utraque specie}), and the radical Taborites, so named from a steep mountain which their blind but brave and victorious leader, Ziska (d. 1424), fortified and called Mount Tabor. The Calixtines accepted the compromise of communion in both kinds, which the Council of Basle offered to them (1433), and mostly returned to the Roman Church. The Taborites rejected all compromise with the hated papal Antichrist, and demanded a thorough reformation, but they were defeated by the allied Romanists and Calixtines near Prague, 1434, and subdued by George Podiebrad, 1453.

THE BOHEMIAN BRETHERN.

From this time the Taborites disappeared as a party, but from their remnants arose, about 1457, a new and a more important sect, the Unitas Fratrum (\emph{Jednota bratrská}), as they called themselves, or the Bohemian Brethren.\footnote{This name applies also to the members who emigrated to Moravia, Saxony, and Poland; but the name \emph{Moravian Brethren} does not occur until the 18th century, when Zinzendorf incorporated into his own society (the Moravians, properly so called) the last survivors of the Bohemian brotherhood, who had come from Moravia to Saxony. See Gindely, Vol. I. p. 36. They were also called \emph{Waldenses,} and in derision \emph{Picards} (probably the same as \emph{Beghards}) and \emph{Grubenheimer, Pit-dwellers} (because they held divine service in pits and caves).} They adhered to the rigid discipline of the Taborites, but were free from their fanaticism and violence. They endeavored to reproduce, in peaceful retirement from the world, the simplicity and spirituality of the Apostolic Church as they understood it. They held to the Bohemian version of the Bible revised by Hus\footnote{Another Bohemian version or revision of the New Testament was made from the Greek by Blahoslav, a member of the Unitas Fratrum and the author of a Bohemian grammar (d. 1571).} as their only standard of faith and conduct. They rejected worldly amusements, oaths, military service, and capital punishment; they opposed the secular power of the clergy, and denounced the Pope of Rome as Antichrist. At first they received the sacraments from Calixtine and Romish priests who joined them.

In 1467 they effected an independent organization at a synodical meeting held in the village of Lhota, which was attended by about fifty members, priests and laymen, scholars and peasants, under the lead of Michael, formerly a Catholic priest. After praying and fasting, they elected by lot (Acts i. 26) three priests out of their number, and laid hands on them. Then they were all solemnly rebaptized. But not satisfied with this independent reconstruction of the Church, they sought regular ordination from a Waldensian bishop, Stephen of Austria, who was reported to have been ordained by a Roman bishop in 1434, and who afterwards suffered martyrdom in Vienna. Stephen ordained Michael; Michael ordained Matthias of Kunwald, and then, laying down his dignity, asked to be ordained afresh by Matthias, who was the first of the three elected by lot, and significantly bore the name of the supplementary apostle. This shows the vacillation of the Brethren between Presbyterianism and Episcopacy, as well as between radical independency and historical conservatism.\footnote{Gindely reports this from the scanty and conflicting sources, and adds the remark (Vol. I. p. 37): '\emph{Es zeigt das Schwanken des Gemüths und den Zweifel an die Berechtigung der gethanen Schritte, dass die Brüder in ihren Schriften gleich nach der Wahl jede Differenz zwischen priesterlicher and bischöflicher Würde verwarfen, mil ängstlicher Gewissenhaftigkeit aber bei sich die letztere einführten.}'} But they retained, or meant to retain, an unbroken succession of the episcopate, and transmitted it afterwards to the Moravian Church.\footnote{The last bishop of the old Unitas Fratrum was John Amos Comenius (or Komensky, a Czech, born in Moravia, 1592, died at Amsterdam, 1671), who acquired great celebrity by his new method of instruction by pictures and illustrations, and by his \emph{Janua Linguarum reserrata} and his \emph{Orbis pictus.} His nephew, D. E. Jablonsky, was elected and ordained bishop by a Synod of Bohemian Brethren in Poland, 1698, and he ordained David Nitschmann, the first bishop of the Moravians, 1735. See E. von Schweinitz, \emph{The Moravian Episcopate} (Bethlehem, Pa., 1865; comp. his art. \emph{Moravian Church,} in Johnson's \emph{Univ. Cyclop.} Vol. III.), and Benham, \emph{Origin and Episcopate of the Bohemian Brethren} (Lond. 1867). The Moravian episcopate depends on the Bohemian, and the Bohemian on the Waldensian episcopate, which in the thirteenth century did not claim to rest on apostolic succession. Comp. the quotations in Gieseler, \emph{Kirchengesh.} Vol. II. Pt. II. pp. 640, 641.}

The Brethren were cruelly persecuted; many were tortured and burned; others fled to neighboring Moravia, where for a short season they were unmolested. In the beginning of the sixteenth century they numbered in Bohemia about 200,000 members with 400 parishes. They had three printing establishments in 1519, while the Romanists had only one, and the Utraquists two. They made valuable contributions to evangelical hymnology. Their most fruitful author was Lucas of Prague (d. 1528), who did more for the organization of the society than its founder Gregor, and wrote over eighty books.\footnote{Gindely, Vol. I. p. 2OO, and Von Zezschwitz, \emph{Lukas von Prag,} in Herzog's \emph{Encyklop.,} Supplem. Vol. XX. pp 23 sqq., 31. Gindely, however, places no high estimate on the writings of Lucas, and charges him with great obscurity. They are mostly extant in manuscript.}

THE WALDENSES.

Literature.

I. The Waldensian MSS., mostly in the libraries of Geneva, Cambridge, Dublin, and Strasburg. The older prints are not reliable. See a description of these MSS. in Herzog, \emph{Die romanischen Waldenser,} pp. 46 sqq. The Morland MSS. of Cambridge were brought to light again by Henry Bradshaw, 1862.

II. The accounts of mediæval Catholic writers: Bernard Abbas Fontis Calidi (Fonte Claude, d. 1193); Alanus de Insulis (d. 1202); Stephanus de Borbone (Etienne de Bourbon, d. 1225); Yvonet (1275); Rainerius (1250); Pseudo-Rainerius; Moneta of Cremona; Gualter Mapes, of Oxford.

Roman Catholic historians are apt to confound the Waldenses with the heretical Albigenses and Cathari, and include them in the same condemnation; while some of the older Protestant historians reverse the process to clear the Albigenses of the charge of Manicheism.

III. Historical works, mostly in the interest of the Waldenses:

J. P. Perrin:  \emph{Histoire des Vaudois.} Geneva, 1619. English translation with additions by R. Baird and S. Miller. Philadelphia, 1847.

Pierre Gilles:  \emph{Histoire ecclésiastique des églises réformées—autrefois appellées églises Voudoises.} Geneva, 1655.

Jean Leger (pastor and moderator of the Waldensian churches, afterwards of a Walloon church at Leyden): \emph{Histoire générale des églises évangéliques des vallées de Piémont ou Vaudoises.} Leyden, 1669, 2 vols. fol. A German translation by Von Schweinitz. Breslau, 1750.

S. Morland:  \emph{History of the Evangelical Churches of the Valleys of Piedmont.} London, 1658. Morland was sent by Cromwell to Piedmont; he brought back a number of Waldensian MSS., and deposited them in Cambridge.

Jacques Brez (Waldensian): \emph{Histoire des Vaudois.} Paris, Lausanne, and Utrecht, 1796.

S. R. Maitland:  \emph{Tracts and Documents illustrative of the History of the Doctrines and Rites of the Ancient Albigenses and Waldenses.} London, 1832.

Ant. Monastier:  \emph{Histoire de l’église Vaudoise.} Paris and Toulouse, 1847, 2 vols.

Alexis Muston (Waldensian): \emph{Histoire des Vaudois.} Paris, 1834. The same: \emph{L’Israel des Alpes, première histoire complète des Vaudois.} Paris, 1851, 4 vols.

Chr. U. Hahn:  \emph{Geschichte der Waldenser.} Stuttgart, 1847. (The second volume of his learned \emph{Geschichte der Ketzer im Mittelalter.}) Contains many valuable documents.

A. W. Dieckhoff:  \emph{Die Waldenser im Mittelalter.} Göttingen, 1851. Marks an epoch in the critical sifting of the documents, but is too negative, and unjust to the Waldenses.

Herzog:  \emph{Die romanischen Waldenser. }Halle, 1853. Also his valuable art. \emph{Waldenser} in his \emph{Real-Encyklop.} Vol. XVII. pp. 502 sqq. Based upon a careful examination of the Waldensian MSS.

C. A. G. von Zezschwitz:  \emph{Die Katechismen der Waldenser und Böhmischen Brüder als Documente ihres wechselseitigen Lehraustausches. Kritische Textausgabe,} etc. Erlangen, 1863. Compare his \emph{System der christl. kirchl. Katechetik,} Leipz. 1863, Vol. I. pp. 548 sqq.

Palacky:  \emph{Verhältniss der Waldenser zu den böhmischen Secten.} Prag, 1869. (38 pp.)

Edmund de Schweinitz:  \emph{The Catechism of the Bohemian Brethren. Translated from the Old German.} Bethlehem, Pa., 1869.

G. Lechler:  \emph{Johann von Wiclif und die Vorgeschichte der Reformation.} Leipz. 1873, Vol. I. pp. 46–63.

F. Wagenmann:  \emph{Waldenser,} in Schmidt's \emph{Encyklop. des gesammten Erziehungs- und Unterrichtswesens,} Vol. X. (1875), pp. 259–274.

Soon after their organization the Brethren came into friendly contact with the older and like-minded Waldenses (Vaudois), so called from their founder, Peter Waldo, or Waldus, a lay evangelist of Lyons (about 1170), who gave his rich possessions to the poor. They called themselves originally \emph{the Poor of Lyons,} who by voluntary poverty and celibacy aimed at evangelical perfection.\footnote{The Dominican Stephen of Borbone says: '\emph{Incepit hæc secta circa annum ab incarnatione Domini }1170 . . . Waldenses  \emph{dicti sunt a primo huius hæresis auctore, qui nominatus fuit Waldensis. Dicuntur etiam Pauperes de Lugduno quia ibi inceperunt in professione paupertatis.}' They were also called \emph{Leonistæ,} from \emph{Leona,} Lyons; \emph{Sabatati,} from their wooden sandals (\emph{sabot}); and \emph{Humiliati,} from their humility.} The early confessional and catechetical books of the two sects are closely connected. The Brethren derived, as already noted, their episcopate from the Waldenses, and in 1497 they sent two delegates, Lucas of Prague and Thomas of Landskron (Germanus), to France and Italy, who reported that the Waldenses in those countries were far advanced in the knowledge of Scripture truth, while elsewhere they found nothing but false doctrine, superstition, loose discipline, and corrupt morals.\footnote{Joachim Camerarius, in his \emph{Historica narratio de Fratrum orthod. ecclesiis in Bohemia }(ed. by his grandson, Heidelb. 1605), gives a full account of two deputations of the Brethren to the Waldenses, one in 1467, and the other in 1497. See Herzog, pp. 290 sqq., and Gindely. Vol. I. pp. 88 sq.} On the other hand, many of the exiled Waldenses, who spread in every direction,\footnote{Pseudo-Rainerius: '\emph{fere nulla est terra, in qua hæc secta non sit.}'} emigrated to Bohemia, attracted by the religious commotions of that country, and coalesced with the Brethren into one community.

The Bohemian Brethren and the Waldenses made a near approach to evangelical Protestantism, and are the only mediæval sects which have maintained their existence to this day. But we must distinguish between their position before and their position after the Reformation, which marks an important epoch in their creed. Much confusion (as Gieseler observes) has been introduced into their history both by friend and foe.

The Waldenses formed at first no separate church, but an \emph{ecclesiola in ecclesia,} a pious lay community of Bible-readers. They were well-versed in Scripture, and maintained its supremacy over the traditions of men; they preached the gospel to the poor, allowing women also to preach; and gradually rejected the papal hierarchy, purgatory, prayers for the dead, the worship of saints and relics, the mass, transubstantiation, the oath, and capital punishment. Being excommunicated by Lucius III. (1184) and other popes as schismatics and heretics, they seceded and became a persecuted church. They had a clergy of their own with bishops, priests, and deacons. The origin and succession of their orders are involved in obscurity. They survived the fierce persecutions in France and the valleys of Piedmont, and extended their influence through emigrants to other countries, kindling a zeal for the study of the Scriptures in the vernacular, and strengthening the opposition to the papal Church. When they heard the glad tidings of the Reformation, they sent a deputation—Morel and Masson—to Œcolampadius, Bucer, and other reformers, in 1530, and derived from them clearer views of the distinction between canonical and apocryphal books, justification by faith, election and free-will, the marriage of the clergy, and the nature and number of sacraments. At a synod in the valley of Angrogne, Sept. 12–18, 1532, which was attended also by Farel and two other Reformed preachers of French Switzerland, the Reformation was adopted by a large majority, and subsequently carried out. Since that time the Waldenses became and remained a regular branch of the Reformed Church.\footnote{Herzog, pp. 378 sqq.}

In the course of time the consciousness of this change was obscured, and in their polemic zeal against Romanism they traced the Reformed doctrines to their fathers, who certainly prepared the way for them. Their manuscripts were interpolated and assigned to a much earlier date.\footnote{Leger dates, without any proof, the \emph{Nobla Leyczon} and the Waldensian Catechism from the year 1100; the Confession of Faith, the tracts on Purgatory and the Invocation of Saints, from 1120; the book on Antichrist from 1126.} Some of their historians even constructed an imaginary Waldensian succession of pure evangelical catholicity up to the apostolic age, in opposition to the papal succession of an apostate pseudo-catholicity, which they dated from the fictitious donation of Constantine to Pope Sylvester and the consequent secularization of the Church. This is the Protestant counterpart of the Romish caricatures of the Reformation, and deserves equal condemnation in the name of common honesty and historical truth.

A critical examination and comparison of the Waldensian manuscripts and the reports of the conferences with the Reformers have exposed these literary frauds, and produced at first a reaction against the Waldenses and in favor of the Bohemian Brethren, from whom some of their books were supposed to be derived. But on still further examination it appears that there was a mutual exchange of views and writings between the two, and that the assertions of some later Bohemian Brethren concerning their independence are as little to be trusted and as clearly unfounded as the claims of the Waldenses. Their oldest writings, from the twelfth to the fourteenth century, were popular translations of the Scriptures and extracts from the fathers, followed by more extended works, such as \emph{La Nobla Leyczon}\footnote{Given in the original by Herzog, pp. 444–457, from the Geneva MS., with the variations of the Dublin text. Herzog assigns it to the year 1400. Ebrard, \emph{Ueber das Alter der Nobla Leyczon,} in the \emph{Zeitschrift für histor. Theologie,} 1864, and in his \emph{Kirchengesch. }Vol. II. p. 193, traces it to the beginning of the thirteenth century, and defends the date of the Geneva MS., that the work was written fully eleven hundred years after St. John wrote, 'It is the last time' (1 John ii. 18), i.e., about 1200.} (i.e., \emph{lectio,} a didactic poem on Bible history and an exhortation to repentance), the \emph{Cantica,} an allegorical exposition, or application rather, of the Song of Songs, and several poems and ascetic tracts. The second class embraces the writings of the fifteenth century (on Purgatory, the Invocation of Saints, and the Sacraments), which are more or less dependent on the \emph{Confessio Taboritarum }(1433), and other Hussite documents.\footnote{See the comparison in Dieckhoff, pp. 377 sqq.} The third class was not composed or put into its present shape till after the adoption of the Reformation in 1532. Their chief confession is based upon the Gallican (1559), and was issued during the fearful massacre of 1655.\footnote{See Vol. III. pp. 757 sqq.}

The indebtedness of the Waldenses to the Reformation for a purer creed does not deprive them of a claim to the deep sympathy of all Protestant Christians, which in the period of their fiercest persecution in Piedmont (1655) provoked the threat of Cromwell to make the thunder of English cannon resound in the castle of St. Angelo, and inspired the sublime sonnet of Milton—

'Avenge, O Lord, thy slaughtered saints, whose bones

Lie scattered on the Alpine mountains cold;

Even them who kept thy truth so pure of old,

When all our fathers worshiped stocks and stones.

Forget not: in thy book record their groans,

Who were thy sheep, and in their ancient fold

Slain by the bloody Piedmontese, that rolled

Mother with infant down the rocks. Their moans

The vales redoubled to the hills, and they

To heaven. Their martyred blood and ashes sow

O'er all the Italian fields, where still doth sway

The triple tyrant; that from these may grow

A hundredfold, who having learnt thy way

Early may fly the Babylonian woe.'

The last lines sound like a prophecy; for since the day of liberty dawned on Italy (in 1848), that venerable martyr church has, from its mountain retreats in Piedmont, with youthful vigor established missions in nearly all the cities of the peninsula.

THE WALDENSIAN CATECHISM (1489) AND THE BOHEMIAN CATECHISM (1521).

The doctrinal affinity of the Waldenses and the Bohemian Brethren appears especially in their Catechisms, which are the most important of all their writings before the Reformation, and which prove their zeal for Christian education on the basis of the Scriptures. They bear such a striking resemblance to each other that the one must be in part a copy from the other. The Waldensian Catechism has a better claim to originality, and, although not nearly as old as was formerly supposed,\footnote{Leger, Monastier, and Hahn trace it to the beginning of the twelfth century.} must have been written before 1500; while the Bohemian, in the form in which it was presented to Luther, first appeared in print in 1521 or 1522, and was probably the work of Lucas of Prague (d. 1528), who had visited the Waldenses in Italy and France (1489). But both rest probably on older sources. Palacky brought to light (1869) a similar Catechism, which he derives from Hus before 1414.\footnote{Dieckhoff (pp. 98–115), from an imperfect knowledge of the Waldensian Catechism (as given by Perrin and Leger), maintained the priority of the Bohemian Catechism, and charged the Waldenses with gross plagiarism. Dr. Herzog (pp. 324 sq.) inclined to the same opinion, but with some qualification, and first edited the original text of the Waldensian Catechism from the Dublin MSS. in the Romance language (pp. 438–444). Since then Prof. Von Zezschwitz, of Erlangen, has published (1863) both Catechisms in their authentic form, with an elaborate argument for the priority of the Waldensian from internal evidence and from its affinity with other undoubted Waldensian documents. Ebrard (Vol. II. p. 491) assents to this view, and says: 'The Waldensian Catechism is thoroughly and characteristically Waldensian.' But Palacky traces both to a Bohemian Catechism (of about 4 pages) which he found in the imperial library of Vienna, and published, with a Latin version, in his \emph{Docmnenta }relating to Hus (pp. 703, 708). The authorship of Hus, however, is a mere conjecture ('\emph{cuius autor Hus esse videtur}'). The resemblance extends only to a few questions, and does not settle the point of priority; for Palacky himself admits that the Waldenses were in Prague as early as 1408, and known to Hus. 'The Hussites,' he says (\emph{Das Verhältniss der Waldenser,} etc., p. 20), were both disciples and teachers of the Waldenses, but more the latter than the former.'}

The Waldensian Catechism, called 'The Smaller Questions,'\footnote{\emph{Las interrogations menors.} The more extensive work on Antichrist was likewise arranged in questions and answers.} intended for children, is a remarkable production for an age of prevailing popular superstition and ignorance. It consists of fifty-seven questions by the teacher (\emph{lo barba,} i.e., \emph{uncle}), and as many answers by the pupil (\emph{l’enfant}). It embodies the Apostles' Creed, the Lord's Prayer, and the Ten Commandments, and is divided into three divisions—Faith (Ques. 6), Hope (Ques. 32), and Love (Ques. 47). This division was suggested by St. Paul (1 Cor. xiii. 13) and Augustine (\emph{Enchiridion}), and is followed also in the Greek Catechism of Mogila and the Russian Catechism of Philaret. Under the head of Faith we have a practical exposition of the Apostles' Creed and the Ten Commandments, showing their subjective bearing on a living faith. In the Second Part (Ques. 32), Love is defined to be a gift of the Holy Spirit and an intimate union of the human will with the divine will. In the Third Part (Ques. 48), Hope is defined to be a certain expectation of grace and future glory. The Catechism is directed against the idolatry and superstition of the anti-Christian Church, but the opposition is indirect and moderate. The characteristic Waldensian features are the distinction between a living and a dead faith (Ques. 8); the six evangelical commandments (Ques. 21); the seven gifts of the Holy Spirit (Ques. 23); the distinction between the true or essential (invisible) Church (\emph{la gleisa de la part de la substancia}), which consists of all the elect of God in Christ, known only to him, and the outward or institutional (visible) church (\emph{de la part de li menisteri}), i.e., the ministers and the people subject to them (Ques. 35); and the rigid exposition of the second commandment against all forms of idolatry (Ques. 29). Of the sacraments it is said (Ques. 46): 'Two are absolutely necessary for all; the rest are less necessary.' This clearly indicates that the Catechism was written before the Reformation period, when the Waldenses rejected all but two sacraments.

The Bohemian Catechism is longer, having seventy-five questions and answers. It follows the ``Waldensian in the general arrangement and first part, and introduces also (like the Greek catechisms) the Beatitudes from the Sermon on the Mount (Ques. 31); it has more to say of idolatry, the worship of Mary, the saints and martyrs, and especially on the Lord's Supper; but these additions lack perspicuity, and are too long for the use of children.

The following specimen will give an idea of these Catechisms, and the relation they sustain to each other and to the Catechism ascribed to Hus:

The Waldensian Catechism.

The Waldensian Catechism.

The Bohemian Catechism.

\emph{Las interrogacions menors.}

\emph{Translated.}

(\emph{von Zezschwitz,} p.41)

1. Si tu fosses demandà qui sies-tu? Respont:

1. \emph{If thou art asked, Who art thou? Answer:}

1. \emph{Was bistu? Antwort:}

\emph{Di.}\footnote{\par That is, \emph{Discipulus.} In other copies, \emph{L’enfant.}} \emph{Yo soy creatura de Dio racional e mortal.}

I am a creature of God, rational and mortal.

A. Ein vernunfftige schopfung Gottes vnd ein tötliche.

2. Dio perque te ha creà?

2. \emph{For what end has God made you?}

2. \emph{Warumb beschüff dich Gott?}

\emph{Di. Afin que yo conoissa lui meseyme e cola e havent la soa gracia meseyme sia salvà.}

That I may know and serve him, and be saved by his grace.

A. Das ich in solt kennen un liephaben vnd habende die liebe gottes das ich selig wurdt.

3. En que ista la toa salù?

3. \emph{On what rests thy salvation?}

3. \emph{Warauff steht dein seligkayt?}

\emph{Di. En tres vertùs substantials de necessità pertenent a salù.}

On three fundamental virtues, which are necessary to salvation.

Auff dreyen göttlichen tugenden.

4. Quals son aquellas?

4. \emph{Which are they?}

4. \emph{Welche seints?}

\emph{Di. Fè, sperancza e carità.}

Faith, Hope, and Love.

A. Der glaub, die lieb, die hofnung.

5. Per que cosa provarès aiczo?

5. \emph{How do you prove this?}

5. \emph{Bewer das.}

\emph{Di. L’apostol scriv. } 1 Cor. xiii.: \emph{aquestas cosas permanon, fè, sperancza e carità}

The Apostle writes, 1 Cor. xiii., 'Now abideth faith, hope, love, these three; but the greatest of these is love.'

A. S. Paul' spricht, utzundt bleyben vns dize drey tugendt, der glaub, die lieb vnd die hofnung, vnd das gröst ausz den ist die lieb.

6. [Qual es la \emph{prumiera vertù substancial?}

6. \emph{Which is the first fundamental virtue?}

6. \emph{Welches ist die erst grundtfest deiner seligkayt?}

\emph{Di. La fè. Car l’apostol di: non possibla cosa es placzer a Dio senza la fè. Mas a l’appropiant a Dio conven creyre, car el es e serè reguiardonador de li cresent en si.}]

Faith; for the Apostle says, 'It is impossible to please God without faith: for he that cometh to God must believe that he is, and that he is a rewarder of them that diligently seek him [Heb. xi. 6].

A. Der glaub.

7. Qual cosa es la fè?

7. \emph{What is faith?}

7. \emph{Bewer das.}

\emph{Di. Segond l’apostol } Heb. xi. \emph{es subsistencia de las cosas de sperar e argument de las non appareissent.}\footnote{Hus begins with Ques. 7 (\emph{Quid est fides? Respondet S. Paulas in Ep. ad Hebr.,} etc.), and gives the substance of Ques. 6, but omits Ques. 1–5, and has no trace of a threefold division.}

According to the Apostle, Heb. xi., faith is the substance of things hoped for, the evidence of things not seen.\footnote{Hus begins with Ques. 7 (\emph{Quid est fides? Respondet S. Paulas in Ep. ad Hebr.,} etc.), and gives the substance of Ques. 6, but omits Ques. 1–5, and has no trace of a threefold division.}

A. S. Paul' sagt zu den Juden, es ist vnmüglich Gott zugefallen on den glauben, dann d’zünhenen\footnote{That is, \emph{hinzunahen.}} will zu Gott, der musz gelauben das Gott sey, auch das er ein belöner sey der die in suchen.\footnote{Hus begins with Ques. 7 (\emph{Quid est fides? Respondet S. Paulas in Ep. ad Hebr.,} etc.), and gives the substance of Ques. 6, but omits Ques. 1–5, and has no trace of a threefold division.}

8. De quanta maniera es la fè?\footnote{The Waldensian Catechism begins with the subjective faith, the Bohemian Catechism (Ques. 1O) with the objective faith, as laid down in the Creed. Hus agrees with the former.}

8. \emph{How many kinds of faith are there?}\footnote{The Waldensian Catechism begins with the subjective faith, the Bohemian Catechism (Ques. 1O) with the objective faith, as laid down in the Creed. Hus agrees with the former.}

8. \emph{Was ist der glaub?}

\emph{Di. De doas manieras, czo es viva e morta.}

Two kinds, a living faith and a dead faith.

A. S. Paulus sagt, der glaub ist ein grundfest der ding welcher man hat zuversicht, vnd ein bewerung der vnsichtigen.

[Hus (third Ques.): \emph{Duplex est fides, altera viva, altera mortua.}]

9. Qual cosa e fè viva?

9. \emph{What is living faith?}

9. \emph{Welches glaubens bistu?}

\emph{Di. Lo es aquella, laqual obra per carità, testificant l’apostol} Gal. v.; [\emph{czo es l’observancza de li comandament de Dio. Fè viva es creyre en Dio, czo es amar luy meseyme e gardar li seo comandament.}]

It is faith active in love (as the Apostle testifies, Gal. v. 6), that is by keeping God's commandments. Living faith is to believe in God, that is, to love him and to keep his commandments.

A. Des gemainen christenlichen.

10. Qual cosa es fè morta?

10. \emph{What is dead faith?}

10. \emph{Welches ist der?}

\emph{Di. Segond Sanct Jaques, la fè, s’ilh non ha obras, es morta en si meseyme; e dereco, la fè es ociosa sencza las obras. O fè morta es creire esser Dio, creyre a Dio, creyre de Dio, e non creire en Dio.}\footnote{The distinction between \emph{credere Deum, credere Deo,} and \emph{credere in Deum} often occurs in the writings of Hus and in the Catechism ascribed to him (Palacky, p. 710).}

According to St. James, faith which has no works is dead in itself; faith is idle without works. Or dead faith is to believe that God is, to believe about God, of God, but not to believe in God.\footnote{The distinction between \emph{credere Deum, credere Deo,} and \emph{credere in Deum} often occurs in the writings of Hus and in the Catechism ascribed to him (Palacky, p. 710).}

A. Ich gelaub in Gott vatter almechtigen, etc.

[The Apostles' Creed in full.]

11. De laqual fè siès-tu?

11. \emph{What is your faith?}

11. \emph{Welcher unterschaid ist diser glaube?}

\emph{Di. De la vera fè catholica e apostolica.}\footnote{This is fuller than 'the common Christian faith' in the Bohemian Catechism (Ques. 9).}

The true catholic and apostolic faith.\footnote{This is fuller than 'the common Christian faith' in the Bohemian Catechism (Ques. 9).}

Das ein glaub ist lebendig, der ander tod.

12. Qual es aquella?

12. \emph{Which is that?}

12. \emph{Was ist der tod glauben?}

\emph{Di. Lo es aquella, la qual al conselh de li apostol es departià en docze articles.}

It is the one which at the Council of the Apostles was divided into twelve articles.\footnote{\par According to the mediæval tradition. Hus puts the names of the apostles before each article, and adds the damnatory clause of the Athanasian Creed.}

A. Es ist zu glauben Gott den herrn zu sein, Gott dem herren, vnd von Gott dem herrn, aber nicht in Gott den herrn.\footnote{The distinction between \emph{credere Deum, credere Deo,} and \emph{credere in Deum} often occurs in the writings of Hus and in the Catechism ascribed to him (Palacky, p. 710).}

13. Qual es aquella?

13. \emph{Which is it?}

13. \emph{Was ist der lebendig glauben?}

\emph{Di. Yo creo en Dio lo payre tot poissent.}

I believe in God the Father Almighty, etc.

A. Es ist zu glaubn in Gott den vater, den sun, den heylig geyst.

[Now follows the Apostles' Creed in full.]

\xsubsection{The Bohemian Confessions after the Reformation. A.D. 1535 and 1575.}

§ 73. The Bohemian Confessions after the Reformation. A.D. 1535 and 1575.

THE REFORMATION AND COUNTER-REFORMATION IN BOHEMIA.

Comp. the work of Pescheck, quoted p. 565; and Reuss:  \emph{La Destruction du Protestantisme en Bohème.} Strasburg, 1867.

The Reformation rekindled the fire of the Husite movement, and made rapid progress within and without the Catholic Church. The Bohemian Brethren sent, from 1520, several delegations to Wittenberg to confer with Luther. They received new light in doctrine, but painfully missed discipline in the churches of Germany. Luther was at first displeased with their figurative theory of the Lord's Supper, their views of justification, and the celibacy of the clergy, and induced them to conform them to his teaching, but afterwards he treated them with a degree of indulgence and forbearance that contrasts favorably with his uncompromising antagonism to the Zwinglians. Nevertheless, the Bohemian Brethren, like the Waldenses, ultimately passed in a body to the Reformed communion, with which they had more sympathy in matters of doctrine and discipline.\footnote{They wrote afterwards to Beza (Dec. 3, 1575): '\emph{Lutherus nostra sic fuit interpretatus, quasi ipsius sententiæ sint consentanea, sua quidem ille culpa, non nostra.}' Zezschwitz, p. 153; Ebrard, Vol. III. p. 400. They had sent a deputation to Bucer and Calvin at Strasburg in 1510, who were well received.} Besides them we find in Bohemia, after the Reformation, three Protestant parties, Utraquists, Lutherans, and Calvinists.

There was at one time, during the reign of Maximilian II., a fair prospect of the conversion of the whole Bohemian nation, as also of the German provinces of Austria, to Protestantism; but before the work was consolidated, the Jesuits, backed by the whole power of the Hapsburg dynasty, inaugurated a counter-reformation and a series of cruel persecutions which crushed the evangelical faith, and turned that kingdom into a second Spain. The bloody drama of the Thirty-Years' War began at Prague (1618). Emperor Ferdinand II. (1619–1637), a fanatical pupil of the Jesuits, fulfilled his terrible vow to exterminate heresy by all possible means, though he should have to reign over a desert. The execution of twenty-seven of the most distinguished Protestants, in June, 1621, was the signal for this war of extermination. The richest families were deprived of their property. Protestant worship was forbidden. Protestant churches, schools, and hospitals were razed to the ground, or turned into Jesuit churches and colleges. All Protestant preachers, professors, and school-teachers were ordered, in 1624, to leave the country within a week, under pain of death. Bohemian and German Bibles and all Bohemian works published after 1414, being suspected of heresy, were destroyed in immense quantities on marketplaces and beneath the gallows. One Jesuit, Anton Koniasch (1637) boasted that he had burned over 60,000 books. Thus the whole Czech literature and civilization was overwhelmed with ruin, and ignorance as dark as midnight spread over the land.\footnote{See, on this wholesale destruction of books, Pescheck's \emph{Geschichte der Gegenreformation in Böhmen,} Vol. II. pp. 93 sqq. Bohemian works published from 1414 to 1635 are exceedingly rare, or are to be found only outside of Bohemia in the libraries at Herrnhut, Dresden, and Leipzig.} Protestants were forbidden the rights of citizens; they could not carry on a trade, nor contract marriage, nor make a will. Even light and air were denied them. 'More than thirty thousand Bohemian families, and among them five hundred belonging to the aristocracy, went into banishment. Exiled Bohemians were to be found in every country of Europe, and were not wanting in any of the armies that fought against Austria. Those who could or would not emigrate held to their faith in secret. Against them dragonades were employed. Detachments of soldiers were sent into the various districts to torment the heretics till they were converted. The ``Converters'' (\emph{Seligmacher}) went thus throughout all Bohemia, plundering and murdering. ... A desert was created; the land was crushed for a generation. Before the war Bohemia had 4,000,000 inhabitants, and in 1648 there were but 700,000 or 800,000. These figures appear preposterous, but they are certified by Bohemian historians.'\footnote{Heusser, \emph{The Period of the Reformation,} English translation, New York, 1874, p. 426. Dr. Döllinger, in his concluding address at the Bonn Union Conference in August, 1875, speaking of the suppression of the Reformation in Austria, made the following remarks: '\emph{Nach römischer Lehre ist eine katholische Regierung verpflichtet, die Andersgläubigen zu unterdrücken. Die Päpste haben die Habsburger durch die Jesuiten stets zur Befolgung dieser Lehre angehalten. In der zweiten Hälfte des sechszehnten Jahrhunderts war die Bevölkerung in einigen überwiegend deutschen Erbstaaten fast zu neun Zehntel protestantisch. Durch das System der Zwangsbekehrung und der Austreibung der Protestanten wurde am Ende des }16. \emph{Jahrhunderts and im }17. \emph{der römische Katholicismus wieder herrschend. Die wenigen Schriftsteller, welche Oesterreich im} 17\emph{ten Jahrhundert hatte, klagen einmüthig über den Schaden, den die Protestanten-Austreibung dem Wohlstand Oesterreichs gebracht. Man darf sagen, es macht sich noch heute fühlbar, dass damals der beste Theil der städtischen Bevölkerung vertrieben wurde. Sine grosse geistige Versumpfung ist die Folge der engen Verbindung der Habsburger Dynastie mit der Curie gewesen. Ich sage: der Habsburgischen Dynastie; die jetzige Dynastie ist die lothringische, aus welcher ganz andere Regenten hervorgegangen sind. Ihr gehört Joseph II. an, aber auch die andern Kaiser dieser Dynastie haben nicht ihre Unterthanen um der Religion willen unterdrückt. Oesterreich leidet noch jetzt an den schlimmen Folgen früherer Missregierungen, aber es ist ein Staat, der noch eine Zukunft hat, und sein neues Emporblühen ist von grosser Wichtigkeit für Europa. Wenn wir den Satz des Herrn; an ihren Früchten sollt ihr sie erkennen, auf das Papalsystem anwenden, so können wir nur ein hartes Urtheil über dasselbe fällen. Das jetzige Verhalten des römischen Stuhles zeigt aber, dass er aus der Weltgeschichte nichts gelernt hat, dass sie ihm ein mit sieben Siegeln verschlossenes Buch ist.}'}

The exiled Bohemian Brethren became the nucleus of the Moravian Brotherhood (1722), and in this noble little Church, so distinguished for its missionary zeal, they continue to this day. Their last and worthy Bishop, Amos Comenius, died an exile in Holland, 1671, with the hope of the future revival of his persecuted Church, which was fulfilled through the labors of Count Zinzendorf. But even in Bohemia Protestantism could not be utterly annihilated. It began again to raise its feeble head when Emperor Joseph II. issued the Edict of Toleration (1781). The recent revival of Czech patriotism and literature came to its aid. The fifth centenary of the birth of Hus was celebrated at Prague, 1869, and his works and letters were published. In 1875 there were forty-six Reformed congregations in Bohemia and twenty-two in Moravia. The number of Lutheran congregations is smaller, and they belong almost entirely to the German part of the population.

THE BOHEMIAN CONFESSION OF 1535.

The Latin text in the \emph{Corpus et Syntagma Conf.,} and in Niemeyer, pp. 771–818; the German text in Böckel, pp. 780–830.

The Bohemian Brethren surpass all Churches in the number of their confessions of faith, which amount to no less than thirty-four from 1467 to 1671, in the Bohemian, Latin, and German languages.\footnote{Gindely enumerates them in \emph{Fontes,} etc., pp. 453 sqq. Comp. Zezschwitz, in Herzog's \emph{Real-Encyklop. }Vol. XX. p. 31.} But they were all superseded by two, respectively called the First and the Second Bohemian Confessions.

The first of these confessions was prepared, after the example of the Lutherans at the Diet of Augsburg, in proof of their orthodoxy, signed by the noblemen belonging to the Unitas, and laid by a deputation before King Ferdinand at Vienna, Nov. 14, 1535, who promised to take it into consideration.\footnote{\emph{Confessio Fidei ac Religionis, Baronum ac Nobilium Regni Bohemiæ, serenissimo ac invictissimo Romanorum, Bohemiæ,} etc., \emph{Regi, Viennæ Austriæ, sub anno Domini }1535 \emph{oblata.}'} It was written in Latin by an unknown author, probably by John Augusta, Senior of the Brethren, and, after the death of Lucas of Prague, their most influential leader (d. 1572), and with his aid it was translated into German.\footnote{So Gindely, Vol. I. p. 233, 238. Niemeyer (Proleg. p. lxxxvi.) asserts: '\emph{Prodiit primum lingua Bohemica, deinde Latine reddita Vitembergæ publici juris facta est.}' But Gindely is a much better authority in Bohemian matters.}

The confession consists of a long apologetic preface against the charges of heresy and immorality, and of twenty articles. It closely resembles in form and contents the Augsburg Confession. In Art. XII., on Baptism, it is stated that the Brethren had formerly rebaptized converts, but that they had given up this practice as useless. Infant baptism is acknowledged (Matt. xix. 14). The doctrine of the Lord's Supper (Art. XIII.) is accommodated to the Lutheran theory, though framed somewhat vaguely.\footnote{'\emph{Docent etiam, quod his Christi verbis, quibus ipse panem corpus suum, et vinum speciatim sanguinem suum esse pronunciat, nemo de suo quidquam affingat, admisceat aut detrahat, sed simpliciter his Christi verbis, neque ad dexteram neque ad sinistram declinando credat.}'}

The Bohemians sent the confession with a deputation to the Reformers at Wittenberg (1536). Luther disapproved the articles on celibacy and justification, but after the Brethren had made some corrections he published the document, at their request and expense, in 1538, with a favorable preface. In later editions the Bohemians made many changes.\footnote{See Niemeyer, \emph{Proleg.} p. lxxxvii.}

THE SECOND BOHEMIAN CONFESSION. A.D. 1575.

The Latin text in Niemeyer, pp. 819–851; the German text in Böckel, pp. 827–849.

The historical notices I have chiefly derived from Pescheck's  \emph{Geschichte der Gegenreformation in Böhmen,} 2d ed. Vol. I. pp. 103 sqq., and from Gindely's  \emph{Geschichte der Böhmischen Brüder,} Vol. II. pp. 141 sqq.

The mild and liberal Emperor Maximilian II. (1564–1576) was kindly disposed towards his Protestant subjects, and had a certain degree of sympathy with their creed. While holding a diet at Prague he allowed the non-Catholic Bohemians to compose a united confession of their faith. The Utraquists, Lutherans, Calvinists, and Bohemian Brethren laid aside their disputes and agreed upon a moderate doctrinal statement, which is more particularly called the Bohemian Confession.\footnote{\par '\emph{Confessio Bohemica, hoc est, Confessio sanctæ et Christianæ fidei, omnium trium Ordinum Regni Bohemiæ, Corpus et Sanguinem Domini nostri Jesu Christi in Cœna sub utraque specie accipientium},' etc.} It was prepared in the Bohemian language by two divines—Dr. Paul Pressius and M. Krispin\footnote{A new chapter in the history of Bohemia—now called Czechoslovakia—and of the followers of Huss was opened with the war, 1914. After nearly 400 years of subjection to Austria, the country gained its freedom and was made a republic, with T. G. Masaryk, who had spent some time in the U. S., as president, 1918. Religious liberty was established, the statue of Mary, erected on the central square of Prag, pulled down, and a massive monument dedicated to Huss with an address by Masaryk. The 'Czechoslovak National Reformed Catholic Church,' now constituting 4 per cent, of the population, was organized, 1920, and occupies a position midway between Romanism and Protestantism, leaning to the Greek Church. Its first bishop received consecration from a Serbian bishop. It uses the national tongue in worship and practises clerical marriage and the elective principle in the choice of its bishops. The Hussites, adopting the name 'The Evangelical Union of Bohemian Brethren,' have doubled in numbers and carry on a theological seminary in Prag with a full faculty, partly supported by the state and occupying quarters in the university buildings. In 1925 the papal nuncio, as a protest at the part taken by Masaryk and the government at a Huss celebration, left Prag. Papal relations were re-established, 1927. Masaryk holds the degree of doctor of theology from the theological seminary in Prag for a thesis on Huss.—Ed.}—and adopted with some changes by the Diet of Prague. It was presented to Maximilian, May 17, 1575. He gave the delegates the verbal promise of protection in their faith and worship. It was afterwards presented to Maximilian's son and successor, Rudolph II., 1608, who, under the political pressure of the times, in an imperial letter, or charter, granted the Protestant Bohemians equal rights with the Roman Catholics, a separate consistory at Prague, and the control of the university (1609). But these concessions were of short duration. Emperor Matthias violated the compact, and Ferdinand II. annulled it by his Edict of Restitution (1629), which gave the Romanists full power to suppress Protestantism.

The Second Bohemian Confession consists of an address to Maximilian II. and twenty-five articles on the holy Scriptures, on God, the Holy Trinity, the fall and original sin, free-will, the law, justification, faith and good works, the Church, the sacraments, etc. It is in essential agreement with the Augsburg Confession and the older Bohemian Confession. The doctrine of the Lord's Supper is conformed to the later Melanchthonian view. A German translation was transmitted to the divines at Wittenberg, and approved by them Nov. 3, 1575. A Latin translation appeared in 1619.

\xsubsection{The Reformation in Poland and the Consensus of Sendomir. A.D. 1570.}

§ 74. The Reformation in Poland and the Consensus of Sendomir. A.D. 1570.

Literature.

\emph{Consensus Sendomiriensis,} in Niemeyer, pp. 561 sqq. The German text in Beck, Vol. II. pp. 87 sqq.

Joannis a Lasco:  \emph{Opera tam edita quam inedita recensuit vitam autoris enarravit} A. Kuyper.. Amstel. 1866, 2 Tom. The first volume contains his dogmatic and polemic writings, including the \emph{Responsio adv. Hosium} (1559); the second his Confession, Catechisms, and Letters, including a few from Poland, 1556–59 (Vol. II. pp. 746–765). His Letters were previously published by Gerdesius, in his \emph{Scrinium antiquarium,} Groning. 1750.

Dan. Ern. Jablonski:  \emph{Historia consensus Sendomiriensis inter evangelicos regni Poloniæ et M.D. Lithuaniæ in synodo generali evangelicorum utriusque partis Sendomiriæ A.D.} 1570 \emph{die} 14 \emph{Aprilis initi. }Berolini, 1731.

C. G. von Friese:  \emph{Reformationsgeschichte von Polen und Lithauen.} Breslau, 1786, 3 vols.

Valerian Krasinski (an exiled Polish Count): \emph{Historical Sketch of the Rise, Progress, and Decline of the Reformation in Poland.} London, 1838 and 1840, 2 vols. German translation by W. Ad. Lindau. Leipz. 1841. Krasinski: \emph{Sketch of the Religious History of the Slavonic Nations.} Edinburgh, 1851. The same in French (\emph{Histoire religieuse des peuples slaves}), Paris, 1853, with an introduction by Merle d'Aubigné.

G. W. Th. Fischer:  \emph{Versuch einer Geschichte der Reformation in Polen.} Grätz, 1855–56, 2 vols.

P. Bartels:  \emph{Johannes a Lasco.} Elberfeld, 1860. In Vol. IX. of \emph{Leben der Väter der reform. Kirche.}

Dr.  Erbkam: Art. \emph{Sendomir,} in Herzog's \emph{Real-Encykl.} Vol. XXI. pp. 24–45. Dr. Erdmann: Art. \emph{Polen,} ibid. Vol. XII. pp. 1 sqq.

The history of the Reformation in Poland is as sad as that in Bohemia. It started with fair prospects of success, but was suppressed by the counter-reformation under the energetic and unscrupulous leadership of the Jesuits, who took advantage of the dissensions among Protestants, the weakness of the court, and the fickleness of the nobility, obtained the control of the education of the aristocracy and clergy, and ultimately brought that unfortunate kingdom to the brink of internal ruin before its political dismemberment by the surrounding powers.

POLAND IN THE SIXTEEN CENTURY.

Poland became a mighty kingdom by the union with Lithuania (1386) and the successful wars with the Teutonic order in Prussia. In the middle of the sixteenth century it extended from the shores of the Baltic to the Black Sea, and embraced Great Poland (Posen), Little Poland (Warsaw), Lithuania, Samogitia (Wilna), Courland, Livonia, Esthland, Podlesia, Volhynia, Podolia, Ukraine, and the Prussian territories of Dantzic, Culm, and Ermeland. The population was Slavonic, with a large number of Germans and Jews. It originally received Christianity from the Greek Church, through Bohemia, but, owing to its close connection with the German empire, it became, like Bohemia, Roman Catholic during the tenth century. The government was in the hands of the nobility, which controlled the king. The power of the Church was restricted to spiritual affairs, and weakened by the immorality of the clergy.

THE REFORMATION.

Poland never showed special devotion to the Roman See, and during the Council of Constance manifested some sympathy with the reform of Hus. Waldenses, Bohemians, and all classes of Protestants, even Socinians and Anabaptists, found hospitable shelter.

The Lutheran Reformation was introduced by Polish students returning from Wittenberg, and by Lutheran tutors employed in the families of the nobles. It triumphed in the German cities of Dantzic (1525) and Thorn (1530).

Among the Slavonic population and the higher nobility, and in the University of Cracow, Calvinism made rapid progress. It was patronized by Prince Nicholas Radziwill, the Chancellor of Poland under King Sigismund Augustus II. (1548–1572). The king himself corresponded with Calvin, and read his 'Institutes' with great zeal. Calvin dedicated to him his Commentary on the Epistle to the Hebrews, and in some remarkable letters solemnly urged him to use the favorable opportunity for the introduction of the pure doctrine and worship of Christ before the door might be forever closed. In a large kingdom with strongly feudal institutions he would allow, for the sake of unity and order, and after the model of the ancient Church, the episcopal organization, with an archbishop and a regular succession; but he thought that under the circumstances the Reformation could not be introduced without some irregularity, since the papal bishops had become the open enemies of the gospel. He became at last discouraged by the indecision of the king, and lost confidence in the sincerity of the nobles. His fears were only too well realized.\footnote{On Calvin's relation to Poland, see Stähelin, \emph{Joh. Calvin,} Vol. II. pp. 22 sqq.}

Another powerful element were the Bohemian Brethren, who, driven from their native land in 1548, emigrated in large numbers and organized forty congregations in Great Poland.\footnote{Vergerius wrote, 1557, to Stanislaus Ostrorog: '\emph{Esse jam in Polonia circiter XL ad eorum normam institutas ecclesias, quæ sane florent, multo autem plures propediem instituendas.}'} They were well received, and, by the affinity of race and language, their purity, simplicity, and strict discipline, they made a deep impression on the Slavonic Poles. The Brethren united with the Calvinists at the first general Protestant Synod held at Kosminek, 1555. The latter adopted the confession, liturgy, and episcopal government of the former. This step was highly approved by Calvin, who wrote to a Polish nobleman, Stanislaus Krasinski: 'From a union with the Waldenses [as the Brethren were sometimes called] I hope the best, not only because God blesses every act of a holy union of the members of Christ, but also because at the present crisis the experience of the Waldenses, who are so well drilled in the service of the Lord, will be of no small benefit to you.' He also advocated union with the adherents of the Augsburg Confession as this was understood and explained by its author. He was invited by the nobility to Poland, but could not leave Geneva.

JOHN A LASCO.

In Calvin's place appeared, by his advice and probably at the invitation of the king, John a Lasco, or Laski, a Polish nobleman, distinguished among the Reformers of the second rank. Born at Warsaw, 1499, and educated for the priesthood by his uncle, the Archbishop of Gnesen and Primas of Poland, he made a literary journey to Holland and Switzerland, and became personally acquainted with Zwingli at Zurich (1524) and with Erasmus at Basle (1525), who shook his faith in the Roman Church.\footnote{Erasmus spoke of Laski in the highest terms, and sold him his library for three hundred crowns, with the privilege of retaining it till his death. Krasinski, l.c. p. 98 (German ed.).} On his return to Poland he endeavored to introduce a moderate reformation, but the country was not prepared for it. He declined an offer to a bishopric, and sacrificed bright prospects to his conviction, preferring to be in a foreign land 'a poor servant of Christ crucified for him.' He labored several years as Reformed pastor in Emden, East Friesland, until the Interim troubles drove him and his friends to England. He organized in London three congregations of Dutch, German, French, and Italian emigrants (\emph{ecclesiæ peregrinorum}) on a Presbyterian and voluntary basis, under the protection of Archbishop Cranmer and Edward VI. The persecution of Queen Mary forced him again to wander in exile. When he landed with a hundred and seventy-five members of his flock in Denmark, 1553, he was refused shelter in cold winter because he could not subscribe to the Lutheran doctrine of the real presence. He fully experienced the force of his motto, 'The pious have no home on earth, for they seek heaven.' After laboring a short time in a congregation of English and other pilgrims in Frankfort-on-the Main, he accepted the invitation to his native country in 1556, and was made General Superintendent of Little Poland. Here he prepared, with the aid of other scholars, an admirable Polish translation of the Scriptures, published after his death, organized Reformed Churches (which increased in his time to the number of one hundred and twenty-two), and confirmed the union of the Calvinists with the Bohemian Brethren, although he himself preferred the Presbyterian polity with lay representation to the Bohemian episcopacy, and differed from their view of the Lord's Supper and other articles of their confession. He died Jan. 7, 1560, in the midst of work and care.\footnote{He wrote to Calvin, Feb. 19, 1557 (\emph{Opera,} Vol. II. p. 746): '\emph{Ita nunc obruor curis ac negotiis, mi Calvine! ut nihil possim scribere. Hinc hostes, illinc falsi fratres nos adoriuntur, ut non sit quies ulla, sed et pios multos habemus, sit Deo gratia! qui nobis sunt et adiumento et consolationi.}'}

PETER PAUL VERGERIO.\footnote{See Chr. H. Sixt: \emph{Petrus Paulas Vergerius, . . . eine reformationsgeschichtliche Monographie} (Braunschweig, 1855), pp. 391 sqq. and 437 sqq. Comp. also Herzog's art. \emph{Vergerius,} in his \emph{Real-Encykl.} Vol. XVII. pp. 65 sqq.}

During the same period Poland was twice visited (1557 and 1559) by another remarkable man among the secondary reformers— Peter Paul Vergerio (1498–1565), formerly papal nuncio to Germany and Bishop of Capo d’Istria.\footnote{See Chr. H. Sixt: \emph{Petrus Paulas Vergerius, . . . eine reformationsgeschichtliche Monographie} (Braunschweig, 1855), pp. 391 sqq. and 437 sqq. Comp. also Herzog's art. \emph{Vergerius,} in his \emph{Real-Encykl.} Vol. XVII. pp. 65 sqq.} In the attempt to refute the Lutheran writings he had become a Protestant, introduced the Reformation in the Italian parts of the Grisons (Valtellina, Poschiavo, and Bregaglia), and then took up his residence in Tübingen under the protection of Duke Christopher of Würtemberg, writing many books and making important missionary journeys. He was well received in Poland by Prince Radziwill and the king. He associated mainly with Lutherans and the Bohemian Brethren, but labored for the cause of union, like Laski.\footnote{He thought at one time of joining the \emph{Unitas Fratrum,} being disgusted with the renewal of the sacramental war. Even Melanchthon once expressed a similar desire, '\emph{in Valdensium ecclesiis me inserere et in illis mori; placent enim mihi summopere.}' See his letter to V. Dietrich, quoted by Herzog, p. 71.}

He aided the Reformation by his able pen, and the Roman historian Raynaldus says that 'this wretched heretic led many weak Catholics into the camp of Satan.' But his stay in Poland was too short to leave permanent results.

THE PAPAL REACTION AND TRIUMPH.

In the mean time the Roman Catholic party, under the leadership of Cardinal Stanislaus Hosius, Bishop of Ermeland (d. 1579), was very active. Pope Paul IV. sent a nuncio, Lipomani, to Poland, and urged the king to banish Laski and Vergerio from the country, and to suppress, with every power at his command, the rising heresy, if he would save his honor, his crown, and his soul. The weak king vacillated between the advice of Calvin and the threats of the Pope, and did nothing. He allowed the glorious opportunity to pass, and died in 1572, the last of the House of Jagellon. The nobles were likewise undecided, and many of them were carried away by the Unitarian heresy which began to spread in Poland in 1558.

During the interregnum which followed the death of Sigmund Augustus, the nobles, before electing a new king, concluded in 1573 a patriotic treaty of peace for the protection of religious freedom, under the name of \emph{Pax Dissidentium}—that is, of the Roman Catholic and the three evangelical Churches.\footnote{The Roman Catholics objected to being called \emph{Dissidentes,} and were opposed to the whole treaty.} They required Duke Henry of Anjou, the brother of the King of France and a violent enemy of the Huguenots, to accept the treaty as a condition of the crown, hoping to break it afterwards. On being peremptorily told by the Great Marshal, in the midst of the act of coronation, '\emph{Si non jurabis non regnabis}' he took the oath in spite of the remonstrance of the Romish party; but he left Poland in 1574, being called to the throne of France after the death of his brother, Charles IX. His Protestant successor, Stephen Bathori of Transylvania (1575–86), took the same oath, but afterwards joined the Roman Church and opened the door to the Jesuits This was the turning-point.

Under Sigmund III.—a Swedish prince, who had been educated and converted by Jesuits, and was elected king in 1587—there began a series of vexations and oppressions of the Protestants which gradually reduced them to a poor remnant, except in the Prussian part of Poland where the German element prevailed. Even Laski's relations and the four sons of Radziwill returned to the Roman Church; one of these sons became a cardinal; another made a pilgrimage to Jerusalem, and spent five thousand ducats for the purchase and destruction of Polish Bibles which his father had published (1563) at his expense.\footnote{Krasinski, p. 297.} Hence the great scarcity of this work. It was an essential part of the Jesuit counter-reformation to burn the whole Protestant literature, and thus, to suppress all independent thought. In this it succeeded only too well. The Polish nation, after the light of the gospel was extinguished, hastened step by step to its internal and external ruin.

THE CONSENSUS OF SENDOMIR.

After the death of Laski (1560) and Prince Radziwill (1567) the Protestants had no commanding leader, and felt the more the necessity of some union for their own safety. An organic union would have been the best, and would perhaps have made them strong enough to carry the king and the nobles with them. But for such a step they were not prepared. Instead of this the Lutherans (influenced by the liberal advice of the Melanchthonian divines of Wittenberg), the Calvinists, and the Bohemian Brethren effected a confederate union at the Synod of Sendomir,\footnote{A town on the Vistula in Little Poland. Krasinski and Gindely call it Sandomir.} April 14, 1570, and expressed it in the \emph{Consensus Sendomiriensis,} the only important confessional document of the evangelical Churches in Poland. It was published by authority, in Latin and Polish, in 1586, with a preface signed by Erasmus Gliczner, Lutheran Superintendent of Great Poland, in the name of the ministers of the Augsburg Confession, by John Laurentius, Superintendent of the Bohemian Brethren in Great Poland, and by Paulus Gilovius, Superintendent of the Reformed Churches in Little Poland.\footnote{The full title is '\emph{Consensus in fide et religione Christiana inter Ecclesias Evangelicas Majoris et Minoris Poloniæ, Magnique Ducatus Lithuaniæ et cæterarum ejus regni provinciarum, primo Sendomiriæ Anno MDLXX. in Synodo generali sancitus, et deinceps in aliis, ac demum in Wlodislaviensi generali Synodo Anno MDLXXXIII. confirmatus, et Serenissimis Poloniæ Regibus, Augusto, Henrico ac Stephano oblatus, nunc autem ex decreto Synodico in publicum typis editus. Anno Christi MDLXXXVI.}' This edition contains the supplementary resolutions of the Synods of Posen (1570), Cracow (1573), Petricow (1578), and Vladislav (1583). It was reprinted at Thorn, 1592 and 1596 (with the \emph{Acta et conclusiones synodi generalis Thoruniensis anni }1595); at Heidelberg, 1605; at Geneva, in the \emph{Corpus et Syntagma Conf.,} 1612 and 1654 (from the Heidelberg edition); at Frankfort-on-the-Oder, 1704 (with a Preface and German translation of Dr. Sam. Strimesius); and at Berlin, 1731, in Jablonski's \emph{Historia cons. Send.} Niemeyer (1840) gives the Latin text from the edition of Thorn, with all the supplements (pp. 551–591). Böckel excludes the Consensus (as not being strictly Reformed) from his collection. Beck gives the German text, but without the additions; and so also Dr. Nitzsch, in his \emph{Urkundenbuch der Evangelischen Union }(Bonn, 1853), pp. 72 sqq.}

The Consensus sets forth that the three orthodox evangelical Churches are agreed in the doctrines of God, the Holy Trinity, the Incarnation, the person of Christ, justification by faith, and other fundamental articles, as taught in the Augsburg, the Bohemian, and Helvetic Confessions, against papists, sectarians, and all enemies of the gospel; that in the unfortunate sacramentarian controversy they adopt that explanation of the words of institution which distinguishes (with Irenæus) between the earthly form and the heavenly substance in the Lord's Supper, and regards the visible elements not as mere signs, but as conveying to the believer truly through faith that which they represent.\footnote{Niemeyer, p. 554: '\emph{Convenimus in sententia verborum Domini nostri Jesu Christi, ut illa orthodoxe intellecta sunt a patribus, ac imprimis Irenæo, qui duabus rebus, scilicet terrena et cœlesti, hoc mysterium constare dixit; neque elementa signave nuda et vacua illa esse asserimus, sed simul reipsa } credentibus  \emph{exhibere et præstare } fide,  \emph{quod significant. Denique ut expressius clariusque loquamur, convenimus, ut credamus et confiteamur, } substantialem præsentiam Christi [not corporis et sanguinis Christi], \emph{non significari duntaxat, sed vere in cœna eo} [sc. Christo] \emph{vescentibus repræsentari, distribui, et exhiberi } corpus et sanguinem Domini  \emph{symbolis adjectis ipsi rei minime nudis, secundum Sacramentorum naturam.}' The Lutheran members demanded the phrase '\emph{præsentiam } corporis  \emph{Christi}' for '\emph{præsentiam} Christi,' and the insertion of the entire article of the Saxon Confession on the Lord's Supper. The first request was denied by the Calvinists and Bohemian Brethren; the second was granted, because the Saxon Confession uses the words '\emph{in hac communione vere et substantialiter adesse } Christum' (not corpus \emph{Christi}). See Gindely, \emph{Gesch. der Böhm. Brüder,} Vol. II. p. 86.}

Then follows a long extract on the sacraments from the Repetition of the Augsburg Confession, or Saxon Confession, which Melanchthon prepared in 1551 for the Council of Trent.

The Consensus thus adopts the later Melanchthonian or Calvinistic theory; it avoids the characteristic Lutheran terms (\emph{manducatio oralis,} etc.), and demands faith as the medium of receiving the matter represented by the elements. The doctrine of predestination was not touched, as there seems to have been no controversy about it.

In conclusion the Consensus acknowledges the orthodoxy and Christian character of the three parties, and pledges them to cultivate peace and charity, and to avoid strife and dissension, which greatly hinder the progress of the gospel. They should seal this compact by exchange of pulpits and of delegates to general synods, and by frequent sacramental intercommunion; each denomination retaining its peculiarities in worship and discipline which (according to the Augsburg and the Saxon Confessions) are consistent with the unity of the Church.

Then follow the signatures of noblemen and ministers.

Great joy was felt at this happy result, and was expressed by mutual congratulations and united praise of God.

A few weeks afterwards, May 20, 1570, a synodical meeting was held at Posen in the same spirit of union, and twenty brief supplementary articles were adopted for the purpose of confirming and preserving the Consensus.\footnote{\emph{Consignatio observationum necessariarum ad confirmandum et conservandum mutuum Consensum Sendomiriæ Anno DN. MDLXX. die }14 \emph{April, in vera religione Christiana initum inter Ministros Augustanæ Confessionis et Fratrum Bohemorum, Posaniæ eodem anno, Maii }20 \emph{facta, et a Ministris utriusque cœtus approbata ac recepta.} Printed in the \emph{Corpus et Syntagma Conf.,} and in Niemeyer, pp. 561–565.} One of the articles forbids polemics in the pulpit. When the people, who stood outside of the house where the meeting was held, heard the happy conclusion, they joined in the singing of the \emph{Te Deum,} with tears of joy and gratitude to God. The union was sealed on the following Sunday by two united services in the Lutheran church and in the Bohemian chapel.

The Consensus was again confirmed by the general synods at Cracow, 1573; Petricow, 1578; Vladislav, 1583; and Thorn, 1595. The last was the largest synod ever held in Poland.\footnote{See the Acts of these synods relating to the Consensus and to matters of discipline, in Niemeyer, pp. 565–591.}

The Lutherans who adhered to the Formula of Concord (1580) withdrew from the Consensus. But the spirit of union which produced it passed into the three Brandenburg Confessions of the seventeenth century, and revived in the Evangelical Union of Prussia.\footnote{See above, pp. 545 sqq. Comp. also Nitzsch, \emph{Urkundenbuch der Evangelischen Union,} pp. 80 sqq.}

\xsubsection{The Reformation in Hungary and the Confession of Czenger. A.D. 1557.}

§ 75. The Reformation in Hungary and the Confession of Czenger.

Literature.

I. The Latin text of the \emph{Confessio Czengerina,} or \emph{Hungarica,} in the \emph{Corpus et Syntagma Conf.,} and in Niemeyer, pp. 539–550; the German text in Böckel, pp. 851–863.

II. P. Ember (Reform.): \emph{Historia ecclesiæ reform. in Hungaria et Transylvania} (ed. Lampe). Utrecht, 1728.

Ribini (Luth.): \emph{Memorabilia Aug. Conf. in regno Hungariæ.} 1787, 2 vols.

\emph{Geschichte der evang. Kirche in Ungarn vom Anfang der Reformation bis }1850 [by Bauhofer, not named]. \emph{Mit einer Einleitung von Merle d’Aubigné.} Berlin, 1854.

Gieseler: \emph{Church History,} Vol. IV. pp. 258 sqq. (Am. ed.).

Baur: \emph{Geschichte der christl. Kirche,} Vol. IV. (1863), pp. 214 sqq., 552 sqq.

Ebrard: \emph{Kirchen- und Dogmengeschichte,} Vol. III. (1866), pp. 415–432.

E. L. Th. Henke (d. 1872): \emph{Neuere Kirchengeschichte} (ed. by W. Gass). Halle, 1874, Vol. I. pp. 352 sqq.

Burgovszky: Art. \emph{Ungarn,} in Herzog's \emph{Real-Encykl.} Vol. XVI. pp. 636 sqq.

Hungary, an extensive and fertile country on the banks of the lower Danube, once an independent kingdom, then united with the empire of Austria, and containing a mixed population of Magyars, Germans, Slowaks, Ruthenians, Croats, Serbs, etc., received the first seeds of the Christian religion from Constantinople; but the real apostle of the Hungarians was Stephen I. (979–1038), a king and a saint, who by persuasion and violence overthrew heathenism and barbarism, gave rich endowments to the churches and clergy, and brought his country into close contact with the Roman Church and the German Empire.

THE REFORMATION.

The way for the Reformation was prepared by Waldenses and Bohemian Brethren who sought refuge in Hungary from persecution. The writings of Luther found ready access among the German population, and were read with avidity, especially the one on the Babylonian Captivity of the Church. Many young Hungarians, among them Matthias Dévay (De Vay), called 'the Hungarian Luther,'\footnote{Dévay lived in the home of Luther, who calls him '\emph{vir honestus, gravis et eruditus.}' He sympathized, however, with Melanchthon in the eucharistic controversy, and inclined to the Calvinistic view, so as to cause complaint on the part of the strict Lutherans in Hungary (1544). See Luther's \emph{Letters,} Vol. V. p. 644 (ed. De Wette), and Henke, p. 355.} and Leonard Stöckel, studied at Wittenberg; others, as John Honter, at Basle; and on their return they introduced the new doctrines at Ofen, Cronstadt, and other cities, without any compulsion or aid from the government. It was a spontaneous movement of the people. Even some bishops and other dignitaries of the Roman Church became Protestants from conviction.

In 1545 a meeting of twenty-nine ministers at Erdöd adopted a creed of twelve articles in essential agreement with the Augsburg Confession. Another Lutheran synod at Medwisch (Medias), in 1548, drew up the \emph{Confessio Pentapolitana,} which represented five free cities in Upper Hungary, and was declared legal in 1555. The Saxon or German population of Hungary and Transylvania remained mostly Lutheran.

On the other hand, the majority of the Magyars or Hungarians proper (the ruling race in that country) were more influenced by the Latin writings of Melanchthon and Calvin than by the German of Luther, and during the violent eucharistic controversies in Germany embraced the Calvinistic creed, which they formally adopted at the Synod of Czenger, 1557, and which they nominally profess to this day.\footnote{We say nominally, for both the Reformed and Lutheran Churches of Hungary have been much affected by rationalism. This applies, however, to nearly all the State Churches of the Continent.} A large number of Magyar pastors left the Lutheran Confession and embraced Calvinism in 1563. The Presbyterian polity and discipline were introduced by the Synods of Tarczal, Göntz, and Debreczin. Thus the separation of the two evangelical Churches was completed.

Protestantism made rapid progress under Maximilian II. At the close of the sixteenth century the larger part of the people and the whole nobility, with the exception of three magnates, had accepted the Reformation. It gave a vigorous impulse to national life and literary activity. 'It is astonishing to see the amount of religious information which was then spread among the citizens and the lower classes, and the fertility of the press in places where now not even an almanac is printed.'\footnote{Burgovszky, l.c. p. 643.}

But under the reign of Rudolph II., King of Hungary from 1572 to 1608, began the counter-reformation of the Jesuits (among whom Peter Pázmány, a nobleman of Calvinistic parents, was the most successful in making converts), and a series of cruel persecutions by the Hapsburg rulers, urged on by the Popes, which continued for nearly two centuries, amid reactions, rebellions, civil wars, and wars with the Turks. A Jesuitical formula for the conversion of Hungarian Protestants pronounces awful curses on the evangelical faith, with the promise to persecute it by the sword. Whether genuine or not, it shows the intense bitterness of the conflict.\footnote{See above, p. 92, note 2.} General Caraffa, a cruel papist, erected in the market-place at Eperjes a bloody scaffold, or 'slaughter-bank,' where for several months daily tortures and executions by fire and sword took place (1657).\footnote{Sismondi and Merle d’Aubigné (l.c. p. ix.) state that the persecutions of the Hungarian Protestants surpassed in cruelty the persecutions of the Huguenots under Louis XIV.}

Protestantism survived these trials. Joseph II., by his famous Edict of Toleration, Oct. 29, 1781, secured to the followers of the Augsburg and Helvetic Confessions liberty of conscience and public worship. His brother and successor, Leopold, confirmed it in 1791. The remaining restrictions were removed in 1848. The present number of Protestants in Hungary is about three millions, or one fifth of the whole population (which in 1869 amounted to fifteen millions and a half). The Lutheran Confession prevails among the German population; the followers of the Reformed or Helvetic Confession are twice as numerous, and are mostly Magyars.

THE HUNGARIAN CONFESSION.

The Hungarian Confession, or Confessio Czengerina, was prepared and adopted at a Reformed Synod held at Czenger in 1557 or 1558,\footnote{The date is uncertain.} and printed in 1570 at Debreczin.\footnote{Debreczin is a royal free city in the northeastern part of the Hungarian Lowland, with about fifty thousand inhabitants, and contains the principal Calvinistic college of the kingdom. In 1849 it was the seat of the revolutionary government of Kossuth, and the independence of Hungary was there declared in the Reformed Church.}

It treats, in brief articles or propositions, of the Triune God, of Jesus Christ, the Holy Spirit, the Scripture designations of the Holy Spirit, the rules for explaining the phrases concerning God, the law and the gospel, the rights and sacraments of the Church, Christian liberty, election, the cause of sin, and the only mediator Jesus Christ. It is preceded by a strong Biblical argument against the anti-Trinitarians and Socinians, who had spread in Transylvania. It vehemently rejects the Romish transubstantiation and the Lutheran '\emph{sarcophagia},'\footnote{'\emph{Damnamus Papisticum delirium . . . primo panem transsubstantiari, et offerri in missa: deinde sola accidentia panis manere. . . . Ita et eorum insaniam damnamus, qui asserunt Sarcophagiam, id est, ore corporali sumi corpus Christi naturale, sanguinolentum, sine ulla mutatione et transsubstantiatione.}'—Niemeyer, pp. 544 sq. The severe judgment of the Lutheran doctrine was a retaliation for the condemnation of Zwingli and Calvin as sacramentarians by a Lutheran Synod of Hermanstadt. Ebrard, Vol. III. p. 424.} but also the 'sacramentarian' view of a purely symbolical presence, and teaches that Christ is truly though spiritually present, and communicates himself in the Lord's Supper as the living bread and the celestial drink, with all his gifts, to the believer.\footnote{'\emph{Rejicimus et eorum delirium, qui Cœnam Domini } vacuum signum,  \emph{vel Christi absentis tantum } memoriam  \emph{his signis recoli docent. Nam sicut Christus est } Amen, testis fidelis, verax, veritas et vita  . . . \emph{ita Cœna Domini est præsentis et infiniti æternique Filii Dei unigeniti a Patre memoria: qui se et sua bona, carnem suam et sanguinem suum, id est, panem vivum et potum cœlestem, Spiritus Sancti ope per verbum promissionis gratiæ, offert et exhibet electis fide vera evangelium Christi apprehendentibus.}'—Page 545.} It defends infant baptism against the Anabaptists. It teaches a free election, but is silent about reprobation, and denies that God is the author of sin. Later synods professed more clearly the doctrine of predestination and the perseverance of saints.

This Confession presents some original and vigorous features, but has only a secondary historical importance. It was practically superseded by the Second Helvetic Confession of 1566, which is far superior, and was subscribed by the entire Reformed clergy of Hungary convened at Debreczin in 1567. The Heidelberg Catechism was also introduced.

\xsection{V. the Anglican Articles of Religion.}

\xsubsection{The English Reformation.}

§ 76. The English Reformation.

Literature.

I. Works on the Thirty-nine Articles.

(a) Historical.

Charles Hardwick (B.D., Archdeacon of Ely, and Christian Advocate in the University of Cambridge, d. 1859): \emph{A History of the Articles of Religion; to which is added a Series of Documents from A.D.} 1536 \emph{to A.D.} 1615, \emph{together with Illustrations from Contemporary Sources.} Cambridge, 1851 (reprinted in Philadelphia, 1852); second edition, thoroughly revised, Cambridge, 1859 (pp. 399).

(b) Commentaries.

Thomas R. Jones:  \emph{An Exposition of the Thirty-nine Articles by the Reformers; being Extracts from the Works of Latimer, Ridley, Cranmer, Hooper, Jewell, Philpot, Pilkington, Coverdale, Becon, Bradford, Sandys, Grindal, Whitgift,} etc. London, 1849.

Thomas Rogers (Chaplain to Archbishop Bancroft): \emph{The Catholic Doctrine of the Church of England, an Exposition of the Thirty-nine Articles.} London, 1579, 1585, 1607, and other editions (under various titles). Newly edited by \emph{J. J. S. Perowne,} for 'The Parker Society,' Cambridge, 1854. This is the oldest commentary, and was countenanced by Bancroft, to whom it was dedicated.

Gilbert Burnet (Bishop of Salisbury; b. 1643, d. 1715): \emph{An Exposition of the Thirty-nine Articles of the Church of England. }Oxford, 1814 (Clarendon Press), and other editions. Revised, with notes, by James R. Page.

Richard Laurence, L.L.D. (formerly Reg. Prof. of Hebrew in Oxford): \emph{An Attempt to illustrate those Articles of the Church of England which the Calvinists improperly consider as Calvinistical. }In eight sermons (Bampton Lectures for 1834). Oxford, third edition, 1838.

Edward Harold Browne (b. 1811, Bishop of Winchester since 1873, formerly of Ely): \emph{An Exposition of the Thirty-nine Articles, Historical and Doctrinal. }London, 1850–53, in two vols.; since often republished in one vol. (ninth edition, 1871); Amer. edition, with notes by Bishop Williams of Connecticut, New York, 1865.

A. P. Forbes (Bishop of Brechin): \emph{An Explanation of the Thirty-nine Articles, with an Epistle dedicatory to the Rev. E. B. Pusey, D.D.} Oxford and London, 1867. (High Church.)

E. W. Jelf (Canon of Christ Church, Oxford): \emph{The Thirty-nine Articles of the Church of England explained in a Series of Lectures.} Edited by J. R. King. London, 1873.

II. History of the Reformation in England.

(a) Documents and Contemporary Sources.

Works of the English Reformers, published by 'The Parker Society,' Cambridge, 1841–54, fifty-four vols. Contains the writings of Cranmer, Ridley, Latimer, Hooper, Sandys, Coverdale, Jewell, Grindal, Whitgift, the Zurich Letters, etc.

The State Calendars, now being published under the direction of the Master of the Rolls.

John Foxe (one of the Marian exiles, d. 1587): \emph{Acts and Monuments of the Church, or Book of Martyrs.} London, 1563, and often in three or more volumes. Not accurate, but full of facts told in a forcible style.

Wilkins:  \emph{Concilia Magnæ Brittaniæ et Hiberniæ} (446–1717). Four vols. folio. 1736 sq.

E. Cardwell:  \emph{Documentary Annals of the Church of England} (1546–1716), Oxford, 1844, 2 vols.; \emph{Synodalia} (1547–1717), Oxford, 1842, 2 vols.; \emph{The Reformation of the Laws in the Reigns of Henry VIII., Edward VI., and Elizabeth,} Oxford, 1850.

(b) Historical Works.

John Strype (a most laborious and valuable contributor to the Church history and biography of the English Reformation period; b. 1643, d. 1737): \emph{Ecclesiastical Memorials . . . of the Church of England under King Henry VIII., Edward VI., and Queen Mary }(London, 1725–37; Oxford, 1822, 3 vols.); \emph{Annals of the Reformation . . . in the Church of England during Queen Elizabeth's Happy Reign }(London, 1738; Oxford, 1824, 4 vols.; \emph{Memorials of Archbishops Cranmer} (2 vols.), \emph{Parker} (3 vols.), \emph{Grindal }(1 vol.), \emph{Whitgift} (3 vols.). See his \emph{Complete Works,} Oxford, 1822–40, in twenty-seven vols.

Gilbert Burnet:  \emph{The History of the Reformation of the Church of England.} London, 1679 sqq., 7 vols., and other editions. New edition by Pocock.

C. Hardwick:  \emph{History of the Christian Church during the Reformation,} third edition (by W. Stubbs). London, 1873, pp. 165–249.

Fred. Seebohm:  \emph{The Oxford Reformers, Colet, Erasmus, and More.} London, 1869. The same: \emph{The Era of the Protestant Revolution. }1874.

The Church Histories of England and of the English Reformation by J. Collier (non-Juror), Dodd (Rom. Cath.), Thos. Fuller (Royalist; \emph{Church History of Great Britain until} 1658 and \emph{The Worthies of England}), Neal (\emph{History of the Puritans}), Heylin,  Soames,  Massingbeard,  Short,  Blunt,  Waddington,  Weber,  d’Aubigné,  Fisher.

Also the secular Histories of England by Hume,  Macaulay (the introductory chapter), Hallam (\emph{Constitut. Hist}.), Lingard (Rom. Cath.), Knight,  Froude,  Ranke,  Green, in the sections on the Reformation period.

The last and, in its final results, the most important chapter in the history of the reformation was acted in that remarkable island which has become the chief stronghold of Protestantism in Europe, the ruler of the waves, and the pioneer of modern Christian civilization and constitutional liberty. The Anglo-Saxon race is intrusted by Providence with the sceptre of empire in its eastward and westward coarse. The defeat of the Armada was that turning-point in history when the dominion in which the sun never sets passed from Roman Catholic Spain to Protestant England.

The Reformation in Britain, favored by insular independence, was a national political as well as ecclesiastical movement, and carried with it Church and State, rulers and subjects; while on the Continent it encountered a powerful opposition and Jesuitical reaction. It began with outward changes, and was controlled by princes, bishops, and statesmen rather than by scholars and divines; while in other countries the reform proceeded from the inner life of religion and the profound study of the Scriptures. Good and bad men, from pure and low motives, took part in the work, but were overruled by a higher power for a noble end.\footnote{Robert Southey (\emph{Life of Wesley,} Vol. I. p. 266, Harpers' edition) says: 'In England the best people and the worst combined in bringing about the Reformation, and in its progress it bore evident marks of both.'} England produced no reformers of such towering genius, learning, and heroism as Luther and Calvin, but a large number of learned and able prelates and statesmen, and a noble army of martyrs worthily led by Cranmer, Latimer, Ridley, Hooper, and Rogers. It displayed less theological depth and originality than Germany and Switzerland, where the ideas and principles of the Reformation were wrought out, but a greater power of practical organization. It gave the new ideas a larger field of action and application to all the ramifications of society and all departments of literature, which entered upon its golden age in the reign of Elizabeth, and which, in wealth of genius and in veneration for the truths of Christianity, far surpassed that of any other nation.\footnote{Fisher (\emph{The Reformation,} p. 533): 'The boldness and independence of the Elizabethan writers, their fearless and earnest pursuit of truth, and their solemn sense of religion, apart from all asceticism and superstition, are among the effects of the Reformation. This is equally true of them as it is of Milton and of the greatest of their successors. Nothing save the impulse which Protestantism gave to the English mind, and the intellectual ferment which was engendered by it, will account for the literary phenomena of the Elizabethan times.' Even that brilliant and racy French critic, Taine, must acknowledge the constant influence of 'the grave and grand idea of religion, of faith and prayer,' upon such writers as Bacon, Raleigh, Burton, and Sir Thomas Browne.} Although at first despotic and intolerant, English Protestantism by its subsequent development became the guardian of civil and religious liberty. The fierce struggle between 'the old and new learning' lasted for more than a century, and passed through a baptism of blood which purified and fertilized the soil of England and became the seed of new colonies and empires beyond the sea.

The British Reformation is full of romantic interest, and developed a great variety of strongly marked characters, who still excite the passions, prejudices, and contradictory judgments of writers and readers. It is a succession of tragedies; it abounds in actions and reactions, in crimes and punishments, in changes of fortune, in men and women elevated to the pinnacle of power and happiness and hurled to the abyss of disgrace and misfortune. It furnishes a striking illustration of the truth that the history of the Church, as well as of the world, is a judgment of the Church. This idea of righteous retribution imparts a thrilling moral effect to the tragedies of Shakspere, who lived at the close of these shifting scenes, and gathered from them his marvelous knowledge of human nature, in all its phases and conditions, such as no poet ancient or modern ever possessed.

The richest fruit of the British Reformation is the translation of the Bible—the work of three generations, the best ever made, and to this day the chief nursery of piety among the Protestant denominations of the English-speaking race; and next to it that noble responsive liturgy which animates and regulates the devotions of the Episcopal communion on land and sea. These two works are truly national institutions, and command a veneration and affection above all other books, not only by their sacred contents, but also by their classical diction, which sounds in the ear like solemn music from a higher and better world.

EPOCHS OF THE ENGLISH REFORMATION.

The history of the English Reformation naturally divides itself into four periods:

1. From 1527 to 1547. The abolition of the authority of the Roman See over England and the dissolution of the monasteries by Henry VIII. This was chiefly a destructive process and a political change of the supreme governing power of the Church, prompted by unworthy personal motives, but it prepared the way for the religious reformation under the following reign. The despotic and licentious monarch, whom Leo X. rewarded for his book against Luther with the title 'Defender of the Faith,' remained a Catholic in belief and sentiment till his death; he merely substituted king-worship for pope-worship, a domestic tyranny for a foreign one, by cutting off the papal tiara from the episcopal hierarchy and placing his own crown on the bleeding neck; but he could not have effected so great a revolution without the sanction of Parliament and a strong clerical and popular current towards ecclesiastical independence and reform, which showed itself even before his breach with Rome, and became dominant under his successor.

2. From 1547 to 1553. The introduction of the Reformation in doctrine and worship under Edward VI., Henry's only son, and the commencing conflict between the semi-Catholic and the Puritan tendencies. The ruling genius of this period was Archbishop Cranmer, the Melanchthon of England, who by cautious trimming and facile subservience to Henry had saved the cause of the Reformation through the trials of a despotic reign for better times.

3. From 1553 to 1558. The papal reaction under Henry's oldest daughter, Mary Tudor, that 'unhappiest of queens and wives and women.'\footnote{Tennyson, in \emph{Queen Mary,} act v. scene 2.} She had more Spanish than English blood in her veins, and revenged the injustice done to her mother, Catharine of Aragon. Her short but bloody reign was the period of Protestant martyrdom, which fertilized the soil of England, and of the exile of about eight hundred Englishmen, who were received with open arms on the Continent, and who brought back clearer and stronger views of the Reformation. The violent restoration of the old system intensified the hatred of Popery, and forever connected it in the English mind with persecution and bloodshed, with national humiliation and disgrace. 'The tale of Protestant sufferings was told with wonderful pathos and picturesqueness by John Foxe, an exile during the persecution, and his ``Book of Martyrs,'' which was (under the following reign) set up by royal order in the churches for public reading, passed from the churches to the shelves of every English household.'

4. From 1558 to 1603. The permanent establishment of the Reformed Church of England in opposition both to Roman Catholic and to Puritan dissent during the long, brilliant, and successful reign of Queen Elizabeth.

This masculine woman, the last and the greatest of the Tudors, inherited the virtues and vices of her Catholic father (Henry VIII.) and her Protestant mother (Anne Boleyn).\footnote{Her character is admirably drawn by Froude, and by the latest historian of England, J. R. Green, \emph{A Short History of the English People }(London, 1875), pp. 362–370.} She was endowed with rare gifts by nature, and favored with, the best education; she was brave and bold, yet prudent and cautious; fond of show, jewelry and dress, yet parsimonious and mean; coldly intellectual, high-tempered, capricious, haughty, selfish, and vain, and well versed in the low arts of intrigue and dissimulation. She trusted more in time and her good fortune than in Almighty God. She was destitute of religious enthusiasm, and managed the Church question from a purely political point of view. She dropped the blasphemous title 'Head of the Church of England,' and was content to be the supreme 'Governor' of the same.\footnote{Parliament, in the act of supremacy (1534), declared King Henry, his heirs and successors, to be 'the only supreme head, on earth, of the Church of England, called the \emph{Anglicana Ecclesia.}' For denying this royal supremacy in spiritual matters, More and Fisher suffered martyrdom. The thirty-seventh of the Elizabethan Articles modifies it considerably, but still claims for 'the Queen's Majesty the chief power in this Realm of England, . . . unto whom the chief government of all estates, whether they be ecclesiastical or civil, in all causes doth appertain,' etc. Elizabeth disclaimed the sacerdotal character which her father had assumed, but retained and exercised the vast power of appointing her prelates, summoning and dissolving convocations, sanctioning creeds and canons, and punishing heresies and all manner of abuses with the civil sword.} But with this limitation the royal supremacy was the chief article in her creed, and she made her bishops feel her power. 'Proud prelate,' she wrote to the Bishop of Ely, when he resisted the spoliation of his see in favor of one of her favorites, 'you know what you were before I made you what you are! If you do not immediately comply with my request, by God! I will unfrock you.' As a matter of taste she liked crucifixes, images, and the gorgeous display of the Roman hierarchy and ritual; and, being proud of her own virginity, she disliked the marriage of the clergy; she insulted the worthy wife of Archbishop Parker by refusing to call her 'Madam,' the usual address to married ladies. But she had the sagacity to perceive that her true interests were identified with the cause of Protestantism, and she maintained it with a strong arm, aided by the ablest council and the national sentiment, against the excommunication of the Pope, the assaults of Spain, and the intrigues of the Jesuits at home. This is the basis of the popularity which she enjoyed as a ruler with all classes of her subjects except the Romanists.

Her ecclesiastical policy at home was a system of compromise in the interest of outward uniformity. It was fortified by a penal code which may be explained though not justified by the political necessities and the general intolerance of the times, but which was nevertheless cruel and abominable, and has been gradually swept away by the progress of a nobler and more enlightened policy of religious liberty.

As in the case of her predecessors, we should remember that the policy of Elizabeth was merely the outward frame which surrounds the true inward history of the religious movement of her age. The doctrinal reformation with which we are concerned was begun in the second and completed in the fourth period.

With the reign of Elizabeth ended the great conflict with Rome. It was followed by the internal conflict between Puritanism and Episcopacy, which, after a temporary triumph of the former under Cromwell, resulted in the re-establishment of the Episcopal Church and the expulsion of Puritanism (1662), until another revolution (1688) brought on the final downfall of the treacherous Stuarts and the toleration of the Dissenters, who thereafter represented, in separate organizations, the left or radical wing of English Protestantism.

\xsubsection{The Doctrinal Position of the Anglican Church and Her Relation to Other Churches.}

§ 77. The Doctrinal Position of the Anglican Church and her Relation to other Churches.

The Reformed Church of England occupies an independent position between Romanism on the one hand, and Lutheranism and Calvinism on the other, with strong affinities and antagonisms in both directions. She nursed at her breasts Calvinistic Puritans, Arminian Methodists, liberal Latitudinarians, and Romanizing Tractarians and Ritualists. This comprehensiveness of the Church as a whole is quite consistent with the narrowness and exclusiveness of particular parties. It repels and attracts; it caused the large secessions which occurred at critical junctures in the seventeenth, eighteenth, and nineteenth centuries, but it also explains the individual accessions which she continually though quietly receives from other Churches.

The English mind is not theorizing and speculative, but eminently practical and conservative; it follows more the power of habit than the logic of thought; it takes things as they are, makes haste slowly, mends abuses cautiously, and aims at the attainable rather than the ideal. The Reformation in England was less controlled by theology than on the Continent, and more complicated with ecclesiastical and political issues. Anglican theology is as much embodied in the episcopal polity and the liturgical worship as in the doctrinal standards. The Book of Common Prayer is catholic, though purged of superstitious elements; the Articles of Religion are evangelical and moderately Calvinistic.\footnote{The ingenious and sophistical attempt of Dr. John Henry Newman, in his famous \emph{Tract Number Ninety} (Oxford, 1841), to un-Protestantize the Thirty-nine Articles, has been best refuted by his own subsequent transition to Rome. As a specimen of this non-natural interpretation we mention what he says on Art. XI., which teaches as 'a most wholesome doctrine' 'that we are justified by faith only.' This means that faith is the sole \emph{internal} instrument of justification, while baptism is the sole \emph{outward} means and instrument; it does not interfere with the doctrine that good \emph{works} are also a means of justification (pp. 21 sqq.). That is, the Protestant doctrine of justification by faith \emph{alone,} which the Council of Trent condemned, is identical with the Romish doctrine of justification by faith \emph{and works,} which the same Council taught. A more learned and elaborate work, which minimizes the Protestantism of the Articles and makes them bear a catholic sense, is the \emph{Explanation} by the late Bishop Forbes of Brechin, above quoted.} The hierarchical, sacerdotal, and sacramental systems of religion are congenial and logically inseparable; they moderate and check the Protestant tendency, and if unduly pressed they become Romanizing. In great minds we often find great antagonisms or opposite truths dwelling together unreconciled; while partisans look only at one side. Augustine, Luther, and even the more logical Calvin, believed in divine sovereignty and human responsibility, free election and sacramental grace, and combined reverence for Church authority with independence of private judgment. The English Church leaves room for catholic and evangelical, mediæval and modern ideas, without an attempt to harmonize them; but her parties are one-sided, and differ as widely as separate denominations, though subject to the same bishop and worshiping at the same altar. She is composite and eclectic in her character, like the English language; she has more outward uniformity than inward unity; she is fixed in her organic structure, but elastic in doctrinal opinion, and has successively allowed opposite schools of theology to grow up which claim to be equally loyal to her genius and institutions. She has lost in England by those periodical separations which followed her great religious movements (the Puritan, the Methodist, the Anglo-Catholic) nearly one half of the nation she once exclusively controlled; yet she remains to this day the richest and strongest national Church in Protestant Christendom, and exercises more power over England than Lutheranism does over Germany or Calvinism over Switzerland and Holland. In the United States the Protestant Episcopal Church is numerically much smaller than most of the denominations which in England were cast out or voluntarily went out from the established Church as Non-conformists and Dissenters; but she will always occupy a commanding position among the higher classes and in large cities, because she represents the noble institutions and literature of the aristocratic, conservative, and venerable Church of England.

THE MELANCHTHONIAN INFLUENCE.

Germany received Roman Catholic Christianity from England through Winfrid or Boniface, and in turn gave to England the first impulse of the evangelical Reformation. The writings of Luther were read with avidity by students in Oxford and Cambridge as early as 1527. Cranmer spent some time in Germany, and was connected with it by domestic ties.\footnote{His second wife, whom he secretly married in 1532, before his elevation to the primacy (March, 1533), was a niece of the Lutheran divine Osiander at Nürnberg, who subsequently excited a violent controversy about the doctrine of justification.} Henry VIII. never overcame his intense dislike of Luther, kindled by their unfortunate controversy on the seven sacraments, and strengthened by Luther's breach with Erasmus; but he respected Melanchthon for his learning and wisdom, and invited him to assist in reforming the English Church.\footnote{Melanchthon was twice called to England in 1534 ('\emph{Ego jam alteris literis in Angliam vocor}'). In 1535 he dedicated an edition of his \emph{Loci} to Henry, at the request of Barnes, who thought it would promote the progress of the Reformation. Henry renewed the invitation in 1538, and requested the Elector of Saxony to send '\emph{Dominum Philippum Melancthonem, in cuius excellenti eruditione et sano judicio a bonis omnibus multa spes reposita est},' together with some other learned men, to England. Under Edward VI. Melanchthon was called again, and in 1553 he was appointed Professor of Divinity in Cambridge, but he never saw England. See Laurence, l.c. pp. 198 sqq.; Hardwick, \emph{Hist. of the Art.} pp. 52 sqq.; C. Schmidt, \emph{Phil. Mel.} pp. 283–289.} He entered into negotiations with the Wittenberg divines and the Lutheran princes of the Smalcald League, but chiefly from political motives and without effect.

Under Edward VI. the influence of the Melanchthonian theology, as embodied in the Augsburg Confession (1530) and the Suabian Confession (1552), became more apparent, and can be clearly traced in Cranmer's earlier writings, in some of the Articles of Religion, and in those parts of the Book of Common Prayer which were borrowed from the 'Consultation' of Archbishop Hermann of Cologne, compiled by Bucer and Melanchthon (1543). Hence the English Church has been called sometimes by Lutheran divines an \emph{Ecclesia Lutheranizans.}

But the peculiar views of Luther on the real presence and the ubiquity of Christ's body found no congenial soil in England. Cranmer himself abandoned them as early as Dec. 14, 1548, when a public discussion was held in London on the eucharist; and he adopted, together with Ridley, the Calvinistic doctrine of a virtual presence and communication of Christ's glorified humanity. He held that 'Christ is figuratively in the bread and wine, and spiritually in them that worthily eat the bread and drink the wine; but, on the other hand, contended that our blessed Lord is really, carnally, and corporally in heaven alone, from whence he shall come to judge the quick and the dead.'\footnote{So his ultimate doctrine is correctly stated by Hardwick, \emph{History of the Reformation,} p. 209. Cranmer wrote very extensively on the eucharist, and especially against the Romish mass. See the first volume of the Parker Society's edition of his \emph{Works.} His change of view is due to the influence of the book of Ratramnus (Bertram) against transubstantiation, the tract of Bullinger on the eucharist, and the personal influence of Ridley, Peter Martyr, and Bucer. Bishop Browne says (on Art. XXVIII. Sect. I. p. 711 of the Am. ed.): 'Both Cranmer and Ridley, to whom we are chiefly indebted for our formularies, maintained the doctrine nearly identical with that maintained by Calvin, and before him by Bertram. . . . These sentiments of our Reformers were undoubtedly embodied in our Liturgy and Articles. . . . In the main, Calvin, Melanchthon in his later views, and the Anglican divines were at one.' John Knox entirely agreed with Cranmer in the Reformed doctrine of the eucharist, and he objected only to the kneeling posture, which led to the insertion of a special rubric in the Prayer-Book. See Lorimer, \emph{John Knox in England,} pp. 49 and 145.} This doctrinal change was embodied (1552) in the revision of the first Prayer-Book of Edward VI.; the prayer of oblation was converted into a thanksgiving, and the old formula of distribution, which was compatible even with a belief in transubstantiation ('The Body of our Lord Jesus Christ,' etc.), was replaced by another which a Zwinglian may approve ('Take and eat this in remembrance that Christ died for thee,' etc.). In the Elizabethan Service-Book the two formulas were combined (the second being an explanation of the first), and have ever since continued in use.

In the violent controversies which agitated Germany after Luther's death, and which led to the Formula of Concord, England sided with the milder Melanchthonian school. Queen Elizabeth made an effort to prevent the adoption of the Formula and the condemnation of the Reformed doctrines.\footnote{See above, p. 335.}

THE ZWINGLIAN AND CALVINISTIC INFLUENCE.

The doctrines of Zurich and Geneva began to spread in England under the reign of Edward VI. Calvin, whose books were prohibited by Henry VIII. (in 1542), corresponded freely with the Duke of Somerset (Oct. 22, 1548), Edward VI., and Cranmer, and urged a more thorough reformation of doctrine and discipline, and a better education of the clergy, but left episcopacy untouched, which he was willing to tolerate in England as well as in the kingdom of Poland.\footnote{Stähelin, Vol. II. pp. 51 sqq., discusses at length Calvin's correspondence with England. Hardwick speaks of 'the \emph{obtrusive} letters of Calvin;' but his counsel was solicited from every direction. In the controversy of the English exiles at Frankfort both parties (Cox and Knox) appealed to the Genevan Reformer for advice. Cranmer requested him to write often to King Edward. See Calvin to Farel, June 15, 1551 (\emph{Opera,} Vol. XIV. fol. 133): '\emph{Cantuariensis nihil me utilius facturum admonuit, quam si ad Regem sæpius scriberem. Hoc mihi longe gratius, quam si ingenti pecuniæ summa ditatus forem.}' Viret informed Farel in the same year and month (ibid. fol. 131), that the king sent to Calvin '\emph{coronatos centum et libellum a se conscriptum gallice in papatum, cuius censuram a Calvino exigit. . . . Accepit Calvinus a multis Angliæ proceribus multas literas plenas humanitatis. Omnes testantur se ejus ingenio et laboribus valde oblectari. Hortantur ut sæpe scribat. Protector scripsit nominatim.}'} His controversy with Pighius about predestination excited considerable sympathy in England (1552), and his doctrine of the eucharist gained ground more rapidly. Cranmer called to his aid prominent Reformed and Unionistic divines, such as Peter Martyr, Ochino, Laski, Bucer, and Fagius, and gave them high positions in Oxford, Cambridge, and London. It is characteristic of his catholicity of spirit that in 1548 he conceived the plan of inviting Melanchthon of Wittenberg, Bullinger of Zurich, Calvin of Geneva, Bucer of Strasburg, Peter Martyr, Laski, and others to Lambeth for the purpose of drawing up a union creed for all evangelical Churches.\footnote{Strype's \emph{Memorials of Cranmer,} Vol. I. p. 584; Hardwick, \emph{History of the Reformation,} p. 212.} John Hooper, who had resided two years at Zurich, was made Bishop of Gloucester (1551), although he went even beyond Bullinger and Calvin in matters of clerical vestments and ceremonies, and may be called a forerunner of Puritanism. He died heroically for his faith under Mary (1555). John Knox was elected one of the chaplains of Edward VI., and was offered the bishopric of Rochester, which he declined. He exerted considerable influence, and would no doubt have retained it under Elizabeth, had he not (together with his teacher and friend, Calvin) incurred her personal dislike by his trumpet-blast 'against the monstrous regimen of women,' which was provoked by the fatal misgovernment of her sister.\footnote{The influence of Knox upon the English Reformation has been more fully brought to light from the Knox Papers in Dr. Williams's library at London by Dr. Peter Lorimer, in \emph{John Knox and the Church of England }(London, 1875), pp. 98 sqq.}

Under the reign of Mary the English exiles formed the closest ties of personal and theological friendship with the Reformers of Switzerland, and on their return under Queen Elizabeth they took the lead in the restoration and reconstruction of the Reformed Church of England. Bishop Jewel, the final reviser of the Thirty-nine Articles, wrote to Peter Martyr at Zurich (Feb. 7, 1562): 'As to matters of doctrine, we have pared every thing away to the very quick, and do not differ from you by a nail's breadth; for as to the ubiquitarian [i.e., the Lutheran] theory there is no danger in this country. Opinions of that kind can only gain admittance where the stones have sense.'\footnote{Zurich Letters, second series, I. 100. Prof. Fisher, in quoting this passage, adds the just remark (\emph{The Reformation,} p. 341): 'There is no need in bringing further evidence on this point, since the Articles themselves explicitly assert the Calvinistic view [on the Lord's Supper]. In speaking of the English Reformers as Calvinists, it is not implied that they derived their opinions from Calvin exclusively, or received them on his authority. They were able and learned men, and explored the Scriptures and the patristic writers for themselves. Yet no name was held in higher honor among them than that of the Genevan Reformer.'}

Bullinger's 'Decades' were for some time the manual of the clergy. Afterwards Calvin's 'Institutes' became the text-book of theology in Oxford and Cambridge.\footnote{When Robert Sanderson (Professor of Theology in Oxford, 1642, afterwards Bishop of Lincoln, d. 1663) began to study theology in Oxford about 1606, he was recommended, as was usual at that time, to read Calvin's \emph{Institutes,} 'as the best and perfectest system of divinity, and the fittest to be laid as the ground-work in the study of this profession.' Blunt, \emph{Dictionary of Sects,} etc., p. 97. Comp. Hooker's judgment below, p. 607.} Even his Catechism was ordered to be used by statute in the universities (1587). Next to him his friend and successor, Beza, was for many years the highest theological authority. The University of Cambridge, in thanking him for the valuable gift of Codex D of the New Testament, in 1581, acknowledges its preference for him and John Calvin above any men that ever lived since the days of the Apostles.\footnote{\par '\emph{Nam hoc scito, post unicæ scripturæ sacratissimam cognitionem, nullos unquam ex omni memoria temporum scriptores extitisse, quos memorabili viro Joanni Calvino tibique præferamus.}' See Scrivener's \emph{Codex Bezæ,} Introd. p. vi., and his \emph{Introd. to the Critic. of the New Testament,} second edition, 1874, p. 112. Scrivener regards this veneration as an ill omen 'for the peace of the English Church.'} Beza's editions of the Greek Testament, his elegant Latin translation, and exegetical notes were in general use in England during the reigns of Elizabeth and James, and were made the chief basis not only of the Geneva Bible (1560), but also of the revision of the Bishops' Bible under King James (1611).\footnote{See my tract on the \emph{Revision of the English Version of the New Testament,} pp. 28, 29, and Westcott's \emph{History of the English Bible,} pp. 294 sq. A number of errors in the English Version, as well as excellences, can be traced to Beza.}

It is not too much to say that the ruling theology of the Church of England in the latter half of the sixteenth and the beginning of the seventeenth century was Calvinistic.\footnote{Macaulay (in his introductory chapter, p. 39, Boston edition) says: 'The English Reformers were eager to go as far as their brethren on the Continent. They unanimously condemned as anti-Christian numerous dogmas and practices to which Henry had stubbornly adhered, and which Elizabeth reluctantly abandoned. Many felt a strong repugnance even to things indifferent, which had formed part of the polity or ritual of the mystical Babylon.'} The best proof of this is furnished by the 'Zurich Letters,'\footnote{So called because they are mostly derived from the extensive Simler Collection of Zurich, where the Marian exiles, as Bishop Burnet says, 'were entertained both by the magistrates and the ministers—Bullinger, Gualter, Weidner, Simler, Lavater, Gesner, and all the rest of that body—with a tenderness and affection that engaged them to the end of their lives to make the greatest acknowledgments possible for it.' The correspondence was published by the Parker Society (Cambridge, 1842–47, in four vols.), in two series, the first of which covers the reigns of Henry VIII., Edward VI., and Mary; the second and more important the reign of Elizabeth (1558–1602). They include letters of most of the English Reformers and leading bishops and divines to the Swiss Reformers, with their answers, and are noble monuments of Christian and theological friendship.} extending over the whole period of the Reformation, the Elizabethan Articles, the Second Book of Homilies (chiefly composed by Bishop Jewel), the Lambeth Articles, the Irish Articles, and the report of the delegation of King James to the Calvinistic Synod of Dort.\footnote{\emph{The Suffrage of the Divines of Great Britain concerning the Articles of the Synod of Dort signed by them in the Year }1619. London, 1624. There is, however, at the close of this document (p. 176) a wholesome warning 'concerning the mystery of \emph{reprobation},' that it be 'handled sparingly and prudently,' and that 'those fearful opinions, and such as have no ground in the Scriptures, be carefully avoided, which tend rather unto desperation than edification, and do bring upon some of the Reformed Churches a grievous scandal.'}

EPISCOPACY.

This theological sympathy between the English and the Continental Churches extended also to the principles of Church government, which was regarded as a matter of secondary importance, and subject to change, like rites and ceremonies, 'according to the diversities of countries, times, and men's manners, so that nothing be ordained against God's Word' (Art. XXXIV.). The difference was simply this: the English Reformers, being themselves bishops, retained episcopacy as an ancient institution of the Church catholic, but fully admitted (with the most learned fathers and schoolmen, sustained by modern commentators and historians) the original identity of the offices of bishop and presbyter; while the German and Swiss Reformers, being only presbyters or laymen, and opposed by their bishops, fell back from necessity rather than choice upon the parity of ministers, without thereby denying the human right and relative importance or expediency of episcopacy as a superintendency over equals in rank. The more rigid among the Puritans departed from both by attaching primary importance to matters of discipline and ritual, and denouncing every form of government and public worship that was not expressly sanctioned in the New Testament.

The most learned English divines before the period of the Restoration, such as Cranmer, Jewel, Hooker, Field, Ussher, Hall, and Stillingfleet, did not hold the theory of an exclusive \emph{jure divino} episcopacy, and fully recognized the validity of presbyterian ordination. They preferred and defended episcopacy as the most ancient and general form of government, best adapted for the maintenance of order and unity; in one word, as necessary for the well-being, but not for the being of the Church. Cranmer invited the co-operation of Lutherans and Calvinists even in the most important work of framing the Articles of Religion and revising the Liturgy, without questioning their ordination; his own views of episcopacy were so low that he declared 'election or appointment thereto sufficient' without consecration, and he was so thoroughly Erastian that after the death of Henry he and his suffragans took out fresh commissions from the new king.\footnote{In accordance with an act of the thirty-seventh year of Henry VIII., which declares that 'Archbishops and the other ecclesiastical persons had no manner of jurisdiction ecclesiastical but by, under, and from his Royal Majesty; and that his Royal Majesty was the only supreme head of the Church of England and Ireland, to whom, by holy Scripture, all authority and power was wholly given,' etc.} His three successors in the primacy (Parker, Grindal, and Whitgift) did not differ from him in principle. 'Archbishop Grindal,' says Macaulay, 'long hesitated about accepting a mitre, from dislike of what he regarded as the mummery of consecration. Bishop Parkhurst uttered a fervent prayer that the Church of England would propose to herself the Church of Zurich as the absolute pattern of a Christian community. Bishop Ponet was of opinion that the word \emph{bishop }should be abandoned to the Papists, and that the chief officers of the purified Church should be called \emph{superintendents.}' The nineteenth of the Elizabethan Articles, which treats of the visible Church, says nothing of episcopacy as a mark of the Church. The statute of the thirteenth year of Elizabeth, cap. 12, permits ministers of the Scotch and other foreign Churches to exercise their ministry in England without re-ordination. After the union with Scotland the English sovereign represented in his official character the national Churches of the two countries, and when in Scotland, Queen Victoria takes the communion from the hands of a Presbyterian parson. Prominent clergymen of the Church of England, such as Travers (Provost of Trinity College, Dublin), Whittingham (Dean of Durham), Cartwright (Professor of Divinity in Cambridge, afterwards Master of Warwick Hospital), and John Morrison (from Scotland), had received only Presbyterian ordination in foreign Churches. Similar instances of Scotch, French, and Dutch Reformed ministers who were received simply on subscribing the Articles occurred down to the civil war. The English delegates to the Synod of Dort, which was presided over by a presbyter, were high dignitaries and doctors of divinity, one of them (Carleton) a bishop, and two others (Davenant and Hall) were afterwards raised to bishoprics. Archbishop Ussher, the greatest English divine of his age, who in eighteen years had mastered the whole mass of patristic literature, defended episcopacy only as a presidency of one presbyter over his peers, and declared that when abroad he would take the holy communion from a Dutch Reformed or French minister as readily as from an Episcopalian clergyman at home.

But the reigns of James and Charles I. form the transition. In the heat of the Puritan controversy both parties took extreme ground, Presbyterians and Independents as well as Episcopalians, and claimed exclusive Scripture authority and divine right for their form of government. Truth and error were mixed on both sides; for the primitive government was neither Episcopalian nor Presbyterian nor Independent, but apostolic; and the Apostles, as inspired and infallible teachers and rulers of the whole Church of all ages, have and can have no successors, as Christ himself can have none.

The doctrine of the divine and exclusive right of episcopacy was first intimated, in self-defense, by Bishop Bancroft, of London (in a sermon, 1589), then taught and rigidly enforced by Archbishop Laud (1633–1645), the most un-Protestant of English prelates,\footnote{Laud made such a near approach to Rome that he was offered a cardinal's hat (Aug. 1633). When he first maintained, in his exercise for Bachelor of Divinity, in 1604, the doctrine that there could be no true Church without a bishop, he was reproved by the authorities at Oxford, because he 'cast a bone of contention between the Church of England and the Reformed on the Continent.' But when he was in power he spared no effort to force his theory upon reluctant Puritans in England and Presbyterians in Scotland.} and was apparently sanctioned in 1662 by the Act of Uniformity, which forbade any person to hold a benefice or to administer the sacraments before he be ordained a priest by Episcopal ordination. By this cruel Act two thousand ministers, including some of the ablest and most worthy men in England, were expelled from office and driven into non-conformity.

Notwithstanding this change, the Church of England has never officially and expressly pronounced on the validity or invalidity of non-episcopal orders in other Churches; she only maintains that no one shall officiate in her pulpits and at her altars who has not received episcopal ordination according to the direction of the Prayer-book.\footnote{The facts above stated are acknowledged by the best authorities of the Church of England of all parties, such as Strype, Burnet, Lathbury, Keble, and by secular historians such as Hallam and Macaulay. See a calm and thorough argument of Prof. G. P. Fisher, \emph{The Relation of the Church of England to the other Protestant Churches,} in the 'New-Englander' for January, 1874, pp. 121–172. This article grew out of a newspaper controversy in the \emph{New York Tribune,} occasioned by the secession of Bishop Cummins after the General Conference of the Evangelical Alliance at New York, October, 1873. This inter-denominational Conference had the express sanction of the Archbishop of Canterbury in a letter addressed to the Dean of Canterbury, one of the prominent delegates. See \emph{Proceedings} (published N. Y., 1874), p. 720. Comp. also Dr. Washburn, \emph{Relation of the Episcopal Church to other Christian Bodies,} N. Y., 1874.}

RICHARD HOOKER.

The truest representative of the conservative and comprehensive genius of Anglicanism in doctrine and polity, towards the close of the Elizabethan period, is the 'judicious Hooker' (1553–1600), who to this day retains the respect of all parties. In his great work on the 'Laws of Ecclesiastical Polity' he went to the root of the rising controversy between Episcopacy and Puritanism, by representing the Church as a legislative body which had the power to make and unmake institutions and rites not affecting the doctrines of salvation laid down in the Scriptures and œcumenical creeds.

He defends episcopacy, but without invalidating other forms of government, or unchurching other Churches. He highly commends Calvin's 'Institutes' and 'Commentaries,' and calls him 'incomparably the wisest man that ever the French Church did enjoy.'\footnote{He also says: 'Of what account the Master of Sentences [Peter Lombard] was in the Church of Rome, the same and more amongst the preachers of Reformed Churches Calvin had purchased; so that the perfectest divines were judged they which were skillfulest in Calvin's writings; his books almost the very canon to judge both doctrine and discipline by.' See Hooker's lengthy account of Calvin's life and labors in the Preface to his work on the \emph{Laws of Ecclesiastical Polity,} Vol. I. pp. 158–174, edition of Dr. John Keble.} He generally agrees with his theology, at least as far as it is Augustinian, and he clearly adopts his view of the eucharist—namely, as he expresses it, that 'Christ is, \emph{personally} present, albeit a part of Christ be \emph{corporally} absent,' and 'that the real presence is not to be sought for in the sacrament (i.e., in the elements), but in the worthy receiver of the sacrament.' But he keeps clear of the logical sharpness and rigor of Calvinism, and subjects it to the higher test of the fathers and the early Church.\footnote{Dr. Keble, who was a High Anglican or Anglo-Catholic of the Oxford school, says in the Preface to his edition (p. xcix.): 'With regard to the points usually called Calvinistic, Hooker undoubtedly favored the tone and language, which has since come to be characteristic of that school, commonly adopted by those theologians to whom his education led him as guides and models on occasions where no part of Calvinism comes expressly into debate. It is possible that this may cause him to appear, to less profound readers, a more decided partisan of Calvin than he really was. At least it is certain that on the following subjects he was himself decidedly in favor of very considerable modifications of the Genevan theology.' Keble then contrasts the strict Calvinism of the Lambeth Articles with the cautious predestinarianism of Hooker as expressed in a fragment which teaches eternal election and the final perseverance of the foreknown elect, without mentioning reprobation, and makes condemnation depend on 'the foresight of sin as the cause.' Judas went to his place, which was 'of his own proper procurement. Devils were not ordained of God for hell-fire, but hell-fire for them; and for men so far as it was foreseen that men would be like them.' There are, however, as Keble himself admits, passages in Hooker which are more strongly Calvinistic, especially on the doctrine of the perseverance of saints, which he considers hardly consistent with his doctrine of universal baptismal grace. But both these doctrines were held by Augustine likewise, from whom Hooker borrowed them.}

His respect for antiquity and his churchly conservatism gained ground after his death in the conflict with Puritanism; and when the Synod of Dort narrowed the Calvinism of the Reformation to a five-angular scholastic scheme, Arminian doctrines, in connection with High-Church principles, spread rapidly in the Church of England. She became, as a body, more and more exclusive, and broke off the theological interchange and fraternal fellowship with non-episcopal Churches. But we hope the time is coming when the Christian communion which characterized her formative period will be revived under a higher and more permanent form.

Note.—My friend, the Rev. Dr. E. A. Washburn, of New York, an Episcopalian divine of rare culture and liberality of spirit, has kindly furnished the following contribution to this chapter, which will give the reader a broad inside view of Anglicanism under the various phases of its historic development:

'The doctrinal system of the English Church, in its relation to other Reformed communions, especially needs a historic treatment; and the want of this has led to grave mistakes, alike by Protestant critics and Anglo-Catholic defenders. It was one in its positive principles, as opposed to the dogmatic falsehoods of Rome, with the great bodies of the Continental Reformation; yet it grew as a national Church, while those were more fully shaped by the theology of their leaders—Luther, Calvin, and Zwingli. This fact is the key of its history. England felt the same influences, religious and social, that awakened Europe, but its ideas were not borrowed from abroad; it only completed the growth begun in the day of Wyclif. Its earliest step was thus a national, one. Nor was this, as has been proved by its latest historians from the record, the act of Henry VIII.; for before his quarrel the Parliament annulled forever, by its own decree, the supremacy of Rome. It could not be expected that during his reign the standard of doctrine should be greatly changed; and it should be remembered, that Luther himself renounced only by degrees the idea of Papal authority. The ``Articles devised to establish Christian Quietness,'' probably the original of the later Cotton MSS., and the ``Institution of a Christian Man'' following it in 1537, show that the dogma of the mass, the seven sacraments, intercessory prayers for the dead, and reverence of the Virgin and saints as mediators, remained. It is worth noting, however, that the ``Erudition'' in 1543 gives signs of change, as the ``corporal'' presence is there only the ``very body,'' and the idea of special intercession is modified to prayer ``for the universal congregation of Christian people, quick and dead.'' But the next reign proves that the act of national freedom held in solution the whole result. Ultramontanism meant then, as now, not only the feudal headship of Rome, but its scholastic and priestly system. The Reformation, ripened in the minds of Cranmer, Latimer, Ridley, and other devout thinkers, bore its fruit in the revised Liturgy and Articles; nor can any thing be clearer than the doctrinal standard of the Church, if we trace it with just historic criticism to the time when these were fixed.

'The Articles ask our first study. It is plain that the foundation-truths of the Reformation—justification by faith, the supremacy and sufficiency of written Scripture, the fallibility of even general councils—are its basis. Yet it is just as plain that in regard of the specific points of theology, which were the root of discord in the Continental Churches, as election, predestination, reprobation, perseverance, and the rest, these Articles speak in a much more moderate tone. It is from a narrow study of that age that they have been called articles of compromise between a Calvinistic and Arminian party. There were some of extreme views, as the Lambeth Articles prove, but they did not represent the body. The English Reformers had been bred, like the great Genevan, in the school of the greater Augustine; and his richer, more ethical spirit appears in not only the Articles, but in the writings of well-nigh all from Hooper or Whitgift to Hooker. There was the friendliest intercourse between them and the divines of the Continent. Melanchthon, Calvin, Bucer were consulted in their common work. But the unity of the national Church, not the system of a school, was uppermost; and we may write the character of them all in the words of the biographer of Field, that ``in points of extreme difficulty he did not think fit to be so positive in defining as to turn matters of opinion into matters of faith.''

'We may thus learn the structure of the liturgical system. The English Reformers aimed not to create a new, but to reform the historic Church; and therefore they kept the ritual with the episcopate, because they were institutions rooted in the soil. They did not unchurch the bodies of the Continent, which grew under quite other conditions. No theory of an exclusive Anglicanism, as based on the episcopate and general councils, was held by them. Such a view is wholly contradictory to their own Articles. But the historic character of the Church gave it a positive relation to the past; and they sought to adhere to primitive usage as the basis of historic unity. In this revision, therefore, they weeded out all Romish errors, the mass, the five added sacraments, the legends of saints, and superstitious rites; but they kept the ancient Apostles' Creed and the Nicene in the forefront of the service, the sacramental offices, the festivals and fasts relating to Christ or Apostles with whatever they thought pure. Such a work could not be perfect, and it is false either to think it so or to judge it save by its time. There are archaic forms in these offices which retain some ideas of a scholastic theology. The view of regeneration in the baptismal service, decried to-day as Romish, can be found by any scholar in Melanchthon or in Bullinger's Decades. We may see in some of the phrases of the communion office the idea of more than a purely spiritual participation, yet the view is almost identical with that of Calvin. The dogma of the mass had been renounced, but the Aristotelian notions of spirit and body were still embodied in the philosophy of the time. The absolution in the office for the sick, and like features, have been magnified into ``Romanizing germs'' on one side and Catholic verities on another. The Athanasian Creed, revered by all the Reformers, was retained, yet not as that of Nice in the body of the worship; and it was wisely excluded by the American revisers, as the English Church will by-and-by displace it, because a better criticism shows it to be the metaphysical deposit of a later time, un-catholic in descent or structure. Such is the rule by which we are to know the unity of the English system. The satire, so often repeated since Chatham, that the Church has a ``Popish Liturgy and Calvinistic Articles,'' is as ignorant as it is unjust. All liturgical formularies need revision; but such a task must be judged by the standard of the Articles, the whole tenor of the Prayer-book, and the known principles of the men. In the same way we learn their view of the Episcopate. Not one leading divine from Hooper to Hooker claimed any ground beyond the fact of primitive and historic usage; and Whitgift, the typical High-Churchman of the Elizabethan time, in reply to the charge of Cartwright against prelacy as unscriptural, took the ground that to hold it ``of necessity to have the same kind of government as in the Apostles' time, and expressed in Scripture,'' is a ``rotten pillar.'' The Puritan of that day was as narrow as the narrow Churchman of our own.

'This historic sketch of the English Reformation explains its whole character. It had in it varied elements, but by no means contradictory. Had not other influences dwarfed its design, it would have done much to harmonize the communions of Protestantism, to blend the new life with a sober reverence for the historic past. Lutheranism and Calvinism did each its part in the development of a profound theology. The English Church had a more comprehensive doctrine and a more conservative order. It placed the simple Apostles' Creed above all theological confessions as its basis, and a practical system above the subtleties of controversy. But its defect lay in the policy which sought uniformity instead of a large unity; and the loss of the conscientious men who left the national Church gave its ecclesiastical element an undue growth. Yet it has retained throughout much of its comprehensiveness. It has had many schools of thought, but none has ruled it. Calvinism, although shorn of its early strength, has had always adherents, from the saintly Leighton to Toplady and Venn. The Arminian doctrine entered early from Holland, and in the visit of the divines sent by James to the Synod of Dort, among whom were Hall and Davenant, we have the early traces of the change. Davenant was nominally against the Remonstrants, but the ``Suffrages'' prove already the milder tone of the English theology. It is with Laud that the system gained strong ground, yet it never led to such quarrels as in the land of Grotius; it represented the growing dislike of a harsh supralapsarianism and the mild spirit of scholars like Jeremy Taylor. The criticism has often been made that Arminianism is more akin to a High-Church system, because it teaches that divine grace is conditioned by works; but if so, perhaps it shows, as in the case of Jansenism, that a metaphysical creed, in losing sight of the moral side of its own truth, will always drive men to its opposite. The English theology of the next period has the like variety. It had its divines of rich learning—Bramhall, Cosin, and others—inclined to a stricter view of the sacraments and ministry than the Reformers; yet it is mere exaggeration to call them the Anglo-Catholic fathers, as if they were the exponents of the whole Church. They belong to one school of their time. Nor is it a less mistake to judge from their opposition, as members of the national Church, to the Dissenters, that they unchurched the Continental Protestants. Bramhall held an episcopate to be of the \emph{Ecclesia integra,} not \emph{vera}; and Morton, while bitter towards the Presbyterians, is ``not so uncharitable'' towards foreign Reformed bodies ``as to censure them for no Churches, for that which is their infelicity, not their fault.'' Chillingworth and Hales are leaders in this period of a more liberal thought. The Cambridge school, which a modern critic calls the herald of broad Churchmanship, begins here with Smith and Whichcote. The theology of England passed into a still more comprehensive growth. Its larger conflict with Deism took it out of the guerrilla war of the past into the field of Biblical criticism, Christian evidence, and history. No party wholly represents it. Such different minds as Tillotson and Waterland, Cudworth and Paley, Arnold and Keble have been of the same communion. Its successive movements have stirred, yet not rent it. The Methodist revival came from the Arminian Wesley, and the wave of spiritual life left its true influence, although a cold establishment policy ignored it. The evangelical movement was Calvinistic, yet it was mainly the protest of devout men like Wilberforce against formalism, and did little for theological growth. Our time has been busy with the Oxford divinity, which has sought to build a theory of Anglo-Catholicism on the basis of an exclusive episcopal succession, a Nicene authority concurrent with Scripture, and a priesthood dispensing grace through the sacraments. It will end as the theory of a passing school. Our sketch will show on what grounds we judge it a contradiction to the standards of the body, the \emph{consensus} of its fathers down to Hooker, and an utter misstatement of the historic position of the Church of England. It may be hoped that the long strife will lead to a better understanding of its relation to other Reformed communions, and to its place in the common work for the unity of Christendom.'

\xsubsection{The Doctrinal Formulas of Henry VIII.}

§ 78. The Doctrinal Formulas of Henry VIII.

THE TEN ARTICLES.

The first doctrinal deliverance of the Church of England after the rupture with Rome is contained in the Ten Articles of 1536, devised by Henry VIII. (who styles himself in the preface 'by the grace of God king of England and of France, defender of the faith, lord of Ireland, and in earth supreme head of the Church of England'), and approved by convocation.\footnote{First printed by Thomas Berthelet, under the title 'Articles | devised by the Kinges Highnes Majestie, | to stablyshe Christen quietnes and unitie | amonge us, | and | to avoyde contentious opinions, | which articles be also approved | by the consent and determination of the hole | clergie of this realme. | Anno M.D.XXXVI.' They are given by Fuller, Burnet, (Addenda), Collier, and Hardwick (Appendix I). In the Cotton MS. the title is, 'Articles about Religion, set out by the \emph{Convocation,} and published by the King's authority.' It is impossible to determine how far the Articles are the product of the king (who in his own conceit was fully equal to any task in theology as well as Church government), and how far the product of his bishops and other clergy. See Hardwick, pp. 40 sqq.} They are essentially Romish, with the Pope left out in the cold. They can not even be called a compromise between the advocates of the 'old learning,' headed by Gardiner (Bishop of Winchester from 1531), and of the 'new learning,' headed by Cranmer (Archbishop of Canterbury from March, 1533). Their chief object, according to the preface, was to secure by royal authority unity and concord in religious opinions, and to 'repress' and 'utterly extinguish' all dissent and discord touching the same. They were, in the language of Foxe, intended for 'weaklings newly weaned from their mother's milk of Rome.' They assert (1) the binding authority of the Bible, the three œcumenical creeds, and the first four œcumenical councils; (2) the necessity of baptism for salvation, even in the case of infants;\footnote{Art. II. says that 'infants ought to be baptized;' that, dying in infancy, they 'shall undoubtedly be saved thereby, and \emph{else not};' that the opinions of Anabaptists and Pelagians are 'detestable heresies, and utterly to be condemned.'} (3) the sacrament of penance, with confession and absolution, which are declared 'expedient and necessary;' (4) the substantial, real, corporal presence of Christ's body and blood under the form of bread and wine in the eucharist; (5) justification by faith, joined with charity and obedience; (6) the use of images in churches; (7) the honoring of saints and the Virgin Mary; (8) the invocation of saints; (9) the observance of various rites and ceremonies as good and laudable, such as clerical vestments, sprinkling of holy water, bearing of candles on Candlemas-day, giving of ashes on Ash-Wednesday; (10) the doctrine of purgatory, and prayers for the dead in purgatory.

THE BISHOPS' BOOK AND THE KING'S BOOK.\footnote{Printed in \emph{Formularies of Faith put forth by Authority during the Reign of Henry VIII.} Oxford, 1825.}

These Articles were virtually, though not legally, superseded by the 'Bishops' Book,' or the 'Institution of a Christian Man,' drawn up by a Committee of Prelates, 1537, but never sanctioned by the king. It contains an Exposition of the Creed, the Seven Sacraments, the Ten Commandments, the Lord's Prayer, the Ave Maria, and a discussion of the disputed doctrines of justification and purgatory, and the human origin of the papacy. It marks a little progress, which must be traced to the influence of Cranmer and Ridley, but it was superseded by a reactionary revision called the 'King's Book,' or the 'Necessary Doctrine and Erudition for any Christian Man,' sanctioned by Convocation, and set forth by royal mandate in 1543, when Gardiner and the Romish party were in the ascendant.

THE THIRTEEN ARTICLES.

During the negotiations with the Lutheran divines (1535–1538), held partly at Wittenberg, partly at Lambeth, an agreement consisting of Thirteen Articles was drawn up in Latin, at London, in the summer of 1538, which did not receive the sanction of the king, but was made use of in the following reign as a basis of several of the Forty-two Articles. They have been recently discovered in their collected form, by Dr. Jenkyns, among the manuscripts of Archbishop Cranmer in the State Paper Office.\footnote{They are printed in Jenkyns's \emph{Remains of Cranmer} (1833), Vol. IV. pp. 273 sqq.; in Cox's (Parker Soc.) edition of \emph{Cranmer's Works} (1846), Vol. II. pp. 472–480; and in Hardwick's \emph{History of the Articles,} Append. II. pp. 261–273. Six of these thirteen Articles were previously published by Strype and Burnet, but with a false date (1540) and considerable variations.} They treat of the Divine Unity and Trinity, Original Sin, the Two Natures of Christ, Justification, the Church, Baptism, the Eucharist, Penitence, the Use of the Sacraments, the Ministers of the Church, Ecclesiastical Rites, Civil Affairs, the Resurrection and Final Judgment. They are based upon the Augsburg Confession, some passages being almost literally copied from the same.\footnote{See the comparison in Hardwick, pp. 62 sqq.}

THE SIX ARTICLES.

The Thirteen Articles remained a dead letter in the reign of Henry. He broke off all connection with the Lutherans, and issued in 1539, under the influence of Gardiner and the Romish party, and in spite of the protest of Cranmer, the monstrous statute of the Six Articles, ' for the abolishing of Diversity of Opinions.' They are justly called the 'bloody' Articles, and a 'whip with six strings.' They bore severely not only upon the views of the Anabaptists and all radical Protestants, who in derision were called 'Gospellers,' but also upon the previous negotiations with the Lutherans. After the burning of some Dissenters the Articles were somewhat checked in their operation, but remained legally in force till the death of the king, who grew more and more despotic, and prohibited (in 1542) Tyndale's 'false translation' of the Bible, and even the reading of the New Testament in English to all women, artificers, laborers, and husbandmen.

The Six Articles imposed upon all Englishmen a belief (1) in transubstantiation, (2) the needlessness of communion in both kinds, (3) in clerical celibacy, (4) the obligation of vows of chastity or widowhood, (5) the necessity of private masses, (6) auricular confession. Here we have some of the most obnoxious features of Romanism. Whoever denied transubstantiation was to be burned at the stake; dissent from any of the other Articles was to be punished by imprisonment, confiscation of goods, or death, according to the degree of guilt.

\xsubsection{The Edwardine Articles. A.D. 1553.}

§ 79. The Edwardine Articles. A.D. 1553.

With the accession of Edward VI. (Jan. 28, 1547) Cranmer and the reform party gained the controlling influence. The Six Articles were abolished. The First Prayer-Book of Edward VI. was prepared and set forth (1549), and a few years afterwards the Second, with sundry changes (1552).

The reformation of worship was followed by that of doctrine. For some time Cranmer entertained the noble but premature idea of framing, with the aid of the German and Swiss Reformers, an evangelical catholic creed, which should embrace 'all the heads of ecclesiastical doctrine,' especially an adjustment of the controversy on the eucharist, and serve as a protest to the Council of Trent, and as a bond of union among the Protestant Churches.\footnote{See Cranmer's letters of invitation to Calvin, Bullinger, and Melanchthon, in Cox's edition of \emph{Cranmer's Works,} Vol. II. pp. 431–433.}

This project was reluctantly abandoned in favor of a purely English formula of public doctrine, the Forty-two Articles of Religion. They were begun by Cranmer in 1549, subjected to several revisions, completed in November, 1552, and published in 1553, together with a short Catechism, by 'royal authority,' and with the approval of 'a Synod (Convocation) at London.'\footnote{'\emph{Articuli de quibus in Synodo Londinensi, A.D. M.D.LII. ad tollendam opinionum dissensionem et consensum veræ religionis firmandum, inter Episcopos et alios Eruditos Viros convenerat.}' 'Articles agreed on by the Bishopes, and other learned menne in the Synode at London, in the yere of our Lorde Godde, M.D.LII., for the auoiding of controuersie in opinions, and the establishment of a godlie concorde, in certeine matters of Religion.' They are printed in Hardwick, Append. III. pp. 277–333, in Latin and English, and in parallel columns with the Elizabethan Articles. The Latin text is also given by Niemeyer, pp. 592–600. On minor points concerning their origin, comp. Hardwick, pp. 73 sqq.} It is, however, a matter of dispute whether they received the formal sanction of Convocation, or were circulated on the sole authority of the royal council during the brief reign of Edward (who died July 6, 1553).\footnote{Palmer, Burnet, and others maintain the latter; Hardwick (p. 107), the former.} The chief title to the authorship of the Articles, as well as of the revised Liturgy, belongs to Cranmer; it is impossible to determine how much is due to his fellow-Reformers—'bishops and other learned men'—and the foreign divines then residing in England, to whom the drafts were submitted, or whose advice was solicited.\footnote{John Knox and the other royal chaplains were also consulted; see Lorimer, 1.c. pp. 126 sqq. Knox did not object to the doctrines of the Articles, but to the rubric on kneeling in the eucharistic service of the Liturgy, and his opposition led to the 'Declaration on Kneeling,' which is a strong protest against ubiquitarianism and any idolatrous veneration of the sacramental elements. It was inserted as a rubric by order of Council in 1552, was omitted to 1559, and restored in 1662.}

The Edwardine Articles are essentially the same as the Thirty-nine, with the exception of a few (three of them borrowed from the Augsburg Confession), which were omitted in the Elizabethan revision—namely, one on the blasphemy against the Holy Ghost (Art. XVI.); one on the obligation of keeping the moral commandments—against antinomianism—(XIX.); one on the resurrection of the dead (XXXIX.); one on the state of the soul after death—against the Anabaptist notion of the psychopannychia—(XL.); one against the millenarians (XLI.);\footnote{'\emph{Qui Millenariorum fabulam revocare conantur, sacris literis adversantur, et in Judaica deliramenta sese præcipitant} (cast themselves headlong into a Juishe dotage).' Comp. the Augsburg Confession, Art. XVII., where the Anabaptists and others are condemned for teaching the final salvation of condemned men and devils, and the Jewish opinions of the millennium.} and one against the doctrine of universal salvation (XLII.).\footnote{\par '\emph{Hi quoque damnatione digni sunt, qui conantur hodie perniciosam opinionem instaurare, quod omnes, quantumvis impii, servandi sunt tandem, cum definito tempore a justitia divina pænas de admissis flagitiis luerunt.}'} A clause in the article on Christ's descent into Hades (Art. III.),\footnote{'\emph{Nam corpus} [\emph{Christi}] \emph{usque ad resurrectionem in sepulchro jacuit, Spiritus ab illo emissus} (his ghost departing from him) \emph{cum spiritibus qui in carcere sive in inferno detinebantur, fuit, illisque prædicavit, quemadmodum testatur Petri locus.} (\emph{At suo ad inferos descensu nullos a carceribus aut tormentis liberavit Christus Dominus.})'} and a strong protest against the ubiquity of Christ's body, and 'the real and bodily presence of Christ's flesh and blood in the sacrament of the Lord's Supper' (in Art. XXIX.), were likewise omitted.

\xsubsection{The Elizabethan Articles. A.D. 1563 and 1571.}

§ 80. The Elizabethan Articles. A.D. 1563 and 1571.

After the temporary suppression of Protestantism under Queen Mary, the Reformed hierarchy, Liturgy, and Articles of Religion were permanently restored, with a number of changes, by Queen Elizabeth.

In 1559, Archbishop Parker, with the other prelates, set forth, as a provisional test of orthodoxy, Eleven Articles, taken in part from those of 1553, but differing in form and avoiding controverted topics.\footnote{They are printed by Hardwick in Append. IV. pp. 337–339.} They were superseded by the Thirty-nine Articles.

THE LATIN EDITION, 1563.

At the first meeting of the two Convocations, which were summoned by Elizabeth in January, 1563, Parker submitted a revision of the Latin Articles of 1553, prepared by him with the aid of Bishop Cox of Ely, Bishop Guest of Rochester, and others, who had already taken an active part in the revision of the Prayer-book.\footnote{A manuscript copy of this revision, with numerous corrections and autograph signatures of Matthæus Cantuar.' (Parker), and other prelates (including some of the northern province), is preserved among the Parker MSS. in Corpus Christi College, Cambridge, and was published by Dr. Lamb in 1829. The handwriting (as Mr. Lewis, the librarian, informed me when there on a visit in July, 1875) is probably Jocelin's, the secretary of Parker. The copy contains also the older Articles Nos. 40–42, but marked by a red line as to be omitted. This copy is probably the same which Parker submitted to Convocation, but it presents several variations (especially in Art. XX.) from the copy of the Convocation records. Comp. Hardwick, pp. 125 and 135 sqq.} After an examination by both houses, the Articles, reduced to the number of thirty-nine, were ratified and signed by the Bishops and the members of the lower house, and published by the royal press, 1563.

It is stated that Elizabeth 'diligently read and sifted' the document before giving her assent. To her influence must probably be traced two characteristic changes of the printed copy as compared with the Parker MS.—namely, the insertion of the famous clause in Art. XX., affirming the authority of the Church in matters of faith—and the omission of Art. XXIX., which denies that the unworthy communicants partake of the body and blood of Christ.\footnote{Hardwick, pp. 143 sqq.} The latter Article, however, was restored by the Bishops, May 11, 1571, and appears in all the printed copies since that time, both English and Latin.

THE ENGLISH EDITION, 1571.

The authorized English text was adopted by Convocation in 1571, and issued under the editorial care of Bishop Jewel of Salesbury. It presents sundry variations from the Latin edition of 1563. Both editions are considered equally authoritative and mutually explanatory.\footnote{This is the view of Burnet and Waterland, adopted by Hardwick, p. 158. Waterland says \{\emph{Works,} Vol. II. pp. 316, 317): 'As to the Articles, English and Latin, I may just observe for the sake of such readers as are less acquainted with these things: \emph{first,} that the Articles were passed, recorded, and ratified in the year 1562 [1563], and \emph{in Latin only}. \emph{Secondly,} that those Latin Articles were revised and corrected by the convocation of 1571. \emph{Thirdly,} that an authentic English translation was then made of the Latin Articles by the same convocation, and the Latin and English adjusted as nearly as possible. \emph{Fourthly,} that the Articles thus perfected \emph{in both languages} were published the same year, and by the royal authority. \emph{Fifthly,} subscription was required the same year to the English Articles, called the Articles of 1562, by the famous act of the 13th of Elizabeth.—These things considered, I might justly say with Bishop Burnet, that the Latin and English are both \emph{equally authentical.} Thus much, however, I may certainly infer, that if in any places the English version be ambiguous, where the Latin original is clear and determinate, the Latin ought to fix the more doubtful sense of the other (as also \emph{vice versa}), it being evident that the Convocation, Queen, and Parliament intended the same sense in both.'}

THE ROYAL DECLARATION OF 1628.

After the Synod of Dort, to which James I. sent a strong delegation, the Arminian controversy spread in England, and caused such an agitation that the king, who, according to his own estimate and that of his flatterers, was equal to Solomon in wisdom, ordered Archbishop Abbot (Aug. 4, 1622) to prohibit the lower clergy from preaching on the five points.\footnote{One of the directions reads: 'That no preacher of what title soever, under the degree of a Bishop, or Dean at least, do from henceforth presume to preach in any popular auditory the deep points of predestination, election, reprobation, or the universality, efficacy, resistibility or irresistibility of divine grace; but leave those themes to be handled by learned men, and that moderately and modestly, by way of use and application, rather than by way of positive doctrine, as being fitter for the schools and Universities than for simple auditories.'—Wilkins, Vol. IV. p. 465; Hardwick, p. 202.} Charles I., in concert with Archbishop Laud (who sympathized with Arminianism), issued a Proclamation (1626) of similar import, deploring the prevalence of theological dissension, and threatening to visit with severe penalties those clergymen who should raise, publish, or maintain opinions not clearly warranted by the formularies of the Church.

As this proclamation did not silence the controversy, Charles was advised by Laud to order the republication of the thirty-nine Articles with a Preface regulating the interpretation of the same. This Preface, called 'His Majesty's Declaration,' was issued in 1628, and has ever since accompanied the English editions of the Articles.\footnote{It disappeared, of course, in the American editions. It is printed in Vol. III. p. 486.} Its object was to check Calvinism (although it is not named), and the quinquarticular controversy ('all further curious search' on 'those curious points in which the present differences lie'), and to restrict theological opinions to the 'literal and grammatical sense' of the Articles.\footnote{'No man shall either print or preach or draw the Article' [the previous sentence speaks of the \emph{Articles} generally, perhaps Art. XVII. on predestination is meant particularly] 'aside any way, but shall submit to it in the plain and full meaning thereof; and shall not put his own sense or comment to be the meaning of the Article, but shall take it in the literal and grammatical sense.' In a 'Declaration' of Charles on the dissolution of Parliament (March 10, 1628), he says, concerning his intention in issuing the Declaration before the Articles: 'We did tie and restrain all opinions to the sense of these Articles that nothing might be left to fancies and invocations' [probably an error for 'innovations']. 'For we call God to record, before whom we stand, that it is, and always hath been, our chief heart's desire, to be found worthy of that title, which we account the most glorious in all our crown, \emph{Defender of the Faith.}'—Hardwick, p. 206.} It was greeted by Arminians and High-Churchmen, who praise its moderation,\footnote{Hardwick says (p. 205): 'A document more sober and conciliatory could not well have been devised.' Bishop Forbes goes further, and thinks that it was 'the enunciation of the Catholic sense of the Articles,' and that Newman's \emph{Tract XC.} and Pusey's \emph{Irenicon} are 'legitimate outcomes of the King's Declaration' (1.c. Vol. I. p. xi.).} but was resisted by Calvinists and the Puritan party then prevailing in the House of Commons, which declared its determination to suppress both 'Popery and Arminianism.'\footnote{The House passed the following vote and manifesto on the royal Declaration: 'We, the Commons in Parliament assembled, do claim, protest, and avow for truth, the sense of the Articles of Religion which were established by Parliament in the thirteenth year of our late Queen Elizabeth, which by the public act of the Church of England, and by the general and current expositions of the writers of our Church, have been delivered unto us. And we reject the sense of the Jesuits and Arminians, and all others, wherein they differ from us.'—Hardwick, p. 206.} The subsequent history of England has shown how little royal and parliamentary proclamations and prohibitions avail against the irresistible force of ideas and the progress of theology.

SUBSCRIPTION.

Queen Elizabeth was at first opposed to any action of Parliament on questions of religious doctrine, which she regarded as the highest department of her own royal supremacy; but in May, 1571, she was forced by her council, in view of popish agitations, to give her assent to a bill of Parliament which required all priests and teachers of religion to subscribe the Thirty-nine Articles.\footnote{Stat. 13 Eliz. c. 12. It enacts 'by the authority of the present Parliament, that every person under the degree of a bishop, which doth or shall pretend to be a priest or minister of God's holy Word and Sacraments, by reason of any other form of institution, consecration, or ordering, than the form set forth by Parliament in the time of the late King of most worthy memory, King Edward the Sixth, or now used, . . . shall . . . declare his assent, and subscribe to all the Articles of Religion, which only concern the confession of the true Christian faith and the doctrine of the sacraments comprised in a book entitled \emph{Articles,} . . . put forth by the Queen's authority.' The subscription to the Articles was urged by the Puritanic party in Parliament in opposition to Romanism. See Hardwick, pp. 150 sq. The wording of the statute was made use of to confine assent to the \emph{doctrinal} Articles ('which \emph{only} concern,' etc.), and to relieve the conscience of the Puritans who objected to the royal supremacy, the surplice, and other 'defiled robes of Antichrist.'}

Subscription was first rigidly enforced by Archbishop Whitgift (in 1584, which is noted as 'the woful year of subscription'), and by Bancroft (1604).

This test of orthodoxy was even applied to academical students. At Oxford a decree of Convocation, in 1573, required students to subscribe before taking their degrees, and in 1576 this requirement was extended to students above sixteen years of age on their admission. At Cambridge the law was less rigid.

The Act of Uniformity under Charles II. imposed with more stringency than ever subscription on the clergy and every head of a college. But the Toleration Act of William and Mary gave some relief by exempting dissenting ministers from subscribing to Arts. XXXIV—XXXVI. and a portion of XXVII. Subsequent attempts to relax or abolish subscription resulted at last in the University Tests Act of 1871, by which 'no one, at Oxford, Cambridge, or Durham, in order to take a degree, \emph{except in divinity,} or to exercise any right of graduates, can be required to make any profession of faith.'\footnote{The various acts enforcing and relaxing subscription are conveniently collected in the \emph{Prayer-Book Interleaved,} London, 7th ed. 1873, pp. 360 sqq. See also chap. xi. of Hardwick's \emph{History of the Articles.}}

RELATION TO THE EDWARDINE ARTICLES.

The Elizabethan Articles differ from the Edwardine Articles, besides minor verbal alterations—

(1.) In the omission of seven Articles (Edwardine X., XVI., XIX., XXXIX. to XLII.). The last four of them reject certain Anabaptist doctrines, which had in the mean time disappeared or lost their importance.\footnote{See p. 615.} Art. XIX. of the old series, touching the obligation of the moral law, was transferred in substance to Art. VII. of the new series.

(2.) In the addition of four Articles, viz.: On the Holy Ghost (Eliz. V.); on good works (XII.); on the participation of the wicked in the eucharist (XXIX.); on communion in both kinds (XXX.).

(3.) In the partial curtailment or amplification of seventeen Articles. Among the amplifications are to be noticed the list of Canonical and Apocryphal Books (VI.), and of the Homilies (XXXV.); the restriction of the number of sacraments to two (XXV.); the condemnation of transubstantiation, and the declaration of the spiritual nature of Christ's presence (XXVIII.); the disapproval of worship in a foreign tongue (XXIV.); the more complete approval of infant baptism (XXVII.), and clerical marriage (XXXII.).

The difference of the two series, and their relation to the Thirteen Articles, will be more readily seen from the following table:

Thirteen Articles.1538

Forty-two Articles.1553

Thirty-nine Articles.1571

1. De Unitate Dei et Trinitate Personarum.

1. Of faith in the holie Trinitie.

1. Of Faith in the Holy Trinity.

2. De Peccato Originali.

2. That the worde, or Sonne of God, was made a very man.

2. Of Christ the Son of God, which was made very man.

3. De duabus Christi Naturis.

4. De Justificatione.

3. Of the goying doune of Christe into Helle.

3. Of the Going down of Christ into Hell.

5. De Ecclesia.

6. De Baptismo.

4. The Resurrection of Christe.

4. Of the Resurrection of Christ.

7. De Eucharistia.

8. De Pœnitentia.

5. Of the Holy Ghost.

9. De Sacramentorum Usu.

5. The doctrine of holie Scripture is sufficient to Saluation.

6. Of the Sufficiency of the Holy Scripture for Salvation.

10. De Ministris Ecclesiæ.

11. De Ritibus Ecclesiasticis.

6. The olde Testamente is not to be refused.

7. Of the Old Testament.

7. The three Credes.

8. Of the Three Creeds.

12. De Rebus Civilibus.

8. Of originall or birthe sinne.

9. Of Original or Birth Sin.

13. De Corporum Resurrectione et Judicio Extremo.

9. Of free wille.

10. Of Free Will.

10. Of Grace.

11. Of the Justification of manne.

11. Of the Justification of man.

[This order follows, as far as it goes, the order of the doctrinal articles of the Augsburg Confession.]

12. Of Good Works.

12. Workes before Justification.

13. Of Works before Justification.

13. Workes of Supererogation.

14. Of Works of Supererogation.

14. No man is without sinne, but Christe alone.

15. Of Christ alone without sin.

15. Of sinne against the holie Ghoste.

16. Of Sin after Baptism.

16. Blasphemie against the holie Ghoste.

17. Of predestination and election.

17. Of Predestination and Election.

18. We must truste to obteine eternal salvation onely by the name of Christ.

18. Of obtaining Salvation by the name of Christ.

19. All men are bound to kepe the moral commaundementes of the Lawe.

20. Of the Church.

19. Of the Church.

21. Of the aucthoritie of the Churche.

20. Of the Authority of the Church.

22. Of the aucthoritie of General Counsailes.

21. Of the Authority of General Councils.

23. Of Purgatorie.

22. Of Purgatory.

24. No manne maie minister in the Congregation except he be called.

23. Of Ministering in the Congregation.

25. Menne must speake in the Congregation in soche toung as the people understandeth.

24. Of Speaking in the Congregation in such a tongue as the people understandeth.

26. Of the Sacramentes.

25. Of the Sacraments.

27. The wickednesse of the Ministres dooeth not take awaie the effectuall operation of Goddes ordinances.

26. Of the Unworthiness of Ministers which hinder not the effect of the Sacraments.

28. Of Baptisme.

27. Of Baptism.

29. Of the Lordes Supper.

28. Of the Lord's Supper.

29. Of the Wicked which eat not the Body of Christ in the use of the Lord's Supper.

30. Of Both Kinds.

30. Of the perfeicte oblacion of Christe made upon the crosse.

31. Of the one Oblation of Christ finished upon the cross.

31. The state of single life is commaunded to no man by the worde of God.

32. Of the Marriage of Priests,

32. Excommunicate persones are to bee auoided.

33. Of Excommunicate Persons, how they are to be avoided.

33. Tradicions of the Churche.

34. Of the Traditions of the Church.

34. Homelies.

35. Of Homilies.

35. Of the booke of Praiers and Ceremonies of the Churche of England.

36. Of Consecrating of Bishops and Ministers.

36. Of Ciuile Magistrates.

37. Of Civil Magistrates.

37. Christien mennes gooddes are not commune.

38. Of Christian men's goods, which are not common.

38. Christien menne maie take an oath.

39. Of a Christian man's oath.

39. The Resurrection of the dead is not yeat brought to passe.

The Ratification.

40. The soulles of them that departe this life doe neither die with the bodies nor sleep idlie.

41. Heretickes called Millenarii.

42. All men shall not bee saved at the length.

\xsubsection{The Interpretation of the Articles.}

§ 81. The Interpretation of the Articles.

The theological interpretation of the Articles by English writers has been mostly conducted in a controversial rather than an historical spirit, and accommodated to a particular school or party. Moderate High-Churchmen and Arminians, who dislike Calvinism, represent them as purely Lutheran;\footnote{\par So Archbishop Laurence, of Cashel, and Hardwick, in their learned works on the Articles.} Anglo-Catholics and Tractarians, who abhor both Lutheranism and Calvinism, endeavor to conform them as much as possible to the contemporary decrees of the Council of Trent;\footnote{Newman, Pusey, Forbes. Archbishop Laud had prepared the way for this Romanizing interpretation.} Calvinistic and evangelical Low-Churchmen find in them substantially their own creed.\footnote{Even the Puritans accepted the doctrinal Articles, and the Westminster Assembly first made them the basis of its Calvinistic Confession.} Continental historians, both Protestant and Catholic, rank the Church of England among the Reformed Churches as distinct from the Lutheran, and her Articles are found in every collection of Reformed Confessions.\footnote{From the \emph{Corpus et Syntagma} down to the collections of Niemeyer and Böckel. The Roman Catholic Möhler likewise numbers the Articles among the Reformed (Calvinistic) Confessions, \emph{Symbolik,} p. 22. On the other hand, the Articles have no place in any collection of Lutheran symbols; still less, of course, could they be included among Greek or Latin symbols.}

The Articles must be understood in their natural grammatical and historical sense, from the stand-point and genius of the Reformation, the public and private writings of their compilers and earliest expounders. In doubtful cases we may consult the Homilies, the Catechism, the several revisions of the Prayer-book, the Canons, and other contemporary documents bearing on the reformation of doctrine and discipline in the Church of England.

In a preceding section we have endeavored to give the historical key for the understanding of the doctrinal character of the English Articles. A closer examination will lead us to the following conclusions:

1. The Articles are \emph{Catholic} in the œcumenical doctrines of the Holy Trinity and the Incarnation, like all the Protestant Confessions of the Reformation period; and they state those doctrines partly in the very words of two \emph{Lutheran} documents, viz., the Augsburg Confession and the Würtemberg Confession.

2. They are \emph{Augustinian} in the anthropological and soteriological doctrines of free-will, sin, and grace: herein likewise agreeing with the Continental Reformers, especially the Lutheran.

3. They are \emph{Protestant} and \emph{evangelical} in rejecting the peculiar errors and abuses of Rome, and in teaching those doctrines of Scripture and tradition, justification by faith, faith and good works, the Church, and the number of sacraments, which Luther, Zwingli, and Calvin held in common.

4. They are \emph{Reformed} or \emph{moderately Calvinistic} in the two doctrines of Predestination and the Lord's Supper, in which the Lutheran and Reformed Churches differed; although the chief Reformed Confessions were framed after the Articles.

5. They are \emph{Erastian} in the political sections, teaching the closest union of Church and State, and the royal supremacy in matters ecclesiastical as well as civil; with the difference, however, that the Elizabethan revision dropped the title of the king as 'supreme head in earth,' and excluded the ministry of the Word and Sacraments from the 'chief government' of the English Church claimed by the crown.\footnote{The modification of the royal supremacy in Art. XXXVII., as compared with Art XXXVI. of Edward, was intended to meet the scruples of Romanists and Calvinists. Nevertheless this article, and the two acts of supremacy and uniformity, form the basis of that restrictive code of laws which pressed so heavily for more than two centuries upon the consciences of Roman Catholic and Protestant dissenters. Comp. the third chapter of Hallam's \emph{Constitutional History of England} (Harper's ed. pp. 71 sqq.).} All the Reformation Churches were more or less intolerant, and enforced uniformity of belief as far as they had the power; but the Calvinists and Puritans were more careful of the rights of the Church over against the State than the Lutherans.

6. Art. XXXV., referring to the Prayer-book and the consecration of archbishops, bishops, priests, and deacons, is purely Anglican and Episcopalian, and excited the opposition of the Puritans.

We have now to furnish the proof as far as the doctrinal articles are concerned.

THE ARTICLES AND THE AUGSBURG CONFESSION.

The Edwardine Articles were based in part, as already observed, upon a previous draft of Thirteen Articles, which was the joint product of German and English divines, and based upon the doctrinal Articles of the Augsburg Confession. Some passages were transferred verbatim from the Lutheran document to the Thirteen Articles, and from these to the Forty-two (1553), and were retained in the Elizabethan revision (1563 and 1571). This will appear from the following comparison. The corresponding words are printed in italics.

Augsburg Confession. 1530. Art. I. De Deo.

Thirteen Articles. 1538. Art. I. De Unitate Dei et Trinitate Personarum.

Thirty-nine Articles. 1563. Art. I. De Fide in Sacrosanctum Trinitatem.

Ecclesiæ magno consensu apud nos docent, Decretum Nicænæ Synodi, de unitate essentiæ divinæ et de tribus personis, verum et sine ulla dubitatione credendum esse. Videlicet, quod sit una essentia divina, quæ et appellatur et est \emph{Deus, æternus, incorporeus impartibilis, immensa potentia, sapientia, bonitate, creator et conservator omnium} rerum, \emph{visibilium} et \emph{invisibilium;} et tamen \emph{tres sint personæ, ejusdem essentiæ} et \emph{potentiæ,} et coæternæ, \emph{Pater, Filius et Spiritus Sanctus.} Et nomine personæ utuntur ea significatione, qua usi sunt in hac causa scriptores ecclesiastici, ut significet non partem aut qualitatem in alio, sed quod proprie subsistit.

De Unitate Essentiæ Divinæ et de Tribus Personis, censemus decretum Nicenæ Synodi verum, et sine ulla dubitatione credendum esse, videlicet, quod sit una Essentia Divina, quæ et appellatur et est \emph{Deus, æternus, incorporeus, impartibilis, immensa potentia, sapientia, bonitate, creator et conservator omnium} rerum \emph{visibilium} et \emph{invisibilium,} et tamen \emph{tres sint personæ ejusdem essentiæ} et \emph{potentiæ,} et coæternæ, \emph{Pater, Filius, et Spiritus Sanctus; }et nomine personæ utimur ea significatione qua usi sunt in hac causa scriptores ecclesiastici, ut significet non partem aut qualitatem in alio, sed quod proprie subsistit.

Unus est vivus et verus \emph{Deus æternus, incorporeus impartibilis,} impassibilis, \emph{immensæ potentiæ, sapientiæ} ac \emph{bonitatis: creator et conservator omnium} tum \emph{visibilium} tum \emph{invisibilium.} Et in unitate huius divinæ naturæ \emph{tres sunt Personæ ejusdem essentiæ, potentiæ,} ac æternitatis, \emph{Pater, Filius, et Spiritus Sanctus.}\footnote{The same passage occurs in the \emph{Reformatio Legum ecclesiasticarum} (De Summa Trinitate, c. 2), a work prepared by a committee consisting of Cranmer, Peter Martyr, and six others, 1551. It was edited by Cardwell, Oxford, 1850, and serves as a commentary on the Articles. See Hardwick, pp. 82 and 371.}

Damnant omnes hæreses, contra hunc articulum exortas, ut Manichæos, qui duo principia ponebant, Bonum et Malum; item Valentinianos, Arianos, Eunomianos, Mahometistas,

Damnamus omnes hæreses contra hunc articulum exortas, ut Manichæos, qui duo principia ponebant, Bonum et Malum: item Valentinianos, Arianos, Eunomianos, Mahometistas,

et omnes horum similes. Damnant et Samosatenos, veteres et neotericos, qui, cum tantum unam personam esse contendant, de Verbo et de Spiritu Sancto astute et impie rhetoricantur, quod non sint personæ distinctæ, sed quod Verbum significet verbum vocale, et Spiritus motum in rebus creatum.

et omnes horum similes. Damnamus et Samosatenos, veteres et neotericos, qui cum tantum unam personam esse contendant, de Verbo et Spiritu Sancto astute et impie rhetoricantur, quod non sint personæ distinctæ, sed quod Verbum significet verbum vocale, et Spiritus motum in rebus creatum.

Art. III. De Filio Dei.

Art. III. De Duabus Christi Naturis.

Art. II. Verbum Dei verum hominem esse factum.

Item, docent, quod \emph{Verbum,} hoc est, \emph{Filius} Dei, \emph{assumpserit humanam naturam in utero beatæ} Mariæ \emph{virginis, ut} sint \emph{duæ naturæ, divina et humana, in unitate personæ inseparabiliter conjunctæ, unus Christus, vere Deus et vere homo,} natus ex virgine Maria, \emph{vere passus, crucifixus, mortuus, et sepultus, ut reconciliaret nobis Patrem,} et \emph{hostia esset non tantum pro culpa originis,} sed \emph{etiam pro omnibus actualibus hominum peccatis.}

Item docemus, quod \emph{Verbum,} hoc est \emph{Filius} Dei, \emph{assumpserit humanam naturam in utero beatæ} Mariæ \emph{virginis, ut} sint \emph{duæ naturæ, divina et humana, in unitate personæ inseparabiliter conjunctæ, unus Christus, vere Deus, et vere homo,} natus ex virgine Maria, \emph{vere passus, crucifixus, mortuus, et sepultus, ut reconciliaret nobis Patrem,} et \emph{hostia esset non tantum pro culpa originis,} sed \emph{etiam pro omnibus actualibus hominum peccatis.}

\emph{Filius,} qui est \emph{Verbum} Patris ab æterno a Patre genitus verus et æternus Deus, ac Patri consubstantialis, \emph{in utero Beatæ virginis} ex illius substantia \emph{naturam humanam assumpsit:} ita \emph{ut duæ naturæ, divina et humana} integre atque perfecte \emph{in unitate personæ,} fuerint \emph{inseparabiliter coniunctæ:} ex quibus est \emph{unus Christus, verus Deus et verus homo:} qui \emph{vere passus} est, \emph{crucifixus, mortuus, et sepultus, ut Patrem nobis reconciliaret, esset}que \emph{hostia non tantum pro culpa originis,} verum \emph{etiam pro omnibus actualibus hominum peccatis.}

Idem descendit ad inferos, et vere resurrexit tertia die, deinde ascendit ad cœlos, ut sedeat ad dexteram Patris, et perpetuo regnet et dominetur omnibus creaturis, sanctificet credentes in ipsum, misso in corda eorum Spiritu Sancto, qui regat, consoletur ac vivificet eos, ac defendat adversus diabolum et vim peccati.

Item descendit ad inferos, et vere resurrexit tertia die, deinde ascendit ad cœlos, ut sedeat ad dexteram Patris et perpetuo regnet et dominetur omnibus creaturis, sanctificet credentes in ipsum, misso in corde eorum Spiritu Sancto, qui regat, consoletur, ac vivificet eos, ac defendat adversus diabolum et vim peccati.

Idem Christus palam est rediturus, ut judicet vivos et mortuos, etc., juxta Symbolum Apostolorum.

Idem Christus palam est rediturus ut judicet vivos et mortuos, etc., juxta Symbolum Apostolorum.

Art. IV. De Justificatione.

Art. IV. De Justificatione.

Art. XI. De Hominis Iustificatione.

Item docent, quod homines non possint justificari \emph{coram Deo} propriis viribus, \emph{meritis} aut \emph{operibus,} sed gratis justificentur \emph{propter Christum per fidem,} cum credunt se in gratiam recipi, et peccata remitti propter Christum, qui sua morte pro nostris peccatis satisfecit. Hanc fidem imputat Deus pro justitia coram ipso. Rom. III. et IV.

[Art. IV. of the Augsburg Confession is enlarged, and Art. V. added. In this case the English Articles do not give the language, but the sense of the Lutheran symbols, with the unmistakeable 'sola fide,' which was Luther's watchword.]

Tantum \emph{propter} meritum Domini ac Servatoris nostri Iesu \emph{Christi, per fidem,} non propter \emph{opera} et \emph{merita} nostra, iusti \emph{coram Deo} reputamur. Quare sola fide nos \emph{iustificari,} doctrina est saluberrima, ac consolationis plenissima: ut in Homilia de Iustificatione hominis fusius explicatur.

Art. VII. De Ecclesia.

Art. V. De Ecclesia.

Art. XIX. De Ecclesia.

Item docent, quod una Sancta Ecclesia pepetuo mansura sit. \emph{Est} autem \emph{Ecclesia} congregatio Sanctorum [Versammlung aller \emph{Gläubigen}], in qua Evangelium recte [\emph{rein}] docetur, et \emph{recte} [laut des Evangelii] \emph{administrantur Sacramenta.} Et ad veram unitatem Ecclesiæ satis est consentire de doctrina Evangelii et administratione Sacramentorum. Nec necesse est ubique esse similes traditiones humanas, seu ritus aut ceremonias, ab hominibus institutas. Sicut inquit Paulus (Eph. iv. 5, 6): Una fides, unum Baptisma, unus Deus et Pater omnium, etc.

[This Article is much enlarged, and makes an important distinction between the Church as the 'congregatio omnium sanctorum et \emph{ fidelium},' (the invisible Church), which is the mystical body of Christ, and the Church as the 'congregatio omnium hominum qui baptizati sunt' (the visible Church).]

\emph{Ecclesia} Christi visibilis, \emph{est} cœtus \emph{fidelium, in quo} verbum Dei purum prædicatur, et \emph{ sacromenta,} quoad ea quæ necessario exiguntur, iuxta Christi institutum \emph{recte administrantur.}\footnote{\par The silence of this Article concerning the episcopal succession was observed by Joliffe, prebendary at Worcester, who added among the marks of the Church, '\emph{legitima et continua successio vicariorum Christi.}'}  Sicut erravit Ecclesia Hierosolymitana, Alexandrina et Antiochena: ita et erravit Ecclesia Romana, non solum quoad agenda et cæremoniarum ritus, verum in his etiam quæ credenda sunt. [Compare Art. XXXIII., which treats of ecclesiastical traditions, and corresponds in sentiment to the second clause in Art. VII. of the Augsburg Confession.]

Art. XIII. De Usu Sacramentorum.

Art. IX. De Sacramentorum Usu.

Art. XXV. De Sacramentis.

De usu Sacramentorum docent, quod \emph{Sacramenta instituta,} sint, \emph{non} modo ut

Docemus, quod \emph{Sacramenta} quæ per verbum Dei \emph{instituta} sunt, \emph{non tantum}

\emph{Sacramenta} a Christo \emph{instituta non tantum sunt notæ professionis Christianorum},

sint \emph{notæ professionis} inter homines, \emph{sed} magis ut sint \emph{signa} et \emph{testimonia voluntatis Dei} erga \emph{nos, ad excitandam et confirmandam fidem,} in his, qui utuntur, proposita. Itaque utendum est Sacramentis ita, ut fides accedat, quæ credat promissionibus, quæ per Sacramenta exhibentur et ostenduntur. Damnant igitur illos, qui docent, quod Sacramenta ex opere operato justificent, nec docent fidem requiri in usu Sacramentorum, quæ credat remitti peccata.

\emph{sint notæ professionis} inter \emph{Christianos, sed} magis \emph{certa quædam testimonia et efficacia signa gratiæ,} et \emph{bonæ voluntatis Dei} erga \emph{nos, per quæ} Deus \emph{invisibiliter operatur in nobis,} et suam gratiam in nos invisibiliter diffundit, siquidem ea rite susceperimus; quodque per ea \emph{excitatur} et \emph{confirmatur fides} in his qui eis utuntur. Porro docemus, quod ita utendum sit sacramentis, ut in adultis, præter veram contritionem, necessario etiam debeat accedere fides, quæ credat præsentibus promissionibus, quæ per sacramenta ostenduntur, exhibentur, et præstantur. Neque, etc.

\emph{sed certa quædam} potius \emph{testimonia,} et \emph{efficacia signa gratiæ} atque \emph{bonæ} in \emph{nos voluntatis Dei, per quæ invisibiliter} ipse \emph{in nobis operatur,} nostramque \emph{fidem} in se, non solum \emph{excitat,} verum etiam \emph{confirmat.}

Besides these passages, there is a close resemblance in thought, though not in language, in the statements of the doctrine of original sin,\footnote{Conf. Aug. Art. II., English Art. IX., from Augustine.} and of the possibility of falling after justification.\footnote{Conf. Aug. Art. XII. ('\emph{Damnant Anabaptistas qui negant semel justificatos posse amittere Spiritum Sanctum},' etc.), English Art. XVI.} Several of the Edwardine Articles, also, which were omitted in the Elizabethan revision, were suggested by Art XVII. of the Augsburg Confession, which is directed against the Anabaptists.

THE ARTICLES AND THE WÜRTEMBERG CONFESSION.

In the Elizabethan revision of the Articles another Lutheran Confession was used (in Arts. II., V., VI., X., XI., and XX.)—namely, the \emph{Confessio Würtembergica,} drawn up by the Suabian Reformer, Brentius (at a time when he was still in full harmony with Melanchthon), in the name of Duke Christopher of Würtemberg (1551), and presented by his delegates to the Council of Trent (Jan. 24, 1552).\footnote{Printed in the \emph{Corpus et Syntagma Conf.,} and in Dr. Heppe's \emph{Bekenntniss-Schriften der altprotestantischen Kirche Deutschlands,} Cassel, 1855, pp. 491–554. See above, § 47, pp. 343 sq. Archbishop Laurence (\emph{Bampton Lectures,} pp. 40 and 233 sqq.) first discovered and pointed out this resemblance. Hardwick (pp. 126 sqq.) and the 'Interleaved Prayer-Book' speak of the Confession of Brentius alternately as the 'Saxon' Confession, and the 'Würtemberg' (or Wirtemburg!) Confession, as if the Saxon city of Wittenberg and the Duchy (now Kingdom) of Würtemberg were one and the same. The 'Saxon Confession,' so called, or the 'Repetition of the Augsburg Confession,' is a different document, written about the same time and for the same purpose by Melanchthon, in behalf of the Wittenberg and other Saxon divines. See above, p. 340, and the Oxford \emph{Sylloge,} which incorporates the Saxon but not the Würtemberg Confession.} Soon after the accession of Elizabeth the negotiations with the German Lutherans (which had been broken off in 1538) were resumed, with a view to join the Smalcaldian League, but led to no definite result. It was probably during these negotiations that the Würtemberg Confession became known in England; and as it had acquired a public notoriety by its presentation at Trent, and was a restatement of the Augsburg Confession adapted to the new condition of things, it was very natural that it should be compared in the revision of the Articles. Melanchthon's 'Saxon Repetition of the Augsburg Confession' would indeed have answered the same purpose equally well, but perhaps it was not known in time.

Confessio Würtembergica, 1552.

Thirty-nine Articles, 1563.

Art. II. De Filio Dei (Heppe, p.492).

Art. II. Verbum Dei verum hominem esse factum.

Credimus et confitemur Filium Dei, Dominum nostrum Jesum Christum, ab æterno a Patre suo genitum, verum et æternum Deum, Patri suo consubstantialem, et in plenitudine temporia factum hominem, etc.

Ab æterno a Patre genitus, verus et æternus Deus, ac Patri consubstantialis.

Art. III. De Spiritu Sancto (Heppe, p. 493).

Art. V. De Spiritu Sancto.

Credimus et confitemur Spiritum Sanctum ab æterno procedere a Deo Patre et Filio, et esse ejusdem cum Patre et Filio essentiæ, majestatis, et gloriæ, verum ac æternum Deum.

Spiritus Sanctus, a Patre et Filio procedens, ejusdem est cum Patre et Filio essentiæ, majestatis, et gloriæ, verus ac æternus Deus.

Art. XXX. De Sacra Scriptura (Heppe, p. 540).

Art. VI. Divinæ Scripturæ doctrina sufficit ad salutem.

Sacram Scripturam vocamus eos Canonicos libros veteris et novi Testamenti, de quorum authoritate in Ecclesia nunquam dubitatum est.

. . . Sacræ Scripturæ nomine eos Canonicos libros veteris et novi Testamenti intelligimus, de quorum auctoritate in Ecclesia nunquam dubitatum est.

Art. IV. De Peccato (Heppe, p.498).

Art. X. De Libero Arbitrio.

Quod autem nonnulli affirmant homini post lapsum tantam animi integritatem relictam, ut possit sese, naturalibus suis viribus et bonis operibus, ad fidem et invocationem Dei convertere ac præparare, haud obscure pugnat cum Apostolica doctrina, et cum vero Ecclesiæ Catholicæ consensu.

Ea est hominis post lapsum Adæ conditio, ut sese, naturalibus suis viribus et bonis operibus, ad fidem et invocationem Dei convertere ac præparare non possit. [The next clause, 'Quare absque gratia Dei,' etc., is taken almost verbatim from Augustine, \emph{ De gratia et lib. arbitrio,} c. 17 (al. 33).]

Art. V. De Justificatione (Heppe, p. 495).

Art. XI. De Hominis Justificatione.

Homo enim fit Deo acceptus, et reputatur coram eo justus, propter solum Filium Dei, Dominum nostrum Jesum Christum, per fidem.

Tantum propter meritum Domini ac Servatoris nostri Jesu Christi, per fidem, non propter opera et merita nostra, justi coram Deo reputamur.

Art. VIII. De Evangelio Christi (Heppe, p. 500.

Nec veteris nec novi Testamenti hominibus contingat æterna salus propter meritum operum Legis, sed tantum propter meritum Domini nostri Jesu Christi, per fidem.

Art. VII. De Bonis Operibus (Heppe, p. 499).

Art. XII. De Bonis Operibus.

Non est autem sentiendum, quod iis bonis operibus, quæ per nos facimus, in judicio Dei, ubi agitur de expiatione peccatorum, et placatione divinæ iræ, ac merito æternæ salutis, confidendem sit. Omnia enim bona opera, quæ nos facimus, sunt imperfecta, nec possunt severitatem divini judicii ferre.

Bona opera, quæ sunt fructus fidei, et justificatos sequuntur, quanquam peccata nostra expiare, et divini judicii severitatem ferre non possunt, Deo tamen, grata sunt et accepta in Christo. . . .

Art. XXXII. De Ecclesia (Heppe, p. 544).

Art. XX. De Ecclesiæ Autoritate.

Credimus et confitemur, quod una sit Sancta Catholica et Apostolica Ecclesia, juxta symbolum Apostolorum et Nicænum. . . .

Habet Ecclesia ritus sive ceremonias statuendi jus, et in fidei controversiis auctoritatem, quamvis Ecclesiæ non licet quicquam instituere, quod verbo Dei scripto adversetur nec unum Scripturæ locum sic exponere potest ut alteri contradicat

Quod hæc Ecclesia habeat jus judicandi de omnibus doctrinis, juxta illud, \emph{Probate spiritus, num ex Deo sint.}

Quod hæc Ecclesia habeat jus interpretandæ Scripturæ.

THE ARTICLES AND THE REFORMED CONFESSIONS.

We now proceed to those doctrines in which the Lutheran and the Reformed Churches differed and finally separated—namely, the doctrines of predestination and the eucharistic presence. Here we find the English Articles on the Reformed side. The authors and revisers formed their views on these subjects partly from an independent study of the Scriptures and Augustine, partly from contact with the Swiss divines.

The principal Reformed Confessions were indeed published at a later date—the Gallican Confession in 1559; the Belgic in 1561; the Heidelberg Catechism in 1563; the Second Helvetic Confession in 1566. But Zwingli's and Bullinger's works, Calvin's Institutes (1536), and his Tract on the Lord's Supper (1541), the Zurich Consensus (1549), and the Geneva Consensus (1552), must have been more or less known in England. Bishop Hooper had become a thorough disciple of Bullinger by a long residence in Zurich before the accession of Edward VI., and was consulted on the Articles. Cranmer (as previously mentioned) embraced, with Ridley, the Reformed doctrine of the Lord's Supper as early as 1548; he corresponded with the Swiss Reformers, as well as with Melanchthon, and invited them (March 1552) to England to frame a general creed; and he was in intimate personal connection with Bucer, Peter Martyr, John Laski, and Knox at the time he framed the Articles.\footnote{One of the last letters of Cranmer was written from his prison, 1555, to Peter Martyr, who was a decided Calvinist. See \emph{Zurich Letters,} First Series, Vol. I. p. 29.} From the same period we have a remarkable witness to the influence of Calvin's tracts in defense of the doctrine of predestination.\footnote{See above, p. 474.} Bartholomew Traheron, then Dean of Chichester, and librarian to King Edward, wrote to Bullinger from London, Sept. 10, 1552, as follows:\footnote{\emph{Zurich Letters,} First Series, Vol. I. p. 325.} 'I am exceedingly desirous to know what you and the other very learned men who live at Zurich think respecting the predestination and providence of God. If you ask the reason, there are certain individuals here who lived among you some time, and who assert that you lean too much to Melanchthon's views.\footnote{From this we might infer that Melanchthon's influence, in consequence of his abandonment of absolute predestinarianism, was declining in England, while Calvin's was increasing.} But the \emph{greater} \emph{number among us,} of whom I own myself to be one, embrace the opinion of John Calvin as being perspicuous, and most agreeable to holy Scripture. And we truly thank God that that excellent treatise of the very learned and excellent John Calvin against Pighius and one Georgius Siculus should have come forth at the very time when the question began to be agitated among us.\footnote{He means the \emph{Consensus Genevensis de æterna Dei prædestinatione,} which appeared in 1552, and acquired semi-symbolical authority in Geneva. Calvin had also previously (1543) written a tract against Pighius on the doctrine of free-will, and dedicated it to Melanchthon, who gratefully acknowledged the compliment, but modestly intimated his dissent and his inability to harmonize the all-ruling providence of God with the action of the human will. See Stähelin, \emph{Calv.} Vol. I. p. 241.} For we confess that he has thrown, much light upon the subject, or rather so handled it as that we have never before seen any thing more learned or more plain. We are anxious, however, to know what are your opinions, to which we justly allow much weight. We certainly hope that you differ in no respect from his excellent and most learned opinion. At least you will please to point out what you approve in that treatise, or think defective, or reject altogether, if indeed you do reject any part of it, which we shall not easily believe.' To this letter Bullinger replied at length, but not to the satisfaction of the Dean, who wrote to him again, June 3, 1553, as follows:\footnote{\emph{Zurich Letters,} First Series, Vol. I. p. 327. Bullinger's tract \emph{De providentia,} which was occasioned by Traheron, is still extant in MS. in Zurich, and is fully noticed by Schweizer. See above, p. 475.} 'You do not approve of Calvin, when he states that God not only foresaw the fall of the first man, and in him the ruin of his posterity, but that he also at his own pleasure arranged it. But unless we allow this, we shall certainly take away both the providence and the wisdom of God altogether. I do not indeed perceive how this sentence of Solomon contains any thing less than this: ``The Lord hath made all things for himself; yea, even the wicked for the day of evil'' (Prov.xvi.4). And that of Paul: ``Of him and through him, and to him are all things'' (Rom. xi. 36). I pass over other expressions which the most learned Calvin employs, because they occur everywhere in the holy Scriptures.'

The Elizabethan revision was the work of the Marian exiles, who felt themselves in complete theological harmony with the Swiss divines, especially with Bullinger of Zurich, who represented an improved type of Zwinglianism, and agreed with Calvin on the subject of the Lord's Supper (as expressed in the Consensus Tigurinus, 1549), but was more moderate and guarded on the subject of predestination.\footnote{\par On Bullinger's intimate personal relations with English divines, which began before the reign of Edward and continued till his death (1575), compare Pestalozzi's \emph{Heinrich Bullinger,} pp. 441 sqq.} His writings seem to have been better known and exerted more influence in the earlier part of Elizabeth's reign than those of Calvin, which were more congenial to the Scotch mind; but they became all-powerful even in England towards the close of the sixteenth century.

On this point we have the explicit testimonies of the very men who were the chief assistants of Archbishop Parker in the revision of the Articles. Bishop Horn, of Winchester, wrote to Henry Bullinger, Dec. 13, 1563, soon after the adoption of the Latin revision: 'We have throughout England the same ecclesiastical doctrine as yourselves. . . . The people of England entertain on these points' [the sacraments, and 'against the ubiquitarianism of Brentius'] 'the same opinions as you do at Zurich.'\footnote{\par \emph{Zurich Letters,} Second Series, Vol. I. (A.D. 1558–1579), p. 135.} Bishop Grindal, of London, afterwards (1575) the successor of Parker in the primacy, wrote to Bullinger, Aug. 27, 1566: 'We, who are now bishops, most fully agree in the pure doctrines of the gospel with your churches, and with the Confession you have lately set forth' [i.e., the Second Helvetic Confession, which appeared in the same year]. 'And we do not regret our resolution; for in the mean time, the Lord giving the increase, our churches are enlarged and established, which under other circumstances would have become a prey to the Ecebolians, Lutherans, and semi-papists.'\footnote{Ibid. p. 169. Ecebolus was a sophist of Constantinople in the fourth century, who followed the Emperor Julian in his apostasy.} In a letter to Calvin, dated June 19, 1563, Grindal says: 'As you and Bullinger are almost the only chief pillars remaining, we desire to enjoy you both (if it please God) as long as possible. I purposely omit mention of Brentius, who having undertaken the advocacy of the very worst of causes' [ubiquitarianism], 'seems no longer to acknowledge us as brethren.'\footnote{\par Ibid. Vol. II. p. 97. Brentius advocated the absolute ubiquity of Christ's body, and fiercely attacked the Reformed in several tracts, from 1560 to 1564 (ten years after he wrote the Würtemberg Confession). He was answered by Bullinger and Peter Martyr. See above, p. 290.} The letters of Bishop Cox, of Ely, to Bullinger and Peter Martyr, though not so explicit, breathe the same spirit of grateful respect and affection. The strong testimony of Bishop Jewel of Salesbury, the final reviser of the English text and chief author of the Second Book of Homilies, we have already quoted.\footnote{See his letter to his revered teacher, Peter Martyr, p. 603. Grindal called him after his death (Sept. 22, 1571), 'the jewel and singular ornament of the Church, as his name implies.'—\emph{Zurich Letters,} Second Series, Vol. I. p. 260. An adversary, Moren, said of him : 'I should love thee, Jewel, if thou wert not a Zwinglian; in thy faith I hold thee an heretic, but surely in thy life thou art an angel.' Queen Elizabeth ordered a copy of Jewel's 'Apology of the Church of England' (1562) to be chained in every parish church.}

PREDESTINATION AND ELECTION.

On the premundane mystery of predestination, which no system of philosophy or theology can satisfactorily solve in this world, and which ought to be approached with profound reverence and humility, all the Reformers, in their private writings, followed originally the teaching of the great Augustine and the greater St. Paul; meaning thereby to cut human merit and pride at the roots, and to give all the glory of our salvation to God alone. But the Lutheran symbols (with the exception of the later Formula of Concord) are silent on the subject, while most of the Reformed standards, under the influence of Calvin, give it a prominent place. The English Articles handle it with much wisdom and moderation, dwelling exclusively on the election of saints or predestination to life. We give the XVIIth Article in its original form with the later amendments; the clauses which were omitted in the Elizabethan revision are printed in \emph{italics,} the words which were inserted or substituted are inclosed in brackets.

Art. XVIII.

OF PREDESTINATION AND ELECTION.

Predestination to Life is the everlasting purpose of God, whereby (before the foundations of the world were laid) he hath constantly decreed by his counsel secret to us, to deliver from curse and damnation those whom he hath chosen [in Christ]\footnote{The insertion 'in Christ' is Scriptural and in accordance with all the Reformed Confessions. There is no election out of Christ or apart from Christ.} out of mankind, and to bring them by Christ to everlasting salvation, as vessels made to honor. Wherefore, \emph{such as have} [they which be endued with] so excellent a benefit of God \emph{given unto them,} be called according to God's purpose by his Spirit working in due season: they through grace obey the calling: they be justified freely: they be made sons [of God] by adoption: they be made like the image of \emph{God's} [his] only begotten Son Jesus Christ: they walk religiously in good works, and at length, by God's mercy, they attain to everlasting felicity.

As the godly consideration of Predestination, and our Election in Christ, is full of sweet, pleasant, and unspeakable comfort to godly persons, and such as feel in themselves the working of the Spirit of Christ, mortifying the works of the flesh, and their earthly members, and drawing up their mind to high and heavenly things, as well because it doth greatly establish and confirm their faith of eternal salvation to be enjoyed through Christ, as because it doth fervently kindle their love towards God: so, for curious and carnal persons, lacking the Spirit of Christ, to have continually before their eyes the sentence of God's Predestination, is a most dangerous downfall, whereby the Devil \emph{may} [doth] thrust them either into desperation, or into wretchlessness of most unclean living, no less perilous than desperation.

Furthermore, \emph{although the Decrees of Predestination are unknown unto us, yet} we must receive God's promises in such wise, as they be generally set forth to us in holy Scripture; and, in our doings, that Will of God is to be followed, which we have expressly declared unto us in the Word of God.

This Article can not be derived from the Augsburg Confession, nor from the Thirteen Articles, nor from the Würtemberg Confession—for they omit the subject of predestination altogether\footnote{With the exception of an incidental allusion to the absolute freedom of divine grace in the Augsburg Confession, Art. V., De Ministerio: '\emph{Per verbum et sacramenta tamquam per instrumenta donatur Spiritus Sanctus, qui fidem efficit, } ubi et quando visum est Deo, \emph{in iis qui audiunt evangelium.}' Compare with this the expression of the Form. Concordiæ (Sol. decl. Art. II. de lib. arbitr. p. 673): '\emph{Trahit Deus hominem, } quem convertere decrevit.' It is significant that in the altered edition of 1540 Melanchthon omitted the words \emph{ubi et quando visum est Deo},' as also the words '\emph{non adjuvante Deo}' in Art. XIX. The brevity of allusion shows that even in 1530, although still holding to the Augustinian scheme, he laid less stress on it than in the first edition of his \emph{Loci.} This appears also from a letter to Brentius, Sept. 30, 1531 (\emph{Corp. Ref.} Vol. II. p. 547), where Melanchthon says: '\emph{Sed ego in tota Apologia fugi illam longam et inexplicabilem disputationem de prædestinatione. Ubique sic loquor, quasi prædestinatio sequatur nostram fidem et opera.}'}—nor from Melanchthon's private writings, for he abandoned his former views, and suggested the synergistic theory as early as 1535, and more fully in 1548.\footnote{See above, pp. 262 sqq., and Schweizer, \emph{Centraldogmen,} Vol. I. p. 384. There is not a trace of synergism in the XVIIth Art, and Art. X. expressly denies the freedom of will, while Melanchthon asserts it in the later editions of his \emph{Loci} ('\emph{Liberum arbitrium esse in homine facultatem applicandi se ad gratiam}'). Laurence (p. 179) and Hardwick (p. 383) derive the last clause about the 'general' promises and the 'revealed will' from Melanchthon, but the same sentiments are found in Calvin, Bullinger, and the Reformed Confessions. See below.} It can not be naturally understood in any other than an Augustinian or moderately Calvinistic sense. It does not, indeed, go as far as the Lambeth Articles (1595), which the stronger Calvinism of the rising generation thought necessary to add as an explanation. It omits the knotty points; it is cautiously framed and guarded against abuse.\footnote{This element of caution and modesty is well expressed by Bishop Ridley: 'In these matters [of God's election] I am so fearful that I dare not speak further, yea, almost none otherwise than the very text doth, as it were, lead me by the hand.' Ridley's \emph{Works} (Parker ed.), p. 368. He thus wrote in a letter of sympathy to his friend and chaplain, Bradford, who in prison, at London, had a dispute with a certain 'free-willer,' Henry Hart, and wrote an excellent 'Defense of Election.' This treatise was approved by his fellow-prisoners, and shows what an unspeakable comfort they derived from this doctrine. See \emph{The Writings of John Bradford, Martyr,} 1555 (Parker Soc. ed.), pp. 307 sqq.} But it very clearly teaches a free eternal election in Christ, which carries with it, by way of execution in time, the certainty of the call, justification, adoption, sanctification, and final glorification (Rom. viii. 29, 30).

This is all that is essential, and a matter of dogma in the Reformed Churches; the rest of what is technically called Calvinism, in distinction from Arminianism, is logical inference, and belongs to the theology of the school. It should be remembered that all the Reformed Confessions (even the Canons of Dort, the Westminster Confession, and the Helvetic Consensus Formula) keep within the limits of infralapsarianism, which puts the fall under a \emph{permissive} decree, and makes man alone responsible for sin and condemnation; the most authoritative, as the Helvetic Confession of Bullinger, the Heidelberg Catechism, and the Brandenburg Confessions (also the Scotch Confession of 1560) teach only the positive and comforting part of predestination, and ignore or deny a separate decree of reprobation; thus taking the ground practically that all that are saved are saved by the free grace of God, while all that are lost are lost by their own guilt. They also teach that God's promises and Christ's redemption are general, and that we must abide by the \emph{revealed} will of God, which sincerely offers the gospel salvation to all who repent and believe.\footnote{Conf. Helv. post., cap. X.: '\emph{Bene sperandum est de omnibus. Vestrum non est de his curiosius inquirere. . . . Audienda est prædicatio evangelii, eique credendum est, et pro indubitato habendum, si credis ac sis in Christo, electum te esse.} . . . ``\emph{Venite ad me omnes},'' etc. . . . ``\emph{Sic Deus dilexit mundum},'' etc. . . . ``\emph{Non est voluntas Patris, ut quisque de his pusillis pereat.}'' . . . \emph{Promissiones Dei sunt universales fidelibus}' (not \emph{electis}), etc. Heidelb. Cat., Qu. 37: 'Christ bore the wrath of God against the sin of the whole human race (1 Pet. ii. 24; 1 John ii. 2, etc.).' Conf. Belg., Art. XIII.: '\emph{Sufficit nobis ea duntaxat discere quæ ipse verbo suo nos docet, neque hos fines transilire fas esse ducimus.}' Calvin himself often warns against idle curiosity and speculation on the secret will of God, and exhorts men to abide by the revealed will of God. See the passages quoted by Stähelin, Vol. II. p. 279. Comp. the remarks of Dr. Julius Müller on the Reformed Confessions concerning predestination, in his work, \emph{Die evang. Union} (1854), p. 214, and his \emph{Dogmat. Abhandlungen} (1870), p. 194.}

The remarks of the Article about the 'sweet, pleasant, and unspeakable comfort' of our election in Christ, and the caution against abuse by carnal persons, are consistent only with the Calvinistic interpretation, and wholly inapplicable to Arminian views, which are neither comfortable nor dangerous, and have never thrust any man 'into desperation, or into wretchlessness of most unclean living.'\footnote{Dr. Cunningham (\emph{The Reformers and the Theology of the Reformation,} p. 194), says: 'It is only the Calvinistic, and not the Arminian doctrine that suggests or requires such guards or caveats; and it is plainly impossible that such a statement could ever have occurred to the compilers of the Articles as proper and necessary, unless they had been distinctly aware that they had just laid down a statement which at least included the Calvinistic doctrine.'}

The view here taken is confirmed by the contemporary testimonies already quoted, and by the first learned commentator of the Articles, Thomas Rogers, who was chaplain to Archbishop Bancroft, and did not sympathize with the Puritan party. He draws the following propositions from the XVIIth Article, and fortifies them with abundant Scripture passages:\footnote{\emph{The Catholic Doctrine of the Church of England,} etc., first published, London, 1586, Parker Society ed. (by J. J. S. Perowne), 1854, p. 143. This important work has not been even alluded to by any writer I have consulted on the subject.}

'1. There is a predestination of men unto everlasting life.

'2. Predestination hath been from everlasting.

'3. They who are predestinate unto salvation can not perish.

'4. Not all men, but certain, are predestinate to be saved.

'5. In Christ Jesus, of the mere will and purpose of God, some are elected, and not others, unto salvation.

'6. They who are elected unto salvation, if they come unto years of discretion, are called both outwardly by the Word and inwardly by the Spirit of God.

'7. The predestinate are both justified by faith, sanctified by the Holy Ghost, and shall be glorified in the life to come.

'8. The consideration of predestination is to the godly-wise most comfortable, but to curious and carnal persons very dangerous.

'9. The general promises of God, set forth in the holy Scriptures, are to be embraced of us.

'10. In our actions, the Word of God, which is his revealed will, must be our direction.'

To this theological comment I add the judgment of an impartial and well-informed secular historian. Henry Hallam\footnote{\emph{Constit. History of England,} ch. vii. p. 230 (Amer. ed.).} declares that the Articles on predestination, original sin, and total depravity, 'after making every allowance for want of precision, are totally irreconcilable with the scheme usually denominated Arminian.' He justly appeals in confirmation of this judgment to contemporary and other early authorities, and adds: 'Whatever doubts may be raised as to the Calvinism of Cranmer and Ridley, there can surely be no room for any as to the chiefs of the Anglican Church under Elizabeth. We find explicit proofs that Jewel, Nowell, Sandys, and Cox professed to concur with the Reformers of Zurich and Geneva in every point of doctrine. The works of Calvin and Bullinger became the text-books in the English universities. Those who did not hold the predestinarian theory were branded with reproach by the name of Free-willers and Pelagians; and when the opposite tenets came to be advanced, as they were at Cambridge about 1590, a clamor was raised as if some unusual heresy had been broached.'

The Arminian interpretation of the Article under consideration is an anachronism and a failure. The Lutheran interpretation is more plausible, but true only so far as the Lutheran system is itself Augustinian. The Tractarian interpretation, which identifies eternal election with ecclesiastical calling, and the elect with the baptized, is contrary both to the spirit and letter of the Article. It must in all fairness be admitted that Art. XVII., in connection with Arts. X. and XIII., implies the infralapsarian scheme, and that the Lambeth Articles are not a reaction, but a legitimate though one-sided development.

Note.—The anti-Calvinistic interpretation began after the Synod of Dort with Archbishop Laud, or his biographer, Peter Heylin (in his \emph{Historia Quinqu-Articularis,} London, 1660, which was answered and refuted by Henry Hickman, in his \emph{Historia Quinqu-Articularis Exarticulata,} 1673). It was maintained, with hesitation, by Waterland (1721), more decidedly by Dr. Winchester, d. 1780 (\emph{Dissertation on the XVIIth Article,} new ed. London, 1808); by Dean Kipling (\emph{The Articles of the Church of England proved not to be Calvinistic,} Cambridge, 1802); by Bishop Tomline, d. 1827 (\emph{A Refutation of Calvinism,} London, 1811); and, with considerable learning, by Archbishop Laurence, d. 1839 (\emph{Bampt. Lect.,} Lect. VII. and VIII., Oxford, 1834, 3d ed. 1838), and by Hardwick (\emph{Hist. of the Articles}).

Laurence and Hardwick, as already remarked, trace Article XVII to Lutheran sources, but they overlook the difference between the Lutheran system (which admits the Augustinian premises, and even the doctrine of unconditional election of grace—see the \emph{formula of Concord,} ch. xi.) and the Arminian system (which denies the Augustinian anthropology, and makes both election and reprobation conditional), and show more dislike than real knowledge of Calvin. It is little less than a caricature when Laurence says of Calvin that his 'love of hypothesis' was superior to his great talent and piety (p. 43); that his '\emph{vanity} induced him to frame a peculiar system of his own' (pp. 262, 263), and that 'no man, perhaps, was ever less scrupulous in the adoption of general expressions, and no man adopted them with more mental reservations' (p. 375). Principal Cunningham has exposed this unfairness (\emph{The Reformers and the Theology of the Reformers,} 1866, pp. 179 sqq.).

Bishop Burnet (who was an Arminian and Latitudinarian) and Bishop Browne (a moderate High-Churchman) hesitate between the Augustinian and the Arminian interpretation. Burnet, after calmly reviewing the different theories of predestination, says (p. 236, Oxford ed.): 'It is not to be denied, but that the Article seems to be framed according to \emph{St. Austin's} doctrine: it supposes men to be under a \emph{curse and damnation,} antecedently to predestination, from which they are delivered by it; so it is directly against the supralapsarian doctrine; nor does the Article make any mention of reprobation—no, not in a hint; no definition is made concerning it. The Article does also seem to assert the efficacy of grace—that in which the knot of the whole difficulty lies is not defined; that is, whether God's eternal purpose or decree was made according to what he foresaw his creatures would do, or purely upon an absolute will, in order to his own glory. It is very probable that those who penned it meant that the decree was absolute; but yet since they have not said it, those who subscribe the Articles do not seem to be bound to any thing that is not expressed in them; and, therefore, since the Remonstrants do not deny but that God having foreseen what all mankind would, according to all the different circumstances in which they should be put, do or not do, he upon that did by a firm and eternal decree lay that whole design in all its branches, which he executes in time; they may subscribe this Article without renouncing their opinion as to this matter. On the other hand, the Calvinists have less occasion for scruple, since the Article \emph{does seem more plainly to favor them.} The three cautions that are added to it do likewise intimate that St. Austin's doctrine was designed to be settled by the Articles for \emph{the danger of men's having the sentence of God's predestination always before their eyes, which may occasion either desperation on the one hand, or the wretchedness of most unclean living on the other,} belongs only to that side; since these mischiefs do not arise out of the other hypothesis. The other two, of taking \emph{the promises of God in the sense in which they are set forth to us in holy Scriptures,} and \emph{of following that will of God that is expressly declared to us in the Word of God,} relate very visibly to the same opinion.

Bishop Browne, after a long discussion, comes to the conclusion (p. 425) that 'the Article was designedly drawn up in guarded and general terms, on purpose to comprehend all persons of tolerably sober views. . . . I am strongly disposed to believe that Cranmer's own opinions were certainly neither Arminian nor Calvinistic, nor probably even Augustinian; yet I can hardly think that he would have so worded this Article had he intended to declare very decidedly against either explanation of the doctrine of election.'

Bishop Forbes, a Tractarian, admits the Article to be 'Augustinian, but not Calvinistic' (p. 252), and identifies the baptized with the elect, saying (p. 254), 'God's predestination is bestowed on every baptized Christian. . . . The fact of God bringing men to baptism is synonymous with his choosing them in Christ out of mankind.'

John Wesley, unable to reconcile Art. XVII. with his Arminianism, omitted it altogether from his revision of the Articles.

BAPTISMAL REGENERATION AND FALL FROM GRACE.

The Articles teach also the possibility of falling away from grace (XVI.) and the doctrine of general baptismal regeneration (XXVII.). This seems to exclude an absolute decree of election 'to everlasting life,' which involves final perseverance as a necessary means to a certain end. Hence the attempts to explain away either the one or the other in order to save the logical consistency of the formulary.\footnote{Dr. Goode, in his learned work, \emph{The Doctrine of the Church of England as to the Effects of Baptism in the case of Infants} (1849), labors to show that inasmuch as the founders of the Church of England were Calvinists, they can not have held the Tractarian doctrine of baptismal regeneration, which is incompatible with Calvinism. Archdeacon Wilberforce, who afterwards seceded to Rome, showed, in his \emph{Doctrine of Holy Baptism} (London, 1849), in opposition to Goode, that the formularies of the Church of England do clearly teach baptismal regeneration. J. B. Mozley, B.D., Fellow of Magdalen College, Oxford, in his able work on \emph{The Primitive Doctrine of Baptismal Regeneration} (London, 1856), takes a middle ground, viz., that the Church of England imposes the doctrine 'that God gives regenerating grace to the whole body of the baptized,' and tolerates the doctrine 'that God gives grace sufficient for salvation only to some of this body' and 'that these two positions can not really be in collision with each other, though apparently they are.' Mozley grapples with the difficulties of the problem, but has after all not succeeded in making it clear.}

In Article XVI. there is no real difficulty. It is directed against the Anabaptists, who 'say they can no more sin,' and the modern Novatians, who 'deny the place of forgiveness to such as truly repent,' and accords with a similar article in the Augsburg Confession.\footnote{Comp. Augs. Conf., Art. XII.: '\emph{Damnant Anabaptistas, qui negant semel justificatos posse amittere Spiritum Sanctum. . . . Damnantur et Novatiani qui nolebant absolvere lapsos post baptismum redeuntes ad pœnitentiam.}' Also Bullinger's Confes. Helv., cap. XIV.: '\emph{Damnamus et veteres et novos Novatianos, atque Catharos.}'} It simply teaches the possibility of a temporary fall of the baptized and regenerated, but not a \emph{total} and \emph{final} fall of the elect, as is clear from the addition, 'and by the grace of God we may arise again and amend our lives.' This is quite consistent with Augustinianism, and even with the most rigorous form of Calvinism.\footnote{See the defense of this Article by Dean Bridges, of Sarum, quoted by Hardwick, p. 211.}

On the subject of baptism the Anglican Church agrees much more with the Lutheran than with the Calvinistic creed. She retained the Catholic doctrine of baptismal regeneration, but rejected the \emph{opus operatum} theory, and the doctrine that baptism destroys the nature of original sin as well as its guilt. Baptismal regeneration is taught indefinitely in Article XXVII.,\footnote{'Baptism is . . . a sign of regeneration or new birth, whereby, as by an instrument, they that receive baptism rightly, are grafted into the Church.' The language of this Article bears a Reformed or Calvinistic interpretation. Bishop Hooper and several of the Marian exiles were Zwinglians, but the views of Cranmer and Ridley, in their private writings, on the effects of baptism and baptismal grace, agree substantially with those of Luther. See Browne on \emph{Art. XXVII.} pp. 668 sq.; the passages collected by Jones, \emph{Expos. of the Art.} pp. 157 sqq.; also Hardwick, pp. 393–395.} more plainly in the Catechism,\footnote{The second question: 'Who gave you this name? \emph{Ans.} My godfather and godmother in baptism, wherein I was made a member of Christ, the child of God, and an inheritor of the kingdom of heaven.'} and in the baptismal service of the Liturgy, which pronounces every child after baptism to be regenerated.\footnote{After the public baptism of infants, the priest shall say: 'Seeing now, dearly beloved brethren, that this child is regenerate, and grafted into the body of Christ's Church, let us give thanks to Almighty God for these benefits,' etc. And in the prayer which follows: 'We yield thee hearty thanks, most merciful Father, that it hath pleased thee to regenerate this infant with thy Holy Spirit, to receive him for thine own child by adoption, and to incorporate him into thy holy Church.' The same prayer is prescribed for the office of private baptism of infants. The baptismal service is derived from the Sarum Manual and from the 'Consultation' of Archbishop Hermann of Cologne, which was borrowed from Luther's \emph{Taufbüchlein.} See Daniel, \emph{Cod. Liturg. Eccl. Luth.} p. 185, and Procter, \emph{History of the Book of Common Prayer,} p. 371, 11th ed. (1874). Among the eight particulars in the Prayer-Book, which Baxter and his Nonconformist brethren objected to as sinful, the fourth was 'that ministers be forced to pronounce all baptized infants to be regenerate by the Holy Ghost, whether they be the children of Christians or not' (Procter, p. 133). The last clause intimates that baptized children of Christian parents were regarded by them as regenerate.}

This doctrine seems to be contradicted by the undeniable fact that multitudes of baptized persons in all churches, especially in those where infant baptism is indiscriminately practiced, show no signs of a holy life or real change of heart, and belie their baptismal engagements.

To remove this difficulty, some Anglicans take the language of the baptismal service, not in a real and literal, but in a hypothetical or charitably presumptive meaning.\footnote{So Mozley, who endeavors to fasten this meaning upon the fathers, and the standard Anglican writers, including Hooker. But the strong language of the Greek and Latin fathers, who almost identify baptism with regeneration, and seem to know no other regeneration but that by baptism (which they call \textgreek{ἀναγέννησις, παλιγγενεσία, θεογένεσις, φωτισμός, } \emph{regeneratio, secunda nativitas, renascentia, illuminatio}), must be understood chiefly of adult baptism, which in the first four centuries of the Church was the rule, while infant baptism was the exception, and which was administered to such only as had passed through a course of catechetical instruction, and professed repentance and faith in Christ. The same is true of the passages of the New Testament on baptism.} Others make a distinction between baptismal or ecclesiastical regeneration (\emph{i.e.,} incorporation into the visible Church) and moral or spiritual regeneration (which includes renovation and conversion). Still others distinguish between the regenerate and the elect, and thus harmonize Art. XXVII. with Art. XVII. Augustine regards the elect as an inner circle of the baptized; and holds that, in addition to the baptismal grace of regeneration, the elect receive from God the gift of perseverance to the end, which puts into execution the eternal and unchangeable decree of election. The reason why God grants this grace to some and withholds it from others is unknown to us, and must be traced to his inscrutable wisdom. 'Both the grace of the beginning,' he says, 'and the grace of persevering to the end is not given according to our merits, but according to a most secret, just, wise, and beneficent will.' 'Wonderful indeed, very wonderful, that to some of his own sons, whom he has regenerated, and to whom he has given faith, hope, and charity, God does not give perseverance.'\footnote{\par See his tract \emph{ De dono perseverantiæ,} and Mozley's \emph{Treatise on the Augustinian Doctrine of Predestination} (Lond. 1855), pp. 191 sqq., and the \emph{Primitive Doctrine of Baptismal Regeneration,} pp. 113 sqq. Mozley thinks that Augustine means by baptismal regeneration only capacity for goodness and holiness. Browne (on Art. XXVII.) presents a somewhat different view, viz., that Augustine uses the term regeneration sometimes in a wider, sometimes in a stricter and deeper sense. 'At one time he speaks of all the baptized as regenerate in Christ, and made children of God by virtue of that sacrament; at another time he speaks of baptismal grace as rather enabling them to become, than as actually constituting them God's children; and says that, in the higher and stricter sense, persons are not to be called sons of God unless they have the grace of perseverance, and walk in the love of God' (p. 660). There is no doubt that Augustine wished to adhere to the traditional orthodox view of baptism, and yet he could not help seeing that his new doctrine of predestination required a modification, which, however, he did not fully and clearly carry out.}

Here is a point where Calvin differs from Augustine, at least in logic, although they agree in the result—namely, the non-salvation of the non-elect, whether baptized or not. Calvin likewise brings baptism into close connection with regeneration,\footnote{\par This is undoubtedly the case in the New Testament wherever Christian baptism is mentioned: John iii. 5; Acts ii. 38;  Rom. vi. 3, 4; Gal. iii. 27; Col. ii. 12; Eph. v. 26; Tit. iii. 5; 1 Pet. iii. 21. Calvin's exposition of some of these passages in his commentaries should be compared with his teaching in the 'Institutes.'} but he draws a sharper distinction between the outward visible sign and seal (Rom. iv. 11) and the inner invisible grace; he takes moreover a higher view of regeneration as a thorough moral renovation, and identifies the truly regenerate with the elect. He consequently restricts the regenerating efficacy of the Spirit to the elect, and makes it so far independent of the sacramental act that it need not always coincide with it, but may precede or follow the same. Thus the Westminster Confession calls baptism 'a sign and a seal of the covenant of grace, of his [the baptized person's] ingrafting into Christ, of regeneration, of remission of sins, and of his giving up unto God through Jesus Christ, to walk in newness of life.' But it adds that 'grace and salvation are not so inseparably annexed unto it [baptism], as that no person can be regenerated or saved without it (Rom. iv. 11;  Acts x. 2, 4, 22, 31, 45, 47); or, that all that are baptized are undoubtedly regenerated (Acts viii. 13, 23). The efficacy of baptism is not tied to that moment of time wherein it is administered (John iii. 8): yet, notwithstanding by the right use of this ordinance, the grace promised is not only offered, but really exhibited and conferred by the Holy Ghost to such (whether of age, or infants) as that grace belongeth unto, according to the counsel of God's own will, in his appointed time (Gal. iii. 27; Tit. iii. 5;  Eph. v. 25, 26;  Acts ii. 38, 41).'\footnote{Chap. xxviii. 1, 5, 6.}

The objection to the Calvinistic view is that it resolves the baptism of the non-elect into an empty ceremony (not to say solemn mockery); while the Augustinian view turns the baptismal regeneration of the non-elect into a failure. The former sacrifices the universality of baptismal grace to the particularism of election, the latter sacrifices the higher view of regeneration to the claims of baptism. The real difficulty of both theories lies in the logical incompatibility of a \emph{limited} election and a \emph{universal} baptismal grace. The predestinarian system and the sacramental system are two distinct lines of thought, which neither Augustine nor Calvin have been able satisfactorily to adjust and to harmonize.

NECESSITY OF BAPTISM.

As to the necessity of baptism for salvation, the Anglican Church at first followed, but afterwards softened the rigor of the Augustinian and Roman Catholic doctrine, which excludes even unbaptized infants dying in infancy from heaven, and assigns them to the \emph{ limbus infantum,} on the borders of hell. In the second of the Ten Articles of Henry VIII. (1536), it is asserted that 'infants and children dying in infancy shall undoubtedly be saved thereby [by baptism], \emph{and else not.}' In the first revision of the Liturgy, the introductory prayer that the child may be received by baptism into the ark of Christ's Church contains the exclusive clause 'and so saved from perishing.'\footnote{Borrowed from the Lutheran service composed by Melanchthon and Bucer for Cologne: 'That being separated from the number of the ungodly, he may be kept safe in the holy ark of thy Church (\emph{in sancta Ecclesiæ, tuæ Arca tutus servari possit}).' See Laurence, p. 71; Procter, p. 374. The Augsburg Confession (Art. IX., Latin ed.) teaches \emph{quod baptismus sit necessarius ad salutem}, and condemns the Anabaptists for teaching that infants may be saved without baptism.} But in the revision of 1552 this clause was omitted; for Cranmer, who framed the Liturgy, had in the mean time changed his opinion, as we may infer from the treatise upon the 'Reformation of Ecclesiastical Laws,' composed under his superintendency, where the 'scrupulous superstition' of the necessity of infant baptism for infant salvation is rejected.\footnote{\emph{Reformat. Leg., De Baptismo:} '\emph{Illorum etiam videri debet scrupulosa superstitio, qui Dei gratiam et Spiritum Sanctum tantopere cum sacramentorum elementis colligant, ut plane affirment, nullum Christianorum infantem salutem esse consecuturum, qui prius morte fuerit occupatus, quam ad Baptismum adduci potuerit; quod longe secus habere judicamus.}'} This change must be traced to the influence of Zwingli and Bullinger, who first boldly asserted that all infants dying before committing actual sin, whether baptized or not, whether of Christian or heathen parents, are saved in consequence of the universal merit of Christ ('\emph{propter remedium per Christum exhibitum}'), which holds good until rejected by unbelief.\footnote{See above, p. 378. Zwingli was not quite so positive about the salvation of heathen children, but he declared it at least '\emph{probabilius ut gentium liberi per Christum salventur quam ut damnentur.}' Bullinger held the same view, though not so clearly expressed. See the passages quoted by Laurence, pp. 266, 267, who agrees on this subject with the Zurich Reformers.} Calvin likewise taught the possibility of salvation without baptism, but confined it to the elect. Thomas Becon (chaplain to Cranmer, and one of the six preachers of Canterbury Cathedral, died 1567) is very explicit on this subject. As many Jewish children, he says, were saved without circumcision, so many Christian children, and even Turks and heathens, may be spiritually baptized and saved without water baptism. 'Besides all these things, what shall we say of God's election? Can the lack of outward baptism destroy and make of none effect the election of God; so that when God hath chosen to everlasting salvation, the want of an external sign shall cast down into everlasting damnation? . . . As many people are saved which never received the sacrament of the body and blood of Christ, so likewise are many saved though they were never outwardly baptized with water; forasmuch as the regeneration of the Christian consisteth rather in the spirit than in the flesh. This text, therefore, of Christ, ``Except a man be born of water,'' etc., is to be understood of such as may conveniently be baptized, and yet, notwithstanding, contemptuously refuse baptism, and despise the ordinance of Christ.'\footnote{Quoted by Jones, 1.c. pp. 167 sq.} Bishop Jewel says: 'The grace of God is not tied to any sacraments. He is able to work salvation both with them and without them.'\footnote{\par Ibid. p. 171.} Hooker is much more cautious and churchly. 'Predestination,' he says, 'bringeth not to life, without the grace of external vocation, wherein our baptism is implied, . . . which both declareth and maketh us Christians. In which respect we justly hold it to be the door of our actual entrance into God's house; the first apparent beginning of life; a seal, perhaps, to the grace of election, before received (Calvin, \emph{Instit.} iv. 15, 22), but to our sanctification here a step that hath not any before it. . . . If Christ himself which giveth salvation do require baptism (Mark xvi. 16), it is not for us that look for salvation to sound and examine him, whether unbaptized men may be saved, but seriously to do that which is required, and religiously to fear the danger which may grow by want thereof.' Yet, touching infants who die unbaptized, he inclines, at least in regard to the offspring of Christian parents, to a charitable presumption of 'the great likelihood of their salvation,' for the reasons that 'grace is not absolutely tied unto sacraments;' that 'God bindeth no man unto things altogether impossible;' that 'there is in their Christian parents, and in the Church of God, a presumed desire that the sacrament of baptism might be given them;' and that 'the seed of faithful parentage is holy from the very birth (1 Cor. vii. 14).'\footnote{\emph{Eccles. Polity,} Book V. ch. 60 (Vol. II. pp. 341, 342, 346, 347, Keble's ed.).}

The Anglican Church, then, as far as we may infer from her authoritative declarations, makes certain the salvation of all baptized infants dying in infancy, and leaves the possibility of salvation without baptism an open question, with a strong leaning towards the liberal view. The Roman Church makes infant salvation without baptism impossible; the Lutheran Church makes it at least improbable; the Calvinistic Churches make it certain in the case of all the elect, without regard to age, and decidedly incline to the charitable belief that all children dying in infancy belong to the number of the elect.

The doctrine of the absolute necessity of baptism for salvation has always been based upon two declarations of our Lord,  Mark xvi. 16, and  John iii. 5 (on the assumption that 'water' refers to baptism). But in the first passage our Lord, after declaring that faith followed by baptism saves, states the negative without adding, \emph{and is not baptized;} intimating by this omission, that only the want of faith or the refusal of the gospel, not the want of baptism, condemns. In the discourse with Nicodemus he does not say that water baptism is regeneration, nor that every one that is born of water is also born of the Spirit (which was certainly not the case with Simon Magus, who, notwithstanding his baptism, remained 'in the gall of bitterness and the bond of iniquity'); he simply lays down two conditions for entering into the kingdom of God, and puts the emphasis on being born of the Spirit. This is evident from the fact that in that discourse 'water' is mentioned but once, but the Spirit four times. The most that can be inferred from the two passages is the ordinary necessity of baptism where it can be had—that is, within the limits of the Christian Church. We are bound to God's ordinances, but God's Spirit is free and 'bloweth where it listeth.' We should never forget that the same Lord was the special friend of children, and declared them to belong to the kingdom of heaven, without any reference to baptism or circumcision, adding these significant words, 'It is not the will of your Father who is in heaven that one of these little ones should perish' (Matt. xviii. 14).

THE LORD'S SUPPER.

If the Articles on Predestination and Baptism leave room for different interpretations, there can be no reasonable doubt about the meaning of Art. XXVIII. on the Lord's Supper. It clearly teaches the Reformed doctrine of the \emph{spiritual} presence and \emph{spiritual} eating by \emph{faith} only, in opposition both to transubstantiation and consubstantiation, which imply a \emph{corporal} presence and an \emph{oral} manducation by \emph{all} communicants, both good and bad, although with opposite effects.

The wide departure from the Lutheran formularies, otherwise so freely consulted, may be seen from the following comparison:

Augsburg Confession. 1530. Art. X.

Thirteen Articles. 1538. Art. VII.

Thirty-nine Articles. 1563 and 1571. Art. XXVIII.

\emph{De cœna Domini docent, quod corpus et sanguis Christi } vere adsint,  \emph{et distribuantur }  vescentibus  \emph{in cæna Domini; et improbant secus docentes.}

\emph{De Eucharistia constanter credimus et docemus, quod in sacramento corporis et sanguinis Domini } vere, substantialiter,\footnote{The term \emph{substantialiter} is borrowed from the Apology of the Augsburg Conf., Art. X.}  et realiter adsint  \emph{corpus et sanguis Christi } sub speciebus panis et vini.\footnote{\emph{Sub speciebus panis et vini,} from the German edition of the Augsburg Conf. (\emph{unter Gestalt des Brotes und Weines}).} \emph{Et quod sub ejusdem speciebus vere et realiter exhibentur et } distribuuntur  \emph{illis qui sacramentum accipiunt, sive bonis } sive malis.

\emph{Corpus Christi datur, accipitur, et manducatur } tantum cœlesti et spirituali ratione (only after an heavenly and spiritual manner). \emph{Medium, autem quo Corpus Christi accipitur et manducatur in cœna, } fides  \emph{est} (and the mean whereby the body of Christ is received and eaten in the Supper, is faith).

The clause here quoted from the Elizabethan revision was wanting in the Edwardine Articles, and was inserted on motion of Bishop Guest of Rochester.\footnote{This is inferred from a letter to Cecil, Dec. 22, 1566, where Guest justifies the use of the word 'only' by saying that he did not intend to exclude 'the presence of Christ's body from the sacrament, but only the grossness and sensibleness in the receiving thereof.' Hardwick, p. 130.} Both series contain the assertion that the bread which we break is a communion of the body of Christ 'to such as \emph{rightly, worthily,} and \emph{with faith} receive the same,' which was meant to exclude the oral manducation. Both strongly condemn transubstantiation. The Edwardine Articles protest also against the Lutheran hypothesis of the ubiquity of Christ's body.\footnote{\par 'Forasmuch as the truth of man's nature requireth that the body of one and the self-same man can not be at one time in diverse places, but must needs be in some one certain place: therefore the body of Christ can not be present at one time in many and diverse places. And because (as holy Scripture doth teach) Christ was taken up into heaven, and there shall continue unto the end of the world, a faithful man ought not either to believe or openly to confess the real and bodily presence (as they term it) of Christ's flesh and blood, in the sacrament of the Lord's Supper.'} This same protest against ubiquity is found substantially in the Parker MS. of the Latin revision of 1563, but it was struck out in the Convocation.\footnote{Hardwick regards this omission as a protest against Zwinglianism. But the leading Elizabethan bishops, especially Horn, Jewel, and Grindal, assure Bullinger and Peter Martyr of their full agreement with them against the ubiquitarian hypothesis, which was at that time defended by Brentius and Andreae, and opposed by the Swiss. See pp. 603 and 632.} Instead of it a new Article was added in the English revision of 1571, denying that the unworthy partake of Christ in the communion.\footnote{\par Art. XXIX. 'Of the wicked which do not eat the body of Christ in the use of the Lord's Supper. The wicked, and such as be void of a lively faith, although they do carnally and visibly press with their teeth (as St. Augustine saith) the sacrament [i.e., the sacramental \emph{sign}] of the body and blood of Christ: yet in no way are they partakers of Christ, but rather to their condemnation do eat and drink the sign or sacrament of so great a thing.' This Article is wanting in the Latin edition of 1563, having probably been withdrawn from the Convocation records in compliance with the desire of the Queen and her council to deal gently with the adherents of the 'old learning' (whether Romish or Lutheran); but it was inserted in the Latin editions after the year 1571. See Hardwick, pp. 144 and 315.}

The Catechism likewise limits the reception of Christ's body and blood to the 'faithful,' and declares the benefit of the Lord's Supper to be 'the strengthening and refreshing of our \emph{souls.}' The communion service does not rise above this view, and the distribution formula, inserted in the revision of 1552, expresses the commemorative theory. The rubric on kneeling, at the close of the service, which was inserted in the second Prayer-Book of Edward VI. (1552) by Cranmer, through the influence of Hooper and Knox (one of the royal chaplains),\footnote{See the lengthy discussion of this subject in Lorimer's \emph{John Knox,} pp. 100–136.} then omitted in Elizabeth's reign from regard to the Catholics, but which was again restored in the Reign of Charles II. (1662) to conciliate the Puritans, explains the kneeling at the communion not to mean an adoration of the sacramental bread and wine, or any corporal presence of Christ's natural flesh and blood. 'For the natural body and blood of Christ are in heaven, and not here; it being against the truth of Christ's natural body to be at one time in more places than one.' This is a plain declaration against consubstantiation and ubiquity.

Before the Articles were framed a public disputation on the eucharistic presence was held before the royal commissioners at the University of Oxford, May, 1549, in which Peter Martyr, then professor of theology, defended the figurative interpretation of the words, 'This \emph{is} my body,' and the commemorative character of the ordinance. The acts of the disputation were published by Cranmer, with a preface and discourse of Peter Martyr.\footnote{\emph{Tractatio de sacramento Eucharistiæ habita in celeberrima Universitate Oxoniensi. Ad hæc: Disputatio de eodem sacramento in eadem Universitate habita.} London, 1549; also in Zurich, 1552, and an English translation, 1583. See an account in Dr. C. Schmidt, \emph{Peter Martyr Vermigli, Leben und ausgewählte Schriften} (Elberfeld, 1858), pp. 91–100, 105.} In June of the same year a disputation on the same subject, in which Bucer took part, was held in the University of Cambridge.\footnote{\par Schmidt, p. 106. Ridley's \emph{Works,} pp. 171 sqq.}

Cranmer, after holding first to transubstantiation, then to consubstantiation, adopted at last the Calvinistic theory of a spiritual real presence and a spiritual reception by faith only, and embodied it in the Articles and the second revision of the Liturgy.\footnote{See above, p. 601. Cranmer admits the charge of his opponents, Bishop Gardiner and Dr. Smith, that he was upon this point first a Papist, then a Lutheran, and at last a Zwinglian. 'After it hath pleased God,' he says, 'to show unto me, by his holy Word, a more perfect knowledge of his Son Jesus Christ, from time to time as I grew in knowledge of him, by little and little I put away my former ignorance. And as God of his mercy gave me light, so through his grace I opened mine eyes to receive it, and did not willfully repugn unto God and remain in darkness. And I trust in God's mercy and pardon for my former errors, because I erred but of frailness and ignorance.' Answer to Smith's Preface, \emph{Works,} Vol. I. p. 374.} He openly confessed this change at a public disputation held in London, Dec. 14, 1548, in the Parliament house, 'in the presence of almost all the nobility of England.'\footnote{Of this recantation Bartholomew Traheron wrote to Bullinger from London, Dec. 31, 1548, as follows: 'I can not refrain, my excellent Bullinger, from acquainting you with circumstances that have lately given us the greatest pleasure, that you and your fellow-ministers may participate in our enjoyment. On the 14th of December, if I mistake not, a disputation was held at London concerning the eucharist, in the presence of almost all the nobility of England. The argument was sharply contested by the Bishops. The Archbishop of Canterbury, contrary to general expectation, most openly, firmly, and learnedly maintained your opinion upon this subject. His arguments were as follows: The body of Christ was taken up from us into heaven. Christ has left the world. ``Ye have the poor always with you, but me ye have not always,'' etc. Next followed the Bishop of Rochester [Ridley], who handled the subject with so much eloquence, perspicuity, erudition, and power, as to stop the mouth of that most zealous papist, the Bishop of Worcester [Heath]. The truth never obtained a more brilliant victory among us. I perceive that it is all over with Lutheranism, now that those who were considered its principal and almost only supporters have altogether come over to our side. We are much indebted to the Lord who provides for us also in this particular.' In a postscript to this letter, John of Ulmis adds: 'The foolish Bishops have made a marvelous recantation.' The same 'notable disputation of the sacrament' is mentioned in King Edward's Journal as having taken place in the Parliament house. See \emph{Zurich Letters,} 1537–1558, pp. 322, 323.} He wrote an elaborate exposition and defense of his final view against the attacks of Gardiner.\footnote{\par \emph{An Answer unto a Crafty and Sophistical Cavillation, devised by Stephen Gardiner, Doctor of Law, late Bishop of Winchester, against the True and Godly Doctrine of the most holy Sacrament of the Body and Blood of our Saviour Jesus Christ} (1550). The sacramental writings of Cranmer fill the first volume of the Parker Society's edition of his works (Cambridge, 1844).} He does not allude to Calvin's writings on the eucharist, although he can hardly have been ignorant of them, but quotes largely from Augustine, Tertullian, Origen, Theodoret, and other fathers who seem to favor a figurative interpretation, and approvingly mentions Bertram, Berengarius, and Wycliff among mediæval divines, and Bucer, Peter Martyr, Zwingli, Œcolampadius among the Reformers, as teaching substantially the same doctrine.\footnote{\emph{Works,} Vol. I. pp. 14, 173, 196, 225, 374.} He also expressed his unqualified approbation of Bullinger's 'Tract on the Sacraments,' which was by his desire republished in England (1551) by John à Lasco, to whom he remarked that 'nothing of Bullinger's required to be read and examined previously.'\footnote{See a letter of John à Lasco to Bullinger, dated London, April 10, 1551; Cardwell's \emph{Liturgies of Edward VI.} (Preface), and Lorimer's \emph{John Knox,} p. 49.} But he traced his change directly to Bishop Ridley,\footnote{See a letter of John à Lasco to Bullinger, dated London, April 10, 1551; Cardwell's \emph{Liturgies of Edward VI.} (Preface), and Lorimer's \emph{John Knox,} p. 49.} and Ridley derived his view not so much from Swiss sources as from Bertram (Ratramnus), who, in the middle of the ninth century, wrote with great ability against the magical transubstantiation theory of Paschasius Radbertus, and in favor of a spiritual and dynamic presence.\footnote{Bishop Browne correctly says (p. 710): 'Ridley, indeed, refused to take the credit of converting Cranmer, but Cranmer himself always acknowledged his obligations to Ridley.' In his last examination at Oxford, before Bishop Brooks of Gloucester (Sept., 1555), Cranmer said that 'Doctor Ridley, by sundry persuasions and authorities drew me quite from my opinion' (on the real presence). \emph{Works,} Vol. II. p. 218. Brooks on the same occasion remarked: 'Latimer leaneth to Cranmer. Cranmer to Ridley, and Ridley to the singularity of his own wit;' to which Ridley replied, that this was 'most untrue, in that he was but a young scholar in comparison of Master Cranmer.' Ridley's \emph{Works,} pp. 283, 284.} Cranmer's last utterances on this subject, shortly before his condemnation and martyrdom, were made in the Oxford disputations with the Romanists to which he, with Ridley and Latimer, was summoned from prison, April (and again in September), 1555. He declared there that Christ's 'true body is truly present to them that truly receive him, but spiritually. And so it is taken after a spiritual sort. . . . If ye understand by this word ``\emph{really},'' \emph{re ipsa,} i.e., \emph{in very deed and effectually,} so Christ, by the grace and efficacy of his passion, is in deed and truly present to all his true and holy members. But if ye understand by this word ``\emph{really}'' \emph{corporaliter,} i.e., \emph{corporally,} so that by the body of Christ is understanded a natural body and organical, so the first proposition doth vary, not only from usual speech and phrase of Scripture, but also is clean contrary to the holy Word of God and Christian profession: when as both the Scripture doth testify by these words, and also the Catholic Church hath professed from the beginning, Christ to have left the world, and to sit at the right hand of the Father till he come unto judgment.'\footnote{\par \emph{Works,} Vol. I. pp. 394, 395.}

We add the last confessions of the other two English Reformers at their examination in Oxford.

Bishop Latimer declared 'that there is none other presence of Christ required than a spiritual presence; and this presence is sufficient for Christian man, as the presence by the which we both abide in Christ, and Christ in us to the obtaining of eternal life, if we persevere in his true gospel.'\footnote{Jones, l.c. p. 176, where also the passages of the leading divines and bishops of the Elizabethan age on the subject of the Lord's Supper are collected.}

Bishop Ridley said: 'I worship Christ in the sacrament, but not because he is included in the sacrament: like as I worship Christ also in the Scriptures, not because he is really included in them. . . . The body of Christ is present in the sacrament, but yet sacramentally and spiritually (according to his grace) giving life, and in that respect really, that is, according to his benediction, giving life.The true Church of Christ doth acknowledge a presence of Christ's body in the Lord's Supper to be communicated to the godly by grace, and spiritually, as I have often showed, and by a sacramental signification, but not by the corporal presence of the body of his flesh.\footnote{Ridley's \emph{Works,} pp. 235 sq. Jewel expresses the same views very fully in his controversy with Harding, \emph{Works,} Vol. I. pp. 448 sqq. (Parker Soc. ed. 1845). Bishop Browne (p. 715) says that all the great luminaries of the Church of England (naming Mede, Andrewes, Hooker, Taylor, Hammond, Cosin, Bramhall, Ussher, Pearson, Patrick, Bull, Beveridge, Wake, Waterland) agree with the doctrine of the formularies in denying a corporal and acknowledging a spiritual feeding in the Supper of the Lord.}

REVISION OF THE ARTICLES.

The Thirty-nine Articles have remained unchanged in England since the reign of Elizabeth. The objections of Nonconformists to some of the Articles (XXIV., XXV., the affirmative clause of XX., and a portion of XXVII) have been removed since 1688 by relaxation and exemption; and the difficulties arising from the development of theological schools with widely divergent tendencies, within the bosom of the Church of England itself, have been met by liberal decisions allowing a great latitude of interpretation.

During the reign of William III., in 1689, a thorough revision of the Book of Common Prayer was undertaken and actually made in the interest of an agreement with Protestant Dissenters, by an able royal commission of ten bishops and twenty divines, including the well-known names of Stillingfleet, Patrick, Tillotson, Sharp, Hall, Beveridge, and Tenison. But the revision has never been acted upon, and was superseded by the toleration granted to Dissenters. The alterations did not extend to the Articles directly, but embraced some doctrinal features in the liturgical services—namely, the change of the word \emph{Priest} to 'Presbyter' or 'Minister;' \emph{Sunday} to 'Lord's Day;' the omission of the \emph{Apocryphal Lessons} in the calendar of Saints' days, for which chapters from Proverbs and Ecclesiastes were substituted, a concession to conscientious scruples against kneeling in receiving the sacrament, and an addition to the rubric before the Athanasian Creed, stating that 'the \emph{condemning clauses} are to be understood as relating only to those who obstinately deny the substance of the Christian faith.\footnote{A revision of the Book of Common Prayer was adopted by the National Church Assembly, July, 1927, the vote being 34 to 4 bishops, 255 to 37 clergymen, 230 to 92 laymen, but rejected by the House of Commons, Dec., 1927, by a vote of 238 to 205. A second revision was rejected by the Commons, June 14, 1928, by an increased majority, 266 to 220. The revision seemed to permit the reservation of the sacrament and introduced after the consecration of the elements the \emph{epiclesis} of the Greek Church, stating the change of the bread and wine. The Revised Book is issued by the S. P. C. K.—Ed.}

\xsubsection{American Revision of the Thirty-nine Articles, A.D. 1801.}

§ 82. American Revision of the Thirty-nine Articles, A.D. 1801.

William White, D.D. (first Bishop of the Protestant Episcopal Church in the diocese of Pennsylvania; d. 1836): \emph{Memoirs of the Protestant Episcopal Church in the United States of America.} New York, 1820; 3d ed. by De Costa, 1880.

William Stevens Perry, D.D. (Secretary of the House of Clerical and Lay Deputies of the General Convention of the Protestant Episcopal Church in the United States): \emph{A Hand-book of the General Convention of the Protestant Episcopal Church, giving its History and Constitution,} 1785—1874. New York, 1874. The same: \emph{Journals of the General Convention,} etc., 1785—1835. Claremont, N. H., 1874.

Also Samuel Wilberforce (late Bishop of Oxford): \emph{A History of the Protestant Episcopal Church in America} (1844); Caswall: \emph{History of the American Church} (2d ed. 1851); and Procter: \emph{A History of the Book of Common Prayer,} pp. 162 sqq. (11th ed. 1874).

For the colonial history, comp. the \emph{Historical Collections relating to the American Colonial Church,} ed. by Dr. Perry. Hartford, 1871 sqq. 3 vols. 4to.

The members of the Church of England in the American Colonies, from the first settlement of Virginia (1607) till after the War of the Revolution, belonged to the diocese of the Bishop of London, who never visited the country, and could exercise but an imperfect supervision. Several attempts were made, by the friends of the Church, to establish colonial bishoprics, but failed.

The separation from the crown of England necessitated an independent organization, which assumed the title of The Protestant Episcopal Church in the United States of America. The first steps towards such an organization were taken by a meeting of clergy and laity in New Brunswick, New Jersey, May 11, 1784, and by another and larger one, held in New York, Oct. 6 and 7, of the same year. The first General Convention, consisting of sixteen clerical and twenty-six lay deputies, assembled in Philadelphia, Sept. 27 and 28, 1785, Dr. White presiding, adopted a constitution and such changes in the Book of Common Prayer as were deemed necessary to conform it 'to the American Revolution and the Constitutions of the respective States,' and petitioned the English hierarchy to consecrate such bishops for the independent Church as may be elected by the separate dioceses.\footnote{Shortly before the Convention, Bishop Seabury, of Connecticut, had received consecration at Aberdeen from three Bishops of Scotland (Nov. 14, 1784), but he did not attend the Convention, and was opposed from High-Church principles to the introduction of lay representation and the limitation of the power of the episcopate.} The revised provisional Liturgy was rather hastily prepared and published, 1786. It is called the 'Proposed Book.'\footnote{It is sometimes also called 'Bishop White's Prayer-Book,' who was the chairman of the committee of revision, Dr. William Smith, of Maryland, and Dr. Wharton, of Delaware, being the other members. Smith is made chiefly responsible for the changes by Perry, p. 23. The book was printed in Philadelphia, 1786, in London, 1789, and again (with omission of the amended Articles of Religion) in New York, Dec., 1873, for provisional use in the new Reformed Episcopal Church,' which has since adopted a new revision.} It contains, besides many necessary ritual changes and improvements, Twenty Articles of Religion, based upon the Thirty-nine Articles, but differing widely from them, being a mutilation rather than an improvement.\footnote{Given by Perry, \emph{Hand-book,} pp. 34–39, from the original MSS. in the Convention archives. He calls the Proposed Book a 'hasty, crude, and unsatisfactory compilation, which failed utterly to establish itself in the American Church' (p. 42).} The alterations and omissions were made in the interest of an unchurchly latitudinarianism which then prevailed. The Nicene Creed and the Athanasian Creed, which Art. VIII. of the English series acknowledges, were entirely omitted in Art. IV. of the new series; the Apostles' Creed was retained, but without the clause 'He descended into hell.'

The book failed to give general satisfaction at home or abroad. The English Archbishops demanded the restoration of the three Œcumenical Creeds in their integrity.\footnote{See their letter in Perry, pp. 50–55.}

The General Convention held at Wilmington, Del., Oct. 11, 1786, complied with this request so far as the Nicene Creed and the discretionary use of the clause of the descent in the Apostles' Creed were concerned.\footnote{In the first edition of the new Prayer-Book, 1790, the objectionable clause was printed in italics, and put in parentheses. But the General Convention of 1792 left it discretionary to use it, or to omit it, or to substitute for it the words, '\emph{He went into the place of departed spirits},' as being equivalent to the words in the Creed.} The omission of the Athanasian Creed was adhered to,\footnote{Bishop Seabury was very zealous for the Athanasian Creed; and in the Convention of 1789 the House of Bishops agreed to its permissory use, but the House of Deputies 'would not allow of the Creed in any shape.' Bishop White favored a compromise—viz., to leave it in the Prayer-Book as a doctrinal document, but not to read it in public worship. See his \emph{Memoirs,} pp. 149, 150, and a letter of White, quoted by Perry, p. 76.} and subsequently acquiesced in by the English Bishops. The obstacle of the oath of allegiance required in England having been removed by act of Parliament, the Rev. Drs. White, of Pennsylvania, and Provoost, of New York, received the long-sought 'Apostolical succession,' in the chapel of Lambeth Palace, Feb. 4, 1787. At one time this result seemed so doubtful that steps were taken to secure ordination, with a broken succession, from the Lutheran bishops of Denmark, and the consent of the Danish government had actually been obtained, when the difficulties in England were removed.

In the Special Convention of Philadelphia, June, 1799 (the General Convention having been prevented in the preceding year by an epidemic), a new revision of the Articles of Religion, reduced to seventeen, was considered, but not finally acted upon by the House of Deputies, and was printed as an Appendix to the Journal of that House.\footnote{Perry, pp. 90–95.} But it gave no satisfaction, and shared the same fate with the previous draft of twenty Articles.

Finally, the General Convention held at Trenton, New Jersey, Sept. 3–12, 1801, settled the question by adopting the Thirty-nine Articles in the form which they have since retained in the American Episcopal Church, and are incorporated in its editions of the Prayer-Book.\footnote{See Vol. III. pp. 477 sqq., where they are given in full.} The only doctrinal difference is the omission of the Athanasian Creed from Art. VIII.; the remaining changes are political, and adapted to the separation of Church and State. Otherwise even 'the obsolete diction' is retained. The following is the action of this Convention:\footnote{Perry, pp. 99–101.}

'Resolutions of the Bishops, the Clergy, and the Laity of the Protestant Episcopal Church in the United States of America, in Convention, in the city of Trenton, the 12th day of September, in the year of our Lord 1801, respecting the Articles of Religion.

'The Articles of Religion are hereby ordered to be set forth with the following directions, to be observed in all future editions of the same; that is to say—

'The following to be the title, viz.:

'''Articles of Religion, as established by the Bishops, the Clergy, and the Laity of the Protestant Episcopal Church in the United States of America, in Convention, on the 12th day of September, in the year of our Lord 1801.''

'The Articles to stand as in the Book of Common Prayer of the Church of England, with the following alterations and omissions, viz.:

'In the 8th Article, the word ``three'' in the title, and the words ``three—Athanasius' creed'' in the Article, to be omitted, and the Article to read thus:

'''Art. VIII. Of the Creeds.

'''The Nicene Creed, and that which is commonly called the Apostles' Creed, ought thoroughly to be received and believed, for they may be proved by most certain warrants of Holy Scripture.''

'Under the title ``Article 21,'' the following note to be inserted, namely,

'''The 21st of the former Articles is omitted, because it is partly of a local and civil nature, and is provided for, as to the remaining parts of it, in other Articles.''

'The 35th Article to be inserted with the following note, namely,

'''This Article is received in this Church, so far as it declares the Books of Homilies to be an explication of Christian doctrine, and instructive in piety and morals. But all references to the constitution and laws of England are considered as inapplicable to the circumstances of this Church: which also suspends the order for the reading of said homilies in churches until a revision of them may conveniently be made, for the clearing of them, as well from obsolete words and phrases, as from the local references.''

'The 36th Article, entitled ``Of Consecration of Bishops and Ministers,'' to read thus:

'''The Book of Consecration of Bishops, and ordering of Priests and Deacons, as set forth by the General Convention of this Church in 1792, doth contain all things necessary to such consecration and ordering: neither hath it any thing that, of itself, is superstitious and ungodly. And, therefore, whosoever are consecrated or ordered according to said form, we decree all such to be rightly, orderly, and lawfully consecrated and ordered.''

'The 37th Article to be omitted, and the following substituted in its place:

'''Art. XXXVII. Of the Power of the Civil Magistrate.

'''The power of the civil magistrate extendeth to all men, as well Clergy as Laity, in all things temporal; but hath no authority in things purely spiritual. And we hold it to be the duty of all men who are professors of the gospel, to pay respectful obedience to the civil authority, regularly and legitimately constituted.''\footnote{This Art. appears as the last in the XVII. Articles of 1799.}

'Adopted by the House of Bishops.

WILLIAM WHITE, D.D., Presiding Bishop.

'Adopted by the House of Clerical and Lay Deputies.

ABRAHAM BEACH, D.D., President.

On the nature and aim of this action, Bishop White remarks:\footnote{\emph{Memoirs,} p. 33.}

'The object kept in view, in all the consultations held, and the determinations formed, was the perpetuating of the Episcopal Church, on the ground of the general principles which she had inherited from the Church of England; and of not departing from them, except so far as either local circumstances required, or some very important cause rendered proper. To those acquainted with the system of the Church of England, it must be evident that the object here stated was accomplished on the ratification of the Articles.'

The only change in the Prayer-Book which has a doctrinal bearing, besides the omission of the Athanasian Creed, is the insertion of the Prayer of Oblation and Invocation from the Scotch (and the First Edwardine) Prayer-Book, through the influence of Bishop Seabury, who had been consecrated in the Scotch Episcopal Church.

\xsubsection{The Catechisms of the Church of England. A.D. 1549 and 1662.}

§ 83. The Catechisms of the Church of England. A.D. 1549 and 1662.

Literature.

The Church Catechism is contained in all the English and American editions of the Book of Common Prayer, between the baptismal and the confirmation services, and is printed in this work with the American emendations, Vol. III. pp. 517 sqq. The authentic text of the final revision of 1662 is in the corrected copy of the \emph{Black-Letter Prayer-Book,} which was attached to the Act of Uniformity, and has been republished in fac-simile, Lond. 1871. (It was supposed to be lost, when in 1867 it was discovered in the Library of the House of Lords.)

Archibald John Stephens:  \emph{The Book of Common Prayer, with notes legal and historical.} Lond. 1854 Vol. III. pp. 1449–1477.

Francis Procter:  \emph{A History of the Book of Common Prayer,} 11th ed. Lond. 1874, ch. V. sect. i (pp. 397 sqq.).

See other works on the Anglican Liturgy, noticed by Procter, p. viii.

EARLIER CATECHISMS.

The English Church followed the example of the Lutheran and Reformed Churches on the Continent in providing for regular catechetical instruction. English versions and expositions of the Lord's Prayer, the Creed, and the Ten Commandments, with some prayers, were known before the Reformation, and constituted the 'Prymer,' which is commonly mentioned in the fifteenth century as a well-known book of private devotion.\footnote{The earliest known copy, belonging to the latter part of the 14th century, has been published by Maskel in \emph{Monumenta ritualia Ecclesiæ Anglicanæ,} Vol. II. It contains Matins and Hours of our Lady; Evensong and Compline; the seven Penitential Psalms; the Psalmi graduales (Psa. CXX.–CXXXIV.); the Litany; Placebo (Vespers); Dirge (the office for the departed); the Psalms of Commendation; Pater noster; Ave Maria; Creed; Ten Commandments; the seven deadly sins. See Procter, p. 15.} In 1545 Henry VIII. set forth a Primer which was 'to be taught, learned, and read, and none other to be used throughout all his dominions.\footnote{It contained, besides the contents of the older Primers, the Salutation of the Angel, the Passion of our Lord, and several prayers. See Procter, p. 15, and Barton, \emph{Three Primers,} pp. 437 sqq.} During his reign the curates were frequently enjoined to teach the people the Lord's Prayer, the Creed, and the Ten Commandments, sentence by sentence, on Sundays and Holydays, and to make all persons recite them when they came to confession.

CRANMER'S CATECHISM.

'Cranmer's Catechism,' which appeared with his sanction in 1548, was for the most part a translation of the Latin Catechism of Justus Jonas, and retains the Catholic and Lutheran consolidation of the first and second commandments, and the sacrament of penance or absolution; but it was soon superseded.\footnote{So Hardwick (\emph{Hist. of the Reform.} p. 194) and other Episcopal writers. This matter needs further investigation. The very existence of a Catechism of Jonas is doubted by Langemack and Mönckeberg, who have written with authority on Luther's Catechism. But it is a fact that Luther, before be prepared his own Catechisms (1529), charged with this task his colleagues and friends Justus Jonas and Agricola of Eisleben (who afterwards became the leader of Antinomian views in opposition to Luther), for he wrote to Hausmann, Feb. 2, 1525: '\emph{Jonæ et Eislebio mandatas est catechismus puerorum parandus}' (De Wette, Vol. II. p. 621). This is probably the Catechism which appeared in the same year in a Latin translation anonymously under the title '\emph{Quo pacto statim a primis annis, pueri debeant in Christianismo institui. Libellus perutilis.}' At the close: '\emph{Impressum Wittembergæ per Georgium Rhaw. An.} 1525.' The original German edition has not been traced, but Dr. Schneider has discovered a copy of an improved German edition, under the title '\emph{Ein Buchlein fur die kinder gebessert und gemehret. Der Leyen Biblia. Wittemberg,} 1528,' and has reproduced it in the appendix to his critical edition of Luther's Small Catechism, 1853. He leaves it, however, uncertain whether it was composed by Jonas. Comp. his Introduction, pp. xx sqq. It consists of a brief exposition of the Ten Commandments, the Creed, the Lord's Prayer, the Sacrament of Baptism, and the Lord's Supper, with an addition on Confession; and so far it anticipates the order of Luther's Catechism. This must be the basis of Cranmer's Catechism; but as the Parker Soc. edition of his works gives only his dedicatory Preface to King Edward (Vol. II. p. 418), I can not verify the identity. It seems strange that Cranmer did not translate rather the far more perfect Catechism of Luther. The reason was, no doubt, his personal acquaintance with the author's son, Justus Jonas, jun., who was recommended to him by Melanchthon, was very kindly treated by him, and seems to have been the chief medium of his communication with the German Lutherans. See Strype's \emph{Memoir of Cranmer,} Vol. II. p. 581; Laurence, p. 17; and Cranmer's \emph{Works,} Vol. II. p. 425.} Cranmer changed about that time his view of the real presence.

THE CATECHISM OF THE PRAYER-BOOK.

When the Reformation was positively introduced under Edward VI., and the Book of Public Worship was prepared, a Catechism was embodied in it, to insure general instruction in the elements of the Christian religion. In the Prayer-Books of Edward VI. (1549, 1552) and Elizabeth (1559) this Catechism bears the title 'Confirmation, wherein is contained a Catechism for Children.'

This work has undergone, with other parts of the Prayer-Book, sundry alterations. The commandments were given, first very briefly (as in King Henry's Primer), then in full with a Preface in the edition of 1552. The explanation of the sacraments was added in 1604 by royal authority, in compliance with the wish of the Puritans expressed at the Hampton Court Conference,\footnote{Dr. Reynolds said at that Conference: 'The Catechism in the Common Prayer-Book is too brief, and that of Mr. Nowell (late Dean of St. Paul's) too long for novices to learn by heart. I request, therefore, that one uniform Catechism may be made, and none other generally received.' To this King James replied: 'I think the doctor's request very reasonable, yet so that the Catechism may be made in the fewest and plainest affirmative terms that may be,—not like the many ignorant Catechisms in Scotland, set out by every one who was the \emph{son of a good man.}'—Fuller's \emph{Church History of Britain,} Vol. V. p. 284.} and is attributed to Bishop Overall, then Dean of St. Paul's. In the last revision of the Prayer-Book, in 1661, the title was changed into 'A Catechism,' and two emendations were introduced in the answer on Baptism, as follows:

Earlier Editions.

Edition of 1661 (1662)

What is the outward visible sign or form in Baptism?

What is the outward visible sign or form in Baptism?

Water; wherein the person baptized \emph{is dipped or sprinkled with it,} in the name, etc.

Water; wherein the person is \emph{baptized,} in the name, etc.

Why then are infants baptized when by reason of their tender age they can not perform them [repentance and faith]?

Why then are infants baptized, when by reason of their tender age they can not perform them?

\emph{Yes; they do perform them by their Sureties, who promise and vow them both in their names: which} when they come to age themselves are bound to perform.

\emph{Because they promise them both by their Sureties; which promise,} when they come to age, themselves are bound to perform.

In the explanation of the Commandments the words 'the King \emph{and his Ministers}' were so changed as to read 'the King \emph{and all that are put in authority under him.}'

This Catechism is a considerable improvement on the mediæval primers, but very meagre if we compare it with the Catechisms of Luther, Calvin, and other Continental Reformers.

The Nonconformist ministers at the Savoy Conference (April, 1661), in reviewing the whole Liturgy, objected to the first three questions of the Catechism, and desired a full exposition of the Lord's Prayer, the Creed, and the Commandments, and additional questions on the nature of faith, repentance, the two covenants, justification, adoption, regeneration, and sanctification. These censures were not heeded.\footnote{Dr. Shields, in his edition of the \emph{Book of Common Prayer as amended by the Savoy Conference} (Phila. 1867), has inserted the Shorter Westminster Catechism in the place of the Anglican Catechism. But it does not harmonize with the genius of the Prayer-Book.}

The American Episcopal Church adopted, with the body of the Book of Common Prayer, the Catechism also, substituting 'the civil authority' for 'the King,' and omitting several directions in the appended rubrics.

Outside of the Anglican communion the Catechism is used only by the Irvingites who more nearly approach that Church, especially in their liturgy, than any other.

LARGER CATECHISMS.

The need of a fuller Catechism for a more advanced age was felt in the Church of England. Such a one was prepared by Poynet, Bishop of Winchester, and published, together with the Forty-two Articles, in Latin and English, in 1553,\footnote{Both editions are reprinted by the Parker Society in \emph{Liturgies,} etc., \emph{of Edward VI.}} apparently with the approval of Cranmer and the Convocation.\footnote{'\emph{Catechismus brevis Christianæ disciplinæ summam continens:}' '\emph{A short Catechism, or plain instruction, containing the sum of Christian learning, set forth by the King's Majesty's authority, for all schoolmasters to teach.}' The authority of this Catechism was afterwards disputed. See Hardwick, \emph{Hist. of the Articles,} p. 109.} On the basis of this, Dean Nowell, of St. Paul's, prepared another in 1562, which was amended, but not formally approved by Convocation (Nov. 11, 1562), and published (1570) in several forms—larger, middle, and smaller. The smaller differs but slightly from that in the Prayer-Book.\footnote{The larger Catechism appeared first in Latin under the title '\emph{Catechismus, sive prima institutio disciplinaque pietatis Christianæ, latine explicata.} Reprinted in Bishop Randolph's \emph{Enchirid. Theolog.} See Churton's \emph{Life of Nowell,} pp. 183 sq., and Lathbury, \emph{History of Convoc.} pp. 167 sq.}

Besides these English productions, the Catechisms of Œcolampadius, Leo Judæ, and especially those of Calvin and Bullinger were extensively used, even in the Universities, during the reign of Elizabeth.\footnote{Procter says (p. 400): 'Even in 1578, when the exclusive use of Nowell's Catechism had been enjoined in the canons of 1571, those of Calvin, Bullinger, and others were still ordered by statute to be used in the University of Oxford.'}

\xsubsection{The Lambeth Articles, A.D. 1595.}

§ 84. THE LAMBETH ARTICLES, A.D. 1595.

Literature.

\emph{Articuli Lambethani.} London, 1651. Appended to Ellis's \emph{Artic. XXXIX. Eccl. Angl. Defensio;} reprinted 1720.

Peter Heylin (Arminian): \emph{Historia Quinqu-Articularis.} London, 1660. Chaps. xx.-xxii. Also his \emph{History of the Presbyterians.}

Strype: \emph{Life and Acts of John Whitgift,} Vols. II. and III. (Oxford ed. 1822).

Thomas Fuller:  \emph{Church History of Britain,} Vol. V. pp. 219–227 (Oxford ed. of 1845).

R. Hooker's \emph{Works,} ed. Keble, Vol. I. p. cii.; Vol. II. p. 752.

Collier: \emph{An Ecclesiastical History of Great Britain,} Vol. VII. pp. 184–195.

Neale: \emph{History of the Puritans,} Vol. I. pp. 208 sqq. (Harper's ed.).

Hardwick: \emph{History of the Articles of Religion,} chap. vii. pp. 162–180, 343–347.

The Lambeth Articles are printed in Vol. III. p. 523, and also in Strype, Fuller, Collier, and Hardwick, l.c.

The Lambeth Articles have never had full symbolical authority in the Church of England, but they are of historical interest as showing the ascendency of the predestinarian system of Calvin in the last decade of the sixteenth century.\footnote{Fuller says (Vol. V. p. 227): 'All that I will say of the credit of these Articles is this: that as medals of gold and silver, though they will not pass in payment for current coin, because not stamped with the King's inscription, yet they will go with goldsmiths for as much as they are in weight; so, though these Articles want authentic reputation to pass for provincial acts, as lacking sufficient authority, yet will they be readily received of orthodox Christians for as far as their own purity bears conformity to God's Word. . . . Their testimony is an infallible evidence what was the general and received doctrine of England in that age about the forenamed controversies.'}

As Calvin became more fully known in England, he acquired an authority over the leading divines and the Universities almost as great is that of St. Augustine during the reign of Edward VI., or, in the language of Hooker, as that of the 'Master of Sentences' in the palmy days of scholasticism, 'so that the perfectest divines were judged they which were skillfullest in Calvin's writings.' Hardwick, speaking of the latter part of the Elizabethan period, admits that 'during an interval of nearly thirty years the extreme opinions of the school of Calvin, not excluding his theory of irrespective reprobation, were predominant in almost every town and parish.' The stern, bold, uncompromising predestinarianism of the Geneva Reformer seemed to furnish the best antidote to the twin errors of Pelagianism and Popery. The Puritan party without an exception, and the great majority of the conforming clergy, understood the Articles of Religion as teaching his doctrines of free-will, election, and perseverance; but some of them thought them not strong enough.

The University of Cambridge was a stronghold of the Calvinistic system. It was taught there by Thomas Cartwright, the Margaret Professor of Divinity (who, however, was deposed in 1571 for Puritanic sentiments—d. 1603); William Perkins, Fellow and Tutor of Christ's College (d. 1602);\footnote{He wrote the \emph{Golden Chain,} or \emph{Armilla aurea} (1592), which contains a very clear, logical exposition of the predestinarian order of the causes of salvation and damnation. His works were published in 3 vols. London, 1616–18.} and especially by Dr. William Whitaker (Whittaker), the Regius Professor of Divinity (d. 1595).\footnote{He wrote the best defense of the Protestant doctrine of the Scriptures against Bellarmine and Stapleton. His works were published in Latin at Geneva (1610), 2 vols., and in part republished by the Parker Society, Cambridge, 1849.}

But in the same University there arose an opposition which created great stir. It began with Baro (Baron), a French refugee, who, by the favor of Burghley, was promoted to the Margaret Professorship of Divinity (1574). He inferred from the history of the Ninevites that God predestinated all men to eternal life, but on condition of their faith and perseverance.\footnote{\emph{Prælect. in Ionam Prophetam,} London, 1579, and \emph{Concio ad Clerum,} preached in 1595. See the Letter of the heads of Cambridge, March 8, 1595, to Secretary Lord Burghley (Cecil), Chancellor of the University, in Collier, Vol. VII. p. 193.} For this opinion, which he more fully explained in a sermon, he was cited before Dr. Goade, the Vice-Chancellor of the University; and although the proceedings were stopped by the interposition of Burghley, he retired to London (1596), where he died a few years afterwards. The same cause was taken up more vigorously by William Barrett, a fellow of Caius College, who, in a 'concio ad clerum,' preached in Great St. Mary's Church, April 29, 1595, indulged in a virulent attack on the honored names of Calvin, Beza, Peter Martyr, and Zanchius, and their doctrine of irrespective predestination.

The academic controversy was carried by both parties first to the Vice-Chancellor and heads of Colleges, and then to Archbishop Whitgift, of Canterbury. Whitgift, a High-Churchman and an enemy of Puritanism, seemed at first inclined to take part with Barrett, but yielded to the pressure of the University. Barrett was obliged to admit his ignorance and mistake, and to modify his dogmatic statements. He left England and joined the Church of Rome.

To settle this controversy, and to prevent future trouble, the heads of the University sent Dr. Whitaker and Dr. Tyndal (Dean of Ely) to London, to confer with the Archbishop and other learned divines. The result was the adoption of Nine Articles, at Lambeth, Nov. 20, 1595.\footnote{This is the correct date, given by Strype from the authentic MS. copy which is headed, '\emph{Articuli approbati a reverendissimis dominis D. D. Ioanne archiepiscopo Cantuariensi, et Richardo episcopo Londinensi, et aliis Theologis, Lambethæ, Novembris} 20, \emph{anno} 1595.' Heylin and Collier give the 10th of November.} They contain a clear and strong enunciation of the predestinarian system, by teaching—

1. The eternal election of some to life, and the reprobation of others to death.

2. The moving cause of predestination to life is not the foreknowledge of faith and good works, but only the good pleasure of God.

3. The number of the elect is unalterably fixed.

4. Those who are not predestinated to life shall necessarily be damned for their sins.

5. The true faith of the elect never fails finally nor totally.

6. A true believer, or one furnished with justifying faith, has a full assurance and certainty of remission and everlasting salvation in Christ.

7. Saving grace is not communicated to all men.

8. No man can come to the Son unless the Father shall draw him, but all men are not drawn by the Father.

9. It is not in every one's will and power to be saved.

The Articles were drawn up by Whitaker (who died soon afterwards), and somewhat modified by the Bishops to make them less objectionable to anti-Calvinists. Thus the fifth Article originally stated that true faith could not totally and finally fail 'in those who had once been partakers of it;' while in the revision the words 'in the elect' (i.e., a special class of the regenerated) were substituted.\footnote{See the original draft and the comments thereon, in Hardwick, p. 345, where we find the remark: '\emph{In autographo Whitakeri verba erant, ``in iis qui semel ejus participes fuerunt;'' pro quibus a Lambethanis substituta sunt} ``\emph{in electis,'' sensu plane alio, et ad mentem Augustini; cum in autographo sint ad mentem Calvini. Augustinus enim opinatus est,} ``\emph{veram fidem quæ per dilectionem operatur, per quam contingit adoptio, justificatio et sanctificatio, posse et intercidi et amitti: fidem vero esse commune donum electis et reprobis, sed perseverantiam electis propriam}:'' \emph{Calvinus autem,} ``\emph{veram et justificantem fidem solis salvandis et electis contingere.}'''} The Articles thus amended were signed by Archbishop Whitgift, Dr. Richard Fletcher,\footnote{Not Richard Bancroft, as Fuller states; for Bancroft was not made Bishop of London till 1597.} Bishop of London, Dr. Richard Vaughan, Bishop elect of Bangor, and others. They were also sent to Dr. Hutton, Archbishop of York, and Dr. Young, Bishop of Rochester. Hutton indorsed the first Article with '\emph{verissimum},' and approved the rest with the remark that they could all be plainly collected or fairly deduced from the Scriptures and the writings of St. Augustine.

Whitgift sent the Lambeth Articles to the University of Cambridge (Nov. 24), not as new laws and decrees, but as an explanation of certain points already established by the laws of the land. But inasmuch as they had not the Queen's sanction (though he states that the Queen was fully persuaded of the truth of them, which is inconsistent with her conduct), they should be used privately and with discretion.\footnote{Heylin endeavors to relieve Whitgift from the odium of signing the Lambeth Articles by casting doubt on his honesty. Whitgift sided with Hooker against Travers, and entertained Dr. Harsnet in his family, who derided the doctrine of unconditional reprobation in a sermon at St. Paul's Cross (1584). See Collier, pp. 186, 189. But while he may have been opposed to strict Calvinism, as he certainly was to Puritanism, he seems to have been in full accord with the Augustinian infralapsarianism.}

Queen Elizabeth, who had no special liking for Calvinism and dogmatic controversies, was displeased with the calling of a Synod without her authority, which subjected the Lambeth divines to prosecution.\footnote{Fuller (Vol. V. p. 222) relates that the Queen, in her laconic style, reminded the Primate, half in jest, that by his unauthorized call of a council he had 'incurred the guilt of \emph{præmunire.}'} She commanded the Archbishop to recall and suppress those Articles without delay. At the Hampton Court Conference of King James and several prelates with the leaders of the Puritans (Jan., 1604), Dr. Reynolds made the request that 'the nine orthodoxal assertions concluded on at Lambeth might be inserted into the Book of Articles.'\footnote{See Fuller, who gives a minute account of this famous Conference, Vol. V. p. 275.} It is stated that they were exhibited at the Synod of Dort by the English deputies, as the judgment of their Church on the Arminian controversy. But the anti-Calvinistic reaction under the Stuarts gradually deprived them of their force in England, while in Ireland they obtained for some time a semi-symbolical authority.

It is interesting to compare with the Lambeth Articles a brief predestinarian document of Calvin, recently discovered by the Strasburg editors of his works,\footnote{It is printed in Vol. III. pp. 524 sq. of this work.} and a fragment of Hooker on free-will, predestination, and perseverance. The former is stronger, the latter is milder, and presents the following slight modification of those Articles:\footnote{Hooker's \emph{Works,} ed. Keble, Vol. II. pp. 752 sq.}

'It followeth therefore [says Hooker, at the close of his fragment]—

'1. That God hath predestinated certain men, not all men.

'2. That the cause moving him hereunto was not the foresight or any virtue in us at all.

'3. That to him the number of his elect is definitely known.

'4. That it can not be but their sins must condemn them to whom the purpose of his saving mercy doth not extend.

'5. That to God's foreknown elect final continuance of grace is given.

[Art. 6 of the Lambeth series is omitted by Hooker.]

'6. [7.] That inward grace whereby to be saved is deservedly not given unto all men.

'7. [8.] That no man cometh unto Christ whom God by the inward grace of his Spirit draweth not.

'8. [9.] And that it is not in every, no, not in any man's own mere ability, freedom, and power, to be saved, no man's salvation being possible without grace. Howbeit, God is no favorer of sloth; and therefore there can be no such absolute decree touching man's salvation as on our part includeth no necessity of care and travail, but shall certainly take effect, whether we ourselves do wake or sleep.'

\xsubsection{The Irish Articles. A.D. 1615.}

§ 85. THE IRISH ARTICLES. A.D. 1615.

Literature.

\emph{Works of the Most Rev.} James Ussher, D.D.,  \emph{Lord Archbishop of Armagh, and Primate of all Ireland. With a Life of the Author, and an Account of his Writings.} By  Charles Richard Elrington, D.D. Dublin, 1847, 16 Vols. See Vol. I. pp. 38 sqq. and Appendix IV.

Ch. Hardwick:  \emph{A History of the Articles of Religion,} pp. 181 sqq., 351 sqq.

James Seaton Reid, D.D.:  \emph{History of the Presbyterian Church in Ireland.} Belfast, 1834, 3 vols.

W. D. Killen, D.D. (Presb. Prof. of Eccles. Hist. at Belfast): \emph{The Ecclesiastical History of Ireland from the Earliest Period to the Present Time.} London, 1875, 2 vols. (Vol. I. pp. 492 sqq.; Vol. II. pp. 17 sqq.)

The Irish Articles are printed in Vol. III. pp. 526 sqq. of this work, in Dr. Elrington's \emph{Life of Ussher} (Vol. I. Append. IV.), in Hardwick (Append. VI.), and in Killen (Vol. I. Append. III.).

The Protestant clergy in Ireland accepted the English Prayer-Book in 1560. Whether the Elizabethan Articles of Religion were also adopted is uncertain.\footnote{Archbishop Ussher, in a sermon preached before the English House of Commons, 1621, declared: 'We all agree that the Scriptures of God are the perfect rule of our faith: we all consent in the main grounds of religion drawn from thence; \emph{we all subscribe to the Articles of Doctrine} agreed upon in the Synod of the year 1562.' But he must have understood this in the general sense of assent, as he was addressing laymen who never subscribed the Articles. Elrington, p. 43, and Hardwick, p. 182. The Irish Church adopted, in 1566 (1567), a 'Brief Declaration' in XII. Articles of Religion; but these are substantially the same as the XI. Articles prepared by Archbishop Parker, 1559 or 1560, and provisionally used in England till 1563. In Ireland they continued in force till 1615. See Elrington, Append.; Hardwick, pp. 122, 337; and Killen, Vol. I. pp. 395, 515–520.} At all events, they did not fully satisfy the rigorous Calvinism which came to prevail there for a period even more extensively than in England, and which found an advocate in an Irish scholar and prelate of commanding character and learning.

The first Convocation of the Irish Protestant clergy, which took place after the model of the English Convocation, adopted a doctrinal formula of its own, under the title 'Articles of Religion, agreed upon by the Archbishops and Bishops, and the rest of the clergy of Ireland, in the Convocation holden at Dublin in the year of our Lord God 1615, for the avoiding of diversities of opinions, and the establishing of consent touching true religion.'

They were drawn up by James Ussher,\footnote{He and his family spell the name with double s (Latin, \emph{Usserius}), but it is often spelled \emph{Usher.}} head of the theological faculty and Vice-Chancellor of Trinity College, Dublin, afterwards Archbishop of Armagh, and Primate of all Ireland. He was born in 1580, died 1656, and was buried in Westminster Abbey by order of Cromwell. He was the greatest theological and antiquarian scholar of the Episcopal Church of his age, and was highly esteemed by Churchmen and Puritans, being a connecting link between the contending parties. He was elected into the Westminster Assembly of Divines, but the King's prohibition and his loyalty to the cause of the crown and episcopacy forbade him to attend. He had an extraordinary familiarity with Biblical and patristic literature, and, together with his friend Vossius of Holland, he laid the foundation for a critical investigation of the œcumenical creeds. Whether formally commissioned by the Convocation or not, he must, from his position, have had the principal share in the preparation of those Articles. They are 'in strict conformity with the opinions he entertained at that period of his life.'\footnote{Dr. Elrington, \emph{Life of J. Ussher,} pp. 43, 44. Comp. also the 'Body of Divinity,' which was published in Ussher's name during the sessions of the Westminster Assembly, and which, he admitted to have compiled, in early life, from the writings of others.}

By a decree of the Synod appended to the Dublin Articles, they were to be a rule of public doctrine, and any minister who should publicly teach any doctrine contrary to them, and after due admonition should refuse to conform, was to be 'silenced and deprived of all spiritual promotions.' The Viceroy of Ireland, in the name of King James, gave his approval. James, with all his high notions of episcopacy and hatred of Puritanism, was a Calvinist in theology, and countenanced the Synod of Dort. It is stated that the adoption of this Confession induced Calvinistic ministers of Scotland to settle in Ireland.\footnote{Killen, Vol. I. p. 495.}

But in the reign of Charles I. and his adviser, Archbishop Laud, a reaction set in against Calvinism. An Irish Convocation in 1635, under the lead of the Earl of Strafford, Lord-Lieutenant of Ireland, and his chaplain, John Bramhall (one of the ablest High-Church Episcopalians, who was made Bishop of Londonderry, 1634, and Archbishop of Armagh, 1661—died, 1663), adopted the Thirty-nine Articles 'for the manifestation of agreement with the Church of England in the confession of the same Christian faith and the doctrine of the sacraments.' This act was intended quietly to set aside the Irish Articles; and hence they were ignored in the canons adopted by that convocation.\footnote{Killen, Vol. II. p. 23: 'The silence of the canons in respect to the Calvinistic formulary, now nearly twenty years in use, was fatal to its claims, and thus it was quietly superseded. Heylin errs in stating (\emph{Life of Laud}) that the Dublin Articles were actually 'called in.'} Ussher, however, who continued to adhere to Calvinism, though on terms of friendship with Laud, required subscription to both series, and in a contemporary letter to Dr. Ward he says: 'The Articles of Religion agreed upon in our former Synod, anno 1615, \emph{we let stand} as we did before. But for the manifestation of our agreement with the Church of England, we have received and approved your Articles also, concluded in the year 1562, as you may see in the first of our Canons.'\footnote{Elrington, \emph{Life,} p. 176.} After the Restoration the Dublin Articles seem to have been lost sight of, and no mention was made of them when, in the beginning of the nineteenth century, the English and Irish establishments were consolidated into 'the United Church of England and Ireland.'\footnote{Hardwick, p. 190.}

The Irish Articles are one hundred and four in number, arranged under nineteen heads. They are a clear and succinct system of divinity, in full harmony with Calvinism, excepting the doctrine of the ecclesiastical supremacy of the crown (which is retained from the English Articles). They incorporate the substance of the Thirty-nine Articles and the Lambeth Articles, but are more systematic and complete. They teach absolute predestination and perseverance, denounce the Pope as Antichrist, inculcate the Puritan view of Sabbath observance, and make no mention of three orders in the ministry, nor of the necessity of episcopal ordination. In all these particulars they prepared the way for the doctrinal standards of the Westminster Assembly. They were the chief basis of the Westminster Confession, as is evident from the general order, the headings of chapters and subdivisions, and the almost literal agreement of language in the statement of several of the most important doctrines.\footnote{This agreement has been proved by Professor Mitchell, D.D., of St. Andrews, in his tract \emph{The Westminster Confession of Faith,} 3d ed., Edinburgh, 1867, and in the Introduction to his edition of the \emph{Minutes of the Westminster Assembly,} 1874, pp. xlvi. sqq. We shall return to the subject more fully in the section on the Westminster Confession.}

\xsubsection{The Articles of the Reformed Episcopal Church. A.D. 1875.}

§ 86. The Articles of the Reformed Episcopal Church. A.D. 1875.

Literature.

I. Articles of Religion of the Reformed Episcopal Church,  \emph{as adopted by the General Council of the Reformed Episcopal Church, on the} 18\emph{th day of May, in the year of our Lord} 1875. New York (38 Bible House), 1876. They are printed in the last section of the third volume of this work.

II. The Book of Common Prayer of the Reformed Episcopal Church. \emph{Adopted and set forth for use by the Second General Council of the said Church, held in the City of New York, May,} 1874. Philadelphia (James A. Moore), 1874. (This took the place of the 'Proposed Book' of 1785, republished for provisional use in Dec., 1873.)

III. \emph{Journal of the First General Council of the Reformed Episcopal Church, held in New York, Dec.} 2, 1873. New York, 1873.

\emph{Journal of the Proceedings of the Second General Council of the Ref. Epis. Church, held in New York.} Philadelphia, 1874.

\emph{Journal of the Proceedings of the Third General Council of the Ref. Epis. Church, held in Chicago, Illinois, May} 12 \emph{to May} 18, 1875. Philadelphia, 1875.

IV. Bishop  George David Cummins:  \emph{Primitive Episcopacy: A Return to the Old Paths of Scripture and the Early Church. A Sermon preached in Chicago, Dec.} 14, 1873, \emph{at the Consecration of the Rev. Charles Edward Cheney, D.D., as a, Bishop in the Ref. Epis. Church.} New York, 1874.—By the same: \emph{The Lord's Table, and not the Altar.} New York, 1875.

Bishop  Chas. Edw. Cheney:  \emph{The Evangelical Ideal of a Visible Church} (a sermon). Philadelphia, 1874.

James A. Latane:  \emph{Letter of Resignation to Bishop Johns of Virginia.} Wheeling, Va., 1874.

Bishop  W. R. Nicholson:  \emph{Reasons why I became a Reformed Episcopalian.} Philadelphia, 1875.

Benj. Aycrigg:  \emph{Memoirs of the Ref. Epis. Church, and of the Prot. Epis. Church.} N.Y., 1875; 2d ed., 1889.

Before closing this section we must notice a recent American reconstruction of the English Articles of Religion, which goes much farther than the revision of the Protestant Episcopal Church, and is disowned by it, but must still be considered as an offshoot from the same root. We mean the 'Articles of Religion' set forth in 1875 by the Reformed Episcopal Church.

ORIGIN.

This body seceded from the Protestant Episcopal Church in the United States under the lead of the Rev. Dr. George David Cummins, Assistant Bishop of the Diocese of Kentucky (d. 1876). The reason of his sudden and unexpected resignation was his dissatisfaction with High-Church ritualism and exclusiveness, and his despair of checking their progress within the regular Episcopal Church. The occasion was the manifestation of this exclusiveness in a public protest of the Bishop of the diocese of New York against the General Conference of the Evangelical Alliance in Oct., 1873, and against the interdenominational communion services, in which Bishop Cummins, together with the Dean of Canterbury (with the full approval of the Archbishop of Canterbury), had taken a prominent part\footnote{In his letter of resignation to Bishop B. B. Smith, of Kentucky, dated Nov. 10, 1873, Cummins alludes to those solemn services, and adds: 'As I can not surrender the right and privilege thus to meet my fellow-Christians of other Churches around the table of our dear Lord, I must take my place where I can do so without alienating those of my own household of faith. I therefore leave the communion in which I have labored in the sacred ministry for over twenty-eight years, and transfer my work and office to another sphere of labor.'} He compared his conduct with the Old Catholic reaction against modern Romanism.\footnote{There is, however, this material difference, that the Episcopal Church as a body has not altered her creed, nor added new dogmas, as the Roman Church has done in the Vatican Council.} He desired simply to organize the theology and polity of the Low-Church party on the historic basis of the American Episcopal Church itself in its initial stage, as represented by Bishop White and the first bishops of Virginia and New York. Hence his return to the 'Proposed Book' of 1785, and to the labors of the Royal Commission in 1689.

The resignation of Bishop Cummins was followed by his canonical deposition. The majority of his brethren preferred to fight the battle within the old Church, or quietly to wait for a favorable reaction, and strongly disapproved of his course.\footnote{Though a gentleman of unblemished moral character, he was publicly charged by one of his evangelical fellow-bishops with the threefold crime of breaking his ordination vows, creating a schism, and consecrating, single-handed, a deposed clergyman (Dr. Cheney, of Chicago) to the episcopate. The last act was considered the crowning offense; for thereby he destroyed the monopoly of the apostolic succession, which, in the estimation of many modern Episcopalians, is the article of a standing or falling Church.} Others deprecated from principle the multiplication of denominations, and feared that the new sect might become narrower than the old. Still others, though unwilling to share the risk and responsibility, wished it well, in the hope that it might administer a wholesome rebuke to the hierarchical spirit. A small number of Low-Church clergymen and laymen followed his example. A new ecclesiastical organization, under the name of the Reformed Episcopal Church, was effected at a council held in the Young Men's Christian Association building, at New York, Dec. 2, 1873.\footnote{It has since grown steadily, though by no means rapidly. It numbers now (1884) ten bishops, ninety-eight presbyters, and about as many congregations in the United States, Canada, British Columbia, Bermuda Islands, and England. The number of communicants is about 7000. See art. \emph{Episcopal Church, Reformed,} by Rev. W. T. Sabine, in Schaff-Herzog \emph{Encycl.}} It set forth the following

DECLARATION OF PRINCIPLES:

I. The Reformed Episcopal Church, holding 'the faith once delivered unto the saints,' declares its belief in the Holy Scriptures of the Old and New Testaments as the Word of God, and the sole Rule of Faith and Practice; in the Creed 'commonly called the Apostles' Creed;' in the divine institution of the Sacraments of Baptism and the Lord's Supper; and in the doctrines of grace substantially as they are set forth in the Thirty-nine Articles of Religion.

II. This Church recognizes and adheres to Episcopacy, not as of divine right, but as a very ancient and desirable form of church polity.

III. This Church, retaining a Liturgy which shall not be imperative or repressive of freedom in prayer, accepts the Book of Common Prayer, as it was revised, proposed, and recommended for use by the General Convention of the Protestant Episcopal Church, A.D. 1785, reserving full liberty to alter, abridge, enlarge, and amend the same, as may seem most conducive to the edification of the people, 'provided that the substance of the faith be kept entire.'

IV. This Church condemns and rejects the following erroneous and strange doctrines as contrary to God's Word:

\emph{First,} That the Church of Christ exists only in one order or form of ecclesiastical polity.

\emph{Second,} That Christian ministers are 'priests' in another sense than that in which all believers are 'a royal priesthood.'

\emph{Third,} That the Lord's Table is an altar on which the oblation of the Body and Blood of Christ is offered anew to the Father.

\emph{Fourth,} That the Presence of Christ in the Lord's Supper is a presence in the elements of Bread and Wine.

\emph{Fifth,} That Regeneration is inseparably connected with Baptism.

The next work was the revision of the Liturgy on the basis of the 'Proposed Book' of 1785, by the Second Council, held at New York, 1874. The Apostles' Creed and the Nicene Creed were retained, but the clause 'He descended into hell' was stricken out from the former. In the baptismal service, thanksgiving for the regeneration of the child was omitted. Throughout the book the words 'minister' and 'Lord's table' were substituted for 'priest' and 'altar'—a change which had been proposed long before by the English commission of 1689.

THE ARTICLES OF RELIGION.

A considerable number of the Western members of this new denomination were in favor of adopting simply the Apostles' Creed and the Nine Articles of the Evangelical Alliance. But the majority insisted on retaining the Thirty-nine Articles with a few changes. The revision was intrusted to a Committee of Doctrine and Worship, consisting of Rev. W. R. Nicholson, D.D. (since consecrated Bishop, March, 1876), Rev. B. B. Leacock, D.D., Rev. Joseph D. Wilson, and some laymen. The report of the committee was amended and adopted at the Third General Council, held in Chicago, May 12–18, 1875.

The Articles of Religion are thirty-five in number. They follow the order of the Thirty-nine Articles, and adhere to them in language and sentiment much more closely than the Twenty Articles of the 'Proposed Book' of 1785 and the Seventeen Articles of the Episcopal Convention of 1799. Articles 1 and 2, of the Trinity and Incarnation, are retained with slight verbal alterations. Art. 3, of the descent of Christ into Hades, is omitted. Art. 3, of the Resurrection 'and the Second Coming' of Christ, Art. 4, of the Holy Ghost, and Art. 5, of the Holy Scriptures, are enlarged. Art. 8, of the old series, concerning the three creeds, is omitted; but in Art. 22 the Nicene Creed and the Apostles' Creed are acknowledged. The Articles of free-will, justification, and good works are retained, with some enlargements on justification by faith alone (which Bishop Cummins regards with Luther as the article of a standing or falling Church). Art. 18 is an abridgment of Art. 17, but affirms, together with predestination and election, also the doctrine of human freedom and responsibility, without attempting a reconciliation. The Articles of the Church and Church Authority are enlarged, but not altered in sense. Art. 24 wholly rejects the doctrine of 'Apostolic Succession' as 'unscriptural and productive of great mischief;' adding, 'This Church values its historic ministry, but recognizes and honors as equally valid the ministry of other Churches, even as God the Holy Ghost has accompanied their work with demonstration and power.' Baptism is declared to be only 'a sign of regeneration ' (not an instrument). Art. 27 rejects consubstantiation as well as transubstantiation, as 'equally productive of idolatrous errors and practices,' but otherwise agrees with Art. 28 of the old series. Arts. 31 and 32 reject purgatory, the worship of saints and images, confession for absolution, and other Romish practices. Art. 34, of the power of tie civil authority, is the same as Art. 37 of the Protestant Episcopal Church (retained from the draft of 1799), except that the words 'as well \emph{clergy} as \emph{laity}' are exchanged for 'as well \emph{ministers} as \emph{people.}'

\xsection{The Presbyterian Confessions of Scotland.}

\xsubsection{The Reformation in Scotland.}

§ 87. The Reformation in Scotland.

Literature.

I. Confessions.

[Wm. Dunlop]:  \emph{A Collection of Confessions of Faith, Catechisms, Directories, Books of Discipline, etc., of publick Authority in the Church of Scotland.} Edinburgh, 1719–22, 2 vols.

Horatius Bonar:  \emph{Catechisms of the Scottish Reformation. With Preface and Notes.} London, 1866.

Alexander Taylor Innes (Solicitor before the Supreme Court of Scotland): \emph{The Law of Creeds in Scotland. A Treatise on the Legal Relation of Churches in Scotland, established and not established, to their Doctrinal Standards.} Edinburgh, 1867 (pp. 495).

II. History of the Reformation and Church in Scotland.

Wodrow Society's Publications: 24 vols. 8vo. Lond. 1842 sqq. Comprising Knox's \emph{Works,} Calderwood's \emph{History of the Kirk of Scotland, Autobiography of Robert Blair} (1593–1636), Scott's \emph{Apologetical Narration} (1560–1633), Twedie's \emph{Select Biographies, The Wodrow Correspondence,} and other works. (The Wodrow Society was founded in 1841, in honor of Robert Wodrow, an indefatigable Scotch Presbyterian historian, b. 1679, d. 1734, for the publication of the early standard writings of the Reformed Church of Scotland.)

Spottiswoode Society's Publications. 16 vols. 8vo. Edinburgh, 1844 sqq. Comprising Keith's \emph{History} (to 1568), the Spottiswoode's \emph{History} and \emph{Miscellany,} etc.

John Knox (1505–1572): \emph{Historie of the Reformation of Religioun in Scotland} (till 1567). Edinburgh, 1584; London, 1664; better ed. by McGavin, Glasgow, 1831. Best ed. in complete \emph{Works,} edited by \emph{David Laing,} Edinburgh, 1846–64. 6 vols. (The first two vols. contain the \emph{History of the Reformation,} including the Scotch Conf. of Faith and the Book of Discipline.)

George Buchanan (1506–1682): \emph{Rerum Scoticarum Historia.} Edinburgh, 1582; Aberdeen, 1762; in English, 1690.

John Spottiswoode:  \emph{History of the Church and State of Scotland} (from 203 to the death of James VI.). London, 1668; 4th ed. 1677: ed. by the Spottiswoode Society, Edinburgh, 1847–51, in 3 vols. (John Spotswood, or Spottiswoode, was b. 1565; Archbishop of Glasgow, 1603, and then of St. Andrew's, 1615, and Chancellor of Scotland, 1635; the first in the succession of the modified Scotch episcopacy introduced by James; was obliged to retire to England, and died in London, 1639.)

David Calderwood (a learned and zealous defender of the Presbyterian Church of Scotland, d. 1650): \emph{The History of the Kirk of Scotland.} (London, 1678.) New ed. by Thomas Thomson. Edinburgh, 1842–49, 8 vols. (Wodrow Soc.)

Sir James Balfour (King-at-arms to Charles I. and II.): \emph{Historical Works published from the Original MSS.} Edinburgh, 1824, 4 vols. (Contains the Annals and Memorials of Church and State in Scotland, from 1057 to 1652.)

Robt. Keith (Primus Bishop of the Scotch Episcopal Church, Bishop of Fife, d. 1757): \emph{History of the Affairs of Church and State in Scotland, from the Beginning of the Reformation to the Retreat of Queen Mary into England,} 1568. Edinburgh, 1734, fol. (reprinted by the Spottiswoode Soc. in 2 vols. 8vo). By the same: \emph{An Historical Catalogue of the Scottish Bishops down to the year} 1688. New ed. by M. Russell. Edinburgh, 1824.

Gilbert Stuart (d. 1786): \emph{History of the Reformation of Religion in Scotland }(1517–1561). London, 1780 and 1796. By the same: \emph{History of Scotland from the Establishment of the Reformation till the Death of Queen Mary.} London, 1783, 1784, 2 vols. (In vindication of Queen Mary.)

George Cook:  \emph{History of the Reformation in Scotland.} Edinburgh, 2d ed. 1819, 2 vols. By the same: \emph{History of the Church of Scotland, from the Reformation to the Revolution.} Edinburgh, 1815, 2d ed. 1819, 3 vols.

Thomas M'Crie (d. 1835): \emph{Life of John Knox.} Edinburgh, 1811, 2 vols. 5th ed. 1831, and often; Philadelphia, 1845; \emph{Works} of M'Crie, 1858. By the same: \emph{Life of Andrew Melville.} London, 1819; 1847, 2 vols.

Thomas M'Crie, Jun.: \emph{Sketches of Scottish Church History.} 2d ed. 1843.

Prince  Alex. Labanoff:  \emph{Lettres, Instructions, et Mémoirs de Marie Stuart.} London, 1844, 7 vols.

Thomas Stephen:  \emph{History of the Church of Scotland from the Reformation to the Present Time.} London, 1843–45, 4 vols.

W. M. Hetherington (Free Church): \emph{History of the Church of Scotland till }1843. 4th ed. Edinburgh, 1853 (also New York, 1845), 2 vols.

Gen.  Von Rudloff:  \emph{Geschichte der Reformation in Schottland.} Berlin, 1847–49, 2 vols. 2d ed. 1854.

G. Weber:  \emph{Gesch. der akatholischen Kirchen u. Secten in Grossbritannien.} Leipzig, 1845 and 1853 (Vol. I. pp. 607–652; Vol. II. pp. 461–660).

John Cunningham (Presbyt.): \emph{Church History of Scotland to the Present Time.} 1859. 2 vols.

John Lee:  \emph{Lectures on the History of the Church of Scotland.} Edinburgh, 1860, 2 vols.

George Grubb (Liberal Episcopalian): \emph{Ecclesiastical History of Scotland.} London, 1861, 4 vols.

A. Teulet:  \emph{Relation politiques de la France et de l’Espagne avec l’Ecosse, en} 16\emph{me siècle.} Paris, 1862, 5 vols.

Fr. Brandes:  \emph{John Knox, der Reformator Schottlands.} Elberfeld, 1862. (The 10th vol. of Fathers and Founders of the Reformed Church.)

Merle d’Aubigne (d. 1872): \emph{History of the Reformation in Europe, in the Time of Calvin.} Vol. VI. (1876), chaps. i.–xv. (to 1546). Comp. also his \emph{Three Centuries of Struggle} (1850).

Dean Stanley (Broad-Church Episcopalian): \emph{Lectures on the History of the Church of Scotland, delivered in Edinburgh in} 1872 (with a sermon on the Eleventh Commandment, preached in Greyfriars' Church). London and New York, 1872.

Prof. R. Rainy (Free-Church Presbyterian): \emph{Three Lectures on the Church of Scotland} (against Stanley's praise of Moderatism). Edinburgh, 1872.

Geo. P. Fisher:  \emph{History of the Reformation,} pp. 351 sqq. (New York, 1873).

Peter Lorimer, D.D. (Prof. in the English Presbyterian College, London): \emph{Patrick Hamilton} (London, 1857); \emph{The Scottish Reformation} (1860); \emph{John Knox and the Church of England} (London, 1875).

Compare also the general and secular Histories of Scotland by Robertson (1759 and often, 2 vols.); Pinkerton (1814, 2 vols.); P. F. Tytler (1828–43, 9 vols., new ed. 1866, 10 vols.); John Hill Burton (from Agricola's Invasion to the Revolution of 1688. London, 1867–70, 7 vols.—From 1689 to 1748. 1870, 2 vols.); the chapters relating to Scotland in the Histories of England by Hume, Lingard (Rom. Cath.), Knight, Ranke, Froude.

The Reformation in Scotland was far more consistent and radical than in England, and resulted in the establishment of Calvinistic Presbyterianism under the sole headship of Christ. While in England politics controlled religion, in Scotland religion controlled politics. The leading figure was a plain presbyter, a man as bold, fearless, and uncompromising as Cranmer was timid, cautious, and conservative. In England the crown and the bishops favored the Reformation, in Scotland they opposed it; but Scotch royalty was a mere shadow compared with the English, and was, during that crisis, represented by a woman as blundering and unfortunate as Elizabeth was sagacious and successful. George Buchanan, the Erasmus of Scotland, the classical tutor of Mary and her son James, maintained, as the Scotch doctrine, that governments existed for the sake of the governed, which in England was regarded at that time as the sum of all heresy and rebellion.\footnote{His book, \emph{De jure regni apud Scotos} (1569), was burned at Oxford in 1683, together with the works of Milton.} When James became king of England, he blessed God's gracious goodness for bringing him 'into the promised land, where religion is purely professed, where he could sit amongst grave, learned, and reverend men; not as before, elsewhere, a king without state, without honor, without order, where beardless boys would brave him to the face.'\footnote{So he addressed the English prelates at the Hampton Court Conference. Fuller, \emph{Church History of Britain,} Vol. V. pp. 267 sq.}

The Scotch Reformation was carried on, agreeably to the character of the people of that age and country, with strong passion and violence, and in close connection with a political revolution; but it elevated Scotland at last to a very high degree of religious, moral, and intellectual eminence, which contrasts most favorably with its own mediæval condition, as well as with the present aspect of Southern Roman Catholic countries, once far superior to it in point of civilization and religion.\footnote{Thomas Carlyle calls the Scotch Reformation 'a resurrection from death to life. It was not a smooth business; but it was welcome surely, and cheap at that price; had it been far rougher, on the whole, cheap at any price, as life is. The people began to \emph{live}; they needed first of all to do that, at what cost and costs soever. Scotch literature and thought, Scotch industry; James Watt, David Hume, Walter Scott, Robert Burns: I find Knox and the Reformation acting in the heart's core of every one of these persons and phenomena; I find that without the Reformation they would not have been. Or what of Scotland? The Puritanism of Scotland became that of England, of New England. A tumult in the High Church of Edinburgh spread into a universal battle and struggle over all these realms; then came out, after fifty years' struggling, what we call the glorious Revolution, a Habeas-Corpus Act. Free Parliaments, and much else!'—\emph{Heroes,} Lect. IV.}

In the middle of the sixteenth century the Scotch were still a semi-barbarous though brave and energetic race. Their character and previous history are as wild and romantic as their lochs, mountains, and rapids, and show an exuberance of animal life, full of blazing passions and violent commotions, but without ideas and progress. The kings of the house of Stuart were in constant conflict with a restless and rebellious nobility and the true interests of the common people. The history of that ill-fated dynasty, from its fabulous patriarch Banquo, in the eleventh century, down to the execution of Queen Mary (1587), and the final expulsion of her descendants from England (1688), is a series of tragedies foreshadowed in Shakspere's 'Macbeth,' where crimes and retributions come whirling along like the rushing of a furious tempest. The powerful and fierce nobility were given to the chase and the practice of arms, to rapine and murder. Their dress was that of the camp or stable; they lived in narrow towers, built for defense, without regard to comfort or beauty. They regarded each other as rivals, the king as but the highest of their own order, and the people as mere serfs, who lived scattered under the shadow of castles and convents. The patriarchal or clan system which prevailed in the Highlands, and the feudal system which the Norman barons superinduced in the south, kept the nation divided into a number of jealous and conflicting sections, and made the land a scene of chronic strife and anarchy.

In this unsettled state of society morals and religion could not flourish. The Church kept alive the faith in the verities of the supernatural world, restrained passion and crime, distributed the consolations of religion from the cradle to the grave, and built such monuments as the Cathedral of Glasgow and the Abbey of Melrose; but it left the people in ignorance and superstition. It owned the full half of all the wealth of the nation from times when land was poor and cheap, and it had the controlling influence in the privy council, the parliament, and over the people. But this very wealth and political power became a source of corruption, which rose to a fearful height before the Reformation. The law of celibacy was practically annulled, and the clergy were shamefully dissolute and disgracefully ignorant. Some priests are said to have regarded Luther as the author of the New Testament. The bishops and abbots, by frequently assisting the king against the nobles, and rivaling with them in secular pomp and influence, excited their envy and hatred, which hastened their ruin.

Owing to its remoteness, poverty, and inhospitable climate, Scotland was more free than England from the interference of the pope and his Italian creatures. But this independence was rather a disadvantage, for without preventing the progress of the native corruptions, it kept off the civilizing influences of the Continent, and removed the check upon the despotism of the king. James III. usurped the right of filling the episcopal vacancies without the previous election of the chapters and the papal sanction, and consulted his temporal interest more than that of religion. Simony of the most shameful kind became the order of the day. James V. (1528–42) provided for his illegitimate children by making them abbots and priors of Holyrood House, Kelso, Melrose, Coldingham, and St. Andrew's. Most of the higher dignities of the Church were in the hands of the royal favorites and younger sons of the nobility, who were sometimes not ordained, nor even of age, but who drew, nevertheless, the income of the cathedrals and abbeys, and disgraced the holy office. 'By this fraudulent and sacrilegious dealing'—says an impartial old authority—'the rents and benefices of the Church became the patrimony of private families, and persons in no ecclesiastical orders, and even boys too, were, by the presentation of our kings and the provision of the popes, set over the episcopal sees themselves. The natural result of this was that by far too many of these prelates, being neither bred up in letters, nor having in them any virtuous dispositions, did not only live irregularly themselves, but through neglect of their charge did likewise introduce by degrees such a deluge of ignorance and vice among the clergy and all ranks of men that the state of the Church seemed to call loudly for a reformation of both.'

The first impulse to the Reformation in Scotland came from Lutheran writings and from copies of Tyndale's New Testament. The first preachers and martyrs of Protestantism were Patrick Hamilton, who had studied in Wittenberg and Marburg, and was burned (1528), George Wishart, who shared the same fate (1546), and the aged Walter Mill, who predicted from the flames (Aug. 28, 1558), 'A hundred better men shall rise out of the ashes of my bones, and I shall be the last to suffer death in Scotland for this cause.'

In the mean time God had prepared the right man for this crisis.

\xsubsection{John Knox.}

§ 88. John Knox.

Literature.

Besides the works of Knox, the excellent biography of M'Crie, and Lorimer's monograph quoted in the preceding section, comp. Froude's Lecture on \emph{The Influence of the Reformation on the Scottish Character,} 1865 (in \emph{Short Studies on Great Subjects,} Vol. I. pp. 128 sqq.), and an exceedingly characteristic essay of Thomas Carlyle on the \emph{Portraits of John Knox,} which first appeared in \emph{Fraser's Magazine} for April, 1875, and then as an appendix to his \emph{Early Kings of Norway.} London, 1875 (pp. 209–307), and New York (Harper's ed. pp. 173–257). Brandes follows M'Crie very closely. Laing, in the first vol. of his edition of Knox's \emph{History of the Reformation} (pp. xiii.–xliv.), gives a convenient chronological summary of the chief events of his life.

John Knox (1505–1572), the Luther of Scotland, was educated in the University of Glasgow, and ordained to the Romish priesthood (1530), but became a convert to Protestantism (1545, the year of Wishart's martyrdom\footnote{This is the date given by Laing, while M'Crie assigns Knox's conversion to the year 1542.} through the study of the Bible and the writings of Augustine and Jerome. He went at once to the extreme of opposition, as is often the case with strong and determined characters of the Pauline type. He abhorred the mass as an 'abominable idolatry and profanation of the Lord's Supper,' and popery as the great anti-Christian apostasy and Babylonish harlot predicted in the Bible.\footnote{His first Protestant sermon in the parish church at St. Andrew's was on Dan. vii., to prove that the pope was the last beast, the man of sin, the Antichrist. Some of the hearers said: 'Others hewed at the branches of papistry, but he struck at the root to destroy the whole.' Calderwood, Vol. I. p.230; Knox's \emph{Works,} Vol. I. p. 192.}

After preaching awhile to the Protestant soldiers in the garrison of St. Andrew's, he was taken prisoner by the French fleet (1547), and made a galley-slave for nineteen months, 'going in irons, miserably entreated and sore troubled by corporal infirmity.' Regardless of danger, he remained true to his faith. When called upon to kiss an image of the Holy Virgin, he declared that it was 'no mother of God, but a painted piece of wood, fit for swimming rather than being worshiped;' and he flung the picture into the river Loire.

On obtaining his liberty, he labored five years (1549–1554) in England as a pioneer of English Puritanism. He preached in Berwick, on the borders of Scotland, in Newcastle, and in London. He was elected one of the six chaplains of Edward VI. (1551), was consulted about the Articles of Religion and the revision of the Liturgy, and was offered the bishopric of Rochester, which he declined from opposition to the large extent of dioceses, the secular business, vestments, and 'other popish fooleries remaining.'\footnote{His labors in England, and the reasons for his \emph{nolo episcopari}, are fully described by Dr. Lorimer, in part from unpublished sources.}

After the accession of Bloody Mary he fled among the last, at the urgent request of friends, to the Continent, and spent five years (from January, 1554, to January, 1559, interrupted by a journey to Scotland, November, 1555, to July, 1556), at Frankfort-on-the-Main, and especially at Geneva. Here he found 'the most perfect school of Christ that ever was since the days of the Apostles.' Though four years older, he sat an admiring pupil at the feet of John Calvin, and became more Calvinistic than the great Reformer. He preached to a flock of English exiles, took part in the Geneva version of the Bible, and aided by his pen the cause of evangelical religion in England and Scotland.

The accession of Queen Elizabeth opened the way for his final return and crowning work, although she refused him passage through her dominion, and never forgave him his '\emph{blast}' at the dignity and ruling capacity of her sex.\footnote{Before his return, while the fires of Smithfield were still burning, he had published anonymously his 'First Blast of the Trumpet against the Monstrous Regiment [i.e., regimen or government] of Women,' 1558, which was aimed at the misgovernment of Mary Tudor and Mary of Guise. This singular and characteristic but unfortunate book begins with the sentence, 'To promote a woman to bear rule, superiority, dominion, or empire, above any realm, nation, or city, is repugnant to nature, contumely to God, a thing most contrarious to his revealed will and approved ordinance, and, finally, it is a subversion of all equity and justice.' He appealed to the creation, to the Jews, to St. Paul, to ancient philosophers and legislators, to the fathers, to the Salic and French law. His error was that from some bad examples he drew sweeping conclusions, which were soon confirmed by Mary Stuart, but disproved by Elizabeth (as they are in our day by the reign of Victoria). No wonder that Queen Mary and Queen Elizabeth were incensed at what they regarded a personal insult. Knox himself foresaw the bad consequences, and expected to be called 'a sower of sedition, and one day perchance to be attainted for treason,' but he was too manly to retract, and retained his opinion to the last, but, not wishing to obstruct the path of Elizabeth, he never published the intended \emph{Second} and \emph{Third Blast.} See M'Crie's \emph{J. Knox,} pp. 141–147 (Philadelphia ed.), and Carlyle, l.c. pp. 230 sqq.}

The remaining twelve years of his life were devoted to the fierce struggle and triumph of the Reformation in his native land, which he has himself so vividly, truthfully, and unselfishly described in his \emph{History.}\footnote{Knox wrote four Books of his \emph{History of the Reformation,} down to 1564, at the request of his friends. The Fifth Book is not found in any MS. copy, and was first published by David Buchanan in 1644; it relates the affairs of the most controverted period in Scottish history, from Sept., 1564, to Aug., 1567, when Queen Mary abdicated. Laing thinks that it is mostly derived from Knox's papers by some unknown hand (\emph{Works,} Vol. II. p. 468). Carlyle regrets that this 'hasty and strangely interesting, impressive, and peculiar \emph{History} has not been rendered far more extensively legible to serious mankind at large.' Laing has added a vocabulary.} Shortly before his death he heard the news of the terrible massacre of the Huguenots on St. Bartholomew's night, and summoning up the remainder of his broken strength, he thundered from the pulpit in Edinburgh his indignation and the vengeance of God against 'that cruel murderer and false traitor, the King of France' (Charles IX.). His last sermons were on our Lord's crucifixion, a theme on which he wished to close his ministry. He presided at the installation of Lawson as his colleague and successor, and made an impressive address and prayer. As he left the church a crowd of people lined the street and followed him to his house to take farewell of their pastor. He found his last comfort in the sacerdotal prayer, the fifty-third chapter of Isaiah, and some psalms, 'hearing' what was read, and 'understanding far better.' He died, weary of life and longing for heaven, in the sixty-seventh year of his age, in peace, without a struggle, lamented by the clergy, the nobles, and the people (Nov. 24, 1572). He could conscientiously say on his death-bed, before God and his holy angels, that he never made merchandise of religion, never studied to please men, never indulged his private passions, but faithfully used his talents for the edification of the Church over which he was called to watch. He was buried in the graveyard of St. Giles's; no monument was erected; a plain stone with his name marks the spot.

Knox was the greatest of Scotchmen, as Luther the greatest of Germans. He was the incarnation of all the noble and rugged energies of his nation and age, and devoted them to the single aim of a thorough reformation in doctrine, worship, and discipline, on the basis of the Word of God.\footnote{Thomas Carlyle, himself a typical Scotchman, calls Knox 'the most Scottish of Scots, and to this day typical of all the qualities which belong nationally to the very choicest Scotsmen we have known, or had clear record of: utmost sharpness of discernment and discrimination, courage enough, and, what is still better, no particular consciousness of courage, but a readiness in all simplicity to do and dare whatsoever is commanded by the inward voice of native manhood; on the whole, a beautiful and simple but complete incompatibility with whatsoever is false in word or conduct; inexorable contempt and detestation of what in modern speech is called \emph{humbug,} . . . a most clear-cut, hardy, distinct, and effective man; fearing God, and without any other fear.' He severely characterizes the patriarchal, long-bearded, but stolid picture of Knox in Beza's \emph{Icones} (Geneva, 1580), and in Laing's edition, and represents the 'Somerville portrait,' with a sharp, stern face, high forehead, pointed beard, and large white collar, as the only probable likeness of the great Reformer.} In genius, learning, wealth of ideas, and extent of influence, he was inferior to Luther and Calvin, but in boldness, strength, and purity of character, fully their equal.\footnote{M'Crie (p. 355) well compares him with the three leading Reformers: 'Knox bore a striking resemblance to Luther in personal intrepidity and in popular eloquence. He approached nearest to Calvin in his religious sentiments, in the severity of his manners, and in a certain impressive air of melancholy which pervaded his character. And he resembled Zwinglius in his ardent attachment to the principles of civil liberty, and in combining his exertions for the reformation of the Church with uniform endeavors to improve the political state of the people. Not that I would place our Reformer on a level with this illustrious triumvirate. There is a splendor which surrounds the great German Reformer, partly arising from the intrinsic heroism of his character, and partly reflected from the interesting situation in which his long and doubtful struggle with the Court of Rome placed him in the eyes of Europe, which removes him at a distance from all who started in the same glorious career. The Genevese Reformer surpassed Knox in the extent of his theological learning, and in the unrivaled solidity and clearness of his judgment. And the Reformer of Switzerland, though inferior to him in masculine elocution and in daring courage, excelled him in self-command, in prudence, and in that species of eloquence which steals into the heart, convinces without irritating, and governs without assuming the tone of authority. But although ``he attained not to the first three,'' I know not, among all the eminent men who appeared at that period, any name which is so well entitled to be placed next to theirs as that of Knox, whether we consider the talents with which he was endowed, or the important services which he performed.'} He was the most heroic man of a heroic race. His fear of God made him fearless of man. Endowed with a vigorous and original intellect, he was eminently a man of action, with the pulpit for his throne and the word for his sword. A statesman as well as a theologian, he possessed rare political sagacity and intuitive knowledge of men. Next to Calvin, he is the chief founder of the Presbyterian polity, which has proved its vitality and efficiency for more than three centuries. Like St. Paul and Calvin, he was small in person and feeble in body, but irresistible in moral force.\footnote{\par '\emph{Haud scio an unquam majus ingenium in fragili et imbecillo corpusculo collocarit.}' Principal Smeton, as quoted by M'Crie, p. 355.} 'He put more life into his hearers from the pulpit in an hour than six hundred trumpets.'\footnote{So the English embassador, Sir Nicholas Throckmorton, wrote to Cecil.} When old and decrepit, leaning on his staff and the arm of his faithful servant, he had to be lifted to the pulpit; but before the close he became so animated and vigorous that he seemed 'likely to ding the pulpit in blads [to beat it in pieces] and flie out of it.'\footnote{Thus his eloquence was described, in 1571, by James Melville, then a student and constant hearer of Knox. A lively Frenchman, in the \emph{Journal des Debats,} gave the following amusing version of this account: 'A Presbyterian fanatic named Knox, . . . old and broken down, . . . began his sermon in a feeble voice and slow action; but soon heating himself by the force of his passion and hatred, he bestirred himself like a madman; \emph{he broke his pulpit, and jumped into the midst of his hearers} (\emph{sautoit au milien des auditeurs}).' M'Crie, p. 325.} Well did the Earl of Morton, the newly elected regent, characterize him over his open grave in that sentence which has since been accepted as the best motto of his life: 'Here lies he who never feared the face of man.'\footnote{Or, in the less graceful but more expressive original phrase, as given by James Melville (the only authority for it), 'He neither feared nor flattered any flesh.'} And in a different spirit, James VI paid the same tribute to his fearless character, when with uplifted hands he thanked God that the three surviving bairns of Knox were all lasses; 'for if they had been three lads,' he said to Mrs. Welch, 'I could never have bruiked [enjoyed] my three kingdoms in peace.'\footnote{Mrs. Welch was a daughter of Knox, and gained admission to the King, in London, 1622, to ask his permission for the return of her sick husband (a worthy Presbyterian minister, who had been exiled for his resistance to the re-establishment of episcopacy) to his native Scotland. James at last yielded on condition that she should persuade him to submit to the bishops; but the lady, lifting up her apron and holding it towards the King, replied, in the genuine spirit of her father, 'Please your Majesty, I'd rather kep [receive] his head there.' Mr. Welch died in London soon after this singular conversation; his widow returned to Ayr, and survived him three years, 'a spouse and daughter worthy of such a husband and such a father.' M'Crie, p. 362. Knox was twice married and had two sons by his first wife, Marjory Bowes, of London, and three daughters by his second wife, Margaret Stewart, of a high noble family in Scotland. The sons were educated at Cambridge, but died young, without issue.}

Knox had the stern and uncompromising spirit of a Hebrew prophet. He confronted Queen Mary as Elijah confronted Jezebel, unmoved by her beauty, her smiles, or her tears. He himself relates the four or five interviews he had with that graceful, accomplished, fascinating, but ill-fated lady, whose charms and misfortunes still excite fresh feelings of sympathy in every human heart. It is difficult to imagine a more striking contrast: Knox the right man in the right place, Mary the wrong woman in the wrong place; he intensely Scotch in character and aim, she thoroughly French by education and taste; he in the vigor of manhood, she in the bloom of youth and beauty; he terribly in earnest, she gay and frivolous; he a believer in God's sovereignty and the people's right and duty to disobey and depose treacherous princes, she a believer in her own absolute right to rule and the subject's duty of passive obedience; he abhorring her religion as idolatry and her policy as ruin to Scotland, she fearing him as a rude fanatic, an impertinent rebel and sorcerer in league with Beelzebub.\footnote{Carlyle thus speaks of this remarkable chapter in the Scotch Reformation: 'The interviews of Knox with the Queen are what one would most like to produce to readers; but unfortunately they are of a tone which, explain as we might, not one reader in a thousand could be made to sympathize with or do justice to in behalf of Knox. The treatment which that young, beautiful, and high chief personage in Scotland receives from the rigorous Knox, would to most modern men seem irreverent, cruel, almost barbarous. Here more than elsewhere Knox proves himself,—here more than any where bound to do it,—the Hebrew Prophet in complete perfection; refuses to soften any expression or to call any thing by its milder name, or in short for one moment to forget that the Eternal God and His Word are great, and that all else is little, or is nothing; nay, if it set itself against the Most High and His Word, is the one frightful thing that this world exhibits. He is never in the least ill-tempered with her Majesty; but she can not move him from that fixed centre of all his thoughts and actions: Do the will of God, and tremble at nothing; do against the will of God, and know that, in the Immensity and the Eternity around you, there is nothing but matter of terror. Nothing can move Knox here or elsewhere from that standing-ground; no consideration of Queen's sceptres and armies and authorities of men is of any efficacy or dignity whatever in comparison; and becomes not beautiful, but horrible, when it sets itself against the Most High.'} We must not judge from his conversations with the Queen that he was a woman-hater: he respected right women in their proper sphere, as he was respected by them, and his correspondence reveals a vain of tenderness and kindly genial humor beneath his severity.\footnote{See his letters of comfort to Mrs. Bowes, his mother-in-law, who suffered much from religious melancholy, in \emph{Works} by Laing, Vol. III. pp. 337–343, and Vol. VI. p. 513; also in Lorimer, pp. 39 sqq.} But in this case he sacrificed all personal considerations to what he believed to be his paramount duty to God and his Church.

The pulpit proved mightier than the throne. The suicidal blunders of the Queen, who had more trouble from her three husbands—two of them handsome but heartless and worthless ruffians and murderers—than her grand-uncle Henry VIII. had from his six wives, are the best vindication of the national policy, if not the personal conduct, of the Reformer. Had Mary's popish policy triumphed, there would have been an end to Protestantism and liberty in Scotland, and probably in England too; while Knox, fighting intolerance with intolerance, laid the solid foundation for future liberty. He felt at that turning-point of history that, what is comparatively harmless now, 'one mass was more dangerous to Scotland than an army of ten thousand enemies.'\footnote{Froude says: 'Toleration is a good thing in its place; but you can not tolerate what will not tolerate you, and is trying to cut your throat. . . . The Covenanters fought the fight and won the victory, and then, and not till then, came the David Humes with their essays on miracles, and the Adam Smiths with their political economies, and steam-engines, and railroads, and philosophical institutions, and all the other blessed or unblessed fruits of liberty' 1.c. pp. 148, 149).}

If Knox lacked the sweet and lovely traits of Christian character, it should be remembered that God wisely distributes his gifts. Neither the polished culture of Erasmus, nor the gentle spirit of Melanchthon, nor the cautious measures of Cranmer could have accomplished the mighty change in Scotland. Knox was, beyond a doubt, the providential man for his country. Scotland alone could produce a Knox, and Knox alone could reform Scotland. If any man ever lived to some purpose, and left the indelible impress of his character upon posterity, it was John Knox. His is to this day the best known and the most popular name in Scotland. Such fearless and faithful heroes are among the best gifts of God to the world.

We need not wonder that Knox, like the other Reformers, was pursued by malignant calumny during his life, and even charged with unnatural crimes, which would make him ridiculous as well as hideous. But those who knew him best esteemed him most. Bannatyne, his faithful clerk, calls him, in his journal, 'the light of Scotland, the comfort of the Church, the mirror of godliness, the pattern of all true ministers in purity of life, soundness of doctrine, and boldness in reproving wickedness.' James Melville, who heard his last sermons, speaks of him as 'that most notable prophet and apostle' of Scotland.\footnote{Beza also calls him '\emph{Scotorum apostolum.}'} Posterity has judged differently, according to the religious stand-point. To some he still appears as a semi-barbarous fanatic, a dangerous heretic, or at best as a 'holy savage;' while Froude regards him as 'the grandest figure in the entire history of the British Reformation,' and Carlyle as 'more than a man of genius—a heaven-inspired prophet and heroic leader of men.'

\xsubsection{The Scotch Confession of Faith. A.D. 1560.}

§ 89. The Scotch Confession of Faith. A.D. 1560.

Literature.

The Scotch Confession in the original Scotch dialect, together with the authorized Latin version of Patrick Adamson (1572), is printed in Vol. III. pp. 427–470, from Dunlop's \emph{Collection,} Vol. II. pp. 13 sqq. It appeared at Edinburgh, 1561 (Robert Lekprevik), without the marginal Scripture references, in the Minutes of Parliament, in Knox's \emph{History of the Reformation} (Vol. II. pp. 93 sqq.; Laing's ed.), in Calderwood's \emph{History of the Kirk of Scotland} (Vol. II. pp. 16 sqq.; Thomson's ed. for the Wodrow Soc.), in Edward Irving's reprint of the Conf. and the Book of Discipline (1831), and (abridged) in Innes, \emph{Law of Creeds} (pp. 38 sqq.). In the \emph{Writings of John Knox,} by the Presbyterian Board of Publication, Phila., 1842, pp. 237 sqq., it is given in modern English.

A Latin version (less correct and elegant than that of Adamson) appeared in the \emph{Corpus et Syntagma Conf.,} 1612 and 1654, and is reproduced in Niemeyer's \emph{Collectio,} pp. 340 sqq. Niemeyer's critical notice in the \emph{Prolog.,} p. li., is very brief and meagre. For a German translation, see Böckel, pp. 645 sqq.

The supplementary Scotch Confession of 158O is printed in Vol. III. pp. 470–475.

ORIGIN OF THE SCOTCH CONFESSION.

'The Creed of Scotland and the Church of Scotland emerge into history so nearly at the same moment that it is difficult to say which has the precedence even in order of time. It is at least equally difficult to say which is first in respect of authority; and, indeed; the question whether the Church is founded upon the creed or the creed upon the Church appears to be at the root of most of the legal difficulties that lie before us.'\footnote{Innes, \emph{The Law of Creeds in Scotland,} p. 4.}

The Reformed Church of Scotland was not legally recognized and established by Parliament till 1567, seven years after the Scotch Confession was adopted and the first General Assembly was held; but it existed in fact, under royal protest, as a voluntary body from December 3, 1557, when a number of Protestant nobles and gentlemen signed, at Edinburgh, a 'Covenant' to maintain, nourish, and defend to the death 'the whole \emph{Congregation of Christ,} and every member thereof.' This was one of those religious bonds or mutual agreements by which the confederation of Protestants of Scotland was so often ratified to secure common privileges. The term \emph{Congregation} (\textgreek{ἐκκλησία,} \emph{ecclesia}), which afterwards was exchanged for \emph{Kirk} (\textgreek{κυριακόν}), then signified the true Church of Christ in opposition to the apostate Papal Church, and the leaders were called the 'Lords of the Congregation.' For a few years the Liturgy of Edward VI. and the 'Order of Geneva' seem to have been used, but there is no record of any formal approval of a doctrinal standard before 1560.\footnote{'The Confession of Faith of the English Congregation at Geneva,' 1558, consists only of four articles: 1, of God the Father; 2, of Jesus Christ; 3, of the Holy Ghost; 4, of the Church and the Communion of Saints. It was probably drawn up by Knox. Chaps. 1 and 4 have some resemblance to the corresponding articles of the Scotch Confession. It is reprinted in Dunlop's \emph{Collection,} Vol. II. pp. 3–12. The editor says that it was 'received and approved by the Church of Scotland in the beginning of the Reformation.'}

On the first of August, 1560, after the death of the Queen Regent, Mary of Guise, and the expulsion of the French troops, but before the arrival of Queen Mary, the Scotch Parliament convened at Edinburgh to settle the new state of things in this transition period. It proved to be the most important meeting in the history of that kingdom. The Church question came up on a petition to abolish popery, to restore the purity of worship and discipline, and to devote the ecclesiastical revenues to the support of a pious clergy, the promotion of learning, and the relief of the poor. In answer to the first request, the Parliament directed the Protestant ministers to draw up a Confession of Faith. This was done hastily, though not without mature preparation, in four days, by John Knox and his compeers.\footnote{Knox reports (Vol. II. p. 128): 'Commission and charge was given to Mr. John Winram, sub-prior of St. Andrew's, Mr. John Spottiswoode, John Hillock, Mr. John Douglas, rector of St. Andrew's, Mr. John Rowe, and John Knox, to draw in a volume the policy and discipline of the Kirk, \emph{as well as they had done the doctrine.}' Thus six Johns composed both the Confession of Faith and the Book of Discipline, which breathe the spirit of the Church militant, and are Pauline rather than Johannean. Knox was no doubt the chief author of both. He had experience in the preparation of such documents, as he was consulted about the Edwardine Articles of Religion, prepared the Confession for the English congregations in Geneva, and must have been familiar with the Swiss Confessions.} The document was read twice, article by article, and ratified by the three estates, August 17, 1560, 'as a doctrine grounded upon the infallible Word of God.' Every member was requested to vote. The papal bishops were charged to object and refute, but they kept silence. The temporal lords all voted for the Confession except three, the Earl of Athole, Lord Somerville, and Lord Borthwick, who declared as their only reason of dissent, 'We will beleve as our fathers belevet.'\footnote{Knox, \emph{Hist.} Vol. II. p. 121; Calderwood, Vol. II. p. 37.}

Randolph, the English envoy, wrote to Cecil two days afterwards: 'I never heard matters of so great importance neither sooner dispatched, nor with better will agreed unto. . . . The rest of the Lords, with common consent and as glad a will as ever I heard men speak, allowed the same. . . . Many offered to shed their blood in defense of the same. The old Lord Lindsay, as grave and godly a man as ever I saw, said, ``I have lived many years; I am the oldest in this company of my sort; now that it hath pleased God to let me see this day, where so many nobles and others have allowed so worthy a work, I will say with Simeon, \emph{Nunc dimittis.}'''\footnote{Knox, \emph{Works,} Vol. VI. pp. 116–118: Innes, p. 11.}

The adoption of the Confession was followed (Aug. 24, 1560) by acts abolishing the mass, the jurisdiction of the pope, and rescinding all the laws formerly made in support of the Roman Catholic Church and against the Reformed religion. A messenger was dispatched with the Confession to Queen Mary, in Paris, to secure her ratification, but was not graciously received. Her heart's design was to restore in due time her own religion.

In December of the same year the first General Assembly convened, and approved the Book of Discipline, prepared by the same authors. It was submitted to the state authority, but this refused to ratify it.\footnote{See 'The Booke of the Universall Kirk of Scotland,' containing the earliest records of the Minutes of the Assembly, published in one volume, 1839; Calderwood, Vol. II. pp. 44 sqq.; Innes, pp. 21 sqq.}

Seven years afterwards (1567), the Parliament formally established the Reformed Church, by declaring the ministers of the blessed Evangel and the people of the realm professing Christ \emph{according to the Confession of Faith} 'to be the only true and holy Kirk of Jesus Christ within this realm.' Subscription was required from all ministers first in 1572.\footnote{Innes, pp. 30 and 49.} From that time till the Revolution of 1688 this native Confession was the only legally recognized doctrinal standard of both the Presbyterian and Episcopal Churches in Scotland. The Covenanters, however, during the Commonwealth, adopted the Westminster standards, which have thrown the older Confession into the shade. Besides, the General Assembly approved and recommended also the Second Helvetic Confession, which Beza transmitted to Scotland (1566), Calvin's Catechism, and the Heidelberg Catechism, but no subscription to these foreign confessions was ever exacted.

CONTENTS.

The Scotch Confession consists of twenty-five Articles, and a short Preface, which breathes the spirit of true confessors ready for martyrdom. It begins: 'Long have we thirsted, dear brethren, to have notified unto the world the sum of that doctrine which we profess, and for the which we have sustained infamy and danger;' and it ends with the words: 'We firmly purpose to abide to the end in the confession of this our faith.' But the authors are far from claiming infallibility for their own statements of the truth, and subject them to improvement and correction from the Holy Scriptures.\footnote{'We protest that if any one will note in this our Confession any article or sentence repugnant to God's Holy Word, that it would please him of his gentleness and for Christian charity's sake, to admonish us of the same in writing; and we, upon our honor and fidelity, by God's grace, do promise unto him satisfaction from the mouth of God (that is, from his Holy Scriptures), or \emph{else reformation of that which he shall prove to be amiss.}' Dean Stanley, in quoting this passage from the Preface (\emph{Lectures,} etc. p. 113), says that it is the only Protestant Confession which, far in advance of its age, acknowledges its own fallibility. But the First Confession of Basle (1534) does the same in express words in the closing article (see Niemeyer, \emph{Collect.} pp. 84 and 104); and the changes of the Augsburg Confession (Art. X.), and of the English Articles, imply the recognition of their imperfection on the part of the authors. The 19th Article, in declaring that all Churches have erred in matters of faith, could certainly not intend to exempt the Church of England and her formularies.} In harmony with this, the 20th Article denies the infallibility of general councils, 'some of which have manifestly erred, and that in matters of great weight and importance.'

The Confession covers the œcumenical and evangelical doctrines, beginning with God and ending with the Church, the Sacraments, and the Civil Magistrate. It exhibits a clear, fresh, and forcible summary of the orthodox Reformed faith, as then held in common by the Protestants of England, Switzerland, France, and Holland. Though decidedly Calvinistic, it is yet free from the scholastic technicalities and angular statements of the Calvinism of a later generation. The doctrine of the Sacraments is similar to and rather stronger than that of the Thirty-nine Articles.\footnote{Tytler (\emph{History of Scotland,} Vol. III. p. 129, ed. of 1872) observes: 'It is worthy of remark that in these holy mysteries of our faith this Confession, drawn up by the primitive Scotch Reformers, keeps in some points at a greater distance from the rationalizing of ultra-Protestantism than the Articles of Edward.' On Knox's view of the eucharist, see Lorimer, 1.c. pp. 129 and 131. He held the Calvinistic view before he came to Geneva, and while still a disciple of Wishart, who learned it from his intercourse with the Swiss Churches to 1540, and translated the First Helvetic Confession of 1536 into English.} The Church is declared to be uninterruptedly one from the beginning to the end of the world,' one company and multitude of men chosen of God, who rightly worship and embrace him by true faith in Christ Jesus, who is the \emph{only Head} of the same Church, which also is the body and spouse of Christ Jesus; which Church is catholic, that is, universal, because it containeth the elect of all ages, all realms, nations, and tongues, who have communion with God the Father, and with his Son, through the sanctification of the Holy Spirit.' But this Church is put in strong contrast with the false and apostate Church of the Papacy, and distinguished from it by three marks— namely, the pure preaching of the gospel, the right administration of the sacraments, and the exercise of ecclesiastical discipline. The first two are mentioned in the Augsburg Confession and the English Articles; the third is peculiarly Calvinistic and Presbyterian.

But no particular form of Church government or worship is laid down in this Confession as binding, and freedom is allowed in ceremonies.\footnote{Art. XX.: 'In the Church, as in the house of God, it becometh all things to be done decently and in order: not that we think that one policy, and one order of ceremonies can be appointed for all ages, times, and places; for as ceremonies, such as men have devised, are but temporal, so may and ought they to be changed, when they rather foster superstition than edify the Church using the same.'} Knox himself prepared, after the Geneva model, a liturgy, or Book of Common Order, which was indorsed by the General Assembly (Dec. 26, 1564), and used in Scotland for a long time.\footnote{It has been republished by the Rev. John Cumming, London, 1840. Cumming says (p. v.): 'The Scotch Church never objected to a written liturgy in her public worship, provided there was room left in the service for extemporaneous service.' John Knox's Liturgy was never formally abolished, but, like the Scotch Confession, it was silently superseded by the Westminster standards.} The exclusive theory of a \emph{jure divino} Presbyterianism dates not so much from Knox as from Andrew Melville, and the aversion to forms of prayer was a reaction against the attempt of Laud to force a foreign episcopacy and liturgy upon the reluctant Scotch.

Edward Irving, himself one of the purest and noblest sons of Scotland, who for several years thrilled the English metropolis with his pentecostal gift of tongues, and to whom Thomas Carlyle, the friend of his youth, paid such a touching tribute, was in the habit of reading the Scotch Confession twice in the year to his congregation, and bestowed this encomium upon it:\footnote{\emph{Collected Writings of Edward Irving,} London, 1864, Vol. I. p. 601, quoted by Innes, p. 55.} 'This document is the pillar of the Reformation Church of Scotland, which hath derived little help from the Westminster Confession of Faith: whereas these twenty-five articles, ratified in the Parliament of Scotland in the year 1560, not only at that time united the states of the kingdom in one firm band against the Papacy, but also rallied the people at sundry times of trouble and distress for a whole century thereafter, and it may be said even until the Revolution, when the Church came into that haven of rest which has proved far more pernicious to her than all the storms she ever passed through; for, though the Westminster Confession was adopted as a platform of communion with the English Presbyterians in the year 1647, it exerted little or no influence upon our Church, and was hardly felt as an operative principle either of good or evil, until the Revolution of 1688; so that the Scottish Confession was the banner of the Church in all her wrestlings and conflicts, the Westminster Confession but as the camp colors which she hath used during her days of peace—the one for battle, the other for fair appearance and good order. This document consisteth of twenty-five articles, and is written in a most honest, straightforward, manly style, without compliment or flattery, without affectation of logical precision and learned accuracy, as if it came fresh from the heart of laborious workmen, all the day long busy with the preaching of the truth, and sitting down at night to embody the heads of what was continually taught. There is a freshness of life about it which no frequency of reading wears off.'

\xsubsection{The Scotch Covenants and the Scotch Kirk.}

§ 90. The Scotch Covenants and the Scotch Kirk.

Literature.

The Covenants are added to some Scotch editions of the Westminster Standards. The Solemn League and Covenant was often separately printed.

James Aikman:  \emph{An Historical Account of Covenanting in Scotland, from the first Band in Mearns,} 1556, \emph{to the Signature of the Grand National Covenant,} 1638. Edinburgh, 1848 (82 pp.).

National Covenants or politico-religious agreements for the maintenance and defense of certain principles and privileges are a peculiar and prominent feature in the history of the Kirk of Scotland. They were copied from Jewish precedents.\footnote{Josh. xxiv. 25: 'So Joshua made a covenant with the people that day, and set them a statute and an ordinance at Shechem;' 2 Kings xi. 17: 'And Jehoiada made a covenant between the Lord and the king and the people, that they should be the Lord's people;' also Isa. xliv. 5.} They originated in critical periods, when the sacred rights and convictions of the people were in imminent danger, and when the religious and national sentiments were inseparably blended. They are not properly confessions of faith, but closely connected with them, and must therefore be noticed here. They are solemn pledges to defend the doctrines and polity of the Reformed Kirk against all hostile attempts from within or from without, and to die rather than surrender.\footnote{Dr. M'Crie says of the Scotch Covenants (p. 120): 'Although they have been condemned as unwarranted in a religious point of view, and dangerous in a political, yet are they completely defensible upon the principles both of conscience and policy. A mutual agreement, compact, or covenant, is virtually implied in the constitution of every society, civil or religious; and the dictates of natural law conspire with the declarations of revelation in sanctioning the warrantableness and propriety of explicit engagements, about any lawful and important matter, and of ratifying these, if circumstances shall require it, by formal subscription, and by a solemn appeal to the searcher of hearts. By strengthening the motives to fidelity and constancy, and thus producing mutual confidence among those who are embarked in the same cause, they have proved eminently beneficial in the reformation of churches and nations, and in securing the religious and political privileges of men. The misapplication of them, when employed in a bad cause and for mischievous ends, can be no argument against their use in a legitimate way, and for laudable purposes. And the reasoning employed to prove that such covenants should not be entered into without the permission of rulers would lead to the conclusion that subjects ought never to profess a religion to which their superiors are hostile, nor make any attempts to obtain the reform of abuses, or the redress of grievances, without the consent and approbation of those who are interested in their support.' From Scotland the custom of covenanting passed to the Puritans in England and New England, and remains to-day in the shape of solemn engagements assumed by individual Christians when they enter into full communion with a church. Such covenants take the place of confirmation vows customary in the Lutheran and Anglican Churches.}

The earlier Covenants were safeguards against popery, the later against episcopacy. In the ecclesiastical history of Scotland since the Reformation we may distinguish three main periods: the period of anti-popery (1560 to 1590), the period of anti-prelacy (until 1690), and the period of anti-patronage (until 1875).

The first Covenants were made for mutual protection against the Romanists by a number of Protestant nobles and gentlemen, at Mearns, 1556, at Edinburgh, Dec. 3, 1557, at Perth, Dec. 31, 1559, before the Reformed Kirk was properly organized.

THE NATIONAL COVENANT, 1581 AND 1638.

Far more important is the 'National Covenant,' or the 'Second Scotch Confession,' also called the 'King's Confession,' and the 'Negative Confession.' It was drawn up in English and Latin by the Rev. John Craig, a noble, well-educated, and devoted man, a colleague of Knox and author of two Catechisms, who, after an eventful and romantic career, died in 1600, in the eighty-eight year of his life. It is a solemn indorsement of the Confession of Faith of 1560, with the strongest possible protest against 'all kind of papistry in general and particular heads,' especially against the 'usurped tyranny of the Roman Antichrist upon the Scriptures of God, upon the Kirk, the civil magistrate, and consciences of men; all his tyrannous laws made upon indifferent things, against our Christian liberty; . . . his five bastard sacraments, with all his rites, ceremonies, and false doctrine added to the ministration of the true sacraments without the Word of God; his cruel judgment against infants departing without the sacrament;\footnote{This is the first \emph{confessional} declaration against the damnation, and, by implication, in favor of the salvation, of unbaptized infants; and agrees with the \emph{private} opinion previously expressed by Zwingli and Bullinger.} his absolute necessity of baptism; his blasphemous opinion of transubstantiation; his devilish mass; his blasphemous priesthood; his profane sacrifice for sins of the dead and the quick; . . . his worldly monarchy and wicked hierarchy; his three solemn vows; his erroneous and bloody decrees made at Trent, with all the subscribers and approvers of that cruel and bloody band conjured against the Kirk of God.' No other Protestant Confession is so fiercely anti-popish.

This Covenant was subscribed by King James VI., his household, and a number of nobles and ministers, at Edinburgh, Jan. 28, 1581 (or 1580, old style\footnote{'They did not begin the year in Scotland, at that time, till the 25th of March.'—Dunlop's \emph{Collection,} Vol. II. p. 101.}); then by the Assembly and by persons of all ranks in March, 1581; again in 1590, together with a 'General Band for Maintenance of the True Religion and the King's Person or Estate;' it was solemnly renewed, with additions, in 1638 and 1639; ratified by an Act of Parliament in 1640, and signed by King Charles II., in exile, at Spey, June 23, 1650, and again when he was crowned at Scone, in Scotland, Jan. 1, 1651.\footnote{See the text in Vol. III. p. 480; and in Calderwood, Vol. III. p. 502. Calderwood thinks (p. 505) that this confession, under the name of 'wicked hierarchy,' condemns \emph{episcopal} government; but it is evident from the context that the \emph{papal} hierarchy is meant.}

The renewal of the Covenant in 1638, which is more particularly called the National Covenant, marks the Second Reformation. It includes the old Covenant of 1581, the Acts of Parliament condemning popery, and a protest against the government of the Kirk by bishops, all those measures of King Charles I. which 'do sensibly tend to the re-establishment of the Popish religion and tyranny, and to the subversion and ruin of the true Reformed religion, and of our liberties, laws, and estates.' The additions were prepared by Alexander Henderson and Johnston of Warriston, to meet a great crisis.\footnote{See the additions in Dunlop's \emph{Collection,} Vol. II. pp. 125–137, also the Acts of the Assemblies of Glasgow, 1638, and Edinburgh, 1639, pp. 114 sqq.}

The introduction of the semi-presbyterian mongrel episcopacy of James was comparatively harmless. But when his son Charles and his spiritual adviser, Archbishop Laud, in criminal ignorance or contempt of public feeling, attempted to force upon the Scots the royal supremacy, with a Romanizing hierarchy and liturgy, it produced a revolution and civil war which extended to England, and culminated in the temporary triumph of Puritanism. Macaulay traces the freedom of England to this 'act of insane bigotry.' In 1633 Laud displayed the most elaborate pomp of ceremonial worship in Holyrood Chapel to impress the descendants of John Knox! His new service-book differed from the English in a marked tendency to popery. When it was first introduced, July 23, 1637, in the cathedral church of St. Giles, in the presence of the privy council, the two archbishops of Scotland, several bishops, and the city magistrates, a poor old woman, named Jenny Geddes, confounding 'colic' and 'collect,' indignantly exclaimed, 'Villain, dost thou say mass at my lug,' and hurled her famous stool at the head of the unfortunate dean, who read 'the black, popish, and superstitious book.' Instantly all was uproar and confusion all over the city. The people shouted through the streets, 'A pope, a pope! Antichrist! The sword of the Lord and Gideon!' The unpremeditated riot extended into a popular revolution. The result was the overthrow of the artificial scheme which bigotry and tyranny had concocted.\footnote{'Never,' says Dean Stanley (p. 82), 'except in the days of the French Revolution, did a popular tumult lead to such important results. The stool which was on that occasion flung at the head of the Dean of Edinburgh extinguished the English Liturgy entirely in Scotland for the seventeenth century, to a great extent even till the nineteenth, and gave to the civil war of England an impulse which only ended in the overthrow of the Church and monarchy.'}

The renewal of the Covenant took place in Greyfriars' Church, in Edinburgh, the 28th of February, 1638, and was a most solemn and extraordinary scene. No less than sixty thousand people flocked to the city from all parts of the kingdom. The dense crowd which filled the church and adjoining graveyard listened with breathless attention to the prayers, the addresses, and the reading of the Covenant. The aged Earl of Sutherland first signed his name with trembling hand upon the parchment roll. Name followed name in swift succession. 'Some wept aloud; some burst into a shout of exultation; some, after their names, added the words, \emph{till death;} and some, opening a vein, subscribed with their own warm blood. As the space became filled, they wrote their names in a contracted form, limiting them at last to the initial letters, till not a spot remained on which another letter could be inscribed. . . . Never, except among God's peculiar people, the Jews, did any national transaction equal in moral and religious sublimity that which was displayed by Scotland on the great day of her sacred National Covenant.'\footnote{Hetherington, \emph{History of the Church of Scotland,} p. 91 (3d ed.).}

Similar scenes were repeated throughout the Northern Kingdom. Noblemen and gentlemen carried copies of the Covenant in their pockets and portmanteaus, soliciting subscriptions. Women sat in church day and night, from Friday till Sunday, to receive the communion with it. To refuse signature seemed to some denial of Christianity itself.\footnote{For fuller particulars, see Baillie's \emph{Letters,} Vol. I., Rothes's \emph{Relation,} Aiton's \emph{Life of Henderson,} Burton (Vol. VI. p. 442). Accounts from the episcopal side, in Thomas Stephen's \emph{History of the Church of Scotland,} Vol. I. pp. 552 sqq.; Stanley, 1.c. pp. 80 sqq.}

THE SOLEMN LEAGUE AND COVENANT, 1643.

'The Solemn League and Covenant for Reformation and Defense of Religion, the Honor and Happiness of the King, and the Peace and Safety of the Three Kingdoms of Scotland, England, and Ireland,' is the last and the most important of these national compacts which grew out of the Reformation. It has neither the doctrinal import nor the ring and fervor of the National Covenant of 1580 and 1638, but it had a wider scope and greater effect. It is anti-episcopal as well as anti-papal. It is the connecting link between Scotch Presbyterianism and English Puritanism, between the General Assembly and the Westminster Assembly, between the Scotch Parliament and the Long Parliament. It aimed to secure uniformity of religion in the united realms, while the National Covenant, like the Confession of 1560, was purely Scotch, and never exceeded its original boundary.\footnote{It is surprising that these two Covenants should be confounded by such a scholar as Dean Stanley, in his eloquent description of it, in \emph{Lectures on the Church of Scotland,} pp. 83–86 (Am. ed.). Dean Hook makes the same mistake—\emph{Life of Laud,} p. 267.}

We present first the text in full:\footnote{From 'The [Westminster] Confession of Faith, the Larger and Shorter Catechisms, together with the Sum of Saving Knowledge, Covenants, National, and Solemn League,' etc. Edinburgh, 1788, pp. 501 sqq. Masson, in his \emph{Life of Milton,} Vol. III. p. 13, gives the essential parts of the National Covenant. Fuller inserts it in full, Vol. VI. p. 255, and compares it (p. 259) to 'the superstitious and cruel Six Articles enacted by King Henry VIII.' Comp. Baillie's \emph{Letters,} Vol. II. pp. 81–90; the Acts of the General Assembly for 1643; Stoughton, \emph{The Church of the Civil Wars,} pp. 293 and 320; Masson, 1.c. Vol. III. pp. 6–15; Hetherington, l.c. pp. 110 sqq.}

'We Noblemen, Barons, Knights, Gentlemen, Citizens, Burgesses, Ministers of the Gospel, and Commons of all sorts, in the kingdoms of Scotland, England, and Ireland, by the providence of God, living under one King, and being of one reformed religion, having before our eyes the glory of God and the advancement of the kingdom of our Lord and Saviour Jesus Christ, the honor and happiness of the King's Majesty and his posterity, and the true public liberty, safety, and peace of the kingdoms, wherein every one's private condition is included: And calling to mind the treacherous and bloody plots, conspiracies, attempts, and practices of the enemies of God, against the true religion and professors thereof in all places, especially in these three kingdoms, ever since the reformation of religion; and how much their rage, power, and presumption are of late and at this time increased and exercised, whereof the deplorable state of the Church and kingdom of Ireland, the distressed estate of the Church and kingdom of England, and the dangerous estate of the Church and kingdom of Scotland are present and public testimonies; we have now at last (after other means of supplication, remonstrance, protestation, and sufferings, for the preservation of ourselves and our religion from utter ruin and destruction, according to the commendable practice of these kingdoms in former times, and the example of God's people in other nations), after mature deliberation, resolved and determined to enter into a mutual and Solemn League and Covenant, wherein we all subscribe, and each one of us for himself, with our hands lifted up to the most High God, do swear,

'I. That we shall sincerely, really, and constantly, through the grace of God, endeavor, in our several places and callings, the preservation of the reformed religion in the Church of Scotland, in doctrine, worship, discipline, and government, against our common enemies; the reformation of religion in the kingdoms of England and Ireland, in doctrine, worship, discipline, and government, according to the Word of God and the example of the best Reformed Churches; and shall endeavor to bring the Churches of God in the three kingdoms to the nearest conjunction and uniformity in religion, confession of faith, form of Church government, directory for worship and catechising; that we, and our posterity after us, may, as brethren, live in faith and love, and the Lord may delight to dwell in the midst of us.

'II. That we shall, in like manner, without respect of persons, endeavor the extirpation of Popery, Prelacy (that is, Church government by Archbishops, Bishops, their Chancellors and Commissaries, Deans, Deans and Chapters, Archdeacons, and all other ecclesiastical Officers depending on that hierarchy), superstition, heresy, schism, profaneness, and whatsoever shall be found to be contrary to sound doctrine and the power of godliness; lest we partake in other men's sins, and thereby be in danger to receive of their plagues; and that the Lord may be one, and his name one, in the three kingdoms.

'III. We shall, with the same sincerity, reality, and constancy, in our several vocations, endeavor, with our estates and lives, mutually to preserve the rights and privileges of the Parliaments, and the liberties of the kingdoms; and to preserve and defend the King's Majesty's person and authority, in the preservation and defense of the true religion and liberties of the kingdoms; that the world may bear witness with our consciences of our loyalty, and that we have no thoughts or intentions to diminish his Majesty's just power and greatness.

'IV. We shall also, with all faithfulness, endeavor the discovery of all such as have been or shall be incendiaries, malignants, or evil instruments, by hindering the reformation of religion, dividing the King from his people, or one of the kingdoms from another, or making any faction or parties amongst the people, contrary to this League and Covenant; that they may be brought to public trial, and receive condign punishment, as the degree of their offenses shall require or deserve, or the supreme judicatories of both kingdoms respectively, or others having power from them for that effect, shall judge convenient.

'V. And whereas the happiness of a blessed peace between these kingdoms, denied in former times to our progenitors, is, by the good providence of God, granted unto us, and hath been lately concluded and settled by both Parliaments; we shall each one of us, according to our place and interest, endeavor that they may remain conjoined in a firm peace and union to all posterity, and that justice may be done upon the willful opposers thereof, in manner expressed in the precedent article.

'VI. We shall also, according to our places and callings, in this common cause of religion, liberty, and peace of the kingdoms, assist and defend all those that enter into this League and Covenant in the maintaining and pursuing thereof; and shall not suffer ourselves, directly or indirectly, by whatsoever combination, persuasion, or terror, to be divided and withdrawn from this blessed union and conjunction, whether to make defection to the contrary part, or to give ourselves to a detestable indifferency or neutrality in this cause which so much concerneth the glory of God, the good of the kingdom, and honor of the king; but shall, all the days of our lives, zealously and constantly continue therein against all opposition, and promote the same according to our power against all lets and impediments whatsoever; and what we are not able ourselves to suppress or overcome we shall reveal and make known, that it may be timely prevented or removed: all which we shall do as in the sight of God.

'And, because these kingdoms are guilty of many sins and provocations against God and his Son Jesus Christ, as is too manifest by our present distresses and dangers, the fruits thereof, we profess and declare, before God and the world, our unfeigned desire to be humbled for our own sins, and for the sins of these kingdoms; especially, that we have not as we ought valued the inestimable benefit of the gospel; that we have not labored for the purity and power thereof; and that we have not endeavored to receive Christ in our hearts nor to walk worthy of him in our lives; which are the causes of other sins and transgressions so much abounding amongst us; and our true and unfeigned purpose, desire, and endeavor for ourselves, and all others under our power and charge, both in public and in private, in all duties we owe to God and man, to amend our lives, and each one to go before another in the example of a real reformation; that the Lord may turn away his wrath and heavy indignation, and establish these Churches and kingdoms in truth and peace.

'And this Covenant we make in the presence of Almighty God, the searcher of all hearts, with a true intention to perform the same, as we shall answer at that great day when the secrets of all hearts shall be disclosed; most humbly beseeching the Lord to strengthen us by his Holy Spirit for this end, and to bless our desires and proceedings with such success as may be deliverance and safety to his people and encouragement to other Christian Churches, groaning under, or in danger of, the yoke of anti-Christian tyranny, to join in the same or like association and covenant, to the glory of God, the enlargement of the kingdom of Jesus Christ, and the peace and tranquillity of Christian kingdoms and commonwealths.'

The immediate origin of this international politico-religious Covenant was the combined application of the English Parliament, then at war with King Charles I., and the Westminster Assembly of Divines, then sitting under its authority, for the effectual aid of the Scots, who occupied a position of neutrality. Six commissioners—four from the Parliament (Sir William Armyn, Sir Harry Vane the younger, Mr. Hatcher, and Mr. Darley) and two from the Westminster Assembly (Stephen Marshall and Philip Nye)—appeared with official and private letters before the Scotch Convention of Estates and the General Assembly at Edinburgh, in August, 1643. The English desired a civil league; the Scotch were for a religious covenant, and made the latter a condition of the former. Alexander Henderson, a highly esteemed minister at Edinburgh, Rector of the University (since 1640), and then for the third time Moderator of the General Assembly, was intrusted with the preparation of the document. He had drawn up a part of the National Covenant five years before. The English suggested some modifications which gave greater prominence to the political feature. The draft was unanimously and enthusiastically adopted by the General Assembly and the Scottish Convention, Aug. 17, 1643. The people, who had not forgotten the Covenant of 1638, manifested their most hearty approval, and went into the new engagement with the 'perfervidum ingenium Scotorum.'

The Solemn League and Covenant became a signal of war and victory in the history of Puritanism. It was followed by the appointment of Scotch commissioners to the Westminster Assembly, who took a leading part in the preparation of the Westminster standards of doctrine, worship, and discipline. It was debated for three or four days in that Assembly, and approved, with a few verbal alterations, by all the members except the Episcopalians. On the 21st of September Parliament ordered it to be published and subscribed throughout England. On the 25th of September the members of the House of Commons (two hundred and twenty-eight) and the divines of the Assembly set the example in St. Margaret's Church,\footnote{It is still used as a place of worship on special occasions by the Houses of Parliament.} beneath the shadow of Westminster Abbey. It was one of the strangest and most solemn events in the history of England. It reminds one of the formation of the Swiss Confederacy on the green meadow at Grütli. After prayer and addresses by White of Dorchester, Philip Nye, and Henderson, the Covenant was read, article by article, from the pulpit, and every member, standing up and lifting his right hand to heaven, took the pledge, and then signed his name on the rolls of parchment. The House of Lords followed a few weeks afterwards (Oct. 15). The same solemn scene was re-enacted in almost every English town and parish where the authority of Parliament prevailed. Cromwell among the Commons, and probably, also, Milton as a householder, signed the document, though Cromwell afterwards made war on the Scots, and Milton came to the conclusion that 'new Presbyter is but old Priest writ large.' In vain did the King, from his head-quarters in Oxford, forbid the League (Oct. 9), as 'a traitorous and seditious combination against himself and the established religion of his kingdom.' It became the shibboleth of Puritan religion and patriotism. There were, however, some exceptions. England, after all, was not so zealous for Presbyterianism as Scotland, and not used to covenanting. Richard Baxter raised his voice against the indiscriminate enforcement of the Covenant, and prevented its being taken in Kidderminster and the neighborhood.\footnote{Marsden (\emph{History of the Later Puritans,} p. 77): 'Such is the weight of character: one country clergyman prevailed against the rulers of two kingdoms.'}

From England the tide flowed back to Scotland, and Scotland, stimulated by the example, outran the neighboring country in zeal for the League. On the 13th of October, 1643, most of the nobles, including eighteen members of the Privy Council, solemnly signed it in Edinburgh, and from that day on for months there was 'a general swearing to the Covenant' by the people of Scotland, as by the Parliamentarians in England, from district to district, from city to city, from village to village, from parish to parish.\footnote{Stoughton, Vol. I. p. 294; Masson, Vol. III. pp. 12, 13.}

'O'er hill and dale the summons flew,

Nor rest nor pause the herald knew.

Each valley, each sequester'd glen,

Mustered its little horde of men,

That met, as torrents from the height,

In Highland dales, when streams unite,

Still gathering as they pour along,

A voice more loud, a tide more strong.'

On the 29th of November, 1643, the two countries entered into a treaty, by which the Scots promised to furnish an army for the war, the expenses to be refunded after the conclusion of peace. The Scots felt that they were playing the part of the good Samaritan towards the neighbor who had fallen among thieves. 'Surely,' says Baillie, 'it was a great act of faith in God, and huge courage and unheard-of compassion' on the part of the Scotch nation, 'to hazard their own peace and venture their lives and all, for to save a people so irrecoverably ruined, both in their own and in all the world's eyes.'

The united army fought under the banner of the Anglo-Scotch Covenant against royal and episcopal tyranny, and for the establishment of presbyterian uniformity. The negative end was gained, the positive failed. 'Trusting in God and keeping their powder dry,' the Puritans overthrew both monarchy and prelacy, but only to be overthrown in turn by the Nemesis of history. No human power could bring the two kingdoms under one creed and one form of government and worship. Presbyterian uniformity in England was as preposterous as Episcopal uniformity in Scotland.

The Solemn League and Covenant was weakened by the quarrel between the Presbyterians and Independents, and was virtually broken with the destruction of the monarchy and the execution of Charles I. (1648).\footnote{The Westminster Assembly, or what was left of it, sympathized with Presbyterian Scotland in loyalty to the monarchy, and unanimously signified its desire for the King's release. Forty-seven ministers, meeting at Sion College, signed a document addressed to Fairfax, in which they protested most earnestly in the name of religion and the Solemn League and Covenant against the military usurpation and the violence intended to the King's person. Masson, Vol. III. p. 716; Stoughton, Vol. I. p. 529.} The English army put down the Covenant which the Scotch army had set up. After the Restoration it became an object of intense hatred, and was publicly burned by the common hangman in Westminster Hall by order of Parliament (1661). Charles II., who had twice sworn both to the Solemn League and to the National Covenant as a part of his coronation oath in Scotland (June 23, 1650, and Jan. 1, 1651), broke his oath as soon as he ascended the English throne, and established the royal Supremacy and Episcopacy even in Presbyterian Scotland (1662). But the Covenanters fought for the institutions of their fathers with the heroic spirit of martyrdom through all those troubled times,

'Whose memory rings through Scotland to this hour.'

THE SCOTCH KIRK.

After severe struggles Prelacy was again overthrown and Presbyterianism permanently re-established in Scotland by Parliament in 1690, though with a degree of dependence on the state which kept up a constant irritation, and which led from time to time to new secessions.

These secessions from the Established Kirk, down to the great exodus of the Free Church in 1843, were no new departures, but, like the sects in Russia, returns to the old landmarks. The system of Calvinistic Presbyterianism which the great Reformer had established in Geneva found in Scotland a larger and more congenial field of action, and became there more free and independent of the civil power. It was wrought into the bone and sinew of the nation which seems to be predestinated for such a manly, sturdy, God-fearing, solid, persevering type of Christianity. Romanism in the Highlands is only an unsubdued remnant of the Middle Ages, lately reinforced by Irish emigrants to the large cities. Episcopacy is an English exotic for Scotchmen educated in England and associated with the English aristocracy. The body of the people are Presbyterian to the back-bone. The differences between the Established Kirk, the United Presbyterians, the Free Church, and the smaller secession bodies seem insignificant to an outside observer, and turn on questions of psalmody, patronage, and relation to the civil government. The vital doctrines and principles are held in common by all. Differences of opinion, which in other countries constitute merely theological schools or parties in one and the same denomination, give rise in Scotland to separate ecclesiastical organizations. The scrupulous conscientiousness and stubbornness which clothe minor questions with the dignity and grandeur of fundamental principles, and are made to justify separation and schism, are the shadow of a virtue. Scotland is an unconquerable fort of orthodox Protestantism. In no other country and Church do we find such fidelity and tenacity; such unswerving devotion to the genius of the Reformation; such union of metaphysical subtlety with religious fervor and impetuosity; such general interest in ecclesiastical councils and enterprizes; such jealousy for the rights and self-government of the Church; such loyalty to a particular denomination combined with a generous interest in Christ's kingdom at large; such reverence for God's holy Word and holy day, that after the hard and honest toil of the week lights up the poorest man's cottage on 'Saturday night.'

The history of Christianity, since the days of the apostles, furnishes no brighter chapter of heroic and successful sacrifices for the idea of the sole headship of Christ, and the honor and independence of his Church, than the Free-Church movement, whose leaders—Chalmers, Welsh, Candlish, Cunningham, Duncan, Fairbairn, Guthrie, Buchanan, Arnot—have now one by one taken their place among the great and good men of the past, but will continue to live in the memory of a grateful people. Dr. Norman Macleod, himself one of the noblest of Scotchmen, who was a member of the disruption Assembly of 1843, and found it harder to stay in the Established Church as 'a restorer of the breach' than to go out of it amid the huzzas of popular enthusiasm, honored himself as much as Dr. Chalmers, his teacher, when he spoke of him after his death as a man 'whose noble character, lofty enthusiasm, and patriotic views will rear themselves before the eyes of posterity like Alpine peaks, long after the narrow valleys which have for a brief period divided us are lost in the far distance of past history.'\footnote{\emph{Memoir of Norman Macleod, by his Brother,} 1876, Vol. I. p. 263 (N. Y. ed.).} In securing liberty for itself, the Free Church conferred a blessing upon the mother Church by rousing it to greater activity, and setting in motion an agitation which resulted in the total abolition of the Law of Patronage by Act of Parliament (1875).

\xsubsection{The Scotch Catechisms.}

§ 91. The Scotch Catechisms.

Catechetical instruction became soon after the Reformation, and remains to this day, one of the fundamental institutions of Presbyterian Scotland, and accounts largely for the general diffusion of religious information among the people.

The First Book of Discipline, adopted in 1560, prescribes public catechising of the children before the people on Sunday afternoon. The General Assembly of 1570 ordered ministers and elders to give to all the children within their parishes three courses of religious instruction—when they were nine, twelve, and fourteen years of age. Later assemblies enacted similar laws, and enjoined it also upon the heads of families to catechise their children and servants. The Assembly of 1649 renewed the act of the Assembly of 1639 'for a day of weeklie catechising, to be constantly observed in every kirk.'\footnote{\emph{Book of Discipline,} ch. xi. sect. 3; \emph{Buik of Universal Kirk,} p. 121 (Peterkin's edition); Horatius Bonar, \emph{Catechisms of the Scottish Reformation} (London, 1866), Preface, p. xxxvii.}

The older Catechisms, both domestic and foreign, contain the same system of doctrine in a fresher though less logical form than the Westminster standards, by which they were superseded after the middle of the seventeenth century. 'Our Scottish Catechisms,' says Dr. Bonar, the hymnist,' though gray with the antiquity of three centuries, are not yet out of date. They still read well, both as to style and substance; it would be hard to amend them, or to substitute something better in their place. Like some of our old church-bells, they have retained for centuries their sweetness and amplitude of tone unimpaired. It may be questioned whether the Church has gained any thing by the exchange of the Reformation standards for those of the seventeenth century. . . . In the Reformation we find doctrine, life, action nobly blended. Between these there was harmony, not antagonism; for antagonism in such cases can only arise when the parts are disproportionately mingled. Subsequently the balance was not preserved: the purely dogmatical preponderated. This was an evil, yet an evil not so easily avoided as some think; for, as the amount of error flung upon society increased, the necessity for encountering it increased also; controversy spread, dialectics rose into repute, and the dogmatical threatened to stifle or dispossess the vital.'\footnote{L.c. p. viii.}

FOREIGN CATECHISMS.

The Catechism of Calvin and the Palatinate or Heidelberg Catechism were approved by the Church of Scotland, and much used in the sixteenth and seventeenth centuries.\footnote{See both in Dunlop's and Bonar's \emph{Collections.} Comp. above, pp. 467 and 537 sq.}

An English edition of the former by the translators of the Geneva Bible appeared first at Geneva, 1556, for the use of the congregation of exiles, of which Knox was pastor, and then at Edinburgh, 1564. The latter was printed in Edinburgh, 1591, 1615, and 1621.

NATIVE CATECHISMS.

The number of these must have been very large. King James remarked at the Hampton Court Conference that in Scotland every son of a good woman thought himself competent to write a Catechism. We mention only those which had ecclesiastical sanction:

1. Two Catechisms of John Craig (1512–1600), an eminent minister at Aberdeen, and then at Edinburgh.\footnote{Both in Bonar, pp. 187–285. The Shorter Catechism is also printed in Dunlop's \emph{Collection,} Vol. II. pp. 365–377.} He was the author of the Second Scotch Confession.\footnote{See p. 686; Calderwood, Vol. III. p. 354; M'Crie, \emph{J. Knox,} pp. 236 sqq.}

The Larger Catechism of Craig was first printed in Edinburgh, by Henrie Charteris, in 1581, and in London, 1589. The General Assembly of 1590 indorsed it, and ordered an abridgment by the author, which was approved and published in 1591. In this shorter form it was generally used till superseded by the Westminster Catechism. The author says in the Preface (dated July 20, 1581): 'First, I have abstained from all curious and hard questions; and, next, I have brought the questions and the answers to as few words as I could, and that for the ease of children and common people, who can not understand nor gather the substance of a long question or a long answer confirmed with many reasons.' The work begins with some historical questions, and then explains the Apostles' Creed, the Ten Commandments, and the Lord's Prayer, and ends with the means of grace and the way of salvation. The questions and answers are short, and of almost equal length. We give some specimens from the larger work; which is little known:

\emph{First Questions.}

\emph{Ques.} Who made man and woman?

\emph{Ans}. The eternal God of his goodness.

\emph{Ques.} Whereof made he them?

\emph{Ans.} Of an earthly body and an heavenly spirit.

\emph{Ques.} To whose image made he them?

\emph{Ans.} To his own image.

\emph{Ques.} What is the image of God?

\emph{Ans.} Perfect uprightness in body and soul.

\emph{Ques.} To what end were they made?

\emph{Ans.} To acknowledge and serve their Maker.

\emph{Ques.} How should they have served him?

\emph{Ans.} According to his holy will.

\emph{Ques.} How did they know his will?

\emph{Ans.} By his Works, Word, and Sacraments.

\emph{Ques.} What liberty had they to obey his will?

\emph{Ans.} They had free will to obey and disobey.

\emph{Of the Sacraments.}

\emph{Ques.} What is a Sacrament?

\emph{Ans.} A sensible sign and seal of God's favor offered and given to us.

\emph{Ques.} To what end are the Sacraments given?

\emph{Ans.} To nourish our faith in the promise of God.

\emph{Ques.} How can sensible signs do this?

\emph{Ans.} They have this office of God, not of themselves.

\emph{Ques.} How do the Sacraments differ from the Word?

\emph{Ans.} They speak to the eye, and the Word to the ear.

\emph{Ques.} Speak they other things than the Word?

\emph{Ans.} No, but the same diversely.

\emph{Ques.} But the word doth teach us sufficiently?

\emph{Ans.} Yet the Sacraments with the Word do it more effectually.

\emph{Ques.} What, then, are the Sacraments to the Word?

\emph{Ans.} They are sure and authentic seals given by God.

\emph{Ques.} May the Sacraments be without the Word?

\emph{Ans.} No, for the Word is their life.

\emph{Ques.} May the Word be fruitful without the Sacraments?

\emph{Ans.} Yes, no doubt, but it worketh more plenteously with them.

\emph{Ques.} What is the cause of that?

\emph{Ans.} Because more senses are moved to the comfort of our faith.

\emph{Baptism. }

\emph{Ques.} What is the signification of baptism?

\emph{Ans. }Remission of our sins and regeneration.

\emph{Ques. }What similitude hath baptism with remission of sins?

\emph{Ans. }As washing cleanseth the body, so Christ's blood our souls.

\emph{Ques. }Wherein doth this cleansing stand?

\emph{Ans. }In putting away of sin, and imputation of justice.

\emph{Ques. }Wherein standeth our regeneration?

\emph{Ans. }In mortification and newness of life.

\emph{Ques. }How are these things sealed up in baptism?

\emph{Ans. }By laying on of water.

\emph{Ques. }What doth the laying on of the water signify?

\emph{Ans. }Our dying to sin and rising to righteousness.

\emph{Ques. }Doth the external washing work these things?

\emph{Ans. }No, it is the work of God's Holy Spirit only.

\emph{Ques. }Then the sacrament is a bare figure?

\emph{Ans. }No, but it hath the verity joined with it.

\emph{Ques. }Do all men receive these graces with the Sacraments?

\emph{Ans. }No, but only the faithful.

\emph{The Lord's Supper.}

\emph{Ques. }What signifieth the Lord's Supper to us?

\emph{Ans. }That our souls are fed with the body and blood of Christ.

\emph{Ques. }Why is this represented by bread and wine?

\emph{Ans. }Because what the one doth to the body, the same doth the other to the soul spiritually.

\emph{Ques. }But our bodies are joined corporally with the elements, or outward signs?

\emph{Ans. }Even so our souls be joined spiritually with Christ his body.

\emph{Ques. }What need is there of this union with him?

\emph{Ans. }Otherwise we can not enjoy his benefits.

\emph{Ques. }Declare that in the Sacrament?

\emph{Ans. }As we see the elements given to feed our bodies, even so we see by faith Christ gave his body to us to feed our souls.

\emph{Ques. }Did he not give it upon the Cross for us?

\emph{Ans. }Yes, and here he giveth the same body to be our spiritual food, which we receive and feed on by faith.

\emph{Ques. }How receive we his body and blood?

\emph{Ans. }By our own lively faith only.

\emph{Ques. }What followeth upon this receiving by faith?

\emph{Ans. }That Christ dwelleth in us, and we in him.

\emph{Ques. }Then we receive only the tokens, and not his body?

\emph{Ans. }We receive his very substantial body and blood by faith.

\emph{Ques. }How can that be proved?

\emph{Ans. }By the truth of his Word, and nature of a Sacrament.

\emph{Ques. }But his natural body is in heaven?

\emph{Ans. }I no doubt, but yet we receive it in earth by faith.

\emph{Ques. }How can that be?

\emph{Ans. }By the wonderful working of the Holy Spirit.

\emph{Cause and Progress of Salvation.}

\emph{Ques. }Out of what fountain doth this our stability flow?

\emph{Ans}. Out of God's eternal and constant [unchanging] election in Christ.

\emph{Ques. }By what way cometh this election to us?

\emph{Ans. }By his effectual calling in due time.

\emph{Ques. }What worketh this effectual calling in us?

\emph{Ans. }The obedience of faith.

\emph{Ques. }What thing doth faith work?

\emph{Ans. }Our perpetual and inseparable union with Christ.

\emph{Ques. }What worketh this union with Christ?

\emph{Ans. }A mutual communion with him and his graces.

\emph{Ques. }What worketh this communion?

\emph{Ans. }Remission of sins and imputation of justice.

\emph{Ques. }What worketh remission of sins and imputation of justice?

\emph{Ans}. Peace of conscience and continual sanctification.

\emph{Ques. }What worketh sanctification?

\emph{Ans. }The hatred of sin and love of godliness.

2. A Latin Catechism, entitled \emph{Rudimenta Pietatis} and \emph{Summula Catechismi,} for the use of grammar schools.\footnote{In Dunlop's \emph{Collection,} Vol. II. pp. 378–382, and in Bonar, pp. 289–293.} It is ascribed to Andrew Simpson, who was master of the grammar school at Perth, and the first Protestant minister at Dunbar. It was used in the high-school at Edinburgh down to 1710.

Besides this, the Latin editions of the Heidelberg Catechism and Calvin's Catechism (translated by Patrick Adamson) were also in use.

3. The Catechism of John Davidson, minister at Salt-Preston, approved by the Provincial Assembly of Lowthiane and Tweddale, 1599.\footnote{Bonar, p. 324.}

4. A metrical Catechism by the Wedderburns in the time of Knox.\footnote{Bonar, p. 301.} The sentiment is better than the poetry. The Reformation in Scotland, as well as in France and Holland, called forth metrical versions of the Psalms, while in Germany it produced original hymns. The gospel was sung as well as preached into the hearts of the common people. But a Catechism is for instruction, and requires plain, clear, precise statements for common comprehension.

\xsection{The Westminster Standards.}

\xsubsection{The Puritan Conflict.}

§ 92. The Puritan Conflict.

Literature.

1. Sources.

1. The Parliamentary Acts, the Minutes and Standards of the Westminster Assembly, the royal Proclamations, Cromwell's Letters, Milton's state papers, and other public documents. See the \emph{State Calendars}; Rushworth's \emph{Collection} (1616–1648); Cardwell's \emph{Documentary Annals of the Church of England} (1546–1716); Camden's \emph{Annals of James I.} (with the king's own works); Winwood's \emph{Memorials of State}; and the literature mentioned in § 93 and § 94.

2. The private writings of the Episcopal and Puritan divines during the reigns of Elizabeth and the Stuarts, too numerous even to classify. Much material for history may be drawn from the works of Archbishop Laud (b. 1573, beheaded 1645), especially his \emph{Diary} (in the first vol. of his \emph{Remains}, publ. by H. Wharton, 1695–1700, in 2 vols. fol., and in the \emph{Anglo-Catholic Library}, Oxford, 1847–1850, 5 vols.), and of Richard Baxter (1615–1691), especially in the \emph{Narrative of his Life and Times} (publ. by Sylvester, 1696, under the title \emph{Reliquiæ Baxterianæ}, in 1 vol. fol., and by Dr. Calamy, 1713, in 4 vols., and in ed. of his \emph{Practical Works}, Lond. 1830, 23 vols. Baxter's numerous controversial tracts have never been collected, and have gone, with his medical prescriptions, to 'everlasting rest,' but his practical works will last). Mrs. Lucy Hutchinson's \emph{Memoirs of }(her husband) \emph{Colonel Hutchinson, with Original Anecdotes of many of his most Distinguished Contemporaries, and a Summary Review of Public Affairs} (publ. from MS. 7th ed. Lond. 1848), present an admirable picture of the inner and private life of the Puritans.

3. Innumerable controversial pamphlets and tracts for the times, which did the work of the newspapers of to-day. From 1640 to 1660 no less than 30,000 pamphlets on Church government alone are said to have appeared. Milton's tracts surpass all others in eloquence and force.

2. Historical.

Thomas Fuller (1608–1661, Prebendary of Sarum): \emph{The Church History of Britain, from the Birth of Christ until the Year} 1648. Ed. of Brewer, Oxford, 1845, in 6 vols. (Vols. V. and VI.).

Clarendon (1608–1674, Royalist and Episcopalian): \emph{History of the Rebellion.} Oxford ed. 1839 and 1849, 7 vols.

Daniel Neal (1678–1743, Independent): \emph{History of the Puritans, or Protestant Nonconformists, from the Reformation in }1517 \emph{to the Revolution in }1688. Lond. 1732; Toulmin's ed. 1793, 5 vols.; Choules's ed. New York (Harpers), 1858, in 2 vols.

J. B. Marsden (Vicar of Great Missenden): \emph{The History of the Early Puritans, from the Reformation to the Opening of the Civil War in }1642. Lond. 1850, 2d ed. 1853. By the same: \emph{The History of the Later Puritans, from the Opening of the Civil War in }1642 \emph{to the Ejection of the Nonconforming Clergy in }1662. Lond. 1852.

Hallam: \emph{Constitutional History of England}, 5th ed. ch. vii.–xi.

Th. Carlyle:  \emph{Life and Letters of Cromwell.} Lond. and New York, 1845, 2 vols. ('Edited with the care of an antiquarian and the genius of a poet.'—Green, \emph{Hist. of the English People}, p. 530.)

Guizot's French works on \emph{Charles I. }(1625–1649, 2 vols.), \emph{Cromwell }(1649–1658), the \emph{Re-establishment of the Stuarts }(1658–1660, 2 vols.), \emph{Monk} (1660, transl. by Scoble, 1851), the \emph{English Revolution of }1640 (transl. by Hazlitt, Lond. 1856).

Samuel Hopkins: \emph{The Puritans during the Reigns of Edward VI. and Queen Elizabeth. }Boston, 1859–61, 3 vols.

Principal Tulloch (Scotch Presbyt.): \emph{English Puritanism and its Leaders: Cromwell, Milton, Baxter, Bunyan.} Lond. 1861.

Dr. John Stoughton (Independent): \emph{Ecclesiastical History of England }(during the Civil Wars, the Commonwealth, and the Restoration). Lond. 1867–1875, 5 vols. By the same: \emph{Church and State Two Hundred Tears ago. A History of Ecclesiastical Affairs in England from }1660 to 1663. Lond. 1862. By the same: \emph{Spiritual Heroes; or, Sketches of the Puritans} (Ch. VI., The Westminster Assembly, pp. 120 sqq.). Lond. 1850.

Dr. David Masson (Prof. of Rhetoric and English Lit. in the Univ. of Edinb.): \emph{The Life of John Milton: Narrated in connection with the Political, Ecclesiastical, and Literary History of his Times.} Lond. 1859 80, 6 vols. See Vol. II. (1871), Books III. and IV., and Vol. III. (1873), Books I., II., and III.

On the early history of New England Puritanism, see the well-known works of Palfrey, Bancroft, Felt; and Leonard Bacon's  \emph{Genesis of the New England Churches} (New York, 1874)

PROTESTANTISM AND CIVIL WARS.

The Reformation has often been charged by Roman Catholic writers with being the mother of the bloody civil wars which grew out of the close union of Church and State, and which devastated Europe for more than a century. But the fault is primarily on the side of Rome. Exclusiveness and intolerance are fundamental principles of her creed, and persecution her consistent practice wherever she has the power. In Italy and Spain Protestantism was strangled in its cradle. In Bohemia, Hungary, and Poland it was reduced to a struggling minority by the civil sword and the Jesuit intrigues. In France it barely escaped annihilation in the massacre of the night of St. Bartholomew, which the pope hailed with a \emph{Te Deum;} and after fighting its way to the throne, and acquiring the limited toleration of the Edict of Nantes, it was again persecuted almost to extermination by the most Catholic King Louis XIV. In Switzerland the war between the Catholic and Reformed Cantons, in which Zwingli fell, fixed the boundaries of the two religions on a basis of equality. Germany had to pass through the fearful ordeal of the Thirty-Years' War, which destroyed nearly one half of its population, but ended, in spite of the protest of the pope, with the legal recognition of the Lutheran and Reformed Confessions by the Treaty of Westphalia in 1648. The United Provinces of Holland came out victorious from the long and bloody struggle with the tyranny and bigotry of Spain. Scotland fought persistently and successfully against popery and prelacy. England, after the permanent establishment of the Reformation under Elizabeth, was shaken to the base by an internal conflict, not between Protestants and foreign Romanists, but between Protestants and native Romanizers, ultra-Protestant Puritans and semi-Catholic Churchmen.

This conflict marks the most important period in the Church history of that island; it called forth on both sides its deepest moral and religious forces; it made England at last the stronghold of constitutional liberty in Europe, and laid the foundations for a Protestant republic in America. The Puritans were the pioneers in this struggle in Old England, and the fathers of New England beyond the sea. As the blood of martyrs is the seed of the Church, so freedom is the sweet fruit of bitter persecution.

CHARACTER OF PURITANISM.

Puritanism—an honorable name, etymologically and historically, though originally given in reproach,\footnote{The name \emph{Puritans} (from \emph{pure}, as \emph{Catharists} from \textgreek{καθαρός}), or \emph{Precisians}, occurs first about 1564 or 1566, and was employed to brand those who were opposed to the use of priestly vestments, as the cap, surplice, and the tippet (but not the gown, which the Puritans and Presbyterians retained, as well as the Continental Protestant ministers). Shakspere uses the term half a dozen times, and always reproachfully (see Clarke's \emph{Shaksp. Concordance} and Schmidt's \emph{Shaksp. Lexicon}, s.v.). In the good sense, it denotes those who went back to the purity and simplicity of apostolic Christianity in faith and morals. Neal defines a Puritan to be 'a man of severe morals, a Calvinist in doctrine, and a Nonconformist to the ceremonies and discipline of the Church, though not totally separated from it'} like Pietism and Methodism—aimed at a radical purification and reconstruction of Church and State on the sole basis of the Word of God, without regard to the traditions of men. It was a second Reformation, as bold and earnest as the first, but less profound and comprehensive, and more radical in its antagonism to the mediæval Church. It was a revolution, and ran into the excesses of a revolution, which called forth, by the natural law of reaction, the opposite excesses of a reactionary restoration; but it differs from more recent revolutions by the predominance of the religious motive and aim. The English Puritans, the Scotch Covenanters, and the French Huguenots were alike spiritual descendants of Calvin, and represent, with different national characteristics, the same heroic faith and severe discipline. They were alike animated by the fear of God, which made them strong and free. They bowed reverently before his holy Word, but before no human authority. In their eyes God alone was great.

The Puritans were no separate organization or sect, but the advanced wing of the national Church of England, and at one time they became the national Church itself, treating their opponents as Nonconformists, as they had been treated by them before, and as they were treated afterwards in turn. Conformity and Nonconformity were relative terms, which each party construed in its own way and for its own advantage. The Puritan ministers were educated at Oxford and Cambridge, and had bishops, deans, and professors of theology among their leaders and sympathizers. Their intention was not to secede, but simply to reform still further the national Church in the interest of primitive purity and simplicity by legislative and executive sovereignty. The tyrannical measures of the ruling party drove them to greater opposition, and a large portion of them into complete independency and the advocacy of toleration and freedom. But originally they were as intolerant and exclusive as their opponents. The common error of both was that they held to a close union of Church and State, and aimed at one national Church, to which all citizens must conform.

ORIGIN AND PROGRESS OF THE CONTROVERSY.

'Nonconformity,' says Thomas Fuller in his quaint and pithy way, 'was conceived in the days of King Edward, born in the reign of Queen Mary (but beyond the sea, at Frankfort-on-the-Main), nursed and weaned in the reign of Elizabeth, grew up a youth or tall stripling under King James, and shot up under Charles I. to the full strength and stature of a man able not only to cope with, but to conquer the hierarchy, its adversary.'

The open conflict between Puritanism and High-Churchism dates from the closing years of the sixteenth century, but its roots may be traced to the beginning of the Reformation, which embraced two distinct tendencies—one semi-Catholic, conservative and aristocratic; the other anti-Catholic, radical and democratic.

The aristocratic politico-ecclesiastical movement, headed by the monarch and the bishops, grew out of the mediæval conflict of the English crown and Parliament with the foreign papacy, and effected under Henry VIII. the national independence of the English Church, and under Edward VI. a positive though limited reformation in doctrine and ritual.

The democratic religious movement, which sprang from the desire of the people after salvation and unobstructed communion with God and the Bible, had its forerunners in Wycliffe and the Lollards, and was nurtured by Tyndale's English Testament, the writings of the Continental Reformers, and the personal contact of the Marian exiles with Bullinger and Calvin. At first it was nearly crushed under Henry VIII., who would not even tolerate the circulation of the English Bible; but it gained considerable influence under Edward VI., passed through a baptism of blood under Mary, and became a strong party under Elizabeth. It included a number of bishops, pervaded the universities, and was backed by the sympathies of the common people as they were gradually weaned from the traditions of popery.

Under Edward VI. the martyr-bishop Hooper, of Gloucester, a friend of Bullinger, and one of the fathers of Puritanism, opened the ritualistic controversy by refusing to be consecrated in the sacerdotal garments, and to take the customary episcopal oath, which included an appeal to the saints. He was quieted by the representations of the young king, of Bucer, and Peter Martyr, who regarded those externals as things indifferent; but he continued to strive after 'an entire purification of the Church from the very foundation.'

Under Queen Mary the conflict continued in the prisons and around the fires of Smithfield, and was transferred to the Continent with the English exiles, such as Jewel, Grindal, Sandys, Pilkington, Parkhurst, Humphrey, Sampson, Whittingham, Coverdale, Cox, Nowel, Foxe, Horn, and Knox. It produced an actual split in the congregation at Frankfort-on-the-Main. There it turned on the question of the Prayer-Book of Edward VI., whether it should be adhered to, or reformed still further after the model of the simpler worship of Zurich and Geneva. The episcopal and liturgical party was led by Dr. Cox (afterwards bishop of Ely), and formed the majority; the Puritan party was headed by John Knox, who was required to leave, and organized another congregation of exiles at Geneva.

After the accession of Elizabeth both parties flocked back to their native land, and forgot the controversy for a while in the common zeal for the re-establishment of Protestantism. As long as the ruling powers favored the Reformation the Puritans were satisfied, and heartily co-operated in every step. Though badly treated by the proud queen, they were to the last among her most loyal subjects, and prayed even in their dungeons for her welfare. They overlooked her faults for her virtues. They were the strongest supporters of the government and the crown against popish plots and foreign aggression, and helped to defeat the Spanish Armada, whose 'proud shipwrecks' were scattered over 'the Northern Ocean even to the frozen Thule.' But when the anti-Romish current stopped, and the Church of England seemed to settle down in a system of compromise between Rome and Geneva, fortified and hedged in by a cruel penal code against every dissent, the radicals assumed an antagonistic attitude of nonconformity against the rigorous enforcement of conformity, and stood up for the rights of conscience and the progress of ecclesiastical reform.

The controversy was renewed in different ways, between Cartwright and Whitgift, and between Travers and Hooker. In both cases the combatants were unequally matched: Cartwright, the father of Presbyterianism, was a much abler man than Archbishop Whitgift, the father of High-Church episcopacy; while Hooker, the Master of the Temple, far excelled Travers, the Lecturer at the Temple, in learning and depth. Here the question was chiefly whether the Scriptures as interpreted by private judgment, or the Scriptures as interpreted by the fathers of the primitive Church, should be the rule of faith and discipline. With this was connected another question—whether the Roman Church had lost the character of a Christian Church, and was therefore to be wholly disowned, or whether she was still a true though corrupt Church, with valid ordinances, coming down through an unbroken historical succession. The Puritans advocated Scripture Christianity \emph{versus} historical Christianity, Hooker historical Christianity as consistent with Scripture Christianity. But in substance of doctrine both parties were Augustinians and Calvinists, with this difference, that the Puritans were high Calvinists, the Churchmen low Calvinists. Whitgift advocated even the Lambeth Articles, and Hooker adopted them with some modifications. Arminianism did not make its appearance in England till the close of the reign of James.

THE HAMPTON COURT CONFERENCE.

The accession of James I. (1603-1625) marks a new epoch. He was no ordinary man. His learning ranged from the mysteries of predestination to witchcraft and tobacco; he had considerable shrewdness, mother-wit, ready repartee, and uncommon sense, but little common-sense, and no personal dignity nor moral courage; he was given to profanity, intemperance, and dissimulation. His courtiers and bishops lauded him as the Solomon of his age, but Henry IV. of France characterized him better as 'the wisest fool in Christendom.' He was brought up in the school of Scotch Presbyterianism, subscribed the Scotch Confession, and once said of the Anglican liturgy that 'it is an ill-said mass in English.' But the Stuart blood was in him, and when he arrived in England he felt relieved of his tormentors, who 'pulled his sleeve as they administered their blunt rebukes to him,' and was delighted by the adulation of prelates who had much higher notions of royalty than Scotch presbyters.

He lost no time in showing his true character. He answered the famous Millenary (or Millemanus) petition, signed by nearly a thousand Puritan ministers, and asking for the reform of certain abuses and offenses in worship and discipline,\footnote{Fuller, Vol. V. pp. 305–309. The petition was dated January 14, 1603 (old style), but was presented April 4. The real number of signers was only 825.} by the imprisonment of ten petitioners on the ground that their act tended to sedition and treason, although it contained no demand inconsistent with the established Church. Thus the opportunity for effecting a compromise was lost. He agreed, however, to a Conference, which suited his ambition for the display of his learning and wit in debate.

The Conference was held January 14, 16, and 18, 1604 (old style, 1603), at Hampton Court. The persons summoned were nine bishops, headed by Archbishop Whitgift of Canterbury and Bishop Bancroft of London, and eight deans, on the part of the Conformists, and four of the most learned and moderate Puritan divines, under the lead of Dr. John Reynolds, President of Corpus Christi College, Oxford.\footnote{Fuller (Vol. V. pp. 378, 379) speaks in very high terms of Reynolds, who was so unceremoniously snubbed by Bishop Bancroft. He praises his memory, which was 'little less than marvelous,' and 'a faithful index,' as his reason was 'a solid judex of what he read,' and his humility, which 'set a lustre on all; communicative of what he knew to any that desired information herein, like a tree loaded with fruit, bowing down its branches to all that desired to ease it of the burden thereof, deserving this epitaph,  \par '\emph{Incertum est utrum doctior an melior.'} \par He associates him with Bishop Jewel and Richard Hooker, all born in Devonshire, and educated at Corpus Christi College, and says, ' No one county in England have three such men (contemporary at large), in what college soever they were bred; no college in England bred such three men, in what county soever they were born.' John Reynolds was at first a zealous papist and turned an eminent protestant; while his brother William was as earnest a protestant, and became by their mutual disputation an inveterate papist, which gave occasion to the distich: \par '\emph{Quod genus hoc pugnæ est? ubi victus gaudet uterque,} \par \emph{Et simul alteruter se superasse dolet.'} \par 'What war is this? when conquer'd both are glad, \par And either to have conquer'd other sad.'} The King himself acted both as moderator and judge, and took the leading part in the discussion. He laid down his famous pet-principle (which he called his 'aphorism'), 'No bishop, no king;'\footnote{He also said to Dr. Reynolds: 'If you aim at a Scotch presbytery, it agreeth as well with monarchy as God and the devil. Then Jack, and Tom, and Will, and Dick shall meet and censure me and my council. Therefore I reiterate my former speech, \emph{Le roy s’avisera.}'} and, after browbeating the Puritans, used as his final argument, 'I will make them conform themselves, or else I will harry them out of the land, or else do worse.'

Archbishop Whitgift was so profoundly impressed with the King's theological wisdom that be said, 'Undoubtedly your Majesty speaks by the special assistance of God's Spirit;' and Bishop Bancroft, of London (who first proclaimed the doctrine of a \emph{jure divino} episcopacy), thanked God on his knees that of his singular mercy he had given to them 'such a king, as since Christ's time the like hath not been.' The same haughty prelate rudely interrupted Dr. Reynolds, one of the most learned men in England, saying, 'May your Majesty be pleased that the ancient canon be remembered—\emph{Schismatici contra episcopos non sunt audiendi}; and there is another decree of a very ancient council, that no man should be admitted to speak against that whereunto he hath formerly subscribed. And as for you, Doctor Reynolds, and your associates, how much are ye bound to his Majesty's clemency, permitting you, contrary to the statute \emph{primo Elizabethæ}, so freely to speak against the liturgy or discipline established.'

Fuller remarks 'that the King in this famous Conference went beyond himself, that the Bishop of London (when not in a passion) appeared even with himself, and that Dr. Reynolds fell much beneath himself.' The Nonconformists justly complained that the King invited their divines, not to have their scruples satisfied, but his pleasure propounded—not to hear what they had to say, but to inform them what he would do. Hallam, viewing the Conference calmly from his stand-point of constitutional history, says: 'In the accounts that we read of this meeting we are alternately struck with wonder at the indecent and partial behavior of the King and at the baseness of the bishops, mixed, according to the custom of servile natures, with insolence toward their opponents. It was easy for a monarch and eighteen churchmen to claim the victory, be the merits of their dispute what they might, over four abashed and intimidated adversaries.'\footnote{The accounts of the Hampton Court Conference are mostly derived from the partial report of Dr. William Barlow, Dean of Chester, who was present. It appeared in 1604, and again in 1638. See Fuller, Vol. V. pp. 266–303; Cardwell, \emph{Hist. of Conferences}, p. 121; Procter, \emph{Hist. of the Book of Common Prayer}, p. 88; Marsden, \emph{Early Puritans}, p. 255.}

The Conference, however, had one good and most important result—the revision of our English Bible. The revision was suggested and urged by Dr. Reynolds, who was subsequently appointed one of the revisers,\footnote{He was assigned to the company which was charged with the translation of the writings of the greater and lesser Prophets. But he died in 1607, before the completion of the work.} and it was ordered to be executed by King James, from whom it has its name.\footnote{The discussion bearing upon this subject is likewise characteristic of the King, the Bishop, and the Puritan, and may be added here (from Fuller, Vol. V. pp. 284, 285):  \par '\emph{Dr. Reynolds.} ``May your Majesty be pleased that the Bible be new translated, such as are extant not answering the original.'' And he instanced in three particulars:     \emph{In the Original} \emph{Ill Translated}   'Gal. iv. 25  \textgreek{συστοιχεῖ} Bordereth.   Psalm cv. 28. They were not disobedient. They were not obedient.   Psalm cvi. 30. Phinehas executed judgment. Phinehas prayed.   \par '\emph{Bishop of London.} ``If every man's humor might be followed, there would be no end of translating.'' \par '\emph{His Majesty.} ``I profess I could never yet see a Bible well translated in English; but I think that of all, that of Geneva is the worst. I wish some special pains were taken for an uniform translation; which should be done by the best learned in both universities, then reviewed by the bishops, presented to the privy council, lastly ratified by royal authority to be read in the whole Church, and no other.'' \par '\emph{Bishop of London.} ``But it is fit that no marginal notes should be added thereunto.'' \par '\emph{His Majesty.} ``That caveat is well put in; for in the Geneva translation some notes are partial, untrue, seditious, and savoring of traitorous conceits: as when, from Exodus i. 19, disobedience to kings is allowed in a marginal note; and, 2 Chron. xv. 16, King Asa taxed in the note for only deposing his mother for idolatry, and not killing her. To conclude this point, let errors in matters of faith be amended, and indifferent things be interpreted, and a gloss added unto them; for, as Bartolus de Regno saith, 'Better a king with some weakness than still a change;' so rather a Church with some faults than an innovation. And surely, if these were the greatest matters that grieved you, I need not have been troubled with such importunate complaints.'''}

With all his high notions about royalty, James had not the moral courage to carry them into full practice, and with all his high notions about episcopacy, he had no sympathy with Arminianism, but actually countenanced the Calvinistic Presbyterian Synod of Dort, and sent five delegates to it, among them a bishop. In both these respects Charles went as far beyond James as Laud went beyond Whitgift and Bancroft.

KING CHARLES AND ARCHBISHOP LAUD.

The antagonism was intensified and brought to a bloody issue under Charles I. (1625-1649) and William Laud. They belong to the most lauded and the most abused persons in history, and have been set down by opposite partisans among the saints and among the monsters. They were neither. They were good men in private life, but bad men in public. They might have been as respected and useful in a humble station, or in another age or country, as they were hateful and hurtful at the helm of government in Protestant England. It was their misfortune rather than their crime that they were utterly at war with the progressive spirit of their age. Both were learned, cultured, devout gentlemen and churchmen, but narrow, pedantic, reactionary, haughty aristocrats. The one was constitutionally a tyrant, the other constitutionally a pope or an inquisitor-general. They fairly represented in congenial alliance the principle and practice of political and ecclesiastical absolutism, and the sovereign contempt for the rights of the people, whose sole duty in their opinion was passive obedience. Kingcraft and priestcraft based upon divine right was their common shibboleth. By their suicidal follies they destroyed the very system which they so long defended with a rod of iron, and thus they became the benefactors of Protestantism, which they labored to destroy. Both died as martyrs of despotism, and their last days were their best. 'Nothing in life became them as the leaving it.'

Charles wanted to rule without a Parliament; he did so, in fact, for more than eleven years, and the four Parliaments which he was compelled to convoke he soon arbitrarily dissolved (1625, 1626, 1629, and 1640). He preferred ship-money to legal taxation. He made himself intolerable by his duplicity and treachery. 'Faithlessness was the chief cause of his disasters, and is the chief stain on his memory. He was in truth impelled by an incurable propensity to dark and crooked ways. It may seem strange that his conscience, which on occasions of little moment was sufficiently sensitive, should never have reproached him with this great vice. But there is reason to believe that he was perfidious, not only from constitution and from habit, but also on principle. He seems to have learned from theologians whom he most esteemed that between him and his subjects there could be nothing of the nature of mutual contract; and that he could not, even if he would, divest himself of his despotic authority; and that in every promise which he made there was an implied reservation that such promise might be broken in case of necessity, and that of the necessity he was the sole judge.'\footnote{Macaulay, chap. i. p. 65 (Boston ed.). I add the admirable description of Charles by Mrs. Lucy Hutchinson, in the \emph{Memoirs} of her husband (Bohn's ed. p. 84): 'King Charles was temperate, chaste, and serious; so that the fools and bawds, mimics and catamites, of the former court, grew out of fashion; and the nobility and courtiers, who did not quite abandon their debaucheries, yet so reverenced the king as to retire into corners to practice them. Men of learning and ingenuity in all arts were in esteem, and received encouragement from the king, who was a most excellent judge and a great lover of paintings, carvings, gravings, and many other ingenuities, less offensive than the bawdry and profane abusive wit which was the only exercise of the other court. But, as in the primitive times, it is observed that the best emperors were some of them stirred up by Satan to be the bitterest persecutors of the Church, so this king was a worse encroacher upon the civil and spiritual liberties of his people by far than his father. He married a Papist, a French lady, of a haughty spirit, and a great wit and beauty, to whom he became a most uxorious husband. By this means the court was replenished with Papists, and many who hoped to advance themselves by the change turned to that religion. All the Papists in the kingdom were favored, and, by the king's example, matched into the best families; the Puritans were more than ever discountenanced and persecuted, insomuch that many of them chose to abandon their native country, and leave their dearest relations, to retire into any foreign soil or plantation where they might, amidst all outward inconveniences, enjoy the free exercise of God's worship. Such as could not flee were tormented in the bishops' courts, fined, whipped, pilloried, imprisoned, and suffered to enjoy no rest, so that death was better than life to them; and notwithstanding their patient sufferance of all these things, yet was not the king satisfied till the whole land was reduced to perfect slavery. The example of the French king was propounded to him, and he thought himself no monarch so long as his will was confined to the bounds of any law; but knowing that the people of England were not pliable to an arbitrary rule, he plotted to subdue them to his yoke by a foreign force, and till he could effect it, made no conscience of granting any thing to the people, which he resolved should not oblige him longer than it served his turn; for he was a prince that had nothing of faith or truth, justice or generosity, in him. He was the most obstinate person in his self-will that ever was, and so bent upon being an absolute, uncontrollable sovereign that he was resolved either to be such a king or none. His firm adherence to prelacy was not for conscience of one religion more than another, for it was his principle that an honest man might be saved in any profession; but he had a mistaken principle that kingly government in the State could not stand without episcopal government in the Church; and, therefore, as the bishops flattered him with preaching up his sovereign prerogative, and inveighing against the Puritans as factious and disloyal, so he protected them in their pomp and pride, and insolent practices against all the godly and sober people of the land.'}

William Laud\footnote{Born at Reading, Oct. 7, 1573; ordained 1601; Bishop of St. David's, 1621; of Bath and Wells, 1626; of London, 1628; Chancellor of Oxford University, 1630; Archbishop of Canterbury, 1633; impeached of high-treason, 1641; beheaded Jan. 10, 1645.} rose, like Cardinal Wolsey, by his abilities and the royal favor from humble origin to the highest positions in Church and State. He began his career of innovation early at Oxford, and asserted in his exercise for the degree of Bachelor of Divinity (1604) the absolute necessity of baptism for salvation, and the necessity of diocesan episcopacy, not only for the well-being, but for the very existence of the Church. This position exposed him to the charge of heresy, and no one would speak to him in the street. Under James he was kept back,\footnote{'Because,' as King James said, in keen discernment of his character, 'he hath a restless spirit, and can not see when matters are well, but loves to toss and change, and to bring things to a pitch of reformation, floating in his own brain, which may endanger the steadfastness of that which is in a good pass.' He restrained his early plans 'to make that stubborn [Scotch] Kirk stoop to the English pattern,' for 'he knows not the stomach of that people.'} but under Charles he rose rapidly, and after the death of Abbot, who was a Puritan, he succeeded him in the primacy of the English Church. When he crossed the Thames to take possession of Lambeth, he met with an ominous accident, which he relates in his \emph{Diary} (Sept. 18, 1633). The overloaded ferry-boat upset, and his coach sank to the bottom of the river, but he was saved as by water, and 'lost neither man nor horse.'

Laud was of small stature\footnote{\par He was called 'the little Archbishop.'} and narrow mind, but strong will and working-power, hot and irascible in temper, ungracious and unpopular in manner, ignorant of human nature, a zealous ritualist, a pedantic disciplinarian, and an overbearing priest. He was indefatigable and punctilious in the discharge of his innumerable duties as archbishop and prime minister, member of the courts of Star-Chamber and High-Commission, of the committee of trade, the foreign committee, and as lord of the treasury. He was for a number of years almost omnipotent and omnipresent in three kingdoms, looking after every appointment and every executive detail in Church and State.\footnote{'His influence extended every where, over every body, and every thing, small as well as great—like the trunk of an elephant, as well suited to pick up a pin as to tear down a tree.'—Stoughton, Vol. I. p. 33.}

His chief zeal was directed to the establishment of absolute outward uniformity in religion as he understood it, without regard to the rights of conscience and private judgment. His religion consisted of High-Church Episcopalianism and Arminianism in the nearest possible approach to Rome, which he admired and loved, and the furthest possible distance from Geneva, which he hated and abhorred.\footnote{I must add, however, that in his book against Fisher the Jesuit there are a few favorable allusions to Calvin as a theologian, especially to his doctrine of the spiritual real presence.} But while Arminianism in Holland was a protestant growth, and identified with the cause of liberal progress, Laud made it subservient to his intolerant High-Churchism, and liked it for its affinity with the Semipelagianism of the Greek fathers. To enforce this Semipelagian High-Churchism, and to secure absolute uniformity in the outward service of God in the three kingdoms, was the highest aim of his administration, to which he bent every energy. He could not conceive spiritual unity without external uniformity. This was his fundamental error. In a characteristic sermon which he preached at Westminster before Parliament, March 17, 1628, on unity in Church and State (Eph. iv. 3), he says: 'Unity \emph{of any kind} will do much good; but the best is safest, and that is unity of the Spirit. . . . The way to keep unity both in Church and State is for the governors to carry a watchful eye over all such as are discovered or feared to have private ends. . . . Provide for the keeping of unity, and . . . God will bless you with the success of this day. For this day, the seventeenth of March, Julius Cæsar overthrew Sextus Pompeius. . . . And this very day, too, Frederick II. entered Jerusalem, and recovered whatsoever Saladin had taken from the Christians. But I must tell you, these emperors and their forces were great keepers of unity.'\footnote{\emph{Works} (Oxf. 1847), Vol. I. pp. 161, 167, 180, 181.}

In the same year he caused the Royal Declaration to be added to the Thirty-nine Articles to check their Calvinistic interpretation.\footnote{That Laud is the author of this Declaration was charged by Prynne, and is proved by the Oxford editor of his \emph{Works}, Vol. I. pp. 153 sq. Comp. above, p. 617.} From the same motive he displaced, through the agency of Wentworth and Bramhall, the Calvinistic Irish Articles, and neutralized the influence of Archbishop Ussher in Ireland. But the height of his folly, and the beginning of his fall, was the enforcement of his episcopal and ritualistic scheme upon Presbyterian Scotland in criminal defiance of the will of the people and the law of the land. This brought on the Scotch Covenant and hastened the Civil War.

In England he filled all vacancies with Churchmen and Arminians of his own stamp. He kept (as he himself informs us in his \emph{Diary}) a ledger for the guidance of his royal master in the distribution of patronage: those marked by the letter \emph{O} (Orthodox) were recommended to all favors, those marked \emph{P} (Puritans) were excluded from all favors. Bishop Morely, on being asked what the Arminians held, wittily and truthfully replied, 'The best bishoprics and deaneries in England.' He expelled or silenced the Puritans, and shut up every unauthorized meeting-house. 'Even the devotions of private families could not escape the vigilance of his spies.' In his eyes the Puritans were but a miserable 'fraction' of fanatics and rebels, a public nuisance which must be crushed at any price. He made the congregations of French and Dutch refugees conform or leave the land, and forbade the English ambassador in Paris to attend the service of the Huguenots. He restrained the press and the importation of foreign books, especially the favorite Geneva translation of the Bible prepared by the Marian exiles. He stopped several ships in the Thames which were to carry persecuted and disheartened Puritans to New England, and thus tried to prevent Providence from writing the American chapter in history. In this way Oliver Cromwell is said to have been kept at home, that in due time he might overthrow the monarchy.

With equal rigor Laud enforced his ritualism, which was to him not only a desirable matter of taste and propriety, but also an essential element of reverence and piety. He took special care and showed great liberality for the restoration of cathedrals and the full cathedral service with the most pompous ceremonial; he made it a point of vital importance that the communion-tables be removed from the centre of the church to the east end of the chancel, elevated above the level of the pavement, placed altar-ways, railed in, and approached always with the prescribed bows and genuflexions\footnote{He informed the king of 'a very ill accident which happened at Taplow, by reason of not having the communion-table, railed in, that it might be kept from profanations. For in the sermon time a dog came to the table and took the loaf of bread prepared for the Holy Sacrament in his mouth, and ran away with it. Some of the parishioners took the same from the dog and set it again on the table. After sermon the minister could not think fit to consecrate this bread, and other fit for the Sacrament was not to be had in that town, so there was no Communion.'—\emph{Works}, Vol. V. p 367. This brings to mind the grave and curious disputes of the mediæval schoolmen on the question what effect the consecrated wafer would have upon a mouse or a rat.} He called the altar 'the greatest place of God's residence on earth,' and magnified it above the pulpit, because on the altar was Christ's body, which was more than his Word; but he denied the charge of transubstantiation. He introduced pictures, images, crucifixes, candles, and brought put every worn-out relic from the ecclesiastical wardrobe of the Middle Ages. Being himself unmarried, he preferred celibates in the priesthood. In the University of Oxford, to which he was a munificent benefactor, he was addressed as His Holiness, and Most Holy father.

No wonder that he was charged with the intention to reintroduce popery into England. The popular mind, especially in times of excitement, takes no notice of minor shades of distinction, and knows only friend and foe. Laud, no doubt, did the pope's work effectually, but he did it unintentionally. He loved the Roman Church much better than the Protestant sects, but he loved the Anglican Church more. He once dreamed, as he tells us, 'that he was reconciled to the Church of Rome,' but was much troubled by it.\footnote{\emph{Diary}, March 8, 1626 (\emph{Works}, Vol. III. p. 201).} He was twice offered, by some unnamed agent, a cardinal's hat, but promptly declined it.\footnote{He relates, in his \emph{Diary}, Aug. 4, 1633 (on the day of Archbishop Abbot's death), that 'there came \emph{one} to me, seriously, . . . and offered me to be a Cardinal. I went presently to the King and acquainted him both with the thing and the person.' On the 17th of August, having in the mean time (Aug. 6) been appointed Archbishop of Canterbury, he had a second offer of a red hat, and again answered 'that something dwelt within him which would not suffer that till Rome were other than it is' (\emph{Works}, Vol. III. p. 219). In his Marginal Notes on Prynne's \emph{Breviate} (p. 266), he adds that his 'conscience' also went against this. But it is by no means certain or even probable that the pope himself (as Fuller states without proof) authorized such an offer. It may have been a trap laid for Laud on the eve of his elevation to the primacy. Lingard, the Roman Catholic historian of England, says that Laud was 'in bad repute in Rome' (Vol. X. p. 139), and Dean Hook, his Anglo-Catholic biographer, asserts that he was 'dreaded and hated at Rome,' and that his death was greeted there with joy (\emph{Life of L.} p. 233). Lingard adds, however, that 'in the solitude of his cell, and with the prospect of the block before his eyes, Laud began to think more favorably of the Catholic [Roman] Church,' and he shows that Rosetti inquired of Cardinal Barberini whether, if Laud should escape from the Tower, the pope would afford him an asylum in Rome with a pension of 1000 crowns. But this is inconsistent with Laud's last defense. He was then over seventy, and anxious to die.} He preferred to be an independent pope in England, and aped the Roman original as well as he could, with more or less show of real or imaginary opposition that springs from rivalry and affinity. Neal says that he was not 'an absolute papist,' but 'ambitious of being the sovereign patriarch of three kingdoms.'\footnote{\emph{Hist. of the Puritans}, Vol. I. p. 280.} From his 'Conference' with Fisher the Jesuit, which is by far his ablest and most learned performance, it is very evident that he differed from Rome on several points of doctrine and practice, such as the invocation of Mary and the saints, the worship of images, transubstantiation, the sacrifice of the mass, works of supererogation, the temporal power of the pope, and the infallibility of councils; and that his mind, though clear and acute, was not sufficiently logical to admit the ultimate conclusions of some of his own premises.\footnote{The Conference with Fisher (whose real name was Piersey or Percy) took place, by command and in the presence of King James, May 24, 1622, and was edited, with final corrections and additions, by Laud himself in 1639. It was republished 1673 and 1686, and by the Oxford University Press 1839, with an Introduction by Edward Cardwell. It is also included in Vol. II. of the Oxf. ed. of his \emph{Works.} Laud thought that his way of defense was the only one by which the Church of England could justify her separation from the Church of Rome. He bequeathed £100 for a Latin translation of this book.} He regarded the Reformation merely as an incident in the history of the English Church, and rejected only such doctrines of Romanism as he was unable to find in the Bible and the early fathers. In his long and manly defense before the House of Lords he claimed to have converted several persons (Chillingworth among them) from popery, but frankly admitted that 'the Roman Church never erred in fundamentals, for fundamentals are in the Creed, and she denies it not. Were she not a true Church, it were hard with the Church of England, since from her the English bishops derive their apostolic succession. She is, therefore, a true but not an orthodox Church. Salvation may be found in her communion; and her religion and ours are one in the great essentials. I am not bound to believe each detached phrase in the Homilies, and I do not think they assert the pope to be Antichrist; yet it can not be proved that I ever denied him to be so. As to the charge of unchurching foreign Protestants, I certainly said generally, according to St. Jerome, ``No bishop, no Church;'' and the preface of the book of ordination sets forth that the three orders came from the apostles.' In his last will and testament he says: 'For my faith, I die as I have lived, in the true orthodox profession of the Catholic faith of Christ, foreshadowed by the prophets and preached to the world by Christ himself, his blessed apostles, and their successors; and a true member of his Catholic Church within the communion of a living part thereof, the present Church of England, as it stands established by law.'

In one word, Laud was a typical Anglo-Catholic, who unchurched all non-episcopal Churches, and regarded the Anglican Church as an independent sister of the Latin and Greek communions, and as the guardian of the whole truth as against the 'sects,' and of nothing but the truth as against Rome. The Anglo-Catholicism of the nineteenth century is simply a revival of Laud's system divested of its hateful tyranny and political ambition and entanglements. Dr. Pusey, the father of modern Anglo-Catholicism, is superior to Archbishop Laud in learning, spirituality and charity, but in their theology and logic there is no difference. [**placement of this note is uncertain**] \footnote{The \emph{Works} of Laud embrace five volumes in the Oxford 'Library of Anglo-Catholic Theology.' His seven sermons preached on great state occasions abound with his high notions of royalty, episcopacy, and uniformity, but do not rise above mediocrity. His \emph{Diary}—the chief source of his autobiography—though not 'contemptible' (as Hallam characterizes it), is dry and pedantic, and notices trifling incidents as important occurrences, \emph{e.g.}, the bad state of the weather, his numerous dreams, the marriage of K. C. with a minister's widow, the particular posture of the Elector of the Palatinate at communion 'upon a stool by the wall before the traverse, and with another and a cushion before him to kneel at' (Dec. 25, 1635), and his unfortunate affairs with 'E. B.' (of which he deeply repented; see his \emph{Devot.} Vol. III. p. 81). His \emph{Devotions} are made up mostly of passages of the Psalms and the fathers, and reveal the best side of his private character. His last prayer, as he kneeled by the block to receive the fatal stroke, is the crown of his prayers, and worth quoting: 'Lord, I am coming as fast as I can. I know I must pass through the shadow of death before I can come to see Thee. But it is but \emph{umbra mortis}, a mere shadow of death, a little darkness upon nature; but Thou, by Thy merits and passion, hast broken through the jaws of death. So, Lord, receive my soul, and have mercy upon me; and bless this kingdom with peace and plenty, and with brotherly love and charity, that there may not be this effusion of Christian blood amongst them, for Jesus Christ His sake, if it be Thy will.' The opinions on Laud are mostly tinctured by party spirit. His friend Clarendon says, 'His learning, piety, and virtue have been attained by very few, and the greatest of his infirmities are common to all, even the best of men.' Prynne, who lost his two ears by Laud's influence, calls him the most execrable traitor and apostate that the English soil ever bred ('Canterbury's Doome'). His biographers, Peter Heylin (\emph{Cyprianus Anglicanus}, Lond. 1671), John Parker Lawson (\emph{The Life and Times of William Laud}, Lond. 1829, 2 vols.), and Dr. Hook (in the \emph{Lives of the Archbishops of Canterbury}, Vol. XI. Lond. 1875), are vindicators of his character and policy. May, Hallam, Macaulay, Lingard, Green, Häusser, and Stoughton (Vol. I. pp. 402 sq.) condemn his public acts, but give him credit for his private virtues. May (\emph{History of Parliament}, approvingly quoted by Hallam, chap. viii. Charles I.) says: 'Laud was of an active, or, rather, of a restless mind; more ambitious to undertake than politic to carry on; of a disposition too fierce and cruel for his coat. He had few vulgar and private vices, as being neither taxed of covetousness, intemperance, nor incontinence; and, in a word, a man not altogether so bad in his personal character as unfit for the state of England.'}

THE STAR-CHAMBER AND THE HIGH-COMMISSION COURT.

The two chief instruments of this royal episcopal tyranny were the Star-Chamber and the High-Commission Court—two kinds of inquisition—the first political, the second ecclesiastical, with an unlimited jurisdiction over all sorts of misdemeanors, and with the power to inflict the penalties of deprivation, imprisonment, fines, whipping, branding, cutting ears, and slitting noses.

Freedom of speech and the press, which is now among the fundamental and inalienable rights of every Anglo-Saxon citizen, was punished as a crime against society. Prynne, a graduate of Oxford, and a learned barrister of Lincoln's Inn, who published an unreadable book (\emph{Histrio-Mastix, the Player' Scourge, or Actors' Tragedie, divided into Two Parts}) against theatres, masquerades, dancing, and women actors, with reflections upon the frivolities of the queen, was condemned by the Star-Chamber to be expelled from Oxford and Lincoln's Inn, to be fined £5000, to stand in the pillory at Westminster and Cheapside, to have his ears cut off, his cheeks and forehead branded with hot irons, and to be imprisoned for life. His huge quarto volume of 1006 pages, with quotations from as many authors, was burned under his nose, so that he was nearly suffocated with the smoke. Leighton, a Scotchman (father of the saintly archbishop), Bastwick, a learned  physician, and Henry Burton, a B.D. of Oxford, and rector of a church in London, were treated with similar cruelty for abusing in printed pamphlets the established hierarchy. No doubt their language was violent and coarse,\footnote{Burton called the bishops \emph{step}-fathers, \emph{cater}-pillars, limbs of the beast, blind watchmen, dumb dogs, new Babel-builders, antichristian mushrumps, etc. Prynne called them 'silk and satin divines,' and said that 'Christ himself was a Puritan, and that, therefore, all men should become Puritans.' But their opponents could be equally abusive. Lord Cottington, one of Prynne's judges, said that, in writing the \emph{Histrio-Mastix}, 'either the devil had assisted Prynne or Prynne the devil.' Another judge, the Earl of Dorset, called him '\emph{omnium malorum nequissimum.} '} but torture and mutilation are barbarous and revolting. And yet Laud not only thanked the lords of the Star-Chamber for their 'just and honorable sentence upon these men,' but regretted, in a letter to Strafford, that he could not resort to more 'thorough' measures.

THE CIVIL WAR AND THE COMMONWEALTH.

The excesses of despotism, sacerdotalism, ceremonialism, intolerance, and cruelty exhausted at last the patience of a noble, freedom-loving people, and kindled the blazing war-torch which burned to the ground the throne and the temple. The indignant nation rose in its majesty, and asserted its inherent and constitutional rights.

The storm burst forth from the North. The Scots compelled the King to abandon his schemes of innovation, and to admit that prelacy was contrary to Scripture. In England the memorable Long Parliament organized the opposition, and assumed the defense of constitutional liberty against royal absolutism. It met Nov. 3, 1640, and continued till April 20, 1653, when it was dissolved by Cromwell to give way to military despotism. The war between the Parliament and the King broke out in August, 1642. For several months the Cavaliers fought more bravely and successfully than the undisciplined forces of the Roundheads; but the fortunes of war changed when Oliver Cromwell, a country gentleman, bred to peaceful pursuits, appeared at the head of his Ironsides, whom he selected from the ranks of the Puritans. It was an army such as England never saw before or since—an army which feared God and hated the pope; which believed in the divine decrees and practiced perseverance of saints; which fought for religion; which allowed no oath, no drunkenness, no gambling in the camp; which sacredly respected private property and the honor of woman; which went praying and psalm-singing into the field of battle, and never returned from it without the laurels of victory. And when these warriors were disbanded at the Restoration, they astonished the royalists by quietly taking their place among the most industrious, thrifty, and useful citizens.\footnote{One of the noblest specimens of a Puritan officer was Col. Hutchinson, whose character and life have been so admirably described by his widow (pp. 24 sqq. Bohn's ed.).}

During the reign of the Long Parliament the Star-Chamber and the High-Commission Court were ignominiously and forever swept out of existence amid the execrations of the people. The episcopal hierarchy and the Liturgy were overthrown (Sept. 10, 1642); about two thousand royalist ministers, many of them noted for incapacity, idleness, and immorality, others highly distinguished for scholarship and piety—as Hammond, Sanderson, Pocock, Byron Walton, Hall, Prideaux, Pearson—were ejected as royalists from their benefices and given over to poverty and misery, though one fifth of the revenues of the sequestered livings was reserved for the sufferers.\footnote{Comp. Marsden, \emph{The Later Puritans}, pp. 40 sqq. Baxter himself allows that 'some able, godly preachers were cast out for the war alone.' Among these was also the excellent Thomas Fuller, the author of the incomparable books on \emph{Church History} and the \emph{Worthies of England}, although in the days of Laud he had been stigmatized as a Puritan in doctrine.} This summary and cruel act provoked retaliation, which in due time came with increased severity. The leaders of despotism—the Earl of Strafford (May 12, 1641), Archbishop Laud (Jan. 10, 1645), and at last the King himself (Jan. 30, 1649)—were condemned to death on the block, and thus surrounded by the halo of martyrdom. Their blood was the seed of the Restoration. The execution of Charles especially was in the eyes of the great majority of the English and Scotch people a crime and a blunder, and set in motion the reaction in favor of monarchy and episcopacy.

At first, however, Cromwell's genius and resolution crushed every opposition in England, Ireland, and Scotland. On the ruins of the monarchy and of Parliament itself he raised a military government which inspired respect and fear at home and abroad, and raised England to the front rank of Protestant powers, but which created no affection and love except among his invincible army. The man of blood and iron, the ablest ruler that England ever had, died at the height of his power, on the anniversary of his victories at Dunbar and Worcester (Sept. 3), and was buried with great pomp among the legitimate kings of England in Westminster Abbey (Nov. 23, 1658).'\footnote{On his last days and utterances, see the \emph{Mercurius Politicus} for Sept. 2–9,1658, and Stoughton, \emph{The Church of the Commonwealth}, p. 511. Macaulay pays the following tribute to Cromwell's foreign policy: 'The Protector's foreign policy at the same time extorted the ungracious approbation of those who most detested him. The Cavaliers could scarcely refrain from wishing that one who had done so much to raise the fame of the nation had been a legitimate king; and the Republicans were forced to own that the tyrant suffered none but himself to wrong his country, and that, if he had robbed her of liberty, he had at least given her glory in exchange. After half a century, during which England had been of scarcely more weight in European politics than Venice or Saxony, she at once became the most formidable power in the world, dictated terms of peace to the United Provinces, avenged the common injuries of Christendom on the pirates of Barbary, vanquished the Spaniards by land and sea, seized one of the finest West India islands, and acquired on the Flemish coast a fortress which consoled the national pride for the loss of Calais. She was supreme on the ocean. She was the head of the Protestant interest. All the Reformed Churches scattered over Roman Catholic kingdoms acknowledged Cromwell as their guardian. The Huguenots of Languedoc. the shepherds who, in the hamlets of the Alps, professed a Protestantism older than that of Augsburg, were secured from oppression by the mere terror of that great name. The pope himself was forced to preach humanity and moderation to popish princes. For a voice which seldom threatened in vain had declared that, unless favor were shown to the people of God, the English guns should be heard in the Castle of Saint Angelo. In truth, there was nothing which Cromwell had, for his own sake and that of his family, so much reason to desire as a general religious war in Europe. In such a war he must have been the captain of the Protestant armies. The heart of England would have been with him. His victories would have been hailed with a unanimous enthusiasm unknown in the country since the rout of the Armada, and would have effaced the stain which one act, condemned by the general voice of the nation, has left on his splendid fame. Unhappily for him. he had no opportunity of displaying his admirable military talents except against the inhabitants of the British Isles.'—\emph{History of England}, ch. i. Carlyle says that Cromwell was the best thing that England ever did.}

THE RESTORATION.

The Puritan Commonwealth was but a brilliant military episode, and died with its founder. His son Richard, amiable, good-natured, weak and incompetent, succeeded him without opposition, but resigned a few months after (April 22, 1659). The army, which under its great commander had ruled the divided nation, was now divided, while the national sentiment in the three kingdoms became united, and demanded the restoration of the old dynasty as the safest way to escape the dangers of military despotism. Puritanism represented only a minority of the English people, and the majority of this minority were royalists. The Presbyterians, who were in the saddle during the interregnum, were specially active for the unconditional recall of the treacherous Stuarts. The event was brought about by the cautious and dexterous management of General Monk, a man of expediency, who had successively served under Charles I. and Cromwell, and worshiped with Episcopalians, Presbyterians, and Independents, and at last returned to the Episcopal Church. Charles II., 'who never said a foolish thing, and never did a wise one,' was received with such general enthusiasm on his triumphal march from Dover to London that he wondered where his enemies were, or whether he ever had any. The revolution of national sentiment was complete. The people seemed as happy as a set of unruly children released from the discipline of the school.\footnote{'Almost all the gentry of all parts went—some to fetch him over, some to meet him at the sea-side, some to fetch him into London, into which he entered on the 29th day of May, with a universal joy and triumph, even to his own amazement; who, when he saw all the nobility and gentry of the land flowing in to him, asked where were his enemies. For he saw nothing but prostrates, expressing all the love that could make a prince happy. Indeed, it was a wonder in that day to see the mutability of some, and the hypocrisy of others, and the servile flattery of all. Monk, like his better genius, conducted him and was adored like one that had brought all the glory and felicity of mankind home with this prince.'—\emph{Memoirs of the Life of Col. Hutchinson}, p. 402.}

The restoration of the monarchy was followed by the restoration of Episcopacy and the Liturgy with an exclusiveness that did not belong to it before. The Savoy Conference between twenty-one Episcopalians and an equal number of Presbyterians (April 15 till July 25, 1661) utterly failed, and left both parties more exasperated and irreconcilable than before. The Churchmen, once more masters of the situation, refused to make any concessions and changes.\footnote{The fullest account of the conference held in the Savoy Hospital, London, is given by Baxter, who was a member, in his Autobiography. Comp. Neal, Cardwell, Stoughton (\emph{Restor.} Vol. I. p. 157), Hallam (Ch. XI. Charles II.), and Procter (\emph{History of the Book of Common Prayer}, p. 113). Hallam casts the chief blame on the Churchmen, who had it in their power to heal the division and to retain or to expel a vast number of worthy clergymen. But both parties lacked the right temper, and smarted under the fresh recollection of past grievances. Baxter embodied the changes desired by the Puritans in his Liturgy, the hasty work of a fortnight, which was never used, but republished by Prof. Shields of Princeton, Philadelphia, 1867.} Thus another opportunity of comprehension was lost. In the revision of the Liturgy, which was completed by Convocation at the close of the same year (Dec., 1661), approved by the King, and ratified by Act of Parliament (April, 1662), not the slightest regard was paid to Presbyterian objections, reasonable or unreasonable, although about six hundred alterations were made; on the contrary, all the ritualistic and sacerdotal features complained of were retained and even increased.\footnote{Procter (p. 141): 'Some changes were made, in order to avoid the appearance of favoring the Presbyterian form of Church government; thus, \emph{church}, or \emph{people}, was substituted for \emph{congregation}, and \emph{ministers} in for of \emph{the congregation; priests and deacons} were especially named instead of \emph{pastors and ministers.}' The Apocryphal lessons were retained, and the legend of Bel and the Dragon (omitted in 1604) was again introduced in the Calendar of Daily Lessons, to show contempt for the Puritan scruples. In the Litany the words 'rebellion' and 'schism' were added to the petition against 'sedition.'} The Act of Uniformity, which received the royal assent May 19, 1662, and took effect on the ominous St. Bartholomew's Day, Aug. 24, 1662 (involuntarily calling to mind the massacre of the Huguenots), required not only from ministers, but also from all schoolmasters, absolute conformity to the revised Liturgy and episcopal ordination, or reordination. By this cruel act more than two thousand Puritan rectors and vicars—that is, about one fifth of the English clergy, including such men as Baxter, Howe, Poole, Owen, Goodwin, Bates, Manton, Caryl—were ejected and exposed to poverty, public insult, fines, and imprisonment for no other crime than obeying God rather than men. A proposition in the House of Commons to allow these heroes of conscience one fifth of their income, as the Long Parliament had done in the removal of royalist clergymen, was lost by a vote of ninety-four to eighty-seven.\footnote{Dr. Stoughton, a well-informed and impartial historian, gives it as the result of his careful inquiry that the persecution and sufferings of the Episcopalians under the Long Parliament and the Commonwealth are not to be compared with the persecution of the Nonconformists under Charles I. and Charles II. (\emph{Ch. of the Commonwealth}, p. 346). Hallam is of the same opinion. Richard Baxter, one of the ejected ministers, gives a sad account of their sufferings: 'Many hundreds of these, with their wives and children, had neither house nor bread. . . . Their congregations had enough to do, besides a small maintenance, to help them out of prisons, or to maintain them there. Though they were as frugal as possible, they could hardly live; some lived on little more than brown bread and water; many had but eight or ten pounds a year to maintain a family, so that a piece of flesh has not come to one of their tables in six weeks' time; their allowance could scarce afford them bread and cheese. One went to plow six days and preached on the Lord's day. Another was forced to cut tobacco for a livelihood. . . . Many of the ministers, being afraid to lay down their ministry after they had been ordained to it, preached to such as would hear them in fields and private houses, till they were apprehended and cast into gaols, where many of them perished' (quoted by Green, p. 612). Baxter himself was repeatedly imprisoned, although he was a royalist and openly opposed Cromwell's rule. For many details of suffering, see Orme's \emph{Life of Baxter }(Lond. 1830), pp. 229 sqq.}

Even the dead were not spared by the spirit of 'mean revenge.' The magnates of the Commonwealth, twenty-one in number (including Dr. Twisse, the Prolocutor of the Westminster Assembly), who had been buried in Westminster Abbey since 1641, were exhumed and reinterred in a pit (Sept. 12, 1661). Seven only were exempt; among them Archbishop Ussher, who had been buried there at Cromwell's express desire, and at a cost of £200 paid by him. Cromwell himself, Ireton, and Bradshaw were dug up Jan. 29, 1661, next day dragged to Tyburn, hanged (with their faces turned to Whitehall), decapitated, and buried under the gallows. Cromwell's head was planted on the top of Westminster Hall.\footnote{Stanley's \emph{Hist. Memorials of Westminster Abbey}, pp. 191 sq., 247, 320 (3d ed. Lond. 1869).}

The Puritans were now a target of hatred and ridicule as well as persecution. They were assailed from the pulpit, the stage, and the press by cavaliers, prelatists, and libertines as a set of hypocritical Pharisees and crazy fanatics, noted for their love of Jewish names, their lank hair, their sour faces, their deep groans, their long prayers and sermons, their bigotry and cant.\footnote{Butler's \emph{Hudibras} fairly reflects the prevailing sentiment of the Restoration period about the Puritans. He caricatures them in his mock-heroic style (Part I. Canto I. vers. 192 sqq.) as  \par 'That stubborn crew \par Of errant saints, whom all men grant \par To be the true Church militant: \par Such as do build their faith upon \par The holy text of pike and gun; \par Decide all controversy by \par Infallible artillery; \par And prove their doctrine orthodox \par by apostolic blows and knocks; \par Call fire, and sword, and desolation \par A godly thorough Reformation, \par Which always must be carried on, \par And still be doing, never done, \par As if religion were intended \par For nothing else but to be mended.'} And yet the same Puritanism, blind, despised, forsaken, or languishing in prison, produced some of the noblest works, which can never die. It was not dead—it was merely musing and dreaming, and waiting for a resurrection in a nobler form. Milton's 'Paradise Lost' (1667) and Bunyan's 'Pilgrim's Progress' (1678) are the shining lights which illuminate the darkness of that disgraceful period.\footnote{'Puritanism,' says an Oxford historian, 'ceased from the long attempt to build up a kingdom of God by force and violence, and fell back on its truer work of building up a kingdom of righteousness in the hearts and consciences of men. It was from the moment of its seeming fall that its real victory began. As soon as the wild orgy of the Restoration was over, men began to see that nothing that was really worthy in the work of Puritanism had been undone. The revels of Whitehall, the skepticism and debauchery of courtiers, the corruption of statesmen, left the mass of Englishmen what Puritanism had made them—serious, earnest, sober in life and conduct, firm in their love of Protestantism and of freedom. In the Revolution of 1688 Puritanism did the work of civil liberty, which it had failed to do in that of 1642. It wrought out through Wesley and the revival of the eighteenth century the work of religious reform which its earlier efforts had only thrown back for a hundred years. Slowly, but steadily, it introduced its own seriousness and purity into English society, English literature, English politics. The whole history of English progress, since the Restoration, on its moral and spiritual sides, has been the history of Puritanism.'—J. R. Green's \emph{Short History of the English People}, p. 586.}

With the Restoration rushed in a flood of frivolity and immorality; the King himself setting the example by his shameless adulteries, which he blazoned to the world by raising his numerous mistresses and bastards to the rank and wealth of the nobility of proud old England. 'The violent return to the senses,' says a French writer who has not the slightest sympathy with Puritanism, 'drowned morality. Virtue had the semblance of Puritanism. Duty and fanaticism became mingled in a common reproach. In this great reaction, devotion and honesty, swept away together, left to mankind but the wreck and the mire. The more excellent parts of human nature disappeared; there remained but the animal, without bridle or guide, urged by his desires beyond justice and shame.'\footnote{Taine's \emph{History of English Literature}, vol. i. p. 461 (Am. ed.).}

THE REVOLUTION.

Bad as was Charles II. (1660-1685), his brother, James II. (1685-1688), was worse. He seemed to combine the vices of the Stuarts without their redeeming traits. Charles, indifferent to religion and defiant to virtue during his life, sent on his death-bed for a Romish priest to give him absolution for his debaucheries. James openly professed his conversion to Romanism, filled in defiance of law the highest posts in the army and the cabinet with Romanists, and opened negotiations with Pope Innocent XI. At the same time he persecuted with heartless cruelty the Protestant Dissenters, and outraged justice by a series of judicial murders which have made the name of Chief Justice Jeffreys as infamous as Nero's.

At last the patience of the English people was again exhausted, the incurable race of the Stuarts, unwilling to learn and to forget any thing, was forever hurled from the throne, and the Prince of Orange, who had married Mary, the eldest daughter of James, was invited to rule England as William III.

THE RESULT.

The Revolution of 1688 was a political triumph of Puritanism, and secured to the nation constitutional liberty and the Protestant religion. The Episcopal Church remained the established national Church, but the Act of Toleration of 1689 guaranteed liberty and legal protection to such Nonconformists as could subscribe thirty-five and a half of the Thirty-nine Articles of Religion, omitting those to which the Puritans had conscientious scruples. Though very limited, this Act marked a great progress. It broke up the reign of intolerance, and virtually destroyed the principle of uniformity. The Act of Uniformity of 1662 was intended for the whole kingdom, and proceeded on the theory of an ecclesiastical incorporation of all Englishmen; now it was confined to the patronized State Church. It recognized none but the Episcopal form of worship, and treated non-Episcopalians as disloyal subjects, as culprits and felons; now other Protestant Christians—Presbyterians, Independents, Baptists, and even Quakers—were placed under the protection of the law, and permitted to build chapels and to maintain pastors at their own expense. The fact was recognized that a man may be a good citizen and a Christian without conforming to the State religion. Uniformity had proved an intolerable tyranny, and had failed. Comprehension of different denominations under one national Church, though favored by William, seemed impracticable. Limited toleration opened the way for full liberty and equality of Christian denominations before the law; and from the soil of liberty there will spring up a truer and deeper union than can be secured by any compulsion in the domain of conscience, which belongs to God alone.

Puritanism did not struggle in vain. Though it failed as a national movement, owing to its one-sidedness and want of catholicity, it accomplished much. It produced statesmen like Hampden, soldiers like Cromwell, poets like Milton, preachers like Howe, theologians like Owen, dreamers like Bunyan, hymnists like Watts, commentators like Henry, and saints like Baxter, who though dead yet speak. It lives on as a powerful moral element in the English nation, in the English Church, in English society, in English literature. It has won the esteem of the descendants of its enemies. In our day the Duke of Bedford erected a statue to Bunyan (1874) in the place where he had suffered in prison for twelve years; and Episcopalians, Presbyterians, and Independents united in a similar tribute of justice and gratitude to the memory of Baxter at Kidderminster (1875), where he is again pointing his uplifted arm to the saints' everlasting rest. The liberal-minded and large-hearted dean of Westminster represented the nobler part of the English people when he canonized those great and good men in his memorial discourses at the unveiling of their statues. Puritanism lives moreover in New England, which was born of the persecutions and trials of its fathers and founders in old England, and gave birth to a republic truer, mightier, and more enduring than the ephemeral military commonwealth of Cromwell. It will continue to preserve and spread all over the Saxon world the love of purity, simplicity, spirituality, practical energy, liberty, and progress in the Christian Church.

On the other hand, it is for the children of the Puritans to honor the shining lights of the Church of England who stood by her in the days of her trial and persecution. That man is to be pitied indeed who would allow the theological passions of an intolerant age to blind his mind to the learning, the genius, and the piety of Ussher, Andrewes, Hall, Pearson, Prideaux, Jeremy Taylor, Barrow, and Leighton, whom God has enriched with his gifts for the benefit of all denominations. It is good for the Church of England—it is good for the whole Christian world—that she survived the fierce conflict of the seventeenth century and the indifferentism of the eighteenth to take care of venerable cathedrals, deaneries, cloisters, universities, and libraries, to cultivate the study of the fathers and schoolmen, to maintain the importance of historical continuity and connection with Christian antiquity, to satisfy the taste for stability, dignity, and propriety in the house of God, and to administer to the spiritual wants of the aristocracy and peasantry, and all those who can worship God most acceptably in the solemn prayers of her liturgy, which, with all its defects, must be pronounced the best ever used in divine service.

While the fierce conflict about religion was raging, there were prophetic men of moderation and comprehension on both sides—

'Whose dying pens did write of Christian union,

How Church with Church might safely keep communion;

Who finding discords daily to increase,

Because they could not live, would die, in peace.'

In a sermon before the House of Commons, under the arched roof of Westminster Abbey, Richard Baxter uttered this sentence: 'Men that differ about bishops, ceremonies, and forms of prayer, may be all true Christians, and dear to one another and to Christ, if they be practically agreed in the life of godliness, and join in a holy, heavenly conversation. But if you agree in all your opinions and formalities, and yet were never sanctified by the truth, you do but agree to delude your souls, and neither of you will be saved for all your agreement.'\footnote{\emph{Vain Religion of the Formal Hypocrite.} Baxter's \emph{Works}, Vol. XVII. p. 80. Quoted by Stoughton, p. 195. The sermon was preached Apr. 30, 1660, just before the recall of Charles II. See Orme. \emph{Life of Baxter}, p. 160.}

This is a noble Christian sentiment, echoing the words of a greater man than Baxter: 'In Christ Jesus neither circumcision availeth any thing, nor uncircumcision,'—we may add, neither surplice nor gown, neither kneeling nor standing, neither episcopacy nor presbytery nor independency—' but a new creature.'\footnote{Gal. vi. 15.}

\xsubsection{The Westminster Assembly.}

§ 93. The Westminster Assembly.

Literature.

I. Original Sources.

The  Westminster Standards—see § 94.

Minutes of the Sessions of the Westminster Assembly of Divines (from Nov. 1644 to March, 1649). \emph{From Transcripts of the Originals procured by a Committee of the General Assembly of the Church of Scotland, ed. by the} Rev. Alex. F. Mitchell, D.D., \emph{and the} Rev. John Struthers, LL.D. Edinb. and Lond. 1874. (The MS. Minutes of the Westm. Assembly from 1643 to 1652, formerly supposed to have been lost in the London fire of 1666, were recently discovered in Dr. Williams's library, Grafton St., London, and form 3 vols. of foolscap fol. They are mostly in the handwriting of Adoniram Byfield, one of the scribes of the Assembly. A complete copy was made for the General Assembly of the Church of Scotland, and is preserved in Edinburgh. They are, upon the whole, rather meagre, and give only the results, with brief extracts from the speeches, without the arguments.)

Robert Baillie (Principal of the University of Glasgow, and one of the Scotch delegates to the Assembly of Westminster, b. 1599, d. 1662): \emph{Letters and Journals ed. from the author's MSS. by David Laing, Esq.} Edinb. 1841–42, 3 vols. (These Letters and Journals extend from Jan. 1637 to May, 1662, and exhibit in a lively and graphic manner 'the stirring scenes of a great national drama,' with the hopes and fears of the time. Vol. II. and part of Vol. III. bear upon the Westm. Assembly.)

John Lightfoot, D.D. (Master of Catharine Hall and Vice-Chancellor of Cambridge, one of the members of the Westm. Assembly, b. 1602, d. 1675): \emph{Journal of the Proceedings of the Assembly of Divines from Jan.} 1, 1643 \emph{to Dec.} 31, 1644. In Vol. XIII. pp. 1–344 of his \emph{Whole Works, ed. by John Rogers Pitman} (Lond. 1825, in 13 vols.).

George Gillespie (the youngest of the Scotch Commissioners to the Assembly, d. 1648): \emph{Notes of Debates and Proceedings of the Westminster Assembly}, ed. from the MSS. by David Meek, Edinb. 1846. Comp. also Gillespie's \emph{Aaron's Rod Blossoming} (a very able defense of Presbyterianism against Independency and Erastianism), Lond. 1646, republ. with his other works and a memoir of his life by Hetherington, Edinb. 1844–46, 2 vols.

Journals of the House of Lords and the House of Commons from 1643 to 1649.

John Rushworth (assistant clerk and messenger of the Long Parliament, and afterwards a member of the House of Commons, d. 1690): \emph{Historical Collections of remarkable Proceedings in Parliament.} Lond. 1721, 7 vols.

(The 'fourteen or fifteen octavo vols.' of daily proceedings which Dr. Thomas Goodwin, the eminent Independent member of the Assembly, is reported by his son to have written 'with his own hand,' have never been published or identified. They must not be confounded with the \emph{three folio} vols. of official minutes in Dr. Williams's library.)

Historical.

The respective sections in Fuller (Vol. VI. pp. 247 sqq.), Neal (Part III. chaps. 2, 4, 6, 8, 10), Stoughton (Vol. I. pp. 271, 327, 448 sqq.), Masson (\emph{Life of Milton}, Vols. II. and III.), and other works mentioned in § 92.

W. M. Hetherington: \emph{History of the Westminster Assembly of Divines.} Edinb. 1843; New York, 1844.

James Reid: \emph{Memoirs of the Lives and Writings of those eminent Divines who convened in the famous Assembly at Westminster.} Paisley, 1811 and 1815, 2 vols.

Gen. von Rudloff: \emph{Die Westminster Synode}, 1643–1649. In Niedner's \emph{Zeitschrift für die histor. Theologie} for 1850, pp. 238–296. (The best account of the Assembly in the German language.)

P. Schaff: Art. \emph{Westminster Synode}, etc., in Herzog's \emph{Real-Encykl.} Vol. XVIII. pp. 52 sqq., and Art. on the same subject in his \emph{Relig. Encycl.} N. Y. 1884, Vol. III. pp. 2499 sqq.

Thos. M'Crie: \emph{Annals of English Presbytery from the Earliest to the Present Time.} Lond. 1872.

J. B. Bittinger: \emph{The Formation of our Standards}, in the 'Presbyterian Quarterly and Princeton Review' for July, 1876, pp. 387 sqq.

C. A. Briggs: Art. \emph{Documentary History of the Westminster Assembly, in Pres. Rev.} for 1889, pp. 127–164.

Alexander F. Mitchell, D.D. (Prof. of Ch. Hist at St. Andrews, and ed. of the Minutes of the Assembly): \emph{The Westminster Assembly: its History and Standards.} London, 1883. (519 pages.)

IMPORTANCE OF THE ASSEMBLY.

It was after such antecedents, and in such surroundings, that the Westminster Assembly of Divines was called to legislate for Christian doctrine, worship, and discipline in three kingdoms. It forms the most important chapter in the ecclesiastical history of England during the seventeenth century. Whether we look at the extent or ability of its labors, or its influence upon future generations, it stands first among Protestant Councils. The Synod of Dort was indeed fully equal to it in learning and moral weight, and was more general in its composition, since it embraced delegates from nearly all Reformed Churches; while the Westminster Assembly was purely English and Scotch, and its standards even to-day are little known on the Continent of Europe.\footnote{It is characteristic that Dr. Niemeyer published his collection of Reformed Confessions, the most complete we have, at first without the Westminster Standards, being unable to find a copy, and issued them afterwards in a supplement. Dr. Winer barely mentions the Westminster Confession in his \emph{Symbolik}, and never quotes from it. If German Church historians (including Gieseler) were to be judged by their knowledge of English and American affairs, they would lose much of the esteem in which they are justly held. What lies \emph{westward} is a \emph{terra incognita} to most of them. They are much more at home in the by-ways of the remote past than in the living Church of the present, outside of Germany.} But the doctrinal legislation of the Synod of Dort was confined to the five points at issue between Calvinism and Arminianism; the Assembly of Westminster embraced the whole field of theology, from the eternal decrees of God to the final judgment. The Canons of Dort have lost their hold upon the mother country; the Confession and Shorter Catechism of Westminster are as much used now in Anglo-Presbyterian Churches as ever, and have more vitality and influence than any other Calvinistic Confession.

It is not surprising that an intense partisan like Clarendon should disparage this Assembly.\footnote{Clarendon, who hated Presbyterianism as a plebeian religion unfit for a gentleman, disposes of the Westminster Assembly in a few summary and contemptuous sentences: 'Of about one hundred and twenty members,' he says, 'of which the Assembly was to consist, a few very reverend and worthy persons were inserted; yet of the whole number there were not above twenty who were not declared and avowed enemies of the doctrine or discipline of the Church of England; some were infamous in their lives and conversations, and most of them of very mean parts in learning, if not of scandalous ignorance; and of no other reputation but of malice to the Church of England.' These charges are utterly without foundation, and belong to the many misrepresentations and falsehoods which disfigure his otherwise classical \emph{History of the Rebellion.} The number of members was 151.} Milton's censure is neutralized by his praise, for, although he hated presbytery only less than episcopacy, he called the Assembly a 'select assembly,' 'a learned and memorable synod,' in which 'piety, learning, and prudence were housed.' This was two years after the Assembly had met, when its character was fully shown. He afterwards changed his mind, chiefly for a personal reason—in consequence of the deservedly bad reception of his unfortunate book on 'Divorce,' which he had dedicated in complimentary terms to this very Assembly and to the Long Parliament.\footnote{In his \emph{Fragments of a History of England} (1670), Milton speaks both of the Long Parliament and the Assembly in vindictive scorn, and calls the latter 'a certain number of divines neither chosen by any rule or custom ecclesiastical, nor eminent for either piety or knowledge above others left out; only as each member of Parliament, in his private fancy, thought fit, so elected one by one.' He charges them with inconsistency in becoming pluralists and nonresidents, and with intolerance, as if 'the spiritual power of their ministry were less available than bodily compulsion,' and the authority of the magistrate 'a stronger means to subdue and bring in conscience than evangelical persuasion.' On his unhappy marriage and his tracts on Divorce growing out of it, see Masson, Vol. III. pp. 42 sqq.}

Richard Baxter, who was not a member of the Assembly, but knew it well, and was a better judge of its theological and religious character than either Clarendon or Milton, pays it this just tribute: 'The divines there congregated were men of eminent learning, godliness, ministerial abilities, and fidelity; and being not worthy to be one of them myself, I may the more freely speak the truth, even in the face of malice and envy, that, as far as I am able to judge by the information of all history of that kind, and by any other evidences left us, the Christian world, since the days of the apostles, had never a synod of more excellent divines (taking one thing with another) than this and the Synod of Dort.' He adds, however, 'Yet, highly as I honor the men, I am not of their mind in every part of the government which they have set up. Some words in their Catechism I wish had been more clear; and, above all, I wish that the Parliament, and their more skillful hand, had done more than was done to heal our breaches, and had hit upon the right way, either to unite with the Episcopalians and Independents, or, at least, had pitched on the terms that are fit for universal concord, and left all to come in upon those terms that would.'\footnote{\emph{Life and Times}, Pt. I. p. 73. Comp. Orme's \emph{Life of Baxter}, p. 69.}

Hallam censures the Assembly for its intolerant principles, but admits that it was 'perhaps equal in learning, good sense, and other merits to any Lower House of Convocation that ever made a figure in England.' One of the best-informed German historians says of the Assembly: 'A more zealous, intelligent, and learned body of divines seldom ever met in Christendom.'\footnote{General Rudloff, in his article above quoted, p. 263.}

The chief fault of the Assembly was that it clung to the idea of a national State Church, with a uniform system of doctrine, worship, and discipline, to which every man, woman, and child in three kingdoms should conform. But this was the error of the age; and it was only after a series of failures and persecutions that the idea of religious freedom took root in English soil.

APPOINTMENT OF THE ASSEMBLY.

Soon after the opening of the Long Parliament the convening of a conference of divines for the settlement of the theological and ecclesiastical part of the great conflict suggested itself to the minds of leading men. The first bill of Parliament to that effect was conceived in a spirit hostile to the Episcopal hierarchy, but rather friendly to the ancient liturgy, and was passed Oct. 15, 1642, but failed for the want of royal assent.

As the king's concurrence became hopeless, Parliament issued on its own responsibility an ordinance, June 12, 1643, commanding that an assembly of divines should be convened at Westminster, in London, on the first day of July following, to effect a more perfect reformation of the Church of England in its liturgy, discipline, and government on the basis of the Word of God, and thus to bring it into nearer agreement with the Church of Scotland and the Reformed Churches on the Continent. Presbyterianism was not mentioned, but pretty plainly pointed at. The Assembly was to consist of one hundred and fifty-one members in all, viz., thirty lay assessors (ten Lords and twenty Commoners), who were named first,\footnote{'There must be some laymen in the Synod to overlook the clergy, lest they spoil the civil work; just as when the good woman puts a cat into the milk-house to kill a mouse, she sends her maid to look after the cat, lest the cat eat up the cream.'—Selden, \emph{Table-Talk}, p. 169. (Quoted by Stoughton and Stanley.)} and included such eminent scholars, lawyers, and statesmen as John Selden, John Pym, Boulstrode Whitelocke, Oliver St. John, Sir Benjamin Rudyard, and Sir Henry Vane, and of one hundred and twenty-one divines, who were selected from the different counties, chiefly from among the Presbyterians, with a few of the most influential Episcopalians and Independents. Forty members constituted a quorum.

The Assembly was thus created by State authority. In like manner, the ancient œcumenical councils were called by emperors, and the Synod of Dort by the government of the United Provinces. The English Convocations also can not meet, nor make canons, nor discuss topics without royal license. The twenty-first of the Thirty-nine Articles forbids the calling of General Councils except 'by the command and will of princes.' Parliament now exercised the privilege of the crown, and usurped the ecclesiastical supremacy. It nominated all the members, with the exception of the Scotch commissioners, who were appointed by the General Assembly, and were admitted by Parliament. It fixed the time and place of meeting, it prescribed the work, and it paid the expenses (allowing to each member four shillings a day); it even chose the prolocutor and scribes, filled the vacancies, and reserved to its own authority all final decision; reducing thus the Assembly to an advisory council. Hence even the Westminster Confession was presented to Parliament simply as a 'humble Advice.' But with all its horror of ecclesiastical despotism, engendered by the misgovernment of Laud, the Long Parliament was the most religious political assembly that ever met in or out of England, and was thoroughly controlled by the stern spirit of Puritanism. Once constituted, the Assembly was not interfered with, and enjoyed the fullest freedom of debate. Its standards were wholly the work of competent divines, and received the full and independent assent of ecclesiastical bodies.

The king by proclamation prohibited the meeting of the Assembly, and threatened those who disobeyed his order with the loss of all their ecclesiastical livings and promotions. This unfortunately prevented the attendance of loyal Episcopalians.

COMPOSITION AND PARTIES.

It was the intention of Parliament to comprehend within the Assembly representatives of all the leading parties of the English Church with the exception of that of Archbishop Laud, whose exclusive High-Churchism and despotism had been the chief cause of the troubles in Church and State, and made co-operation impossible.\footnote{Laud says of the Assembly: 'The greatest part of them were Brownists, or Independents, or New England ministers, if not worse; or at best enemies to the doctrine and discipline of the Church of England,' The facts are, that the Independents were a small minority, and that New England was not represented at all.} The selection was upon the whole judicious, though some of the ablest and soundest Puritan divines, as Richard Baxter and John Owen, were omitted. Scotland came in afterwards, but in time to be of essential service and to give the Assembly a strong Presbyterian preponderance. The Colonial Churches of New England were invited by a letter from members of Parliament (Sept., 1642) to send the Rev. John Cotton, Thomas Hooker, and John Davenport as delegates; but they declined, because compliance would subject them to all the laws that might be made, and might prove prejudicial to them. Hooker, of Hartford, 'liked not the business,' and deemed it his duty rather to stay in quiet and obscurity with his people in Connecticut than to go three thousand miles to plead for Independency with Presbyterians in England. Davenport could not obtain leave from his congregation at New Haven. Cotton, of Boston, would not go alone.\footnote{Masson, \emph{Life of Milton}, Vol. II. p. 605; Bancroft, \emph{History of the United States of America} (Centennial ed. 1876), Vol. I. pp. 331, 332.}

The Assembly itself, by direction of Parliament, addressed fraternal letters to the Belgic, French, Helvetic, and other Reformed Churches (Nov. 30, 1643), and received favorable replies, especially from Holland, Switzerland, and the Huguenot congregation in Paris.\footnote{See the correspondence in Neal, Vol. I. pp. 470 sqq. (Harper's ed.).} Hesse Cassel advised against meddling with the bishops. The king issued a counter manifesto from Oxford, May 14, 1644, in Latin and English, to all foreign Protestants, and denied the charge of designing to introduce popery.\footnote{Neal, Vol. I. p. 472.}

As to doctrine, there was no serious difference among the members. They all held the Calvinistic system with more or less rigor. There were no Arminians, Pelagians, or Antinomians among them.

But in regard to Church government and discipline the Assembly was by no means a unit, although the Scotch Presbyterian polity ultimately prevailed, and became for a brief season, by act of Parliament, even the established form of government in England. The most frequent and earnest debates were on this point rather than on doctrine and worship. This conflict prevented the Assembly, says Neal (an Independent), from 'laying the top stone of the building, so that it fell to pieces before it was perfected.' Hereafter the common name of Puritans gave way to the party names of Presbyterians and Independents.

We may arrange the members of the Assembly under four sections:\footnote{Comp. the full accounts in Neal, Part III. ch. iv. (Vol. I. pp. 488 sqq.), Hetherington, Stoughton, and Masson.}

1. The Episcopalians. Parliament elected four prelates, viz.: James Ussher (Archbishop of Armagh and Bishop of Carlisle), Brownrigg (Bishop of Exeter), Westfield (Bishop of Bristol), Prideaux (Bishop of Worcester);\footnote{Prideaux's name seems to have been omitted in the final ordinance of June, 1643.} and five doctors of divinity, viz.: Drs. Featley (Provost of Chelsea College), Hammond (Canon of Christ's Church, Oxford), Holdsworth (Master of Emmanuel College, Cambridge), Sanderson (afterwards Bishop of Lincoln), and Morley (afterwards Bishop of Winchester). An excellent selection. But with one or two exceptions they never attended, and could not do so without disloyalty and disobedience to the king; besides, they objected to the company with an overwhelming number of Puritans, and a council not elected by the clergy and mixed with laymen. Ussher is said to have attended once, but on no good authority; he was present, however, in spirit, and great respect was paid to his theology by the Assembly.\footnote{Ussher was a second time appointed by the House of Commons a member of the Assembly when he came to London in 1647, and on his petition received permission to preach in Lincoln's Inn.—\emph{Journals of the House of Commons}, Vol. V. p. 423 (quoted by Dr. Mitchell).} Brownrigg sent in an excuse for non-attendance. Westfield was present, at least, at the first meeting. Dr. Featley, a learned Calvinist in doctrine, and a violent polemic against the Baptists, was the only Episcopalian who attended regularly and took a prominent part in the proceedings until, after the adoption of the Scotch Covenant, he was expelled by Parliament for revealing, contrary to pledge, the secrets of the Assembly in a letter to Ussher, then in the king's headquarters at Oxford, and was committed to prison (Sept. 30, 1643). This act of severity is strongly condemned by Baxter. Here ends the connection of Episcopacy with the Assembly.

Before this time Parliament had been seriously agitated by the Episcopal question. As early as Nov. 13, 1640, the 'Root and Branch' party sent in a petition signed by 15,000 Londoners for the total overthrow of the Episcopal hierarchy, while 700 clerical petitioners prayed merely for a reduction and modification of the same. Radicalism triumphed at last under the pressure of political necessity and the popular indignation created by Laud's heartless tyranny. First the bishops were excluded from the House of Lords (Feb. 5, 1642), with the reluctant assent of the king; and then the hierarchy itself was decreed out of existence (Sept. 10, 1642), the bill to take effect Nov. 5, 1643,\footnote{'An act for the utter abolishing and taking away of all archbishops, bishops, their chancellors and commissaries,' etc. Clarendon says that marvelous art was used, and that the majority of the Commons were really against the bill; but the writer of the 'Parliamentary Chronicle' says that it passed \emph{unanimously}, and was celebrated by bonfires and the ringing of bells all over London.—Neal, Vol. I. p. 421. Hallam also follows the latter account.} but the ordinances to carry this measure into full effect were not-passed till Oct. 9 and Nov. 16, 1646.\footnote{Neal, Vol. II. pp. 35 sq.} The old building was destroyed before a new building was agreed upon. This was the very question to be decided by the Assembly; hence the interval between the law and its execution. For nearly twenty years the Episcopal Church, though not legally abolished, from want of royal assent, was an \emph{ecclesia pressa el illicita} on her own soil.

Among the scores or hundreds of pamphlets which appeared in this war upon the bishops, the five anti-Episcopal treatises of John Milton were the most violent and effective. He attacked the English hierarchy, especially as it had developed itself under the Stuarts, with a force and majesty of prose which is unsurpassed even by his poetry. He went so far as to call Lucifer 'the first prelate-angel,' and treats Ussher with lofty contempt as a mere antiquarian or dryasdust. 'He rolls,' says his biographer, 'and thunders charge after charge; he tasks all his genius for epithets and expressions of scorn; he says things of bishops, archbishops, the English Liturgy, and some of the dearest forms of the English Church, the like of which could hardly be uttered now in any assembly of Englishmen without hissing and execration.'\footnote{Masson, Vol. II. p. 245. Comp. pp. 356 sqq., and the just estimate of Stoughton, \emph{The Ch. of the Civil Wars.} p. 129.}

2. The Presbyterians formed the great majority and gained strength as the Assembly advanced. Their Church polity is based upon the two principles of ministerial parity, as to ordination and rank (or the original identity of presbyters and bishops), and the self-government of the Church by representative judicatories composed of clerical and lay members. It was essentially the scheme of Calvin as it prevailed in the Reformed Churches on the Continent, and was established in Scotland.

The Scots seemed to be predestinated for Calvinistic Presbyterianism by an effective decree of Providence. The hostility of their bishops to the Reformation, and the repeated attempts of the Stuarts to force English institutions upon them, filled the nation with an intense aversion to Episcopacy and liturgical worship. Bishop Bancroft, of London, the first real High-Church Episcopalian, called English Presbyterianism an 'English Scottizing for discipline.'

In England, on the contrary, Episcopacy and the Prayer-Book were identified with the Reformation and Protestant martyrdom, and hence were rooted in the affections of the people. Besides, the early bishops were in fraternal correspondence with the Swiss Churches. But in the latter part of Elizabeth's reign, when Episcopacy took exclusive ground and rigorously enforced uniformity against all dissent, Presbyterianism began to raise its head under the lead of two eminent Calvinists, Thomas Cartwright (1535–1603), Professor of Theology in Cambridge, and Walter Travers (d. 1624), Preacher in the Temple, London, afterwards Provost of Trinity College, Dublin. The former was in conflict with the High-Churchism of Archbishop Whitgift;\footnote{Even Whitgift, however, did not go to the extreme of \emph{jure divino} Episcopacy, but admitted that the Scripture has not set down 'any one certain form of Church government to be perpetual.' Cartwright, on the other hand, was an able and earnest, but radical Presbyterian, and with Calvin and Beza advocated the death penalty for heretics.} the latter with the moderate Churchism of Richard Hooker, who was far his superior in ability, and whom he himself esteemed as 'a holy man.' The first English presbytery within the prelatic Church, as an \emph{ecclesiola} in \emph{ecclesia}, was formed at Wandsworth, in Surrey, in 1572, and Cartwright drew up for it a 'Directory of Church-Government,' or 'Book of Discipline,' in 1583, which is said to have been subscribed by as many as five hundred clergymen, and which was printed by authority of Parliament in 1644.\footnote{A fac-simile of this \emph{Directory} was reproduced in London, 1872 (James Nesbit \& Co.), for the tercentenary celebration of the Presbytery at Wandsworth, with an introduction by Prof. Lorimer. On Cartwright and the Elizabethan Presbyterianism, comp. Masson, \emph{Life of Milton}, Vol. II. pp. 581 sqq., and M'Crie, \emph{Annals of English Presbytery}, pp. 87–131.}

This anomalous organization was stamped out by authority, but the recollection of it continued through the reigns of James and Charles, and gathered strength with the rising Conflict.

The Westminster divines, with the exception of the Scotch Commissioners and two French Reformed pastors of London,\footnote{Samuel de la Place and Jean de la March.} were in Episcopal orders, and graduates of Oxford and Cambridge, and therefore as a body not opposed to Episcopacy as such. A goodly number inclined to Ussher's scheme of a 'reduced' or limited Episcopacy, \emph{i.e.}, a common government of the Church by presbyters under the supervision of the bishop as \emph{primus inter pares.}\footnote{\par \emph{The Reduction of Episcopacy unto the Form of Synodical Government received in the Ancient Church}, written in 1641, but not fully published till 1658, and brought forward again after the Restoration; in Ussher's \emph{Works} by Elrington, Vol. XII. Comp. Masson, Vol. II. p. 230.}

Had the moderate Episcopalians attended, the result would probably have been a compromise between Episcopacy and Presbytery. But the logic of events which involved Parliament in open war with the stubborn king, and necessitated the calling in of the aid of Presbyterian Scotland, changed the aspect of affairs. The subscription of the 'Solemn League and Covenant' (Sept., 1643) bound both the Parliament and the Assembly to the preservation of the doctrine, worship, and discipline of the Church of Scotland and the extirpation of popery and prelacy (\emph{i.e.}, the government of the Church by archbishops and bishops).

There were, however, two classes of Presbyterians, corresponding to the Low and High Church Episcopalians. The liberal party maintained that the Presbyterian form of government was based on human right, and 'lawful and agreeable to the Word of God,' but subject to change according to the wants of the Church. The high and exclusive Presbyterians of the school of Andrew Melville maintained that it was based on divine right, and 'expressly instituted or commanded' in the New Testament as the only normal and unchangeable form of Church polity. Twisse, Gataker, Reynolds, Palmer, and many others advocated the \emph{jus humanum} of Presbytery, all the Scotch Commissioners and the five 'Smectymnuans,'\footnote{The Smectymnuans were Stephen Marshall, Edmund Calamy, Thomas Young (the chief author), Matthew Newcomen, and William Spurstow. The oddity and ugliness of the title, composed of the initials of each author, helped the circulation and provoked witty rhymes, such as  \par 'The Sadducees would raise the question, \par Who must be \emph{Smec} at the resurrection.'} so called from their famous tract \emph{Smectymnuus}, in reply to Bishop Hall's defense of Episcopacy (March, 1641), advocated the \emph{jus divinum.} The latter triumphed, but for the sake of union they had to forego some details of their theory.\footnote{One of the dividing questions was that of ruling elders. 'Sundry of the ablest,' says Baillie (Vol. II. pp. 110 sq.), 'were flat against the institution of any such officer by divine right, such as Dr. Smith, Dr. Temple, Mr. Gataker, Mr. Vines, Mr. Price, Mr. Hall, and many more, besides the Independents, who truly spake much and exceedingly well. The most of the Synod was in our opinion, and reasoned bravely for it; such as Mr. Seaman, Mr. Walker, Mr. Marshall, Mr. Newcomen, Mr. Young, Mr. Calamy. Sundry times Mr. Henderson, Mr. Rutherford, Mr. Gillespie, all three spoke exceedingly well. When all were tired, it came to the question. There was no doubt but we would have carried it by far most voices; yet because the opposites were men very considerable, above all gracious little Palmer, we agreed upon a committee to satisfy, if it were possible, the dissenters.' He afterwards expresses the hope that the advance of the Scotch army 'will much assist our arguments.'}

The sequel, however, proved that Presbyterianism, so congenial to Scottish soil, was an artificial plant in England. Milton's prophetic words were fulfilled: 'Woe be to you, Presbyterians especially, if ever any of Charles's race recovers the English sceptre! Believe me, you shall pay all the reckoning.' Independency has ultimately far outgrown Presbytery, and is preferred by the English mind because it comes nearer to Episcopacy in making each pastor a bishop in his own congregation. Baxter says that Ussher agreed with the Independents in this, 'that every bishop was independent, and that synods and councils were not so much for government as concord.'\footnote{Quoted by Neal, Vol. I. p. 493.} If Presbyterianism has recently taken a new start and made great progress in London and other cities of England, it is owing mostly to the immigration of energetic and liberal Scotchmen and the high character of its leading ministers.

3. The Independents, called 'the five dissenting brethren' by the Presbyterians. They were, led by Dr. Thomas Goodwin and Rev. Philip Nye.\footnote{The others were Jeremiah Burroughs, William Bridge, and Sydrach Simpson. These five were the signers of the 'Apologetic Narration.' Afterwards William Carter, William Greenhill, John Bond (perhaps also Anthony Burgess), joined them. Baillie (Vol. II. p. 110) counts ten or eleven, including Carter, Caryl, Philips, and Sterry. Among its lay-assessors lord Viscount Say and Seale and Sir Harry Vane sympathized with the Independents. Neal says: 'Their numbers were small at first, though they increased prodigiously and grew to a considerable figure under the protectorship of Oliver Cromwell.'} Though small in number (twelve at the most), they were strong in ability, learning, and weight of character, and possessed the confidence of the rising Cromwell and the army, as well as the distant colonies in New England. Some of them had been driven to Holland by the persecution of Laud and Wren, and had administered to congregations of their expatriated countrymen, which occupied a middle ground between Brownism and Presbytery, after the model of John Robinson's pilgrims in Leyden. They were allowed the use of the Reformed churches, with liberty to ring the bell for service. After their return they advocated congregational independency and toleration, which the Presbyterians abhorred.\footnote{Baillie declares 'liberty of conscience and toleration of all or any religion' (as advocated by Roger Williams against John Cotton) to be 'so prodigious an impiety that this religious Parliament can not but abhor the very naming of it.\emph{'—Tracts on Liberty of Conscience} [published by the Hansard Knollys Society), p. 270, note. But Baillie was opposed to the employment of 'secular violence' in dealing with heretics. See M'Crie, p. 191.} The Independents maintained that a Christian congregation should consist of converted believers, and govern itself according to Christ's law, without being subject to the jurisdiction of presbyteries and synods, and that such a congregation had even a right to ordain its own minister. They fought the Presbyterians at every step on the questions of ruling elders, ordination, jurisdiction of presbyteries and synods, toleration, and threatened at times to break up the harmony of the meeting.

The longest debate, called 'the Grand Debate,' which lasted thirty days, was on the divine right of presbytery. And yet the two parties had great respect for each other. 'I wish,' said Gillespie, in the heat of the controversy, 'the dissenting brethren prove to be as unwilling to divide from us as we have been unwilling to divide from them. I wish that, instead of toleration, there may be a mutual endeavor for a happy accommodation,'\footnote{\emph{Minutes}, p. 28.}

The Independents appealed, rather inconsistently, to Cæsar, and addressed 'An Apologetic Narration to Parliament' (Dec., 1643). Under the Protectorate of Cromwell they became the ruling party, and had great political influence; but after the Restoration they resolved to seek for toleration outside of the National Church rather than for comprehension within it. New England was their Eldorado.\footnote{On the Independent controversy, see Baillie, Gillespie, and Masson (Vol. III. pp. 18 sqq.).}

4. The  Erastians\footnote{So called from the Swiss professor and physician, Erastus, properly Liebler, or Lieber, who wrote against Bullinger and Beza, and died at Basle, 1583.} maintained the ecclesiastical supremacy of the civil government in all matters of discipline, and made the Church a department of the State. They held that clergymen were merely teachers, not rulers, and that the power of the keys belonged to the secular magistrate. They hoped in this way to secure national unity and to prevent an \emph{imperium in imperio} and all priestly tyranny over conscience; but in fact they simply substituted a political for an ecclesiastical despotism, a cæsaropapacy for a hierarchical papacy. They were willing to submit to a \emph{jure humano} Presbyterianism, but they denied that any particular form of Church government was prescribed in the New Testament, and claimed for the State the right to establish such a form as might be most expedient.

The advocates of Erastianism in the Assembly were Selden, Lightfoot, and Coleman, all distinguished for Hebrew learning, which they used to good advantage. They appealed to the example of Moses and the kings of Israel, and the institutions of the Synagogue. They were backed by the lawyers among the lay-assessors and by the House of Commons, most of whom were (according to Baillie) 'downright Erastians.' The Assembly itself owed its existence to an act of Erastianism.

In strong opposition to them the Presbyterians maintained that the Lord Jesus, as sole King and Head of his Church, has appointed a spiritual government with distinct officers.

The controversy was ably conducted on both sides, and, we may say, exhausted.\footnote{The chief books on the Erastian side are Selden's \emph{De Synedriis} and Lightfoot's \emph{Journal;} on the Presbyterian side, Gillespie's \emph{Aaron's Rod Blossoming, or, the Divine Ordinance of Church-Government Vindicated} (dedicated to the Westminster Assembly; a very learned book of 590 pages), and Rutherford's \emph{Divine Right of Church Government} (both published in London, 1646). The Erastian controversy was afterwards transferred to Scotland, and led to several secessions. Comp. Principal Cunningham's Essay on the Erastian controversy in his \emph{Historical Theology}, Vol. II. pp. 557–588.}

The Independents and Erastians withdrew before the final adoption of the Book of Discipline, and left the field to the Presbyterians. The Presbyterian Church polity was at length established by the English Parliament, which ordained, June 29, 1647, that 'all parishes within England and Wales be brought under the government of congregational, classical, provincial, and national churches, according to the form of Presbyterial government agreed upon by the Assembly of Divines at Westminster.' Provinces were to take the place of dioceses, and were again divided into classes or presbyteries, and these were to elect representatives to a national assembly. But Parliament retained an Erastian power in its own hand, and would not permit even exclusion from the Lord's table without allowing to the offender recourse to the civil courts. Presbyterianism was nominally the established religion, but only in two provinces, London and Lancashire, was it fairly established, until its overthrow by the Restoration.\footnote{See M'Crie, pp. 189 sqq.}

THE LEADING MEMBERS.

Among the 121 divines of the Assembly there was a goodly portion of worthy and distinguished men who had suffered privation and exile under the misgovernment of Laud, who jeopardized their livings by accepting the appointment, notwithstanding the threats of the king, and who had the courage, after the Restoration, to sacrifice all earthly comforts to their conscientious convictions. Not a few of them combined rare learning, eloquence, and piety in beautiful harmony. 'The Westminster divines,' says Dr. Stoughton, 'had learning—Scriptural, patristic, scholastic, and modern—enough and to spare: all solid, substantial, and ready for use. Moreover, in the perception and advocacy of what is most characteristic and fundamental in the gospel of Jesus Christ they were as a body considerably in advance of some who could put in a claim to equal and perhaps higher scholarship.'\footnote{\emph{Church of the Civil Wars}, p. 453.}

It is sufficient for our purpose to mention the most eminent of the Westminster divines.\footnote{For a fall list of members, with biographical notices, the reader is referred to D. Masson, \emph{Life of John Milton}, Vol. II. pp. 516–524, where they are arranged in alphabetical order; and to Dr. Mitchell, in his Introduction to the \emph{Minutes}, pp. lxxxi.–lxxxiv., where they are given in the order of the ordinance of Parliament calling the Assembly (dated June 12, 1643), with some twenty members subsequently added to fill vacancies. Meek gives various lists in his edition of Gillespie's \emph{Notes.} Neal's list has several errors. Much information on the leading members may be gathered from Baillie's \emph{Journals}, Fuller's \emph{Church History} and \emph{Worthies of England}, Anthony Wood's \emph{Athenæ et Fasti Oxonienses}, Neal's \emph{History of the Puritans}, Stoughton's historical works, and Masson's \emph{Milton.} Reid gives biographical sketches of the Westminster divines in alphabetical order, with lists of their works.}

William Twisse, D.D. (Oxon.), Rector of Newbury, Prolocutor or Moderator by appointment of Parliament till his death (July, 1646). He was of German descent, about sixty-nine years of age, noted as a high Calvinist of the supralapsarian school, full of learning and subtle speculative genius, but 'merely bookish,' as Baillie says, and poorly fitted to guide a delicate assembly. Bishop Hall calls him 'a man so eminent in school-divinity that the Jesuits shrunk under his strength.' Thomas Fuller says: \footnote{\emph{Worthies of England}, Vol. I. p. 93. Dr. Owen, though he wrote against him, called him, 'the veteran leader, so well trained in the scholastic field; this great man; the very learned and illustrious Twisse.' M'Crie describes him as 'a venerable man, verging on seventy years of age, with a long, pale countenance, an imposing beard, lofty brow, and meditative eye; the whole contour indicating a life spent in severe and painful study' (\emph{Annals of the English Presbytery}, p. 145). The last words of Twisse were, 'Now at length I shall have leisure to follow my studies to all eternity.'} 'His plain preaching was good, solid disputing better, pious living best of all good.'

Charles Herle (d. 1659), an Oxford scholar, and Rector of Winwick in Lancashire, succeeded Twisse as Prolocutor. He was a moderate Presbyterian, and, in the language of Fuller, 'so much Christian, scholar, and gentleman that he could unite in affection with those who were disjoined in judgment from him.' He wrote against independency, but remarked in the Preface: 'The difference between us is not so great; at most it does but ruffle a little the fringe, not any way rend the garment of Christ.'\footnote{'The presence of such a man in the chair is sufficient to redeem the Assembly from the charge of illiberality or vulgar fanaticism.'—M'Crie, p. 151.}

John White (Oxon., d. 1648) and Dr. Cornelius Burgess (Oxon., d. 1665), the two Assessors, enjoyed general esteem. White was surnamed 'the patriarch of Dorchester,' but he 'would willingly contribute his shot of facetiousness on any just occasion' (Fuller). He was the great-grandfather of the Wesleys on the maternal side. Burgess was 'very active and sharp,' bold and fearless, an eminent debater and valiant defender of Presbyterianism and royalty.

Dr. Arrowsmith, head of St. John's College, Cambridge, 'a man with a glass eye,' having lost one by an arrow-shot, a 'learned divine' and 'elegant Latinist,' and long remembered in Cambridge for his 'sweet and admirable temper,' and Dr. Tuckney (d. 1670), Vice-Chancellor of the University, an inspiring teacher and bountiful friend of the poor, must be mentioned together as the chief composers of the Larger and Shorter Catechisms. They were both friends of the broad-minded Whichcote, who calls Arrowsmith 'the companion of his special thought.'\footnote{Tulloch, \emph{Rat. Theol. in England}, Vol. II. (the Cambridge Platonists), pp. 56 sq.} Dr. Tuckney, when requested by some members of Parliament to pay special regard to piety in his elections in Cambridge, made the sensible reply: 'No man has a greater respect than I have for the truly godly; but I am determined to choose none but scholars. They may deceive me in their godliness—they can not in their scholarship.' He is said to be the author of the exposition of the Ten Commandments in the Larger Catechism.

Edmund Calamy, B.D. (Cantab.), one of the four representatives of the London clergy, was a very popular preacher and a leader in the Presbyterian party. 'He was the first openly to avow and defend the Presbyterian government before a committee of Parliament; and though tempted afterwards with a bishopric, he continued stanch to his principles to his dying day.'\footnote{M'Crie, p. 155.} He died soon after the great fire in London (1666). His grandson, of the same name, was still more celebrated.

Joseph Caryl, M.A. (Oxon., 1602–1673), was a moderate Independent, a distinguished preacher, and 'a man of great learning, piety, and modesty' (Neal). He became afterwards one of Cromwell's Triers, was ejected in 1662, and lived privately, preaching to his congregation as the times would permit. He is chiefly known as the indefatigable author of a commentary on Job, in twelve volumes, 4to (Lond. 1648–1666), which is an excellent school of its chief topic, the virtue of patience.\footnote{Another edition in two large folio vols. was published in 1676 sq. Darling calls this exposition 'a most elaborate, learned, judicious, and pious work.'}

Thomas Coleman (Oxon.) was called 'Rabbi Coleman' for his profound Hebrew learning. Baillie describes him as half-scholar and half-fool, and of small estimation. He died during the heat of the Erastian debate (1647).

Thomas Gataker, B.D. (Cantab., d. 1654, \emph{aet.} eighty), a devourer of books, and equally esteemed for learning, piety, and sound doctrine. He refused various offers of preferment.

Thomas Goodwin, D.D. (Cantab., d. 1680, \emph{aet.} eighty), one of the two •patriarchs of English Independency,' Philip Nye being the other. He was Vicar of Trinity Church, Cambridge, relinquished his preferments in 1634, was pastor of a congregation of English exiles at Arnheim, Holland, then in London,\footnote{He founded a Congregational church in London in 1640, which continues to this day, and has recently (under the pastorate of Dr. Joseph Parker) erected the City Temple, with a memorial tablet to Goodwin in the vestibule.} and afterwards President of Magdalen College, in Oxford, till the Restoration, when he resigned. He was the favorite minister of Cromwell, eloquent in the pulpit, orthodox in doctrine, and exemplary in life, but 'tinctured with a shade of gloom and austerity' (M'Crie). 'Though less celebrated than Owen, his great attainments in scholarship and the range and variety of his thoughts astonish us when we read his writings, showing how familiar he was with all forms of theological speculation, ancient and modern' (Stoughton).\footnote{His austerity gave rise to the story related by Addison, in the \emph{Spectator}, that Dr. Goodwin, 'with half-a-dozen night-caps on his head and religious horror in his countenance,' overawed and terrified an applicant for examination in Oxford by asking him in a sepulchral voice, 'Are you prepared for death?' His works were published in London, 1681–1704, in 5 vols.}

Dr. Joshua Hoyle (Oxon., d. 1654), Divinity Professor in Dublin, afterwards Master of University College, Oxford, was the only Irish divine of the Assembly, 'a master of the Greek and Latin fathers,' who 'reigned both in the chair and in the pulpit.'

John Lightfoot, D.D. (Cantab.), the greatest rabbinical scholar of his age, whose \emph{Horæ Hebraicæ et Talmudicæ} are still familiarly quoted in illustration of the New Testament. His \emph{Journal} is one of the chief sources for the history of the Assembly, especially for exegetical and antiquarian aspects of the Erastian controversy. In 1649 he became Master of Catharine Hall, Cambridge, and retained his post till he died, 1675, aged seventy-three.

Stephen Marshall, B.D. (Cantab.), Lecturer at St. Margaret's, Westminster, was 'the best preacher in England' (Baillie), a fearless leader in the political strife, a great favorite in the Assembly, 'their trumpet, by whom they sounded their solemn fasts' (Fuller). One of his royalist enemies called him 'the Geneva bull, a factions and rebellions divine.' He was buried in Westminster Abbey, 1655, but disinterred with the other Puritans after the Restoration.

Philip Nye (Oxon., d. 1672), minister of Kimbolton, who had been in exile with his friend Goodwin, took a leading part, as a Commissioner of Parliament, in soliciting the assistance of the Scots, and securing subscription to the Covenant; but he conceived a dislike to their Church polity and gave them a world of trouble. He kept them for three weeks debating on the superior propriety, as he contended, of having the elements handed to the communicants in their own seats instead of calling them out to the table. He was a stanch Independent, a keen debater, and a 'great politician, of uncommon depth, and seldom if ever outreached' (Neal). He was one of the Triers under Cromwell, and the leader of the Congregational Savoy Conference. After the Restoration lie declined tempting offers, and preached privately to a congregation of Dissenters till he died, seventy-six years of age.

Herbert Palmer, B.D. (Cantab.), Vicar of Ashwell, afterwards Master of Queen's College, Cambridge, was a little man with a childlike look, but very graceful and accomplished, a fluent orator in French as well as English, and a model pastor. He spent his fortune in works of charity, and his delicate frame in the cure of souls. He had scruples about the divine right of ruling elders, but became a convert to Presbyterianism. He is the real author of the 'Christian Paradoxes,' which have so long been attributed to Lord Bacon.\footnote{This fact has recently been discovered by Rev. A. B. Grosart (1864). See Masson, Vol. II. p. 520.}

Dr. Edward Reynolds (Oxon., d. 1676), 'the pride and glory of the Presbyterian party' (Wood), was very learned, eloquent, cautious, but lacking backbone. He accepted from Charles II. the bishopric of Norwich (Jan., 1660), owing, it was said, to the influence of 'a covetous and politic consort' (Wood); but 'he carried the wounds of the Church in his heart and in his bowels to the grave with him.'

Sir Francis Rous (or Rowse, b. 1579, d. 1659), 'an old, most honest' member of Parliament, afterwards a member of Cromwell's Privy Council, was one of the twenty Commoners who were deputed to the Assembly. He innocently acquired an immortal fame by his literal versification of the Psalms, which was first printed in 1643, then revised, and is used to this day in Scotland and in many Presbyterian congregations in America in preference to all other versions and hymns.\footnote{See Baillie, Vol. II. p. 120; Vol. III. pp. 532 sqq.; and the \emph{Minutes }of the Westminster Assembly, pp. 131, 163, 418.}

Lazarus Seaman, B.D. (Cantab., 1667), one of the four representatives of the London clergy, a very active member and reputed as an Orientalist, who always carried with him a small Hebrew Bible without points. He is described as 'an invincible disputant' and 'a person of most deep, piercing, and eagle-eyed judgment in all points of controversial divinity, in which he had few equals, if any superiors.' He became Master of Peterhouse, Cambridge, but was ejected after the Restoration.

John Selden (1584–1654), one of the lay assessors, and a scholar and wit of European reputation.\footnote{\emph{Opera omnia}, ed. Dav. Wilkins, London, 1726, 3 vols. in folio.} His scholarship was almost universal, but lay chiefly in languages, law, and antiquities (hence '\emph{antiquariorum coryphæus}'). For a long time he took an active part in the debates, and often perplexed the divines by raising scruples. He liked to correct their 'little English pocket Bibles' from the Greek and Hebrew. Not especially fond of the flesh of the Scriptures, he cast the 'bones' at them 'to break their teeth therewith' (Fuller). He was an Erastian and a clergy-hater, but on his death-bed he declared that 'out of the numberless volumes he had read, nothing stuck so close to his heart, or gave him such solid satisfaction, as the single passage of Paul, 'The grace of God that bringeth salvation hath appeared unto all men.'

Richard Vines, Master of Pembroke Hall, Cambridge (d. 1656), 'an excellent preacher and very powerful in debate, and much respected on all accounts' (Masson).

Thomas Young, Master of Jesus College, Cambridge, a Scotchman by birth, Milton's preceptor, and the chief of the five 'Smectymnuans.'

THE SCOTCH COMMISSIONERS.

After the adoption of the international League and Covenant, Scotland sent five clerical and three lay commissioners who admirably represented their Church and country. They formed a group by themselves at the right hand of the Prolocutor. They were the only delegates who were elected by proper ecclesiastical authority, viz., the General Assembly of their Church (Aug. 19, 1643), at the express request of the English Parliament; they declined being considered members in the ordinary sense, but they were allowed by warrant of Parliament to be present and to debate, and practically they exerted an influence disproportionate to their number. They arrived in London in September, fresh from the battle 'with lordly bishops, popish ceremonies, and royal mandates,' and full of the '\emph{perfervidum ingenium Scotorum.}'

Alexander Henderson, Rector of the University of Edinburgh since 1640, sixty years of age, ranks next to John Knox and Andrew Melville in the history of Scotch Presbyterianism, and was the author of the 'Solemn League and Covenant,' which linked the Scottish and English nations in a civil and religious alliance for the Reformed religion and civil liberty. Being unmarried, he gave himself entirely to the Assembly from Aug., 1643, to Aug., 1646. He has heretofore been too much ignored. 'My researches,' says Masson,\footnote{Vol. III. p. 16.} 'have more and more convinced me that he was, all in all, one of the ablest and best men of his age in Britain, and the greatest, the wisest, and most liberal of the Scottish Presbyterians. They all had to consult him; in every strait and conflict he had to be appealed to, and came in at the last as the man of supereminent composure, comprehensiveness, and breadth of brow. Although the Scottish Presbyterian rule was that no churchman should have authority in State affairs, it had to be practically waived in his case; he was a cabinet minister without office.'

Robert Baillie (b. 1599, d. 1662), Professor of Divinity and Principal of the University of Glasgow, did not speak much, but was a regular attendant for fully three years, a shrewd observer, and has been called the Boswell of the Assembly and 'the pleasantest of letter gossips.' His 'Letters and Journals' (not properly edited until 1842) are among the most graphic books of contemporary memoir to be found in any language. His faculty of narration in his pithy native Scotch is nothing short of genius. Whenever we have an account from Baillie of any thing he saw or was present at, it is worth all accounts put together for accuracy and vividness; so in his accounts of Strafford's trial, and so in his account of his first impressions of the Westminster Assembly' (Masson).

George Gillespie, minister of Edinburgh (d. 1648), Was only thirty-one years of age when he entered the Assembly, the youngest, and yet one of the brightest stars, 'the prince of disputants, who with the fire of youth had the wisdom of age.' He first attracted public attention in his twenty-fourth year by 'A Dispute against the English-Popish Ceremonies obtruded upon the Church of Scotland' (1637), which helped the revolt against Laud's innovations. He took a leading part in the debates of the Assembly against Erastianism and Independency. According to Scotch tradition he once made even Selden reel and say, 'That young man, by his single speech, has swept away the labors of ten years of my life.' This is probably a patriotic exaggeration. The excessive ardor and activity of his mind wore out his frame, and he returned from the Assembly to die in his native land.

Samuel Rutherford (1600–1661), Professor of Divinity and Principal of St. Mary's College in St. Andrews, was one of the most fervid and popular preachers in Scotland, and highly esteemed for his learning and piety. 'The characteristics of his mind were clearness of intellect, warmth and earnestness of affection, and loftiness and spirituality of devotional feeling.' His book, 'Lex Rex,' is considered one of the best expositions of the principles of civil and religious liberty; and his glowing letters of comfort from his prison in Aberdeen (which he called 'Christ's Palace') show him to be 'the true saint and martyr of the Covenant.'

Rev. Robert Douglas never sat. Among the lay commissioners, John Lord Maitland (afterwards Earl of Lauderdale) distinguished himself first by his zeal for the Scotch Covenanters, and afterwards by his apostasy and cruelty against them. Sir Archibald Johnstone, of Warristone, was from 1637 a leader among the Scotch Covenanters, a great lawyer, and a devout Christian, who, as Bishop Burnet, his nephew, narrates, often prayed in his family two hours at a time with unexhausted copiousness. The Marquis of Argyle also, who afterwards suffered death for his loyalty to the Scotch Kirk, sat for some time as an elder in the Assembly.

OPENING OF THE ASSEMBLY.

The Assembly was opened on Saturday, July 1, 1643, in the grand national Abbey of Westminster, in the presence of both Houses of Parliament and a large congregation, by a sermon of Dr. Twisse on John xiv. 18: 'I will not leave you comfortless; I will come unto you'—a text which was deemed 'pertinent to these times of sorrow, anguish, and misery, to raise up the drooping spirits of the people of God who lie under the pressure of Popish wars and combustions.'\footnote{From the Parliamentarian newspaper No. 25, for July 3–10, 1643, quoted by Mitchel, p. xi. Lightfoot reports in his Journal (p. 3) that 'a great congregation' was present besides the members of the Assembly and of Parliament.}

After service the members of the Assembly, 'three score and nine'\footnote{This is about the average attendance of the Lower House of the Convocation of Canterbury,—Stanley, \emph{Memorials of Westminster Abbey}, p. 507.} (twenty-nine more than the required quorum), repaired for organization to the Chapel of Henry VII., that 'most gorgeous of sepulchres,' where the Upper House of Convocation used to meet. The mediæval architecture formed a striking contrast to the Puritan simplicity of worship and dress. The divines appeared in black coats or cloaks, skull-caps, and Geneva bands in imitation of the foreign Protestants,\footnote{Neal and Stoughton.} with the exception of a few Royalists and Episcopalians, who in their canonical gowns seemed 'the only non-Conformists.'\footnote{Fuller.} Add to this apparel their solemn looks, the peaked beards and mustaches, and the broad double ruff around the neck, and we have a spectacle of a synod differing as much from a modern Presbyterian Assembly as from an Episcopal Convocation or a Roman Catholic Council.\footnote{M'Crie and Mitchell compare it to a synod of Huguenots as pictured on the title-page of the first volume of Quick's \emph{Synodicon.} But there the Frenchmen wear broad-brimmed hats.}

Every member had to take the following vow (which was read in the Assembly every Monday morning):

'I do seriously promise and vow, in the presence of almighty God, that in this Assembly, whereof I am a member, I will maintain nothing in point of doctrine but what I believe to be most agreeable to the Word of God; nor in point of discipline, but what may make most for God's glory and the peace and good of his Church.'

THE ASSEMBLY IN THE JERUSALEM CHAMBER.

For several weeks the meetings were held in the Chapel of Henry VII. But when extreme cold weather set in at the close of September, the Assembly repaired to the 'Jerusalem Chamber,' in the Deanery of Westminster.\footnote{The origin of the name is uncertain. Some derive it from the tapestries or pictures of Jerusalem on the wall. Dr. Stoughton, who is well informed in English history and archaeology, informs me (by letter of May 4, 1876) that it probably arose 'from the fact of its adjoining the sanctuary, the place of peace;' and he quotes a passage from the account of King John's death: '\emph{Nec providet quod est Romæ ecclesia Jerusalem dicta, id est, visio pacis; quia quicunque illuc confugerit, cuiuscunque criminis obnoxius, subsidium invenit}' (William of Malmesbury, \emph{De gestis Angl.} Lib. II. p. 67).} 'What place more proper for the building of Sion,' asks Fuller, 'than the Chamber of Jerusalem, the fairest of the Dean's lodgings, where King Henry IV. died, and where these divines did daily meet together?'\footnote{\emph{Church Hist.} Vol. VI. p. 253.}

This large and venerable hall, furnished with a long table and chairs, and ornamented with tapestry (pictures of the Circumcision, the Adoration of the Magi, and the Passage through the Wilderness), was originally the withdrawing-room of the abbot, and has become famous in romance and history as the cradle of many memorable schemes and events from the Reformation down to the present time.

There, before the fire of the hearth—then a rare luxury in England—King Henry IV., who intended to make a pilgrimage to Jerusalem, died of a hideous leprosy (March 20, 1413). When informed of the name of the chamber, he exclaimed,

'Laud be to God! even there my life must end.

It hath been prophesied to me many years

I would not die but in Jerusalem;

Which vainly I supposed the Holy Land.

But bear me to that chamber; there I'll lie:

In that Jerusalem shall Harry die.'\footnote{Shakspere, Second Part of King Henry IV., act iv. sc. 4.}

There Sir Thomas More was confined (1534), and urged by the abbot to acknowledge the king's ecclesiastical supremacy; and there probably he wrote his appeal to a general council which never met, but may yet meet at some future day.

There, under the genial warmth of the fire which had attracted the dying king, the grave Puritan Assembly prepared its standards of doctrine, worship, and discipline, to be disowned by England, but honored by Scotland and America.

There the most distinguished Biblical scholars of the Church of England, in fraternal co-operation with scholars of Dissenting denominations, both nobly forgetting old feuds and jealousies, are now engaged in the truly catholic and peaceful work of revising the common version of the Bible for the general benefit of English-speaking Christendom.\footnote{For a fuller description of the Jerusalem Chamber, see Dean Stanley's \emph{Memorials of Westminster Abbey}, pp. 417 sqq. I may be permitted to add from personal experience an interesting recent incident in the history of that chamber. At the kind invitation of the Dean of Westminster, the delegates to the International Council of Presbyterian Churches, then meeting in London for the formation of a Presbyterian Alliance, repaired to the Jerusalem Chamber on Thursday afternoon, July 22, 1875, and, standing around the long table, were instructed and entertained by the Dean, who, modestly taking 'the Moderator's chair,' gave them a graphic historical description of the chamber, interspersed with humorous remarks and extracts from Baillie. He dwelt mainly on the Westminster Assembly, promising, in his broad-Church liberality, at some future time to honor that Assembly by a picture on the northern wall. Dr. McCosh, as Moderator of the Presbyterian Council, proposed a vote of thanks for the courtesy and kindness of the Dean, which was, of course, unanimously and heartily given. The writer of this expressed the hope that the Jerusalem Chamber may yet serve a still nobler purpose than any in the past, namely, the reunion of Christendom on the basis of God's revealed truth in the Bible; and he alluded to the fact that the Dean had recently (in the 'Contemporary Review,' and in an address at Saint Andrews) paid a high compliment to the Westminster Confession by declaring its first chapter, on the Holy Scriptures, to be one of the best, if not the very best symbolical statement ever made.}

BAILLIE'S DESCRIPTION OF THE ASSEMBLY.

The Assembly in actual session in this famous locality, and its order of business, can not be better described than in the graphic language of one of the Scotch Commissioners:

'The like of that Assembly, 'says Professor Baillie,\footnote{In a letter to his cousin, William Spang, dated London, Dec. 7, 1643. See \emph{Letters and Journals}, Vol. II. pp. 107–109. I have retained the Scotch words, but modernized the spelling. Extracts from this letter are quoted by Neal, Hetherington, Stanley, Stoughton, Mitchell.} 'I did never see, and, as we hear say, the like was never in England, nor any where is shortly like to be. They did sit in Henry the Seventh's Chapel, in the place of the Convocation; but since the weather grew cold, they did go to Jerusalem Chamber, a fair room in the Abbey of Westminster, about the bounds of the College forehall, but wider. At the one end nearest the door and on both sides are stages of seats as in the new Assembly-House at Edinburgh, but not so high, for there will be room but for five or six score. At the upmost end there is one chair set on a frame, a foot from the earth, for the Mr. Prolocutor Dr. Twisse. Before it, on the ground, stand two chairs for the two Mr. Assessors, Dr. Burgess and Mr. White. Before these two chairs, through the length of the room, stands a table, at which sit the two scribes, Mr. Byfield and Dr. Roborough. The house is all well hung and has a good fire, which are some dainties at London. Foranent [in front of] the table, upon the Prolocutor's right hand, there are three or four ranks of forms. On the lowest we five do sit. Upon the other, at our backs, the members of Parliament deputed to the Assembly. On the forms foranent us, on the Prolocutor's left hand, going from the upper end of the house to the chimney, and at the other end of the house, and backside of the table, till it comes about to our seats, are four or five stages of forms, whereupon their divines sit as they please, albeit commonly they keep the same place. From the chimney to the door there are no seats, but a void for passage. The Lords of Parliament use to sit on chairs in that void, about the fire. We meet every day of the week but Saturday. We sit commonly from nine to one or two [in the] afternoon. The Prolocutor at the beginning and end has a short prayer. The man, as the world knows, is very learned in the questions he has studied, and very good, beloved by all, and highly esteemed; but merely bookish, and not much, as it seems, acquainted with conceived prayer, [and] among the unfittest of all the company for any action; so after the prayer he sits mute. It was the canny conveyance of those who guide most matters for their own interest to plant such a man of purpose in the chair. One of the Assessors, our good friend Mr. White, has keeped in of the gout since our coming; the other, Dr. Burgess, a very active and sharp man, supplies, so far as is decent, the Prolocutor's place.

'Ordinarily there will be present above threescore of their divines. These are divided into three committees, in one whereof every man is a member; no man is excluded who pleases to come to any of the three. Every committee, as the Parliament gives order in writing to take any purpose into consideration, takes a portion, and in their afternoon meeting prepares matters for the Assembly, sets down their mind in distinct propositions, [and] backs their propositions with texts of Scripture. After the prayer, Mr. Byfield, the scribe, reads the proposition and Scriptures, whereupon the Assembly debates in a most grave and orderly way. No man is called up to speak; but who stands up of his own accord, he speaks so long as he will without interruption. If two or three stand up at once, then the divines confusedly call on his name whom they desire to hear first: on whom the loudest and maniest [most] voices call, he speaks. No man speaks to any but to the Prolocutor. They harangue long and very learnedly. They study the questions well beforehand, and prepare their speeches; but withal the men are exceeding prompt and well-spoken. I do marvel at the very accurate and extemporal replies that many of them usually do make. When, upon every proposition by itself, and on every text of Scripture that is brought to confirm it, every man who will has said his whole mind, and the replies, and duplies, and triplies are heard, then the most part calls ``To the question.'' Byfield, the scribe, rises from the table and comes to the Prolocutor's chair, who, from the scribe's book, reads the proposition, and says, ``As many as are of opinion that the question is well stated in the proposition, let them say I;'' when I is heard, he says, ``As many as think otherwise, say No.'' If the difference of I's and No's be clear, as usually it is, then the question is ordered by the scribes, and they go on to debate the first Scripture alleged for proof of the proposition. If the sound of I and No be near equal, then says the Prolocutor, ``As many as say I, stand up;'' while they stand, the scribe and others number them in their mind; when they sit down the No's are bidden to stand, and they likewise are numbered. This way is clear enough, and saves a great deal of time, which we spend in reading our catalogue. When a question is once ordered, there is no more debate of that matter; but if a man will vaige,'\footnote{Probably 'wander' (from 'vague').} he is quickly taken up by Mr. Assessor, or many others, confusedly crying, ``Speak to order, to order.'' No man contradicts another expressly by name, but most discreetly speaks to the Prolocutor, and at most holds on the general—The reverend brother, who lately or last spoke, on this hand, on that side, above, or below.

'I thought meet once for all to give you a taste of the outward form of their Assembly. They follow the way of their Parliament. Much of their way is good, and worthy of our imitation: only their longsomeness is woeful at this time, when their Church and Kingdom lies under a most lamentable anarchy and confusion. They see the hurt of their length, but can not get it helped; for being to establish a new Platform of worship and discipline to their nation for all time to come, they think they can not be answerable if solidly and at leisure they do not examine every point thereof.'

DEVOTIONAL EXERCISES.

With theological discussion the Assembly combined devotional exercises, and observed with Parliament regular and occasional fasts which are characteristic of the Puritan piety of that age. At the joint meeting of the Parliament and the Assembly in St. Margaret's Church, for the signing of the Covenant (Monday, Sept. 25, 1643), Mr. White 'prayed near upon an hour,' Mr. Nye 'made an exhortation of another hour long,' Mr. Henderson 'did the like;' then there was the reading of the Covenant, a prayer by Dr. Yonge, 'another psalm by Mr. Wilson,' and a concluding prayer, when they 'adjourned till Thursday morning, because of the fast.'\footnote{Lightfoot, \emph{Journal}, p. 16.}

Baillie describes the fast observed May 17, 1644, at the request of General Essex before his march into the field, as 'the \emph{sweetest} day' he saw in England, although it lasted eight hours, from nine to five, without interruption. 'After Dr. Twisse,' he writes, 'had begun with a brief prayer, Mr. Marshall prayed large two hours, most divinely, confessing the sins of the members of the Assembly in a wonderfully pathetic and prudent way. After, Mr. Arrowsmith preached one hour; then a psalm; thereafter, Mr. Vines prayed near two hours, and Mr. Palmer preached one hour, and Mr. Seaman prayed near two hours; then a psalm. After, Mr. Henderson brought them to a short, sweet conference of the heart confessed in the Assembly, and other seen faults\footnote{Probably a misprint for 'heart-confessed and other seen faults in the Assembly.'} to be remedied, and the convenience to preach against sects, especially Anabaptists and Antinomians. Dr. Twisse closed with a short prayer and blessing. God was so evidently in all this exercise that we expect certainly a blessing both in our matter of the Assembly and whole kingdom.'\footnote{\emph{Letters and Journals}, Vol. II. pp. 184 sq.}

We can not read such accounts without amazement at the devotional fervor and endurance of the Puritan divines. And yet, if we consider the length of their prayers and sermons, their austerity in society, dress and manner, their peculiar phraseology and cant, their aversion to the fine arts and public amusements, however innocent, we need not be surprised at the popular rebound to the opposite extreme under the frivolous and licentious Charles II. 'All that was beautiful in Church music, architecture, or ornament, and in personal elegance and refinement, was rigidly proscribed. Even poetry was at a discount; Milton himself, in his lifetime, in more senses than one, ``sung darkling;'' and the literary style, of the day, unlike either that of the foregoing or the subsequent age, was harsh, stiff, and void of elegance. Even the typography of the period is peculiarly grim and unseemly.'\footnote{M'Crie, \emph{Annals of English Presb.} p. 173. The last remark applies also to the early editions of the Westminster standards and controversial pamphlets.}

It should not be forgotten, however, that there are times when aesthetics must give way to more important matters, and that radical extremes are unavoidable in critical periods. The Catholic Church itself, in the first three centuries, passed through the gloom of the catacombs, and, in its ascetic abhorrence of heathen art and beauty, strangely misconceived even our blessed Lord's personal appearance as homely and repulsive in the days of his humiliation. Tertullian, in his way, went farther than the Puritans.

DURATION AND CLOSE.

The Assembly occupied about five years and six months for the completion of its proper work—the standards of doctrine, worship, and discipline—and held no less than 1163 regular sessions from July 1, 1643, till February 22, 1649, when it ought to have adjourned \emph{sine die.} It met every day, except Saturday and Sunday, from nine o'clock till one or two—the afternoons being left to committees. After Nov. 9, 1647, we find no mention of the Scotch Commissioners. But the Assembly continued to drag out a shadowy existence, with scanty and irregular attendance, as a standing committee for the examination and ordination of candidates for the ministry, meeting every Thursday,\footnote{The sessions held after Feb. 22, 1649 (1648), are not numbered. The last regular meetings were likewise devoted merely to executive business. See \emph{Minutes}, p. 539.} till March 25, 1652, when it informally broke up before the dissolution of the 'Rump' Parliament by Oliver Cromwell (April 19, 1653). 'It dwindled away by degrees, though never legally dissolved,' says Fuller. It vanished with the Long Parliament which gave it birth.

\xsubsection{The Westminster Confession.}

§ 94. The Westminster Confession.

I. Standard Editions.

1. \emph{English.}

The editio princeps, without Scripture texts, was printed, but not published, Dec. 7, 1646, at London, under the title, '\emph{The Humble | Advice | of the | Assembly | of | Divines, | Now by authority of Parliament | sitting at Westminster, | concerning | a Confession of Faith, | presented by them lately to both Houses | of Parliament.} | ... London. Printed for the Company of Stationers.' 1647.

A second edition (of 600 copies) was printed in London, under the same title, 'with the Quotations and Texts of Scripture annexed,' by order of Parliament, dated April 29, 1647.

The first Edinburgh ed. is a reprint of the second London ed. in somewhat different type. Only 300 copies were printed, Aug. 9, 1647, for the use of the General Assembly. See fac-simile in Vol. III. p. 598.

The typography and paper of these early editions are very poor. After the adoption, innumerable editions appeared under the proper title, 'Confession of Faith.' The earliest \emph{small} ed. of \emph{Edinb.} appeared 1650; the earliest \emph{small} ed. in Lond., 1648 or 1649. See \emph{Minutes}, p. 418, note 4.

The edition which was adopted by the English Parliament, with some changes (similar to those afterwards made in the Savoy Declaration), bears a different title, viz.: Articles | \emph{of | Christian Religion, | Approved and Passed by both Houses | of } Parliament, | \emph{After Advice had with the Assembly | of } | Divines | \emph{by | Authority of Parliament sitting at | Westminster.} | London: | . . . June 27, 1648.

Copies of the earliest and other rare editions I found and compared in the British Museum, in the Libraries of Edinburgh, the Free Church College and the Advocates' Libraries, and that of Union Theol. Seminary in New York. The texts vary but slightly. I used also a London ed. of 1658 (pp. 108), which is a little superior in typography, and still bears the title \emph{Humble Advice}, etc. It has the Scripture proofs printed out in full.

Prof. Mitchell proposes to publish, with other documents, 'a careful collation of the earlier editions of the Confession' (\emph{Minutes}, p. 546).

A very good edition of the Westminster Confession of Faith and Catechisms, together with the Covenants (National and Solemn League), the acts of Parliament and the General Assembly relative to and approving of the same, was printed by authority at Edinburgh (University Press), 1858 (pp. 561).

The American editions differ from the English and Scotch in Chaps. XXIII. and XXXI., and in the close of XX. The changes are given in Vol. III. pp. 600 sqq.

2. \emph{Latin.}

\emph{Confessio Fidei in Conventu theologorum authoritate Parliamenti Anglicani indicto elaborata; eidem Parliamento postmodum exhibita; quin et ab eodem, deinque ab Ecclesia Scoticana cognita et approbata; una cum Catechismo duplici, majori, minorique; e sermone Anglicano summa cum fide in Latinum versa.} Cantabrigiæ, 1656 and 1659, small 8vo (229 pp.). Other eds., Edinb. 1670, 1694, 1708, 1711; Glasgow, 1660; in the appendix to Niemeyer's \emph{Collectio Conf.} 1840. See Vol. III. pp. 600 sqq. The translation is good, but the translator is not named, nor could I ascertain his name from the librarians in Edinburgh and London, not even from the learned Mr. David Laing and Dr. Mitchell. The initials below the preface are 'G. D.' (perhaps G. Dillingham, D.D., of Emanuel College, Cambridge; others surmised G. Duport, of Cambridge).

3. \emph{German.}

A German translation appeared as early as 1648. A new one in Böckel's  \emph{Bekenntniss-Schriften der evang. reform. Kirche}, pp. 683 sqq. (under the title \emph{Das puritanische Glaubensbekenntniss}). Another version is published by the Presbyterian Board in Philadelphia.

Historical.

See Literature on Westminster Assembly, § 93.

Dr.  Alex. F. Mitchell (Prof. of Ch. Hist, in St. Andrews): \emph{The Westminster Confession of Faith: a Contribution to the Study of its Historical Relations and to the Defence of its Teaching.} Edinb. 3d ed. 1867. Comp. his valuable Introduction to the \emph{Minutes}, 1874.

Alex. Taylor Innes: \emph{The Law of Creeds in Scotland.} Edinburgh, 1867.

Explanatory and Apologetic.

\emph{Truth's Victory over Error; or, an Abridgment of the chief Controversies in Religion}, etc. [By  David Dickson.] Edinb. (1649), 1684; Glasgow, 1725. A catechetical exposition of the Westm. Conf.

\emph{A Brief Sum of Christian Doctrine contained in Holy Scripture, and holden forth in the Confession of faith and Catechisms of the Westminster Assembly}, etc. [Drawn up by David Dickson.] Edinb. 1693.

Robert Shaw (Minister of the Free Church at Whitburn): \emph{An Exposition of the Confession of Faith of the Westminster Assembly of Divines.} With an Introduction by W. M. Hetherington. Edinb. 1845.

Archibald Alexander Hodge, D.D. (Prof. of Theol. in Allegheny Seminary): \emph{A Commentary on the Confession of Faith.} Philad. 1869 (Presbyt. Board).

Critical and Polemical.

W. Parker: \emph{The late Assembly of Divines' Conf. of Faith Examined, wherein many of their Excesses and Defects, of their Confusions and Disorders, of their Errors and Contradictions, are presented.} Lond. 1651.

James Stark: \emph{The Westminster Confession of Faith critically Compared with the Holy Scripture and found wanting.} Lond. 1863. A candid but captious critique of all the chapters.

Joseph Taylor Goodsir: \emph{The Westminster Confession of Faith Examined on the Basis of the other Protestant Confessions.} Lond. 1868. Directed chiefly against Ch. XI., on Justification by Faith.

A. M. Fairbairn: \emph{The Westminster Confession of Faith and Scotch Theology.} An article in the 'Contemporary Review,' answered by Prof. Mitchell in the Introduction to \emph{Minutes of the Westminster Assembly.}

William Marshall: \emph{The Principles of the Westminster Standards Persecuting.} Edinb. 1873.

REVISION OF THE ENGLISH ARTICLES.

The Assembly was at first employed for ten weeks on a revision of the Thirty-nine Articles of the Church of England, being directed by an order of Parliament (July 5, 1643) 'to free and vindicate the doctrine of them from all aspersions and false interpretations.' The Puritans regarded the doctrinal Articles as sound and orthodox in substance and spirit, but capable of improvement in the line marked out by the Lambeth Articles and the Irish Articles; in other words, they desired to make them more explicitly Calvinistic.

Fifteen of these Articles, including the most important doctrines, were thus revised, and provided with Scripture proofs.\footnote{The revised Fifteen Articles have been reprinted from the copy as approved by Parliament, in Hall's \emph{Harmony of Protestant Confessions;} in Appendix No. VII. to Neal's \emph{History of the Puritans;} in Stoughton, \emph{Church of the Commonwealth}, Append. pp. 228 sqq.} Very few changes were made. Art. I., on the Trinity, was left untouched. In Art. II., on the Son of God, the word 'all' before 'actual sins of men' is missing, which, if not an oversight, was a misimprovement in the interest of Calvinistic particularism.\footnote{The 'all' was in the original edition of 1563 and the edition of 1628, but is missing in the edition of 1630 and other English editions, and also in the American Episcopal revision; see Vol. III. p. 478.} In Art. III. the unhistorical interpretation of Christ's descent into Hades, which makes it a mere repetition of the preceding clause in the Creed, is put in. In Art. VI. the allusion to the Apocrypha is omitted. The remaining Articles are retained with some verbal improvements, except Art. VIII. of the three Creeds, which is omitted in almost all the printed copies. But in the original copy which the Assembly sent to Parliament, Art. VIII. was retained with a slight verbal change,\footnote{'The three creeds that \emph{go under the name of the} Nicene Creed, Athanasius' Creed,' etc., instead of 'The three Creeds, Nicene Creed, Athanasian Creed,' etc. Ussher and Vossius had proved the post-Athanasian origin of the creed which bears his name. Lightfoot (\emph{Journal}, p. 10) notices, probably from an earlier stage of the debate, another change, viz.: 'for that the \emph{matter} of them [for \emph{they}] may be proved by most certain warrants of Holy Scripture.' He adds that 'at last it was concluded that the Creeds should be printed at the end of the Thirty-nine Articles.' Comp. Mitchell, in \emph{Minutes}, p. 542.} and omitted in the copy which Parliament sent to the King at the Isle of Wight. The Assembly certainly had no objection to the doctrine of the œcumenical creeds, and teaches it in its own standards. And yet the omission of all allusion to them in the Confession of Faith is so far characteristic as it reveals a difference of stand-point. The Puritan Assembly was unwilling to adopt any rule of faith except the Scripture explained by itself; while the Episcopal Church was reformed on the basis of the Scripture as interpreted by the ancient Church, or at all events with respectful reference to primitive creeds and canons.

The work of revision was suspended by an order of Parliament, Oct. 12, 1643, requiring the Assembly to enter upon the work of Church government, and then given up in consequence of an order 'to frame a Confession of Faith for the three kingdoms, according to the Solemn League and Covenant.' The framing of the Westminster Confession is therefore due to Scotch influence and the adoption of the Solemn League and Covenant.\footnote{See this important document and its history above, pp. 689 sqq. Marsden says (\emph{Later Puritans}, p. 90): 'The taking of the Covenant in Scotland was perhaps the most solemn scene in the history of nations. The forced imposition of it in England was an insult and a burlesque.' Fuller refutes it at length from his English and Episcopal stand-point (\emph{Church Hist.} Vol. VI. pp. 259 sqq.). It certainly turned out to be a blunder in England, but it was a sublime blunder for a noble end, and not without important results, among which is the one mentioned in the text.}

This was a wise conclusion. The alteration or reconstruction of an established creed (except in minor particulars) is in itself a difficult and ungrateful task, and more apt to produce confusion than harmony, as is shown by the history of the Nicene Creed and the Augsburg Confession.

PREPARATION OF THE CONFESSION.

The first appointment of a Committee to prepare matter for a joint Confession of Faith was made Aug. 20, 1644, and embraced, besides the Commissioners of the Church of Scotland, the following Englishmen: Dr. Gouge, Mr. Gataker, Mr. Arrowsmith, Dr. Temple, Mr. Burroughs, Mr. Burges, Mr. Vines, Dr. Goodwin, and Dr. Hoyle. The chairman, Dr. William Gouge, a graduate of Cambridge, was Minister of Blackfriars, London (from 1608), and stood in high veneration among the Puritans, there being 'scarce a lord or lady or citizen of quality in or about the city that were piously inclined but they sought his acquaintance.'\footnote{Masson, Vol. II. p. 518. Gouge's Commentary on Hebrews was republished, 1866, at Edinburgh, in 3 vols., with a memoir, in which he is called 'the father of the London ministers and the oracle of his time' (p. xii.).} He died Dec. 12, 1653, seventy-nine years of age. The Committee was enlarged Sept. 4, 1644, by adding Messrs. Palmer, Newcomen, Herle, Reynolds, Wilson, Tuckney, Smith, Young, Ley, and Sedgwicke.\footnote{See excerpts from Vol. II. of the MS. Minutes, in Mitchell's ed. of \emph{Minutes} (which begin Nov. 18, 1644), p. lxxxvi.}

This Committee, it seems, prepared the material and reported in the 434th session, May 12, 1645, when a smaller Committee was appointed to digest the material into a formal draught. The members were taken from the old Committee, with Dr. Gouge as chairman. The Scotch Commissioners were to be again consulted.\footnote{\emph{Minutes}, p. 91.} On July 7th, 1645, Dr. Temple made a report of a part of the Confession touching the Holy Scripture, which was read and debated.\footnote{\emph{Ibid.} p. 110.} The following day, Reynolds, Herle, and Newcomen, to whom were afterwards added Tuckney and Whitaker, were appointed a Committee 'to take care of the wording of the Confession, as it is voted in the Assembly from time to time, and report to the Assembly when they think fit there should be any alteration in the words,' after first consulting 'with the Scotch Commissioners or any one of them.'\footnote{\emph{Minutes}, p. 110.} In the 470th session, July 16, 1645, the heads of the Confession were distributed among three large committees to be elaborated and prepared for more formal discussion.\footnote{\emph{Ibid.} p. 114.} The chapters were reported, read, and debated, section by section, and sometimes word by word.

The sub-committees sat two days every week, and reported as they progressed. On Sept. 25, 1646, the title was fixed ('The Humble Advice,' etc.) and the first nineteen chapters were sent up to the House of Commons at their request. A few days afterwards (Oct. 1) a duplicate was sent to the House of Lords.\footnote{\emph{Ibid.} pp. 290, 291; \emph{Journals of the H. of Commons}, Vol. IV. p. 677; and \emph{the H. of Lords}, Vol. VIII. pp. 505, 588.} The House of Lords passed these chapters, after a third reading, unanimously (Nov. 6). The House of Commons delayed definite action till the whole was presented. In the 752d Session, Dec. 4, 1646, the Confession was completed and presented to both Houses of Parliament in a copy transcribed with great pains by Dr. Burgess, for which he received a vote of thanks from the Assembly.\footnote{\emph{Minutes}, p. 308; \emph{Journals of the H. of Commons}, Vol. IV. p. 739; \emph{of the Lords}, Vol. VIII. p. 597.}

The Confession was thus prepared in two years and three months, amid many interruptions by discussions on the Catechism and on discipline. No other symbolical book cost so much time and labor, except the Tridentine and Vatican Decrees, and perhaps the Lutheran Formula of Concord. Besides the chairman, Drs. Tuckney, Arrowsmith, Reynolds (afterwards bishop), Temple, Hoyle, Palmer, Herle, and the Scotch divines seem to have been the chief authors of the work.

The Confession was first printed Dec., 1646, or Jan., 1647, for the exclusive use of Parliament and the Assembly, without the Scripture proofs. The House of Commons, not satisfied, expressly requested the Assembly to send them the Scripture texts (April 22, 1647), which was promptly done (April 29).\footnote{\emph{Journals of the House of Commons}, Vol. V. p. 151; \emph{Minutes}, p. 352. Baillie (in a letter to Spang, Jan. 20, 1647, Vol. III. p. 2) ascribes this request of Parliament to the 'retarding party,' and as a change of tactics of the opponents, and remarks that the Assembly omitted the Scripture proofs at first 'only to eschew the offense of the House, whose practice hitherto has been to enact nothing of religion on divine right or Scriptural ground, but upon their own authority alone.'} Whereupon the House of Commons ordered 'that six hundred copies, and no more, of the Advice of the Assembly of Divines concerning the Confession of Faith, with the quotations and texts of Scripture annexed, presented to this House, and likewise six hundred copies of the Proceedings of the Assembly of Divines upon the Nine-and-thirty Articles of the Church of England, be forthwith printed for the service of both Houses and of the Assembly of Divines; and the printer is enjoined at his peril not to print more than six hundred copies of each, or to divulge or publish any of them.'\footnote{\emph{Journals}, Vol. V. p. 156, and \emph{Minutes}, p. 354.} At the same time a vote of thanks to the Assembly was passed 'for their great pains in these services.' This second edition. appeared May, 1647, and contains the received and ecclesiastically authorized text. It must not be confounded with the revised text of Parliament.

THE ACTION OF THE PARLIAMENT.

The House of Commons began, May 19, 1647, the consideration of the 'Humble Advice,' chapter by chapter, resumed it in October, and completed it March 22, 1648. It made some alterations in the governmental chapters, and gave the document the title, 'Articles of Christian Religion approved and passed by both Houses of Parliament, after Advice had with the Assembly of Divines by authority of Parliament sitting at Westminster.'\footnote{The original title, 'A Confession of Faith,' was voted down by sixty-one to forty-one.—\emph{Minutes}, p. 415.}

The House of Lords agreed to all the alterations, excepting to that on marriage, June 3, 1648. Whereupon the House of Commons, on the 20th of June, ordered 'that the Articles of Christian Religion sent from the Lords with some alterations, the which were this day read, and upon the question agreed unto, be forthwith printed and published.' The next day it was resolved 'that the texts of Scripture be printed with the Articles of Faith.'

A copy of the authorized edition of these Articles is preserved in the British Museum. It differs from the Assembly's Confession by the omission of the entire Ch. XXX. (on Church Censures) and Ch. XXXI. (on Synods and Councils), and parts of Ch. XX. (§ 4) and Ch. XXIV. (§§ 5, 6, and part of 4).

When, after Cromwell's death, the Long Parliament was restored in 1659, it adopted the Confession with the exception of Ch. XXX. and Ch. XXXI., and requested Dr. Reynolds, Mr. Calamy, and Mr. Manton to superintend the publication (March 5, 1660).\footnote{\emph{Journals of the House of Commons}, Vol. VII. p. 862; Mitchell, in \emph{Minutes}, p. 417. Mitchell gives no information of copies of this edition.}

The English Parliament thus twice indorsed the Westminster Confession as to its doctrinal articles, but retained an Erastian control over matters of discipline. With the restoration of the monarchy the Confession shared the fate of Presbyterianism in England.

THE ACTION OF THE GENERAL ASSEMBLY OF SCOTLAND.

The Confession was at once brought to Scotland, and most favorably received.\footnote{Baillie brought a copy of the first edition, without proofs, in January (\emph{Letters}, Vol. III. p. 2); Gillespie probably a copy of the second ed., with proofs, in July, when he returned. The Assembly ordered an edition of 300 copies to be printed at Edinburgh, for the use of the members.—\emph{Minutes}, p. 419.} The General Assembly at Edinburgh, Aug. 27, 1647, after careful examination, adopted it in full as it came from the hands of the Westminster divines, declaring it 'to be most agreeable to the Word of God, and in nothing contrary to the received doctrine, worship, discipline, and government of this Kirk,' and thankfully acknowledging the great mercy of the Lord, 'in that so excellent a Confession of Faith is prepared, and thus far agreed upon in both kingdoms.' The Scotch Parliament indorsed this action, Feb. 7, 1649.

Thus the Confession, as well as the two Catechisms, received the full sanction of the highest ecclesiastical and civil authorities of Scotland. But the royal sanction was not obtained till 1690, under William and Mary.\footnote{See the Acts of the Scotch Assembly and Parliament, and of the English Parliament, in \emph{Minutes}, pp. 419 sqq.; in the Edinb. ed. of the Conf., 1855; and in Innes, \emph{The Law of Creeds}, pp. 95 sqq.}

It is a very remarkable fact that this Confession failed in its native land, and succeeded in foreign lands. The product of English Puritans became the highest standard of doctrine for Scotch and American Presbyterians, and supplanted the older Confession of their own Reformers. The Shorter Catechism, however, was for a long time extensively used in England.

Another remarkable fact is that the English authors, with their sad experience of the laws of uniformity, never intended to make their Confession binding upon the conscience as a document for subscription, while the Scots adopted it at once.\footnote{Dr. Tuckney, one of the chief authors of the Confession and Catechisms, says: 'For the matter of imposing upon I am not guilty. In the Assembly I gave my vote with others that the Confession of Faith put out by authority should not be required to be either sworn or subscribed to—our having been burnt in the hand in that kind before; but [only] so as not to be publicly preached or written against' (quoted by M'Crie, \emph{Annals}, p. 221). Baxter, also, while highly recommending the Westminster Standards, expressed the hope that 'the Assembly intended not all that long Confession and those Catechisms to be imposed as a test of Christian communion, nor to disown all that scrupled every word in it [them]. If they did, I could not have commended it for any such use, though it be useful for the instruction of families' (Sylvester's \emph{Life of Baxter}, p. 122, quoted by M'Crie, p. 222).} Dr. M'Crie accounts for this difference partly 'by national idiosyncrasies, partly by the extreme desire of the Scots to obtain that ``covenanted uniformity'' for which England was not prepared, but which Scotland, with a Church fully organized and a Parliament favorably disposed, regarded as the sheet-anchor of her safety, and to which afterwards, as a sacred engagement, she resolutely clung, in hope and against hope, in days of darkness and storms. In England Presbytery had yet to be organized, and at every step it encountered conflicting and neutralizing influences.'

\xsubsection{Analysis of the Westminster Confession.}

§ 95. Analysis of the Westminster Confession.

SOURCES.

The Westminster Confession sets forth the Calvinistic system in its scholastic maturity after it had passed through the sharp conflict with Arminianism in Holland, and as it had shaped itself in the minds of Scotch Presbyterians and English Puritans during their conflict with High-Church prelacy. The leading ideas, with the exception of the theory of the Christian Sabbath, were of Continental growth, but the form was entirely English.

The framers of the Confession were no doubt quite familiar with Continental theology; Latin was then still the theological language; the Arminian controversy had excited the greatest attention in England, and agitated the pulpit and the press for years; the English Church was well represented at the Synod of Dort; several divines of the Assembly had spent some time in Holland, where they found a hospitable refuge from persecution under Charles I., and were treated with great respect by the Dutch ministers and divines.\footnote{Dr. M'Crie (\emph{Annals}, p. 177) asserts without proof that the 'Westm. Conf. bears unmistakably the stamp of the Dutch theology in the sharp distinctions, logical forms, and judicial terms into which the reformed doctrine had gradually moulded itself under the red heat of the Arminian and Socinian controversies.' This is an error if we look to the direct source. See below.}

But while the Confession had the benefit of the Continental theology, and embodied the results of the Arminian controversy, it was not framed on the model of any Continental Confession, nor of the earlier Scottish Confessions, notwithstanding the presence and influence of the Commissioners from the Church of Scotland. On the contrary, it kept in the track of the English Articles of Religion, which the Assembly was at first directed to revise, and with which it was essentially agreed. It wished to carry on that line of development which was begun, several years before the Arminian controversy, by the framers of the Lambeth Articles (1595), and which was continued by Archbishop Ussher in the Irish Articles (1615).\footnote{See pp. 658 and 662.} It is a Calvinistic completion and sharper logical statement of the doctrinal system of the Thirty-nine Articles, which stopped with the less definite Augustinian scheme, and left a considerable margin for different interpretations. In point of theological ability and fullness it is far superior to its predecessors.

The Westminster Confession agrees more particularly with the Articles which were adopted by the Protestant Church in Ireland, but afterwards set aside by Archbishop Laud through the Earl of Strafford. This is manifest in the order and arrangement, in the titles of chapters, in phraseology, and especially in the most characteristic features of Calvin's theology—the doctrine of Predestination and of the Sacraments. The resemblance is so striking that it must have been intended for the purpose of showing the essential agreement of the Assembly with the doctrinal standards of the English and Irish Reformation. Ussher himself had pursued the same course and incorporated in his work the substance of the English Articles and the full text of the Lambeth Articles. He was a doctrinal Puritan, and although he declined the invitation to a seat in the Assembly, he was highly esteemed by the members for his learning, orthodoxy, and piety. His friend, Dr. Hoyle, Professor of Divinity at Dublin, belonged to the committee which framed the Confession.\footnote{This agreement was first brought to light and set forth in detail by Prof. Mitchell, of St. Andrews, in the pamphlet above quoted, and also in the Introduction to the \emph{Minutes}, p. xlvii.}

The following tables will illustrate the relation of the Westminster Confession to the preceding standards of the English and Irish Church.

WESTMINSTER CONFESSION. 1647.

IRISH ARTICLES. 1615.

Chapter I.—Of Holy Scripture.

Of Holy Scripture.

VII. All things in Scripture are not alike plain in themselves, nor alike clear unto all; yet \emph{those things which are necessary to be known}, believed, and observed \emph{for salvation} are so clearly propounded and opened in some place of Scripture or other, that not only \emph{the learned, but the unlearned}, in a due use of the ordinary means, may attain unto a sufficient understanding of them.

5. Although there be some hard things in the Scripture, . . . yet \emph{all things necessary to be known unto everlasting salvation} are clearly delivered therein; and nothing of that kind is spoken under dark mysteries in one place which is not in other places spoken more familiarly and plainly, to the capacity \emph{both of learned and unlearned.}

Chapter II.—Of God and of the Holy Trinity.

Of Faith in the Holy Trinity.

I. There is \emph{but one only living and true God}, who is infinite in being and perfection, a most pure spirit, invisible, \emph{without body, parts, or passions}, etc.

8. There is \emph{but one living and true God, everlasting, without body, parts, or passions}, of infinite power, wisdom, and goodness, etc.

III. \emph{In the unity of the Godhead there be three persons, of one substance, power, and eternity}—God the Father, God the Son, and God the Holy Ghost.

And \emph{in unity of this Godhead, there be three persons of one} and the same \emph{substance, power, and eternity}—the Father the Son, and the Holy Ghost. [English Art. I.]

Chapter III.—Of God's Eternal Decree.

Of God's Eternal Decree and Predestination.

I. \emph{God from all eternity did}, by the most wise and holy counsel of his own will, freely and unchangeably \emph{ordain whatsoever comes to pass; yet so as thereby} neither is God the author of sin, \emph{nor is violence offered to the will of the creatures}, nor is \emph{the liberty or contingency of second causes taken away, but rather established.}

11. \emph{God, from all eternity, did}, by his unchangeable counsel, \emph{ordain whatsoever} in time should \emph{come to pass: yet so as thereby no violence is offered to the wills of the reasonable creatures}, and neither \emph{the liberty nor the contingency of the second causes is taken away, but established rather.}

III. By the decree of \emph{God}, for the manifestation of his glory, \emph{some men and angels are predestinated unto} everlasting \emph{life}, and others \emph{foreordained} to everlasting \emph{death.}

12. By the same eternal counsel \emph{God hath predestinated some unto life}, and \emph{reprobated} some \emph{unto death:} of both which there is a \emph{certain number} known only to God, which \emph{can neither be increased nor diminished.} Lambeth Art. I. and III.]

IV. These angels and men, thus predestinated and foreordained, are particularly and unchangeably designed; and their \emph{number} is so \emph{certain} and definite that it \emph{can not be either increased} or \emph{diminished.}

V. Those of mankind that are \emph{predestinated unto life}, God, \emph{before the foundation of the world was laid}, according to his eternal and immutable \emph{purpose}, and the \emph{secret counsel} and \emph{good pleasure} of his will, \emph{hath chosen in Christ unto everlasting} glory, out of his mere free grace and love, \emph{without any foresight of faith or good works}, or \emph{perseverance} in either of them, \emph{or any} other \emph{thing} in the creature, as

13. \emph{Predestination to life} is the everlasting \emph{purpose} of God, whereby, \emph{before the foundations of the world were laid}, he hath constantly decreed in his \emph{secret counsel} to deliver from curse and damnation those whom he \emph{hath chosen in Christ} out of mankind, and to bring them by Christ \emph{unto everlasting} salvation, as vessels made to honor.

14. The \emph{cause moving God} to \emph{predestinate}

conditions, or causes \emph{moving him} thereunto; and all to the praise of his glorious grace.

\emph{unto life} is \emph{not the foreseeing of faith, or perseverance, or good works, or of any thing} which is in the person predestinated, but only the \emph{good pleasure} of God himself. For all things being ordained for the manifestation of his glory, and his glory being to appear both in the works of his mercy and of his justice, it seemed good to his heavenly wisdom to choose out a certain number towards whom he would extend his undeserved mercy, \emph{leaving the rest} to be spectacles of \emph{his justice.}

VI. As God hath appointed the elect unto glory, so hath he, by the eternal and most free purpose of his will, foreordained all the means thereunto. Wherefore they who are elected, being fallen in Adam, are redeemed by Christ; are effectually \emph{called} to faith in Christ by \emph{his Spirit working in due season;} are \emph{justified, adopted}, sanctified, and kept by his power through faith unto salvation. Neither are any other redeemed by Christ, effectually called, justified, adopted, sanctified, and saved, but the elect only.

15. Such as are predestinated unto life, be \emph{called} according unto God's purpose (\emph{his Spirit working in due season}) and through grace they obey the calling, they be \emph{justified} freely, they be made sons of God \emph{by adoption}, they be made like the image of his only-begotten Son Jesus Christ, they walk religiously in good works, and at length, by God's mercy they attain to everlasting felicity. But such as are not predestinated to salvation shall finally be condemned for their sins. [English Art. XVII.; Lambeth Art. II.]

VII. The \emph{rest} of mankind God was pleased, according to the unsearchable counsel of his own will, whereby he extendeth or withholdeth mercy as he pleaseth, for the glory of his sovereign power over his creatures, \emph{to pass by}, and to ordain them to dishonor and wrath for their sin, to the praise of \emph{his} glorious \emph{justice.} [Comp. Irish Art. § 14: '\emph{leaving the rest} to be spectacles \emph{of his justice.}']

32. None can come unto Christ unless it be given unto him, and unless the Father draw him. And all men are not so drawn by the Father that they may come unto the Son. Neither is there such a sufficient measure of grace vouchsafed unto every man whereby he is enabled to come unto everlasting life. [Lambeth Art. VII., VIII., IX.]

VIII. The doctrine of this high mystery of predestination is to be handled with special prudence and care, that men attending to the will of God revealed in his Word, and yielding obedience thereunto, may, from the certainty of their effectual vocation, be assured of their eternal election.

17. We must receive God's promises in such wise as they be generally set forth unto us in Holy Scripture; and in our doings, that will of God is to be followed which we have expressly declared unto us in the Word of God. [English Art. XVII.]

So shall this doctrine afford matter of praise, reverence, and admiration of God, and of humility, diligence, and abundant consolation, to all that sincerely obey the gospel.

16. The godlike consideration of predestination and our election in Christ is full of sweet, pleasant, and unspeakable comfort to godly persons, etc. [English Art. XVII.]

Chapter V.—Of Providence.

Of the Fall of Man, etc.

IV. [His providence] extendeth itself even to the first fall, and all other sins of angels and men, and that \emph{not by a bare permission}, but such as has joined with it a most wise and powerful bounding, and otherwise \emph{ordering and governing} of them in a manifold dispensation to his own holy ends: yet so as the sinfulness thereof proceedeth only from the creature and \emph{not from God, who,} being most holy and righteous, \emph{neither} is \emph{nor can be the author} or approver \emph{of sin.}

28. God is \emph{not the author of sin;} howbeit he \emph{doth not only permit}, but also by his providence \emph{govern and order} the same, guiding it in such sort by his infinite wisdom as it turneth to the manifestation of his own glory, and to the good of his elect.

Chapter VI.—Of the Fall of Man, of Sin, etc.

Of Original Sin.

V. \emph{This corruption of nature}, during this life, \emph{doth remain in those that are regenerated:}

24. \emph{This corruption of nature doth remain} even \emph{in those that are regenerated;} . . . And

and although it be through Christ pardoned and mortified, yet both itself and all the motions thereof are truly and properly sin.

howsoever for Christ's sake there be no condemnation to such as are regenerate and do believe, yet doth the apostle acknowledge that in itself this concupiscence hath the nature of sin. [English Art. IX.]

Chapter VIII.—Of Christ the Mediator.

Of Christ, the Mediator of the Second Covenant.

II. The Son of God, the second person in the Trinity, being \emph{very and eternal God, of one substance} and equal \emph{with the Father}, did, when the fullness of time was come, \emph{take} upon him \emph{man's nature}, with all the essential properties and common infirmities thereof, yet without sin: being conceived by the power of the Holy Ghost, \emph{in the womb of the Virgin Mary, of her substance. So that two whole, perfect}, and distinct \emph{natures, the Godhead and} the \emph{manhood, were inseparably joined} together in \emph{one person}, without conversion, composition, or confusion. Which person is \emph{very God}, and \emph{very man}, yet \emph{one Christ;} the only Mediator between God and man.

29. The \emph{Son}, which is the Word of the Father, begotten from everlasting of the Father, \emph{the true and eternal God, of one substance with the Father, took man's nature in the womb of the blessed Virgin, of her substance: so that two whole} and \emph{perfect natures}, that is to say, \emph{the Godhead and manhood, were inseparably joined in one person}, making \emph{one Christ very God and very Man.} [English Art. II.]

Chapter XVI.—Of Good Works.

Of Sanctification and Good Works.

I. Good works are only such as God \emph{hath commanded in his holy Word}, and not such as, \emph{without the warrant thereof, are devised by men, out of blind zeal}, or upon any pretense of good intention.

42. The works which God would have his people to walk in are \emph{such as he hath commanded in his Holy Scripture}, and not such works \emph{as men have devised} out of their own brain, \emph{of a blind zeal} and devotion, \emph{without the warrant} of the Word of God.

Chapter XVII.—Of the Perseverance of the Saints.

Of Justification and Faith.

I. They whom God hath accepted in his Beloved, effectually called, and sanctified by his Spirit, can \emph{neither totally nor finally fall away} from the state of grace; but shall certainly persevere therein to the end, and be eternally saved.

38. A true, lively, justifying faith, and the sanctifying Spirit of God, is not extinguished, nor vanisheth \emph{away}, in the regenerate, \emph{either finally or totally.} [Lambeth Art. V.]

Chapter XXI.—Of Religious Worship and the Sabbath Day.

Of the Service of God.

II. \emph{Religious worship is to be given to God} the Father, Son, and Holy Ghost, and to him \emph{alone.}

54. \emph{All religious worship ought to be given to God alone.}

VIII. This Sabbath is then kept holy unto the Lord, when men . . . do not only observe an holy rest all the day from their own works, words, and thoughts about their worldly employments and recreations, but also are taken up the \emph{whole time} in the \emph{public and private exercises} of his worship, and in the duties of necessity and mercy.

56. The first day of the week, which is the Lord's day, is \emph{wholly} to be dedicated unto the service of God; and therefore we are bound therein to rest from our common and daily business, and to bestow that leisure upon holy \emph{exercises, both public and private.}

Chapter XXIII.—Of the Civil Magistrate.

Of the Civil Magistrate.

III. The Civil Magistrate may not assume to himself \emph{the administration of the Word and}

58. . . . Neither do we give unto him hereby the \emph{administration of the Word and sacraments},

\emph{sacraments, or the power of the keys of} the kingdom of heaven.

\emph{or the power of the keys}, etc. [See English Art. XXXVII.]

Chapter XXV.—Of the Church.

Of the Church, etc.

I. The \emph{Catholic} or \emph{Universal Church}, which is invisible, consists of the whole number of the elect \emph{that have been, are, or shall be gathered} into one, \emph{under Christ, the head thereof;} and is the spouse, the body, the fullness of him who filleth all in all.

68. There is but one \emph{Catholic Church}, out of which there is no salvation: containing the \emph{universal} company of all the saints \emph{that ever were, are, or shall be gathered} together \emph{in one} body, under one \emph{head, Christ} Jesus.

Chapter XXVIII.—Of Baptism.

Of Baptism.

I. Baptism is a sacrament of the New Testament, ordained by Jesus Christ, \emph{not only} for the solemn \emph{admission} of the party baptized \emph{into the} visible \emph{Church;} but also to be unto him a \emph{sign} and \emph{seal} of the covenant of grace, of his ingrafting into Christ, of \emph{regeneration}, of remission of sins, and of his giving up unto God, through Jesus Christ to walk in newness of life.

89. Baptism is \emph{not only} an outward \emph{sign} of our profession, . . . but much more a sacrament of our \emph{admission into the Church}, sealing unto us our \emph{new birth} (and consequently our justification, adoption, and sanctification) by the communion which we have with Jesus Christ. [English Art. XXVII.]

Chapter XXIX.—Of the Lord's Supper

Of the Lord's Supper

I. The sacrament of his body and blood . . . for the perpetual remembrance of the sacrifice of himself in his death, the \emph{sealing} all the benefits thereof unto true believers, their \emph{spiritual nourishment and growth in him.}

92. The Lord's Supper is not only a sign, but much more a sacrament of our preservation in the Church, \emph{sealing} unto us our \emph{spiritual nourishment} and continual \emph{growth in Christ.} [English Art. XXVIII.]

VII. Worthy receivers, outwardly partaking of the visible elements in this sacrament, do then also inwardly by \emph{faith, really} and indeed, yet not carnally and corporally, but spiritually, receive and feed upon Christ crucified, and all benefits of his death: \emph{the body and blood of Christ} being then not corporally or carnally in, with, or under the bread and wine, yet as \emph{really}, but spiritually, present to the faith of believers in that ordinance, as the elements themselves are to the outward senses.

94. But in the inward and spiritual part the same body and blood is \emph{really} and substantially presented unto all those who have grace to receive the Son of God, even to all those that \emph{ believe} in his name. And unto such as in this manner do worthily and with \emph{faith} repair unto the Lord's table, \emph{the body and blood of Christ} is not only signified and offered, but also truly exhibited and communicated.

VIII. Although ignorant and wicked men receive the outward elements in this sacrament, yet they receive not the thing signified thereby; but by their unworthy coming thereunto are guilty of the body and blood of the Lord, to their own damnation. Wherefore, all ignorant and ungodly persons, as they are unfit to enjoy communion with him, so are they unworthy of the Lord's table, and can not, without great sin against Christ, while they remain such, partake of these holy mysteries, or be admitted thereto.

96. The wicked, and such as want a lively faith, although they do carnally and visibly (as St. Augustine speaketh) press with their teeth the sacrament of the body and blood of Christ, yet in no wise are they made partakers' of Christ; but rather to their condemnation do eat and drink the sign or sacrament of so great a thing. [English Art. XXIX.]

CONTENTS.

Neal says: 'Though all the divines were in the anti-Arminian scheme, yet some had a greater latitude than others. I find in my MS. the dissent of several members against some expressions relating to reprobation, to the imputation of the active as well as passive obedience of Christ, and to several passages in the chapter on liberty of conscience and Church discipline; but the Confession, as far as related to articles of faith, passed the Assembly and Parliament by a very great majority.'\footnote{Vol. II. p. 41.} Neal does not specify the differences to which he alludes. Since the publication of the \emph{Minutes} we are enabled to ascertain them, at least to some extent, from the meagre and broken reports of debates on election and reprobation, on the fall of Adam, on the Covenants, on providence, free-will, creation, justification, sanctification, the sacraments, and other topics. In most cases the fact is simply mentioned that there was a debate; in others brief extracts of speeches are given which reveal minor differences of views, though not of parties, or even of schools. The debates on Church government were much more serious and heated. The harmony of so many scholars from all parts of England and Scotland, on a whole scheme of divinity, is truly surprising, and accounts for their sanguine hopes of securing a doctrinal uniformity in the three kingdoms.

The Confession consists of thirty-three chapters, which cover, in natural order, all the leading articles of the Christian faith from the creation to the final judgment. It exhibits the consensus of the Reformed Churches on the Continent and in England and Scotland, which was one of the objects of Parliament intrusted to the Assembly.

BIBLIOGRAPHY.

Following the precedent of most of the Continental Reformed Confessions and the Irish Articles, the Westminster formulary properly begins with the Bible, on which all our theology must be based, and sets forth its divine inspiration, authority, and sufficiency as an infallible rule of faith and practice, in opposition both to Romanism, which elevates ecclesiastical tradition to the dignity of a joint rule of faith, and to Rationalism, which teaches the sufficiency of natural reason. It excludes the Jewish Apocrypha entirely from the Canon, while in the English and Irish Articles they are at least enumerated, though distinguished from the canonical books.\footnote{The Lutheran symbols make no such distinction and give no list of the canonical books. They have no separate article on the Scriptures at all, beyond the important statement in the introduction to the Formula of Concord.} The Confession gives to reason, or the light of nature, its proper place, distinguishes between the original Scripture and the translations, maintains the true exegetical principle of the self-interpretation of Scripture in the light of the Spirit that inspired it, and carefully avoids committing itself to any mechanical or magical or any other particular theory concerning the mode and degrees of inspiration, or obstructing the investigation of critical questions concerning the text and the authorship (as distinct from the canonicity) of the several books.\footnote{Thus we find that the Epistle to the Hebrews is named separately, and not included in '\emph{fourteen} Epistles of Paul,' as in the Belgic Confession. Canonicity is not necessarily dependent on a traditional view of authorship or genuineness.} It rests the authority of the Bible on its own intrinsic excellence and the internal testimony of the Spirit rather than the external testimony of the Church, however valuable this is as a continuous witness.\footnote{Ch. I. 5: 'We may be moved and induced by the testimony of the Church to an high and reverent esteem of the Holy Scripture, and the heavenliness of the matter, the efficacy of the doctrine, the majesty of the style, the consent of all the parts, the scope of the whole (which is to give all glory to God), the full discovery it makes of the only way of man's salvation, the many other incomparable excellencies, and the entire perfection thereof, are arguments whereby it doth abundantly evidence itself to be the Word of God; yet, notwithstanding, our full persuasion and assurance of the infallible truth and divine authority thereof is from the inward work of the Holy Spirit, bearing witness by and with the Word in our hearts.'}

No other Protestant symbol has such a clear, judicious, concise, and exhaustive statement of this fundamental article of Protestantism. It has been pronounced equal in ability to the Tridentine decree on justification.\footnote{While arguing \emph{against} creeds and councils, Dean Stanley (in the \emph{Contemp. Rev.} for Aug. 1874, p. 499) writes: 'Is there any single theological question which any council or synod has argued and decided with an ability equal to that of any of the great theologians, lay or clerical? The \emph{nearest approaches} to it are the chapters on Justification in the Decrees of Trent, and on the \emph{Bible} in the \emph{Westminster Confession.}' Comp. also the remarks of Dr. Mitchell, Introd. to \emph{Minutes}, p. xlix.} It may more aptly be compared to the Tridentine decree on Scripture and tradition (Sess. IV.) and the recent Vatican decree on the dogmatic constitution of the Catholic faith (Sess. III.), as far as this relates to reason and revelation, and may be regarded as the best Protestant counterpart of the Roman Catholic doctrine of the rule of faith. The Confession plants itself exclusively on the Bible platform, without in the least depreciating the invaluable aid of human learning—patristic, scholastic, and modern  —in its own proper place, as a means to an end and an aid in ascertaining the true sense of the mind of the Holy Spirit, who through his own inspired Word must alternately decide all questions of the Christian faith and duty. It is clear that Protestantism must sink or swim with this principle. Criticism, philosophy, and science may sweep away human traditions, confessions, creeds, and other outworks, but they can never destroy the fortress of God's Word, which liveth and abideth forever.

THEOLOGY AND CHRISTOLOGY.

Ch. II., 'Of the Trinity,' and Ch. XVIII., 'Of Christ the Mediator,' contain one of the best statements of the Nicene doctrine of the Trinity and of the Chalcedonian Christology, as held by all orthodox Churches. On these articles the evangelical Protestant Confessions are entirely agreed.

PREDESTINATION.

Ch. III., 'Of God's Eternal Decree,'\footnote{The English and Scotch editions use the singular, some American editions the plural (as in the Catechisms). There was a dispute in the Assembly about \emph{decree} and \emph{decrees.} Several members were opposed to dividing the one, all-comprehending decree of God. Seaman said: 'All the odious doctrine of the Arminians is from their distinguishing of the decrees, but our divines say they are one and the same decree.' Reynolds differed. See \emph{Minutes}, p. 151. But both Catechisms in all editions have \emph{decrees} (comprehended under the one \emph{purpose} of God; see Shorter Catechism, Quest. 7).} Ch. V., 'Of Providence,' Ch. IX., 'Of Free Will,' and Ch. XVIII., 'Of the Perseverance of the Saints,' are closely connected. They present a logical chain of ideas which make up what is technically called 'the Calvinistic system,' as developed first by Calvin himself against Romanism, then in Holland and England against Arminianism.

This system had at that time a powerful hold upon the serious religious minds in England and Scotland, including many leading Episcopal divines (not of the Laudian type) who otherwise had no sympathy with Puritanism, and ridiculed it with bitter sarcasm, like Dr. South. Even the authorized English version of the Bible (1611) has been charged by Arminians with a Calvinistic bias, while Calvinists have never complained of any defect in this respect.\footnote{The charge derives some plausibility from the fact that the supralapsarian Beza, by his Greek Testament and his Latin translation and notes, exerted a marked influence on the translators. It is supported chiefly by three passages. In Matt. xx. 23, the words '\emph{it shall be given}' are unnecessarily inserted (after the precedent of the Geneva version). In Acts ii. 47, we read, 'The Lord added to the Church \emph{such as should be saved},' instead of 'such as were being saved, or in the way of salvation' (\textgreek{τοὺς σωζομένους,} not \textgreek{τοὺς σωθησομένους}). In Heb. x. 38—'Now the just shall live by faith; but if \emph{any man} draw back, my soul shall have no pleasure in him'—\emph{any man} is inserted, with Beza ('si \emph{quis} se subduxerit'), to distinguish the subject of \textgreek{ὑποστείληται} from the \textgreek{δίκαιος} of the first clause, and to evade an argument against the perseverance of saints. But the case here is doubtful.} The only question in the Assembly was as to the logical extent to which they should carry the doctrine of predestination in a confessional statement. The more consistent and rigorous scheme of supralapsarianism had its advocates in Westminster as well as in Dort, and was favored by Dr. Twisse, the Prolocutor, who followed Beza and Gomarus to the giddy abyss of including the fall itself in the absolute eternal decree as a necessary means for the manifestation of God's justice; but the infralapsarian (or sublapsarian) scheme of Augustine decidedly triumphed. Supralapsarianism has always remained only a private speculation.

The Westminster Confession goes, indeed, beyond the two Helvetic Confessions, the Heidelberg Catechism, the Scotch Confession, and the Thirty-nine Articles; but it goes not a whit further than the Canons of Dort (which had the approval of the delegates of King James), the Lambeth Articles, and the Irish Articles.\footnote{See the comparative table, pp. 762, 768. Ussher adhered to his views on predestination, which he had expressed in the Irish Articles. In his 'Method of the Christian Religion,' written in his youth, but revised and republished shortly before his death, he has even a stronger passage on reprobation than the Westminster Confession, viz., 'Did God, then, before he made man, determine to save some and reject others? A. Yes, surely; before they had done either good or evil, God in his eternal counsel set some apart upon whom he would in time show the riches of his mercy, and determined to withhold the same from others, upon whom he would show the severity of his justice.' See Vol. XI. of his \emph{Works;} and Mitchell, p. liv. note.} It teaches really no more on predestination than the great Catholic Augustine had taught in the fourth century, as well as two archbishops of Canterbury—Anselm in the eleventh, and Bradwardine in the fourteenth century.\footnote{Bradwardine's treatise, \emph{De causa Dei adversus Pelagium}, which leads even to supralapsarianism, was republished in London in 1618 by Archbishop Abbot, the Calvinistic predecessor of the anti Calvinistic Laud.} It gives, however, a clearer logical shape and greater prominence to the doctrine in the system by placing it among the first articles. It puts the fall with its sinful consequences only under a \emph{permissive} (as distinct from a \emph{causal} or \emph{effective}) decree, and emphatically exempts God from all authorship of sin.\footnote{Ch. V. 4: 'God, being most holy and righteous, \emph{neither is nor can be} the author or approver of sin.'} It does not teach the horrible and blasphemous doctrine (so often unjustly and unscrupulously charged upon Calvinism) that God from eternity foreordained men for sin and damnation; but it does teach that out of the fallen mass of corruption God elected a definite number of men to salvation and 'passed by' the rest, leaving them to the just punishment of their sins.

This is severe and harsh enough, but very different from a decree of eternal \emph{reprobation}, which term nowhere occurs in the Confession. The difference is made more clear from the debates in the 'Minutes.' Several prominent members, as Calamy, Arrowsmith, Vines, Seaman, who took part in the preparation of the doctrinal standards, sympathized with the hypothetical universalism of the Saumur school (Cameron and Amyrauld) and with the moderate position of Davenant and the English delegates to the Synod of Dort. They expressed this sympathy on the floor of the Assembly, as well as on other occasions. They believed in a special \emph{effective} election and final perseverance of the elect (as a necessary means to a certain end), but they held at the same time that God sincerely \emph{intends} to save \emph{all} men; that Christ \emph{intended} to die, and \emph{actually} died, for \emph{all} men; and that the difference is not in the intention and offer on the part of God, but in the acceptance and appropriation on the part of men.\footnote{Calamy said, in a sermon before the House of Commons: 'It is most certain that God is not the cause of any man's damnation. He found us sinners in Adam, but made none sinners.' In the debate on redemption in the Assembly, he stated: 'I am far from universal redemption in the Arminian sense, but I hold with our divines in the Synod of Dort that Christ did pay a price for all, [with] absolute intention for the elect, [with] conditional intention for the reprobate in case they do believe; that all men should be \emph{salvabiles, non obstante lapsu Adami;} that Jesus Christ did not only die sufficiently for all, but God did \emph{intend}, in giving of Christ, and Christ in giving himself did \emph{intend}, to put \emph{all} men in a state of salvation in case they do obey.' . . . 'This universality of redemption does neither intrude upon either doctrine of special election or special grace' (\emph{Minutes}, p. 152). 'The difference is not in the offer, but in the application. For the word \emph{world} [in John iii. 16] signifies the \emph{whole} world' (p. 156). 'It can not be meant of the elect because of that \emph{whosoever believeth}, and Mark xvi., ``Preach the Gospel to \emph{every creature''' }(p. 154). 'In the point of election I am for special election, and for reprobation I am for \emph{massa corrupta;} . . . there is \emph{ea administratio} of grace to the reprobate that they do \emph{willfully damn themselves}' (p. 153). Seaman said: 'All in the first Adam were made liable to damnation, so all are liable to salvation in the second Adam. Every man was \emph{damnnabilis}, so is every man \emph{salvabilis}' (p. 154). Dr. Mitchell (pp. lvi. sqq.) shows that Arrowsmith, Gataker, and other members of the Assembly, in their private writings, agreed with Calamy. His interpretation of \textgreek{κόσμος,} in John iii. 16, is indeed the only tenable one, and seems to be favored by the exegetical tact of Calvin himself (\emph{in loc.}), for Calvin the exegete is more fair and free than Calvin the theologian. Dr. Arrowsmith, who was a member of the Committees on the Confession and on the Catechisms, in his explanation of Rom. ix. 22, 23, justly presses the important difference between the passive \textgreek{κατηρτισμένα} and the active \textgreek{προητοίμασεν} 'I desire,' he says, 'to have it punctually observed that the vessels of wrath are only said to be fitted to destruction, without naming by whom—God, Satan, or themselves; whereas, on the other side, God himself is expressly said to have prepared his chosen vessels of mercy unto glory. Which was purposely done (as I humbly conceive) to intimate a remarkable difference between election and preterition, in that election is a proper cause not only of salvation itself, but of all the graces which have any causal tendency thereunto, and therefore God is said to prepare his elect to glory; whereas negative reprobation is no proper cause either of damnation itself or of the sin that bringeth it, but an antecedent only; wherefore the non-elect are indeed said to be fitted to that destruction which their sins in conclusion bring upon them, but not by God. I call it a remarkable difference, because where it is once rightly apprehended and truly believed, it sufficeth to stop the mouth of one of those greatest calumnies and odiums which are usually cast upon our doctrine of predestination, viz., that God made sundry of his creatures on purpose to damn them—a thing which the rhetoric of our adversaries is wont to blow up to the highest pitch of aggravation. But it is soon blown away by such as can tell them, in the words of the excellent Dr. Davenant, ``It is true that the elect are severally created to the end and intent that they may be glorified together with their head, Christ Jesus; but for the non-elect, we can not truly say that they are created to the end that they may be tormented with the devil and his angels. No man is created by God with a nature and quality fitting him to damnation. Yea, neither in the state of his innocency nor in the state of the fall and his corruption doth he receive any thing from God which is a proper and fit means of bringing him to his damnation.'''—\emph{Chain of Principles}, pp. 335, 336, etc., edition 1659 (quoted by Mitchell, p. lxi.).}

Another important and modifying feature is that the Confession, far from teaching fatalism or necessitarianism, expressly recognizes the freedom of will, and embraces in the divine decrees 'the liberty or contingency of second causes' (Ch. III., 1).\footnote{Comp. Ch. IX. 1: 'God hath endued the will of man with that natural liberty that it is neither forced, nor by any absolute necessity of nature determined, to good or evil (Matt. xvii. 12; Deut. xxx. 19).} Herein it agrees with Ussher, Bullinger, and Calvin himself, and favorably differs from the Lutheran Formula of Concord, which (following the strong expressions of Luther and Flacius) unphilosophically represents the human will before conversion to be as passive as a dead log or stone. The Confession makes no attempt to solve the apparent contradiction between divine sovereignty and human freedom, but it at least recognizes both sides of the problem, and gives a basis for the assertion that God's absolute decrees have no causal effect upon the sinful actions of men, for which they alone are responsible.

With the Calvinistic particularism the limitation of redemption\footnote{The term \emph{atonement} is not used in the Confession. The English Bible exceptionally renders Rom. v. 11, \textgreek{καταλλαγή} (\emph{reconciliation}), by \emph{atonement}, which in its old sense (=at-one-ment) means \emph{reconciliation}, but is now equivalent to \emph{expiation, satisfaction} (\textgreek{ἱλασμός}). \emph{Redemption} (\textgreek{ἀπολύτρωσις}) is a wider term. This distinction should be kept in view in the explanation of the Confession.} is closely connected. The difference is chiefly one of logical consistency. It refers to the efficiency of redemption or its actual application. All were agreed as to its absolute sufficiency or its infinite intrinsic value. All could subscribe the formula that Christ died \emph{sufficienter pro omnibus, efficaciter pro electis.} Dr. Reynolds, who seems to have defended the more rigorous view, said in the debate: 'The Synod intended no more than to declare the sufficiency of the death of Christ; it is \emph{pretium in se}, of sufficient value to all—nay, ten thousand worlds.'\footnote{\emph{Minutes}, p. 153. The ablest modern defendants of a limited atonement, Drs. Cunningham and Hodge (see his \emph{Theology}, Vol. II. pp. 544 sqq.), are as emphatic on the absolute \emph{sufficiency} as Reynolds. Their arguments are chiefly logical; but logic depends on the premises, and is a two-edged sword which may be turned against them as well. For if the atonement be limited in \emph{design}, it must be limited in the \emph{offer;} or if unlimited in offer, the offer made to the non-elect must be \emph{insincere} and \emph{hypocritical}, which is inconsistent with the truthfulness and goodness of God. Every Calvinist preaches on the assumption that the offer of salvation is truly and sincerely extended to \emph{all} his hearers, and that it is their \emph{own} fault if they are not saved.}

Nevertheless, behind the logical question is the far more important theological and practical question concerning the extent of the divine \emph{intention or purpose}, viz., whether this is to be measured by God's love and the intrinsic value of Christ's merits, or by the actual result. On this question there was a difference of opinion among the divines, as the 'Minutes' show, and this difference seems to have been left open by the framers of the Confession. On the one hand, the closing sentences of Ch. III. 6 ('neither are any other \emph{redeemed} by Christ, effectually called, justified, adopted, sanctified, \emph{and} saved, but the elect only'), and Ch. VIII. 8 ('To all those for whom Christ hath purchased \emph{redemption}, he doth certainly and effectually apply and communicate the same'), favor a limited redemption, unless the word \emph{redeemed} be understood in a narrower sense, so as to be equivalent to \emph{saved}, and to imply the subjective application or actual execution.\footnote{Compare the remarks of Mitchell, p. lvii., who considers the language of the Confession in Ch. III. compatible with the liberal view, while the other passage, strictly construed, excludes it, unless 'redemption' be there taken in the sense of Baxter, as meaning 'that special redemption proper to the elect which was accompanied with an intention of actual application of the saving benefits in time.' The difference of views came up again in the debate on the 68th question of the Larger Catechism. See \emph{Minutes}, pp. 369, 392, 393.} On the other hand, Ch. VII. 3 teaches that under the covenant of grace the Lord 'freely offereth unto \emph{sinners} life and salvation by Jesus Christ, requiring of them faith in him, that they may be saved; and promising to give unto all those that are \emph{ordained unto life} his Holy Spirit, to make them willing and able to believe.' This looks like a compromise between conditional universalism taught in the first clause, and particular election taught in the second. This is in substance the theory of the school of Saumur, which was first broached by a Scotch divine, Cameron (d.1625), and more fully developed by his pupil Amyrault, between A.D. 1630 and 1650, and which was afterwards condemned in the Helvetic Consensus Formula (1675).\footnote{See pp. 480 sqq.}

ANTHROPOLOGY.

Chapters VI. to IX. present the usual doctrines of the Evangelical Reformed (Augustinian) anthropology, with the new feature of the \emph{Covenants.} The doctrine of covenants belongs to a different scheme of theology from that of the divine decrees. It is biblical and historical rather than scholastic and predestinarian. It views man from the start as a free responsible agent, not as a machine for the execution of absolute divine decrees.

Ch. VII. distinguishes two covenants of God with man, the covenant of \emph{works} made with Adam and his posterity on condition of perfect and personal obedience, and a covenant of \emph{grace} made in Christ with believers, offering free salvation on condition of faith in him. The covenant of grace again is administered under two dispensations, the law and the gospel. In the Old Testament it was administered by promises, sacrifices, circumcision, the paschal lamb, and other types and ordinances which forshadowed the future Saviour. Under the New Testament the covenant of grace is dispensed through the preaching of the Word and the administration of the Sacraments. There are therefore not two covenants of grace differing in substance, but one and the same under various dispensations.

The exegetical arguments for the covenant of works are derived chiefly from Gal. iii. 10, 12, 21; Rom. iii. 20; x. 5; but these passages refer to the covenant of the law of Moses, not to a covenant in the primitive state, and lead rather to a distinction between the covenant of the law (which, however, was also a covenant of promise) and the covenant of the gospel (the fulfillment of the law and promise).\footnote{Later federalists based the primitive covenant of works on Hos. vi. 7. See p. 484.}

The doctrine of covenants is usually traced to Dutch origin; but it was inaugurated after the middle of the sixteenth century by Caspar Olevianus (d. 1587), one of the authors of the Heidelberg Catechism, in a work on 'the Nature of God's Covenant of Mercy with the Elect,' on the basis of Jer. xxxviii. 31–34; Heb. viii. 8–12.\footnote{\par \emph{De substantia fœderis gratuiti}, etc. See a German version in Sudhoff's \emph{Olevianus und Ursinus} (Elberfeld, 1857), pp. 573 sqq.} Dr. Mitchell says that the Confession teaches no more on this subject than had been taught before by Rollock in Scotland and Cartwright in England. It is not probable, though not impossible, that the more fully developed theory of the covenants by John Coccejus was already known in England at the time when the Confession was framed. Coccejus likewise distinguishes the \emph{fœdus operum} or \emph{naturæ} in the state of innocence, and a \emph{fœdus gratiæ}, after the fall, but he views the latter under \emph{three} stages, the patriarchal or Abrahamic (\emph{œconomia ante legem}), the Mosaic (\emph{œconomia sub lege} and the Christian (\emph{œconomia post legem}).\footnote{\par Coccejus, or Koch, was at first Professor in Bremen (his native place), then at Franeker, 1636, and last at Leyden, 1649, where he died, 1669. His chief work, \emph{Summa doctrinæ de fœdere et testamento Dei}, appeared in 1648 (a year after the Westminster Conf.) and again in 1653. It was the first attempt of a biblical and exegetical theology in distinction from the scholastic orthodoxy which then prevailed in Holland. Coccejus was denounced by the orthodox as a Judaizing and Pelagianizing heretic. Comp. the article \emph{Coccejus and his School}, by Dr. Ebrard. in Herzog's \emph{Real-Encykl.} Vol. II. pp. 742 sqq.}

SOTERIOLOGY.

Chapters X. to XVIII. contain the best confessional statement of the evangelical doctrines of justification, adoption, sanctification, saving faith, good works, and assurance of salvation. The statement of justification by faith is as guarded and discriminating on the Protestant side of the question as the Tridentine statement of justification by faith and works is on the Roman Catholic side.

ECCLESIOLOGY.

Chapters XXV. and XXVI. In the doctrine of the Church the Protestant distinction between the invisible and visible Church is first clearly formulated, and the purest Churches under heaven are admitted to be 'subject to mixture and error.' Christ is declared to be the only head of the Church—a most important principle, for which the Church of Scotland has contended faithfully against the encroachments of the civil power through years of trial and persecution. On the subject of the independence and self-government of the Church in her own proper sphere, the Presbyterian Church of Scotland (as also the Dissenting Churches in England, and all American Churches) are immeasurably in advance of all the Protestant Churches on the Continent, and even of the Church of England, which is still dependent on the crown and the will of a Parliament composed of professors of all religions and no religion.

But while the Confession claims full freedom for the Church in the management of her own affairs, it claims no authority or superiority over the State like the hierarchical principle. It declares the Pope of Rome, who pretends to be the supreme head of the Church on earth, to be 'that Antichrist, that man of sin and son of perdition that exalteth himself in the Church against Christ and all that is called God' (2 Thess. ii. 3, 4, 8, 9).\footnote{\par This statement, which is made also in other Protestant Confessions and in the Irish Articles (No. 80; see Vol. III. p. 540), does not unchurch the Church of Rome, or declare her ordinances invalid; for Antichrist sits in the temple of God, and there is a material difference between the papacy and the Roman Catholic Church, as there is between the Jewish hierarchy and the people of Israel.}

The chapter on the Communion of Saints urges the duty of cherishing and promoting union and harmony with all Christians of whatever part of the visible Church.\footnote{Presbyterians therefore act in perfect consistency with their Confession if they take a leading part in all Bible Societies, Tract Societies, the Evangelical Alliance, and other catholic societies. They are among the most liberal of orthodox denominations in the support of these societies.}

THE SACRAMENTS.

The doctrine of the Sacraments in general, and Baptism, and the Lord's Supper in particular, in Chs. XXVII.-XXIX., is the Calvinistic theory which we have already discussed elsewhere.\footnote{See pp. 281, 376, 455, 601, 639, 641, 645.} It is the same which is taught in all the Reformed Confessions—Continental, Anglican, and Scotch. This is admitted by candid scholars. 'On the doctrine of the sacraments,' says Marsden, an English Episcopalian, 'we do not perceive a shade of difference from the teaching of the Church of England.'\footnote{\emph{History of the Later Puritans}, p. 84. He then quotes the questions of the Shorter Catechism on the Sacraments.} And Dr. Mitchell, a Scotch Presbyterian, says: 'The teaching of the Confession on the Lord's Supper is the teaching of Cranmer, Latimer, and Ridley, of Hooker, Ussher, and many others, . . . as well as of Knox, who from his long residence in England, and with English exiles on the Continent, had thoroughly caught up their warm and catholic utterances. This teaching is as far removed from the bare remembrance theory attributed to the early Swiss Reformers as from the consubstantiation of Luther and the local or supra-local presence contended for by Roman Catholics and Anglo-Catholics. It is so spiritual, yet so really satisfying, that even some High-Churchmen have owned that it would be difficult to find a better directory in the study of questions relating to this sacrament than is supplied in the Confession of Faith.'\footnote{Introduction to \emph{Minutes}, p. lxviii.}

THE CHRISTIAN SABBATH.

Ch. XXI., ' Of Religious Worship and the Sabbath Day,' must be mentioned as (next to the Irish Articles) the first symbolical indorsement of what may be called the Puritan theory of the Christian Sabbath which was not taught by the Reformers and the Continental Confessions, but which has taken deep root in England, Scotland, and the United States, and has become the basis of a far stricter observance of the Lord's day than exists in any other country. This observance is one of the most prominent national and social features of Anglo-American Christianity, and at once strikes the attention of every traveler.\footnote{The most recent manifestation of the national American sentiment was the closing of the Centennial Exhibition in Philadelphia (1876) on the Lord's day.}

The way was gradually prepared for it. Calvin's view of the authority of the fourth commandment was stricter than Luther's, Knox's view stricter than Calvin's, and the Puritan view stricter than Knox's.\footnote{There is a tradition that Knox once called on Calvin on Sunday, and found him enjoying the recreation of bowling on a green. Knox himself on one occasion had one or two friends taking supper with him on Sunday night, and no doubt considered this innocent (see Randolph's letter to Cecil, Nov. 30, 1562, quoted by Hessey, \emph{Bampton Lectures on Sunday}, Lond. 1860, p. 270). On the other hand, it is a fact that the designation of 'Sabbath' for Sunday, and the enumeration of 'the breaking of the Sabbath' among the grosser sins, originated with Knox, or at all events in Scotland at his time. The \emph{First Book of Discipline}, which was drawn up by Knox and five other ministers, abolishes Christmas, Circumcision, and Epiphany, 'because they have no assurance in God's Word,' but enjoins the observance of Sunday in these words: 'The Sabbath must be kept strictly in all towns, both forenoon and afternoon, for hearing of the Word; at afternoon upon the Sabbath, the Catechism shall be taught, the children examined, and the baptism ministered. Public prayers shall be used upon the Sabbath, as well afternoon as before, when sermons can not be had.' The third General Assembly resolved, July 4, 1562, to petition the queen for the punishing of Sabbath-breaking and all the vices which are 'commanded to be punished by the law of God, and yet not by the law of the realm.' Similar acts occur in the Assemblies of 1575, 1590, and 1596. See Gilfillan's work on the \emph{Sabbath}, and Appendix D to Mitchell's tract on the \emph{Westminster Confession}, pp. 53 sqq.} The Prayer-Book of the Church of England, by incorporating the responsive reading of the Decalogue in the regular service, kept alive in the minds of the people the perpetual obligation of the fourth commandment, and helped to create a public sentiment within the Church of England favorable to the Puritan theory, although practically great desecration prevailed during Elizabeth's reign. The 'judicious' Hooker, who was no Puritan, says: 'We are bound to account the sanctification of one day in seven a duty which God's immutable \emph{law} doth exact \emph{forever.}'\footnote{\emph{Eccles. Polity}, Bk. V. ch. 70, sec. 9. The fifth book came out in 1597, two years after Bownd's book. Ussher, Leighton, Pearson, Beveridge, Cecil, and other leading divines of the Church of England take the same ground on the perpetuity of the fourth commandment, and so far agree with the Puritan theory. But the Puritan practice in Scotland and New England often runs into Judaizing excesses.}

Towards the close of Elizabeth's reign the Sabbath question assumed the importance and dignity of a national movement, and of a practical reformation which traveled from England to Scotland and from both countries to North America. The chief impulse to this movement was given in 1595 by Dr. Nicolas Bownd (or Bound),\footnote{\par He was a graduate of Cambridge, was suspended with others in 1583 for some act of non-conformity, and died in 1607. Isaac Walton states (in his \emph{Life of Hooker}) that he was offered by Whitgift the mastership of the Temple, but this seems inconsistent with the Archbishop's hostility to his book. Bownd wrote also \emph{The Holy Exercise of Fasting} (1604); \emph{A Storehouse of Comfort for the Afflicted} (1604); and a sermon on the \emph{Unbelief of Thomas for the Comfort of all who desire to believe, which armeth us against Despair in the Hour of Death} (1608). There is a biographical sketch of Bownd in Brook's \emph{Lives of the Puritans}, Vol. II. pp. 171–176.} a learned Puritan clergyman of Norton in Suffolk. He is not the originator, but the systematizer or first clear expounder of the Puritan theory of the Christian Sabbath, namely, that the Sabbath or weekly day of holy rest is a primitive institution of the benevolent Creator for the benefit of man, and that the fourth commandment as to its substance (that is, the keeping holy one day out of seven) is as perpetual in design and as binding upon the Christians as any other of the Ten Commandments, of which Christ said that not 'one jot or one tittle' shall pass away till all be fulfilled.\footnote{The first edition of Bownd's book appeared in 1595, and was dedicated to the Earl of Essex (see the title in Vol. V. p. 211 of Fuller's \emph{Church History}, Brewer's ed.). The second and enlarged edition of 1606 was dedicated to the Bishop of Norwich and the Dean of Ely, and bears the following characteristic title (which somewhat differs from the title of the first): '\emph{Sabbathum Veteris et Novi Testamenti: or, The True Doctrine of the Sabbath, held and practised of the Church of God, both before and under the Law, and in the time of the Gospel: Plainly laid forth and soundly proved by testimonies both of Holy Scripture and also of old and new Ecclesiastical Writers, Fathers and Councils, and Laws of all sorts, both civil, canon, and common. Declaring first from what things God would have us straitly to rest upon the Lord's day, and then by what means we ought publicly and privately to sanctify the same. Together with the sundry Abuses of men in both these kinds, and how they ought to be reformed. Divided into two Books by} Nicolas Bownd, \emph{Doctor of Divinity; and now by him the second time perused, and enlarged with an Interpretation of sundry points belonging to the Sabbath, and a more ample proof of such things as have been gainsaid or doubted of by some divines of our time, and a more full Answer unto certain objections made against the same: with some other things not impertinent to this argument.}' London, 1606, 4to, pp. 479. Having been unable to obtain this rare work, I copied the title from Robert Cox, \emph{The Literature of the Sabbath Question} (in 2 vols. Edinb. 1865), Vol. I. p. 145. There is a copy in the Bodleian Library, and another in the library of the University of Edinburgh. Cox himself is opposed to the Puritan theory, and holds the Church of England responsible for originating it by requiring the fourth commandment to be read and responded to in the Liturgy. Of Bownd's book he says: 'In the treatise bearing this long title the Sabbatarian opinions of the Puritans, which afterwards found more precise expression in the Westminster Confession and Catechisms, and are now maintained by the Evangelical sects in this country, were for the first time broadly and prominently asserted in Christendom.' Fuller gives a full account of the contents, Vol. V. pp. 211 sqq. His editor, Brewer, says that Bownd's book 'is written in a truly Christian spirit, and ought by no means to be considered as the fruit of Puritan principles.' The accounts of Collier (\emph{Eccl. Hist.} Vol. VII. pp. 182 sqq.), Neal (Vol. I. pp. 208 sq.), and Hesse (\emph{Sunday}, pp. 276 sqq.) are drawn from Fuller.}

The work in which this theory was ably and earnestly vindicated proved to be a tract for the times. Heylin, a High-Church opponent, says 'that in a very little time it grew the most bewitching error, the most popular deceit that had ever been set on foot in the Church of England.'\footnote{Quoted by Hessey, p. 281.} Fuller dates from it 'the more solemn and strict observance of the Lord's day,' and gives the following description of the effect produced by it:

'It is almost incredible how taking this doctrine was, partly because of its own purity, and partly for the eminent piety of such persons as maintained it, so that the Lord's day, especially in corporations, began to be precisely kept, people becoming a law to themselves, forbearing such sports as [were] yet by statute permitted; yea, many rejoicing at their own restraint therein. On this day the stoutest fencer laid down the buckler, the most skilful archer unbent his bow, counting all shooting besides the mark; May-games and Morris-dances grew out of request, and good reason that bells should be silenced from gingling about men's legs, if their very ringing in steeples were adjudged unlawful; some of them were ashamed of their former pleasures, like children which, grown bigger, blushing themselves out of their rattles and whistles. Others forbore them for fear of their superiors, and many left them off out of a politic compliance, lest otherwise they should be accounted licentious.

'Yet learned men were much divided in their judgments about these Sabbatarian doctrines. Some embraced them as ancient truths consonant to Scripture, long disused and neglected, now seasonably revived for the increase of piety. Others conceived them grounded on a wrong bottom, but because they tended to the manifest advance of religion it was pity to oppose them, seeing none have just reason to complain being deceived into their own good. But a third sort flatly fell out with these positions, as galling men's necks with a Jewish yoke, against the liberty of Christians: that Christ, as Lord of the Sabbath, had removed the rigor thereof, and allowed men lawful recreations; that this doctrine put an unequal lustre on the Sunday, on set purpose to eclipse all other holy days, to the derogation of the authority of the Church; that the strict observance was set up out of faction to be a character of difference, to brand all for libertines who did not entertain it.'\footnote{Vol. V. pp. 214 sqq.}

The Puritan Sabbath theory was denounced and assailed by the rising school of High-Churchism as a Sabbatarian heresy and a cunningly concealed attack on the authority of the Church of England, by substituting the Jewish Sabbath for the Christian Sunday and all the Church festivals.\footnote{The chief writers against the Puritan theory were Thomas Rogers, Bancroft's chaplain (in his \emph{Preface to the Articles}); and afterwards Bishop White of Ely (\emph{A Treatise of the Sabbath-Day . . . against Sabbatarian Novelty}, Lond. 1635); Peter Heylin, Laud's chaplain (\emph{The History of the Sabbath}, Lond. 2d ed. 1636); and Dr. John Pocklington (\emph{Sunday no Sabbath}, Lond. 1636). See extracts from their works by Cox, 1.c. Vol. I. pp. 166 sqq. White and Heylin wrote at the request of Laud. Bishop Prideaux (1622), Bishop Cosin (1635), and Dr. Young (1639) took a more moderate view. Richard Baxter (1671), though strongly leaning to the Puritanic side, tried to mediate between the strict Sabbath theory and the ecclesiastical Sunday theory, and maintained the joyous rather than the penitential character of the Lord's day. See Hessey, pp. 288 sq.} Attempts were made by Archbishop Whitgift in 1599, and by Chief Justice Popham in 1600, to suppress Bownd's book and to destroy all the copies, but 'the more it was called in the more it was called on;' its price was doubled, and 'though the book's wings were clipped from flying abroad in print, it ran the faster from friend to friend in transcribed copies, and the Lord's day, in most places, was most strictly observed. The more liberty people were offered the less they used it. . . . It was sport for them to refrain from sports. . . . Scarce any comment, catechism, or controversy was set forth by the stricter divines, wherein this doctrine (the diamond in this ring) was not largely pressed and proved; so that, as one saith, the Sabbath itself had no rest.'\footnote{Fuller, pp. 218, 219.}

At last King James I. brought his royal authority to bear against the Puritan Sabbatarianism so called, and issued the famous 'Book of Sports,' May 24, 1618, which was afterwards republished, with an additional order, by his son, Charles I., no doubt by advice of Archbishop Laud, Oct. 18, 1633.\footnote{Of the first edition no copy is known to exist. The second edition, of which a copy is preserved in the British Museum, bears the title: 'The Kings | \emph{Maiesties | Declaration to | His Subjects}, | Concerning | \emph{lawfull } Sports  \emph{to | bee vsed. | Imprinted at } London  \emph{by | Robert Barker, Printer to the Kings | most Excellent Maiesties And by | the Assignes of John Bill.} | M.DC.XXXIII.' 4to, 24 pp. This edition has been reprinted on tinted paper, in exact imitation of the original, at London (Bernard Quaritch), 15 Piccadilly, 1860. The Long Parliament, in 1643, ordered the book to be burned by the common hangman, in Cheapside and other places.} This curious production formally authorizes and commends the desecration of the evening of the Lord's day by dancing, leaping, fencing, and other 'lawful recreations,' on condition of observing the earlier part by strict outward conformity to the worship of the Church of England.\footnote{'Our expresse pleasure therefore is, that. . . no lawfull Recreation shall bee barred to Our good People, which shall not tend to the breach of Our aforesayd Lawes, and Canons of Our Church: which to expresse more particularly, Our pleasure is, That the Bishop, and all other inferiour Churchmen, and Churchwardens, shall for their parts bee carefull and diligent, both to instruct the ignorant, and conuince and reforme them that are mis-led in Religion, presenting them that will not conforme themselues, but obstinately stand out to Our Judges and Iustices: Whom We likewise command to put the Law in due execution against them. \par  'Our pleasure likewise is, That the Bishop of that Diocesse take the like straight order with all the Puritanes and Precisians within the same, either constraining them to conforme themselues, or to leaue the Country according to the Lawes of Our Kingdome, and Canons of Our Church, and so to strike equally on both hands, against the contemners of Our Authority, and aduersaries of Our Church. And as for Our good peoples lawfull Recreation, Our pleasure likewise is, That after the end of Diuine Seruice, Our good people be not disturbed, letted, or discouraged from any lawfull recreation, Such as dauncing, either men or women, Archery for men, leaping, vaulting, or any other such harmelesse Recreation, nor from hauing of May-Games, Whitson Ales, and Morris-dances, and the setting vp of May-poles \& other sports therewith vsed, so as the same be had in due \& conuenient time, without impediment or neglect of Diuine Seruice.'—\emph{Book of Sports}, pp. 8 sqq.} The professed object of this indulgence to the common people was to check the progress of the Papists and Puritans (or 'Precisians'), and to make 'the bodies more able for war' when his majesty should have 'occasion to use them.' The court set the example of desecration by balls, masquerades, and plays on Sunday evening; and the rustics repaired from the house of worship to the ale-house or the village green to dance around the Maypole and to shoot at butts. To complete the folly, King James ordered the book to be read in every parish church, and threatened clergymen who refused to do so with severe punishment. King Charles repeated the order. But in both cases it became the source of great trouble and confusion.\footnote{Fuller says (Vol. V. p. 452): 'When this declaration was brought abroad, it is not so hard to believe as sad to recount what grief and distraction thereby was occasioned in many honest men's hearts.'} Several bishops disapproved of it. Archbishop Abbot (the Puritan predecessor of Laud) flatly forbade it to be read at Croydon. The Lord Mayor of London commanded the king's own carriages to be stopped as they were passing through the city on a Sunday. James raged and swore, and countermanded the prohibition. The Lord Mayor yielded, with this answer: 'While I was in my power I did my duty, but that being taken away, it is my duty to obey.' Some clergymen, after reading the book from the pulpit, followed it up by a sermon against it, or by reading the fourth commandment—'Remember the Sabbath day to keep it holy'—and added, 'This is the law of God, the other the injunction of man.' Those who refused to read the royal Book of Sports were suspended from office and benefice, or even excommunicated by Laud and his sympathizing fellow-bishops.\footnote{Prynne says: 'How many hundred godly ministers have been suspended from their ministry, sequestered, driven from their livings, excommunicated, prosecuted in the High Commission, and forced to leave the kingdom, for not publishing this declaration, is experimentally known to all men.' For particulars, see Neal, Vol. I. pp. 312 sqq.} Many left England, and joined

'The pilgrim bands, who crossed the sea to keep

Their Sabbaths in the eye of God alone,

In his wide temple of the wilderness.''

This persecution of conscientious ministers for obeying God rather than men gave moral strength to the cause of Sabbath observance, and rooted it deeper in the affections of the people. It was one of the potent causes which overwhelmed Charles and Laud in common ruin. The sober and serious part of the nation were struck with a kind of horror that they should be invited by the highest authorities in Church and State to destroy the effect of public worship by a desecration of a portion of the day consecrated to religion.

On the Sunday question Puritanism achieved at last a permanent triumph, and left its trace upon the Church of England and Scotland, which reappeared after the licentious period of the Restoration. For, although the Church of England, as a body, never committed itself to the Puritan Sabbath theory, it adopted at least the practice of a much stricter observance than had previously obtained under Elizabeth and the Stuarts, and would never exchange it for the Continental laxity, with its disastrous effects upon the attendance at public worship and the morals of the people.

The Westminster Confession, without entering into details or sanctioning the incidental excesses of the Puritan practice, represents the Christian rest-day under its threefold aspect: (1) as a divine law of nature (\emph{jus divinum naturale}), rooted in the constitution of man, and hence instituted (together with marriage) at the creation, in the state of innocence, for the perpetual benefit of body and soul; (2) as a positive moral law (\emph{jus divinum positivum}), given through Moses, with reference to the primitive institution ('Remember') and to the typical redemption of Israel from bondage; (3) as the commemoration of the new creation and finished redemption by the resurrection of Christ; hence the change from the last to the first day of the week, and its designation 'the Lord's day' (\emph{dies Dominica}). And it requires the day to be wholly devoted to the exercises of public and private worship and the duties of necessity and mercy.

To this doctrine and practice the Presbyterian, Congregational, and other Churches in Scotland, England, and America have faithfully adhered to this day. Yea, twenty-seven years before it was formulated by the learned divines of Westminster, the Pilgrim Fathers of America had transplanted both theory and practice first to Holland, and, finding them unsafe there, to the wild soil of New England. Two days after their landing from the \emph{Mayflower} (Dec. 22, 1620), forgetting the pressing necessities of physical food and shelter, the dreary cold of winter, the danger threatening from wild beasts and roaming savages, they celebrated their first Sunday in America on a barren rock and under the stormy sky of heaven, and, in the exercise of the general priesthood of believers, they offered the sacrifices of contrite hearts and the praises of devout lips to their God and Saviour, on his own appointed day of holy rest; not dreaming that they were the bearers of the hopes and destinies of a mighty future and the founders of a republic stretching across a continent and embracing millions of intelligent Christian freemen.\footnote{\par Comp. my essay on the \emph{Anglo-American Sabbath.} New York, 1863.}

The political articles of the Confession touching the power of the civil magistrate and the relation of Church and State will be discussed hereafter (§ 97) in connection with the subject of religious toleration and the changes which have been introduced in later editions.

\xsubsection{The Westminster Catechisms.}

§ 96. The Westminster Catechisms.

Editions.

\emph{The Humble} | Advice | \emph{of the} | Assembly | \emph{of} | Divines, | \emph{Now by Authority of Parliament | sitting at} | Westminster; | \emph{Concerning} | A Larger Catechism: | \emph{Presented by them lately to both Houses of Parliament.} | Printed at London [Oct. 1647, without Scripture proofs], and reprinted at Edinburgh, by Evan Tyler, Printer to the King's most Excellent Majestie, 1647 [Dec.]. The Edinburgh reprint has fifty-six pages, and no Scripture proofs. See fac-simile in Vol. III. p. 674. Of the London \emph{editio princeps}, six hundred copies were printed, but not published, by order of Parliament, for its own use. Of the Edinburgh \emph{editio princeps}, eight hundred copies were ordered by the General Assembly, Dec. 23, 1647. The second ed., which appeared in London [after April 14, 1648], contains the proofs from Scripture.

The Shorter Catechism appeared under the same title (except \emph{Shorter} for \emph{Larger}) a little later [after Nov. 25, 1647], by order of Parliament. Mr. John Laing, the obliging librarian of the Free Church College in Edinburgh, informs me that both Catechisms appeared in one vol. of seventy-nine pages, at Edinburgh, Dec. 23,1647, with a general title and a separate title for each. A statement to the same effect I see in the Advertisement to Dunlop's \emph{Collection of Confessions}, Vol. I. p. clviii., with the additional remark that this edition was sent to the Presbyteries for examination.

The Larger and Shorter Catechisms often appeared in connection with the Westminster Confession, and exist in innumerable English and American editions, especially the Shorter. The textual variations are insignificant, except that the American (General Assembly's) editions of the Larger Catechism omit the words 'tolerating a false religion' in the answer to Question 109.

I have made use of the first Edinb. ed., and a large London ed. of 1658, which contains the Conf. and both Catechisms under their original (three separate) titles (\emph{The Humble Advice}, etc.), with the Scripture proofs in full. Opposite the special title of the Shorter Catechism is the order of Parliament, dated 'Die Lunæ 15. Septemb., 1648,' directing that the Shorter Catechism 'be forthwith printed and published, wherein Mr. Henry Roborough and Mr. Adoniram Byfield, Scribes of the Assembly of Divines, are requested to use all possible care and diligence.'

The Catechisms have been translated into many languages, especially the Shorter. A Latin version appeared, together with the version of the Confession, in Cambridge, 1656, as has been noted above, p. 753. The Latin text of the Shorter Catechism is printed in Vol. III. pp. 676 sqq. For a German version of both, see Böckel, pp. 716 sqq. A Greek version of the Shorter Catechism (with the Latin), by John Harmar (Regius Professor of Greek In Oxford), was published at London, 1660; a new one by Robert Young (\textgreek{ἡ κατήχησις συντομωτέρα}), Edinburgh, 1854. A Hebrew version by G. Seaman, M.D. (London, 1689), and another by H. S. McKee (Edinb. 1854; Dublin, 1864). Also Syriac, Arabic, modern Greek, Portuguese, Welsh, and other versions.

The largest number of editions and translations are to be found, as far as I know, in the British Museum.

Expositions.

Thomas Lye  (Minister in London, d. 1684): \emph{An Explanation of the Shorter Catechism.} London, 1676.

Hugh Binning (d. 1653, Prof. of Moral Philos., Glasgow): \emph{The Common Principles of the Christian Religion. . . . A Practical Catechism.} 1671.

Thomas Vincent (Minister in London, d. 1671): \emph{An Explanation of the Assembly's Shorter Catechism.} London, 1708; Edinb. 1799; Presbyterian Board of Publication, Philadelphia.

Thomas Watson (Minister in London, d. 1690): \emph{A Body of Practical Divinity, consisting of above} 176 \emph{Sermons on the Shorter Catechism.} 5th ed. Glasgow, 1797; Lond. 1807; Glasgow, 1838; N. Y. 1836.

John Flavel, (b. 1627, d. 1691): \emph{Exposition of the Catechism.} 1692. In his \emph{Whole Works}, 2 vols. fol. 1701, 7th ed. Edinb. 1762; and in 6 vols. London, 1820.

Thomas Ridgley (b. 1667, d. 1734): \emph{A Body of Divinity . . . Being the Substance of Lectures on the Assembly's Larger Catechism}, London, 1731–33, 2 vols. fol.; an ed. in 4 vols. 8vo, 1814; Edinb. 1845, 2 vols. 8vo; New York, 1855.

Samuel Willard (b. 1640, d. 1707): \emph{A Body of Divinity in} 250 \emph{Lectures on the Assembly's Catechism.} 1 vol. fol. Boston, 1726.

John Willison (Minister of Dundee from 1718 to 1750): \emph{An Example of Plain Catechising upon the Assembly's Shorter Catechism.} Edinb. 1737; 2d ed. Glasgow, 1764.

Fisher's Catechism: \emph{The Westminster Assembly's Shorter Catechism Explained}, by way of question and answer. \emph{By some Ministers of the Gospel.} The authors are Ralph Erskine (d. 1752), Ebenezer Erskine (d. 1754), and  James Fisher (d. Sept. 28, 1775, Secession Minister at Greyfriars, Glasgow). Fisher prepared the second part alone, and issued the third ed. Glasgow, 1753. Hence the whole work is called by his name. 14th ed. Edinb. 1800; 17th ed. Glasgow, 1813; also by the Board of Publication, Philadelphia.

John Brown (Minister at Haddington from 1751 to 1787): \emph{Easy Explication of the Assembly's Shorter Catechism.} 8th ed. Edinb. 1812; 9th ed. Montrose, 1822.

Henry Belfrage (d. 1835): \emph{A Practical Exposition of the Assembly's Shorter Catechism, exhibiting a System of Theology in a Popular Form.} Edinb. 2d ed. 1834. 2 vols.

Alex. Mair (d.1751): \emph{A Brief Explication of the Assembly's Shorter Catechism.} New ed. Montrose, 1837.

Alex. Smith Paterson: \emph{A Concise System of Theology: being the Shorter Catechism Analyzed and Explained.} Edinb. 1841; 2d ed. 1844.

Ashbel Green, D.D. (President of Princeton College from 1812 to 1822; d. 1848): \emph{Lectures on the Shorter Catechism.} Phila. 1841, 2 vols., Presbyt. Board of Publ.

Jonathan Cross: \emph{Illustrations of the Shorter Catechism. Proof-texts, Exposition, and Anecdotes.} 2 vols. 18mo. Presbyt. Board of Publ.

Edwin Hall, D.D.: \emph{The Shorter Catechism of the Westminster Assembly, with Analysis and Scripture Proofs.} Presbyt. Board of Publ.

James R. Boyd, D.D.: \emph{The Westminster Shorter Catechism; with Analysis, Proofs, Explanations, and Illustrative Anecdotes.} 18mo. Presbyt. Board of Publ.

\emph{The Bellefonte Series of Tracts} on the Answers to the Shorter Catechism, written by numerous Presbyterian ministers, and edited by the Rev. Wm. T. Wylie. Bellefonte, Pa. 1875.

PREPARATION AND ADOPTION.

Simultaneously with the Confession, the Assembly prepared first one, and afterwards two Catechisms: a larger one for public exposition in the pulpit, according to the custom of the Reformed Churches on the Continent, and a smaller one for the instruction of children, a clear and condensed summary of the former.\footnote{The first Catechism of the Assembly, according to Baillie, was nearly agreed on at the end of 1644, but was never published. Perhaps it was the same which is partially inserted in the \emph{Minutes;} or it may have been the MS. Catechism of Sam. Rutherford, which is preserved in the University library at Edinburgh. In the 774th session, Jan. 14, 1647 (old style, 1646), the Assembly ordered 'that the Committee for the Catechism do prepare a draught of \emph{two} Catechisms, one more large and another more brief, in which they are to have an eye to the Confession of Faith, and to the matter of the Catechism already begun' (\emph{Minutes}, p. 321).} Both are amply provided with Scripture proofs. The questions of Church polity and discipline are properly omitted.

The Catechisms were finished and presented to Parliament for examination and approval in the autumn of 1647.\footnote{Both Catechisms were first presented to Parliament \emph{without Scripture proofs}, the Larger before Oct. 25, 1647, the Shorter on Nov. 25, 1647 (\emph{Minutes}, pp. 485, 486, 492), and were forthwith printed in London and Edinburgh. The Catechisms \emph{with Scripture proofs} were presented to Parliament on or before April 14, 1648 (\emph{Minutes}, p. 511).} Parliament ordered six hundred copies to be printed, and then examined and approved the Catechisms, with some slight exceptions (Sept. 15, 1648). The General Assembly at Edinburgh adopted the Larger Catechism, July 20, 1648, and the Shorter Catechism, July 28, declaring both to be 'agreeable to the Word of God, and in nothing contrary to the received doctrine, worship, discipline, and government of this Kirk.' These acts were approved by the Scottish Parliament, Feb. 7, 1649, but repealed under Charles II. in 1661. When the Scottish Parliament, in 1690, established Presbyterian government in Scotland, and ratified the Westminster Confession of Faith, no express mention was made of the Catechisms, but both continued in ecclesiastical use, and the Shorter Catechism was often earnestly enjoined upon ministers, teachers, and parents by the General Assembly.\footnote{Mitchell, \emph{Minutes}, p. 515. note. Innes (\emph{Law of Creeds}, p. 195) says: 'The Shorter Catechism has been for many generations the real creed of Scotland, so far as the mass of the people is concerned.'}

GENERAL CHARACTER.

The two Catechisms are, in the language of a Scotch divine, 'inimitable as theological summaries; though, when it is considered that to comprehend them would imply an acquaintance with the whole circle of dogmatic and controversial divinity, it may be doubted whether either of them is adapted to the capacity of childhood. . . . Experience has shown that few who have been carefully instructed in our Shorter Catechism have failed to discover the advantage of becoming acquainted in early life, even as a task, with that admirable ``form of sound words.'''\footnote{M'Crie, \emph{Annals}, pp. 177 sq. Neal (Vol. II. p. 42) judges similarly. 'The Larger Catechism,' he says, 'is a comprehensive system of divinity, and the smaller a very accurate summary, though it has by some been thought a little too long, and in some particulars too abstruse for the capacities of children.' Baillie was of the same opinion (\emph{Letters}, III. 59).}

Both Catechisms have the peculiarity that each answer embodies the question, and thus forms a complete proposition or sentence in itself.

Both depart from the catechetical tradition by omitting the Apostles' Creed, which in other orthodox Catechisms is the common historical basis of the exposition of the Articles of Faith. It is, however, annexed to the Shorter Catechism,' not as though it were composed by the Apostles or ought to be esteemed canonical Scripture, as the Ten Commandments and the Lord's Prayer, but because it is a brief sum of the Christian faith, agreeable to the Word of God, and anciently received in the Churches of Christ.' A note is attached to the article on the descent into Hell (better, \emph{Hades} or \emph{Sheol}), to the effect that it simply means Christ 'continued in the state of the dead and under the power of death until the third day.' This explanation (like that of Calvin and the Heidelberg Catechism) misses the true sense of the descent, and ignores its peculiar significance in the work of redemption for the world of the departed (comp. Luke xxiii. 43; Acts ii. 31; Eph. iv. 8, 9; 1 Cor. xv. 55, 57; 1 Pet. iii. 18, 19; iv. 6; Rev. i. 18). The eschatology of the Reformation standards is silent or defective on the middle state, and most Protestant versions of the Bible confound \emph{Hell} and \emph{Hades}, which represent separate and distinct though cognate ideas.

THE LARGER CATECHISM.

The Larger Catechism occupied, as the \emph{Minutes} show, a good deal of the Assembly's attention during the year 1647, and was discussed question by question. It was prepared before the Shorter.\footnote{This appears from the \emph{Minutes}, p. 410. The report on the Shorter Catechism was first called for in the 896th session, Aug. 9, 1647. Mr. Palmer reported, and Messrs. Calamy and Gower were added to the Committee. The opposite view is clearly wrong, though advocated by Neal (Vol. II. p. 42), and even quite recently by Dr. M'Crie, who says (\emph{Annals}, p. 177): 'The Larger Catechism was not prepared till some time \emph{after the Shorter}, of which it was evidently intended to form an amplification and exposition.'} It is chiefly the work of Dr. Anthony Tuckney, Professor of Divinity and Vice-Chancellor at Cambridge.\footnote{It is based in part on Ussher's catechetical \emph{Body of Divinity}, perhaps also on the concise theological compendium of John Wolleb, Antistes at Basle (1626).} It is a masterpiece of catechetical skill, superior to any similar work, and exhibits in popular form a complete system of divinity, like the Roman Catechism and the Longer Russian Catechism of Philaret. It also serves in part as a valuable commentary or supplement to the Confession, especially on the ethical part of our religion. But it is over-minute in the specification of what God has commanded and forbidden in the Ten Commandments, and loses itself in a wilderness of details.\footnote{Take for example Question 113:  \par \emph{What are the sins forbidden in the third commandment?} \par 'The sins forbidden in the third commandment are, the not using of God's name as is required; and the abuse of it in an ignorant, vain, irreverent, profane, superstitious, or wicked mentioning, or otherwise using his titles, attributes, ordinances, or works, by blasphemy, perjury; all sinful cursings, oaths, vows, and lots; violating our oaths and vows, if lawful; and fulfilling them, if of things unlawful; murmuring and quarreling at, curious prying into, and misapplying of God's decrees and providences; misinterpreting, misapplying, or any way perverting the Word, or any part of it, to profane jests, curious or unprofitable questions, vain janglings, or the maintaining of false doctrines; abusing it, the creatures, or any thing contained under the name of God, to charms or sinful lusts and practices; the maligning, scorning, reviling, or any wise opposing God's truth, grace, and ways; making profession of religion in hypocrisy or for sinister ends; being ashamed of it, or a shame to it, by uncomformable, unwise, unfruitful, and offensive walking or backsliding from it.'}

THE SHORTER CATECHISM.

Dr. Tuckney was also the convener of the Committee which prepared the Shorter Catechism, but its concise and severely logical answers are traced to the Rev. John Wallis, M.A., an eminent mathematician, who as a young man fresh from Cambridge was appointed an amanuensis of the Assembly.\footnote{In the \emph{Minutes}, p. 488, Wallis is mentioned in connection with the Shorter Catechism. He published an exposition of it.} He afterwards became Professor of Geometry at Oxford and one of the founders of the Royal Society. He was probably the last survivor of the Westminster divines, for he died 1703, \emph{aet.} eighty-eight.\footnote{Masson's \emph{Milton}, Vol. II. p. 515.} Gillespie's name is traditionally connected with the question 'What is God?' He is said to have answered it in prayer, apparently without meditation, when the Assembly were in suspense for words to define the Being of beings. But the Scotch Commissioners had little to do with the Shorter Catechism, as most of them had left before it was discussed in the Assembly.\footnote{The Scotch Commissioners took leave Dec. 25, 1646. The last mention of them is Nov. 9, 1647, when Rutherford took his leave.—\emph{Minutes}, pp. 471, 487. Dr. Mitchell informs me that the fourth question is probably derived from 'A Compendious Catechism' (by J. F.), printed at London in April, 1645: 'God is a Spirit, One, Almighty, Eternal, Infinite, Unchangeable Being, Absolutely Holy, Wise, Just, and Good.'}

The Shorter Catechism is one of the three typical Catechisms of Protestantism which are likely to last to the end of time. It is fully equal to Luther's and to the Heidelberg Catechism in ability and influence, it far surpasses them in clearness and careful wording, and is better adapted to the Scotch and Anglo-American mind, but it lacks their genial warmth, freshness, and childlike simplicity.\footnote{For a fuller comparison, see pp. 543–545.} It substitutes a logical scheme for the historical order of the Apostles' Creed. It deals in dogmas rather than facts. It addresses the disciple as an interested outsider rather than as a church-member growing up in the nurture of the Lord. Its mathematical precision in definitions, some of which are almost perfect,\footnote{For example, Questions 4, 21, 92.} though above the capacity of the child, is a good preparation for the study of theology. Its use among three denominations (Presbyterians, Congregationalists, and Regular Baptists) proves its solid worth. Baxter called it 'the best Catechism he ever saw, a most excellent sum of the Christian faith and doctrine, and a fit test to try the orthodoxy of teachers.' Thomas Carlyle, in speaking against modern materialism, made this confession (1876): 'The older I grow—and I now stand upon the brink of eternity—the more comes back to me the first sentence in the Catechism which I learned when a child, and the fuller and deeper its meaning becomes: ``What is the chief end of man? To glorify God, and to enjoy him forever.'''

\xsubsection{Criticism of the Westminster System of Doctrine.}

§ 97. Criticism of the Westminster System of Doctrine.

The Westminster Confession, together with the Catechisms, is the fullest and ripest symbolical statement of the Calvinistic system of doctrine. In theological ability and merit it is equal to the best works of the kind, and is not surpassed by the Lutheran Formula of Concord or the Roman Decrees of the Councils of Trent and the Vatican. Its intrinsic worth alone can explain the fact that it has supplanted the older Scottish standards of John Knox and John Craig in the land of their birth, and that it was adopted by three distinct denominations: by the Presbyterians in full, and by the Congregationalists and the Regular Baptists with some slight modifications. Of these the Congregationalists had but a small though very able representation in the Westminster Assembly, the Baptists none at all. It has at this day as much vitality as any of the Protestant symbols and more vitality than most of them. It materially aids in shaping theological thought and religious activity as far as the English tongue prevails. Altogether it represents the most vigorous and yet moderate form of Calvinism, which has found (like Christianity itself) a more congenial and permanent home in the Anglo-Saxon race than in the land of its birth.

The doctrines of the Confession are stated with unusual care, logical precision, clearness, caution, and circumspection, and with an eye to all their various aspects and mutual relations. Where they seem to conflict or can not be harmonized by our finite intelligence—as absolute sovereignty and free agency, the fall of Adam and personal guilt, the infinite divinity and the finite humanity of Christ—both truths are set forth, and room is left for explanations and adjustments by scientific theology within the general limits of the system. The important difference between a public confession of faith and a private system of theology was at least distinctly recognized in principle, although (as we shall see presently) not always consistently carried out.\footnote{In the debate on predestination Dr. Reynolds wisely said, 'Let us not put disputes and scholastic things into a confession of faith.'—\emph{Minutes}, p. 151.}

The style of the Confession and Catechisms is clear, strong, dignified, and well adapted to the grave subject. The selection of Scripture proofs is careful and judicious, and reveals a close familiarity with the sacred writings.

The merits of the Westminster standards have been admitted not only by Presbyterians,\footnote{Principal Baillie wrote (Jan. 26, 1647, \emph{Letters}, Vol. III. p. 2): 'The Confession is much cried up by all, \emph{even many of our greatest opposites}, as the best confession yet extant.' The moderate and judicious Richard Baxter esteemed the Westminster Confession and Catechisms the best books in his library next to the Bible, and says (in his \emph{Confession}, ch. i. § 5): 'I have perused oft the Confession of the Assembly, and verily judge it the most excellent, for fullness and exactness, that I have ever read from any Church; and though the truths therein, being of several degrees of evidence and necessity, I do not hold them with equal clearness, confidence, or certainty; and though some few points in it are beyond my reach, yet I have observed nothing in it contrary to my judgment, if I may be allowed those expositions following.' The saintly Archbishop Leighton, though he left the Church for which his father had suffered such cruelties from Laud, taught the doctrine of the Confession to the end of his life.} but also by liberal Episcopalians,\footnote{J. B. Marsden (\emph{The History of the Later Puritans}, 1852, pp. 80, 81), while judging severely of the Assembly on account of its treatment of Episcopacy, thinks the Westminster Confession inferior to none of the Protestant Confessions except in originality, and adds: 'It does not, however, detract from the real merit of these later divines, that they availed themselves of the labors of the Reformation; or that Bullinger and Calvin, especially the latter, should have left them little to accomplish, except in the way of arrangement and compression. The Westminster Confession should be read by those who can not encounter the more ponderous volumes of the great masters from which it is derived. It is in many respects an admirable summary of Christian faith and practice. None can lay it down with a mean opinion of the Westminster divines. The style is pure and good, the proofs are selected with admirable skill, the arguments are always clear, the subjects well distributed, and sufficiently comprehensive to form at least the outline of a perfect system of divinity.' It is but just to add that Marsden goes on to censure what he calls its 'rigid ultra-Calvinism, which has always repelled the great majority of English Christians.' Dean Stanley, who has no theological sympathy with the Westminster Confession, says that of all Protestant Confessions 'it far more nearly approaches the full proportions of a theological treatise, and exhibits \emph{far more depth of theological insight, than any other},' He adds, however, that 'it reflects also far more than any other the minute hair-splitting and straw-dividing distinctions which had reached their height in the Puritanical theology of that age, and which in sermons ran into the sixteenthly, seventeenthly sections that so exercised the soul of Dugald Dalgetty as he waited for the conclusion of the discourse in the chapel of Inverary Castle. It accordingly furnished the food for which the somewhat hard and logical intellect of Scotland had a special appetite' (\emph{Lectures on the History of the Church of Scotland}, delivered in 1872, Am. ed. p. 88). In another place Stanley calls the Westminster formulary 'that famous Confession of Faith which, alone within these islands, was imposed by law on the whole kingdom; and which, alone of all Protestant Confessions, still, in spite of its sternness and narrowness, retains a hold on the minds of its adherents, to which its fervor and its logical coherence in some measure entitle it' (\emph{Memorials of Westminster Abbey}, p. 513).} and even by Methodists, who entirely dissent from its theology.\footnote{Dr. Currey, for many years editor of the 'Methodist Advocate,' of New York, in an editorial on Creeds (Aug. 6, 1874), calls 'the Westminster Confession of Faith the \emph{ablest, clearest, and most comprehensive system of Christian doctrine ever framed.} That venerable instrument purposely embodies in its unity the dogma of absolute predestination, which necessarily becomes the corner-stone of the edifice, so giving it shape and character. But, despite that capital fault, it is not only a wonderful monument of the intellectual greatness of its framers, but a comprehensive embodiment of nearly all the precious truths of the gospel. If set forth without ecclesiastical authority, for the edification of believers, it would, despite its faults, be a work of inestimable worth; but enforced by such authority, and imposed upon men's consciences, it is a yoke and a chain and a cage of iron. And yet this is the accepted formula of faith of nearly all the Calvinistic Churches of America. Even the Congregationalists in National Council, at Plymouth Rock, only a few years ago, reaffirmed their acceptance of it.'}

DEFECTS.

The Westminster standards, like all human productions, including the translations of the Bible itself, have imperfections.

The great revival of the sixteenth century was followed in the Reformed and Lutheran Churches by a dry scholasticism which was more biblical and evangelical than the mediæval scholasticism, but shared with it the defects of a one-sided intellectualism to the exclusion of the mystic and emotional types of Christianity. Scholasticism in the technical sense—whether Roman Catholic or Protestant—is the product of the devout understanding rather than the glowing heart, and approaches the deepest mysteries of faith, such as the Trinity, the Incarnation, the eternal decrees of election and reprobation of men and angels, with profound reverence indeed, yet with a boldness and assurance as if they were mathematical problems or subjects of anatomical dissection.\footnote{Dr. Wallis, the mathematician, who is said to be the chief author of the sharp definitions of the Shorter Catechism (see p. 786), wrote towards the close of the seventeenth century a pamphlet in defense of the doctrine of the Trinity against rising Unitarianism, where he compares the Almighty to a \emph{cube} with its length, breadth, and height infinitely extended, \emph{longum, latum, profundum}, which are the equal sides of one substance, and fairly resemble the Father, Son, and Holy Ghost. He finds nothing mysterious in this doctrine. 'It is,' he says, 'but this, that there be three \emph{somewhats}, which are but one God, and these \emph{somewhats} are called Persons.' Quoted by Stoughton, \emph{The Church of the Revolution}, p. 213.} It shows usually a marvelous dexterity in analysis, division, subdivision, distinction, and definition, but it lacks the intuition into the hidden depths and transcending heights where the antagonisms of partial truths meet in unity.

The Westminster standards do not go so far in this direction as the Canons of Dort or the Helvetic Consensus Formula, but certainly further than the Reformation symbols, which are less logical and precise, and more fresh and elastic. They reflect the hard severity of Puritanism. They embody too much metaphysical divinity, and overstep the limits which divide a public confession of faith from a scientific treatise of theology. It would be impossible nowadays to pass such an elaborate system through any Protestant ecclesiastical body with a view to impose it upon all teachers of religion. The Confession, however, as already mentioned, was not intended as a yoke by the English framers, nor has subscription ever been required to all its details, but only to the general scheme. The Bible is expressly declared by Calvinists to be 'the only infallible rule of faith and practice,' and the Confession is adopted 'as containing the system of doctrine taught in the holy Scriptures.'\footnote{This is the American formula of subscription required from ministers. On the Scottish subscription formulas, see Innes, pp. 66, 81, 84, 103, 453.}

The chief characteristics of Calvinistic scholasticism as it prevailed in the seventeenth century are that it starts from God's sovereignty and justice rather than from God's love and mercy, and that it makes the predestinarian scheme to control the historical and christological scheme. This brings us to the most assailable point in the Westminster Confession and Larger Catechism, the abstract doctrine of eternal decrees, which will always repel a large portion of evangelical Christendom. We believe that the divine-human person and work of Christ furnish the true key to the full understanding of the plan of salvation and the solid platform for the ultimate agreement of all evangelical creeds.

PRETERITION OF THE REST OF MANKIND.

Absolute predestinarianism is the strength and the weakness of Calvinism. The positive decree of eternal election is its impregnable fort, the negative decree of eternal reprobation its Achilles' heel. Predestination to holiness and happiness, being a gracious purpose of God's love, is full of 'sweet, pleasant, and unspeakable comfort to godly persons,'\footnote{Articles of the Church of England, Art. XVII.} and affords 'matter of praise, reverence, and admiration of God, and of humility, diligence, and abundant consolation to all that sincerely obey the gospel.'\footnote{Westm. Conf. Ch. VIII. § 8. This last section is the best in the whole chapter.} Predestination to death and damnation, being a judicial decree of God's wrath on account of Adam's fall, is—whether true or false—a '\emph{decretum horribile}' (as Calvin himself significantly calls it, in view of the apparent ruin of whole nations with their offspring), and ought never to be put into a creed or confession of the Church, but should be left to the theology of the school. Hence it is wisely omitted by the Heidelberg Catechism, the Helvetic Confessions, the Thirty-nine Articles, and other Reformed symbols. Even the old Scotch Confession of John Knox does not mention it, and the Second Scotch Confession expressly rejects, as an antichristian error, the horrible popish doctrine of the damnation of unbaptized infants.

The Westminster Confession, it is true, carefully avoids the term reprobation, and substitutes for it the milder idea of preterition. It uses the verb \emph{predestinate} only with reference to eternal life, while the lost are spoken of as being \emph{ordained} or judicially condemned to death. Yet it makes the dogmatic assertion that 'God was pleased, according to the unsearchable counsel of his own will, whereby he extendeth or withholdeth mercy as he pleaseth, for the glory of his sovereign power over his creatures, to \emph{pass by} the rest of mankind, and to \emph{ordain} them to dishonor and wrath for their sin, to the praise of his glorious justice.'\footnote{Ch. III. 7. This seventh section is the one dark spot in the Confession, and mars its beauty and usefulness. Comp. Larger Catechism, Quest. 13: 'God hath passed by and foreordained the rest to dishonor and wrath to be for their own sin inflicted, to the praise of the glory of his justice.' The Shorter Catechism (Quest. 7) wisely omits the negative part of predestination.} Now there are indeed passages in the Old and New Testaments, especially the ninth chapter of the Epistle to the Romans, which seem to bear out this statement,\footnote{Matt. xi. 25 ('Thou hast hid these things,' etc.); Rom. ix. 17, 18, 21, 22; 2 Tim. ii. 20; Jude 4; 1 Pet. ii. 8—all quoted in the Confession. The ninth chapter of Romans is the exegetical bulwark of the doctrine of reprobation; but it must be explained in connection with the tenth chapter, which brings out the unbelief of the creature as the cause, and with the eleventh chapter, which opens the prospect of a glorious solution of the problem in the conversion of the fullness of the Gentiles and the people of Israel, and ends with the grand declaration that 'God hath shut up all unto disobedience, that \emph{he might have mercy upon all.}' We have no more right to limit the \emph{all} in the second clause than in the first. Comp. the parallelism in Rom. v. 12 sqq.} but they must be interpreted in the light of the biblical idea of a God of infinite love and mercy, and in connection with other passages which in their obvious and natural sense declare that God sincerely desires all men to repent and be saved, that Christ is the Saviour of the world, that he is the propitiation not only for our sins, 'but also for the sins of the whole world,' and that he condemns no one absolutely and finally except for unbelief—that is, for the willful rejection of the gospel salvation.\footnote{John i. 29; iii. 16; iv. 24; 1 John ii. 2; iii. 8, 16; iv. 14; 1 Tim. ii. 4; Titus ii. 11; 2 Pet. iii. 9; Mark xvi. 16.} This fundamental doctrine of God's universal love and abundant provision for the salvation of all mankind should be put into a confession of faith rather than the doctrine of reprobation or preterition, which is, to say the least, as objectionable in such a document as the damning clauses in the Athanasian Creed.

The exegetical and theological adjustment of this whole subject of predestination, and of the unequal distribution and partial withholding of the favors of Providence and the means of grace in this world, is involved in insurmountable difficulties, and the contemplation of it should, make us cautious and charitable. A few general remarks may tend to set the problem in its true light, and to open the prospect of at least a partial solution.\footnote{Comp. our remarks, pp. 451 sqq.}

It must in fairness be admitted that the Calvinistic system only traces undeniable facts to their first ante-mundane cause in the inscrutable counsel of God. It draws the legitimate logical conclusions from such anthropological and eschatological premises as are acknowledged by all other orthodox Churches, Greek, Roman, Lutheran, and Reformed. They all teach the condemnation of the human race in consequence of Adam's fall, and confine the opportunity and possibility of salvation from sin and perdition to this present life,\footnote{The Roman Catholic doctrine, of purgatory is no exception, for this is confined to members of the Catholic Church who were converted in this life but need further purification before they can enter heaven. The Roman creed is more pronounced than the Greek and the Protestant on the impossibility of salvation outside of the visible Church on earth.} And yet every body must admit that the vast majority of mankind, no worse by nature than the rest, and without personal guilt, are born and grow up in heathen darkness, out of the reach of the means of grace, and are thus, as far as we know, actually 'passed by' in this world. No orthodox system can logically reconcile this stubborn and awful fact with the universal love and impartial justice of God.

The only solution seems to lie either in the Quaker doctrine of universal light—that is, an uncovenanted offer of salvation to all men in this earthly life—or in an extension of the period of saving grace beyond death till the final judgment for those (and for those only) who never had an opportunity in this world to accept or to reject the gospel salvation. But the former view implies a depreciation of the visible Church, the ministry of the gospel, and the sacraments; the latter would require a liberal reconstruction of the traditional doctrine of the middle state such as no orthodox Church, in the absence of clear Scripture light on this mysterious subject, and in view of probable abuse, would be willing to admit in its confessional teaching, even if theological exegesis should be able to produce a better agreement than now exists on certain disputed passages of the New Testament and the doctrine of Hades.

So far, then, the only difference is that, while the other orthodox Confessions conceal the real difficulty, Calvinism reveals it, and thus brings it nearer to a solution.

Moreover, the Calvinistic system, by detaching election from the absolute necessity of water-baptism, has a positive advantage over the Augustinian system, and is really more liberal. All the creeds which teach baptismal regeneration as an indispensable prerequisite of salvation virtually exclude the overwhelming majority of mankind—whole nations, with untold millions of infants dying in infancy—from the kingdom of heaven, whether they expressly say so or not. The Christian heart of the great African father shrunk from this fearful but inevitable conclusion of his logical head, and tried to mitigate it by making a distinction between positive damnation or actual suffering, and negative damnation or absence of bliss, and by subjecting unbaptized infants to the latter only. And this is the doctrine of Roman Catholic divines. The Calvinistic theory affords a more substantial relief, and allows, after the precedent of Zwingli and Bullinger, and in accordance with the analogy of Melchisedek, Job, and other exceptional cases of true piety under the Jewish dispensation, an indefinite extension of God's saving grace beyond the limits of the visible Church and the ordinary means of grace. It leaves room for the charitable hope of the salvation of all infants dying in infancy, and of those adults who, without an historical knowledge of Christ, live up to the light of nature and Providence, and die with a humble and penitent longing after salvation—that is, in a frame of mind like that of Cornelius when he sent for St. Peter.\footnote{See above, p. 378.} This was, indeed, not the professed Calvinism of Calvin and Beza, nor of the divines of Dort and Westminster, nor of the older divines of New England;\footnote{The Rev. Michael Wigglesworth, of Malden, Mass., a graduate and tutor of Harvard College (d. 1705), published a popular poem, \emph{The Day of Doom} (1662; 6th ed. 1715; reprinted as a curiosity by the Amer. News Company, New York, 1867), in which God reasons on the judgment-day with reprobate infants, who 'from the womb unto the tomb were straightway carried,' about the justice of their eternal damnation; and in consideration of their lesser guilt, assigns them (like St. Augustine) 'the easiest room in hell!'} but it is consistent with the Calvinistic scheme, which never presumed to fix the limits of divine election, and with a liberal interpretation of the Westminster Confession, which expressly acknowledges that elect infants and elect adults are regenerated and saved by Christ without being outwardly called by the gospel.\footnote{Ch. X. 3: 'Elect infants dying in infancy are regenerated and saved by Christ through the Spirit, who worketh when and where and how he pleaseth. So are all other elect persons who are incapable of being outwardly called by the ministry of the Word.' The Confession nowhere speaks of reprobate infants, and the existence of such is not \emph{necessarily} implied by way of distinction, although it \emph{probably} was in the minds of the framers as their private opinion, which they wisely withheld from the Confession. I think the interpretation of Dr. A. A. Hodge, of Allegheny, in his Commentary on this section (p. 240), is fairly admissible: 'The Confession affirms what is certainly revealed, and leaves that which revelation has not decided to remain without the suggestion of a positive opinion upon one side or the other.'. He agrees, as to the salvation of \emph{all} infants dying in infancy, with his father, who asserts that 'he never saw a Calvinistic theologian who held the doctrine of infant damnation in any sense' (\emph{System. Theol.}, Vol. III. p. 605).}

Modern Calvinism, at least in America, has decidedly taken a liberal view of this subject, and freely admits at least the probability of the universal salvation of infants, and hence the salvation of the greater part of the human race. Christianity can not be a failure in any sense—it must be a triumphant success, which is guaranteed from eternity by the infinite goodness and wisdom of God.\footnote{Dr. Hodge, of Princeton, is of the opinion, which would be preposterous in the Augustinian and Roman Catholic system, that the number of those who are ultimately lost is 'very inconsiderable as compared with the whole number of the saved.' This is the closing sentence of his \emph{System. Theol.}, Vol. III. p. 879. That the number of the saved will far exceed the number of the lost may be fairly inferred from the \textgreek{πολλῷ μᾶλλον} of Paul (Rom. v. 15, 17); but this inference can not well be harmonized with the declaration of our Lord, Matt. vii. 14, that but few enter the strait gate, unless we assume the universal salvation of infants, and look forward to great progress of the gospel in the future.}

But whatever may be the theoretical solution of this deep and dark mystery, there is a practical platform on which evangelical Christians can agree, namely, that all men who are and will be saved are saved by the free grace of God, without any merit of their own (faith itself being a gift of grace); while all who are lost are lost by their own guilt. It has often been said that pious Calvinists preach like Arminians, and pious Arminians pray like Calvinists. In this both may be inconsistent, but it is a happy and a useful inconsistency. The Calvinistic Whitefield was as zealous and successful in converting souls as the Arminian Wesley, and Wesley was as fervent and prevailing in prayer as Whitefield. They parted in this world, but they have long since been reconciled in heaven, where they see the whole truth face to face. We must work as if all depended on our efforts, and we must pray as if all depended on God. This is the holy paradox of St. Paul, who exhorts the Philippians to work out their own salvation with fear and trembling, for the very reason that it is God who worketh effectively in them both to will and to work of his own good pleasure. God's work in us and for us is the basis and encouragement of our work in him and for him.

INTOLERANCE.\footnote{\par On the subject of Toleration and Persecution, with special reference to England, the reader may profitably consult a series of \emph{Tracts on Liberty of Conscience and Persecution}, 1614–1661, edited by Edward B. Underhill for the Hansard Knollys Society, London, 1846; W. E. H. Lecky, \emph{History of Rationalism in Europe.} (4th edition, London, 1870; New York edition, 1875, in 2 vols.), ch. iv.; Masson, \emph{Life of Milton}, Vol. III. pp. 87 sqq., 383 sqq.; Stoughton, \emph{The Church of the Revolution} (London, 1874), ch. iv. pp. 114 sqq.; and Marshall's book quoted on p. 754.}

The principle of intolerance has been charged upon Chaps. XXIII. (Of the Civil Magistrate), XXX. (Of Church Censures), XXXI. (Of Synods and Councils), and the last clause of Ch. XX. (Of Christian Liberty, viz., the words 'and by the power of the civil magistrate'). The same charge applies to a few words in the 109th question of the Larger Catechism, where 'tolerating a false religion' is included among the sins forbidden in the Second Commandment with reference to some passages of the Old Testament and of the Book of Revelation (ii. 2, 16, 20; xvii. 16, 17).

There is no doubt that these passages assume a professedly Christian government, or the union of Church and State as it had come to be established in all Christian countries since the days of Constantine, and as it was acknowledged at that time by Protestants as well as Roman Catholics.\footnote{The first dissenting voices came from Anabaptists and Socinians, and from Castellio, who had nothing to gain and every thing to lose from the existing alliance of government and religion.} It is on this ground that the Confession claims for the civil magistrate (of whatever form of government) the right and duty not only legally to protect, but also to support the Christian Church, and to \emph{prohibit} or \emph{punish} heresy, idolatry, and blasphemy.

The power to coerce and punish implies the \emph{principle} of intolerance and the \emph{right} of persecution in some form or other, though this right may never be exercised. For just as far as a civil government is identified with a particular Church, an offense against that Church becomes an offense against the State, and subject to its penal code. All acts of uniformity in religion are necessarily exclusive, and must prohibit the public manifestations of dissent, whatever may be the private thoughts and sentiments, which no human government can reach.

It is a fact, moreover, that the Westminster Assembly was called for the purpose of legislating for the faith, government, and worship of three kingdoms, and that by adopting the Solemn League and Covenant it was pledged for the extirpation of popery and prelacy and all heresy.\footnote{And yet, in the face of this fact and the whole history of the seventeenth century, Dr. Hetherington (in his Introduction to Shaw's \emph{Exposition of the Confession of Faith}, pp. xxviii.) broadly denies any taint of intolerance in the Confession.}

The few Independents demanded a limited toleration, and were backed by Cromwell and his army, which was full of Independents, Baptists, Antinomians, Socinians, New Lights, Familists, Millenarians, and other 'proud, self-conceited, hot-headed sectaries' (as Baxter calls them). All these sectaries, who sprung up during the great religious excitement of the age, but mostly subsided soon afterwards, were of course tolerationists in their own interest. But for this very reason the prevailing sentiment in the Assembly was stoutly opposed to toleration, as the great Diana of the Independents and supposed mother and nurse of all sorts of heresies and blasphemies threatening the overthrow of religion and society.\footnote{Thomas Edwards, a zealous Presbyterian minister at London, published in 1645 a treatise of 60 pages, dedicated to Parliament, under the title, \emph{Gangræna; or, a Catalogue and Discovery of many of the Errors, Heresies, Blasphemies, and Pernicious Practices of the Sectaries of this Time}, in which he collects no less than one hundred and seventy-six miscellaneous 'errors, heresies, and blasphemies,' and enumerates sixteen heretical sects—namely: 1, Independents; 2, Brownists; 3, Millenaries; 4, Antinomians; 5, Anabaptists; 6, Arminians; 7, Libertines; 8, Familists; 9, Enthusiasts; 10, Seekers: 11, Perfectists: 12, Socinians; 13, Arians; 14, Antitrinitarians; 15, Antiscripturists; 16, Skeptics. 'The industrious writer,' says Neal, 'might have enlarged his catalogue with Papists, Prelatists, Deists, Ranters, Behemenists, etc., etc., or, if he had pleased, a less number might have served his turn, for very few of these sectaries were collected into societies; but his business was to blacken the adversaries of Presbyterian uniformity, that the Parliament might crush them by sanguinary methods.' See an account of this book in Neal, Part III. ch. vii. (Vol. II. p. 37), and Masson, Vol. III. pp. 143 sqq.} The Scottish delegation was a unit on the subject, and Baillie wrote a \emph{Dissuasive from the Errors of the Time} (1645) against toleration, and attacked it in his \emph{Letters.}\footnote{Innes (\emph{Law of Creeds}, pp. 243 and 244) says: 'Toleration was long unknown in the law, as in the history, of Scotland. The intense sentiment of national unity was strongly against it. The nation was one, and the Church became one. The Church claimed to be the Church of Christ in the realm, exclusively and of divine right. . . . The Scottish commissioners went to the Westminster Assembly to work out the ``covenanted uniformity in religion,'' and the new doctrine of the ``toleration of sects'' which met them there they most earnestly resisted.} Innumerable pamphlets were published on both sides. The advocates of toleration were defeated, and could only exact from the Assembly the important declaration that God alone is Lord of the conscience.

And yet, if we judge the Westminster standards from the stand point of the seventeenth century, and compare them with similar documents, they must be pronounced moderate.

1. They go no further on the subject of intolerance than the Belgic Confession,\footnote{Art. 36. See Vol. III. p. 432.} the Gallican Confession,\footnote{Art. 39. See Vol. III. p. 372.} the English Articles,\footnote{\par Art. 37. See Vol. III. p. 512.} and the Irish Articles.\footnote{No. 70. See Vol. III. p. 540.} They teach less than is implied in the Anglican doctrine of the royal supremacy, which puts the religion of a whole nation in the hands of the temporal sovereign, and which was employed for the severest measures against all dissenters, Roman Catholic and Protestant.

2. The Presbyterians, during the fifteen years of their domination,\footnote{We exempt the five years of Cromwell's Protectorate (1653–1658), during which the Independents were in the ascendency.} used their power very moderately, with the exception of a wholesale ejectment of a large number of prelatists from office (allowing them, however, one fifth of their income). This was a folly and a crime (viewed from our standpoint), but not nearly as cruel as the hanging and burning, the imprisonment, torture, and mutilation so freely exercised against themselves and other non-conformists before 1640 and after 1661. During the disgraceful period of the Restoration, which they unwisely brought about without exacting any pledges from the faithless Stuart, they suffered for their loyalty to the Westminster standards as much hardship and displayed as much heroism, both in England and Scotland, as any Church or sect in Christendom ever did.\footnote{A recent able writer, who has no sympathy whatever with the faith of Presbyterians, thus describes their persecutions under the Stuarts: 'In Scotland, during almost the whole period that the Stuarts were on the throne of England, a persecution rivaling in atrocity almost any on record was directed by the English government, at the instigation of the Scotch bishops, and with the approbation of the English Church, against all who repudiated episcopacy. If a conventicle was held in a house, the preacher was liable to be put to death. If it was held in the open air, both minister and people incurred the same fate. The Presbyterians were hunted like criminals over the mountains; their ears were torn from the roots; they were branded with hot irons; their fingers were wrenched asunder by the thumbkins; the bones of their legs were shattered in the boots; women were scourged publicly through the streets; multitudes were transported to the Barbadoes; an infuriated soldiery was let loose upon them, and encouraged to exercise all their ingenuity in torturing them.' (Lecky, l.c. Vol. II. p. 48, Amer. ed.)}

3. The Confession expresses for the first time among the confessions of faith, whether consistently or not, the true \emph{principle} of religious liberty, which was made the basis of the Act of Toleration, in the noble sentiment of Ch. XX. 2: '\emph{God alone is Lord of the conscience} (James iv. 12; Rom. xiv. 4), \emph{and hath left it free from the doctrines and commandments of men}, which are in any thing contrary to his Word, or beside it, in matters of faith or worship (Acts iv. 19; v. 29; 1 Cor. vii. 23; Matt. xxiii. 8-10; xxv. 9; 2 Cor. 1, 24). So that to believe such doctrines or to obey such commandments out of conscience is to betray true liberty of conscience; and the requiring of an implicit faith, and an absolute and blind obedience, is to destroy liberty of conscience, and reason also' (Isa. viii. 20; Acts xvii. 11).

4. The objectionable clauses in the Confession and Larger Catechism have been mildly interpreted and so modified by the Presbyterian Churches in Europe as to disclaim persecuting sentiments.\footnote{The Established Church of Scotland, the Original Secession Church, the English Presbyterian Church, and the Irish Presbyterian Church adhere to the 'whole doctrine' of the Westminster Confession, with a slight qualification of Ch. XXXI. 2. The Reformed Presbyterian Church does the same, but declares in its Testimony that it is 'not pledged to defend every sentiment or expression,' and asserts that 'to employ civil coercion of any kind for the purpose of inducing men to renounce an erroneous creed, or to espouse and profess a sound Scriptural one, is incompatible with the nature of true religion, and must ever prove ineffectual in practice.' The United Presbyterian Church introduces into its Formula of subscription this clause: 'It being understood that you are not required to approve of any thing in these documents which teaches, or is supposed to teach, compulsory or persecuting and intolerant, principles in religion.' The Free Church of Scotland meets the difficulty by a questionable exegesis, declaring (in an 'Act anent Questions and Formula,' June 1, 1846): ' The General Assembly, in passing this Act, think it right to declare that, while the Church firmly maintains the same Scriptural principles as to the duties of nations and their rulers in reference to true religion and the Church of Christ, for which she has hitherto contended, she disclaims intolerant or persecuting principles, and does not regard her Confession of Faith, or any portion thereof, when fairly interpreted, as favoring intolerance or persecution, or consider that her office-bearers, by subscribing it, profess any principles inconsistent with liberty of conscience and the right of private judgment.' See Innes, \emph{The Law of Creeds}, pp. 453, 461, 463.} The Presbyterian Churches in the United States have taken the more frank and effective course of an entire reconstruction of those chapters, so as to make them expressly teach the principle of religious freedom, and claim no favor from the civil magistrate but that protection which it owes to the lives, liberties, and constitutional rights of all its citizens.\footnote{See next section.}

GENERAL REMARKS ON THE PROGRESS OF RELIGIOUS LIBERTY.

The question in the sixteenth and seventeenth centuries was about toleration and persecution. But religious freedom requires much more, and is now regarded as one of the fundamental and most precious rights of men, which must be sacredly protected in its public exercise by the civil government, within the limits of order, peace, and public morals. This liberty is the final result and gain of ages of intolerance and persecution.

The history of religious persecution is the darkest chapter in Church history—we may call it the devil's chapter—and the darkest part in it is the persecution of Christians by Christians. It is, however, relieved by the counter-manifestation of the heroic virtues of Christian martyrdom and the slow but steady progress of liberty through streams of martyr blood.

All Christian Churches, except a few denominations of recent date which never had a chance, have more or less persecuted when in power, and must plead guilty. The difference is only one of degree. The Episcopalians were less intolerant than the Roman Catholics, the Presbyterians less intolerant than the Episcopalians, the Independents less intolerant (in theory) than the Presbyterians. But they were all intolerant. Even the Independents of Old England, with the great Cromwell and the great Milton as their leaders, excluded Romanists, Prelatists (i.e., Episcopalians), and Unitarians from their programme of toleration,\footnote{Milton, the independent of Independents and the boldest as well as most eloquent champion of civil and religious liberty in the seventeenth century, was unwilling to tolerate Romanists, because he regarded them as idolaters and as enemies of freedom. See his \emph{Areopagitica}, of which Lecky (Vol. II. p. 80) says that it is as glorious a monument of the genius of Milton as his \emph{Paradise Lost}, and that it 'probably represents the very highest point that English eloquence has attained.'} and, strange to say, when in power in New England, they expelled Baptists and hanged Quakers on the virgin soil of Massachusetts before and after the Westminster Assembly. On the other hand, however, there is not a Christian Church or sect that has not complained of intolerance and injustice under persecution, and that has not furnished some bold advocates of toleration and freedom, from Tertullian and Lactantius down to Roger Williams and William Penn. This is the redeeming feature in this fearful picture, and must not be overlooked in making up a just estimate.

It is therefore the greatest possible injustice to charge the persecutions to Christianity, which breathes the very opposite spirit of forbearance, forgiveness, love, and liberality; which teaches us to suffer wrong rather than to inflict wrong; and which, by restoring the divine image in man, and lifting him up to the sphere of spiritual freedom, is really the pure source of all that is truly valuable in our modern ideas of civil and religious liberty. Whatever may be said of the severity of the Mosaic legislation, which assumes the union of the civil and ecclesiastical power, Christ and the Apostles, both by precept and example, strictly prohibit the use of carnal means for the promotion of the kingdom of heaven, which is spiritual in its origin, character, and aim. The reminiscence of this spirit lingered in the Church through the darkest ages in the maxim \emph{Ecclesia non sitit sanguinem.}

It is also wrong to derive intolerance from the strength and intensity of religious conviction—although this undoubtedly may come in as an additional stimulus—and to trace toleration to skepticism and unbelief.\footnote{This is the theory of Lecky.} For who had stronger convictions than St. Paul? His Jewish conviction or pharisaical fanaticism made him a bitter persecutor, but his Christian conviction inspired his seraphic description of love (1 Cor. xiii.) and strengthened him for martyrdom. On the other hand, the Deist philosopher, Hobbes, by giving the civil power an absolute right to determine the religion of a nation, taught the extreme doctrine of persecution; and the reign of terror in France proves that infidelity may be as fanatical and intolerant as the strongest faith, and may instigate the most horrible of persecutions.

Intolerance is rooted in the selfishness and ambition of human nature and in the spirit of sectarian exclusiveness, which assumes that we and the sect to which we belong have the monopoly of truth and orthodoxy, and that all who dissent from us must be in error. Persecution follows as a legitimate consequence of this selfishness and bigotry wherever the intolerant party has the power to persecute.

The Roman Church, wherever she controls the civil government, can not consistently tolerate, much less legally recognize, any form of worship besides her own, because she identifies herself with the infallible Church of Christ, out of which there is no salvation, and regards all who dissent from her as damnable schismatics and heretics.\footnote{The limited toleration in some Roman Catholic countries exists in spite of Romanism, and the liberal opinions and Christian feelings of individual Catholics have no influence on the system, which is the same as ever, as may be inferred from the papal Syllabus of 1864, and from the recent papal protest against even the minimum of religious toleration in Spain (1876). In Protestant countries the Roman Church claims as much liberty as she can get, and advocates toleration in her own interest, but would deny it to others as soon as she attained to power.} Protestants, who began with the assertion of private judgment against the authority of Rome, and complained bitterly of her persecuting spirit, are inconsistent and more inexcusable if they refuse the same right to others and persecute them for its exercise. For a long time, however, Protestantism clung to the traditional idea of uniformity in religion, and this was the source of untold suffering, especially in England, until it became manifest beyond a doubt that doctrinal and ceremonial uniformity was an impossibility in a nation of intelligent freemen. The Toleration Act of May 24, 1689, for the relief of Dissenters, marks the transition. Since that time religious persecution by the civil power has ceased in the Anglo-Saxon race, and the principle of religious liberty has gradually become a settled conviction of the most advanced sections of the Christian world.

For this change of public sentiment the chief merit is due to the English Non-conformists, who in the school of persecution became advocates of toleration, especially to the Baptists and Quakers, who made religious liberty (within the limits of the golden rule) an article of their creed, so that they could not consistently persecute even if they should ever have a chance to do so.\footnote{See the 'Fourteenth Proposition' of Barclay, adopted by the Quakers: 'Since God hath assumed to himself the power and dominion of the conscience, who alone can rightly instruct and govern it, therefore it is not lawful for any whatsoever, by virtue of any authority or principality they bear in the government of this world, to force the consciences of others; and therefore all killing, banishing, fining, imprisoning, and other such things, which men are afflicted with, for the alone exercise of their conscience, or difference in worship or opinion, proceedeth from the spirit of Cain, the murderer, and is contrary to the truth; provided always that no man, under the pretense of conscience, prejudice his neighbor in his life or estate, or do any thing destructive to, or inconsistent with, human society; in which case the law is for the transgressor, and justice to be administered upon all, without respect of persons.' This was published in 1675. Bossuet, therefore, was imperfectly informed when at the close of the seventeenth century (1688) he mentioned the Anabaptists and Socinians as the \emph{only} Christians who did not admit the power of the civil sword '\emph{dans les matières de la religion et de la conscience}' (\emph{Hist. des Variations}, LIV. X. 56).} It was next promoted by the eloquent advocacy of toleration in the writings of Chillingworth,\footnote{\emph{The Religion of Protestants a Safe Way to Salvation}, 1637 (or 1638; dedicated in a most humble preface to King Charles I.; 3d ed. 1664; 10th ed. 1742; reprinted in the first two vols. of the Oxford ed. of Chillingworth's \emph{Works}, 1838, in 3 vols.). This book is a vindication of Protestantism and of the author's return to it, and proclaims that the Bible, the whole Bible, and nothing but the Bible, is the religion of Protestants, and that no Church of one denomination is infallible. At Chillingworth's burial, in Jan., 1644, Dr. Cheynell, who had shown him great kindness during his sickness, flung this book into the grave, with the words, 'Get thee gone, thou cursed book; go rot with thy author.' Chillingworth, however, had no idea of civil liberty, and wrote as an extreme royalist on the \emph{Unlawfulness of Resisting the Lawful Prince, although most Impious, Tyrannical, and Idolatrous.}} Jeremy Taylor,\footnote{\emph{Liberty of Prophesying}, written in exile (1647), and unfortunately retracted in part after the Restoration by the author himself, who declared it to have been a \emph{ruse de guerre.} Coleridge regards this weakness as almost the only stain on Taylor's character.} and other Anglican divines of the latitudinarian school; further, by the mingling of creeds and sects in the same country where persecution failed of its aim; and, lastly, by the skeptical philosophy and the religious indifferentism of the eighteenth century, which, however, has repeatedly shown itself most intolerant of all forms of positive belief, and can therefore be no more trusted than the bigotry of superstition. Religious freedom is best guaranteed by an enlightened Christian civilization, a liberal culture, a large-hearted Christian charity, a comprehensive view of truth, a free social intercourse of various denominations, and a wise separation of civil and ecclesiastical government.

During the last stages of the age of persecution Providence began to prepare in the colonies of North America the widest field and the proper social basis for the full exercise of religious liberty and equality by bringing together under one government the persecuted of all nations and sects, so that the enjoyment of the liberty of each depends upon and is guaranteed by the recognition and protection of the liberty of all the rest.

\xsubsection{The Westminster Standards in America.}

§ 98. The Westminster Standards in America.

With the Puritan emigration from England and the Presbyterian emigration from Scotland and the North of Ireland, the Westminster standards were planted on the virgin soil of America long before the Declaration of Independence. The most popular is the Shorter Catechism, which has undergone no change except a very slight one among the Cumberland Presbyterians.\footnote{See next section.}

THE CONGREGATIONAL CHURCHES OF NEW ENGLAND.

The Confession of Faith was first adopted 'for substance of doctrine,' but without the principles of Presbyterian discipline, by the Congregational Synod of Cambridge, in the Colony of Massachusetts, A.D. 1648, one year after its issue in England; then, in the Savoy recension, by the Synod of Boston, Mass., May 12, 1680; and again, in the same form, by the Congregational churches of Connecticut at a Synod of Saybrook, Sept. 9, 1708.

The Smaller Catechism was formerly used as a school-book in New England, but has been thrust into the background by the modern prejudice against catechisms and by a flood of more entertaining but less solid Sunday-school literature.

THE PRESBYTERIAN CHURCHES.

The various Presbyterian bodies of English and Scotch descent used at first all the Westminster standards without alteration. The Presbytery of Philadelphia, the oldest in America, was organized in 1706, the Synod of Philadelphia in 1717, and the Synod of New York in 1743. The Synod of Philadelphia, Sept. 19, 1729, adopted the Confession with a liberal construction, in these words:

'Although the Synod do not claim or pretend to any authority of imposing our faith upon other men's consciences, but do profess our just dissatisfaction with and abhorrence of such impositions, and do utterly disclaim all legislative power and authority in the Church, being willing to receive one another as Christ has received us to the glory of God, and admit to fellowship in sacred ordinances all such as we have grounds to believe Christ will at last admit to the kingdom of heaven: yet we are undoubtedly obliged to take care that the faith once delivered to the saints be kept pure and uncorrupt among us, and so handed down to our posterity.

'\emph{And} [we] \emph{do therefore agree that all the ministers of this Synod, or that shall hereafter be admitted to this Synod, shall declare their agreement in and approbation of the Confession of Faith, with the Larger and Shorter Catechisms of the Assembly of Divines at Westminster, as being, in all the essential and necessary articles, good forms of sound words and systems of Christian doctrine, and do also adopt the said Confession and Catechisms as the confession of our faith.}

'And we do also agree that all the Presbyteries within our bounds shall always take care not to admit any candidate of the ministry into the exercise of the sacred function but what declares his agreement in opinion with all the essential and necessary articles of said Confession, either by subscribing the said Confession of Faith and Catechisms, or by a verbal declaration of his assent thereto, as such minister or candidate shall think best. And in case any minister of this Synod, or any candidate for the ministry, shall have any scruple with respect to any article or articles of said Confession or Catechisms, he shall, at the time of his making said declaration, declare his sentiments to the Presbytery or Synod, who shall, notwithstanding, admit him to the exercise of the ministry within our bounds, and to ministerial communion, if the Synod or Presbytery shall judge his scruple or mistake to be only about articles not essential and necessary in doctrine, worship, or government. But if the Synod or Presbytery shall judge such ministers or candidates erroneous in essential and necessary articles of faith, the Synod or Presbytery shall declare them incapable of communion with them. And the Synod do solemnly agree that none of us will traduce or use any opprobrious terms of those that differ from us in these extra-essential and not-necessary points of doctrine, but treat them with the same friendship, kindness, and brotherly love as if they had not differed from us in such sentiments.'\footnote{Minutes of the Synod of Philadelphia, as published in the \emph{Records of the Presbyterian Church in the United States of America} (embracing the Minutes of the Presbytery of Philadelphia, and of the Synods of New York and Philadelphia, from 1706 to 1788). Philad. Presbyt. Board of Public. 1841, p. 92. See also W. E. Moore's  \emph{Presbyterian Digest: a Compend of the Acts and Deliverances of the General Assembly of the Presbyterian Church in the United States of America} (Philad. Presbyt. Board), second ed. 1873, pp.45 sq.}

In the afternoon session the scruples about adopting these standards were solved, and the Synod unanimously declared that they do not receive 'some clauses in the twentieth and twenty-third chapters in any such sense as to suppose the civil magistrate hath a controlling power over Synods with respect to the exercise of their ministerial authority, or power \emph{to persecute any for their religion}, or in any sense contrary to the Protestant succession to the throne of Great Britain.'

This supplementary action foreshadows the changes which were afterwards made.

When the Synods of Philadelphia and New York united in one body at Philadelphia, May 29, 1758, they adopted, as the first article of the plan of union, the following:

'Both Synods having always approved and received the Westminster Confession of Faith and Larger and Shorter Catechisms, as an orthodox and excellent system of Christian doctrine, founded on the Word of God, we do still receive the same as the confession of out faith; and also adhere to the plan of worship, government, and discipline contained in the Westminster Directory, strictly enjoining it on all our members and probationers for the ministry, that they preach and teach according to the form of sound words in said Confession and Catechisms, and avoid and oppose all errors contrary thereto.'\footnote{See Minutes of the Synod of 1758 as published in the \emph{Records of the Presbyterian Church}, p. 286. Also Moore's \emph{Digest}, p. 48; and Gillett, \emph{Hist. of the Presbyt. Ch. in the U. S. of America}, Vol. I. p. 138.}

THE AMERICAN REVISION.

After the Revolutionary War the united Synod of Philadelphia and New York, which met at Philadelphia, May 28, 1787, appointed a committee to prepare an alteration in the Confession of Faith, Ch. XX. (closing paragraph), Ch. XXIII., 3, and Ch. XXXI., 1, 2, in consequence of the new relation of Church and State.\footnote{See \emph{Records of the Presbyterian Church}, p. 539, where we find the following minute, dated May 28, 1787: 'The Synod took into consideration the last paragraph of the twentieth chapter of the Westminster Confession of Faith, the third paragraph of the twenty-third chapter, and the first paragraph of the thirty-first chapter; and having made some alterations, agreed that the said paragraphs, as now altered, be printed for consideration, together with the draught of a plan of government and discipline. The Synod also appointed the above-named committee to revise the Westminster Directory for public worship, and to have it, when thus revised, printed, together with the draught, for consideration. And the Synod agreed that when the above proposed alterations in the Confession of Faith shall have been finally determined on by the body, and the Directory shall have been revised as above directed, and adopted by the Synod, the said Confession thus altered, and Directory thus revised and adopted, shall be styled, ``The Confession of Faith and Directory for Public Worship of the Presbyterian Church in the United States of America.'''}

The changes proposed were adopted by the joint Synod at a subsequent meeting in Philadelphia, May 28, 1788, in the following action:

'The Synod having fully considered the draught of the form of government and discipline, did, on a review of the whole, and hereby do ratify and adopt the same, as now altered and amended, as the Constitution of the Presbyterian Church in America, and order the same to be considered and strictly observed as the rule of their proceedings by all the inferior judicatories belonging to the body. And they order that a correct copy be printed, and that the \emph{Westminster Confession of faith, as now altered}, be printed in full along with it, as making a part of the Constitution.

'\emph{Resolved}, That the true intent and meaning of the above ratification by the Synod is, that the Form of Government and Discipline, and the \emph{Confession of Faith}, as now ratified, is to continue to be our constitution and the confession of our faith and practice unalterable, unless two thirds of the Presbyteries under the care of the General Assembly shall propose alterations or amendments, and such alterations or amendments shall be agreed to and enacted by the General Assembly.'\footnote{\emph{Records of the Presbyterian Church}, p. 546; Moore's \emph{Digest}, p. 51.}

On the day following (May 29) the Synod 'took into consideration the Westminster Larger and Shorter Catechisms, and having made a small amendment of the Larger, did approve, and do hereby approve and ratify the said Catechisms, as now agreed on, as the Catechisms of the Presbyterian Church in the said United States.' At the same time it was ordered that all these standards, as altered and adapted to the wants of the American churches, be printed and bound up in one volume.\footnote{\emph{Records}, p. 547; Moore's \emph{Digest}, p. 52. The first edition of the new book appeared Philad. 1789, under the title: '\emph{The Constitution of the Presbyterian Church in the United States of America, containing the Confession of Faith, the Catechisms, the Government and Discipline, and the Directory of the Worship of God, ratified and adopted by the Synod of New York and Philadelphia, May} 28, 1788. The Assembly of 1792 ordered a new edition with the Scripture texts annexed, and appointed a committee for the purpose. This edition was adopted by the Assembly in 1794 (Moore's \emph{Digest}, p. 52).}

The changes consist in the omission of those sentences which imply the union of Church and State, or the principle of ecclesiastical establishments, making it the duty of the civil magistrate not only to protect, but also to support religion, and giving to the magistrate power to call and ratify ecclesiastical synods and councils, and to punish heretics. Instead of this, the American revision confines the duty of the civil magistrate to the legal protection of religion in its public exercise, without distinction of Christian creeds or organizations. It thus professes the principle of religious liberty and equality of all denominations before the law. This principle has been faithfully and consistently adhered to by the large body of the Presbyterian Church in America, and has become the common law of the land. To facilitate the comparison we present the respective sections in parallel columns:

Original Text.

American Text.

Ch. XXIII. 3.—Of the Civil Magistrate.

Ch. XXIII. 3.—Of the Civil Magistrate.

The civil magistrate, may not assume to himself the administration of the Word and Sacraments, or the power of the keys of the kingdom of heaven;\footnote{2 Chron. xxvi. 18; Matt. xviii. 17; Matt. xvi. 19; 1 Cor. xii. 28, 29; Eph. iv. 7, 12; 1 Cor. iv. 1, 2; Rom. x. 15; Heb. v. 4.} yet he hath authority, and it is his duty to take order, that unity and peace be preserved in the Church, that the truth of God be kept pure and entire, that all blasphemies and heresies be suppressed, all corruptions and abuses in worship and discipline prevented or reformed; and all the ordinances of God duly settled, administered,

Civil magistrates may not assume to themselves the administration of the Word and Sacraments;\footnote{2 Chron. xxvi. 18}or the power of the keys of the kingdom of heaven;\footnote{Matt. xvi. 19; 2 Cor. iv. 1, 2} or, in the least, interfere in matters of faith.\footnote{John xviii. 36; Mal. ii. 7; Acts v. 29.} Yet, as nursing fathers, it is the duty of civil magistrates to protect the Church of our common Lord, without giving the preference to any denomination of Christians above the rest, in such a manner that all ecclesiastical persons whatever shall enjoy the full, free, and unquestioned liberty

and observed.\footnote{Isa. xlix. 23; Psa. cxxii. 9; Ezra vii. 23–28; Lev. xxiv. 16; Deut. xiii. 5, 6, 12; 1 Kings xviii. 4; 1 Chron. xiii. 1–9; 2 Kings xxiii. 1–26; 2 Chron. xxxiv. 33; xv. 12, 13.} For the better effecting whereof he hath power to call synods, to be present at them, and to provide that whatsoever is transacted in them be according to the mind of God.\footnote{2 Chron. xv. 8–17; xxix. 30; Matt. ii. 4, 5.}

of discharging every part of their sacred functions without violence or danger.\footnote{Isa. xlix. 23.} And as Jesus Christ hath appointed a regular government and discipline in his Church, no law of any commonwealth should interfere with, let, or hinder the due exercise thereof among the voluntary members of any denomination of Christians, according to their own profession and belief.\footnote{Psa. cv. 15; Acts xviii. 14, 15, 16.} It is the duty of civil magistrates to protect the person and good name of all their people, in such an effectual manner as that no person be suffered, either upon pretense of religion or infidelity, to offer any indignity, violence, abuse, or injury to any other person whatsoever; and to take order that all religious and ecclesiastical assemblies be held without molestation or disturbance.\footnote{2 Sam. xxiii. 13; 1 Tim. ii. 1; Rom. xiii. 4.}

Ch. XXXI.—Of Synods and Councils.

Ch. XXXI.—Of Synods and Councils.

For the better government and further edification of the Church, there ought to be such assemblies as are commonly called synods or councils.\footnote{Acts xv. 2, 4, 6.}

For the better government and further edification of the Church, there ought to be such assemblies as are commonly called synods or councils.\footnote{Acts xv. 2, 4, 6.} And it belongeth to the overseers and other rulers of the particular churches, by virtue of their office, and the power which Christ hath given them for edification, and not for destruction, to appoint such assemblies; and to convene together in them, as often as they shall judge it expedient for the good of the Church.\footnote{Acts xv. 22, 23, 25.}

II. As magistrates may lawfully call a synod of ministers and other fit persons to consult and advise with about matters of religion:\footnote{Isa. xlix. 23; 1 Tim. ii. 1, 2; 2 Chron. xix. 8–12; xxix. and xxx.; Matt. ii. 4, 5; Prov. xi. 14.} so, if magistrates be open enemies to the Church, the ministers of Christ, of themselves, by virtue of their office; or they, with other fit persons, upon delegation from their churches, may meet together in such assemblies.\footnote{Acts xv. 2, 4, 22, 23, 25.}

In Ch. XX., § 4, the last sentence, 'and by the power of the civil magistrate,' was omitted, so as to read, 'they [the offenders] may lawfully be called to account, and proceeded against by the censures of the Church.'

The only change made in the Larger Catechism was the striking out of the words 'tolerating a false religion,' among the sins forbidden in the Second Commandment (Quest. 109).

The example set by the Presbyterian Church in the United States was afterwards (1801) followed by the Protestant Episcopal Church in the revision of the political sections of the Thirty-nine Articles of Religion.

PRESBYTERIAN REUNION.\footnote{For the documentary history of this remarkable movement, see the Minutes of the two General Assemblies for 1867–69, and of the reunited Assembly from 1870 to 1872; also the new edition of Moore's \emph{Presbyterian Digest} (1873), pp. 57–106; and the Memorial Volume on \emph{Presbyterian Reunion}, New York, 1870.}

The division of the Presbyterian Church into Old School and New School, which took place at Philadelphia, June 8, 1837, arose chiefly from contentions in consequence of the Plan of Union formed in 1801 between the General Assembly and the Congregational Association of Connecticut, and involved two different constructions of the doctrinal standards—the one more strict and conservative, the other more liberal and progressive—but did not affect the organic law of the Church.\footnote{For the documentary history of the separation of the Presbyterian Church and the 'Exscinding Acts' of the Old-School Assembly, see Baird's \emph{Collection} (O. S.), pp. 710 sqq., and the first edition of Moore's \emph{New Digest} (N. S.), pp. 456 sqq. In the new edition of Moore's \emph{Digest} (1873), the chapter on the division is omitted, and the documents on the reunion inserted instead.} The Old School, it is true, charged the New School with sixteen Pelagian and Arminian errors, which had their origin in recent developments of New England theology; but the New School met the charge with the 'Auburn Declaration' (Aug. 1837), which denied those errors and adopted sixteen 'true doctrines' in essential harmony with the Calvinistic anthropology and soteriology. This Declaration must be regarded as expressing the belief of the New-School body at that time, whatever the views of individual members may have been.\footnote{The sixteen errors charged are found in Baird's \emph{Collection}, pp. 711 and 745 sqq., together with the reply of the New School, which was afterwards, in Aug. of the same year, adopted by a convention of 98 commissioned ministers and 58 laymen (besides 24 corresponding members) at Auburn, N. Y., and is hence called the 'Auburn Declaration.' The latter is also embodied in the third volume of this work, p. 771. On its history, comp. Dr. Morris, in the \emph{Amer. Presbyt. Review}, for January, 1876.}

In the preparatory steps towards a reunion of these two bodies after a separation of thirty-two years, the question of the doctrinal basis took a prominent part. It was proposed that 'in the United Church the Westminster Confession of Faith shall be received and adopted as containing the system of doctrine taught in the Holy Scriptures.' It is characteristic of the excellent temper and spirit of concession which prevailed on both sides, that at the 'Presbyterian National Union Convention,' held in November, 1867, at Philadelphia, Dr. Henry B. Smith, of the Union Theological Seminary, New York, a prominent leader of the New School, proposed a defining clause, to satisfy the demands of Old School orthodoxy;\footnote{The 'Smith amendment' was in these words: 'It being understood that this Confession is received in its proper historical, that is, the Calvinistic or Reformed, sense.' This would exclude, of course, Antinomianism and Fatalism on the one hand, and Arminianism and Pelagianism on the other.} while the Rev. Dr. Gurley, pastor of an Old-School church in Washington City, proposed an additional clause to guarantee the New School liberty of interpretation.\footnote{The 'Gurley amendment' was in these words: 'It is also understood that various methods of viewing, stating, explaining, and illustrating the doctrines of the Confession, which do not impair the integrity of the Reformed or Calvinistic system, are to be freely allowed in the United Church, as they have hitherto been allowed in the separate Churches.'} The amendments were received unanimously, with great joy and gratitude.

But after further consideration it was found best to drop both these amendments, and when the reunion was consummated by the two assemblies at Pittsburgh, Pa., Nov. 10, 1869, the following article was unanimously adopted:

'The reunion shall be effected on the doctrinal and ecclesiastical basis of our common Standards; the Scriptures of the Old and New Testaments shall be acknowledged to be the inspired Word of God, and the only infallible rule of faith and practice; the Confession of Faith shall continue to be sincerely received and adopted, as containing the system of doctrine taught in the Holy Scriptures; and the government and discipline of the Presbyterian Church in the United States shall be approved as containing the principles and rules of our polity.'

Thus the Presbyterian Church in the United States of America, which had been unfortunately separated by a \emph{permissive} decree of God, was happily and, we trust, forever reunited by an \emph{efficient} and \emph{gracious} decree of God.\footnote{See the address of Dr. Masgrave at the meeting in Pittsburgh, \emph{Memorial Volume}, p. 388.}

OTHER PRESBYTERIAN CHURCHES IN THE UNITED STATES.

In addition to this large Presbyterian Church, there are in the United States a number of smaller ones having distinctively a Scottish origin. Of these and of their relation to the Westminster standards the Rev. G. D. Mathews, of New York, from his own familiar acquaintance with the Presbyterian Churches in Scotland and the United States, kindly furnishes for this work the following account:

'Among the emigrants into this country in the last century were many who had been connected with the Associate Church of Scotland. The fathers of that Church, the Erskines, objected not so much to the constitution of the Established Church as to its administration, especially in reference to patronage and to Church discipline. In 1753 the American \emph{Associate Church} was organized as a Presbytery subordinate to the Antiburgher Synod of Scotland, equalling if not surpassing the mother Church in its rigid adherence to the doctrinal system of the Westminster standards. Its zeal for these, indeed, served to deepen its opposition to the Scottish Establishment as a Church that had become unfaithful to its religious profession.

'In 1774 a \emph{Reformed Presbyterian Presbytery} was constituted in America by followers of Cargill, Cameron, and Renwick. These held that the Church of Scotland had marred its standing as a true Church of Christ by entering into union with an immoral government—the government of Great Britain being of this character because not based on Scriptural principles. Of this latter position the proof was alleged to lie in its disregard, as shown by the national acceptance of Episcopacy at the Restoration in 1660, and again at the Revolution in 1688, of that Solemn League and Covenant which had been sworn to in 1643, a Covenant whose engagements were affirmed to be binding on the people of the British Empire until fulfilled. An additional proof lay in the absence from its constitution of any acknowledgment of God as the Author of its existence and the source of its authority, of Jesus Christ as its Ruler, and of the Bible as the supreme law of its conduct.

'Notwithstanding some actual differences, the force of circumstances brought these Churches together, so that in 1782 they became united under the name of the \emph{Associate Reformed Church}—minorities on both sides refusing to enter the union, and thus perpetuating their respective Churches. In 1799 the Associate Reformed Church issued an edition of the Westminster Confession containing the following changes from the original documents:

Chap. XX. 4.— . . . faith, worship, conversation, (insert) \emph{or the order which Christ hath established in his Church, they may be lawfully called to account, and proceeded against by the censures of the Church; and in proportion as their erroneous opinions or practices, either in their own nature or in the manner of publishing or maintaining them, are destructive to the external peace of the Church and of civil society, they may also be proceeded against} by the power of the civil magistrate.

Chap. XXII. 3.— . . . the keys of the kingdom of heaven. (Add) \emph{Yet, as the gospel revelation lays indispensable obligations upon all classes of people who are favored with it, magistrates, as such, are bound to execute their respective offices in a subserviency thereto, administering government on Christian principles, and ruling in the fear of God, according to the directions of his Word; as those who shall give an account to the Lord Jesus, whom God hath appointed to be the Judge of the world.}

\emph{Hence magistrates, as such, in a Christian country are bound to promote the Christian religion, is the most valuable interest of their subjects, by all such means as are not inconsistent with civil rights, and do not imply an interference with the policy of the Church, which is the free and independent kingdom of the Redeemer, nor an assumption of dominion over conscience.}

Chap. XXXI. 2.—(Substitute.) \emph{The ministers of Christ themselves, and by virtue of their office; or they with other fit persons, upon delegation from their churches, have the exclusive right to appoint, adjourn, or dissolve such synods or councils; though in extraordinary cases it may be proper for magistrates to desire the calling of a synod of ministers and other fit persons, to consult and advise with about matters of religion; and in such cases it is the duly of churches to comply with their desire.}

'In the Larger Catechism, under the things forbidden by the Second Commandment, the word \emph{authorizing} was substituted for ``\emph{tolerating} a false religion.''

'In 1858 the \emph{Associate Church}, which had by this time grown considerably, joined with the \emph{Associate Reformed Church}, when the name \emph{United Presbyterian Church} was assumed and the Westminster Confession again altered. The edition used by this Church differs from the original in the following passages:

Chap. XX. 4.— . . . hath established in the Church, they (add) \emph{ought to be called to account, and proceeded against by the censures of the Church, if they belong to her communion, and thus be amenable to her own spiritual authority. And as the civil magistrate is the minister of God for good to the virtuous and a revenger to execute wrath upon him that doeth evil, he is therefore bound to suppress individuals and combinations, whatever may be their avowed objects, whether political or religious, whose principles and practices, openly propagated and maintained, are calculated to subvert the foundations of properly constituted society.}

Chap. XXIII. 3.— . . . kingdom of heaven, (add) \emph{or in the least interfere to regulate matter's of faith and worship. As nursing fathers, magistrates are bound to administer their government according to the revealed principles of Christianity, and to improve the opportunities which their high station and extensive influence afford in promoting the Christian religion as their own most valuable interest and the good of the people demand, by all such means as do not imply any infringement of the inherent rights of the Church, or any assumption of dominion over the consciences of men. They ought not to punish any as heretics or schismatics. No authoritative judgment concerning matters of religion is competent to them, as their authority extends only to the external works or practices of their subjects as citizens, and not as Christians. It is their duty to protect the Church in such a manner that all ecclesiastical persons shall enjoy the free, full, and unquestioned liberty of discharging every part of their sacred functions without violence or danger. They should enact no law which would in any way interfere with or hinder the due exercise of government and discipline established by Jesus Christ in his Church. It is their duty also to protect the person, good name, estate, natural and civil rights of all their subjects in such a way that no person be suffered, upon any pretense, to violate them; and to take order that all religious and ecclesiastical assemblies be held without molestation or disturbance. God alone being Lord of the conscience, the civil magistrate may not compel any under his civil authority to worship God contrary to the dictates of their own consciences; yet it is competent in him to restrain such opinions and to punish such practices as tend to subvert the foundations of civil society and violate the common rights of men.}

Chap. XXXI. 2.—(Substitute.) \emph{We declare that as the Church of Jesus Christ is a kingdom distinct from and independent of the state, having a government, laws, office-bearers, and all spiritual power peculiar to herself for her own edification; so it belongs exclusively to the ministers of Christ, together with other fit persons, upon delegation from their churches, by virtue of their office and the intrinsic power committed unto them, to appoint their own assemblies, and to convene together in them as often as they should judge it expedient for the good of the Church.}

'In the question of the Larger Catechism, changed in 1799, the original word \emph{tolerating} was restored.

'At no period has the \emph{Associate Church}, which still exists, altered the language of the Confession. It has refrained from doing this, ``judging it to be improper for one ecclesiastical body to alter any deed of another, making it rather express their own views than those of the body by whom it was originally framed, for hereby the sentiments of one body may be unfairly palmed upon another.'' Any obscurity or error in the Confession should be remedied by the emitting of a \emph{Testimony}, in which there could be given a full and accurate statement of the particular truth in question. In 1784, therefore, the Associate Church issued such a Testimony, in which (Articles 15–19), speaking of the civil magistrate, it affirmed that the magistrate, as such, is no ruler in the Church; that he should not grant any privileges to those whom he judges professors of the true religion which may hurt others in their natural rights; that his whole duty, as a magistrate, respects men, not as Christians, but as members of civil society; that any \emph{de facto} government governing orderly is that ordinance of God which must be obeyed, and that with any such government Christians may lawfully co-operate.

'\emph{The Reformed Presbyterian Church} has also retained the Westminster Confession unaltered. Adhering to its teaching on the Civil Magistrate, as this was received by the Church of Scotland in the Adopting Act of 1647, it issued in 1806 a Testimony, in which it declared that civil government is a natural institution, but that, to be a lawful one, so that a Christian man may take part in it, God must be acknowledged in its constitution as the fountain of all power and authority, and that Christian rulers, appointed to office according to a righteous civil constitution, have authority from God to rule, in subserviency to the kingdom of Christ. The absence from the American national constitution of any such acknowledgment renders that covenant unscriptural and immoral, and so precludes Christian men from becoming identified with its administration. Another reason for this political dissent is the doctrine of the binding obligation of the Scottish Covenants.

'A difference of opinion that had gradually risen within this Church as to the extent of this precluding led to the formation, in 1833, of the \emph{Synod of the Reformed Presbyterian Church} holding the extremest view of political dissent, and of the \emph{General Synod} of the same Church, permitting its members to exercise the political franchise.

'As regards the doctrinal articles of the Confession, all these Churches are \emph{Calvino Calviniores.}'

\xsubsection{The Westminster Standards in the Cumberland Presbyterian Church.}

§ 99. The Westminster Standards in the Cumberland Presbyterian Church.

Sources.

I. On the part of the Cumberland Presbyterian Church:

\emph{The Confession of Faith of the Cumberland Presbyterian Church in the United States of America. Revised and adopted by the General Assembly, at Princeton, Ky., May}, 1829. Nashville, Tennessee (Board of Publ. of the C. P. Ch.), 1875 (pp. 286). The same book contains also the Shorter Catechism, the Form of Government and Discipline, the Directory of Worship, and Manual.

The history of the origin of the schism is contained in the \emph{Circular Letter} of the late Cumberland Presbytery; the \emph{Reply} to a Pastoral Letter of West Tennessee Presbytery.

II. On the part of the Presbyterian Church

Samuel Baird: \emph{Collection of the Acts, Deliverances, and Testimonies of the Presbyterian Church.} Philad. (Presbyt. Board), 1855; second ed. 1859, pp. 640 sqq. Contains the official acts of the General Assembly on the origin and disorders of the Cumberland Presbytery.

Wm. E. Moore: \emph{A New Digest of the Acts and Deliverances of the General Assembly of the Presbyterian Church in the United States of America.} Philadelphia, 1861, p. 95 (on the validity of the Cumberland Presbyterian ordinances), and p. 448 (on terms of correspondence).

Robert Davidson: \emph{History of the Presbyterian Church in the State of Kentucky.} New York, 1847 (Ch. ix. pp. 223 sqq., 'The Cumberland Presbyterian Schism').

Historical and Doctrinal.

James Smith: \emph{History of the Christian Church, including a History of the Cumberland Presbyterian Church.} Nashville, 1835.

H. B. Crisman: \emph{Origin and Doctrines of the Cumberland Presbyterian Church.} 1856, new ed. Nashville, Tenn. 1875.

Richard Beard (D.D. and Prof. of Syst. Theol. in Cumberland University, Lebanon, Tennessee): \emph{Why am I a Cumberland Presbyterian?} Nashville, Tenn. 1872. By the same: \emph{Lectures on Systematic Theology}, 3 vols. Nashville (Board of Publ.). Comp. his Art. in Johnson's \emph{Universal Cyclop.} 1876, Vol. I.

F. R. Cossitt: \emph{Life and Times of Rev. Finis Ewing.} Louisville, 1853.

HISTORICAL.

The Cumberland Presbyterian Church in the United States of America, so called from its birth-place, the 'Cumberland Country' in Kentucky and Tennessee, took its rise in an extensive revival of religion which began in the southwestern part of Kentucky in 1797, and reached its height in 1800 and 1801, among a population mostly of Scotch-Irish descent. Methodist ministers took part in it. This revival called for a larger number of ministerial laborers than could be supplied in the regular way by the few Presbyterian institutions of learning then existing. Hence the Presbytery of Cumberland ('at the recommendation of the Rev. Mr. Rice, the oldest Presbyterian minister then residing in Kentucky') licensed and ordained a number of pious men without a liberal education, and allowed them, in subscribing the Westminster Confession, to express their dissent from what they called the doctrine of 'fatality,' i.e., the doctrine of absolute decrees. The Synod of Kentucky demanded a re-examination of these ministers and candidates; this being refused, it dissolved the Cumberland Presbytery in 1806. The General Assembly confirmed the action, but ultimately recognized the Cumberland Presbyterians as an independent organization, and entered into terms of correspondence with them as with other evangelical denominations.\footnote{In 1825 the General Assembly declared that the ministrations of the Cumberland Presbyterians 'are to be viewed in the same light with those of other denominations' (Baird's \emph{Collection}, p. 646). In 1849 the General Assembly of the New School entered into correspondence with them, and passed this resolution: 'The General Assembly of each Church shall appoint and receive delegates from the General Assembly of the other Church, who shall be possessed of all the powers and privileges of other members of such Assemblies, except that of voting' (\emph{Minutes}, p. 184; Moore, p. 448). The Rev. Dr. Alexander J. Baird appeared as a delegate of the Cumberland Presbyterian Church before the United General Assembly in Baltimore, 1873, and was cordially received (\emph{Minutes of the General Assembly of the Presbyt. Church} for 1873, p. 485). In the following year the General Assembly at St. Louis sent a salutation to the Cumberland Presbyterian Assembly then in session at Springfield, Mo., with the words: 'Serving the same Lord, we are one in him. May he dwell in us.' To this the Cumberland Assembly responded in the same fraternal spirit (\emph{Minutes} for 1874, pp. 18 and 20). A committee of conference on union was also appointed, but was discharged by the General Assembly of 1875 (\emph{Minutes}, p. 480).}

The dissenters organized an independent 'Cumberland Presbytery,' February 4, 1810, consisting of four regularly ordained ministers, six licentiates, and seven candidates. The presbytery grew into the Cumberland Synod in 1818, and this adopted a Confession, Catechism, and Form of Church Government. The Confession was the work of a committee of which the Rev. Finis Ewing was the leading spirit. The Cumberland Synod was divided into three (1828), and a General Assembly was formed, which held its first session in May, 1829. This Body subjected the Confession of Faith to a final revision. 'In so doing, the Synod and General Assembly only exercised an undeniable right, allowed by the God of the Bible and secured by the civil constitution; and discharged what they conceived to be a duty to the Church and the world. . . . Let the work be tried neither by tradition nor the fathers, but by the holy Scriptures.'\footnote{Preface to the Confession.}

The Cumberland Church has since spread rapidly, and extends now from Western Pennsylvania to Texas and California. It furnishes the proof that people may be good Presbyterians without being Calvinists.

THE CUMBERLAND PRESBYTERIAN CONFESSION.

The Cumberland Presbyterians differ from the regular Presbyterians in two points—the education for the ministry and the doctrine of predestination. They adopt and use the Westminster Confession in full, with the American amendments in Chs. XXIII. and XXXI., and slight verbal changes, but they depart from it in rejecting the unconditional election and reprobation as taught in Ch. III.\footnote{See the changes in Vol. III. p. 771.} They retain, however, substantially Ch. XVII. on perseverance, although perseverance presupposes unconditional election, and is inconsistent with conditional election. The Cumberland Confession teaches on the one hand conditional election and unlimited atonement, and on the other the final perseverance of the saints. It is an eclectic compromise between Calvinism and Arminianism; it is half Calvinistic and half Arminian, and makes no attempt to harmonize these antagonistic elements. 'Cumberland Presbyterians,' says one of their writers, 'believe as firmly as Arminians do that salvation, in all cases, is conditional. But they believe that every genuine saint will comply with the conditions; and thus salvation becomes certain to saints. It is uncertain to sinners because it is doubtful whether they will comply with the conditions; but certain to saints because it is certain that they will comply with the conditions—``My sheep hear my voice, and they follow me.'''\footnote{Crisman, 1.c. p. 158. Comp. art. of Prof. R. Beard, 1.c.: 'Its theology is Calvinistic, with the exception of the offensive doctrine of predestination so expressed as to seem to embody the old pagan dogma of \emph{necessity} or \emph{fatality.}'} The same writer answers the usual objections to the doctrine of perseverance (the fall of Adam and the angels, of Solomon and Peter, the warnings and exhortations of Scripture, the alleged inconsistency of the doctrine with free agency and the duty of watchfulness), and urges nine reasons against the Arminian view of falling from grace.\footnote{The difficulties of this great problem of predestination have been discussed more fully in § 97, pp. 791 sqq.}

Another departure connected with the former is the affirmation of the salvation of all infants dying in infancy. The old Confession says. Ch. X. 3: '\emph{Elect} infants, dying in infancy, are regenerated and saved by Christ through the Spirit, who worketh when and where and how he pleaseth.' This seems naturally (though not necessarily) to imply the existence of \emph{reprobate} infants who are not saved. To avoid this interpretation, the Cumberland Confession substitutes \emph{all} for \emph{elect}, and thus positively teaches universal infant salvation. In this point it has anticipated what seems now to be the general sentiment among American Presbyterians, who harmonize it with the Westminster Confession either by interpreting that \emph{all} infants dying in infancy are \emph{elect}, or that it confines itself to state as an article of faith what is clearly warranted in Scripture, and leaves the rest to private opinion.

The Shorter Catechism of the Assembly has been changed by the Cumberland Presbyterians in Question 7 as follows:

westminster catechism.

cumberland catechism.

\emph{What are the decrees of God?}

\emph{What are the decrees of God?}

The decrees of God are his \emph{eternal} purpose according to the counsel of his will, whereby, for his own glory, he hath foreordained whatsoever comes to pass.

The decrees of God are his purpose according to the counsel of his \emph{own} will, whereby he hath foreordained \emph{to bring} to pass what \emph{shall be} for his own glory: \emph{sin not being for God's glory, therefore he hath not decreed it.}

In Question 20 the words 'God did provide salvation \emph{for all mankind}' are substituted for 'God, having elected \emph{some} to everlasting life,' and the phraseology is otherwise changed. In Question 31, for the phrase 'What is effectual calling?' is substituted 'What is the work of the Spirit?'

[Note.—In 1906, the Cumberland Presbyterian Church was ``reunited'' with the Presbyterian Church, U. S. A., accepting the Westminster Confession as revised, 1902. A dissenting element retained the old name and has perpetuated the organization with a membership, 1929, of 64,081. At the time of the union, 1906, the Cumberland Church reported 200,000 members in 114 presbyteries.—Ed.]

\xchapter{Chapter 8. The Creeds of Modern Evangelical Denominations.}

\xsection{General Survey.}

§ 100.  General Survey.

With the Westminster standards the creed-making period of the Reformed Churches was brought to a close. Calvinism found in them its clearest and fullest exposition. The Helvetic Consensus Formula (1675) was only a weak symbolical after-birth, called forth by the Saumur controversies on the extent of divine election and the inspiration of Hebrew vowel-points. The creative power of Lutheran symbolism had exhausted itself much earlier in the Formula of Concord (1577), and was followed by a period of scholastic analysis and demonstration of the Lutheran system as embodied in its authoritative confessions. The prevailing tendency in these Churches is to greater confessional freedom and catholic expansion rather than sectarian contraction. While the Roman Catholic Church in our age has narrowed its creed by adding two new dogmas of wide range and import, and has doomed to silence every dissent from the infallible decisions of the Vatican, like a machine that is worked by a single motive force, and makes resistance impossible, the Protestant Churches would simplify and liberalize their elaborate standards of former days rather than increase their bulk and tighten their authority. The spirit of the age refuses to be bound by rigorous formulas, and demands greater latitude for private opinion and theological science.

We might therefore close our history of creeds at this point. But evangelical Protestantism extends far beyond the boundaries of Lutheranism and Calvinism.

Since the middle of the seventeenth century there arose, mainly from the fruitful soil of the Reformed Church in England, first amid much persecution, then under the partial protection of the Toleration Act of 1689, a number of distinct ecclesiastical organizations, which, while holding fast to the articles of the œcumenical faith of orthodox Christendom, and the evangelical principles of the Protestant Reformation, differ on minor points of doctrine, worship, and discipline. They have passed through the bloody baptism of persecution as much as the older Churches of the Reformation, and by their fruits they have fully earned a title to an honorable standing in the family of Christian Churches.

The most important among these modern denominations are the Congregationalists, Baptists, and Quakers, who rose in the seventeenth century, and the Methodists and Moravians, who date from the middle of the eighteenth century. They originated in England, with the exception of the Moravians (who are of Bohemian and German descent), and found from the start a fruitful and congenial soil in the American colonies, which offered an hospitable asylum to all who suffered from religious persecution. The Congregationalists had established flourishing colonies in Massachusetts and Connecticut before they were even tolerated in the mother country. Roger Williams, the patriarch of the American Baptists, though of English birth and training, made Rhode Island his permanent home. The fathers and founders of the Society of Friends—Fox and Penn; of Methodism— Wesley and Whitefield; of the Moravian Church—Zinzendorf, Spangenberg, Nitschmann—visited America repeatedly, and with such success that they gave to their denominations an Anglo-American stamp. Two of these denominations, the Methodists and Baptists, have in the United States during the nineteenth century numerically far outgrown the older Protestant Churches, and are full of aggressive zeal and energy, both at home and in distant missionary fields.\footnote{\par The following comparative table of ministers and churches in 1776 and 1876 gives at least an approximate idea of the growth of churches in the United States during its first centennial:     Statistics of 1776 (or 1780-90)  Statistics of 1876    Denominations. Ministers. Churches. Denominations. Ministers. Churches.   Baptists. . 722 872 Baptists. . 13,779 22,924   Congregationalists.  575 700 Congregationalists.  3,333 3,509   Episcopalians.  150 200 Episcopalians.  3,216 4,000     (No bishop.)     (61 bishops)     Friends (Quakers).  400 500 Friends (Quakers).  865 885   Lutherans (1786).  25 60 Lutherans.  2,662 4,623   Methodists.  24 . . . . Methodists.  20,453 40,000   Moravians.  12(?) 8(?) Moravians.  75 75   Presbyterians (1788).  177 419 Presbyterians.  4,744 5,077   Reformed, Dutch.  40 100 Reformed, Dutch.  546 506   Reformed, German.  12 60 Reformed, German.  644 1,353   Roman Catholics.  26(?) 52(?) Roman Catholics.  5,141 5,046           (56 bishops)}

On the Continent of Europe these Anglo-American denominations till quite recently were little known, and were even persecuted as intruders and unchurchly sects. National State Churches will allow the  widest latitude of theological speculation within the limits of outward conformity rather than grant freedom of public worship to dissenting organizations, however orthodox.\footnote{Under the disparaging name of \emph{sects} the Methodists and Baptists, and other denominations figure usually in German works on Symbolics that recognize only three \emph{Churches} or \emph{Confessions}—the Catholic (Greek and Roman), the Lutheran, and the Reformed (Calvinistic). The late Professor Marheineke, one of the chief writers on Symbolics, after explaining to his catechumens of Trinity Parish, in Berlin, that there are three Churches in Christendom, asked a pupil, 'To what Church do you belong?' and received the answer, 'To Trinity Church.' The science of Symbolics, or Comparative Theology, has thus far been almost exclusively cultivated in Germany, but should be reconstructed on a much more liberal scale in England and America, where all denominations meet in daily intercourse and on terms of equal rights.}

The nineteenth century has given birth in England to the Irvingites and Darbyites, and in America to the Cumberland Presbyterians,  Reformed Episcopalians, and other organizations, which more or less depart from the older Protestant confessions, but adhere to the supernatural revelation in the Bible and the fundamental articles of general orthodoxy.\footnote{Some of these have already been considered, the Cumberland Presbyterians in connection with the Westminster Confession, the Reformed Episcopalians in connection with the history of the Thirty-nine Articles.}

The creeds of these modern Protestant denominations (if we except the Savoy Declaration of 1658 and the Baptist Confession of 1688, which contain the body of the Westminster Confession) are thin, meagre, and indefinite as compared with the older confessions, which grew out of the profound theological controversies of the sixteenth century. They contain much less theology; they confine themselves to a popular statement of the chief articles of faith for practical use, and leave a large margin for the exercise of private judgment. In this respect they mark a return to the brevity and simplicity of the primitive baptismal creeds and rules of faith. The authority of creeds, moreover, is lowered, and the absolute supremacy and sufficiency of the Scriptures is emphasized.

In the present age there is, especially in America, a growing tendency towards a liberal recognition and a closer approach of the various evangelical denominations in the form of a free union and co-operation in the common work of the Master, without interfering with the inner organization and peculiar mission of each. This union tendency manifests itself from different starting-points and in different directions, now in the form of voluntary associations (such as Bible and Tract Societies, Young Men's Christian Associations, the Evangelical Alliance, the German Church Diet), now in the form of ecclesiastical confederations (Pan-Anglican Council, Presbyterian Alliance, Anglo-Greek Committees, the Bonn Conferences), now in the form of organic union (the evangelical Union of Lutherans and Reformed Churches in Prussia and other German States, Presbyterian Reunion of Old and New School). The same tendency calls forth efforts, feeble as yet, to formulate the essential consensus of the creeds of congenial sections of Christendom. The old motto, \emph{in necessariis unitas, in dubiis libertas, in omnibus caritas,} is struggling to become a practical reality; the age of separation and division is passing away, and the age of the reunion of divided Christendom is beginning to dawn, and to gather the corps of Christ's army, so long engaged in internal war, against the common foe Antichrist.

\xsection{The Congregationalists.}

§ 101. The Congregationalists.

Literature.

I. English Congregationalism.

See the sources of the Westminster Assembly, and the historical works of Neal, Stoughton, and others mentioned in §§ 92, 93, and 94.

John Robinson (Pastor of the Pilgrim Fathers in Leyden, d. 1626): \emph{Works, with Memoir by Robert Ashton.} London, 1851, 3 vols.

\emph{The Grand Debate concerning Presbytery and Episcopacy} in the Westminster Assembly (Lond. 1652).

The works of Drs. Goodwin,  Owen,  Howe, and other patriarchs of Independency.

Benjamin Brook:  \emph{The Lives of the Puritans from Queen Elizabeth to} 1662. London, 1813, 3 vols.

Benjamin Hanbury:  \emph{Historical Memorials relating to the Independents or Congregationalists, from their Rise to the Restoration of the Monarchy, A.D.} 1660. London (Congreg. Union of England and Wales), 1839–1844, 3 vols.

Jos. Fletcher:  \emph{History of Independency in England since the Reformation.} London, 1847–1849, 4 vols.

George Punchard (of Boston): \emph{History of Congregationalism from about A.D.} 250 \emph{to the Present Time.} 2d ed. rewritten and enlarged, New York and Boston (Hurd \& Houghton), 1865–81, 5 vols. (The first two vols. are irrelevant.)

John Waddington:  \emph{Congregational History}, 1200–1567. London, 1869–78, 4 vols. Second volume from 1567 to 1700, Lond. 1874. (See a searching and damaging review of this work by Dr. Dexter in the ``Congreg. Quarterly'' for July, 1874, Vol. XVI. pp. 420 sqq.)

Herbert S. Skeats:  \emph{A History of the Free Churches of England from} l688 \emph{to} l851. London, 1867; 2d. ed. 1860.

II. American Congregationalism.

(1) \emph{Sources.}

The works of John Robinson, above quoted, especially his \emph{Justification of Separation from the Church of England} (1610, printed in 1639).

John Cotton (of Boston, England, and then of Boston, Mass.): \emph{The Way of the Churches of Christ in New England. Or the Way of Churches Walking in Brotherly Equality or Co-ordination, without Subjection of one Church to another. Measured by the Golden Reed of the Sanctuary.} London, 1645. By the same: \emph{The Way of Congregational Churches cleared} (against Baillie and Rutherford). London, 1648.

Thomas Hooker (of Hartford, Conn.): \emph{A Survey of the Summe of Church Discipline.} London, 1648.

Robinson, Cotton, and Hooker are the connecting links between English Independency and American Congregationalism. Their rare pamphlets (wretchedly printed, like most works during the period of the civil wars, from want of good type and paper) are mostly found in the Congregational Library at Boston, and ought to be republished in collected form.

Alexander Young:  \emph{Chronicles of the Pilgrim Fathers of the Colony of Plymouth, from} 1602 \emph{to} 1628. Boston, 1841.

Alexander Young:  \emph{Chronicles of the first Planters of the Colony of Massachusetts Bay. From} 1623 \emph{to} 1636. Boston, 1846.

George B. Cheever:  \emph{The Journal of the Pilgrims at Plymouth, in New England, in} 1620; \emph{reprinted from the original volume, with illustrations.} New York, 1848.

Nathanael Morton (Secretary to the Court for the Jurisdiction of New Plymouth): \emph{New England's Memorial.} Boston, 1855 (6th ed. Congreg. Board of Publication). Reprints of Memorial of 1669, Bradford's History of Plymouth Colony, etc.

(2) \emph{Histories.}

Benjamin Trumbull, D.D.:  \emph{A Complete History of Connecticut, Civil and Ecclesiastical, from the Emigration of its first Planters, from England, in the year }1630, \emph{to the year }1764. New Haven, 1818, 2 vols.

Leonard Bacon:  \emph{Thirteen Historical Discourses, on the Completion of Two Hundred Years from the Beginning of the First Church in Sew Haven.} New Haven, 1839.

Joseph B. Felt:  \emph{The Ecclesiastical History of New England; comprising not only Religious, but also Moral and other Relations.} Boston, Mass. (Congregational Library Association), 1855–1862, 2 vols.

Joseph S. Clark:  \emph{A Historical Sketch of the Congregational Churches in Massachusetts from }1620 \emph{to }1858. Boston, 1858.

\emph{Memorial of the Semi-Centennial Celebration of the Founding of the Theological Seminary at Andover.} Andover, Mass. 1859.

\emph{Contributions to the Ecclesiastical History of Connecticut; prepared under the Direction of the General Association to Commemorate the Completion of One Hundred and Fifty Years since its First Annual Assembly.} New Haven (publ. by Wm. L. Kingsley), 1861.

Daniel Appleton White:  \emph{New England Congregationalism in its Origin and Purity; Illustrated by the Foundation and Early Records of the First Church in Salem} [Mass.]. Salem, 1861. Comp. \emph{Reply} to the above, by Joseph B. Felt. Salem, 1861.

The first vols. of  G. Bancroft's  \emph{History of the United States} (begun in 1834); last ed. 1876, 6 vols.

John Gorham Palfrey:  \emph{History of New England.} Boston, 1859–1874, 4 vols.

Leonard Bacon:  \emph{The Genesis of the New England Churches.} New York, 1874.

Henry Martyn Dexter:  \emph{As to Roger Williams and his 'Banishment' from the Massachusetts Plantation; with a few further Words concerning the Baptists, the Quakers, and Religious Liberty.} Boston, 1876 (Congregational Publishing Society). A vindication of the Massachusetts Colony against the charge of intolerance.

Numerous essays and reviews relating to the Congregational polity and doctrine and the history of Congregational Churches may be found in the volumes of the following periodicals:

\emph{American Quarterly Register.} Boston, Mass. 1827–1843, 15 vols.

\emph{The Christian Spectator.} 1st series monthly; 2d series quarterly. New Haven, 1819–1838, 20 vols.

\emph{The New-Englander}, quarterly (continued). New Haven, 1843–1876, 34 vols.

\emph{The Congregational Quarterly} (continued). Boston, Mass. 1st series, 1859–1868, 10 vols.; 2d series, 1869–1876, 8 vols.

\emph{The Congregational Year-Book.} New York, 1854–1859, 5 vols.

Other light is thrown on the Congregational history and polity by \emph{Results of Councils}, many of which, in cases of peculiar interest, have been published in pamphlet form.

(3) \emph{Congregational Polity.}

\emph{Congregational Order. The Ancient Platforms of the Congregational Churches of New England, with a Digest of Rules and Usages in Connecticut. Publ. by direction of the General Association of Connecticut.} Middletown, Conn. 1843. [Edited by Leonard Bacon,  David D. Field,  Timothy P. Gillet.]

Thomas C. Upham:  \emph{Ratio Disciplinæ; or, The Constitution of the Congregational Churches, Examined and Deduced from Early Congregational Writers, and other Ecclesiastical Authorities, and from Usage.} 2d edition. Portland, 1844.

Preston Cummings:  \emph{A Dictionary of Congregational Usages and Principles according to Ancient and Modern Authors; to which are added brief Notices of some of the Principal Writers, Assemblies, and Treatises referred to in the Compilation.} Boston, 1852.

George Punchard:  \emph{A View of Congregationalism, its Principles and Doctrines; the Testimony of Ecclesiastical History in its Favor, its Practice, and its Advantages.} [1st edition, 1840.] Third edition, revised and enlarged. Boston (Congreg. Board of Publication), 1856.

Henry Martyn Dexter:  \emph{Congregationalism: What it is; Whence it is; How it Works; Why it is Better than any other Form of Church Government.} Boston, 1865; 5th ed. revised, 1879.

Congregationalism has its name from the prominence it gives to the particular congregation as distinct from the general Church.\footnote{This term is preferable to \emph{Independency.} In England both terms are used synonymously. The American Congregationalists rather disclaim the designation Independents, except for a small portion of their ancestors, namely, the 'Pilgrim Fathers' of Plymouth. See below.} It aims to establish a congregation of real believers or converts, and it declares such a congregation to be independent of outward jurisdiction, whether it be that of a king or a bishop or a presbytery. Under the first aspect it has several precedents; under the latter aspect it forms a new chapter in Church history, or at least it carries the protest against foreign jurisdiction a great deal farther than the Reformers, who protested against the tyrannical authority of the papacy, but recognized some governmental jurisdiction over local congregations.

CONGREGATIONS IN THE APOSTOLIC AGE.

In the New Testament the word \emph{church} or \emph{congregation}\footnote{\textgreek{ἐκκλησία,} from \textgreek{ἐκκαλέω, } \emph{to call out}, means (like \texthebrew{קָהָל}) any public assembly, but especially a religious assembly.} denotes sometimes the Church universal, the whole body of Christian believers spread throughout the world;\footnote{Matt. xvi. 18; Acts xx. 28; Gal. i. 13; Eph. i. 22, etc.} sometimes a particular congregation at Jerusalem, Antioch, Corinth, Rome, or any other place.\footnote{Matt. xviii. 17; Acts v. 11; viii. 3; xv. 41 (in the plural, \textgreek{αἱ ἐκκλησίαι}); Gal. i. 22; Rom. xvi. 4, 5, etc.} The congregations are related to the Church as members to the body. The denominational and sectarian use of the word is foreign to the Scriptures, which know of no sect but the sect called Christians.\footnote{Comp. Acts xi. 26; xxvi. 28; 1 Pet. iv. 16. There were parties or sects among the Christians at Corinth which assumed apostolic designations, but Paul rebuked them (1 Cor. i, 10–13; iii. 3, 4). The tribes of Israel may be quoted as a Jewish precedent of the divisions in Christendom, but they formed one nation.} Denominations or Confessions are the growth of history and adaptations of Christianity to the differences of race, nationality, and psychological constitution; and after fulfilling their mission they will, as to their human imperfections and antagonisms, disappear in the one kingdom of Christ, which, however, in the beauty of its living unity and harmony, will include an endless variety.

An organized local congregation in the apostolic age was a company of saints,\footnote{\textgreek{ἐκκλησίαι τῶν ἁγίων, } 1 Cor. xiv. 33.} or a self-supporting and self-governing society of Christian believers, with their offspring, voluntarily associated for purposes of worship, growth in holiness, and the promotion of Christ's kingdom. The Apostolic churches were not free from imperfection and corruption, but they were separated from the surrounding world of unbelievers, and constantly reminded of their high and holy calling.

THE ANTE-NICENE CHURCHES.

In the ante-Nicene age a distinction was made between the church of \emph{believers} or communicant members and the church of \emph{catechumens} or hearers who were in course of preparation for membership, but not allowed to partake of the communion.\footnote{Comp. the modern American distinction between church proper and congregation.} Public worship was accordingly divided into the service of the faithful (\emph{missa fidelium}) and the service of the catechumens (\emph{missa catechumenorum}).

MIXTURE OF THE CHURCH WITH THE WORLD.

With the union of Church and State since Constantine the original idea of a church of real believers was gradually lost, and became identical with a parish which embraced all nominal Christians in a particular place or district. Baptism, confirmation, and attendance at communion were made obligatory upon all residents, whether converted or not, and every citizen was supposed to be a Christian.\footnote{The Jews—like the 'untaxed Indians' in the United States—were excluded from the rights of citizenship, and as unmercifully persecuted during the Middle Ages as the Christians were persecuted by the Jews in the apostolic age.} The distinction between the Church and the world was well-nigh obliterated, and the Church at large became a secular empire with an Italian sovereign at its head. Hence the complaint of Dante (in Milton's rendering):

'Ah! Constantine, of how much ill was cause,

Not thy conversion, but those rich domains

That the first wealthy Pope received of thee!'

ATTEMPTS TO RESTORE THE PURITY OF THE CHURCH.

Monasticism was an attempt in the Catholic Church itself to save the purity of the congregation by founding convents and nunneries secluded not only from the world, but also from all ties of domestic and social life. It drained the Church of many of its best elements, and left the mass more corrupt.

The Bohemian Brethren and the Waldenses introduced strict congregational discipline in opposition to the ruling Church.

The Reformers of the sixteenth century deplored the want of truly Christian congregations after the apostolic model, and wished to revive them, but Luther and Zwingli gave it up in despair from the want of material for congregational self-government (which can never be developed without an opportunity and actual experiment).

Calvin was more in earnest, and astonished the world by founding in Geneva a flourishing Christian commonwealth of the strictest discipline, such as had not been seen since the age of the Apostles. But it was based on a close union of the civil and ecclesiastical power, which destroyed the voluntary feature, and ended at last in the same confusion of the Church and the world.

The Anabaptists and Mennonites emphasized the voluntary principle and the necessity of discipline, but they injured their cause by fanatical excesses.

The German Pietists of the school of Spener and Francke realized their idea of \emph{ecclesiolæ} in \emph{ecclesia}, or select congenial circles within the outward organization of the promiscuous national Church, from which they never separated. Wesley did originally the same thing, but his movement resulted in a new denomination.

The Moravians went farther, and established separate Christian colonies, which in the period of rationalism and infidelity were like beacon-lights in the surrounding darkness.

ENGLISH AND AMERICAN CONGREGATIONALISM.

English and American Congregationalism, or Congregationalism as a distinct denomination, arose among the Puritans during the latter part of the reign of Queen Elizabeth. It was at first identified with the name of the Rev. Robert Browne, and called Brownism; but, being an unworthy representative and an apostate from his principles, he was disowned.\footnote{Robert Browne, a clergyman of the Established Church and a restless agitator, urged a reformation 'without tarrying for any,' a complete separation from the national Church as an anti-Christian institution, and the formation of independent Christian societies. After suffering persecution and exile (he was imprisoned about thirty times), he returned to the Ministry of the national Church, where he led an idle and dissolute life till his death, in 1630, at the age of eighty years.} It had other and more worthy pioneers, such as Barrowe, Greenwood, Johnson, Ainsworth, Penry, and especially John Robinson.\footnote{See on these early witnesses and martyrs of Independency, Hanbury (Vol. I. chaps. ii.-xxvi.), Brook (Vol. III.), and Punchard (Vol. III.).} The Independents were, like every new sect, persecuted under the reigns of James and Charles I., and obliged to seek shelter first in Holland and then in the wilderness of New England.

But with the opening of the Long Parliament, which promised to inaugurate a jubilee to all tender consciences, they began to breathe freely, and hastened to return from exile; 'for,' says Fuller, 'only England is England indeed, though some parts of Holland may be like unto it.'\footnote{Vol. VI. p. 280.} They had a considerable share in the labors of the Westminster Assembly of Divines, especially through Dr. Goodwill and Rev. Philip Nye, who are styled the 'patriarchs' of orthodox Independency. They became the ruling political and religious power in England during the short protectorate of Cromwell, and furnished the majority to his ecclesiastical commission, called the Triers. After the Restoration they were again persecuted, being held chiefly responsible for the execution of King Charles and the overthrow of the monarchy. In 1689 they acquired toleration, and are now one of the most intelligent, active, and influential among the Dissenting bodies in England.

The classical soil of Congregationalism is New England, where it established 'a Church without a bishop and a State without a king.' From New England it spread into the far West, to the shores of the Pacific Ocean, and exerted a powerful influence upon other Churches. Puritan Congregationalism is the father of New England and one of the grandfathers of the American Republic, and it need not be ashamed of its children.\footnote{I beg leave to quote from an essay which I wrote and published in the midst of our civil war (1863), when New England was most unpopular, the following tribute to its influence upon American history: 'It seems superfluous, even in these days of sectional prejudice, party animosity, and slander, to say one word in praise of New England. Facts and institutions always speak best for themselves. We might say with Daniel Webster, giving his famous eulogy on Massachusetts a more general application to her five sister States: ``There they stand: look at them, and judge for yourselves. There is their history—the world knows it by heart: the past at least is secure.'' The rapid rise and progress of that rocky and barren country called New England is one of the marvels of modern history. In the short period of two centuries and a half it has attained the height of modern civilization which it required other countries more than a thousand years to reach. Naturally the poorest part of the United States, it has become the intellectual garden, the busy workshop, and the thinking brain of this vast republic. In general wealth and prosperity, in energy and enterprise, in love of freedom and respect for law, in the diffusion of intelligence and education, in letters and arts, in virtue and religion, in every essential feature of national power and greatness, the people of the six New England States, and more particularly of Massachusetts, need not fear a comparison with the most favored nation on the globe. But the power and influence of New England, owing to the enterprising and restless character of its population, extends far beyond its own limits, and is almost omnipresent in the United States. The twenty thousand Puritans who emigrated from England within the course of twenty years, from 1620 to 1640, and received but few accessions until the modern flood of mixed European immigration set in, have grown into a race of several millions, diffused themselves more or less into every State of the Union, and take a leading part in the organization and development of every new State of the great West to the shores of the Pacific. Their principles have acted like leaven upon American society; their influence reaches into all the ramifications of our commerce, manufactures, politics, literature, and religion; there is hardly a Protestant Church or Sabbath-school in the land, from Boston to San Francisco, which does not feel, directly or indirectly, positively or negatively, the intellectual and moral power that constantly emanates from the classical soil of Puritan Christianity.'} It lacks a proper appreciation of historical Christianity and its claims upon our regard and obedience; but by bringing to light the manhood and freedom of the Christian people, and the rights and privileges of individual congregations, it marks a real progress in the development of Protestantism, and has leavened other Protestant denominations in America; for here congregations justly claim and exercise a much larger share, and have consequently a much deeper interest in the management of their own affairs than in the State Churches of Europe. The Congregational system implies, of course, the power of self-government and a living faith in Christ, without which it would be no government at all. It moreover requires the cementing power of fellowship.

INDEPENDENCY AND FELLOWSHIP.

Anglo-American Congregationalism has two tap roots, independency and fellowship, on the basis of the Puritan or Calvinistic faith. It succeeds in the measure of its ability to adjust and harmonize them. It is a compromise between pure Independency and Presbyterianism. It must die without freedom, and it can not live without authority, Independency without fellowship is ecclesiastical atomism; fellowship without Independency leads to Presbyterianism or Episcopacy.\footnote{Dr. Emmons, one of the leaders of New England Congregationalism, is credited with this memorable \emph{dictum}: 'Associationism leads to Consociationism; Consociationism leads to Presbyterianism; Presbyterianism leads to Episcopacy; Episcopacy leads to Roman Catholicism; and Roman Catholicism is an ultimate fact' (Prof. Park, in \emph{Memoir of Emmons}, p. 163). But there would be equal force in the opposite reasoning from Independency to anarchy, and from anarchy to dissolution. Independents have a right to protest against tyranny, whether exercised by bishops or presbyters ('priests writ large'); but there are Lord Brethren as well as Lord Bishops, and the tyranny of a congregation over a minister, or of a majority over a minority, is as bad as any other kind of tyranny.}

It starts from the idea of an apostolic congregation as an organized brotherhood of converted believers in Christ. This was the common ground of the Westminster divines.\footnote{'The Form of Presbyterial Church Government agreed upon by the Assembly of Divines at Westminster,' and adopted by the General Assembly of Scotland in 1645, thus defines a local Church: 'Particular churches in the primitive times were made up of visible saints, viz., such as, being of age, professed faith in Christ and obedience unto Christ, according to the rules of faith and life taught by Christ and his apostles, and of their children.' The Form of Government ratified by the General Assembly of the Presbyterian Church in the United States in May, 1821, gives this definition (Ch. II. 4): 'A particular church consists of a number of professing Christians, with their offspring, voluntarily associated together for divine worship and godly living, agreeably to the Holy Scriptures, and submitting to a certain form of government.'} But they parted on the question of jurisdiction and the relation of the local congregation to the Church general. The Independents denied the authority of presbyteries and synods, and maintained that each congregation properly constituted is directly dependent on Christ, and subject to his law, and his law \emph{only.} The whole power of the keys is vested in these individual churches.

At the same time, however, it is admitted and demanded that there should be a free fraternal intercommunion between them, with the rights and duties of advice, reproof, and co-operation in every Christian work.

This fellowship manifests itself in the forms of Councils, Associations (in Massachusetts), Consociations (in Connecticut), on a larger scale in 'the Congregational Union of England and Wales,' and 'the National Council of the Congregational Churches in the United States.' It is this fellowship which gives Congregationalism the character of a denomination among other denominations. But the principle of congregational sovereignty is guarded by denying to those general meetings any legislative authority, and reducing them simply to advisory bodies.\footnote{The most serious conflict between the principles of Independency and Fellowship in recent times has grown out of the unhappy Beecher trial, which has shaken American Congregationalism to the very base. See Proceedings of the two Councils held in Brooklyn in 1874 and 1876, which represent both sides of the question (Dr. Storrs's and Mr. Beecher's), though presided over by the same Nestor of American Congregationalism (Dr. Leonard Bacon).}

There were from the start two tendencies among Congregationalists—the extreme Independents or Separatists, of whom the 'Pilgrim Fathers' are the noblest representatives, and the more churchly Independents, who remained in the English Church, and who established on a Calvinistic theocratic basis the Commonwealth of Massachusetts. John Robinson, the Moses of American Independency, who accompanied his flock to the deck of the Speedwell, but never saw the promised land himself, was a separatist from the Church of England, though he disowned Brownism with its extravagances. His colony at Plymouth were Separatists. The settlers of Boston, Salem, Hartford, and New Haven, on the other hand, were simply Nonconformists within the Church of England. Their ministers—John Cotton, Richard Mather, Thomas Hooker, John Davenport, Samuel Stone, and others—were trained in the English Universities, mostly in Cambridge,\footnote{Masson (\emph{Life of Milton}, Vol. II. p. 563) says that of seventeen noted ministers who emigrated to New England, fourteen were bred in Cambridge, and only three (Davenport, Mather, and Williams) at Oxford. R. Williams was probably likewise a Cambridge graduate. It was therefore natural that the first college in New England should be called after Cambridge.} and had received Episcopal ordination. They rejected the term Independents, and inconsistently relapsed into the old notion of uniformity in religion, with an outburst of the dark spirit of persecution. But this was only temporary. American Congregationalism at present is a compromise between the two tendencies, and vacillates between them, leaning sometimes to the one, sometimes to the other side.

CONGREGATIONALISM AND CREEDS.

The effect of the Congregational polity upon creeds is to weaken the authority of general creeds and to strengthen the authority of particular creeds. The principle of fellowship requires a general creed, but it is reduced to a mere \emph{declaration} of the common faith prevailing among Congregationalists at a given time, instead of a binding formula of subscription. The principle of independency calls for as many particular creeds as there are congregations. Each congregation, being a complete self-governing body, has the right to frame its own creed, to change it \emph{ad libitum}, and to require assent to it not only from the minister, but from every applicant for membership. Hence there are a great many creeds among American Congregationalists which have purely local authority; but they must be in essential harmony with the prevailing faith of the body, or the congregations professing them forfeit the privileges of fellowship. They must flow from the same system of doctrine, as many little streams flow from the same fountain.

In this multiplication of local creeds Congregationalism far outstrips the practice of the ante-Nicene age, where we find varying yet essentially concordant rules of faith in Jerusalem, Cæsarea, Antioch, Aquileja, Carthage, Rome.

With these local creeds are connected 'covenants' or pledges of members to live conformably to the law of God and the faith and discipline of the Church. A covenant is the ethical application of the dogmatic creed.

In the theory of creeds and covenants, as on the whole subject of Church polity, the Regular or Calvinistic Baptists entirely agree with the Congregationalists.

\xsection{English Congregational Creeds.}

§ 102. English Congregational Creeds.

Literature.

\emph{A } | Declaration  | \emph{of the } | Faith  \emph{and } Order  | \emph{Owned and practised in the }| Congregational Churches  | \emph{in } | England;  | \emph{Agreed upon and consented unto } | \emph{by their } | Elders and Messengers  | \emph{in }| \emph{their } Meeting  \emph{at the } Savoy,  | \emph{Octob.} 12, 1658. | London | Printed for D.L. And are to be sold in Paul's Churchyard, Fleet | Street, and Westminster Hall, 1659.

A Latin edition appeared in 1662 at Utrecht, under the title, \emph{Confessio nuper edita Independentium seu Congregationalium in Anglia.}

The Preface, the Platform, and those doctrinal articles which differ from the Westminster Confession are printed in Vol. III. pp. 707 sqq., from the first London edition. The Savoy Declaration, without the Preface, is also given by Hanbury,  \emph{Memorials}, Vol. III. pp. 517 sqq.; and by Dr. A. H. Quint, in the 'Congregational Quarterly' for July and October, 1866 (Vol. VIII. pp. 241–267 and 341–344).

On the Savoy meeting, comp. Hanbury,  \emph{Memorials}, Vol. III. pp. 515 sqq.

THE SAVOY DECLARATION. A.D. 1658.

We now proceed to the general creeds or declarations of faith which have been approved by the Congregational Churches in England and America. They agree substantially with the Westminster Confession, or the Calvinistic system of doctrine, but differ from Presbyterianism by rejecting the legislative and judicial authority of presbyteries and synods, and by maintaining the independence of the local churches. In the course of time the rigor of old Calvinism has relaxed, both in England and America. 'New England theology,' as it is called, attempts to find a \emph{via media} between Calvinism and Arminianism in anthropology and soteriology. But the old standards still remain unrepealed.

The first and fundamental Congregational confession of faith and platform of polity is the Savoy Declaration, so called from the place where it was composed and adopted.\footnote{The Savoy, in the Strand, London, is remarkable for its historical associations. The palace, on the banks of the Thames, was built by Peter, Earl of Savoy and Richmond, in 1245; enlarged and beautified by Henry, Duke of Lancaster, 1328. King John II., of France, while a prisoner in England, resided there (1357–63). It was burned in Wat Tyler's insurrection, 1381; rebuilt and endowed as a hospital by Henry VII., 1505. It was the city residence of the Bishop of London. The royal chapel was burned down in 1864, but beautifully restored by Queen Victoria, and reopened Nov. 26, 1865. The Congregational meeting of 1658 must not be confounded with the 'Savoy Conference' between Episcopalians and Presbyterians which was held there from April 15 to July 25, 1661.}

The position of the Congregationalists during the short period of their ascendency under Cromwell's Protectorate (1653–1658) was rather anomalous. They were by no means so strongly committed to the voluntary principle and against a national Church as to refuse appointments in the universities and parish churches, with the tithes and other emoluments connected therewith. Dr. Goodwin was President of Magdalen College, Cambridge; Dr. Owen, Dean of Christ Church and Vice-Chancellor at Oxford; Philip Nye, Rector of St. Bartholomew's, London; Joseph Caryl, Rector of St. Mary Magnus; William Greenhill, incumbent of the village of Stepney; William Bridge, town lecturer at Yarmouth; John Howe, parish minister at Torrington, and afterwards court chaplain to Cromwell until his death.\footnote{Comp. Stoughton, \emph{Church of the Commonwealth}, ch. ix. pp. 207 sqq. A number of the Baptists likewise accepted preferments under the Protectorate. See ib. p. 242, and Ivimey's list of Baptists who were ejected at the Restoration, \emph{History of Baptists}, Vol. I. p. 328.} Cromwell himself had no idea of disconnecting the government from religion. Christianity was fully recognized under his rule as part and parcel of the law of the land. It accompanied with its solemn worship the ordinary business of Parliament. Public fasts were frequently appointed by the Protector (to which the Presbyterians objected as an Erastian intrusion), and lasted usually from nine in the morning until four in the afternoon. The rights of patronage were not disturbed; the tithes and other provisions for the support of the clergy and the repair of churches were continued. A commission of Triers, or judicial examiners, one fourth of whom were laymen, was appointed to test the fitness of clerical applicants and to remove unworthy incumbents, and Church boards of gentry and clergy were set up in every county for the supervision of ecclesiastical affairs. The Triers took the place of the late Westminster Assembly in its administrative work, but were less numerous, and included Independents, Presbyterians, and Baptists. Dr. Owen, Goodwin, and Manton belonged to them, besides others of less wisdom and charity. They were subject to a certain Erastian control by the Protector and his Council of State, but left to decide each case according to their best judgment, without imposing any creed or canon or statute. The plan seems to have worked well, and furnished the country, as Baxter says, who was no friend of Cromwell, with 'able, serious preachers, who lived a godly life, of what tolerable opinion soever they were.' Cromwell's Protectorate was too short to develop a full system of ecclesiastical polity. It was a government of experiments in accommodation to existing circumstances. Upon the whole, it was more tolerant than any previous reign, but only to Puritanism and such Protestant sects as recognized the Scriptures and the fundamentals of the Christian faith; while it was intolerant to Romanists, Socinians, and Episcopal royalists, who endangered his government. In his foreign policy Cromwell was the boldest protector of Protestantism and religious liberty that England has ever produced.\footnote{Comp. Stoughton, 1.c. pp. 81 sqq. Green (\emph{History of the English People}, p. 573) judges upon the whole quite favorably of Cromwell's ecclesiastical polity: 'In England, Cromwell dealt with the Royalists as irreconcilable enemies; but in every other respect he carried out fairly his pledge of ``healing and settling.'' . . . From the Church, which was thus reorganized, all power of interference with faiths differing from its own was resolutely withheld. Cromwell remained true to his great cause of religious liberty. Even the Quaker, rejected by all other Christian bodies as an anarchist and blasphemer, found sympathy and protection in Cromwell. The Jews had been excluded from England since the reign of Edward the First; and a prayer which they now presented for leave to return was refused by the commission of merchants and divines to whom the Protector referred it for consideration. But the refusal was quietly passed over, and the connivance of Cromwell in the settlement of a few Hebrews in London and Oxford was so clearly understood that no one ventured to interfere with them.'}

Under these favorable circumstances, and in view of the successful establishment of an exclusively Congregational commonwealth by their transatlantic brethren, the Independents might think of repeating in a milder form the experiment of the Westminster Assembly to secure at least a certain degree of religious uniformity in England, with a limited amount of toleration to orthodox dissenters. Their great protector did not seem to favor such a scheme, but shortly before his death he reluctantly gave his consent to 'the humble petition and advice' of influential members of Parliament to issue a confession of faith for the whole kingdom, yet 'without compelling the people thereto by penalties,' and to extend liberty to all Christian professions, except 'popery or prelacy,' or such as 'publish horrid blasphemies or practice or hold forth licentiousness or profaneness under the profession of Christ.' A notice from the clerk of the Council of State summoned the Congregational churches, in and near London, to a meeting in the Savoy, but it was not held till twenty-six days after Cromwell's death. About two hundred delegates from one hundred and twenty congregations attended the Conference, which lasted from Sept. 29 till Oct. 12, 1658. They agreed unanimously upon the Confession and Order of Discipline. It was regarded by them, in the language of the Preface, 'as a great and special work of the Holy Ghost that so numerous a company of ministers and other principal brethren should so \emph{readily, speedily}, and \emph{jointly} give up themselves unto such a whole body of truths that are after godliness.'

The Savoy Declaration is the work of a committee, consisting of Drs. Goodwill, Owen, Nye, Bridge, Caryl, and Greenhill, who had been members of the Westminster Assembly, with the exception of Dr. Owen. It contains a lengthy Preface (fourteen pages), the Westminster Confession of Faith with sundry changes (twenty-two pages), and a Platform of Church Polity (five pages).

1. The Preface is prolix and indifferently written, but deserves notice for inaugurating a more liberal view of the authority of creeds and the toleration of other creeds. The chief ideas are these: To confess our faith is an indispensable duty we owe to God as much as prayer. Public confessions are a means of expressing the common faith, but ought not to be enforced. 'Whatever is of force or constraint in matters of this nature causes them to degenerate from the name and nature of \emph{Confessions}, and turns them into \emph{Exactions} and \emph{Impositions of Faith.}' With this we should acknowledge 'the great principle that among all Christian States and Churches there ought to be vouchsafed a forbearance and mutual indulgence unto saints of all persuasions that keep unto and hold fast the necessary foundations of faith and holiness, in all other matters extra-fundamental, whether of faith or order.'

This was a considerable step beyond the prevailing notion of uniformity, although it falls far short of the modern theory of religious liberty. The Preface goes on to guard itself against the charge of indifference or carelessness.

2. The Declaration of Faith. This is a slight modification of the Westminster Confession. 'To this Confession,' the Preface states, 'we fully assent, as do our brethren of New England and the churches also of Scotland, as each in their general synods have testified. A few things we have added for obviating some erroneous opinions, and made other additions and alterations in method here and there, and some clearer explanations as we found occasion.' The Declaration is divided into thirty-two chapters, in the same order as the Westminster Confession, which has thirty-three chapters. In the exceptions taken the Savoy Council followed the example set by the Long Parliament in its edition of the Westminster Confession. The only important changes refer to matters of Church government and discipline. Chaps. XXX., 'Of Church Censures,' and XXXI., 'Of Synods and Councils,' are omitted altogether. Chaps. XXIII. (XXIV.), 'Of the Civil Magistrates,' XXIV. (XXV.), 'Of Marriage and Divorce,' and XXVI., 'Of the Church,' are modified. Chap. XX., 'Of the Gospel,' in the Savoy Declaration, is inserted, and hence the difference in the numbering of the remaining chapters. The change in Chap. XXIV. is a decided improvement, if we judge it from the American theory of Church and State. A similar and more thorough change was subsequently made by the American Presbyterians in the Westminster Confession.

3. The Declaration of 'the Institution of Churches and the Order appointed in them by Jesus Christ' contains the principles of the Congregational Church polity which we have already explained. Similar Platforms of Discipline, as they are called, have been issued from time to time by the American Congregationalists—at Cambridge, 1648, at Saybrook, 1708, and at Boston, 1865.

THE DECLARATION OF 1833.

This is a popular abridgment of the older confessions, and presents a milder form of Calvinism. It was prepared in 1833 by the Rev. Dr. Redford, of Worcester, and other members of a committee of the 'Congregational Union of England and Wales,' which was organized in 1831. It is annually printed in the 'Congregational Year-Book,' but it disclaims any authority as a standard of subscription.\footnote{See Vol. III. pp. 730 sqq.}

Note.—The Rev. Dr. John Stoughton, of London, a leading divine and historian among the English Independents, has kindly supplied me with the following statement concerning the prevailing sentiment of that body on the authority of creeds, a statement which applies largely to American Congregationalists in the present age:

'Looking at the principles of Congregationalism, which involve the repudiation of all human authority in matters of religion, it is impossible to believe that persons holding those principles can consistently regard any ecclesiastical creed or symbol in the same way in which Catholics, whether Roman or Anglican, regard the creeds of the ancient Church. There is a strong feeling among English Congregationalists against the use of such documents for the purpose of defining the limits of religious communion, or for the purpose of checking the exercise of sober, free inquiry; and there is also a widely spread conviction that it is impossible to reduce the expression of Christian belief to a series of logical propositions, so as to preserve and represent the full spirit of gospel truth. No doubt there may be heard in some circles a great deal of loose conversation seeming to indicate such a repugnance to the employment of creeds as would imply a dislike to any formal definition of Christian doctrine whatever; but I apprehend that the prevailing sentiment relative to this subject among our ministers and churches does not go beyond the point just indicated. Many consider that while creeds are objectionable as tests and imperfect as confessions, yet they may have a certain value as manifestoes of conviction on the part of religious communities.

'The Westminster Assembly's Catechism never had the authority in Congregational churches which from the beginning it possessed in the Presbyterian Church of Scotland, and its use in schools and families for educational purposes, once very common, has diminished of late years to a very low degree. The Savoy Declaration, which perhaps never had much weight with Congregationalists, is a document now little known, except by historical students. The Declaration of 1833 was prepared by a committee of the Congregational Union, of which the Rev. Dr. Redford, of Worcester, was a member. He, I believe, drew up the Articles, and it was only in accordance with his well-known character as a zealous antagonist of human authority in religion that he introduced the following passages in the preliminary notes:

'''It is not designed, in the following summary, to do more than to state the leading doctrines of faith and order maintained by Congregational churches in general.

'''It is not intended that the following statement should be put forth with any authority, or as a standard to which assent should be required.

'''Disallowing the utility of creeds and articles of religion as a bond of union, and protesting against subscription to any human formularies as a term of communion, Congregationalists are yet willing to declare, for general information, what is commonly believed among them, reserving to every one the most perfect liberty of conscience.''

'It would be well to insert a statement made to me by one who from his official position has the best means of ascertaining the state of opinion in our churches:

'''I do not believe that the Declaration of 1833 could now with success be submitted for adoption to an Assembly of the Congregational Union; in part, because not a few would dispute its position, and in part because many more—I believe the majority—without objecting on strictly doctrinal grounds, would object on grounds of policy.''

'I may add to this, in the words of the Dean of Westminster, who wrote them on the authority of ``a respected Congregational minister,'' that, beyond care in the matter of ordination, ``no measures are adopted or felt to be either desirable or necessary for preserving uniformity of doctrine, excepting only that the trust-deeds of most of their places of worship contain a reference to leading points of doctrine to which the minister may be required to express his assent. In practice this is merely a provision against any decided departure from the faith as commonly received among us, the trustees of the property having it in their power to refuse the use of the building to any minister whose teaching may be contrary to the doctrines contained in the deed. Such cases, however, are extremely rare.''

'In some cases trust-deeds make reference to the Declaration of 1833, as containing the doctrines to be taught in substance within the places of worship secured by such deeds; but in most cases a brief schedule of doctrines is employed, of which the following is an example:

'''1. The divine and special inspiration of the holy Scriptures of the Old and New Testament, and their supreme authority in faith and practice.

'''2. The unity of God. The Deity of the Father, of the Son, and of the Holy Ghost.

'''3. The depravity of man, and the absolute necessity of the Holy Spirit's agency in man's regeneration and sanctification.

'''4. The incarnation of the Son of God, in the person of the Lord Jesus Christ; the universal sufficiency of the atonement by his death; and the free justification of sinners by faith alone in him.

'''5. Salvation by grace, and the duty of all who hear the gospel to believe in Christ.

'''6. The resurrection of the dead and the final judgment, when the wicked 'shall go away into everlasting punishment, but the righteous into life eternal.'''

'The Secretary of our Chapel Building Society informs me that ``one reason for the disuse of the Declaration may be its length, and the circumstance that, to put it beyond question that document is meant, it has been thought it would be needful to embody it in the deed, which would add to the cost.''

'It has been remarked, on the authority of one already cited, ``that, notwithstanding the absence of tests, there is among Independents a marked uniformity of opinion on all important points.'' Perhaps this statement, still true on the whole, would require more qualification than it did some years ago. There are among us a few men of mental vigor who have departed very considerably from the published creeds of Congregationalism. There may be a larger number whose opinions are of an Arminian cast; but, again to use language supplied by a friend, in whom I place confidence as to this subject: ``It would still be fair, I think, to describe our ministry as moderately Calvinistic. An immense majority of the ministers are so. An impression to the contrary has, I am aware, become prevalent; but that is owing, I believe, to the fact that the greater number of the men who have departed from the Calvinistic type hold prominent positions, and have 'the habit of the pen.''' It is a difficult and delicate task to report the state of large religious communities among whose members there exist some diversities of opinion. One person biased by his own predilections will give one account, and another person under an influence of the same kind will give another.

'In what I have said I have endeavored to be as impartial as possible; and, to give the more weight to my statements, I have sought the assistance of official brethren who have wider means of information than I possess, and who may look at things from points of view not exactly identical with my own.'

\xsection{American Congregational Creeds.}

§ 103. American Congregational Creeds.

Literature.

Special essays relating to the creeds and Church order of American Congregationalists.

\emph{The Formation of Creeds.} Article by the Rev. Joseph P. Thompson in the 'New-Englander,' Vol. IV. pp. 265–274. 1846.

\emph{Congregationalism and Symbolism.} Article by the Rev. Wm. G. T. Shedd in the 'Bibliotheca Sacra,' Vol. XV. pp. 661–690. 1858. (An argument showing the need of a more positive creed for Congregationalism.)

\emph{Confessions of Faith.} Article by the Rev. Edward W. Gilman in the 'Congregational Quarterly,' Vol. IV. pp. 179–191. 1862.

\emph{Declaration of Faith and the Confession.} Article by the Rev. Edward A. Lawrence. Ib. Vol. VIII. pp. 173–190. 1866.

\emph{Ancient Confessions of Faith and Family Covenants.} By E. W. G. Ib. Vol. XI. pp. 516–527. 1869.

\emph{The National Council} (of 1871). Article by Dr. A. H. Quint in the 'Cong. Quarterly,' Vol. XIV. pp. 61, 80. 1872.

The American Congregationalists have from time to time adopted the Westminster standards of doctrine, with the exception of the sections relating to synodical Church government. Formerly the Assembly's Shorter Catechism was taught in all the schools of New England; but of late years those standards have gone much out of use, though they have never been disowned.

THE SYNOD OF CAMBRIDGE, 1648.\footnote{'The Congregational Order' above quoted contains the Cambridge Platform and the Saybrook Platform, together with the 'Saybrook Confession of Faith,' \emph{i.e.}, the Savoy Confession as previously adopted by the Synod of Boston.}

The 'Elders and Messengers of the churches assembled in the Synod at Cambridge, in New England,' in June, 1648, adopted the Westminster Confession one year after its publication, in these words: 'This Synod having perused and considered with much gladness of heart, and thankfulness to God, the Confession of Faith published of late by the reverend Assembly in England, do judge it to be very holy, orthodox, and judicious in all matters of faith; and do therefore freely and fully consent thereunto, for the substance thereof. Only in those things which have respect to Church government and discipline [in some sections of Chaps. XXV., XXX., and XXXI.] we refer ourselves to the Platform of Church Discipline agreed upon by this present assembly; and do therefore think it meet that this Confession of Faith should be commended to the churches of Christ among us, and to the honored court, as worthy of their consideration and acceptance. Howbeit, we may not conceal, that the doctrine of \emph{vocation}, expressed in Chap. X., § 1, and summarily repeated in Chap. XIII., § 1, passed not without some debate. Yet considering that the term of \emph{vocation} and others by which it is described are capable of a large or more strict sense or use, and that it is not intended to bind apprehensions precisely in point of order or method, there hath been a general condescendency thereunto. Now by this our professed consent and free concurrence with them in all the doctrinals of religion, we hope it may appear to the world that as we are a remnant of the people of the same nation with them, so we are professors of the same common faith, and fellow-heirs of the same common salvation.'

The Cambridge Synod thus anticipated by ten years the work of the Savoy Conference (1658).

The Cambridge Platform, which is said to be the work of the Rev. Richard Mather, sets forth in substance the same principles of independent Church government and discipline as the Savoy Declaration.

THE SYNOD OF BOSTON, 1680.

The Synod of Elders and Messengers of the New England Congregational churches, held in Boston, Mass., May 12, 1680, adopted and published the Savoy recension of the Westminster Confession, together with the Cambridge Platform. It says, in the preface to its Declaration:

'That which was consented unto by the Elders and Messengers of the Congregational churches in England, who met at the Savoy (being for the most part, some small variations excepted, the same with that which was agreed upon first by the Assembly at Westminster, and was approved of by the Synod at Cambridge, in New England, anno 1648, as also by a General Assembly in Scotland), was twice publicly read, examined, and approved of: that little variation which we have made from the one, in compliance with the other, may be seen by those who please to compare them. But we have (for the main) chosen to express ourselves in the words of those reverend Assemblies, that so we might not only with one heart, but with one mouth, glorify God and our Lord Jesus Christ.'\footnote{The changes are very slight, and in part restorations of the Westminster text. They are noted by Dr. Quint in the 'Congregational Quarterly' for July, 1866, p. 266.}

THE SYNOD OF SAYBROOK, 1708.

The Elders and Messengers of the churches in the Colony of Connecticut assembled at Saybrook, Sept. 9, 1708, agreed that the Boston Confession should 'be recommended to the honorable general assembly of this Colony, at the next session, for their public testimony thereunto, as the faith of the churches of the Colony.' They also accepted 'the Heads of Agreement assented to [in 1692] by the united ministers [of England], formerly called Presbyterian and Congregational,' and so virtually gave indorsement to three creeds as essentially teaching the same system—the doctrinal part of the Articles of the Church of England, the Westminster Confession or Catechisms, and the Confession agreed on at the Savoy.

THE NATIONAL COUNCIL OF BOSTON, 1865.\footnote{\emph{Debates and Proceedings of the National Council of the Congregational Churches, held at Boston, Mass., June} 14–24, 1865. \emph{From the Phonographic Report by J. M. W. Yerrinton and Henry M. Parkhurst.} Boston, Amer. Cong. Association, 1866 (ed. under the care of the Rev. A. H. Quint and the Rev. Isaac P. Langworthy), pp. 95–98, 344–347, 361–363, 401–403.}

The National Council of Congregational churches of the United States, held in the Old South Meeting-house of the city of Boston after the close of the Civil War (which suggested this Council), in the year 1865 (June 14–24), adopted a 'Declaration of Faith.' This Declaration passed through three transformations:

The first draft was prepared by a committee consisting of three divines (two progressive, one conservative), viz., Dr. Joseph P. Thompson (then Pastor of the Church of the Tabernacle, New York), Dr. Edward A. Lawrence (Prof. in the Theol. Seminary of East Windsor [now at Hartford], Conn.), and Dr. George P. Fisher (Prof. of Ecclesiastical History in Yale College). The Committee declined to give a formulated statement of doctrines, but characterized, in a comprehensive way, the doctrines held in common by the Congregational churches, and referred to the ancient Confessions of Westminster and Savoy, as sufficiently answering the end of a substantial unity in doctrine. This draft was read, discussed, and referred to a larger committee.

The second draft was presented by the Rev. J. O. Fiske, of Bath, Maine, and in conformity with the usage of the councils at Cambridge, 1648, at Boston, 1680, and at Saybrook, 1708, expresses adherence to the Westminster and Savoy Confessions for 'substance of doctrine' and the system of truths commonly known as 'Calvinism,' and emphasizes in opposition to modern infidelity the doctrine of the Trinity, the incarnation, the atonement, and other fundamental articles of the common Christian faith.

The third draft was read by the Rev. Alonzo H. Quint, by direction of the business committee, at a meeting of the Council held June 22d, on Burial Hill, Plymouth, on the spot where the first meeting-house of the Pilgrims stood, and which Dr. Bacon declared to be to Congregationalists 'the holiest spot of all the earth.' This paper was substantially approved and referred to a committee of revision to improve the form. This committee reported, Friday, June 23, through the Rev. Dr. Stearns, President of Amherst College, a number of slight verbal alterations. In this improved form the Declaration was twice read 'in a distinct and impressive manner,' and after prayer by the Rev. Dr. Ray Palmer, of New York, unanimously adopted by rising. The singing of Dr. Palmer's well-known hymn, 'My faith looks up to thee,' and the old doxology, 'To God the Father, God the Son,' concluded the solemnity.\footnote{The Boston Declaration is printed in Vol. III. p. 734.}

The same Council adopted a new Platform of Discipline, called the Boston Platform of 1865, and published by the Congregational Board. This virtually supersedes the Cambridge and Saybrook Platforms.

THE OBERLIN NATIONAL COUNCIL, 1871.

The Oberlin Council of 1871 is the first of a regular triennial series of National Councils of the Congregational churches in the United States.\footnote{Formerly General Councils or Synods were held only occasionally (1637, 1646, 1648, 1662, 1680, 1708, 1852, 1865), when some controversy or matter of special concern to all the churches seemed to justify them.} It adopted a constitution, one paragraph of which briefly refers to the rule of faith in a very general way.\footnote{Printed in Vol. III. p. 737.}

Note.—Besides the creeds of General Councils, there are in use among American Congregationalists a great number and variety of creeds, concerning which the Rev. Edward W. Gilman, D.D. (Secretary of the American Bible Society) kindly furnishes the following information:

'1. State Associations and Conferences..

'The usage is various. The General Association of Massachusetts, founded in 1803, accepts as a basis of union ``the doctrines of Christianity as they are generally expressed in the Assembly's Shorter Catechism.'' So do the General Convention of Vermont, founded 1796, and the General Association of New Hampshire, founded 1747. The General Association of New York, founded 1834, has separate Articles of Faith. So has the General Association of Illinois. The General Conferences of Maine and Connecticut have no express doctrinal basis.

'2. County Consociations (of twenty or thirty churches).

' The Lincoln and Kennebec Consociation (Maine), 1808, recommended to its constituent churches Articles of ``Union, Faith, and Practice.'' The Northwestern Consociation (Vermont), 1818, recommended to its churches a uniform Confession and Covenant. The Litchfield South Consociation (Conn.), 1828, prepared a Confession and Covenant for the general use of its churches. The New Haven West Consociation (Conn.) admits only churches which accept the doctrinal part of the Saybrook Platform.

'3. Institutions of Learning.

'The Hollis Professor of Divinity in Harvard College must ``declare it as his belief that the Scriptures of the Old and New Testament are the only perfect rule of faith and practice,'' and the first incumbent (1722), being examined by the Corporation, declared his assent to the Confession of Faith in the Assembly's Catechism and to the doctrinal Articles of the Church of England. Assent to the Westminster Confession or the Saybrook Platform was required of Professors in Yale College from 1753 to 1823. In the Theological Institution at Andover both Visitors and Professors are required to subscribe a Declaration of Faith drawn up by the founders in 1808, and to renew this declaration at intervals of five years.

'4. Local Churches.

'The types are various, and while each church is at liberty to construct and alter its own formulas, certain tendencies towards uniformity of usage, at different periods, are noticeable.

'(a) Individual Professions. Such were those made by John Cotton, at Charlestown, in 1630, and by John Davenport, at New Haven, in 1639. (See the latter in \emph{Ancient Waymarks}, published at New Haven in 1853. See also \emph{Cong. Quarterly}, 1869, Vol. XI. p. 517.)

'(b) Brief general references, either to the holy Scriptures as the only rule of belief and duty, or to the Westminster Catechism or the Boston (\emph{i.e.}, Savoy) Confession, as agreeable to the Scriptures. This usage came in at an early day, and was current at the beginning of this century.

'(c) Articles of Faith, embracing in theological phraseology the outlines of a system of divinity. After the year 1800 these came into general use as formulas for the reception of members, and great reliance was placed upon them as helps in maintaining the purity of the churches against the inroads of false doctrine. Candidates for admission to Church privileges were required to give their assent to the several propositions, which thus in many cases were made tests of worthiness. Dr. Samuel Worcester (Fitchburg, 1798) and Dr. E. D. Griffin, the first pastor of the Park Street Church, Boston (1811), had much to do in shaping the practice of the churches from their day to the present time. Formulas of this class have, however, been subjected to various modifications, by way of accommodation to individual opinions, or for the sake of denying current error, or of emphasizing truths peculiar to the Calvinistic system, but especially in order to secure brevity in the Church service. In this way it has unfortunately sometimes happened that doctrines fundamental to Christianity have failed to find a place in the formal Confession of Faith.

'(d) Creeds divested of theological terms, and clothed in language so clear and simple and general as to prevent no Christian from giving them his prompt and hearty assent. The revisions of the last twenty years have been looking in this direction, and churches are beginning to be formed with no other symbol of faith than the Apostles' Creed.'

\xsection{The Anabaptists and Mennonites.}

§ 104. The Anabaptists and Mennonites.

Literature.

I. On the Anabaptists.

The writings of Luther, Melanchthon, Zwingli, Calvin, Bullinger, and other Reformers and older divines against the Anabaptists are polemical.

H. W. Erbkam:  \emph{Geschichte der Protest. Sekten im Zeitalter der Reformation.} Hamburg und Gotha, 1848, pp. 479 sqq.

Cornelius:  \emph{Geschichte des Münsterischen Aufruhrs.} Leipz. 1856 and 1860, 2 vols.

Karl Hase:  \emph{Das Reich der Wiedertäufer. Neue Propheten.} 2d edition, Leipz. 1860.

Bouterweck:  \emph{Zur Liter. und Geschichte der Wiedertäufer.} Bonn, 1865.

Gerh. Uhlhorn:  \emph{Die Wiedertäufer in Münster}, in his \emph{Vermischte Vorträge.} Stuttgart, 1875, pp. 193 sqq.

Comp. also  Schreiber's  \emph{Biography of Hübmaier, }in his \emph{Taschenbuch f. Geschichte und Alterthum in Süddeutschland}, 1839 and 1840.

II. On the Mennonites.

Menno Simons:  \emph{Fundamentum}, 1539, 1558, etc.; \emph{Opera}, Amst. 1646, 4to; \emph{Opera omnia theologica}, Amst. 1681, in 1 vol. fol. (Both editions in Dutch.)

Herm. Schyn:  \emph{Historia Christianorum qui in Belgio fœderato Mennonitæ appellantur.} Amst. 1723. The same in Dutch, with additions by Gerardus Maatschoen, Amst. 1743–1745, in 3 vols. 12mo. By the same: \emph{Histor. Mennonit. plenior Deductio.} 1729.

S. Blaupot Ten Cate:  \emph{Geschiedenis der Doopsgezinden.} Amsterdam, 1839–47. 5 vols. 8vo.

Cramer:  \emph{The Life of Menno Sim.} Amst. 1837 (Dutch)

Harder:  \emph{Leben Menno Simons.} Königsberg, 1846.

Roosen:  \emph{Menno Simons.} Leipz. 1848.

Erbkam:  \emph{Geschichte der Protest. Sekten}, pp. 480, 571.

Gieseler:  \emph{Kirchengeschichte}, Vol. III. Part II. pp. 90 sqq.

Henke:  \emph{Neuere Kirchengeschichte} (\emph{herausgegeben von Dr. Gass}). Halle, 1874, Vol. I. pp. 414 sqq.

The various branches of the Baptist family of Christians\footnote{Mennonites, Calvinistic Baptists, Arminian Baptists, Dunkers, River Brethren, Seventh-Day Baptists, Six-Principle Baptists, Disciples or Campbellites. The last are very numerous in the West; they reject all creeds on principle.} differ very widely, and have little or no connection except that they agree in rejecting infant baptism and in requiring a personal and voluntary profession of faith in Christ as a necessary condition of baptism. Most of them agree also in opposition to sprinkling, or any other mode of baptism but that by total immersion of the body in water. The largest and most respectable denomination of Baptists took its rise in the great religious commotion of England during the seventeenth century, and differed from the Puritans only in the doctrine of baptism and in the steadfast advocacy of religious freedom. But the Baptist movement began a century earlier on the Continent, and this first stage must at least be briefly noticed.

THE ANABAPTISTS.

The early history of the Anabaptists exhibits a strange chaos of peaceful reforms and violent revolutions—separatism, mysticism, millenarianism, spiritualism, contempt of history, ascetic rigor, fanaticism, communism, and some novel speculations concerning the body of Christ as being directly created by God, and different from the flesh and blood of other men. An impartial history, with a careful critical sifting of these incongruous elements, is still a desideratum.

The modern Anabaptists\footnote{Or Rebaptizers, so called by their opponents because they rebaptized those baptized in infancy, while they themselves denied the validity of infant baptism (some of them \emph{Catholic} baptism in general), and regarded voluntary baptism in years of discretion as the only true baptism. The ancient Anabaptists or Rebaptizers, headed by Cyprian, denied the validity of \emph{heretical} baptism, and carried the principle of Catholic exclusivism to a logical extreme, which the Roman Church has always rejected.} figure prominently in the history of the Reformation, and meet us in Germany, Switzerland, Holland, and England. They were Protestant radicals, who rejected infant baptism as an invention of the Roman Antichrist, and aimed at a thorough reconstruction of the Church. They spread mostly among the laboring classes. Some of their preachers had no regular education, despised human learning, and relied on direct inspiration; but others were learned and eloquent men, as Grebel, Manz, Hetzer, Hübmaier, Denk, Röublin, and Rothmann. They were regarded as a set of dangerous fanatics, who could not be tolerated under a Christian government. Their supposed or real connection with the Peasant War, against the tyranny of landholders (1524), and with the bloody and disastrous excesses at Münster (1534), increased the opposition. Their doctrines were condemned in the Lutheran and Reformed Confessions. The Reformers, even the mildest among them (Melanchthon, Bucer, Bullinger, and Cranmer, as well as Luther, Zwingli, and Calvin), felt that their extermination was necessary for the salvation of the churchly Reformation and social order. And yet they must have known worthy men among them; Calvin himself married the widow of an Anabaptist pastor. Protestant and Roman Catholic magistrates vied with each other in cruelty against them, and put them to death by drowning, hanging, and burning.

But it is the greatest injustice to make the Anabaptists as such responsible for the extravagances that led to the tragedy at Münster. Their original and final tendencies were orderly and peaceful. They disowned the wild fanaticism of Thomas Münzer, John Bockelsohn, and Knipperdolling. They were opposed to war and violence. They were the crude harbingers and martyrs of some truths which have germinated in other ages. They upheld the necessity of discipline and congregational organization on the basis of personal faith in Christ, instead of carnal descent and parochial boundaries. They attacked the doctrine of the eternal damnation of unbaptized infants, and the equally horrible doctrine of persecution. Balthasar Hübmaier (Hübmör, or, as he was called by a Latin name, Pacimontanus), the ablest and most learned among the Anabaptists, a pupil of Dr. Eck (Luther's opponent), and for some time Professor of Catholic Theology at Ingolstadt, then a zealous and eloquent Protestant preacher, was perhaps the first who taught the principle of universal religious liberty, on the ground that Christ came not to kill and to burn, but to save, and condemned the employment of force in his kingdom. He held that those only are heretics who willfully and wickedly oppose the holy Scriptures; and even these ought to be treated by no other than moral means of persuasion and instruction.\footnote{\emph{Von Ketzern und ihren Verbrennern.} A very rare book.} He was burned at the stake in Vienna, March 10,1528, and died with pious joy; his wife, who encouraged him in his martyr spirit, was three days afterwards drowned in the Danube.

THE MENNONITES.

Menno Simons, a converted Roman Catholic priest, collected the scattered remnant of the Anabaptists into a well-organized, peaceful, and industrious community in Holland and on the borders of Germany (1536). He gave them a strict system of discipline, and endeavored to revive the idea of a pure apostolic congregation consisting of true believers unmixed with the world. He labored in constant peril of life with untiring patience till his death, Jan. 13, 1561. 'For eighteen years,' he says, 'with my poor feeble wife and little children, has it behooved me to bear great and various anxieties, sufferings, griefs, afflictions, miseries, and persecutions, and in every place to find a bare existence, in fear and danger of my life. While some preachers are reclining on their soft beds and downy pillows, we oft are hidden in the caves of the earth; while they are celebrating the nuptial or natal days of their children with feasts and pipes, and rejoicing with the timbrel and the harp, we are looking anxiously about, fearing the barking of the dogs, lest persecutors should be suddenly at the door; while they are saluted by all around as doctors, masters, lords, we are compelled to hear ourselves called Anabaptists, ale-house preachers, seducers, heretics, and to be hailed in the devil's name. In a word, while they for their ministry are remunerated with annual stipends and prosperous days, our wages are the fire, the sword, the death.'\footnote{Schyn, \emph{Plenior Deduct}, p. 133 (quoted in Introd. to Baptist \emph{Tracts on Liberty of Conscience}, p. lxxxii.).}

His followers were called \emph{Mennonites} after his death.\footnote{Or \emph{Doopsgezinden}, \emph{i.e.}, Dippers. In Menno's writings they are called \emph{Gemeente Gods, ellendige, weerloze Christenen, broeders}, etc., but never \emph{Mennonites}, See Gieseler, Vol. III. Pt. II. p. 92.} They acquired at last toleration, first in Holland from Prince William of Orange, 1572, and full liberty in 1626. They spread to the Palatinate, Switzerland, Eastern Prussia, and by emigration to South Russia, Pennsylvania, and other parts of North America. Quite recently several hundred families left their Russian settlements for America because the privilege of exemption from military service was withdrawn. They are a small, quiet, peaceful, industrious, and moral community, like the Quakers. Their historian, Schyn, labors to show that they have no connection whatever with the fanatical and revolutionary Anabaptists of Münster.

The Mennonites were divided during the lifetime of Menno into two parties on questions of discipline: 1, the 'coarse' Mennonites (\emph{die Groben}), or Waterlanders, who were more numerous, and flourished in the Waterland district of North Holland; 2, the 'refined' Mennonites (\emph{die Feinen}), who were chiefly Flemings, Frieslanders, and Germans. The latter adhered to the strict discipline of the founder.

The Mennonites acknowledge 'the Confession of Waterland,' which was drawn up by two of their preachers, John Ris (Haus de Rys) and Lubbert Gerardi (Gerritsz), in the Dutch language.\footnote{Schyn gives a Latin translation, in his \emph{Historia Mennonitarum}, pp. 172–220, under the title, \emph{Præcipuorum Christianæ fidei Articulorum brevis Confessio adornata a Joanne Risio et Lubberto Gerardi.} He calls it also \emph{Mennonitarum Confessio}, or \emph{Formula Consensus inter Waterlandos.} He says the confessions of the other branches of the Mennonites agree with it in all fundamental articles. Winer (\emph{Compar. Darstellung}, etc., pp. 24, 25), gives a list of Mennonite Confessions and Catechisms.}

It consists of forty Articles, and teaches, besides the common doctrines of Protestant orthodoxy, the peculiar views of this community. It rejects oaths (Art. XXXVIII., on the ground of Matt. v. 37 and James v. 10), war (XVIII.), and secular office-holding, because it is not commanded by Christ and is inconsistent with true Christian character; but it enjoins obedience to the civil magistrate as a divine appointment wherever it does not contradict the Word of God and interfere with the dictates of conscience (XXXVII.). The Church consists of the faithful and regenerate men scattered over the earth, under Christ the Lord and King (XXIV.). Infant baptism is rejected as unscriptural (XXXI.); but the Mennonites differ from other Baptists by sprinkling.\footnote{One branch of them, the Collegiants or Rhynsburgers, held, however, to the necessity of immersion. They have recently become extinct, having had among them some men of distinction.} On the Lord's Supper they agree with Zwingli. They admit hereditary sin, but deny its guilt (Art. IV.). They hold to conditional election and universal redemption.\footnote{\par Art. VII. derives sin exclusively from the will of man, and teaches that God predestinated and created all men for salvation (\emph{omnes decrevit et creavit ad salutem}), that he provided the remedy for all, that Christ died for all, and saves all who believe and persevere.  \par{} [Note.—McGlothlin gives as the earliest Anabaptist articles of the sixteenth century two brief Swiss statements of 1527 which bear solely on practical questions. Two of the teachings inculcate communism and that the Lord's Supper be celebrated 'as often as the brethren come together.' The articles of the Moravian Anabaptists forbade the Lord's Supper to persons having property.—Ed.]} But during the Arminian controversy a portion sided with the strict Calvinists. They reject also law-suits, revenge, every kind of violence, and worldly amusements. In many respects they are the forerunners of the Quakers quite as much as of the English and American Baptists.

\xsection{The Regular or Calvinistic Baptists.}

§ 105. The Regular or Calvinistic Baptists.

Literature.

\emph{Confessions of Faith and other Public Documents illustrative of the History of the Baptist Churches of England in the Seventeenth Century. Edited for the Hanserd Knollys Society by }  Edward Bean Underhill. London (Haddon Brothers \& Co.), 1854. Contains reprints of seven Baptist Confessions from 1611 to 1688, the Baptist Catechism of Collins, and several letters and other documents from the early history of Baptists in England.

Thos. Crosby:  \emph{The History of the English Baptists, from the Reformation to the Beginning of the Reign of King George I.} London, 1740. 4 vols. Contains important documents, but also many inaccuracies.

Joseph Ivimey:  \emph{History of the English Baptists, including an Investigation of the History of Baptism in England.} London, 1811–23. In 3 vols. 8vo.

Isaac Backus (d. 1806): \emph{History of New England, with especial Reference to the Baptists.} In 3 vols. A new edition, by David Weston, was published by the Backus Historical Society, Newton Centre, Mass. 1871.

David Benedict (Pastor of the Baptist Church in Pawtucket, R. I.): \emph{A General History of the Baptist Denomination in America and other Parts of the World.} Boston, 1813, in 2 vols.; new edition, New York, 1848, in 1 vol. (970 pp.). A chaos of facts.

Francis Wayland:  \emph{Notes on the Principles and Practices of the Baptist Churches.} New York (Sheldon, Blakeman, \& Co.), 1857.

Sewall S. Cutting:  \emph{Historical Vindications; . . . with Appendices containing Historical Notes and Confessions of Faith.} Boston (Gould \& Lincoln), 1859.

J. M. Cramp:  \emph{Baptist History, from the Foundation of the Christian Church to the Close of the Eighteenth Century.} Philadelphia (American Baptist Publication Society), 1868. For popular use.

J. Jackson Goadby:  \emph{Bye-Paths in Baptist History: A Collection of Interesting, Instructive, and Curious Information, not generally known, concerning the Baptist Denomination.} London, 1874 (pp. 375). Chap. VI. treats of Baptist Confessions of Faith.

\emph{The Baptists and the National Centennial: A Record of Christian Work}, 1776–1876. Edited by  Lemuel Moss, D.D. Philadelphia (Baptist Publication Society), 1876. Contains a chapter on 'Doctrinal History and Position,' by Dr. Pepper, pp. 51 sqq.

William R. Williams:  \emph{Lectures on Baptist History.} Philadelphia, 1877.

The English and American Baptists have inherited some of the principles without the eccentricities and excesses of the Continental Anabaptists and Mennonites.\footnote{Their older scholars claim an origin earlier than the Continental or the English Reformation, going back to the Waldenses and Albigenses, and to the Lollard movement following, in Britain, the labors of Wycliff. The tradition of the Holland Mennonites gave them a Waldensian ancestry. But these points are disputed, and no historical connection can be traced.} They are radical but not revolutionary in politics and religion, and as sober, orderly, peaceful, zealous, and devoted as any other class of Christians. They rose simultaneously in England and America during the Puritan conflict, and have become, next to the Methodists, the strongest denomination in the United States.

The great body of Baptists are called Regular or Particular or Calvinistic Baptists, in distinction from the smaller body of General or Arminian or Free-Will Baptists. They are Calvinists in doctrine and Independents in Church polity, but differ from both in their views on the subjects and mode of baptism. They teach that believers only ought to be baptized, that is, dipped or immersed, on a voluntary confession of their faith. They reject infant baptism as an unscriptural innovation and profanation of the sacrament, since an infant can not hear the gospel, nor repent and make a profession of faith. They believe, however, in the salvation of all children dying before the age of responsibility. Baptism in their system has no regenerative and saving efficacy: it is simply an outward sign of grace already bestowed, a public profession of faith in Christ to the world, and an entrance into the privileges and duties of church membership.\footnote{The Campbellites, or Disciples, differ from the other Baptists by identifying baptismal immersion with regeneration, or teaching a concurrence of both acts.} They also opposed from the start national church establishments, and the union of Church and State, which one of their greatest writers (Robert Hall) calls 'little more than a compact between the priest and the magistrate to betray the liberties of mankind, both civil and religious.' They advocate voluntaryism, and make the doctrine of religions freedom, as an inherent and universal right of man, a part of their creed.

THE BAPTISTS IN ENGLAND.

In England the Baptists were for a long time treated with extreme severity on account of their supposed connection with the fanatical fraction of the German and Dutch Anabaptists. A number of them who had fled from Holland were condemned to death or exiled (1535 and 1539). Latimer speaks, in a sermon before Edward VI., of Anabaptists who were burned to death under Henry VIII., in divers towns, and met their fate 'cheerfully and without any fear.'

Under Edward VI. they became numerous in the south of England, especially in Kent and Essex. Two were burned—a Dutchman, named George van Pare, and an English woman, Joan Boucher, usually called Joan of Kent. These were the only executions for heresy during his reign. The young king reluctantly and with tears yielded to Cranmer, who urged on him from the Mosaic law the duty of punishing blasphemy and fundamental heresy. Joan of Kent, besides rejecting infant baptism, was charged with holding the doctrine of some German and Dutch Anabaptists, that Christ's sinless humanity was not taken 'from the substance of the Virgin Mary,' who was a sinner, but was immediately created by God. She resisted every effort of Cranmer to change her views, and preferred martyrdom (May 2, 1550). Several of the Forty-two Edwardine Articles were directed against the Anabaptists.

Under Elizabeth a congregation of Dutch Anabaptists was discovered in London; twenty-seven members were imprisoned, some recanted, some were banished from the kingdom. The two most obstinate, John Wielmaker and Henry Terwoort, were committed to the flames in Smithfield, July 22, 1575, notwithstanding the petition of John Foxe, the martyrologist, who begged the queen to spare them, not indeed from prison or exile (which he deemed a just punishment for heresy), but from being 'roasted alive in fire and flame,' which was 'a hard thing, and more agreeable to the practice of Romanists than to the custom of Evangelicals.'\footnote{See Foxe's letter to Queen Elizabeth, in Latin, in Append. III. to Neal's \emph{History} (Vol. II. p. 439).} These Dutch Anabaptists were charged with 'most damnable and detestable heresies,' such as that Christ took not flesh from the substance of Mary; that infants ought not to be baptized; that it is not lawful for a Christian man to be a magistrate or bear the sword or take an oath. These are evidently doctrines of the Mennonites, afterwards adopted by the Quakers, and now generally tolerated without any injury to society.

During the reigns of James and Charles the Baptists made common cause with the Puritans, especially the Independents, against the prelatical Church, but withdrew more completely from the national worship, and secretly assembled in woods, stables, and barns for religious worship. They began to organize separate congregations (1633), but were punished whenever discovered. Many fled to Holland, and some to America. Their earliest publications were pleas for liberty of conscience.\footnote{See the \emph{Tracts on Liberty of Conscience}, republished for the Hanserd Knollys Society by E. B. Underhill (London, 1846), which contains seven Baptist works on this subject from 1614 to 1661. On Roger Williams, see below.}

With the Long Parliament they acquired a little freedom, though their views were opposed by Presbyterians and Independents, as well as by Episcopalians. They increased rapidly during the civil wars. In 1644 they numbered seven congregations in London, and forty-seven in the country. Cromwell left them unmolested. He had many of them in his army, and some even held positions in his experimental Broad Church.\footnote{Samuel Richardson, a Baptist, who knew him personally, speaks very highly of Cromwell, as a man who 'aimeth at the general good of the nation and just liberty of every man, who is faithful to the saints, who hath owned the poor despised people of God, and advanced many to a better way and means of living.' See \emph{Tracts on Liberty of Conscience}, p. 240.} Milton is claimed by them, on the ground of a passage unfavorable to infant baptism, but with no more justice than Arians, Unitarians, and Quakers may claim him.\footnote{'Milton, it seems, withdrew at last from all Church organizations, regarding them with equal respect and indifference, except the Romanists, whom he excludes from toleration as idolaters and enemies of toleration. With his illustrious friend, the younger Sir Henry Vane, whom, as understanding the true relations of Church and State, he praises in one of his most beautiful sonnets, he joined the 'Seekers,' a body looking for a more perfect Church yet to come. Roger Williams, the friend of both poet and statesman, joined them in his last years in occupying the same ground. In 1673, the year before his death, Milton published a treatise on 'True Religion, Heresy, Schism, Toleration, and the Best Means against the Growth of Popery,' in which he defines heresy to be 'a religion taken up and believed from the traditions of men and additions to the Word of God.' In this sense Popery is the only or the greatest heresy; its very name, Roman Catholic, a contradiction; one of the Pope's bulls as universal particular, or catholic schismatic; while Protestants are free from heresy, which is in the will and choice professedly against the Scriptures. He represents four classes of Protestants—Lutherans, Calvinists, Anabaptists, and Socinians—as agreed in the articles essential to salvation, and says: 'The Lutheran holds consubstantiation; an error, indeed, but not mortal. The Calvinist is taxed with predestination, and to make God the author of sin, not with any dishonorable thought of God, but it may be overzealously asserting his absolute power, not without plea of Scripture. The Anabaptist is accused of denying infants their right to baptism; again, they say they deny nothing but what Scripture denies them. The Arian and Socinian are charged to dispute against the Trinity; they affirm to believe the Father, Son, and Holy Ghost according to Scripture and the Apostolic Creed. As for terms of \emph{trinity, trini-unity, co-essentiality, tri-personality}, and the like, they reject them as scholastic notions, not to be found in Scripture, which, by a general Protestant maxim, is plain and perspicuous abundantly to explain its own meaning in the properest words belonging to so high a matter and so necessary to be known; a mystery indeed in their sophistic subtleties, but in Scripture a plain doctrine. Their other opinions are of less moment. They dispute the satisfaction of Christ, or rather the word \emph{satisfaction}, as not Scriptural, but they acknowledge him both God and their Saviour. The Arminian, lastly, is condemned for setting up free-will against free-grace; but that imputation he disclaims in all his writings, and grounds himself largely upon Scripture only.'}

After the Restoration they were again persecuted by fines, imprisonment, and torture. They suffered more severely than any other Non-conformists, except the Quakers. Among their most distinguished confessors, who spent much time in prison, were Vavasor Powell (d. 1670), Hanserd Knollys (d. 1690),\footnote{Knollys fled to Massachusetts (1638), and preached for some time in the extreme northern part of the colony, but, being exposed to danger as a Baptist and Separatist, he returned to England in 1641. The society for the republication of scarce old Baptist tracts is called after him.} Benjamin Keach, and John Bunyan (d. 1688).

The Act of Toleration (1689) brought relief to the Baptists, and enabled them to build chapels and spread throughout the country. Since then they have become one of the leading branches of Dissenters in England. They have produced some of the most eminent preachers and authors in the English language, such as John Bunyan, Andrew Fuller, Robert Hall, John Foster, Joseph Angus, C. H. Spurgeon.

ROGER WILLIAMS.

Literature.

See Lives of Roger Williams by  Knowles (1834),  Gammell (1845, 1846, 1854), and  Elton (1852); also  Arnold's  \emph{History of Rhode Island} (1860), Vol. I.;  Palfrey's  \emph{History of New England}, Vols. I. and II.;  Bancroft's  \emph{History of the U. S.}, Vol. I.;  Masson,  \emph{Life of Milton}, Vol. II. pp. 560 sqq., 573 sq.;  Allibone,  \emph{Dict. of Brit. and Amer. Authors}, Vol. III. p. 2747;  Dexter,  \emph{As to Roger Williams and his 'Banishment' from the Massachusetts Plantation} (Boston, 1876);  J. L. Diman,  \emph{Monument to R. W. in Providence} (Providence, 1877).

The works of Williams were republished by the Narragansett Club (First Series, Vol. I., Providence, 1866), and by Underhill for the Hanserd Knollys Society (London, 1848).

In America the Baptists trace their origin chiefly but not exclusively to Roger Williams (b. probably in Wales, 1599,\footnote{The accounts of the year of his birth vary from 1598 to 1606. He was a protégé of the celebrated judge, Sir Edward Coke. Historians differ as to whether he was \emph{Rodericus} Williams, from Wales, who entered Jesus College, Oxford, in 1624, or \emph{Rogerus} Williams, whose name appears in the subscription-book of Pembroke College, Cambridge, in 1626. Elton and Masson take the former, Arnold and Dexter the latter view, which better agrees with his Christian name.} d. in Providence, R. I., 1683), the founder of Rhode Island. Originally a clergyman in the Church of England, he became a rigid separatist, a radical come-outer of all Church establishments, an 'arch-individualist,' and an advocate of 'soul-liberty' in the widest acceptation of the term. He was a pious, zealous, unselfish, kind-hearted, but eccentric, 'conscientiously contentions,' and impracticable genius, a real troubler in Israel, who could not get along with any body but himself; and this accounts for his troubles, which, however, were overruled for good. Cotton Mather compared him to a windmill, which, by its rapid motion in consequence of a violent storm, became so intensely heated that it took fire and endangered the whole town.

Pursued out of his land by Bishop Laud, as he says, he emigrated with a heavy heart, in company with his wife Mary, to the colony of Massachusetts, and arrived after a tedious and tempestuous voyage in February, 1631.

He first exercised his ministerial gifts as an assistant to the pastor of Plymouth Colony, and acquired a knowledge of the Indian language. In 1633 he removed to Salem as assistant of Mr. Skelton, and in 1635 he was ordained pastor of Salem Church. But he was even then in open opposition to the prevailing views and customs of the colony, and refused to take the oath of fidelity. Besides this, he was charged with advocating certain opinions supposed to be dangerous, viz., that the magistrate ought not to punish offenses against the first table; that an oath ought not to be tendered to an unregenerate man; that a regenerate man ought not to pray with the unregenerate, though it be his wife or child; that a man ought not to give thanks after the sacrament nor after meat. He was unwilling to retract, and advised his church to withdraw from communion with the other churches of the colony, 'as full of anti-Christian pollution.' For these reasons the court banished Williams (Oct., 1635). The question of toleration was implied in the first charge; he denied the jurisdiction of the civil magistrate over matters of conscience and religion, and defended this principle afterwards in a book, 'The Bloudy Tenent of Persecution for Cause of Conscience,' against John Cotton (1644).\footnote{This book was anonymously published in London, when Williams was there occupied in obtaining a charter for Rhode Island, and is exceedingly rare, only six copies being known to exist; but it has been reprinted from the copy in the Bodleian Library by Edward Bean Underhill, together with the Answer to Cotton's Letter and a Memoir of Williams (London, 1848, pp. 439 and xxxvi.). It is written in a kindly and moderate spirit, free from the controversial bitterness of the age, in the form of a conference between Truth and Peace. Williams begins with this sentence: 'The blood of so many hundred thousand souls of Protestants and Papists, spilt in the wars of present and former ages, for their respective consciences, is not required nor accepted by Jesus Christ, the Prince of Peace.' He maintains that civil government has nothing whatever to do with spiritual matters, over which God alone rules, and that religious liberty should be extended not only to all Christian denominations and sects, but even to 'the most Paganish, Jewish, Turkish, or anti-Christian consciences and worships' (p. 2). John Cotton, his chief opponent, wrote in reply 'The Bloudy Tenent washed, and made white in the Bloud of the Lambe: being discussed and discharged of blood-guiltiness by just Defense' (London, 1647). Williams defended his position in 'The Bloody Tenent yet more Bloody by Mr. Cotton's endeavour to wash it white in the Blood of the Lambe' (London, 1652, 4to, pp. 373). John Cotton (1585–1652), who emigrated to America two years after Williams (1633), was one of the patriarchs of New England, and, together with Hooker and Stone, constituted the 'glorious triumvirate' that supplied the Puritans in the wilderness with their three great necessities—'\emph{Cotton} for their clothing, \emph{Hooker} for their fishing, and \emph{Stone} for their building.'—Cotton Mather's \emph{Magnalia}, Vol. III. p. 20.} His views on baptism were developed afterwards; but they would only have aggravated his case, and in fact his rebaptism brought upon him the sentence of excommunication from the church of Salem, of which he was still nominally a member.\footnote{Dr. Dexter's monograph is a learned and elaborate partisan defense of the action of the young Colony, which, he says, 'was reluctantly compelled to choose between the expulsion of Williams and the immediate risk of social, civil, and religious disorganization' (p. 88). He takes the ground that Williams was banished, not on religious, but on political grounds. But religion and politics were inseparably interwoven in the New England theocracy.}

The banishment was the best thing that could have happened to Williams: it led to the development of his heroic qualities, and gave him a prominent position in American history. He left Salem with a few friends, and made his way in dreary winter through 'a howling wilderness' to the wigwams of his Indian friends, and was sorely tossed in frost and snow among barbarians for fourteen weeks, 'not knowing what bread or bed did mean.' In June, 1636, he founded with five families who adhered to him the town of Providence. He scrupulously bought the land from the Indians, and acted as pastor of this democratic settlement. In 1638 he became a Baptist; he was immersed by Ezekiel Hollyman, and in turn immersed Hollyman and ten others. This was the first Baptist church on the American Continent. But a few months afterwards he renounced his rebaptism on the ground that Hollyman was unbaptized, and therefore unauthorized to administer the rite to him. He remained for the rest of his life a 'Seeker,' cut loose from all existing Church organizations and usages, longing for a true Church of God, but unable to find one on the face of the whole earth. He conceived 'that the apostasy of Antichrist hath so far corrupted all that there can be no recovery out of that apostasy till Christ send forth, new apostles to plant churches anew.'

In 1643 he went to England, and obtained through the Commissioners of Plantation a charter which allowed the planters to rule themselves according to the laws of England, 'so far as the nature of the case would admit.' In 1663 he accepted for the colony another and more successful charter, a patent from the English crown similar to that of Massachusetts, to which he had formerly objected. He kept up friendly relations with the Indians, and twice saved the Massachusetts colony from danger, thus returning good for evil. His fame rests on his advocacy of the sacredness of conscience. Bancroft goes too far when in his eloquent eulogy he calls him 'the first person in modern Christendom who asserted in its plenitude the doctrine of the liberty of conscience, the equality of opinions before the law.' The Anabaptists and Mennonites had done the same a hundred years before. But Williams planted the first civil government on the principle of universal 'soul-liberty,' and was followed by William Penn in his Quaker colony in Pennsylvania. Roger Williams has been called 'that noble confessor of religious liberty, that extraordinary man and most enlightened legislator, who, after suffering persecution from his brethren, persevered, amidst incredible hardships and difficulties, in seeking a place of refuge for the sacred ark of conscience.'\footnote{Mrs. P. S. Elton, in \emph{The Piedmontese Envoy; or, The Men, Manners, and Religion of the Commonwealth: A Tale} (London, 1852), puts this eulogy into the mouth of John Milton; hence it is sometimes falsely quoted as Milton's (Allibone, Vol. III. p. 2747).}

In the other colonies the Baptists were more or less persecuted till the time of the Revolution, but after that they spread with great rapidity.

The American Baptists differ from their English brethren by a stricter discipline and closer communion practice. They are very zealous in missions, education, and other departments of Christian activity. In theology they cultivate especially biblical studies with great success.

BAPTIST CONFESSIONS.

The Baptists, like the Congregationalists, lower the authority of general creeds to mere declarations of faith prevailing at the time in the denomination, to which no one is bound to give assent beyond the pleasure of his conviction; and they multiply the number and elevate the authority of local or congregational creeds and covenants, by which the members of particular congregations voluntarily bind themselves to a certain scheme of doctrine and duty. Notwithstanding the entire absence of centralization in their government, and the unrestrained freedom of private judgment, the Calvinistic Baptists have maintained as great a degree of essential harmony of faith as they themselves deem desirable.

'The Baptist creeds,' says Dr. Joseph Angus, in behalf of English Baptists,\footnote{In a letter to the author.} 'were prepared in the first instance for apologetic and defensive purposes. They merely describe the doctrines held by the bodies from which they emanated. They were never imposed on ministers and members of the churches of either section of the Baptists. Even when adopted, as they sometimes were, by any church, as an expression of its sentiments, all sister churches were left free, and in the particular church a considerable latitude of judgment was allowed in interpreting them. They have never been accepted as tests, and merely represent in a general way the sentiment of the body. In trust deeds or in the rules of associations they never appear. Property in trust is held for the use of evangelical Christians maintaining the doctrines commonly held by Particular (or General) Baptists; sometimes these doctrines are enumerated in the briefest possible way—the trinity, the atonement, etc.—and sometimes they are not enumerated at all. Of course, in the event of an appeal to law, the creeds and confessions would be evidence of the faith of the body. Substantially the two sections of the Baptist body believe as of old. But their confessions are not authoritative except as evidence and in matters of property; while in the interpretation of them it is a principle to allow as much freedom as is consistent with a substantial agreement in the same general truth.'

'Confessions of faith,' says Dr. Osgood, with special reference to the Baptists in the United States,\footnote{Letter to the author.} 'have never been held as tests of orthodoxy, as of any authoritative or binding force; they merely reflect the existing harmony of views and the scriptural interpretations of the churches assenting to them. ``We believe,'' says Wayland, ``in the fullest sense, in the independence of every individual church of Christ. We hold that each several church is a Christian society, on which is conferred by Christ the entire power of self-government. No church has any power over any other church. No minister has any authority in any church except that which has called him to be its pastor. Every church, therefore, when it expresses its own belief, expresses the belief of no other than its own members. If several churches understand the Scriptures in the same way, and all unite in the same confession, then this expresses the opinions and belief of those who profess it. It, however, expresses their belief because all of them, from the study of the Scriptures, understand them in the same manner, and not because any tribunal has imposed such interpretations upon them. We can not acknowledge the authority of any such tribunal. We have no right to delegate such an authority to any man or to any body of men. It is our essential belief that the Scriptures are a revelation from God, given . . . to every individual man. They were given to every individual that he might understand them for himself, and the word that is given him will judge him at the great day. It is hence evident that we can have no standards which claim to be of any authority over us.'''\footnote{F. Wayland, \emph{Principles and Practices of Baptist Churches}, pp. 13, 14.}

I. The Confession of the Seven Churches in London. Dr. Daniel Featley, a prominent Episcopalian of the Puritan party and member of the Westminster Assembly (from which, however, he was expelled for informing the king of its proceedings), had a public disputation with the Baptists in 1644, and published it, with a dedication to the Parliament, under the title, 'The Dippers dipt; or, the Anabaptists Duck'd and Plung'd over Head and Ears at a Disputation in Southwark.'\footnote{London, 3d ed. 1645; 7th ed. 1660. The spirit of this book may be judged from the title and the following passage of the \emph{Epistle Dedicatory:} 'Of all heretics and schismatics, the Anabaptists ought to be most carefully looked into, and severely punished, if not utterly exterminated and banished out of the Church and Kingdom. . . . They preach and print and practice their heretical impieties openly; they hold their conventicles weekly in our chief cities and suburbs thereof, and there prophesy by turns; . . . they flock in great multitudes to their Jordans, and both sexes enter into the river, and are dipt after their manner with a kind of spell, containing the heads of their erroneous tenets. . . , And as they defile our rivers with their impure washings, and our pulpits with their false prophecies and fanatical enthusiasms, so the presses sweat and groan under the load of their blasphemies.'}

This gave rise to a Confession of Faith, on the part of seven London churches, with an Epistle Dedicatory to the two houses of Parliament. It appeared in 1644 (three years before the Westminster Confession), and again with some additions and changes in 1646, under the title, 'A Confession of Faith of Seven Congregations or Churches of Christ in London, which are commonly (but unjustly) called Anabaptists.'\footnote{Printed in Underhill's Collection, pp. 11–48. The title-pages, which are all given by Underhill, slightly differ in the three editions of 1644, '46, and '51. I have before me a copy of the fourth ed., London, 1652, which has been for more than two hundred years in the family of the Rev. Dr. Holme, a Baptist clergyman of New York. It has the same title as the third ed., but, only fifty-one Articles; Art. XXXVIII., on the support of the ministry by the congregation, being omitted.} This document consists of fifty-two (51) Articles, and shows that in all important doctrines and principles, except on the sacraments and Church government, the Baptists agreed with the orthodox Reformed Churches. The concluding paragraph admits the fallibility of human confessions, and the readiness of Baptists to receive further light, but also their determination 'to die a thousand deaths rather than do any thing against the least tittle of the truth of God, or against the light of our own consciences.'

II. The Confession of Somerset, 1656. It was signed by the delegates of sixteen churches of Somerset and the adjoining counties. It consists of forty-six Articles.\footnote{Underhill, pp. 74–106.}

III. The Confession of 1688. This is by far the most important and authoritative. It has superseded the two earlier confessions, and is to this day held in the highest esteem. It appeared first in 1677, at London, under the title, 'A Confession of Faith put forth by the Elders and Brethren of many congregations of Christians baptized upon profession of their faith.' It was reprinted in 1688, 1689, and approved and recommended by the ministers and messengers of above a hundred congregations met in London, July 3–11, 1689.\footnote{The following certificate was prefixed: 'We, the ministers and messengers of, and concerned for, upwards of one hundred congregations in England and Wales, denying \emph{Arminianism}, being met together in London, from the third day of the seventh month to the eleventh of the same, 1689, . . . have thought meet for the satisfaction of all other Christians that differ from us in the point of \emph{baptism}, to recommend to their perusal the \emph{confession of our faith}, . . . which confession we own, as containing the doctrine of our faith and practice; and do desire that the members of our churches respectively do furnish themselves therewith.' Signed by thirty-seven persons in the name of the whole assembly.} It has been often reprinted.\footnote{Editions of 1699, 1719, 1720, etc. An American ed. was issued by Benj. Franklin, and one at Pittsburgh (S. Williams), 1831. It is also reprinted by Crosby, Vol. III. Append. II. pp. 56–111; Underhill, pp. 169–246.} 'It is still generally received by all those congregations that hold the doctrine of personal election and the certainty of the saints' final perseverance.'\footnote{Dr. Angus.} In America it was adopted by the Baptist Association which met in Philadelphia, Sept. 25, 1742, and hence is known also by the name of the Philadelphia Confession.

This Confession consists of thirty-two chapters, beginning with the holy Scriptures and ending with the last judgment. It is simply the Baptist recension of the Westminster Confession, as the Savoy Declaration is the Congregational recension of the same Westminster Confession. It follows the Westminster Confession in sentiment and language, with very few verbal alterations, except in the doctrine of the Church and the Sacraments. The Preface sets forth that the Confession of Westminster is retained in substance for the purpose of showing the agreement of the Baptists with the Presbyterians and Congregationalists 'in all the fundamental Articles of the Christian religion,' and also to convince all that they have 'no itch to clog religion with new words, but do readily acquiesce in that form of sound words which has been, in consent with the holy Scriptures, used by others before us; hereby declaring before God, angels, and men our hearty agreement with them in that wholesome Protestant doctrine which with so clear evidence of Scripture they have asserted.' The Appendix is a defense of the Baptist theory against Pædobaptists.

The Confession differs from that of the Westminster in the chapters on the Church and on the sacraments. It omits the chapter 'Of Church Censuses' (XXX.) and 'Of Synods and Councils.' The chapter 'Of the Church' (XXV.) is adapted to the independent polity; and the chapter 'Of Baptism' is altered to suit the Baptist theory, limiting the right of baptism to those 'who do actually profess repentance towards God, faith in and obedience to our Lord Jesus,' and declaring 'immersion or dipping of the person in water' to be 'necessary to the due administration of this ordinance' (XXIX.). A chapter, 'Of the Gospel and the Extent of Grace thereof,' is inserted from the Savoy Declaration as Ch. XX. (which causes the change of the numbering of the chapters which follow).\footnote{See Vol. III. pp. 738 sqq.}

IV. In 1693 a \emph{Catechism }based on this Confession was drawn up by William Collins, at the request of the General Assembly which met at London in June of that year. It is taken chiefly from the Westminster Shorter Catechism, and follows closely its order and method. It is also called 'Keach's Catechism.' Benjamin Keach was with Collins among the signers of the Confession of 1688, and seems to have had much to do with the work. It is the only Catechism which has found general acceptance among Baptists in England and America.\footnote{Underhill says, p. xv.: 'It is the only Catechism of value among Baptists.' He gives it from the 16th Engl. ed., pp. 247–270, but says nothing of Keach's co-authorship, and ascribes to him another Catechism ('The Child's Instructor: a New and Easy Primer,' 24mo, 1664), for which he was imprisoned under Charles II. The American Baptist Publication Society publishes it under the title, 'The Baptist Catechism commonly called Keach's Catechism; or, A Brief Instruction in the Principles of the Christian Religion, agreeably to the Confession of Faith put forth by upwards of a hundred congregations in Great Britain, July 3, 1689, and adopted by the Philadelphia Baptist Association, Sept. 22, 1742.' Here the name of Collins is omitted. But the Catechism is literally the same as the one in Underhill's Collection.}

During the seventeenth century there were also some private confessions written by John Bunyan, Vavasor Powell, Benjamin Keach, and Elias Keach.

V. The New Hampshire Confession was prepared about 1833 or 1834, by the Rev. J. Newton Brown, of New Hampshire (d. 1868), the editor of a 'Universal Cyclopædia of Religious Knowledge.' It is shorter and simpler than the Confession of 1688, and presents the Calvinistic system in a milder form. It has been accepted by the Baptists of New Hampshire and other Northern and Western States, and is now the most popular creed among American Baptists.\footnote{It is printed in Vol. III. pp. 742 sqq.}

\xsection{Arminian or Free-will Baptists.}

§ 106. Aminian or Free-will Baptists.

IN ENGLAND.

[See Literature on p. 845.]

The General or Arminian Baptists differ from the Particular or Calvinistic Baptists in rejecting unconditional election and the perseverance of saints, and in maintaining the freedom of will and the possibility of falling from grace. So far they followed the Mennonites. They assign greater power to a general assembly of associated churches, and hold three orders—bishops or messengers, pastors or elders, and deacons; while the Particular Baptists, like the Congregationalists, recognize only two—bishops or pastors and deacons (elders being a title applicable to the first or to both).

I. The first Confession of Arminian Baptists was published by English refugees in Holland, under the title, 'A Declaration of Faith of English People remaining at Amsterdam in Holland,' Amsterdam, 1611.\footnote{It is reprinted in Crosby's \emph{History}, Vol. II. Appendix I. pp. 1–9, and in Underhill's Collection, pp. 1–10. A manuscript copy exists in the archives of the Mennonite church at Amsterdam, to which the original subscriptions of forty-two names are appended, preceded by the modest remark, 'We subscribe to the truth of these Articles, desiring further instruction.'} It was drawn up by Smyth and Helwisse. It consists of twenty-seven (26) Articles. The first Article confesses the doctrine of the Trinity in the spurious words of 1 John v. 7. Election is conditioned by foreknown faith, reprobation by foreknown unbelief, and the perseverance of saints is denied.\footnote{Art. V.: 'God before the foundation of the world hath predestinated that all that believe in him shall be saved, and all that believe not shall be damned; all which he knew before. And this is the election and reprobation spoken of in the Scriptures, . . . and not that God hath predestinated men to be wicked, and so be damned, but that men being wicked shall be damned.' Art. VII.: 'Men may fall away from the grace of God, and from the truths which they have received and acknowledged.'} The Church of Christ is defined (Art. X.) to be 'a company of faithful people separated from the world by the Word and Spirit of God, being knit unto the Lord, and one unto another, by baptism, upon their own confession of the faith.' Baptism is confined to adults, but nothing is said of immersion. The duty of obedience to the magistrate is very earnestly enjoined (Art. XXIV.).

II. The 'London Confession' was approved by more than twenty thousand Baptists, and was presented to Charles II., July 26, 1660. It contains twenty-five Articles.\footnote{Underhill, pp. 107–120.}

III. The 'Orthodox Creed' was published in 1678, by the General Baptists of Oxfordshire and the parts adjacent. It makes a near approach to Calvinism, with a view to unite the Protestants in the fundamental articles against the errors of Rome.\footnote{Ibid. pp. 121–168.}

IN AMERICA.

Literature.

I. D. Stewart:  \emph{The History of the Free-will Baptists for Half a Century.} Dover, 1862 sqq. (Vol. I. from 1780 to 1830). Comp. also the \emph{Lives} of Randall, Stinchfield, Colby, Thornton, Marks, Bowles, Phinney, and Elias Smith; the \emph{Records} of \emph{Yearly Meetings} and \emph{Quarterly Meetings, }and sundry articles in the religious periodicals and other publications of the Free-will Baptists issued from their Printing Establishment at Dover, New Hampshire.

The American General Baptists are called Free-will Baptists or Free Baptists. They trace their origin to Benjamin Randall (1749–1808), who was converted by one of the last sermons of Whitefield at Portsmouth, New Hampshire, Sept. 28, 1770. He was at first a Congregationalist, but in 1776 he united himself with a regular Baptist church in South Berwick, Maine, and entered the ministry. In 1780 he organized, in New Durham, New Hampshire, a Baptist church, which became the nucleus of a new denomination, holding the doctrines of conditional election, free will, and open communion. In government it is congregational.

In 1827 the Free-will Baptists organized a General Conference in New England, and opened correspondence with the Arminian Baptists in England and North Carolina.

Their Confession of Faith, together with a directory of discipline, was prepared by order of the General Conference of 1832, approved 1834, revised by a committee in 1848, 1865, and 1868. It is the clearest and ablest exposition of the principles of the Free-will Baptists.\footnote{It is published at Dover. N. H., under the title, \emph{Treatise, on the Faith and Practice of the Free-will Baptists}, and forms a little book of fifty pages. The doctrinal part is printed in Vol. III. pp. 749 sqq.}

\xsection{The Society of Friends, or Quakers.}

§ 107. The Society of Friends, or Quakers.

Literature.

I. Sources.

Geo. Fox (founder of the Society of Friends, d. 1690): \emph{Works} (containing his Journal, Letters, and Exhortations), London, 1694–1706, in 3 vols. fol.; also Philadelphia, in 8 vols. 8vo.

Robert Barclay (the standard divine of the Quakers, d. 1690): \emph{Works}, edited by William Penn, London, 1692, under the title, '\emph{Truth Triumphant through the Spiritual Warfare, Christian Labors and Writings of that Able and Faithful Servant of Jesus Christ, Robert Barclay},' etc. The principal of these works are: \emph{Apologia Theologiæ vere Christianæ}, first in Latin, Amst. 1675; then in English, by the author himself; also in German, Dutch, French, and Spanish. The full title of the English edition is, '\emph{An Apology for the True Christian Divinity, being an Explanation and Vindication of the Principles and Doctrines of the People called Quakers.}' (I have a very elegant copy of the eighth edition, Birmingham, 1765.) \emph{A Catechism and Confession of Faith, approved of and agreed unto by the General Assembly of the Patriarchs, Prophets, and Apostles, Christ himself Chief Speaker in and among them.} (The answers wholly biblical.) 1673. The same, in Latin (\emph{Catechesis et Fidei Confessio}, etc.). Rotterdam, 1676. \emph{Treatise on Christian Discipline}, etc.

William Penn (d. 1718): \emph{A Summary of the History, Doctrine, and Discipline of Friends} (London, 1692); \emph{Brief Account of the Rise and Progress of the People called Friends} (London, 1694); '\emph{Quakerism a New Nickname for Old Christianity;}' '\emph{The Great Case of Liberty of Conscience Debated and Defended},' etc. Some of Penn's tracts were translated into German by Seebohm (Pyrmont, 1792 and 1798).

II. historical.

Gerard Croese:  \emph{History of the Quakers, containing the Lives, Tenets, Sufferings, Trials, Speeches, and Letters of all the most Eminent Quakers from the First Rise of the Sect.} London, 1696, 8vo.

William Sewel (d. 1725): \emph{History of the Rise, Increase, and Progress of the Christian People called Quakers.} London, 1725, fol.; 6th edition, 1834, in 2 vols.; also in Dutch and German. (Charles Lamb pronounced this book 'far more edifying and affecting than any thing of Wesley and his colleagues.')

Joseph Besse:  \emph{Collection of the Sufferings of the People called Quakers, for the Testimony of a Good Conscience.} London, 1753, 2 vols. fol.

John Gough:  \emph{The History of the Quakers.} Dublin, 1789, 4 vols. 8vo.

Sam. M. Janney:  \emph{History of the Friends.} Philadelphia, 1859–1867, 4 vols.

Biographies of G. Fox, by  Jonah Marsh (1848),  S. M. Janney (1853),  W. Tallack (1868).

Biographies of W. Penn, by  Marsiliac (1791),  Clarkson (1813),  Ellis (1852),  Janney (1852),  Hepworth Dixon (1856).

III. Explanatory and Apologetic.

Thos. Clarkson (d. 1846): \emph{A Portraiture of Quakerism.} London, 1806; 2d ed. 1807, 3 vols.

Joseph John Gurney (d. 1847): \emph{Observations on the Distinguishing Views and Practices of the Society of Friends.} 7th edition, London, 1834; 2d American from the 7th London edition, New York, 1869.

Thos. Evans:  \emph{An Exposition of the Faith of the Religious Society of Friends.} Philadelphia, 1828. Approved by the Quakers at a meeting held in Philadelphia, Oct. 19, 1827, and often printed. (Manchester edition, 1867.)

\emph{The Ancient Testimony of the Religious Society of Friends,... revived and given forth by the Yearly Meeting held in Philadelphia in the Fourth Month}, 1843. Philadelphia, at Friends' book-store. A summary of orthodox Quakerism, chiefly from the writings of Barclay.

W. I. Allinson: Art. \emph{Friends}, in M'Clintock and Strong's \emph{Cyclop.}, Vol. III. pp. 667 sqq. (New York, 1870).

\emph{Friends' Review, a Religious, Literary, and Miscellaneous Journal.} Philadelphia, so far twenty-nine vols. till 1876 (edited by Henry Hartshorne).

IV. Polemical and Critical.

For a full account of the literature against the Quakers, see Jos. Smith's  \emph{Bibliotheca anti-Quakeriana; or, A Catalogue of Books adverse to the Society of Friends. Alphabetically arranged. With Biographical Notices of the Authors, together with the Answers which ham been given to some of them by Friends and others.} London, 8vo, pp. 474.

Möhler (R. C.): \emph{Symbolik}, pp. 488–532;  Rud. Hofmann:  \emph{Symbolik, }pp. 514–520;  Schenckenburger,  \emph{Lehrbegriffe der kleineren protest. Kirchenparteien}, pp. 69–102.

HISTORICAL SKETCH.

The Religious Society of Friends, as they call themselves—or Quakers, as they are usually called\footnote{The name 'Friends' designates a democratic brotherhood in Christ. The name 'Quakers' is sometimes wrongly derived from the warning of Fox to the magistrates 'to quake for fear' and 'to tremble at the Word of the Lord' (Isa. lxvi. 2). It comes rather from their own tremulous utterance of emotion in prayer and exhortation. Barclay (\emph{Apology}, p. 310, on Prop. XI.) speaks of the trembling motion of the body under the power of the truth, by which Quakers are exercised as in the day of battle, and says: 'From this the name of \emph{Quakers}, i.e., \emph{Tremblers}, was first reproachfully cast upon us; which, though it be none of our choosing, yet in this respect we are not ashamed of it, but have rather reason to rejoice therefore, even that we are sensible of this power that hath oftentimes laid hold of our adversaries and made them yield unto us.' Allinson says (1.c. p. 668): 'The epithet Quakers was given in derision, because they often trembled under an awful sense of the infinite purity and majesty of God, and this name, rather submitted to than accepted by them, has become general as a designation.'}—originated in the Puritan commotion which roused all the religious energies of England.

It was founded by George Fox (1624–1690), one of the oddest saints in Christendom, a self-taught and half-inspired man of genius, who was called by a higher power from the shepherd's staff to the evangelism of the baptism by fire and by the Spirit. In early youth he felt inclined to ascetic retirement, like the hermits of old. He was a thorough mystic, and desired to get at the naked truth without the obstruction of church, sacrament, ceremonies, theology, and ordinary study, except the Scriptures spiritually understood. He loved to commune with nature and nature's God, to walk in the inward light, to enjoy the indwelling Christ, and to receive inspirations from heaven. He spent much time in fasting and prayer, he wrestled with the devil, and passed through deep mental distress, doubt, and despondency. His moral character was beyond reproach—pure, truthful, unworldly, just, temperate, meek, and gentle, yet bold and utterly regardless of conventional usage and propriety. He began his public testimony in his twenty-third year, and traveled through England, Holland, and the American colonies, preaching and praying with pentecostal fervor and power, revealing hidden truths, boldly attacking pride, formality, and worldliness, and exhorting to repentance, humility, and mercy. He sometimes interrupted the clergymen at public service, and the lawyers in court, and warned them against the wrath to come. He was a stern ascetic, clad in leather, and wearing long hair. He addressed every body 'thou' or 'thee,' and sublimely ignored all worldly honors and dignities.\footnote{'The Lord forbade him,' says Sewel, 'to put off his hat to any man, high or low; he was required to \emph{Thou} and \emph{Thee} every man and woman without distinction, and not to bid people \emph{Good-morrow} or \emph{Good-evening;} neither might he bow or scrape his leg to any one.'} He was nine times thrown into prison for breaches of the peace and blasphemy, and suffered much hardship and indignity with imperturbable temper; but towards the close of his meteoric career he enjoyed comparative rest. His 'Journal' gives an account of his labors, and is, in the language of Sir James Mackintosh, 'one of the most extraordinary and instructive narratives in the world.' Fox was providentially provided with the best aid in founding his society.

Robert Barclay (1648–1690) was the apologist and theologian of the Quakers, the only one known to fame. Descended from a noble family in Scotland, and educated in Paris, he became a convert first to Romanism, then to Quakerism (1667). He had therefore the advantage of an experimental as well as theoretical knowledge of the Scotch Calvinistic and the Roman Catholic creeds. He made various missionary journeys in company with William Penn; he walked through the streets of Aberdeen in sackcloth and ashes, and was several times imprisoned, but spent his last years in peace on his estate of Ury.

William Penn (1644–1718), the statesman and politician of the Quakers, and the founder of Pennsylvania, was the son of an admiral, and enjoyed the favor of James II. (his father's friend), which he used in the cause of justice and mercy.\footnote{The charges of Lord Macaulay against Penn's integrity have been repelled by W. E. Forster (\emph{William Penn and Thomas Babington Macaulay}, 1850) and J. Paget (Edinburgh, 1858).} He himself was expelled for his religion from the University of Oxford and his father's house, and was twice imprisoned, but ably defended the liberty of conscience, and was acquitted. By his influence more than twelve hundred Quakers were set at liberty. In 1680 he obtained from the king, in payment of a claim of £16,000, an extensive tract of land west of the Delaware River, and organized a colony on the basis of perfect freedom of religion (1682). The city of Philadelphia, or brotherly love, became the chief asylum of persecuted Quakers, a century afterwards the cradle of American independence, and in 1876 the theatre of the most remarkable centennial ever celebrated by any nation. Penn was twice in America, but died in England. He made a treaty with the Indians, of which Voltaire said that it was the only treaty never sworn to and never broken. The United States government would have fared better with the aborigines of the country if it had followed the humane example of Roger Williams and William Penn.

The Quakers, during the first forty years of their history, were more severely persecuted than any sect of Christians had ever been, with the exception of the Waldenses, and bore it with unflinching heroism. Their eccentricities and fanatical excesses, their utter disregard for the courtesies and conventionalities of civilized life, their fierce abuse of the national churches (or 'steeple-houses') and clergymen, their opposition to tithes, salary, the oath, and military service, provoked the combined hostility of magistrates, ministers, and people. Their places of worship were invaded by the populace armed with staves, cudgels, and pitchforks; the windows broken by stones and bullets; their religious services rudely interrupted by hallooing and railing; their property destroyed or sold; their persons ridiculed, buffeted, assailed with stones and filth, dragged by the hair through the streets, or thrown into loathsome prisons and punished as heretics and blasphemers.

Cromwell, who had a tender feeling for all 'godly' radicals and enthusiasts, was rather pleased with George Fox, with whom he had an interview (1654); he allowed him to keep on his hat, and to speak about the mysteries of spiritual experience; and, although he disapproved his disorderly conduct, he pressed his hand and said, 'Come again to my house; if thou and I were together but an hour in every day, we should be nearer one to the other.' But Cromwell could not control the local magistrates and the rabble.

Under Charles II. the Quakers fared much worse, and notwithstanding the influence of Penn upon James II., who favored them for political reasons in the interest of the Roman Catholics, they continued to suffer until the Act of Toleration, in 1689, or rather until 1696, when by a special Act of Parliament their solemn affirmation was recognized as equivalent to an oath.

During the period from 1650 to 1689, according to the patient researches of their historian, Joseph Besse, no less than 13,258 Quakers suffered fine, imprisonment, torture, and mutilation in England, Scotland, and Ireland, 219 were banished, and 360 perished in prisons, some almost literally rotting in pestilential cells.

In New England they were not treated any better: 170 instances of hard usage are enumerated, 47 were banished, and 4 hanged (three men and one woman, Mary Dyer). In explanation, though not in justification, of this severity of the Puritan colony towards them, we should remember those offenses against public decency which led some Quaker men and women to invade churches during divine service, and to promenade the streets of Boston, Cambridge, and Salem in sackcloth and ashes, even \emph{in puris naturalibus}, for 'a sign and wonder' (in imitation of Isa. xx. 2, 3), to symbolize the 'naked truth,' and to utter a prophetic 'testimony' against the 'hireling priests,' the tyrannical magistrates, and the wicked and perverse generation, warning them of the impending judgments of the Lord, who would come with fire and sword.\footnote{Palfrey, \emph{History of New England}, Vol. II. pp. 461–485; Dexter, \emph{As to Roger Williams},' etc., pp. 124 sqq. One such case of Oriental teaching by signs occurred also in England, and is mentioned by Fox himself in his \emph{Journal}: 'The Lord made one to go naked amongst you, a figure of thy nakedness, and as a sign, before your destruction cometh, that you might see that you were naked and not covered with the truth.' See Stoughton, \emph{The Church of the Commonwealth}, p. 360.} Even Roger Williams, in his debate with the Quakers at Newport (1672), with all his liberality, condemned such conduct.\footnote{He wrote a curious book, \emph{George Fox digg'd out of his Burrowes}, etc., which was republished by the Narragansett Club, 1872, with an introduction by Prof. Diman. Comp. Dexter, 1.c. p. 138.}

Notwithstanding these persecutions, the Society of Friends spread rapidly, and numbered about 70,000 members towards the close of the seventeenth century. They afterwards diminished in England, but increased in America, though not as much as other denominations. On the Continent they had only a few adherents in Holland and Germany.

The fanatical heat of the martyr period of the Quakers cooled down with the cessation of persecution. They became a sober, quiet, orderly, and peaceful community. The oddities which they still retain are perfectly harmless, and form an interesting chapter in the history of morals. Quakerism is not so much a new theology as a new mode of Christian life, representing the utmost simplicity in opposition to show, ornament, and amusement.

QUAKER CONFESSIONS.

The Quakers are more radical than the Independents and the Baptists. They utterly broke with historical Christianity, and reject its visible ordinances, which the Independents and the Baptists retained. They kept aloof from the Puritans, and would have nothing whatever to do with the national English or any other Church or sect in Christendom. They oppose all outward authority in religion, though it be the letter of the Bible itself.

With such views they can not consistently recognize any binding standards of doctrine which might obstruct the freedom of interpretation of the divine Word under the direct illumination of the Spirit.

Nevertheless, with all their radicalism, the Quakers retained the substance of the Christian faith, and, following the example of the early Christians, they set forth their tenets in a number of apologies against the misrepresentations of their enemies. The first 'Confession and Profession of Faith in God' was published by Richard Farnsworth in 1658. Similar apologetic documents followed in 1659 and 1661 by George Fox the Younger, in 1662 by John Crook, in 1664 by William Smith, in 1668 by William Penn, in 1671 by Whitehead and Perm, in 1698 by Penn and others, in 1671, 1675, and 1682 by George Fox.\footnote{On these earlier confessions, see Evans, pp. xii. sqq.}

The ablest and most authoritative exposition of the belief of the Quakers is the 'Apology' of Robert Barclay, written in his quiet retreat in Ury, Scotland, 1675, and addressed to Charles II. It is his most elaborate work, and is still held in the highest estimation by the orthodox Friends. He pays the school-divinity the compliment that, although it takes up almost a man's whole life-time to learn, it 'brings not a whit nearer to God, neither makes any man less wicked or more righteous.' 'Therefore,' he continues, 'hath God laid aside the wise and the learned and the disputers of this world, and hath chosen a few despicable and unlearned instruments as he did fishermen of old, to publish his pure and naked truth, and to free it of those mists and fogs wherewith the clergy hath clouded it.' Nevertheless, Barclay makes use of a considerable amount of learning—classical, patristic, and modern—for the defense of his views.

The 'Catechism' of Barclay (written in 1673) treats in fourteen chapters of the doctrines of the Christian faith, and answers the questions in the language of the Bible, without addition or comment, evidently for the purpose of showing the entire harmony of the Quakers with the written Word of God. Their distinctive peculiarities are skillfully put into the question, and the Scripture passages are so selected as to confirm them.\footnote{Comp. Ch. XI., concerning Baptism, and Bread and Wine. I will select, as a specimen, the questions on the Lord's Supper:  \par '\emph{Ques.} I perceive there was a baptism of water, which was John's baptism, and is therefore by John himself contradistinguished from Christ's: was there not likewise something of the like nature appointed by Christ to his disciples, of eating bread, and drinking wine, in remembrance of him? \par '\emph{Ans.} For I have received of the Lord, that which also I delivered unto you, That the Lord Jesus, the same night in which he was betrayed, took bread; and when he had given thanks, he brake it, and said, Take, eat; this is my body which is broken for you; this do in remembrance of me. After the same manner also he took the cup, when he had supped, saying, This cup is the new testament in my blood; this do ye, as oft as ye drink it, in remembrance of me. 1 \emph{Cor.} xi. 23–25. \par '\emph{Ques.} How long was this to continue? \par \emph{Ans.} For as often as ye eat this bread, and drink this cup, ye do show the Lord's death till he come. 1 \emph{Cor. }xi. 26. \par '\emph{Ques.} Did Christ promise to come again to his disciples? \par \emph{'Ans.} And I will not leave you comfortless; I will come to you. Jesus answered and said unto him, If a man love me, he will keep my words, and my Father will love him, and we will come unto him, and make our abode with him. \emph{John} xiv. 18, 23. \par '\emph{Ques.} Was this an inward coming? \par \emph{'Ans.} At that day ye shall know that I am in my Father, and you in me, and I in you. \emph{John }xiv. 20. \par '\emph{Ques.} But it would seem this was even practiced by the church of Corinth, after Christ was come inwardly: was it so, that there were certain appointments positively commanded, yea, and zealously and conscientiously practiced by the saints of old, which were not of perpetual continuance, nor yet now needful to be practiced in the Church? \par '\emph{Ans.} If I then your Lord and Master have washed your feet, ye ought also to wash one another's feet. For I have given you an example, that you should do as I have done to you. \emph{John} xiii. 14, 15. \par 'For it seemed good to the Holy Ghost and to us, to lay upon you no greater burthen than these necessary things: that ye abstain from meats offered to idols, and from blood, and from things strangled, and from fornication; from which if ye keep yourselves, ye shall do well: Fare ye well. \emph{Acts} xv. 28, 29. \par 'Is any man sick among you? let him call for the elders of the church, and let them pray over him, anointing him with oil in the name of the Lord. \emph{James} v. 14. \par '\emph{Ques.} These commands are no less positive than the other; yea, some of them are asserted as the very sense of the Holy Ghost, as no less necessary than abstaining from fornication, and yet the generality of Protestants have laid them aside, as not of perpetual continuance: but what other Scriptures are there, to show that it is not necessary for that of bread and wine to continue? \par '\emph{Ans.} For the kingdom of God is not meat and drink; but righteousness and peace, and joy in the Holy Ghost. \emph{Rom.} xiv. 17. \par 'Let no man therefore judge you in meat or drink, or in respect of an holy day, or of the new moon, or of the Sabbath days. Wherefore if ye be dead with Christ from the rudiments of the world, why, as though living in the world, are ye subject to ordinances (touch not, taste not, handle not: which all are to perish with the using), after the commandments and doctrines of men? \emph{Col.} ii. 16, 20–22. \par '\emph{Ques.} These Scriptures are very plain, and say as much for the abolishing of this, as to any necessity, as aught that can be alleged for the former: but what is the bread then, wherewith the saints are to be nourished? \par '\emph{Ans.} Then Jesus said unto them, Verily, verily, I say unto you, Moses gave you not that bread from heaven, but my Father giveth you the true bread from heaven,' etc. \par Then follows the whole section, John vi. 32–35, 48–58.} To the Catechism is added a brief 'Confession of Faith,' in twenty-three Articles, which is almost entirely composed of Scripture passages.

THE DISTINCTIVE PRINCIPLES OF THE FRIENDS.

The Friends are few in number, but honorably distinguished for their philanthropy, their consistent advocacy of religious freedom and the universal rights of men, their zeal in behalf of prison reform, the abolition of slavery and war. In private and social life they excel in simplicity, honesty, neatness, temperance, self-control, industry, and thrift. Their oddities in dress and habits are the shadows of virtues.

In theology and religion they are on the extreme border of Protestant orthodoxy, and reject even a regular ministry and the visible sacraments; yet they strongly believe in the supernatural and the constant presence and power of the Holy Spirit. They hold the essentials of the evangelical faith, the divine inspiration and infallibility of the Scriptures (though they disparage the letter and the human means of interpretation), the doctrine of the Trinity (in substance, though not in name),\footnote{I can not find the term \emph{Trinity} in Fox's \emph{Journal} nor in Barclay's \emph{Apology}, but both teach very clearly that Christ is God, and that the Holy Spirit is God, that all knowledge of the Father comes through the Son, and all knowledge of the Son through the Holy Spirit.} the incarnation, the divinity of Christ, the atonement by his blood, the regeneration and sanctification by the Spirit, everlasting life and everlasting punishment. And while they deny the necessity of water baptism and the Lord's Supper as a participation of the elements of bread and wine, and regard such rites as a relapse into the religion of forms and shadows, they believe in the inward substance or invisible grace of the sacraments, viz., the baptism of the Spirit and fire, and the vital communion with Christ by faith. They belong to the supernaturalistic line of Protestant dissenters, while the Socinians and Unitarians tend in the opposite rationalistic direction.

Several of the peculiar views and practices of the Quakers were anticipated by Carlstadt, the Zwickau Prophets, the Mennonites, and especially by Caspar von Schwenkfeld, a pious and retiring nobleman of Silesia (born 1490, banished 1548, d. 1561 at Ulm). Schwenkfeld embraced and preached the doctrines of the Lutheran Reformation with zeal till 1524, when he adopted, as by a higher revelation, a peculiar view of the Lord's Supper, explaining the words of institution to mean, My body is this bread, \emph{i.e.}, spiritual nourishment for the soul.\footnote{He understood \textgreek{σῶμα} and \textgreek{αἷμα} to be the subject, and \textgreek{τοῦτο} the predicate.} He also taught the deification of Christ's flesh, and opposed bibliolatry and all outward ecclesiasticism. A small remnant of his sect that was banished from Germany still survives in the eastern counties of Pennsylvania.\footnote{See Erbkam, \emph{Geschichte der protest. Sekten im Zeitalter der Reformation}, pp. 357 sqq., and Kadelbach, \emph{Geschichte K. v. Schwenkfeld's}, etc. (Lauban, 1861). The German Catechism of the Schwenkfeldians of Pennsylvania, by Christopher Schultz, Senior (translated by Daniel Rupp, Stippackville, Pa. 1863), teaches Schwenkfeld's peculiar doctrine of the Lord's Supper, but not the deification of Christ's flesh.} There is, however, no historical connection between George Fox and these predecessors. His views were entirely his own. The history of the Roman Catholic Church furnishes a parallel in the quietism of Miguel de Molinos (1627–1698), who taught that Christian perfection consists in the sweet repose of all the mental faculties in God, and in indifference to all the actions of the body. He was condemned as a heretic by Pope Innocent XI. (1687), and shut up for life in a monastic prison.

Quakerism is a system of mystic spiritualism. It is the only organized sect of mystics in England and America. The strong practical common-sense of the English race is constitutionally averse to mystic tendencies. Quakerism is an extreme reaction against ecclesiasticism, sacerdotalism, and sacramentalism. It demonstrates the paramount importance of the spirit in opposition to the worship of the letter; the superiority and independence of the inward and invisible in opposition to the overestimate of the external and visible; and the power of silence against the excess of speech.

Christianity undoubtedly is spirit and life, and may exist under different forms, or if necessary without form, like the spirit in the disembodied state. But the normal condition is a sound spirit in a sound body, and while God is independent of his own ordinances, we are bound to them. The Quakers make the exception the rule, but by the law of reaction formalism takes revenge. Their antiformalism becomes itself a stereotyped form, and their peculiar hats and coats are as distinctive as the clerical surplice and gown. When they leave their Society they usually join the Episcopal Church, the most formal among the Protestant denominations.

THE INNER LIGHT.

The ruling principle of Quakerism is the universal inner light.\footnote{Penn (in the Preface to Fox's \emph{Journal}, p. xiv.) calls it 'the fundamental principle which is as the corner-stone of their fabric, and, to speak eminently and properly, their characteristic or main distinguishing point or principle, viz., the light of Christ within, as God's gift for man's salvation. This is as the root of the goodly tree of doctrines that grew and branched out from it.' Fox's \emph{Journal} is full of it; see the list of passages in Vol. II. pp. 551 sq. of the 6th ed. (Leeds, 1836).} It is also called the seed, the Word of God, the gift of God, the indwe1ling Christ. This is not to be confounded with reason or conscience, or any natural faculty of man.\footnote{Barclay (\emph{Apol.} p. 74) rejects the errors of Pelagians and Socinians, and teaches the corruption of human nature in consequence of the fall, but maintains, in opposition to Augustine, Luther, and Calvin, that God does not impute sin to infants until they commit actual transgression. Gurney says (l.c. p. 6): 'Never did they [the Quakers] dare to consider this light as a part of fallen man's corrupt nature; never did they hesitate to ascribe it to the free and universal grace of God through Christ Jesus our Lord.'} It is supernatural and divine in its origin; it is a direct illumination of the mind and heart by the Spirit of God for the purpose of salvation. It is the light of the Logos, which shines 'in darkness' and 'lighteth every man that cometh into the world,'\footnote{\par John i. 9. The difference in the construction of \textgreek{ἐρχόμενου εἰς τὸν κόσμον} does not affect the universality, which is sufficiently sustained by \textgreek{πάντα ἄνθρωπον} but the question is whether John means the light of reason or the light of grace, and in the latter case whether it is sufficient for salvation or merely preparatory to it. When Fox, on his second visit to Cromwell (in 1656), quoted this passage, he was met with the objection that John meant 'the natural light;' but he 'showed him the contrary—that it was divine and spiritual, proceeding from Christ, the spiritual and heavenly man' (\emph{Journal}, Vol. I. p. 383).} It is Christ himself dwelling in man as the fountain of life, light, and salvation. It is the primary source of all religious truth and knowledge. It opens the sense of spiritual mysteries; it convinces and converts; it gives victory over sin, and brings joy and peace. It is communicated to men without distinction of race or religion or education, not indeed in the same measure, but in a degree sufficient to save them if they obey it, and to condemn them if they reject it. 'The grace of God that bringeth salvation hath appeared to \emph{all} men.'\footnote{Titus ii. 11. Other passages quoted by Quakers for their favorite doctrine are, Gen. vi. 3; Deut. xxx. 14; Rom. x. 3; Luke ii. 10; Rom. ii. 14, 15; Col. i. 23; Eph. v. 13; Acts x. 35.} A day of merciful visitation comes to every human being at least once in his life, and marks a critical turning-point which determines his character in this world and his eternal fate in the world to come. To many the voice from heaven speaks often.

Cornelius was, under the divine influence of that light before the arrival of Peter and the hearing of the gospel. Socrates traced his better impulses to the divine monitor in his breast, who from childhood checked his evil passions without coercion.\footnote{\emph{Apol. Soc.} He calls his \textgreek{δαιμόνιον} (in Jowett's translation) 'a voice which comes to me and always forbids me to do something which I am going to do, but never commands me to do any thing, and which stands in the way of my being a politician.' He goes on to say that in politics he would have perished long ago without doing any good either to the people or to himself. The case of Socrates is not mentioned by Barclay, but by Gurney, p. 42: 'When Socrates, as compared with his fellow-countrymen, attained to an eminent degree of disinterestedness, integrity, justice, and charity; when he obeyed the counsels of that unknown monitor who so frequently checked him in the hour of temptation; when he bore so clear a testimony to virtue as to be persecuted to death for virtue's sake—on what scriptural grounds can any man deny that he was made a partaker, to a certain degree, of a divine influence?'} The savage Indians of North America followed the light when, after having been long engaged in war, they sacrificed a spotless white dog to the Great Spirit and threw their tomahawks into the lake.\footnote{Gurney, p. 42.}

If Christ died for all men, his benefits must in some way be offered to all. He is the personal Light of the whole world, which shines into all parts of the human family backward to Adam and forward to the end of time. As many are sinners without ever having heard of Adam and the fall, so many are partakers of Christ without any external knowledge of him or the Scriptures. Else idiots, infants, and the saints who died before Christ's advent could not be saved. Historical knowledge can not save without experimental knowledge, but experimental knowledge may save without historical knowledge.

The inner light agrees with the teaching of the Bible, though not confined to its letter. It is the true interpreter of the Bible, which without it remains a sealed book. It holds in this respect the same position which the Roman Catholic Church assigns to unwritten tradition, with this important difference, that tradition is an outward, objective authority, and confined to the visible Church, while the inner light is subjective, and shines upon all men.

Quakerism thus boldly breaks through the confines of historical Christianity and the means of grace, indefinitely expands the sphere of revelation, and carries the saving power of Christ, even in this present life, into the regions of heathen darkness. It must consistently regard all virtuous and pious heathen as unconscious Christians, who, like the Athenians of old, 'unknowingly' worship an 'unknown God.' Justin Martyr, the first Christian philosopher, advanced the idea that the 'Logos spermaticos,' \emph{i.e.}, the Eternal Word of God, before his incarnation, scattered the divine seed of truth and righteousness among the Greeks as well as the Jews. Zwingli taught the salvation of many heathen and of all children dying in infancy. But these were isolated private opinions; the doctrinal standards of the orthodox Churches—Greek, Latin, and Protestant—know of no Christ and no salvation outside of Christendom and without the written or preached gospel. The Quakers teach the absolute universality, not indeed of salvation, but of the \emph{offer} and the \emph{opportunity} of salvation.

This doctrine is the corner-stone of their system.\footnote{Hence their name, 'Professors of the Light,' 'Friends of Light,' 'Children of Light.'} It is the source of their democracy, their philanthropy, their concern for the lowest and most neglected classes of society, their opposition to slavery, war, and violence, their meekness under suffering, their calmness and serenity of temper. But the same doctrine explains also their comparative disregard of the \emph{written} Scriptures, the visible Church, the ministry, the means of grace, the forms of worship, and their indifference to heathen missions. There is, however, more recently among orthodox Friends a growing disposition to aid in the circulation of the Bible, the work of foreign missions, and to associate with evangelical Christians of other Churches.

BARCLAY'S THESIS.

Barclay reduces the doctrinal system of the Friends to fifteen propositions or theological theses, which are briefly as follows:\footnote{See them in full, Vol. III. p. 749.}

1. \emph{The Foundation of Knowledge.}—The height of happiness is in the true knowledge of God and of Jesus Christ (John xvii. 3).

2. \emph{Immediate Revelation.}—This comes from the Son of God (Matt. xi. 27) through the testimony of the Spirit.

This is the inner light which has already been sufficiently explained.

3. \emph{The Holy Scriptures.}—They contain the revelations of the Spirit of God to the saints. They are a declaration of the fountain, but not the fountain itself; they are the secondary rule of faith and morals, subordinate to the Spirit from which they derive all their excellency and certainty (John xvi. 13).

4. \emph{The Condition of Man after the Fall.}—All men are by nature fallen, degenerated, and spiritually dead, but hereditary sin is not \emph{imputed} to infants until they make it their own by actual transgression. Socinianism and Pelagianism are rejected, but also the doctrine of the 'Papists and most Protestants,' that a man without the grace of God may be a true minister of the gospel.

5. \emph{Universal Redemption by Christ.}—God wills all men to be saved; Christ died for all men; the light is sent to every man for salvation, if not resisted.

On this point the Quakers side with Lutherans and Arminians against Calvinists, but go far beyond them.

6. Objections to the universality of redemption refuted.

7. \emph{Justification.}—Man is regenerated and justified when he receives the inner light. It is not by our works that we are justified, but by Christ who is both the gift and the giver, and the cause producing the effects in us.

The Quakers closely connect justification with sanctification, and approach the Roman view, with this difference, that they teach justification \emph{in} our works, not \emph{on account} of our works. Penn distinguishes between legal justification, that is, the forgiveness of past sins through Christ, the alone propitiation, and moral justification or sanctification, whereby man is made inwardly just through the cleansing and sanctifying power and Spirit of Christ.

8. \emph{Perfection.}—Man may become free from actual sinning, and so far perfect; yet perfection admits of growth, and there remains a possibility of sinning.\footnote{Penn (Preface to Fox's \emph{Journal}, p. xiv.) says that the Friends 'never held a perfection in wisdom and glory in this life, or from infirmities or death, as some have with a weak or ill mind imagined and insinuated against them.'}

The Methodists have substantially adopted this view, and call it entire consecration or perfect love.

9. \emph{Perseverance.}—Those who resist the light, or disobey it after receiving it, fall away (Heb. vi. 4–6; Tim. i. 6); but it is possible in this life to attain such a stability in the truth from which there can be no total apostasy.

This is a compromise between Calvinism and Arminianism.

10. \emph{The Ministry.}—Those and only those are qualified ministers of the gospel who are illuminated and called by the Spirit, whether male or female, whether learned or unlearned. These ought to preach without hire or bargaining (Matt. x. 8), although they may receive a voluntary temporal support from the people to whom they administer in spiritual things.

11. \emph{Worship.}—It consists 'in the inward and immediate moving and drawing of the Spirit, which is neither limited to places or times or persons.' All other worship which man appoints and can begin and end at his pleasure is superstition, will-worship, and idolatry.

All forms and even sacred music are excluded from the naked spiritualism of Quaker worship. It is simply reverent communion of the soul with God, uttered or silent. I once attended a Quaker meeting in London whose solemn silence was more impressive than many a sermon. I felt the force of the word, 'There was silence in heaven for the space of half an hour.' At another meeting I heard one man and several women exhort and pray in a tremulous voice and with reverential awe, as if in the immediate presence of the great Jehovah. All depends upon the power of the Holy Spirit.

12. \emph{Baptism.}—It is 'a pure and spiritual thing, a baptism of the Spirit and of fire,' by which we are purged from sin (1 Pet. iii. 21; Rom. vi. 4; Col. ii. 12; Gal. iii. 27; John iii. 30). Of this the water-baptism of John was a figure commanded for a time. The baptism of infants is a human tradition, without Scripture precept or practice.

13. \emph{The Communion of the Body and Blood of Christ} is likewise inward and spiritual, of which the breaking of bread at the last Supper was a figure. It was used for a time, for the sake of the weak, even by those who had received the substance, as the washing of feet and the anointing of the sick with oil was practiced; all which are only the shadows of better things. (John vi. 32–35; 1 Cor. x. 16, 17.)

This doctrine of the sacraments is a serious departure from the universal consensus of Christendom and the obvious intention of our Saviour. It can only be accounted for as a protest against the opposite extreme, which substitutes the visible sign for the invisible grace.

14. \emph{The Power of the Civil Magistrate.}—It does not extend over the conscience, which God alone can instruct and govern, provided always that no man under pretense of conscience do any thing destructive to the rights of others and the peace of society. All civil punishments for matters of conscience proceed from the spirit of Cain the murderer.

Here the Quakers, like the Baptists, commit themselves most unequivocally to the doctrine of universal religious liberty as a part of their creed.

15. \emph{Salutations and Recreations.}—Under this head are forbidden the taking off the hat to a man, the bowings and cringings of the body, and 'all the foolish or superstitious formalities' which feed pride and vanity and belong to the vain pomp and glory of this world; also all unprofitable and frivolous plays and recreations which divert the mind from the fear of God, from sobriety and gravity. Penn said of Fox that he was 'civil beyond all forms of breeding.'

The Apology of Barclay is a commentary on these propositions.

Note.—The Hicksites.—In the year 1827 a schism took place among the Friends in Philadelphia, and extended to most of the yearly meetings in America, but had no influence in England. Since then the Quakers are divided into 'orthodox' Quakers and 'Hicksites,' although the latter refuse to be called by any other name but that of 'Friends' or 'Quakers.' The founder of this rupture was Elias Hicks, born in Hempstead, Long Island, March 19, 1768; died in Jericho, N.Y., Feb. 27, 1830.

He took strong ground against slavery, and abstained from all participation in the fruits of slave labor. He was for a long time an acceptable preacher, but early in the present century he advocated radical Unitarian and other heterodox doctrines, which shocked the majority of the Quakers and led to commotion, censure, and schism. The first separation took place in the Yearly Meeting at Philadelphia, and then a similar one in New York, Baltimore, Ohio, and Indiana. Many espoused the cause of Hicks, in the interest of religious liberty and progress, without indorsing his heretical opinions on the articles of the Trinity, the divinity, and the atonement of Christ.

The extreme left of the Hicksites broke off in 1853 in Chester County, Pa., and organized a separate party under the name of \emph{Progressive Friends.} They opened the door to all who recognize the equal brotherhood of the human family, without regard to sex, color, or condition, and engage in works of beneficence and charity. They disclaim all creeds and disciplinary authority, and are opposed to every form of ecclesiasticism.

The Hicksite movement drove the orthodox Quakers more closely to the Scriptures, and called forth several official counter-demonstrations.

On the 'Hicksite' Quakers, see Elias Hicks,  \emph{Journal of his Life and Labors}, and his \emph{Sermons}, Phila. 1828; and Janney (a Hicksite), \emph{History of the Society of Friends}, Vol. IV.

\xsection{The Moravians.}

§ 108. The Moravians.

See the Literature on the Bohemian Brethren, § 75, p. 565, and the Waldenses, p. 568.

Doctrinal and Confessional.

I. Zinzendorf:  \emph{Ein und zwanzig Discourse über die Augsburgische Confession}, 1747–1748 (never published through the trade, and therefore rare). Also the other writings of Zinzendorf, and especially his hymns and spiritual poems, collected and published by Albert Knapp, with a spirited sketch of bis life and character (Stuttg. 1845).

Aug. Gottlieb Spangenberg:  \emph{Idea Fidei Fratrum oder Kurzer Begriff der christlichen Lehre in den evang. Brüdergemeinen.} Barby, 1778, 1782; Gnadau, 1833; English ed. Lond. 1784. Accepted as authority. By the same: \emph{Declaration über die zeither gegen uns ausgegangenen Beschuldigungen.} Berlin, 1772.

Hermann Plitt (Pres. of the Morav. Theol. Seminary in Gnadenfeld): \emph{Evangelische Glaubenslehre nach Schrift und Erfahrung.} Gotha, 1864, 2 vols. Not authoritative. By the same: \emph{Zinzendorf's Theologie.} Gotha, 1869–1874, 3 vols.

The hymns and liturgies of the Moravian Church.

Edm de Schweinitz (Morav. Bishop): \emph{The Moravian Manual.} Publ. by authority of the Synod. 2d enlarged ed. Bethlehem, Pa. 1869.

II. Among the early opponents of the Moravians we mention Fresenius,  Fabricius,  Georgius, and the celebrated commentator, J. A. Bengel (\emph{Abriss der sogen. Brüdergemeinde, in welchem die Lehre und die ganze Sache geprüft, das Gute und Böse dabei unterchieden}, etc. Stuttg. 1751; republ. Berlin, 1859).

III. Modern representations by divines not of the Moravian Church.

Möhler:  \emph{Symbolik}, pp. 533 sqq.; Schneckenburger:  \emph{Vorlesungen über die kleineren protest. Kirchenparteien}, pp. 152–171; R. Hofmann:  \emph{Symbolik}, pp. 533 sqq.

Historical.

I. Biographies of Count Zinzendorf.

Spangenberg:  \emph{Leben des Grafen Zinzendorf.} Barby, 1772–1775, 8 vols. Thorough, reliable, and prolix.

J. G. Müller (brother of the Swiss historian, John von M.): \emph{Bekenntnisse merkmürdiger Männer von sich selbst.} 3 vols. 1775.

L. C. von Schrautenbach:  \emph{Der Graf v. Zinz. und die Brüdergemeinde seiner Zeit, herausgeg. v. F. W. Kölbing.} Gnadau, 1851. Written in 1782, but not for publication, and kept as MS. in the Archives of the Moravian Church till 1851. One of the most interesting works on Zinzendorf, setting forth the philosophy of his religion.

Varnhagen von Ense:  \emph{Leben Zinzendorf's.} Berlin, 1830; 2d ed. 1846. The view of an outsider, similar to Southey's Life of Wesley.

J. W. Verbeck:  \emph{Gr. Zinzendorf's Leben und Charakter.} Gnadau, 1845. An extract from Spangenberg.

F. Bovet:  \emph{Le Comte de Zinzendorf.} Paris, 1860.

G. Burkhardt:  \emph{Zinzendorf und die Brüdergemeinde}, in Herzog's \emph{Real-Encykl.} Vol. XVIII. pp. 508–592 (Gotha, 1864), and published as a separate volume.

II. Histories of the Moravian Church.

Many MS. sources in the Archives of Herrnhut, Saxony, especially the 'Lissa Folios,' relating to the history of the Ancient Bohemian and Moravian Church; the 'Diarium der Gemeinde zu Herrnhut' down to 1736; the journals and letters of Zinzendorf; and the history both of the Ancient and Renewed Church, by John Plitt, from 1722 to 1836, in 9 vols.

The \emph{Büding'sche Sammlung.} Büdingen and Leipzig, 1742–1744, 3 vols. A collection of documents.

The \emph{Barby'sche Sammlung.} Barby, 1760, 2 vols. A continuation of the former.

David Cranz:  \emph{Alte und neue Brüderhistorie} (down to 1769). Barby, 1772; continued by Hegner, in 3 parts, 1791–1816. Engl. transl. by La Trobe, London, 1780.

\emph{Die Gedenktage der erneuerten Brüderkirche} (\emph{Memorial Days of the Renewed Brethren's Church}). Gnadau, 1820.

Bp.  Holmes:  \emph{History of the United Brethren.} Lond. 1825, 2 vols.

A. Bost:  \emph{Histoire de l’Église des Frères de Bohème et Moravie.} Paris, 1844, 2 vols. Abridged English transl. publ. by the Relig. Tract Soc. of London, 1848.

Bp.  E. W. Cröger:  \emph{Geschichte der erneuerten Brüderkirche} (down to 1822). Gnadau, 1852–1854, 3 vols. (The same wrote also a \emph{Geschichte der alten Brüderkirche.} Gnadau, 1865 and 1866, 2 vols.)

Verbeek:  \emph{Geschichte der alten und neuen Brüder-Unität.} Gnadau, 1857.

H. Plitt:  \emph{Die Gemeine Gottes in ihrem Geiste und ihren Formen mit Beziehung auf die Brüdergemeine.} Gotha, 1859.

Dr.  Nitzsch:  \emph{Kirchengeschichtliche Bedeutung der Brüdergemeinde.} Berlin, 1853.

Missionary.

The missionary literature of the Moravians is very large and important, and embraces the works of Cranz on Greenland (1767); Oldendorp (1777) on Danish Missions; Heckewelder (1817) on Indian  Missions; L. Kölbing,  \emph{Uebersicht der Missionsgeschichte der evang. Brüderkirche} (1832 and 1833); Bp. von Schweinitz,  \emph{Life of David Zeisberger} (Phila. 187O). Comp. the \emph{Missionary Manual and Directory of the Unitas Fratrum}, Bethlehem, Pa. 1875.

HISTORICAL SKETCH.

We must distinguish between the old Bohemian and Moravian Brethren who belonged to the Slavonic race, and the new Moravians who are chiefly German or of German descent. The connecting link between the two was the celebrated educator, John Amos Comenius (1592–1671), the Jeremiah of the former, and the John the Baptist of the latter, who, hoping against hope for the resurrection of the Bohemian Unitas Fratrum, nearly crushed to death by persecution, left behind him their order of discipline, and made provision for the ordination of two bishops, that through them the succession might be preserved in a quiescent state, until, in 1735, it was transferred to the renewed Church.

The new Moravian Church\footnote{Also called the Unitas Fratrum, the United Brethren, the Moravian Brethern; in German, Brüdergemeine, or Herrnhuter. They must not be confounded with the Methodist 'United Brethren in the United States,' founded by Rev. William Otterbein in 1800.} took its origin from the remnant (the ' Hidden Seed') of the Bohemian and Moravian Brethern, to whom Count Zinzendorf (1700–1760), under the guidance of a special providence, gave an hospitable refuge on his estates at Berthelsdorf, in Upper Lusatia, Saxony. The asylum was called \emph{Herrnhut} (the Lord's Protection), and became the mother church and the centre of the denomination.

The little colony of immigrants from Moravia soon increased, by the accession of German families of the pietistic school of Spener, to the number of three hundred souls. It was organized on the basis of the \emph{Ratio Disciplinæ} of Comenius. David Nitschmann was consecrated the first bishop by Daniel Ernst Jablonsky (court chaplain in Berlin) and Christian Sitkov, the surviving bishops of the old succession (March 13, 1735). This consecration was performed secretly in the presence of only two members of the Bohemian congregation in Berlin, for the sole purpose of sending ordained ministers to the distant missions and colonies. It was not intended to establish an episcopal form of government, separate and distinct from the national Lutheran Church, but this separation was the natural consequence. The second bishop was Count Zinzendorf himself, who gave up his office at the Saxon court and his worldly prospects to devote himself entirely to the Church of his own planting.\footnote{It is an interesting fact that Frederic William I., king of Prussia, advised Zinzendorf to get the old Moravian Episcopal ordination, and that Zinzendorf conferred on the subject with Bishop Jablonsky, and with his friend, the Archbishop of Canterbury (John Potter).} With all his eccentricities he was one of the purest and most remarkable men in the history of Christianity, a religious and poetic genius, and a true nobleman by nature and divine grace as well as by rank. He had but one all-absorbing passion—Christ and him crucified.\footnote{'\emph{Ich habe nur eine Passion, und die ist Er, nur Er.}'} From his childhood, when he used to write letters to his beloved Saviour, this sacred fire burned in him, and continued to burn till he was called to see him face to face. He early conceived the idea, by planting in the spirit of Spener a true Church in the nominal Church, to reform the Church at home, and to carry the gospel to the heathen. We may call him the German Wesley; he was an organizer like John Wesley, and a true hymnist like his brother Charles. The Oxford Methodists started with a legalistic type of piety, but they received a new inspiration from the childlike, cheerful, serene, and sublime trust in God which characterized the Moravians with whom they came in contact.

The patriarchs of Moravianism—Zinzendorf, Nitschmann, and Spangenberg—like the patriarchs of Methodism, labored in both hemispheres at a time when the stagnant State Churches of Germany and England cared little or nothing for their children in America. They founded Bethlehem (1741) and Nazareth in Pennsylvania, and other colonies which remain to this day. Zinzendorf endeavored to unite the other German denominations and sects in Pennsylvania into one Church, but in vain.\footnote{On the unionistic labors of Count Zinzendorf in Pennsylvania from 1742 to 1748, see an interesting article of the Rev. L. F. Reichel (mostly from unpublished MSS.) in Schaff's \emph{Deutscher Kirchenfreund} for 1849, pp. 93–107.}

The Moravian brotherhood is \emph{par excellence} a missionary society at home and abroad. It has but few regularly organized congregations scattered in Christian lands, but in an age of indifferentism and rationalism they were like cities of refuge and oases in the wilderness, with fresh fountains and green pastures for multitudes who flocked to them for refreshment.\footnote{Hase (\emph{Kirchengeschichte}, p. 636, 9th ed.): '\emph{Die Frömmigkeit ist in Herrnhut eine Manier geworden, aber viele stille oder gebrochene Herzen hatten hier eine Heimath, und der alte Christus in den Zeiten des Unglaubens ein Heiligthum.}'} They are still holding up the model of living congregations of real Christians. Besides, they have mission stations, called Diaspora (1 Pet. i. 1), for those who wish to derive spiritual benefit from them without severing their connection with the established Churches. These half-members may be compared to the Jewish proselytes of the gate as distinguished from the proselytes of righteousness. The Moravians, however, are free from the spirit of proselytism, and endeavor to promote peace and union among the Christians at home. But they are aggressive abroad, and concentrate their energies on foreign missions. Their chief glory lies in the extraordinary zeal and self-denial with which, since 1732, they have labored for the conversion of the most ignorant and degraded heathen in Greenland, Labrador, among the American Indians, and the African negroes and Esquimaux, at a time when orthodox Protestant Christendom had not yet awoke to a sense of its long-neglected duty. To the small band of Moravians belongs the first place of honor in the work of foreign missions.

DISCIPLINE AND WORSHIP.

The Moravian congregations in Germany are select communities of converted Christians, \emph{ecclesiæ in ecelesia}, separate and distinct from the national Churches and the vanities of the surrounding world.\footnote{The Moravian settlements in the United States were organized on the same exclusive principle, but have recently been thrown open to other people.} They have a strict discipline, but they are free from gloomy asceticism, and cherish a cheerful and trustful piety with love for music and social refinement. Their educational institutions attract pupils from all directions.

The form of government is a kind of Episcopal Presbyterianism, under the supreme legislative power of synods, and an executive administration of an elective board of bishops and elders, called the 'Unity's Elders' Conference.' The bishops ordain deacons and presbyters: they represent the whole \emph{Unitas Fratrum}, are official members of the synods, and have usually a seat in the governing boards. They claim an unbroken succession, but lay no stress on it, and fully recognize the validity of Presbyterian orders.

The home churches are divided into three provinces, Continental, British, and American. In 1857 these were declared independent in local and provincial affairs, but they continue to be united in doctrine and the work of foreign missions.

In worship, the Moravians combine liturgical and extemporaneous prayer. At all the liturgical services music forms a prominent feature. Their liturgy and hymn-book are of a superior order. They have greatly enriched the treasures of German hymnology, and produced also one of the best English hymnists in James Montgomery (1771–1854), 'the Cowper of the nineteenth century.' Love-feasts are held preparatory to the communion, in imitation of the ancient Agapæ. Foot-washing was formerly practiced, but has been discontinued since the beginning of the present century. The former use of the lot in connection with marriage has been practically abandoned; and in connection with the appointment of ministers it has been restricted or is left discretional.

DOCTRINES.

The Moravians acknowledge no exclusive and compulsory symbols. They are essentially unionistic, and seek union in harmony of spirit, life, and worship, rather than in a logical statement of doctrine.\footnote{Burkhardt (in Herzog, Vol. XVIII. p. 589) says: '\emph{Die Brüdergemeinde stellt nie ein äusserlich formulirtes Bekenntniss nach aussen hin auf, das sie von anderen evangelischen Glaubensgenossen trennen könnte. Sie wird es und kann es nie thun, denn nicht Abschluss und Scheidung, sondern Union ist ihr Princip. Aber nur jene wahre und positive Union auf Grund der heiligen Schrift und der lebendigen Herzens-Erfahrung, die allein die Herzen vereinigt.}' Bishop Schweinitz says \{\emph{Manual}, p. 95): 'The Renewed Church of the Brethren has no Confession of Faith as such, that is, no document bearing this name.'} Their most authoritative creed is the \emph{Easter Litany}, which dates from 1749, and is still used annually in all Moravian churches, but as an act of worship, not as a formula for subscription.\footnote{See the Moravian Litany in Vol. III. p. 793.} They have always laid the chief stress on the atoning death of Christ, and the personal union of the soul with him, but more in a devotional and practical than doctrinal way. Christ crucified and living in them is the all in all of their religion, their only comfort in life and death; but they have not formulated any particular theory of the atonement or of the \emph{unio mystica.} They prefer the chiaroscuro of mystery and the personal attachment to Christ to all scientific theology.

Historically and nationally, they are more nearly related to the Lutheran denomination than to any other. They sustain to it a relation similar to that which the Wesleyans sustain to the Church of England. They professed from the start their agreement with the Augsburg Confession. Spangenberg, the exponent of their doctrinal system, begins the preface to his \emph{Idea Fidei Fratrum} with the declaration that his book is no new confession, but that the Confessio Augustana of 1530 is and shall remain their confession.

But we should remember that this indorsement of the doctrinal articles of the Augsburg Confession, though no doubt sincere, was partly a matter of policy and necessity to secure toleration in Lutheran countries.\footnote{After ten years' banishment from Saxony, Zinzendorf secured in 1748 recognition of his congregation as \emph{Augsburgische Religionsverwandte} (\emph{Addicti Augustanæ Conf.})—a title under which the Reformed, or Calvinists, were included in the Treaty of Westphalia.} It had no force outside of Germany and Scandinavia, and even there no subscription to this document was ever required.\footnote{\emph{Manual}, p. 95: 'This acknowledgment, according to the declaration of the General Synod, does not bind the conscience of any member, much less is it of any weight in those provinces of the Unity where the Augsburg Confession has no other value than as being the creed of one (the Lutheran) among many Churches enjoying equal rights' (\emph{Synod. Results of} 1857, p. 96).} The Moravians never adopted the other Lutheran symbols, least of all the Formula of Concord, which strict Lutherans regard as a legitimate development of the Augustana. They never wished to be considered, nor were they recognized as Lutherans, but were violently assailed by them for their alleged doctrinal latitudinarianism and various excesses during their early history. Even the Pietists for a period made common cause with their orthodox enemies against the new sect, though less on doctrinal grounds. The Moravians claim to be the legitimate descendants and heirs of the Bohemian Brethren, who were closely connected with the Waldenses, and had their own Confessions and Catechisms before and after the Reformation. They admitted to their communion Lutherans, Pietists, Calvinists, Anglicans, without inquiring into their creed, if only they were devout Christians. In England they were recognized by Parliament, with the concurrence of the bench of bishops, as 'an ancient Episcopal Church' (1749), and allowed to settle in the American colonies. They also freely associated with Wesleyans. They were the advocates of a conservative evangelical union of three chief types of doctrine\footnote{\emph{Lehrtropen} (\textgreek{τρόποι παιδείας}), as Zinzendorf called them. He meant different educational ways of God adapted to the varieties of national and individual character. The Lutheran type prevailed among the Moravians in Saxony, the Reformed in Holland and England. The Moravian type furnished the historical base and a peculiar element in discipline rather than doctrine.}—the old Moravian or Bohemian, the Lutheran, and the Reformed—living in brotherly harmony as a true \emph{unitas fratrum}, and having their common centre in Christ. They rise above the boundaries of nationality and sect, and represent a real catholicity or universalism of creed with Christ as the only fundamental article. 'I know of no other foundation,' says Zinzendorf, 'but Christ, and I can associate with all who build on this foundation.' He was at one time even open to a project of union with the Greek and Latin Churches and all sorts of Christian sects, but he learned that the union here below must be spiritual and inward.

It is a remarkable fact that the great German theologian, Schleiermacher, was cradled in the Moravian community, and conceived there his love for Christian union and personal devotion to Christ, which guided him through the labyrinth of speculation and skepticism, and triumphed on his death-bed. He shook almost every dogma of orthodoxy, and was willing, if necessary; to sacrifice all, if he could only retain a perfect and sinless Saviour.

Zinzendorf's theology and piety passed through a process of development—first a sound evangelical stage (1723–1742), then a period of sickly sentimentalism (1743–1750), and, last, a period of purification and reconstruction (1750–1760).\footnote{See especially Plitt and Burkhardt.} These phases are reflected in the history of his followers. Encouraged by his own unguarded language, in poetry and prose, about the luxurious reveling in the wounds of the Lamb,\footnote{Or 'Lambkin,' \emph{Lämmlein}, as the favorite phrase was. The side-wound was made especially prominent.} and the personal intimacy with the Saviour, they ran into wild and dangerous excesses of an overheated imagination. As is often the case in the history of religious enthusiasm, the spirit was about to end in the flesh.\footnote{Bishop Schweinitz thus describes this period (\emph{Moravian Manual}, pp. 35 sq.): 'The relation between Christ and his Church was described in language more highly figurative, and under images more sensuous, than any thing found even in the Song of Solomon. A mania spread to spiritualize, especially the marriage relation, and to express holy feelings in extravagant terms. Hymns abounded, treating of the passion of Jesus, apostrophizing the wound in his side, degrading sacred things to a level with the worst puerilities, and pouring forth sentimental nonsense like a flood; while services, in themselves devotional and excellent, were changed into occasions for performances more in keeping with the stage of a common theatre than with the sanctity of the house of God. In short, fanaticism rioted among ministers and people, and spread from Herrnhaag and Marienborn to other churches both on the Continent of Europe and in England. Those in America escaped, or were but slightly affected. This continued for about five years, reaching its climax in 1749. It is possible that immoralities of life may have occurred in single instances, although there are no positive proofs of this; the great majority of the Brethren, however, were preserved from such extremes.' Similar antinomian excesses occurred in the Moravian congregations in England (1751), and turned Wesley and Whitefield against their old friends, whom they charged with neglecting to preach the law either as a schoolmaster or as a rule of life, with irreverent sentimentalism and superstitions fopperies. See Tyerman, \emph{Life of John Wesley}, Vol. II. pp. 95 sqq. (Harper's ed.).} But Zinzendorf himself, honestly confessing his share of responsibility, recalled his followers from the abyss to the purity and simplicity of the gospel.

The purified and matured system of the Moravians is best exhibited in Spangenberg's \emph{Idea Fidei}, which occupies a similar position among them as Melanchthon's \emph{Loci} in the Lutheran Church. It is also set forth from time to time in the \emph{Synodical Results.} The Synod of 1869 issued the following summary of such doctrines as are deemed most essential to salvation:

'1. The doctrine of the total depravity of human nature: that there is no health in man, and that the fall absolutely deprived him of the divine image.

'2. The doctrine of the love of God the Father, who has ``chosen us in Christ before the foundation of the world,'' and ``so loved the world that he gave his only-begotten Son, that whosoever believeth in him should not perish, but have everlasting life.''

' 3. The doctrine of the real godhead and the real manhood of Jesus Christ: that God, the Creator of all things, was manifested in the flesh, and has reconciled the world unto himself; and that ``he is before all things, and by him all things consist.''

'4. The doctrine of the atonement and satisfaction of Jesus Christ for us: that he ``was delivered for our offenses, and was raised again for our justification;'' and that in his merits alone we find forgiveness of sins and peace with God.

'5. The doctrine of the Holy Ghost and the operations of his grace: that it is he who works in us the knowledge of sin, faith in Jesus, and the witness that we are children of God.

'6. The doctrine of the fruits of faith: that faith must manifest itself as a living and active principle, by a willing obedience to the commandments of God, prompted by love and gratitude to him who died for us.

'In conformity with these fundamental articles of faith, the great theme of our preaching is Jesus Christ, in whom we have the grace of the Lord, the love of the Father, and the communion of the Holy Ghost. We regard it as the main calling of the Brethren's Church to proclaim the Lord's Death, and to point to him, ``as made of God unto us wisdom, and righteousness, and sanctification, and redemption.'''\footnote{Bishop Schweinitz, in M'Clintock and Strong's \emph{Cyclop.} Vol. VI. p. 587. Comp. his Compend of Doctrine in XVII. Articles, compiled from the authorized publication in the \emph{Moravian Manual}, pp. 95–100. A popular statement is contained in the \emph{Catechism of Christian Doctrine for the Instruction of Youth in the Church of the United Brethren}, and the \emph{Epitome of Christian Doctrine for the Instruction of Candidates for Confirmation} (various editions in German and English).}

\xsection{Methodism.}

§ 109. Methodism.

Literature.

I. Doctrinal Standards.

John Wesley (1703-1791): \emph{ Sermons on Several Occasions}; and \emph{Explanatory Notes on the New Test.} In many eds., London, Bristol, New York, Cincinnati, etc. Best ed. of the Sermons by Thomas Jackson, Lond. 1825, New York, 1875.

Richard Watson (1781-1833): \emph{Theological Institutes: or a View of the Evidences, Doctrines, Morals, and Institutions of Christianity. }First ed. Lond. 1822-28, in 6 parts; best ed., with an Analysis by John M'Clintock, New York, in 2 vols. (29th ed. 1875).

W. B. Pope (Theol. Tutor, Didsbury College, Manchester): \emph{A Compendium of Christian Theology: being Analytical Outlines of a Course of Theological Study, Biblical, Dogmatic, Historical. }London (Wesleyan Conference Office), 1875 (752 pp.). By the same: \emph{The Peculiarities of Methodist Doctrine. }London, 1873.

D. D. Whedon, D.D. (Ed. of the 'Methodist Quarterly Review,' and of a Popular Commentary on the New Test.): \emph{Doctrines of Methodism. }In 'Bibliotheca Sacra' for April, 1862, pp. 241–274. Andover, Mass.

W. F. Warren:  \emph{System. Theologie.} Bremen, 1865, Vol. I.

\emph{The Doctrines and Discipline of the Methodist Episcopal Church. }1872. Ed. by Bishop Harris. New York (Nelson \& Phillips) and Cincinnati (Hitchcock \& Walden).

\emph{Catechisms of the Methodist Episcopal Church. }New York (Nelson \& Phillips). Especially No. 3, which is designed 'for an advanced grade of study.' Approved by the General Conference, 1852. Two German Catechisms by the Rev. Dr. William Nast, 1868.

II. Other Sources for the Doctrines and History of Methodism.

\emph{The Complete Works} of John Wesley (first ed. Bristol, 1771 sqq., in 32 small vols. full of typographical errors; 3d and best ed. with the author's last corrections, ed. by Thomas Jackson, Lond. 1831, 14 vols.; New York, 7 vols.).

\emph{The Poetical Works of John and Charles Wesley.} Ed. by G. Osborn, D.D. Lond. 1872, 13 vols.

\emph{The Works }of John Fletcher (Lond. 1815, 10 vols.; New York, 1831, 4 vols.).

\emph{The Sermons }and \emph{Journals }of  George Whitefield (1756, 1771).

\emph{The Journals }of Bishop  Asbury (new ed. N. Y. 1854, 3 vols.).

III. Biographies.

\emph{John Wesley}, by Coke and Moore (Lond. 1792); by John Hampson (1791, 3 vols.); by Robert Southey (with Notes by Sam. T. Coleridge, 3d ed. Lond. 1846; Amer. ed. with Notes by Coleridge, Alex. Knox, and Daniel Curry, N.Y. 1847, 2 vols.); by  Richard Watson (Lond. 1831; Amer. ed. with Notes by T. O. Summers); by L. Tyerman (Lond. and New York, 1872, 3 vols.); Isaac Taylor:  \emph{Wesley and Methodism }(Lond. and New York, 1855); James H. Rigg:  \emph{The Living Wesley as he was in his Youth and his Prime }(Lond. 1875; New York ed. with Introduction by Dr. Hurst, of Drew Theol. Seminary). Comp. Dr. Rigg's article on the \emph{Churchmanship of John Wesley}, in the 'Contemporary Review' for Sept. 1876.

\emph{Charles Wesley }(1708 to 1788), by  Thomas Jackson (Lond. 1841, 2 vols.).

\emph{George Whitefield }(the founder of 2 Methodism, b. 1714, d. 1770), by  J. Gillie (Lond. 1772, 1813); by  Robert Philip (Lond. 1830; also in German, with a Preface by Tholuck, Leipz. 1834); by  L. Tyerman (London and New York, 1877, 2 vols.; the best).

\emph{The Oxford Methodists: Memoirs of Clayton, Ingham, Gambold, Hervey, and Broughton. }By L. Tyerman. London and New York, 1873.

\emph{Early Methodist Preachers.} Ed. by Thomas Jackson (Lond. 1839, 2 vols.).

IV. General Histories of Methodism.

Dr.  Abel Stevens (\emph{History of Methodism}, New York and Lond. 1858–61, 3 vols.; \emph{History of the Methodist Episcopal Church}, N. Y. 1866–67, 4 vols.; \emph{Centenary of American Methodism}, N. Y. 1865); Dr.  George Smith (Lond. 1857–62, 3 vols.: illustrated popular edition, 1864), and a number of other works. For a concise summary, see Stevens's art. 'Methodism,' in Johnson's 'Univers. Cyclop.' Vol. III. (1876). Also for popular use,  James Porter:  \emph{The Revised Compendium of Methodism.} New York, 1875.  Jacoby:  \emph{Geschichte des Methodismus. }Bremen, 1870.

Comp. \emph{The Wesleyan Methodist Magazine. }London (Wesleyan Conference Office), 1778 to 1876 (xcix. vols.).

\emph{The Methodist Quarterly Review.} New York (Nelson \& Phillips), Vols. LVIII. till 1876.

M'Clintock and Strong's  \emph{Cyclopædia} (New York, 1867–81, 10 vols. (three supplementary vols. promised), is edited by Methodists, and pays special attention to Methodist and Arminian articles.

V. Bibliographical, Critical, and Polemical.

For the anti-Methodist literature, see  H. C. Decanver:  \emph{Catalogue of Works in Refutation of Methodism from its Origin, in }1729, \emph{to the Present Time}, Phila. (John Penington), 1846. Contains in alphabetical order the titles of 227 books and sermons against Methodism, most of which are forgotten.

G. Osborn:  \emph{Outlines of Wesleyan Bibliography.} London, 1869.

M. Schneckenburger:  \emph{Lehrbegriffe der kleineren protest. Kirchenparteien.} 1863, pp. 103–151.

Joh. Jüngst:  \emph{Amerikanischer Methodismus in Deutschland und R. Pearsall Smith. }Gotha, 1875. By the same: \emph{Wesen und Berechtigung des Methodismus. }Gotha, 1876.

CHARACTER OF METHODISM.

Methodism is the most successful of all the younger offshoots of the Reformation. In one short century it has become one of the largest denominations in England, and the largest in the United States, with missionary stations encircling the globe.

The founders were admirably qualified for their work, and as well fitted together as the Reformers. John Wesley was one of the greatest preachers and organizers, and in the abundance of his labors perhaps the most apostolic man that England ever produced. As a revivalist of practical religion he may be called the English Spener, as an organizer the Protestant Ignatius Loyola. His brother Charles occupies, next to Watts, the first place in English hymnology, and sang Methodism into the hearts of the people. Whitefield, the orator and evangelist, kindled a sacred fire in two hemispheres which burns to this day. Their common, single, and sole purpose was to convert sinners from the service of Satan to the service of God, by means of incessant preaching, praying, and working. For this end they were willing to spend and be spent, to be ridiculed, reviled, pelted and hooted by mobs, maltreated by superiors, and driven from the church into the street; for this they would in another age have suffered torture, mutilation, and death itself as cheerfully as the Puritans did before them. The practical activity of these great and good men was equaled only by that of the Reformers in the theoretic sphere. During the fifty years of his itinerant ministry, John Wesley traveled 'a quarter of a million of miles, and preached more than forty thousand sermons.'\footnote{Tyerman, \emph{John Wesley}, Vol. III. p. 658 (Harper's ed.). Dr. Rigg (\emph{The Living Wesley}, Hurst's ed. p. 208) remarks that Wesley rode ordinarily sixty miles a day, and not seldom eighty and ninety miles, besides preaching twice or thrice.} Charles Wesley composed over six thousand religious poems,\footnote{\par Osborn's edition contains 7600 poems of Wesley, including those of John, who composed all the translations from the German.} in the study, in the pulpit, on horseback, in bed, and in his dying hour.\footnote{When hardly able to articulate any more, he dictated to his wife these lines:  \par 'In age and feebleness extreme, \par Who shall a helpless worm redeem? \par Jesus, my only hope thou art, \par Strength of my failing flesh and heart; \par Oh could I catch a smile from thee, \par And drop into eternity!'} Whitefield, besides traveling through England, Ireland, and Scotland, made seven evangelistic voyages to America, turning the ship into a church, and 'preached in four-and-thirty years upwards of eighteen thousand sermons, many of them to enormous crowds, and in the teeth of brutal persecution.'\footnote{\par Tyerman, Vol. III. p. 78.} A day before his death he preached his last sermon of nearly two hours' length in the open air, 'weary \emph{in }the work, but not \emph{of} the work' of his Lord. Fletcher labored in a more restricted sphere, as Vicar of Madely, but just as faithfully and devotedly, visiting his people and the poor ignorant colliers early and late, in rain and snow, studying intensely, living all the while on bread and cheese or fruit, and exhibiting an angelic type of character, so that Wesley, from a personal acquaintance of more than thirty years, gave him the testimony that 'he never heard him speak an improper word or saw him do an improper action,' and that he never knew a man 'so inwardly and outwardly devoted to God, so unblamable in every respect.'\footnote{See Wesley's Funeral Sermon on the death of John W. Fletcher, who was a French Swiss by birth (de la Fléchière), born at Nyon, Canton de Vand, 1729, educated at Geneva, died at Madeley, 1785. His chief works is \emph{Checks to Antinomianism}, against Calvinism.} The pioneers of American Methodism were animated by the same zeal. Bishop Asbury, 'in the forty-five years of his American ministry, preached about 16,500 sermons, or at least one a day, and traveled about 270,000 miles, or 6000 a year, and presided in no less than 224 annual conferences, and ordained more than 4000 preachers.'\footnote{Stevens, \emph{Centenary of American Methodism} (N. Y. 1865), p. 94.} He was ordained bishop (1784) when the number of American Methodists fell below 15,000, and he died (1816) when it exceeded 211,000, with more than 700 itinerant preachers.

Methodism owes its success to this untiring zeal in preaching the gospel of the new birth and a 'full and free salvation' to the common people, in churches, chapels, and the open air, and to its peculiar methods and institutions—itinerancy, missionary bishops, presiding elders, lay helpers, class-meetings, camp-meetings, conferences, and systematic collections. Methodism, as Dr. Chalmers characterized it, is 'Christianity in earnest.' It works powerfully upon the feelings; it inspires preachers and members with enthusiasm; it gives every man and woman too a distinct vocation and responsibility; it 'keeps all at work and always at it,' according to Wesley's motto; it knows nothing of churches without ministers, or ministers without charges, as long as there are sinners to be converted in any corner of the globe. Methodism is better organized than any other Protestant denomination, and resembles in this respect the Church of Rome and its great monastic orders. It is a powerful rival of that Church. It has an efficient machinery with an abundance of steam, and is admirably adapted for pioneer work in a new country like America. It is a well-disciplined army of conquest, though not so good an army of occupation, since it allows so many 'to fall away from grace,' not only temporarily, but even 'totally and finally.' Till 1872 the laity was excluded from participation in Church government (and is so still in England), but was compensated by a large liberty in the sphere of worship, in class-meetings, band-meetings, love-feasts, which tend to develop the social and emotional element in religion.

METHODISM AND PURITANISM.

Methodism forms the third great wave of the Evangelical Protestant movement in England, and represents the idea of revival. The Reformation destroyed the power of the papacy. Puritanism aimed at a more thorough Reformation in Church and State, and controlled for a time the civil and religious life of the nation. Methodism kept aloof from politics, and confined itself to the sphere of practical religion. Puritanism was animated by the genius of Calvinism; Methodism, in its main current, by the genius of Arminianism. Both made a deep and lasting impression upon the national Church from which they proceeded, and moulded the character of American Christianity. The Methodist revival checked the progress of skepticism and infidelity which had begun to set in with deism. It brought the life and light of the gospel to the most neglected classes of society.

If evangelical Christianity to-day has a stronger hold on the Anglo-Saxon race in both hemispheres than on any other nation, it is chiefly due to the influence of Puritanism and Methodism.

RELATION TO THE CHURCH OF ENGLAND.

Methodism is a daughter of the Church of England, and was nursed in the same University of Oxford which, a century later, gave rise to the Tractarian school in the opposite direction towards Rome. The 'Holy Club' of the fourteen Oxford students associated for prayer, holy living, and working, began, like Dr. Pusey and his friends, with a revival of earnest, ascetic, and ritualistic High-Churchism, and received the name 'Methodists' for its punctual and methodical habits of devotion. Wesley was at first so exclusive an Episcopalian that he shrank from street-preaching and lay-preaching, and, at least on one occasion, even rebaptized Dissenters. But his contact with the simple-hearted, trustful, and happy German Moravians (Peter Böhler, Nitschmann, and Spangenberg) whom he met on his voyage across the Atlantic, in the Colony of Georgia, and after his return, led to his second 'conversion,' which took place May 24, 1738, and imparted to his piety a cheerfully evangelical and, we may say, a liberal Broad-Church character.\footnote{'At the first,' says Dr. Rigg ('Contemporary Review' for 1876, pp. 656 sq.), 'with Wesley faith had meant the intellectual acceptance of the creeds, together with the submission of the will to the laws and services of the Church. . . . Until he met with Böhler, he had not embraced, scarcely, it would seem, had conceived the idea of faith as being, in its main element, personal trust and self-surrender, as having for its central object the atonement of Jesus Christ, and as inspired and sustained by the supernatural aid and concurrence of the Holy Spirit. . . . Wesley confessed that Böhler's teaching was true gospel teaching. . . . Here ended his High-Church stage of life. Here began his work as an evangelist and Church revivalist. All dates from his final acceptance of Böhler's teaching as to the nature of faith.' Dr. Stevens says (\emph{Centenary}, p. 31): 'Methodism is indebted to Moravianism for not only some of the most important features of its moral discipline, but for the personal conversion of both the Wesleys.' But Wesley was converted before as much so as Luther was when he entered the convent of Erfurt several years before he experienced his second or evangelical conversion to the doctrine of justification by faith alone. On the other hand, some of the Oxford Tractarians were converted over again, or backward, when they joined the Church of Rome.}

He now entered upon his independent evangelistic career, yet with no idea of forming a separate denomination. His object was simply to revive experimental piety within the limits of the Anglican Church, as Spener and Francke had done before within the Lutheran Confession in Germany. Although badly treated by bishops and other clergy, he had no quarrel with the authorities in Church or State, but only with sin and Satan. His aim was to build the city of God and to save souls within the establishment, if possible; without it, if necessary. He performed indeed some uncanonical acts which led ultimately to secession, but he did it from necessity, not from choice. He never made common cause with Dissenters. He lived and died in the Church of his fathers. His brother Charles was even more conservative, and took great offense at his violation of the canons.

Had the Church of England been as wise and politic as the Church of Rome, she would have encouraged and utilized the great revival of the eighteenth century for the spread of vital Christianity at home and abroad, and might have made the Wesleyan society an advocate of her own interests as powerful as the order of the Jesuits is of the Papacy. Now, after a century of marvelous success, the founder of Methodism is better appreciated, and has been assigned (1876) a place of honor among England's mighty dead in Westminster Abbey.

The English Wesleyans continue to hold a middle position between the Established Church and the Dissenters proper, but tend latterly more to Free-Churchism.

AMERICAN METHODISM.

In the United States the Methodists were made an independent organization with an episcopal form of government by Wesley's own act. As a Tory and a believer in political non-resistance, he at first wrote against the American 'rebellion,' but accepted the providential result; and, considering himself as a 'Scriptural Episcopos,' he ordained, on the second day of September, 1784, two presbyters (Richard Whatcoat and Thomas Vasey) and one superintendent or bishop, viz., the Rev. Thomas Coke, LL.D. (a presbyter of the Church of England), for his American mission, which then embraced 83 traveling preachers and 14,988 members.\footnote{The first Methodist society in America was formed in 1766, in the city of New York, among a few Irish emigrants, by Philip Embury, a local preacher, and by his cousin, Mrs. Barbara Heck, a true 'mother in Israel.' Hence Methodism celebrated its centenary in 1866 with great festivities.} This was a bold and an irregular act, but a master-stroke of policy, justified by necessity and abundant success.\footnote{He also ordained a few presbyters for Scotland and England to assist him in administering the sacraments, on the plea that the regular clergy often refused to admit his people to the Lord's table. At the Conference of 1788 he consecrated (according to Samuel Bradburn's statement) one of his preachers as a superintendent or bishop. He had long before been convinced by Stillingfleet's 'Irenicon' and Lord King's 'Primitive Church' that bishops and presbyters were originally one order, and that diocesan episcopacy was not founded on divine right. In a letter to his brother Charles (1785) he calls the uninterrupted episcopal succession 'a fable which no man ever did or can prove.'—Rigg, 1.c. p. 669. For a full discussion of Wesley's ordination acts, see Stevens, \emph{History of Methodism}, Vol. II. pp. 209 sqq., and Tyerman, \emph{John Wesley.} Vol. III. pp. 426 sqq.}

Bishop Coke, assisted by the Rev. P. W. Otterbein, of the German Reformed Church, ordained, according to Wesley's direction, Francis Asbury to the office of joint superintendent, and twelve others to the office of presbyters, at the first General Conference held in Baltimore (Dec. 27, 1784). These were the first Protestant bishops in America, with the exception of Dr. Samuel Seabury, who was consecrated a few weeks before (Nov. 14, 1784), at Aberdeen, as bishop of the Protestant Episcopal diocese in Connecticut.\footnote{Bishop White, of Pennsylvania, was not consecrated by the Archbishop of Canterbury until Feb. 4, 1787, the consecration being delayed and nearly frustrated by certain impediments.} In a short time the society, thus fully organized, overtook older denominations, and kept pace with the rapid progress of the young republic.

The separation from the mother Church of England was complete, but her blood still flows in the veins of Methodism and shows itself in a half-way assent to her doctrinal standards (as far as they admit of an Arminian interpretation), to her liturgy (as far as it does not encourage sacerdotalism and ritualism or interfere with the freedom of worship), and to her episcopacy (as based upon expediency, and not on the divine right of succession).

BRANCHES OF METHODISM.

The Methodist Christians in England and America are divided into a number of distinct ecclesiastical organizations—the 'Wesleyans,' the 'Methodist Episcopal Church,' the 'Primitive Methodists,' the 'Primitive Wesleyans of Ireland,' the 'Bandroom Methodists,' the 'Methodist Protestant Church,' the 'Welsh Calvinistic Methodists,' the 'Free Methodist Church,' the 'African (Bethel and Zion) Methodist Episcopal Church,' etc. To the Methodist family belong also the 'Evangelical Association' (or 'Albright's Brethren,' so called from Jacob Albright, a Pennsylvania German, who founded this society in 1800), and the 'United Brethren in Christ' (founded by Philip William Otterbein, a German Reformed minister, d. in Baltimore, 1813).

The great parent body, however, are the Wesleyans in England and the Methodist Episcopal Church in the United States. They far outnumber all the other branches put together. The Methodist Episcopal Church was divided in 1844 on the question of slavery into 'the Methodist Episcopal Church' (North), and 'the Methodist Episcopal Church, South,' but measures have been inaugurated (1876) for reuniting them. Similar schisms for the same cause rent other Churches before the civil war, but have been healed or will be healed, since the war has removed the difficulty. The Roman Catholic, and next to it the Protestant Episcopal Church, owing to their conservatism, were least affected by the disturbing question of slavery, and remained intact.

The differences between the various branches of Methodism refer to the episcopate, the relative powers of the bishops and the general conference, lay representation, and other matters of government and discipline which do not come within the scope of this work. The doctrinal creed is the same in all, with the exception of the Whitefieldian Methodists, who are Calvinists, while all the rest are Arminians.

Note.—The \emph{Cyclopædia} of M'Clintock and Strong, Vol. VI. p. 159, gives the following list of Methodist denominations, with the date of their organization and estimate of their ministers and church members in 1872:

GREAT BRITAIN AND IRELAND.

Denominations

Date of  Organization.

Number of  Ministers.

Number of  Church Members

Wesleyan Methodists

1739

3,157

557,995

Welsh Calvinistic Methodists

(1745)

207

58,577

New Connection Methodists

1797

260

35,706

Primitive Methodists

1810

943

161,229

Primitive (Ireland) Methodists

1816

85

14,247

Bible Christians

1815

254

26,241

United Methodist Free Churches

1828–49

312

68,062

Wesleyan Reform Union

1849

20

9,393

Totals

5,238

931,450

AMERICA.

Denominations

Date of  Organization.

Number of  Ministers.

Number of  Church Members

Methodist Episcopal Church (in 1872)

1784

10,742

1,458,441

Methodist Church (Non-Episcopal)

1866

624

75,000

United Brethern

1800

. . . .

. . . .

Evangelical Association (Albrights)

1800

632

78,716

African Methodist Episcopal

1816

600

20,000

African Methodist Episcopal (Zion)

1819

694

164,000

Canada Wesleyans

1828

. . . .

69,597

Eastern British American Wesleyan Methodists

1854 ?

147

16,118

Methodist Episcopal Church of Canada

1828

228

21,103

Methodist Protestants, South

1830

2,858

600,900

Free Methodists

1860

about 90

6,000

Primitive Methodists

about 20

2,000

Totals

17,308

2,591,875

* This does not include the colored membership now separately organized as the \emph{Colored Methodist Episcopal Church, South.}

\xsection{Methodists Creeds.}

§ 110. Methodist Creeds.

The American Methodists have three classes of doctrinal standards.

1. The Twenty-five Articles of Religion.\footnote{\par See Vol. III. pp. 766 sqq. Comp. also Emory, \emph{History of the Discipline}, ch. i. § 2; Comfort, \emph{Exposition of the Articles} (New York, 1847); Jimeson, \emph{Notes on the Twenty-five Articles} (Cincinnati, 1853).} They were prepared by John Wesley, from the Thirty-nine Articles of the Church of England (together with an abridgment of the Book of Common Prayer), for the American Methodists, and were adopted by the Conference in Baltimore, 1784, with the exception of Article XXIII., which recognizes the United States as 'a sovereign and independent nation,' and which was adopted in 1804. These articles are now unalterably fixed, and can neither be revoked nor changed.\footnote{'The General Conference shall not revoke, alter, or change our Articles of Religion, nor establish any new standards or rules of doctrine contrary to our present existing and established standards of doctrine.' This article can not be amended (\emph{Discipline}, p. 51). The General Conference is the highest of the five judicatories, and the only legislative body of the Methodist Episcopal Church.}

2. John Wesley's Sermons and Notes on the New Testament. They are legally binding only on the British Wesleyans, but they are in fact as highly esteemed and as much used by American Methodists, and constitute the life of the denomination. When eighty-one years of age (Feb. 28, 1784), Wesley, in his famous Deed of Declaration, which is called the Magna Charta of Methodism, bequeathed the property and government of all his chapels in the United Kingdom (then 359 in number) to the 'Legal Hundred,' \emph{i.e.}, a conference of one hundred of his traveling preachers and their successors, on condition that they should accept as their basis of doctrine his Notes on the New Testament and the four volumes of Sermons which had been published by him or in his name in or before 1771.\footnote{Tyerman, Vol. III. pp. 417 sqq.} These sermons are fifty-eight in number, and cover the common faith and duties of Christians,\footnote{Thirteen discourses are on the Sermon on the Mount, chiefly ethical; two are funeral discourses (on the death of Whitefield and Fletcher); one on the cause and cure of earth-quakes; one on the use of money.} but contain at the same time the doctrines which constitute the distinctive creed of Methodism.\footnote{On Salvation by Faith; Scriptural Christianity; Original Sin; Justification by Faith; Free Grace; the Witness of the Spirit (three sermons); on Christian Perfection. It is singular there is not one sermon on the Freedom of the Will.} The Notes on the New Testament are for the most part a popular version of Bengel's \emph{Gnomon.}

3. The Book of Discipline and several Catechisms, one published in 1852, another in 1868 (by Dr. Nast), are at least secondary standards for the American Methodists.

The distinctive features of the Methodist creed are not logically formulated, like those of the Lutheran and Reformed Churches. It allows a liberal margin for further theological development. John Wesley, though himself an able logician and dialectician, sought Christianity more in practical principles and sanctified affections than in orthodox formulas, and laid greater stress on the œcumenical consensus which unites than on the sectarian dissensus which divides the Christians. The General Rules, or recognized terms of membership, for the origina1 Methodist 'societies' (not churches), are ethical and practical, and contain not a single article of doctrine. They require 'a desire to flee the wrath to come and be saved from sin,' and to avoid certain specific vices.

Nevertheless Methodists claim to have more doctrinal harmony than many denominations which impose a minute creed. There is a Methodist system of doctrine and a Methodist theology, however elastic they may be. But there is a difference of opinion among their standard writers as to the degree of originality and completeness of this system and its relation to other confessions. We may distinguish an American and an English view on the subject.

An ingenious attempt has recently been made to raise the Methodist creed to the importance and dignity of a fourth confession or symbolical system alongside of the Roman Catholic, the Lutheran, and the Calvinistic, and far above them. According to Dr. Warren, Catholicism makes salvation dependent upon a meritorious co-operation of man with God, and is essentially pagan; Calvinism makes salvation depend exclusively on the eternal decree and free grace of God, and views Christianity from the stand-point of the Old Testament; Lutheranism derives salvation from the personal relation of man to the means of grace (the Word and Sacraments), and views Christianity from the stand-point of justification by faith alone; Methodism makes salvation exclusively dependent upon man's own free relation to the illuminating, renewing, and sanctifying influences of the Holy Spirit, and represents the stand-point of Christian perfection. Calvin retains the Christians under the dispensation of the Father, Luther under the dispensation of the Son, Wesley leads them into the dispensation of the Spirit. The first confines salvation to the favorite number of the elect; the second binds it to the baptismal font, the altar, and the pulpit; the third offers it freely to all. Calvin's ideal Christian is a servant of God, Luther's a child of God, Wesley's a perfect man in the full stature of Christ.\footnote{\par \emph{Syst. Theol.} Vol. I. pp. 90, 99, 119, 140, 149, 166. Dr. Warren (who is now President of the Methodist University in Boston) wrote this able book (which is as yet, 1876, unfinished) while in Germany, and under the stimulus of the generalizing theories of some German divines. Zinzendorf had made a somewhat similar distinction between the Lutheran, Reformed, and Moravian types of doctrine (\emph{Lehrtropen}), but comprehended them all in his brotherhood. James Martineau, from the Unitarian point of view, represents Luther, Calvin, and Wesley as the representatives of the orthodox gospel in three dialects (\emph{Studies of Christianity}, London, 1873, pp. 399 sq.).}

English Methodists claim for their system a humbler position, and represent it, in accordance with the intention of the founders, as a liberal evangelical modification of the Anglican creed, with some original doctrines to which they attach great importance.\footnote{\par Professor William B. Pope, of Didsbury College, Manchester, one of the leading Wesleyan divines, makes the following statement concerning the creed of the English Methodists (in the Introduction to his translation of Winer's \emph{Comparative View of the Doctrines and Confessions of the various Communities of Christendom}, Edinb. 1873, pp. lxxvi.–lxxviii.): 'It may be said that English Methodism has no distinct articles of faith. At the same time it is undoubtedly true that no community in Christendom is more effectually hedged about by confessional obligations and restraints. Reference has been made to the distinction of creeds, confessions, and standards. Methodism combines the three in its doctrinal constitution after a manner on the whole peculiar to itself. Materially if not formally, virtually if not actually, implicitly if not avowedly, its theology is bound by the ancient œcumenical Creeds, by the Articles of the English Church, and by comprehensive standards of its own, the peculiarity of its maintenance of these respectively having been determined by the specific circumstances of its origin and consolidation—circumstances with which it is not our business here to enter. In common with most Christian Churches it holds fast the Catholic Symbols; the Apostolical and Nicene are extensively used in the Liturgy, and the Athanasian, not so used, is accepted so far as concerns its doctrinal type. The doctrine of the Articles of the Church of England is the doctrine of Methodism. This assertion must, of course, be taken broadly, as subject to many qualifications. For instance, the Connection has never avowed the Articles as its Confession of Faith; some of those Articles have no meaning for it in its present constitution; others of them are tolerated in their vague and doubtful bearing, rather than accepted as definitions; and, finally, many Methodists would prefer to disown any relation to them of any kind. Still the verdict of the historical theologian, who takes a comprehensive view of the estate of Christendom, in regard to the history and development of Christian truth, would locate the Methodist community under the Thirty-nine Articles. He would draw his inference from the posture towards them of the early founders of the system; and he would not fail to mark that the American branch of the family, which has spread simultaneously with its European branch, has retained the Articles of the English Church, with some necessary modifications, as the basis of its Confession of Faith. Setting aside the Articles that have to do with discipline rather than doctrine, the Methodists universally hold the remainder as tenaciously as any of those who sign them, and with as much consistency as the great mass of English divines who have given them an Arminian interpretation. That is to say, where they diverge in doctrine from the Westminster Confession, Methodism holds to them; while this Confession rather expresses their views on Presbyterian Church government. It may suffice to say generally on this subject, that so far as concerns the present volume [of Winer], every quotation from the English Articles may stand, if justly interpreted, as a representative of the Methodist Confession.  \par 'Finally, we have the Methodist Standards, belonging to it as a society within a Church, which entirely regulate the faith of the community, but are binding only upon its ministers. Those Standards are to be found in certain rather extensive theological writings which have none of the features of a Confession of Faith, and are never subscribed or accepted as such. More particularly, they are some Sermons and Expository Notes of John Wesley; more generally, these and other writings, catechisms, and early precedents of doctrinal definition; taken as a whole, they indicate a standard of experimental and practical theology to which the teaching and preaching of its ministers are universally conformed. What that standard prescribes in detail it would be impossible to define here. . . . Suffice that the Methodist doctrine is what is generally termed Arminian, as it regards the relation of the human race to redemption; that it lays great stress upon the personal assurance which seals the personal religion of the believer; and that it includes a strong testimony to the office of the Holy Spirit in the entire renewal of the soul in holiness, as one of the provisions of the covenant of grace upon earth. It may be added, though only as an historical fact, that a rigorous maintenance of this common standard of evangelical doctrine has been attended by the preservation of a remarkable unity of doctrine throughout this large communion.' \par Dr. Whedon, the editor of the 'Methodist Quarterly Review,' in a notice of Pope's \emph{Winer} (October No., 1873, pp. 680 sqq.), enters 'his firm, fraternal protest against being recorded before the eyes of the world as training under the Thirty-nine Articles of the Church of England,' and says, 'The entire body of Methodists of the United States no more hold the Thirty-nine Articles, doctrinally, than they do the Westminster Confession. They reject a large share of both for the same reason, namely, that they are, in their proper interpretation, Calvinistic. Nor does this Confession express their views on Presbyterian Church government: for the Confession affirms the divine obligation of Presbyterianism, and the large body of American Methodists believe in the right of a voluntary episcopacy.'}

\xsection{Analysis of Arminian Methodism.}

§ 111. Analysis of Arminian Methodism

THE SEMI-ANGLICAN DOCTRINES.

The Twenty-five Articles represent the doctrines which Methodism holds in common with other evangelical Churches, especially with the Church of England. They are an abridgment of the Thirty-nine Articles of Religion, with a view to simplify and to liberalize them. Wesley omitted the political articles, which apply only to England, and those articles which are strongly Augustinian, especially Article 17, of Predestination (which teaches unconditional election to salvation and the perseverance of the elect), Art. 13, of Works before Justification (which are said to have the nature of sin), and Art. 8 (which indorses the three Creeds). On the other hand, Art. 10, of Free Will, which teaches (with Augustine, Luther, and Calvin) the natural inability of man to do good works without the grace of God, is literally retained (Meth. Art. 8).

Minor doctrinal changes were made in Art 2 (Art. 2), where the clauses 'begotten from everlasting of the Father,' and 'of her [the Virgin's] substance,' are omitted (either as doubtful or lying outside of a creed);\footnote{Emory, in his \emph{History of the Discipline}, inserts the clause, 'begotten of everlasting from the Father,' as adopted in 1784, but omitted in 1786 and in later editions, perhaps by typographical error. A Methodist correspondent (Rev. D. A. Whedon) suggests to me that Wesley may have made a distinction between the eternal \emph{Sonship} and the eternal \emph{Generation}, and may have maintained the former, but questioned the latter as referring to the manner rather than the fact. Prof. Pope, the latest Methodist writer on Dogmatics, avoids this question as belonging to the transcendental mysteries (\emph{Christ. Theol.} p. 272).} in Art. 9 (7), where the last clauses, which affirm the continuance of original sin in the regenerate, are left out (as inconsistent with Wesley's view of perfection); in Art. 16 (12), where 'sin after justification' is substituted for 'sin after baptism' (to avoid the doctrine of baptismal regeneration); in Art. 25 (16), of the Sacraments, where the words 'sure witnesses and effectual,' before 'signs of grace,' are stricken out (which betrays a lowering of the doctrine of the Sacraments); in Art. 34 (22), where 'traditions of the Church' are changed into 'Rites and Ceremonies.'

These omissions and changes are significant, and entirely consistent with Methodism, but they are negative rather than positive. Wesley eliminated the latent Calvinism from the Thirty-nine Articles, but did not put in his Arminianism, nor his peculiar doctrines of the Witness of the Spirit and Christian Perfection, leaving them to be derived from other documents of his own composition.

THE ARMINIAN DOCTRINES.

The five points in which Arminius differed from the Calvinistic system are clearly and prominently brought out in Wesley's writings, though mostly in the form of popular and practical exposition and exhortation. He put the name of Arminius on his periodical organ, and struck the keynote to the Arminian tone of Methodist preaching. The Arminian features of Methodism are, freedom of the will (taken in the sense of \emph{liberum arbitrium}, or power of contrary choice) as necessary to responsibility; self-limitation of divine sovereignty in its exercise and dealings with free agents; foreknowledge as preceding and conditioning foreordination; universality of redemption; resistibility of divine grace; possibility of total and final apostasy from the state of regeneration and sanctification.

Calvinism and Methodism agree in teaching man's salvation by God's free grace, in opposition to Pelagianism and Semipelagianism. But Calvinism traces salvation to the eternal purpose of God, and confines it to the elect; Methodism makes it dependent on man's free acceptance of that grace which is offered alike to all and on the same terms. Calvinism emphasizes the divine side, Methodism the human.\footnote{Dr. Warren, 1.c. p. 140, states the difference in an extreme form, which would subject Methodism to the charge of downright Pelagianism: '\emph{Nach der Methodistischen Auffassung des Heilsverhältnisses Gottes und des Menschen hängt das Heil oder Nicht-Heil eines jeden Menschen lediglich von seinem eigenen freien Verhalten gegenüber den erleuchtenden, erneuernden und heiligenden Einwirkungen des heiligen Geistes ab. Verhält man sich gegenüber diesen Einwirkungen empfänglich, so wird man hier, und einst dort, selig werden; verschliesst man sein Herz gegen dieselben, so wird man hier, und auf ewig im Tode verbleiben. Mit dieser Grundanschauung hängen alle sonstigen Eigenthümlichkeiten des Methodismus, wie z. B. seine eigenthümliche Freiheitslehre, seine Betonung der Wirksamkeit des heiligen Geistes, seine Lehre von der christlichen Vollkommenheit, und dergleichen, eng zusammen. Seinem innersten Geist und Wesen nach ist er eine Auffassung des Christenthums vom Standpunkte der christlichen Vollkommenheit oder der völligen Liebe.}'} Herein Methodism entirely agrees with Arminianism, and is even more emphatically opposed to the doctrines of absolute predestination, limited atonement, and the perseverance of saints than Arminius was, who left the last point undecided.

Wesley began the thunder against the imaginary horrors and blasphemies of Calvinism which has since resounded from innumerable Methodist pulpits. He defines predestination to be 'an eternal, unchangeable, irresistible decree of God, by virtue of which one part of mankind are infallibly saved, and the rest infallibly damned; it being impossible that any of the former should be damned, or that any of the latter should be saved;' and then he goes on to show that this doctrine makes all preaching useless; that it makes void the ordinance of God; that it tends directly to destroy holiness, meekness, and love, the comfort and happiness of religion, zeal for good works, and the whole Christian revelation; that it turns God into a hypocrite and deceiver; that it overturns his justice, mercy, and truth, and represents him 'as worse than the devil, more false, more cruel, and more unjust.' 'This,' he says, 'is the blasphemy clearly contained in the horrible decree of predestination, and for this I abhor it (however I love the persons who assert it).' To this decree he sets over the other decree, 'I will set before the sons of men life and death, blessing and cursing; and the soul that chooseth life shall live, as the soul that chooseth death shall die.' The elect are all those who 'suffer Christ to make them alive.'\footnote{Sermon liv., on Free Grace (Rom. viii. 32), preached at Bristol. It follows immediately after the eulogistic funeral discourse on the Calvinistic Whitefield. His brother Charles wrote a polemical poem on 'The Horrible Decree,' in which his poetic genius left him, as may be inferred from the following specimens:  \par 'O Horrible Decree, \par Worthy of whence it came! \par Forgive their hellish blasphemy, \par Who charge it on the Lamb.' \par 'To limit Thee they dare, \par Blaspheme Thee to Thy face, \par Deny their fellow-worms a share \par In Thy redeeming grace.' \par In another poem, on 'Predestination,' he prays: \par 'Increase (if that can be) \par The perfect hate I feel \par To Satan's Horrible Decree, \par That genuine child of hell; \par Which feigns thee to \emph{pass by} \par The most of Adam's race, \par And leave them in their blood to die, \par Shut out from saving grace.' \par How infinitely superior to these polemical effusions is his genuine hymn: \par 'Jesus, lover of my soul,' \par which a Calvinist may sing as heartily as a pious Methodist will join in his antagonist's (Toplady's): \par 'Rock of Ages, cleft for me.'}

The vehemence of this opposition to the doctrine of predestination must be explained in part from the subjective and emotional nature of Methodist piety, which exposes it much more to an antinomian abuse of this doctrine than is the case with the calm intellectual tendency, of Calvinism.

On the other hand, however, the 'evangelical' Arminianism of Wesley, as it is called, differs from the Dutch Arminianism, as developed by Episcopius and Limborch, and inclines as much towards Augustinianism as Arminianism inclines towards Pelagianism. In this respect it resembles somewhat the Lutheran anthropology of the Formula of Concord, though it differs altogether from its christology and sacramentalism.

1. Methodism holds a much stronger view of original sin than Arminianism, and regards it not simply as a disease or weakness,\footnote{Episcopius calls the \emph{peccatum originis} an \emph{infirmitas} or \emph{calamitas}or \emph{malum}, but not a \emph{malum culpæ} and \emph{malum pœnæ.} Limborch calls it \emph{malum naturale}, not \emph{peccatum nostri respectu.} See Winer, \emph{Comp. Symb.} pp. 60 sqq.} but as a total depravity that unfits man altogether for co-operation with the grace of God towards conversion. Wesley, Fletcher, and Watson describe this natural corruption in consequence of Adam's fall in the darkest colors, almost surpassing the descriptions of Augustine, Luther, and Calvin; but they deny the personal responsibility of Adam's posterity for his fall or the doctrine of original guilt; and herein they agree with the Arminians and the Quakers.

2. Methodism teaches the freedom of will as a gift of prevenient grace, which is given to every man as a check and antidote to original sin; while Arminianism, with its milder view of the fall, allows man a certain freedom of will in a weakened state as an inherent and inherited power of nature.

3. Methodism lays greater stress on the subjective experience of conversion and regeneration. Its preaching is essentially radical evangelistic revival preaching, which rouses the sinner to a sense of his danger, and the paramount necessity of an immediate, sudden, and radical change of heart and life.

THE ORIGINAL DOCTRINES OF METHODISM.

To these modifications of Arminianism must be added a few doctrines which Methodism claims as its own contributions to the better understanding of the Christian system.

1. The doctrine of the \emph{universality of divine grace}, not only in its intention, but in its \emph{actual offer.} Herein Methodism resembles the Quaker doctrine of universal light. It is assumed—on the ground of Paul's parallel between the first and second Adam (Rom. v.)—that all men are born into an order of saving grace, as well as into an order of sin. Adam brought a universal seed of death, but Christ brought a universal seed of life, which is available for all who do not reject it.\footnote{'No man living,' says Wesley, 'is without some preventing grace, and every degree of grace is a degree of life. There is a measure of free will supernaturally restored to every man, together with that supernatural light which enlightens every man that cometh into this world.' 'That by the offense of one, judgment came upon all men (all born into the world) unto condemnation, is an undoubted truth, and affects every infant as well as every adult person. But it is equally true that by the righteousness of One, the free gift came upon all men (all born into the world—infants and adults) unto justification.' D. D. Whedon (\emph{Biblioth. Sacra}, 1862, p. 258): 'Under the redemptive system, the man is born into the world, from Adam, a depraved being. It is as a depraved being that he becomes an Ego. But instantly after, in the order of nature, he is met by the provisions of atonement.' 'Every human being,' says Warren, 'has a measure of grace (unless he has cast it away), and those who faithfully use this intrusted gift will be accepted of God in the day of judgment, whether Jew or Greek, Christian or heathen. In virtue of Christ's mediation between God and the fallen race, all men since the first promise, Gen. iii. 15, are under an economy of grace, and the only difference between them as subjects of the moral government of God is that, while all have grace and light enough to attain salvation, some, over and above this, have more and others less' (Vol. I. pp. 146 sq.). Pope (pp. 239–248) distinguishes this doctrine from the Augustinian, Pelagian, Semipelagian, Tridentine, Lutheran, Calvinistic, and Arminian, and says that there is no doctrine which 'so irresistibly and universally appeals for its confirmation to the common conscience and judgment of mankind.'} For by virtue of the universal atonement, man, though born in sin, is held guiltless until he arrives at the point of personal responsibility.

While Romanism and Lutheranism save those only who are brought into contact with the Church and the Sacraments, Calvinism those only who are elect from eternity, Methodism brings the opportunity of salvation to all men in this present life, though in different forms and degrees, so that they are actually saved if they do not incur the guilt of rejecting salvation by unbelief. Hence all children are saved if they die before they commit actual sin. Though born in sin, they are not held guilty before the age of responsible agency. They are saved by the same power of the universal atonement which saves adults; though there is a difference of opinion as to the regeneration of infants before death.\footnote{Dr. D. D. Whedon (\emph{Biblioth. Sacra}, 1862, p. 258) remarks on this point: 'That the dying infant is saved, and saved by the atonement, we all agree. But his precise condition, as affected by the atonement, while a living infant, seems to be a somewhat undecided matter. Probably a large majority of the Methodist Episcopal Church have, for some time past, held, without much discussion, that the living infant was both unjustified and unregenerate, and yet upon his death he obtained both blessings. This making death the condition of justification and regeneration appears to many hardly logical, and not without danger. Mr. Wesley's earlier expressions of opinion indicated a holding of the churchly doctrine of baptismal regeneration in infancy. His later indications of opinion indicate that he held all infants to be members of the kingdom of heaven; and he also held that regeneration is a condition to membership in the kingdom of heaven; but he does not expressly draw the inference that all infants are regenerate. Fletcher maintained the doctrine both of infant justification and regeneration. Dr. Fisk held to infant justification. Our baptismal service first declares, in its Scripture lesson of infants, that ``of such is the kingdom of God,'' and yet declares ``that none can enter into the kingdom of God unless he be regenerate.'' But neither here is the inference expressly drawn. The subject is a matter of calm discussion, and perhaps the number of those holding the doctrine of infant regeneration has decidedly increased.'} On the same ground all heathen may be saved who do not neglect their opportunities. Ability and opportunity are the measure of responsibility, and God requires no more from man than he empowers him to perform. Christ's atonement covers the deficiency of ability in the case of infants, and the deficiency of opportunity in the case of the heathen.

Fletcher distinguishes three dispensations in this general economy of grace: the dispensation of the Father, embracing the heathen and Mohammedans, who know God only from his general revelation in nature, providence, and the conscience; the dispensation of the Son, for those who live within the limits of Christendom and the reach of the gospel; and the dispensation of the Holy Spirit, for those who have an experimental knowledge of the regenerating and sanctifying Spirit. Wesley, Watson, and Pope teach essentially the same view of the universality of grace.

2. The next distinctive doctrine of Methodism is the \emph{Witness of the Spirit} or the \emph{assurance} of salvation (Rom. viii. 15, 16). It is a double and concurrent witness of God's Spirit and of our spirit concerning our justification. The former is objective and divine, and antecedes; the latter is subjective and human, and follows. The Holy Spirit bears testimony to our spirit that by faith we are the children of God. This testimony is immediate and direct, and follows the work of justification and regeneration. On the ground of this testimony the believer feels \emph{assured} of his \emph{present} acceptance with God, and has a \emph{hope} of his \emph{final} salvation, but he is at the same time guarded against carnal security by the fear of a total and final fall from grace. Hence there are so many backsliders, who constitute a special class among Methodists.\footnote{Comp. the three sermons of Wesley on the Witness of the Spirit (x.–xii.), Vol. I. pp. 85 sqq. He traced this doctrine to his contact with some Moravians on his voyage to Georgia (1735), whose childlike trust and serene cheerfulness led him to exclaim: 'I, who went to America to convert others, was never myself converted to God.' He meant conversion from legal bondage to evangelical freedom and a sense of assurance of pardon. He subsequently visited Count Zinzendorf and the Moravians in Germany to study their discipline (1739). Watson (Vol. II. p. 271) distinguishes four views on the testimony of the Spirit, and thus states his own, which agrees with Wesley's: 'It is twofold; a direct testimony or ``inward impression on the soul, whereby the Spirit of God witnesses to my spirit that I am a child of God; that Christ hath loved me, and given himself for me, that I, even I, am reconciled to God'' (Wesley's \emph{Sermons}); and an indirect testimony, arising from the work of the Spirit in the heart and life, which St. Paul calls the testimony of our own spirit; for this is inferred from his expression, ``And the Spirit beareth witness with our spirit, etc.'' This testimony of our own spirit, or indirect testimony of the Holy Spirit by and through our own spirit, is considered confirmatory of the first testimony.' Pope (p. 465): 'Assurance is the fruit, not the essence of faith. . . .Perfect faith must be assured of its object. . . . The internal assurance of faith is a privilege that all may claim and expect; seasons of darkness and depression and uncertainty are only the trial of that faith of assurance.'}

Herein the Methodist doctrine differs from the Calvinistic doctrine of assurance which is based, not on subjective feeling, but on the divine promises and the unchangeable decree of God's election, and which covers not only the present state, but the whole process to its final completion, conditioned by the perseverance of saints as the final test of genuine conversion.\footnote{The Westminster Confession, Ch. XVIII., says that true believers 'may in this life be certainly assured that they are in a state of grace, and may rejoice in hope of the glory of God, which hope shall never make them ashamed.' This assurance is 'founded upon the divine truth of the promises of salvation, the testimony of the Spirit witnessing with our spirit that we are the children of God.' It is not of 'the essence of faith,' and may be 'shaken, diminished, and intermitted,' yet revived again in due time and keep us from utter despair.}

3. The last and crowning doctrine of Methodism, in which the Quakers likewise preceded it, is \emph{Perfectionism.} It is regarded as a mighty stimulus to progressive holiness, and forms the counterpart of the doctrine of apostasy, which acts as a warning against backsliding. It is derived from such passages as Matt. v. 48; Phil. iii. 15; Heb. vi. 1; x. 14; 1 John iii. 6; v. 18. Methodist perfection is not a sinless perfection or faultlessness, which Wesley denied,\footnote{In his sermons on Temptation, Vol. II. p. 215, and on Perfection, Vol. I. p. 356; Vol. II. p. 168: 'The highest perfection,' he says, 'which man can attain while the soul dwells in the body, does not exclude ignorance and error and a thousand infirmities.'} but a sort of imperfect perfection, from which it is possible to fall again temporarily or forever.\footnote{Meth. Catech. No. 3, p. 37: 'It is the privilege of every believer to be wholly sanctified, and to love God with all his heart in the present life; but at \emph{every stage of Christian experience there is danger of falling from grace}, which danger is to be guarded against by watchfulness, prayer, and a life of faith in the Son of God.'} It is entire sanctification or perfect love (1 John ii. 5; iv. 12), which every Christian may and ought to attain in this present life. From this state all voluntary transgressions or sinful volitions are excluded, though involuntary infirmities may and do remain; in this state all the normal qualities are possessed and enjoyed in their fullness. As to the attainment of perfection, it comes according to the prevailing view from gradual growth in grace, according to others by a special act of faith.\footnote{Wesley has two sermons on Christian Perfection, one on Phil. iii. 12 (Vol. I. p. 355), and one on Heb. vi. 1 (Vol. II. p. 167). He distinguishes, (1) angelic, (2) Adamic, (3) absolute perfection—all of which he denies to man in his present state—and (4) the relative perfection, which he claims for him under the gospel dispensation, namely, perfect love to God. From 1 John iii. 6 and v. 18, he reasons, 'A Christian is so far perfect as not to commit sin' (Vol. I. p. 365). He affirms that several persons have enjoyed this blessing of freedom from sin without interruption for many years, and not a few unto their death, as they have declared with their last breath (Vol. II. p. 174). Pope says (p. 527): 'The Spirit is imparted in this fullness for the perfect consecration of the soul to the Triune God: this is called the love of God perfected in us. The commandment requires from us in return the perfect love of the soul to God and man; and this perfection, promised to faith working by love, is abundantly attested as the possible and attained experience of Christians.' Pope distinguishes the Methodist theory of perfection from the ascetic, the fanatical, the Pelagian, the mystical, the Romanist, the imputationist (modern Calvinistic), and the Arminian (p. 535); and he mentions five characteristic marks of the Methodist doctrine, the chief of which is entire consecration to God in perfect love (p. 540).}

\xsection{Calvinistic Methodism.}

§ 112. Calvinistic Methodism.

WHITEFIELD.

George Whitefield labored with Wesley until 1741, when they parted on the question of predestination and free will; the former taking the Calvinistic, the latter, with his brother and the majority of Methodists, the Arminian side, and henceforth they pursued different paths, like Paul and Barnabas. Personally they became cordial friends again, and their friendship continued until death. This should not be forgotten when we read the bitter predestinarian controversy which their friends and followers carried on and renewed from time to time. When Whitefield heard of the dangerous illness of Wesley, who had already written his own epitaph, he sent him an affectionate letter (Dec. 3, 1753), saying, 'I pity myself and the Church, but not you. A radiant throne awaits you, and ere long you will enter into your Master's joy.'\footnote{See the whole letter in Tyerman, \emph{J. Wesley}, Vol. II. p. 175.} When Whitefield died in Newburyport (Sept. 30, 1770), Wesley preached his funeral sermon (Nov. 18) at Whitefield's Chapel in Tottenham Court Road and at the Tabernacle, near Moorfields, on the text Numb. xxiii. 10, 'Let me die the death of the righteous, and let my last end be like his!' Without alluding to their temporary separation, he speaks of him in the highest terms as an eminent instrument of God, who in the business of salvation put Christ; as high as possible, and man as low as possible, and who brought a larger number of sinners from darkness to the light than any other man. He praises his 'unparalleled zeal, his indefatigable activity, his tenderness of heart towards the afflicted, and charitableness to the poor, his deep gratitude, his most generous aud tender friendship, his modesty, frankness, patience, courage, and steadfastness to the end.'\footnote{\emph{Sermon LIII.} Vol. I. pp. 470 sqq.}

Whitefield was free from sectarian spirit and cared little for organization. His sole purpose was to convert sinners to Christ, and to revive Churches to new zeal and energy.\footnote{In this unselfish zeal he has a worthy successor in our day in Mr. Moody.} His labors were crowned with signal success. The day of judgment alone will reveal the number of his converts, and the amount of good which he kindled by his flaming sermons among Episcopalians, Presbyterians, Congregationalists, and other denominations, as well as among the crowds of ungodly people who were attracted by his eloquence.\footnote{'Whitefield's preaching was such as England never heard before—theatrical, extravagant, often commonplace, but hushing all criticism by its intense reality, its earnestness of belief, its deep, tremulous sympathy with the sin and sorrow of mankind. It was no common enthusiast who could wring gold from the close-fisted Franklin and admiration from the fastidious Horace Walpole, or who could look down from the top of a green knoll at Kingswood on twenty thousand colliers, grimy from the Bristol coal-pits, and see as he preached the tears ``making white channels down their blackened cheeks.'''—Green, \emph{History of the English People}, p. 718 (Engl. ed.). Dr. Abel Stevens, an Arminian Methodist, calls Whitefield 'the most eloquent, the most flaming preacher that the Christian Church has known since its apostolic age, whose eloquence sanctified, wakened the whole British empire' (\emph{Centenary of Amer. Methodism}, p. 24).}

But although most of his converts fell in with existing denominations, a considerable number of them formed three separate organizations. One of them, called 'the Whitefield Methodists,' were lost among the Independents. The other two still remain.

THE COUNTESS OF HUNTINGDON'S CONNECTION.

Selina, Countess of Huntingdon (1707–1791), a lady of true nobility of heart and intellect as well as rank, devoted, after the death of her husband and four children, her time and fortune to the spread of vital religion among the nobility and the court as well as the common people. She purchased halls and theatres in London, Bristol, and Dublin, built over sixty chapels, supported ministers, founded a college at Trevecca, in Wales, and stirred up others to similar liberality. She dispensed with her luxurious equipage and sold even her jewels for the benefit of this work. She took Whitefield, with whose Calvinism she sympathized, under her special patronage, and made him her chaplain, and exercised a sort of leadership over his congregations. Hence they became known as the 'Countess of (or Lady) Huntingdon's Connection.

Whitefield bequeathed to the Countess his benevolent institutions and lands in Georgia, and this resulted in a mission to America.

The ministers of this connection are almost identical in doctrine and Church polity with the Independents, but in public worship they use to some extent the Anglican Liturgy. Their principal institution is Chesunt College, in Herts.

THE WELSH CALVINISTIC METHODISTS.

Literature

\emph{The History, Constitution, Rules of Discipline, and Confession of Faith of the Calvinistic Methodists in Wales. Drawn up by their own Associated Ministers.} Third ed. Mold, 1840.

John Hughes:  \emph{History of Welsh Methodism}, (in Welsh). Liverpool, 1856, 3 vols.

William Williams:  \emph{Welsh Calvinistic Methodism. An Historical Sketch} (in English). London, 1872.

Whitefield's preaching through Wales, and the kindred labors of Howell Harris, of Trevecca, Griffith Jones, Daniel Rowlands, Howell Davies, and William Williams—most of them clergymen of the Established Church who joined the Methodists—produced a powerful and extensive revival, and resulted in a new connection in 1743, and more fully in 1785, when the Rev. Thomas Charles, of Bala, one of the most zealous and useful preachers of his day, joined it.\footnote{Charles graduated at Oxford as A.B. in 1778, labored seven years as a clergyman of the Established Church, united himself with the Calvinistic Methodists in 1785, and drew up in 1790 a series of Rules for conducting Associations or Quarterly Meetings. He was one of the founders of the British and Foreign Bible Society.}

For many years the Welsh Methodists existed without a settled form of government or doctrinal confession.

In 1823 it was unanimously agreed at the Associations of Aberystwith and Bala to issue in the Welsh language such a document, together with a sketch of the origin and early history of the denomination. An English edition was published in 1827.

The Confession of Faith consists of forty-four chapters, and accords substantially in spirit and arrangement with the Westminster Confession, though it is far inferior to it in ability and accuracy.

The articles in which it differs from the Wesleyan scheme are Arts. V., XII., and XXXIV., which are as follows:

V.—\emph{Of the Decree of God.}

God, from eternity, according to the counsel of his own will, and for the manifestation and exaltation of his glorious attributes, decreed all things which he should perform in time and to eternity, in the creation and governing of his creatures, and in the salvation of sinners of the human race; yet in such a manner that he is neither the Author of sin, nor does he force the will of his creatures in the fulfillment of his decree; and this decree of God is not depended on any thing in a creature, nor yet on the foreknowledge of God; but rather God knows that such and such circumstances will take place, because he has ordained that it should be so. God's decree is infinitely wise, perfectly righteous, and existing from eternity; it is a free, an ample, a secret, gracious, holy, good, an unchangeable and effectual decree.

XII.—\emph{Of the Election of Grace.}

God from eternity elected and ordained Christ to be a Covenant Head, a Mediator, and a Surety to his Church; to redeem and to save it. God also elected in Christ a countless multitude out of every tribe, tongue, people, and nation, to holiness and everlasting life; and every means were employed to effect this purpose most securely. This election is eternal, righteous, sovereign, unconditional, peculiar or personal, and unchangeable. It wrongs none, though God has justly left some without being elected, yet he has not wronged them: they are in the same condition as if there had been no election; and had there been no election, no flesh had been saved.

XXXIV.—\emph{Of Perseverance in Grace.}

Those whom God has made acceptable in the Beloved, whom he has effectually called, and whom the Spirit sanctifies, can not completely and forever fall from a state of grace, but they shall assuredly be supported unto the end, and they shall be saved. Their perseverance depends not on their own will, but on the unchangeableness of the purpose of God, the election of grace, the power of the Father's love, the sufficiency of the propitiation of Christ, the success of his intercession, union with him, the indwelling of the Spirit within them, the seed of God implanted in their souls, the nature and strength of the covenant, and the promise and oath of God. Founded on these things, perseverance is certain and unfailing. Though they may, through the temptations of Satan and the world, the great power of their indwelling corruption, and the neglect of using the means for their support, fall into sins, and remain in them for some time, and thus displease God, grieve the Holy Spirit, injure their grace, lose their comfort, harden their hearts, sting their consciences, draw a temporal judgment upon themselves, harm others, and disgrace the cause of God, yet they shall be kept by the power of God through faith to salvation, though their falls will be felt most bitterly by them.

Those who continue to live quietly in sin, and comfort themselves that they are in a gracious state, show evident signs that they are self-deceivers. For by perseverance in grace is not meant the continuing to enjoy and to inherit external Gospel privileges merely; but a continuance in holiness, diligence, watchfulness, a holy temper and walk, and a scrupulous observance of every duty. There is nothing more opposed to sin than a perseverance in grace; and whosoever shall thus continue in grace to the end shall be saved.

\xsection{The Catholic Apostolic Church (Called Irvingites).}

§ 113. The Catholic Apostolic Church (called Irvingites.)

Literature.

I. Sources.

Edward Irving:  \emph{Works}, collected and edited by his nephew, the Rev. G. Carlyle. London, 1864–65, 5 vols.

Michael Hohl:  \emph{Bruchstücke aus dem Leben und den Schriften E. Irving's. }St. Gallen, 1839; 2d ed. 1850.

Mrs.  M. O. W. Oliphant:  \emph{The Life of Edward Irving, Minister of the National Scotch Church, London, illustrated by his Journals and Correspondence. }London and New York (Harpers), 1862.

\emph{A Testimony to the King of England, and another to the Bishops of England. }London, 1836. (Anonymous. Prepared by the Apostles.)

\emph{A Testimony addressed to all Patriarchs, Archbishops, and Bishops, and the Reigning Sovereigns of Christendom. }1837. (Anonymous.)

\emph{Liturgy and other Divine Offices of the Church. }London, 1842. Drawn up by the 'Apostles,' and enlarged from time to time.

C. M. Carré:  \emph{The First and Last Days of the Church of Christ.} London, 1851.

\emph{Readings upon the Liturgy.} (By one of the Apostles.) London, 1852.

\emph{The Catechism. }(The English Episcopal Catechism enlarged.)

Thomas Carlyle (one of the Apostles): \emph{The Door of Hope for Britain, }and \emph{The Door of Hope for Christendom. }London, 1853. By the same: \emph{Apostles Given, Lost, Restored: Pleadings with my Mother.}

Rev.  William Dow (one of the Apostles, originally a Scotch Presbyt.): \emph{First Principles of the Doctrine of Christ. }Edinb. 1856. By the same: \emph{A Series of Discourses on practical and Doctrinal Subjects. }Edinb. 1853; 2d series, Edinb. 1860.

Rev.  J. S. Davenport:  \emph{Edward Irving and the Catholic Apostolic Church. }New York, 1863. By the same: \emph{Christian Unity and its Recovery. }New York, 1866. By the same: \emph{Letter to Bishop Whitehouse: The Church and the Episcopate. }Montreal, 1873.

W. W. Andrews:  \emph{The True Constitution of the Church and its Restoration. }New York, 1854. By the same: \emph{Review of Mrs. Oliphant's Life of E. Irving, }in the 'New-Englander' for July and Oct. 1863. By the same: \emph{The Catholic Apostolic Church, its History, Organization, Doctrine, and Worship, }in the 'Bibliotheca Sacra' for Jan. and April, 1866. Andover, Mass. By the same: \emph{The True Marks of the Church. }Hartford, 1867.

Rev.  Nicholas Armstrong (one of the Apostles): \emph{Sermons on Various Subjects. }2d ed. London, 1870. By the same: \emph{Homilies on the Epistles and Gospels.} London, 1870.

Rev.  T. Groser:  \emph{Sermons, }1st and 2d series. London, 1871 and 1874.

\emph{Apostles' Doctrine and Fellowship. }Anonymous. London, 1871.

\emph{The Purpose of God in Creation and Redemption.} Anonymous. 4th ed. Edinburgh, 1874.

\emph{Readings for the Sundays and Holydays of the Church's Year.} Anonymous. London, 1875.

\emph{The Dispensation of the Parousia.} Hartford, 1876.

Various writings of Henry Drummond (one of the Apostles), Chas. Böhm,  C. Rothe,  A. Köppen,  Ernst Gaab,  Rosstäuscher (author of an essay 'On the Gift of Tongues,' and a history of the movement under the title \emph{Der Aufbau der Kirche Christi auf den ursprünglichen Grundlagen}), and especially  H. W. J. Thiersch (the Tertullian of this modern Montanism, and its most learned minister in Germany, who wrote \emph{Lectures on Catholicism and Protestantism}, 1848, on the \emph{Canon of the N. T.}, 1845, on the \emph{Church in the Apostolic Age, }1852, and other excellent works).

II. Criticisms.

De Quincey, in \emph{Literary Reminiscences}, Vol. II.

Thomas Carlyle, in 'Fraser's Magazine' for Jan. 1835.

Articles on Irving in 'Edinburgh Review' for Oct. 1862; 'North British Review' for Aug. 1862; 'Blackwood's Magazine' for Nov. 1858, and June, 1862; 'London Quarterly Review' for Oct. 1862; 'Methodist Quarterly Review,' Jan. 1849, 1863.

Philip Schaff:  \emph{Der Irvingismus und die Kirchenfrage}, in his 'Deutscher Kirchenfreund,' Jahrg. III. 1850, pp. 49 sqq, 81 sqq. 161 sqq. 223 sqq. Mercersburg, Pa.

G. W. Lehmann:  \emph{Ueber die Irvingianer.} Hamburg, 1853.

Comp.  J. L. Jacobi:  \emph{Die Lehre der Irvingiten oder der sogenannten apostolischen Gemeinde verglichen mit der heiligen Schrift. }Berlin, 1868.

EDWARD IRVING.

Edward Irving, the herald and pioneer of the 'Catholic Apostolic Church,' was born at Annan, in Scotland, 1792, and died in the vigor of manhood at Glasgow, Dec. 8, 1834, where he lies buried in the crypt of the cathedral. He belonged to the Presbyterian Church, and for several years (1819–1822) labored in Glasgow as the assistant of the great and good Dr. Chalmers.

In 1822 he accepted a call to the Caledonian Chapel, Hatton Garden, London, and at once became the most powerful and popular preacher of the metropolis. He was at that time overflowing with bodily and spiritual life and energy. He excelled in the noblest gifts of eloquence, cultivated on the models of Hooker and Jeremy Taylor. Lofty thoughts clothed in gorgeous, semi-poetic language, devotional fervor, a solemn manner, a sonorous voice, a quaint antique style, a broad Scotch accent, an imposing figure, bushy hair flowing down in ringlets, a beaming face (which reminded Sir Walter Scott of that of the Saviour on Italian pictures), all combined to attract large and intelligent audiences, and to secure their closest attention, as if they listened to a messenger from the presence of the great Jehovah. De Quincey judged him to be, more than any man he ever saw, 'a son of thunder, and unquestionably by many degrees the greatest orator of our times.' He attracted people from all classes—noblemen, statesmen, and authors. When on a visit to Edinburgh and Glasgow, he roused the population at sunrise from their beds to hear his discourses. He shook the kingdom with his eloquence.

While he ruled like a monarch from his pulpit, he was a docile pupil of Coleridge, and received from the suggestive conversations of the old sage seeds of truth which seriously modified his Scotch Calvinistic creed. He now made more account of the incarnation and the true humanity of Christ, maintaining that he assumed our \emph{fallen}, i.e., temptable, mortal, corruptible nature, yet without sin itself, into complete fellowship with his divine person. This exposed him to the charge of denying the sinlessness of our Saviour, which was far from his thoughts. He also gave a large place to the hope of the glorious return of Christ, and the revival of the miraculous gifts of the Spirit in the Church.

In these views he was greatly strengthened by the sudden reappearance of what he believed to be the supernatural gifts of tongues, prophesying, and healing. These manifestations first occurred in the spring of 1830 in the west of Scotland, on the shores of the Clyde, among some pious Presbyterian men and women, who believed that their organs of speech were made use of by the Spirit of God for the utterance of his thoughts and intentions. Several persons from London, on hearing of these things, visited Scotland, and, on their return, held prayer-meetings in private houses, attended by devout members of different denominations. They united in supplications for the restoration of spiritual gifts. In April, 1831, the same manifestations took place among members of the Church of England and friends of Irving in London. The 'prophesyings' were addressed to the audience in intelligible English, and resembled the solemn exhortations of Quakers moved by the Spirit. The speaking in tongues consisted of soliloquies of the speaker, or dialogues between him and God which no one could understand. The burden of the prophetic utterances was the judgments impending on the apostate Church, the speedy coming of Christ, and the duty of preparing his way.\footnote{See \emph{A Brief Account of a Visit to some of the Brethren in the West of Scotland}, London, 1831 (J. Nisbet); Robert Baxter (first a believer in the divine origin and then in the satanic origin of these gifts): \emph{Narrative of Facts characteristic of the Supernatural Manifestations in Members of Mr. Irving's Congregation and other Individuals, in England and Scotland, and formerly in the Writer himself}, Lond. (Nisbit), 1833; Hohl, 1.c. (quoted in my \emph{Hist. of the Apost. Ch.} § 55, p. 198). Comp. also Stanley, \emph{Comment. on the Epp. to the Corinthians}, 4th ed. London, 1876, pp. 250 sqq.}

Similar manifestations of ecstatic utterances in seasons of powerful religious excitement appeared among the Montanists in the second century, the persecuted Protestants in France, called the 'Prophets of Cevennes,' and among the early Quakers.

These extraordinary proceedings naturally led to a rupture between Irving and the Presbytery of London (1832). He was turned out of the church built for him in Regent Square, and ultimately deposed from the ministry of the Church of Scotland by the Presbytery of Annan (1833), from which he had received his first license to preach.

On being driven from Regent Square, he was followed by the larger part of his congregation to Newman Street; and the following year, when his Presbyterian orders had been taken from him, he humbly submitted to reordination by one whom he received as an apostle. He never rose beyond the position of an 'angel,' or pastor, in the new Church, and, after less than two years of great labors and sufferings, passed from this world of trial into the regions of light.

He is little mentioned in the writings of his followers, and is regarded by them merely as a forerunner or John the Baptist, not as the founder of their community. His brilliant meteoric career, lofty character, and sad end created profound interest and sympathy. Dr. Chalmers, on hearing of his death, said that 'he was one in whom the graces of the humble Christian were joined to the virtues of the old Roman.' Thomas Carlyle, his countryman and early friend, thus characterizes Irving: 'He was appointed a Christian priest, and strove with the whole force that was in him to be it. I call him, upon the whole, the best man I have ever, after trial enough, found in this world, or now hope to find.'\footnote{When he adds, 'Oh foulest Circean draught, thou poison of popular applause! madness is in thee, and death; thy end is Bedlam and the grave,' he seems to cast a reflection on Irving's character which is not justified by facts; for Mrs. Oliphant's \emph{Life }shows him to have willingly sacrificed popularity to his convictions.}

THE CATHOLIC APOSTOLIC CHURCH.

This remarkable man, whose purity and piety can be as little doubted as his genius and eloquence, whatever may be thought of his soundness and judgment, gave the strongest if not the first impulse to the religious movement which, since its organization, is usually called after his name, but which calls itself 'The Catholic Apostolic Church.'\footnote{'They do not lay claim to the name \emph{Catholic Apostolic }as exclusively their own, but they use it as a proper designation of the one body of Christ, of which they are an organic part, and they refuse to be called by any other. They do this on the ground that it is wrong to affix to the Church the name of an eminent leader, like Luther or Calvin or Wesley; or one founded upon some feature of Church polity, such as Episcopal, Presbyterian, or Congregational; or one derived from some peculiar doctrine or rite, as Baptist or Free-will Baptist; or one expressing geographical limitations, such as Roman, Greek, Anglican, or Moravian. The essential characteristic of a thing should be expressed by its name, and the Church has for its three chief features, \emph{Unity}, as the only organism of which Christ is head; \emph{Catholicity}, as having a universal mission; and \emph{Apostolicity}, as sent by Christ into the world, even as he was sent by the Father. It is a significant fact that this name, adopted in the Nicene Creed, has practically every where been changed, as into the \emph{Roman }Catholic, the \emph{Greek }Orthodox, the \emph{Protestant }Episcopal, or something still narrower and more sectarian.'—W. W. Andrews, in \emph{Biblioth. Sacra}, 1.c.} It took full shape and form after his death, as it claims, under supernatural direction. It is one of the unsolved enigmas of Church history: it combines a high order of piety and humility of individual members with astounding assumptions, which, if well founded, would require the submission of all Christendom to the authority of its inspired apostles.

The modern 'Apostolic' Church believes and teaches that the Lord, who will soon appear in glory, has graciously restored, or at least begun to restore his one true Church, by reviving the primitive supernatural offices and gifts, which formed the bridal outfit of the apostolic age, but were soon afterwards lost or marred by the ingratitude and unbelief of Christendom. It claims to have apostles, prophets, and evangelists for the general care of the Church, and angels (or bishops), presbyters (or priests), and deacons for the care of particular congregations. All officers are called by the Holy Ghost through the voice of the prophets, except the deacons, who are chosen by the congregation as its representatives. They form a more complete hierarchy than that of the Episcopal or even the Greek and Roman Churches, whose bishops never claimed to be inspired apostles, but only successors of the apostles.

If the twelve modern apostles were truly called by Christ and endowed with all the powers and functions of that unique office, men will naturally look for sufficient evidence of the fact. But nine of these apostles died before 1876, and their vacancies have not been filled, nor are they expected to be filled. The Church, then, is relapsing into the same destitute condition which, according to their own theory, preceded this movement.\footnote{From a conversation with a learned minister of that Church, to whom I mentioned this difficulty, I infer that he at least—I do not know how many more—regards its testimony as a partial failure, or merely as a temporary provision, to be superseded by a better one. Another writes to me in answer to the same question: 'We are quite ready to admit failure, great failure, so far as to the present effects of the movement upon Christendom. But \emph{intrinsically}, and in relation to God's plans, we do not think it a failure.'} Their only hope is in the speedy return of our Lord.

To this apostolic hierarchy corresponds a highly ritualistic worship, with a solemn liturgy, based upon the Anglican and ancient Greek liturgies, and with an elaborate symbolism, derived from a fanciful interpretation of the Jewish tabernacle as a type of the worship of the Christian Church in the wilderness.

In this hierarchical constitution and ritualistic worship consists the chief peculiarity of this community. Its ministers and members have accordingly a very high idea of the Church and of the Sacraments. They are strict believers in baptismal regeneration and the real presence, though neither in the Roman nor the Lutheran sense. They reject transubstantiation and consubstantiation as well as the merely symbolical presence, and hold to the spiritual real presence of Calvin, but combine with it the view of Irenæus and other early fathers, that the elements, after being consecrated by the invocation of the Holy Ghost, have a heavenly and spiritual, as well as a material character, and are antitypes of the body and blood of Christ. They regard the eucharist as the centre of Christian worship, and not only as a sacrament, but also as a sacrifice in the patristic sense of a thank-offering, and they connect with it a commemoration of the departed. They are, upon the whole, the highest of High-Churchmen. They are in this respect the very antipodes of the Plymouth Brethren, the lowest of Low-Churchmen and the most independent of Independents, although both agree in their antagonism to the historical Churches and their expectation of the speedy coming of the Lord.

Yet, on the other hand, the Irvingites are unquestionably Protestant, and accept the positive results of the Reformation. They reject the Pope, not indeed as the Antichrist or 'the man of sin,' who will be revealed in the last times as the outgrowth of unbelief and lawlessness, but as an antichristian usurper of supreme authority in the Church. In their general belief they are as orthodox as any other denomination. They receive the whole Scriptures with devout reverence as their supreme guide. They lay stress on the œcumenical creeds, and embody them in their liturgical services. In catechetical instruction they use the Anglican Catechism, with an additional part inculcating their peculiar views about the constitution and order of the Christian Church. They manifest a catholic spirit, and sustain, as individuals, fraternal relations with members of other denominations. Upon the whole, they have most sympathy with the Episcopal Church, from which they received the majority of their original members. Of their apostles, eight were Anglicans (including two clergymen and two members of Parliament), three Presbyterians, and one Independent. Their main strength is in London, where they have seven churches, after the model of the seven churches in Asia Minor. They have also congregations in many of the principal cities in England and Scotland, and in some parts of the Continent of Europe, especially North Germany; while in Roman Catholic countries and in America they have made little or no progress.

The Irvingite movement has directed the attention of many serious minds to a deeper study of the supernatural order and outfit of the Apostolic Church, the divisions and reunion of Christendom, and the eschatological questions connected with the second advent.

STATEMENT OF THE REV. W. W. ANDREWS.

With these remarks we introduce a fuller inside account of the Catholic Apostolic Church, which was kindly prepared for this work by the Rev. W. W. Andrews, of Wethersfield, Conn. He has been thoroughly acquainted with the movement from the beginning, and is highly esteemed by all who know him as a Christian gentleman and scholar:

'The body of Christians who call themselves by the name of the Catholic Apostolic Church, not as exclusively their own, but because it is the proper designation of the one Catholic Church, is distinguished from all other Christian communions by the claim to the possession of gifts and ministries which, after having been long lost or suspended in their exercise, they believe to be now again restored to prepare the way for the coming and kingdom of the Lord.

History.

'The history of this religious movement can be given in few words. About the beginning of the second quarter of the present century, there was much prayer in many countries, but especially in Great Britain, for the outpouring of the Holy Ghost; and early in the year 1830 supernatural manifestations occurred in several parts of Scotland, in devout members of the Presbyterian Church, in the form of tongues, prophesyings, and healings. The following year similar manifestations took place in London, first in members of the Church of England, and afterwards among other religious bodies.

'Towards the end of the year 1832, by which time the supernatural character and divine origin of these spiritual phenomena had been abundantly attested, and a considerable number of persons had become believers, another and most important step was taken in the restoration of the apostolic office. The will of God that certain men should serve him as apostles was made known through supernatural utterances of the Holy Ghost by prophets, as when, at Antioch, he said, ``Separate me Barnabas and Saul for the work whereunto I have called them.'' The apostolate to the Gentiles, begun in the calling of Paul, but then left unfinished, the Lord now, at the end of the dispensation, set his hand to restore; and by the middle of the year 1835 the full number was completed, and they entered as a twelvefold Apostolic College on the work of caring for the whole Christian Church. As Great Britain had been chosen of God to be the centre of this catholic movement, one of the first duties laid upon the restored apostles was the preparing of a Testimony to the Bishops of the Church of England and Ireland, and of another to the King's Privy Council, in which they pointed out the sins and perils of those lands, and testified to the coming of the Lord as the only hope of mankind, and to the work of the Holy Ghost as the necessary means of preparation.

'A year or two later, they addressed a more full and complete testimony, of the same general character, to all the Rulers in Church and State throughout Christendom. They did this, because it was their duty, from the nature of their office, to seek the blessing of the whole flock of God. Apostles alone have universal jurisdiction, as they alone receive their commission directly from the Lord; and it belonged to them, when restored towards the close of the long history of the Church, to take up those questions in respect to doctrine, organization, and worship which had broken the unity of Christendom; and having examined the creeds and rites and usages of every part, to separate the evil from the good, and to stamp with their apostolic authority every fragment of divine truth and order which had been preserved. This they have been doing for more than forty years, and the results to which they have arrived may be thus briefly stated.

Doctrines.

'They hold the holy Scriptures of the Old and New Testaments (the Old Testament as received from the Jews, without the Apocrypha) to contain the sum and substance of all divine revelations, and therefore to be the supreme and infallible standard of doctrine.

'But they also believe that Christ's promise to be with his Church to the end of the world has not been made void, and that the Holy Spirit has borne a living witness to the one faith in all generations; and they have adopted the three great creeds commonly called the Apostles', the Nicene, and the Athanasian, as expressing more clearly than any others the belief of the Universal Church. The great doctrines of the holy Trinity, the incarnation, the atoning death and bodily resurrection of the Lord, his ascension and high priestly work in heaven, the descent of the Holy Ghost to draw men to Christ, and to regenerate, sanctify, and endow with heavenly gifts them that believe, together with the second personal coming of the Lord to judge the quick and the dead, and to administer eternal retributions, they hold in their plain and obvious import, in harmony with the whole Orthodox Church, Greek, Roman, and Protestant.

'These creeds they have appointed to be used in divine worship: the Apostles', at the daily morning and evening services; the Nicene, in the ordinary celebrations of the eucharist on the Lord's Day; and the Athanasian, four times in the year, at the great festivals of Christmas, Easter, Pentecost, and All-Saints. They use the Nicene Creed in the form in which the Western Church receives it, retaining the \emph{Filioque}, but not condemning the Eastern Church for using it in the form in which it was left by the Council of Constantinople.

'In respect to the great central truth of the incarnation, the key to all the purposes and works of God, they teach that the second Person in the adorable Godhead, the only and eternally begotten Son, became man by assuming our entire humanity—body, soul, and spirit —under the conditions of the fall, but without sin, through the overshadowing of the Holy Ghost. They reject, therefore, the dogma of the immaculate conception of the mother of the Lord as against the truth of holy Scripture, which declares the whole human race to have been involved in the fall of the first Adam. They teach that by being born of a mother of the fallen race, he took the common nature of man, with all its infirmities, burdens, and liabilities, so that he could be tempted in all points like as we are, and be dealt with in all things by the Father as the representative of mankind. But they also make prominent the work of the Holy Ghost in effecting the incarnation, holding that it was through his presence and power that the Son of God was conceived of the Virgin Mary, and afterwards anointed for his public ministry; so that while it was a divine person who became incarnate, he had no advantage of his Godhead in his earthly life, but did every thing as man upheld, guided, and energized by the Holy Ghost.

'They hold, with the Church of England, and all the great leaders of the Reformation, that the death of the Lord Jesus Christ was ``a full, perfect, and sufficient sacrifice, oblation, and satisfaction for the sins of the whole world'' (and not merely for those of the elect); and on this ground they stand aloof both from the rationalism which denies its vicarious and expiatory nature, and from the Roman doctrine of the mass, which teaches that the sacrifice of the cross needs to be supplemented by the sacrifices of the eucharist, in which the Lamb of God is continually immolated afresh.

'But they go beyond the theology of the Reformation in respect to the Church, which they look upon as the fruit of the death and resurrection of Christ, and of the descent of the Holy Ghost which followed his ascension; and as differing, therefore, fundamentally in its spiritual essence and prerogatives from all the companies of the faithful in the preceding dispensations. They believe that in rising from the dead he became the fountain of a new life, the head of a redeemed humanity, of which those who believe in him are made partakers by the operations of the Holy Ghost working in and through the ordinances of his Church. The Sacraments of Baptism and the Lord's Supper are the divinely appointed means of conveying and nourishing this new life of his resurrection, by the implanting and energizing of which the whole multitude of the faithful are made to be the One Body of Christ.

'As to the structure and endowments of the Church, they hold that its original constitution contains the abiding law for all generations. The fourfold ministry of apostles, prophets, evangelists, and pastors, first fulfilled by the Lord himself when upon earth, was continued in his Church after he had gone into heaven, because it was the necessary instrumentality of conveying his manifold grace and blessing, and of bringing his Body to the stature of his fullness (Eph. iv. 11–16). The Holy Ghost was given to be the permanent possession of his people; and the apostles reject the common distinction between ordinary and extraordinary gifts as wholly unscriptural, and as restraining the manifestations of the Spirit. They lay great stress upon the connection of the descent of the Comforter with the glorifying of the Lord Jesus (John vii. 39), and teach that the object of his mission was to reveal the glory and manifest the energies of the Man whom God had exalted from the weakness and dishonor of the grave to his own right hand. Supernatural gifts and miraculous workings are therefore in accordance with the nature of the dispensation, which began with the resurrection of the Lord, and is to end with the resurrection of his saints.

Worship.

'Its chief features are the celebration of the Eucharist on every Lord's day; services at six in the morning and five in the evening of every day in the year, requiring for their complete fulfillment the three ministries of angel, priests, and deacons; the observance of the great feasts of the Church, excluding those in honor of particular saints; and a monthly service by the seven churches in London gathered into one as a symbol of the Universal Church, which is also observed in all the congregations throughout the world.

'The holy Eucharist is made to be the centre of worship, of which Christ, the great High-Priest in the heavens, is the leader, and the Mosaic ritual the shadow and type. The showing to the Father of that one sacrifice of the cross, which is the basis of all intercession, is effected by the Lord himself, by his own bodily presence in heaven; and the Church is enabled to do the same upon the earth by means of that sacrament in which he places in her hands the symbols and spiritual reality of his body and blood. The eucharist is regarded as the antitype of the priestly act of Melchizedek in bringing forth bread and wine to Abraham, the father of the faithful, from whom he received the tenth of all; and in the offertory, both the tithes and the offerings of the people are brought up and presented to God as an act of worship.

'As the death of the cross was itself the fulfillment of all the bloody sacrifices of the Law, the commemoration of it in the holy Supper becomes the distinguishing Christian rite, from which all other acts of worship, especially the daily morning and evening services—the antitype of the daily services of the Tabernacle—derive their life and power. All the purest and most catholic parts of all the rituals of Christendom have been gathered up and woven together, to form, with such additions as the present exigencies of the Church demand, a comprehensive and organic system of worship, at once purely Scriptural, and embodying the richest liturgical treasures of the past. Among the errors and superstitions which have been weeded out are transubstantiation, the worship of the Virgin Mary and of saints and angels, the use of images and pictures, and prayers for deliverance from purgatorial fires. But in rejecting the corruption of the truth, the truth itself has not been cast away; and the doctrine of the real presence (as a spiritual mystery involving no physical change of the elements), the thankful and reverential mention of the Mother of the Lord (``And with the holy angels, and with thy Church in all generations, we call her blessed''), and continual supplications and intercessions in behalf of the faithful departed, that they ``may rest in the peace of God, and awake to a joyful resurrection,'' all have place in the services appointed by the apostles.

Organization and Unity of the Church.

'The unity of the Church is held as a fundamental fact, resulting from the acts and operations of God, and not from the agreements and confederacies of men. There is one Body of Christ, embracing all who have been baptized in the name of the Father, and of the Son, and of the Holy Ghost (though, like the unfruitful branches of the vine, many may at last be cut out and cast away); and, in the absence of the Head, the harmonious intercommunion of the members is secured by the inworking of the One Spirit, and by a ministry proceeding immediately from the Head, and having jurisdiction over all the parts. The distinction between the Church Universal and the local or particular churches which compose it, is sharply drawn in the organization which has been developed under the rule of the apostles. The apostles themselves are the great Catholic ministry, through which guidance and blessing are conveyed to the whole body, and they are assisted in their work by prophets, evangelists, and pastors.

'But each particular church, when fully organized, is under the rule of an angel, or chief pastor or bishop, with presbyters and deacons helping him in their subordinate places. It is his office to stand continually at his own altar at the head of his flock, carrying on the worship of God, cherishing and directing the gifts of the Holy Ghost, and exercising the pastoral charge over all the souls committed to his care. The threefold ministry of Episcopacy (and, in a lower form, of Presbyterianism) is here united with the central authority which Rome has wrongfully sought to attain by exalting her bishop to the place of universal headship.

The Second Coming.

'In respect to eschatology, they hold, with the Church of the first three centuries, that the second coming of the Lord precedes and introduces the millennium; at the beginning of which the first resurrection takes place, and at the close the general resurrection, with the final judgment and its eternal retributions to the righteous and the wicked. This period of a thousand years will be marked by the presence of the Lord and his risen and translated saints upon or in near proximity to the earth, then freed, at least partially, from the curse; by the re-establishment of the tribes of Israel in their own land, in fulfillment of the promises to their fathers, with Jerusalem rebuilt, to be the metropolitan centre of blessing to all nations; and by the bringing of all the families of mankind into the obedience and order and blessedness of the kingdom of God.

'The restoration of the primitive gifts and ministries, like the ministries of Noah and of John the Baptist at the close of the antediluvian and Jewish dispensations, is to prepare for the ushering in of this next stage of God's actings. The order of events is to be as follows: The immediate and special work of the apostles is to gather and make ready a company of first-fruits, described (Rev. vii. 1–8) as sealed with the seal of the living God—the gift of the Holy Ghost bestowed by the hands of the apostles (Eph. i. 13; Acts xix. 1–5)—and as organized after a twelvefold law, of which the type was given in the structure of the twelve tribes of Israel. They are sealed while the angels are holding back the winds of judgment, before the great tribulation (Rev. vii. 14) is let loose upon the earth, that in them the Lord's words may be fulfilled, and they be counted worthy to escape all the things that are coming to pass, and to stand before the Son of Man (Luke xxi. 36).

'But the taking away of the first-fruits is only the first stage of the mighty work to be done in the bringing of this dispensation to a close. It is to be followed by the revelation of the Man of Sin, the infidel Antichrist, who will be successfully resisted for a time by the two witnesses (Rev. xi. 3–12), but will at length prevail over them, and for a short time rule the nations with the tyrannizing power and lurid splendors of the pit. In the midst of the terrors of that great tribulation the harvest will be reaped, and all the faithful gathered into the garner of the great Husbandman; and thereupon will be the vintage of wrath (Rev. xiv. 15–20), and the Lord will come forth to tread the wine-press of his Father's indignation, and to cast the beast and the false prophet into the lake of fire.

Prospects.

'They regard the failure of their labors to gather the Churches of Christendom into their communion as being after the analogy of the failures at the close of all preceding dispensations, and as furnishing no argument against the reality of their divine mission.

'The apostles do not, therefore, expect to have a large following at this stage of God's work. As a sheaf of first-fruits to the harvest, such will be the relation of the few who receive their testimony to the great multitude who will be saved out of the fiery trial of the time of the Antichrist. Nor does their faith fail because many of their brethren have been taken away by death, and it has pleased God to leave their places unfilled; for they look upon this as an indication that their present work is nearly finished, and that the Lord will soon take those who shall be found ready, to stand with him upon Mount Zion, safe in his hiding-place, while he pours out the vials of his wrath upon the earth. It would seem that the two apostolates at the beginning and the end of the dispensation form the company of the four-and-twenty elders who sit on thrones around the throne of the great King (Rev. iv. 4), partakers of his dominion, and associated with him in his work of judgment and rule.

Relation to other Churches.

'This brief statement of the position and doctrines of the ``Catholic Apostolic Church'' shows the grounds of their refusal to be called by any other name than belongs to the whole community of the baptized. They are a part of the one Church, differing from their brethren in being gathered under the proper ministries of the Church universal, and in being organized according to the original law of the Church as defined by St. Paul when speaking of the Body of Christ (1 Cor. xii.). They hold the one faith, the one hope, and the one baptism; and, without departing from the exact and literal teachings of the New Testament, they have added to these the larger statements of truth which have been the fruits of God's presence with his Church through all her generations.

'Having its origin among the Protestant Churches, and retaining all the great truths pertaining to the cross of Christ, for which the Reformation was a noble and successful struggle, this Catholic work has laid under contribution the rich stores of the Greek and Roman communions, and is leading the Church on into still deeper knowledge of the purposes of God contained in holy Scripture, by means of the living ministers of Christ and the revelations of the Holy Ghost, to the end of preparing her as a bride for the marriage of the Lamb.'

\xsection{The Articles of the Evangelical Alliance.}

§ 114. The Articles of the Evangelical Alliance.

Literature.

\emph{Report of the Proceedings of the Conference, held at Freemasons' Hall, London, from August} 19\emph{th to September }2\emph{d inclusive}, 1846. \emph{Published by Order of the Conference. }London (Partridge \& Oakey, Paternoster Row), 1847.

Comp. also the \emph{Proceedings of the Seven General Conferences of the Alliance}, held at \emph{London}, 1851, \emph{Paris}, 1855, \emph{Berlin}, 1857, \emph{Geneva}, 1861, \emph{Amsterdam}, 1867, \emph{New York}, 1873, and \emph{Basle}, 1879, all published in English, some also in the German, French, Dutch, and other languages.

The General Conference of New York, the first held on American soil, was the most important, and its proceedings (published by Harper \& Brothers, N. Y. 1874) form an interesting panoramic view of the intellectual and spiritual state of the Christian world at that time.

CHARACTER AND AIM OF THE ALLIANCE.

The 'Evangelical Alliance' is not an ecclesiastical organization, and has, therefore, no authority to issue and enforce an ecclesiastical creed or confession of faith. It is a voluntary society for the manifestation and promotion of Christian union, and for the protection of religious liberty. Its object is not to bring about an organic union of Churches, nor a confederation of independent Churches, but to exhibit, and to strengthen union and co-operation among individual members of different Protestant denominations without interfering with their respective creeds and internal affairs. It aims to realize the idea of such a Christian union as is consistent with denominational distinctions and varieties in doctrine, worship, and government. It may ultimately lead to a closer approximation of the Churches themselves, but it may and does exist without ecclesiastical union; and ecclesiastical union would be worthless without Christian union. It is remarkable that our Lord, in his sacerdotal prayer, which is the magna charta of Christian union, makes no reference to the Church or to any outward organization. The communion of saints has its source and centre in their union with Christ, and this reflects his union with the Father.

The Alliance extends to all nationalities and languages, but is confined, so far, to Christians who hold what is understood to be the Scriptural or evangelical system of faith as professed by the Churches of the Reformation and their legitimate descendants. It thus embraces Episcopalians, Lutherans, Presbyterians, Congregationalists, Baptists, Methodists, Moravians, and other orthodox Protestants, but it excludes Roman and Greek Catholics on the one hand, and the antitrinitarian Protestants on the other. The Quakers, though unwisely excluded by Art. IX., are in full sympathy with one of the two chief objects of the Alliance—the advocacy of religious liberty.

THE CONFERENCE OF 1846.

The call to the London Conference of 1846 for the formation of an Evangelical Alliance against infidelity was sufficiently liberal to encourage all orthodox Protestants to attend without doing any violence to their confessional conscience. But the High-Church elements, from aversion to miscellaneous ecclesiastical company, kept aloof, and left the enterprise in the hands of the evangelical Low-Church and Broad-Church ranks of Protestantism. The meeting was overwhelmingly English, and controlled by Episcopalians, Scotch Presbyterians, and English Dissenters. Next to them, America was best represented, and exerted the most influence. The delegation from the Continent was numerically small, but highly respectable. The whole number of attendants was over eight hundred ministers and laymen, from about fifty distinct ecclesiastical organizations of Protestant Christendom, among them many scholars and ministers of the highest Christian standing in their respective Churches and countries. Those who took the most active part in the proceedings were Sir Culling Eardley Smith (President), E. Bickersteth, B. W. Noel, W. M. Bunting, J. Angell James, Dr. Steane, Wm. Arthur, T. Binney, O. Winslow, Andrew Reed, of England; Norman Macleod, W. Cunningham, W. Arnot, R. Buchanan, James Begg, James Henderson, Ralph Wardlaw, of Scotland; Drs. Samuel H. Cox, Lyman Beecher, W. Patton, Robert Baird, Thomas Skinner, E. W. Kirk, S. S. Schmucker, of the United States; Drs. Tholuck, W. Hoffmann, E. Kuntze, of Germany; Adolphe Monod, Georges Fisch, La Harpe, of France and Switzerland. The meeting was one of unusual enthusiasm and interest. One of its most eloquent speakers, Dr. Samuel H. Cox, of New York, characterized it as an assembly

'Such as earth saw never,

Such as Heaven stoops down to see.'

The late Dr. Norman Macleod wrote during the meeting, in a private letter recently brought to light:\footnote{\emph{Memoir, by his Brother}, 1876, Vol. I. p. 260 (N.Y. ed.). The letter to his sister dated Aug. 4, 1846, should be dated Aug. 24.} 'I have just time to say that our Alliance goes on nobly. There are one thousand members met from all the world, and the prayers and praises would melt your heart. Wardlaw, Bickersteth, and Tholuck say that in their whole experience they never beheld any thing like it. . . . It is much more like heaven than any thing I ever experienced on earth.'

THE DOCTRINAL BASIS.

The part of the proceedings with which we are concerned here is the attempt made to set forth the doctrinal consensus of evangelical Christendom as a basis for the promotion of Christian union and religious liberty.

The Rev. Edward Bickersteth, Rector of Walton, Herts, and one of the leaders of the evangelical party in the Established Church of England, moved the adoption of the doctrinal basis, and Dr. S. H. Cox, a Presbyterian of New York, supported it in a stirring speech, on the third day (Aug. 21). After considerable discussion and some unessential modifications, the basis was adopted on the fifth day (Aug. 24), \emph{nemine contradicente;} the vast majority raising their hands in approval, the rest abstaining from voting. The chairman then gave out the hymn,

'All hail the great Immanuel's name,

Let angels prostrate fall.'

It 'was sung by the Conference with a depth of devotional feeling which, even during the meetings of the Conference, had never been surpassed.'\footnote{\emph{Proceedings}, p. 193.}

The doctrinal basis is expressly declared 'not to be a creed or confession in any formal or ecclesiastical sense, but simply an indication of the class of persons whom it is desirable to embrace within the Alliance.' It consists of nine articles: (1) the divine inspiration and supreme authority of the Holy Scriptures; (2) the right and duty of private judgment in the interpretation of the Holy Scriptures; (3) the unity and trinity of the Godhead; (4) the total depravity of man in consequence of the fall; (5) the incarnation of the Son of God, his atonement, and his mediatorial intercession and reign; (6) justification by faith alone; (7) the work of the Holy Spirit in conversion and sanctification; (8) the immortality of the soul, the resurrection of the body, the judgment of the world by Jesus Christ, with the eternal blessedness of the righteous and the eternal punishment of the wicked; (9) the divine institution of the Christian ministry, and the perpetuity of Baptism and the Lord's Supper.

The basis is merely a skeleton: it affirms 'what are usually understood to be evangelical views' on the nine articles enumerated. To give an explicit statement of these views would require a high order of theological wisdom and circumspection. For the practical purpose of the Alliance, the doctrinal basis has upon the whole proved sufficient, though some would have it more strict, others more liberal, since it excludes the orthodox Quakers. It has been variously modified and liberalized by branch Alliances in calling General Conferences. The American branch, at its organization in New York, Jan., 1867, adopted it with a qualifying preamble, subordinating it to the more general consensus of Christendom, and allowing considerable latitude in its construction.\footnote{See Vol. III. p. 821.}

\xsection{The Consensus and Dissensus of Creeds.}

§ 115. The Consensus and Dissensus of Creeds.

Philip Schaff:  \emph{The Antagonisms of Creeds}, in the ``Contemporary Review,'' London, Oct 1876, (Vol. XI. pp. 836–850). \emph{The Consensus of the Reformed Confessions}, In the ``Proceedings of the First Gen. Pres. Council,'' Edinburgh, 1877; separately issued, New York, 1877.

The Creeds of orthodox Christendom have passed before us. A concluding summary of the points of agreement and disagreement will aid the reader in forming an intelligent judgment on the possibility, nature, and extent of an ultimate adjustment of the doctrinal antagonisms which are embodied and perpetuated in the symbols of the historic Churches. The argumentation from Scripture, tradition, and reason belongs to the science of Symbolics.

A. The Catholic Consensus of Greek, Latin, and Evangelical Christendom.

The Consensus is contained in the Scriptures, and in the œcumenical Creeds which all orthodox Churches adopt. It may be more fully and clearly specified as follows:

I.—RULE OF FAITH AND PRACTICE.

The Divine Inspiration and Authority of the Canonical Scriptures in matters of faith and morals. (Against Rationalism.)

II.—THEOLOGY.

1. The Unity of the Divine essence. (Against Atheism, Dualism, Polytheism.)

2. The Trinity of the Divine Persons.

Father, Son, and Holy Ghost, the Maker, Redeemer, and Sanctifier.

(Against Arianism, Socinianism, Unitarianism.)

3. The Divine perfections.

Omnipotence, omnipresence, omniscience, wisdom, holiness, justice, love, and mercy.

4. Creation of the world by the will of God out of nothing for his glory and the happiness of his creatures. (Against Materialism, Pantheism, Atheism.)

5. Government of the world by Divine Providence.

III.—ANTHROPOLOGY.

1. Original innocence.

Man made in the image of God, with reason and freedom, pure and holy; yet needing probation, and liable to fall.

2. Fall: sin and death.

Natural depravity and guilt; necessity and possibility of salvation. (Against Pelagianism and Manichæism.)

3. Redemption by Christ.

IV.—CHRISTOLOGY.

1. The Incarnation of the eternal Logos or second Person in the Holy Trinity.

2. The Divine-human constitution of the Person of Christ.

3. The life of Christ.

His superhuman conception; his sinless perfection; his crucifixion, death, and burial; resurrection and ascension; sitting at the right hand of God; return to judgment.

4. Christ our Prophet, Priest, and King forever.

5. The mediatorial work of Christ, or the atonement.

'He died for our sins, and rose for our justification.'

V.—PNEUMATOLOGY.

1. The Divine Personality of the Holy Spirit.

2. His eternal Procession (\textgreek{ἐκπόρευσις, } \emph{processio}) from the Father, and his historic Mission (\textgreek{πέμψις, } \emph{missio}) by the Father and the Son.

3. His Divine work of regeneration and sanctification.

VI.—SOTERIOLOGY.

1. Eternal predestination or election of believers to salvation.

2. Call by the gospel.

3. Regeneration and conversion. Necessity of repentance and faith.

4. Justification and sanctification.

Forgiveness of sins and necessity of a holy lire.

5. Glorification of believers.

VII.—ECCLESIOLOGY AND SACRAMENTOLOGY.

1. Divine origin and constitution of the Catholic Church of Christ.

2. The essential attributes of the Church universal.

Unity, catholicity, holiness, and indestructibility of the Church. Church militant and Church triumphant.

3. The ministry of the gospel.

4. The preaching of the gospel.

5. Sacraments: visible signs, seals, and means of grace.

6. Baptism for the remission of sins.

7. The Lord's Supper for the commemoration of the atoning death of Christ.

VIII.—ESCHATOLOGY.

1. Death in consequence of sin.

2. Immortality of the soul.

3. The final coming of Christ.

4. General resurrection.

5. Judgment of the world by our Lord Jesus Christ.

6. Heaven and Hell.

The eternal blessedness of saints, and the eternal punishment of the wicked.

7. God all in all (1 Cor. xv. 28).

B. Consensus and Dissensus of the Greek and Roman Churches.

(\emph{a}) CONSENSUS.

I. The articles of the œcumenical Creeds, excepting the \emph{Filioque} of the Latin recension of the Nicene Creed and the \emph{et filio} of the Athanasian Creed.

II. Most of the post-œcumenical doctrines, which are not contained in the œcumenical Creeds, and from which Protestants dissent, viz.:

1. The authority of ecclesiastical tradition, as a joint rule of faith with the Scriptures.

2. The worship (\textgreek{τιμητικὴ προσκύνησις}) of the Virgin Mary, the Saints, their pictures (not statues), and relics.

3. The infallibility of the Church—that is, the teaching hierarchy (\emph{ecclesia docens}). The Roman Church lodges infallibility in the papal monarchy, the Greek Church in the (seven) œcumenical Councils, and the patriarchal oligarchy as a whole.\footnote{We say as a \emph{whole}; for the Greek Church does not claim infallibility for any individual patriarch, and has herself condemned, in œcumenical Synods, as heretics not only Pope Honorius, of Rome, but also several of her own patriarchs, e.g., Nestorius, of Constantinople; Dioscurus, of Alexandria; Peter the Fuller, of Antioch; Sallustius, of Jerusalem; Cyril Lucar, of Constantinople.}

4. Justification by faith and works, as joint conditions.

5. The Seven Sacraments or Mysteries, with minor differences as to confirmation and unction.

6. Baptismal regeneration (in an unqualified sense), and the necessity of water-baptism for salvation.

7. Priestly absolution by divine authority.

8. Transubstantiation (\textgreek{μετουσίωσις}), and the adoration of the consecrated elements.

9. The sacrifice of the Mass for the living and the dead.

This forms the centre of worship. Preaching is subordinate.

10. Prayers for the departed.

On the authority of the Apocryphal books of the Old Testament, transubstantiation, Purgatory, and a few other points, the Greek doctrine is not so clearly developed and formulated; but, upon the whole, much nearer the Roman view than the Protestant.

As to the popular use of the Bible, there is this important difference, that the Greek Church has never prohibited it, like the Roman, and that the Russian Church has recently favored it, and thus opened the way for a wholesome progress and possible reformation

(\emph{b}) DISSENSUS.

I. The eternal Procession of the Holy Ghost from the Son (\emph{Filioque}): denied by the Greek, taught by the Latin Church.

II. The papal supremacy and infallibility: rejected by the Greek Church as an antichristian usurpation, asserted by the Latin Church as its corner-stone.

III. The immaculate conception of the Virgin Mary: proclaimed as a dogma by the Pope, 1854:\footnote{The Greek Archbishop Lykurgos, of Syra and Tenos (d. 1876), declared, while in England, in a conference with the Bishop of Ely, Feb. 4, 1870: 'The Orthodox Church considers the immaculate conception to be blasphemous. It destroys the doctrine of the Incarnation.' But in practice the worship of the blessed Virgin is carried as far in the Greek Church as in the Latin.}

IV. The marriage of the lower clergy: allowed by the Greek, forbidden by the Latin Church.

V. Withdrawal of the eucharistic cup from the laity.

VI. A number of rites and ceremonies.

Greek rites: threefold baptismal immersion, instead of pouring or sprinkling; use of leavened, instead of unleavened, bread in the eucharist; the invocation of the Holy Ghost for the benediction of the sacred elements; infant communion; anointing baptized infants; the repetition of holy unction (\textgreek{τὸ εὐχέλαιον}) in sickness.

C. Consensus and Dissensus of the Greek Church and the Evangelical Churches.

(\emph{a}) CONSENSUS.

I. They believe the Scriptures and the doctrines of the œcumenical Creeds. (See A.)

II. They reject:

1. The supremacy and infallibility of the Pope.

2. The immaculate conception of the Virgin Mary.

3. The withdrawal of the cup from the laity.

4. The enforced celibacy of priests and deacons.

(The Greek Church, however, prohibits the second marriage of the lower clergy, and requires the celibacy of the bishops.)

(\emph{b}) DISSENSUS.

I. The double Procession of the Holy Spirit.\footnote{In this doctrine the Protestant Confessions side with the Latin Church, or at least they do not oppose it. The eternal procession of the Spirit was no topic of controversy in the period of the Reformation, and may be regarded as an open question subject to further exegetical and theological investigation. A number of Episcopalians in England and America would be willing to expunge the \emph{Filioque} from the Nicene Creed, or to compromise with the Orientals on the single procession of the Spirit from the Father \emph{through the Son.} See the Theses of the Bonn Conference of 1875, at the close of Vol. II.}

II. In the post-œcumenical doctrines mentioned sub B. (\emph{a}), II., the Greek Church sides with Rome against Protestantism.

D. Consensus and Dissensus of the Roman Catholic and the Evangelical Protestant Churches.

(\emph{a}) CONSENSUS. (See sub A.)

(\emph{b}) DISSENSUS.

I. Scripture and Tradition, as a rule of faith.

Roman Catholic doctrine:

The necessity of ecclesiastical tradition (culminating in the infallible decisions of the papal see), as a joint rule of faith and as the sole interpreter of Scripture.

Protestant doctrine:

The absolute supremacy and sufficiency of the Scriptures as a guide to salvation.

II. Other differences concerning the Scriptures.

1. Extent of the Canon:

The Apocrypha of the Old Testament are included in the Roman, excluded from the Protestant Canon.

2. Authority of the Latin Vulgate:

Put on a par with the Hebrew and Greek Scriptures by Rome; while Protestantism claims divine authority only for the original Scriptures of the inspired authors.

3. Popular use and circulation of the Bible:

Discouraged (and relatively forbidden) by Rome; encouraged by Protestantism, which goes hand in hand with the Word of God, and must stand or fall with it.

III. Objects of Worship.

Roman Catholic doctrine:

1. God (\emph{latria});

2. The Virgin Mary (\emph{hyperdulia});

3. Angels and Saints (\emph{dulia});

4. Images and Relics of Saints. Protestant doctrine:

God alone. All other worship is gross or refined idolatry.

The Roman Catholic Christian approaches Christ through human mediators, and virtually substitutes the worship of Mary for the worship of Christ; the Protestant approaches Christ directly, and prays to him as his only and all-sufficient High-Priest and Intercessor with the Father.

IV. Primitive State.

Difference (asserted by Roman Catholics, denied by Protestants) between the image of God (\emph{imago}, \textgreek{εἰκών, } \texthebrew{צֶלֶם}), \emph{i.e.}, the natural perfection of the first man as a rational and free being, and similitude of God (\emph{similitudo, } \textgreek{ὁμοίωσις, } \texthebrew{דְמדּת}), \emph{i.e.}, supernatural endowment of man with righteousness and holiness together with the immortality of the body.

V. Original Sin.

Roman Catholic doctrine:

Original sin is a negative defect (\emph{carentia justitiæ originalis}), or the loss of the similitude—not of the image—of God, and is entirely removed by baptism.

Protestant doctrine:

Original sin is a positive corruption and total depravity, involving the loss of (spiritual) freedom, and retains the character of sin after baptism.

VI. Justification by faith and good works (Roman Catholic);—or by faith alone (Protestant).

1. Different conceptions of justification (\textgreek{δικαίωσις, } \emph{justificatio}): a gradual process of making the sinner righteous (identical with sanctification);—or a judicial and declaratory act of God (acquittal of the penitent sinner on the ground of Christ's merits and on condition of faith in Christ), followed by sanctification.

2. Different conceptions of faith: intellectual assent and submission to divine authority;—or personal trust in Christ and living union with him.

3. Different position assigned to works: condition of justification;—or evidence of justification.

4. Assurance of justification and salvation: denied (except on the ground of a special revelation) by Roman Catholics; asserted by Protestants (though in different degrees).

Paul and James. Basis of reconciliation: faith operative in love.\footnote{Gal. v. 6, (\textgreek{πιστις δἰ ἀγάπης ἐνεργουμένη,} is to be explained as the dynamic middle, not as the passive, 'completed in love' (the \emph{fides formata} of Roman Catholic commentators).}

VII. Good works of believers.

The meritoriousness of good works (\emph{meritum ex congruo} and \emph{meritum ex condigno}): Works of supererogation, not commanded, but recommended (\emph{consilia evangelica}), with corresponding extra merits, which constitute a treasury at the disposal of the Pope for the dispensation of indulgences.

Here is the root of the ascetic and monastic system (\emph{vota monastica}: voluntary obedience, poverty, and celibacy), and the chief difference between Roman Catholic and Evangelical ethics.

VIII. The Church.

1. Identification of the Church of Christ with the Church of Rome—the fundamental error (the \textgreek{πρῶτον ψεῦδος}) of the papacy.

2. Distinction of the invisible Church (one and universal under the sole headship of Christ), and the visible Church (existing in many organizations or denominations): asserted by Protestants; denied by Roman Catholics.

3. Different conception and application of the attributes of the Church; unity, holiness, catholicity, apostolicity, indefectibility, infallibility, and exclusiveness, especially the last (\emph{extra ecclesiam nulla salus}, which is made to mean \emph{extra ecclesiam Romanam}).

IX. The Pope.

The infallible head of the Universal Church, the Vicar of Christ on earth, by virtue of his office as the successor of Peter.

This is the cardinal doctrine of Romanism, but rejected by Greeks and Protestants as an antichristian usurpation of the prerogative of Christ.

X. Sacraments in general.

1. Definition: visible signs of invisible grace instituted by the express command of Christ in the New Testament (Protestant);—or simply by the authority of the Church (Roman Catholic).

2. Number: seven (Roman Catholic);—or two (Baptism and the Lord's Supper).

3. Effect: \emph{ex opere operato} (\emph{i.e.}, by virtue of the objective act);—or through faith (as the subjective condition).

XI. Baptism.

Its effect on original sin; its relation to regeneration; its necessity for salvation; and several ritual differences.

XII. The Eucharist. Romanism holds, Protestantism denies:

1. Transubstantiation and the adoration of the elements.

2. The withdrawal of the cup from the laity.

3. The Eucharist as a sacrifice, \emph{i.e.}, an actual though unbloody repetition of Christ's sacrifice on the cross by the priest for the sins of the living and the dead (the souls in purgatory).

The celebration of the Mass is the centre of Roman Catholic worship.

XIII. The other five Sacraments: Confirmation, Penance, Matrimony, Ordination, Extreme Unction.

Maintained by Rome as sacraments proper; rejected by Protestants, or admitted only as semi- or quasi-sacramental acts.

1. Confirmation.

Retained by the Lutheran, Anglican, and the German Reformed Churches (as supplementary to infant baptism after a course of catechetical instruction). Rejected by other Protestant Churches, in which a voluntary union with the Church by a public profession of faith takes the place of confirmation.

2. Penance (\emph{sacramentum pænitentiæ}).

Auricular confession and priestly absolution; satisfaction for venial sins; indulgences. The Lutheran (and Anglican) standards approve private confession to the \emph{minister;} other Churches leave it entirely optional; all Protestants deny the efficacy of priestly absolution except as an official declaration of God's forgiving mercy to the penitent

3. Ordination.

A separate priesthood and clerical celibacy (Roman Catholic); the general priesthood of the laity and the right of the laity to participate in Church government (Protestant).

4. Matrimony.

Differences in matrimonial legislation, mixed marriages, and divorce.

5. Extreme unction.

Rejected by Protestants, who in James v. 14 emphasize the praying rather than 'the anointing with oil' (a physical remedy).

XIV. Purgatory.

A temporary middle place and state (until the final judgment) between heaven and hell for the purification of imperfect Christians, which may be advanced by prayers and masses in their behalf.

Protestantism holds that there are only two conditions in the other world, but with various degrees of bliss or misery.

The indulgences closely connected with purgatory were the first occasion, though not the cause, of the Reformation.

E. Doctrinal Differences among Evangelical Protestants.

I.—LUTHERANISM AND CALVINISM.

1. Baptismal Regeneration.

Baptism a means of regeneration (as concurrent with the sacramental act), and hence necessary for salvation;—or only a sign and seal of regeneration (whether concurrent or preceding or succeeding, according to God's free pleasure).

2. The Encharistic presence.

Corporeal real presence (in, with, and under the elements) for all communicants;—or spiritual real (dynamic and effective) presence for believers only.

3. Christological.

The extent of the \emph{communicatio idiomatum}.\footnote{That is, whether it includes also the \emph{genus majestaticum}, or the communication of the attributes of the divine nature to the human nature of Christ—affirmed by the Lutheran symbols, denied by the Reformed. See pp.319 sqq.} The ubiquity of Christ's body: asserted by the Lutheran Church (as a dogmatic support to its doctrine of the eucharistic multipresence); denied by the Reformed (as inconsistent with the limitations of humanity and the fact of Christ's ascension to heaven).

4. Predestination and the perseverance of saints.

No difference between Luther and Calvin, who were both Augustinians, but between their followers. (Synergism of Melanchthon in his later period. Semi-Augustinianism of the Formula of Concord. Extreme Calvinism of the Synod of Dort.)

II.—CALVINISM AND ARMINIANISM.

1. Election: unconditional;—or conditional.

2. Extent of redemption: limited to the elect;—or unlimited to all men.

3 and 4. Nature of faith and grace: irresistible;—or resistible.

5. Perseverance of saints;—or the possibility of total and final apostasy.

III.—CONGREGATIONALISM.

1. Conception of a Christian congregation or local church: a self-governing body of converted believers voluntarily associated for spiritual ends.

2. Independence of such a church of foreign jurisdiction.

3. Duty of voluntary fellowship with other churches.

IV.—BAPTIST DOCTRINES.

1. Congregationalism as sub III.

2. Baptism.

(\emph{a}) Its subjects: only responsible converts on the ground of a voluntary profession of their faith.

(\emph{b}) Its mode: total immersion of the body.

3. Universal liberty of conscience as a sphere over which civil government has no control. ('Soul-liberty.')\footnote{President Anderson, of Rochester University (article \emph{Baptists} in Johnson's \emph{Cyclopædia}, Vol. I. p. 383), enumerates four distinctive doctrinal principles of the Baptists: (1) immersion; (2) believers only to constitute a visible church; (3) responsible converts only entitled to baptism; (4) separation of Church and State, and independence of each individual church as a body of baptized believers of any other body, whether ecclesiastical or political. But the second article is held also by the Congregationalists, and the fourth can not be called an article of faith.}

V.—QUAKER DOCTRINES.

1. Universal diffusion of the inner light for the salvation of men.

2. Immediate revelation superior to, though concordant with, the outward testimony of the Scriptures.

3. The ministry of the gospel depending on inspiration, and not confined to a class or sex.

4. The sacraments are spiritual acts, not visible rites and ceremonies, as under the old dispensation.

5. Worship is purely inward, and depends upon the immediate moving of the Holy Spirit.

6. Universal religious liberty.

VI.—METHODIST DOCTRINES.

1. Universal offer of salvation in different dispensations.

2. Witness of the Spirit, or assurance of present acceptance with God.

3. Christian perfection, or perfect sanctification.

F. Orthodox Protestantism and Heterodox Protestantism.

I. Socinianism (Unitarianism). Denies the following œcumenical doctrines:

1. The Trinity.

2. The Incarnation and eternal Divinity of Christ.

3. Original sin and guilt.

4. The vicarious atonement.

II. Universalism departs from the orthodox doctrines of the—

1. Nature and extent of sin and its consequences.

2. Endless punishment. (Difference between Restorationism and Universalism proper).

III. Swedenborgianism asserts:

1. A new revelation and a new Church (the New Jerusalem).

2. Intercourse with the spirit world.

3. It limits the number of the canonical Scriptures.

4. It claims to unlock the deeper inner sense of the Scriptures.

5. It dissents from the evangelical doctrines of the tripersonality of the Godhead, the incarnation, the atonement, justification, the Church, the sacraments, and the resurrection.

\xsection{The Disciples of Christ.}

§ 116. The Disciples of Christ

Literature.

Richardson:  \emph{Memoirs of A. Campbell}, new ed., Cinti., 1888—\emph{A Debate between Rev. A. Campbell and Rev. N. L. Rice on Christian Baptism and Eccles. Creeds as Terms of Communion}, Nov. 15–Dec. 2, 1843, Lexington, 1844, pp. 912.—  I. Errett:  \emph{Our Position}, a tract, Cinti., 1901.—  J. H. Garrison:  \emph{The Old Faith Restated}, St. Louis, 1891.—  B. B. Tyler:  \emph{Hist. of the Disciples of Christ}, Am. Ch. Hist. Series, N.Y., 1894.— F. D. Powers: Art. in Schaff-Herzog Enc., III., 443 sq.—  J. S. Lamar:  \emph{Life of I. Errett}, 2 vols., Cinti., 1894.— W. T. Moore:  \emph{Compar. Hist. of the Disciples of Christ}, N.Y., 1909, pp. 830.— Summerbell:  \emph{The Christians and the Disciples}, Dayton, 1906.— Errett Gates:  \emph{The Disciples of Christ}, in ``Story of the Chh.,'' Chaps. XI–XIV., N.Y., 1905.— P. Ainslie:  \emph{The Message of the Disciples for the Union of the Church}, N.Y., 1913, pp. 210. Gives the ``Declaration and Address'' in full. Also \emph{The Scandal of Christianity}, N. Y., 1929.—  J. K. Kellems:  \emph{A. Campbell and the Disciples}, N.Y., 1930, pp. 409.

A large and influential Christian body whose historic position has been antagonism to all Church creeds containing articles formulated by non-biblical writers is the group known as The Disciples of Christ. In the number of its members, the fifth Protestant ecclesiastical body in the United States, it arose early in the nineteenth century in Western Pennsylvania under the leadership of Thomas Campbell, a Seceder minister from North Ireland and his son, Alexander Campbell. Born in Ireland, Sept. 12, 1788, Alexander came to America, 1809, settled in the Western part of Pennsylvania and died March 4, 1866, in Bethany, West Virginia, where he had established a college. Seeking relief from the restraints of ecclesiastical formularies and a return to the so-called implicity and ordinances of ``original Christianity,'' Thomas, then a Presbyterian minister, formed ``the Christian Association of Washington,'' Pennsylvania, and issued, 1809, a ``Declaration and Address'' to which the Disciples go back as the justification and basis of their existence as a distinct group. In 1811, he and his followers joined themselves in an independent organization at Brush Run, Pa. Two years later, the organization united with the Redstone Baptist Association. A division arising in this body over the principles of the ``Church Reformation,'' as the movement led by the Campbells was called, the followers of the Campbells constituted themselves an independent body, 1827. Four years later, this body was enlarged by the accession of a number of churches which followed Rev. Barton W. Stone, once a member of the Presbytery of Lexington, Kentucky. The members of the Kentucky churches preferred to call themselves by the simple name ``Christians'' and for this reason the ``Disciples'' have often gone by that name.

Alexander Campbell, the real founder of the new movement, was ordained 1812 and immersed a few months later. He was a man of intellectual vigor and independence of thought, of positiveness of conviction and statement, and became abundant in labors. His views were set forth not only in the pulpit but through the columns of two periodicals, the Christian Baptist and the Millennial Harbinger, and in public discussions on the platform. These discussions, which aroused wide attention in Southern Ohio and the South West, were carried on with Robert Owen, 1829, and Archbishop Purcell of the Roman Catholic Church, 1837, both in Cincinnati, and with the noted Presbyterian polemic theologian, Dr. Nathan L. Rice, in Lexington, Ky., 1843. The last discussion, which lasted sixteen days, had Henry Clay as its chairman. Campbell was accused of being ``contentious'' and ``disputatious'' while his skill as a debater was generally recognized.

The distinctive tenets of the Disciples, as set forth in the Declaration and Address, are the ``alone-sufficiency and all-sufficiency of the Bible''—to use the language of Alexander Campbell—and the unadulterated evil of Christian creeds and the denominational divisions in Christendom. With other Christians, they hold to the doctrines of the trinity, original sin, Christ's atonement and resurrection, the necessity of repentance and regeneration, and the two future states. They reject—to follow the Declaration and Address—human opinions and the inventions of men as having any authority in the Church of God and profess ``to stand upon the ground on which the Church stood at the beginning . . . and to take up things just as the Apostles left them.'' The Society, so it was affirmed, ``was formed for the sole purpose of promoting simple evangelical Christianity and to promote only such measures as, reduce to practice the original form of Christianity expressly exhibited in the sacred page. . . . Everything not taught and enjoined in the Bible is of no authority and nothing is to be made a term of communion among Christians which is not as old as the New Testament.'' Creeds of human composition are a calamity and have been the cause of the historic divisions in the Church and ``divisions among Christians is a horrid evil. Although the Church of Christ must necessarily exist in particular and distinctive societies, locally separate one from another, yet ought there to be no schisms, no uncharitable divisions, among them.'' Further, by the Declaration and Address, the society rejected the application to itself of the name ``Church'' and its members desired to be regarded ``merely as voluntary advocates of Church reformation; and, tired of the jarrings and janglings of a party spirit, to restore unity, peace and purity to the whole Church of God,''—a large, if not ambitious, and certainly most laudable purpose. In accordance with these affirmations, whereby the Scriptures are not only treated as the sole standard of Christian teaching but its language the sole organ through which it is to be conveyed and all human formularies intended to state and summarize those teachings are set aside as evil in their consequences, the Disciples have been inclined to look upon themselves as pioneers in the movement of Christian unity and Church union. For twenty years the body has had within itself an Association for the Promotion of Christian Unity. In spite of their formal declarations and wishes, the Disciples are a distinct ecclesiastical body, following the congregational scheme of church government and observing usages distinguishing them from many Christian bodies such as the weekly observance of the Lord's Supper and baptism by immersion, though denying baptismal regeneration. Whether such usages, the doctrines explicitly held in common with all Christians and their opposition to creeds, constitute a confessional formulary or not, seems to admit of two answers. The large volume of Kellems, 1930, ascribes to Alexander Campbell a doctrinal system no less positive than are the systems of other ecclesiastical bodies. By their evangelical activity and warm manifestation of union among themselves, the Disciples of Christ have not only reached large bodies of people but won the fellowship of other Christian communions and have shown that a Church's efficiency and solidarity does not necessarily depend upon an explicit formulary of human composition. The most prominent personage among the Disciples since Alexander Campbell has been James A. Garfield, President of the United States.

\xsection{The Universalists.}

§ 117. The Universalists.

The Universalist Churches of the United States, starting in New England, have modified the orthodox Christian system as expressed in the historic creeds but while differing among themselves, they retain reverence for Christ as a divine teacher, the belief in the immortality of the soul and suitable awards after death for conduct in this world. A movement towards the union of Congregational and Universalist churches has had advocates as in California, 1930. The New England Convention of Universalist Churches, meeting in Winchester, New Hampshire, 1803, adopted a Profession of Belief in three articles. Eddy, in his \emph{Hist. of Universalism}, says that ``while the Profession was sufficiently definite to exclude the possibility of mistaking its most prominent thought, the reconciliation of all souls to God, it was sufficiently liberal to be acceptable alike to Trinitarians and to Unitarians, to the believer in future punishment and to the believer that the consequences of sin are confined to this life.'' In 1899 the General Convention, meeting in Boston, added to the Winchester Profession, clauses giving ``The Conditions of Fellowship.'' At a meeting of the Convention in Winchester, 1903, celebrating the adoption of the Profession, its three articles were spoken of ``as the first explicit statement in a creed of what is known as liberal Christianity.''

``The Profession of Belief and Conditions of Fellowship are as follows''\footnote{Minutes of the Winchester Convention, Washington, 1929. In 1878 the Universalists of Boston and vicinity put forth a statement of belief in 9 articles which was not adopted by the convention. See R. Eddy: \emph{Hist. of Universalism}, in Am. Church Hist, series, X; 255–507.—Ed.}

Art. I. We believe that the Holy Scriptures of the Old and New Testaments contain a revelation of the character of God and of the duty, interest and final destination of mankind. Art. II. We believe that there is one God, whose nature is Love revealed in one Lord Jesus Christ, by one Holy Spirit of Grace, who will finally restore the whole family of mankind to holiness and happiness. Art. III. We believe that holiness and true happiness are inseparably connected, and that believers ought to be careful to maintain order and practice good works; for these things are good and profitable unto men.

The Conditions of Fellowship shall be as follows: 1. The acceptance of the essential principles of the Universalist Faith to wit: 1 The Universal Fatherhood of God; 2 The Spiritual authority and leadership of His Son, Jesus Christ; 3 The trustworthiness of the Bible as containing a revelation from God; 4 The certainty of just retribution for sin; 5 The final harmony of all souls with God.

The Winchester Profession is commended as containing these principles but neither this nor any other precise form of words is required as a condition of fellowship, provided always that the principles above stated be professed.

2. The acknowledgment of the authority of the General Convention and assent to its laws.

\xsection{The Unitarians.}

§ 118. The Unitarians.

The Unitarian Churches are a free association of societies which lay stress upon practical aims and require no subscription to a doctrinal formula. The nearest approach to such a formula is the declaration written by James Freeman Clarke and published by the Unitarian Sunday School Union, affirming ``belief in the fatherhood of God, the brotherhood of man, the leadership of Jesus, salvation by character, the progress of mankind upward and onward forever.'' Founded in 1825, the American Unitarian Association announced its object to be ``to diffuse the knowledge and promote the interest of pure Christianity.'' The older Unitarians, represented by Channing, while they rejected binding creedal formulas, the doctrines of the Trinity and total depravity and what was called the bleak Calvinism of New England, held to the exaltation of Christ and the immortality of the soul. Twenty of the original congregations of Massachusetts, including the Plymouth church, allied themselves with the movement. Humanistic efforts have been emphasized. In Boston, which became the home of literary culture, many of its representatives broke loose entirely from the historic New England ecclesiastical system and abandoned Christianity as a supernatural revelation for the philosophy known as transcendentalism. By the middle of the 19th century, Unitarianism had become synonymous with religious liberalism, basing its conclusions in part upon the results of German rationalistic criticism. Outside of Massachusetts, its following has been small and its churches include groups which join with the name Unitarian extreme tenets of religious liberalism. President Eliot of Harvard pronounced ``independent thought the chief feature of Unitarianism.''

In 1865, largely under the influence of Dr. Bellows of New York City, ``the National Conference of Unitarian and Other Christian Churches'' was organized and passed the following resolution:

``To secure the largest unity of the spirit and the widest practical co-operation of our body, it is hereby understood that all the resolutions and declarations of this convention are expressions only of its majority, committing in no degree those who object to them, claiming no other than a moral authority over the members of the convention, or the churches represented here, and are all dependent wholly for their effect upon the consent they command on their own merits from the churches here represented or belonging within the circles of our special fellowship.''

The words used in the preamble ``The obligations of all disciples of the Lord Jesus Christ'' led to a warm discussion and the formation of ``The Free Religious Association.'' At a meeting of the National Conference, 1894, the following preamble concerning faith and fellowship was adopted and has been interpreted in some sections to include in the fellowship of the Unitarian churches, members of the Brahmo Somaj of India and all others who ``sympathize with the spirit and practical aims'' of the Unitarians:\footnote{Batchelor in \emph{Christian Register}, 1906, p. 202, 203. See J. H. Allen: \emph{The Unitarian Movement since the Reformation}, in Am. Ch. Hist. Series, X, pp. 1–249, N.Y., 1894; G. W. Cooke: \emph{Unitarianism in America}, Boston, 1902. E. Emerton: \emph{Unitar. Thought}, N.Y., 1911.}

``The Conference of Unitarian and other Christian Churches was formed in the year 1865, with the purpose of strengthening the churches and societies which should unite in it for more and better work for the kingdom of God. These churches accept the religion of Jesus, holding, in accordance with his teaching, that practical religion is summed up in love to God and love to man. The Conference recognizes the fact that its constituency is Congregational in tradition and polity. Therefore, it declares that nothing in this constitution is to be construed as an authoritative test; and we cordially invite to our working fellowship any who, while differing from us in belief, are in general sympathy with our spirit and our practical aims.''
